% Kôsaku Yosida: Functional Analysis (6th ed., Springer, Grundlehren 123)
% 按 1995 Springer 重印版 (第6版) 重排, 版式复刻原书
\documentclass{yosida}

% XeLaTeX text, keeping Latin Modern and the existing AMS math fonts.
\usepackage{lmodern}
\usepackage{fontspec}
\setmainfont{Latin Modern Roman}[
  Ligatures=TeX,
  SmallCapsFont=lmromancaps10-regular.otf,
  BoldFeatures={SmallCapsFont=lmromancaps10-regular.otf,
    SmallCapsFeatures={FakeBold=1.5}}]
\usepackage{textcomp}
\usepackage{multicol}
\usepackage{indentfirst}
\usepackage{enumerate}
\usepackage{etoolbox}
\AtBeginEnvironment{enumerate}{\everymath{}}
\everymath{\displaystyle}
% ---------- 定理环境在 yosida.cls 中以 amsthm 自定义样式声明 ----------

% ---------- 记号宏 ----------
% 原书各种常见记号统一由下列宏给出, 正文一律使用宏而非临时写法。
\DeclareMathOperator{\supp}{supp}          % 支集 support
\DeclareMathOperator{\dist}{dist}          % 距离 distance
\DeclareMathOperator{\Conv}{Conv}          % 凸包 convex hull
\DeclareMathOperator{\tr}{tr}              % 迹 trace
\DeclareMathOperator{\LIM}{LIM}            % 广义极限 (Banach limit)
\DeclareMathOperator*{\vraimax}{vrai\,max} % 本质上确界 (essential supremum)
\DeclareMathOperator{\sgn}{sgn}            % 符号函数
\DeclareMathOperator{\diam}{diam}          % 直径
% 微分符号: \mathop{} 让它在前面留出一个细空格, \! 取消紧跟其后的空隙,
% d 用直立体(原书如此), 全书统一用 \dd。
\providecommand{\dd}{\mathop{}\!\mathrm{d}}
% 实部/虚部记号沿用 LaTeX 内建的 \Re (ℜ) 与 \Im (ℑ); 0 处 OCR 残写法
% (Re( / Im( / "R e (" ) 已统一为这两个宏。

% ---------- 超链接 (规范 §8) ----------
\usepackage[dvipsnames]{xcolor}
\usepackage[
  colorlinks=true,
  pageanchor=false,
  linkcolor=blue,
  citecolor=red,
  urlcolor=MidnightBlue,
  linktoc=all,
  bookmarksnumbered=true,
  bookmarksopen=true,
  bookmarksopenlevel=1]{hyperref}
\usepackage{bookmark}
\bookmarksetup{numbered,open,openlevel=1}
% Explicit foreground destinations avoid a discarded zero-size DVI page-anchor
% box on dense math pages. Index targets retain the exact page.N naming contract.
\makeatletter
\AddToHook{shipout/foreground}{%
  \begingroup
    \edef\yos@shipoutpage{\thepage}%
    \ifx\yos@shipoutpage\@empty\else
      \xdef\@currentHpage{page.\yos@shipoutpage}%
      \special{pdf:dest (page.\yos@shipoutpage) [@thispage /XYZ 70 806 null]}%
    \fi
  \endgroup}
\makeatother

% Source Supplementary Notes footnotes use stars; numbered notes restart by chapter.
\makeatletter
\@addtoreset{footnote}{yoschap}
\newcommand{\yosstarfootnote}[1]{%
  \begingroup\renewcommand{\thefootnote}{*}\footnote[0]{#1}\endgroup}
% The source already prints a star in the heading. Create its footnote body
% and a unique hyperlink destination without a second mark in the paragraph.
\newcommand{\yosstarfootnotetext}[1]{%
  \begingroup\renewcommand{\thefootnote}{*}%
  \stepcounter{Hfootnote}%
  \global\let\Hy@saved@currentHref\@currentHref
  \hyper@makecurrent{Hfootnote}%
  \global\let\Hy@footnote@currentHref\@currentHref
  \global\let\@currentHref\Hy@saved@currentHref
  \footnotetext[0]{#1}\endgroup}
% A printed paragraph condition has a target without consuming an equation number.
\newcommand{\yosparagraphtag}[1]{%
  \phantomsection\def\@currentlabel{#1}%
  \unskip\nobreak\hfill\mbox{(#1)}\par}
\newcommand{\yosnamedtag}[1]{\phantomsection\def\@currentlabel{#1}}
\makeatother

% ---------- 参考文献 (biblatex, 编号按原书次序) ----------
\usepackage[backend=biber,style=numeric,sorting=none,
  isbn=false,url=false,doi=false,eprint=false]{biblatex}
\DeclareSortingTemplate{yosidasource}{\sort{\field{userc}}}
\ExecuteBibliographyOptions{sorting=yosidasource}
\addbibresource{references.bib}
% 姓前名后 + 小型大写, 复刻原书著者样式
\DeclareNameAlias{default}{family-given}
\renewcommand*{\mkbibnamefamily}{\textsc}
% 文献条目照录原书文字(置于 note 字段), 精简输出驱动
\DeclareFieldFormat{note}{#1}
% Source citations restart within each primary author. Biber keys remain unique.
\DeclareFieldFormat{labelnumber}{\thefield{usera}}
\DeclareFieldFormat{labelnumberwidth}{\mkbibbrackets{\thefield{usera}}}
\DeclareBibliographyDriver{misc}{%
  \usebibmacro{bibindex}%
  \usebibmacro{begentry}%
  \usebibmacro{author}%
  \newunit\newblock
  \printfield{note}%
  \usebibmacro{finentry}}

\begin{document}

% 各章/节内编号(公式按节从1重排)后, 给 hyperref 提供唯一的锚点名,
% 避免不同章节的同号公式/定理共用同一 PDF 目标 (链接跳错页)。
% \tag 不步进 equation 计数器 (插入编号如 (9') 不应打乱后续编号), 因此带
% \tag 的公式会与相邻公式的锚点重名; 这里把源码行号并入公式锚点使其唯一
% (书中不存在同一行有两个公式环境的情形), 从而不必改动 hyperref 的默认
% 页码锚点 (索引的 \hyperlink{page.N} 依赖它)。
\AtBeginDocument{%
  \def\theHequation{\arabic{yoschap}.\arabic{yossec}.\arabic{equation}.\the\inputlineno}%
  \def\theHtheorem{\arabic{yoschap}.\arabic{yossec}.\arabic{theorem}}%
  \def\theHproposition{\arabic{yoschap}.\arabic{yossec}.\arabic{proposition}}%
  \def\theHlemma{\arabic{yoschap}.\arabic{yossec}.\arabic{lemma}}%
  \def\theHcorollary{\arabic{yoschap}.\arabic{yossec}.\arabic{corollary}}%
  \def\theHdefinition{\arabic{yoschap}.\arabic{yossec}.\arabic{definition}}%
  \def\theHexample{\arabic{yoschap}.\arabic{yossec}.\arabic{example}}}

% ==== 标题页 ====
\begin{titlepage}
\thispagestyle{front}\pagenumbering{gobble}
\vspace*{30mm}
\begin{center}
{\Large K\^osaku Yosida\par}
\vspace{12mm}
{\huge\bfseries Functional Analysis\par}
\vspace{10mm}
{\large Sixth Edition\par}
\vspace{14mm}
{\large Springer-Verlag\par Berlin Heidelberg New York\par}
{\large 1980\par}
\end{center}
\end{titlepage}
% ==== 序言 ====
\pagenumbering{roman}\setcounter{page}{5}
\vspace{6mm}\begin{center}\bfseries Preface to the First Edition\end{center}\vspace{2mm}
The present book is based on lectures given by the author at the University of Tokyo during the past ten years. It is intended as a textbook to be studied by students on their own or to be used in a course on Functional Analysis, i.e., the general theory of linear operators in function spaces together with salient features of its application to diverse fields of modern and classical analysis. 

Necessary prerequisites for the reading of this book are summarized, with or without proof, in Chapter~\hyperref[chap:0]{$0$} under titles: Set Theory, Topological Spaces, Measure Spaces and Linear Spaces. Then, starting with the chapter on Semi-norms, a general theory of Banach and Hilbert spaces is presented in connection with the theory of generalized functions of S. L. \textsc{Sobolev} and L. \textsc{Schwartz}. While the book is primarily addressed to graduate students, it is hoped it might prove useful to research mathematicians, both pure and applied. The reader may pass, e.g., from Chapter~\hyperref[chap:IX]{IX} (Analytical Theory of Semi-groups) directly to Chapter~\hyperref[chap:XIII]{XIII} (Ergodic Theory and Diffusion Theory) and to Chapter~\hyperref[chap:XIV]{XIV} (Integration of the Equation of Evolution). Such materials as ``Weak Topologies and Duality in Locally Convex Spaces'' and ``Nuclear Spaces'' are presented in the form of the appendices to Chapter~\hyperref[chap:V]{V} and Chapter~\hyperref[chap:X]{X}, respectively. These might be skipped for the first reading by those who are interested rather in the application of linear operators. 

In the preparation of the present book, the author has received valuable advice and criticism from many friends. Especially, Mrs. K. \textsc{Hille} has kindly read through the manuscript as well as the galley and page proofs. Without her painstaking help, this book could not have been printed in the present style in the language which was not spoken to the author in the cradle. The author owes very much to his old friends, Professor E. \textsc{Hille} and Professor S. \textsc{Kakutani} of Yale University and Professor R. S. \textsc{Phillips} of Stanford University for the chance to stay in their universities in $1962$, which enabled him to polish the greater part of the manuscript of this book, availing himself of their valuable advice. Professor S. \textsc{Ito} and Dr. H. \textsc{Komatsu} of the University of Tokyo kindly assisted the author in reading various parts of the galley proof, correcting errors and improving the presentation. To all of them, the author expresses his warmest gratitude. 

Thanks are also due to Professor F. K. \textsc{Schmidt} of Heidelberg University and to Professor T. \textsc{Kato} of the University of California at Berkeley who constantly encouraged the author to write up the present book. Finally, the author wishes to express his appreciation to Springer-Verlag for their most efficient handling of the publication of this book. 

Tokyo, September $1964$ 

Kôsaku Yosida 

\vspace{6mm}\begin{center}\bfseries Preface to the Second Edition\end{center}\vspace{2mm}
In the preparation of this edition, the author is indebted to Mr. \textsc{Floret} of Heidelberg who kindly did the task of enlarging the Index to make the book more useful. The errors in the second printing are corrected thanks to the remarks of many friends. In order to make the book more up-to-date, Section~\hyperref[sec:XIV.4]{$4$} of Chapter~\hyperref[chap:XIV]{XIV} has been rewritten entirely for this new edition. 

Tokyo, September $1967$ 

Kôsaku Yosida 

\vspace{6mm}\begin{center}\bfseries Preface to the Third Edition\end{center}\vspace{2mm}
A new Section ($9$. Abstract Potential Operators and Semi-groups) pertaining to G. \textsc{Hunt}'s theory of potentials is inserted in Chapter~\hyperref[chap:XIII]{XIII} of this edition. The errors in the second edition are corrected thanks to kind remarks of many friends, especially of Mr. \textsc{Klaus-Dieter Bier-Stedt}. 

Kyoto, April $1971$ 

Kôsaku Yosida 

\vspace{6mm}\begin{center}\bfseries Preface to the Fourth Edition\end{center}\vspace{2mm}
Two new Sections ``$6$. Non-linear Evolution Equations $1$ (The Kōmura-Kato Approach)'' and ``$7$. Non-linear Evolution Equations $2$ (The Approach Through The Crandall-Liggett Convergence Theorem)'' are added to the last Chapter~\hyperref[chap:XIV]{XIV} of this edition. The author is grateful to Professor Y. Kōmura for his careful reading of the manuscript. 

Tokyo, April $1974$ 

Kôsaku Yosida 

\vspace{6mm}\begin{center}\bfseries Preface to the Fifth Edition\end{center}\vspace{2mm}
Taking advantage of this opportunity, supplementary notes are added at the end of this new edition and additional references to books have been entered in the bibliography. The notes are divided into two categories. The first category comprises two topics: the one is concerned with the time reversibility of Markov processes with invariant measures, and the other is concerned with the uniqueness of the solution of time dependent linear evolution equations. The second category comprises those lists of recently published books dealing respectively with Sobolev Spaces, Trace Operators or Generalized Boundary Values, Distributions and Hyperfunctions, Contraction Operators in Hilbert Spaces, Choquet's Refinement of the Krein-Milman Theorem and Linear as well as Nonlinear Evolution Equations. 

A number of minor errors and a serious one on page $459$ in the fourth edition have been corrected. The author wishes to thank many friends who kindly brought these errors to his attention. 

Kamakura, August $1977$ 

Kôsaku Yosida 

\vspace{6mm}\begin{center}\bfseries Preface to the Sixth Edition\end{center}\vspace{2mm}
Two major changes are made to this edition. The first is the rewriting of the Chapter~\hyperref[chap:VI]{VI},$6$ to give a simplified presentation of Mikusiński's Operational Calculus in such a way that this presentation does not appeal to Titchmarsh's theorem. The second is the rewriting of the Lemma together with its Proof in the Chapter~\hyperref[chap:XII]{XII},$5$ concerning the Representation of Vector Lattices. This rewriting is motivated by a letter of Professor E. Coimbra of Universidad Nova de Lisboa kindly suggesting the author's careless phrasing in the above Lemma of the preceding edition. 

A number of misprints in the fifth edition have been corrected thanks to kind remarks of many friends. 

Kamakura, June $1980$ 

Kôsaku Yosida 



% 目录 (默认 book 样式; 编号文本由 \yoschapter/\yossection 写入条目)
\setcounter{page}{9}
\begingroup
\yostocentryfalse
\yoschapterstar{Contents}\label{chap:toc}
\endgroup
\yostoc

% 正文
\clearpage
\pagenumbering{arabic}

% ==== 0 ====
\yoschapter{0}{Preliminaries}
\label{chap:0}

It is the purpose of this chapter to explain certain notions and theorems used throughout the present book. These are related to Set Theory, Topological Spaces, Measure Spaces and Linear Spaces. 


\yossection{1}{Set Theory}
\label{sec:0.1}

Sets. $x \in X$ means that $x$ is a member or element of the set $X; x \notin X$ means that $x$ is not a member of the set $X$. We denote the set consisting of all $x$ possessing the property $P$ by $\{x; P\}$. Thus $\{y; y = x\}$ is the set $\{x\}$ consisting of a single element $x$. The void set is the set with no members, and will be denoted by $\emptyset$. If every element of a set $X$ is also an element of a set $Y$, then $X$ is said to be a subset of $Y$ and this fact will be denoted by $X \subseteq Y$, or $Y \supseteq X$. If $\mathfrak{X}$ is a set whose elements are sets $X$, then the set of all $x$ such that $x \in X$ for some $X \in \mathfrak{X}$ is called the union of sets $X$ in $\mathfrak{X}$; this union will be denoted by $\bigcup\limits_{X \in \mathfrak{X}} X$. The intersection of the sets $X$ in $\mathfrak{X}$ is the set of all $x$ which are elements of every $X \in \mathfrak{X}$; this intersection will be denoted by $\bigcap\limits_{X \in \mathfrak{X}} X$. Two sets are disjoint if their intersection is void. A family of sets is disjoint if every pair of distinct sets in the family is disjoint. If a sequence $\{X_n\}_{n=1,2,\ldots}$ of sets is a disjoint family, then the union $\bigcup\limits_{n=1}^{\infty} X_n$ may be written in the form of a sum $\sum\limits_{n=1}^{\infty} X_n$. 

Mappings. The term mapping, function and transformation will be used synonymously. The symbol $f: X \to Y$ will mean that $f$ is a single-valued function whose domain is $X$ and whose range is contained in $Y$; for every $x \in X$, the function $f$ assigns a uniquely determined element $f(x) = y \in Y$. For two mappings $f: X \to Y$ and $g: Y \to Z$, we can define their composite mapping $gf: X \to Z$ by $(gf)(x) = g(f(x))$. The symbol $f(M)$ denotes the set $\{f(x); x \in M\}$ and $f(M)$ is called the image of $M$ under the mapping $f$. The symbol $f^{-1}(N)$ denotes the set $\{x; f(x) \in N\}$ and $f^{-1}(N)$ is called the inverse image of $N$ under the mapping $f$. It is clear that 
$Y_{1} = f(f^{-1}(Y_{1}))$ for all $Y_{1} \subseteq f(X)$, and $X_{1} \subseteq f^{-1}(f(X_{1}))$ for all $X_{1} \subseteq X$.

If $f: X \to Y$, and for each $y \in f(X)$ there is only one $x \in X$ with $f(x) = y$, then $f$ is said to have an inverse (mapping) or to be one-to-one. The inverse mapping then has the domain $f(X)$ and range $X$; it is defined by the equation $x = f^{-1}(y) = f^{-1}(\{y\})$. 

The domain and the range of a mapping $f$ will be denoted by $D(f)$ and $R(f)$, respectively. Thus, if $f$ has an inverse then 
\[
f ^{-1} (f (x)) = x \text{ for all } x \in D (f), \text{ and } f (f ^{-1} (y)) = y \text{ for all } y \in R (f).
\]

The function $f$ is said to map $X$ onto $Y$ if $f(X) = Y$ and into $Y$ if $f(X) \subseteq Y$. The function $f$ is said to be an extension of the function $g$ and $g$ a restriction of $f$ if $D(f)$ contains $D(g)$, and $f(x) = g(x)$ for all $x$ in $D(g)$. 

\yossub{Zorn's Lemma}

\begin{definitionu} Let $P$ be a set of elements $a, b, \ldots$. Suppose there is a binary relation defined between certain pairs $(a, b)$ of elements of $P$, expressed by $a \prec b$, with the properties: 
$a\prec a$; if $a\prec b$ and $b\prec a$, then $a=b$; if $a\prec b$ and $b\prec c$, then $a\prec c$ (transitivity).
\end{definitionu}

Then $P$ is said to be partially ordered (or semi-ordered) by the relation $\prec$. 

\begin{examplen} If $P$ is the set of all subsets of a given set $X$, then the set inclusion relation $(A \subseteq B)$ gives a partial ordering of $P$. The set of all complex numbers $z = x + iy$, $w = u + iv$, $\ldots$ is partially ordered by defining $z \prec w$ to mean $x \leqq u$ and $y \leqq v$.

\end{examplen}
\begin{definitionu} Let $P$ be a partially ordered set with elements $a, b, \ldots$. If $a < c$ and $b < c$, we call $c$ an upper bound for $a$ and $b$. If furthermore $c < d$ whenever $d$ is an upper bound for $a$ and $b$, we call $c$ the least upper bound or the supremum of $a$ and $b$, and write $c = \sup (a, b)$ or $a \vee b$. This element of $P$ is unique if it exists. In a similar way we define the greatest lower bound or the infimum of $a$ and $b$, and denote it by $\inf (a, b)$ or $a \wedge b$. If $a \vee b$ and $a \wedge b$ exist for every pair $(a, b)$ in a partially ordered set $P$, $P$ is called a lattice. 
\end{definitionu}

\begin{examplen} The totality of subsets $M$ of a fixed set $B$ is a lattice by the partial ordering $M_{1} \prec M_{2}$ defined by the set inclusion relation $M_{1} \subseteq M_{2}$.

\end{examplen}
\begin{definitionu} $A$ partially ordered set $P$ is said to be linearly ordered (or totally ordered) if for every pair $(a, b)$ in $P$, either $a < b$ or $b < a$ holds. A subset of a partially ordered set is itself partially ordered by the relation which partially orders $P$; the subset might turn out to be linearly ordered by this relation. If $P$ is partially ordered and $S$ is a subset of $P$, an $m \in P$ is called an upper bound of $S$ if $s < m$ for every $s \in S$. An $m \in P$ is said to be maximal if $p \in P$ and $m < p$ together imply $m = p$. 
\end{definitionu}

\begin{lemman}[Zorn's Lemma] Let $P$ be a non-empty partially ordered set with the property that every linearly ordered subset of $P$ has an upper bound in $P$. Then $P$ contains at least one maximal element. 
\end{lemman}

It is known that Zorn's lemma is equivalent to Zermelo's axiom of choice in set theory. 


\yossection{2}{Topological Spaces}
\label{sec:0.2}

\yossub{Open Sets and Closed Sets}

\begin{definitionu} $A$ system $\tau$ of subsets of a set $X$ defines a topology in $X$ if $\tau$ contains the void set, the set $X$ itself, the union of every one of its subsystems, and the intersection of every one of its finite subsystems. The sets in $\tau$ are called the open sets of the topological space $(X,\tau)$; we shall often omit $\tau$ and refer to $X$ as a topological space. Unless otherwise stated, we shall assume that a topological space $X$ satisfies Hausdorff's axiom of separation: 

For every pair $(x_{1}, x_{2})$ of distinct points $x_{1}, x_{2}$ of $X$, there exist disjoint open sets $G_{1}, G_{2}$ such that $x_{1} \in G_{1}, x_{2} \in G_{2}$. 
\end{definitionu}

A neighbourhood of the point $x$ of $X$ is a set containing an open set which contains $x$. A neighbourhood of the subset $M$ of $X$ is a set which is a neighbourhood of every point of $M$. A point $x$ of $X$ is an accumulation point or limit point of a subset $M$ of $X$ if every neighbourhood of $x$ contains at least one point $m \in M$ different from $x$. 

\begin{definitionu} Any subset $M$ of a topological space $X$ becomes a topological space by calling ``open'' the subsets of $M$ which are of the form $M \cap G$ where $G'$'s are open sets of $X$. The induced topology of $M$ is called the relative topology of $M$ as a subset of the topological space $X$. 
\end{definitionu}

\begin{definitionu} $A$ set $M$ of a topological space $X$ is \emph{closed} if it contains all its accumulation points. It is easy to see that $M$ is closed iff\footnote{\hspace{0.35em}\emph{iff} is the abbreviation for ``if and only if''.} its \emph{complement} $M^{C}=X-M$ is open. Here $A$-$B$ denotes the totality of points $x\in A$ not contained in $B$. If $M\subseteq X$, the intersection of all closed subsets of $X$ which contain $M$ is called the \emph{closure} of $M$ and will be denoted by $M^{a}$ (the superscript ``a'' stands for the first letter of the German: abgeschlossene Hülle). 
\end{definitionu}

Clearly $M^a$ is closed and $M \subseteq M^a$; it is easy to see that $M = M^a$ iff $M$ is closed. 

\yossub{Metric Spaces}

\begin{definitionu} If $X$, $Y$ are sets, we denote by $X \times Y$ the set of all ordered pairs $(x, y)$ where $x \in X$ and $y \in Y$; $X \times Y$ will be called the \emph{Cartesian product} of $X$ and $Y$. $X$ is called a metric space if there is defined a function $d$ with domain $X \times X$ and range in the real number field $\mathbb{R}^{1}$ such that 
$d(x_1,x_2)\geqq0$, and $d(x_1,x_2)=0$ iff $x_1=x_2$; $d(x_1,x_2)=d(x_2,x_1)$; $d(x_1,x_3)\leqq d(x_1,x_2)+d(x_2,x_3)$ (the triangle inequality).
$d$ is called the metric or the distance function of $X$. With each point $x_{0}$ in a metric space $X$ and each positive number $r$, we associate the set $S(x_{0}; r) = \{x \in X; d(x, x_{0}) < r\}$ and call it the open sphere with centre $x_{0}$ and radius $r$. Let us call ``open'' the set $M$ of a metric space $X$ iff, for every point $x_{0} \in M$, $M$ contains a sphere with centre $x_{0}$. Then the totality of such ``open'' sets satisfies the axiom of open sets in the definition of the topological space. 
\end{definitionu}

Hence a metric space $X$ is a topological space. It is easy to see that a point $x_{0}$ of $X$ is an accumulation point of $M$ iff, to every $\varepsilon > 0$, there exists at least one point $m \neq x_{0}$ of $M$ such that $d(m, x_{0}) < \varepsilon$. The $n$-dimensional Euclidean space $\mathbb{R}^{n}$ is a metric space by 
$d(x,y) = \left(\sum\limits_{i=1}^{n}(x_i- y_i)^2\right)^{1/2}$, where $x = (x_1, \ldots, x_n)$ and $y = (y_1, \ldots, y_n)$. 

\yossub{Continuous Mappings}

\begin{definitionu} Let $f: X \to Y$ be a mapping defined on a topological space $X$ into a topological space $Y$. $f$ is called ccontinuous at a point $x_0 \in X$ if to every neighbourhood $U$ of $f(x_0)$ there corresponds a neighbourhood $V$ of $x_0$ such that $f(V) \subseteq U$. The mapping $f$ is said to be continuous if it is continuous at every point of its domain $D(f) = X$. 
\end{definitionu}

\begin{theoremu} Let $X$, $Y$ be topological spaces and $f$ a mapping defined on $X$ into $Y$. Then $f$ is continuous iff the inverse image under $f$ of every open set of $Y$ is an open set of $X$. 
\end{theoremu}

\begin{proof}
 If $f$ is continuous and $U$ an open set of $Y$, then $V = f^{-1}(U)$ is a neighbourhood of every point $x_0 \in X$ such that $f(x_0) \in U$, that is, $V$ is a neighbourhood of every point $x_0$ of $V$. Thus $V$ is an open set of $X$. Let, conversely, for every open set $U \ni f(x_0)$ of $Y$, the set $V = f^{-1}(U)$ be an open set of $X$. Then, by the definition, $f$ is continuous at $x_0 \in X$. 
\end{proof}

\yossub{Compactness}

\begin{definitionu} $A$ system of sets $G_{\alpha}, \alpha \in A$, is called a covering of the set $X$ if $X$ is contained as a subset of the union $\cup_{\alpha \in A} G_{\alpha}$. A subset $M$ of a topological space $X$ is called compact if every system of open sets of $X$ which covers $M$ contains a finite subsystem also covering $M$. 
\end{definitionu}

In view of the preceding theorem, a continuous image of a compact set is also compact. 

\begin{proposition}\label{prop:0.2.Proposition.1} Compact subsets of a topological space are necessarily closed. 
\end{proposition}

\begin{proof}
 Let there be an accumulation point $x_0$ of a compact set $M$ of a topological space $X$ such that $x_0 \in M$. By Hausdorff's axiom of separation, there exist, for any point $m \in M$, disjoint open sets $G_{m,x_0}$ and $G_{x_0,m}$ of $X$ such that $m \in G_{m,x_0}$, $x_0 \in G_{x_0,m}$. The system $\{G_{m,x_0}; m \in M\}$ surely covers $M$. By the compactness of $M$, there exists a finite subsystem $\{G_{m_i,x_0}; i = 1,2,\ldots,n\}$ which covers $M$. Then $\bigcap\limits_{i=1}^{n} G_{x_0,m_i}$ does not intersect $M$. But, since $x_0$ is an accumulation point of $M$, the open set $\bigcap\limits_{i=1}^{n} G_{x_0,m_i} \ni x_0$ must contain a point $m \in M$ distinct from $x_0$. This is a contradiction, and $M$ must be closed. 
\end{proof}

\begin{proposition}\label{prop:0.2.Proposition.2} $A$ closed subset $M_{1}$ of a compact set $M$ of a topological space $X$ is compact. 
\end{proposition}

\begin{proof}
 Let $\{G_{\alpha}\}$ be any system of open sets of $X$ which covers $M_{1}$. $M_{1}$ being closed, $M_{1}^{C}=X-M_{1}$ is an open set of $X$. Since $M_{1}\subseteq M$, the system of open sets $\{G_{\alpha}\}$ plus $M_{1}^{C}$ covers $M$, and since $M$ is compact, a properly chosen finite subsystem $\{G_{\alpha_{i}};i=1,2,\ldots,n\}$ plus $M_{1}^{C}$ surely covers $M$. Thus $\{G_{\alpha_{i}};i=1,2,\ldots,n\}$ covers $M_{1}$. 
\end{proof}

\begin{definitionu} $A$ subset of a topological space is called relatively compact if its closure is compact. A topological space is said to be locally compact if each point of the space has a compact neighbourhood. 
\end{definitionu}

\begin{theoremu} Any locally compact space $X$ can be embedded in another compact space $Y$, having just one more point than $X$, in such a way that the relative topology of $X$ as a subset of $Y$ is just the original topology of $X$. This $Y$ is called a one point compactification of $X$. 
\end{theoremu}

\begin{proof}
 Let $y$ be any element distinct from the points of $X$. Let $\{U\}$ be the class of all open sets in $X$ such that $U^{C} = X- U$ is compact. We remark that $X$ itself $\in \{U\}$. Let $Y$ be the set consisting of the points of $X$ and the point $y$. A set in $Y$ will be called open if either \begin{enumerate}[(i)]
\item it does not contain $y$ and is open as a subset of $X$
\item it does contain $y$ and its intersection with $X$ is a member of $\{U\}$.
\end{enumerate}
 It is easy to see that $Y$ thus obtained is a topological space, and that the relative topology of $X$ coincides with its original topology.

Suppose $\{V\}$ be a family of open sets which covers $Y$. Then there must be some member of $\{V\}$ of the form $U_0 \cup \{y\}$, where $U_0 \in \{U\}$. By the definition of $\{U\}$, $U_0^C$ is compact as a subset of $X$. It is covered by the system of sets $V \cap X$ with $V \in \{V\}$. Thus some finite subsystem: $V_1 \cap X, V_2 \cap X, \ldots, V_n \cap X$ covers $U_0^C$. Consequently, $V_1, V_2, \ldots, V_n$ and $U_0 \cup \{y\}$ cover $Y$, proving that $Y$ is compact. 
\end{proof}

\yossub{Tychonov's Theorem}

\begin{definitionu} Corresponding to each $\alpha$ of an index set $A$, let there be given a topological space $X_{\alpha}$. The Cartesian product $\prod\limits_{\alpha \in A} X_{\alpha}$ is, by definition, the set of all functions $f$ with domain $A$ such that $f(\alpha) \in X_{\alpha}$ for every $\alpha \in A$. We shall write $f = \prod\limits_{\alpha \in A} f(\alpha)$ and call $f(\alpha)$ the $\alpha$-th coordinate of $f$. When $A$ consists of integers $(1, 2, \ldots, n)$, $\prod\limits_{k=1}^{n} X_k$ is usually denoted by $X_1 \times X_2 \times \cdots \times X_n$. We introduce a (weak) topology in the product space $\prod\limits_{\alpha \in A} X_{\alpha}$ by calling ``open'' the sets of the form $\prod\limits_{\alpha \in A} G_{\alpha}$, where the open set $G_{\alpha}$ of $X_{\alpha}$ coincides with $X_{\alpha}$ for all but a finite set of $\alpha$. 
\end{definitionu}

\begin{theoremn}[Tychonov's Theorem] The Cartesian product $X = \prod\limits_{\alpha \in A} X_{\alpha}$ of a system of compact topological spaces $X_{\alpha}$ is also compact. 
\end{theoremn}

\begin{remarku} As is well known, a closed bounded set on the real line $\mathbb{R}^{1}$ is compact with respect to the topology defined by the distance $d(x,y)=|x-y|$ (the Bolzano-Weierstrass theorem). By the way, a subset $M$ of a metric space is said to be bounded, if $M$ is contained in some sphere $S(x_{0},r)$ of the space. Tychonov's theorem implies, in particular, that a parallelopiped: 
\[
- \infty < a _{i} \leqq x _{i} \leqq B_{i} < \infty \quad (i = 1, 2, \dots , n)
\]
\end{remarku}

of the $n$-dimensional Euclidean space $\mathbb{R}^{n}$ is compact. From this we see that $\mathbb{R}^{n}$ is locally compact. 

\begin{proof}[Proof of Tychonov's Theorem.]
A system of sets has the finite intersection property if its every finite subsystem has a non-void intersection. It is easy to see, by taking the complement of the open sets of a covering, that a topological space $X$ is compact iff, for every system $\{M_{\alpha}; \alpha \in A\}$ of its closed subsets with finite intersection property, the intersection $\bigcap\limits_{\alpha \in A} M_{\alpha}^{a}$ is non-void.

Let now a system $\{S\}$ of subsets $S$ of $X = \prod\limits_{\alpha \in A} X_{\alpha}$ have the finite intersection property. Let $\{N\}$ be a system of subsets $N$ of $X$ with the following properties: 

\begin{enumerate}[(i)]
\item $\{S\}$ is a subsystem of $\{N\}$, 

\item $\{N\}$ has the finite intersection property, 

\item $\{N\}$ is maximal in the sense that it is not a proper subsystem of other systems having the finite intersection property and containing $\{S\}$ as its subsystem. 

\end{enumerate}

The existence of such a maximal system $\{N\}$ can be proved by Zorn's lemma or transfinite induction. 

For any set $N$ of $\{N\}$ we define the set $N_{\alpha} = \{f(\alpha); f \in N\} \subseteq X_{\alpha}$. We denote then by $\{N_{\alpha}\}$ the system $\{N_{\alpha}; N \in \{N\}\}$. Like $\{N\}$, $\{N_{\alpha}\}$ enjoys the finite intersection property. Thus, by the compactness of $X_{\alpha}$, there exists at least one point $p_{\alpha} \in X_{\alpha}$ such that $p_{\alpha} \in \bigcap\limits_{N \in \{N\}} N_{\alpha}^{a}$. We have to show that the point $p = \prod\limits_{\alpha \in A} p_{\alpha}$ belongs to the set $\bigcap\limits_{N \in \{N\}} N^{a}$. 

But since $p_{\alpha_0}$ belongs to the intersection $\bigcap\limits_{N \in \{N\}} N_{\alpha_0}^a$, any open set $G_{\alpha_0}$ of $X_{\alpha_0^-}$ which contains $p_{\alpha_0}$ intersects every $N_{\alpha_0} \in \{N_{\alpha_0}\}$. Therefore the open set 
\[
G ^ {(\alpha_{0})} = \{x; x = \prod\limits_{\alpha \in A} x _{\alpha} \text{ with } x _ {\alpha_{0}} \in G _ {\alpha_{0}} \}
\]
of $X$ must intersect every $N$ of $\{N\}$. By the maximality condition (iii) of $\{N\}$, $G^{(\alpha_0)}$ must belong to $\{N\}$. Thus the intersection of a finite number of sets of the form $G^{(\alpha_0)}$ with $\alpha_0 \in A$ must also belong to $\{N\}$ and so such a set intersect every set $N \in \{N\}$. Any open set of $X$ containing $p$ being defined as a set containing such an intersection, we see that $p = \prod\limits_{\alpha \in A} p_\alpha$ must belong to the intersection $\bigcap\limits_{N \in \{N\}} N^a$. 
\end{proof}

\yossub{Urysohn's Theorem}

\begin{propositionu} $A$ compact space $X$ is normal in the sense that, for any disjoint closed sets $F_{1}$ and $F_{2}$ of $X$, there exist disjoint open sets $G_{1}$ and $G_{2}$ such that $F_{1} \subseteq G_{1}, F_{2} \subseteq G_{2}$. 
\end{propositionu}

\begin{proof}
 For any pair $(x, y)$ of points such that $x \in F_1$, $y \in F_2$, there exist disjoint open sets $G(x, y)$ and $G(y, x)$ such that $x \in G(x, y)$, $y \in G(y, x)$. $F_2$ being compact as a closed subset of the compact space $X$, we can, for fixed $x$, cover $F_2$ by a finite number of open sets $G(y_1, x)$, $G(y_2, x), \ldots, G(y_{n(x)}, x)$. Set 
\[
G _{x} = \bigcup\limits_{j=1} ^{n(x)} G (y _{j}, x) \quad \text{ and } \quad G (x) = \bigcap\limits_{j=1} ^{n(x)} G (x, y _{j}). \]

Then the disjoint open sets $G_{x}$ and $G(x)$ are such that $F_{2} \subseteq G_{x}, x \in G(x) \cdot F_{1}$ being compact as a closed subset of the compact space $X$, we can cover $F_{1}$ by a finite number of open sets $G(x_{1}), G(x_{2}), \ldots, G(x_{k})$. Then 
\[
G _{1} = \bigcup\limits_{j=1} ^{k} G (x _{j}) \quad \text{ and } \quad G _{2} = \bigcap\limits_{j=1} ^{k} G _ {x _{j}}
\]
satisfy the condition of the proposition. 
\end{proof}

\begin{corollaryu} $A$ compact space $X$ is regular in the sense that, for any non-void open set $G_1'$ of $X$, there exists a non-void open set $G_2'$ such that $(G_2')^a \subseteq G_1'$. 
\end{corollaryu}

\begin{proof}
 Take $F_{1} = (G_{1}')^{C}$ and $F_{2} = \{x\}$ where $x \in G_{1}'$. We can then take for $G_{2}'$ the open set $G_{2}$ obtained in the preceding proposition. 
\end{proof}

\begin{theoremn}[Urysohn's Theorem] Let $A, B$ be disjoint closed sets in a normal space $X$. Then there exists a real-valued continuous function $f(t)$ on $X$ such that 
\[
0 \leqq f (t) \leqq 1 \text{ on } X, \text{ and } f (t) = 0 \text{ on } A, f (t) = 1 \text{ on } B.
\]
\end{theoremn}

\begin{proof}
 We assign to each rational number $r = k / 2^n$ ( $k = 0, 1, \ldots, 2^n$), an open set $G(r)$ such that \begin{enumerate}[(i)]
\item $A \subseteq G(0)$, $B = G(1)^c$
\item $G(r)^a \subseteq G(r')$ whenever $r < r'$.
\end{enumerate}
 The proof is obtained by induction with respect to $n$.

For $n = 0$, there exist, by the normality of the space $X$, disjoint open sets $G_0$ and $G_1$ with $A \subseteq G_0$, $B \subseteq G_1$. We have only to set $G_0 = G(0)$. Suppose that $G(r)'s$ have been constructed for $r$ of the form $k/2^{n-1}$ in such a way that condition (ii) is satisfied. Next let $k$ be an odd integer $> 0$. Then, since $(k-1)/2^n$ and $(k+1)/2^n$ are of the form $k'/2^{n-1}$ with $0 \leqq k' \leqq 2^{n-1}$, we have $G((k-1)/2^n)^a \subseteq G((k+1)/2^n)$. Hence, by the normality of the space $X$, there exists an open set $G$ which satisfies $G((k-1)/2^n)^a \subseteq G$, $G^a \subseteq G((k+1)/2^n)$. If we set $G(k/2^n) = G$, the induction is completed. 

Define $f(t)$ by 
\[
f (t) = 0 \text{ on } G (0), \text{ and } f (t) = \sup\limits _{t \in G (r)} r \text{ whenever } t \in G (0) ^{C}.
\]

Then, by (i), $f(t) = 0$ on $A$ and $f(t) = 1$ on $B$. We have to prove the continuity of $f$. For any $t_0 \in X$ and positive integer $n$, we take $r$ with $f(t_0) < r < f(t_0) + 2^{-n-1}$. Set $G = G(r) \cap G(r- 2^{-n})^{aC}$ (we set, for convention, $G(s) = \emptyset$ if $s < 0$ and $G(s) = X$ if $s > 1$). The open set $G$ contains $t_0$. For, $f(t_0) < r$ implies $t_0 \in G(r)$, and $(r- 2^{-n-1}) < f(t_0)$ implies $t_0 \in G(r- 2^{-n-1})^C \subseteq G(r- 2^{-n})^{aC}$. Now $t \in G$ implies $t \in G(r)$ and so $f(t) \leqq r$; similarly $t \in G$ implies $t \in G(r- 2^{-n})^{aC} \subseteq G(r- 2^{-n})^C$ so that $r- 2^{-n} \leqq f(t)$. Therefore we have proved that $|f(t)- f(t_0)| \leqq 1/2^n$ whenever $t \in G$. 
\end{proof}
\yossub{The Stone-Weierstrass Theorem}

\begin{theoremn}[Weierstrass' Polynomial Approximation Theorem] Let $f(x)$ be a real-valued (or complex-valued) continuous function on the closed interval [$0$, $1$]. Then there exists a sequence of polynomials $P_{n}(x)$ which converges, as $n \to \infty$, to $f(x)$ uniformly on [$0$, $1$]. According to S. \textsc{Bernstein}, we may take 
\begin{equation}\label{eq:0.2.1}
P _{n} (x) = \sum\limits_{p=0} ^{n} n C _{p} f (p / n) x ^{p} (1- x) ^{n-p}.
\end{equation}
\end{theoremn}

\begin{proof}
 Differentiating $(x + y)^n = \sum\limits_{p=0}^{n} {}_n C_p x^p y^{n-p}$ with respect to $x$ and multiplying by $x$, we obtain $nx(x + y)^{n-1} = \sum\limits_{p=0}^{n} p {}_n C_p x^p y^{n-p}$. Similarly, by differentiating the first expression twice with respect to $x$ and multiplying by $x^2$, we obtain $n(n-1)x^2(x + y)^{n-2} = \sum\limits_{p=0}^{n} p(p-1) {}_n C_p x^p y^{n-p}$. Thus, if we set 
\begin{equation}\label{eq:0.2.2}
r _{p} (x) = _{n} C _{p} x ^{p} (1- x) ^{n-p},
\end{equation}
we have 
\begin{equation}\label{eq:0.2.3}
\sum\limits_{p=0} ^{n} r _{p} (x) = 1, \sum\limits_{p=0} ^{n} p r _{p} (x) = n x, \sum\limits_{p=0} ^{n} p (p- 1) r _{p} (x) = n (n- 1) x ^{2}.
\end{equation}

Hence 
\begin{equation}\label{eq:0.2.4}
\begin{aligned} \sum\limits_{p=0} ^{n} (p- n x) ^{2} r _{p} (x) & = n ^{2} x ^{2} \sum\limits_{p=0} ^{n} r _{p} (x)- 2 n x \sum\limits_{p=0} ^{n} p r _{p} (x) + \sum\limits_{p=0} ^{n} p ^{2} r _{p} (x) \\ & = n ^{2} x ^{2}- 2 n x \cdot n x + (n x + n (n- 1) x ^{2}) \\ & = n x (1- x). \end{aligned}
\end{equation}

We may assume that $|f(x)| \leqq M < \infty$ on [$0$, $1$]. By the uniform continuity of $f(x)$, there exists, for any $\varepsilon > 0$, a $\delta > 0$ such that 
\begin{equation}\label{eq:0.2.5}
\left| f (x)- f \left(x ^{\prime}\right) \right| < \varepsilon \text{ whenever } \left| x- x ^{\prime} \right| < \delta .
\end{equation}

We have, by \eqref{eq:0.2.3}, 
\[
\begin{aligned}
&\left| f (x)- \sum\limits_{p=0} ^{n} f (p / n) r _{p} (x) \right| = \left| \sum\limits_{p=0} ^{n} (f (x)- f (p / n)) r _{p} (x) \right| \leqq \left| \sum\limits_{| p- n x | \leqq \delta n} \right| + \left| \sum\limits_{| p- n x | > \delta n} \right|.
\end{aligned}
\]

For the first term on the right, we have, by $r_{p}(x) \geqq 0$ and \eqref{eq:0.2.3}, 
\[
\left| \sum\limits_{| p- n x | \leqq \delta n} \right| \leqq \varepsilon \sum\limits_{p=0} ^{n} r _{p} (x) = \varepsilon .
\]

For the second term on the right, we have, by \eqref{eq:0.2.4} and $|f(x)| \leqq M$, 
\[ \begin{aligned}\left| \sum\limits_{| p- n x | > \delta n} \right|&\leqq 2 M \sum\limits_{| p- n x | > \delta n} r _{p} (x) \leqq \frac {2 M}{n ^{2} \delta^{2}} \sum\limits_{p=0} ^{n} (p- n x) ^{2} r _{p} (x)\\&= \frac {2 M x (1- x)}{n \delta^{2}} \leqq \frac {M}{2 \delta^{2} n} \rightarrow 0 (\text{ as } n \rightarrow \infty).\end{aligned}
\]
\end{proof}
\begin{theoremn}[The Stone-Weierstrass Theorem] Let $X$ be a compact space and $C(X)$ the totality of real-valued continuous functions on $X$. Let a subset $B$ of $C(X)$ satisfy the three conditions: \begin{enumerate}[(i)]
\item if $f, g \in B$, then the function product $f \cdot g$ and linear combinations $\alpha f + \beta g$, with real coefficients $\alpha, \beta$, belong to $B$,
\item the constant function $1$ belongs to $B$
\item the uniform limit $f_{\infty}$ of any sequence $\{f_n\}$ of functions $\in B$ also belongs to $B$.
\end{enumerate}
 Then $B = C(X)$ iff $B$ separates the points of $X$, i.e. iff, for every pair $(s_1, s_2)$ of distinct points of $X$, there exists a function $x$ in $B$ which satisfies $x(s_1) \neq x(s_2)$. 
\end{theoremn}

\begin{proof}
 The necessity is clear, since a compact space is normal and so, by Urysohn's theorem, there exists a real-valued continuous function $x$ such that $x(s_1) \neq x(s_2)$. 

To prove the sufficiency, we introduce the lattice notations: 
\[
(f \vee g) (s) = \max (f (s), g (s)), (f \wedge g) (s) = \min (f (s), g (s)), | f | (s) = | f (s) |.
\]

By the preceding theorem, there is a sequence $\{P_n\}$ of polynomials such that 
\[
\left| \left| t \right|- P _{n} (t) \right| < 1 / n \quad \text{ for } \quad- n \leqq t \leqq n.
\]

Hence $\left|\left|f(s)\right|- P_n(f(s))\right| < 1/n$ if $-n \leqq f(s) \leqq n$. This proves, by (iii), that $|f| \in B$ if $f \in B$, because any function $f(s) \in B \subseteq C(X)$ is bounded on the compact space $X$. Thus, by 
\[
f \vee g = \frac {f + g}{2} + \frac {| f- g |}{2} \quad \text{ and } \quad f \wedge g = \frac {f + g}{2}- \frac {| f- g |}{2},
\]
we see that $B$ is closed under the lattice operations $\vee$ and $\wedge$. 

Let $h \in C(X)$ and $s_1, s_2 \in X$ be arbitrarily given such that $s_1 \neq s_2$. Then we can find an $f_{s_1s_2} \in B$ with $f_{s_1s_2}(s_1) = h(s_1)$ and $f_{s_1s_2}(s_2) = h(s_2)$. To see this, let $g \in B$ be such that $g(s_1) \neq g(s_2)$, and take real numbers $\alpha$ and $\beta$ so that $f_{s_1s_2} = \alpha g + \beta$ satisfies the conditions: $f_{s_1s_2}(s_1) = h(s_1)$ and $f_{s_1s_2}(s_2) = h(s_2)$. 

Given $\varepsilon > 0$ and a point $t \in X$. Then, for each $s \in X$, there is a neighbourhood $U(s)$ of $s$ such that $f_{st}(u) > h(u)- \varepsilon$ whenever $u \in U(s)$. Let $U(s_1), U(s_2), \ldots, U(s_n)$ cover the compact space $X$ and define 
\[
f _{t} = f _ {s _{1} t} \vee \dots \vee f _ {s _{n} t}.
\]

Then $f_{t} \in B$ and $f_{t}(u) > h(u)- \varepsilon$ for all $u \in X$. We have, by $f_{s,t}(t) = h(t)$, $f_{t}(t) = h(t)$. Hence there is a neighbourhood $V(t)$ of $t$ such that $f_{t}(u) < h(u) + \varepsilon$ whenever $u \in V(t)$. Let $V(t_{1}), V(t_{2}), \ldots, V(t_{k})$ cover the compact space $X$, and define 
\[
f = f _ {t _{1}} \wedge \dots \wedge f _ {t _{k}}.
\]

Then $f \in B$ and $f(u) > h(u)- \varepsilon$ for all $u \in X$, because $f_{t_{j}}(u) > h(u)- \varepsilon$ for $u \in X$. Moreover, we have, for an arbitrary point $u \in X$, say $u \in V(t_i)$, $f(u) \leqq f_{t_{i}}(u) < h(u) + \varepsilon$. 

Therefore we have proved that $|f(u)- h(u)| < \varepsilon$ on $X$. 

We have incidentally proved the following two corollaries. 
\end{proof}

\begin{corollary}[\textsc{Kakutani}-\textsc{Krein}]\label{cor:0.2.Corollary.1} Let $X$ be a compact space and $C(X)$ the totality of real-valued continuous functions on $X$. Let a subset $B$ of $C(X)$ satisfy the conditions: \begin{enumerate}[(i)]
\item if $f, g \in B$, then $f \vee g, f \wedge g$ and the linear combinations $\alpha f + \beta g$, with real coefficients $\alpha, \beta$, belong to $B$,
\item the constant function $1$ belongs to $B$
\item the uniform limit $f_{\infty}$ of any sequence $\{f_n\}$ of functions $\in B$ also belongs to $B$.
\end{enumerate}
 Then $B = C(X)$ iff $B$ separates the points of $X$. 
\end{corollary}

\begin{corollary}\label{cor:0.2.Corollary.2} Let $X$ be a compact space and $C(X)$ be the totality of complex-valued continuous functions on $X$. Let a subset $B$ of $C(X)$ satisfy the conditions: \begin{enumerate}[(i)]
\item if $f, g \in B$, then the function product $f \cdot g$ and the linear combinations $\alpha f + \beta g$, with complex coefficients $\alpha, \beta$, belong to $B$,
\item the constant function $1$ belongs to $B$
\item the uniform limit $f_{\infty}$ of any sequence $\{f_n\}$ of functions $\in B$ also belongs to $B$.
\end{enumerate}
 Then $B = C(X)$ iff $B$ satisfies the conditions: \begin{enumerate}[(i)]
\setcounter{enumi}{3}
\item $B$ separates points of $X$
\item if $f(s) \in B$, then its complex conjugate function $\bar{f}(s)$ also belongs to $B$.
\end{enumerate}
\end{corollary}

\begin{theoremn}[Weierstrass' Trigonometric Approximation Theorem] Let $X$ be the circumference of the unit circle of $\mathbb{R}^{2}$. It is a compact space by the usual topology, and a complex-valued continuous function on $X$ is represented by a continuous function $f(x),-\infty < x < \infty$, of period $2\pi$. If we take, in the above Corollary~\ref{cor:0.2.Corollary.2}, for $B$ the set of all functions representable by linear combinations, with complex coefficients of the trigonometric functions 
\[
e ^{inx} (n = 0, \pm 1, \pm 2, \dots)
\]

and by those functions obtainable as the uniform limit of such linear combinations, we obtain Weierstrass' trigonometric approximation theorem: Any complex-valued continuous function $f(x)$ with period $2\pi$ can be approximated uniformly by a sequence of trigonometric polynomials of the form $\sum\limits_{n} c_n e^{inx}$. 
\end{theoremn}

\yossub{Completeness}

A sequence $\{x_{n}\}$ of elements in a metric space $X$ converges to a limit point $x\in X$ iff $\lim\limits_{n\to \infty}d(x_n,x) = 0$. By the triangle inequality $\tilde{d}(x_n,x_m)\leqq d(x_n,x) + d(x,x_m)$, we see that a convergent sequence $\{x_{n}\}$ in $X$ satisfies Cauchy's convergence condition 
\begin{equation}\label{eq:0.2.1a}
\lim\limits _{n, m \rightarrow \infty} d \left(x _{n}, x _{m}\right) = 0.
\end{equation}

\begin{definitionu} Any sequence $\{x_{n}\}$ in a metric space $X$ satisfying the above condition \eqref{eq:0.2.1a} is called a Cauchy sequence. A metric space $X$ is said to be complete if every Cauchy sequence in it converges to a limit point $\in X$. 
\end{definitionu}

It is easy to see, by the triangle inequality, that the limit point of $\{x_{n}\}$, if it exists, is uniquely determined. 

\begin{definitionu} $A$ subset $M$ of a topological space $X$ is said to be non-dense in $X$ if the closure $M^{a}$ does not contain a non-void open set of $X$. $M$ is called dense in $X$ if $M^{a} = X$. $M$ is said to be of the first category if $M$ is expressible as the union of a countable number of sets each of which is non-dense in $X$; otherwise $M$ is said to be of the second category. 
\end{definitionu}

\yossub{Baire's Category Argument}

\begin{theoremn}[The Baire-Hausdorff Theorem] $A$ non-void complete metric space is of the second category. 
\end{theoremn}

\begin{proof}
 Let $\{M_{n}\}$ be a sequence of closed sets whose union is a complete metric space $X$. Assuming that no $M_{n}$ contains a non-void open set, we shall derive a contradiction. Thus $M_{1}^{C}$ is open and $M_{1}^{Ca}=X$, hence $M_{1}^{C}$ contains a closed sphere $S_{1}=\{x;d(x_{1},x)\leqq r_{1}\}$ whose centre $x_{1}$ may be taken arbitrarily near to any point of $X$. We may assume that $0<r_{1}<1/2$. By the same argument, the open set $M_{2}^{C}$ contains a closed sphere 

$S_{2} = \{x; d(x_{2}, x) \leqq r_{2}\}$ contained in $S_{1}$ and such that $0 < r_{2} < 1/2^{2}$. By repeating the same argument, we obtain a sequence $\{S_{n}\}$ of closed spheres $S_{n} = \{x; d(x_{n}, x) \leqq r_{n}\}$ with the properties: 
\[
0 < r _{n} < 1 / 2 ^{n}, S _{n+1} \subseteq S _{n}, S _{n} \cap M _{n} = \emptyset \quad (n = 1, 2, \dots). \]

The sequence $\{x_{n}\}$ of the centres forms a Cauchy sequence, since, for any $n < m$, $x_{m} \in S_{n}$ so that $d(x_{n}, x_{m}) \leqq r_{n} < 1/2^{n}$. Let $x_{\infty} \in X$ be the limit point of this sequence $\{x_{n}\}$. The completeness of $X$ guarantees the existence of such a limit point $x_{\infty}$. By $d(x_{n}, x_{\infty}) \leqq d(x_{n}, x_{m}) + d(x_{m}, x_{\infty}) \leqq r_{n} + d(x_{m}, x_{\infty}) \to r_{n}$ (as $m \to \infty$), we see that $x_{\infty} \in S_{n}$ for every $n$. Hence $x_{\infty}$ is in none of the sets $M_{n}$, and hence $x_{\infty}$ is not in the union $\bigcup\limits_{n=1}^{\infty} M_{n} = X$, contrary to $x_{\infty} \in X$. 
\end{proof}

\begin{theoremn}[Baire's Theorem 1] Let $M$ be a set of the first category in a compact topological space $X$. Then the complement $M^C = X- M$ is dense in $X$. 
\end{theoremn}

\begin{proof}
 We have to show that, for any non-void open set $G$, $M^C$ intersects $G$. Let $M = \bigcup\limits_{n=1}^{\infty} M_n$ where each $M_n$ is a non-dense closed set. Since $M_1 = M_1^a$ is non-dense, the open set $M_1^C$ intersects $G$. Since $X$ is regular as a compact space, there exists a non-void open set $G_1$ such that $G_1^a \subseteq G \cap M_1^C$. Similarly, we can choose a non-void open set $G_2$ such that $G_2^a \subseteq G_1 \cap M_2^C$. Repeating the process, we obtain a sequence of non-void open sets $\{G_n\}$ such that 
\[
G _{n+1} ^{a} \subseteq G _{n} \cap M _{n+1} ^{C} \quad (n = 1, 2, \dots). \]
The sequence of closed sets $\{G_n^a\}$ enjoys, by the monotony in $n$, the finite intersection property. Since $X$ is compact, there is an $x \in X$ such that $x \in \bigcap\limits_{n=1}^{\infty} G_n^a$. $x \in G_1^a$ implies $x \in G$, and from $x \in G_{n+1}^a \subseteq G_n \cap M_{n+1}^C$ ( $n = 0, 1, 2, \ldots; G_0 = G$), we obtain $x \in \bigcap\limits_{n=1}^{\infty} M_n^c = M^c$. Therefore we have proved that $G \cap M^c$ is non-void. 
\end{proof}

\begin{theoremn}[Baire's Theorem 2] Let $\{x_{n}(t)\}$ be a sequence of real-valued continuous functions defined on a topological space $X$. Suppose that a finite limit: 
\[
\lim\limits _{n \rightarrow \infty} x _{n} (t) = x (t)
\]

exists at every point $t$ of $X$. Then the set of points at which the function $x$ is discontinuous constitutes a set of the first category. 
\end{theoremn}

\begin{proof}
 We denote, for any set $M$ of $X$, by $M^{i}$ the union of all the open sets contained in $M$; $M^{i}$ will be called the interior of $M$. 

 Put $P_{m} (\varepsilon) = \{t \in X; | x (t)- x _{m} (t) | \leqq \varepsilon , \varepsilon > 0 \}, G (\varepsilon) = \bigcup\limits_{m=1} ^{\infty} P _{m} ^{i} (\varepsilon). $ Then we can prove that $C = \bigcap\limits_{n=1}^{\infty} G(1/n)$ coincides with the set of all points at which $x(t)$ is continuous. Suppose $x(t)$ is continuous at $t = t_0$. 

We shall show that $t_0 \in \bigcap\limits_{n=1}^{\infty} G(1/n)$. Since $\lim\limits_{n \to \infty} x_n(t) = x(t)$, there exists an $m$ such that $|x(t_0)- x_m(t_0)| \leqq \varepsilon/3$. By the continuity of $x(t)$ and $x_m(t)$ at $t = t_0$, there exists an open set $U_{t_0} \ni t_0$ such that $|x(t)- x(t_0)| \leqq \varepsilon/3$ $|x_m(t)- x_m(t_0)| \leqq \varepsilon/3$ whenever $t \in U_{t_0}$. Thus $t \in U_{t_0}$ implies $|x(t)- x_m(t)| \leqq |x(t)- x(t_0)| + |x(t_0)- x_m(t_0)| + |x_m(t_0)- x_m(t)| < \varepsilon$, which proves that $t_0 \in P_m^i(\varepsilon)$ and so $t_0 \in G(\varepsilon)$. Since $\varepsilon > 0$ was arbitrary, we must have $t_0 \in \bigcap\limits_{n=1}^{\infty} G(1/n)$. 

Let, conversely, $t_0 \in \bigcap\limits_{n=1}^{\infty} G(1/n)$. Then, for any $\varepsilon > 0$, $t_0 \in G(\varepsilon/3)$ and so there exists an $m$ such that $t_0 \in P_m^i(\varepsilon/3)$. Thus there is an open set $U_{t_0} \ni t_0$ such that $t \in U_{t_0}$ implies $|x(t)- x_m(t)| \leqq \varepsilon/3$. Hence, by the continuity of $x_m(t)$ and the arbitrariness of $\varepsilon > 0$, $x(t)$ must be continuous at $t = t_0$. 

After these preparations, we put 
\[
F _{m} (\varepsilon) = \{t \in X; | x _{m} (t)- x _{m+k} (t) | \leqq \varepsilon \quad (k = 1, 2, \dots) \}.
\]
This is a closed set by the continuity of the $x_{n}(t)$'s. We have $X = \bigcup\limits_{m=1}^{\infty} F_{m}(\varepsilon)$ by $\lim\limits_{n \to \infty} x_{n}(t) = x(t)$. Again by $\lim\limits_{n \to \infty} x_{n}(t) = x(t)$, we have $F_{m}(\varepsilon) \subseteq P_{m}(\varepsilon)$. Thus $F_{m}^{i}(\varepsilon) \subseteq P_{m}^{i}(\varepsilon)$ and so $\bigcup\limits_{m=1}^{\infty} F_{m}^{i}(\varepsilon) \subseteq G(\varepsilon)$. On the other hand, for any closed set $F$, $(F- F^{i})$ is a non-dense closed set. Thus $X- \bigcup\limits_{m=1}^{\infty} F_{m}^{i}(\varepsilon) = \bigcup\limits_{m=1}^{\infty} (F_{m}(\varepsilon)- F_{m}^{i}(\varepsilon))$ is a set of the first category. Thus its subset $G(\varepsilon)^{C} = X- G(\varepsilon)$ is also a set of the first category. Therefore the set of all the points of discontinuity of the function $x(t)$, which is expressible as $X- \bigcap\limits_{n=1}^{\infty} G(1/n) = \bigcup\limits_{n=1}^{\infty} G(1/n)^{C}$, is a set of the first category. 
\end{proof}

\begin{theoremu} $A$ subset $M$ of a complete metric space $X$ is relatively compact iff it is totally bounded in the sense that, for every $\varepsilon > 0$, there exists a finite system of points $m_{1}, m_{2}, \ldots, m_{n}$ of $M$ such that every point $m$ of $M$ has a distance $< \varepsilon$ from at least one of $m_{1}, m_{2}, \ldots, m_{n}$. In other words, $M$ is totally bounded if, for every $\varepsilon > 0$, $M$ can be covered by a finite system of spheres of radii $< \varepsilon$ and centres $\in M$. 
\end{theoremu}

\begin{proof}
 Suppose $M$ is not totally bounded. Then there exist a positive number $\varepsilon$ and an infinite sequence $\{m_{n}\}$ of points $\in M$ such that $d(m_{i}, m_{j}) \geqq \varepsilon$ for $i \neq j$. Then, if we cover the compact set $M^{a}$ by a system of open spheres of radii $< \varepsilon$, no finite subsystem of this system can cover $M^{a}$. For, this subsystem cannot cover the infinite subset $\{m_{i}\} \subseteq M \subseteq M^{a}$. Thus a relatively compact subset of $X$ must be totally bounded. 

Suppose, conversely, that $M$ is totally bounded as a subset of a complete metric space $X$. Then the closure $M^a$ is complete and is totally bounded with $M$. We have to show that $M^a$ is compact. To this purpose, we shall first show that any infinite sequence $\{p_n\}$ of $M^a$ contains a subsequence $\{p_{(n')}\}$ which converges to a point of $M^a$. Because of the total boundedness of $M$, there exist, for any $\varepsilon > 0$, a point $q_e \in M^a$ and a subsequence $\{p_{n'}\}$ of $\{p_n\}$ such that $d(p_{n'}, q_e) < \varepsilon / 2$ for $n = 1, 2, \ldots$; consequently, $d(p_{n'}, p_{m'}) \leqq d(p_{n'}, q_e) + d(q_e, p_{m'}) < \varepsilon$ for $n, m = 1, 2, \ldots$. We set $\varepsilon = 1$ and obtain the sequence $\{p_{i'}\}$, and then apply the same reasoning as above with $\varepsilon = 2^{-1}$ to this sequence $\{p_{i'}\}$. We thus obtain a subsequence $\{p_{n''}\}$ of $\{p_{n'}\}$ such that 
\[
d \left(p _ {n ^{\prime}}, p _ {m ^{\prime}}\right) < 1, d \left(p _ {n ^{\prime \prime}}, p _ {m ^{\prime \prime}}\right) < 1 / 2 \quad (n, m = 1, 2, \dots).
\]

By repeating the process, we obtain a subsequence $\{\phi_{n^{(k+1)}}\}$ of the sequence $\{\phi_{n^{(k)}}\}$ such that 
\[
d \left(\phi_ {n ^{(k+1)}}, \phi_ {m ^{(k+1)}}\right) < 1 / 2 ^{k} \quad (n, m = 1, 2, \dots).
\]

Then the subsequence $\{\phi_{(n)'}\}$ of the original sequence $\{\phi_n\}$, defined by the diagonal method: 
\[
\phi_ {(n) ^{\prime}} = \phi_ {n ^{(n)}},
\]
surely satisfies $\lim\limits_{n,m\to \infty}d(p_{(n)'}p_{(m)'}) = 0$. Hence, by the completeness of $M^a$, there must exist a point $p\in M^a$ such that $\lim\limits_{n\to \infty}d(p_{(n)'}p) = 0$. 

We next show that the set $M^{a}$ is compact. We remark that there exists a countable family $\{F\}$ of open sets $F$ of $X$ such that, if $U$ is any open set of $X$ and $x \in U \cap M^{a}$, there is a set $F \in \{F\}$ for which $x \in F \subseteq U$. This may be proved as follows. $M^{a}$ being totally bounded, $M^{a}$ can be covered, for any $\varepsilon > 0$, by a finite system of open spheres of radii $\varepsilon$ and centres $\in M^{a}$. Letting $\varepsilon = 1, 1/2, 1/3, \ldots$ and collecting the countable family of the corresponding finite systems of open spheres, we obtain the desired family $\{F\}$ of open sets. 

Let now $\{U\}$ be any open covering of $M^a$. Let $\{F^*\}$ be the subfamily of the family $\{F\}$ defined as follows: $F \subseteq \{F^*\}$ iff $F \subseteq \{F\}$ and there is some $U \in \{U\}$ with $F \subseteq U$. By the property of $\{F\}$ and the fact that $\{U\}$ covers $M^a$, we see that this countable family $\{F^*\}$ of open sets covers $M^a$. Now let $\{U^*\}$ be a subfamily of $\{U\}$ obtained by selecting just one $U \in \{U\}$ such that $F \subseteq U$, for each $F \in \{F^*\}$. Then $\{U^*\}$ is a countable family of open sets which covers $M^a$. We have to show that some finite subfamily of $\{U^*\}$ covers $M^a$. Let the sets in $\{U^*\}$ be indexed as $U_1, U_2, \ldots$. Suppose that, for each $n$, the finite union $\bigcup\limits_{j=1}^{n} U_j$ fails to cover $M^a$. Then there is some point $x_n \in \left(M- \bigcup\limits_{k=1}^{n} U_k\right)$. By what was proved above, the sequence $\{x_{n}\}$ contains a subsequence $\{x_{(n)}\}$ which converges to a point, say $x_{\infty}$, in $M^{a}$. Then $x_{\infty} \in U_{N}$ for some index $N$, and so $x_{n} \in U_{N}$ for infinitely many values of $n$, in particular for an $n > N$. This contradicts the fact that $x_{n}$ was chosen so that $x_{n} \in \left(M- \bigcup\limits_{k=1}^{n} U_{k}\right)$. Hence we have proved that $M^{a}$ is compact. 
\end{proof}


\yossection{3}{Measure Spaces}
\label{sec:0.3}

\yossub{Measures}

\begin{definitionu} Let $S$ be a set. A pair $(S, \mathfrak{B})$ is called a \emph{$\sigma$-ring} or a \emph{$\sigma$-additive family} of sets $\subseteq S$ if $\mathfrak{B}$ is a family of subsets of $S$ such that 
\begin{align}
S& \in \mathfrak{B}, \label{eq:0.3.1} \\
B& \in \mathfrak{B} \text{ implies } B ^{C} = (S- B) \in \mathfrak{B}, \label{eq:0.3.2} \\
B_{j}& \in \mathfrak{B} (j = 1, 2, \dots) \text{ implies that } \bigcup\limits_{j=1} ^{\infty} B_{j} \in \mathfrak{B} (\sigma \text{-additivity }). \label{eq:0.3.3}
\end{align}
\end{definitionu}

Let $(S, \mathfrak{B})$ be a $\sigma$-ring of sets $\subseteq S$. Then a triple $(S, \mathfrak{B}, m)$ is called a \emph{measure space} if $m$ is a non-negative, $\sigma$-additive measure defined on $\mathfrak{B}$: 
\begin{equation}
 m(B) \geqq 0 \text{ for every } B \in \mathfrak{B}, \label{eq:0.3.4} \\
\end{equation}
\begin{equation}\label{eq:0.3.5}
\parbox[b]{0.92\linewidth}{\centering
$m\left(\sum\limits_{j=1}^{\infty}B_j\right)=\sum\limits_{j=1}^{\infty}m(B_j)$ for any disjoint sequence $\{B_j\}$ of sets $\in\mathfrak{B}$\par
(\emph{countable- or $\sigma$-additivity of $m$}),}
\end{equation}
\par\medskip\refstepcounter{equation}\label{eq:0.3.6}
\noindent
\parbox[b]{\linewidth}{\setlength{\parindent}{1.2em}\setlength{\parfillskip}{0pt}\indent
$S$ is expressible as a countable union of sets $B_j\in\mathfrak{B}$ such that $m(B_j)<\infty$ $(j=1,2,\ldots)$ (\emph{$\sigma$-finiteness} of the measure space $(S,\mathfrak{B},m)$).\hfill(\theequation)}
\par\medskip
This value $m(B)$ is called the \emph{$m$-measure} of the set $B$. 

\yossub{Measurable Functions}

\begin{definitionu} $A$ real- (or complex-) valued function $x(s)$ defined on $S$ is said to be \emph{$\mathfrak{B}$-measurable} or, in short, \emph{measurable} if the following condition is satisfied: 

\par\medskip\refstepcounter{equation}\label{eq:0.3.7}
\noindent
\parbox[b]{\linewidth}{\setlength{\parindent}{1.2em}\setlength{\parfillskip}{0pt}\indent
For any open set $G$ of the real line $\mathbb{R}^{1}$ (or complex plane $\mathbb{C}^{1}$), the set $\{s; x(s)\in G\}$ belongs to $\mathfrak{B}$.\hfill(\theequation)}
\par\medskip
It is permitted that $x(s)$ takes the value $\infty$.
\end{definitionu}

\begin{definitionu} $A$ property $P$ pertaining to points $s$ of $S$ is said to hold \emph{$m$-almost everywhere} or, in short \emph{$m$-a.e.}, if it holds except for those $s$ which form a set $\in \mathfrak{B}$ of $m$-measure zero. 
\end{definitionu}

A real- (or complex-) valued function $x(s)$ defined $m$-a.e. on $S$ and satisfying condition \eqref{eq:0.3.7} shall be called a $\mathfrak{B}$-measurable function defined $m$-a.e. on $S$ or, in short, a $\mathfrak{B}$-measurable function. 

\begin{theoremn}[Egorov's Theorem] If $B$ is a $\mathfrak{B}$-measurable set with $m(B) < \infty$ and if $\{f_n(s)\}$ is a sequence of $\mathfrak{B}$-measurable functions, finite $m$-a.e. on $B$, that converges $m$-a.e. on $B$ to a finite $\mathfrak{B}$-measurable function $f(s)$, then there exists, for each $\varepsilon > 0$, a subset $E$ of $B$ such that $m(B- E) \leqq \varepsilon$ and on $E$ the convergence of $f_n(s)$ to $f(s)$ is uniform. 
\end{theoremn}

\begin{proof}
 By removing from $B$, if necessary, a set of $m$-measure zero, we may suppose that on $B$, the functions $f_{n}(s)$ are everywhere finite, and converge to $f(s)$ on $B$. 


The set $B_{n} = \bigcap\limits_{k=n+1}^{\infty}\{s\in B; |f(s)- f_{k}(s)| < \varepsilon\}$ is $\mathfrak{B}$-measurable and $B_{n} \subseteq B_{k}$ if $n < k$. Since $\lim\limits_{n\to\infty}f_{n}(s) = f(s)$ on $B$, we have $B = \bigcup\limits_{n=1}^{\infty}B_{n}$. Thus, by the $\sigma$-additivity of the measure $m$, we have 
\[
\begin{aligned}
&m(B)  = m \{B_{1} + (B_{2}- B_{1}) + (B_{3}- B_{2}) + \dots \} = m(B_{1}) + m(B_{2}- B_{1}) + m(B_{3}- B_{2}) + \dots \\
&= m(B_{1}) + (m(B_{2})- m(B_{1})) + (m(B_{3})- m(B_{2})) + \dots = \lim\limits _{n \to \infty} m(B_{n}).
\end{aligned}
\]

Hence $\lim\limits_{n\to \infty}m(B- B_n) = 0$, and therefore, from a sufficiently large $k_{0}$ on, $m(B- B_k) < \eta$ where $\eta$ is any given positive number. 

Thus there exist, for any positive integer $k$, a set $C_k \subseteq B$ such that $m(C_k) \leqq \varepsilon / 2^k$ and an index $N_k$ such that 
\[
\left| f (s)- f _{n} (s) \right| < 1 / 2 ^{k} \text{ for } n > N _{k} \text{ and for } s \in B- C _{k}.
\]

Let us set $E = B- \bigcup\limits_{k=1}^{\infty} C_k$. Then we find 
\[
m(B- E) \leqq \sum\limits_{k=1} ^{\infty} m(C _{k}) \leqq \sum\limits_{k=1} ^{\infty} \varepsilon / 2 ^{k} = \varepsilon ,
\]
and the sequence $f_{n}(s)$ converges uniformly on $E$.
\end{proof}

\yossub{Integrals}

\begin{definitionu} $A$ real- (or complex-) valued function $x(s)$ defined on $S$ is said to be finitely-valued if it is a finite non-zero constant on each of a finite number, say $n$, of disjoint $\mathfrak{B}$-measurable sets $B_j$ and equal to zero on $S- \bigcup\limits_{j=1}^{n} B_j$. Let the value of $x(s)$ on $B_j$ be denoted by $x_j$. Then $x(s)$ is $m$-integrable or, in short, integrable over $S$ if $\sum\limits_{j=1}^{n} |x_j| m(B_j) < \infty$, and the value $\sum\limits_{j=1}^{n} x_j m(B_j)$ is defined as the integral of $x(s)$ over $S$ with respect to the measure $m$; the integral will be denoted by $\int_S x(s) m(\dd s)$ 
or, in short, by $\int_{S} x(s)$ or simply by $\int x(s)$ when no confusion can be expected. A real- (or complex-) valued function $x(s)$ defined $m$-a.e. on $S$ is said to be $m$-integrable or, in short, integrable over $S$ if there exists a sequence $\{x_n(s)\}$ of finitely-valued integrable functions converging to $x(s)$ $m$-a.e. and such that 
\[
\lim\limits _{n, k \rightarrow \infty} \int_{S} | x _{n} (s)- x _{k} (s) | m(\dd s) = 0.
\]
\end{definitionu}

It is then proved that a finite $\lim\limits_{n\to\infty}\int_{S}x_{n}(s)m(\dd s)$ exists and the value of this limit is independent of the choice of the approximating sequence $\{x_{n}(s)\}$. The value of the integral $\int_{S}x(s)m(\dd s)$ over $S$ with respect to the measure $m$ is, by definition, given by $\lim\limits_{n\to\infty}\int_{S}x_{n}(s)m(\dd s)$. We shall sometimes abbreviate the notation $\int_{S}x(s)m(\dd s)$ to $\int x(s)m(\dd s)$ or to $\int x(s)$. 

\yossub{Properties of the Integral}

\begin{enumerate}[i)]
 \item If $x(s)$ and $y(s)$ are integrable, then $\alpha x(s) + \beta y(s)$ is integrable and $\int_{S} (\alpha x(s) + \beta y(s)) m(\dd s) = \alpha \int_{S} x(s) m(\dd s) + \beta \int_{S} y(s) m(\dd s).$ 
 \item $x(s)$ is integrable iff $|x(s)|$ is integrable. 
 \item If $x(s)$ is integrable and $x(s) \geqq 0$ a.e., then $\int_{S} x(s) m(\dd s) \geqq 0$, and the equality sign holds iff $x(s) = 0$ a.e. 
 \item If $x(s)$ is integrable, then the function $X(B) = \int_{B} x(s) m(\dd s)$ is $\sigma$-additive, that is, $X\left(\sum\limits_{j=1}^{\infty} B_j\right) = \sum\limits_{j=1}^{\infty} X(B_j)$ for any disjoint sequence $\{B_j\}$ of sets $\in \mathfrak{B}$. Here $\int_{B} x(s) m(\dd s) = \int_{S} C_B(s) x(s) m(\dd s)$, where $C_B(s)$ is the defining function of the set $B$, that is, 
 \[
 C _{B} (s) = 1 \text{ for } s \in B \text{ and } C _{B} (s) = 0 \text{ for } s \in S- B.
 \]
 \item $X(B)$ in iv) is absolutely continuous with respect to $m$ in the sense that $m(B) = 0$ implies $X(B) = 0$. This condition is equivalent to the condition that $\lim\limits_{m(B) \to 0} X(B) = 0$ uniformly in $B \in \mathfrak{B}$. 
\end{enumerate}

\begin{lemman}[The Lebesgue-Fatou Lemma] Let $\{x_{n}(s)\}$ be a sequence of real-valued integrable functions. If there exists a real-valued integrable function $x(s)$ such that $x(s) \geqq x_{n}(s)$ a.e. for $n = 1, 2, \ldots$ (or $x(s) \leqq x_{n}(s)$ a.e. for $n = 1, 2, \ldots$), then 
\begin{gather*}\int_{S} \left(\varlimsup\limits_{n \rightarrow \infty} x _{n} (s)\right) m(d s) \geqq \varlimsup\limits_{n \rightarrow \infty} \int_{S} x _{n} (s) m(d s) \\\left(\text{ or } \int_{S} \left(\varliminf\limits_{n \to \infty} x _{n} (s)\right) m(d s) \leqq \varliminf\limits_{n \to \infty} \int_{S} x _{n} (s) m(d s)\right), 
\end{gather*}

under the convention that if $\varlimsup\limits_{n\to\infty}x_n(s)$ (or $\varliminf\limits_{n\to\infty}x_n(s)$) is not integrable, we understand that $\varlimsup\limits_{n\to\infty}\int_{S}x_n(s)m(\dd s)=-\infty$ (or $\varliminf\limits_{n\to\infty}\int_{S}x_n(s)m(\dd s)=\infty$). 
\end{lemman}

\begin{definitionu} Let $(S, \mathfrak{B}, m)$ and $(S', \mathfrak{B}', m')$ be two measure spaces. We denote by $B \times B'$ the smallest $\sigma$-ring of subsets of $S \times S'$ which contains all the sets of the form $B \times B'$, where $B \in B$, $B' \in B'$. It is proved that there exists a uniquely determined $\sigma$-finite, $\sigma$-additive and non-negative measure $m \times m'$ defined on $B \times B'$ such that 
\[
(m \times m ^{\prime}) (B \times B ^{\prime}) = m(B) m ^{\prime} (B ^{\prime}).
\]
$m \times m'$ is called the product measure of $m$ and $m'$. We may define the $\mathfrak{B} \times \mathfrak{B}'$-measurable functions $x(s, s')$ defined on $S \times S'$, and the $m \times m'$-integrable functions $x(s, s')$. The value of the integral over $S \times S'$ of an $m \times m'$-integrable function $x(s, s')$ will be denoted by 
\[
\iint_ {S \times S ^{\prime}} x (s, s ^{\prime}) (m \times m ^{\prime}) (\dd s \dd s ^{\prime}) \quad \text{ or } \quad \iint_ {S \times S ^{\prime}} x (s, s ^{\prime}) m(\dd s) m ^{\prime} (\dd s ^{\prime}).
\]
\end{definitionu}

\begin{theoremn}[The Fubini-Tonelli Theorem] $A$ $\mathfrak{B} \times \mathfrak{B}'$-measurable function $x(s, s')$ is $m \times m'$-integrable over $S \times S'$ iff at least one of the iterated integrals 
\[
\int_ {S ^{\prime}} \left\{\int_{S} | x (s, s ^{\prime}) | m(\dd s) \right\} m ^{\prime} (\dd s ^{\prime}) \quad \text{ and } \quad \int_{S} \left\{\int_ {S ^{\prime}} | x (s, s ^{\prime}) | m ^{\prime} (\dd s ^{\prime}) \right\} m(\dd s)
\]

is finite; and in such a case we have 
\end{theoremn}
\[
\begin{aligned} \iint_ {S \times S ^{\prime}} x (s, s ^{\prime}) m(\dd s) m ^{\prime} (\dd s ^{\prime}) & = \int_ {S ^{\prime}} \left\{\int_{S} x (s, s ^{\prime}) m(\dd s) \right\} m ^{\prime} (\dd s ^{\prime}) \\ & = \int_{S} \left\{\int_ {S ^{\prime}} x (s, s ^{\prime}) m ^{\prime} (\dd s ^{\prime}) \right\} m(\dd s). \end{aligned}
\]

\yossub{Topological Measures}

\begin{definitionu} Let $S$ be a locally compact space, e.g., an $n$-dimensional Euclidean space $\mathbb{R}^{n}$ or a closed subset of $\mathbb{R}^{n}$. The Baire subsets of $S$ are the members of the smallest $\sigma$-ring of subsets of $S$ which contains every compact $G_{\delta}$-set, i.e., every compact set of $S$ which is the intersection of a countable number of open sets of $S$. The Borel subsets of $S$ are the members of the smallest $\sigma$-ring of subsets of $S$ which contains every compact set of $S$. 
\end{definitionu}

If $S$ is a closed subset of a Euclidean space $\mathbb{R}^{n}$, the Baire and the Borel subsets of $S$ coincide, because in $\mathbb{R}^{n}$ every compact (closed bounded) set is a $G_\delta$-set. If, in particular, $S$ is a real line $\mathbb{R}^{1}$ or a closed interval on $\mathbb{R}^{1}$, the Baire (= Borel) subsets of $S$ may also be defined as the members of the smallset $\sigma$-ring of subsets of $S$ which contains half open intervals $(a, b]$. 

\begin{definitionu} Let $S$ be a locally compact space. Then a non-negative Baire (Borel) measure on $S$ is a $\sigma$-additive measure defined for every Baire (Borel) subset of $S$ such that the measure of every compact set is finite. The Borel measure $m$ is called regular if for each Borel set $B$ we have 
\[
m(B) = \inf\limits _{U \supseteq B} m(U)
\]
where the infimum is taken over all open sets $U$ containing $B$. We may also define the regularity for Baire measures in a similar way, but it turns out that a Baire measure is always regular. It is also proved that each Baire measure has a uniquely determined extension to a regular Borel measure. Thus we shall discuss only Baire measures. 
\end{definitionu}

\begin{definitionu} $A$ complex-valued function $f(s)$ defined on a locally compact space $S$ is a Baire function on $S$ if $f^{-1}(B)$ is a Baire set of $S$ for every Baire set $B$ in the complex plane $\mathbb{C}^{1}$. Every continuous function is a Baire function if $S$ is a countable union of compact sets. A Baire function is measurable with respect to the $\sigma$-ring of all Baire sets of $S$. 
\end{definitionu}

\yossub{The Lebesgue Measure}

\begin{definitionu} Suppose $S$ is the real line $\mathbb{R}^{1}$ or a closed interval of $\mathbb{R}^{1}$. Let $F(x)$ be a monotone non-decreasing function on $S$ which is continuous from the right: $F(x) = \inf\limits_{x<y} F(y)$. Define a function $m$ on half closed intervals $(a, b]$ by $m((a, b]) = F(b)- F(a)$. This $m$ has a uniquely determined extension to a non-negative Baire measure on $S$. The extended measure $m$ is finite, i.e., $m(S) < \infty$ iff $F$ is bounded. If $m$ is the Baire measure induced by the function $F(s) = s$, then $m$ is called the Lebesgue measure. The Lebesgue measure in $\mathbb{R}^{n}$ is obtained from the $n$-tuple of the one-dimensional Lebesgue measures through the process of forming the product measure. 
\end{definitionu}

Concerning the Lebesgue measure and the corresponding Lebesgue integral, we have the following two important theorems: 

\begin{theorem}\label{thm:0.3.Theorem.1} Let $M$ be a Baire set in $\mathbb{R}^{n}$ whose Lebesgue measure $|M|$ is finite. Then, if we denote by $B \ominus C$ the symmetric difference of the set $B$ and $C: B \ominus C = B \cup C- B \cap C$, we have 
\[\lim\limits_{|h|\to 0}|(M + h)\ominus M| = 0, \text{ where } M + h = \{x\in \mathbb{R}^{n}; x = m + h,m\in M\}.\] 
Here $m + h = (m_{1} + h_{1},\dots ,m_{n} + h_{n})$ for $m = (m_1,\dots ,m_n)$ $h = (h_1,\dots ,h_n)$ and $|h| = \left(\sum\limits_{j=1}^{n}h_{j}\right)^{1/2}$.
\end{theorem}

\begin{theorem}\label{thm:0.3.Theorem.2} Let $G$ be an open set of $\mathbb{R}^{n}$. For any Lebesgue integrable function $f(x)$ in $G$ and $\varepsilon > 0$, there exists a continuous function $C_{\varepsilon}(x)$ in $G$ such that $\{x \in G; C_{\varepsilon}(x) \neq 0\}^a$ is a compact subset of $G$ and 
\[
\int_{G} | f (x)- C _{\varepsilon} (x) | \dd x < \varepsilon .
\]
\end{theorem}


\begin{remarku} Let $m$ be a Baire measure on a locally compact space $S$. A subset $Z$ of $S$ is called a set of $m$-measure zero if, for each $\varepsilon > 0$, there is a Baire set $B$ containing $Z$ with $m(B) < \varepsilon$. One can extend $m$ to the class of $m$-measurable sets, such a set being one which differs from a Baire set by a set of $m$-measure zero. Any property pertaining to a set of $m$-measure zero is said to hold $m$-almost everywhere ($m$-a.e.). One can also extend integrability to a function which coincides $m$-a.e. with a Baire function. 
\end{remarku}


\yossection{4}{Linear Spaces}
\label{sec:0.4}

\yossub{Linear Spaces}

\begin{definitionu} $A$ set $X$ is called a linear space over a field $K$ if the following conditions are satisfied: 
\begin{equation}\label{eq:0.4.1}
\parbox{0.90\linewidth}{$X$ is an abelian group (written additively),}
\end{equation}
A scalar multiplication is defined: to every element $x \in X$ and each $\alpha \in K$ there is associated an element of $X$, denoted by $\alpha x$, such that we have 
\begin{equation}\label{eq:0.4.2}
\begin{aligned}
&\alpha (x + y) = \alpha x + \alpha y (\alpha \in K; x, y \in X), (\alpha + \beta) x = \alpha x + \beta x (\alpha , \beta \in K; x \in X),
\\
&(\alpha \beta) x = \alpha (\beta x) (\alpha , \beta \in K; x \in X), 1 \cdot x = x (1 \text{ is the unit element of the field } K).
\end{aligned}
\end{equation}
\end{definitionu}

In the sequel we consider linear spaces only over the real number field $\mathbb{R}^{1}$ or the complex number field $\mathbb{C}^{1}$. A linear space will be said to be real or complex according as the field $K$ of coefficients is the real number field $\mathbb{R}^{1}$ or the complex number field $\mathbb{C}^{1}$. Thus, in what follows, we mean by a linear space a real or complex linear space. We shall denote by Greek letters the elements of the field of coefficients and by Roman letters the elements of $X$. The zero of $X$ (= the unit element of the additively written abelian group $X$) and the number zero will be denoted by the same letter $0$, since it does not cause inconvenience as $0 \cdot x = (\alpha- \alpha)x = \alpha x- \alpha x = 0$. The inverse element of the additively written abelian group $X$ will be denoted by $-x$; it is easy to see that $-x = (-1)x$. 

\begin{definitionu} The elements of a linear space $X$ are called vectors (of $X$). The vectors $x_{1}, x_{2}, \ldots, x_{n}$ of $X$ are said to be linearly independent if the equation $\sum\limits_{j=1}^{n} \alpha_{j} x_{j} = 0$ implies $\alpha_{1} = \alpha_{2} = \cdots = 0$. They are linearly dependent if such an equation holds where at least one coefficient is different from $0$. If $X$ contains $n$ linearly independent vectors, but every system of $(n + 1)$ vectors is linearly dependent, then $X$ is said to be of $n$-dimension. If the number of linearly independent vectors is not finite, then $X$ is said to be of infinite dimension. Any set of $n$ linearly independent vectors in an $n$-dimensional linear space constitutes a basis for $X$ 
and each vector $x$ of $X$ has a unique representation of the form $x = \sum\limits_{j=1}^{n} \alpha_j y_j$ in terms of the basis $y_1, y_2, \ldots, y_n$. A subset $M$ of a linear space $X$ is called a linear subspace or, in short, a subspace, if whenever $x, y \in M$, the linear combinations $\alpha x + \beta y$ also belong to $M$. $M$ is thus a linear space over the same coefficient field as $X$. 
\end{definitionu}

\yossub{Linear Operators and Linear Functionals}

\begin{definitionu} Let $X$, $Y$ be linear spaces over the same coefficient field $K$. A mapping $T: x \to y = T(x) = Tx$ defined on a linear subspace $D$ of $X$ and taking values in $Y$ is said to be linear, if 
\[
T \left(\alpha x _{1} + \beta x _{2}\right) = \alpha \left(T x _{1}\right) + \beta \left(T x _{2}\right).
\]
\end{definitionu}

The definition implies, in particular, 
\[
T \cdot 0 = 0, T (- x) =- (T x).
\]

We denote 
\[
D = D (T), \{y \in Y; y = T x, x \in D (T) \} = R (T), \{x \in D (T); T x = 0 \} = N (T)
\]
and call them the domain, the range and the null space of $T$, respectively. $T$ is called a linear operator or linear transformation on $D(T) \subseteq X$ into $Y$, or somewhat vaguely, a linear operator from $X$ into $Y$. If the range $R(T)$ is contained in the scalar field $K$, then $T$ is called a linear functional on $D(T)$. If a linear operator $T$ gives a one-to-one map of $D(T)$ onto $R(T)$, then the inverse map $T^{-1}$ gives a linear operator on $R(T)$ onto $D(T)$: 
\[
T ^{-1} T x = x \text{ for } x \in D (T) \quad \text{ and } \quad T T ^{-1} y = y \text{ for } y \in R (T).
\]
$T^{-1}$ is the inverse operator or, in short, the inverse of $T$. By virtue of $T(x_{1}- x_{2}) = Tx_{1}- Tx_{2}$, we have the following. 

\begin{propositionu} $A$ linear operator $T$ admits the inverse $T^{-1}$ iff $Tx = 0$ implies $x = 0$. 
\end{propositionu}

\begin{definitionu} Let $T_{1}$ and $T_{2}$ be linear operators with domains $D(T_{1})$ and $D(T_{2})$ both contained in a linear space $X$, and ranges $R(T_{1})$ and $R(T_{2})$ both contained in a linear space $Y$. Then $T_{1} = T_{2}$ iff $D(T_{1}) = D(T_{2})$ and $T_{1}x = T_{2}x$ for all $x \in D(T_{1}) = D(T_{2})$. If $D(T_{1}) \subseteq D(T_{2})$ and $T_{1}x = T_{2}x$ for all $x \in D(T_{1})$, then $T_{2}$ is called an extension of $T_{1}$, and $T_{1}$ a restriction of $T_{2}$; we shall then write $T_{1} \subseteq T_{2}$. 
\end{definitionu}

Convention. The value $T(x)$ of a linear functional $T$ at a point $x \in D(T)$ will sometimes be denoted by $\langle x, T \rangle$, i.e. 
\[
T (x) = \langle x, T \rangle .
\]

\yossub{Factor Spaces}

\begin{propositionu} Let $M$ be a linear subspace in a linear space $X$. We say that two vectors $x_{1}, x_{2} \in X$ are equivalent modulo $M$ if $(x_{1}- x_{2}) \in M$ and write this fact symbolically by $x_{1} \equiv x_{2} \pmod{M}$. Then we have: 

\begin{enumerate}[(i)]
\item $x \equiv x \pmod{M}$,
\item if $x_{1} \equiv x_{2} \pmod{M}$, then $x_{2} \equiv x_{1} \pmod{M}$,
\item if $x_{1} \equiv x_{2} \pmod{M}$ and $x_{2} \equiv x_{3} \pmod{M}$, then $x_{1} \equiv x_{3} \pmod{M}$.
\end{enumerate}
\end{propositionu}
\begin{proof}
 \begin{enumerate}[(i)]
\item is clear since $x- x = 0 \in M$.
\item If $(x_{1}- x_{2}) \in M$, then $(x_{2}- x_{1}) =-(x_{1}- x_{2}) \in M$.
\item If $(x_{1}- x_{2}) \in M$ and $(x_{2}- x_{3}) \in M$, then $(x_{1}- x_{3}) = (x_{1}- x_{2}) + (x_{2}- x_{3}) \in M$. \qedhere
\end{enumerate}
\end{proof}

We shall denote the set of all vectors $\in X$ equivalent modulo $M$ to a fixed vector $x$ by $\xi_{x}$. Then, in virtue of properties (ii) and (iii), all vectors in $\xi_{x}$ are mutually equivalent modulo $M$. $\xi_{x}$ is called a class of equivalent (modulo $M$) vectors, and each vector in $\xi_{x}$ is called a representative of the class $\xi_{x}$. Thus a class is completely determined by any one of its representatives, that is, $y \in \xi_{x}$ implies that $\xi_{y} = \xi_{x}$. Hence, two classes $\xi_{x}, \xi_{y}$ are either disjoint (when $y \in \xi_{x}$) or coincide (when $y \in \xi_{x}$). Thus the entire space $X$ decomposes into classes $\xi_{x}$ of mutually equivalent (modulo $M$) vectors. 

\begin{theoremu} We can consider the above introduced classes (modulo $M$) as vectors in a new linear space where the operation of addition of classes and the multiplication of a class by a scalar will be defined through 
\[
\xi_{x} + \xi_{y} = \xi_{x+y}, \alpha \xi_{x} = \xi_{\alpha x}.
\]
\end{theoremu}

\begin{proof}
 The above definitions do not depend upon the choice of representatives $x, y$ of the classes $\xi_x, \xi_y$ respectively. In fact, if $(x_1- x) \in M$, $(y_1- y) \in M$, then 
\[
\begin{gathered}
(x _{1} + y _{1})- (x + y) = (x _{1}- x) + (y _{1}- y) \in M, (\alpha x _{1}- \alpha x) = \alpha (x _{1}- x) \in M.
\end{gathered}
\]

We have thus proved $\xi_{x_{1}+y_{1}}=\xi_{x+y}$ and $\xi_{ax_{1}}=\xi_{ax}$, and the above definitions of the class addition and the scalar multiplication of the classes are justified. 
\end{proof}

\begin{definitionu} The linear space obtained in this way is called the factor space of $X$ modulo $M$ and is denoted by $X/M$. 
\end{definitionu}

\yossub{References}

Topological Space: P. \textsc{Alexandrov}-H. \textsc{Hopf} \cite{alexandrov6}, N. \textsc{Bourbaki} \cite{bourbaki26}, J. L. \textsc{Kelley} \cite{kelley146}. 

Measure Space: P. R. \textsc{Halmos} \cite{halmos92}, S. \textsc{Saks} \cite{saks280}. 
   % 0. Preliminaries
% ==== I ====
\yoschapter{I}{Semi-norms}
\label{chap:I}

The semi-norm of a vector in a linear space gives a kind of length for the vector. To introduce a topology in a linear space of infinite dimension suitable for applications to classical and modern analysis, it is sometimes necessary to make use of a system of an infinite number of semi-norms. It is one of the merits of the Bourbaki group that they stressed the importance, in functional analysis, of locally convex spaces which are defined through a system of semi-norms satisfying the axiom of separation. If the system reduces to a single semi-norm, the corresponding linear space is called a \emph{normed linear space}. If, furthermore, the space is complete with respect to the topology defined by this semi-norm, it is called a \emph{Banach space}. The notion of complete normed linear spaces was introduced around $1922$ by S. \textsc{Banach} and N. \textsc{Wiener} independently of each other. A modification of the norm, the \emph{quasi-norm} in the present book, was introduced by M. \textsc{Fréchet}. A particular kind of limit, the \emph{inductive limit}, of locally convex spaces is suitable for discussing the generalized functions or the distributions introduced by L. \textsc{Schwartz}, as a systematic development of S. L. \textsc{Sobolev}'s generalization of the notion of functions. 


\yossection{1}{Semi-norms and Locally Convex Linear Topological Spaces}
\label{sec:I.1}

As was stated in the above introduction, the notion of semi-norm is of fundamental importance in discussing linear topological spaces. We shall begin with the definition of the semi-norm. 

\begin{definition}\label{def:I.1.Definition.1} $A$ real-valued function $p(x)$ defined on a linear space $X$ is called a semi-norm on $X$, if the following conditions are satisfied: 
 \begin{align}
 P(x + y) &\leqq P(x) + P(y) \text{ (subadditivity) }, \label{eq:I.1.1} \\
 P(\alpha x) &= | \alpha | P(x). \label{eq:I.1.2}
 \end{align}
\end{definition}


\begin{example}\label{ex:I.1.Example.1} The $n$-dimensional Euclidean space $\mathbb{R}^{n}$ of points $x = (x_{1}, \ldots, x_{n})$ with coordinates $x_{1}, x_{2}, \ldots, x_{n}$ is an $n$-dimensional linear space by the operations: 
 \[
\begin{gathered}
x + y = (x_{1} + y_{1}, x_{2} + y_{2}, \dots , x_{n} + y_{n}), \alpha x = (\alpha x_{1}, \alpha x_{2}, \dots , \alpha x_{n}).
\end{gathered}
\]
\end{example}


In this case $p(x) = \max\limits_{1 \leqq j \leqq n} |x_j|$ is a semi-norm. As will be proved later, $p(x) = \left(\sum\limits_{j=1}^{n} |x_j|^q\right)^{1/q}$ with $q \geqq 1$ is also a semi-norm on $\mathbb{R}^{n}$. 

\begin{proposition}\label{prop:I.1.Proposition.1} $A$ semi-norm $p(x)$ satisfies 
\begin{align}
P(0)& = 0, \label{eq:I.1.3} \\
P\left(x_{1}- x_{2}\right)& \geqq | P\left(x_{1}\right)- P\left(x_{2}\right) |, \text{ in particular, } P(x) \geqq 0. \label{eq:I.1.4}
\end{align}
\end{proposition}

\begin{proof}
 $p(0) = p(0 \cdot x) = 0 \cdot p(x) = 0$. We have, by the subadditivity, $p(x_1- x_2) + p(x_2) \geqq p(x_1)$ and hence $p(x_1- x_2) \geqq p(x_1)- p(x_2)$. Thus $p(x_1- x_2) = |-1| \cdot p(x_2- x_1) \geqq p(x_2)- p(x_1)$ and so we obtain \eqref{eq:I.1.4}. 
\end{proof}

\begin{proposition}\label{prop:I.1.Proposition.2} Let $p(x)$ be a semi-norm on $X$, and $c$ any positive number. Then the set $M = \{x \in X; p(x) \leqq c\}$ enjoys the properties: 
\begin{equation}\label{eq:I.1.5}
M\ni0,
\end{equation}
\begin{equation}\label{eq:I.1.6}
\begin{minipage}{0.92\linewidth}
\everymath{}\let\limits\nolimits
$M$ is convex: $x,y\in M$ and $0<\alpha<1$ imply $\alpha x+(1-\alpha)y\in M$.
\end{minipage}
\end{equation}
\begin{equation}\label{eq:I.1.7}
\begin{minipage}{0.92\linewidth}
\everymath{}\let\limits\nolimits
$M$ is balanced (équilibré in Bourbaki's terminology): $x\in M$ and $|\alpha|\leqq1$ imply $\alpha x\in M$.
\end{minipage}
\end{equation}
\begin{equation}\label{eq:I.1.8}
\begin{minipage}{0.92\linewidth}
\everymath{}\let\limits\nolimits
$M$ is absorbing: for any $x\in X$, there exists $\alpha>0$ such that $\alpha^{-1}x\in M$.
\end{minipage}
\end{equation}
\begin{equation}\label{eq:I.1.9}
P(x)=\inf_{\alpha>0,\,\alpha^{-1}x\in M}\alpha c
\qquad\text{(inf = infimum = the greatest lower bound).}
\end{equation}
\end{proposition}

\begin{proof}
 \eqref{eq:I.1.5} is clear from \eqref{eq:I.1.3}. \eqref{eq:I.1.7} and \eqref{eq:I.1.8} are proved by \eqref{eq:I.1.2}. \eqref{eq:I.1.6} is proved by the subadditivity \eqref{eq:I.1.1} and \eqref{eq:I.1.2}. \eqref{eq:I.1.9} is proved by observing the equivalence of the three propositions below: 
\[
[ \alpha^{-1} x \in M ] \rightleftarrows [ P(\alpha^{-1} x) \leqq c ] \rightleftarrows [ P(x) \leqq \alpha c ]. \qedhere\]
\end{proof}

\begin{definition}\label{def:I.1.Definition.2} The functional 
\begin{equation}\tag{9'}\label{eq:I.1.9'}
p_{M} (x) = \inf\limits_ {\alpha > 0, \alpha^{-1} x \in M} \alpha
\end{equation}
is called the Minkowski functional of the convex, balanced and absorbing set $M$ of $X$. 
\end{definition}

\begin{proposition}\label{prop:I.1.Proposition.3} Let a family $\{p_{\gamma}(x); \gamma \in \Gamma\}$ of semi-norms of a linear space $X$ satisfy the axiom of separation: 
\begin{equation}
\label{eq:I.1.10}
\begin{minipage}{0.92\linewidth}
\everymath{}\let\limits\nolimits
For any $x_0\neq0,$ there exists $P_{\gamma_0}(x)$ in the family such that $P_{\gamma_0}(x_0)\neq0.$
\end{minipage}
\end{equation}
\end{proposition}

Take any finite system of semi-norms of the family, say $p_{\gamma_1}(x), p_{\gamma_2}(x), \ldots, \ldots, p_{\gamma_n}(x)$ and any system of $n$ positive numbers $\varepsilon_1, \varepsilon_2, \ldots, \varepsilon_n$, and set 
\begin{equation}\label{eq:I.1.11}
U = \{x \in X; P_ {\gamma_{j}} (x) \leqq \varepsilon_{j} \quad (j = 1, 2, \dots , n) \}.
\end{equation}

$U$ is a convex, balanced and absorbing set. Consider such a set $U$ as a neighbourhood of the vector $0$ of $X$, and define a neighbourhood of any vector $x_{0}$ by the set of the form 
\begin{equation}\label{eq:I.1.12}
x_{0} + U = \{y \in X; y = x_{0} + u, u \in U \}.
\end{equation}

Consider a subset $G$ of $X$ which contains a neighbourhood of each of its point. Then the totality $\{G\}$ of such subsets $G$ satisfies the axiom of open sets, given in Chapter~\ref{chap:0}, Preliminaries, \ref{sec:0.2}. 

\begin{proof}
 We first show that the set $G_0$ of the form $G_0 = \{x \in X; p_\gamma(x) < c\}$ is open. For, let $x_0 \in G_0$ and $p_\gamma(x_0) = \beta < c$. Then the neighbourhood of $x_0, x_0 + U$ where $U = \{x \in X; p_\gamma(x) \leqq 2^{-1}(c- \beta)\}$, is contained in $G_0$, because $u \in U$ implies $p_\gamma(x_0 + u) \leqq p_\gamma(x_0) + p_\gamma(u) < \beta + (c- \beta) = c$. 

Hence, for any point $x_0 \in X$, there is an open set $x_0 + G_0$ which contains $x_0$. It is clear, by the above definition of open sets, that the union of open sets and the intersection of a finite number of open sets are also open. 

Therefore we have only to prove Hausdorff's axiom of separation: If $x_{1} \neq x_{2}$, then there exist disjoint open sets $G_{1}$ and $G_{2}$ such that 
\begin{equation}\label{eq:I.1.13}
x_{1} \in G_{1}, x_{2} \in G_{2}.
\end{equation}

In view of definition \eqref{eq:I.1.12} of the neighbourhood of a general point $x_0$, it will be sufficient to prove \eqref{eq:I.1.13} for the case $x_1 = 0$, $x_2 \neq 0$. We choose, by \eqref{eq:I.1.10}, $p_{\gamma_1}(x)$ such that $p_{\gamma_2}(x_2) = \alpha > 0$. Then $G_1 = \{x \in X; p_{\gamma_1}(x) < \alpha / 2\}$ is open, as proved above. Surely $G_1 \ni 0 = x_1$. We have to show that $G_1$ and $G_2 = x_2 + G_1$ have no point in common. Assume the contrary and let there exist a $y \in G_1 \cap G_2$. $y \in G_2$ implies $y = x_2 + g = x_2- (-g)$ with some $g \in G_1$ and so, by \eqref{eq:I.1.4}, $p_{\gamma_1}(y) \geqq p_{\gamma_2}(x_2)- p(-g) \geqq \alpha- 2^{-1}\alpha = \alpha / 2$, because $-g$ belongs to $G_1$ with $g$. This contradicts the inequality $p_{\gamma_1}(y) < \alpha / 2$ implied by $y \in G_1$. 
\end{proof}

\begin{proposition}\label{prop:I.1.Proposition.4} By the above definition of open sets, $X$ is a linear topological space, that is, $X$ is a linear space and at the same time a topological space such that the two mappings $X \times X \to X : (x, y) \to x + y$ and $K \times X \to X : (\alpha, x) \to \alpha x$ are both continuous. Moreover, each semi-norm $p_{y}(x)$ is a continuous function on $X$. 
\end{proposition}

\begin{proof}
 For any neighbourhood $U$ of $0$, there exists a neighbourhood $V$ of $0$ such that 
\[
V + V = \{w \in X; w = v_{1} + v_{2} \text{ where } v_{1}, v_{2} \in V \} \subseteq U,
\]
since the semi-norm is subadditive. Hence, by writing 
\[
(x + y)- (x_{0} + y_{0}) = (x- x_{0}) + (y- y_{0}),
\]
we see that the mapping $(x, y) \to x + y$ is continuous at $x = x_0$, $y = y_0$. For any neighbourhood $U$ of $0$ and any scalar $\alpha \neq 0$, the set $\alpha U = \{x \in X; x = \alpha u, u \in U\}$ is also a neighbourhood of $0$. Thus, by writing 
\[
\alpha x- \alpha_{0} x_{0} = \alpha (x- x_{0}) + (\alpha- \alpha_{0}) x_{0},
\]
we see by \eqref{eq:I.1.2} that $(\alpha, x) \to \alpha x$ is continuous at $\alpha = \alpha_0$, $x = x_0$. 

The continuity of the semi-norm $p_{\gamma}(x)$ at the point $x = x_{0}$ is proved by $|p_{\gamma}(x)- p_{\gamma}(x_{0})| \leqq p_{\gamma}(x- x_{0})$. 
\end{proof}

\begin{definition}\label{def:I.1.Definition.3} $A$ linear topological space $X$ is called a locally convex, linear topological space, or, in short, a locally convex space, if any of its open sets $\ni0$ contains a convex, balanced and absorbing open set. 
\end{definition}

\begin{proposition}\label{prop:I.1.Proposition.5} The Minkowski functional $p_{M}(x)$ of the convex, balanced and absorbing subset $M$ of a linear space $X$ is a semi-norm on $X$. 
\end{proposition}

\begin{proof}
 By the convexity of $M$, the inclusions \[
 x / (P_{M} (x) + \varepsilon) \in M, y / (P_{M} (y) + \varepsilon) \in M \text{ for any } \varepsilon > 0
 \]
 imply 
 \[
 \frac {P_{M} (x) + \varepsilon}{P_{M} (x) + P_{M} (y) + 2 \varepsilon} \cdot \frac {x}{P_{M} (x) + \varepsilon} + \frac {P_{M} (y) + \varepsilon}{P_{M} (x) + P_{M} (y) + 2 \varepsilon} \cdot \frac {y}{P_{M} (y) + \varepsilon} \in M
 \]
 and so $p_{M}(x+y)\leqq p_{M}(x)+p_{M}(y)+2\varepsilon$. Since $\varepsilon>0$ was arbitrary, we obtain the subadditivity of $p_{M}(x)$. Similarly we obtain $p_{M}(\alpha x)=|\alpha|p_{M}(x)$ since $M$ is balanced. 
\end{proof}

We have thus proved 


\begin{theoremu} $A$ linear space $X$, topologized as above by a family of semi-norms $p_{\gamma}(x)$ satisfying the axiom of separation \eqref{eq:I.1.10}, is a locally convex space in which each semi-norm $p_{\gamma}(x)$ is continuous. Conversely, any locally convex space is nothing but the linear topological space, topologized as above through the family of semi-norms obtained as the Minkowski functionals of convex balanced and absorbing open sets of $X$. 
\end{theoremu}

\begin{definition}\label{def:I.1.Definition.4} Let $f(x)$ be a complex-valued function defined in an open set $\Omega$ of $\mathbb{R}^{n}$. By the support (or carrier) of $f$, denoted by $\supp(f)$, we mean the smallest closed set (of the topological space $\Omega$) containing the set $\{x \in \Omega; f(x) \neq 0\}$. It may equivalently be defined as the smallest closed set of $\Omega$ outside which $f$ vanishes identically. 
\end{definition}

\begin{definition}\label{def:I.1.Definition.5} By $C^k(\Omega)$, $0 \leqq k \leqq \infty$, we denote the set of all complex-valued functions defined in $\Omega$ which have continuous partial derivatives of order up to and including $k$ (of order $< \infty$ if $k = \infty$). By $C_0^k(\Omega)$, we denote the set of all functions $\in C^k(\Omega)$ with compact support, i.e., those functions $\in C^k(\Omega)$ whose supports are compact subsets of $\Omega$. A classical example of a function $\in C_0^\infty(\mathbb{R}^{n})$ is given by 
\begin{gather}f (x) = \exp\left((| x | ^{2}- 1) ^{-1}\right) \text{ for } | x | = | (x_{1}, \dots , x_{n}) | = \left(\sum\limits_{j=1} ^{n} x_{j} ^{2}\right) ^{1/2} < 1, = 0 \text{ for } | x | \geqq 1. \label{eq:I.1.14}
\\\text{ The Space } \mathfrak {E} ^{k} (\Omega) \notag
\end{gather}
$C^k (\Omega)$ is a linear space by 
\[
\left(f_{1} + f_{2}\right) (x) = f_{1} (x) + f_{2} (x), (\alpha f) (x) = \alpha f (x).
\]
\end{definition}

For any compact subset $K$ of $\Omega$ and any non-negative integer $m \leqq k$ ( $m < \infty$ when $k = \infty$), we define the semi-norm 
\[
p_{K,m} (f) = \sup\limits_{| s | \leqq m, x \in K} | D ^{s} f (x) |, f \in C ^{k} (\Omega),
\]
where $\sup= \text{supremum} =$ the least upper bound and 
\[
\begin{aligned}
&D ^{s} f (x) = \frac {\partial^ {s_{1}} + s_{2} + \cdots + s_{n}}{\partial x_{1} ^ {s_{1}} \partial x_{2} ^ {s_{2}} \cdots \partial x_{n} ^ {s_{n}}} f (x_{1}, x_{2}, \dots , x_{n}), | s | = | (s_{1}, s_{2}, \dots , s_{n}) | = \sum\limits_{j=1} ^{n} s_{j}.
\end{aligned}
\]

Then $C^k(\Omega)$ is a locally convex space by the family of these semi-norms. We denote this locally convex space by $\mathfrak{E}^k(\Omega)$. The convergence $\lim\limits_{h\to \infty}f_h = f$ in this space $\mathfrak{E}^k(\Omega)$ is exactly the uniform convergence $\lim\limits_{h\to \infty}D^s f_h(x) = D^s f(x)$ on every compact subset $K$ of $\Omega$, for each $s$ with $|s| \leqq k (|s| < \infty \text{ if } k = \infty)$. We often write $\mathfrak{E}(\Omega)$ for $\mathfrak{E}^\infty (\Omega)$. 

\begin{proposition}\label{prop:I.1.Proposition.6} $\mathfrak{E}^k (\Omega)$ is a metric space. 
\end{proposition}

\begin{proof}
 Let $K_{1} \subseteq K_{2} \subseteq \cdots \subseteq K_{n} \subseteq \cdots$ be a monotone increasing sequence of compact subsets of $\Omega$ such that $\Omega = \bigcup\limits_{n=1}^{\infty} K_{n}$. Define, for each positive integer $h$, the distance 
\[
d_{h} (f, g) = \sum\limits_{m=0} ^{k} 2 ^{-m} P_ {K_{h}, m} (f- g) \cdot (1 + P_ {K_{h}, m} (f- g)) ^{-1}. \]

Then the convergence $\lim\limits_{s\to \infty}f_s = f$ in $\mathfrak{E}_k(\Omega)$ is defined by the distance 
\[
d (f, g) = \sum\limits_{h=1} ^{\infty} 2 ^{-h} d_{h} (f, g) \cdot (1 + d_{h} (f, g)) ^{-1}.
\]

We have to show that $d_h(f, g)$ and $d(f, g)$ satisfy the triangle inequality. The triangle inequality for $d_h(f, g)$ is proved as follows: by the subadditivity of the semi-norm $p_{K_h,m}(f)$, we easily see that $d_h(f, g)$ satisfies the triangle inequality $d_h(f, g) \leqq d_h(f, k) + d_h(k, g)$, if we can prove the inequality 
\[
\begin{aligned}
&| \alpha- \beta | \cdot (1 + | \alpha- \beta |) ^{-1}  \leqq | \alpha- \gamma | (1 + | \alpha- \gamma |) ^{-1} + | \gamma- \beta | (1 + | \gamma- \beta |) ^{-1}
\end{aligned}
\]
for complex numbers $\alpha, \beta$ and $\gamma$; the last inequality is clear from the inequality valid for any system of non-negative numbers $\alpha, \beta$ and $\gamma$: 
\[
(\alpha + \beta) (1 + \alpha + \beta) ^{-1} \leqq \alpha (1 + \alpha) ^{-1} + \beta (1 + \beta) ^{-1}.
\]
The triangle inequality for $d(f, g)$ may be proved similarly. 
\end{proof}

\begin{definition}\label{def:I.1.Definition.6} 
 Let $X$ be a linear space. Let a family $\{X_{\alpha}\}$ of linear subspaces $X_{\alpha}$ of $X$ be such that $X$ is the union of $X_{\alpha}$'s. Suppose that each $X_{\alpha}$ is a locally convex linear topological space such that, if $X_{\alpha_{1}} \subseteq X_{\alpha_{2}}$, then the topology of $X_{\alpha_{1}}$ is identical with the relative topology of $X_{\alpha_{1}}$ as a subset of $X_{\alpha_{2}}$. We shall call ``open'' every convex balanced and absorbing set $U$ of $X$ iff the intersection $U \cap X_{\alpha}$ is an open set of $X_{\alpha}$ containing the zero vector $0$ of $X_{\alpha}$, for all $X_{\alpha}$. If $X$ is a locally convex linear topological space whose topology is defined in the stated way, then $X$ is called the (strict) inductive limit of $X_{\alpha}$'s. 
\end{definition}

\begin{remarku} Take, from each $X_{\alpha}$, a convex balanced neighbourhood $U_{\alpha}$ of $0$ of $X_{\alpha}$. Then the convex closure $U$ of the union $V = \bigcup\limits_{\alpha} U_{\alpha}$, i.e., 
\[
\begin{aligned}
U& = \left\{u \in X; u = \sum\limits_{j=1} ^{n} \beta_{j} v_{j}, v_{j} \in V, \beta_{j} \geqq 0 (j = 1, 2, \dots , n), \sum\limits_{j=1} ^{n} \beta_{j} = 1 \right.
\\
&\quad \text{ with arbitrary finite } n \Bigg \}
\end{aligned}
\]
\end{remarku}

surely satisfies the condition that it is convex balanced and absorbing in such a way that $U \cap X_{\alpha}$ is a convex balanced neighbourhood of $0$ of $X_{\alpha}$, for all $X_{\alpha}$. The set of all such $U$'s corresponding to an arbitrary choice of $U_{\alpha}$'s is a fundamental system of neighbourhoods of $0$ of the (strict) inductive limit $X$ of $X_{\alpha}$'s, i.e., every neighbourhood of $0$ of the (strict) inductive limit $X$ of $X_{\alpha}$'s contains one of the $U$'s obtained above. This fact justifies the above definition of the (strict) inductive limit. 

\yossub{The Space $\mathfrak{D}(\Omega)$}

$C_0^\infty (\Omega)$ is a linear space by 
\[
\left(f_{1} + f_{2}\right) (x) = f_{1} (x) + f_{2} (x), (\alpha f) (x) = \alpha f (x).
\]

For any compact subset $K$ of $\Omega$, let $\mathfrak{D}_K(\Omega)$ be the set of all functions $f \in C_0^\infty(\Omega)$ such that $\supp(f) \subseteq K$. Define a family of semi-norms on $\mathfrak{D}_K(\Omega)$ by 
\[
p_{K,m} (f) = \sup\limits_{| s | \leqq m, x \in K} | D ^{s} f (x) |, \text{ where } m < \infty .
\]
$\mathfrak{D}_K(\Omega)$ is a locally convex linear topological space, and, if $K_{1} \subseteq K_{2}$, then it follows that the topology of $\mathfrak{D}_{K_1}(\Omega)$ is identical with the relative topology of $\mathfrak{D}_{K_1}(\Omega)$ as a subset of $\mathfrak{D}_{K_2}(\Omega)$. Then the (strict) inductive limit of $\mathfrak{D}_K(\Omega)$'s, where $K$ ranges over all compact subsets of $\Omega$, is a locally convex, linear topological space. Topologized in this way, $C_0^\infty(\Omega)$ will be denoted by $\mathfrak{D}(\Omega)$. It is to be remarked that, 
\[
P(f) = \sup\limits_{x \in \Omega} | f (x) |
\]
is one of the semi-norms which defines the topology of $\mathfrak{D}(\Omega)$. For, if we set $U = \{f \in C_0^\infty(\Omega); p(f) \leqq 1\}$, then the intersection $U \cap \mathfrak{D}_K(\Omega)$ is given by $U_K = \{f \in \mathfrak{D}_K(\Omega); p_K(f) = \sup\limits_{x \in K} |f(x)| \leqq 1\}$. 

\begin{proposition}\label{prop:I.1.Proposition.7} The convergence $\lim\limits_{h\to \infty}f_h = 0$ in $\mathfrak{D}(\Omega)$ means that the following two conditions are satisfied: \begin{enumerate}[(i)]
\item there exists a compact subset $K$ of $\Omega$ such that $\supp(f_h)\subseteq K(h = 1,2,\ldots)$
\item for any differential operator $D^{s}$, the sequence $\{D^s f_h(x)\}$ converges to $0$ uniformly on $K$.
\end{enumerate}
\end{proposition}

\begin{proof}
 We have only to prove (i). Assume the contrary, and let there exist a sequence $\{x^{(k)}\}$ of points $\in\Omega$ having no accumulation points in $\Omega$ and a subsequence $\{f_{h_{j}}(x)\}$ of $\{f_{h}(x)\}$ such that $f_{h_{j}}(x^{(j)})\neq0$. Then the semi-norm 
\[
P(f) = \sum\limits_{k=1}^{\infty} 2\sup\limits_{x\in K_k-K_{k-1}}
\left|f(x)/f_{h_k}\left(x^{(k)}\right)\right|.
\]
\begingroup\everymath{}
Here the monotone increasing sequence of compact subsets $K_j$ of $\Omega$
satisfies $\bigcup_{j=1}^{\infty}K_j=\Omega$ and
$x^{(k)}\in K_k-K_{k-1}$ ($k=1,2,\ldots$), with $K_0=\emptyset$.
This semi-norm defines a neighbourhood $U = \{f \in C_0^\infty(\Omega); p(f) \leqq 1\}$ of $0$ of $\mathfrak{D}(\Omega)$. However, none of the $f_{h_{k}}$'s is contained in $U$. \qedhere\par\endgroup 
\end{proof}

\begin{corollaryu} The convergence $\lim\limits_{h\to \infty}f_h = f$ in $\mathfrak{D}(\Omega)$ means that the following two conditions are satisfied: \begin{enumerate}[(i)]
\item there exists a compact subset $K$ of $\Omega$ such that $\supp(f_h)\subseteq K$ ( $h = 1,2,\ldots$)
\item for any differential operator $D^{s}$, the sequence $D^{s}f_{h}(x)$ converges to $D^{s}f(x)$ uniformly on $K$.
\end{enumerate}
\end{corollaryu}

\begin{proposition}[$A$ theorem of approximation]\label{prop:I.1.Proposition.8} Any continuous function $f \in C_0^0(\mathbb{R}^{n})$ can be approximated by functions of $C_0^\infty(\mathbb{R}^{n})$ uniformly on $\mathbb{R}^{n}$. 
\end{proposition}

\begin{proof}
 Let $\theta_{1}(x)$ be the function introduced in \eqref{eq:I.1.14} and put 
\begin{equation}\label{eq:I.1.15}
\begin{aligned}
&\theta_{a} (x) = h_{a} ^{-1} \theta_{1} (x / a), \text{ where } a > 0 \text{ and } h_{a} > 0 \text{ are such that } \int_ {\mathbb{R}^{n}} \theta_{a} (x) \dd x = 1.
\end{aligned}
\end{equation}

We then define the regularization $f_{a}$ of $f$: 
\begin{equation}\label{eq:I.1.16}
\begin{aligned}
f_{a} (x) &= \int_ {\mathbb{R}^{n}} f (x- y) \theta_{a} (y) \dd y = \int_ {\mathbb{R}^{n}} f (y) \theta_{a} (x- y) \dd y, \text{ where }
\\
x- y &= (x_{1}- y_{1}, x_{2}- y_{2}, \ldots , x_{n}- y_{n}).
\end{aligned}
\end{equation}

The integral is convergent since $f$ and $\theta_{a}$ have compact support. Moreover, since 
\[
f_{a} (x) = \int_ {\supp (f)} f (y) \theta_{a} (x- y) \dd y,
\]
the support of $f_{a}$ may be taken to be contained in any neighbourhood of the $\supp(f)$ if we take $a > 0$ sufficiently small. Next, by differentiating under the integral sign, we have 
\begin{equation}\label{eq:I.1.17}
D ^{s} f_{a} (x) = D_{x} ^{s} f_{a} (x) = \int_ {\mathbb{R}^{n}} f (y) D_{x} ^{s} \theta_{a} (x- y) \dd y,
\end{equation}
and so $f_{a}$ is in $C_0^\infty (\mathbb{R}^{n})$. Finally we have by $\int \limits_{\mathbb{R}^{n}}\theta_a(\dot{x}-y)\dd y = 1$, 
\[
\begin{aligned}
&| f_{a} (x)- f (x) |  \leqq \int_ {\mathbb{R}^{n}} | f (y)- f (x) | \theta_{a} (x- y) \dd y \leqq \int_{| f (y)- f (x) | \leqq \varepsilon} | f (y)- f (x) | \theta_{a} (x- y) \dd y \\
&+ \int_{| f (y)- f (x) | > \varepsilon} | f (y)- f (x) | \theta_{a} (x- y) \dd y.
\end{aligned}
\]

The first term on the right is $\leqq\varepsilon$; and the second term on the right equals $0$ for sufficiently small $a>0$, because, by the uniform continuity of the function $f$ with compact support, there exists an $a>0$ such that $|f(y)-f(x)|>\varepsilon$ implies $|y-x|>a$. We have thus proved our Proposition. 
\end{proof}


\yossection{2}{Norms and Quasi-norms}
\label{sec:I.2}

\begin{definition}\label{def:I.2.Definition.1} $A$ locally convex space is called a normed linear space, if its topology is defined by just one semi-norm. 
\end{definition}

Thus a linear space $X$ is called a normed linear space, if for every $x \in X$, there is associated a real number $\|x\|$, the norm of the vector $x$, such that 
\begin{align}
\left\| x \right\|& \geqq 0 \text{ and } \left\| x \right\| = 0 \text{ iff } x = 0, \label{eq:I.2.1} \\
\left\| x + y \right\|& \leqq \left\| x \right\| + \left\| y \right\| (\text{ triangle inequality }), \label{eq:I.2.2} \\
\left\| \alpha x \right\|& = | \alpha | \cdot \left\| x \right\|. \label{eq:I.2.3}
\end{align}

The topology of a normed linear space $X$ is thus defined by the distance 
\begin{equation}\label{eq:I.2.4}
d (x, y) = \| x- y \|.
\end{equation}

In fact, $d(x, y)$ satisfies the axiom of distance: 
\begin{gather}d (x, y) \geqq 0 \text{ and } d (x, y) = 0 \text{ iff } x = y, d (x, y) \leqq d (x, z) + d (z, y) (\text{ triangle inequality }), \notag \\d (x, y) = d (y, x). \notag
\end{gather}

For, $d(x, y) = \|x- y\| = \|y- x\| = d(y, x)$ and $d(x, y) = \|x- y\| = \|x- z + z- y\| \leqq \|x- z\| + \|z- y\| = d(x, z) + d(z, y)$ by \eqref{eq:I.2.1}, \eqref{eq:I.2.2}, 

\eqref{eq:I.2.3} and \eqref{eq:I.2.4}. 

The convergence $\lim\limits_{n\to\infty}d(x_n,x)=0$ in a normed linear space $X$ will be denoted by $\slim\limits_{n\to\infty}x_n=x$ or simply by $x_n\to x$, and we say that the sequence $\{x_n\}$ converges strongly to $x$. The adjective ``strong'' is introduced to distinguish it from the ``weak'' convergence to be introduced later. 

\begin{proposition}\label{prop:I.2.Proposition.1} In a normed linear space $X$, we have 
\begin{align}
\lim\limits_{n \rightarrow \infty} \| x_{n} \|& = \| x \| \text{ if } \slim_{n \rightarrow \infty} x_{n} = x, \label{eq:I.2.5} \\
\slim_{n \to \infty} \alpha_{n} x_{n}& = \alpha x \text{ if } \lim\limits_{n \to \infty} \alpha_{n} = \alpha \text{ and } \slim_{n \to \infty} x_{n} = x, \label{eq:I.2.6} \\
\slim_{n \to \infty} (x_{n} + y_{n})& = x + y \text{ if } \slim_{n \to \infty} x_{n} = x \text{ and } \slim_{n \to \infty} y_{n} = y. \label{eq:I.2.7}
\end{align}
\end{proposition}

\begin{proof}
 \eqref{eq:I.2.5}, \eqref{eq:I.2.6} and \eqref{eq:I.2.7} are already proved, since $X$ is a locally convex space topologized by just one semi-norm $p(x) = \|x\|$. However, we shall give a direct proof as follows. As a semi-norm, we have 
\begin{equation}\label{eq:I.2.8}
\left\| x- y \right\| \geqq \left| \left\| x \right\|- \left\| y \right\| \right|
\end{equation}
and hence \eqref{eq:I.2.5} is clear. \eqref{eq:I.2.7} is proved by $\|(x + y)- (x_n + y_n)\| = \|(x- x_n) + (y- y_n)\| \leqq \|x- x_n\| + \|y- y_n\|$. From $\|\alpha x- \alpha_n x_n\| \leqq \|\alpha x- \alpha_n x\| + \|\alpha_n x- \alpha_n x_n\| \leqq |\alpha- \alpha_n| \cdot \|x\| + |\alpha_n| \cdot \|x- x_n\|$ and the boundedness of the sequence $\{\alpha_n\}$ we obtain \eqref{eq:I.2.6}. 
\end{proof}

\begin{definition}\label{def:I.2.Definition.2} $A$ linear space $X$ is called a \emph{quasi-normed linear space}, if, for every $x \in X$, there is associated a real number $\|x\|$, the \emph{quasi-norm} of the vector $x$, which satisfies \eqref{eq:I.2.1}, \eqref{eq:I.2.2} and 
\begin{equation}\tag{3'}\label{eq:I.2.3'}
\left\|- x \right\| = \left\| x \right\|, \lim\limits_ {\alpha_{n} \rightarrow 0} \left\| \alpha_{n} x \right\| = 0 \text{ and } \lim\limits_ {\left\| x_{n} \right\|\rightarrow 0} \left\| \alpha x_{n} \right\| = 0.
\end{equation}
\end{definition}

\begin{proposition}\label{prop:I.2.Proposition.2} In a quasi-normed linear space $X$, we have \eqref{eq:I.2.5}, \eqref{eq:I.2.6} and \eqref{eq:I.2.7}. 
\end{proposition}

\begin{proof}
 We need only prove \eqref{eq:I.2.6}. The proof in the preceding Proposition shows that we have to prove 
\begin{equation}
\tag{9}\label{eq:I.2.9}
\begin{minipage}{0.92\linewidth}
\everymath{}\let\limits\nolimits
$\lim\limits_{n \rightarrow \infty} \| x_{n} \| = 0$ implies that $\lim\limits_{n \rightarrow \infty} \| \alpha x_{n} \| = 0$ uniformly in $\alpha$ on any bounded set of $\alpha.$
\end{minipage}
\end{equation}

The following proof of \eqref{eq:I.2.9} is due to S. \textsc{Kakutani} (unpublished). Consider the functional $p_n(\alpha) = \|\alpha x_n\|$ defined on the linear space $\mathbb{R}^{1}$ of real numbers normed by the absolute value. By the triangle inequality of $p_n(\alpha)$ and \eqref{eq:I.2.3}, $p_n(\alpha)$ is continuous on $\mathbb{R}^{1}$. Hence, from $\lim\limits_{n\to \infty}p_n(\alpha) = 0$ implied by \eqref{eq:I.2.3} and Egorov's theorem (Chapter~\ref{chap:0}, Preliminaries, \ref{sec:0.3}. Measure Spaces), we see that there exists a Baire measurable set $A$ on the real line $\mathbb{R}^{1}$ with the property: 
\begin{equation}\tag{10}\label{eq:I.2.10}
\parbox[b]{0.92\linewidth}{\setlength{\parindent}{1.2em}\indent
the Lebesgue measure $|A|$ of $A$ is $>0$ and $\lim\limits_{n\to\infty}p_n(\alpha)=0$ uniformly on $A$.}
\end{equation}


Since the Lebesgue measure on the real line is continuous with respect to translations, we have, denoting by $B \ominus C$ the symmetric difference $B \cup C- B \cap C$, 
\[
\left| (A + \sigma) \ominus A \right|\rightarrow 0 \quad \text{ as } \quad \sigma \rightarrow 0.
\]

Thus there exists a positive number $\sigma_0$ such that 
$|\sigma| \leqq \sigma_{0}$ implies $|(A + \sigma) \ominus A| < |A|/2$, in particular, $|(A + \sigma) \cap A| > 0$. Hence, for any real number $\sigma$ with $|\sigma| \leqq \sigma_{0}$, there is a representation 
\[
\sigma = \alpha- \alpha^{\prime} \text{ with } \alpha \in A, \alpha^{\prime} \in A.
\]

Therefore, by $p_n(\sigma) = p_n(\alpha- \alpha') \leqq p_n(\alpha) + p_n(\alpha')$, we see that 
\[
\lim\limits_{n \rightarrow \infty} P_{n} (\sigma) = 0 \text{ uniformly in } \sigma \text{ when } | \sigma | \leqq \sigma_{0}.
\]

Let $M$ be any positive number. Then, taking a positive integer $k \geqq M / \sigma_0$ and remembering $\phi_n(k\sigma) \leqq k\phi_n(\sigma)$, we see that \eqref{eq:I.2.9} is true for $|\alpha| \leqq M$. 
\end{proof}

\begin{remarku} The above proof may naturally be modified so as to apply to complex quasi-normed linear spaces $X$ as well. 
\end{remarku}

As in the case of normed linear spaces, the convergence $\lim\limits_{n\to\infty}\left\|x-x_{n}\right|=0$ in a quasi-normed linear space will be denoted by $\slim\limits_{n\to\infty}x_{n}=x$, or simply by $x_{n}\to x$; we shall then say that the sequence $\{x_{n}\}$ converges strongly to $x$. 

\begin{examplen} Let the topology of a locally convex space $X$ be defined by a countable number of semi-norms $p_{n}(x)$ ( $n = 1, 2, \ldots$). Then $X$ is a quasi-normed linear space by the quasi-norm
\[
\left\| x \right\| = \sum\limits_{n=1} ^{\infty} 2 ^{-n} P_{n} (x) (1 + P_{n} (x)) ^{-1}.
\]

For, the convergence $\lim\limits_{h\to \infty}p_n(x_h) = 0$ ( $n = 1,2,\ldots$) is equivalent to $s\text{- }\lim\limits_{h\to \infty}x_h = 0$ with respect to the quasi-norm $\|x\|$ above. 


\end{examplen}
\yossection{3}{Examples of Normed Linear Spaces}
\label{sec:I.3}

\begin{example}\label{ex:I.3.Example.1} $C(S)$. Let $S$ be a topological space. Consider the set $C(S)$ of all real-valued (or complex-valued), bounded continuous functions $x(s)$ defined on $S$. $C(S)$ is a normed linear space by 
\[
(x + y) (s) = x (s) + y (s), (\alpha x) (s) = \alpha x (s), \| x \| = \sup\limits_{s \in S} | x (s) |.
\]
\end{example}

In $C(S)$, \slim $x_{n} = x$ means the uniform convergence of the functions $x_{n}(s)$ to $x(s)$. 

\begin{example}\label{ex:I.3.Example.2} $L^p(S, \mathfrak{B}, m)$, or, in short, $L^p(S)$ ($1 \leqq p < \infty$). Let $L^p(S)$ be the set of all real-valued (or complex-valued) $\mathfrak{B}$-measurable functions $x(s)$ defined $m$-a.e. on $S$ such that $|x(s)|^p$ is $m$-integrable over $S$. $L^p(S)$ is a linear space by 
\[
(x + y) (s) = x (s) + y (s), (\alpha x) (s) = \alpha x (s).
\]
\end{example}

For, $(x(s) + y(s))$ belongs to $L^p(S)$ if $x(s)$ and $y(s)$ both belong to $L^p(S)$, as may be seen from the inequality $|x(s) + y(s)|^p\leqq 2^p(|x(s)|^p+ |y(s)|^p)$. We define the norm in $L^p(S)$ by 
\begin{equation}\label{eq:I.3.1}
\left\| x \right\| = \left(\int_{S} | x (s) | ^{p} m (\dd s)\right) ^{1/p}.
\end{equation}

The subadditivity 
\begin{equation}\label{eq:I.3.2}
\begin{aligned}
&\left(\int_{S} | x (s) + y (s) | ^{p} m (\dd s)\right) ^{1/p}  \leqq \left(\int_{S} | x (s) | ^{p} m (\dd s)\right) ^{1/p} + \left(\int_{S} | y (s) | ^{p} m (\dd s)\right) ^{1/p},
\end{aligned}
\end{equation}
called Minkowski's inequality, is clear for the case $p = 1$. To prove the general case $1 < p < \infty$, we need 

\begin{lemma}\label{lem:I.3.Lemma.1} Let $1 < p < \infty$ and let the conjugate exponent $p'$ of $p$ be defined by 
\begin{equation}\label{eq:I.3.3}
\frac {1}{p} + \frac {1}{p^{\prime}} = 1.
\end{equation}
\end{lemma}

Then, for any pair of non-negative numbers $a$ and $b$, we have 
\begin{equation}\label{eq:I.3.4}
a b \leqq \frac {a ^{p}}{p} + \frac {b ^ {p^{\prime}}}{p^{\prime}},
\end{equation}
where the equality is satisfied iff $a = b^{1/(p-1)}$. 

\begin{proof}
 The minimum of the function $f(c) = \frac{c^p}{p} + \frac{1}{p'}- c$ for $c \geqq 0$ is attained only at $c = 1$, and the minimum value is $0$. By taking $c = ab^{-1/(p-1)}$ we see that the Lemma is true. 
\end{proof}

\begin{proof}[Proof of \eqref{eq:I.3.2}.]
 We first prove Hölder's inequality 
\begin{equation}\label{eq:I.3.5}
\int | x (s) y (s) | \leqq \left(\int | x (s) | ^{p}\right) ^{1/p} \cdot \left(\int | y (s) | ^ {p^{\prime}}\right) ^ {1 / p^{\prime}}
\end{equation}

(for convenience, we write $\int z(s)$ for $\int_{S}z(s)m(\dd s)$). 


To this end, we assume that $A = \left( \int |x(s)|^p\right)^{1/p}$ and $B = \left( \int |y(s)|^{p'} \right)^{1/p'}$ are both $\neq 0$, since otherwise $x(s)y(s) = 0$ a.e. and so \eqref{eq:I.3.5} would be true. Now, by taking $a = |x(s)|/A$ and $b = |y(s)|/B$ in \eqref{eq:I.3.4} and integrating, we obtain 
\[
\frac {\int | x (s) y (s) |}{A B} \leqq \frac {1}{p} \frac {A ^{p}}{A ^{p}} + \frac {1}{p^{\prime}} \frac {B ^ {p^{\prime}}}{B ^ {p^{\prime}}} = 1 \text{ which implies } \eqref{eq:I.3.5}.
\]

Next, by \eqref{eq:I.3.5}, we have 
\[
\begin{aligned}
&\int | x (s) + y (s) | ^{p} \leqq \int | x (s) + y (s) | ^{p-1} \cdot | x (s) | + \int | x (s) + y (s) | ^{p-1} \cdot | y (s) |
\\
&\leqq \left(\int | x (s) + y (s) | ^ {p^{\prime} (p- 1)}\right) ^ {1 / p^{\prime}} \left(\int | x (s) | ^{p}\right) ^{1/p}
\\
&\quad + \left(\int | x (s) + y (s) | ^ {p^{\prime} (p- 1)}\right) ^ {1 / p^{\prime}} \left(\int | y (s) | ^{p}\right) ^{1/p},
\end{aligned}
\]
which proves \eqref{eq:I.3.2} by $p'(p - 1) = p$.
\end{proof}

\begin{remarkn}[Remark $1$] The equality sign in \eqref{eq:I.3.2} holds iff there exists a non-negative constant $c$ such that $x(s) = cy(s)$ $m$-a.e. (or $y(s) = cx(s)$ $m$-a.e.). This is implied from the fact that, by Lemma~\ref{lem:I.3.Lemma.1}, the equality sign in Hölder's inequality \eqref{eq:I.3.5} holds iff $|x(s)| = c \cdot |y(s)|^{1/p-1}$ (or $|y(s)| = c \cdot |x(s)|^{1/p-1}$) are satisfied $m$-a.e. 
\end{remarkn}

\begin{remarkn}[Remark $2$] The condition $\|x\| = \left(\int |x(s)|^p\right)^{1/p} = 0$ is equivalent to the condition that $x(s) = 0$ $m$-a.e. We shall thus consider two functions of $L^p(S)$ as equivalent if they are equal $m$-a.e. By this convention, $L^p(S)$ becomes a normed linear space. The limit relation $\slim\limits_{n \to \infty} x_n = x$ in $L^p(S)$ is sometimes called the mean convergence of $p$-th order of the sequence of functions $x_n(s)$ to the function $x(s)$. 
\end{remarkn}

\begin{example}\label{ex:I.3.Example.3} $L^{\infty}(S)$. A $\mathfrak{B}$-measurable function $x(s)$ defined on $S$ is said to be essentially bounded if there exists a constant $\alpha$ such that $|x(s)| \leqq \alpha m$-a.e. The infimum of such constants $\alpha$ is denoted by 
\[
\vraimax_{s \in S} | x (s) | \text{ or essential } \sup\limits_{s \in S} | x (s) |.
\]
$L^{\infty}(S, \mathfrak{B}, m)$ or, in short, $L^{\infty}(S)$ is the set of all $\mathfrak{B}$-measurable, essentially bounded functions defined $m$-a.e. on $S$. It is a normed linear space by 
\[
(x + y) (s) = x (s) + y (s), (\alpha x) (s) = \alpha x (s), \| x \| = \vraimax_{s \in S} | x (s) |,
\]
under the convention that we consider two functions of $L^{\infty}(S)$ as equivalent if they are equal $m$-a.e. 
\end{example}

\begin{theorem}\label{thm:I.3.Theorem.1} Let the total measure $m(S)$ of $S$ be finite. Then we have 
\begin{equation}\label{eq:I.3.6}
\lim\limits_{p\rightarrow \infty} \left(\int_{S} | x (s) | ^{p} m (\dd s)\right) ^{1/p} = \vraimax_{s \in S} | x (s) | \quad \text{ for } \quad x (s) \in L ^{\infty} (S)
\end{equation}
\end{theorem}

\begin{proof}
 It is clear that
\begin{equation}\left(\int\limits_{S}\left|x(s)\right|^{p}m(\dd s)\right)^{1/p} \leqq m(S)^{1/p}\vraimax\limits_{s\in S}\left|x(s)\right|, \varlimsup\limits_{p\to\infty}\left(\int\limits_{S}\left|x(s)\right|^{p}\right)^{1/p} \leqq \vraimax\limits_{s\in S}\left|x(s)\right|. \notag\end{equation}
By the definition of the vrai max, there exists, for any $\varepsilon>0$, a set $B$ of $m$-measure $>0$ at each point of which $\left|x(s)\right|\geqq\vraimax\limits_{s\in S}\left|x(s)\right|-\varepsilon$. Hence
\begin{align}
\left(\int\limits_{S}\left|x(s)\right|^{p}m(\dd s)\right)^{1/p}& \geqq m(B)^{1/p}\left(\vraimax\limits_{s\in S}\left|x(s)\right|-\varepsilon\right), \notag
\end{align}
and therefore $\lim\limits_{p\to\infty}\left(\int\limits_{S}\left|x(s)\right|^{p}\right)^{1/p}\geqq\vraimax\limits_{s\in S}\left|x(s)\right|-\varepsilon$, and so \eqref{eq:I.3.6} is true. 
\end{proof}

\begin{example}\label{ex:I.3.Example.4} Let, in particular, $S$ be a discrete topological space consisting of countable points denoted by $1$, $2$, \ldots; the term discrete means that each point of $S = \{1, 2, \ldots\}$ is itself open in $S$. Then as linear subspaces of $C(\{1, 2, \ldots\})$, we define $(c_{0})$, $(c)$ and $(l^{p})$, $1 \leqq p< \infty$. 

$(c_{0})$: Consider a bounded sequence of real or complex numbers $\{\xi_{n}\}$. Such a sequence $\{\xi_{n}\}$ defines a function $x(n)=\xi_{n}$ defined and continuous on the discrete space $S=\{1,2,\ldots\}$; we shall call $x=\{\xi_{n}\}$ a vector with components $\xi_{n}$. The set of all vectors $x = \{\xi_{n}\}$ such that $\lim\limits_{n\to \infty}\xi_n = 0$ constitutes a normed linear space $(c_0)$ by the norm 
\[
\left\| x \right\| = \sup\limits_{n} | x (n) | = \sup\limits_{n} \left| \xi_{n} \right|.
\]
\end{example}

$(c)$: The set of all vectors $x = \{\xi_n\}$ such that finite $\lim\limits_{n \to \infty} \xi_n$ exist, constitutes a normed linear space $(c)$ by the norm $\|x\| = \sup\limits_n |x(n)| = \sup|\xi_n|$. 

$(l^{p}), 1 \leqq p< \infty$: The set of all vectors $x = \{\xi_{n}\}$ such that $\sum\limits_{n=1}^{\infty} |\xi_{n}|^{p} < \infty$ constitutes a normed linear space $(l^{p})$ by the norm $\|x\| = \left( \sum\limits_{n=1}^{\infty} |\xi_{n}|^{p} \right)^{1/p}$. As an abstract linear space, it is a linear subspace of $C(\{1, 2, \ldots\})$. It is also a special case of $L^{p}(S, \mathfrak{B}, m)$ in which $m(\{1\}) = m(\{2\}) = \cdots = 1$.

$(l^{\infty}) = (m)$: As in the case of $L^{\infty}(S)$, we shall denote by $(l^{\infty})$ the linear space $C(\{1, 2, \ldots\})$, normed by $\|x\| = \sup\limits_{n} |x(n)| = \sup\limits_{n} |\xi_n|$. $(l^{\infty})$ is also denoted by $(m)$. 

The Space of Measures. Let $\mathfrak{B}$ be a $\sigma$-ring of subsets of $S$. Consider the set $A(S, \mathfrak{B})$ of all real- (or complex-) valued functions $\varphi(B)$ defined on $\mathfrak{B}$ such that 
\begin{equation}\label{eq:I.3.7}
\left| \varphi (B) \right| \neq \infty \text{ for every } B \in \mathfrak {B},
\end{equation}
\begin{equation}\label{eq:I.3.8}
\varphi \left( \sum\limits_{j=1} ^{\infty} B_j \right) = \sum\limits_{j=1} ^{\infty} \varphi (B_j) \text{ for any disjoint sequence } \{B_j \} \text{ of sets } \in \mathfrak{B}.
\end{equation}


$A(S, \mathfrak{B})$ will be called the space of signed (or complex) measures defined on $(S, \mathfrak{B})$. 

\begin{lemma}\label{lem:I.3.Lemma.2} Let $\varphi \in A(S, \mathfrak{B})$ be real-valued. Then the total variation of $\varphi$ on $S$ defined by 
\begin{equation}\label{eq:I.3.9}
V (\varphi ; S) = \overline {{V}} (\varphi ; S) + | \underline {{V}} (\varphi ; S) |
\end{equation}
is finite; here the positive variation and the negative variation of $\varphi$ over $B\in\mathfrak{B}$ are given respectively by 
\begin{equation}\label{eq:I.3.10}
\overline {{V}} (\varphi ; B) = \sup\limits_ {B_{1} \subseteq B} \varphi (B_{1}) \text{ and } \underline {{V}} (\varphi ; B) = \inf\limits_ {B_{1} \subseteq B} \varphi (B_{1}).
\end{equation}
\end{lemma}

\begin{proof}
 Since $\varphi(\emptyset) = 0$, we have $V(\varphi; B) \geqq 0 \geqq \underline{V}(\varphi; B)$. Suppose that $V(\varphi; S) = \infty$. Then there exists a decreasing sequence $\{B_n\}$ of sets $\in \mathfrak{B}$ such that 
\[
V (\varphi ; B_{n}) = \infty , | \varphi (B_{n}) | \geqq n- 1. \]

The proof is obtained by induction. Let us choose $B_{1}=S$ and assume that the sets $B_{2}, B_{3}, \ldots, B_{k}$ have been defined so as to satisfy the above conditions. By the first condition with $n=k$, there exists a set $B \in B$ such that $B \subseteq B_k$, $|\varphi(B)| \geqq |\varphi(B_k)| + k$. We have only to set $B_{k+1} = B$ in the case $V(\varphi; B) = \infty$ and $B_{k+1} = B_k- B$ in the case $V(\varphi; B) < \infty$. For, in the latter case, we must have $V(\varphi; B_k- B) = \infty$ and $|\varphi(B_k- B)| \geqq |\varphi(B)|- |\varphi(B_k)| \geqq k$ which completes the induction. 

By the decreasing property of the sequence $\{B_n\}$, we have 
\[
\begin{aligned}
&S- \underset {n = 1} {\overset {\infty} {\cap}} B_{n}  = \underset {n = 1} {\overset {\infty} {\sum}} (S- B_{n}) = (S- B_{1}) + (B_{1}- B_{2}) + (B_{2}- B_{3}) + \dots + (B_{n}- B_{n+1}) + \dots
\end{aligned}
\]
so that, by the countable additivity of $\varphi$, 
\[
\begin{aligned} \varphi \left(S- \bigcap\limits_{n=1} ^{\infty} B_{n}\right) & = \varphi (S- B_{1}) + \varphi (B_{1}- B_{2}) + \varphi (B_{2}- B_{3}) + \dots \\ & = [ \varphi (S)- \varphi (B_{1}) ] + [ \varphi (B_{1})- \varphi (B_{2}) ] \\ & \quad + [ \varphi (B_{2})- \varphi (B_{3}) ] + \dots \\ & = \varphi (S)- \lim\limits_{n \to \infty} \varphi (B_{n}) = \infty \text{ or }- \infty , \end{aligned}
\]
which is a contradiction of \eqref{eq:I.3.7}. 
\end{proof}
\begin{theorem}[Jordan's decomposition]\label{thm:I.3.Theorem.2} Let $\varphi \in A(S, \mathfrak{B})$ be real-valued. Then the positive variation $\overline{V}(\varphi; B)$, the negative variation $\underline{V}(\varphi; B)$ and the total variation $V(\varphi; B)$ are countably additive on $B$. Moreover, we have the Jordan decomposition 
\begin{equation}\label{eq:I.3.11}
\varphi (B) = \bar {V} (\varphi ; B) + \underline {{{V}}} (\varphi ; B) \text{ for any } B \in \mathfrak {B}.
\end{equation}
\end{theorem}

\begin{proof}
 Let $\{B_n\}$ be a sequence of disjoint sets $\in \mathfrak{B}$. For any set $B \in \mathfrak{B}$ such that $B \subseteq \sum\limits_{n=1}^{\infty} B_n$, we have $\varphi(B) = \sum\limits_{n=1}^{\infty} \varphi(B \cap B_n) \leqq \sum\limits_{n=1}^{\infty} \overline{V}(\varphi; B_n)$ and hence $\overline{V}\left(\varphi; \sum\limits_{n=1}^{\infty} B_n\right) \leqq \sum\limits_{n=1}^{\infty} \overline{V}(\varphi; B_n)$. On the other hand, if $C_n \in \mathfrak{B}$ is a subset of $B_n$ ( $n = 1, 2, \ldots$), then we have $\overline{V}\left(\varphi; \sum\limits_{n=1}^{\infty} B_n\right) \geqq \varphi\left(\sum\limits_{n=1}^{\infty} C_n\right) = \sum\limits_{n=1}^{\infty} \varphi(C_n)$ and so $\overline{V}\left(\varphi; \sum\limits_{n=1}^{\infty} B_n\right) \geqq \sum\limits_{n=1}^{\infty} \overline{V}(\varphi; B_n)$. Hence we have proved the countable additivity of $\overline{V}(\varphi; B)$ and those of $\underline{V}(\varphi; B)$ and of $V(\varphi; B)$ may be proved similarly. 


To establish \eqref{eq:I.3.11}, we observe that, for every $C \in \mathfrak{B}$ with $C \subseteq B$, we have $\varphi(C) = \varphi(B)- \varphi(B- C) \leqq \varphi(B)- \underline{V}(\varphi; B)$ and so $\overline{V}(\varphi; B) \leqq \varphi(B)- \underline{V}(\varphi; B)$. Similarly we obtain $\underline{V}(\varphi; \overline{B}) \geqq \varphi(B)- \overline{V}(\varphi; B)$. These inequalities together give \eqref{eq:I.3.11}.
\end{proof}

\begin{theorem}[Hahn's decomposition]\label{thm:I.3.Theorem.3} Let $\varphi \in A(S, \mathfrak{B})$ be a signed measure. Then there exists a set $P\in \mathfrak{B}$ such that 

$\varphi(B) \geqq 0$ for every $B \in \mathfrak{B}$ with $B \subseteq P$, 
$\varphi(B) \leqq 0$ for every $B \in \mathfrak{B}$ with $B \subseteq P^{C} = S- P$. 
\end{theorem}

The decomposition $S = P+ (S- P)$ is called the Hahn decomposition of $S$ pertaining to $\varphi$. 

\begin{proof}
 For each positive integer $n$ we choose a set $B_n \in \mathfrak{B}$ such that $\varphi(B_n) \geqq \overline{V}(\varphi; S)- 2^{-n}$. Hence by \eqref{eq:I.3.11}, we have 
\begin{equation}\label{eq:I.3.12}
\underline {{{V}}} (\varphi ; B_{n}) \geqq- 2 ^{-n} \quad \text{ and } \quad \overline {{{V}}} (\varphi ; S- B_{n}) \leqq 2 ^{-n}. \end{equation}

The latter inequality is obtained from $\overline{V} (\varphi ;S- B_n) = \overline{V} (\varphi ;S)- \overline{V} (\varphi ;B_n)$ and $\overline{V} (\varphi ;B_n)\geqq \varphi (B_n)$. We then put 
\[
P= \varliminf\limits_{n \to \infty} B_{n} = \bigcup\limits_{k=1} ^{\infty} \bigcap\limits_{n=k} ^{\infty} B_{n}.
\]

Then $S- P= \varlimsup\limits_{n\to \infty}(S- B_n) = \bigcap\limits_{k=1}^{\infty}\bigcup\limits_{n=k}^{\infty}(S- B_n)\subseteq \bigcup\limits_{n=k}^{\infty}(S- B_n)$ every $k$, and therefore, by the $\sigma$-additivity of $\overline{V} (\varphi ;B)$, 
\[
\overline {{{V}}} (\varphi ; S- P) \leqq \sum\limits_{n=k} ^{\infty} \overline {{{V}}} (\varphi ; S- B_{n}) \leqq 2 ^{-(k-1)},
\]
which gives $\overline{V}(\varphi; S- P) = 0$. On the other hand, the negative variation $\underline{V}(\varphi; B)$ is a non-positive measure and so, by \eqref{eq:I.3.12} and similarly as above, 
\[
\left| \underline {{V}} (\varphi ; P) \right| \leqq \varliminf\limits_{n \rightarrow \infty} \left| \underline {{V}} (\varphi ; B_{n}) \right| = 0,
\]
which gives $\underline{V} (\varphi ;P) = 0$. The proof is thus completed. 
\end{proof}

\begin{corollaryu} The total variation $V(\varphi; S)$ of a signed measure $\varphi$ is defined by 
\begin{equation}\label{eq:I.3.13}
V (\varphi ; S) = \sup\limits_{\sup| x (s) | \leqq 1} \left| \int_{S} x (s) \varphi (\dd s) \right|
\end{equation}
where $x(s)$ ranges through $\mathfrak{B}$-measurable functions defined on $S$ such that $\sup\limits_{s}|x(s)| \leqq 1$. 
\end{corollaryu}

\begin{proof}
 If we take $x(s) = 1$ or =-$1$ according as $s \in P$ or $s \in S- P$, then the right hand side of \eqref{eq:I.3.13} gives $V(\varphi; S)$. On the other hand, it is easy to see that 
\[
\left| \int_{S} x (s) \varphi (\dd s) \right| \leqq \sup\limits_{s} | x (s) | \cdot \int_{S} V (\varphi ; \dd s) = \sup\limits_{s} | x (s) | \cdot V (\varphi ; S)
\]
and hence \eqref{eq:I.3.13} is proved. 
\end{proof}

\begin{example}\label{ex:I.3.Example.5} $A(S, \mathfrak{B})$. The space $A(S, \mathfrak{B})$ of signed measures $\varphi$ on $\mathfrak{B}$ is a real linear space by 
\[
\left(\alpha_{1} \varphi_{1} + \alpha_{2} \varphi_{2}\right) (B) = \alpha_{1} \varphi_{1} (B) + \alpha_{2} \varphi_{2} (B), B \in \mathfrak {B}.
\]
\end{example}

It is a normed linear space by the norm 
\begin{equation}\label{eq:I.3.14}
\left\| \varphi \right\| = V (\varphi ; S) = \sup\limits_{\sup| x (s) | \leqq 1} \left| \int_{S} x (s) \varphi (\dd s) \right|.
\end{equation}

\begin{example}\label{ex:I.3.Example.6} The space $A(S, \mathfrak{B})$ of complex measures $\varphi$ is a complex linear space by 

$(\alpha_{1}\varphi_{1} + \alpha_{2}\varphi_{2})(B) = \alpha_{1}\varphi_{1}(B) + \alpha_{2}\varphi_{2}(B), B \in \mathfrak{B}$ with complex $\alpha_{1}, \alpha_{2}$. 
\end{example}

It is a normed linear space by the norm 
\begin{equation}\label{eq:I.3.15}
\left\| \varphi \right\| = \sup\limits_{\sup| x (s) | \leqq 1} \left| \int_{S} x (s) \varphi (\dd s) \right|,
\end{equation}
where complex-valued $B$-measurable functions $x(s)$ defined on $S$ are taken into account. We shall call the right hand value of \eqref{eq:I.3.15} the total variation of $\varphi$ on $S$ and denote it by $V(\varphi; S)$. 


\yossection{4}{Examples of Quasi-normed Linear Spaces}
\label{sec:I.4}

\begin{example}\label{ex:I.4.Example.1} $\mathfrak{E}^k (\Omega)$. The linear space $\mathfrak{E}^k (\Omega)$, introduced in Chapter~\ref{chap:I}, \ref{sec:I.1}, is a quasi-normed linear space by the quasi-norm $\|x\| = d(x,0)$, where the distance $d(x,y)$ is as defined there. 
\end{example}

\begin{example}\label{ex:I.4.Example.2} $M(S, \mathfrak{B}, m)$. Let $m(S) < \infty$ and let $M(S, \mathfrak{B}, m)$ be the set of all complex-valued $\mathfrak{B}$-measurable functions $x(s)$ defined on $S$ and such that $|x(s)| < \infty$ $m$-a.e. Then $M(S, \mathfrak{B}, m)$ is a quasi-normed linear space by the algebraic operations 
\[
(x + y) (s) = x (s) + y (s), (\alpha x) (s) = \alpha x (s)
\]
and (under the convention that $x = y$ iff $x(s) = y(s)$ $m$-a.e.) 
\begin{equation}\label{eq:I.4.1}
\left\| x \right\| = \int_{S} | x (s) | (1 + | x (s) |) ^{-1} m (\dd s).
\end{equation}
\end{example}

The triangle inequality for the quasi-norm $\|x\|$ is clear from 
\[
\frac {| \alpha + \beta |}{1 + | \alpha + \beta |} \leqq \frac {| \alpha | + | \beta |}{1 + | \alpha | + | \beta |} \leqq \frac {| \alpha |}{1 + | \alpha |} + \frac {| \beta |}{1 + | \beta |}.
\]

The mapping $\{\alpha, x\} \to \alpha x$ is continuous by the following 

\begin{propositionu} The convergence $\slim\limits_{n \to \infty} x_n = x$ in $M(S, \mathfrak{B}, m)$ is equivalent to the asymptotic convergence (or the convergence in measure) in $S$ of the sequence of functions $\{x_n(s)\}$ to $x(s)$: 
\begin{equation}\label{eq:I.4.2}
\text{ For any } \varepsilon > 0, \lim\limits_{n \to \infty} m \{s \in S; | x (s)- x_{n} (s) | \geqq \varepsilon \} = 0.
\end{equation}
\end{propositionu}

\begin{proof}
 Clear from the inequality 
\[
\frac {\delta}{1 + \delta} m (B_{\delta}) \leqq \| x \| \leqq m (B_{\delta}) + \frac {\delta}{1 + \delta} m (S- B_{\delta}), B_{\delta} = \{s \in S; | x (s) | \geqq \delta \}. \]
\end{proof}

\begin{remarku} It is easy to see that the topology of $M(S, \mathfrak{B}, m)$ may also be defined by the quasi-norm 
\begin{equation}\tag{1'}\label{eq:I.4.1'}
\left\| x \right\| = \inf\limits_{\varepsilon > 0} \tan^{-1} [ \varepsilon + m \{s \in S; | x (s) | \geqq \varepsilon \} ].
\end{equation}
\end{remarku}

\begin{example}\label{ex:I.4.Example.3} $\mathfrak{D}_K(\Omega)$. The linear space $\mathfrak{D}_K(\Omega)$, introduced in Chapter~\ref{chap:I}, \ref{sec:I.1}, is a quasi-normed linear space by the quasi-norm $\|x\| = d(x,0)$, where the distance $d(x,y)$ is defined in Chapter~\ref{chap:I}, \ref{sec:I.1}. 
\end{example}


\yossection{5}{Pre-Hilbert Spaces}
\label{sec:I.5}

\begin{definition}\label{def:I.5.Definition.1} $A$ real or complex normed linear space $X$ is called a pre-Hilbert space if its norm satisfies the condition 
\begin{equation}\label{eq:I.5.1}
\left\| x + y \right\| ^{2} + \left\| x- y \right\| ^{2} = 2 \left(\left\| x \right\| ^{2} + \left\| y \right\| ^{2}\right).
\end{equation}
\end{definition}

\begin{theorem}[M. \textsc{Fréchet}-J. \textsc{von} \textsc{Neumann}-P. \textsc{Jordan}]\label{thm:I.5.Theorem.1} We define, in a real pre-Hilbert space $X$, 
\begin{equation}\label{eq:I.5.2}
(x, y) = 4 ^{-1} (\| x + y \| ^{2}- \| x- y \| ^{2}).
\end{equation}
\end{theorem}

Then we have the properties: 
\begin{align}
(\alpha x, y)& = \alpha (x, y) \quad (\alpha \in \mathbb{R}^{1}), \label{eq:I.5.3} \\
(x + y, z)& = (x, z) + (y, z), \label{eq:I.5.4} \\
(x, y)& = (y, x), \label{eq:I.5.5} \\
(x, x)& = \| x \| ^{2}. \label{eq:I.5.6}
\end{align}

\begin{proof}
 \eqref{eq:I.5.5} and \eqref{eq:I.5.6} are clear. We have, from \eqref{eq:I.5.1} and \eqref{eq:I.5.2}, 
\begin{align}
(x, z) + (y, z)& = 4 ^{-1} \left(\left\| x + z \right\| ^{2}- \left\| x- z \right\| ^{2} + \left\| y + z \right\| ^{2}- \left\| y- z \right\| ^{2}\right) \notag \\
&= 2 ^{-1} \left(\left\| \frac {x + y}{2} + z \right\| ^{2}- \left\| \frac {x + y}{2}- z \right\| ^{2}\right) \notag
\\
&= 2 \left(\frac {x + y}{2}, z\right).  \label{eq:I.5.7}
\end{align}

If we take $y = 0$, we obtain $(x, z) = 2\left(\frac{x}{2}, z\right)$, because $(0, z) = 0$ by \eqref{eq:I.5.2}. Hence, by \eqref{eq:I.5.7}, we obtain \eqref{eq:I.5.4}. Thus we see that \eqref{eq:I.5.3} holds for rational numbers $\alpha$ of the form $\alpha = m / 2^n$. In a normed linear space, $\|\alpha x + y\|$ and $\|\alpha x- y\|$ are continuous in $\alpha$. Hence, by \eqref{eq:I.5.2}, $(\alpha x, y)$ is continuous in $\alpha$. Therefore \eqref{eq:I.5.3} is proved for every real number $\alpha$. 
\end{proof}

\begin{corollaryn}[J. von \textsc{Neumann}-P. \textsc{Jordan}] We define, in a complex normed linear space $X$ satisfying \eqref{eq:I.5.1},
\begin{equation}(x, y) = (x, y)_{1} + i (x, i ^{*} y)_{1}, \text{ where } i = \sqrt {- 1}, (x, y)_{1} = 4 ^{-1} (\| x + y \| ^{2}- \| x- y \| ^{2}). \label{eq:I.5.8}\end{equation}

Then, we have \eqref{eq:I.5.4}, \eqref{eq:I.5.6} and 
\begin{align}
(\alpha x, y)& = \alpha (x, y) \quad (\alpha \in \mathbb{C}^{1}), \tag{3'}\label{eq:I.5.3'} \\
(x, y)& = \overline {{(y , x)}} (\text{ complex- conjugate number }). \tag{5'}\label{eq:I.5.5'}
\end{align}

\end{corollaryn}
\begin{proof}
 $X$ is also a real pre-Hilbert space and so \eqref{eq:I.5.4} and \eqref{eq:I.5.3} with real $\alpha$ hold good. We have, by \eqref{eq:I.5.8}, $(y, x)_1 = (x, y)_1$, $(ix, iy)_1 = (x, y)_1$ and hence $(y, ix)_1 = (-iiy, ix)_1 =-(iy, x)_1 =-(x, iy)_1$. Therefore 
\[
(y, x) = (y, x)_{1} + i (y, i x)_{1} = (x, y)_{1}- i (x, i y)_{1} = \overline {{(x , y)}}. \]

Similarly, we have 
\[
(i x, y) = (i x, y)_{1} + i (i x, i y)_{1} =- (x, i y)_{1} + i (x, y)_{1} = i (x, y),
\]
and therefore we have proved \eqref{eq:I.5.3}. Finally we have \eqref{eq:I.5.6}, because 
\[
(x, x)_{1} = \| x \| ^{2} \text{ and } (x, i x)_{1} = 4 ^{-1} (| 1 + i | ^{2}- | 1- i | ^{2}) \| x \| ^{2} = 0.
\]
\end{proof}

\begin{theorem}\label{thm:I.5.Theorem.2} $A$ (real or) complex linear space $X$ is a (real or) complex pre-Hilbert space, if to every pair of elements $x, y \in X$ there is associated a (real or) complex number $(x, y)$ satisfying $\eqref{eq:I.5.3}$, $\eqref{eq:I.5.4}$, $\eqref{eq:I.5.5}$ and 
\begin{equation}\label{eq:I.5.9}
(x, x) \geqq 0, \text{ and } (x, x) = 0 \text{ iff } x = 0.
\end{equation}
\end{theorem}

\begin{proof}
 For any real number $\alpha$, we have, by \eqref{eq:I.5.3}, \eqref{eq:I.5.4} and \eqref{eq:I.5.5} 
\[
\begin{aligned}
(x + \alpha (x, y) y, x + \alpha (x, y) y) &= \| x \| ^{2} + 2 \alpha | (x, y) | ^{2}
\\
\quad + \alpha^{2} | (x, y) | ^{2} \| y \| ^{2} \geqq 0, \text{ where } \| x \| &= (x, x) ^{1/2}
\end{aligned}
\]
so that we have $|(x, y)|^4- \|x\|^2 |(x, y)|^2 \|y\|^2 \leqq 0$. Hence we obtain Schwarz' inequality 
\begin{equation}\label{eq:I.5.10}
\left| (x, y) \right| \leqq \left\| x \right\| \cdot \left\| y \right\|,
\end{equation}
where the equality is satisfied iff $x$ and $y$ are linearly dependent. 

The latter part of \eqref{eq:I.5.10} is clear from the latter part of \eqref{eq:I.5.9}. 

We have, by \eqref{eq:I.5.10}, the triangle inequality for $\|x\|$: 
\[
\begin{aligned}
&{\| x + y \| ^{2} = (x + y, x + y) = \| x \| ^{2} + (x, y) + (y, x) + \| y \| ^{2}} {\qquad \leqq (\| x \| + \| y \|) ^{2}.}
\end{aligned}
\]

Finally, the equality \eqref{eq:I.5.1} is verified easily. 
\end{proof}

\begin{definition}\label{def:I.5.Definition.2} The number $(x, y)$ introduced above is called the scalar product (or inner product) of two vectors $x$ and $y$ of the pre-Hilbert space $X$. 
\end{definition}

\begin{example}\label{ex:I.5.Example.1} $L^2(S, \mathfrak{B}, m)$ is a pre-Hilbert space in which the scalar product is given by $(x, y) = \int x(s) \overline{y(s)} m(\dd s)$. 
\end{example}

\begin{example}\label{ex:I.5.Example.2} The normed linear space $(l^2)$ is a pre-Hilbert space in which the scalar product is given by $(\{\xi_n\}, \{\eta_n\}) = \sum\limits_{n=1}^{\infty} \xi_n \overline{\eta}_n$. 
\end{example}

\begin{example}\label{ex:I.5.Example.3} Let $\Omega$ be an open domain of $\mathbb{R}^{n}$ and $0 \leqq k < \infty$. Then the totality of functions $f \in C^k(\Omega)$ for which 
\begin{equation}\label{eq:I.5.11}
\left\| f \right\|_{k} = \left(\sum\limits_{| j | \leqq k} \int_{\Omega} \left| D ^{j} f (x) \right| ^{2} \dd x\right) ^{1/2} < \infty , \text{ where } \dd x = \dd x_{1} \dd x_{2} \dots \dd x_{n}
\end{equation}
constitutes a pre-Hilbert space $\hat{H}^k (\Omega)$ by the scalar product 
\begin{equation}\label{eq:I.5.12}
(f, g)_{k} = \sum\limits_{| j | \leqq k} \int_{\Omega} D ^{j} f (x) \cdot D ^{j} \overline {{g (x)}} \dd x.
\end{equation}
\end{example}

\begin{example}\label{ex:I.5.Example.4} Let $\Omega$ be an open domain of $\mathbb{R}^{n}$ and $0 \leqq k < \infty$. Then $C_0^k(\Omega)$ is a pre-Hilbert space by the scalar product \eqref{eq:I.5.12} and the norm \eqref{eq:I.5.11}. We shall denote this pre-Hilbert space by $\hat{H}_0^k(\Omega)$. 
\end{example}

\begin{example}\label{ex:I.5.Example.5} Let $G$ be a bounded open domain of the complex $z$-plane. Let $A^{2}(G)$ be the set of all holomorphic functions $f(z)$ defined in $G$ and such that 
\begin{equation}\label{eq:I.5.13}
\left\| f \right\| = \left(\iint_{G} | f (z) | ^{2} \dd x \dd y\right) ^{1/2} < \infty , (z = x + i y).
\end{equation}
\end{example}

Then $A^2(G)$ is a pre-Hilbert space by the norm \eqref{eq:I.5.13}, the scalar product 
\begin{equation}\label{eq:I.5.14}
(f, g) = \iint_{G} f (z) \overline {{g (z)}} \dd x \dd y
\end{equation}
and the algebraic operations 
\[
(f + g) (z) = f (z) + g (z), (\alpha f) (z) = \alpha f (z).
\]

\begin{example}\label{ex:I.5.Example.6} Hardy-Lebesgue class $H- L^2$. Let $H- L^2$ be the set of all functions $f(z)$ which are holomorphic in the unit disk $\{z; |z| < 1\}$ of the complex $z$-plane and such that 
\begin{equation}\label{eq:I.5.15}
\sup\limits_{0<r<1} \left(\int_{0} ^{2 \pi} | f (r e ^{i \theta}) | ^{2} \dd \theta\right) < \infty .
\end{equation}
\end{example}

Then, if $f(z) = \sum\limits_{n=0}^{\infty} c_n z^n$ is the Taylor expansion of $f$, 
\[
\begin{aligned}
&F (r)  = \frac {1}{2 \pi} \int_{0} ^{2 \pi} | f (r e ^{i \theta}) | ^{2} \dd \theta = \frac {1}{2 \pi} \sum\limits_{n,m=0} ^{\infty} \int_{0} ^{2 \pi} c_{n} \overline {{c}}_{m} r ^{n+m} e ^{i (n- m) \theta} \dd \theta = \sum\limits_{n=0} ^{\infty} | c_{n} | ^{2} r ^{2n}
\end{aligned}
\]
is monotone increasing in $r$, $0 < r < 1$, and bounded from above. Thus it is easy to see that 
\begin{equation}\label{eq:I.5.16}
\left\| f \right\| = \sup\limits_{0<r<1} \left[ \frac {1}{2 \pi} \left(\int_{0} ^{2 \pi} | f (r e ^{i \theta}) | ^{2} \dd \theta\right) \right] ^{1/2} = \left(\sum\limits_{n=0} ^{\infty} \left| c_{n} \right| ^{2}\right) ^{1/2}
\end{equation}
is a norm which satisfies condition \eqref{eq:I.5.1}, since $(l^2)$ is a pre-Hilbert space. 

\begin{remarku} Let a sequence $\{c_n\} \in (l^2)$ be given, and consider 
\[
f (z) = f \left(r e ^{i \theta}\right) = \sum\limits_{n=0} ^{\infty} c_{n} z ^{n} = \sum\limits_{n=0} ^{\infty} c_{n} r ^{n} e ^{i n \theta}, | z | < 1.
\]
\end{remarku}

By Schwarz' inequality, we have 
\[
\left| \sum\limits_{n=k} ^{\infty} c_{n} z ^{n} \right| \leqq \left(\sum\limits_{n=k} ^{\infty} \left| c_{n} \right| ^{2}\right) ^{1/2} \left(\sum\limits_{n=k} ^{\infty} r ^{2n}\right) ^{1/2},
\]
and so $\sum\limits_{n=0}^{\infty} c_n z^n$ is uniformly convergent in any disk $|z| \leqq \varrho$ with $0 < \varrho < 1$. Thus $f(z)$ is a holomorphic function in the unit disk $|z| < 1$ such that \eqref{eq:I.5.15} holds good, that is, $f(z)$ belongs to the class $H- L^2$. 

Therefore we have proved 

\begin{theorem}\label{thm:I.5.Theorem.3} The Hardy-Lebesgue class $H- L^2$ is in one-to-one correspondence with the pre-Hilbert space $(l^2)$ as follows: 
\[
H \cdot L ^{2} \ni f (z) = \sum\limits_{n=0} ^{\infty} c_{n} z ^{n} \leftrightarrow \{c_{n} \} \in (l ^{2})
\]
in such a way that 
\begin{align}
f (z)& = \sum\limits_{n=0} ^{\infty} c_{n} z ^{n} \leftrightarrow \{c_{n} \}, g (z) = \sum\limits_{n=0} ^{\infty} d_{n} z ^{n} \leftrightarrow \{d_{n} \} \text{ imply } \notag \\
f (z) + g (z)& \leftrightarrow \left\{c_{n} + d_{n} \right\}, \alpha f (z) \leftrightarrow \left\{\alpha c_{n} \right\} \text{ and } \| f \| = \left(\sum\limits_{n=0} ^{\infty} | c_{n} | ^{2}\right) ^{1/2}. \notag
\end{align}
\end{theorem}

Hence, as a pre-Hilbert space, $H- L^2$ is isomorphic with $(l^2)$. 


\yossection{6}{Continuity of Linear Operators}
\label{sec:I.6}

\begin{proposition}\label{prop:I.6.Proposition.1} Let $X$ and $Y$ be linear topological spaces over the same scalar field $K$. Then a linear operator $T$ on $D(T) \subseteq X$ into $Y$ is continuous everywhere on $D(T)$ iff it is continuous at the zero vector $x = 0$. 
\end{proposition}

\begin{proof}
 Clear from the linearity of the operator $T$ and $T \cdot 0 = 0$. 
\end{proof}

\begin{theorem}\label{thm:I.6.Theorem.1} Let $X$, $Y$ be locally convex spaces, and $\{p\}$, $\{q\}$ be the systems of semi-norms respectively defining the topologies of $X$ and $Y$. Then a linear operator $T$ on $D(T) \subseteq X$ into $Y$ is continuous iff, for every semi-norm $q \in \{q\}$, there exist a semi-norm $P\in \{p\}$ and a positive number $\beta$ such that 
\begin{equation}\label{eq:I.6.1}
q (T x) \leqq \beta P(x) \quad \text{ for all } \quad x \in D (T).
\end{equation}
\end{theorem}

\begin{proof}
 The condition is sufficient. For, by $T \cdot 0 = 0$, the condition implies that $T$ is continuous at the point $x = 0 \in D(T)$ and so $T$ is continuous everywhere on $D(T)$. 


The condition is necessary. The continuity of $T$ at $x = 0$ implies that, for every semi-norm $q \in \{q\}$ and every positive number $\varepsilon$, there exist a semi-norm $P\in \{p\}$ and a positive number $\delta$ such that 
\[
x \in D (T) \text{ and } P(x) \leqq \delta \text{ imply } q (T x) \leqq \varepsilon .
\]

Let $x$ be an arbitrary point of $D(T)$, and let us take a positive number $\lambda$ such that $\lambda p(x) \leqq \delta$. Then we have $p(\lambda x) \leqq \delta$, $\lambda x \in D(T)$ and so $q(T(\lambda x)) \leqq \varepsilon$. Thus $q(Tx) \leqq \varepsilon / \lambda$. Hence, if $p(x) = 0$, we can take $\lambda$ arbitrarily large and so $q(Tx) = 0$; and if $p(x) \neq 0$, we can take $\lambda = \delta / p(x)$ and so, in any case, we have $q(Tx) \leqq \beta p(x)$ with $\beta = \varepsilon / \delta$.
\end{proof}

\begin{corollary}\label{cor:I.6.Corollary.1} Let $X$ be a locally convex space, and $f$ a linear functional on $D(f) \subseteq X$. Then $f$ is continuous iff there exist a semi-norm $P$ from the system $\{p\}$ of semi-norms defining the topology of $X$ and a positive number $\beta$ such that 
\begin{equation}\label{eq:I.6.2}
\left| f (x) \right| \leqq \beta P(x) \text{ for all } x \in D (f).
\end{equation}
\end{corollary}

\begin{proof}
 For, the absolute value $|\alpha|$ itself constitutes a system of semi-norms defining the topology of the real or complex number field. 
\end{proof}

\begin{corollary}\label{cor:I.6.Corollary.2} Let $X$, $Y$ be normed linear spaces. Then a linear operator $T$ on $D(T) \subseteq X$ into $Y$ is continuous iff there exists a positive constant $\beta$ such that 
\begin{equation}\label{eq:I.6.3}
\left\| T x \right\| \leqq \beta \left\| x \right\| \text{ for all } x \in D (T).
\end{equation}
\end{corollary}

\begin{corollary}\label{cor:I.6.Corollary.3} Let $X$, $Y$ be normed linear spaces. Then a linear operator $T$ on $D(T) \subseteq X$ into $Y$ admits a continuous inverse $T^{-1}$ iff there exists a positive constant $\gamma$ such that 
\begin{equation}\label{eq:I.6.4}
\left\| T x \right\| \geqq \gamma \left\| x \right\| \text{ for every } x \in D (T).
\end{equation}
\end{corollary}

\begin{proof}
 By \eqref{eq:I.6.4}, $Tx = 0$ implies $x = 0$ and so the inverse $T^{-1}$ exists. The continuity of $T^{-1}$ is proved by \eqref{eq:I.6.4} and the preceding Corollary~\ref{cor:I.6.Corollary.2}. 
\end{proof}

\begin{definition}\label{def:I.6.Definition.1} Let $T$ be a continuous linear operator on a normed linear space $X$ into a normed linear space $Y$. We define 
\begin{equation}\label{eq:I.6.5}
\left\| T \right\| = \inf\limits_{\beta \in B} \beta , \text{ where } B = \{\beta ; \left\| T x \right\| \leqq \beta \left\| x \right\| \text{ for all } x \in X \}.
\end{equation}
\end{definition}

By virtue of the preceding Corollary~\ref{cor:I.6.Corollary.2} and the linearity of $T$, it is easy to see that 
\begin{equation}\label{eq:I.6.6}
\left\| T \right\| = \sup\limits_{\left\| x \right\| \leqq 1} \left\| T x \right\| = \sup\limits_{\left\| x \right\| = 1} \left\| T x \right\|.
\end{equation}
$\|T\|$ is called the norm of $T$. A continuous linear operator on a normed linear space $X$ into $Y$ is called a bounded linear operator on $X$ into $Y$, since, for such an operator $T$, the norm $\|Tx\|$ is bounded when $x$ ranges over the unit disk or the unit sphere $\{x \in X; \|x\| \leqq 1\}$ of $X$. 

\begin{definition}\label{def:I.6.Definition.2} Let $T$ and $S$ be linear operators such that 
\[
D (T) \text{ and } D (S) \subseteq X, \text{ and } R (T) \text{ and } R (S) \subseteq Y.
\]
\end{definition}

Then the sum $T + S$ and the scalar multiple $\alpha T$ are defined respectively by 
\[
(T + S) (x) = T x + S x \quad \text{ for } \quad x \in D (T) \cap D (S), (\alpha T) (x) = \alpha (T x)
\]

Let $T$ be a linear operator on $D(T) \subseteq X$ into $Y$, and $S$ a linear operator on $D(S) \subseteq Y$ into $Z$. Then the product $ST$ is defined by 
\[
(S T) x = S (T x) \quad \text{ for } \quad x \in \{x; x \in D (T) \quad \text{ and } \quad T x \in D (S) \}.
\]
$T + S, \alpha T$ and $ST$ are linear operators. 

\begin{remarku} $ST$ and $TS$ do not necessarily coincide even if $X = Y = Z$. An example is given by $Tx = tx(t)$, $Sx(t) = \sqrt{(-1)}\frac{d}{\dd t} x(t)$ considered as linear operators from $L^2(\mathbb{R}^{1})$ into $L^2(\mathbb{R}^{1})$. In this example, we have the commutation relation $(ST- TS)x(t) = \sqrt{-1} x(t)$. 
\end{remarku}

\begin{proposition}\label{prop:I.6.Proposition.2} If $T$ and $S$ are bounded linear operators on a normed linear space $X$ into a normed linear space $Y$, then 
\begin{equation}\label{eq:I.6.7}
\left\| T + S \right\| \leqq \left\| T \right\| + \left\| S \right\|, \left\| \alpha T \right\| = | \alpha | \| T \|.
\end{equation}
\end{proposition}

If $T$ is a bounded linear operator on a normed linear space $X$ into a normed linear space $Y$, and $S$ a bounded linear operator on $Y$ into a normed linear space $Z$, then 
\begin{equation}\label{eq:I.6.8}
\left\| S T \right\| \leqq \left\| S \right\| \cdot \left\| T \right\|.
\end{equation}

\begin{proof}
 We prove the last inequality; \eqref{eq:I.6.7} may be proved similarly. $\|STx\| \leqq \|S\| \|Tx\| \leqq \|S\| \|T\| \|x\|$ and so $\|ST\| \leqq \|S\| \|T\|$. 
\end{proof}

\begin{corollaryu} If $T$ is a bounded linear operator on a normed linear space $X$ into $X$, then 
\begin{equation}\label{eq:I.6.9}
\left\| T ^{n} \right\| \leqq \left\| T \right\| ^{n},
\end{equation}
where $T^{n}$ is defined inductively by $T^{n}=TT^{n-1}$ ( $n=1,2,\ldots;T^{0}=I$ which maps every $x$ onto $x$ itself, i.e., $Ix=x$, and $I$ is called the identity operator). 
\end{corollaryu}


\yossection{7}{Bounded Sets and Bornologic Spaces}
\label{sec:I.7}

\begin{definition}\label{def:I.7.Definition.1} $A$ subset $B$ in a linear topological space $X$ is said to be bounded if it is absorbed by any neighbourhood $U$ of $0$, i.e., if there exists a positive constant $\alpha$ such that $\alpha^{-1}B \subseteq U$. Here $\alpha^{-1}B = \{x \in X; x = \alpha^{-1}b, b \in B\}$. 
\end{definition}

\begin{propositionu} Let $X$, $Y$ be linear topological spaces. Then a continuous linear operator on $X$ into $Y$ maps every bounded set of $X$ onto a bounded set of $Y$. 
\end{propositionu}

\begin{proof}
 Let $B$ be a bounded set of $X$, and $V$ a neighbourhood of $0$ of $Y$. By the continuity of $T$, there exists a neighbourhood $U$ of $0$ of $X$ such that $T \cdot U = \{Tu; u \in U\} \subseteq V$. Let $\alpha > 0$ be such that $B \subseteq \alpha U$. Then $T \cdot B \subseteq T(\alpha U) = \alpha(T \cdot U) \subseteq \alpha V$. This proves that $T \cdot B$ is a bounded set of $Y$. 
\end{proof}

\begin{definition}\label{def:I.7.Definition.2} $A$ locally convex space $X$ is called bornologic if it satisfies the condition: 

\begin{equation}\label{eq:I.7.1}
\begin{minipage}[b]{0.92\linewidth}
If a balanced convex set $M$ of $X$ absorbs every bounded set of $X$, then $M$ is a neighbourhood of $0$ of $X$.
\end{minipage}
\end{equation}
\end{definition}


\begin{theorem}\label{thm:I.7.Theorem.1} $A$ locally convex space $X$ is bornologic iff every semi-norm on $X$, which is bounded on every bounded set, is continuous. 
\end{theorem}

\begin{proof}
 We first remark that a semi-norm $p(x)$ on $X$ is continuous iff it is continuous at $x = 0$. This we see from the subadditivity of the semi-norm: $p(x- y) \geqq |p(x)- p(y)|$ (Chapter~\ref{chap:I}, \ref{sec:I.1}, \eqref{eq:I.1.4}). 

Necessity. Let a semi-norm $p(x)$ on $X$ be bounded on every bounded set of $X$. The set $M = \{x \in X; p(x) \leqq 1\}$ is convex and balanced. If $B$ is a bounded set of $X$, then $\sup\limits_{b \in B} p(b) = \alpha < \infty$ and therefore $B \subseteq \alpha M$. Since, by the assumption, $X$ is bornologic, $M$ must be a neighbourhood of $0$. Thus we see that $P$ is continuous at $x = 0$. 

Sufficiency. Let $M$ be a convex, balanced set of $X$ which absorbs every bounded set of $X$. Let $P$ be the Minkowski functional of $M$. Then $P$ is bounded on every bounded set, since $M$ absorbs, by the assumption, every bounded set. Hence, by the hypothesis, $p(x)$ is continuous. Thus $M_{1} = \{x \in X; p(x) < 1/2\}$ is an open set $\ni 0$ contained in $M$. This proves that $M$ is a neighbourhood of $0$. 
\end{proof}

\begin{example}\label{ex:I.7.Example.1} Normed linear spaces are bornologic. 
\end{example}

\begin{proof}
 Let $X$ be a normed linear space. Then the unit disk $S = \{x \in X; \| x \| \leqq 1\}$ of $X$ is a bounded set of $X$. Let a semi-norm $p(x)$ on $X$ be bounded on $S$, i.e., $\sup\limits_{x \in S} p(x) = \alpha < \infty$. Then, for any $y \neq 0$, 
\[
P(y) = P\left(\left\| y \right\| \cdot \frac {y}{\left\| y \right\|}\right) = \left\| y \right\| P\left(\frac {y}{\left\| y \right\|}\right) \leqq \alpha \left\| y \right\|. \]

Thus $p$ is continuous at $y = 0$ and so continuous at every point of $X$. 
\end{proof}

\begin{remarku} As will be seen later, the quasi-normed linear space $M(S, \mathfrak{B})$ is not locally convex. Thus a quasi-normed linear space is not necessarily bornologic. However we can prove 
\end{remarku}

\begin{theorem}\label{thm:I.7.Theorem.2} $A$ linear operator $T$ on one quasi-normed linear space into another such space is continuous iff $T$ maps bounded sets into bounded sets. 
\end{theorem}

\begin{proof}
 As was proved in Chapter~\ref{chap:I}, \ref{sec:I.2}, Proposition~\hyperref[prop:I.2.Proposition.2]{$2$}, a quasi-normed linear space is a linear topological space. Hence the ``only if'' part is already proved above in the Proposition. We shall prove the ``if'' part. 

Let $T$ maPbounded sets into bounded sets. Suppose that $\slim\limits_{k\to \infty}x_k = 0$. Then $\lim\limits_{k\to \infty}\| x_k\| = 0$ and so there exists a sequence of integers $\{n_k\}$ such that $\lim\limits_{k\to \infty}\| x_k\| = 0$, while $\lim\limits_{k\to \infty}\| x_k\| = 0$. 
\[
\lim\limits_{k \rightarrow \infty} n_{k} = \infty \quad \text{ while } \quad \lim\limits_{k \rightarrow \infty} n_{k} \| x_{k} \| = 0.
\]

We may take, for instance, $n_k$ as follows: 
$n_{k}  =$ the largest integer $\leqq \| x_{k} \| ^{-1/2}$ if $x_{k} \neq 0, = k$ if $x_{k} = 0.$


Now we have $\left\|n_{k}x_{k}\right\|=\left\|x_{k}+x_{k}+\cdots+x_{k}\right\|\leqq n_{k}\left\|x_{k}\right\|$ so that $\slim\limits_{k\to\infty}n_{k}x_{k}=0$. But, in a quasi-normed linear space, the sequence $\{n_{k}x_{k}\}$, which converges to $0$, is bounded. Thus, by the hypothesis, $\{T(n_{k}x_{k})\}=\{n_{k}Tx_{k}\}$ is a bounded sequence. Therefore 
\[
\slim_{k \to \infty} T x_{k} = \slim_{k \to \infty} n_{k} ^{-1} (T (n_{k} x_{k})) = 0,
\]
and so $T$ is continuous at $x = 0$ and hence is continuous everywhere. 
\end{proof}

\begin{theorem}\label{thm:I.7.Theorem.3} Let $X$ be bornologic. If a linear operator $T$ on $X$ into a locally convex linear topological space $Y$ maps every bounded set into a bounded set, then $T$ is continuous. 
\end{theorem}

\begin{proof}
 Let $V$ be a convex balanced neighbourhood of $0$ of $Y$. Let $p$ be the Minkowski functional of $V$. Consider $q(x) = p(Tx)$. $q$ is a semi-norm on $X$ which is bounded on every bounded set of $X$, because every bounded set of $Y$ is absorbed by the neighbourhood $V$ of $0$. Since $X$ is bornologic, $q$ is continuous. Thus the set $\{x \in X; Tx \in V^a\} = \{x \in X; q(x) \leqq 1\}$ is a neighbourhood of $0$ of $X$. This proves that $T$ is continuous. 
\end{proof}


\yossection{8}{Generalized Functions and Generalized Derivatives}
\label{sec:I.8}

A continuous linear functional defined on the locally convex linear topological space $\mathfrak{D}(\Omega)$, introduced in Chapter~\ref{chap:I}, \ref{sec:I.1}, is the ``distribution'' or the ``generalized function'' of L. \textsc{Schwartz}. To discuss the generalized functions, we shall begin with the proof of 

\begin{theorem}\label{thm:I.8.Theorem.1} Let $B$ be a bounded set of $\mathfrak{D}(\Omega)$. Then there exists a compact subset $K$ of $\Omega$ such that 
\begin{align}
\supp (\varphi)& \subseteq K \text{ for every } \varphi \in B , \label{eq:I.8.1} \\
\sup\limits_{x \in K, \varphi \in B} | D ^{j} \varphi (x) |& < \infty \text{ for every differential operator } D ^{j}. \label{eq:I.8.2}
\end{align}
\end{theorem}

\begin{proof}
 Suppose that there exist a sequence of functions $\{\varphi_i\} \subseteq B$ and a sequence of points $\{p_i\}$ such that (i): $\{p_i\}$ has no accumulation point in $\Omega$, and (ii): $\varphi_i(p_i) \neq 0$ ( $i = 1, 2, \ldots$). Then
\[
P(\varphi) = \sum\limits_{i=1} ^{\infty} i | \varphi (P_{i}) | / | \varphi_{i} (P_{i}) |
\]
is a continuous semi-norm on every $\mathfrak{D}_K(\Omega)$, defined in Chapter~\ref{chap:I}, \ref{sec:I.1}. Hence, for any $\varepsilon > 0$, the set $\{\varphi \in \mathfrak{D}_K(\Omega); p(\varphi) \leqq \varepsilon\}$ is a neighbourhood of $0$ of $\mathfrak{D}_K(\Omega)$. Since $\mathfrak{D}(\Omega)$ is the inductive limit of $\mathfrak{D}_K(\Omega)$'s, we see that $\{\varphi \in \mathfrak{D}(\Omega); p(\varphi) \leqq \varepsilon\}$ is also a neighbourhood of $0$ of $\mathfrak{D}(\Omega)$. Thus $p$ is continuous at $0$ of $\mathfrak{D}(\Omega)$ and so is continuous on $\mathfrak{D}(\Omega)$. Hence $p$ must be bounded on the bounded set $B$ of $\mathfrak{D}(\Omega)$. However, $p(\varphi_i) \geqq i (i = 1, 2, \ldots)$. This proves that we must have \eqref{eq:I.8.1}. 

We next assume that \eqref{eq:I.8.1} is satisfied, and suppose \eqref{eq:I.8.2} is not satisfied. Then there exist a differential operator $D^{j_0}$ and a sequence of functions $\{\varphi_i\} \subseteq B$ such that $\sup\limits_{x \in K} |D^{j_0} \varphi_i(x)| > i$ ( $i = 1, 2, \ldots$). Thus, if we set 
\[
\phi (\varphi) = \sup\limits_{x \in K} | D ^ {j_{0}} \varphi (x) | \quad \text{ for } \quad \varphi \in \mathfrak {D}_{K} (\Omega),
\]
$p(\varphi)$ is a continuous semi-norm on $\mathfrak{D}_K(\Omega)$ and $p(\varphi_i) > i (i = 1,2,\ldots)$. Hence $\{\varphi_i\} \subseteq B$ cannot be bounded in $\mathfrak{D}_K(\Omega)$, and a fortiori in $\mathfrak{D}(\Omega)$. This contradiction proves that \eqref{eq:I.8.2} must be true. 
\end{proof}
\begin{theorem}\label{thm:I.8.Theorem.2} The space $\mathfrak{D}(\Omega)$ is bornologic. 
\end{theorem}

\begin{proof}
 Let $q(\varphi)$ be a semi-norm on $\mathfrak{D}(\Omega)$ which is bounded on every bounded set of $\mathfrak{D}(\Omega)$. In view of Theorem~\ref{thm:I.7.Theorem.1} in Chapter~\ref{chap:I}, \ref{sec:I.7}, we have only to show that $q$ is continuous on $\mathfrak{D}(\Omega)$. To this purpose, we show that $q$ is continuous on the space $\mathfrak{D}_K(\Omega)$ where $K$ is any compact subset of $\Omega$. Since $\mathfrak{D}(\Omega)$ is the inductive limit of $\mathfrak{D}_K(\Omega)'s$, we then see that $q$ is continuous on $\mathfrak{D}(\Omega)$. 

But $q$ is continuous on every $\mathfrak{D}_K(\Omega)$. For, by hypothesis, $q$ is bounded on every bounded set of the quasi-normed linear space $\mathfrak{D}_K(\Omega)$, and so, by Theorem~\ref{thm:I.8.Theorem.2} of the preceding section, $q$ is continuous on $\mathfrak{D}_K(\Omega)$. Hence $q$ must be continuous on $\mathfrak{D}(\Omega)$. 
\end{proof}

We are now ready to define the generalized functions. 

\begin{definition}\label{def:I.8.Definition.1} $A$ linear functional $T$ defined and continuous on $\mathfrak{D}(\Omega)$ is called a generalized function, or an ideal function or a distribution in $\Omega$; and the value $T(\varphi)$ is called the value of the generalized function $T$ at the testing function $\varphi \in \mathfrak{D}(\Omega)$. 
\end{definition}

By virtue of Theorem~\ref{thm:I.7.Theorem.1} in Chapter~\ref{chap:I}, \ref{sec:I.7} and the preceding Theorem~\ref{thm:I.8.Theorem.2}, we have 

\begin{proposition}\label{prop:I.8.Proposition.1} $A$ linear functional $T$ defined on $\mathfrak{D}(\Omega)$ is a generalized function in $\Omega$ iff it is bounded on every bounded set of $\mathfrak{D}(\Omega)$, that is, iff $T$ is bounded on every set $B \in \mathfrak{D}(\Omega)$ satisfying the two conditions \eqref{eq:I.8.1} and \eqref{eq:I.8.2}. 
\end{proposition}

\begin{proof}
 Clear from the fact that $T(\varphi)$ is continuous iff the semi-norm $|T(\varphi)|$ is continuous. 
\end{proof}

\begin{corollaryu} $A$ linear functional $T$ defined on $C_{0}^{\infty}(\Omega)$ is a generalized function in $\Omega$ iff it satisfies the condition: 

To every compact subset $K$ of $\Omega$, there correspond a 
positive constant $C$ and a positive integer $k$ such that 
\begin{equation}\label{eq:I.8.3}
\left| T (\varphi) \right| \leqq C \sup\limits_{| j | \leqq k, x \in K} \left| D ^{j} \varphi (x) \right| \quad \text{ whenever } \quad \varphi \in \mathfrak {D}_{K} (\varphi).
\end{equation}
\end{corollaryu}

\begin{proof}
 By the continuity of $T$ on the inductive limit $\mathfrak{D}(\Omega)$ of the $\mathfrak{D}_K(\Omega)'s$, we see that $T$ must be continuous on every $\mathfrak{D}_K(\Omega)$. Hence the necessity of condition \eqref{eq:I.8.3} is clear. The sufficiency of the condition \eqref{eq:I.8.3} is also clear, since it implies that $T$ is bounded on every bounded set of $\mathfrak{D}(\Omega)$. 
\end{proof}

\begin{remarku} The above Corollary is very convenient for all applications, since it serves as a useful definition of the generalized functions. 
\end{remarku}

\begin{example}\label{ex:I.8.Example.1} Let a complex-valued function $f(x)$ defined a.e. in $\Omega$ be locally integrable in $\Omega$ with respect to the Lebesgue measure $\dd x = \dd x_1 \dd x_2 \cdots \dd x_n$ in $\mathbb{R}^{n}$, in the sense that, for any compact subset $K$ of $\Omega$, $\int |f(x)| \dd x < \infty$. Then 
\begin{equation}\label{eq:I.8.4}
T_{f} (\varphi) = \int_{\Omega} f (x) \varphi (x) \dd x, \varphi \in \mathfrak {D} (\Omega),
\end{equation}
defines a generalized function $T_{f}$ in $\Omega$. 
\end{example}

\begin{example}\label{ex:I.8.Example.2} Let $m(B)$ be a $\sigma$-finite, $\sigma$-additive and complex-valued measure defined on Baire subsets $B$ of an open set $\Omega$ of $\mathbb{R}^{n}$. Then 
\begin{equation}\label{eq:I.8.5}
T_{m} (\varphi) = \int_{\Omega} \varphi (x) m (\dd x), \varphi \in \mathfrak {D} (\Omega),
\end{equation}
defines a generalized function $T_{m}$ in $\Omega$. 
\end{example}

\begin{example}\label{ex:I.8.Example.3} As a special case of Example~\hyperref[ex:I.8.Example.2]{$2$}, 
\begin{equation}\label{eq:I.8.6}
T_{\delta_{p}}(\varphi) = \varphi (p), \text{ where } p \text{ is a fixed point of } \Omega, \varphi \in \mathfrak{D}(\Omega)
\end{equation}
defines a generalized function $T_{\delta_p}$ in $\Omega$. It is called the Dirac distribution concentrated at the point $p\in \Omega$. In the particular case $p= 0$, the origin of $\mathbb{R}^{n}$, we shall write $T_{\delta}$ or $\delta$ for $T_{\delta_0}$. 
\end{example}

\begin{definition}\label{def:I.8.Definition.2} The set of all generalized functions in $\Omega$ will be denoted by $\mathfrak{D}(\Omega)'$. It is a linear space by 
\begin{equation}\label{eq:I.8.7}
(T + S) (\varphi) = T (\varphi) + S (\varphi), (\alpha T) (\varphi) = \alpha T (\varphi),
\end{equation}
and we call $\mathfrak{D}(\Omega)'$ the space of the generalized functions in $\Omega$ or the dual space of $\mathfrak{D}(\Omega)$. 
\end{definition}

\begin{remarku} Two distributions $T_{f_1}$ and $T_{f_2}$ are equal as functionals $(T_{f_1}(\varphi) = T_{f_2}(\varphi)$ for every $\varphi \in \mathfrak{D}(\Omega)$) iff $f_1(x) = f_2(x)$ a.e. If this fact is proved, then the set of all locally integrable functions in $\Omega$ is, by $f \leftrightarrow T_f$, in a one-one correspondence with a subset of $\mathfrak{D}(\Omega)'$ in such a way that $(f_1 \text{ and } f_2 \text{ being considered equivalent iff } f_1(x) = f_2(x) \text{ a.e. })$ 
\begin{equation}\tag{7'}\label{eq:I.8.7'}
T_ {f_{1}} + T_ {f_{2}} = T_ {f_{1} + f_{2}}, \alpha T_{f} = T_{\alpha f}.
\end{equation}
\end{remarku}

In this sense, the notion of the generalized function is, in fact, a generalization of the notion of the locally integrable function. To prove the above assertion, we have only to prove that a locally integrable function $f$ is $= 0$ a.e. in an open set $\Omega$ of $\mathbb{R}^{n}$ if $\int_{B} f(x) \varphi(x) \dd x = 0$ for all $\varphi \in C_0^\infty(\Omega)$. By introducing the Baire measure $\mu(B) = \int_{B} f(x) \dd x$, the latter condition implies that $\int_{\Omega} \varphi(x) \mu(\dd x) = 0$ for all $\varphi \in C_0^\infty(\Omega)$, which further implies that $\int_{\Omega} \varphi(x) \mu(\dd x) = 0$ for all $\varphi \in C_0^0(\Omega)$, by virtue of Proposition~\ref{prop:I.1.Proposition.8} in Chapter~\ref{chap:I}, \ref{sec:I.1}. Let $B$ be a compact $G_\delta$-set in $\Omega$: $B = \bigcap\limits_{n=1}^{\infty} G_n$, where $G_n$ is an open relatively compact set in $\Omega$. By applying Urysohn's theorem in Chapter~\ref{chap:0}, \ref{sec:0.2}, there exists a continuous function $f_n(x)$ such that 
\[
\begin{aligned} 0 \leqq f_{n} (x) \leqq 1 & \text{ for } x \in \Omega , f_{n} (x) = 1 \text{ for } x \in G_{n+2} ^{a} \text{ and } f_{n} (x) = 0 \\ & \text{ for } x \in G_{n} ^{a}- G_{n+1} (n = 1, 2, \dots), \end{aligned}
\]
assuming that $\{G_n\}$ is a monotone decreasing sequence of open relatively compact sets of $\Omega$ such that $G_{n+2}^a \subseteq G_{n+1}$. Setting $\varphi = f_n$ and letting $n \to \infty$, we see that $\mu(B) = 0$ for all compact $G_\delta$-sets $B$ of $\Omega$. The Baire sets of $\Omega$ are the members of the smallest $\sigma$-ring containing compact $G_\delta$-sets of $\Omega$, we see, by the $\sigma$-additivity of the Baire measure $\mu$, that $\mu$ vanishes for every Baire set of $\Omega$. Hence the density $f$ of this measure $\mu$ must vanish a.e. in $\Omega$. 

We can define the notion of differentiation of generalized functions through 

\begin{proposition}\label{prop:I.8.Proposition.2} If $T$ is a generalized function in $\Omega$, then 
\begin{equation}\label{eq:I.8.8}
S (\varphi) =- T \left(\frac {\partial \varphi}{\partial x_{1}}\right), \quad \varphi \in \mathfrak {D} (\Omega),
\end{equation}
defines another generalized function $S$ in $\Omega$. 
\end{proposition}

\begin{proof}
 $S$ is a linear functional on $\mathfrak{D}(\Omega)$ which is bounded on every bounded set of $\mathfrak{D}(\Omega)$. 
\end{proof}

\begin{definition}\label{def:I.8.Definition.3} The generalized functional $S$ defined by \eqref{eq:I.8.8} is called the generalized derivative or the distributional derivative of $T$ (with respect to $x_{1}$), and we write 
\begin{equation}\label{eq:I.8.9}
S = \frac {\partial}{\partial x_{1}} T,
\end{equation}
so that we have 
\begin{equation}\label{eq:I.8.10}
\frac {\partial}{\partial x_{1}} T (\varphi) =- T \left(\frac {\partial \varphi}{\partial x_{1}}\right).
\end{equation}
\end{definition}

\begin{remarku} The above notion is an extension of the usual notion of the derivative. For, if the function $f$ is continuously differentiable with respect to $x_{1}$, then we have 
\[
\begin{aligned} \frac {\partial}{\partial x_{1}} T_{f} (\varphi) & =- T_{f} \left(\frac {\partial \varphi}{\partial x_{1}}\right) =- \int \dots \int_{\Omega} f (x) \frac {\partial \varphi}{\partial x_{1}} \dd x_{1} \dots \dd x_{n} \\ & = \int \dots \int_{\Omega} \frac {\partial}{\partial x_{1}} f (x) \cdot \varphi (x) \dd x_{1} \dots \dd x_{n} = T_ {\partial f / \partial x_{1}} (\varphi), \end{aligned}
\]
\end{remarku}

as may be seen by partial integration observing that $\varphi(x)$ vanishes identically outside some compact subset of $\Omega$. 

\begin{corollaryu} $A$ generalized function $T$ in $\Omega$ is infinitely differentiable in the sense of distributions defined above and 
\begin{equation}\label{eq:I.8.11}
(D ^{j} T) (\varphi) = (- 1) ^{|j|} T (D ^{j} \varphi), \text{ where } | j | = \sum\limits_{i=1} ^{n} j_{i}, D ^{j} = \frac {\partial^{|j|}}{\partial x_{1} ^ {j_{1}} \dots \partial x_{n} ^ {j_{n}}}.
\end{equation}
\end{corollaryu}

\begin{example}\label{ex:I.8.Example.1a} The Heaviside function $H(x)$ is defined by 
\begin{equation}\label{eq:I.8.12}
H (x) = 1 \text{ or } = 0 \text{ according as } x \geqq 0 \text{ or } x < 0.
\end{equation}
\end{example}

Then we have 
\begin{equation}\tag{12'}\label{eq:I.8.12'}
\frac {d}{\dd x} T_{H} = T_ {\delta_{0}},
\end{equation}
where $T_{\delta_0}$ is the Dirac distribution concentrated at the origin $0$ of $\mathbb{R}^{1}$. In fact, we have, for any $\varphi \in \mathfrak{D}(\mathbb{R}^{1})$, 
\[
\left(\frac {d}{\dd x} T_{H}\right) (\varphi) =- \int_{- \infty} ^{\infty} H (x) \varphi^{\prime} (x) \dd x =- \int_{0} ^{\infty} \varphi^{\prime} (x) \dd x =- [ \varphi (x) ]_{0} ^{\infty} = \varphi (0).
\]

\begin{example}\label{ex:I.8.Example.2a} Let $f(x)$ have a bounded and continuous derivative in the open set $\mathbb{R}^{1}- \bigcup\limits_{j=1}^{k} x_j$ of $\mathbb{R}^{1}$. Let $s_j = f(x_j + 0)- f(x_j- 0)$ be the saltus or the jump of $f(x)$ at $x = x_j$. Since 
\[
\left(\frac {d}{\dd x} T_{f}\right) (\varphi) =- \int_{- \infty} ^{\infty} f (x) \varphi^{\prime} (x) \dd x = \sum\limits_{j} \varphi (x_{j}) s_{j} + \int_{- \infty} ^{\infty} f ^{\prime} (x) \varphi (x) \dd x,
\]
we have 
\begin{equation}\tag{12''}\label{eq:I.8.12a}
\frac {d}{\dd x} T_{f} = T_ {f ^{\prime}} + \sum\limits_{j} s_{j} \delta_{xj},
\end{equation}
where $\delta_{x_{j}}$ is defined by \eqref{eq:I.8.6}. 
\end{example}

\begin{example}\label{ex:I.8.Example.3a} Let $f(x) = f(x_1, x_2, \ldots, x_n)$ be a continuously differentiable function on a closed bounded domain $\Omega \subseteq \mathbb{R}^{n}$ having a smooth boundary $S$. Define $f$ to be $0$ outside $\Omega$. By partial integration, we have 
\[
\begin{aligned}
&\left(\frac {\partial}{\partial x_{j}} T_{f}\right) (\varphi)  =- \int_{\Omega} f (x) \frac {\partial}{\partial x_{j}} \varphi (x) \dd x = \int_{S} f (x) \varphi (x) \cos (\nu , x_{j}) \dd S + \int_{\Omega} \frac {\partial f}{\partial x_{j}} \varphi (x) \dd x,
\end{aligned}
\]
where $\pmb{\nu}$ is the inner normal to $S$, $(\pmb{\nu}, x_j) = (x_j, \pmb{\nu})$ is the angle between $\pmb{\nu}$ and the positive $x_j$-axis and $dS$ is the surface element. We have thus 
\begin{equation}\tag{12''}\label{eq:I.8.12''}
\frac {\partial}{\partial x_{j}} T_{f} = T_ {\partial f / \partial x_{j}} + T_{S}, \text{ where } T_{S} (\varphi) = \int_{S} f (x) \cos (\nu , x_{j}) \varphi (x) \dd S.
\end{equation}
\end{example}

\begin{corollaryu} If $f(x) = f(x_1, x_2, \ldots, x_n)$ is $C^2$ on $\Omega$ and is $0$ outside, then, from \eqref{eq:I.8.12''} and $\frac{\partial}{\partial \nu} = \sum\limits_j \frac{\partial}{\partial x_j} \cos (x_j, \nu)$ we obtain Green's 
integral theorem 
\begin{equation}\tag{12'''}\label{eq:I.8.12''a}
\left(\Delta T_{f}\right) (\varphi) = T_{\Delta f} (\varphi) + \int_{S} \frac {\partial f}{\partial \nu} \varphi (x) \dd S- \int_{S} f (x) \frac {\partial \varphi}{\partial \nu} \dd S,
\end{equation}
where $\Delta$ is the Laplacian $\sum\limits_{j=1}^{n} \partial^2 / \partial x_j^2$. 
\end{corollaryu}

\begin{proposition}\label{prop:I.8.Proposition.3} If $T$ is a generalized function in $\Omega$ and $f \in C^{\infty}(\Omega)$, then 
\begin{equation}\label{eq:I.8.13}
S (\varphi) = T (f \varphi), \varphi \in \mathfrak {D} (\Omega),
\end{equation}
defines another generalized function $S$ in $\Omega$. 
\end{proposition}

\begin{proof}
 $S$ is a linear functional on $\mathfrak{D}(\Omega)$ which is bounded on every bounded set of $\mathfrak{D}(\Omega)$. This we see by applying Leibniz' formula to $f\varphi$. 
\end{proof}

\begin{definition}\label{def:I.8.Definition.4} The generalized function $S$ defined by \eqref{eq:I.8.13} is called the product of the function $f$ and the generalized function $T$. 
\end{definition}

Leibniz' Formula. We have, denoting $S$ in \eqref{eq:I.8.13} by $fT$, 
\begin{equation}\label{eq:I.8.14}
\frac {\partial}{\partial x_{j}} (f T) = \frac {\partial f}{\partial x_{j}} T + f \frac {\partial T}{\partial x_{j}}
\end{equation}
because we have 
\[
- T \left(f \frac {\partial \varphi}{\partial x_{j}}\right) = T \left(\frac {\partial f}{\partial x_{j}} \varphi\right)- T \left(\frac {\partial}{\partial x_{j}} (f \varphi)\right)
\]
by Leibniz' formula for $\partial (f\varphi) / \partial x_j$. This formula is generalized as follows. 

Let $P(\xi)$ be a polynomial in $\xi_1, \xi_2, \ldots, \xi_n$, and consider a linear partial differential operator $P(D)$ with constant coefficients, obtained by replacing $\xi_j$ by $i^{-1} \partial/\partial x_j$. The introduction of the imaginary coefficient $i^{-1}$ is suitable for the symbolism in the Fourier transform theory in Chapter~\ref{chap:VI}. 

\begin{theorem}[Generalized Leibniz' Formula of L. \textsc{Hörmander}]\label{thm:I.8.Theorem.3} We have 
\begin{equation}\label{eq:I.8.15}
P(D) (f T) = \sum\limits_{s} \left(D_{s} f\right) \frac {1}{s !} P^{(s)} (D) T,
\end{equation}
where, for $s = (s_1, s_2, \ldots, s_n)$, 
\begin{equation}\label{eq:I.8.16}
\begin{yoslines}
&\displaystyle P^{(s)} (\xi) = \frac {\partial^ {s_{1} + s_{2} + \dots + s_{n}}}{\partial \xi_{1} ^ {s_{1}} \partial \xi_{2} ^ {s_{2}} \cdots \partial \xi_{n} ^ {s_{n}}} P(\xi) = D_{\xi} ^{s} P(\xi), \displaystyle D_{s} = \left(\frac {1}{i} \frac {\partial}{\partial x_{1}}\right) ^ {s_{1}} \left(\frac {1}{i} \frac {\partial}{\partial x_{2}}\right) ^ {s_{2}} \dots \left(\frac {1}{i} \frac {\partial}{\partial x_{n}}\right) ^ {s_{n}},
\end{yoslines}
\end{equation}
and 
\begin{equation}\label{eq:I.8.17}
s! = s_{1}! s_{2}! \dots s_{n}! .
\end{equation}
\end{theorem}

\begin{proof}
 Repeated application of \eqref{eq:I.8.14} gives an identity of the form 
 \begin{equation}\label{eq:I.8.18}
 P(D) (f T) = \sum\limits_{s} \left(D_{s} f\right) Q_{s} (D) T,
 \end{equation}
where $Q_{s}(\xi)$'s are polynomials. Since \eqref{eq:I.8.18} is an identity, we may substitute 
\[
f (x) = e ^{i \langle x, \xi \rangle} \text{ and } T = e ^{i \langle x, \eta \rangle}, \text{ where } \langle x, \xi \rangle = \sum\limits_{j=1} ^{n} x_{j} \xi_{j}
\]
in \eqref{eq:I.8.18}. Thus, by the symbolism 
\begin{equation}\label{eq:I.8.19}
P(D) e ^{i \langle x, \xi \rangle} = P(\xi) e ^{i \langle x, \xi \rangle},
\end{equation}
we obtain 
\[
P(\xi + \eta) = \sum\limits_{s} \xi^{s} Q_{s} (\eta), \text{ where } \xi^{s} = \xi_{1} ^ {s_{1}} \xi_{2} ^ {s_{2}} \dots \xi_{n} ^ {s_{n}}.
\]

On the other hand, we have, by Taylor's formula, 
\[
P(\xi + \eta) = \sum\limits_{s} \frac {1}{s !} \xi^{s} P^{(s)} (\eta).
\]

 Thus we obtain 
 \[
 Q_{s} (\eta) = \frac {1}{s !} P^{(s)} (\eta).
 \]
\end{proof}

\yossection{9}{\texorpdfstring{$B$}{B}-spaces and \texorpdfstring{$F$}{F}-spaces}
\label{sec:I.9}

In a quasi-normed linear space $X$, $\lim\limits_{n\to \infty}\|x_n- x\| = 0$ implies, by the triangle inequality $\|x_n- x_m\|\leqq \|x_n- x\| + \|x- x_m\|$, that $\{x_n\}$ is a Cauchy sequence, i.e., $\{x_n\}$ satisfies Cauchy's convergence condition 
\begin{equation}\label{eq:I.9.1}
\lim\limits_{n, m \rightarrow \infty} \| x_{n}- x_{m} \| = 0.
\end{equation}

\begin{definition}\label{def:I.9.Definition.1} $A$ quasi-normed (or normed) linear space $X$ is called an $F$-space (or a $B$-space) if it is complete, i.e., if every Cauchy sequence $\{x_{n}\}$ of $X$ converges strongly to a point $x_{\infty}$ of $X$: 
\begin{equation}\label{eq:I.9.2}
\lim\limits_{n \rightarrow \infty} \| x_{n}- x_{\infty} \| = 0.
\end{equation}
\end{definition}

Such a limit $x_{\infty}$, if it exists, is uniquely determined because of the triangle inequality $\|x- x' \| \leqq \|x- x_n\| + \|x_n- x'\|$. A complete pre-Hilbert space is called a Hilbert space. 

\begin{remarku} The names $F$-space and $B$-space are abbreviations of Fréchet space and Banach space, respectively. It is to be noted that \textsc{Bourbaki} uses the term Fréchet spaces for locally convex spaces which are quasi-normed and complete. 
\end{remarku}

\begin{proposition}\label{prop:I.9.Proposition.1} Let $\Omega$ be an open set of $\mathbb{R}^{n}$, and denote by $\mathfrak{E}(\Omega) = C^\infty(\Omega)$ the locally convex space, quasi-normed as in Proposition~\ref{prop:I.1.Proposition.6} in Chapter~\ref{chap:I}, \ref{sec:I.1}. This $\mathfrak{E}(\Omega)$ is an $F$-space. 
\end{proposition}

\begin{proof}
 The condition $\lim\limits_{n,m\to \infty}\|f_n- f_m\| = 0$ in $\mathfrak{E}(\Omega)$ means that, for any compact subset $K$ of $\Omega$ and for any differential operator $D^{\alpha}$, the sequence $\{D^{\alpha}f_n(x)\}$ of functions converges, as $n\to \infty$, uniformly on $K$. Hence there exists a function $f(x)\in C^{\infty}(\Omega)$ such that $\lim\limits_{n\to \infty}D^{\alpha}f_n(x) =$ 

$D^{\alpha}f(x)$ uniformly on $K$. $D^{\alpha}$ and $K$ being arbitrary, this means that $\lim\limits_{n\to \infty}\| f_n- f\| = 0$ in $\mathfrak{E}(\Omega)$. 
\end{proof}

\begin{proposition}\label{prop:I.9.Proposition.2} $L^p(S) = L^p(S, \mathfrak{B}, m)$ is a $B$-space. In particular, $L^2(S)$ and $(l^2)$ are Hilbert spaces. 
\end{proposition}

\begin{proof}
 Let $\lim\limits_{n,m\to\infty}\left\|x_{n}-x_{m}\right\|=0$ in $L^{p}(S)$. Then we can choose a subsequence $\{x_{n_{k}}\}$ such that $\sum\limits_{k}\left\|x_{n_{k+1}}-x_{n_{k}}\right\|<\infty$. Applying the triangle inequality and the Lebesgue-Fatou Lemma to the sequence of functions 
\[
y_{t} (s) = \left| x_ {n_{1}} (s) \right| + \sum\limits_{k=1} ^{t} \left| x_ {n_{k+1}} (s)- x_ {n_{k}} (s) \right| \in L ^{p} (S),
\]
we see that 
\[
\int_{\Omega} \left(\lim\limits_{t \rightarrow \infty} y_{t} (s) ^{p}\right) m (\dd s) \leqq \varliminf\limits_{t \rightarrow \infty} \| y_{t} \| ^{p} \leqq \left(\| x_ {n_{1}} \| + \sum\limits_{k=1} ^{\infty} \| x_ {n_{k+1}}- x_ {n_{k}} \|\right) ^{p}. \]

Thus a finite $\lim\limits_{t\to \infty}y_t(s)$ exists a.e. Hence a finite $\lim\limits_{t\to \infty}x_{n_{t+1}}(s) = x_{\infty}(s)$ exists a.e. and $x_{\infty}(s)\in L^{p}(S)$, since $|x_{n_{t+1}}(s)|\leqq \lim\limits_{t\to \infty}y_t(s)\in L^p(S)$. Applying again the Lebesgue-Fatou Lemma, we obtain 
\[
\left\| x_{\infty}- x_ {n_{k}} \right\| ^{p} = \int_{\Omega} \left(\lim\limits_{t \rightarrow \infty} \left| x_ {n_{t}} (s)- x_ {n_{k}} (s) \right| ^{p}\right) m (\dd s) \leqq \left(\sum\limits_{l=k} ^{\infty} \left\| x_ {n_{l+1}}- x_ {n_{l}} \right\|\right) ^{p}.
\]

Therefore $\lim\limits_{k\to \infty}\|x_{\infty}- x_{n_{k}}\| = 0$, and hence, by the triangle inequality and Cauchy's convergence condition $\lim\limits_{n,m\to \infty}\|x_n- x_m\| = 0$, we obtain 
\[
\varlimsup\limits_{n \rightarrow \infty} \| x_{\infty}- x_{n} \| \leqq \varlimsup\limits_{k \rightarrow \infty} \| x_{\infty}- x_ {n_{k}} \| + \varlimsup\limits_{k, n \rightarrow \infty} \| x_ {n_{k}}- x_{n} \| = 0.
\]
\end{proof}

Incidentally we have proved the following important 

\begin{corollaryu} $A$ sequence $\{x_{n}\} \in L^{p}(S)$ which satisfies Cauchy's convergence condition \eqref{eq:I.9.1} contains a subsequence $\{x_{n_{k}}\}$ such that 
a finite $\lim\limits_{b\to \infty}x_{n_{k}}(s) = x_\infty (s)$ exists a.e., $x_{\infty}(s)\in L^{p}(S)$ and 
\begin{equation}\label{eq:I.9.3}
\slim\limits_{n \rightarrow \infty} x_{n} = x_{\infty}.
\end{equation}
\end{corollaryu}

\begin{remarku} In the above Proposition and the Corollary, we have assumed in the proof that $1 \leqq p< \infty$. However the results are also valid for the case $P= \infty$, and the proof is somewhat simpler than for the case $1 \leqq p< \infty$. The reader should carry out the proof. 
\end{remarku}

\begin{proposition}\label{prop:I.9.Proposition.3} The space $A^2(G)$ is a Hilbert space. 
\end{proposition}

\begin{proof}
 Let $\{f_n(z)\}$ be a Cauchy sequence of $A^2(G)$. Since $A^2(G)$ is a linear subspace of the Hilbert space $L^2(G)$, there exists a subsequence $\{f_{n_{k}}(z)\}$ such that 

a finite $\lim\limits_{k\to \infty}f_{n_{k}}(z) = f_\infty (z)$ exists a.e., $f_{\infty}\in L^{2}(G)$ and 
\[
\lim\limits_{n \rightarrow \infty} \int_{G} | f_{\infty} (z)- f_{n} (z) | ^{2} \dd x \dd y = 0. \]

We have to show that $f_{\infty}(z)$ is holomorphic in $G$. To do this, let the sphere $|z- z_0| \leqq \varrho$ be contained in $G$. The Taylor expansion $f_n(z)- f_m(z) = \sum\limits_{j=0}^{\infty} c_j (z- z_0)^j$ implies 
\begin{equation}\label{eq:I.9.4}
\begin{aligned}
&\left\| f_{n}- f_{m} \right\| ^{2} \geqq \int_ {\left| z- z_{0} \right| \leqq \varrho} \left| f_{n} (z)- f_{m} (z) \right| ^{2} \dd x \dd y = \int_{0} ^{\varrho} \left(\int_{0} ^{2 \pi} \sum\limits_{j} c_{j} r ^{j} e ^{i j \theta} \sum\limits_{k} \bar {c}_{k} r ^{k} e ^{- i k \theta} \dd \theta\right) r \dd r
\\
&= \sum\limits_{j} 2 \pi \int_{0} ^{\varrho} \left| c_{j} \right| ^{2} r ^{2j+1} \dd r = 2 \pi \sum\limits_{j=0} ^{\infty} \varrho^{2j+2} \left| c_{j} \right| ^{2} (2 j + 2) ^{-1} \geqq \pi \left| c_{0} \right| ^{2} \varrho^{2}
\\
&= \pi \varrho^{2} \left| f_{n} (z_{0})- f_{m} (z_{0}) \right| ^{2}.
\end{aligned}
\end{equation}

Thus the sequence $\{f_n(z)\}$ itself converges uniformly on any closed sphere contained in $G$. $f_n(z)$'s being holomorphic in $G$, we see that $f_\infty(z) = \lim\limits_{n \to \infty} f_n(z)$ must be holomorphic in $G$. 
\end{proof}

\begin{proposition}\label{prop:I.9.Proposition.4} $M(S, \mathfrak{B}, m)$ with $m(S) < \infty$ is an $F$-space. 
\end{proposition}

\begin{proof}
 Let $\{x_{n}\}$ be a Cauchy sequence in $M(S, \mathfrak{B}, m)$. Since the convergence in $M(S, \mathfrak{B}, m)$ is the asymptotic convergence, we can choose a sub-sequence $\{x_{n_{k}}(s)\}$ of $\{x_{n}(s)\}$ such that 
\[
m \left(B_{k}\right) \leqq 2 ^{-k} \quad \text{ for } \quad B_{k} = \{s \in S; 2 ^{-k} \leqq \left| x_ {n_{k+1}} (s)- x_ {n_{k}} (s) \right| \}. \]

The sequence $x_{n_{k}}(s) = x_{n_1}(s) + \sum\limits_{j=1}^{k-1} (x_{n_{j+1}}(s)- x_{n_{j}}(s))$ ( $k = 1, 2, \ldots$) is $s$-convergent to a function $\in M(S, \mathfrak{B}, m)$, because, if $s \in \bigcup\limits_{j=t}^{\infty} B_j$, we have $\sum\limits_{j=t}^{\infty} |x_{n_{j+1}}(s)- x_{n_{j}}(s)| \leqq \sum\limits_{j=t}^{\infty} 2^{-j} \leqq 2^{1-t}$ and $m\left(\bigcup\limits_{j=t}^{\infty} B_j\right) \leqq \sum\limits_{j=t}^{\infty} m(B_j) \leqq \sum\limits_{j=t}^{\infty} 2^{-j} \leqq 2^{1-t}$; consequently we see, by letting $t \to \infty$, that the sequence $\{x_{n_{k}}(s)\}$ converges $m$-a.e. to a function $x_\infty(s) \in M(S, \mathfrak{B}, m)$. Hence $\lim\limits_{k \to \infty} \|x_{n_{k}}- x_\infty\| = 0$ and so, by $\lim\limits_{n, m \to \infty} \|x_n- x_m\| = 0$, we obtain $\lim\limits_{n \to \infty} \|x_n- x_\infty\| = 0$. 
\end{proof}

The Space $(s)$. The set $(s)$ of all sequences $\{\xi_{n}\}$ of numbers quasi-normed by 
\[
\left\| \left\{\xi_{n} \right\} \right\| = \sum\limits_{j=1} ^{\infty} 2 ^{-j} \left| \xi_{j} \right| / (1 + \left| \xi_{j} \right|)
\]
constitutes an $F$-space by $\{\xi_n\} + \{\eta_n\} = \{\xi_n + \eta_n\}$, $\alpha\{\xi_n\} = \{\alpha\xi_n\}$. The proof of the completeness of $(s)$ may be obtained as in the case of $M(S, \mathfrak{B}, m)$. The quasi-norm 
$\|\{\xi_{n}\}\|=\inf\limits_{\varepsilon>0}\tan^{-1}\left\{\varepsilon+\text{ the number of }\xi_{n}\text{ 's which satisfy }|\xi_{n}|>\varepsilon\right\}$ also gives an equivalent topology of $(s)$. 

\begin{remarku} It is clear that $C(S), (c_0)$ and $(c)$ are $B$-spaces. The completeness of the space $(l^p)$ is a consequence of that of $L^p(S)$. Hence, by Theorem~\ref{thm:I.5.Theorem.3} in Chapter~\ref{chap:I}, \ref{sec:I.5}, the space $H- L^2$ is a Hilbert space with $(l^2)$. 
\end{remarku}

Soboley Spaces $W^{k,p}(\Omega)$. Let $\Omega$ be an open set of $\mathbb{R}^{n}$, and $k$ a positive integer. For $1 \leqq p< \infty$, we denote by $W^{k,p}(\Omega)$ the set of all complex-valued functions $f(x) = f(x_1, x_2, \ldots, x_n)$ defined in $\Omega$ such that $f$ and its distributional derivatives $D^s f$ of order $|s| = \sum\limits_{j=1}^{n} |s_j| \leqq k$ all belong to $L^p(\Omega)$. $W^{k,p}(\Omega)$ is a normed linear space by 
\[
\begin{aligned} (f_{1} + f_{2}) (x) & = f_{1} (x) + f_{2} (x), (\alpha f) (x) = \alpha f (x) \text{ and } \\ \| f \|_{k,p} & = \left(\sum\limits_{| s | \leqq k} \int_{\Omega} | D ^{s} f (x) | ^{p} \dd x\right) ^{1/p}, \dd x = \dd x_{1} \dd x_{2} \dots \dd x_{n} \end{aligned}
\]
under the convention that we consider two functions $f_{1}$ and $f_{2}$ as the same vector of $W^{k,p}(\Omega)$ if $f_{1}(x) = f_{2}(x)$ a.e. in $\Omega$. It is easy to see that $W^{k,2}(\Omega)$ is a pre-Hilbert space by the scalar product 
\[
(f, g)_{k,2} = \left(\sum\limits_{| s | \leqq k} \int_{\Omega} D ^{s} f (x) \overline {{D ^{s} g (x)}} \dd x\right).
\]

\begin{proposition}\label{prop:I.9.Proposition.5} The space $W^{k,p}(\Omega)$ is a $B$-space. In particular, $W^k (\Omega) = W^{k,2}(\Omega)$ is a Hilbert space by the norm $\|f\|_k = \|f\|_{k,2}$ and the scalar product $(f,g)_k = (f,g)_{k,2}$. 
\end{proposition}

\begin{proof}
 Let $\{f_h\}$ be a Cauchy sequence in $W^{h,p}(\Omega)$. Then, for any differential operator $D^s$ with $|s|\leqq k$, the sequence $\{D^s f_h\}$ is a Cauchy sequence in $L^p(\Omega)$ and so, by the completeness of $L^p(\Omega)$, there exist functions $f^{(s)}\in L^{p}(\Omega)$ ( $|s|\leqq k$) such that $\lim\limits_{h\to \infty}\int_{\Omega}|D^{s}f_{h}(x)- f^{(s)}(x)|^{p}\dd x = 0$. By virtue of Hölder's inequality in Chapter~\ref{chap:I}, \ref{sec:I.3}, applied to compact sets of $\Omega$, we easily see that $f_{h}$ is locally integrable in $\Omega$. Hence, for any function $\varphi \in C_0^\infty (\Omega)$, 
\[
T_ {D ^{s} f_{h}} (\varphi) = \int_{\Omega} D ^{s} f_{h} (x) \cdot \varphi (x) \dd x = (- 1) ^{|s|} \int_{\Omega} f_{h} (x) D ^{s} \varphi (x) \dd x,
\]
and so, again applying Hölder's inequality, we obtain, by $\lim\limits_{h\to \infty}\int_{\Omega}|f_h(x)- f^{(0)}(x)|^p\dd x = 0,$ 
\[
\lim\limits_{h \rightarrow \infty} T_ {D ^{s} f_{h}} (\varphi) = (- 1) ^{|s|} T_ {f ^{(0)}} (D ^{s} \varphi) = D ^{s} T_ {f ^{(0)}} (\varphi). \]

Similarly we have, by $\lim\limits_{h\to \infty}\int_{\Omega}|D^{s}f_{h}(x)- f^{(s)}(x)|^{p}\dd x = 0$ 
\[
\lim\limits_{h \rightarrow \infty} T_ {D ^{s} f_{h}} (\varphi) = T_ {f ^{(s)}} (\varphi).
\]

Hence we must have $D^{s}T_{f^{(0)}} = T_{f^{(s)}}$, that is, the distributional derivative $D^{s}f^{(0)}$ equals $f^{(s)}$. This proves that $\lim\limits_{h\to \infty}\| f_h- f^{(0)}\|_{k,p} = 0$ and $W^{k,p}(\Omega)$ is complete. 
\end{proof}


\yossection{10}{The Completion}
\label{sec:I.10}

The completeness of an $F$-space (and a $B$-space) will play an important role in functional analysis in the sense that we can apply to such spaces Baire's category arguments given in Chapter~\ref{chap:0}, Preliminaries. The following theorem of completion will be of frequent use in this book. 

\begin{theoremn}[of completion] Let $X$ be a quasi-normed linear space which is not complete. Then $X$ is isomorphic and isometric to a dense linear subspace of an $F$-space $\tilde{X}$, i.e., there exists a one-to-one correspondence $x \leftrightarrow \tilde{x}$ of $X$ onto a dense linear subspace of $\tilde{X}$ such that
\begin{equation}\label{eq:I.10.1}
\left(\widetilde {x + y}\right) = \tilde {x} + \tilde {y}, (\widetilde {\alpha x}) = \alpha \tilde {x}, \| \tilde {x} \| = \| x \|.
\end{equation}

The space $\tilde{X}$ is uniquely determined up to isometric isomorphism. If $X$ is itself a normed linear space, then $\tilde{X}$ is a $B$-space. 

\end{theoremn}
\begin{proof}
 The proof proceeds as in Cantor's construction of real numbers from rational numbers. 


The set of all Cauchy sequences $\{x_{n}\}$ of $X$ can be classified according to the equivalence $\{x_{n}\} \sim \{y_{n}\}$ which means that $\lim\limits_{n\to \infty}\|x_n- y_n\| = 0$. We denote by $\{x_{n}\}'$ the class containing $\{x_{n}\}$. Then the set $\tilde{X}$ of all such classes $\tilde{x} = \{x_{n}\}'$ is a linear space by 
\[
\{x_{n} \} ^{\prime} + \{y_{n} \} ^{\prime} = \{x_{n} + y_{n} \} ^{\prime}, \alpha \{x_{n} \} ^{\prime} = \{\alpha x_{n} \} ^{\prime}.
\]

We have $\|x_n\|- \|x_m\| \leqq \|x_n- x_m\|$ and hence $\lim\limits_{n \to \infty} \|x_n\|$ exists. We put 
\[
\left\|\left\{x_{n} \right\} ^{\prime} \right\| = \lim\limits_{n \rightarrow \infty} \left\| x_{n} \right\|.
\]

It is easy to see that these definitions of the vector sum $\{x_{n}\}^{\prime} + \{y_{n}\}^{\prime}$, the scalar multiplication $\alpha\{x_{n}\}^{\prime}$ and the norm $\|\{x_{n}\}^{\prime}\|$ do not depend on the particular representations for the classes $\{x_{n}\}^{\prime}, \{y_{n}\}^{\prime}$, respectively. For example, if $\{x_{n}\} \sim \{x_{n}^{\prime}\}$, then 
\[
\lim\limits_{n \rightarrow \infty} \| x_{n} \| \leqq \lim\limits_{n \rightarrow \infty} \| x_{n} ^{\prime} \| + \lim\limits_{n \rightarrow \infty} \| x_{n} ^{\prime}- x_{n} \| \leqq \lim\limits_{n \rightarrow \infty} \| x_{n} ^{\prime} \|
\]
and similarly $\lim\limits_{n\to \infty}\|x_n'\|\leqq \lim\limits_{n\to \infty}\|x_n\|$ so that we have $\|\{x_n\}'\| = \|\{x_n'\}''\|$ 

To prove that $\|\{x_n\}'\|$ is a quasi-norm, we have to show that 
\[
\lim\limits_{\alpha \rightarrow 0} \| \alpha \{x_{n} \} ^{\prime} \| = 0 \quad \text{ and } \quad \lim\limits_ {\| \{x_{n} \} ^{\prime} \| \rightarrow 0} \| \alpha \{x_{n} \} ^{\prime} \| = 0.
\]

The former is equivalent to $\lim\limits_{\alpha\to0}\lim\limits_{n\to\infty}\|\alpha x_{n}\|=0$ and the latter is equivalent to $\lim\limits_{n\to\infty}\|\alpha x_{n}\|=0$. And these are true because $\|\alpha x\|$ is continuous in both variables $\alpha$ and $x$. 

To prove the completeness of $\tilde{X}$, let $\{\tilde{x}_k\} = \{\{x_n^{(k)}\}\}$ be a Cauchy sequence of $\tilde{X}$. For each $k$, we can choose $n_k$ such that 
\begin{equation}\label{eq:I.10.2}
\left\| x_{m} ^{(k)}- x_ {n_{k}} ^{(k)} \right\| < k ^{-1} \quad \text{ if } \quad m > n_{k}.
\end{equation}

Then we can show that the sequence $\{\tilde{x}_k\}$ converges to the class containing the Cauchy sequence of $X$: 
\begin{equation}\label{eq:I.10.3}
\left\{x_ {n_{1}} ^{(1)}, x_ {n_{2}} ^{(2)}, \dots , x_ {n_{k}} ^{(k)}, \dots \right\}.
\end{equation}

To this purpose, we denote by $\bar{x}_{n_{k}}^{(k)}$ the class containing 
\begin{equation}\label{eq:I.10.4}
\left\{x_ {n_{k}} ^{(k)}, x_ {n_{k}} ^{(k)}, \ldots , x_ {n_{k}} ^{(k)}, \ldots \right\}.
\end{equation}

Then, by \eqref{eq:I.10.2}, 
\begin{equation}\label{eq:I.10.5}
\left\| \tilde {x}_{k}- \bar {x}_ {n_{k}} ^{(k)} \right\| = \lim\limits_{m \rightarrow \infty} \left\| x_{m} ^{(k)}- x_ {n_{k}} ^{(k)} \right\| \leqq k ^{-1}
\end{equation}
and hence 
\[
\begin{aligned} \| x_ {n_{k}} ^{(k)}- x_ {n_{m}} ^{(m)} \| = & \| \bar {x}_ {n_{k}} ^{(k)}- \bar {x}_ {n_{m}} ^{(m)} \| \leqq \| \bar {x}_ {n_{k}} ^{(k)}- \tilde {x}_{k} \| + \| \tilde {x}_{k}- \tilde {x}_{m} \| + \| \tilde {x}_{m}- \bar {x}_ {n_{m}} ^{(m)} \| \\ & \leqq \| \tilde {x}_{k}- \tilde {x}_{m} \| + k ^{-1} + m ^{-1}. \end{aligned}
\]

Thus \eqref{eq:I.10.3} is a Cauchy sequence of $X$. Let $\tilde{x}$ be the class containing \eqref{eq:I.10.3}. Then, by \eqref{eq:I.10.5}, 
\[
\left\| \tilde {x}- \tilde {x}_{k} \right\| \leqq \left\| \tilde {x}- \bar {x}_ {n_{k}} ^{(k)} \right\| + \left\| \bar {x}_ {n_{k}} ^{(k)}- \tilde {x}_{k} \right\| \leqq \left\| \tilde {x}- \bar {x}_ {n_{k}} ^{(k)} \right\| + k ^{-1}.
\]

Since, as shown above, 
\[
\left\| \tilde {x}- \bar {x}_ {n_{k}} ^{(k)} \right\| \leqq \lim\limits_{p\rightarrow \infty} \left\| x_ {n_{p}} ^{(p)}- x_ {n_{k}} ^{(k)} \right\| \leqq \lim\limits_{p\rightarrow \infty} \left\| \tilde {x}_{p}- \tilde {x}_{k} \right\| + k ^{-1}
\]
we prove that $\lim\limits_{k\to \infty}\|\tilde{x}-\bar{x}_{n_{k}}^{(k)}\| = 0$, and so $\lim\limits_{k\to \infty}\|\tilde{x}-\tilde{x}_k\| = 0.$ 

The above proof shows that the correspondence 
\[
X \ni x \leftrightarrow \tilde {x} = \{x, x, \dots , x, \dots \} ^{\prime} = \bar {x}
\]
is surely isomorphic and isometric, and the image of $X$ in $\tilde{X}$ by this correspondence is dense in $\tilde{X}$. The last part of the Theorem is clear.
\end{proof}

\begin{examplen} Let $\Omega$ be an open set of $\mathbb{R}^{n}$ and $k < \infty$. The completion of the space $C_0^k(\Omega)$ normed by 
\[
\left\| f \right\|_{k} = \left(\sum\limits_{| s | \leqq k} \int_{\Omega} \left| D ^{s} f (x) \right| ^{2} \dd x\right) ^{1/2}
\]
will be denoted by $H_0^k(\Omega)$; thus $H_0^k(\Omega)$ is the completion of the pre-Hilbert space $\hat{H}_0^k(\Omega)$ defined in Chapter~\ref{chap:I}, \ref{sec:I.5}, Example~\hyperref[ex:I.5.Example.4]{$4$}. Therefore $H_0^k(\Omega)$ is a Hilbert space. The completion of the pre-Hilbert space $\hat{H}^k(\Omega)$ in Chapter~\ref{chap:I}, \ref{sec:I.5}, Example~\hyperref[ex:I.5.Example.3]{$3$} will similarly be denoted by $H^k(\Omega)$. 
\end{examplen}

The elements of $H_0^k(\Omega)$ are obtained concretely as follows: Let $\{f_h\}$ be a Cauchy sequence of $C_0^k(\Omega)$ with regard to the norm $\|f\|_k$. Then, by the completeness of the space $L^2(\Omega)$, we see that there exist functions $f^{(s)}(x) \in L^2(\Omega)$ with $|s| = \sum\limits_{j=1}^{n} s_j \leqq k$ such that 
\[
\lim\limits_{h \rightarrow \infty} \int_{\Omega} | f ^{(s)} (x)- D ^{s} f_{h} (x) | ^{2} \dd x = 0 \quad (\dd x = \dd x_{1} \dd x_{2} \dots \dd x_{n}).
\]

Since the scalar product is continuous in the norm of $L^2(\Omega)$, we see, for any test function $\varphi(x) \in C_0^\infty(\Omega)$, that 
\[
\begin{aligned} T_ {f ^{(s)}} (\varphi) & = \lim\limits_{h \to \infty} \langle D ^{s} f_{h}, \varphi \rangle = \lim\limits_{h \to \infty} (- 1) ^{|s|} T_ {f_{h}} (D ^{s} \varphi) \\ & = (- 1) ^{|s|} \lim\limits_{h \to \infty} \langle f_{h}, D ^{s} \varphi \rangle = (- 1) ^{|s|} \langle f ^{(0)}, D ^{s} \varphi \rangle = (D ^{s} T_ {f ^{(0)}}) (\varphi). \end{aligned}
\]

Therefore we see that $f^{(s)} \in L^2(\Omega)$ is, when considered as a generalized function, the distributional derivative of $f^{(0)}: f^{(s)} = D^s f^{(0)}$. 

We have thus proved that the Hilbert space $H_{0}^{k}(\Omega)$ is a linear subspace of the Hilbert space $W^{k}(\Omega)$, the Sobolev space. In general $H_{0}^{k}(\Omega)$ is a proper subspace of $W^{k}(\Omega)$. However, we can prove 

\begin{propositionu} $H_0^k (\mathbb{R}^{n}) = W^k (\mathbb{R}^{n})$. 
\end{propositionu}

\begin{proof}
 We know that the space $W^{k}(\mathbb{R}^{n})$ is the space of all functions $f(x) \in L^{2}(\mathbb{R}^{n})$ such that the distributional derivatives $D^{s}f(x)$ with $|s| = \sum\limits_{j=1}^{n}s_{j} \leqq k$ all belong to $L^{2}(\mathbb{R}^{n})$, and the norm in $W^{k}(\mathbb{R}^{n})$ is given by 
\[
\left\| f \right\|_{k} = \left(\sum\limits_{| s | \leqq k} \left| D ^{s} f (x) \right| ^{2} \dd x\right) ^{1/2}. \]

Let $f \in W^{k}(\mathbb{R}^{n})$ and define $f_{N}$ by 
\[
f_{N} (x) = \alpha_{N} (x) f (x),
\]
where the function $\alpha_{N}(x) \in C_{0}^{\infty}(\mathbb{R}^{n})$ ( $N = 1, 2, \ldots$) is such that 
\[
\alpha_{N} (x) = 1 \text{ for } | x | \leqq N \text{ and } \sup\limits_ {x \in \mathbb{R}^{n}; | s | \leqq k; N = 1, 2, \dots} | D ^{s} \alpha_{N} (x) | < \infty .
\]

Then by Leibniz' formula, we have 
\begin{equation}D ^{s} f (x)- D ^{s} f_{N} (x) = 0 \text{ for } | x | \leqq N, D ^{t} \alpha_{N} (x) \cdot D ^{u} f (x) \text{ with } | u | + | t | \leqq k \text{ for } | x | > N. \notag\end{equation}

Hence, by $D^s f \in L^2(\mathbb{R}^{n})$ for $|s| \leqq k$, we see that $\lim\limits_{N \to \infty} \|D^s f_N- D^s f\|_0 = 0$ and so $\lim\limits_{N \to \infty} \|f_N- f\|_k = 0$. 

Therefore, it will be sufficient to show that, for any $f \in W^{k}(\mathbb{R}^{n})$ with compact support, there exists a sequence $\{f_{a}(x)\} \subseteq C_{0}^{\infty}(\mathbb{R}^{n})$ such that $\lim\limits_{a \to \infty} \|f_{a}- f\|_{h} = 0$. To this purpose, consider the regularization of $f$ (see \eqref{eq:I.1.16} in Chapter~\ref{chap:I}, \ref{sec:I.1}): 
\[
f_{a} (x) = \int_ {\mathbb{R}^{n}} f (y) \theta_{a} (x- y) \dd y, a > 0.
\]

By differentiation, we have 
\[
\begin{aligned}
D ^{s} f_{a} (x) &= \int_ {\mathbb{R}^{n}} f (y) D_{x} ^{s} \theta_{a} (x- y) \dd y = (- 1) ^{|s|} \int_ {\mathbb{R}^{n}} f (y) D_{y} ^{s} \theta_{a} (x- y) \dd y
\\
\qquad &= (D ^{s} T_{f}) (\theta_{a,x}) \text{ (where } \theta_{a,x} (y) = \theta_{a} (x- y))
\\
\qquad &= \int_ {\mathbb{R}^{n}} D ^{s} f (y) \cdot \theta_{a} (x- y) \dd y \quad (\text{ for } | s | \leqq k).
\end{aligned}
\]

Hence, by Schwarz' inequality, 
\[
\begin{aligned}
& \int_ {\mathbb{R}^{n}} | D ^{s} f_{a} (x)- D ^{s} f (x) | ^{2} \dd x
\\
& \leqq \left(\int_ {\mathbb{R}^{n}} \theta_{a} (x- y) \dd y\right) \int_ {\mathbb{R}^{n}} \left[ \int_ {\mathbb{R}^{n}} | D_{y} ^{s} f (y)- D_{x} ^{s} f (x) | ^{2} \theta_{a} (x- y) \dd y \right] \dd x
\\
&= \int_{| \varepsilon | \leqq a} \left[ \int_ {\mathbb{R}^{n}} | D_{y} ^{s} f (y)- D_{y} ^{s} f (y + \varepsilon) | ^{2} \dd y \right] \theta_{a} (\varepsilon) \dd \varepsilon , \text{ where }
\\
y + \varepsilon &= (y_{1} + \varepsilon_{1}, y_{2} + \varepsilon_{2}, \ldots , y_{n} + \varepsilon_{n}).
\end{aligned}
\]

We know that the inner integral on the extreme right tends to $0$ as $\varepsilon \to 0$ (see Theorem~\ref{thm:0.3.Theorem.1} in Chapter~\ref{chap:0}, \ref{sec:0.3}), and hence $\lim\limits_{a\to 0}\int \limits_{\mathbb{R}^{n}}|D^s f_a(x)- D^s f(x)|^2 \dd x = 0$. Thus $\lim\limits_{a\to 0}\|f_a- f\|_k = 0$. Therefore, the completion $H_0^k (\mathbb{R}^{n})$ of $C_0^k (\mathbb{R}^{n})$ with regard to the norm $\|\|_k$ is identical with the space $W^{k}(\mathbb{R}^{n})$. 
\end{proof}

\begin{corollaryu} $H_0^k (\mathbb{R}^{n}) = H^k (\mathbb{R}^{n}) = W^k (\mathbb{R}^{n})$. 
\end{corollaryu}


\yossection{11}{Factor Spaces of a \texorpdfstring{$B$}{B}-space}
\label{sec:I.11}

Suppose that $X$ is a normed linear space and that $M$ is a closed linear subspace in $X$. We consider the factor space $X/M$, i.e., the space whose elements are classes modulo $M$. In virtue of the fact that $M$ is closed, all these classes $\xi$ are closed in $X$. 

\begin{propositionu} If we define 
\begin{equation}\label{eq:I.11.1}
\left\| \xi \right\| = \inf\limits_{x \in \xi} \left\| x \right\|,
\end{equation}
then all the axioms concerning the norm are satisfied by $\|\xi \|$ 
\end{propositionu}

\begin{proof}
 If $\xi = 0$, then $\xi$ coincides with $M$ and contains the zero vector of $X$; consequently, it follows from \eqref{eq:I.11.1} that $\|\xi\| = 0$. Suppose, conversely, that $\|\xi\| = 0$. It follows then from \eqref{eq:I.11.1} that the class contains a sequence $\{x_n\}$ for which we have $\lim\limits_{n \to \infty} \|x_n\| = 0$, and hence the zero vector of $X$ belongs to the closed set $\xi$ of $X$. This proves that $\xi = M$ and hence is the zero vector in $X/M$. 


Next suppose $\xi, \eta \in X / M$. By definition \eqref{eq:I.11.1}, there exists for any $\varepsilon > 0$, vectors $x \in \xi, y \in \eta$ such that 
\[
\left\| x \right\| \leqq \left\| \xi \right\| + \varepsilon , \left\| y \right\| \leqq \left\| \eta \right\| + \varepsilon .
\]

Hence $\|x + y\| \leqq \|x\| + \|y\| \leqq \|\xi\| + \|\eta\| + 2\varepsilon$. On the other hand, $(x + y) \in (\xi + \eta)$, and therefore $\|\xi + \eta\| \leqq \|x + y\|$ by \eqref{eq:I.11.1}. Consequently, we have $\|\xi + \eta\| \leqq \|\xi\| + \|\eta\| + 2\varepsilon$ and so we obtain the triangle inequality $\|\xi + \eta\| \leqq \|\xi\| + \|\eta\|$. 

Finally it is clear that the axiom $\|\alpha \xi \| = |\alpha \||\xi \|$ holds good.
\end{proof}

\begin{definitionu} The space $X$/$M$, normed by \eqref{eq:I.11.1}, is called a normed factor space. 
\end{definitionu}

\begin{theoremu} If $X$ is a $B$-space and $M$ a closed linear subspace in $X$, then the normed factor space $X$/$M$ is also a $B$-space. 
\end{theoremu}

\begin{proof}
 Suppose $\{\xi_n\}$ is a Cauchy sequence in $X / M$. Then $\{\xi_n\}$ contains a subsequence $\{\xi_{n_{k}}\}$ such that $\| \xi_{n_{k+1}}- \xi_{n_{k}} \| < 2^{-k-2}$. Further, by definition \eqref{eq:I.11.1} of the norm in $X / M$, one can choose in every class $(\xi_{n_{k+1}}- \xi_{n_{k}})$ a vector $y_k$ such that 
\[
\left\| y_{k} \right\| < \left\| \xi_ {n_{k+1}}- \xi_ {n_{k}} \right\| + 2 ^{-k-2} < 2 ^{-k-1}. \]

Let $x_{n_1} \in \xi_{n_1}$. The series $x_{n_1} + y_1 + y_2 + \cdots$ converges in norm and consequently, in virtue of the completeness of $X$, it converges to an element $x$ of $X$. Let $\xi$ be the class containing $x$. We shall prove that $\xi = \slim\limits_{n \to \infty} \xi_n$. 


Denote by $s_{k}$ the partial sum $x_{n_{1}} + y_{1} + y_{2} + \cdots + y_{k}$ of the above series. Then $\lim\limits_{k \to \infty} \|x- s_{k}\| = 0$. On the other hand, it follows from the relations $x_{n_{1}} \in \xi_{n_{1}}, y_{p} \in (\xi_{n_{p+1}}- \xi_{n_{p}})$ that $s_{k} \in \xi_{n_{k+1}}$, and so, by \eqref{eq:I.11.1}, 
\[
\left\| \xi- \xi_ {n_{k}} \right\| \leqq \left\| x- s_{k} \right\|\rightarrow 0 \quad \text{ as } \quad k \rightarrow \infty .
\]

Therefore, from the inequality $\|\xi- \xi_n\| \leqq \|\xi- \xi_{n_{k}}\| + \|\xi_{n_{k}}- \xi_n\|$ and the fact that $\{\xi_n\}$ is a Cauchy sequence, we conclude that $\lim\limits_{n \to \infty} \|\xi- \xi_n\| = 0$.
\end{proof}

\yossection{12}{The Partition of Unity}
\label{sec:I.12}

To discuss the support of a generalized function, in the next section, we shall prepare the notion and existence of the partition of unity. 

\begin{propositionu} Let $G$ be an open set of $\mathbb{R}^{n}$. Let a family $\{U\}$ of open subsets $U$ of $G$ constitute an open base of $G$: any open subset of $G$ is representable as the union of open sets belonging to the family $\{U\}$. Then there exists a countable system of open sets of the family $\{U\}$ with the properties: 
\begin{equation}\label{eq:I.12.1}
\parbox{0.90\linewidth}{the union of open sets of this system equals $G$,}
\end{equation}
\begin{equation}\label{eq:I.12.2}
\parbox{0.90\linewidth}{any compact subset of $G$ meets (has a non-void intersection with) only a finite number of open sets of this system.}
\end{equation}
\end{propositionu}

\begin{definition}\label{def:I.12.Definition.1} The above system of open sets is said to constitute a scattered open covering of $G$ subordinate to $\{U\}$. 
\end{definition}

\begin{proof}[Proof of the Proposition.]
$G$ is representable as the union of a countable number of compact subsets. For example, we may take the system of all closed spheres contained in $G$ such that the centres are of rational coordinates and the radii of rational numbers.

Hence we see that there exists a sequence of compact subsets $K_r$ such that \begin{enumerate}[(i)]
\item $K_r \subseteq K_{r+1}$ ( $r = 1, 2, \ldots$),
\item $G$ is the union of $K_r$'s
\item each $K_r$ is contained in the interior of $K_{r+1}$
\end{enumerate}
. Set
$U_{r} = (\text{ the interior of } K_{r+1})- K_{r-2}$ and $V_{r} = K_{r}- (\text{ the interior of } K_{r-1})$, where for convention we set $K_{0} = K_{-1} =$ the void set. Then $U_{r}$ is open and $V_{r}$ is compact such that $G = \bigcup\limits_{r=1}^{\infty} V_{r}$. For any point $x \in V_{r}$, take an open set $U(x; r) \in \{U\}$ such that $x \in U(x; r) \subseteq U_{r}$. Since $V_{r}$ is compact, there exists a finite system of points $x^{(1)}, x^{(2)}, \ldots, x^{(h_{r})}$ such that $V_{r} \subseteq \bigcup\limits_{i=1}^{h_{r}} U(x^{(i)}; r)$. Then, since any compact set of $G$ meets only a finite number of $U_{r}$'s, it is easy to see that the system of open sets $U(x^{(i)}; r)$ ( $r = 1, 2, \ldots; 1 \leqq i \leqq h_{r}$) is a scattered open covering of $G$ subordinate to $\{U\}$. 
\end{proof}

\begin{theoremn}[the partition of unity] Let $G$ be an open set of $\mathbb{R}^{n}$, and let a family of open sets $\{G_i; i \in I\}$ cover $G$, i.e., $G = \bigcup\limits_{i \in I} G_i$. Then there exists a system of functions $\{\alpha_j(x); j \in J\}$ of $C_0^\infty(\mathbb{R}^{n})$ such that
\begin{equation}\label{eq:I.12.3}
\parbox{0.90\linewidth}{for each $j \in J$, $\supp (\alpha_j)$ is contained in some $G_i$,}
\end{equation}
\begin{align}
\text{ for every } j& \in J, 0 \leqq \alpha_{j} (x) \leqq 1, \label{eq:I.12.4} \\
\sum\limits_{j \in J} \alpha_{j} (x)& \equiv 1 \quad \text{ for } \quad x \in G. \label{eq:I.12.5}
\end{align}

\end{theoremn}
\begin{proof}
 Let $x^{(0)} \in G$ and take a $G_i$ which contains $x^{(0)}$. Let the closed sphere $S(x^{(0)}; r)$ of centre $x^{(0)}$ and radius $r$ be contained in $G_i$. We construct, as in ($14$), Chapter~\ref{chap:I}, \ref{sec:I.1}, a function $\beta_{x^{(0)}}^{(r)}(x) \in C_0^\infty(\mathbb{R}^{n})$ such that $\beta_{x^{(0)}}^{(r)}(x) > 0$ for $|x- x^{(0)}| < r$, $\beta_{x^{(0)}}^{(r)}(x) = 0$ for $|x- x^{(0)}| \geqq r$. 


We put $U_{x^{(0)}}^{(r)} = \{x; \beta_{x^{(0)}}^{(r)}(x) \neq 0\}$. Then $U_{x^{(0)}}^{(r)} \subseteq G_i$ and $\bigcup\limits_{x^{(0)} \in G, r > 0} U_{x^{(0)}}^{(r)} = G$, and, moreover, $\supp\left(\beta_{x^{(0)}}^{(r)}\right)$ is compact. 

There exists, by the Proposition, a scattered open covering $\{U_j; j \in J\}$ subordinate to the open base $\{U_{x^{(0)}}^{(r)}; x^{(0)} \in G, r > 0\}$ of $G$. Let $\beta_j(x)$ be any function of the family $\{\beta_{x^{(0)}}^{(r)}(x)\}$ which is associated with $U_j$. Then, since $\{U_j; j \in J\}$ is a scattered open covering, only a finite number of $\beta_j(x)$'s do not vanish at a fixed point $x$ of $G$. Thus the sum $s(x) = \sum\limits_{j \in J} \beta_j(x)$ is convergent and is $>0$ at every point $x$ of $G$. Hence the functions 
\[
\alpha_{j} (x) = \beta_{j} (x) / s (x) (j \in J)
\]
satisfy the condition of our theorem.
\end{proof}

\begin{definition}\label{def:I.12.Definition.2} The system $\{\alpha_j(x); j \in J\}$ is called a partition of unity subordinate to the covering $\{G_i; i \in I\}$. 
\end{definition}


\yossection{13}{Generalized Functions with Compact Support}
\label{sec:I.13}

\begin{definition}\label{def:I.13.Definition.1} We say that a distribution $T \in \mathfrak{D}(\Omega)'$ vanishes in an open set $U$ of $\Omega$ if $T(\varphi) = 0$ for every $\varphi \in \mathfrak{D}(\Omega)$ with support contained in $U$. The support of $T$, denoted by $\supp(T)$, is defined as the smallest closed set $F$ of $\Omega$ such that $T$ vanishes in $\Omega- F$. 
\end{definition}

To justify the above definition, we have to prove the existence of the largest open set of $\Omega$ in which $T$ vanishes. This is done by the following 

\begin{theorem}\label{thm:I.13.Theorem.1} If a distribution $T \in \mathfrak{D}(\Omega)'$ vanishes in each $U_i$ of a family $\{U_i; i \in I\}$ of open sets of $\Omega$, then $T$ vanishes in $U = \bigcup\limits_{i \in I} U_i$. 
\end{theorem}

\begin{proof}
 Let $\varphi \in \mathfrak{D}(\Omega)$ be a function with $\supp(\varphi) \subseteq U$. We construct a partition of unity $\{\alpha_j(x); j \in J\}$ subordinate to the covering of $\Omega$ consisting of $\{U_i; i \in I\}$ and $\Omega- \supp(\varphi)$. Then $\varphi = \sum\limits_{j \in J} \alpha_j \varphi$ is a finite sum and so $T(\varphi) = \sum\limits_{j \in J} T(\alpha_j \varphi)$. If the $\supp(\alpha_j)$ is contained in some $U_i$, $T(\alpha_j \varphi) = 0$ by the hypothesis; if the $\supp(\alpha_j)$ is contained in $\Omega- \supp(\varphi)$, then $\alpha_j \varphi = 0$ and so $T(\alpha_j \varphi) = 0$. Therefore we have $T(\varphi) = 0$. 
\end{proof}

\begin{proposition}\label{prop:I.13.Proposition.1} $A$ subset $B$ of the space $\mathfrak{E}(\Omega)$ is bounded iff, for any differential operator $D^{j}$ and for any compact subset $K$ of $\Omega$, the set of functions $\{D^{j}f(x); f \in B\}$ is uniformly bounded on $K$. 
\end{proposition}

\begin{proof}
 Clear from the definition of the semi-norms defining the topology of $\mathfrak{E}(\Omega)$. 
\end{proof}

\begin{proposition}\label{prop:I.13.Proposition.2} $A$ linear functional $T$ on $\mathfrak{E}(\Omega)$ is continuous iff $T$ is bounded on every bounded set of $\mathfrak{E}(\Omega)$. 
\end{proposition}

\begin{proof}
 Since $\mathfrak{E}(\Omega)$ is a quasi-normed linear space, the Proposition is a consequence of Theorem~\ref{thm:I.7.Theorem.2} in Chapter~\ref{chap:I}, \ref{sec:I.7}. 
\end{proof}

\begin{proposition}\label{prop:I.13.Proposition.3} $A$ distribution $T \in \mathfrak{D}(\Omega)'$ with compact support can be extended in one and only one way to a continuous linear functional $T_0$ on $\mathfrak{E}(\Omega)$ such that $T_0(f) = 0$ if $f \in \mathfrak{E}(\Omega)$ vanishes in a neighbourhood of $\supp(T)$. 
\end{proposition}

\begin{proof}
 Let us put $\supp(T) = K$ where $K$ is a compact subset of $\Omega$. For any point $x^0 \in K$ and $\varepsilon > 0$, we take a sphere $S(x^0, \varepsilon)$ of centre $x^0$ and radius $\varepsilon$. For any $\varepsilon > 0$ sufficiently small, the compact set $K$ is covered by a finite number of spheres $S(x^0, \varepsilon)$ with $x^0 \in K$. Let $\{\alpha_j(x); j \in J\}$ be the partition of the unity subordinate to this finite system of spheres. Then the function $\psi(x) = \sum\limits_{\supp(\alpha_j) \cap K' \neq \emptyset} \alpha_j(x)$, where $K'$ is a compact neighbourhood of $K$ contained in the interior of the finite system of spheres above, satisfies: 
\[
\psi (x) \in C_{0} ^{\infty} (\Omega) \quad \text{ and } \quad \psi (x) = 1 \text{ in a neighbourhood of } K. \]

We define $T_0(f)$ for $f \in C^\infty(\Omega)$ by $T_0(f) = T(\psi f)$. This definition is independent of the choice of $\psi$. For, if $\psi_1 \in C_0^\infty(\Omega)$ equals $1$ in a neighbourhood of $K$, then, for any $f \in C^\infty(\Omega)$, the function $(\psi- \psi_1)f \in \mathfrak{D}(\Omega)$ vanishes in a neighbourhood of $K$ so that $T(\psi f)- T(\psi_1f) = T((\psi- \psi_1)f) = 0$. 


It is easy to see, by applying Leibniz' formula of differentiation to $\psi f$, that $\{\psi f\}$ ranges over a bounded set of $\mathfrak{D}(\Omega)$ when $\{f\}$ ranges over a bounded set of $\mathfrak{E}(\Omega)$. Thus, since a distribution $T \in \mathfrak{D}(\Omega)'$ is bounded on bounded sets of $\mathfrak{D}(\Omega)$, the functional $T_0$ is bounded on bounded sets of $\mathfrak{E}(\Omega)$. Hence, by the Theorem~\ref{thm:I.7.Theorem.2} in Chapter~\ref{chap:I}, \ref{sec:I.7} mentioned above, $T_0$ is a continuous linear functional on $\mathfrak{E}(\Omega)$. Let $f \in \mathfrak{E}(\Omega)$ vanish in a neighbourhood $U(K)$ of $K$. Then, by choosing a $\psi \in C_0^\infty(\Omega)$ that vanishes in $(\Omega- U(K))$, we see that $T_0(f) = T(\psi f) = 0$.
\end{proof}

\begin{proposition}\label{prop:I.13.Proposition.4} Let $K'$ be the support of $\psi$ in the above definition of $T_0$. Then for some constants $C$ and $k$ 
\[
\left| T_{0} (f) \right| \leqq C \sup\limits_ {| j | \leqq k, x \in K ^{\prime}} \left| D ^{j} f (x) \right| \quad \text{ for all } \quad f \in C ^{\infty} (\Omega).
\]
\end{proposition}

\begin{proof}
 Since $T$ is a continuous linear functional on $\mathfrak{D}(\Omega)$, there exist, for any compact set $K'$ of $\Omega$, constants $C'$ and $k'$ such that 
\[
\left| T (\varphi) \right| \leqq C ^{\prime} \sup\limits_ {| j | \leqq k ^{\prime}, x \in K ^{\prime}} \left| D ^{j} \varphi (x) \right| \quad \text{ for all } \quad \varphi \in \mathfrak {D}_ {K ^{\prime}} (\Omega) \]

(the Corollary of Proposition~\ref{prop:I.8.Proposition.1} in Chapter~\ref{chap:I}, \ref{sec:I.8}). But, for any $g \in C^{\infty}(\Omega)$, we have $\varphi = \psi g \in D_{K'}(\Omega)$. Consequently, we see, by Leibniz' formula of differentiation, that 
\[
\sup\limits_ {| j | \leqq k ^{\prime}, x \in K ^{\prime}} | D ^{j} (\psi g) (x) | \leqq C ^{\prime \prime} \sup\limits_ {| j | \leqq k ^{\prime}, x \in K ^{\prime}} | D ^{j} g (x) |
\]
with a constant $C''$ which is independent of $g$. Setting $g = f$ and $k = k'$, we obtain the Proposition. 
\end{proof}

\begin{proposition}\label{prop:I.13.Proposition.5} Let $S_0$ be a linear functional on $C^\infty(\Omega)$ such that, for some constant $C$ and a positive integer $k$ and compact subset $K$ of $\Omega$, 
\[
\left| S_{0} (f) \right| \leqq C \sup\limits_{| j | \leqq k, x \in K} \left| D ^{j} f (x) \right| \quad \text{ for all } \quad f \in C ^{\infty} (\Omega).
\]
\end{proposition}

Then the restriction of $S_{0}$ to $C_{0}^{\infty}(\Omega)$ is a distribution $T$ with support contained in $K$. 

\begin{proof}
 We observe that $S_0(f) = 0$ if $f$ vanishes identically in a neighbourhood of $K$. Thus, if $\psi \in C_0^\infty(\Omega)$ equals $1$ in a neighbourhood of $K$, then 
\[
S_{0} (f) = S_{0} (\psi f) \quad \text{ for all } \quad f \in C ^{\infty} (\Omega). \]

It is easy to see that if $\{f\}$ ranges over a bounded set of $\mathfrak{D}(\Omega)$, then, in virtue of Leibniz' formula, $\{\psi f\}$ ranges over a set which is contained in a set of the form 
\[
\{g \in C ^{\infty} (\Omega); \sup\limits_{| j | \leqq k, x \in K} | D ^{j} g (x) | = C_{k} < \infty \}.
\]

Thus $S_0(\psi f) = T(f)$ is bounded on bounded sets of $\mathfrak{D}(\Omega)$ so that $T$ is a continuous linear functional on $\mathfrak{D}(\Omega)$. 
\end{proof}

We have thus proved the following 


\begin{theorem}\label{thm:I.13.Theorem.2} The set of all distributions in $\Omega$ with compact support is identical with the space $\mathfrak{E}(\Omega)'$ of all continuous linear functionals on $\mathfrak{E}(\Omega)$, the dual space of $\mathfrak{E}(\Omega)$. A linear functional $T$ on $C^\infty(\Omega)$ belongs to $\mathfrak{E}(\Omega)'$ iff, for some constants $C$ and $k$ and a compact subset $K$ of $\Omega$, 
\[
\left| T (f) \right| \leqq C \sup\limits_{| j | \leqq k, x \in K} \left| D ^{j} f (x) \right| \quad \text{ for all } \quad f \in C ^{\infty} (\Omega).
\]
\end{theorem}

We next prove a theorem which gives the general expression of distributions whose supports reduce to a single point. 

\begin{theorem}\label{thm:I.13.Theorem.3} Let an open set $\Omega$ of $\mathbb{R}^{n}$ contain the origin $0$. Then the only distributions $T \in \mathfrak{D}(\Omega)'$ with supports reduced to the origin $0$ are those which are expressible as finite linear combinations of Dirac's distribution and its derivatives at $0$. 
\end{theorem}

\begin{proof}
 For such a distribution $T$, there exist, by the preceding Theorem~\ref{thm:I.13.Theorem.2}, some constants $C$ and $k$ and a compact subset $K$ of $\Omega$ which contains the origin $0$ in such a way that 
\[
\left| T (f) \right| \leqq C \sup\limits_{| j | \leqq k, x \in K} \left| D ^{j} f (x) \right| \quad \text{ for all } \quad f \in C ^{\infty} (\Omega). \]

We shall prove that the condition 
\[
D ^{j} f (0) = 0 \quad \text{ for all } \quad j \quad \text{ with } | j | \leqq k
\]
implies $T(f) = 0$. To this purpose, we take a function $\psi \in C^{\infty}(\Omega)$ which is equal to $1$ in a neighbourhood of $0$ and put 
\[
f_{\varepsilon} (x) = f (x) \psi (x / \varepsilon).
\]

We have $T(f) = T(f_{\varepsilon})$ since $f = f_{\varepsilon}$ in a neighbourhood of the origin $0$. By Leibniz' formula, the derivative of $f_{\varepsilon}$ of order $\leqq k$ is a linear combination of terms of the form $|\varepsilon|^{-j}D^{i}\psi \cdot D^{i}f$ with $|i| + |j| \leqq k$. Since, by the assumption, $D^{i}f(0) = 0$ for $|i| \leqq k$, we see, by Taylor's formula, that a derivative of order $|s|$ of $f_{\varepsilon}$ is $0(\varepsilon^{k+1-|s|})$ in the support of $\psi (x / \varepsilon)$. Thus, when $\varepsilon \downarrow 0$, the derivatives of $f_{\varepsilon}$ of order $\leqq k$ converge to $0$ uniformly in a neighbourhood of $0$. Hence $T(f) = \lim\limits_{\varepsilon \downarrow 0}T(f_{\varepsilon}) = 0$. 

Now, for a general $f$, we denote by $f_k$ the Taylor's expansion of $f$ up to the order $k$ at the origin. Then, by what we have proved above, 
\[
T (f) = T \left(f_{k}\right) + T \left(f- f_{k}\right) = T \left(f_{k}\right) + 0 = T \left(f_{k}\right).
\]

This shows that $T$ is a linear combination of linear functionals in the derivatives of $f$ at the origin of order $\leqq k$. 
\end{proof}


\yossection{14}{The Direct Product of Generalized Functions}
\label{sec:I.14}

We first prove a theorem of approximation. 

\begin{theorem}\label{thm:I.14.Theorem.1} Let $x = (x_{1}, x_{2}, \ldots, x_{n}) \in \mathbb{R}^{n}$, $y = (y_{1}, y_{2}, \ldots, y_{m}) \in \mathbb{R}^{m}$ and $z = x \times y = (x_{1}, x_{2}, \ldots, x_{n}, y_{1}, y_{2}, \ldots, y_{m}) \in \mathbb{R}^{n+m}$. Then, for any function $\varphi(z) = \varphi(x, y) \in C_{0}^{\infty}(\mathbb{R}^{n+m})$, we can choose functions $u_{ij}(x) \in C_{0}^{\infty}(\mathbb{R}^{n})$ and functions $v_{ij}(y) \in C_{0}^{\infty}(\mathbb{R}^{m})$ such that the sequence of functions 
\begin{equation}\label{eq:I.14.1}
\varphi_{i} (z) = \varphi_{i} (x, y) = \sum\limits_{j=1} ^ {k_{i}} u_{ij} (x) v_{ij} (y)
\end{equation}
tends, as $i \to \infty$, to $\varphi(z) = \varphi(x, y)$ in the topology of $\mathfrak{D}(\mathbb{R}^{n+m})$. 
\end{theorem}

\begin{proof}
 We shall prove Theorem~\ref{thm:I.14.Theorem.1} for the case $n = m = 1$. Consider 
\begin{align}
\Phi (x, y, t)& = \left(2 \sqrt {\pi t}\right) ^{-2} \int_{- \infty} ^{\infty} \int_{- \infty} ^{\infty} \varphi (\xi , \eta) \exp\left(- ((x- \xi) ^{2} + (y- \eta) ^{2}) / 4 t\right) \dd \xi \dd \eta , \notag \\
t& > 0; \Phi (x, y, 0) = \varphi (x, y).  \label{eq:I.14.2}
\end{align}

We have, by the change of variables $\xi_{1} = (\xi- x)/2\sqrt{t}$, $\eta_{1} = (\eta- y)/2\sqrt{t}$, 
\[
\Phi (x, y, t) = (\sqrt {\pi}) ^{-2} \int_{- \infty} ^{\infty} \int_{- \infty} ^{\infty} \varphi (x + 2 \xi_{1} \sqrt {t}, y + 2 \eta_{1} \sqrt {t}) e ^ {- \xi_{1} ^{2}- \eta_{1} ^{2}} \dd \xi_{1} \dd \eta_{1}.
\]

Hence, by $\int_{-\infty}^{\infty}\int_{-\infty}^{\infty}e^{-\xi_1^2-\eta_1^2}\dd\xi_1\dd\eta_1 = \pi ,$ 
\[
\begin{aligned}
&\left| \Phi (x, y, t)- \varphi (x, y) \right| \leqq \left(\sqrt {\pi}\right) ^{-2} \int_{- \infty} ^{\infty} \int_{- \infty} ^{\infty} \left| \varphi (x + 2 \xi \sqrt {t}, y + 2 \eta \sqrt {t})- \varphi (x, y) \right| \times e ^ {- \xi^{2}- \eta^{2}} \dd \xi \dd \eta
\\
&\leqq \pi^{-1} \left\{\iint_ {\xi^{2} + \eta^{2} \geqq T ^{2}} + \iint_ {\xi^{2} + \eta^{2} < T ^{2}} \right\}.
\end{aligned}
\]

Since the function $\varphi$ is bounded and $e^{-\xi^2- \eta^2}$ is integrable in $\mathbb{R}^{2}$, we see that the first term on the right tends to zero as $T \uparrow \infty$. The second term on the right tends, for fixed $T > 0$, to zero as $t \downarrow 0$. Hence we have proved that $\lim\limits_{t \downarrow 0} \Phi(x, y, t) = \varphi(x, y)$ uniformly in $(x, y)$. 


Next, since $\supp(\varphi)$ is compact, we see, by partial integration, 
\[
\begin{aligned} \frac {\partial^{m+k} \Phi (x , y , t)}{\partial x ^{m} \partial y ^{k}} & = \iint_{- \infty} ^{\infty} (2 \sqrt {\pi t}) ^{-2} \frac {\partial^{m+k} \varphi (\xi , \eta)}{\partial \xi^{m} \partial \eta^{k}} e ^ {- [ (x- \xi) ^{2} + (y- \eta) ^{2} ] / 4 t} \dd \xi \dd \eta , t > 0, \\ & = \frac {\partial^{m+k} \varphi (x , y)}{\partial x ^{m} \partial y ^{k}}, t = 0. \end{aligned}
\]

Thus we see as above that 
\begin{equation}\label{eq:I.14.3}
\lim\limits_{t \downarrow 0} \frac {\partial^{m+k} \Phi (x , y , t)}{\partial x ^{m} \partial y ^{k}} = \frac {\partial^{m+k} \varphi (x , y)}{\partial x ^{m} \partial y ^{k}} \quad \text{ uniformly in } (x, y).
\end{equation}

It is easy to see that $\Phi(x, y, t)$ for $t > 0$ given by \eqref{eq:I.14.2} may be extended as a holomorphic function of the complex variables $x$ and $y$ for $|x| < \infty$, 
$|y|<\infty$. Hence, for any given $\gamma>0$, the function $\Phi(x,y,t)$ for fixed $t>0$ may be expanded into Taylor's series 
\[
\Phi (x, y, t) = \sum\limits_{m=0} ^{\infty} \sum\limits_{s=0} ^{m} c_{s} (t) x ^{s} y ^{m-s}
\]
which is absolutely and uniformly convergent for $|x| \leqq \gamma$, $|y| \leqq \gamma$ and may be differentiated term by term: 
\[
\frac {\partial^{m+k} \Phi (x , y , t)}{\partial x ^{m} \partial y ^{k}} = \sum\limits_ {m_{1} = 0} ^{\infty} \sum\limits_{s=0} ^ {m_{1}} c_{s} (t) \frac {\partial^{m+k} x ^{s} y ^ {m_{1}- s}}{\partial x ^{m} \partial y ^{k}}.
\]

Let $\{t_i\}$ be a sequence of positive numbers such that $t_i \downarrow 0$. By the above we can choose, for each $t_i$, a polynomial section $P_i(x, y)$ of the series $\sum\limits_{m=0}^{\infty} \sum\limits_{s=0}^{m} c_s(t) x^s y^{m-s}$ such that 
\[
\lim\limits_{i \to \infty} P_{i} (x, y) = \varphi (x, y) \text{ in the topology of } \mathfrak {E} (\mathbb{R}^{2}),
\]
that is, for any compact subset $K$ of $\mathbb{R}^{2}$, $\lim\limits_{i \to \infty} D^s P_i(x, y) = D^s \varphi(x, y)$ uniformly on $K$ for every differential operator $D^s$. Let us take $\varrho(x) \in C_0^\infty(\mathbb{R}^{1})$ and $\sigma(y) \in C_0^\infty(\mathbb{R}^{1})$ such that $\varrho(x) \sigma(y) = 1$ on the supP$(\varphi(x, y))$. Then we easily see that $\varphi_i(x, y) = \varrho(x) \sigma(y) P_i(x, y)$ satisfies the condition of Theorem~\ref{thm:I.14.Theorem.1}.
\end{proof}

\begin{remarku} We shall denote by $\mathfrak{D}(\mathbb{R}^{n}) \times \mathfrak{D}(\mathbb{R}^{m})$ the totality of functions $\in \mathfrak{D}(\mathbb{R}^{n+m})$ which are expressible as 
\[
\sum\limits_{j=1} ^{k} \varphi_{j} (x) \psi_{j} (y) \quad \text{ with } \quad \varphi_{j} (x) \in \mathfrak {D} (\mathbb{R}^{n}), \psi_{j} (y) \in \mathfrak {D} (\mathbb{R}^{m}).
\]
\end{remarku}

The above Theorem~\ref{thm:I.14.Theorem.1} says that $\mathfrak{D}(\mathbb{R}^{n}) \times \mathfrak{D}(\mathbb{R}^{m})$ is dense in $\mathfrak{D}(\mathbb{R}^{n+m})$ in the topology of $\mathfrak{D}(\mathbb{R}^{n+m})$. The linear subspace $\mathfrak{D}(\mathbb{R}^{n}) \times \mathfrak{D}(\mathbb{R}^{m})$ of $\mathfrak{D}(\mathbb{R}^{n+m})$ equipped with the relative topology is called the direct product of $\mathfrak{D}(\mathbb{R}^{n})$ and $\mathfrak{D}(\mathbb{R}^{m})$. 

We are now able to define the direct product of distributions. To indicate explicitly the independent variables $x = (x_1, x_2, \ldots, x_n)$ of the function $\varphi(x) \in \mathfrak{D}(\mathbb{R}^{n})$, we shall write $(\mathfrak{D}_x)$ for $\mathfrak{D}(\mathbb{R}^{n})$. We also write $(\mathfrak{D}_y)$ for $\mathfrak{D}(\mathbb{R}^{m})$ consisting of the functions $\psi(y), y = (y_1, y_2, \ldots, y_m)$. Likewise we shall write $(\mathfrak{D}_{x \times y})$ for $\mathfrak{D}(\mathbb{R}^{n+m})$ consisting of the functions $\chi(x, y)$. We shall accordingly write $T_{(x)}$ for the distribution $T \in \mathfrak{D}(\mathbb{R}^{n})' = (\mathfrak{D}_x)'$ in order to show that $T$ is to be applied to functions $\varphi(x)$ of $x$. 

\begin{theorem}\label{thm:I.14.Theorem.2} Let $T_{(x)} \in (\mathfrak{D}_x)'$, $S_{(y)} \in (\mathfrak{D}_y)'$. Then we can define in one and only one way a distribution $W = W_{(x \times y)} \in (\mathfrak{D}_{x \times y})'$ such that 
\begin{align}
W (u (x) v (y))& = T_{(x)} (u (x)) S_{(y)} (v (y)) \quad \text{ for } \quad u \in (\mathfrak {D}_{x}), v \in (\mathfrak {D}_{y}), \label{eq:I.14.4} \\
W (\varphi (x, y))& = S_{(y)} \left(T_{(x)} (\varphi (x, y))\right) = T_{(x)} \left(S_{(y)} (\varphi (x, y))\right) \quad \text{ for } \quad \varphi \in (\mathfrak {D}_{x \times y}) \label{eq:I.14.5}
\end{align}
\end{theorem}

\begin{remarku} The distribution $W$ is called the direct product or the tensor product of $T_{(x)}$ and $S_{(y)}$, and we shall write 
\begin{equation}\label{eq:I.14.6}
W = T_{(x)} \times S_{(y)} = S_{(y)} \times T_{(x)}.
\end{equation}
\end{remarku}

\begin{proof}[Proof of Theorem~\ref{thm:I.14.Theorem.2}.]
Let $\mathfrak{B} = \{\varphi(x, y)\}$ be a bounded set of the space $(\mathfrak{D}_{x \times y})$. For fixed $y^{(0)}$, the set $\{\varphi(x, y^{(0)}); \varphi \in \mathfrak{B}\}$ is a bounded set of $(\mathfrak{D}_x)$. We shall show that
\begin{equation}\label{eq:I.14.7}
\{\psi (y ^{(0)}); \psi (y ^{(0)}) = T_{(x)} (\varphi (x, y ^{(0)})), \varphi \in \mathfrak {B} \}
\end{equation}
is a bounded set of $(\mathfrak{D}_{y^{(0)}})$. The proof is given as follows. 

Since $\mathfrak{B}$ is a bounded set of $(\mathfrak{D}_{x\times y})$, there exist a compact set $K_{x}\subset \mathbb{R}^{n}$ and a compact set $K_{y}\subset \mathbb{R}^{m}$ such that 
$\supp(\varphi) \subseteq \{(x, y) \in \mathbb{R}^{n+m}; x \in K_x, y \in K_y\}$ whenever $\varphi \in \mathfrak{B}$. 

Hence $y^{(0)} \in K_y$ implies $\varphi(x, y^{(0)}) = 0$ and $\psi(y^{(0)}) = T_{(x)}(\varphi(x, y^{(0)})) = 0$. Thus 
\begin{equation}\label{eq:I.14.8}
\supp (\psi) \subseteq K_{y} \quad \text{ whenever } \quad \varphi \in \mathfrak {B}.
\end{equation}

We have to show that, for any differential operator $D_{y}$ in $\mathbb{R}^{m}$, 
\begin{equation}\label{eq:I.14.9}
\sup\limits_{y, \psi} | D_{y} \psi (y) | < \infty \text{ where } \psi (y) = T_{(x)} (\varphi (x, y)), \varphi \in \mathfrak {B}.
\end{equation}

To prove this, we take, e.g., $D_{y_1} = \partial/\partial y_1$. Then, by the linearity of $T_{(x)}$, 
\begin{gather}\quad \frac {\psi \left(y_{1} + h , y_{2} , \dots , y_{m}\right)- \psi \left(y_{1} , y_{2} , \dots , y_{m}\right)}{h} \notag \\= T_{(x)} \left\{\frac {\varphi (x , y_{1} + h , y_{2} , \dots , y_{m})- \varphi (x , y_{1} , y_{2} , \dots , y_{m})}{h} \right\}. \notag
\end{gather}

When $\varphi$ ranges over $\mathfrak{B}$, the functions $\left\{\dfrac{\varphi (x, y_{1} + h, y_{2}, \dots , y_{m}) - \varphi (x, y_{1}, y_{2}, \dots , y_{m})}{h}\right\}$ of $x$, with parameters $y \in \mathbb{R}^{m}$ and $h$ such that $|h| \leqq 1$, constitute a bounded set of $(\mathfrak{D}_x)$. This we see from the fact that $\mathfrak{B}$ is a bounded set of $(\mathfrak{D}_{x \times y})$. Hence we see, by letting $h \to 0$ and remembering Proposition~\ref{prop:I.8.Proposition.1} in Chapter~\ref{chap:I}, \ref{sec:I.8}, that \eqref{eq:I.14.9} is true. 

Therefore, by the same Proposition~\ref{prop:I.8.Proposition.1}, we see that the set of values 
\begin{equation}\label{eq:I.14.10}
\{S_{(y)} \left(T_{(x)} (\varphi (x, y))\right); \varphi \in \mathfrak {B} \}
\end{equation}
is bounded. Consequently, the same Proposition~\ref{prop:I.8.Proposition.1} shows that we have defined a distribution $W^{(1)} \in (\mathfrak{D}_{x \times y})'$ through 
\begin{equation}\label{eq:I.14.11}
W ^{(1)} (\varphi) = S_{(y)} \left(T_{(x)} (\varphi (x, y))\right).
\end{equation}

Similarly we define a distribution $W^{(2)} \in (\mathfrak{D}_{x \times y})'$ through 
\begin{equation}\label{eq:I.14.12}
W ^{(2)} (\varphi) = T_{(x)} \left(S_{(y)} (\varphi (x, y))\right).
\end{equation}

Clearly we have, for $u \in (\mathfrak{D}_x)$ and $v \in (\mathfrak{D}_y)$, 
\begin{equation}\label{eq:I.14.13}
W ^{(1)} (u (x) v (y)) = W ^{(2)} (u (x) v (y)) = T_{(x)} (u (x)) \cdot S_{(y)} (v (y)).
\end{equation}

Therefore, by the preceding Theorem~\ref{thm:I.14.Theorem.1} and the continuity of the distributions $W^{(1)}$ and $W^{(2)}$, we obtain $W^{(1)} = W^{(2)}$. This proves our Theorem~\ref{thm:I.14.Theorem.2} by setting $W = W^{(1)} = W^{(2)}$. 
\end{proof}

\yossub{References for Chapter I}

For locally convex linear topological spaces and Banach spaces, see N. \textsc{Bourbaki} \cite{bourbaki27}, A. \textsc{Grothendieck} \cite{grothendieck87}, G. \textsc{Köthe} \cite{konishi164}, S. \textsc{Banach} \cite{banach9}, N. \textsc{Dunford}-J. \textsc{Schwartz} \cite{dunford44} and E. \textsc{Hille}-R. S. \textsc{Phillips} \cite{hille97}. For generalized functions, see L. \textsc{Schwartz} \cite{schwartz294}, I. M. \textsc{Gelfand}-G. E. \textsc{Šilov} \cite{gelfand76}, L. \textsc{Hörmander} \cite{hormander118} and A. \textsc{Friedman} \cite{friedman68}.\yosstarfootnote{See also Supplementary Notes, p.~\pageref{notes:I-VI}.} 
   % I
% ==== II ====
\yoschapter{II}{Applications of the Baire-Hausdorff Theorem}
\label{chap:II}

The completeness of a $B$-space (or an $F$-space) enables us to apply the Baire-Hausdorff theorem in Chapter~\ref{chap:0}, and we obtain such basic principles in functional analysis as the uniform boundedness theorem, the resonance theorem, the open mapping theorem and the closed graph theorem. These theorems are essentially due to S. \textsc{Banach} \cite{banach9}. The termwise differentiability of generalized functions is a consequence of the uniform boundedness theorem. 


\yossection{1}{The Uniform Boundedness Theorem and the Resonance Theorem}
\label{sec:II.1}

\begin{theorem}[the uniform boundedness theorem]\label{thm:II.1.Theorem.1} Let $X$ be a linear topological space which is not expressible as a countable union of closed non-dense subsets. Let a family $\{T_{a}; a \in A\}$ of continuous mappings be defined on $X$ into a quasi-normed linear space $Y$. We assume that, for any $a \in A$ and $x, y \in X$, 

$\|T_a(x+y)\|\leqq\|T_a x\|+\|T_a y\|$ and $\|T_a(\alpha x)\|=\|\alpha T_a x\|$ for $\alpha\geqq0$. If the set $\{T_a x;a\in A\}$ is bounded at each $x\in X$, then $\slim\limits_{x\to 0}T_a x=0$ uniformly in $a\in A$. 
\end{theorem}

\begin{proof}
 For a given $\varepsilon > 0$ and for each positive integer $n$, consider $X_{n} = \{x \in X; \sup\limits_{a \in A} \{\|n^{-1}T_{a}x\| + \|n^{-1}T_{a}(-x)\|\} \leqq \varepsilon\}$. Each set $X_{n}$ is closed by the continuity of $T_{a}$. By the assumption of the boundedness of $\{\|T_{a}x\|; a \in A\}$, we have $X = \bigcup\limits_{n=1}^{\infty} X_{n}$. Hence, by the hypothesis on $X$, some $X_{n_{0}}$ must contain a neighbourhood $U = x_{0} + V$ of some point $x_{0} \in X$, where $V$ is a neighbourhood of $0$ of $X$ such that $V = -V$. Thus $x \in V$ implies $\sup\limits_{a \in A} \|n_{0}^{-1}T_{a}(x_{0} + x)\| \leqq \varepsilon$. Therefore we have 
\[
\begin{aligned}
\left\| T _{a} (n _{0} ^{-1} x) \right\|& = \left\| T _{a} (n _{0} ^{-1} (x _{0} + x - x _{0})) \right\| \leqq \left\| n _{0} ^{-1} T _{a} (x _{0} + x) \right\|
\\
\quad + \left\| n _{0} ^{-1} T _{a} (- x _{0}) \right\|& \leqq 2 \varepsilon \text{ for } x \in V, a \in A. \end{aligned}
\]

Thus the Theorem is proved, because the scalar multiplication $\alpha x$ in a linear topological space is continuous in both variables $\alpha$ and $x$. 
\end{proof}

\begin{corollary}[the resonance theorem]\label{cor:II.1.Corollary.1} Let $\{T_{a}; a \in A\}$ be a family of bounded linear operators defined on a $B$-space $X$ into a normed linear space $Y$. Then the boundedness of $\{\| T_{a}x\|; a \in A\}$ at each $x \in X$ implies the boundedness of $\{ \| T_{a}\|; a \in A\}$. 
\end{corollary}

\begin{proof}
 By the uniform boundedness theorem, there exists, for any $\varepsilon > 0$, a $\delta > 0$ such that $\|x\| \leqq \delta$ implies $\sup\limits_{a \in A} \|T_a x\| \leqq \varepsilon$. Thus $\sup\limits_{a \in A} \|T_a\| \leqq \varepsilon / \delta$. 
\end{proof}

\begin{corollary}\label{cor:II.1.Corollary.2} Let $\{T_n\}$ be a sequence of bounded linear operators defined on a $B$-space $X$ into a normed linear space $Y$. Suppose that $\slim\limits_{n\to \infty}T_n x = Tx$ exists for each $x\in X$. Then $T$ is also a bounded linear operator on $X$ into $Y$ and we have 
\begin{equation}\label{eq:II.1.1}
\left\| T \right\| \leqq \varliminf\limits_{n \rightarrow \infty} \left\| T _{n} \right\|.
\end{equation}
\end{corollary}

\begin{proof}
 The boundedness of the sequence $\{|T_{n}x|\}$ for each $x\in X$ is implied by the continuity of the norm. Hence, by the preceding Corollary, $\sup\limits_{n\geqq1}\|T_{n}\|<\infty$, and so $\|T_{n}x\|\leqq\sup\limits_{n\geqq1}\|T_{n}\|\cdot\|x\|$ ( $n=1,2,\ldots$). 

Therefore, again by the continuity of the norm, we obtain 
\[
\left\| T x \right\| = \lim\limits _{n \rightarrow \infty} \left\| T _{n} x \right\| \leqq \underline {{\lim\limits _{n \rightarrow \infty}}} \left\| T _{n} \right\| \cdot \left\| x \right\|,
\]
which is precisely the inequality \eqref{eq:II.1.1}. Finally it is clear that $T$ is linear. 
\end{proof}

\begin{definitionu} The operator $T$ obtained above is called the strong limit of the sequence $\{T_n\}$ and we shall write $T = \slim T_n$. 
\end{definitionu}

We next prove an existence theorem for the bounded inverse of a bounded linear operator. 

\begin{theorem}[C. \textsc{Neumann}]\label{thm:II.1.Theorem.2} Let $T$ be a bounded linear operator on a $B$-space $X$ into $X$. Suppose that $\|I - T\| < 1$, where $I$ is the identity operator: $I \cdot x = x$. Then $T$ has a unique bounded linear inverse $T^{-1}$ which is given by C. \textsc{Neumann}'s series 
\begin{equation}\label{eq:II.1.2}
T ^{-1} x = \slim _{n \rightarrow \infty} (I + (I - T) + (I - T) ^{2} + \dots + (I - T) ^{n}) x, x \in X.
\end{equation}
\end{theorem}

\begin{proof}
 For any $x \in X$, we have 
\[ \begin{aligned}
&\left\| \sum\limits_{n=0} ^{k} (I - T) ^{n} x \right\|  \leqq \sum\limits_{n=0} ^{k} \| (I - T) ^{n} x \| \leqq \sum\limits_{n=0} ^{k} \| (I - T) ^{n} \| | x \| \leqq \sum\limits_{n=0} ^{\infty} \| I - T \| ^{n} \| x \|.
\end{aligned}
\]

The right hand side is convergent by $\|I - T\| < 1$. Hence, by the completeness of $X$, the bounded linear operator $\slim\limits_{k \to \infty} \sum\limits_{n=0}^{k} (I - T)^n$ is defined. 


It is the inverse of $T$ as may be seen from 
\[
\begin{aligned}T \cdot \slim _{k \rightarrow \infty} \sum\limits_{n=0} ^{k} (I - T) ^{n} x&= \slim _{k \rightarrow \infty} (I - (I - T)) \left(\sum\limits_{n=0} ^{k} (I - T) ^{n} x\right)\\&= x - \slim _{k \rightarrow \infty} (I - T) ^{k+1} x = x,\end{aligned}
\]
and the similar equation $s = \lim\limits_{k \to \infty} \left( \sum\limits_{n=0}^{k} (I - T)^n \right) Tx = x$.
\end{proof}

\yossection{2}{The Vitali-Hahn-Saks Theorem}
\label{sec:II.2}

This theorem is concerned with a convergent sequence of measures, and makes use of the following 

\begin{propositionu} Let $(S, \mathfrak{B}, m)$ be a measure space. Let $\mathfrak{B}_0$ be the set of all $B \in \mathfrak{B}$ such that $m(B) < \infty$. Then by the distance 
\begin{align}
d \left(B _{1}, B _{2}\right)& = m \left(B _{1} \ominus B _{2}\right) \notag \\
B _{1} \ominus B _{2}& = B _{1} \cup B _{2} - B _{1} \cap B _{2}, \label{eq:II.2.1}
\end{align}
$\mathfrak{B}_0$ gives rise to a metric space $(\mathfrak{B}_0)$ by identifying two sets $B_1$ and $B_2$ of $\mathfrak{B}_0$ when $m(B_1 \ominus B_2) = 0$. Thus a point $\overline{B}$ of $(\mathfrak{B}_0)$ is the class of sets $B_1 \in \mathfrak{B}_0$ such that $m(B \ominus B_1) = 0$. Under the above metric \eqref{eq:II.2.1}, $(\mathfrak{B}_0)$ is a complete metric space. 
\end{propositionu}

\begin{proof}
 If we denote by $C_B(s)$ the defining function of the set $B$: 
\[
C _{B} (s) = 1 \text{ or } 0 \text{ according as } s \in B \text{ or } s \in B,
\]
we have 
\begin{equation}\label{eq:II.2.2}
d \left(B _{1}, B _{2}\right) = \int_{S} \left| C _ {B _{1}} (s) - C _ {B _{2}} (s) \right| m (\dd s). \end{equation}

Thus the metric space $(\mathfrak{B}_0)$ may be identified with a subset of the $B$-space $L^1(S, \mathfrak{B}, m)$. Let a sequence $\{C_{B_{n}}(s)\}$ with $B_n \in \mathfrak{B}_0$ satisfy the condition 
\[
\lim\limits _{n, k \rightarrow \infty} d (B _{n}, B _{k}) = \lim\limits _{n, k \rightarrow \infty} \int_{S} | C _ {B _{n}} (s) - C _ {B _{k}} (s) | m (\dd s) = 0.
\]

Then, as in the proof of the completeness of the space $L^1(S, \mathfrak{B}, m)$, we can choose a subsequence $\{C_{B_n'}(s)\}$ such that $\lim\limits_{n' \to \infty} C_{B_n'}(s) = C(s)$ exists $m$-a.e. and $\lim\limits_{n' \to \infty} \int_S |C(s) - C_{B_n'}(s)| m(\dd s) = 0$. Clearly $C(s)$ is the defining function of a set $B_\infty \in \mathfrak{B}_0$, and hence $\lim\limits_{n' \to \infty} d(B_\infty, B_n) = 0$. 

This proves that $(\mathfrak{B}_{0})$ is a complete metric space. 
\end{proof}

\begin{theoremn}[\textsc{Vitali}-\textsc{Hahn}-\textsc{Saks}]\label{thm:II.2.VitaliHahnSaks} Let $(S, \mathfrak{B}, m)$ be a measure space with $m(S) < \infty$, and $\{\lambda_n(B)\}$ a sequence of complex measures such that the total variations $|\lambda_n|(S)$ are finite for $n = 1, 2, \ldots$. Suppose that each $\lambda_n(B)$ is $m$-absolutely continuous and that a finite $\lim\limits_{n \to \infty} \lambda_n(B) = \lambda(B)$ exists for every $B \in \mathfrak{B}$. Then the $m$-absolute continuity of $\lambda_n(B)$ is uniform in $n$, that is, $\lim\limits_{n} m(B) = 0$ implies $\lim\limits_{n} \lambda_n(B) = 0$ uniformly in $n$. And $\lambda(B)$ is $\sigma$-additive on $\mathfrak{B}$.

\end{theoremn}
\begin{proof}
 Each $\lambda_{n}$ defines a one-valued function $\bar{\lambda}_{n}(\overline{B})$ on $(\mathfrak{B}_0)$ by $\bar{\lambda}_{n}(\overline{B}) = \lambda_{n}(B)$, since the value $\lambda_{n}(B)$ is, by the $m$-absolute continuity of $\lambda_{n}(B)$, independent of the choice of the set $B$ from the class $\overline{B}$. The continuity of $\bar{\lambda}_{n}(\overline{B})$ is implied by and implies the $m$-absolute continuity of $\lambda_{n}(B)$. 

Hence, for any $\varepsilon > 0$, the set 
\[
F _{k} (\varepsilon) = \left\{\overline {{B}}; \sup\limits _{n \geqq 1} \left| \bar {\lambda} _{k} (\overline {{B}}) - \bar {\lambda} _{k+n} (\overline {{B}}) \right| \leqq \varepsilon \right\}
\]
is closed in $(\mathfrak{B}_0)$ and, by the hypothesis $\lim\limits_{n\to \infty}\lambda_n(B) = \lambda (B)$, we have $(\mathfrak{B}_0) = \bigcup\limits_{k=1}^{\infty}F_k(\varepsilon)$. The complete metric space $(\mathfrak{B}_0)$ being of the second category, at least one $F_{k_0}(\varepsilon)$ must contain a non-void open set of $(\mathfrak{B}_0)$. This means that there exist a $\overline{B}_0\in (\mathfrak{B}_0)$ and an $\eta >0$ such that 
\[
d (B, B _{0}) < \eta \quad \text{ implies } \quad \sup\limits _{n \geqq 1} | \bar {\lambda} _ {k _{0}} (\overline {{B}}) - \bar {\lambda} _ {k _{0} + n} (\overline {{B}}) | \leqq \varepsilon .
\]

On the other hand, any $B \in \mathfrak{B}_0$ with $m(B) \leqq \eta$ can be represented as $B = B_1 - B_2$ with $d(B_1, B_0) \leqq \eta$, $d(B_2, B_0) \leqq \eta$. We may, for example, take $B_1 = B \cup B_0$, $B_2 = B_0 - B \cap B_0$. Thus, if $m(B) \leqq \eta$ and $k \geqq k_0$, we have 
\[
\begin{aligned}
&\left| \lambda_{k} (B) \right|  \leqq \left| \lambda_ {k _{0}} (B) \right| + \left| \lambda_ {k _{0}} (B) - \lambda_{k} (B) \right| \leqq \left| \lambda_ {k _{0}} (B) \right| + \left| \lambda_ {k _{0}} (B _{1}) - \lambda_{k} (B _{1}) \right| + \left| \lambda_ {k _{0}} (B _{2}) - \lambda_{k} (B _{2}) \right| \\
&\leqq \left| \lambda_ {k _{0}} (B) \right| + 2 \varepsilon .
\end{aligned}
\]

Therefore, by the $m$-absolute continuity of $\lambda_{k_0}$ and the arbitrariness of $\varepsilon > 0$, we see that $m(B) \to 0$ implies $\lambda_n(B) \to 0$ uniformly in $n$. Hence, in particular, $m(B) \to 0$ implies $\lambda(B) \to 0$. On the other hand, it is clear that $\lambda$ is finitely additive, i.e., $\lambda\left(\sum\limits_{j=1}^{n} B_j\right) = \sum\limits_{j=1}^{n} \lambda(B_j)$. Thus, by $\lim\limits_{m(B) \to 0} \lambda(B) = 0$ proved above, we easily see that $\lambda$ is $\sigma$-additive since $m(S) < \infty$. 
\end{proof}

\begin{corollary}\label{cor:II.2.Corollary.1} Let $\{\lambda_n(B)\}$ be a sequence of complex measures on $S$ such that the total variation $|\lambda_n|(S)$ is finite for every $n$. If a finite $\lim\limits_{n\to \infty}\lambda_n(B)$ exists for every $B\in \mathfrak{B}$, then the $\sigma$-additivity of $\{\lambda_n(B)\}$ is uniform in $n$, in the sense that, $\lim\limits_{k\to \infty}\lambda_n(B_k) = 0$ uniformly in $n$ for any decreasing sequence $\{B_k\}$ of sets $\in \mathfrak{B}$ such that $\bigcap\limits_{k=1}^{\infty}B_k = \emptyset$. 
\end{corollary}

\begin{proof}
 Let us consider 
\[
m (B) = \sum\limits_{j=1} ^{\infty} 2 ^{-j} \mu_{j} (B) \quad \text{ where } \quad \mu_{j} (B) = | \lambda_{j} | (B) / | \lambda_{j} | (S). \]

The $\sigma$-additivity of $m$ is a consequence of that of the $\lambda_j$'s, and we have $0 \leqq m(B) \leqq 1$. Each $\mu_j$ and hence each $\lambda_j$ is $m$-absolutely continuous. Thus, by the above Theorem, we have $\lim\limits_{k \to \infty} \lambda_n(B_k) = 0$ uniformly in $n$, because $\lim\limits_{k \to \infty} m(B_k) = 0$. 
\end{proof}

\begin{corollary}\label{cor:II.2.Corollary.2} $\lambda(B)$ in the Theorem is $\sigma$-additive and $m$-absolutely continuous even if $m(S)=\infty$. 
\end{corollary}


\yossection{3}{The Termwise Differentiability of a Sequence of Generalized Functions}
\label{sec:II.3}

The discussion of the convergence of a sequence of generalized functions is very simple. We can in fact prove 

\begin{theoremu} Let $\{T_n\}$ be a sequence of generalized functions $\in \mathfrak{D}(\Omega)'$. Let a finite $\lim\limits_{n\to \infty}T_n(\varphi) = T(\varphi)$ exist for every $\varphi \in \mathfrak{D}(\Omega)$. Then $T$ is also a generalized function $\in \mathfrak{D}(\Omega)'$. We say then that $T$ is the limit in $\mathfrak{D}(\Omega)'$ of the sequence $\{T_n\}$ and write $T = \lim\limits_{n\to \infty}T_n(\mathfrak{D}(\Omega))'$. 
\end{theoremu}

\begin{proof}
 The linearity of the functional $T$ is clear. Let $K$ be any compact subset of $\Omega$. Then each $T_n$ defines a linear functional on the $F$-space $\mathfrak{D}_K(\Omega)$. Moreover, these functionals are continuous. For, they are bounded on every bounded set of $\mathfrak{D}_K(\Omega)$ as may be proved by Proposition~\ref{prop:I.8.Proposition.1} in Chapter~\ref{chap:I}, \ref{sec:I.8}. Thus, by the uniform boundedness theorem, $T$ must be a linear functional on $\mathfrak{D}_K(\Omega)$ which is bounded on every bounded set of $\mathfrak{D}_K(\Omega)$. Hence $T$ is a continuous linear functional on every $\mathfrak{D}_K(\Omega)$. Since $\mathfrak{D}(\Omega)$ is the inductive limit of $\mathfrak{D}_K(\Omega)'s$, $T$ must be a continuous linear functional on $\mathfrak{D}(\Omega)$. 
\end{proof}

\begin{corollaryn}[termwise differentiability theorem] Let $T = \lim\limits_{n\to \infty}T_n(\mathfrak{D}(\Omega)')$. Then, for any differential operator $D^j$, we have $D^{j}T = \lim\limits_{n\to \infty}D^{j}T_{n}(\mathfrak{D}(\Omega))'$.

\end{corollaryn}
\begin{proof}
 $\lim\limits_{n\to \infty}T_n = T(\mathfrak{D}(\Omega)')$ implies that $\lim\limits_{n\to \infty}T_n((-1)^{|j|}D^j\varphi) = T((-1)^{|j|}D^j\varphi)$ for every $\varphi \in \mathfrak{D}(\Omega)$. Thus we have 
\[
(D ^{j} T) (\varphi) = \lim\limits _{n \rightarrow \infty} (D ^{j} T _{n}) (\varphi) \quad \text{ for every } \quad \varphi \in \mathfrak {D} (\Omega). \]
\end{proof}


\yossection{4}{The Principle of the Condensation of Singularities}
\label{sec:II.4}

The Baire-Hausdorff theorem may be applied to prove the existence of a function with various singularities. For instance, we shall prove the existence of a continuous function without a finite derivative. 

\begin{theoremn}[Weierstrass' Theorem] There exists a real-valued continuous function $x(t)$ defined on the interval [$0$, $1$] such that it admits a finite differential quotient $x'(t_0)$ at no point $t_0$ of the interval [$0$, $1$/$2$]. 
\end{theoremn}

\begin{proof}
 A function $x(t)$ admits finite upper- and lower-derivatives on the right at $t = t_{0}$ iff there exists a positive integer $n$ such that 
\[
\sup\limits _ {2 ^{-1} > h > 0} h ^{-1} \left| x \left(t _{0} + h\right) - x \left(t _{0}\right) \right| \leqq n. \]

Let us denote by $M_{n}$ the set of all functions $x(t) \in C[0,1]$ such that $x(t)$ satisfies the above inequality at a certain point $t_0$ of $[0,1/2]$; here the point $t_0$ may vary with the function $x(t)$. We have only to show that each $M_{n}$ is a non-dense subset of $C[0,1]$. For, then, the set $C[0,1] - \bigcup_{n=1}^{\infty} M_{n}$ is non-void since; by the Baire-Hausdorff theorem, the complete metric space $C[0,1]$ is not of the first category. 

Now it is easy to see, by the compactness of the closed interval [$0$, $1$/$2$] and the norm of the space $C[0,1]$, that $M_{n}$ is closed in $C[0,1]$. Next, for any polynomial $z(t)$ and $\varepsilon > 0$, we can find a function $y(t) \in C[0,1] - M_{n}$ such that $\sup\limits_{0 \leqq t \leqq 1} |z(t) - y(t)| = \|z - y\| \leqq \varepsilon$. We may take, for example, a continuous function $y(t)$ represented graphically by a zig-zag line so that the above conditions are satisfied. Hence, by Weierstrass' polynomial approximation theorem, $M_{n}$ is non-dense in $C[0,1]$. 
\end{proof}

S. \textsc{Banach} \cite{banach9} and H. \textsc{Steinhaus} have proved a principle of condensation of singularities given below which is based upon the following 

\begin{theoremn}[S. \textsc{Banach}] Given a sequence of bounded linear operators $\{T_{n}\}$ defined on a $B$-space $X$ into a normed linear space $Y_{n}$. Then the set
\[
B = \left\{x \in X; \varlimsup\limits_{n \rightarrow \infty} \| T _{n} x \| < \infty \right\}
\]
either coincides with $X$ or is a set of the first category of $X$. 

\end{theoremn}
\begin{proof}
 We shall show that $B = X$ under the hypothesis that $B$ is of the second category. By the definition of $B$, we have $\limsup\limits_{k \to \infty} \left\|k^{-1} T_n x \right\| = 0$ whenever $x \in B$. Thus, for any $\varepsilon > 0$, 
\[
B \subseteq \bigcup\limits_{k=1} ^{\infty} B _{k} \quad \text{ where } \quad B _{k} = \left\{x \in X; \sup\limits _{n \geqq 1} \| k ^{-1} T _{n} x \| \leqq \varepsilon \right\}. \]

Each $B_{k}$ is a closed set by the continuity of the $T_{n}$'s. Hence, if $B$ is of the second category, then a certain $B_{k_0}$ contains a sphere of positive radius. That is, there exist an $x_0 \in X$ and a $\delta > 0$ such that $\|x - x_0\| \leqq \delta$ implies $\sup\limits_{n \geqq 1} \|k_0^{-1}T_n x\| \leqq \varepsilon$. Hence, by putting $x - x_0 = y$, we obtain, for $\|y\| \leqq \delta$, $\|k_0^{-1}T_n y\| \leqq \|k_0^{-1}T_n x\| + \|k_0^{-1}T_n x_0\| \leqq 2\varepsilon$. We have thus 
\[
\sup\limits _ {n \geqq 1, \| z \| \leqq k _{0} ^{-1} \delta} \| T _{n} z \| \leqq 2 \varepsilon ,
\]
and so $B = X$. 
\end{proof}

\begin{corollaryn}[the principle of the condensation of singularities] Let $\{T_{p,q}\} (q = 1,2,\ldots)$ be a sequence of bounded linear operators defined on a $B$-space $X$ into a normed linear space $Y_{p}$ ( $p = 1,2,\ldots$). Let, for each $p$, there exist an $x_{p} \in X$ such that $\varlimsup\limits_{q \to \infty} \|T_{p,q} x_{p}\| = \infty$. Then the set
\[
B = \left\{x \in X; \varlimsup\limits_{q \rightarrow \infty} \| T _{p,q} x \| = \infty \quad \text{ for all } \quad p = 1, 2, \dots \right\}
\]
is of the second category. 

\end{corollaryn}
\begin{proof}
 For each $p$, the set $B_{p} = \{x \in X; \varlimsup\limits_{q \to \infty} \|T_{p,q} x\| < \infty\}$ is, by the preceding Theorem and the hypothesis, of the first category. Thus $B = X - \bigcup\limits_{p=1}^{\infty} B_{p}$ must be of the second category. 
\end{proof}

The above Corollary gives a general method of finding functions with many singularities. 

\begin{examplen} There exists a real-valued continuous function $x(t)$ of period $2\pi$ such that the partial sum of its Fourier expansion:
\begin{equation}\label{eq:II.4.1}
f _{q} (x; t) = \sum\limits_{k=0} ^{q} \left(a _{k} \cos k t + b _{k} \sin k t\right) = \frac {1}{\pi} \int_{- \pi} ^{\pi} x (s) K _{q} (s, t) \dd s,
\end{equation}
where $K_{q}(s,t) = \sin ((q + 2^{-1})(s - t)) / 2\sin 2^{-1}(s - t)$, 


satisfies the condition that 
\begin{equation}\label{eq:II.4.2}
\varlimsup\limits_{q \to \infty} | f _{q} (x; t) | = \infty \text{ on a set } P \subseteq [ 0, 2 \pi ] \text{ which is of the }
\end{equation}
power of the continuum. 

Moreover, the set $P$ may be taken so as to contain any countable sequence $\{t_j\} \subseteq [0, 2\pi]$. 

\end{examplen}
\begin{proof}
 The totality of real-valued continuous functions $x(t)$ of period $2\pi$ constitutes a real $B$-space $C_{2\pi}$ by the norm $\|x\| = \sup\limits_{0 \leqq t \leqq 2\pi} |x(t)|$. As may be seen from \eqref{eq:II.4.1}, $f_q(x; t)$ is, for a given $t_0 \in [0, 2\pi]$, a bounded linear functional on $C_{2\pi}$. Moreover, the norm of this functional $f_q(x; t_0)$ is given by 
\begin{equation}\label{eq:II.4.3}
\begin{aligned}
&\frac {1}{\pi} \int_{- \pi} ^{\pi} | K _{q} (s, t _{0}) | \dd s = \text{ the total variation of the function } \frac {1}{\pi} \int_{- \pi} ^{t} K _{q} (s, t _{0}) \dd s.
\end{aligned}
\end{equation}

It is also easy to see that, for fixed $t_0$, \eqref{eq:II.4.3} tends, as $q \to \infty$, to $\infty$. 

Therefore, if we take a countable dense sequence $\{t_j\} \subseteq [0, 2\pi]$, then, by the preceding Corollary, the set 
\[
B = \left\{x \in C _{2 \pi}; \varlimsup\limits_{q \rightarrow \infty} | f _{q} (x; t) | = \infty \quad \text{ for } \quad t = t _{1}, t _{2}, \dots \right\}
\]
is of the second category. Hence, by the completeness of the space $C_{2\pi}$, the set $B$ is non-void. We shall show that, for any $x \in B$, the set 
\[
P = \left\{t \in [ 0, 2 \pi ]; \varlimsup\limits_{q \rightarrow \infty} | f _{q} (x; t) | = \infty \right\}
\]
is of the power of the continuum. To this purpose, set 
\[
F _{m,q} = \{t \in [ 0, 2 \pi ]; | f _{q} (x; t) | \leqq m \}, F _{m} = \bigcap\limits_{q=1} ^{\infty} F _{m,q}.
\]

By the continuity of $x(t)$ and the trigonometric functions, we see that the set $F_{m,q}$ and hence the set $F_{m}$ are closed sets of $[0, 2\pi]$. If we can show that $\bigcup\limits_{m=1}^{\infty} F_{m}$ is a set of the first category of $[0, 2\pi]$, then the set $P = \left([0, 2\pi] - \bigcup\limits_{m=1}^{\infty} F_{m}\right) \ni \{t_j\}$ would be of the second category. Being a set of the second category of $[0, 2\pi]$, $P$ cannot be countable and so, by the continuum hypothesis, $P$ must have the power of the continuum. 


Finally we will prove that each $F_{m}$ is a set of the first category of $[0, 2\pi]$. Suppose some $F_{m_0}$ be of the second category. Then the closed set $F_{m_0}$ of $[0, 2\pi]$ must contain a closed interval $[\alpha, \beta]$ of $[0, 2\pi]$. This implies that $\sup\limits_{q \geqq 1} |f_q(x; t)| \leqq m_0$ for all $t \in [\alpha, \beta]$, contradicting the fact that the set $P$ contains a dense subset $\{t_j\}$ of $[0, 2\pi]$.
\end{proof}

\begin{remarku} We can prove that the set $P$ has the power of the continuum without appealing to the continuum hypothesis. See, for example, F. \textsc{Hausdorff} \cite{halmos94}. 
\end{remarku}


\yossection{5}{The Open Mapping Theorem}
\label{sec:II.5}

\begin{theoremn}[the open mapping theorem of S. \textsc{Banach}] Let $T$ be a continuous linear operator on an $F$-space $X$ onto an $F$-space $Y$. Then $T$ maps every open set of $X$ onto an open set of $Y$.

For the proof, we prepare 

\end{theoremn}
\begin{propositionu} Let $X$, $Y$ be linear topological spaces. Let $T$ be a continuous linear operator on $X$ into $Y$, and suppose that the range $R(T)$ of $T$ is a set of the second category in $Y$. Then, to each neighbourhood $U$ of $0$ of $X$, there corresponds some neighbourhood $V$ of $0$ of $Y$ such that $V \subseteq (TU)^{a}$. 
\end{propositionu}

\begin{proof}
 Let $W$ be a neighbourhood of $0$ of $X$ such that $W = -W$, $W + W \subseteq U$. For every $x \in X$, $x / n \to 0$ as $n \to \infty$, and so $x \in nW$ for large $n$. Hence $X = \bigcup\limits_{n=1}^{\infty}(nW)$, and therefore $R(T) = \bigcup\limits_{n=1}^{\infty} T(nW)$. Since $R(T)$ is of the second category in $Y$, there is some positive integer $n_0$ such that $(T(n_0W))^a$ contains a non-void open set. Since $(T(nW))^a = n(T(W))^a$ and since $n(T(W))^a$ and $(T(W))^a$ are homeomorphic\footnote{A one-to-one mapping $M$ of a topological space $S_{1}$ onto a topological space $S_{2}$ is called a homeomorphic mapping if $M$ and $M^{-1}$ both map open sets onto open sets.} to each other, $(T(W))^a$ also contains a non-void open set. Let $y_0 = Tx_0$, $x_0 \in W$, be a point of this open set. Then, since the mapping $x \to -x_0 + x$ is a homeomorphic mapping, we see that there exists a neighbourhood $V$ of $0$ of $Y$ such that $V \subseteq -y_0 + (T(W))^a$. The elements of $-y_0 + T(W)$ are expressible as $-y_0 + Tw = T(w - x_0)$ with $w \in W$. But $w - x_0 \in W + W \subseteq U$, since $W$ is a neighbourhood of $0$ of $X$ such that $W = -W$ and $W + W \subseteq U$. 

Therefore, $-y_0 + T(W) \subseteq T(U)$ and hence, by taking the closure, $-y_0 + (T(W))^a \subseteq (T(U))^a$ and so $V \subseteq -y_0 + (T(W))^a \subseteq (T(U))^a = (TU)^a$. 
\end{proof}

\begin{proof}[Proof of the Theorem.]
Since $Y$ is a complete metric space, it is of the second category. Thus, by the preceding Proposition, the closure of the image by $T$ of a neighbourhood of $0$ of $X$ contains a neighbourhood of $0$ of $Y$.

Let $X_{\varepsilon}, Y_{\varepsilon}$ denote the spheres in $X, Y$ respectively, with centres at the origins and radii $\varepsilon > 0$. Let us set $\varepsilon_i = \varepsilon / 2^i$ ( $i = 0, 1, 2, \ldots$). Then by what we have stated above, there exists a sequence of positive numbers $\{\eta_i\}$ such that 
\begin{equation}\label{eq:II.5.1}
\lim\limits _{i \rightarrow \infty} \eta_{i} = 0 \quad \text{ and } \quad Y _ {\eta_{i}} \subseteq (T X _ {\varepsilon_{i}}) ^{a} \quad (i = 0, 1, 2, \dots).
\end{equation}

Let $y \in Y_{\eta_0}$ be any point. We shall show that there is an $x \in X_{2\varepsilon_0}$ such that $Tx = y$. From \eqref{eq:II.5.1} with $i = 0$, we see that there is an $x_0 \in X_{\varepsilon_0}$ such that $\|y - Tx_0\| < \eta_1$. Since $(y - Tx_0) \in Y_{\eta_1}$, we see again from \eqref{eq:II.5.1} with $i = 1$ that there is an $x_1 \in X_{\varepsilon_1}$ with $\|y - Tx_0 - Tx_1\| < \eta_2$. Repeating the process, we find a sequence $\{x_i\}$ with $x_i \in X_{\varepsilon_i}$ such that 
\[
\left\| y - T \left(\sum\limits_{i=0} ^{n} x _{i}\right) \right\| < \eta_{n+1} \quad (n = 0, 1, 2, \dots).
\]

We have $\left\| \sum\limits_{k=m+1}^{n} x_k \right\| \leqq \sum\limits_{k=m+1}^{n} \left\| x_k \right\| \leqq \sum\limits_{k=m+1}^{n} \varepsilon_k \leqq \left( \sum\limits_{k=m+1}^{n} 2^{-k} \right) \varepsilon_0$, and so the sequence $\left\{ \sum\limits_{k=0}^{n} x_k \right\}$ is a Cauchy sequence. Hence, by the completeness of $X$, we see that $\slim\limits_{n \to \infty} \sum\limits_{k=0}^{n} x_k = x \in X$ exists. Moreover, we have 
\[
\left\| x \right\| = \lim\limits _{n \rightarrow \infty} \left\| \sum\limits_{k=0} ^{n} x _{k} \right\| \leqq \lim\limits _{n \rightarrow \infty} \sum\limits_{k=0} ^{n} \left\| x _{k} \right\| \leqq \left(\sum\limits_{k=0} ^{\infty} 2 ^{-k}\right) \varepsilon_{0} = 2 \varepsilon_{0}.
\]

We must have $y = Tx$ since $T$ is continuous. Thus we have proved that any sphere $X_{2\varepsilon_0}$ is mapped by $T$ onto a set which contains a sphere $Y_{\eta_0}$. After these preliminaries, let $G$ be a non-void open set in $X$ and $x \in G$. Let $U$ be a neighbourhood of $0$ of $X$ such that $x + U \subseteq G$. Let $V$ be a neighbourhood of $0$ of $Y$ such that $TU \supseteq V$. Then $TG \supseteq T(x + U) = Tx + TU \supseteq Tx + V$. Hence $TG$ contains a neighbourhood of each one of its point. This proves that $T$ maps open set of $X$ onto open sets of $Y$. 
\end{proof}

\begin{theoremn}[Corollary of the Theorem] If a continuous linear operator $T$ on an $F$-space onto an $F$-space gives a one-to-one map, then the inverse operator $T^{-1}$ is also a linear continuous operator. 
\end{theoremn}


\yossection{6}{The Closed Graph Theorem}
\label{sec:II.6}

\begin{definition}\label{def:II.6.Definition.1} Let $X$ and $Y$ be linear topological spaces on the same scalar field. Then the product space $X \times Y$ is a linear space by 
\[
\left\{x _{1}, y _{1} \right\} + \left\{x _{2}, y _{2} \right\} = \left\{x _{1} + x _{2}, y _{1} + y _{2} \right\}, \alpha \{x, y \} = \{\alpha x, \alpha y \}.
\]
\end{definition}

It is also a linear topological space if we call open the sets of the form 
\[
G _{1} \times G _{2} = \left\{\{x, y \}; x \in G _{1}, y \in G _{2} \right\},
\]
where $G_{1}, G_{2}$ are open sets of $X$, $Y$ respectively. If, in particular, $X$ and $Y$ are quasi-normed linear spaces, then $X \times Y$ is also a quasi-normed linear space by the quasi-norm 
\begin{equation}\label{eq:II.6.1}
\left\| \{x, y \} \right\| = \left(\left\| x \right\| ^{2} + \left\| y \right\| ^{2}\right) ^{1/2}.
\end{equation}

\begin{proposition}\label{prop:II.6.Proposition.1} If $X$ and $Y$ are $B$-spaces ( $F$-spaces), then $X \times Y$ is also a $B$-space ( $F$-space). 
\end{proposition}

\begin{proof}
 Clear since $\slim \{x_n, y_n\} = \{x, y\}$ is equivalent to $\slim\limits_{n \to \infty} x_n = x$, $\slim\limits_{n \to \infty} y_n = y$. 
\end{proof}

\begin{definition}\label{def:II.6.Definition.2} The graph $G(T)$ of a linear operator $T$ on $D(T) \subseteq X$ into $Y$ is the set $\{ \{x, Tx\}; x \in D(T) \}$ in the product space $X \times Y$. Let $X, Y$ be linear topological spaces. Then $T$ is called a closed linear operator when its graph $G(T)$ constitutes a closed linear subspace of $X \times Y$. If $X$ and $Y$ are quasi-normed linear spaces, then a linear operator $T$ on $D(T) \subseteq X$ into $Y$ is closed iff the following condition is satisfied: 
\begin{equation}\label{eq:II.6.2}
\begin{aligned}\left\{x _{n} \right\} \subseteq D (T),&\slim _{n \rightarrow \infty} x _{n} = x \quad \text{ and } \quad \slim _{n \rightarrow \infty} T x _{n} = y \quad \text{ imply that }\\&x \in D (T) \quad \text{ and } \quad T x = y.\end{aligned}
\end{equation}
\end{definition}

Thus the notion of a closed linear operator is an extension of the notion of a bounded linear operator. A linear operator $T$ on $D(T) \subseteq X$ into $Y$ is said to be closable or pre-closed if the closure in $X \times Y$ of the graph $G(T)$ is the graph of a linear operator, say $S$, on $D(S) \subseteq X$ into $Y$. 

\begin{proposition}\label{prop:II.6.Proposition.2} If $X, Y$ are quasi-normed linear spaces, then $T$ is closable iff the following condition is satisfied: 
\begin{equation}\label{eq:II.6.3}
\{x _{n} \} \subseteq D (T), \slim _{n \rightarrow \infty} x _{n} = 0 \text{ and } \slim _{n \rightarrow \infty} T x _{n} = y \text{ imply that } y = 0.
\end{equation}
\end{proposition}

\begin{proof}
 The ``only if'' $\mu$ art is clear since the closure in $X \times Y$ of $G(T)$ is a graph $G(S)$ of a linear operator $S$ and so $y = S \cdot 0 = 0$. The ``if'' part is proved as follows. Let us define a linear operator $S$ by the following condition and call $S$ the smallest closed extension of $T$: 

$x \in D(S)$ iff there exists a sequence $\{x_n\} \subseteq D(T)$ such that 
\begin{equation}\slim _{n \to \infty} x _{n} = x \text{ and } \slim _{n \to \infty} T x _{n} = y \text{ exist; and we define } S x = y. \label{eq:II.6.4}\end{equation}

That the value $y$ is defined uniquely by $x$ follows from condition \eqref{eq:II.6.3}. We have only to prove that $S$ is closed. Let $w_{n} \in D(S)$, $\slim\limits_{n \to \infty} w_{n} = w$ and $\slim\limits_{n \to \infty} Sw_{n} = u$. Then there exists a sequence $\{x_{n}\} \subseteq D(T)$ such that $\| w_{n} - x_{n} \| \leqq n^{-1}, \| Sw_{n} - Tx_{n} \| \leqq n^{-1} \quad (n = 1, 2, \ldots)$. Therefore $\slim\limits_{n \to \infty} x_{n} = \slim\limits_{n \to \infty} w_{n} = w$, $\slim\limits_{n \to \infty} Tx_{n} = \slim\limits_{n \to \infty} Sw_{n} = u$, and so $w \in D(S)$, $S \cdot w = u$. 
\end{proof}

An example of a discontinuous but closed operator. Let $X = Y = C[0, 1]$. Let $D$ be the set of all $x(t) \in X$ such that the derivative $x'(t) \in X$; and let $T$ be the linear operator on $D(T) = D$ defined by $Tx = x'$. This $T$ is not continuous, since, for $x_n(t) = t^n$, 
\[
\left\| x _{n} \right\| = 1, \left\| T x _{n} \right\| = \sup\limits _{0 \leqq t \leqq 1} \left| x _{n} ^{\prime} (t) \right| = \sup\limits _{0 \leqq t \leqq 1} \left| n t ^{n-1} \right| = n \quad (n = 1, 2, \dots).
\]

But $T$ is closed. In fact, let $\{x_n\} \subseteq D(T)$, $\slim\limits_{n \to \infty} x_n = x$ and $\slim\limits_{n \to \infty} Tx_n = y$. Then $x_n'(t)$ converges uniformly to $y(t)$, and $x_n(t)$ converges uniformly to $x(t)$. Hence $x(t)$ must be differentiable with continuous derivative $y(t)$. This proves that $x \in D(T)$ and $Tx = y$. 

Examples of closable operators. Let $D_{x}$ be a linear differential operator 
\begin{equation}\label{eq:II.6.5}
D _{x} = \sum\limits_{| j | \leqq k} c _{j} (x) D ^{j}
\end{equation}
with coefficients $c_{j}(x) \in C^{k}(\Omega)$, where $\Omega$ is an open domain of $\mathbb{R}^{n}$. Consider the totality $D$ of functions $f(x) \in L^{2}(\Omega) \cap C^{k}(\Omega)$ such that $D_{x}f(x) \in L^{2}(\Omega)$. We define a linear operator $T$ on $D(T) = D \subseteq L^{2}(\Omega)$ into $L^{2}(\Omega)$ by $(Tf)(x) = D_{x}f(x)$. Then $T$ is closable. For, let $\{f_{h}\} \subseteq D$ be such that $s\text{- }\lim\limits_{h \to \infty} f_{h} = 0$, $s\text{- }\lim\limits_{h \to \infty} D_{x}f_{h} = g$. Then, for any $\varphi(x) \in C_{0}^{k}(\Omega)$, we have, by partial integration, 
\begin{equation}\label{eq:II.6.6}
\int_{\Omega} D _{x} f (x) \cdot \varphi (x) \dd x = \int_{\Omega} f (x) \cdot D _{x} ^{\prime} \varphi (x) \dd x,
\end{equation}
where $D_{x}^{\prime}$ is the differential operator formally adjoint to $D_{x}$: 
\begin{equation}\label{eq:II.6.7}
D _{x} ^{\prime} \varphi (x) = \sum\limits_{| j | \leqq k} (- 1) ^{|j|} D ^{j} \left(c _{j} (x) \varphi (x)\right).
\end{equation}

The formula \eqref{eq:II.6.6} is obtained, since the integrated term in the partial integration vanishes by $\varphi(x) \in C_0^k(\Omega)$. Hence, by the continuity of the scalar product in $L^2(\Omega)$, we obtain, by taking $f = f_h$ and letting $h \to \infty$ in \eqref{eq:II.6.6}, 
\begin{equation}\label{eq:II.6.8}
\int_{\Omega} g (x) \cdot \varphi (x) \dd x = \int_{\Omega} 0 \cdot D _{x} ^{\prime} \varphi (x) \dd x = 0.
\end{equation}

Since $\varphi(x) \in C_0^k(\Omega)$ was arbitrary, we must have $g(x) = 0$ a.e., that is, $g = 0$ in $L^2(\Omega)$. 

\begin{proposition}\label{prop:II.6.Proposition.3} The inverse $T^{-1}$ of a closed linear operator on $D(T) \subseteq X$ into $Y$, if it exists, is a closed linear operator. 
\end{proposition}

\begin{proof}
 The graph of $T^{-1}$ is the set $\left\{\{Tx, x\}; x \in D(T)\right\}$ in the product space $Y \times X$. Thus the Proposition is proved remembering the fact that the mapping $\{x, y\} \to \{y, x\}$ of $X \times Y$ onto $Y \times X$ is a homeomorphic one. 
\end{proof}

We next prove Banach's closed graph theorem: 

\begin{theorem}\label{thm:II.6.Theorem.1} $A$ closed linear operator defined on an $F$-space $X$ into an $F$-space $Y$ is continuous. 
\end{theorem}

\begin{proof}
 The graph $G(T)$ of $T$ is a closed linear subspace of the $F$-space $X \times Y$. Hence, by the completeness of $X \times Y$, $G(T)$ is an $F$-space. The linear mapping $U$ defined by $U\{x, Tx\} = x$ is a one-to-one, continuous linear transformation on the $F$-space $G(T)$ onto the $F$-space $X$. Hence, by virtue of the open mapping theorem, the inverse $U^{-1}$ of $U$ is continuous. The linear operator $V$ defined by $V\{x, Tx\} = Tx$ is clearly continuous on $G(T)$ onto $R(T) \subseteq Y$. Therefore, $T = VU^{-1}$ is continuous on $X$ into $Y$. 
\end{proof}

The following theorem of comparison of two linear operators is due to L. \textsc{Hörmander}: 

\begin{theorem}\label{thm:II.6.Theorem.2} Let $X_{i} (i = 0,1,2; X_{0} = X)$ be $B$-spaces and $T_{i} (i = 1,2)$ be linear operators defined on $D(T_{i}) \subseteq X$ into $X_{i}$. Then, if $T_{1}$ is closed and $T_{2}$ closable in such a way that $D(T_{1}) \subseteq D(T_{2})$, there exists a constant $C$ such that 
\begin{equation}\label{eq:II.6.9}
\left\| T _{2} x \right\| \leqq C \left(\left\| T _{1} x \right\| ^{2} + \left\| x \right\| ^{2}\right) ^{1/2} \text{ for all } x \in D (T _{1}).
\end{equation}
\end{theorem}

\begin{proof}
 The graph $G(T_{1})$ of $T_{1}$ is a closed subspace of $X \times X_{1}$. Hence the mapping 
\begin{equation}\label{eq:II.6.10}
G \left(T _{1}\right) \ni \{x, T _{1} x \} \rightarrow T _{2} x \in X _{2}
\end{equation}
is a linear operator on the $B$-space $G(T_1)$ into the $B$-space $X_2$. We shall prove that this mapping is closed. Suppose that $\{x_n, T_1 x_n\}$ $s$-converges in $G(T_1)$ and that $T_2 x_n$ $s$-converges in $X_2$. Since $T_1$ is closed, there is an element $x \in D(T_1)$ such that $x = \slim\limits_{n \to \infty} x_n$ and $T_1 x = \slim\limits_{n \to \infty} T_1 x_n$. By the hypothesis, we have $x \in D(T_2)$, and, since $T_2$ is closable, the existing $\slim\limits_{n \to \infty} T_2 x_n$ can only be $T_2 x$. Hence the mapping \eqref{eq:II.6.10} is closed, and so, by the closed graph theorem, it must be continuous. This proves \eqref{eq:II.6.9}. 
\end{proof}


\yossection{7}{An Application of the Closed Graph Theorem (Hörmander's Theorem)}
\label{sec:II.7}

Any distribution solution $u \in L^2$ of the Laplace equation 
\[
\Delta u = f \in L ^{2}
\]
is defined by a function which is equal to a $C^{\infty}$ function after correction on a set of measure zero in the domain where $f$ is $C^{\infty}$. This result is known as Weyl's Lemma and plays a fundamental role in the modern treatment of the theory of potentials. See H. \textsc{Weyl} \cite{weyl336}. There is an abundant literature on the extensions of Weyl's Lemma. Of these, the research due to L. \textsc{Hörmander} \cite{hormander113} seems to be the most far reaching. We shall begin with his definition of hypoellipticity. 

\begin{definition}\label{def:II.7.Definition.1} Let $\Omega$ be an open domain of $\mathbb{R}^{n}$. A function $u(x)$, $x \in \Omega$, is said to belong to $L_{\mathrm{loc}}^2(\Omega)$ if $\int_{\Omega'} |u(x)|^2 \dd x < \infty$ for any open subdomain $\Omega'$ with compact closure in $\Omega$. A linear partial differential operator $P(D)$ with constant coefficients: 
\begin{equation}\label{eq:II.7.1}
\begin{yoslines}& P (D) = P \left(\frac {1}{i} \frac {\partial}{\partial x _{1}}, \frac {1}{i} \frac {\partial}{\partial x _{2}}, \dots , \frac {1}{i} \frac {\partial}{\partial x _{n}}\right), \text{ where }\\& P (\xi) = P (\xi_{1}, \xi_{2}, \dots , \xi_{n}) \text{ is a polynomial in } \xi_{1}, \xi_{2}, \dots , \xi_{n},\end{yoslines}
\end{equation}
is said to be hypoelliptic, if every distribution solution $u \in L_{\mathrm{loc}}^{2}(\Omega)$ of $P(D)u = f$ is $C^{\infty}$ after correction on a set of measure zero in the subdomain where $f$ is $C^{\infty}$. 
\end{definition}

\begin{theoremn}[\textsc{Hörmander}] If $P(D)$ is hypoelliptic, then there exists, for any large positive constant $C_{1}$, a positive constant $C_{2}$ such that, for any solution $\zeta = \xi + i\eta$ of the algebraic equation $P(\zeta) = 0$, we have
\begin{equation}\label{eq:II.7.2}
| \zeta | = \left(\sum\limits_{j=1} ^{n} | \zeta_{j} | ^{2}\right) ^{1/2} \leqq C _{1} \quad \text{ if } \quad | \eta | = \left(\sum\limits_{j=1} ^{n} | \eta_{j} | ^{2}\right) ^{1/2} \leqq C _{2}.
\end{equation}

\end{theoremn}
\begin{proof}
 Let $U$ be the totality of distribution solutions $u \in L^2(\Omega')$ of $P(D)u = 0$, that is, the totality of $u \in L^2(\Omega')$ such that 
\begin{equation}\label{eq:II.7.3}
\int_ {\Omega^{\prime}} u \cdot P ^{\prime} (D) \varphi \dd x = 0 \quad \text{ for all } \quad \varphi \in C _{0} ^{\infty} (\Omega^{\prime}),
\end{equation}
where the adjoint differential operator $P'(D)$ of $P(D)$ is defined by 
\begin{equation}\label{eq:II.7.4}
P ^{\prime} (\xi) = P (- \xi_{1}, - \xi_{2}, \dots , - \xi_{n}). \end{equation}

We can prove that $U$ is a closed linear subspace of $L^{2}(\Omega')$. The linearity of $U$ is obvious from the linearity of the differential operator $P(D)$. Let a sequence $\{u_{h}\}$ of $U$ be such that $\slim\limits_{h\to\infty}u_{h}=u$ in $L^{2}(\Omega')$. Then, by the continuity of the scalar product in $L^{2}(\Omega')$, we have 
\[
0 = \lim\limits _{h \rightarrow \infty} \int_ {\Omega^{\prime}} u _{h} \cdot P ^{\prime} (D) \varphi \dd x = \int_ {\Omega^{\prime}} u \cdot P ^{\prime} (D) \varphi \dd x = 0, \text{ i.e., } u \in U.
\]

Thus $U$ is a closed linear subspace of $L^{2}(\Omega')$, and as such a $B$-space. Since $P(D)$ is hypoelliptic, we may suppose that every $u \in U$ is $C^{\infty}$ in $\Omega'$. Let $\Omega_{1}'$ be any open subdomain with compact closure in $\Omega'$. Then, for any $u \in U$, the function $\partial u / \partial x_{k}$ is $C^{\infty}$ in $\Omega'$ ( $k = 1, 2, \ldots, n$). By the argument used in the preceding section, the mapping 
\[
U \ni u \rightarrow \frac {\partial u}{\partial x _{k}} \in L ^{2} (\Omega_{1} ^{\prime}) \quad (k = 1, 2, \dots , n)
\]
is a closed linear operator. Hence, by the closed graph theorem, there exists a positive constant $C$ such that 
\[
\int_ {\Omega_{1} ^{*}} \sum\limits_{k=1} ^{n} \left| \frac {\partial u}{\partial x _{k}} \right| ^{2} \dd x \leqq C \int_ {\Omega^{*}} | u | ^{2} \dd x, \quad \text{ for all } u \in U.
\]

If we apply this inequality to the function $u(x) = e^{i\langle x,\zeta \rangle}$, where $\zeta = \xi + i\eta = (\xi_1 + i\eta_1, \xi_2 + i\eta_2, \ldots, \xi_n + i\eta_n)$ is a solution of $P(\zeta) = 0$ and $\langle x, \zeta \rangle = \sum\limits_{j=1}^{n} \langle x_j, \zeta_j \rangle$, we obtain 
\[
\sum\limits_{k=1} ^{n} \left| \zeta_{k} \right| ^{2} \cdot \int_ {\Omega_{1} ^{\prime}} e ^{-2 \langle x, \eta \rangle} \dd x \leqq C \int_ {\Omega^{\prime}} e ^{-2 \langle x, \eta \rangle} \dd x.
\]

Therefore, when $|\eta|$ is bounded, it follows that $|\zeta|$ must be bounded. 
\end{proof}

\begin{remarku} Later on, we shall prove that condition \eqref{eq:II.7.2} implies the hypoellipticity of $P(D)$. This result is also due to \textsc{Hörmander}. Thus, in particular, we see that Weyl's Lemma is a trivial consequence of \textsc{Hörmander}'s result. In fact, the root of the algebraic equation $-\sum\limits_{j=1}^{n}\zeta_j^2 = 0$ satisfies \eqref{eq:II.7.2}. 
\end{remarku}

\yossub{References for Chapter II}

S. \textsc{Banach} \cite{banach9}, N. \textsc{Bourbaki} \cite{bourbaki27}, N. \textsc{Dunford}-J. \textsc{Schwartz} \cite{dunford44}, E. \textsc{Hille}-R. S. \textsc{Phillips} \cite{hille97} and L. \textsc{Hörmander} \cite{hormander118}. 
   % II
% ==== III ====
\yoschapter{III}{The Orthogonal Projection and F. \textsc{Riesz}' Representation Theorem}
\label{chap:III}


\yossection{1}{The Orthogonal Projection}
\label{sec:III.1}

In a pre-Hilbert space, we can introduce the notion of orthogonality of two vectors. Thanks to this fact, a Hilbert space may be identified with its dual space, i.e., the space of bounded linear functionals. This result is the representation theorem of F. \textsc{Riesz} \cite{rickart272}, and the whole theory of Hilbert spaces is founded on this theorem. 

\begin{definition}\label{def:III.1.Definition.1} Let $x, y$ be vectors of a pre-Hilbert space $X$. We say that $x$ is orthogonal to $y$ and write $x \perp y$ if $(x, y) = 0$; if $x \perp y$ then $y \perp x$, and $x \perp x$ iff $x = 0$. Let $M$ be a subset of a pre-Hilbert space $X$. We denote by $M^{\perp}$ the totality of vectors $\in X$ orthogonal to every vector $m$ of $M$. 
\end{definition}

\begin{theorem}\label{thm:III.1.Theorem.1} Let $M$ be a closed linear subspace of a Hilbert space $X$. Then $M^{\perp}$ is also a closed linear subspace of $X$, and $M^{\perp}$ is called the orthogonal complement of $M$. Any vector $x \in X$ can be decomposed uniquely in the form 
\begin{equation}\label{eq:III.1.1}
x = m + n, \text{ where } m \in M \text{ and } n \in M ^{\perp}.
\end{equation}
\end{theorem}

The element $m$ in \eqref{eq:III.1.1} is called the orthogonal projection of $x$ upon $M$ and will be denoted by $P_M x; P_M$ is called the projection operator or the projector upon $M$. We have thus, remembering that $M \subseteq (M^\perp)^\perp$, 
\begin{equation}\tag{1'}\label{eq:III.1.1'}
x = P _{M} x + P _{M \perp} x, \text{ that is, } I = P _{M} + P _{M \perp}.
\end{equation}

\begin{proof}
 The linearity of $M^{\perp}$ is a consequence of the linearity in $x$ of the scalar product $(x, y)$. $M^{\perp}$ is closed by virtue of the continuity of the scalar product. The uniqueness of the decomposition \eqref{eq:III.1.1} is clear since a vector orthogonal to itself is the zero vector. 

To prove the possibility of the decomposition \eqref{eq:III.1.1}, we may assume that $M \neq X$ and $x \notin M$, for, if $x \in M$, we have the trivial decomposition with $m = x$ and $n = 0$. Thus, since $M$ is closed and $x \notin M$, we have 
\[
d = \inf\limits _{m \in M} \| x - m \| > 0.
\]

Let $\{m_n\} \subseteq M$ be a minimizing sequence, i.e., $\lim\limits_{n\to \infty}\|x - m_n\| = d$. Then $\{m_n\}$ is a Cauchy sequence. For, by $\|a + b\|^2 + \|a - b\|^2 = 2(\|a\|^2 + \|b\|^2)$ valid in any pre-Hilbert space (see \eqref{eq:I.5.1} in Chapter~\ref{chap:I}, \ref{sec:I.5}), we obtain 
\[
\begin{aligned}
&{\| m _{k} - m _{n} \| ^{2} = \| (x - m _{n}) - (x - m _{k}) \| ^{2} = 2 (\| x - m _{n} \| ^{2} + \| x - m _{k} \| ^{2})} \\
&{\qquad - \| 2 x - m _{n} - m _{k} \| ^{2}} {\qquad = 2 (\| x - m _{n} \| ^{2} + \| x - m _{k} \| ^{2}) - 4 \| x - (m _{n} + m _{k}) / 2 \| ^{2}} \\
&{\qquad \leqq 2 (\| x - m _{n} \| ^{2} + \| x - m _{k} \| ^{2} - 4 d ^{2})} {\qquad (\text{ since } (m _{n} + m _{k}) / 2 \in M)} \\
&{\qquad \rightarrow 2 (d ^{2} + d ^{2}) - 4 d ^{2} = 0 \text{ as } k, n \rightarrow \infty .}
\end{aligned}
\]

By the completeness of the Hilbert space $X$, there exists an element $m \in X$ such that $\slim\limits_{n \to \infty} m_n = m$. We have $m \in M$ since $M$ is closed. Also, by the continuity of the norm, we have $\|x - m\| = d$. 

Write $x = m + (x - m)$. Putting $n = x - m$, we have to show that $n \in M^{\perp}$. For any $m' \in M$ and any real number $\alpha$, we have $(m + \alpha m') \in M$ 
and so 
\[
\begin{aligned}
&d ^{2}  \leqq \| x - m - \alpha m ^{\prime} \| ^{2} = (n - \alpha m ^{\prime}, n - \alpha m ^{\prime}) = \| n \| ^{2} - \alpha (n, m ^{\prime}) - \alpha (m ^{\prime}, n) + \alpha^{2} \| m ^{\prime} \| ^{2}.
\end{aligned}
\]

This gives, since $\| n \| = d$, $0 \leqq -2\alpha \operatorname{Re}(n, m') + \alpha^2 \| m' \|^2$ for every real $\alpha$. Hence $\operatorname{Re}(n, m') = 0$ for every $m' \in M$. Replacing $m'$ by $im'$, we obtain $\operatorname{Im}(n, m') = 0$ and so $(n, m') = 0$ for every $m' \in M$. 
\end{proof}

\begin{corollaryu} For a closed linear subspace $M$ of a Hilbert space $X$, we have $M = M^{\perp \perp} = (M^{\perp})^{\perp}$. 
\end{corollaryu}

\begin{theorem}\label{thm:III.1.Theorem.2} The projector $P = P_{M}$ is a bounded linear operator such that 
\begin{align}
P& = P ^{2} \text{(idempotent property of } P), \label{eq:III.1.2} \\
(P x, y)& = (x, P y) \text{ (symmetric property of } P). \label{eq:III.1.3}
\end{align}
\end{theorem}

Conversely, a bounded linear operator $P$ on a Hilbert space $X$ into $X$ satisfying \eqref{eq:III.1.2} and \eqref{eq:III.1.3} is a projector upon $M = R(P)$. 

\begin{proof}
 \eqref{eq:III.1.2} is clear from the definition of the orthogonal projection. We have, by \eqref{eq:III.1.1} and $P_M x \perp P_M \perp y$, 
\[ \begin{aligned}
&(P _{M} x, y)  = (P _{M} x, P _{M} \bot y + P _{M} y) = (P _{M} x, P _{M} y) = (P _{M} x + P _{M} \bot x, P _{M} y) = (x, P _{M} y).
\end{aligned}
\]

Next let $y = x + z$, $x \in M$, $z \in M^{\perp}$ and $w = u + v$, $u \in M$, $v \in M^{\perp}$, then $y + w = (x + u) + (z + v)$ with $(x + u) \in M$, $(z + v) \in M^{\perp}$ and so, by the uniqueness of the decomposition \eqref{eq:III.1.1}, $P_M(y + w) = P_My + P_Mw$; similarly we obtain $P_M(xy) = \alpha P_My$. The boundedness of the operator $P_M$ is proved by 
\[
\begin{aligned} & {\| x \| ^{2} = \| P _{M} x + P _{M \perp} x \| ^{2} = (P _{M} x + P _{M \perp} x, P _{M} x + P _{M \perp} x)} \\ & {\quad = \| P _{M} x \| ^{2} + \| P _{M \perp} x \| ^{2} \geqq \| P _{M} x \| ^{2}.} \end{aligned}
\]

Thus, in particular, we have 
\begin{equation}\label{eq:III.1.4}
\left\| P _{M} \right\| \leqq 1.
\end{equation}

The converse part of the Theorem is proved as follows. The set $M = R(P)$ is a linear subspace, since $P$ is a linear operator. The condition $x \in M$ is equivalent to the existence of a certain $y \in X$ such that $x = P \cdot y$, and this in turn is equivalent, by \eqref{eq:III.1.2}, to $x = Py = P^2y = Px$. Therefore $x \in M$ is equivalent to $x = Px$. $M$ is a closed subspace; for, $x_n \in M$, $\slim x_n = y$ imply, by the continuity of $P$ and $x_n = Px_n$, $\slim x_n = \slim Px_n = Py$ so that $y = Py$. 

We have to show that $P = P_{M}$. If $x \in M$, we have $Px = x = P_{M} \cdot x$; and if $v \in M^{\perp}$, we have $P_{M}y = 0$. Moreover, in the latter case, $(Py, Py) = 6^{*}$ 
$(y, P^2 y) = (y, Py) = 0$ and so $Py = 0$. Therefore we obtain, for any $y \in X$, 
\[
\begin{aligned}
&P y  = P (P _{M} y + P _{M \perp} y) = P P _{M} y + P P _{M \perp} y = P _{M} y + 0, \text{ i.e., } P y = P _{M} y.
\end{aligned}
\]
\end{proof}

Another characterization of the projection operator is given by 

\begin{theorem}\label{thm:III.1.Theorem.3} $A$ bounded linear operator $P$ on a Hilbert space $X$ into $X$ is a projector iff $P$ satisfies $P = P^2$ and $\| P \| \leqq 1$. 
\end{theorem}

\begin{proof}
 We have only to prove the ``if'' part. Set $M = R(P)$ and $N = N(P) = \{y; P y = 0\}$. As in the proof of the preceding Theorem~\ref{thm:III.1.Theorem.2}, $M$ is a closed linear subspace and $x \in M$ is equivalent to $x = Px$. $N$ is also a closed linear subspace in virtue of the continuity of $P$. In the decomposition $x = Px + (I - P)x$, we have $Px \in M$ and $(I - P)x \in N$. The latter assertion is clear from $P(I - P) = P - P^{2} = 0$. 

We thus have to prove that $N = M^{\perp}$. For every $x \in X$, $y = Px - x \in N$ by $P^2 = P$. Hence, if, in particular, $x \in N^{\perp}$, then $Px = x + y$ with $(x, y) = 0$. It follows then that $\|x\|^2 \geqq \|Px\|^2 = \|x\|^2 + \|y\|^2$ so that $y = 0$. Thus we have proved that $x \in N^{\perp}$ implies $x = Px$, that is, $N^{\perp} \subseteq M = R(P)$. Let, conversely, $z \in M = R(P)$, so that $z = Pz$. Then we have the orthogonal decomposition $z = y + x$, $y \in N$, $x \in N^{\perp}$, and so $z = Pz = Py + Px = Px = x$, the last equality being already proved. This shows that $M = R(P) \subseteq N^{\perp}$. We have thus obtained $M = N^{\perp}$, and so, by $N = (N^{\perp})^{\perp}$, $N = M^{\perp}$. 
\end{proof}


\yossection{2}{``Nearly Orthogonal'' Elements}
\label{sec:III.2}

In general, we cannot define the notion of orthogonality in a normed linear space. However, we can prove the 

\begin{theoremn}[F. \textsc{Riesz} {\cite{rickart273}}] Let $X$ be a normed linear space, and $M$ be a closed linear subspace. Suppose $M \neq X$. Then there exists, for any $\varepsilon > 0$ with $0 < \varepsilon < 1$, an $x_{\varepsilon} \in X$ such that
\begin{equation}\label{eq:III.2.1}
\left\| x _{\varepsilon} \right\| = 1 \quad \text{ and } \quad \dist (x _{\varepsilon}, M) = \inf\limits _{m \in M} \left\| x _{\varepsilon} - m \right\| \geqq 1 - \varepsilon .
\end{equation}

The element $x_{\varepsilon}$ is thus ``nearly orthogonal'' to $M$. 

\end{theoremn}
\begin{proof}
 Let $y \in X - M$. Since $M$ is closed, $\dist(y, M) = \inf\limits_{m \in M} \| y - m \| = \alpha > 0$. Thus there exists an $m_{\varepsilon} \in M$ such that $\| y - m_{\varepsilon} \| \leqq \alpha \left(1 + \frac{\varepsilon}{1 - \varepsilon}\right)$. The vector $x_{\varepsilon} = (y - m_{\varepsilon}) / \| y - m_{\varepsilon} \|$ satisfies $\| x_{\varepsilon} \| = 1$ and 
\[ \begin{aligned}
&{\| x _{\varepsilon} - m \| = \| y - m _{\varepsilon} \| ^{-1} \| y - m _{\varepsilon} - \| y - m _{\varepsilon} \| \cdot m \| \geqq \| y - m _{\varepsilon} \| ^{-1} \alpha} {\quad \geqq \left(\frac {1}{1 - \varepsilon}\right) ^{-1} = 1 - \varepsilon .}
\end{aligned}
\]
\end{proof}

\begin{corollary}\label{cor:III.2.Corollary.1} Let there exist a sequence of closed linear subspaces $M_{n}$ of a normed linear space $X$ such that $M_{n} \subseteq M_{n+1}$ and $M_{n} \neq M_{n+1}$ ( $n = 1, 2, \ldots$). Then there exists a sequence $\{y_{n}\}$ such that 
\begin{equation}\label{eq:III.2.2}
y _{n} \in M _{n}, \| y _{n} \| = 1 \quad \text{ and } \quad \dist (y _{n+1}, M _{n}) \geqq 1 / 2 (n = 1, 2, \dots).
\end{equation}
\end{corollary}

\begin{corollary}\label{cor:III.2.Corollary.2} The unit sphere $S = \{x \in X; \| x \| \leqq 1\}$ of a $B$-space $X$ is compact iff $X$ is of finite dimension. 
\end{corollary}

\begin{proof}
 If $x_{1}, x_{2}, \ldots, x_{n}$ be a base for $X$, then the map $(\alpha_{1}, \alpha_{2}, \ldots, \alpha_{n}) \to \sum\limits_{j=1}^{n} \alpha_{j} x_{j}$ of $\mathbb{R}^{n}$ onto $X$ is surely continuous and so it is open by the open mapping theorem. This proves the ``if'' part. The ``only if'' part is proved as follows. Suppose $X$ is not of finite dimension. Then there exists, by the preceding Corollary~\ref{cor:III.2.Corollary.1}, a sequence $\{y_{n}\}$ satisfying the conditions: $\|y_{n}\|=1$ and $\|y_{m}-y_{n}\|\geqq1/2$ for $m>n$. This is a contradiction to the hypothesis that the unit sphere of $X$ is compact. 
\end{proof}


\yossection{3}{The Ascoli-Arzelà Theorem}
\label{sec:III.3}

To give an example of a relatively compact infinite subset of a $B$-space of infinite dimension, we shall prove the 

\begin{theoremn}[\textsc{Ascoli}-\textsc{Arzelà}] Let $S$ be a compact metric space, and $C(S)$ the $B$-space of (real- or) complex-valued continuous functions $x(s)$ normed by $\|x\| = \sup\limits_{s \in S} |x(s)|$. Then a sequence $\{x_n(s)\} \subseteq C(S)$ is relatively compact in $C(S)$ if the following two conditions are satisfied:
\begin{equation}\label{eq:III.3.1}
x _{n} (s) \text{ is equi-bounded (in } n), \text{ i.e., } \sup\limits _{n \geqq 1} \sup\limits _{s \in S} | x _{n} (s) | < \infty ,
\end{equation}
$x_{n}(s)$ is equi-continuous (in $n$), i.e., 
\begin{equation}\label{eq:III.3.2}
\lim\limits _{\delta \downarrow 0} \sup\limits _ {n \geqq 1, \dist (s ^{\prime}, s ^{\prime \prime}) \leqq \delta} | x _{n} (s ^{\prime}) - x _{n} (s ^{\prime \prime}) | = 0.
\end{equation}

\end{theoremn}
\begin{proof}
 A bounded sequence of complex numbers contains a convergent subsequence (the Bolzano-Weierstrass theorem). Hence, for fixed $s$, the sequence $\{x_{n}(s)\}$ contains a convergent subsequence. On the other hand, since the metric space $S$ is compact, there exists a countable dense subset $\{s_{n}\} \subseteq S$ such that, for every $\varepsilon > 0$, there exists a finite subset $\{s_{n}; 1 \leqq n \leqq k(\varepsilon)\}$ of $\{s_{n}\}$ satisfying the condition 
\begin{equation}\label{eq:III.3.3}
\sup\limits _{s \in S} \inf\limits _{1 \leqq j \leqq k (\varepsilon)} \dist (s, s _{j}) \leqq \varepsilon . \end{equation}

The proof of this fact is obtained as follows. Since $S$ is compact, it is totally bounded (see Chapter~\ref{chap:0}, \ref{sec:0.2}). Thus there exists, for any $\delta > 0$, a finite system of points $\in S$ such that any point of $S$ has a distance $\leqq \delta$ from some point of the system. Letting $\delta = 1, 2^{-1}, 3^{-1}, \ldots$ and collecting the corresponding finite systems of points, we obtain a sequence $\{s_n\}$ with the stated properties. 

We then apply the diagonal process of choice to the sequence $\{x_{n}(s)\}$, so that we obtain a subsequence $\{x_{n'}(s)\}$ of $\{x_{n}(s)\}$ which converges for $s = s_1, s_2, \ldots, s_k, \ldots$ simultaneously. By the equi-continuity of $\{x_{n}(s)\}$, there exists, for every $\varepsilon > 0$, a $\delta = \delta(\varepsilon) > 0$ such that $\text{ dis }(s', s'') \leqq \delta$ implies $|x_{n}(s') - x_{n}(s'')| \leqq \varepsilon$ for $n = 1, 2, \ldots$. Hence, for every $s \in S$, there exists a $j$ with $j \leqq k(\varepsilon)$ such that 
\begin{align}
\left| x _ {n ^{\prime}} (s) - x _ {m ^{\prime}} (s) \right|& \leqq \left| x _ {n ^{\prime}} (s) - x _ {n ^{\prime}} (s _{j}) \right| + \left| x _ {n ^{\prime}} (s _{j}) - x _ {m ^{\prime}} (s _{j}) \right| \notag \\
+ \left| x _ {m ^{\prime}} (s _{j}) - x _ {m ^{\prime}} (s) \right|& \leqq 2 \varepsilon + \left| x _ {n ^{\prime}} (s _{j}) - x _ {m ^{\prime}} (s _{j}) \right|. \notag
\end{align}

Thus $\lim\limits_{n,m\to \infty}\max\limits_s|x_{n'}(s) - x_{m'}(s)|\leqq 2\varepsilon$ and so $\lim\limits_{n,m\to \infty}\|x_{n'} - x_{m'}\| = 0.$ 
\end{proof}


\yossection{4}{The Orthogonal Base. Bessel's Inequality and Parseval's Relation}
\label{sec:III.4}

\begin{definition}\label{def:III.4.Definition.1} $A$ set $S$ of vectors in a pre-Hilbert space $X$ is called an orthogonal set, if $x \perp y$ for each pair of distinct vectors $x, y$ of $S$. If, in addition, $\|x\| = 1$ for each $x \in S$, then the set $S$ is called an orthonormal set. An orthonormal set $S$ of a Hilbert space $X$ is called a complete orthonormal system or an orthogonal base of $X$, if no orthonormal set of $X$ contains $S$ as a proper subset. 
\end{definition}

\begin{theorem}\label{thm:III.4.Theorem.1} $A$ Hilbert space $X$ (having a non-zero vector) has at least one complete orthonormal system. Moreover, if $S$ is any orthonormal set in $X$, there is a complete orthonormal system containing $S$ as a subset. 
\end{theorem}

Proof (by Zorn's Lemma). Let $S$ be an orthonormal set in $X$. Such a set surely exists; for instance, if $x \neq 0$, the set consisting only of $x / \|x\|$ is orthonormal. We consider the totality $\{S\}$ of orthonormal sets which contain $S$ as a subset. $\{S\}$ is partially ordered by writing $S_1 \prec S_2$ for the inclusion relation $S_1 \subseteq S_2$. Let $\{S'\}$ be a linearly ordered subsystem of $\{S\}$, then $\bigcup\limits_{S' \in \{S'\}} S'$ is an orthonormal set and an upper bound of $\{S'\}$. Thus, by Zorn's Lemma, there exists a maximal element $S_0$ of $\{S\}$. This orthonormal set $S_0$ contains $S$ and, by the maximality, it must be a complete orthonormal system. 

\begin{theorem}\label{thm:III.4.Theorem.2} Let $S = \{x_{\alpha}; \alpha \in A\}$ be a complete orthonormal system of a Hilbert space $X$. For any $f \in X$, we define its Fourier coefficients (with respect to $S$) 
\begin{equation}\label{eq:III.4.1}
f _{\alpha} = (f, x _{\alpha}).
\end{equation}
\end{theorem}

Then we have Parseval's relation 
\begin{equation}\label{eq:III.4.2}
\left\| f \right\| ^{2} = \sum\limits_{\alpha \in A} \left| f _{\alpha} \right| ^{2}.
\end{equation}

\begin{proof}
 We shall first prove Bessel's inequality 
\begin{equation}\tag{2'}\label{eq:III.4.2'}
\sum\limits_{\alpha \in A} | f _{\alpha} | ^{2} \leqq \| f \| ^{2}. \end{equation}

Let $\alpha_{1},\alpha_{2},\ldots ,\alpha_{n}$ be any finite system of $\alpha$'s. For any finite system of complex numbers $c_{\alpha_1},c_{\alpha_2},\ldots ,c_{\alpha_n}$, we have, by the orthonormality of $\{x_{\alpha}\}$, 
\begin{equation}\label{eq:III.4.3}
\begin{aligned}
&\left\| f - \sum\limits_{j=1} ^{n} c _ {\alpha_{j}} x _ {\alpha_{j}} \right\| ^{2}  = \left(f - \sum\limits_{j=1} ^{n} c _ {\alpha_{j}} x _ {\alpha_{j}}, f - \sum\limits_{j=1} ^{n} c _ {\alpha_{j}} x _ {\alpha_{j}}\right) \\
&= \| f \| ^{2} - \sum\limits_{j=1} ^{n} c _ {\alpha_{j}} \bar {f} _ {\alpha_{j}} - \sum\limits_{j=1} ^{n} \bar {c} _ {\alpha_{j}} f _ {\alpha_{j}} + \sum\limits_{j=1} ^{n} | c _ {\alpha_{j}} | ^{2} = \| f \| ^{2} - \sum\limits_{j=1} ^{n} | f _ {\alpha_{j}} | ^{2} + \sum\limits_{j=1} ^{n} | f _ {\alpha_{j}} - c _ {\alpha_{j}} | ^{2}.
\end{aligned}
\end{equation}

Hence the minimum of $\left\| f - \sum\limits_{j=1}^{n} c_{\alpha_j} x_{\alpha_j}\right\|^2$, for fixed $\alpha_1, \alpha_2, \ldots, \alpha_n$, is attained when $c_{\alpha_j} = f_{\alpha_j}$ ( $j = 1, 2, \ldots, n$). We have thus 
\begin{equation}\label{eq:III.4.4}
\left\| f - \sum\limits_{j=1} ^{n} f _ {\alpha_{j}} x _ {\alpha_{j}} \right\| ^{2} = \| f \| ^{2} - \sum\limits_{j=1} ^{n} | f _ {\alpha_{j}} | ^{2}, \text{ and hence } \sum\limits_{j=1} ^{n} | f _ {\alpha_{j}} | ^{2} \leqq \| f \| ^{2}.
\end{equation}

By the arbitrariness of $\alpha_{1},\alpha_{2},\ldots ,\alpha_{n}$, we see that Bessel's inequality \eqref{eq:III.4.2} is true, and $f_{\alpha}\neq 0$ for at most a countable number of $\alpha$'s, say $\alpha_{1},\alpha_{2},\ldots ,\alpha_{n},\ldots$ We then prove that $f = \slim\limits_{n\to \infty}\sum\limits_{j=1}^{n}f_{\alpha_j}x_{\alpha_j}$. First, the sequence $\left\{\sum\limits_{j=1}^{n}f_{\alpha_j}x_{\alpha_j}\right\}$ is a Cauchy sequence, since, by the orthonormality of $\{x_{\alpha}\}$, 
\[
\left| \sum\limits_{j=k} ^{n} f _ {\alpha_{j}} x _ {\alpha_{j}} \right| ^{2} = \left(\sum\limits_{j=k} ^{n} f _ {\alpha_{j}} x _ {\alpha_{j}}, \sum\limits_{j=k} ^{n} f _ {\alpha_{j}} x _ {\alpha_{j}}\right) = \sum\limits_{j=k} ^{n} | f _ {\alpha_{j}} | ^{2}
\]
which tends, by \eqref{eq:III.4.4} proved above, to $0$ as $k \to \infty$. We set $\slim\limits_{n \to \infty} \sum\limits_{j=1}^{n} f_{\alpha_j} x_{\alpha_j} = f'$, and shall prove that $(f - f')$ is orthogonal to every vector of $S$. By the continuity of the scalar product, we have 
\[
(f - f ^{\prime}, x _ {\alpha_{j}}) = \lim\limits _{n \rightarrow \infty} \left(f - \sum\limits_{k=1} ^{n} f _ {\alpha_{k}} x _ {\alpha_{k}}, x _ {\alpha_{j}}\right) = f _ {\alpha_{j}} - f _ {\alpha_{j}} = 0,
\]
and, when $\alpha \neq \alpha_{j}$ ( $j = 1, 2, \ldots$), 
\[
(f - f ^{\prime}, x _{\alpha}) = \lim\limits _{n \rightarrow \infty} \left(f - \sum\limits_{k=1} ^{n} f _{\alpha k} x _{\alpha k}, x _{\alpha}\right) = 0 - 0 = 0.
\]

Thus, by the completeness of the orthonormal system $S = \{x_{\alpha}\}$, we must have $(f - f') = 0$. Hence, by \eqref{eq:III.4.4} and the continuity of the norm, we have 
\[
0 = \lim\limits _{n \rightarrow \infty} \left\| f - \sum\limits_{j=1} ^{n} f _ {\alpha_{j}} x _ {\alpha_{j}} \right\| ^{2} = \| f \| ^{2} - \lim\limits _{n \rightarrow \infty} \sum\limits_{j=1} ^{n} | f _ {\alpha_{j}} | ^{2} = \| f \| ^{2} - \sum\limits_{\alpha \in A} | f _{\alpha} | ^{2}.
\]
\end{proof}

\begin{corollary}\label{cor:III.4.Corollary.1} We have the Fourier expansion 
\begin{equation}\label{eq:III.4.5}
f = \sum\limits_{j=1} ^{\infty} f _ {\alpha_{j}} x _ {\alpha_{j}} = \slim _{n \rightarrow \infty} \sum\limits_{j=1} ^{n} f _ {\alpha_{j}} x _ {\alpha_{j}}.
\end{equation}
\end{corollary}

\begin{corollary}\label{cor:III.4.Corollary.2} Let $l^2(A)$ be the space $L^2(A, \mathfrak{B}, m)$ where $m(\{\alpha\}) = 1$ for every point $\alpha$ of $A$. Then the Hilbert space $X$ is isometrically isomorphic to the Hilbert space $l^2(A)$ by the correspondence 
\begin{equation}\label{eq:III.4.6}
X \ni f \leftrightarrow \{f _{\alpha} \} \in l ^{2} (A)
\end{equation}
in the sense that 
\begin{equation}\label{eq:III.4.7}
(f + g) \leftrightarrow \left\{f _{\alpha} + g _{\alpha} \right\}, \beta f \leftrightarrow \left\{\beta f _{\alpha} \right\} \text{ and } \| f \| ^{2} = \| \left\{f _{\alpha} \right\} \| ^{2} = \sum\limits_{\alpha \in A} | f _{\alpha} | ^{2}.
\end{equation}
\end{corollary}

\begin{examplen} $\left\{\frac{1}{\sqrt{2\pi}} e^{int}; n = 0, \pm 1, \pm 2, \ldots\right\}$ is a complete orthonormal system in the Hilbert space $L^2(0, 2\pi)$.

\end{examplen}
\begin{proof}
 We have only to prove the completeness of this system. We have, by \eqref{eq:III.4.3}, 
\[
\begin{aligned}
\left\| f - \sum\limits_{j=-n} ^{n} \frac {1}{\sqrt {2 \pi}} c _{j} e ^{ijt} \right\| ^{2} \geqq \left\| f - \sum\limits_{j=-n} ^{n} \frac {1}{\sqrt {2 \pi}} f _{j} e ^{ijt} \right\| ^{2} &= \| f \| ^{2} - \sum\limits_{j=-n} ^{n} \frac {1}{2 \pi} | f _{j} | ^{2},
\\
\text{ where } f _{j} &= (f, e ^{ijt}). \end{aligned}
\]

If $f \in L^2(0, 2\pi)$ is continuous and with period $2\pi$, then the left hand side of the inequality above may be taken arbitrarily small by virtue of Weierstrass' trigonometric approximation theorem (see Chapter~\ref{chap:0}, \ref{sec:0.2}). Thus the set of all the linear combinations $\sum\limits_j c_j e^{ijt}$ is dense, in the sense of the norm, in the subspace of $L^2(0, 2\pi)$ consisting of all continuous functions with period $2\pi$. Such a subspace is also dense, in the sense of the norm, in the space $L^2(0, 2\pi)$. Therefore, any function $f \in L^2(0, 2\pi)$, orthogonal to all the functions of $\left\{\frac{1}{\sqrt{2\pi}} e^{int}\right\}$ must be a zero vector of $L^2(0, 2\pi)$. This proves that our system of functions $\frac{1}{\sqrt{2\pi}} e^{int}$ is a complete orthonormal system of $L^2(0, 2\pi)$. 
\end{proof}


\yossection{5}{E. \textsc{Schmidt}'s Orthogonalization}
\label{sec:III.5}

\begin{theoremn}[E. \textsc{Schmidt}'s orthogonalization] Given a finite or countably infinite sequence $\{x_{j}\}$ of linearly independent vectors of a pre-Hilbert space $X$. Then we can construct an orthonormal set having the same cardinal number as the set $\{x_{j}\}$ and spanning the same linear subspace as $\{x_{j}\}$.

\end{theoremn}
\begin{proof}
 Certainly $x_{1} \neq 0$. We define $y_{1}, y_{2}, \ldots$ and $u_{1}, u_{2}, \ldots$ recurrently as follows: 
\[
\begin{aligned}
&y_{1} = x_{1},  u_{1} = y_{1}/\|y_{1}\|, y_{2} = x_{2} - (x_{2}, u_{1})u_{1},  u_{2} = y_{2}/\|y_{2}\|, \cdots   \ddots \\
&y_{n+1} = x_{n+1} - \sum\limits_{j=1}^{n}(x_{n+1}, u_{j})u_{j},  u_{n+1} = y_{n+1}/\|y_{n+1}\|, \\
&\cdots   \ddots
\end{aligned}
\]

This process terminates if $\{x_j\}$ is a finite set. Otherwise it continues indefinitely. Observe that $y_n \neq 0$, because $x_1, x_2, \ldots, x_n$ are linearly independent. Thus $u_n$ is well defined. It is clear, by induction, that each $u_n$ is a linear combination of $x_1, x_2, \ldots, x_n$ and that each $x_n$ is a linear combination of $u_1, u_2, \ldots, u_n$. Thus the closed linear subspace spanned by the $u$'s is the same as that spanned by the $x$'s. 


We see, by $\|u_{1}\|=1$, that $y_{2}\perp u_{1}$ and hence $u_{2}\perp u_{1}$. Thus, by $\|u_{1}\|=1$, $y_{3}\perp u_{1}$ and hence $u_{3}\perp u_{1}$. Repeating the argument, we see that $u_{1}$ is orthogonal to $u_{2}, u_{3}, \ldots, u_{n}, \ldots$. Thus, by $\|u_{2}\|=1$, we have $y_{3}\perp u_{2}$ and so $u_{3}\perp u_{2}$. Repeating the argument, we finally see that $u_{k}\perp u_{m}$ whenever $k>m$. Therefore $\{u_{j}\}$ constitutes an orthonormal set.
\end{proof}

\begin{corollaryu} Let a Hilbert space $X$ be separable, i.e., let $X$ have a dense subset which is at most countable. Then $X$ has a complete orthonormal system consisting of an at most countable number of elements. 
\end{corollaryu}

\begin{proof}
 Suppose that an at most countable sequence $\{a_j\}$ of vectors $\in X$ be dense in $X$. Let $x_1$ be the first non-zero element in the sequence $\{a_j\}$, $x_2$ the first $a_i$ which is not in the closed subspace spanned by $x_1$, and $x_n$ the first $a_i$ which is not in the closed subspace spanned by $x_1, x_2, \ldots, x_{n-1}$. It is clear that the $a$'s and the $x$'s span the same closed linear subspace. In fact, this closed linear subspace is the whole space $X$, because the set $\{a_j\}$ is dense in $X$. Applying Schmidt's orthogonalization to $\{x_j\}$, we obtain an orthonormal system $\{u_j\}$ which is countable and spans the whole space $X$. 


This system $\{u_{j}\}$ is complete, since otherwise there would exist a non-zero vector orthogonal to every $u_{j}$ and hence orthogonal to the space $X$ spanned by $u_{j}$'s.
\end{proof}

\begin{examplen} Let $S$ be the interval $(a, b)$, and consider the real Hilbert space $L^{2}(S, \mathfrak{B}, m)$, where $B$ is the set of all Baire subsets in $(a, b)$. If we orthogonalize the set of monomials 
\[
1, s, s ^{2}, s ^{3}, \dots , s ^{n}, \dots ,
\]
we get the so-called Tchebyschev system of polynomials 
\[
P _{0} (s) = \text{ constant }, P _{1} (s), P _{2} (s), P _{3} (s), \dots , P _{n} (s), \dots
\]
which satisfies 
\[
\int_{a} ^{b} P _{i} (s) P _{j} (s) m (\dd s) = \delta_{ij} (= 0 \text{ or } 1 \text{ according as } i \neq j \text{ or } i = j).
\]

In the particular case when $a = -1$, $b = 1$ and $m(\dd s) = \dd s$, we obtain the Legendre polynomials; when $a = -\infty$, $b = \infty$ and $m(\dd s) = e^{-s^2} \dd s$, we obtain the Hermite polynomials and finally when $a = 0$, $b = \infty$ and $m(\dd s) = e^{-s} \dd s$, we obtain the Laguerre polynomials. 
\end{examplen}

It is easy to see that, when $-\infty < a < b < \infty$, the orthonormal system $\{P_{j}(s)\}$ is complete. For, we may follow the proof of the completeness of the trigonometric system (see the Example in the preceding section $4$); we shall appeal to Weierstrass' polynomial approximation theorem, in place of Weierstrass' trigonometric approximation theorem. As to the completeness proof of the Hermite or Laguerre polynomials, we refer the reader to G. \textsc{Szegő} \cite{szego313} or to K. \textsc{Yosida} \cite{yosida344}. 


\yossection{6}{F. \textsc{Riesz}' Representation Theorem}
\label{sec:III.6}

\begin{theoremn}[F. \textsc{Riesz}' representation theorem] Let $X$ be a Hilbert space and $f$ a bounded linear functional on $X$. Then there exists a uniquely determined vector $y_f$ of $X$ such that
\begin{equation}\label{eq:III.6.1}
f (x) = \left(x, y _{f}\right) \text{ for all } x \in X, \text{ and } \| f \| = \| y _{f} \|.
\end{equation}

Conversely, any vector $y \in X$ defines a bounded linear functional $f_y$ on $X$ by 
\begin{equation}\label{eq:III.6.2}
f _{y} (x) = (x, y) \text{ for all } x \in X, \text{ and } \| f _{y} \| = \| y \|.
\end{equation}

\end{theoremn}
\begin{proof}
 The uniqueness of $y_{f}$ is clear, since $(x,z)=0$ for all $x\in X$ implies $z=0$. To prove its existence, consider the null space $N=N(f)=\{x\in X;f(x)=0\}$ of $f$. Since $f$ is continuous and linear, $N$ is a closed linear subspace. The Theorem is trivial in the case when $N=X$; we take in this case, $y_{f}=0$. Suppose $N\neq X$. Then there exists a $y_{0}\neq0$ which belongs to $N^{\perp}$ (see Theorem~\ref{thm:III.1.Theorem.1} in Chapter~\ref{chap:III}, \ref{sec:III.1}). Define 
\begin{equation}\label{eq:III.6.3}
y _{f} = (\overline {{f (y _{0})}} / \| y _{0} \| ^{2}) y _{0}. \end{equation}

We will show that this $y_{f}$ meets the condition of the Theorem. First, if $x \in N$, then $f(x) = (x, y_{f})$ since both sides vanish. Next, if $x$ is of the form $x = \alpha y_{0}$, then we have 
\[
(x, y _{f}) = (\alpha y _{0}, y _{f}) = \left(\alpha y _{0}, \frac {\overline {{f (y _{0})}}}{\| y _{0} \| ^{2}} y _{0}\right) = \alpha f (y _{0}) = f (\alpha y _{0}) = f (x).
\]

Since $f(x)$ and $(x, y_{f})$ are both linear in $x$, the equality $f(x) = (x, y_{f})$, $x \in X$, is proved if we have proved that $X$ is spanned by $N$ and $y_{0}$. To show the last assertion, we write, remembering that $f(y_{f}) \neq 0$, 
\[
x = \left(x - \frac {f (x)}{f (y _{f})} y _{f}\right) + \frac {f (x)}{f (y _{f})} y _{f}.
\]

The first term on the right is an element of $N$, since 
\[
f \left(x - \frac {f (x)}{f \left(y _{f}\right)} y _{f}\right) = f (x) - \frac {f (x)}{f \left(y _{f}\right)} f \left(y _{f}\right) = 0.
\]

We have thus proved the representation $f(x) = (x, y_f)$. 

Therefore, we have 
\[
\left\| f \right\| = \sup\limits _{\left\| z \right\| \leqq 1} \left| f (x) \right| = \sup\limits _{\left\| z \right\| \leqq 1} \left| (x, y _{f}) \right| \leqq \sup\limits _{\left\| z \right\| \leqq 1} \left\| x \right\| \cdot \left\| y _{f} \right\| = \left\| y _{f} \right\|,
\]
and also $\|f\| = \sup\limits_{\|x\|\leqq 1}|f(x)| \geqq |f(y_f / \|y_f\|)| = \left(\frac{y_f}{\|y_f\|}, y_f\right) = \|y_f\|$. 

Hence we have proved the equality $\|f\| = \|y_f\|$. 

Finally, the converse part of the Theorem is clear from $|f_y(x)| = |(x, y)| \leqq \|x\| \cdot \|y\|$. 
\end{proof}

\begin{corollary}\label{cor:III.6.Corollary.1} Let $X$ be a Hilbert space. Then the totality $X'$ of bounded linear functionals on $X$ constitutes also a Hilbert space, and there is a norm-preserving, one-to-one correspondence $f \leftrightarrow y_{f}$ between $X'$ and $X$. By this correspondence, $X'$ may be identified with $X$ as an abstract set; but it is not allowed to identify, by this correspondence, $X'$ with $X$ as linear spaces, since the correspondence $f \leftrightarrow y_{f}$ is conjugate linear: 
\begin{equation}\label{eq:III.6.4}
\left(\alpha_{1} f _{1} + \alpha_{2} f _{2}\right) \leftrightarrow \left(\bar {\alpha} _{1} y _ {f _{1}} + \bar {\alpha} _{2} y _ {f _{2}}\right),
\end{equation}
where $\alpha_{1},\alpha_{2}$ are complex numbers. 
\end{corollary}

\begin{proof}
 It is easily verified that $X'$ is made into a Hilbert space by defining its scalar product through $(f_1, f_2) = (\overline{y_{f_1}, y_{f_2}})$, so that the statement of Corollary~\ref{cor:III.6.Corollary.1} is clear. 
\end{proof}

\begin{corollary}\label{cor:III.6.Corollary.2} Any continuous linear functional $T$ on the Hilbert space $X'$ is thus identified with a uniquely determined element $t$ of $X$ as follows: 
\begin{equation}\label{eq:III.6.5}
T (f) = f (t) \quad \text{ for all } \quad f \in X ^{\prime}.
\end{equation}
\end{corollary}

\begin{proof}
 Clear from the fact that the product of two conjugate linear transformations is a linear transformation. 
\end{proof}

\begin{definitionu} The space $X'$ is called the dual space of $X$. We can thus identify a Hilbert space $X$ with its second dual $X'' = (X')'$ in the above sense. This fact will be referred to as the reflexivity of Hilbert spaces. 
\end{definitionu}

\begin{corollary}\label{cor:III.6.Corollary.3} Let $X$ be a Hilbert space, and $X'$ its dual Hilbert space. Then, for any subset $F$ of $X'$ which is dense in the Hilbert space $X'$, we have 
\begin{equation}\label{eq:III.6.6}
\left\| x _{0} \right\| = \sup\limits _{f \in F, \| f \| \leqq 1} \left| f (x _{0}) \right|, x _{0} \in X.
\end{equation}
\end{corollary}

\begin{proof}
 We may assume that $x_0 \neq 0$, otherwise the formula \eqref{eq:III.6.6} is trivial. We have $(x_0, x_0 / \|x_0\|) = \|x_0\|$, and so there exists a bounded linear functional $f_0$ on $X$ such that $\|f_0\| = 1$, $f_0(x_0) = \|x_0\|$. Since $f(x_0) = (x_0, y_f)$ is continuous in $y_f$, and since the correspondence $f \leftrightarrow y_f$ is norm-preserving we see that \eqref{eq:III.6.6} is true by virtue of the denseness of $F$ in $X'$. 
\end{proof}

\begin{remarku} Hilbert's original definition of the ``Hilbert space'' is the space $(l^2)$. See his paper \cite{hilbert96}. It was J. von \textsc{Neumann} \cite{neumann228} who gave an axiomatic definition (see Chapter~\ref{chap:I}, \ref{sec:I.9}) of the Hilbert space assuming that the space is separable. F. \textsc{Riesz} \cite{rickart272} proved the above representation theorem without assuming the separability of the Hilbert space. In this paper, he stressed that the whole theory of the Hilbert space may be founded upon this representation theorem. 
\end{remarku}


\yossection{7}{The Lax-Milgram Theorem}
\label{sec:III.7}

Of recent years, it has been proved that a variant of F. \textsc{Riesz}' representation theorem, formulated by P. \textsc{Lax} and A. N. \textsc{Milgram} \cite{lax171}, is a useful tool for the discussion of the existence of solutions of linear partial differential equations of elliptic type. 

\begin{theoremn}[\textsc{Lax}-\textsc{Milgram}] Let $X$ be a Hilbert space. Let $B(x, y)$ be a complex-valued functional defined on the product Hilbert space $X \times X$ which satisfies the conditions:

Sesqui-linearity, i.e., 
\begin{equation}\label{eq:III.7.1}
\begin{aligned}
B \left(\alpha_{1} x _{1} + \alpha_{2} x _{2}, y\right) &= \alpha_{1} B \left(x _{1}, y\right) + \alpha_{2} B \left(x _{2}, y\right) \text{ and }
\\
B \left(x, \beta_{1} y _{1} + \beta_{2} y _{2}\right) &= \bar {\beta} _{1} B \left(x, y _{1}\right) + \bar {\beta} _{2} \dot {B} \left(x, y _{2}\right),
\end{aligned}
\end{equation}

Boundedness, i.e., there exists a positive constant $\gamma$ such that 
\begin{equation}\label{eq:III.7.2}
\left| B (x, y) \right| \leqq \gamma \left\| x \right\| \cdot \left\| y \right\|,
\end{equation}

Positivity, i.e., there exists a positive constant $\delta$ such that 
\begin{equation}\label{eq:III.7.3}
B (x, x) \geqq \delta \| x \| ^{2}.
\end{equation}

Then there exists a uniquely determined bounded linear operator $S$ with a bounded linear inverse $S^{-1}$ such that 
\begin{equation}\label{eq:III.7.4}
\begin{minipage}[b]{0.92\linewidth}
$(x,y) = B(x,Sy)$ whenever $x$ and $y\in X$, and $\| S\| \leqq \delta^{-1},\| S^{-1}\| \leqq \gamma$.
\end{minipage}
\end{equation}


\end{theoremn}
\begin{proof}
 Let $D$ be the totality of elements $y \in X$ for which there exists an element $y^*$ such that $(x, y) = B(x, y^*)$ for all $x \in X$. $D$ is not void, because $0 \in D$ with $0^* = 0 \cdot y^*$ is uniquely determined by $y$. For, if $w$ be such that $B(x, w) = 0$ for all $x$, then $w = 0$ by $0 = B(w, w) \geqq \delta \|w\|^2$. By the sesqui-linearity of $(x, y)$ and $B(x, y)$, we obtain a linear operator $S$ with domain $D(S) = D$: $Sy = y^*$. $S$ is continuous and $\|Sy\| \leqq \delta^{-1} \|y\|$, $y \in D(S)$, because 
\[
\delta \left\| S y \right\| ^{2} \leqq B (S y, S y) = (S y, y) \leqq \left\| S y \right\| \cdot \left\| y \right\|. \]

Moreover, $D = D(S)$ is a closed linear subspace of $X$. Proof: if $y_{n} \in D(S)$ and $\slim\limits_{n \to \infty} y_{n} = y_{\infty}$, then, by the continuity of $S$ proved above, $\{Sy_{n}\}$ is a Cauchy sequence and so has a limit $z = \slim Sy_{n}$. By the continuity of the scalar product, we have $\lim\limits_{n \to \infty} (x, y_{n}) = (x, y_{\infty})$. We have also, by \eqref{eq:III.7.2}, $\lim\limits_{n\to \infty}B(x,Sy_n) = B(x,z)$. Thus, by $(x,y_{n}) = B(x,Sy_{n})$, we must have $(x,y_{\infty}) = B(x,z)$ which proves that $y_{\infty}\in D$ and $Sy_{\infty} = z$. 

Therefore the first part of the Theorem, that is, the existence of the operator $S$ is proved if we can show that $D(S) = X$. Suppose $D(S) \neq X$. Then there exists a $w_0 \in X$ such that $w_0 \neq 0$ and $w_0 \in D(S)^\perp$. Consider the linear functional $F(z) = B(z, w_0)$ defined on $X$. $F(z)$ is continuous, since $|F(z)| = |B(z, w_0)| \leqq \gamma \|z\| \cdot \|w_0\|$. Thus, by F. \textsc{Riesz}' representation theorem, there exists a $w_0' \in X$ such that $B(z, w_0) = F(z) = (z, w_0')$ for all $z \in X$. This proves that $w_0' \in D(S)$ and $Sw_0' = w_0$. But, by $\delta \|w_0\|^2 \leqq B(w_0, w_0) = (w_0, w_0') = 0$, we obtain $w_0 = 0$ which is a contradiction. 

The inverse $S^{-1}$ exists. For, $Sy = 0$ implies $(x, y) = B(x, Sy) = 0$ for all $x \in X$ and so $y = 0$. As above, we show that, for every $y \in X$ there exists a $y'$ such that $(z, y') = B(z, y)$ for all $z \in X$. Hence $y = Sy'$ and so $S^{-1}$ is an everywhere defined operator, and, by $|(z, S^{-1}y)| = |B(z, y)| \leqq \gamma \|z\| \cdot \|y\|$, we see that $\|S^{-1}\| \leqq \gamma$. 
\end{proof}

Concrete applications of the Lax-Milgram theorem will be given in later chapters. In the following four sections, we shall give some examples of the direct application of F. \textsc{Riesz}' representation theorem. 


\yossection{8}{A Proof of the Lebesgue-Nikodym Theorem}
\label{sec:III.8}

This theorem reads as follows. 

\begin{theoremn}[\textsc{Lebesgue}-\textsc{Nikodym}] Let $(S, \mathfrak{B}, m)$ be a measure space, and $\nu(B)$ be a $\sigma$-finite, $\sigma$-additive and non-negative measure defined on $\mathfrak{B}$. If $\nu$ is $m$-absolutely continuous, then there exists a non-negative, $m$-measurable function $p(s)$ such that
\begin{equation}\label{eq:III.8.1}
\nu (B) = \int_{B} p (s) m (\dd s) \quad \text{ for all } \quad B \in \mathfrak {B} \text{ with } \nu (B) < \infty .
\end{equation}

Moreover, the ``density'' $p(s)$ of $v(B)$ (with respect to $m(B)$) is uniquely determined in the sense that any two of them are equal $m$-a.e. 
\end{theoremn}

Proof (due to J. von \textsc{Neumann} \cite{neumann229}). It is easy to see that $\varrho(B) = m(B) + \nu(B)$ is a $\sigma$-finite, $\sigma$-additive and non-negative measure defined on $\mathfrak{B}$. Let $\{B_n\}$ be a sequence of sets $\in \mathfrak{B}$ such that $S = \bigcup\limits_{n=1}^{\infty} B_n$, $B_n \subseteq B_{n+1}$ and $\varrho(B_n) < \infty$ for $n = 1, 2, \ldots$. If we can prove the theorem for every $B \subseteq B_n$ (for fixed $n$) and obtain the density $\phi_n(s)$, then the Theorem is true. For, we have only to take $\phi(s)$ as follows: 
\begin{equation}p (s) = p _{1} (s) \text{ for } s \in B _{1}, \text{ and } p (s) = p _{n+1} (s) \text{ for } s \in B _{n+1} - B _{n} (n = 1, 2, \dots). \notag\end{equation}

Therefore we may assume that $\varrho(S) < \infty$ without losing the generality. Consider the Hilbert space $L^2(S, \mathfrak{B}, \varrho)$. Then 
\[
f (x) = \int_{S} x (s) \nu (\dd s), x \in L ^{2} (S, \mathfrak {B}, \varrho),
\]
gives a bounded linear functional on $L^{2}(S, \mathfrak{B}, \varrho)$, because 
\[
\begin{aligned}
&| f (x) |  \leqq \int_{S} | x (s) | \nu (\dd s) \leqq \left(\int_{S} | x (s) | ^{2} \nu (\dd s)\right) ^{1/2} \left(\int_{S} 1 \cdot \nu (\dd s)\right) ^{1/2} \leqq \| x \| _{\varrho} \cdot \nu (S) ^{1/2},
\end{aligned}
\]
where $\| x\| _{\varrho} = \left(\int \limits_{S}|x(s)|^2\varrho (\dd s)\right)^{1/2}$. Thus, by F. \textsc{Riesz}' representation theorem, there exists a uniquely determined $y\in L^{2}(S,\mathfrak{B},\varrho)$ such that 
\[
\int_{S} x (s) \nu (\dd s) = \int_{S} x (s) \overline {{y (s)}} \varrho (\dd s) = \int_{S} x (s) \overline {{y (s)}} m (\dd s) + \int_{S} x (s) \overline {{y (s)}} \nu (\dd s)
\]
holds for all $x \in L^{2}(S, \mathfrak{B}, \varrho)$. Taking $x$ as non-negative functions and considering the real part of both sides, we may assume that $y(s)$ is a real-valued function. Hence 
\begin{equation}\label{eq:III.8.2}
\int_{S} x (s) (1 - y (s)) v (\dd s) = \int_{S} x (s) y (s) m (\dd s) \text{ if } x (s) \in L ^{2} (S, \mathfrak {B}, \varrho)
\end{equation}
is non-negative. 

We can prove $0 \leqq y(s) < 1$ $\varrho$-a.e. To this purpose, set $E_{1} = \{s; y(s) < 0\}$ and $E_{2} = \{s; y(s) \geqq 1\}$. If we take the defining function $C_{E_1}(s)$ of $E_{1}$ for $x(s)$ in \eqref{eq:III.8.2}, then the left hand side is $\geqq 0$ and hence $\int_{E_1} y(s) m(\dd s) \geqq 0$. 

Thus we must have $m(E_1) = 0$, and so, by the $m$-absolute continuity of $\nu$, $\nu(E_1) = 0$, $\varrho(E_1) = 0$. We may also prove $\varrho(E_2) = 0$, by taking the defining function $C_{E_1}(s)$ for $x(s)$ in \eqref{eq:III.8.2}. Therefore $0 \leqq y(s) < 1$ $\varrho$-a.e. on $S$. 

Let $x(s)$ be $\mathfrak{B}$-measurable and $\geqq 0 \varrho$-a.e. Then, by $\varrho(S) < \infty$, the ``truncated functions'' $x_{n}(s) = \min(x(s), n)$ belong to $L^{2}(S, \mathfrak{B}, \varrho) (n = 1, 2, \ldots)$, and hence 
\begin{equation}\label{eq:III.8.3}
\int_{S} x _{n} (s) (1 - y (s)) \nu (\dd s) = \int_{S} x _{n} (s) y (s) m (\dd s) \quad (n = 1, 2, \dots).
\end{equation}

Since the integrals increase monotonely as $n$ increases, we have 
\begin{equation}\label{eq:III.8.4}
\lim\limits _{n \rightarrow \infty} \int_{S} x _{n} (s) (1 - y (s)) v (\dd s) = \lim\limits _{n \rightarrow \infty} \int_{S} x _{n} (s) y (s) m (\dd s) = L \leqq \infty .
\end{equation}

Since the integrands are $\geqq 0$ $\varrho$-a.e., we have, by the Lebesgue-Fatou Lemma, 
\begin{align}
L& \geqq \int_{S} \varliminf\limits_{n \rightarrow \infty} (x _{n} (s) (1 - y (s)) v (\dd s) = \int_{S} x (s) (1 - y (s)) v (\dd s), \label{eq:III.8.5} \\
L& \geqq \int_ {\mathcal {S}} \lim\limits _{n \rightarrow \infty} (x _{n} (s) y (s)) m (\dd s) = \int_ {\mathcal {S}} x (s) y m (\dd s), \notag
\end{align}
under the convention that if $x(s) (1 - y(s))$ is not $\nu$-integrable, the right hand side is equal to $\infty$; and the same convention for $x(s) y(s)$. If $x(s) y(s)$ is $m$-integrable, then, by the Lebesgue-Fatou Lemma, 
\begin{equation}\label{eq:III.8.6}
L \leqq \int_{S} \overline {{\lim\limits _{n \rightarrow \infty}}} (x _{n} (s) y (s)) m (\dd s) = \int_{S} x (s) y (s) m (\dd s).
\end{equation}

This formula is true even if $x(s)y(s)$ is not $m$-integrable, under the convention that then $L=\infty$. Under a similar convention, we have 
\begin{equation}\label{eq:III.8.7}
L \leqq \int_{S} x (s) (1 - y (s)) v (\dd s).
\end{equation}

Therefore, we have 
\begin{equation}\label{eq:III.8.8}
\int_{S} x (s) (1 - y (s)) v (\dd s) = \int_{S} x (s) y (s) m (\dd s) \text{ for every } x (s) \text{ which }
\end{equation}
is $\mathfrak{B}$-measurable and $\geqq 0$ $\varrho$-a.e., 
under the convention that, if either side of the equality is $=\infty$ then the other side is also $=\infty$. 

Now we put 
\[
x (s) (1 - y (s)) = z (s), y (s) (1 - y (s)) ^{-1} = p (s).
\]

Then, under the same convention as in \eqref{eq:III.8.8}, we have 
\begin{equation}\label{eq:III.8.9}
\int_{S} z (s) \nu (\dd s) = \int_{S} z (s) p (s) m (\dd s) \text{ if } z (s) \text{ is } \mathfrak {B} \text{-measurable }
\end{equation}
and $\geqq 0$ $\varrho$-a.e. 

If we take the defining function $C_B(s)$ of $B \in \mathfrak{B}$ for $z(s)$, we obtain 
\[
\nu (B) = \int_{B} p (s) m (\dd s) \text{ for all } B \in \mathfrak {B}.
\]

The last part of the Theorem is clear in view of definition \eqref{eq:III.8.1}. 

Reference. For a straightforward proof of the Lebesgue-Nikodym theorem based upon Hahn's decomposition (Theorem~\ref{thm:I.3.Theorem.3} in Chapter~\ref{chap:I}, \ref{sec:I.3}) see K. \textsc{Yosida} \cite{yosida345}. This proof is reproduced in \textsc{Halmos} \cite{halmos92}, p. $128$. See also \textsc{Saks} \cite{saks280} and \textsc{Dunford}-\textsc{Schwartz} \cite{dunford44}. 


\yossection{9}{The Aronszajn-Bergman Reproducing Kernel}
\label{sec:III.9}

Let $A$ be an abstract set, and let a system $X$ of complex-valued functions defined on $A$ constitute a Hilbert space by the scalar product 
\begin{equation}\label{eq:III.9.1}
(f, g) = (f (a), g (a)) _{a}.
\end{equation}

A complex-valued function $K(a, b)$ defined on $A \times A$ is called a reproducing kernel of $X$ if it satisfies the condition: 
\begin{align}
\text{ For any fixed } b, K (a, b)& \in X \text{ as a function of } a, \label{eq:III.9.2} \\
f (b)& = \left(f (a), K (a, b)\right) _{a} \text{ and hence } \overline {{f (b)}} = \left(K (a, b), f (a)\right) _{a}. \label{eq:III.9.3}
\end{align}

As for the existence of reproducing kernels, we have 

\begin{theorem}[N. \textsc{Aronszajn} {\cite{aronszajn7}}, S. \textsc{Bergman} {\cite{bergman15}}]\label{thm:III.9.Theorem.1} $X$ has a reproducing kernel $K$ iff there exists, for any $y_{0} \in A$, a positive constant $C_{y_{0}}$, depending upon $y_{0}$, such that 
\begin{equation}\label{eq:III.9.4}
\left| f \left(y _{0}\right) \right| \leqq C _ {y _{0}} \left\| f \right\| \quad \text{ for all } \quad f \in X.
\end{equation}
\end{theorem}

\begin{proof}
 The ``only if'' part is proved by applying Schwarz' inequality to $f(y_{0}) = (f(x), K(x, y_{0}))_{x}$: 
\begin{equation}\label{eq:III.9.5}
\left| f \left(y _{0}\right) \right| \leqq \left\| f \right\| \cdot \left(K \left(x, y _{0}\right), K \left(x, y _{0}\right)\right) _{x} ^{1/2} = \left\| f \right\| K \left(y _{0}, y _{0}\right) ^{1/2}. \end{equation}

The ``if'' part is proved by applying F. \textsc{Riesz}' representation theorem to the linear functional $F_{y_{0}}(f) = f(y_{0})$ of $f \in X$. Thus, there exists a uniquely determined vector $g_{y_{0}}(x)$ of $X$ such that, for every $f \in X$, 
\[
f (y _{0}) = F _ {y _{0}} (f) = (f (x), g _ {y _{0}} (x)) _{x},
\]
and so $g_{y_0}(x) = K(x, y_0)$ is a reproducing kernel of $X$. The proof shows that the reproducing kernel is uniquely determined. 
\end{proof}

\begin{corollaryu} We have 
\begin{equation}\label{eq:III.9.6}
\sup\limits _{\| f \| \leqq 1} | f (y _{0}) | = K (y _{0}, y _{0}) ^{1/2},
\end{equation}
the supremum being attained by 
\begin{equation}\label{eq:III.9.7}
f _{0} (x) = \varrho K (x, y _{0}) / K (y _{0}, y _{0}) ^{1/2}, | \varrho | = 1.
\end{equation}
\end{corollaryu}

\begin{proof}
 The equality in the Schwarz' inequality \eqref{eq:III.9.5} holds iff $f(x)$ and $K(x, y_0)$ are linearly dependent. From the two conditions $f(x) = \alpha K(x, y_0)$ and $\|f\| = 1$, we obtain 
\[
1 = | \alpha | (K (x, y _{0}), K (x, y _{0})) _{x} ^{1/2} = | \alpha | K (y _{0}, y _{0}) ^{1/2}, \text{ that is, } | \alpha | = K (y _{0}, y _{0}) ^{-1/2}. \]

Hence the equality sign in \eqref{eq:III.9.5} is attained by $f_0(x)$. 
\end{proof}

\begin{examplen} Consider the Hilbert space $A^2(G)$. For any $f \in A^2(G)$ and $z \in G$, we have (see \eqref{eq:I.9.4} in Chapter~\ref{chap:I}, \ref{sec:I.9})
\[
\left| f \left(z _{0}\right) \right| ^{2} \leqq (\pi r ^{2}) ^{-1} \int_ {\left| z - z _{0} \right| \leqq r} | f (z) | ^{2} \dd x \dd y (z = x + i y).
\]

Thus $A^2(G)$ has the reproducing kernel which will be denoted by $K_G(z, z')$. This $K_G(z, z')$ is called Bergman's kernel of the domain $G$ of the complex plane. The following theorem of Bergman illustrates the meaning of $K_G(z, z')$ in the theory of conformal mapping. 

\end{examplen}
\begin{theorem}\label{thm:III.9.Theorem.2} Let $G$ be a simply connected bounded open domain of the complex plane, and $z_{0}$ be any point of $G$. By Riemann's theorem, there exists a uniquely determined regular function $w = f(z; z_{0})$ of $z$ which gives a one-to-one conformal map of the domain $G$ onto the sphere $|w| < \varrho_{G}$ of the complex $w$-plane in such a way that 
\[
f _{0} (z _{0}; z _{0}) = 0, (\dd f _{0} (z; z _{0}) / \dd z) _ {z = z _{0}} = 1.
\]
\end{theorem}

Bergman's kernel $K_{G}(z;z_{0})$ is connected with $f_{0}(z;z_{0})$ by 
\begin{equation}\label{eq:III.9.8}
f _{0} (z; z _{0}) = K _{G} \left(z _{0}; z _{0}\right) ^{-1} \int_ {z _{0}} ^{z} K _{G} (t; z _{0}) \dd t,
\end{equation}
where the integral is taken along any rectifiable curve lying in $G$ and connecting $z_{0}$ with $z$. 

$A_{1}^{2}(G) = \{f(z); f(z)\text{ is holomorphic in } G, f'(z) \in A^{2}(G), f(z_{0}) = 0\text{ and } f'(z_{0}) = 1\}$, and consider, for any $f \in A_{1}^{2}(G)$, the number 
\begin{equation}\label{eq:III.9.9}
\left\| f ^{\prime} \right\| ^{2} = \int_{G} \left| f ^{\prime} (z) \right| ^{2} \dd x \dd y, z = x + i y.
\end{equation}

If we denote by $z = \varphi(w)$ the inverse function of $w = f_0(z; z_0)$, then, for any $f \in A_1^2(G)$, 
\[
\left\| f ^{\prime} \right\| ^{2} = \iint_ {| w | < \varrho_{G}} \left| f ^{\prime} (\varphi (w)) \right| ^{2} \left| \varphi^{\prime} (w) \right| ^{2} \dd u \dd v, w = u + i v.
\]

For, by the Cauchy-Riemann partial differential equations 
\[
x _{u} = y _{v}, x _{v} = - y _{u},
\]
we have 
\[
\begin{aligned}
&\dd x \dd y  = \frac {\partial (x , y)}{\partial (u , v)} \dd u \dd v = (x _{u} y _{v} - y _{u} x _{v}) \dd u \dd v = (x _{u} ^{2} + y _{u} ^{2}) \dd u \dd v = | \varphi^{\prime} (w) | ^{2} \dd u \dd v.
\end{aligned}
\]

Let $f \in A_1^2(G)$, and let $F(w) = f(\varphi(w))$ be expanded into power series: 
\[
F (w) = f (\varphi (w)) = w + \sum\limits_{n=2} ^{\infty} c _{n} w ^{n} \text{ for } | w | < \varrho_{G}.
\]

Then $F'(w) = f'(\varphi(w)) \varphi'(w) = 1 + \sum\limits_{n=2}^{\infty} n c_n w^{n-1}$ and so 
\[
\begin{aligned}
\left\| f ^{\prime} \right\| ^{2} &= \iint_ {| w | < \varrho_{G}} \left| 1 + \sum\limits_{n=2} ^{\infty} n c _{n} w ^{n-1} \right| ^{2} \dd u \dd v
\\
&= \int_{0} ^ {\varrho_{G}} \dd r \left\{\int_{0} ^{2 \pi} \left(r + \sum\limits_{n=2} ^{\infty} n ^{2} \left| c _{n} \right| ^{2} r ^{2n-1}\right) \dd \theta \right\} = \pi \varrho_{G} ^{2} + \sum\limits_{n=2} ^{\infty} \pi n \left| c _{n} \right| ^{2} \varrho_{G} ^{2n}.
\end{aligned}
\]

Therefore minimum $\left\|f^{\prime}\right\|=\sqrt{\pi}\varrho_{G}$, and this minimum is attained iff $F(w)=f(\varphi(w))=w$, that is, iff $f(z)=f_{0}(z;z_{0})$. 

We set, for any $f \in A_1^2(G)$, $g(z) = f(z) / \|f'\|$. Then $\|g'\| = 1$. If we put $\tilde{A}^2(G) = \{g(z); g(z) \text{ is holomorphic in } G, g(z_0) = 0, g'(z_0) > 0 \text{ and } \|g'\| = 1\}$, 
then the above remark shows that 
\[
\underset {g \in \widetilde {A} ^{2} (G)} {\text{ maximum }} g ^{\prime} \left(z _{0}\right) = 1 / \left\| f _{0} ^{\prime} \right\| = (\sqrt {\pi} \varrho_{G}) ^{-1},
\]
and this maximum is attained iff $g(z)$ is equal to 
\[
g _{0} (z) = f _{0} \left(z; z _{0}\right) / \left\| f _{0} ^{\prime} \right\| = f _{0} \left(z; z _{0}\right) / \sqrt {\pi} \varrho_{G}.
\]

Hence, by \eqref{eq:III.9.7}, we obtain 
\[
g _{0} ^{\prime} (z) = (\sqrt {\pi} \varrho_{G}) ^{-1} \frac {\dd f _{0} (z ; z _{0})}{\dd z} = \lambda K _{G} (z; z _{0}) / K _{G} (z _{0}; z _{0}) ^{1/2}, | \lambda | = 1.
\]

Hence, by putting $z = z_0$, we have 
\[
\left(\lambda \sqrt {\pi} \varrho_{G}\right) ^{-1} = K _{G} (z _{0}; z _{0}) / K _{G} (z _{0}; z _{0}) ^{1/2} = K _{G} (z _{0}; z _{0}) ^{1/2},
\]
and so we have proved the formula 
\[
\frac {\dd f _{0} (z ; z _{0})}{\dd z} = K _{G} (z; z _{0}) / K _{G} (z _{0}; z _{0}).
\]


\yossection{10}{The Negative Norm of P. \textsc{Lax}}
\label{sec:III.10}

Let $H_{0}^{s}(\Omega)$ be the completion of the pre-Hilbert space $C_{0}^{\infty}(\Omega)$ endowed with the scalar product $(\varphi, \psi)_{s}$ and the norm $\| \varphi \|_{s}$: 
\begin{equation}\label{eq:III.10.1}
(\varphi , \psi) _{s} = \sum\limits_{| j | \leqq s} \int_{\Omega} D ^{j} \varphi (x) \overline {{D ^{j} \psi (x)}} \dd x, \| \varphi \| _{s} = (\varphi , \varphi) _{s} ^{1/2}.
\end{equation}

Any element $b \in H_0^0(\Omega) = L^2(\Omega)$ defines a continuous linear functional $f_b$ on $H_0^s(\Omega)$ by 
\begin{equation}\label{eq:III.10.2}
f _{b} (w) = (w, b) _ {\mathbf {0}}, w \in H _ {\mathbf {0}} ^{s} (\Omega).
\end{equation}

For, by Schwarz' inequality, we have 
\[
\left| (w, b) _{0} \right| \leqq \left\| w \right\| _{0} \cdot \left\| b \right\| _{0} \leqq \left\| w \right\| _{s} \cdot \left\| b \right\| _{0}.
\]

Therefore, if we define the negative norm of $b \in H_0^0(\Omega) = L^2(\Omega)$ 
\begin{equation}\label{eq:III.10.3}
\left\| b \right\| _{-s} = \sup\limits _ {w \in H _{0} ^{\varepsilon} (\Omega), \| w \| _{s} \leqq 1} \left| f _{b} (w) \right| = \sup\limits _ {w \in H _{0} ^{\varepsilon} (\Omega), \| w \| _{s} \leqq 1} \left| (w, b) _{0} \right|,
\end{equation}
then we have 
\begin{equation}\label{eq:III.10.4}
\left\| b \right\| _{-s} \leqq \left\| b \right\| _{0},
\end{equation}
and so, by $\|b\|_{-s} \geqq |(w/\|w\|_s, b)_0|$, 
\begin{equation}\label{eq:III.10.5}
\left| (w, b) _{0} \right| \leqq \left\| w \right\| _{s} \cdot \left\| b \right\| _{-s}.
\end{equation}

Hence we may write 
\begin{equation}\tag{3'}\label{eq:III.10.3'}
\left\| b \right\| _{-s} = \left\| f _{b} \right\| _{-s} = \sup\limits _ {\left\| w \right\| _{s} \leqq 1} | (w, b) _{0} | \text{ for any } b \in H _{0} ^{0} (\Omega).
\end{equation}

We shall prove 

\begin{theorem}[P. \textsc{Lax} {\cite{lewy172}}]\label{thm:III.10.Theorem.1} The dual space $H_0^s(\Omega)'$ of the space $H_0^s(\Omega)$ may be identified with the completion of the space $H_0^0(\Omega) = L^2(\Omega)$ with respect to the negative norm. 
\end{theorem}

For the proof we prepare 

\begin{propositionu} The totality $F$ of the continuous linear functionals on $H_0^s(\Omega)$ of the form $f_b$ is dense in the Hilbert space $H_0^s(\Omega)'$, which is the dual space of $H_0^s(\Omega)$. 
\end{propositionu}

\begin{proof}
 $F$ is total on $H_0^s(\Omega)$ in the sense that, for a fixed $w \in H_0^s(\Omega)$, the simultaneous vanishing $f_b(w) = 0$, $b \in H_0^0(\Omega)$, takes place only if $w = 0$. This is clear since any element $w \in H_0^s(\Omega)$ is also an element of $H_0^0(\Omega)$. 

Now if $F$ is not dense in the Hilbert space $H_0^s(\Omega)'$, then there exists an element $T \neq 0$ of the second dual space $H_0^s(\Omega)'' = (H_0^s(\Omega)')'$ such that $T(f_b) = 0$ for all $f_b \in F$. By the reflexivity of the Hilbert space $H_0^s(\Omega)$, there exists an element $t \in H_0^s(\Omega)$ such that $T(f) = f(t)$ for all $f \in H_0^s(\Omega)'$. Thus $T(f_b) = f_b(t) = 0$ for all $b \in H_0^0(\Omega)$. This implies, by the totality of $F$ proved above, that $t = 0$, contrary to $T \neq 0$. 
\end{proof}
\begin{corollaryu} We have, dually to \eqref{eq:III.10.3}, 
\begin{equation}\label{eq:III.10.6}
\left\| w \right\| _{s} = \sup\limits _ {b \in H _{0} ^{0} (\Omega), \| b \| _{-s} \leqq 1} \left| (w, b) _{0} \right| \text{ for any } w \in H _{0} ^{s} (\Omega).
\end{equation}
\end{corollaryu}

\begin{proof}
 Clear from Corollary~\ref{cor:III.6.Corollary.3} in Chapter~\ref{chap:III}, \ref{sec:III.6}, because $F = \{f_b; b \in H_0^0(\Omega)\}$ is dense in $H_0^s(\Omega)'$. 
\end{proof}

\begin{proof}[Proof of Theorem~\ref{thm:III.10.Theorem.1}.]
Clear from the facts that \begin{enumerate}[i)]
\item $F$ is dense in the dual space $H_{0}^{s}(\Omega)'$
\item $F$ is in one-to-one correspondence to the set $H_{0}^{0}(\Omega)=L^{2}(\Omega)$ preserving the negative norm, i.e.,
\[
F \ni f _{b} \leftrightarrow b \in H _{0} ^{0} (\Omega) \text{ and } \| f _{b} \| _{-s} = \| b \| _{-s}.
\]
\end{enumerate}


We shall denote by $H_0^{-s}(\Omega)$ the completion of $H_0^0 (\Omega)$ with respect to the negative norm $\|b\|_{-s}$. Thus 
\begin{equation}\label{eq:III.10.7}
H _{0} ^{s} (\Omega) ^{\prime} = H _{0} ^{-s} (\Omega).
\end{equation}

For any continuous linear functional $f$ on $H_0^s(\Omega)$, we shall denote by $\langle w, f \rangle$ the value of $f$ at $w \in H_0^s(\Omega)$. Thus, for any $b \in H_0^0(\Omega)$, we may write 
\begin{equation}\label{eq:III.10.8}
f _{b} (w) = (w, b) _{0} = \langle w, f _{b} \rangle = \langle w, b \rangle , w \in H _{0} ^{s} (\Omega),
\end{equation}
and have the generalized Schwarz' inequality 
\begin{equation}\label{eq:III.10.9}
\left| \langle w, b \rangle \right| \leqq \left\| w \right\| _{s} \left\| b \right\| _{-s},
\end{equation}
which is precisely \eqref{eq:III.10.5}. 

Now we can prove 
\end{proof}

\begin{theorem}[P. \textsc{Lax} {\cite{lewy172}}]\label{thm:III.10.Theorem.2} Any continuous linear functional $g(b)$ on $H_0^{-s}(\Omega)$ can be represented, by a fixed element $w \in H_0^s(\Omega)$, as 
\begin{equation}\label{eq:III.10.10}
g (b) = g _{w} (b) = \overline {{\langle w , b \rangle}}.
\end{equation}
\end{theorem}

We have, in particular, 
\begin{equation}\label{eq:III.10.11}
H _{0} ^{s} (\Omega) ^{\prime} = H _{0} ^{-s} (\Omega), H _{0} ^{-s} (\Omega) ^{\prime} = H _{0} ^{s} (\Omega).
\end{equation}



\begin{proof}
 If $b \in H_0^0(\Omega)$, then $\langle w, b \rangle = f_b(w) = (w, b)_0$. Since $F = \{f_b; b \in H_0^0(\Omega)\}$ is dense in the Hilbert space $H_0^s(\Omega)'$, we know, by \eqref{eq:III.10.9}, that $\overline{\langle w, b \rangle} = (b, w)_0$ with a fixed $w \in H_0^s(\Omega)$, defines a linear functional $g_w$ continuous on a dense subset $F$ of $H_0^s(\Omega)'$. The norm of this functional $g_w$ on $F$ will be denoted by $\|g_w\|_s$. Then, by \eqref{eq:III.10.6}, 
\begin{equation}\label{eq:III.10.12}
\left\| g _{w} \right\| _{s} = \sup\limits _{\left\| b \right\| - s \leqq 1} \left| (b, w) _{0} \right| = \sup\limits _{\left\| b \right\| - s \leqq 1} \left| (w, b) _{0} \right| = \left\| w \right\| _{s}. \end{equation}

We may thus extend, by continuity, the functional $g_{w}$ on $F$ to a continuous linear functional on the completion of $F$ (with respect to the negative norm), that is, $g_{w}$ can be extended to a continuous linear functional on $H_{0}^{s}(\Omega)' = H_{0}^{-s}(\Omega)$; we denote this extension by the same letter $g_{w}$. We have thus 
\begin{equation}\label{eq:III.10.13}
\left\| g _{w} \right\| = \sup\limits _{\left\| b \right\| - s \leqq 1} \left| g _{w} (b) \right| = \left\| w \right\| _{s}.
\end{equation}

Hence, in view of the completeness of the space $H_0^s(\Omega)$, the totality $G$ of the continuous linear functionals $g_w$ on $H_0^{-s}(\Omega)$ may be considered as a closed linear subspace of $H_0^{-s}(\Omega)'$ by the correspondence $g_w \leftrightarrow w$. If this closed subspace $G$ were not dense in $H_0^{-s}(\Omega)'$, then there exists a continuous linear functional $f \neq 0$ on $H_0^{-s}(\Omega)'$ such that $f(g_w) = 0$ for all $g_w \in G$. But, since the Hilbert space $H_0^{-s}(\Omega)$ is reflexive, such a functional $f$ is given by $f(g_w) = g_w(f_0)$, $f_0 \in H_0^{-s}(\Omega)$, and so by \eqref{eq:III.10.13} $f_0$ must be equal to $0$, contrary to the fact $f \neq 0$. Therefore we have proved $H_0^{-s}(\Omega)' = H_0^s(\Omega)$. 
\end{proof}

\begin{remarku} The notion of the negative norm was introduced by P. \textsc{Lax} with the view of applying it to the genuine differentiability of distribution solutions of linear partial differential equations. We shall discuss such differentiability problems in later chapters. It is to be noted that the notion of the negative norm is also introduced naturally through the Fourier transform. This was done by J. \textsc{Leray} \cite{leray175} earlier than \textsc{Lax}. We shall explain the point in a later chapter on the Fourier transform. 
\end{remarku}


\yossection{11}{Local Structures of Generalized Functions}
\label{sec:III.11}

A generalized function is locally the distributional derivative of a function. More precisely, we can prove 

\begin{theoremn}[L. \textsc{Schwartz} {\cite{schwartz294}}] Let $T$ be a generalized function in $\Omega \subseteq \mathbb{R}^{n}$. Then, for any compact subset $K$ of $\Omega$, there exist a positive integer $m_0 = m_0(T, K)$ and a function $f(x) = f(x; T, K, m_0) \in L^2(K)$ such that
\begin{equation}\label{eq:III.11.1}
T (\varphi) = \int_{K} f (x) \frac {\partial^ {n m _{0}} \varphi (x)}{\partial x _{1} ^ {m _{0}} \partial x _{2} ^ {m _{0}} \cdots \partial x _{n} ^ {m _{0}}} \dd x \text{ whenever } \varphi \in \mathfrak {D} _{K} (\Omega)
\end{equation}

\end{theoremn}
\begin{proof}
 By the Corollary in Chapter~\ref{chap:I}, \ref{sec:I.8}, there exist a positive constant $C$ and a positive integer $m$ such that 
\begin{equation}\label{eq:III.11.2}
\left| T (\varphi) \right| \leqq C \sup\limits _{| j | \leqq m, x \in K} \left| D ^{j} \varphi (x) \right| \text{ whenever } \varphi \in \mathfrak {D} _{K} (\Omega). \end{equation}

Thus there exists a positive number $\delta$ such that 
\begin{equation}\label{eq:III.11.3}
\phi_{m} (\varphi) = \sup\limits _{| j | \leqq m, x \in K} | D ^{j} \varphi (x) | \leqq \delta , \varphi \in \mathfrak {D} _{K} (\Omega), \text{ implies } | T (\varphi) | \leqq 1.
\end{equation}

We introduce the notation 
\begin{equation}\label{eq:III.11.4}
\frac {\partial^{s}}{\partial x ^{s}} = \frac {\partial^{sn}}{\partial x _{1} ^{s} \partial x _{2} ^{s} \cdots \partial x _{n} ^{s}},
\end{equation}
and prove that there exists a positive constant $\varepsilon$ such that, for $m_0 = m + 1$, 
\begin{equation}\label{eq:III.11.5}
\int_{\Omega} \left| \partial^ {m _{0}} \varphi (x) / \partial x ^ {m _{0}} \right| ^{2} \dd x \leqq \varepsilon , \varphi \in \mathfrak {D} _{K} (\Omega), \text{ implies } p _{m} (\varphi) \leqq \delta .
\end{equation}

This is proved by repeated application of the following inequality 
\[
\begin{aligned}
\left| \psi (x) \right|& \leqq \underset {K \cap (- \infty , x _{i})} {\int} \left| \partial \psi (x _{1}, \dots , x _{i-1}, y, x _{i+1}, \dots , x _{n}) / \partial y \right| \dd y
\\
&\leqq \left(\underset {K \cap (- \infty , x _{i})} {\int} \dd y\right) ^{1/2} \left(\underset {K \cap (- \infty , x _{i})} {\int} \left| \partial \psi (x _{1}, \dots , x _{i-1}, y, x _{i+1}, \dots , x _{n}) / \partial y \right| ^{2} \dd y\right) ^{1/2}
\\
&= t ^{1/2} \left(\underset {K \cap (- \infty , x _{i})} {\int} \left| \partial \psi (x _{1}, \dots , x _{i-1}, y, x _{i+1}, \dots , x _{n}) / \partial y \right| ^{2} \dd y\right) ^{1/2},
\end{aligned}
\]
where $t$ is the diameter of $K$, i.e., the maximum distance between two points of the compact set $K$. 

Consider the mapping $\varphi(x) \to \psi(x) = \partial^{m_0} \varphi(x)/\partial x^{m_0}$ defined on $\mathfrak{D}_K(\Omega)$ into $\mathfrak{D}_K(\Omega)$. As may be seen by integration, $\psi(x) = 0$ implies $\varphi(x) = 0$. Hence the above mapping is one-to-one. Thus $T(\varphi)$, $\varphi \in \mathfrak{D}_K(\Omega)$, defines a linear functional $S(\psi)$, $\psi(x) = \partial^{m_0} \varphi(x)/\partial x^{m_0}$, by $S(\psi) = T(\varphi)$. By \eqref{eq:III.11.3} and \eqref{eq:III.11.5}, $S$ is a continuous linear functional on the pre-Hilbert space $X$ consisting of such $\psi$'s and topologized by the norm $\|\psi\| = \left( \int_K |\psi(x)|^2 \dd x \right)^{1/2}$. 

Thus there exists, by F. \textsc{Riesz}' representation theorem, a uniquely determined function $f(x)$ from the completion of $X$ such that 
\[
T (\varphi) = S (\psi) = \int_{K} \left(\partial^ {m _{0}} \varphi (x) / \partial x ^ {m _{0}}\right) \cdot f (x) \dd x \text{ for all } \varphi \in \mathfrak {D} _{K} (\Omega).
\]

Actually, the completion of $X$ is contained in $L^{2}(K)$ as a closed linear subspace, and so the Theorem is proved. 
\end{proof}

\yossub{References for Chapter III}

For general account of Hilbert spaces, see N. I. \textsc{Achieser}-I. M. \textsc{Glazman} \cite{achieser1}, N. \textsc{Dunford}-J. \textsc{Schwartz} \cite{dunford45}, B. Sz. \textsc{Nagy} \cite{nagy217}, F. \textsc{Riesz}-B. Sz. \textsc{Nagy} \cite{nagy219} and M. H. \textsc{Stone} \cite{stone311}. 
   % III
% ==== IV ====
\yoschapter{IV}{The Hahn-Banach Theorems}
\label{chap:IV}

In a Hilbert space, we can introduce the notion of orthogonal coordinates through an orthogonal base, and these coordinates are the values of bounded linear functionals defined by the vectors of the base. This suggests that we consider continuous linear functionals, in a linear topological space, as generalized coordinates of the space. To ensure the existence of non-trivial continuous linear functionals in a general locally convex linear topological space, we must rely upon the Hahn-Banach extension theorems. 


\yossection{1}{The Hahn-Banach Extension Theorem in Real Linear Spaces}
\label{sec:IV.1}

\begin{theoremu}[\textsc{Hahn} {\cite{hahn90}}, \textsc{Banach} {\cite{banach9}}] Let $X$ be a real linear space and let $p(x)$ be a real-valued function defined on $X$ satisfying the conditions:
\begin{align}
p (x + y)& \leqq p (x) + p (y) \quad (\text{ subadditivity }), \label{eq:IV.1.1} \\
p (\alpha x)& = \alpha p (x) \quad \text{ for } \quad \alpha \geqq 0. \label{eq:IV.1.2}
\end{align}

Let $M$ be a real linear subspace of $X$ and $f_0$ a real-valued linear functional defined on $M$: 
\begin{equation}\label{eq:IV.1.3}
f _{0} (\alpha x + \beta y) = \alpha f _{0} (x) + \beta f _{0} (y) \quad \text{ for } \quad x, y \in M \text{ and real } \alpha , \beta .
\end{equation}

Let $f_0$ satisfy $f_0(x) \leqq p(x)$ on $M$. Then there exists a real-valued linear functional $F$ defined on $X$ such that \begin{enumerate}[i)]
\item $F$ is an extension of $f_0$, i.e., $F(x) = f_0(x)$ for all $x \in M$
\item $F(x) \leqq p(x)$ on $X$.
\end{enumerate}


\end{theoremu}
\begin{proof}
 Suppose first that $X$ is spanned by $M$ and an element $x_0 \bar{\in} M$, that is, suppose that 
\[
X = \{x = m + \alpha x _{0}; m \in M, \alpha \text{ real } \}. \]

Since $x_0 \bar{\in} M$, the above representation of $x \in X$ in the form $x = m + \alpha x_0$ is unique. It follows that, if, for any real number $c$, we set 
\[
F (x) = F (m + \alpha x _{0}) = f _{0} (m) + \alpha c,
\]
then $F$ is a real linear functional on $X$ which is an extension of $f_0$. We have to choose $c$ such that $F(x) \leqq p(x)$, that is, $f_0(m) + \alpha c \leqq p(m + \alpha x_0)$. This condition is equivalent to the following two conditions: 
\[
\begin{gathered}
f _{0} (m / \alpha) + c \leqq p (x _{0} + m / \alpha) \quad \text{ for } \quad \alpha > 0,
\\
f _{0} (m / (- \alpha)) - c \leqq p (- x _{0} + m / (- \alpha)) \quad \text{ for } \quad \alpha < 0.
\end{gathered}
\]

To satisfy these conditions; we shall choose $c$ such that 
\[
f _{0} \left(m ^{\prime}\right) - p \left(m ^{\prime} - x _{0}\right) \leqq c \leqq p \left(m ^{\prime \prime} + x _{0}\right) - f _{0} \left(m ^{\prime \prime}\right) \text{ for all } m ^{\prime}, m ^{\prime \prime} \in M.
\]

Such a choice of $c$ is possible since 
\[
\begin{aligned} f _{0} (m ^{\prime}) + f _{0} (m ^{\prime \prime}) & = f _{0} (m ^{\prime} + m ^{\prime \prime}) \leqq p (m ^{\prime} + m ^{\prime \prime}) = p (m ^{\prime} - x _{0} + m ^{\prime \prime} + x _{0}) \\ & \leqq p (m ^{\prime} - x _{0}) + p (m ^{\prime \prime} + x _{0}): \end{aligned}
\]
we have only to choose $c$ between the two numbers 
\[
\sup\limits _ {m ^{\prime} \in M} [ f _{0} (m ^{\prime}) - p (m ^{\prime} - x _{0}) ] \quad \text{ and } \quad \inf\limits _ {m ^{\prime \prime} \in M} [ p (m ^{\prime \prime} + x _{0}) - f _{0} (m ^{\prime \prime}) ].
\]

Consider now the family of all real linear extensions $g$ of $f_0$ for which the inequality $g(x) \leqq p(x)$ holds for all $x$ in the domain of $g$. We make this family into a partially ordered family by defining $h > g$ to mean that $h$ is an extension of $g$. Then Zorn's Lemma ensures the existence of a maximal linear extension $g$ of $f_0$ for which the inequality $g(x) \leqq p(x)$ holds for all $x$ in the domain of $g$. We have to prove that the domain $D(g)$ of $g$ coincides with $X$ itself. If it does not, we obtain, taking $D(g)$ as $M$ and $g$ as $f_0$, a proper extension $F$ of $g$ which satisfies $F(x) \leqq p(x)$ for all $x$ in the domain of $F$, contrary to the maximality of the linear extension $g$. 
\end{proof}

\begin{corollaryu} Given a functional $p(x)$ defined on a real linear space $X$ such that \eqref{eq:IV.1.1} and \eqref{eq:IV.1.2} are satisfied. Then there exists a linear functional $f$ defined on $X$ such that 
\begin{equation}\label{eq:IV.1.4}
- p (- x) \leqq f (x) \leqq p (x).
\end{equation}
\end{corollaryu}

\begin{proof}
 Take any point $x_0 \in X$ and define $M = \{x; x = \alpha x_0, \alpha \text{ real }\}$. Set $f_0(\alpha x_0) = \alpha p(x_0)$. Then $f_0$ is a real linear functional defined on $M$. We have $f_0(x) \leqq p(x)$ on $M$. In fact, $\alpha p(x_0) = p(\alpha x_0)$ if $\alpha > 0$, and if $\alpha < 0$, we have $\alpha p(x_0) \leqq -\alpha p(-x_0) = p(\alpha x_0)$ by $0 = p(0) \leqq p(x_0) + p(-x_0)$. Thus there exists a linear functional $f$ defined on $X$ such that $f(x) = f_0(x)$ on $M$ and $f(x) \leqq p(x)$ on $X$. Since $-f(x) = f(-x) \leqq p(-x)$, we obtain $-p(-x) \leqq f(x) \leqq p(x)$. 
\end{proof}


\yossection{2}{The Generalized Limit}
\label{sec:IV.2}

The notion of a sequence $\{x_{n}\}$ of a countable number of elements $x_{n}$ is generalized to the notion of a directed set of elements depending on a parameter which runs through an uncountable set. The notion of the limit of a sequence of elements may be extended to the notion of the generalized limit of a directed set of elements. 

\begin{definitionu} $A$ partially ordered set $A$ of elements $\alpha, \beta, \ldots$ is called a directed set if it satisfies the condition: 

\begin{equation}\label{eq:IV.2.1}
\begin{minipage}[b]{0.92\linewidth}
For any pair $\alpha, \beta$ of elements of $A$, there exists $a\gamma \in A$ such that $\alpha < \gamma$, $\beta < \gamma$.
\end{minipage}
\end{equation}
\end{definitionu}


Let, to each point $\alpha$ of a directed set $A$, there be associated a certain set of real numbers $f(\alpha)$. Thus $f(\alpha)$ is a, not necessarily one-valued, real function defined on the directed set $A$. We write 
\[
\lim\limits _{\alpha \in A} f (\alpha) = a \quad (a \text{ is a real number })
\]
if, for any $\varepsilon > 0$, there exists an $\alpha_0 \in A$ such that $\alpha_0 < \alpha$ implies $|f(\alpha) - a| \leqq \varepsilon$ for all possible values of $f$ at $\alpha$. We say, in such a case, that the value $a$ is the generalized limit or Moore-Smith limit of $f(\alpha)$ through the directed set $A$. 

\begin{examplen} Consider a partition $\Delta$ of the real interval [$0$, $1$]:
\[
0 = t _{0} < t _{1} < \dots < t _{n} = 1.
\]

The totality $P$ of the partition of $[0,1]$ is a directed set $A$ by defining the partial order $\Delta \prec \Delta'$ as follows: If the partition $\Delta'$ is given by $0 = t_0' < t_1' < \dots < t_m' = 1$, then $\Delta \prec \Delta'$ means that $n \leqq m$ and that every $t_i$ is equal to some $t_j'$. Let $x(t)$ be a real-valued continuous function defined on $[0,1]$, and let $f(\Delta)$ be the totality of real numbers of the form 
\[
\sum_{j=0}^{n-1}(t_{j+1}-t_j)x(t_j'),
\]
where $t_j'$ is any point of $[t_j,t_{j+1}]$.

Thus $f(\Delta)$ is the totality of the Riemann sum of the function $x(t)$ pertaining to the partition $\Delta$. The value of the Riemann integral $\int_{0}^{1} x(t) \dd t$ is nothing but the generalized limit of $f(\Delta)$ through $P$. 

As to the existence of a generalized limit, we have the 

\end{examplen}
\begin{theoremn}[S. \textsc{Banach}] Let $x(\alpha)$ be a real-valued bounded function defined on a directed set $A$. The totality of such functions constitutes a real linear space $X$ by
\[
(x + y) (\alpha) = x (\alpha) + y (\alpha), (\beta x) (\alpha) = \beta x (\alpha).
\]

We can then define a linear functional, defined on $X$ and which we shall denote by $\LIM x(\alpha)$, satisfying the condition 
\[
\varliminf\limits_{\alpha \in A} x (\alpha) \leqq \operatorname * {L I M} _{\alpha \in A} x (\alpha) \leqq \varlimsup\limits_{\alpha \in A} x (\alpha),
\]
where 
\[
\varliminf\limits_{\alpha \in A} x (\alpha) = \sup\limits _{\alpha} \inf\limits _{\alpha < \beta} x (\beta), \varlimsup\limits_{\alpha \in A} x (\beta) = \inf\limits _{\alpha} \sup\limits _{\alpha < \beta} x (\beta).
\]

Therefore $\lim\limits_{\alpha \in A}x(\alpha) = \lim\limits_{\alpha \in A}x(\alpha)$ if the latter generalized limit exists. 

\end{theoremn}
\begin{proof}
 We put $p(x) = \varlimsup\limits_{\alpha \in A} x(\alpha)$. It is easy to see that this $p(x)$ satisfies the condition of the Hahn-Banach extension theorem. Hence, there exists a linear functional $f$ defined on $X$ such that $-p(-x) \leqq f(x) \leqq p(x)$ on $X$. We can easily prove that $\varliminf\limits_{\alpha \in A} x(\alpha) = -p(-x)$ so that we obtain the Theorem, by putting $\LIM_{\alpha \in A} x(\alpha) = f(x)$. 
\end{proof}


\yossection{3}{Locally Convex, Complete Linear Topological Spaces}
\label{sec:IV.3}

\begin{definitionu} As in numerical case, we may define a directed set $\{x_{\alpha}\}$ in a linear topological space $X$. $\{x_{\alpha}\}$ is said to converge to an element $x$ of 
\end{definitionu}

$X$, if, for every neighbourhood $U(x)$ of $x$, there exists an index $\alpha_{0}$ such that $x_{\alpha} \in U(x)$ for all indices $\alpha > \alpha_{0}$. A directed set $\{x_{\alpha}\}$ of $X$ is said to be fundamental, if every neighbourhood $U(0)$ of the zero vector $0$ of $X$ is assigned an index $\alpha_{0}$ such that $(x_{\alpha} - x_{\beta}) \in U(0)$ for all indices $\alpha, \beta > \alpha_{0}$. A linear topological space $X$ is said to be complete if every directed fundamental set of $X$ converges to some element $x \in X$ in the sense above. 

\begin{remarku} We can weaken the condition of the completeness, and require only that every sequence of $X$ which is fundamental as a directed set converges to an element $x \in X$; a space $X$ satisfying this condition is said to be sequentially complete. For normed linear spaces, the two definitions of completeness are equivalent. However, in the general case, not every sequentially complete space is complete. 
\end{remarku}

\begin{examplen} Let a sequence $\{f_h(x)\}$ of $\mathfrak{D}(\Omega)$ satisfy the condition $\lim\limits_{h,k\to \infty}(f_h - f_k) = 0$ in $\mathfrak{D}(\Omega)$. That is, by the Corollary of Proposition~\ref{prop:I.1.Proposition.7} in Chapter~\ref{chap:I}, \ref{sec:I.1}, we assume that there exists a compact subset $K$ of $\Omega$ for which $\supp(f_h)\subseteq K$ ( $h = 1,2,\ldots$) and $\lim\limits_{h,k\to \infty}(D^s f_h(x) - D^s f_k(x)) = 0$ uniformly on $K$ for any differential operator $D^{s}$. Then it is easy to see, by applying the Ascoli-Arzelà theorem, that there exists a function $f\in \mathfrak{D}(\Omega)$ for which $\lim\limits_{h\to \infty}D^{s}f_{h}(x) = D^{s}f(x)$ uniformly on $K$ for any differential operator $D^{s}$. Hence $\lim\limits_{h\to \infty}f_{h} = f$ in $\mathfrak{D}(\Omega)$ and so $\mathfrak{D}(\Omega)$ is sequentially complete. Similarly, we can prove that $\mathfrak{E}(\Omega)$ is also sequentially complete. 
\end{examplen}

As in the case of a normed linear space, we can prove the 

\begin{theoremu} Every locally convex linear topological space $X$ can be embedded in a locally convex, complete linear topological space, in which $X$ forms a dense subset. 
\end{theoremu}

We omit the proof. The reader is referred to the literature listed in J. A. \textsc{Dieudonné} \cite{dieudonne39}. Cf. also G. \textsc{Köthe} \cite{konishi164}. 


\yossection{4}{The Hahn-Banach Extension Theorem in Complex Linear Spaces}
\label{sec:IV.4}

\begin{theoremn}[\textsc{Bohnenblust}-\textsc{Sobczyk}] Let $X$ be a complex linear space and $p$ a semi-norm defined on $X$. Let $M$ be a complex linear subspace of $X$ and $f$ a complex linear functional defined on $M$ such that $|f(x)| \leqq p(x)$ on $M$. Then there exists a complex linear functional $F$ defined on $X$ such that \begin{enumerate}[i)]
\item $F$ is an extension of $f$
\item $|F(x)| \leqq p(x)$ on $X$.
\end{enumerate}


\end{theoremn}
\begin{proof}
 We observe that a complex linear space is also a real linear space if the scalar multiplication is restricted to real numbers. If $f(x) = g(x) + ih(x)$, where $g(x)$ and $h(x)$ are the real and imaginary parts of $f(x)$ respectively, then $g$ and $h$ are real linear functionals defined on $M$. 

Thus 
\[
\left| g (x) \right| \leqq \left| f (x) \right| \leqq p (x) \quad \text{ and } \quad \left| h (x) \right| \leqq \left| f (x) \right| \leqq p (x) \text{ on } M.
\]

Since, for each $x \in M$, 
\[
g (i x) + i h (i x) = f (i x) = i f (x) = i (g (x) + i h (x)) = - h (x) + i g (x),
\]
we have 
\[
h (x) = - g (i x) \quad \text{ for all } \quad x \in M.
\]

By the Theorem in Chapter~\ref{chap:IV}, \ref{sec:IV.1}, we can extend $g$ to a real linear functional $G$ defined on $X$ with the property that $G(x) \leqq p(x)$ on $X$. Hence $-G(x) = G(-x) \leqq p(-x) = p(x)$, and so $|G(x)| \leqq p(x)$. We define 
\[
F (x) = G (x) - i G (i x).
\]

Then, by $F(ix) = G(ix) - iG(-x) = G(ix) + iG(x) = iF(x)$, we easily see that $F$ is a complex linear functional defined on $X$. $F$ is an extension of $f$. For, $x \in M$ implies that 
\[
F (x) = G (x) - i G (i x) = g (x) - i g (i x) = g (x) + i h (x) = f (x).
\]

To prove $|F(x)| \leqq p(x)$, we write $F(x) = r e^{-i\theta}$ so that $|F(x)| = e^{i\theta} F(x) = F(e^{i\theta} x)$ is real positive; consequently $|F(x)| = |G(e^{i\theta} x)| \leqq p(e^{i\theta} x) = |e^{i\theta}| p(x) = p(x)$. 
\end{proof}


\yossection{5}{The Hahn-Banach Extension Theorem in Normed Linear Spaces}
\label{sec:IV.5}

\begin{theorem}\label{thm:IV.5.Theorem.1} Let $X$ be a normed linear space, $M$ a linear subspace of $X$ and $f_{1}$ a continuous linear functional defined on $M$. Then there exists a continuous linear functional $f$ defined on $X$ such that \begin{enumerate}[i)]
\item $f$ is an extension of $f_{1}$
\item $\|f_{1}\| = \|f\|$.
\end{enumerate}
\end{theorem}

\begin{proof}
 Set $p(x) = \|f_1\| \cdot \|x\|$. Then $p$ is a continuous semi-norm defined on $X$ such that $|f_1(x)| \leqq p(x)$ on $M$. There exists, by the Theorem in the preceding section $4$, a linear extension $f$ of $f_1$ defined on the whole space $X$ and such that $|f(x)| \leqq p(x)$. Thus $\|f\| \leqq \sup\limits_{\|x\| \leqq 1} p(x) = \|f_1\|$. On the other hand, since $f$ is an extension of $f_1$, we must have $\|f\| \geqq \|f_1\|$ and so we obtain $\|f_1\| = \|f\|$. 
\end{proof}

\yossub{An Application to Moment Problems}

\begin{theorem}\label{thm:IV.5.Theorem.2} Let $X$ be a normed linear space. Given a sequence of elements $\{x_{n}\} \subseteq X$, a sequence of complex numbers $\{\alpha_{n}\}$ and a positive number $\gamma$. Then a necessary and sufficient condition for the existence of a continuous linear functional $f$ on $X$ such that $f(x_{i}) = \alpha_{i} (i = 1, 2, \ldots)$ and $\|f\| \leqq \gamma$ is that the inequalities 
\[
\left| \sum\limits_{i=1} ^{n} \beta_{i} x _{i} \right| \leqq \gamma \left\| \sum\limits_{i=1} ^{n} \beta_{i} x _{i} \right\|
\]
hold for any choice of positive numbers $n$ and complex numbers $\beta_1, \beta_2, \ldots, \beta_n$. 
\end{theorem}

\begin{proof}
 The necessity of this condition is clear from the definition of $\|f\|$. We shall prove the sufficiency. Consider the set 
\[
X _{1} = \left\{z; z = \sum\limits_{i=1} ^{n} \beta_{i} x _{i} \quad \text{ where } n \text{ and } \beta \text{ are arbitrary } \right\}. \]

For two representations $z = \sum\limits_{i=1}^{n} \beta_i x_i = \sum\limits_{i=1}^{m} \beta_{i'} x_{i'}$ of the same element $z \in X_1$, we have, by the condition of the Theorem, 
\[
\left| \sum\limits_{i=1} ^{n} \beta_{i} \alpha_{i} - \sum\limits_{i=1} ^{m} \beta_ {i ^{\prime}} \alpha_ {i ^{\prime}} \right| \leqq \gamma \left\| \sum\limits_{i=1} ^{n} \beta_{i} x _{i} - \sum\limits_{i=1} ^{m} \beta_ {i ^{\prime}} x _ {i ^{\prime}} \right\| = 0.
\]

Thus a continuous linear functional $f_{1}$ is defined on $X_{1}$ by $f_{1}\left(\sum\limits_{i=1}^{n}\beta_{i}x_{i}\right)=\sum\limits_{i=1}^{n}\beta_{i}\alpha_{i}$. We have only to extend, by the preceding Theorem~\ref{thm:IV.5.Theorem.1}, $f_{1}$ to a continuous linear functional $f$ on $X$ with $\|f\|=\|f_{1}\|$. 
\end{proof}

\begin{remarku} As will be shown in section $9$ of this chapter, any continuous linear functional $f$ on $C[0,1]$ is representable as 
\[
f (x) = \int_{0} ^{1} x (t) m (\dd t)
\]
\end{remarku}

with a uniquely determined Baire measure $m$ on the interval [$0$, $1$]. Thus, if we take $x_{j}(t) = t^{j-1} (j = 1, 2, \ldots)$, Theorem~\ref{thm:IV.5.Theorem.2} gives the solvability condition of the moment problem: 
\[
\int_{0} ^{1} t ^{j-1} m (\dd t) = \alpha_{j} (j = 1, 2, \dots).
\]


\yossection{6}{The Existence of Non-trivial Continuous Linear Functionals}
\label{sec:IV.6}

\begin{theorem}\label{thm:IV.6.Theorem.1} Let $X$ be a real or complex linear topological space, $x_{0}$ a point of $X$ and $p(x)$ a continuous semi-norm on $X$. Then there exists a continuous linear functional $F$ on $X$ such that $F(x_{0}) = p(x_{0})$ and $|F(x)| \leqq p(x)$ on $X$. 
\end{theorem}

\begin{proof}
 Let $M$ be the set of all elements $\alpha x_{0}$, and define $f$ on $M$ by $f(\alpha x_{0}) = \alpha p(x_{0})$. Then $f$ is linear on $M$ and $|f(\alpha x_{0})| = |\alpha p(x_{0})| = p(\alpha x_{0})$ there. Thus there exists, by the Theorem in Chapter~\ref{chap:IV}, \ref{sec:IV.4}, an extension $F$ of $f$ such that $|F(x)| \leqq p(x)$ on $X$. Hence $F(x)$ is continuous at $x = 0$ with $p(x)$, and so, by the linearity of $F$, $F(x)$ is continuous at every point of $X$. 
\end{proof}

\begin{corollary}\label{cor:IV.6.Corollary.1} Let $X$ be a locally convex space and $x_{0} \neq 0$ be an element of $X$. Then there exists a continuous semi-norm $p$ on $X$ such that $p(x_{0}) \neq 0$. Thus, by Theorem~\ref{thm:IV.6.Theorem.1}, there exists a continuous linear functional $f_{0}$ on $X$ such that 
\[
f _{0} (x _{0}) = p (x _{0}) \neq 0 \quad \text{ and } \quad | f _{0} (x) | \leqq p (x) \text{ on } X.
\]
\end{corollary}

\begin{corollary}\label{cor:IV.6.Corollary.2} Let $X$ be a normed linear space and $x_{0} \neq 0$ be any element of $X$. Then there exists a continuous linear functional $f_{0}$ on $X$ such that 
\[
f _{0} (x _{0}) = \| x _{0} \| \quad \text{ and } \quad \| f _{0} \| = 1.
\]
\end{corollary}

\begin{proof}
 We take $\|x\|$ for $p(x)$ in Corollary~\ref{cor:IV.6.Corollary.1}. Thus $\|f_0\| \leqq 1$ from $|f_0(x)| \leqq \|x\|$. But, by $f_0(x_0) = \|x_0\|$, we must have the equality $\|f_0\| = 1$. 
\end{proof}

\begin{remarku} As above, we prove the following theorem by the Theorem in Chapter~\ref{chap:IV}, \ref{sec:IV.1}. 
\end{remarku}

\begin{theoremn}[Theorem~\ref{thm:IV.6.Theorem.1}']\yosnamedtag{1'}\label{thm:IV.6.Theorem.1prime} Let $X$ be a real linear topological space, $x_0$ a point of $X$ and $p(x)$ a real continuous functional on $X$ such that 
\[
p (x + y) \leqq p (x) + p (y) \quad \text{ and } \quad p (\alpha x) = \alpha p (x) \quad \text{ for } \quad \alpha \geqq 0.
\]

Then there exists a continuous real linear functional $F$ on $X$ such that $F(x_0) = p(x_0)$ and $-p(-x) \leqq F(x) \leqq p(x)$ on $X$.
\end{theoremn}

\begin{theorem}\label{thm:IV.6.Theorem.2} Let $X$ be a locally convex linear topological space. Let $M$ be a linear subspace of $X$, and $f$ a continuous linear functional on $M$. Then there exists a continuous linear functional $F$ on $X$ which is an extension of $f$. 
\end{theorem}

\begin{proof}
 Since $f$ is continuous on $M$ and $X$ is locally convex, there exists an open, convex, balanced neighbourhood of $0$, say $U$, of $X$ such that $x \in M \cap U$ implies $|f(x)| \leqq 1$. Let $p$ be the Minkowski functional of $U$. Then $p$ is a continuous semi-norm on $X$ and $U = \{x; p(x) < 1\}$. For any $x \in M$ choose $\alpha > 0$ so that $\alpha > p(x)$. Then $p(x/\alpha) < 1$ and so $|f(x/\alpha)| \leqq 1$, that is $|f(x)| \leqq \alpha$. We thus see, by letting $\alpha \downarrow p(x)$, that $|f(x)| \leqq p(x)$ on $M$. Hence, by the Theorem in Chapter~\ref{chap:IV}, \ref{sec:IV.4}, we obtain a continuous linear functional $F$ on $X$ such that $F$ is an extension of $f$ and $|F(x)| \leqq p(x)$ on $X$. 
\end{proof}

\begin{theorem}[S. \textsc{Mazur}]\label{thm:IV.6.Theorem.3} Let $X$ be a real or complex, locally convex linear topological space, and $M$ a closed convex balanced subset of $X$. Then, for any $x_0 \bar{\in} M$, there exists a continuous linear functional $f$ on $X$ such that $f_0(x_0) > 1$ and $|f_0(x)| \leqq 1$ on $M$. 
\end{theorem}

\begin{proof}
 Since $M$ is closed, there exists a convex, balanced neighbourhood $V$ of $0$ such that $M \cap (x_0 + V) = \emptyset$. Since $V$ is balanced and convex, we have $\left(M + \frac{V}{2}\right) \cap \left(x_0 + \frac{V}{2}\right) = \emptyset$. The set $\left(x_0 + \frac{V}{2}\right)$ being a neighbourhood of $x_0$, the closure $U$ of $\left(M + \frac{V}{2}\right)$ does not contain $x_0$. Since $M \ni \complement$, the closed convex balanced set $U$ is a neighbourhood of $0$, because $U$ contains $\frac{V}{2}$ as a subset. Let $p$ be the Minkowski functional of $U$. Since $U$ is closed, we have, for any $x_0 \bar{\in} U$, $p(x_0) > 1$ and $p(x) \leqq 1$, if $x \in U$. 

Thus there exists, by Corollary~\ref{cor:IV.6.Corollary.1} of Theorem~\ref{thm:IV.6.Theorem.1}, a continuous linear functional $f_0$ on $X$ such that $f_0(x_0) = p(x_0) > 1$ and $|f_0(x)| \leqq p(x)$ on $X$. Hence, in particular, $|f_0(x)| \leqq 1$ on $M$. 
\end{proof}

\begin{corollaryu} Let $M$ be a closed linear subspace of a locally convex linear topological space $X$. Then, for any $x_0 \in X - M$, there exists a continuous linear functional $f_0$ on $X$ such that $f_0(x_0) > 1$ and $f_0(x) = 0$ on $M$. Moreover, if $X$ is a normed linear space and if $\text{ dis } (x_0, M) > d$, then we may take $\|f_0\| \leqq 1 / d$. 
\end{corollaryu}

\begin{proof}
 The first part is clear from the linearity of $M$. The second part is proved by taking $U=\{x;\dist(x,M)\leqq d\}$ in the proof of Theorem~\ref{thm:IV.6.Theorem.3}. 
\end{proof}

\begin{remarku} As above, we prove the following theorem by virtue of Theorem~\ref{thm:IV.6.Theorem.1}'. 
\end{remarku}

\begin{theoremn}[Theorem~\ref{thm:IV.6.Theorem.3}' (S. \textsc{Mazur})]\yosnamedtag{3'}\label{thm:IV.6.Theorem.3prime} Let $X$ be a locally convex real linear topological space, and $M$ a closed convex subset of $X$ such that $M \ni 0$. Then, for any $x_0 \bar{\in} M$, there exists a continuous real linear functional $f_0$ on $X$ such that $f_0(x_0) > 1$ and $f_0(x) \leqq 1$ on $M$.
\end{theoremn}

\begin{theorem}[S. \textsc{Mazur}]\label{thm:IV.6.Theorem.4} Let $X$ be a locally convex linear topological space, and $M$ a convex balanced neighbourhood of $0$ of $X$. Then, for any $x_{0} \bar{\in} M$, there exists a continuous linear functional $f_{0}$ on $X$ such that 
\[
f _{0} (x _{0}) \geqq \sup\limits _{x \in M} | f _{0} (x) |.
\]
\end{theorem}

\begin{proof}
 Let $p$ be the Minkowski functional of $M$. Then $p(x_0) \geqq 1$ and $p(x) \leqq 1$ on $M$. $p$ is continuous since $M$ is a neighbourhood of $0$ of $X$. Thus there exists, by Corollary~\ref{cor:IV.6.Corollary.1} of Theorem~\ref{thm:IV.6.Theorem.1}, a continuous linear functional $f_0$ on $X$ such that 
\[
f _{0} (x _{0}) = p (x _{0}) \geqq 1 \quad \text{ and } \quad | f _{0} (x) | \leqq p (x) \leqq 1 \text{ on } M. \]
\end{proof}

\begin{theorem}[E. \textsc{Helly}]\label{thm:IV.6.Theorem.5} Let $X$ be a $B$-space, and $f_{1}, f_{2}, \ldots, f_{n}$ be a finite system of bounded linear functionals on $X$. Given $n$ numbers $\alpha_{1}, \alpha_{2}, \ldots, \alpha_{n}$. Then a necessary and sufficient condition that there exists, for each $\varepsilon > 0$, an element $x_{\varepsilon} \in X$ such that 
\[
f _{i} (x _{\varepsilon}) = \alpha_{i} \quad (i = 1, 2, \dots , n) \quad \text{ and } \quad \| x _{\varepsilon} \| \leqq \gamma + \varepsilon
\]
is that the inequality 
\[
\left| \sum\limits_{i=1} ^{n} \beta_{i} \alpha_{i} \right| \leqq \gamma \left\| \sum\limits_{i=1} ^{n} \beta_{i} f _{i} \right\|
\]
holds for any choice of $n$ numbers $\beta_1, \beta_2, \ldots, \beta_n$. 
\end{theorem}

\begin{proof}
 The necessity is clear from the definition of the norm of a continuous linear functional. We shall prove the sufficiency. We may assume, without losing the generality, that $f$'s are linearly independent; otherwise, we can work with a linearly independent subsystem of $\{f_i\}$ which spans the same subspace as the original system $\{f_i\}$. 


We consider the mapping $x \to \varphi(x) = (f_1(x), f_2(x), \ldots, f_n(x))$ of $X$ onto the Hilbert space $l^2(n)$ consisting of all vectors $x = (\xi_1, \xi_2, \ldots, \xi_n)$ normed by $\|x\| = \left(\sum\limits_{j=1}^{n} |\xi_j|^2\right)^{1/2}$. By the open mapping theorem in Chapter~\ref{chap:II}, \ref{sec:II.5}, the image $\varphi(S_\varepsilon)$ of the sphere $S_\varepsilon = \{x \in X; \|x\| \leqq \gamma + \varepsilon\}$ contains the vector $0$ of $l^2(n)$ as an interior point for every $\varepsilon > 0$. Let us suppose that $(\alpha_1, \alpha_2, \ldots, \alpha_n)$ does not belong to $\varphi(S_\varepsilon)$. Then, by Mazur's theorem given above, there exists a continuous linear functional $F$ on $l^2(n)$ such that 
\[
F \left(\left(\alpha_{1}, \alpha_{2}, \dots , \alpha_{n}\right)\right) \geqq \sup\limits _{\| x \| \leqq \gamma + \varepsilon} | F (\varphi (x)) |.
\]

Since $l^2(n)$ is a Hilbert space, the functional $F$ is given by an element $(\beta_1, \beta_2, \ldots, \beta_n) \in l^2(n)$ in such a way that $F((\alpha_1, \alpha_2, \ldots, \alpha_n)) = \sum\limits_{j=1}^{n} \alpha_j \beta_j$. Thus 
\[
\sum\limits_{j=1} ^{n} \alpha_{j} \beta_{j} \geqq \left| \sum\limits_{j=1} ^{n} f _{j} (x) \beta_{j} \right| \quad \text{ for } \quad \| x \| \leqq \gamma + \varepsilon .
\]

But the supremum of the right hand side for $\| x \| \leqq \gamma + \varepsilon$ is $= (\gamma + \varepsilon) \times \left\| \sum\limits_{j=1}^{n} f_j \beta_j \right\|$, and this contradicts the hypothesis of the Theorem.
\end{proof}

\yossection{7}{Topologies of Linear Maps}
\label{sec:IV.7}

Let $X$, $Y$ be locally convex linear topological spaces on the same scalar field (real or complex number field). We denote by $L(X, Y)$ the totality of continuous linear operators on $X$ into $Y$. $L(X, Y)$ is a linear space by 
\begin{align}
(\alpha T + \beta S) x& = \alpha T x + \beta S x \notag \\
T, S& \in L (X, Y) \notag \\
x& \in X \notag
\end{align}

We shall introduce various topologies for this linear space $L(X, Y)$. 

\begin{enumerate}[i)]
\item Simple Convergence Topology. This is the topology of convergence at each point of $X$ and thus it is defined by the family of semi-norms of the form 
\[
p (T) = p (T; x _{1}, x _{2}, \dots , x _{n}; q) = \sup\limits _{1 \leqq j \leqq n} q (T x _{j}),
\]
where $x_{1}, x_{2}, \ldots, x_{n}$ are an arbitrary finite system of elements of $X$ and $q$ an arbitrary continuous semi-norm on $Y$. $L(X, Y)$ endowed with this topology will be denoted by $L_{s}(X, Y)$. It is clearly a locally convex linear topological space.
\item Bounded Convergence Topology. This is the topology of uniform convergence on bounded sets of $X$. Thus it is defined by the family of 
semi-norms of the form 
\[
\phi (T) = \phi (T; B; q) = \sup\limits _{x \in B} q (T x)
\]
where $B$ is an arbitrary bounded set of $X$ and $q$ an arbitrary continuous semi-norm on $Y$. $L(X, Y)$ endowed with this topology will be denoted by $L_{b}(X, Y)$. It is clearly a locally convex linear topological space.
\end{enumerate}
Since any finite set of $X$ is bounded, the simple convergence topology is weaker than the bounded convergence topology, i.e., open sets of $L_{s}(X, Y)$ are also open sets of $L_{b}(X, Y)$, but not conversely. 

\begin{definition}\label{def:IV.7.Definition.1} If $X$, $Y$ are normed linear spaces, then the topology of $L_{s}(X, Y)$ is usually called the strong topology (of operators); the one of $L_{b}(X, Y)$ is called the uniform topology (of operators). 
\end{definition}

\yossub{Dual Spaces. Weak and Weak* Topologies}

\begin{definitionn}[Definition~\ref{def:IV.7.Definition.1}']\yosnamedtag{1'}\label{def:IV.7.Definition.1prime} In the special case when $Y$ is the real or complex number field topologized in the usual way, $L(X, Y)$ is called the dual space of $X$ and will be denoted by $X'$. Thus $X'$ is the set of all continuous linear functionals on $X$. The simple convergence topology is then called the weak* topology of $X'$; provided with this topology, $X'$ will sometimes be denoted by $X'_{w*}$ and we call it the weak* dual of $X$. The bounded convergence topology for $X'$ is called the strong topology of $X'$; provided with this topology, $X'$ is sometimes denoted by $X'_{s}$ and we call it the strong dual of $X$.
\end{definitionn}

\begin{definition}\label{def:IV.7.Definition.2} For any $x \in X$ and $x' \in X'$, we shall denote by $\langle x, x' \rangle$ or $x'(x)$ the value of the functional $x'$ at the point $x$. Thus the weak* topology of $X'$, i.e., the topology of $X_{w*}'$ is defined by the family of semi-norms of the form 
\[
p \left(x ^{\prime}\right) = p \left(x ^{\prime}; x _{1}, x _{2}, \dots , x _{n}\right) = \sup\limits _{1 \leqq j \leqq n} | \langle x _{j}, x ^{\prime} \rangle |,
\]
where $x_{1}, x_{2}, \ldots, x_{n}$ are an arbitrary finite system of elements of $X$. The strong topology of $X'$, i.e., the topology of $X_{s}'$ is defined by the family of semi-norms of the form 
\[
p \left(x ^{\prime}\right) = p \left(x ^{\prime}; B\right) = \sup\limits _{x \in B} | \langle x, x ^{\prime} \rangle |
\]
where $B$ is an arbitrary bounded set of $X$. 
\end{definition}

\begin{theorem}\label{thm:IV.7.Theorem.1} If $X$ is a normed linear space, then the strong dual space $X_s'$ is a $B$-space with the norm 
\[
\left\| f \right\| = \sup\limits _{\left\| x \right\| \leqq 1} | f (x) |.
\]
\end{theorem}

\begin{proof}
 Let $B$ be any bounded set of $X$. Then $\sup\limits_{b \in B} \|b\| = \beta < \infty$, and hence $\|f\| \leqq \alpha$ implies $p(f; B) = \sup\limits_{b \in B} |f(b)| \leqq \sup\limits_{\|x\| \leqq \beta} |f(x)| \leqq \alpha \beta$. 

On the other hand, the unit sphere $S = \{x; \|x\| \leqq 1\}$ of $X$ is a bounded set, and so $\|f\| = p(f; S)$. This proves that the topology of $X_s'$ is equivalent to the topology defined by the norm $\|f\|$. 

The completeness of $X_s'$ is proved as follows. Let a sequence $\{f_n\}$ of $X_s'$ satisfy $\lim\limits_{n,m \to \infty} \|f_n - f_m\| = 0$. Then, for any $x \in X$, $|f_n(x) - f_m(x)| \leqq \|f_n - f_m\| \cdot \|x\| \to 0$ (as $n, m \to \infty$), and hence a finite $\lim\limits_{n \to \infty} f_n(x) = f(x)$ exists. The linearity of $f$ is clear. The continuity of $f$ is proved by observing that $\lim\limits_{n \to \infty} f_n(x) = f(x)$ uniformly on the unit sphere $S$. Incidentally we have proved that $\lim\limits_{n \to \infty} \|f_n - f\| = 0$. 
\end{proof}

Similarly we can prove 

\begin{theorem}\label{thm:IV.7.Theorem.2} If $X$, $Y$ are normed linear spaces, then the uniform topology (of operators) $L_{b}(X, Y)$ is defined by the operator norm 
\[
\left\| T \right\| = \sup\limits _{\left\| x \right\| \leqq 1} \left\| T x \right\|.
\]
\end{theorem}

\begin{definition}\label{def:IV.7.Definition.3} We define the weak topology of a locally convex linear topological space $X$ by the family of semi-norms of the form 
\[
p (x) = p \left(x; x _{1} ^{\prime}, x _{2} ^{\prime}, \dots , x _{n} ^{\prime}\right) = \sup\limits _{1 \leqq j \leqq n} | \langle x, x _{j} ^{\prime} \rangle |,
\]
where $x_{1}^{\prime}, x_{2}^{\prime}, \ldots, x_{n}^{\prime}$ are an arbitrary finite system of elements of $X^{\prime}$. Endowed with this topology, $X$ is sometimes denoted by $X_{w}$. 
\end{definition}


\yossection{8}{The Embedding of \texorpdfstring{$X$}{X} in its Bidual Space \texorpdfstring{$X''$}{X''}}
\label{sec:IV.8}

We first prove 

\begin{theorem}[S. \textsc{Banach}]\label{thm:IV.8.Theorem.1} Let $X$ be a locally convex linear topological space, and $X'$ its dual space. A linear functional $f(x')$ on $X'$ is of the form 
\[
f (x ^{\prime}) = \langle x _{0}, x ^{\prime} \rangle
\]
with a certain $x_0 \in X$ iff $f(x')$ is continuous in the weak* topology of $X'$. 
\end{theorem}

\begin{proof}
 The ``only if'' part is clear since $\left|\langle x_{0}, x^{\prime}\rangle\right|$ is one of the seminorms defining the weak* topology of $X^{\prime}$. The ``if'' part will be proved as follows. The continuity of $f(x^{\prime})$ in the weak* topology of $X^{\prime}$ implies that there exists a finite system of points $x_{1}, x_{2}, \ldots, x_{n}$ such that $\left|f(x^{\prime})\right| \leqq \sup\limits_{1 \leqq i \leqq n} \left|\langle x_{j}, x^{\prime}\rangle\right|$. Thus 
\[
f _{i} (x ^{\prime}) = \langle x _{i}, x ^{\prime} \rangle = 0 (i = 1, 2, \dots , n) \text{ implies } f (x ^{\prime}) = 0. \]

Consider the linear map $L: X' \to l^2(n)$, defined by 
\[
L \left(x ^{\prime}\right) = \left(f _{1} (x ^{\prime}), f _{2} (x ^{\prime}), \dots , f _{n} (x ^{\prime})\right).
\]
$L(x_1') = L(x_2')$ implies $L(x_1' - x_2') = 0$ so that $f_i(x_1' - x_2') = 0 (i = 1, 2, \ldots, n)$ and hence $f(x_1' - x_2') = 0$. Hence we may define a continuous linear map $F$ defined on the linear subspace $L(X')$ of $l^2(n)$ by 
\[
F \left(L \left(x ^{\prime}\right)\right) = F \left(f _{1} (x ^{\prime}), \dots , f _{n} (x ^{\prime})\right) = f (x ^{\prime}).
\]

This map can be extended to a continuous linear functional defined on the whole space $l^{2}(n)$; the extension is possible since $l^{2}(n)$ is of finite dimension (easier than using the Hahn-Banach extension theorem in infinite dimensional linear spaces). We denote this extension by the same letter $F$. Writing 
\[
(y _{1}, y _{2}, \dots , y _{n}) = \sum\limits_{j=1} ^{n} y _{j} e _{j}, \quad \text{ where } \quad e _{i} = (0, 0, \dots , 0, 1, 0, 0, \dots , 0)
\]
with $1$ at the $j$-th coordinate, 
we easily see that 
\[
F \left(y _{1}, y _{2}, \dots , y _{n}\right) = \sum\limits_{j=1} ^{n} y _{j} \alpha_{j}, \alpha_{j} = F \left(e _{j}\right).
\]

Therefore 
\[
f \left(x ^{\prime}\right) = \sum\limits_{j=1} ^{n} \alpha_{j} f _{j} \left(x ^{\prime}\right) = \sum\limits_{j=1} ^{n} \alpha_{j} \left\langle x _{j}, x ^{\prime} \right\rangle = \left\langle \sum\limits_{j=1} ^{n} \alpha_{j} x _{j}, x ^{\prime} \right\rangle .
\]
\end{proof}
\begin{corollaryu} Each $x_0 \in X$ defines a continuous linear functional $f_0(x')$ on $X_s'$ by $f_0(x') = \langle x_0, x' \rangle$. The mapping 
\[
x _{0} \rightarrow f _{0} = J x _{0}
\]
of $X$ into $(X_s')_s'$ satisfies the conditions 
\[
J \left(x _{1} + x _{2}\right) = J x _{1} + J x _{2}, J (\alpha x) = \alpha J (x).
\]
\end{corollaryu}

\begin{theorem}\label{thm:IV.8.Theorem.2} If $X$ is a normed linear space, then the mapping $J$ is isometric, i.e., $\|Jx\| = \|x\|$. 
\end{theorem}

\begin{proof}
 We have $|f_0(x')| = |\langle x_0, x' \rangle| \leqq \|x_0\| \cdot \|x'\|$ so that $\|f_0\| \leqq \|x_0\|$. On the other hand, if $x_0 \neq 0$, then there exists, by Corollary~\hyperref[cor:IV.6.Corollary.2]{$2$} of Theorem~\ref{thm:IV.6.Theorem.1} in Chapter~\ref{chap:IV}, \ref{sec:IV.6}, an element $x_0' \in X'$ such that $\langle x_0, x_0' \rangle = \|x_0\|$ and $\|x_0'\| = 1$. Hence $f_0(x_0') = \langle x_0, x_0' \rangle = \|x_0\|$ so that $\|f_0\| \geqq \|x_0\|$. We have thus proved $\|Jx\| = \|x\|$. 
\end{proof}

\begin{remarku} As the strong dual space of $X_{s}^{\prime}$, the space $(X_{s}^{\prime})_{s}^{\prime}$ is a $B$-space. Hence a normed linear space $X$ may be considered as a linear subspace of the $B$-space $(X_{s}^{\prime})_{s}^{\prime}$ by the embedding $x \to Jx$. Therefore the strong closure of $JX$ in the $B$-space $(X_{s}^{\prime})_{s}^{\prime}$ gives a concrete construction of the completion of $X$. 
\end{remarku}

\begin{definition}\label{def:IV.8.Definition.4} $A$ normed linear space $X$ is said to be reflexive if $X$ may be identified with its second dual or the bidual $(X_{s}^{\prime})_{s}^{\prime}$ by the correspondence $x \leftrightarrow Jx$ above. We know already (see Chapter~\ref{chap:III}, \ref{sec:III.6}) that a Hilbert space is reflexive. As remarked above, $(X_{s}^{\prime})_{s}^{\prime}$ is a $B$-space and so any reflexive normed linear space must be a $B$-space. 
\end{definition}

\begin{theorem}\label{thm:IV.8.Theorem.3} Let $X$ be a $B$-space and $x_0''$ any bounded linear functional on $X_s'$. Then, for any $\varepsilon > 0$ and any finite system of elements $f_1, f_2, \ldots, f_n$ of $X_s'$, there exists an $x_0 \in X$ such that 
\[
\left\| x _{0} \right\| \leqq \left\| x _{0} ^{\prime \prime} \right\| + \varepsilon \quad \text{ and } \quad f _{j} (x _{0}) = x _{0} ^{\prime \prime} (f _{j}) (j = 1, 2, \dots , n).
\]
\end{theorem}

\begin{proof}
 We apply Helly's theorem $5$ in Chapter~\ref{chap:IV}, \ref{sec:IV.6}. For any system of numbers $\beta_{1}, \beta_{2}, \ldots, \beta_{n}$, we have 
\begin{equation}\left| \sum\limits_{j=1} ^{n} \beta_{j} \alpha_{j} \right| = \left| \sum\limits_{j=1} ^{n} \beta_{j} x _{0} ^{\prime \prime} (f _{j}) \right| = \left| x _{0} ^{\prime \prime} \left(\sum\limits_{j=1} ^{n} \beta_{j} f _{j}\right) \right| \leqq \gamma \cdot \left\| \sum\limits_{j=1} ^{n} \beta_{j} f _{j} \right\|, \text{ wehre } \gamma = \| x _{0} ^{\prime \prime} \|, \alpha_{j} = x _{0} ^{\prime \prime} (f _{j}), \notag\end{equation}
and hence, again by Helly's theorem, we obtain an $x_0 \in X$ with the estimate $\| x_0 \| \leqq \gamma + \varepsilon = \| x_0'' \| + \varepsilon$ and $x_0(f_j) = \alpha_j (j = 1, 2, \ldots, n)$. 
\end{proof}

\begin{corollaryu} The unit sphere $S = \{x \in X; \| x \| \leqq 1\}$ of a $B$-space $X$ is dense in the unit sphere of $(X_s')_s'$ in the weak* topology of $(X_s')'$. 
\end{corollaryu}


\yossection{9}{Examples of Dual Spaces}
\label{sec:IV.9}

\begin{example}\label{ex:IV.9.Example.1} $(c_0)' = (l^1)$. 
\end{example}

To any $f \in (c_0)'$, there corresponds a uniquely determined $y_f = \{\eta_n\} \in (\mathcal{I}^1)$ such that, for all $x = \{\xi_n\} \in (c_0)$, 
\begin{equation}\label{eq:IV.9.1}
\langle x, f \rangle = \sum\limits_{n=1} ^{\infty} \xi_{n} \eta_{n} \text{ and } \| f \| = \| y _{f} \|.
\end{equation}

And conversely, any $y = \{\eta_n\} \in (l^1)$ defines an $f_y \in (c_0)'$ such that, for any $x = \{\xi_n\} \in (c_0)$, 
\begin{equation}\tag{1'}\label{eq:IV.9.1'}
\langle x, f _{y} \rangle = \sum\limits_{n=1} ^{\infty} \xi_{n} \eta_{n} \text{ and } \| f _{y} \| = \| y \|.
\end{equation}

\begin{proof}
 Let us define the unit vector $e_k$ by 
\[
e _{k} = \overbrace {(0 , 0 , \dots , 0} ^{k-1}, 1, 0, 0, \dots) (k = 1, 2, \dots). \]

For any $x = \{\xi_n\} \in (c_0)$ and $f \in (c_0)'$, we have, by $s\text{- }\lim\limits_{k \to \infty} \sum\limits_{n=1}^{k} \xi_n e_n = x$. 
\[
\langle x, f \rangle = \lim\limits _{k \rightarrow \infty} \left\langle \sum\limits_{n=1} ^{k} \xi_{n} e _{n}, f \right\rangle = \lim\limits _{k \rightarrow \infty} \sum\limits_{n=1} ^{k} \xi_{n} \eta_{n}, \eta_{n} = f (e _{n}).
\]

Let $\eta_n = \varepsilon_n |\eta_n|$ for $\eta_n \neq 0$, and $\varepsilon_n = \infty$ for $\eta_n = 0$. Take $x^{(n_0)} = \{\xi_n\} \in (c_0)$ in such a way that $\xi_n = \varepsilon_n^{-1}$ for $n \leqq n_0$, and $\xi_n = 0$ for $n > n_0$. Then $\|x^{(n_0)}\| \leqq 1$ and so $\|f\| = \sup\limits_{\|x\| \leqq 1} |\langle x, f \rangle| \geqq |\langle x^{(n_0)}, f \rangle| = \sum\limits_{n=1}^{n_0} |\eta_n|$. Thus, by letting $n_0 \to \infty$, we see that $y_f = \{\eta_n\} \in (l^1)$ and $\|y_f\| = \sum\limits_{n=1}^{\infty} |\eta_n| \leqq \|f\|$. 

If, conversely, $y = \{\eta_n\} \in (l^1)$, then $\left|\sum\limits_{n=1}^{\infty} \xi_n \eta_n\right| \leqq \|x\| \cdot \|y\|$ for all $x = \{\xi_n\} \in (c_0)$, and so $y$ defines an $f_y \in (c_0)'$ and $\|f_y\| \leqq \|y\|$. 
\end{proof}

\begin{example}\label{ex:IV.9.Example.2} $(c)' = (l^1)$. 
\end{example}

For any $x = \{\xi_n\} \in (c)$, we have the representation 
$x = \xi_0e_0 + s\text{- }\lim\limits_{k\to \infty}\sum\limits_{n=1}^{k}(\xi_n - \xi_0)e_n,$ where $\xi_0 = \lim\limits_{n\to \infty}\xi_n,e_0 = (1,1,1,\ldots)$. Thus, for any $f\in (c)^{\prime}$, we have 
\begin{equation}\label{eq:IV.9.2}
\langle x, f \rangle = \xi_{0} \left\langle e _{0}, f \right\rangle + \lim\limits _{k \rightarrow \infty} \left\langle \sum\limits_{n=1} ^{k} \left(\xi_{n} - \xi_{0}\right) e _{n}, f \right\rangle = \xi_{0} \eta_{0} ^{\prime} + \sum\limits_{n=1} ^{\infty} \left(\xi_{n} - \xi_{0}\right) \eta_{n},
\end{equation}
where $\eta_0' = \langle e_0, f \rangle$ and $\eta_n = \langle e_n, f \rangle$ ( $n = 1, 2, \ldots$). As above, we may take $x^{(n_0)} = \{\xi_n\} \in (c_0) \subseteq (c)$ which satisfies 
\[
\left\| x ^ {\left(n _{0}\right)} \right\| \leqq 1, \xi_{0} = \lim\limits _{n \rightarrow \infty} \xi_{n} = 0 \text{ and } \langle x ^ {\left(n _{0}\right)}, f \rangle = \sum\limits_{n=1} ^ {n _{0}} | \eta_{n} |.
\]

Hence, by $|\langle x^{(n_0)}, f \rangle| \leqq \|x^{(n_0)}\| \|f\|$, we see that $\{\eta_n\}^\infty \in (l^1)$. We set $\eta_0' - \sum\limits_{n=1}^\infty \eta_n = \eta_0$. Then, by \eqref{eq:IV.9.2}, we have 
$\langle x, f \rangle = \xi_0 \eta_0 + \sum\limits_{n=1}^{\infty} \xi_n \eta_n$, where $x = \{\xi_n\} \in (c)$ and $\xi_0 = \lim\limits_{n \to \infty} \xi_n$. \eqref{eq:IV.9.2} Let $\eta_n = \varepsilon_n | \eta_n|$ for $\eta_n \neq 0$, and $\varepsilon_n = \infty$ for $\eta_n = 0$ ( $n = 0, 1, 2, \ldots$). Take $x = \{\xi_n\} \in (c)$ such that 
$\xi_{n} = \varepsilon_{n}^{-1}$ if $n\leqq n_0$, and $\xi_{n} = \varepsilon_{0}^{-1}$ if $n > n_0$ 

Then $\|x\| \leqq 1$, $\xi_0 = \lim\limits_{n \to \infty} \xi_n = \varepsilon_0^{-1}$, and $\langle x, f \rangle = |\eta_0| + \sum\limits_{n=1}^{n_0} |\eta_n| + \varepsilon_0^{-1} \sum\limits_{n=n_0+1}^{\infty} \eta_n$. Hence, we must have $|\eta_0| + \sum\limits_{n=1}^{\infty} |\eta_n| \leqq \|f\|$. 

If, conversely, $y = \{\eta_n\}_0^\infty$ is such that $\|y\| = |\eta_0| + \sum\limits_{n=1}^{\infty} |\eta_n| < \infty$, then 
\[
\eta_{0} \cdot \lim\limits _{n \rightarrow \infty} \xi_{n} + \sum\limits_{n=1} ^{\infty} \xi_{n} \eta_{n}, \text{ where } x = \{\xi_{n} \} _{1} ^{\infty} \in (c),
\]
defines an $f_y \in (c)'$ such that $\|f_y\| \leqq |\eta_0| + \sum\limits_{n=1}^{\infty} |\eta_n|$. 

Therefore, we have proved that $(c)' = (l^1)$ as explained above. 

\begin{example}\label{ex:IV.9.Example.3} $L^p(S, \mathfrak{B}, m)' = L^q(S, \mathfrak{B}, m)$ ( $1 \leqq p < \infty$ and $p^{-1} + q^{-1} = 1$). To any $f \in L^p(S)'$, there corresponds a $y_f \in L^q(S)$ such that 
\begin{equation}\label{eq:IV.9.3}
\langle x, f \rangle = \int_{S} x (s) y _{f} (s) m (\dd s) \text{ for all } x \in L ^{p} (S) \text{ and } \| f \| = \| y _{f} \|,
\end{equation}
and conversely, any $y \in L^{q}(S)$ defines an $f_{y} \in L^{p}(S)'$ such that 
\begin{equation}\tag{3'}\label{eq:IV.9.3'}
\langle x, f _{y} \rangle = \int_{S} x (s) y (s) m (\dd s) \text{ for all } x \in L ^{p} (S) \text{ and } \| f _{y} \| = \| y \|.
\end{equation}
\end{example}

\begin{proof}
 Let $S = \bigcup\limits_{j=1}^{\infty} B_j$ with $0 < m(B_j) < \infty$ and set $B^{(n)} = \bigcup\limits_{j=1}^{n} B_j$. For a fixed $n$, the defining function $C_B(s)$ of the set $B \subseteq B^{(n)}$ is $\in L^p(S)$. Thus the set function $\psi(B) = \langle C_B, f \rangle$ is $\sigma$-additive and $m$-absolutely continuous in $B \subseteq B^{(n)}$. By the differentiation theorem of Lebesgue-Nikodym (see Chapter~\ref{chap:III}, \ref{sec:III.8}), there exists a $y_n(s) \in L^1(B^{(n)}, \mathfrak{B}^{(n)}, m)$ such that 
\[
\psi (B) = \int_{B} y _{n} (s) m (\dd s) \text{ whenever } B \subseteq B ^{(n)},
\]
the family $\mathfrak{B}^{(n)}$ of sets being defined by $\mathfrak{B}^{(n)} = \{B \cap B^{(n)}; B \in \mathfrak{B}\}$. Therefore, by setting $y(s) = y_n(s)$ for $s \in B^{(n)}$, we have 
\[
\langle C _{B}, f \rangle = \int_{B} y (s) m (\dd s) \text{ for } B \in \mathfrak {B} ^{(n)} (n = 1, 2, \dots). \]

Hence, for any finitely-valued function $x$ with support in some $B^{(n)}$, 
\begin{equation}\label{eq:IV.9.4}
\langle x, f \rangle = \int_{S} x (s) y (s) m (\dd s).
\end{equation}

Let $x \in L^{p}(S)$ and put 
\[
\begin{aligned} x _{n} (s) & = x (s) \text{ if } | x (s) | \leqq n \text{ and } s \in B ^{(n)}, \\ & = 0 \text{ otherwise }. \end{aligned}
\]

We decompose the set $\{z; |z| \leqq n\}$ of the complex plane into a finite number of disjoint Baire sets $M_{n,k,t}$ ( $t = 1, 2, \ldots, d_{k,n}$) of diameters $\leqq 1/k$. Set, for the $x_n(s) \in L^\infty(S, \mathfrak{B}, m)$, 
$x_{n,k}(s) =$ a constant $z$, such that $z\in$ (the closure $M_{n,k,t}^{a}$) and $|z| = \inf\limits_{w\in Mn,k,t}|w|$ whenever $x_{n}(s)\in M_{n,k,t}$. 

Then $|x_{n,k}(s)| \leqq |x_n(s)|$ and $\lim\limits_{k \to \infty} x_{n,k}(s) = x_n(s)$ and so, by the Lebesgue-Fatou Lemma, $\slim\limits_{k \to \infty} x_{n,k} = x_n (n = 1, 2, \ldots)$. Thus, again by the Lebesgue-Fatou Lemma, 
\begin{equation}\label{eq:IV.9.5}
\begin{aligned} \langle x _{n}, f \rangle & = \lim\limits _{k \to \infty} \langle x _{n,k}, f \rangle = \lim\limits _{k \to \infty} \int_{S} x _{n,k} (s) y (s) m (\dd s) \\ & = \int_{S} \lim\limits _{k \to \infty} x _{n,k} (s) \cdot y (s) m (\dd s) = \int_{S} x _{n} (s) y (s) m (\dd s). \end{aligned}
\end{equation}

Since \slim $x_{n} = x$, we see that $\langle x, f \rangle = \lim\limits_{n \to \infty} \langle x_{n}, f \rangle$. We put, for any complex number $z$, $a(z) = e^{-i\theta}$ if $z = r e^{i\theta}$ and $a(0) = 0$. Then $\|x\| \geqq \|(|x_{n}| \cdot a(y))\|$ and so 
\[
\left\| f \right\| \left\| x \right\| \geqq \langle \left| x _{n} \right| a (y), f \rangle = \int_{S} \left| x _{n} (s) \right| \cdot | y (s) | m (\dd s).
\]

Thus, by the Lebesgue-Fatou Lemma, $\|f\| \cdot \|x\| \geqq \int_{S} |x(s)\|y(s)| m(\dd s)$ and so the function $x(s)y(s)$ belongs to $\in L^{1}(S)$. Therefore, letting 
$n\to \infty$ in \eqref{eq:IV.9.5}, we obtain 
\[
\langle x, f \rangle = \int_{S} x (s) y (s) m (\dd s) \text{ whenever } x \in L ^{p} (S).
\]

We shall show that $y \in L^q(S)$. To this purpose, set 
\[
\begin{aligned} y _{n} (s) & = y (s) \text{ if } | y (s) | \leqq n \text{ and } s \in B ^{(n)}, \\ & = 0 \text{ otherwise }. \end{aligned}
\]

Then $y_{n} \in L^{q}(S)$ and, as proved above, 
\[
\begin{aligned}
&\left\| f \right\| \cdot \left\| x \right\|  \geqq \langle | x | \cdot a (y), f \rangle = \int_{S} | x (s) | | y (s) | m (\dd s) \geqq \int_{S} | x (s) | | y _{n} (s) | m (\dd s).
\end{aligned}
\]

If we take $x(s) = |y_n(s)|^{q/p}$ and apply Hölder's equality, we obtain 
\[
\int_{S} | x (s) | | y _{n} (s) | m (\dd s) = \left(\int_{S} | x (s) | ^{p} m (\dd s)\right) ^{1/p} \left(\int_{S} | y _{n} (s) | ^{q} m (\dd s)\right) ^{1/q}.
\]

Hence $\|f\| \geqq \|y_n\| = \left( \int_S |y_n(s)|^q m(\dd s) \right)^{1/q}$, with the understanding that, when $p = 1$ we have $\|f\| \geqq \|y_n\| = \text{ essential sup }_{s \in S} |y_n(s)|$. 

Therefore, by letting $n \to \infty$ and applying the Lebesgue-Fatou Lemma, we see that $y \in L^{q}(S)$ and $\| f \| \geqq \| y \|$. On the other hand, any $y \in L^{q}(S)$ defines an $f \in L^{p}(S)'$ by $\langle x, f \rangle = \int_{S} x(s) y(s) m(\dd s)$ as may be seen by Hölder's inequality, and this inequality shows that $\| f \| \leqq \| y \|$. 
\end{proof}

\begin{remarku} We have incidentally proved that 
\[
(l ^{p}) ^{\prime} = (l ^{q}) (1 \leqq p < \infty \text{ and } p ^{-1} + q ^{-1} = 1).
\]
\end{remarku}

\begin{example}\label{ex:IV.9.Example.4} Let the measure space $(S, \mathfrak{B}, m)$ with $m(S) < \infty$ have the property that, for any $B \in \mathfrak{B}$ with $0 < m(B) = \delta < \infty$ and positive integer $n$, there exists a subset $B_n$ of $B$ such that $\delta(n + 1)^{-1} \leqq m(B_n) \leqq \delta n^{-1}$. Then no other continuous linear functionals $\in M(S, \mathfrak{B}, m)$ ' than the zero functional can exist. 
\end{example}

\begin{proof}
 Any $x \in L^{1}(S, \mathfrak{B}, m)$ belongs to $M(S, \mathfrak{B}, m)$ and the topology of $L^{1}(S, \mathfrak{B}, m)$ is stronger than that of $M(S, \mathfrak{B}, m)$. Thus any $f \in M(S, \mathfrak{B}, m)'$, when restricted to the functions of $L^{1}(S, \mathfrak{B}, m)$, defines a continuous linear functional $f_{0} \in L^{1}(S, \mathfrak{B}, m)'$. Thus there exists a $y \in L^{\infty}(S, \mathfrak{B}, m)$ such that 
\[
\langle x, f \rangle = \langle x, f _{0} \rangle = \int_{S} x (s) y (s) m (\dd s) \text{ whenever } x \in L ^{1} (S, \mathfrak {B}, m). \]

Since $L^1(S, \mathfrak{B}, m)$ is dense in $M(S, \mathfrak{B}, m)$ in the topology of $M(S, \mathfrak{B}, m)$, the condition $f \neq 0$ implies that $f_0 \neq 0$. Thus there exists an $\varepsilon > 0$ such that $B = \{s; |y(s)| \geqq \varepsilon\}$ has its measure $m(B) = \delta > 0$. Let $B_n \subseteq B$ be as in the hypothesis, and let $y(s) = r e^{i\theta}$ for $s \in B$. Set $x_{n}(s) = e^{-i\theta}$ for $s \in B_{n}$ and $x_{n}(s) = 0$ otherwise. Then $z_{n}(s) = nx_{n}(s)$ converges to $0$ asymptotically, that is, \slim $z_{n} = 0$ in $M(S, \mathfrak{B}, m)$. But 
\[
\lim\limits _{n \rightarrow \infty} \langle z _{n}, f \rangle = \lim\limits _{n \rightarrow \infty} \langle z _{n}, f _{0} \rangle = \lim\limits _{n \rightarrow \infty} \int_{S} z _{n} (s) y (s) m (\dd s) \geqq \delta \varepsilon > 0,
\]
contrary to the continuity of the functional $f$. 
\end{proof}

\begin{example}\label{ex:IV.9.Example.5} $L^{\infty}(S, \mathfrak{B}, m)'$. 
\end{example}

Let an $f \in L^{\infty}(S, \mathfrak{B}, m)'$ be given and set, for any $B \in \mathfrak{B}$, $f(C_B) = \psi(B)$ where $C_B(s)$ is the defining function of the set $B$. We have then: 
\begin{equation}\label{eq:IV.9.6}
B _{1} \cap B _{2} = \emptyset \text{ implies } \psi (B _{1} + B _{2}) = \psi (B _{1}) + \psi (B _{2}),
\end{equation}
that is $\psi$ is finitely additive, 
the real part $\psi_{1}(B)$ and the imaginary part $\psi_{2}(B)$ of 
$\psi (B)$ are of bounded total variation, that is, $\sup\limits_B |\psi_i(B)|$ 
\begin{equation}\label{eq:IV.9.7}
< \infty (i = 1, 2),
\end{equation}
$\psi$ is $m$-absolutely continuous, that is $m(B) = 0$ implies 
\begin{equation}\label{eq:IV.9.8}
\psi (B) = 0.
\end{equation}

The condition \eqref{eq:IV.9.6} is a consequence of the linearity of $f$, and \eqref{eq:IV.9.7} and \eqref{eq:IV.9.8} are clear from $|\psi(B)| \leqq \|f\| \cdot \|C_B\|$. 

For any $x \in L^{\infty}(S, \mathfrak{B}, m)$, we consider a partition of the sphere $\{z; |z| \leqq \|x\|\}$ of the complex plane into a finite system of disjoint Baire sets $A_{1}, A_{2}, \ldots, A_{n}$ of diameters $\leqq \varepsilon$. If we set $B_{i} = \{s \in S; x(s) \in A_{i}\}$, then, no matter what point $\alpha_{i}$ we choose from $A_{i} (i = 1, 2, \ldots, n)$, we 
\[
\left\| x - \sum\limits_{i=1} ^{n} \alpha_{i} C _ {B _{i}} \right\| \leqq \varepsilon ,
\]
and so 
\[
\left| f (x) - \sum\limits_{i=1} ^{n} \alpha_{i} \psi (B _{i}) \right| \leqq \| f \| \cdot \varepsilon .
\]

Thus, if we let $n \to \infty$ in such a way that $\varepsilon \downarrow 0$, we obtain 
\begin{equation}\label{eq:IV.9.9}
f (x) = \lim \sum_{i=1}^{n} \alpha_{i} \psi (B _{i}),
\end{equation}
independently of the manner of partition $\{z; |z| \leqq \|x\|\} = \sum\limits_{i=1}^{n} A_i$ and choice of points $\alpha$'s. The limit on the right of \eqref{eq:IV.9.9} is called Radon's integral of $x(s)$ with respect to the finitely additive measure $\psi$. Thus 
\begin{equation}\label{eq:IV.9.10}
f(x) = \int_{S} x(s) \psi(\dd s) \text{ (Radon's integral) whenever } x \in L^{\infty}(S, \mathfrak{B}, m),
\end{equation}
and so 
\begin{equation}\label{eq:IV.9.11}
\left\| f \right\| = \sup\limits _ {\text{ ess. } \sup | x (s) | \leqq 1} \left| \int_{S} x (s) \psi (\dd s) \right|.
\end{equation}

Conversely, it is easy to see that any $\psi$ satisfying \eqref{eq:IV.9.6}, \eqref{eq:IV.9.7} and \eqref{eq:IV.9.8} defines an $f \in L^{\infty}(S, \mathfrak{B}, m)'$ through \eqref{eq:IV.9.10} and that \eqref{eq:IV.9.11} is true. 

Therefore, we have proved that $L^{\infty}(S,\mathfrak{B},m)'$ is the space of all set functions $\psi$ satisfying \eqref{eq:IV.9.6}, \eqref{eq:IV.9.7} and \eqref{eq:IV.9.8} and normed by the right hand side of \eqref{eq:IV.9.11}, the so-called total variation of $\psi$. 

\begin{remarku} We have so far proved that $L^p(S, \mathfrak{B}, m)$ is reflexive when $1 < p < \infty$. However, the space $L^1(S, \mathfrak{B}, m)$ is, in general, not reflexive. 
\end{remarku}

\begin{example}\label{ex:IV.9.Example.6} $C(S)'$.
\end{example}

Let $S$ be a compact topological space. Then the dual space $C(S)'$ of the space $C(S)$ of complex-valued continuous functions on $S$ is given as follows. To any $f \in C(S)'$, there corresponds a uniquely determined complex Baire measure $\mu$ on $S$ such that 
\begin{equation}\label{eq:IV.9.12}
f (x) = \int_{S} x (s) \mu (\dd s) \text{ whenever } x \in C (S),
\end{equation}
and hence 
\begin{equation}\label{eq:IV.9.13}
\left\| f \right\| = \sup\limits _ {\sup\limits _{s} | x (s) | \leqq 1} \left| \int_{S} x (s) \mu (\dd s) \right| = \text{ the total variation of } \mu \text{ on } S.
\end{equation}

Conversely, any Baire measure $\mu$ on $S$ such that the right side of \eqref{eq:IV.9.13} is finite, defines a continuous linear functional $f \in C(S)'$ through \eqref{eq:IV.9.12} and we have \eqref{eq:IV.9.13}. Moreover, if we are concerned with a real functional $f$ on a real $B$-space $C(S)$, then the corresponding measure $\mu$ is real-valued; if, moreover $f$ is positive, in the sense that $f(x) \geqq 0$ for non-negative functions $x(s)$, then the corresponding measure $\mu$ is positive, i.e., $\mu(B) \geqq 0$ for every $B \in \mathfrak{B}$. 

\begin{remarku} The result stated above is known as the F. Riesz-A. Markov-S. \textsc{Kakutani} theorem, and is one of the fundamental theorems in topological measures. For the proof, the reader is referred to standard text books on measure theory, e.g., P. R. \textsc{Halmos} \cite{halmos92} and N. \textsc{Dunford}-J. \textsc{Schwartz} \cite{dunford44}. 
\end{remarku}

\yossub{References for Chapter IV}

For the Hahn-Banach theorems and related topics, see \textsc{Banach} \cite{banach9}, \textsc{Bourbaki} \cite{bourbaki27} and \textsc{Köthe} \cite{konishi164}. It was \textsc{Mazur} \cite{mazur192} who noticed the importance of convex sets in normed linear spaces. The proof of Helly's theorem given in this book is due to Y. \textsc{Mimura} (unpublished). 
   % IV
% ==== $V$ ====
\yoschapter{V}{Strong Convergence and Weak Convergence}
\label{chap:V}

In this chapter, we shall be concerned with certain basic facts pertaining to strong-, weak- and weak* convergences, including the comparison of the strong notion with the weak notion, e.g., strong- and weak measurability, and strong- and weak analyticity. We also discuss the integration of $B$-space-valued functions, that is, the theory of Bochner's integrals. The general theory of weak topologies and duality in locally convex linear topological spaces will be given in the Appendix. 


\yossection{1}{The Weak Convergence and The Weak* Convergence}
\label{sec:V.1}

\yossub{Weak Convergence}

\begin{definition}\label{def:V.1.Definition.1} $A$ sequence $\{x_{n}\}$ in a normed linear space $X$ is said to be weakly convergent if a finite $\lim\limits_{n\to \infty}f(x_n)$ exists for each $f\in X_s'$; $\{x_{n}\}$ is said to converge weakly to an element $x_{\infty}\in X$ if $\lim\limits_{n\to \infty}f(x_n) = f(x_{\infty})$ for all $f\in X_s'$. In the latter case, $x_{\infty}$ is uniquely determined, in virtue of the Hahn-Banach theorem (Corollary~\hyperref[cor:IV.6.Corollary.2]{$2$} of Theorem~\ref{thm:IV.6.Theorem.1} in Chapter~\ref{chap:IV}, \ref{sec:IV.6}); we shall write $w\text{- }\lim\limits_{n\to \infty}x_n = x_\infty$ or, in short, $x_{n}\rightarrow x_{\infty}$ weakly. $X$ is said to be sequentially weakly complete if every weakly convergent sequence of $X$ converges weakly to an element of $X$. 
\end{definition}

\begin{examplen} Let $\{x_{n}(s)\}$ be a sequence of equi-bounded continuous functions of $C[0,1]$ which is convergent to a discontinuous function $z(s)$ on $[0,1]$. Then, since $C[0,1]'$ is the space of Baire measures on $[0,1]$ of bounded total variation, we see easily that $\{x_{n}(s)\}$ gives an example of a weakly convergent sequence of $C[0,1]$ which does not converge weakly to an element of $C[0,1]$.

\end{examplen}
\begin{theorem}\label{thm:V.1.Theorem.1} \begin{enumerate}[i)]
\item $s\text{- }\lim\limits_{n \to \infty} x_n = x_\infty$ implies $w\text{- }\lim\limits_{n \to \infty} x_n = x_\infty$, but not conversely.
\item $A$ weakly convergent sequence $\{x_n\}$ is strongly bounded, and, in particular, if $w\text{- }\lim\limits_{n \to \infty} x_n = x_\infty$, then $\{\|x_n\|\}$ is bounded and $\|x_\infty\| \leqq \varliminf\limits_{n \to \infty} \|x_n\|$.
\end{enumerate}
\end{theorem}

\begin{proof}
 \begin{enumerate}[i)]
\item The first part is clear from $|f(x_n) - f(x_\infty)| \leqq \|f\| \cdot \|x_n - x_\infty\|$. The second part is proved by considering the sequence $\{x_n\}$ in the Hilbert space $(l^2)$: 

$x_{n} = \{\xi_{m}^{(n)}\}$ where $\xi_{m}^{(n)} = \delta_{nm}$ (= $1$ or $0$ according as $n = m$ or not). 

For, the value of a continuous linear functional $\in (l^2)'$ at $x = \{\xi_n\}$ is given by $\sum\limits_{n=1}^{\infty} \xi_n \bar{\eta}_n$ with some $\{\eta_n\} \in (l^2)$; consequently, $w - \lim\limits_{n \to \infty} x_n = 0$ but $\{x_n\}$ does not converge strongly to $0$ because $\|x_n\| = 1$ ( $n = 1, 2, \ldots$).
\item Consider the sequence of continuous linear functionals $X_n$ defined on the $B$-space $X_s'$ by $X_n(f) = \langle x_n, f \rangle$, and apply the resonance theorem in Chapter~\ref{chap:II}, \ref{sec:II.1}. \qedhere
\end{enumerate}
\end{proof}

\begin{theorem}[\textsc{Mazur}]\label{thm:V.1.Theorem.2} Let \wlim $_{n \to \infty} x_n = x_\infty$ in a normed linear space $X$. Then there exists, for any $\varepsilon > 0$, a convex combination $\sum\limits_{j=1}^{n} \alpha_j x_j \left( \alpha_j \geqq 0, \sum\limits_{j=1}^{n} \alpha_j = 1 \right)$ of $x_j$'s such that $\left\| x_\infty - \sum\limits_{j=1}^{n} \alpha_j x_j \right\| \leqq \varepsilon$. 
\end{theorem}

\begin{proof}
 Consider the totality $M_1$ of elements of the form $\sum\limits_{j=1}^{n} \alpha_j x_j$ with $\alpha_j \geqq 0$, $\sum\limits_{j=1}^{n} \alpha_j = 1$. We may assume that $0 \in M_1$ by replacing $x_\infty$ and $x_j$ by $(x_\infty - x_1)$ and $(x_j - x_1)$, respectively. Suppose that $\|x_\infty - u\| > \varepsilon$ for every $u \in M_1$. The set $M = \{v \in X; \|v - u\| \leqq \varepsilon/2 \text{ for some } u \in M_1\}$ is a convex neighbourhood of $0$ of $X$ and $\|x_\infty - v\| > \varepsilon/2$ for all $v \in M$. Let $p(y)$ be the Minkowski functional of $M$. Since $x_\infty = \beta^{-1} u_0$ with $p(u_0) = 1$ and $0 < \beta < 1$, we must have $p(x_\infty) = \beta^{-1} > 1$. Consider a real linear subspace $X_1 = \{x \in X; x = \gamma u_0, -\infty < \gamma < \infty\}$ and put $f_1(x) = \gamma$ for $x = \gamma u_0 \in X_1$. This real linear functional $f_1$ on $X_1$ satisfies $f(x) \leqq p(x)$ on $X_1$. Thus, by the Hahn-Banach extension theorem in Chapter~\ref{chap:IV}, \ref{sec:IV.1}, there exists a real linear extension $f$ of $f_1$ defined on the real linear space $X$ and such that $f(x) \leqq p(x)$ on $X$. $M$ being a neighbourhood of $0$, the Minkowski functional $p(x)$ is continuous in $x$. Hence $f$ is a continuous real linear functional defined on the real linear normed space $X$. We have, moreover, 
\[
\sup\limits _ {x \in M _{1}} f (x) \leqq \sup\limits _{x \in M} f (x) \leqq \sup\limits _{x \in M} p (x) = 1 < \beta^{-1} = f (\beta^{-1} u _{0}) = f (x _{\infty}). \]

Therefore, it is easy to see that $x_{\infty}$ cannot be a weak accumulation point of $M_1$, contrary to $x_{\infty} = w - \lim\limits_{n \to \infty} x_n$. 
\end{proof}

\begin{theorem}\label{thm:V.1.Theorem.3} $A$ sequence $\{x_{n}\}$ of a normed linear space $X$ converges weakly to an element $x_{\infty} \in X$ iff the following two conditions are satisfied: \begin{enumerate}[i)]
\item $\sup\limits_{n \geqq 1} \|x_{n}\| < \infty$
\item $\lim\limits_{n \to \infty} f(x_{n}) = f(x_{\infty})$ for every $f$ from any strongly dense subset $D'$ of $X_{s}'$.
\end{enumerate}
\end{theorem}

\begin{proof}
 We have only to prove the sufficiency. For any $g \in X_s'$ and $\varepsilon > 0$, there exists an $f \in D'$ such that $\| g - f \| < \varepsilon$. Thus 
\begin{align}
\left| g \left(x _{n}\right) - g \left(x _{\infty}\right) \right|& \leqq \left| g \left(x _{n}\right) - f \left(x _{n}\right) \right| + \left| f \left(x _{n}\right) - f \left(x _{\infty}\right) \right| + \left| f \left(x _{\infty}\right) - g \left(x _{\infty}\right) \right| \notag \\
&\leqq \varepsilon \| x _{n} \| + | f (x _{n}) - f (x _{\infty}) | + \varepsilon \| x _{\infty} \|, \notag
\end{align}
and hence $\lim |g(x_n) - g(x_\infty)| \leqq 2\varepsilon \sup\limits_{\infty \geqq n \geqq 1} \|x_n\|$. This proves that $\lim\limits_{n \to \infty} g(x_n) = g(x_\infty)$. 
\end{proof}

\begin{theorem}\label{thm:V.1.Theorem.4} $A$ sequence $\{x_{n}\}$ in $L^{1}(S, \mathfrak{B}, m)$ converges weakly to an element $x \in L^{1}(S, \mathfrak{B}, m)$ iff $\{\| x_{n} \|\}$ is bounded and a finite $\lim\limits_{n \to \infty} \int_{B} x_{n}(s) m(\dd s)$ exists for every $B \in \mathfrak{B}$. 
\end{theorem}

\begin{proof}
 The ``only if'' part is clear since the defining function $C_{B}(s)$ of $B \in \mathfrak{B}$ belongs to $L^{\infty}(S, \mathfrak{B}, m) = L^{1}(S, \mathfrak{B}, m)'$. 

The proof of the ``if'' part. The set function $\psi(B)=\lim\limits_{n\to\infty}\int_{B}x_{n}(s)m(d,s)$ $B\in\mathfrak{B}$, is $\sigma$-additive and $m$-absolutely continuous by the Vitali-Hahn-Saks theorem. Hence, by the differentiation theorem of \textsc{Lebesgue}-\textsc{Nikodym}, there exists an $x_{\infty} \in L^{1}(S, \mathfrak{B}, m)$ such that 
\[
\lim\limits _{n \rightarrow \infty} \int_{B} x _{n} (s) m (\dd s) = \int_{B} x _{\infty} (s) m (\dd s) \text{ for all } B \in \mathfrak {B}.
\]

Thus, for any decomposition $S = \sum\limits_{j=1}^{k} B_j$ with $B_j \in \mathfrak{B}$, we have 
\[
\lim\limits _{n \rightarrow \infty} \int_{S} x _{n} (s) y (s) m (\dd s) = \int_{S} x _{\infty} (s) y (s) m (\dd s), y (s) = \sum\limits_{j=1} ^{k} \alpha_{j} C _ {B _{j}} (s).
\]

Since such functions as $y(s)$ constitute a strongly dense subset of the space $L^{\infty}(S, \mathfrak{B}, m) = L^{1}(S, \mathfrak{B}, m)'$, we see that the ``if'' part is true by Theorem~\ref{thm:V.1.Theorem.3}. 
\end{proof}

\begin{theorem}\label{thm:V.1.Theorem.5} Let $\{x_{n}\}$ converge weakly to $x_{\infty}$ in $L^{1}(S, \mathfrak{B}, m)$. Then $\{x_{n}\}$ converges strongly to $x_{\infty}$ iff $\{x_{n}(s)\}$ converges to $x_{\infty}(s)$ in $m$-measure on every $\mathfrak{B}$-measurable set $B$ such that $m(B) < \infty$. 
\end{theorem}

\begin{remarku} $\{x_{n}(s)\}$ is said to converge to $x_{\infty}(s)$ in $m$-measure on $B$, if, for any $\varepsilon > 0$, the $m$-measure of the set $\{s \in B; |x_{n}(s) - x_{\infty}(s)| \geqq \varepsilon\}$ tends to zero as $n \to \infty$ (see the Proposition in Chapter~\ref{chap:I}, \ref{sec:I.4}). The space $(l^{1})$ is an example of $L^{1}(S, \mathfrak{B}, m)$ for which $S = \{1, 2, \ldots\}$ and $m(\{n\}) = 1$ for $n = 1, 2, \ldots$. In this case, we have $(l^{1})' = (l^{\infty})$ so that the weak convergence of $\{x_{n}\}, x_{n} = (\xi_{1}^{(n)}, \xi_{2}^{(n)}, \ldots, \xi_{k}^{(n)}, \ldots)$, to $x_{\infty} = (\xi_{1}^{(\infty)}, \xi_{2}^{(\infty)}, \ldots, \xi_{k}^{(\infty)}, \ldots)$ implies that $\lim\limits_{n \to \infty} \xi_{k}^{(n)} = \xi_{k}^{(\infty)}$ ( $k = 1, 2, \ldots$), as may be seen by taking $f \in (l^{1})'$ in such a way that $f(x) = \langle x, f \rangle = \xi_{k}$ for $x = \{\xi_{j}\} \in (l^{1})$. Thus, in the present case, $\{x_{n}\}$ converges to $x_{\infty}$ in $m$-measure on every $\mathfrak{B}$-measurable set $B$ of finite $m$-measure. In this way we obtain the 
\end{remarku}

\begin{corollaryn}[I. \textsc{Schur}] In the space $(l^1)$, if a sequence $\{x_n\}$ converges weakly to an $x_{\infty} \in (l^1)$, then \slim $x_n = x_{\infty}$.
\end{corollaryn}

\begin{proof}[Proof of Theorem~\ref{thm:V.1.Theorem.5}.]
Since the strong convergence in $L^1(S, \mathfrak{B}, m)$ implies the convergence in $m$-measure, the ``only if'' part is clear. We shall prove the ``if'' part. The sequence $\{x_n - x_\infty\}$ converges weakly to $0$, and so
\begin{equation}\label{eq:V.1.1}
\lim\limits _{n \rightarrow \infty} \int_{B} \left(x _{n} (s) - x _{\infty} (s)\right) m (\dd s) = 0 \text{ for every } B \in \mathfrak {B}.
\end{equation}

Consider the sequence of non-negative measures 
\[
\psi_{n} (B) = \int_{B} | x _{n} (s) - x _{\infty} (s) | m (\dd s), B \in \mathfrak {B}.
\]

Then we have 
\begin{equation}
\label{eq:V.1.2}
\begin{minipage}{0.92\linewidth}
\everymath{}\let\limits\nolimits
$\lim\limits_{k \to \infty} \psi_n(B_k) = 0$ uniformly in $n,$ for any decreasing sequence $\{B_k\}$ of sets $\in \mathfrak{B}$ such that $\bigcap\limits_{k=1}^{\infty} B_k = \emptyset.$
\end{minipage}
\end{equation}


If otherwise, there exists an $\varepsilon > 0$ such that for each $k$ there exists some $n_{k}$ for which $\lim\limits_{k\to \infty}n_k = \infty$ and $\psi_{n_{k}}(B_k) > \varepsilon$. Consequently, we must have 
\[
\int_ {B _{k}} \left| \operatorname{Re} \left(x _ {n _{k}} (s) - x _{\infty} (s)\right) \right| m (\dd s) > \frac {\varepsilon}{\sqrt {2}} \quad \text{ or } \quad \int_ {B _{k}} \left| \operatorname{Im} \left(x _ {n _{k}} (s) - x _{\infty} (s)\right) \right| m (\dd s) > \frac {\varepsilon}{\sqrt {2}},
\]
and so there must exist some $B_{k}^{\prime} \subseteq B_{k}$ such that 
\[
\left| \int_ {B _{k} ^{\prime}} (x _ {n _{k}} (s) - x _{\infty} (s)) m (\dd s) \right| > \varepsilon / 2 \sqrt {2} (k = 1, 2, \dots),
\]
contrary to the fact that, in virtue of \eqref{eq:V.1.1}, the $m$-absolute continuity of the sequence of measures $\varphi_{n}(B) = \int_{B} (x_{n}(s) - x_{\infty}(s)) m(\dd s)$ is uniform in $n$ (see the proof of the Vitali-Hahn-Saks Theorem in Chapter~\ref{chap:II}, \ref{sec:II.2}). 

Next let $B_0$ be any set of $\mathfrak{B}$ such that $m(B_0) < \infty$. We shall show that 
\begin{equation}\label{eq:V.1.3}
\lim\limits _{n \rightarrow \infty} \psi_{n} (B _{0}) = 0.
\end{equation}

Suppose there exist an $\varepsilon > 0$ and a subsequence $\{\psi_{n'}\}$ of $\{\psi_n\}$ such that 
\begin{equation}\label{eq:V.1.4}
\psi_ {n ^{\prime}} (B _{0}) > \varepsilon (n = 1, 2, \dots).
\end{equation}

By the hypothesis that $\{(x_n(s) - x_{\infty}(s))\}$ converges to $0$ in $m$-measure on $B_0$, there exist a subsequence $\{(x_{n''}(s) - x_{\infty}(s))\}$ of $\{(x_{n'}(s) - x_{\infty}(s))\}$ and some sets $B_n'' \subseteq B_0$ such that $m(B_n'') \leqq 2^{-n}$ and $|x_{n''}(s) - x_{\infty}(s)| < \varepsilon / m(B_0)$ on $(B_0 - B_n'')$. We put $B_k = \bigcup\limits_{n=k}^{\infty} B_n''$. Then $\{B_k\}$ is a decreasing sequence such that 
$m\left(\bigcap\limits_{k=1}^{\infty} B_k\right) \leqq \sum\limits_{n=k}^{\infty} m(B_n'') \leqq 2^{-k+1} (k=1, 2, \ldots)$ and so $m\left(\bigcap\limits_{k=1}^{\infty} B_k\right) = 0$. 

Hence, by \eqref{eq:V.1.1} and the Corollary of the Vitali-Hahn-Saks Theorem referred to above, $\lim\limits_{k\to \infty}\psi_n(B_k) = 0$ uniformly in $n$. Therefore 
\[
\psi_ {n ^{\prime \prime}} (B _{0}) \leqq \psi_ {n ^{\prime \prime}} (B _{n} ^{\prime \prime}) + \varepsilon m (B _{0}) ^{-1} \cdot m (B _{0} - B _{n} ^{\prime \prime}) \rightarrow (\leqq \varepsilon) \text{ as } n \rightarrow \infty ,
\]
contrary to \eqref{eq:V.1.4}. This proves \eqref{eq:V.1.3}. 

Now we take a sequence $\{B_k'\}$ of sets $\in \mathfrak{B}$ such that $m(B_k') < \infty$ ( $k = 1, 2, \ldots$) and $S = \bigcup\limits_{k=1}^{\infty} B_k'$. Then 
\[
\int_{S} \left| x _{n} (s) - x _{\infty} (s) \right| m (\dd s) = \int_ {\bigcup\limits_{k=1} ^{t} B _{k} ^{\prime}} + \int_ {S - \bigcup\limits_{k=1} ^{t} B _{k} ^{\prime}}
\]

By \eqref{eq:V.1.3} the first term on the right tends to zero as $n \to \infty$ for fixed $t$, and the second term on the right tends, by \eqref{eq:V.1.2}, to zero as $t \to \infty$ uniformly in $n$. Therefore we have proved that $\slim x_n = x_\infty$ in $L^1(S, \mathfrak{B}, m)$. 

A similar situation in the case of the space $\mathfrak{D}(\Omega)'$ is given by 
\end{proof}

\begin{theorem}\label{thm:V.1.Theorem.6} Let $\{T_n\}$ be a sequence of generalized functions $\in \mathfrak{D}(\Omega)'$. If $\lim\limits_{n\to \infty}T_n = T$ in the weak* topology of $\mathfrak{D}(\Omega)'$, then $\lim\limits_{n\to \infty}T_n = T$ in the strong topology of $\mathfrak{D}(\Omega)'$. 
\end{theorem}

\begin{proof}
 The strong topology of the space $\mathfrak{D}(\Omega)'$ is defined (see Definition~\ref{def:IV.7.Definition.1} in Chapter~\ref{chap:IV}, \ref{sec:IV.7}) through the family of semi-norms 

$p_{B}(T)=\sup\limits_{\varphi\in\mathfrak{B}}|T(\varphi)|$, where $B$ is any bounded set of $D$($\Omega$). 


The weak* topology of the space $\mathfrak{D}(\Omega)'$ is defined through the family of semi-norms 
$p_{F}(T)=\sup\limits_{\varphi\in\mathfrak{F}}|T(\varphi)|$, where $F$ is any finite set of $D$($\Omega$). 

Thus $\lim\limits_{n\to \infty}T_n = T$ in the weak* topology of $\mathfrak{D}(\Omega)'$ is precisely $\lim\limits_{n\to \infty}T_n = T(\mathfrak{D}(\Omega)')$ defined in Chapter~\ref{chap:II}, \ref{sec:II.3}. 

Let $\mathfrak{B}$ be any bounded set of $\mathfrak{D}(\Omega)$. Then there exists a compact subset $K$ in $\Omega$ such that $\supp(\varphi) \subseteq K$ for any $\varphi \in \mathfrak{B}$ and $\sup\limits_{x \in K, \varphi \in \mathfrak{B}} |D^j\varphi(x)|$ 
$<\infty$ for any differential operator $D^{j}$ (Theorem~\ref{thm:I.8.Theorem.1} in Chapter~\ref{chap:I}, \ref{sec:I.8}). Thus, by the Ascoli-Arzelà Theorem, $B$ is relatively compact in $\mathfrak{D}_{K}(\Omega)$. We apply the uniform boundedness theorem to the sequence $\{T_{n}-T\}$ to the effect that, for any $\varepsilon>0$, there exists a neighbourhood $U$ of $0$ of $\mathfrak{D}_{K}(\Omega)$ such that 
\[
\sup\limits _{n; \varphi \in U} | (T _{n} - T) (\varphi) | < \varepsilon .
\]

The compact subset $\mathfrak{B}$ of $\mathfrak{D}_K(\Omega)$ is covered by a finite system of sets of the form $\varphi_i + U$, where $\varphi_i \in \mathfrak{B}$ ( $i = 1, 2, \ldots, k$). Hence 
\[
\begin{aligned}
&\left| (T _{n} - T) (\varphi_{i} + u) \right|  \leqq \left| (T _{n} - T) (\varphi_{i}) \right| + \left| (T _{n} - T) (u) \right| \leqq \left| (T _{n} - T) (\varphi_{i}) \right| + \varepsilon \text{ for any } u \in U.
\end{aligned}
\]

Since $\lim\limits_{n\to \infty}(T_n - T)(\varphi_i) = 0$ for $i = 1,2,\ldots ,k$, we have 
\[
\lim\limits _{n \rightarrow \infty} (T _{n} - T) (\varphi) = 0 \text{ uniformly in } \varphi \in \mathfrak {B}.
\]

This proves our Theorem.
\end{proof}

\begin{theorem}\label{thm:V.1.Theorem.7} $A$ reflexive $B$-space $X$ is sequentially weakly complete. 
\end{theorem}

\begin{proof}
 Let a sequence $\{x_{n}\}$ of $X$ be weakly convergent. Each $x_{n}$ defines a continuous linear functional $X_{n}$ on $X_{s}^{\prime}$ by $X_{n}(x^{\prime}) = \langle x_{n}, x^{\prime} \rangle$. Since $X_{s}^{\prime}$ is a $B$-space (Theorem~\ref{thm:IV.8.Theorem.1} in Chapter~\ref{chap:IV}, \ref{sec:IV.8}), we may apply the resonance theorem. Thus a continuous linear functional on $X_{s}^{\prime}$ is defined by a finite $\lim\limits_{n \to \infty} X_{n}(x^{\prime})$ which exists by hypothesis. Since $X$ is reflexive, there exists an $x_{\infty} \in X$ such that $\langle x_{\infty}, x^{\prime} \rangle = \lim\limits_{n \to \infty} X_{n}(x^{\prime}) = \lim\limits_{n \to \infty} \langle x_{n}, x^{\prime} \rangle$, that is, $x_{\infty} = w - \lim\limits_{n \to \infty} x_{n}$. 
\end{proof}

\begin{theorem}\label{thm:V.1.Theorem.8} Let $X$ be a Hilbert space. If a sequence $\{x_{n}\}$ of $X$ converges weakly to $x_{\infty} \in X$, then $\slim\limits_{n \to \infty} x_{n} = x_{\infty}$ iff $\lim\limits_{n \to \infty} \|x_{n}\| = \|x_{\infty}\|$. 
\end{theorem}

\begin{proof}
 The ``only if'' part is clear from the continuity of the norm. The ``if'' part is clear from the equality 
\[ \begin{aligned}
&{\| x _{n} - x _{\infty} \| ^{2} = (x _{n} - x _{\infty}, x _{n} - x _{\infty})} {\qquad = \| x _{n} \| ^{2} - (x _{n}, x _{\infty}) - (x _{\infty}, x _{n}) + \| x _{\infty} \| ^{2}.}
\end{aligned}
\]

In fact, the limit, as $n \to \infty$, of the right hand side is $\| x_{\infty} \|^2 - \| x_{\infty} \|^2 - \| x_{\infty} \|^2 + \| x_{\infty} \|^2 = 0$. 
\end{proof}
\yossub{Weak* Convergence}

\begin{definition}\label{def:V.1.Definition.2} $A$ sequence $\{f_n\}$ in the dual space $X_s'$ of a normed linear space $X$ is said to be weakly* convergent if a finite $\lim\limits_{n\to \infty}f_n(x)$ exists for every $x\in X$; $\{f_n\}$ is said to converge weakly* to an element $f_\infty \in X_s'$ if $\lim\limits_{n\to \infty}f_n(x) = f_\infty (x)$ for all $x\in X$. In the latter case, we write $w^* -\lim\limits_{n\to \infty}f_n = f_\infty$ or, in short, $f_n\to f_\infty$ weakly*. 
\end{definition}

\begin{theorem}\label{thm:V.1.Theorem.9} \begin{enumerate}[i)]
\item \slim $f_n = f_\infty$ implies \wstarlim $f_n = f_\infty$, but not conversely.
\item If $X$ is a $B$-space, then a weakly* convergent sequence $\{f_n\} \subseteq X_s'$ converges weakly* to an element $f_\infty \in X_s'$ and $\|f_\infty\| \leqq \lim\limits_{n \to \infty} \|f_n\|$.
\end{enumerate}
\end{theorem}

\begin{proof}
 \begin{enumerate}[(i)]
\item The first part is clear from $|f_n(x) - f_\infty(x)| \leqq \|f_n - f_\infty\| \cdot \|x\|$. The second part is proved by the counter example given in the proof of Theorem~\ref{thm:V.1.Theorem.1}.
\item By the resonance theorem, we see that $f_\infty(x) = \lim\limits_{n \to \infty} f_n(x)$ is a continuous linear functional on $X$ and $\|f_\infty\| \leqq \lim\limits_{n \to \infty} \|f_n\|$. \qedhere
\end{enumerate}
\end{proof}

\begin{theorem}\label{thm:V.1.Theorem.10} If $X$ is a $B$-space, then a sequence $\{f_n\} \subseteq X_s'$ converges weakly* to an element $f_\infty \in X_s'$ iff \begin{enumerate}[(i)]
\item $\{\|f_n\|\}$ is bounded
\item $\lim\limits_{n \to \infty} f_n(x) = f_\infty(x)$ on a strongly dense subset of $X$.
\end{enumerate}
\end{theorem}

\begin{proof}
 The proof is similar to that of Theorem~\ref{thm:V.1.Theorem.3}. 
\end{proof}

\yossub{Strong and Weak Closure}

\begin{theorem}\label{thm:V.1.Theorem.11} Let $X$ be a locally convex linear topological space, and $M$ a closed linear subspace of $X$. Then $M$ is closed in the weak topology of $X$. 
\end{theorem}

\begin{proof}
 If otherwise, there exists a point $x_{0} \in X - M$ such that $x_{0}$ is an accumulation point of the set $M$ in the weak topology of $X$. Then, by the Corollary of Theorem~\ref{thm:IV.6.Theorem.3} in Chapter~\ref{chap:IV}, \ref{sec:IV.6}, there exists a continuous linear functional $f_{0}$ on $X$ such that $f_{0}(x_{0}) = 1$ and $f_{0}(x) = 0$ on $M$. Hence $x_{0}$ cannot be an accumulation point of the set $M$ in the weak topology of $X$. 
\end{proof}


\yossection{2}{The Local Sequential Weak Compactness of Reflexive \texorpdfstring{$B$}{B}-spaces. The Uniform Convexity}
\label{sec:V.2}

\begin{theorem}\label{thm:V.2.Theorem.1} Let $X$ be a reflexive $B$-space, and let $\{x_{n}\}$ be any sequence which is norm bounded. Then we can choose a subsequence $\{x_{n'}\}$ which converges weakly to an element of $X$. 
\end{theorem}

We will prove this Theorem under the assumption that $X$ is separable, since concrete function spaces appearing in applications of this Theorem are mostly separable. The general case of a non-separable space will be treated in the Appendix. 

\begin{lemmau} If the strong dual $X_s'$ of a normed linear space $X$ is separable, then so is $X$. 
\end{lemmau}

\begin{proof}
 Let $\{x_n'\}$ be a countable sequence which is strongly dense on the surface of the unit sphere $\{x' \in X_s'; \|x'\| = 1\}$ of $X_s'$. Choose $x_n \in X$ so that $\|x_n\| = 1$ and $|\langle x_n, x_n' \rangle| \geqq 1/2$. Let $M$ be the closed linear subspace of $X$ spanned by the sequence $\{x_n\}$. Suppose $M \neq X$ and $x_0 \in X - M$. By Corollary of Mazur's Theorem~\ref{thm:IV.6.Theorem.3} in Chapter~\ref{chap:IV}, \ref{sec:IV.6}, there exists an $x_0' \in X_s'$ such that $\|x_0'| = 1$, $\langle x_0, x_0' \rangle \neq 0$ and $\langle x, x_0' \rangle = 0$ whenever $x \in M$. Thus $\langle x_n, x_0' \rangle = 0$ ( $n = 1, 2, \ldots$), and so $1/2 \leqq |\langle x_n, x_n' \rangle| \leqq |\langle x_n, x_n' \rangle - \langle x_n, x_0' \rangle| + |\langle x_n, x_0' \rangle|$ which implies that $1/2 \leqq \|x_n\| \|x_n' - x_0'\| = \|x_n' - x_0'\|$. This is a contradiction to the fact that $\{x_n'\}$ is strongly dense on the surface of the unit sphere of $X_s'$. Thus $M = X$, and so linear combinations with rational coefficients of $\{x_n\}$ are dense in $X$. This proves our Lemma. 
\end{proof}

\begin{proof}[Proof of Theorem~\ref{thm:V.2.Theorem.1}.]
As we have remarked above, we assume that $X$ is separable and so $(X_s')_s' = X$ is separable also. By the preceding Lemma, $X_s'$ is separable. Let $\{x_n'\}$ be a countable sequence which is strongly dense in $X_s'$. Since $\{x_n\}$ is norm bounded, the sequence $\{\langle x_n, x_1' \rangle\}$ is bounded. Thus there exists a subsequence $\{x_{n_1}\}$ for which the sequence $\{\langle x_{n_1}, x_1' \rangle\}$ is convergent. Since the sequence $\{\langle x_{n_1}, x_2' \rangle\}$ is bounded, there exists a subsequence $\{x_{n_2}\}$ of $\{x_{n_2}\}$ such that $\{\langle x_{n_2}, x_2' \rangle\}$ is convergent. Proceeding in this way, we can choose a subsequence $\{x_{n_{i+1}}\}$ of the sequence $\{x_{n_{i}}\}$ such that the sequence of numbers $\{\langle x_{n_{i+1}}, x_j' \rangle\}$ converges for $j = 1, 2, \ldots, i + 1$. Hence the diagonal subsequence $\{x_{n_{n}}\}$ of the original sequence $\{x_n\}$ satisfies the condition that the sequence $\{\langle x_{n_{n}}, x_j' \rangle\}$ converges for $j = 1, 2, \ldots$. Thus, by Theorem~\hyperref[thm:V.1.Theorem.3]{$3$} in the preceding section, $\lim\limits_{n \to \infty} \langle x_{n_{n}}, x' \rangle$ exists and is finite for every $x' \in X'$. Hence, by Theorem~\hyperref[thm:V.1.Theorem.7]{$7$} of the preceding section, we see that $w - \lim\limits_{n \to \infty} x_{n_{n}}$ exists.
\end{proof}

\yossub{Milman's Theorem}

We owe to D. P. \textsc{Milman} a theorem that a $B$-space is reflexive when it is uniformly convex in the sense that, for any $\varepsilon > 0$, there exists a $\delta = \delta (\varepsilon) > 0$ such that $\|x\| \leqq 1$, $\|y\| \leqq 1$ and $\|x - y\| \geqq \varepsilon$ implies $\|x + y\| \leqq 2(1 - \delta)$. A pre-Hilbert space is uniformly convex as may be seen from the formula 
\[
\left\| x + y \right\| ^{2} + \left\| x - y \right\| ^{2} = 2 \left(\left\| x \right\| ^{2} + \left\| y \right\| ^{2}\right)
\]
valid in such a space. It is known that, for $1 < p < \infty$, the spaces $L^p$ and $(l^p)$ are uniformly convex (see J. A. \textsc{Clarkson} \cite{clarkson35}). 

\begin{theorem}[\textsc{Milman} {\cite{milman198}}]\label{thm:V.2.Theorem.2} $A$ uniformly convex $B$-space $X$ is reflexive. 
\end{theorem}

Proof (due to S. \textsc{Kakutani}). Given an $x_0'' \in (X_s')_s'$ with $\|x_0''\| = 1$. Then there exists a sequence $\{f_n\} \subseteq X_s'$ with $\|f_n\| = 1$, $x_0''(f_n) \geqq 1 - n^{-1}$ ( $n = 1, 2, \ldots$). By Theorem~\ref{thm:IV.6.Theorem.5} in Chapter~\ref{chap:IV}, \ref{sec:IV.6}, there exists, for every $n$, an $x_n \in X$ such that 
\begin{gather}\quad f _{i} (x _{n})= x _{0} ^{\prime \prime} (f _{i}) (i = 1, 2, \dots , n) \quad \text{ and } \quad \| x _{n} \| \leqq \| x _{0} ^{\prime \prime} \| + n ^{-1} = 1 + n ^{-1}. \text{ Since } \notag
\\1 - n ^{-1} \leqq x _{0} ^{\prime \prime} (f _{n}) = f _{n} (x _{n}) \leqq \| f _{n} \| \| x _{n} \| = \| x _{n} \| \leqq 1 + n ^{-1}, \notag
\end{gather}
we must have $\lim\limits_{n\to \infty}\|x_n\| = 1.$ 

If the sequence $\{x_{n}\}$ does not converge strongly, there exists an $\varepsilon > 0$ and $n_{1} < m_{1} < n_{2} < m_{2} < \cdots < n_{k} < m_{k} < \cdots$ such that $\varepsilon \leqq \|x_{n_{k}} - x_{m_{k}}\|$ ( $k = 1, 2, \ldots$). Thus, by $\lim\limits_{n \to \infty} \|x_n\| = 1$ and the uniform convexity of $X$, we obtain $\varlimsup\limits_{k \to \infty} \|x_{n_{k}} + x_{m_{k}}\| \leqq 2(1 - \delta(\varepsilon)) < 2$. But, since $n_k < m_k$, $f_{n_{k}}(x_{n_{k}}) = f_{n_{k}}(x_{m_{k}}) = x_0''(f_{n_{k}})$ and so 
\[
2 \left(1 - n _{k} ^{-1}\right) \leqq 2 x _{0} ^{\prime \prime} \left(f _ {n _{k}}\right) = f _ {n _{k}} \left(x _ {n _{k}} + x _ {m _{k}}\right) \leqq \left\| f _ {n _{k}} \right\| \cdot \left\| x _ {n _{k}} + x _ {m _{k}} \right\|.
\]

Hence, by $\| f_{n_{k}} \| = 1$, we obtain a contradiction $\varlimsup\limits_{k \to \infty} \| x_{n_{k}} + x_{m_{k}} \| \geqq 2$. 

We have thus proved the existence of $\slim\limits_{n \to \infty} x_n = x_0$, and $x_0$ satisfies 
\begin{equation}\label{eq:V.2.1}
\left\| x _{0} \right\| = 1, f _{i} (x _{0}) = x _{0} ^{\prime \prime} (f _{i}) (i = 1, 2, \dots).
\end{equation}

We show that the solution of the above equations in \eqref{eq:V.2.1} is unique. Otherwise, there exists an $\hat{x}_0 \neq x_0$ which satisfies the same equations. By the uniform convexity, $\|\hat{x}_0 + x_0\| < 2$. We also have $f_i(\hat{x}_0 + x_0) = 2x_0''(f_i)$ ( $i = 1, 2, \ldots$). Thus 
$2(1 - i^{-1}) \leqq 2x_0''(f_i) = f_i(\hat{x}_0 + x_0) \leqq \|f_i\| \|\hat{x}_0 + x_0\| = \|\hat{x}_0 + x_0\|,$ and so $\|\hat{x}_0 + x_0\| \geqq \lim\limits_{i \to \infty} 2(1 - i^{-1}) = 2$ which is a contradiction. 

Finally let $f_0$ be any point of $X_s'$. If we show that $f_0(x_0) = x_0''(f_0)$, then $(X_s')_s' \subseteq X$ and the reflexivity of $X$ is proved. To prove that $f_0(x_0) = x_0''(f_0)$, we take $f_0, f_1, \ldots, f_n, \ldots$ in place of $f_1, f_2, \ldots, f_n, \ldots$ above, and hence we obtain $\hat{x}_0 \in X$ such that 
\[
\left\| \hat {x} _{0} \right\| = 1, f _{i} (\hat {x} _{0}) = x _{0} ^{\prime \prime} (f _{i}) (i = 0, 1, \dots , n, \dots);
\]
we must have $\hat{x}_{0}=x_{0}$ by virtue of the uniqueness just proved above, and so the proof of Theorem~\ref{thm:V.2.Theorem.2} is completed. 


\yossection{3}{Dunford's Theorem and The Gelfand-Mazur Theorem}
\label{sec:V.3}

\begin{definition}\label{def:V.3.Definition.1} Let $Z$ be an open domain of the complex plane. A mapping $x(\zeta)$ defined in $Z$ with values in a $B$-space $X$ is called weakly holomorphic in $\zeta$ in the domain $Z$ if, for each $f \in X'$, the numerical function 
\[
f (x (\zeta)) = \langle x (\zeta), f \rangle
\]
of $\zeta$ is holomorphic in $Z$. 
\end{definition}

\begin{theorem}[N. \textsc{Dunford} {\cite{dunford45}}]\label{thm:V.3.Theorem.1} If $x(\zeta)$ is weakly holomorphic in $Z$, then there exists a mapping $x'(\zeta)$ defined in $Z$ with values in $X$ such that, for each $\zeta_0 \in Z$, we have 
\[
\slim _{h \to 0} h ^{-1} (x (\zeta_{0} + h) - x (\zeta_{0})) = x ^{\prime} (\zeta_{0}).
\]
\end{theorem}

In other words, the weak holomorphic property implies the strong holomorphic property. 

\begin{proof}
 Let $C$ be a rectifiable Jordan curve such that the closed bounded domain $\overline{C}$ enclosed by $C$ lies entirely in $Z$ and $\zeta_0 \in \overline{C} - C$. Let $Z_0$ be any open complex domain $\ni \zeta_0$ such that its closure lies in the interior of $\overline{C}$. Then, by Cauchy's integral representation, we have 
\[
f (x (\zeta_{0})) = \frac {1}{2 \pi i} \int_{C} \frac {f (x (\zeta))}{\zeta - \zeta_{0}} \dd \zeta . \]

Hence, if both $\zeta_0 + h$ and $\zeta_0 + g$ belong to $Z_0$, 
\[
\begin{aligned}
&\quad (h - g) ^{-1} \left\{\frac {f (x (\zeta_{0} + h)) - f (x (\zeta_{0}))}{h} - \frac {f (x (\zeta_{0} + g)) - f (x (\zeta_{0}))}{g} \right\}
\\
&= \frac {1}{2 \pi i} \int_{G} f (x (\zeta)) \left\{\frac {1}{(\zeta - \zeta_{0} - h) (\zeta - \zeta_{0} - g) (\zeta - \zeta_{0})} \right\} \dd \zeta .
\end{aligned}
\]

By the assumption, the distance between $Z_{0}$ and $C$ is positive. Hence, for fixed $f \in X'$, the absolute value of the right hand side is uniformly bounded when $\zeta_{0}, \zeta_{0} + h$ and $\zeta_{0} + g$ range over $Z_{0}$. Thus, by the resonance theorem, we have 
\[
\sup\limits _ {\zeta_{0}, \zeta_{0} + h, \zeta_{0} + g \in Z _{0}} \frac {1}{| h - g |} \left\| \left\{\frac {x (\zeta_{0} + h) - x (\zeta_{0})}{h} - \frac {x (\zeta_{0} + g) - x (\zeta_{0})}{g} \right\} \right\| < \infty .
\]

Therefore, by the completeness of the space $X$, $x(\zeta)$ is strongly differentiable at every $\zeta_0 \in Z$. 
\end{proof}

\begin{corollary}[Cauchy's integral theorem]\label{cor:V.3.Corollary.1} The strong differentiability of $x(\zeta)$ implies its strong continuity in $\zeta$. Thus we can define the curvilinear integral $\int_{C} x(\zeta) \dd\zeta$ with values in $X$. Actually we can prove that $\int_{C} x(\zeta) \dd\zeta = 0$, the zero vector of $X$. 
\end{corollary}

\begin{proof}
 We have, by the continuity and the linearity of $f \in X'$, 
\[
f \left(\int_{C} x (\zeta) \dd \zeta\right) = \int_{C} f (x (\zeta)) \dd \zeta . \]

But the right hand side is zero, because of the ordinary Cauchy integral theorem. Since $f \in X'$ was arbitrary, we must have $\int_{C} x(\zeta) \dd\zeta = 0$ by Corollary~\ref{cor:V.3.Corollary.2} of Theorem~\ref{thm:IV.6.Theorem.1} in Chapter~\ref{chap:IV}, \ref{sec:IV.6}. 
\end{proof}

From the above Corollary, we can derive other Corollaries, as in the ordinary theory of functions of a complex variable. 

\begin{corollary}[Cauchy's integral representation]\label{cor:V.3.Corollary.2} 
\[
x(\zeta_0)=\frac{1}{2\pi i}\int_{\mathcal C}\frac{x(\zeta)}{\zeta-\zeta_0}\dd\zeta
\]
for any interior point $\zeta_0$ of $\overline{\mathcal C}$.
\end{corollary}

\begin{corollary}[Taylor's expansion]\label{cor:V.3.Corollary.3} For any point $\zeta_0$ which is in the interior of the closed domain $\overline{C}$, the Taylor expansion of $x(\zeta)$ at $\zeta = \zeta_0$ converges strongly in the interior of the circle with centre at $\zeta_0$, if this circle does not extend outside of $\overline{C}$: 
\[
\begin{gathered}
x (\zeta) = \sum\limits_{n=0} ^{\infty} (n!) ^{-1} (\zeta - \zeta_{0}) ^{n} x ^{(n)} (\zeta_{0}), \text{ where } x ^{(n)} (\zeta_{0}) = \frac {n !}{2 \pi i} \int_{C} \frac {x (\zeta)}{(\zeta - \zeta_{0}) ^{n+1}} \dd \zeta .
\end{gathered}
\]
\end{corollary}

\begin{corollary}[Liouville's theorem]\label{cor:V.3.Corollary.4} If $x(\zeta)$ is (strongly) holomorphic in the whole finite plane: $|\zeta| < \infty$, and $\sup\limits_{|\zeta| < \infty} |x(\zeta)| < \infty$, then $x(\zeta)$ must reduce to a constant vector $x(0)$. 
\end{corollary}

\begin{proof}
 If we take $\left|\zeta\right|=r$ for the curve $C$, then, as $r\to\infty$, 
\[
\left\| x ^{(n)} (0) \right\| = \frac {n !}{2 \pi} \sup\limits _{| \zeta | < \infty} \left\| x (\zeta) \right\| \int_{C} \frac {| \dd \zeta |}{r ^{n+1}} \rightarrow 0 (n = 1, 2, \dots). \]

Hence the Taylor expansion of $x(\zeta)$ at $\zeta = 0$ reduces to the constant term $x(0)$ only. 
\end{proof}

We shall now apply Corollary~\ref{cor:V.3.Corollary.4} to the proof of the Gelfand-Mazur theorem. We first give 

\begin{definition}\label{def:V.3.Definition.2} $A$ commutative field $X$ over the field of complex numbers is called a normed field, if it is also a $B$-space such that the following conditions are satisfied: 

$\|e\| = 1$, where $e$ is the unit of the multiplication in $X$, 
\begin{equation}\label{eq:V.3.1}
\left\| x y \right\| \leqq \left\| x \right\| \left\| y \right\|, \text{ where } x y \text{ is the multiplication in } X.
\end{equation}
\end{definition}

\begin{theorem}[\textsc{Gelfand} {\cite{gelfand76b}}-\textsc{Mazur} {\cite{mazur191}}]\label{thm:V.3.Theorem.2} $A$ normed field $X$ is isometrically isomorphic to the complex number field. In other words, every element $x$ of $X$ is of the form $x = \xi e$ where $\xi$ is a complex number. 

\end{theorem}

\begin{proof}
 Assume the contrary, and let there exist an $x \in X$ such that $(x - \xi e) \neq 0$ for any complex number $\xi$. Since $X$ is a field, the non-zero element $(x - \xi e)$ has the inverse $(x - \xi e)^{-1} \in X$. 

We shall prove that $(x - \lambda e)^{-1}$ is (strongly) holomorphic in $\lambda$ for $|\lambda| < \infty$. We have, in fact, 
\[
\begin{gathered}
h ^{-1} ((x - (\lambda + h) e) ^{-1} - (x - \lambda e) ^{-1}) = h ^{-1} (x - (\lambda + h) e) ^{-1} \left\{e - (x - (\lambda + h) e) (x - \lambda e) ^{-1} \right\}
\\
= h ^{-1} (x - (\lambda + h) e) ^{-1} \left\{e - e + h (x - \lambda e) ^{-1} \right\} = (x - (\lambda + h) e) ^{-1} (x - \lambda e) ^{-1}.
\end{gathered}
\]

On the other hand, for sufficiently small $|h|$, the series $y^{-1}\left(e + \sum\limits_{n=1}^{\infty}(hy^{-1})^{n}\right)$, where $y = (x - \lambda e)$, converges by \eqref{eq:V.3.1}, and it represents the inverse $(y - he)^{-1} = y^{-1}(e - hy^{-1})^{-1}$, as may be seen by multiplying the series by $(y - he)$. Hence, by the strong continuity in $h$ of the series, we can prove that $(x - \lambda e)^{-1}$ is (strongly) holomorphic in $\lambda$ with the strong derivative $(x - \lambda e)^{-2}$. 

Now, if $|\lambda| \geqq 2\|x\|$, then, as above, $(x - \lambda e)^{-1} = -\lambda^{-1}(e - \lambda^{-1}x)^{-1} = -\lambda^{-1}\left(e + \sum\limits_{n=1}^{\infty}(\lambda^{-1}x)^n\right)$ and so 
\[
\left\| (x - \lambda e) ^{-1} \right\| \leqq \left| \lambda^{-1} \right|\left(1 + \sum\limits_{n=1} ^{\infty} (1 / 2) ^{n}\right)\rightarrow 0 \quad \text{ as } \quad | \lambda | \rightarrow \infty .
\]

Moreover, the function $(x - \lambda e)^{-1}$, being continuous in $\lambda$, is bounded on the compact domain of $\lambda: |\lambda| \leqq 2 \|x\|$. Hence, by Liouville's theorem, $(x - \lambda e)^{-1}$ must reduce to the constant vector $x^{-1} = (x - 0e)^{-1}$. But, since \slim $\frac{|\lambda| \to \infty}{(x - \lambda e)^{-1}} = 0$ as proved above, we have arrived at a contradiction $x^{-1} = 0$, $e = x^{-1}x = 0$. 
\end{proof}


\yossection{4}{The Weak and Strong Measurability. Pettis' Theorem}
\label{sec:V.4}

\begin{definition}\label{def:V.4.Definition.1} Let $(S, \mathfrak{B}, m)$ be a measure space, and $x(s)$ a mapping defined on $S$ with values in a $B$-space $X$. $x(s)$ is called weakly $\mathfrak{B}$-measurable if, for any $f \in X'$, the numerical function $f(x(s)) = \langle x(s), f \rangle$ of $s$ is $\mathfrak{B}$-measurable. $x(s)$ is said to be finitely-valued if it is constant $\neq 0$ on each of a finite number of disjoint $\mathfrak{B}$-measurable sets $B_j$ with $m(B_j) < \infty$ and $x(s) = 0$ on $S - \bigcup\limits_{j} B_j$. $x(s)$ is said to be strongly $\mathfrak{B}$-measurable if there exists a sequence of finitely-valued functions strongly convergent to $x(s)$ $m$-a.e. on $S$. 
\end{definition}

\begin{definition}\label{def:V.4.Definition.2} $x(s)$ is said to be separably-valued if its range $\{x(s); s \in S\}$ is separable. It is $m$-almost separably-valued if there exists a $\mathfrak{B}$-measurable set $B_0$ of $m$-measure zero such that $\{x(s); s \in S - B_0\}$ is separable. 
\end{definition}

\begin{theoremn}[B. J. \textsc{Pettis} {\cite{pettis249}}] $x(s)$ is strongly $B$-measurable iff it is weakly $B$-measurable and $m$-almost separably-valued.

\end{theoremn}
\begin{proof}
 The ``only if'' part is proved as follows. The strong $\mathfrak{B}$-measurability implies the weak $\mathfrak{B}$-measurability, because a finitely-valued function is weakly $\mathfrak{B}$-measurable, and, by the strong $\mathfrak{B}$-measurability of $x(s)$, there exists a sequence of finitely-valued functions $x_{n}(s)$ such that $s\text{- }\lim\limits_{n\to\infty}x_{n}(s)=x(s)$ except on a set $B_{0}\in\mathfrak{B}$ of $m$-measure zero. Thus the union of the ranges of $x_{n}(s)$ ( $n=1,2,\ldots$) is a countable set, and the closure of this set is separable and contains the range $\{x(s); s\in S-B_{0}\}$. 

The proof of the ``if'' part. Without losing the generality, we may assume that the range $\{x(s); s \in S\}$ is itself separable. So we may assume that the space $X$ is itself separable; otherwise, we replace $X$ by the smallest closed linear subspace containing the range of $x(s)$. We first prove that $\|x(s)\|$ is itself $B$-measurable. To this purpose, we shall make use of a lemma, to be proved later, which states that the dual space $X'$ of a separable $B$-space satisfies the condition 
\begin{equation}\label{eq:V.4.1}
\begin{minipage}[b]{0.92\linewidth}
there exists a sequence $\{f_n\} \subseteq X'$ with $\|f_n\| \leqq 1$ such that, for any $f_0 \in X'$ with $\|f_0\| \leqq 1$, we may choose a subsequence $\{f_{n'}\}$ of $\{f_n\}$ for which we have $\lim\limits_{n \to \infty} f_{n'}(x) = f_0(x)$ at every $x \in X$.
\end{minipage}
\end{equation}


Now, for any real number $a$, put 
$A = \{s; \|x(s)\| \leqq a\}$ and $A_f = \{s; |f(x(s))| \leqq a\}$, where $f \in X'$. If we can show that $A = \bigcap\limits_{j=1}^{\infty} A_{f_{j}}$, then, by the weak $B$-measurability of $x(s)$, the function $\|x(s)\|$ is $B$-measurable. It is clear that $A \subseteq_{\|f\| \leqq 1} A_f$. But, by Corollary~\hyperref[cor:IV.6.Corollary.2]{$2$} of Theorem~\ref{thm:IV.6.Theorem.1} in Chapter~\ref{chap:IV}, \ref{sec:IV.6}, there exists, for fixed $s$, an $f_0 \in X'$ with $\|f_0\| = 1$ and $f_0(x(s)) = \|x(s)\|$. Hence the reverse inclusion $A \supseteq_{\|f\| \leqq 1} A_f$ is true and so we have $A = \bigcap\limits_{\|f\| \leqq 1} A_f$. By the Lemma, we obtain $\bigcap\limits_{\|f\| \leqq 1} A_f = \bigcap\limits_{j=1}^{\infty} A_{f_{j}}$ and so $A = \bigcap\limits_{j=1}^{\infty} A_{f_{j}}$. 

Since the range $\{x(s); s \in S\}$ is separable, this range may, for any positive integer $n$, be covered by a countable number of open spheres $S_{j,n}$ ( $j = 1, 2, \ldots$) of radius $\leqq 1/n$. Let the centre of the sphere $S_{j,n}$ be $x_{j,n}$. As proved above, $\|x(s) - x_{j,n}\|$ is $\mathfrak{B}$-measurable in $s$. Hence the set $B_{j,n} = \{s \in S; x(s) \in S_{j,n}\}$ is $\mathfrak{B}$-measurable and $S = \bigcup\limits_{j=1}^{\infty} B_{j,n}$. We set 
\[
x _{n} (s) = x _{i,n} \quad \text{ if } \quad s \in B _{i,n} ^{\prime} = B _{i,n} - \bigcup\limits_{j=1} ^{i-1} B _{j,n}.
\]

Then, by $S = \sum\limits_{i=1}^{\infty} B_{i,n}'$, we have $\|x(s) - x_n(s)\| < 1/n$ for every $s \in S$.  

Since $B_{in}'$ is $\mathfrak{B}$-measurable, it is easy to see that each $x_{n}(s)$ is strongly $\mathfrak{B}$-measurable. Therefore $x(s)$, which is the strong limit of the sequence $\{x_{n}(s)\}$, is also strongly $\mathfrak{B}$-measurable. 
\end{proof}

\begin{proof}[Proof of the Lemma.]
Let a sequence $\{x_{n}\}$ be strongly dense in $X$. Consider a mapping $f \to \varphi_{n}(f) = \{f(x_{1}), f(x_{2}), \ldots, f(x_{n})\}$ of the unit sphere $S' = \{f \in X'; \|f\| \leqq 1\}$ of $X'$ into an $n$-dimensional Hilbert space $l^{2}(n)$ of vectors $(\xi_{1}, \xi_{2}, \ldots, \xi_{n})$ normed by $\left\|\left(\xi_{1}, \xi_{2}, \ldots, \xi_{n}\right)\right\| = \left(\sum\limits_{j=1}^{n} |\xi_{j}|^{2}\right)^{1/2}$. The space $l^{2}(n)$ being separable, there exists, for fixed $n$, a sequence $\{f_{n,k}\} (k = 1, 2, \ldots)$ of $S'$ such that $\{\varphi_{n}(f_{n,k}); k = 1, 2, \ldots\}$ is dense in the image $\varphi_{n}(S')$ of $S'$. 

We thus have proved that, for any $f_0 \in S'$, we can choose a subsequence $\{f_{n,m_n}\} (n = 1, 2, \ldots)$ such that $|f_{n,m_n}(x_i) - f_0(x_i)| < 1/n (i = 1, 2, \ldots, n)$. Hence $\lim\limits_{n \to \infty} f_{n,m_n}(x_i) = f_0(x_i) (i = 1, 2, \ldots)$, and so, by Theorem~\ref{thm:V.1.Theorem.10} in Chapter~\ref{chap:V}, \ref{sec:V.1}, we obtain that $\lim\limits_{n \to \infty} f_{n,m_n}(x) = f_0(x)$ for every $x \in X$. 
\end{proof}

\yossection{5}{Bochner's Integral}
\label{sec:V.5}

Let $x(s)$ be a finitely-valued function defined on a measure space $(S, \mathfrak{B}, m)$ with values in a $B$-space $X$; let $x(s)$ be equal to $x_i \neq 0$ on $B_i \in \mathfrak{B}$ ( $i = 1, 2, \ldots, n$) where $B_i$'s are disjoint and $m(B_i) < \infty$ for $i = 1, 2, \ldots, n$, and moreover, $x(s) = 0$ on $\left(S - \sum\limits_{i=1}^{n} B_i\right)$. Then we can define the $m$-integral $\int_S x(s) m(\dd s)$ of $x(s)$ over $S$ by $\sum\limits_{i=1}^{n} x_i m(B_i)$. By virtue of a limiting procedure, we can define the $m$-integral of more general functions. More precisely, we have the 

\begin{definitionu} $A$ function $x(s)$ defined on a measure space $(S, \mathfrak{B}, m)$ with values in a $B$-space $X$ is said to be Bochner $m$-integrable, if there exists a sequence of finitely-valued functions $\{x_n(s)\}$ which $s$-converges to $x(s)$ $m$-a.e. in such a way that 
\begin{equation}\label{eq:V.5.1}
\lim\limits _{n \rightarrow \infty} \int_{S} \| x (s) - x _{n} (s) \| m (\dd s) = 0.
\end{equation}
\end{definitionu}

For any set $B \in \mathfrak{B}$, the Bochner $m$-integral of $x(s)$ over $B$ is defined by 
\begin{gather}\int_{B} x (s) m (\dd s) = \slim _{n \rightarrow \infty} \int_{S} C _{B} (s) x _{n} (s) m (\dd s), \text{ where } C _{B} \text{ is the } \label{eq:V.5.2} \\\quad \text{ defining function of the set } B. \notag
\end{gather}

To justify the above definition, we have to verify that the $s$-limit on the right of \eqref{eq:V.5.2} exists and that the value of this $s$-limit is independent of the approximating sequence of functions $\{x_{n}(s)\}$. 

\begin{definitionn}[Justification of the Definition] First, $x(s)$ is strongly $\mathfrak{B}$-measurable and consequently the condition \eqref{eq:V.5.1} has a sense, since $\|x(s) - - x_n(s)\|$ is $\mathfrak{B}$-measurable as shown in the proof of Petti's Theorem. From the inequality 
\[
\begin{aligned}
&\left| \int_{B} x _{n} (s) m (\dd s) - \int_{B} x _{k} (s) m (\dd s) \right| = \left| \int_{B} (x _{n} (s) - x _{k} (s)) m (\dd s) \right|
\\
&\leqq \int_{B} \| x _{n} (s) - x _{k} (s) \| m (\dd s) \leqq \int_{S} \| x _{n} (s) - x (s) \| m (\dd s) + \int_{S} \| x (s) - x _{k} (s) \| m (\dd s)
\end{aligned}
\]

and the completeness of the space $X$, we see that $s\text{- }\lim\limits_{n\to \infty}\int_{B}x_n(s)m(\dd s)$ 
exists. It is clear also that this \slim is independent of the approximating sequence, since any two such sequences can be combined into a single approximating sequence. 
\end{definitionn}

\begin{theorem}[S. \textsc{Bochner} {\cite{bochner20}}]\label{thm:V.5.Theorem.1} $A$ strongly $\mathfrak{B}$-measurable function $x(s)$ is Bochner $m$-integrable iff $\| x(s) \|$ is $m$-integrable. 
\end{theorem}

\begin{proof}
 The ``only if'' part. We have $\left\|x(s)\right\| \leqq \left\|x_{n}(s)\right\| + \left\|x(s) - x_{n}(s)\right\|$. By the $m$-integrability of $\left\|x_{n}(s)\right\|$ and the condition \eqref{eq:V.5.1}, it is clear that $\left\|x(s)\right\|$ is $m$-integrable and 
\[
\int_{B} \| x (s) \| m (\dd s) \leqq \int_{B} \| x _{n} (s) \| m (\dd s) + \int_{B} \| x (s) - x _{n} (s) \| m (\dd s). \]

Moreover, since 
\[
\int_{B} \left| \left\| x _{n} (s) \right\| - \left\| x _{k} (s) \right\| \right| m (\dd s) \leqq \int_{B} \left\| x _{n} (s) - x _{k} (s) \right\| m (\dd s),
\]
we see from \eqref{eq:V.5.1} that $\lim\limits_{n\to \infty}\int_{B}\|x_n(s)\|m(\dd s)$ exists so that we have 
\[
\int_{B} \left\| x (s) \right\| m (\dd s) \leqq \varliminf\limits_{n \rightarrow \infty} \int_{B} \left\| x _{n} (s) \right\| m (\dd s).
\]

The ``if'' part. Let $\{x_{n}(s)\}$ be a sequence of finitely-valued functions strongly convergent to $x(s)$ $m$-a.e. Put 
\[
\begin{aligned}
&y _{n} (s) = x _{n} (s)  \text{ if } \| x _{n} (s) \| \leqq \| x (s) \| (1 + 2 ^{-1}), = 0  \text{ if } \| x _{n} (s) \| > \| x (s) \| (1 + 2 ^{-1}).
\end{aligned}
\]

Then the sequence of finitely-valued functions $\{y_n(s)\}$ satisfies $\|y_n(s)\| \leqq \|x(s)\| \cdot (1 + 2^{-1})$ and $\lim\limits_{n \to \infty} \|x(s) - y_n(s)\| = 0$ $m$-a.e. Thus, by the $m$-integrability of $\|x(s)\|$, we may apply the Lebesgue-Fatou Lemma to the functions $\|x(s) - y_n(s)\| \leqq \|x(s)\| (2 + 2^{-1})$ and obtain 
\[
\lim\limits _{n \rightarrow \infty} \int_{S} \| x (s) - y _{n} (s) \| m (\dd s) = 0,
\]
that is, $x(s)$ is Bochner $m$-integrable. 
\end{proof}
\begin{corollary}\label{cor:V.5.Corollary.1} The above proof shows that 
\[
\int_{B} \left\| x (s) \right\| m (\dd s) \geqq \left\| \int_{B} x (s) m (\dd s) \right\|,
\]
and hence $\int_{B} x(s) m(\dd s)$ is $m$-absolutely continuous in the sense that 
\[
\underset {m (B) \rightarrow 0} {\slim} \int_{B} x (s) m (\dd s) = 0.
\]
\end{corollary}

The finite additivity $\int_{\sum\limits_{j=1}^{n} B_j}^n x(s) m(\dd s) = \sum\limits_{j=1}^{n} \int_{B_{j}} x(s) m(\dd s)$ is clear and 
so, by virtue of the $\sigma$-additivity of $\int_{B} \| x(s) \| m(\dd s)$, we see that $\int_{B} x(s) m(\dd s)$ is $\sigma$-additive, i.e., 
\[
\begin{aligned}
&B = \sum\limits_{j=1} ^{\infty} B _{j} \text{ with } m (B _{j}) < \infty \text{ implies } \int_ {\sum\limits_{j=1} ^{\infty} B _{j}} x (s) m (\dd s) = \slim _{n \to \infty} \sum\limits_{j=1} ^{n} \int_ {B _{j}} x (s) m (\dd s).
\end{aligned}
\]

\begin{corollary}\label{cor:V.5.Corollary.2} Let $T$ be a bounded linear operator on a $B$-space $X$ into a $B$-space $Y$. If $x(s)$ is an $X$-valued Bochner $m$-integrable function, then $Tx(s)$ is a $Y$-valued Bochner $m$-integrable function, and 
\[
\int_{B} T x (s) m (\dd s) = T \int_{B} x (s) m (\dd s).
\]
\end{corollary}

\begin{proof}
 Let a sequence of finitely-valued functions $\{y_n(s)\}$ satisfy 
\[
\left\| y _{n} (s) \right\| \leqq \left\| x (s) \right\| (1 + n ^{-1}) \quad \text{ and } \quad \slim _{n \rightarrow \infty} y _{n} (s) = x (s) m\text{-a.e.} \]

Then, by the linearity and the continuity of $T$, we have $\int_{B} T y_n(s) m(\dd s) = T \int_{B} y_n(s) m(\dd s)$. We have, moreover, by the continuity of $T$, 
\[
\begin{aligned}&{\left\| T y _{n} (s) \right\| \leqq \left\| T \right\| \cdot \left\| y _{n} (s) \right\| \leqq \left\| T \right\| \cdot \left\| x (s) \right\| \cdot (1 + n ^{-1}) \quad \text{ and }}\\&{\qquad \slim _{n \rightarrow \infty} T y _{n} (s) = T x (s) m\text{-a.e.}}\end{aligned}
\]

Hence $Tx(s)$ is also Bochner $m$-integrable and 
\[
\begin{aligned}
&\int_{B} T x (s) m (\dd s)  = \slim\limits_{n \to \infty} \int_{B} T y _{n} (s) m (\dd s) = \slim\limits_{n \to \infty} T \int_{B} y _{n} (s) m (\dd s) = T \int_{B} x (s) m (\dd s).
\end{aligned}
\]
\end{proof}

\begin{theorem}[S. \textsc{Bochner} {\cite{bochner20}}]\label{thm:V.5.Theorem.2} Let $S$ be an $n$-dimensional euclidean space, $B$ the family of Baire sets of $S$, and $m(B)$ the Lebesgue measure of $B$. If $x(s)$ is Bochner $m$-integrable, and if $P(s_{0};\alpha)$ is the parallelopiped with centre at $s_{0}\in S$ and side length $2\alpha$, then we have the differentiation theorem 
\[
\slim _{\alpha \downarrow 0} (2 \alpha) ^{-n} \int_ {P (s _{0}; \alpha)} x (s) m (\dd s) = x (s _{0}) \text{ for } m\text{-a.e.} s _{0}.
\]
\end{theorem}

\begin{proof}
 Put 
\[
(2 \alpha) ^{-n} \int_ {P (s _{0}; \alpha)} x (s) m (\dd s) = D (x; s _{0}, \alpha). \]

If $\{x_{n}(s)\}$ is a sequence of finitely-valued functions such that $\|x_{n}(s)\| \leqq \|x(s)\| \cdot (1 + n^{-1})$ and $\slim\limits_{n\to \infty}x_n(s) = x(s)$ $m$-a.e., then 
\[
D (x; s _{0}, \alpha) - x (s _{0}) = D (x - x _{k}; s _{0}, \alpha) + D (x _{k}; s _{0}, \alpha) - x (s _{0}),
\]
and so 
\[
\begin{aligned}
\varlimsup\limits_{\alpha \downarrow 0} \| D (x; s _{0}, \alpha) - x (s _{0}) \|& \leqq \varlimsup\limits_{\alpha \downarrow 0} D (\| x - x _{k} \|; s _{0}, \alpha)
\\
& + \varlimsup\limits_{\alpha \downarrow 0} \| D (x _{k}; s _{0}, \alpha) - x _{k} (s _{0}) \| + \| x _{k} (s _{0}) - x (s _{0}) \|.
\end{aligned}
\]

The first term on the right is, by Lebesgue's theorem of differentiation of numerical functions, equal to $\|x(s_0) - x_k(s_0)\|$ $m$-a.e. The second term on the right is = $0$ $m$-a.e., since $x_k(s)$ is finitely-valued. Hence 
\[
\varlimsup\limits_{\alpha \rightarrow 0} \| D (x; s _{0}, \alpha) - x (s _{0}) \| \leqq 2 \| x _{k} (s _{0}) - x (s _{0}) \| \text{ for } m\text{-a.e.} s _{0}.
\]

Therefore, by letting $k \to \infty$, we obtain Theorem~\ref{thm:V.5.Theorem.2}. 
\end{proof}

\begin{remarku} Contrary to the case of numerical functions, a $B$-space-valued, $\sigma$-additive, $m$-absolutely continuous function need not necessarily be represented as a Bochner $m$-integral. This may be shown by a counter example. 
\end{remarku}

\begin{examplen}[$A$ Counter Example] Let $S = [0,1]$ and $\mathfrak{B}$ the family of Baire sets on $[0,1]$, and $m(B)$ the Lebesgue measure of $B \in \mathfrak{B}$. Consider the totality $m[1/3,2/3]$ of real-valued bounded functions $\xi = \xi(\theta)$ defined on the closed interval $[1/3,2/3]$ and normed by $\|\xi\| = \sup\limits_{\theta} |\xi(\theta)|$. We define an $m[1/3,2/3]$-valued function $x(s) = \xi(\theta; s)$ defined on $[0,1]$ as follows: 

the graph in $s$-$y$ plane of the real-valued function $y = y_{\theta}(s)$, which is the $\theta$-coordinate $\xi(\theta; s)$ of $x(s)$, is the polygonal line connecting the three points $(0, 0)$, $(\theta, 1)$ and $(1, 0)$ in this order. 
\end{examplen}

Then, if $s \neq s'$, we have Lipschitz' condition: 
\[
\left\| (s - s ^{\prime}) ^{-1} (x (s) - x (s ^{\prime})) \right\| = \sup\limits _{\theta} \left| (s - s ^{\prime}) ^{-1} (\xi (\theta ; s) - \xi (\theta ; s ^{\prime})) \right| \leqq 3.
\]

Thus, starting with the interval function $(x(s) - x(s'))$ taking values in $m$[$1$/$3$, $2$/$3$], we can define a $\sigma$-additive, $m$-absolutely continuous set function $x(B)$ defined for Baire set $B$ of [$0$, $1$]. 

If this function $x(B)$ is represented as a Bochner $m$-integral, then, by the preceding Theorem~\ref{thm:V.5.Theorem.2}, the function $x(s)$ must be strongly differentiable with respect to $s$ $m$-a.e. Let the corresponding strong derivative $x'(s)$ be denoted by $\eta(\theta; s)$ which takes values in $m$ [$1$/$3$, $2$/$3$]. Then, for 
every $\theta \in [1 / 3,2 / 3]$ and $m$-a.e. $s$, 
\[
\begin{aligned}
&0 = \lim\limits _{h \to 0} \| h ^{-1} (x (s + h) - x (s)) - x ^{\prime} (s) \|  \geqq \lim\limits _{h \to 0} | h ^{-1} (\xi (\theta ; s + h) - \xi (\theta ; s)) - \eta (\theta ; s) |.
\end{aligned}
\]

This proves that $\xi(\theta; s)$ must be differentiable in $s$ $m$-a.e. for all $\theta \in [1/3, 2/3]$. This is contradictory to the construction of $\xi(\theta; s)$. 

\yossub{References for Chapter V}

S. \textsc{Banach} \cite{banach9}, N. \textsc{Dunford}-J. \textsc{Schwartz} \cite{dunford44} and E. \textsc{Hille}-R. S. \textsc{Phillips} \cite{hille97}. 
   % V
% ==== AV ====
\yoschapterstar{Appendix to Chapter V. Weak Topologies and Duality in Locally Convex Linear Topological Spaces}
\label{chap:AV}

The present book is so designed that the reader may skip this appendix in the first reading and proceed directly to the following chapters. 


\yossection{1}{Polar Sets}
\label{sec:AV.1}

\begin{definitionu} Let $X$ be a locally convex linear topological space. For any set $M \subseteq X$, we define its (right) polar set $M^0$ by 
\begin{equation}\label{eq:AV.1.1}
M ^{0} = \left\{x ^{\prime} \in X ^{\prime}; \sup\limits _{x \in M} | \langle x, x ^{\prime} \rangle | \leqq 1 \right\}.
\end{equation}
\end{definitionu}

Similarly, for any set $M' \subseteq X'$, we define its (left) polar set $^0M'$ by 
\begin{equation}\label{eq:AV.1.2}
{ } ^{0} M ^{ \prime } = \{ x \in X ; \operatorname* { s u p } _ { x ^{ \prime } \in M ^{ \prime } } | \langle x , x ^{ \prime } \rangle | \leqq 1 \} = X \cap ( M ^{ \prime } ) ^{0} ,
\end{equation}
where we consider $X$ to be embedded in its bidual $(X_s')'$. 

A fundamental system of neighbourhoods of $0$ in the weak topology of $X$ is given by the system of sets of the form $^{0}M'$ where $M'$ ranges over arbitrary finite sets of $X'$. A fundamental system of neighbourhoods of $0$ in the weak* topology of $X'$ is given by the system of sets of the form $M^{0}$ where $M$ ranges over arbitrary finite sets of $X$. A fundamental system of neighbourhoods of $0$ in the strong topology of $X'$ is given by the system of sets of the form $M^{0}$ where $M$ ranges over arbitrary bounded sets of $X$. 

\begin{propositionu} $M^{0}$ is a convex, balanced set closed in the weak* topology of $X'$. 
\end{propositionu}

\begin{proof}
 For any fixed $x \in X$, the linear functional $f(x') = \langle x, x' \rangle$ is continuous in the weak* topology of $X'$. Thus $M^0 = \bigcap\limits_{m \in M} \{m\}^0$ is closed in the weak* topology of $X'$. The balanced convexity of $M^0$ is clear. 
\end{proof}

\yossub{An Application of Tychonov's Theorem}

\begin{theorem}\label{thm:AV.1.Theorem.1} Let $X$ be a locally convex linear topological space, and $A$ a convex, balanced neighbourhood of $0$ of $X$. Then $A^{0}$ is compact in the weak* topology of $X'$. 
\end{theorem}

\begin{proof}
 Let $p(x)$ be the Minkowski functional of $A$. Consider, for each $x \in X$, a sphere $S_x = \{z \in C; |z| \leqq p(x)\}$ and the topological product $S = \prod\limits_{x \in X} S_x$. $S$ is compact by Tychonov's theorem. Any element $x' \in X'$ is determined by the set of values $x'(x) = \langle x, x' \rangle$, $x \in X$. Since $x \in (p(x) + \varepsilon)A$ for any $\varepsilon > 0$, we see that $x' \in X'$ implies $\langle x, x' \rangle = \langle (p(x) + \varepsilon)a, x' \rangle$ with a certain $a \in A$. Thus $x \in A^0$ implies that $|x'(x)| \leqq p(x) + \varepsilon$, that is, $x'(x) \in S_x$. Hence we may consider $A^0$ as a subset of $S$. Moreover, it is easy to verify that the topology induced on $A^0$ by the weak* topology of $X'$ is the same as the topology induced on $A^0$ in the Cartesian product topology of $S = \prod\limits_{x \in X} S_x$. 

Hence it is sufficient to prove that $A^0$ is a closed subset of $S$. Suppose $y = \prod\limits_{x \in X} y(x)$ is an element of the weak* closure of $A^0$ in $S$. Consider any $\varepsilon > 0$ and any $x_1, x_2 \in X$. The set of all $u = \prod\limits_{x \in X} u(x) \in S$ such that 
$|u(x_{1})-y(x_{1})|<\varepsilon, |u(x_{2})-y(x_{2})|<\varepsilon$ and $|u(x_{1}+x_{2})-y(x_{1}+x_{2})|<\varepsilon$ is a neighbourhood of $y$ in $S$. This neighbourhood contains some point $x'\in A^{0}$ and, since $x'$ is a continuous linear functional on $X$, we have 
\[
\begin{aligned}
\left| y (x _{1} + x _{2}) - y (x _{1}) - y (x _{2}) \right|& \leqq \left| y (x _{1} + x _{2}) - \langle x _{1} + x _{2}, x ^{\prime} \rangle \right|
\\
\quad + \left| \langle x _{1}, x ^{\prime} \rangle - y (x _{1}) \right| + \left| \langle x _{2}, x ^{\prime} \rangle - y (x _{2}) \right|& < 3 \varepsilon .
\end{aligned}
\]

This proves that $y(x_1 + x_2) = y(x_1) + y(x_2)$. Similarly, we prove that $y(\beta x) = \beta y(x)$, and so $y$ defines a linear functional on $X$. By the fact that $y = \prod\limits_{x \in X} y(x) \in S$, we know that $|y(x)| \leqq p(x)$. Since $p(x)$ is continuous, $y(x)$ is a continuous linear functional, i.e., $y \in X'$. On the other hand, since $y$ is a weak* accumulation point of $A^0$, there exists, for any $\varepsilon > 0$ and $a \in A$, an $x' \in A^0$ such that $|y(a) - \langle a, x' \rangle| \leqq \varepsilon$. Hence $|y(a)| \leqq |\langle a, x' \rangle| + \varepsilon \leqq 1 + \varepsilon$, and so $|y(a)| \leqq 1$, that is, $y \in A^0$. 
\end{proof}

\begin{corollaryu} The unit sphere $S^{*} = \{x' \in X'; \| x' \| \leqq 1\}$ of the dual space $X_{s}'$ of a normed linear space $X$ is compact in the weak* topology of $X'$. 
\end{corollaryu}

\yossub{An Application of Mazur's Theorem}

\begin{theorem}\label{thm:AV.1.Theorem.2} Let $M$ be a convex, balanced closed set of a locally convex linear topological space $X$. Then $M = {}^0 (M^0)$. 
\end{theorem}

\begin{proof}
 It is clear that $M \subseteq {}^0 (M^0)$. If there exists an $x_0 \in {}^0 (M^0) - M$, then, by Mazur's theorem $3$ in Chapter~\ref{chap:IV}, \ref{sec:IV.6}, there exists an $x_0' \in X'$ such that $\langle x_0, x_0' \rangle > 1$ and $|\langle x, x_0' \rangle| \leqq 1$ for all $x \in M$. The last inequality shows that $x_0' \in M^0$ and so $x_0$ cannot belong to ${}^0 (M^0)$. 
\end{proof}


\yossection{2}{Barrel Spaces}
\label{sec:AV.2}

\begin{definitionu} In a locally convex linear topological space $X$, any convex, balanced and absorbing closed set is called a barrel (tonneau in Bourbaki's terminology). $X$ is called a barrel space if each of its barrels is a neighbourhood of $0$. 
\end{definitionu}

\begin{theorem}\label{thm:AV.2.Theorem.1} $A$ locally convex linear topological space $X$ is a barrel space if $X$ is not of the first category. 
\end{theorem}

\begin{proof}
 Let $T$ be a barrel in $X$. Since $T$ is absorbing, $X$ is the union of closed sets $nT = \{nt; t \in T\}$, where $n$ runs over positive integers. Since $X$ is not of the first category, at least one of the $(nT)'s$ contains an interior point. Hence $T$ itself contains an interior point $x_{0}$. If $x_{0} = 0$, $T$ is a neighbourhood of $0$. If $x_{0} \neq 0$, then $-x_{0} \in T$ by the fact that $T$ is balanced. Thus $-x_{0}$ is an interior point of $T$ with $x_{0}$. This proves that the convex set $T$ contains $0 = (x_{0} - x_{0})/2$ as an interior point. 
\end{proof}

\begin{corollary}\label{cor:AV.2.Corollary.1} All locally convex $F$-spaces and, in particular, all $B$-spaces and $\mathfrak{E}(\mathbb{R}^{n})$ are barrel spaces. 
\end{corollary}

\begin{corollary}\label{cor:AV.2.Corollary.2} The metric linear space $\mathfrak{D}_K(\mathbb{R}^{n})$ is a barrel space. 
\end{corollary}

\begin{proof}
 Let $\{\varphi_{k}\}$ be a Cauchy sequence with respect to the distance 
\[
\dist (\varphi , \psi) = \sum\limits_{m=0} ^{\infty} 2 ^{-m} \frac {p _{m} (\varphi - \psi)}{1 + p _{m} (\varphi - \psi)}, \text{ where } p _{m} (\varphi) = \sup\limits _{| j | \leqq m, x \in K} | D ^{j} \varphi (x) |. \]

For any differential operator $D^{j}$, the sequence $\{D^{j}\varphi_{k}(x)\}$ is equi-continuous and equi-bounded, that is, 
\[
\lim\limits _ {| x ^{1} - x ^{2} | \downarrow 0} \sup\limits _{k \geqq 1} \left| D ^{j} \varphi_{k} (x ^{1}) - D ^{j} \varphi_{k} (x ^{2}) \right| = 0 \text{ and } \sup\limits _{x \in K, k \geqq 1} \left| D ^{j} \varphi_{k} (x) \right| < \infty .
\]

This we see from the fact that, for any coordinate $x_s$, $\sup\limits_{x \in K, k \geqq 1} \left| \frac{\partial}{\partial x_s} D^j \varphi_k(x) \right| < \infty$. Hence, by the Ascoli-Arzela theorem, there exists a subsequence of $\{D^j \varphi_{k'}(x)\}$ which converges uniformly on $K$. By the diagonal method, we may choose a subsequence $\{\varphi_{k''}(x)\}$ of $\{\varphi_k(x)\}$ such that, for any differential operator $D^j$, the sequence $\{D^j \varphi_{k''}(x)\}$ converges uniformly on $K$. Thus 
\[
\lim\limits _ {k ^{\prime \prime} \rightarrow \infty} D ^{j} \varphi_ {k ^{\prime \prime}} (x) = D ^{j} \varphi (x) \text{ where } \varphi (x) = \lim\limits _ {k ^{\prime \prime} \rightarrow \infty} \varphi_ {k ^{\prime \prime}} (x),
\]
and these limit relations hold uniformly on $K$. Hence the metric space $\mathfrak{D}_K(\mathbb{R}^{n})$ is complete and so it is not of the first category. 
\end{proof}

\begin{remarku} \begin{enumerate}[(i)]
\item The above proof shows that a bounded set of $\mathfrak{D}(\mathbb{R}^{n})$ is relatively compact in the topology of $\mathfrak{D}(\mathbb{R}^{n})$. For, a bounded set $B$ of $\mathfrak{D}(\mathbb{R}^{n})$ is contained in some $\mathfrak{D}_K(\mathbb{R}^{n})$ where $K$ is a compact set of $R_n^n$, and, moreover, the boundedness condition of $B$ implies the equi-boundedness and equi-continuity of $\{D^j\varphi; \varphi \in B\}$ for every $D^j$.
\item Similarly, we see that any bounded set of $\mathfrak{E}(\mathbb{R}^{n})$ is a relatively compact set of $\mathfrak{E}(\mathbb{R}^{n})$.
\end{enumerate}
\end{remarku}

\begin{corollary}\label{cor:AV.2.Corollary.3} $\mathfrak{D}(\mathbb{R}^{n})$ is a barrel space. 
\end{corollary}

\begin{proof}
 $\mathfrak{D}(\mathbb{R}^{n})$ being an inductive limit of $\{\mathfrak{D}_K(\mathbb{R}^{n})\}$ when $K$ ranges over compact subsets of $\mathbb{R}^{n}$, the present corollary is a consequence of the following 
\end{proof}

\begin{propositionu} Let a locally convex linear topological space $X$ be an inductive limit of its barrel subspaces $X_{\alpha}, \alpha \in A$. Then $X$ itself is a barrel space. 
\end{propositionu}

\begin{proof}
 Let $V$ be a barrel of $X$. By the continuity of the identical mapping $T_{\alpha}: x \to x$ of $X_{\alpha}$ into $X$, the inverse image $T_{\alpha}^{-1}(V) = V \cap X_{\alpha}$ is closed with $V$. Thus $V \cap X_{\alpha}$ is a barrel of $X_{\alpha}$. $X_{\alpha}$ being a barrel space, $V \cap X_{\alpha}$ is a neighbourhood of $0$ of $X_{\alpha}$. $X$ being an inductive limit of $X_{\alpha}$'s, $V$ must be a neighbourhood of $0$ of $X$. 
\end{proof}

\begin{theorem}\label{thm:AV.2.Theorem.2} Let $X$ be a barrel space. Then the mapping $x \to Jx$ of $X$ into $(X_{s}^{\prime})_{s}^{\prime}$, defined in Chapter~\ref{chap:IV}, \ref{sec:IV.8}, is a topological mapping of $X$ onto JX, where the topology of JX is provided with the relative topology of JX as a subset of $(X_{s}^{\prime})_{s}^{\prime}$. 
\end{theorem}

\begin{proof}
 Let $B'$ be a bounded set of $X_s'$. Then the polar set $(B')^0 = \{x'' \in (X_s')'\}; \sup\limits_{x' \in B'} |\langle x', x'' \rangle| \leqq 1\}$ of $B'$ is a neighbourhood of $0$ of $(X_s')_s'$, and it is a convex, balanced and absorbing set closed in $\cdot(X_s')'$. Thus $(B')^0 \cap X = {}^0(B')$ is a convex, balanced and absorbing set of $X$. As a (left) polar set, ${}^0(B')$ is closed in the weak topology of $X$, and hence ${}^0(B')$ is closed in the original topology of $X$. Thus ${}^0(B') = (B')^0 \cap X$ is a barrel of $X$, and so it is a neighbourhood of $0$ of $X$. Therefore, the mapping $x \to Jx$ of $X$ into $(X_s')'$ is continuous, because the topology of $(X_s')'$ is defined by a fundamental system of neighbourhoods of $0$ of the form $(B')^0$, where $B'$ ranges over bounded sets of $X'$. 

Let, conversely, $U$ be a convex, balanced and closed neighbourhood of $0$ of $X$. Then, by the preceding section, $U = {}^0 (U^0)$. Thus $JU = JX \cap (U^0)^0$. On the other hand, $U^0$ is a bounded set of $X_s'$, since, for any bounded set $B$ of $X$, there exists an $\alpha > 0$ such that $\alpha B \subseteq U$ and so $(\alpha B)^0 \supseteq U^0$. Hence $(U^0)^0$ is a neighbourhood of $0$ of $(X_s')_s'$. Thus the image $JU$ of the neighbourhood $U$ of $0$ of $X$ is a neighbourhood of $0$ of $JX$ provided with the relative topology of $JX$ as a subset of $(X_s')_s'$. 
\end{proof}


\yossection{3}{Semi-reflexivity and Reflexivity}
\label{sec:AV.3}

\begin{definition}\label{def:AV.3.Definition.1} $A$ locally convex linear topological space $X$ is called semi-reflexive if every continuous linear functional on $X_s'$ is given by 
\begin{equation}\label{eq:AV.3.1}
\langle x, x ^{\prime} \rangle , \text{ with a certain } x \in X.
\end{equation}
\end{definition}

Thus $X$ is semi-reflexive iff 
\begin{equation}\label{eq:AV.3.2}
X _{w} = (X _{s} ^{\prime}) _ {x ^{*}} ^{\prime}.
\end{equation}

\begin{definition}\label{def:AV.3.Definition.2} $A$ locally convex linear topological space $X$ is called reflexive if 
\begin{equation}\label{eq:AV.3.3}
X = (X _{s} ^{\prime}) _{s} ^{\prime}.
\end{equation}
\end{definition}

By Theorem~\ref{thm:AV.3.Theorem.2} in the preceding section, we have 

\begin{proposition}\label{prop:AV.3.Proposition.1} $A$ semi-reflexive space $X$ is reflexive if $X$ is a barrel space. It is also clear, from Definition~\ref{def:AV.3.Definition.2}, that we have 
\end{proposition}

\begin{proposition}\label{prop:AV.3.Proposition.2} The strong dual of a reflexive space is reflexive. 
\end{proposition}

\begin{theorem}\label{thm:AV.3.Theorem.1} $A$ locally convex linear topological space $X$ is semi-reflexive iff every closed, convex, balanced and bounded set of $X$ is compact in the weak topology of $X$. 
\end{theorem}

\begin{proof}
 Let $X$ be semi-reflexive, and $T$ a closed, convex, balanced and bounded set of $X$. Then, by Theorem~\ref{thm:AV.3.Theorem.2} in Section~\hyperref[sec:AV.1]{$1$} of this Appendix, $T = {}^0 (T^0)$. $T$ being a bounded set of $X$, $T^0$ is a neighbourhood of $0$ of $X_s'$. Thus, by Theorem~\ref{thm:AV.3.Theorem.1} in Section~\hyperref[sec:AV.1]{$1$} of this Appendix, $(T^0)^0$ is compact in the weak* topology of $(X_s')'$. Hence, by the semi-reflexivity of $X$, $T = {}^0 (T^0)$ is compact in the weak topology of $X$. 

We next prove the sufficiency part of Theorem~\ref{thm:AV.3.Theorem.1}. Take any $x'' \in (X_s')'$. The strong continuity of $x''$ on $X_s'$ implies that there exists a bounded set $B$ of $X$ such that 
\[
\left| \langle x ^{\prime}, x ^{\prime \prime} \rangle \right| \leqq 1 \text{ whenever } x ^{\prime} \in B ^{0}, \text{ that is, } x ^{\prime \prime} \in (B ^{0}) ^{0}.
\]

We may assume that $B$ is a convex, balanced and closed set of $X$. Thus, by the hypothesis of Theorem~\ref{thm:AV.3.Theorem.1}, $B$ is a compact set in the weak topology of $X$. Hence $B = B^{wa}$ where $B^{wa}$ denotes the closure of $B$ in the weak topology of $X$. Since $X_w$ is embedded in $(X_s')_{w^*}$ as a linear topological subspace, we must have $(B^0)^0 \supseteq B^{wa} = B$. Therefore we have to show that $x''$ is an accumulation point of $B$ in $(X_s')_{w^*}$. Consider the mapping $x \to \varphi(x) = \{\langle x, x_1', \ldots, \langle x, x_n'\rangle\}$ of $X$ into $l^2(n)$, where $x_1', \ldots, x_n' \in X'$. The image $\varphi(B)$ is convex, balanced and compact, since $B$ is convex, balanced and weakly compact. If $\{\langle x_1', x''\rangle, \ldots, \langle x_n', x''\rangle\}$ does not belong to $\varphi(B)$, then, by Mazur's theorem, there would exist a point $\{c_1, \ldots, c_n\} \in l^2(n)$ such that $\sup\limits_{b \in B} \left| \sum\limits_i c_i \langle b, x_i' \rangle \right| \leqq 1$ and $\left( \sum\limits_i c_i \langle x_i', x''\rangle \right) > 1$, proving that $\sum\limits_i c_i x_i' \in B^0$ and $x''$ cannot belong to $B^{00}$. 
\end{proof}

\begin{theorem}\label{thm:AV.3.Theorem.2} $A$ locally convex linear topological space $X$ is reflexive iff it is a barrel space and every closed, convex, balanced and bounded set of $X$ is compact in the weak topology of $X$. In particular, $\mathfrak{D}(\mathbb{R}^{n})$ and $\mathfrak{E}(\mathbb{R}^{n})$ are reflexive. 
\end{theorem}

\begin{proof}
 The sufficiency part is proved already. We shall prove that the first condition of Theorem~\ref{thm:AV.3.Theorem.2} is necessary. 

Let $T$ be a barrel of $X$. We shall prove that $T$ absorbs any bounded set $B$ of $X$ so that $B^0 \supseteq \alpha T^0, \alpha > 0$. $B^0$ being a neighbourhood of $0$ of $X_s'$, we see that $T^0$ is a bounded set of $X_s'$. By the Proposition and Theorem~\ref{thm:AV.3.Theorem.2} in Section~\hyperref[sec:AV.1]{$1$} of this Appendix, we have $T = {}^0(T^0)$. By the hypothesis that $X$ is reflexive, we have ${}^0(T^0) = (T^0)^0$ and so $T = (T^0)^0$. Therefore we have proved that the barrel $T$ is a neighbourhood of $0$ of $X = (X_s')_s'$. Hence $X$ is a barrel space. 

By hypothesis, the closed, convex and balanced set $K = \Conv\left(\bigcup\limits_{|\alpha| \leqq 1} \alpha B\right)^a$ 
is compact in the weak topology of $X$. Here we denote by $\Conv(N)^a$ the closure in $X$ of the convex closure (see p. $28$) $\Conv(N)$ of $N$. Set $Y = \bigcup\limits_{n=1}^{\infty} nK$ and let $p(x)$ be the Minkowski functional of $K$. Then since $K$ is $w$-compact in $Y$, $p(x)$ defines a norm of $Y$. That is, the system $\{\alpha K\}$ with $\alpha > 0$ defines a fundamental system of neighbourhoods of the normed linear space $Y$, and $Y$ is a $B$-space since $K$ is $w$-compact. Hence $Y$ is a barrel space. On the other hand, since $K$ is a bounded set of $X$, the topology of $Y$ defined by the norm $p(x)$ is stronger than the relative topology of $Y$ as a subset of $X$. As a barrel of $X$, $T$ is closed in $X$. Hence $T \cap Y$ is closed in $Y$ with respect to the topology defined by the norm $p(x)$. Therefore $T \cap Y$ is a barrel of the $B$-space $Y$, and so $T \cap Y$ is a neighbourhood of $0$ of the $B$-space $Y$. We have thus proved that $T \cap Y$ and, a fortiori, $T$ both absorb $K \supseteq B$. 
\end{proof}


\yossection{4}{The Eberlein-Shmulyan Theorem}
\label{sec:AV.4}

This theorem is very important in view of its applications. 

\begin{theoremn}[\textsc{Eberlein}-\textsc{Shmulyan}] A $B$-space $X$ is reflexive iff it is locally sequentially weakly compact; that is, $X$ is reflexive iff every strongly bounded sequence of $X$ contains a subsequence which converges weakly to an element of $X$.

For the proof we need two Lemmas: 

\end{theoremn}
\begin{lemma}\label{lem:AV.4.Lemma.1} If the strong dual $X_s'$ of a $B$-space $X$ is separable, then $X$ itself is separable. 
\end{lemma}

\begin{lemma}[S. \textsc{Banach}]\label{lem:AV.4.Lemma.2} $A$ linear subspace $M'$ of the dual space $X'$ of a $B$-space $X$ is weakly* closed iff $M'$ is boundedly weakly* closed; that is, $M'$ is weakly* closed iff $M'$ contains all weak* accumulation points of every strongly bounded subset of $M'$. 
\end{lemma}

\yoslabel{Lemma}{$1$}\label{lem:AV.4.Lemma.1a}is already proved as a Lemma in Chapter~\ref{chap:V}, \ref{sec:V.2}. For Lemma~\ref{lem:AV.4.Lemma.2}, we have only to prove its ``if'' part. It reads as follows. 

\begin{proof}
 (E. \textsc{Hille}-R. S. \textsc{Phillips} \cite{hille97}). We remark that $M'$ is strongly closed by the hypothesis. Let $x_0' \in M'$. Then we can prove that, for each constant $C$ satisfying the condition $0 < C < \inf\limits_{x' \in M'} \|x' - x_0'\|$, there exists an $x_0 \in X$ with $\|x_0\| \leqq 1 / C$ such that 
\begin{equation}\label{eq:AV.4.1}
\langle x _{0}, x _{0} ^{\prime} \rangle = 1 \quad \text{ and } \quad \langle x _{0}, x ^{\prime} \rangle = 0 \quad \text{ for all } \quad x ^{\prime} \in M ^{\prime}. \end{equation}

Thus the strongly closed set $M'$ must contain all of its weak* accumulation points. 


To prove the existence of $x_0$, we choose an increasing sequence of numbers $\{C_n\}$ such that $C_1 = C$ and $\lim\limits_{n \to \infty} C_n = \infty$. Then there exists a finite subset $\sigma_1$ of the unit sphere $S = \{x \in X; \|x\| \leqq 1\}$ such that $\|x' - x_0'\| \leqq C_2$ and $\sup\limits_{x \in \sigma_1} |\langle x, x' \rangle - \langle x, x_0' \rangle| \leqq C_1$ implies $x' \in M'$. 

If not, there would exists, corresponding to each finite subset $\sigma$ of $S$, an $x_{\sigma}' \in M$ such that 
\[
\left\| x _{\sigma} ^{\prime} - x _{0} ^{\prime} \right\| \leqq C _{2} \quad \text{ and } \quad \sup\limits _{x \in \sigma} | \langle x, x _{\sigma} ^{\prime} \rangle - \langle x, x _{0} ^{\prime} \rangle | \leqq C _{1}.
\]

We order the sets $\sigma$ by inclusion relation and denote the weak* closure of the set $\{x_{\sigma}^{\prime};\sigma^{\prime}\geqq \sigma \}$ by $N_{\sigma}^{\prime}$. It is clear that $N_{\sigma}^{\prime}$ enjoys the finite intersection property. On the other hand, since $M^{\prime}$ is boundedly weakly* closed, the Corollary of Theorem~\hyperref[thm:AV.1.Theorem.1]{$1$} in this Appendix, $1$, implies that the set 
\[
M _ {c ^{\prime}} ^{\prime} = \{x ^{\prime} \in M ^{\prime}; \| x ^{\prime} \| \leqq C ^{\prime} \}
\]
is weakly* compact. Hence $N_{\sigma}' \subseteq M_{c'}'$ for $C' = C_2 + \|x_0'|$, and so there exists an $x_1' \in \bigcap\limits_{\sigma} N_{\sigma}' \subseteq M'$. Hence we have $\sup\limits_{x \in S} |\langle x, x_1' \rangle - \langle x, x_0' \rangle| \leqq C_1$ and so $\|x_1' - x_0'| \leqq C_1$, contrary to the hypothesis that $0 < C_1 < \inf\limits_{x' \in M'} \|x' - x_0'|$. 

By a similar argument, we successively prove the existence of a sequence of finite subsets $\sigma_{1},\sigma_{2},\ldots$ of $S$ such that 
\[
\begin{yoslines}
&\| x ^{\prime} - x _{0} ^{\prime} \| \leqq C _{k} \quad \text{ and }  \sup\limits _ {x \in \sigma_{i}} | \langle x, x ^{\prime} \rangle - \langle x, x _{0} ^{\prime} \rangle | \leqq C _{i} (i = 1, 2, \dots , k - 1) \\
&\text{ do not imply } x ^{\prime} \in M ^{\prime}.
\end{yoslines}
\]

Thus, since $\lim\limits_{i\to \infty}C_i = \infty$, we see that $x^{\prime}\bar{\in} M^{\prime}$ if 
\[
\left| \langle x, x ^{\prime} \rangle - \langle x, x _{0} ^{\prime} \rangle \right| \leqq C \quad \text{ for all } \quad x \in (C / C _{j}) \sigma_{j} \quad (j = 1, 2, \dots).
\]

Let $\{x_{n}\}$ be a sequence which successively exhausts the sets $(C / C_j)\sigma_j$ ( $j = 1,2,\ldots$). Then $\lim\limits_{n\to \infty}x_n = 0$ and so $L(x') = \{\langle x_n,x'\rangle \}$ is a bounded linear transformation of $X_s'$ into the $B$-space $(c_0)$. We know that the point $\{\langle x_n,x_0'\rangle \} \in (c_0)$ lies at a distance $>C$ from the linear subspace $L(M')$. Thus, by Corollary of Theorem~\hyperref[thm:IV.6.Theorem.3]{$3$} in Chapter~\ref{chap:IV}, \ref{sec:IV.6}, there exists a continuous linear functional $\{\alpha_{n}\} \in (c_0)' = (l^1)$ such that 
\begin{equation}\left\| \left\{\alpha_{n} \right\} \right\| = \sum\limits_{n=1} ^{\infty} \left| \alpha_{n} \right| \leqq 1 / C, \sum\limits_{n=1} ^{\infty} \alpha_{n} \langle x _{n}, x _{0} ^{\prime} \rangle = 1 \quad \text{ and } \sum\limits_{n=1} ^{\infty} \alpha_{n} \langle x _{n}, x ^{\prime} \rangle = 0 \quad \text{ for all } \quad x ^{\prime} \in M ^{\prime}. \notag\end{equation}

The element $x_0 = \sum\limits_{n=1}^{\infty} \alpha_n x_n$ clearly satisfies condition \eqref{eq:AV.4.1}.
\end{proof}

\begin{corollaryu} Let $\langle x', x_0'' \rangle = F(x')$ be a linear functional defined on the dual space $X'$ of a $B$-space $X$. If $N(F) = N(x_0'') = \{x' \in X'; F(x') = 0\}$ is weakly* closed, then there exists an element $x_0$ such that 
\begin{equation}\label{eq:AV.4.2}
F (x ^{\prime}) = \langle x ^{\prime}, x _{0} ^{\prime \prime} \rangle = \langle x _{0}, x ^{\prime} \rangle \quad \text{ for all } \quad x ^{\prime} \in X ^{\prime}.
\end{equation}
\end{corollaryu}

\begin{proof}
 We may assume that $N(F) \neq X'$. Otherwise we can take $x_{0} = 0$. Let $x_{0}' \in X'$ be such that $F(x_{0}') = 1$. By \eqref{eq:AV.4.1} of the preceding Lemma~\ref{lem:AV.4.Lemma.2}, there exists an $x_{0} \in X$ such that 
\begin{equation}\label{eq:AV.4.3}
\langle x _{0}, x _{0} ^{\prime} \rangle = 1 \quad \text{ and } \quad \langle x _{0}, x ^{\prime} \rangle = 0 \quad \text{ for all } \quad x ^{\prime} \in N (F). \end{equation}

Hence, for any $x' \in X'$, the functional 
\[
x ^{\prime} - F \left(x ^{\prime}\right) x _{0} ^{\prime} = y ^{\prime} \in X ^{\prime}
\]
satisfies $F(y') = 0$, i.e. $y' \in N(F)$. Therefore, by \eqref{eq:AV.4.3}, we obtain \eqref{eq:AV.4.2}. 
\end{proof}

\begin{proof}[Proof of the Theorem.]
``Only if'' part. Let $\{x_{n}\}$ be a sequence of $X$ such that $\|x_{n}\|=1$. The strong closure $X_{0}$ of the subspace spanned by $\{x_{n}\}$ is a separable $B$-space. Being a $B$-space, $X_{0}$ is a barrel space. We shall show that $X_{0}$ is reflexive. Any strongly closed, bounded set $B_{0}$ of $X_{0}$ is also a strongly closed, bounded set of $X$ and hence $B_{0}$ is compact in the weak topology of $X$ by the reflexivity of $X$. But, as a strongly closed linear subspace of $X$, $X_{0}$ is closed in the weak topology of $X$ (see Theorem~\ref{thm:IV.6.Theorem.3} in Chapter~\ref{chap:IV}, \ref{sec:IV.6}). Hence $B_{0}$ is compact in the weak topology of $X_{0}$. Thus $X_{0}$ is reflexive by Theorem~\hyperref[thm:AV.3.Theorem.2]{$2$} in the preceding section. We have thus $X_{0}=((X_{0})_{s}^{\prime})_{s}^{\prime}$. By Lemma~\ref{lem:AV.4.Lemma.1a} above, $(X_{0})_{s}^{\prime}$ is thus separable. Let $\{x_{n}^{\prime}\}$ be strongly dense in $(X_{0})_{s}^{\prime}$. Then the weak topology of $X_{0}$ is defined by an enumerable sequence of semi-norms $p_{m}(x)=|\langle x,x_{m}^{\prime}\rangle|(m=1,2,\ldots)$. Hence it is easy to see that the sequence $\{x_{n}\}$, which is compact in the weak topology of $X_{0}$, is sequentially weakly compact in $X_{0}$ and in $X$ as well. We have only to choose a subsequence $\{x_{n^{\prime}}\}$ of $\{x_{n}\}$ such that finite $\lim\limits_{n\to\infty}\langle x_{n^{\prime}},x_{m}^{\prime}\rangle$ exists for $m$=$1$,$2$, $\ldots$.

``If'' part. Let $M$ be a bounded set of $X$, and assume that every infinite sequence of $M$ contains a subsequence which converges weakly to an element of $X$. We have to show that the closure $\overline{M}$ of $M$ in $X$ in the weak topology of $X$ is weakly compact in $X$. For, then the barrel space $X$ is reflexive by Theorem~\hyperref[thm:AV.3.Theorem.2]{$2$} in the preceding section. 

Since $X_w \subseteq (X_s')'_{w^*}$, we have $\overline{M} = \overline{\overline{M}} \cap X_w$, where $\overline{\overline{M}}$ is the closure of $\overline{M}$ in the weak* topology of $(X_s')'$. Let $S_r'$ be the sphere of $X_s'$ of radius $r > 0$ and centre $0$. By the correspondence 
\begin{equation}\overline {{\overline {{M}}}} \ni m \leftrightarrow \{\langle x ^{\prime}, m \rangle ; \| x ^{\prime} \| \leqq 1 \} \in \prod\limits_ {x ^{\prime} \in S _{1} ^{\prime}} I _ {x ^{\prime}}, \text{ where } I _ {x ^{\prime}} = \{z; | z | \leqq \sup\limits _ {m \in \overline {{M}}} | \langle x ^{\prime}, m \rangle | \}, \notag\end{equation}
$\overline{M}$ may be identified with a closed subset of the topological product $\prod\limits_{x' \in S_i'} I_{x'}$. By Tychonov's theorem, $\prod\limits_{x' \in S_i'} I_{x'}$ is compact and hence $\overline{\overline{M}}$ is compact in the weak* topology of $(X_s')'$. Hence we have only to show $\overline{\overline{M}} \subseteq X_w$. 

Let $x_0'' \in (X_s)'$ be an accumulation point of the set $\overline{M}$ in the weak* topology of $(X_s)'$. To prove that $x_0'' \in X_w$ we have only to show that the set $N(x_0'') = \{x' \in X'; \langle x', x_0'' \rangle = 0\}$ is weakly* closed. For, then, by the above Corollary, there exists an $x_0 \in X$ such that $\langle x', x_0'' \rangle = \langle x_0, x' \rangle$ for all $x' \in X'$. We shall first show that 
for every finite set $x_1', x_2', \ldots, x_n'$ of $X'$, there exists a $z \in \overline{M}$ such that $\langle x_j', x_0'' \rangle = \langle z, x_j' \rangle (j = 1, 2, \ldots, n)$. 

The proof is as follows. Since $x_0''$ is in the weak* closure of $\overline{M}$, there is an element $z_m \in \overline{M}$ such that 
\begin{equation}\label{eq:AV.4.4}
\left| \left\langle z _{m}, x _{j} ^{\prime} \right\rangle - \left\langle x _{j}, x _{0} ^{\prime \prime} \right\rangle \right| \leqq 1 / m (j = 1, 2, \dots , n).
\end{equation}

By hypothesis there exists a subsequence of $\{z_m\}$ which converges weakly to an element $z \in X$ and so $z \in \overline{M}$, since the sequential weak closure of $\overline{M}$ is contained in $\overline{M}$. We have thus \eqref{eq:AV.4.4}. 

Now, by Lemma~\ref{lem:AV.4.Lemma.2}, $N(x_0'')$ is weakly* closed if, for every $r > 0$, the set $N(x_0'') \cap S_r'$ is weakly* closed. Let $y_0'$ be in the weak* closure of $N(x_0'') \cap S_1'$. We have to show that $y_0' \in N(x_0'') \cap S_1'$. To this purpose, we choose an arbitrary $\varepsilon > 0$ and construct three sequences $\{z_n\} \subseteq \overline{M}$, $\{x_n\} \subseteq M$ and $\{y_n'\} \subseteq N(x_0'') \cap S_1'$ as follows: By \eqref{eq:AV.4.4} we can choose a, $z_1 \in \overline{M}$ such that $\langle z_1, y_0' \rangle = \langle y_0', x_0'' \rangle$. $z_1$ being in the weak closure of $M$, there exists an $x_1 \in M$ such that $|\langle x_1, y_0' \rangle - \langle z_1, y_0' \rangle| \leqq \varepsilon/4$. $y_0'$ being in the weak* closure of $N(x_0'') \cap S_1'$, there exists a $y_1' \in N(x_0'') \cap S_1'$ such that $|\langle x_1, y_1' \rangle - \langle x_1, y_0' \rangle| \leqq \varepsilon/4$. Repeating the argument and remembering \eqref{eq:AV.4.4}, we obtain $\{z_n\} \subseteq \overline{M}$, $\{x_n\} \subseteq M$ and $\{y_n'\} \subseteq N(x_0'') \cap S_1'$ such that 
\begin{equation}\label{eq:AV.4.5}
\begin{yoslines}
&\langle z _{1}, y _{0} ^{\prime} \rangle = \langle y _{0} ^{\prime}, x _{0} ^{\prime \prime} \rangle , \\
&\langle z _{n}, y _{m} ^{\prime} \rangle = \langle y _{m} ^{\prime}, x _{0} ^{\prime \prime} \rangle = 0  (m = 1, 2, \dots , n - 1), \\
&| \langle x _{n}, y _{m} ^{\prime} \rangle - \langle z _{n}, y _{m} ^{\prime} \rangle | \leqq \varepsilon / 4  (m = 0, 1, \dots , n - 1), | \langle x _{i}, y _{n} ^{\prime} \rangle - \langle x _{i}, y _{0} ^{\prime} \rangle | \leqq \varepsilon / 4  (i = 1, 2, \dots , n).
\end{yoslines}
\end{equation}

Thus we have 
\begin{equation}\label{eq:AV.4.6}
\left| \langle y _{0} ^{\prime}, x _{0} ^{\prime \prime} \rangle - \langle x _{i}, y _{n} ^{\prime} \rangle \right| \leqq \varepsilon / 4 + \varepsilon / 4 = \varepsilon / 2 (i = 1, 2, \dots , n).
\end{equation}

Since $\{x_{n}\} \subseteq M$, there exists a subsequence of $\{x_{n}\}$ which converges weakly to an element $x \in \overline{M}$. Without losing the generality, we may assume that the sequence $\{x_{n}\}$ itself converges weakly to $x \in \overline{M}$. Hence, from \eqref{eq:AV.4.5}, $|\langle x, y_{m}' \rangle| \leqq \varepsilon / 4$. From $w - \lim\limits_{n \to \infty} x_{n} = x$ and Mazur's theorem $2$ in Chapter~\ref{chap:V}, \ref{sec:V.1}, there exists a convex combination $u = \sum\limits_{j=1}^{n} \alpha_{j} x_{j} \left( \alpha_{j} \geqq 0, \sum\limits_{j=1}^{n} \alpha_{j} = 1 \right)$ such that $\|x - u\| \leqq \varepsilon / 4$. Therefore, by \eqref{eq:AV.4.6}, 
\[
\left| \langle y _{0} ^{\prime}, x _{0} ^{\prime \prime} \rangle - \langle u, y _{n} ^{\prime} \rangle \right| \leqq \sum\limits_{j=1} ^{n} \alpha_{j} \left| \langle y _{0} ^{\prime}, x _{0} ^{\prime \prime} \rangle - \langle x _{j}, y _{n} ^{\prime} \rangle \right| \leqq \varepsilon / 2,
\]
and hence 
\[
\begin{aligned}
&\left| \langle y _{0} ^{\prime}, x _{0} ^{\prime \prime} \rangle \right|  \leqq \left| \langle y _{0} ^{\prime}, x _{0} ^{\prime \prime} \rangle - \langle u, y _{n} ^{\prime} \rangle \right| + \left| \langle u, y _{n} ^{\prime} \rangle - \langle x, y _{n} ^{\prime} \rangle \right| + \left| \langle x, y _{n} ^{\prime} \rangle \right| \leqq \varepsilon / 2 + \| u - x \| \| y _{n} ^{\prime} \| + \varepsilon / 4 \leqq \varepsilon .
\end{aligned}
\]

As $\varepsilon$ was arbitrary, we see that $\langle y_0', x_0'' \rangle = 0$ and so $y_0' \in N(x_0'')$. Combined with the fact that $S_1'$ is weakly* closed, we finally obtain $y_0' \in N(x_0'') \cap S_1'$. 
\end{proof}

\begin{remarku} As for the weak topologies and duality in $B$-spaces, there is an extensive literature. See, e.g., the references in N. \textsc{Dunford}-J. \textsc{Schwartz} \cite{dunford44}. Sections $1$, $2$ and $3$ of this Appendix are adapted and modified from N. \textsc{Bourbaki} \cite{bourbaki26} and A. \textsc{Grothendieck} \cite{grothendieck87}. It is remarkable that necessary tools for proving this far reaching theorem of \textsc{Eberlein} \cite{dunford52}-\textsc{Shmulyan} \cite{shmulyan302} are found, in one form or other, in the book of S. \textsc{Banach} \cite{banach9}. 
\end{remarku}
   % Appendix to Chapter V
% ==== VI ====
\yoschapter{VI}{Fourier Transform and Differential Equations}
\label{chap:VI}

The Fourier transform is one of the most powerful tools in classical and modern analysis. Its scope has recently been strikingly extended thanks to the introduction of the notion of generalized functions of S. L. \textsc{Sobolev} \cite{sobolev306} and L. \textsc{Schwartz} \cite{schwartz294}. The extension has been applied successfully to the theory of linear partial differential equations by L. \textsc{Ehrenpreis}, B. \textsc{Malgrange} and especially by L. \textsc{Hörmander} \cite{hormander118}. 


\yossection{1}{The Fourier Transform of Rapidly Decreasing Functions}
\label{sec:VI.1}

\begin{definition}\label{def:VI.1.Definition.1} We denote by $\mathfrak{S}(\mathbb{R}^{n})$ the totality of functions $f \in C^{\infty}(\mathbb{R}^{n})$ such that 
\begin{equation}\label{eq:VI.1.1}
\sup\limits _ {x \in \mathbb{R}^{n}} | x ^{\beta} D ^{\alpha} f (x) | < \infty \left(x ^{\beta} = \prod\limits_{j=1} ^{n} x _{j} ^ {\beta_{i}}\right)
\end{equation}
for every $\alpha = (\alpha_{1},\alpha_{2},\ldots ,\alpha_{n})$ and $\beta = (\beta_{1},\beta_{2},\dots ,\beta_{n})$ with non-negative integers $\alpha_{j}$ and $\beta_{k}$. Such functions are called rapidly decreasing (at $\infty$). 
\end{definition}

\begin{examplen} $\exp (-|x|^2)$ and functions $f\in C_0^\infty (\mathbb{R}^{n})$ are rapidly decreasing.

\end{examplen}
\begin{proposition}\label{prop:VI.1.Proposition.1} $\mathfrak{S}(\mathbb{R}^{n})$ is a locally convex linear topological space by the algebraic operation of function sum and multiplication of functions by complex numbers, and by the topology defined by the system of semi-norms of the form 

\begin{equation}\label{eq:VI.1.2}
p(f) = \sup\limits_{x\in \mathbb{R}^{n}}|P(x)D^\alpha f(x)|, \text{ where } P(x) \text{ denotes a polynomial}.
\end{equation}
\end{proposition}

\begin{proposition}\label{prop:VI.1.Proposition.2} $\mathfrak{S}(\mathbb{R}^{n})$ is closed with respect to the application of linear partial differential operators with polynomial coefficients. 
\end{proposition}

\begin{proposition}\label{prop:VI.1.Proposition.3} With respect to the topology of $\mathfrak{S}(\mathbb{R}^{n})$, $C_0^\infty (\mathbb{R}^{n})$ is a dense subset of $\mathfrak{S}(\mathbb{R}^{n})$. 
\end{proposition}

\begin{proof}
 Let $f \in \mathfrak{S}(\mathbb{R}^{n})$ and take $\psi \in C_0^\infty(\mathbb{R}^{n})$ such that $\psi(x) = 1$ when $|x| \leqq 1$. Then, for any $\varepsilon > 0$, $f_{\varepsilon}(x) = f(x) \psi(\varepsilon x) \in C_0^\infty(\mathbb{R}^{n})$. By applying Leibniz' rule of differentiation of the product of functions, we see that 
\[
D ^{\alpha} (f _{e} (x) - f (x)) = D ^{\alpha} \{f (x) (\psi (\varepsilon x) - 1) \}
\]
is a finite linear combination of terms of the form 
$D^{\beta}f(x)\cdot (\varepsilon)^{|\gamma |}\{D^{\gamma}\psi (y)\}_{y = \varepsilon x},$ where $|\beta | + |\gamma | = |\alpha |$ with $|\gamma | > 0,$ 
and the term $D^{\alpha}f(x) \cdot (\psi(\varepsilon x) - 1)$. Thus it is easy to see that $f_{\varepsilon}(x)$ tends to $f(x)$ in the topology of $\mathfrak{S}(\mathbb{R}^{n})$ when $\varepsilon \downarrow 0$. 
\end{proof}

\begin{definition}\label{def:VI.1.Definition.2} For any $f \in \mathfrak{S}(\mathbb{R}^{n})$, define its Fourier transform $\hat{f}$ by 
\begin{equation}\label{eq:VI.1.3}
\hat {f} (\xi) = (2 \pi) ^{-n/2} \int_ {\mathbb{R}^{n}} e ^{- i \langle \xi , x \rangle} f (x) \dd x,
\end{equation}
where $\xi = (\xi_1, \xi_2, \ldots, \xi_n)$, $x = (x_1, x_2, \ldots, x_n)$, $\langle \xi, x \rangle = \sum\limits_{j=1}^{n} \xi_j x_j$ and $\dd x = \dd x_1 \dd x_2 \ldots \dd x_n$. We also define the inverse Fourier transform $\tilde{g}$ of $g \in \mathfrak{S}(\mathbb{R}^{n})$ by 
\begin{equation}\label{eq:VI.1.4}
\widetilde {g} (x) = (2 \pi) ^{-n/2} \int e ^{i \langle x, \xi \rangle} g (\xi) \dd \xi .
\end{equation}
\end{definition}

\begin{proposition}\label{prop:VI.1.Proposition.4} The Fourier transform: $f \to \hat{f}$ maps $\mathfrak{S}(\mathbb{R}^{n})$ linearly and continuously into $\mathfrak{S}(\mathbb{R}^{n})$. The inverse Fourier transform: $g \to \tilde{g}$ also maps $\mathfrak{S}(\mathbb{R}^{n})$ linearly and continuously into $\mathfrak{S}(\mathbb{R}^{n})$. 
\end{proposition}

\begin{proof}
 Differentiating formally under the integral sign, we have 
\begin{equation}\label{eq:VI.1.5}
D ^{\alpha} \hat {f} (\xi) = (2 \pi) ^{-n/2} \int e ^{- i \langle \xi , x \rangle} (- i) ^{| \alpha |} x ^{\alpha} f (x) \dd x. \end{equation}

The formal differentiation is permitted since the right hand side is, by \eqref{eq:VI.1.1}, uniformly convergent in $\xi$. Thus $\hat{f} \in C^{\infty}(\mathbb{R}^{n})$. Similarly, we have, by integration by parts, 
\begin{equation}\label{eq:VI.1.6}
(i) ^{| \beta |} \xi^{\beta} \hat {f} (\xi) = (2 \pi) ^{-n/2} \int e ^{- i \langle \xi , x \rangle} D ^{\beta} f (x) \dd x.
\end{equation}

Thus we have 
\begin{equation}\label{eq:VI.1.7}
(i) ^{| \beta | + | \alpha |} \xi^{\beta} D ^{\alpha} \hat {f} (\xi) = (2 \pi) ^{-n/2} \int e ^{- i \langle \xi , x \rangle} D ^{\beta} (x ^{\alpha} f (x)) \dd x,
\end{equation}
and \eqref{eq:VI.1.7} proves that the mapping $f \to \hat{f}$ is continuous in the topology of $\mathfrak{S}(\mathbb{R}^{n})$. 
\end{proof}

\begin{theorem}[Fourier's integral theorem]\label{thm:VI.1.Theorem.1} Fourier's inversion theorem holds: 
\begin{align}
\tilde {f} (x)& = (2 \pi) ^{-n/2} \int e ^{i \langle x, \xi \rangle} \hat {f} (\xi) \dd \xi = f (x), \text{ i.e., we have } \label{eq:VI.1.8} \\
\tilde {\hat {f}}& = f, \text{ and similarly } \hat {\tilde {f}} = f. \tag{8'}\label{eq:VI.1.8'}
\end{align}
\end{theorem}

Therefore it is easy to see that the Fourier transform maps $\mathfrak{S}(\mathbb{R}^{n})$ onto $\mathfrak{S}(\mathbb{R}^{n})$ linearly and continuously in both directions, and the inverse Fourier transform gives the inverse mapping of the Fourier transform. 

\begin{proof}
 We have 
\begin{equation}\label{eq:VI.1.9}
\int g (\xi) \hat {f} (\xi) e ^{i \langle x, \xi \rangle} \dd \xi = \int \hat {g} (y) f (x + y) \dd y (f \text{ and } g \in \mathfrak {S} (\mathbb{R}^{n})). \end{equation}

In fact, the left hand side is equal to 
\[
\begin{aligned}
&\int g (\xi) \left\{(2 \pi) ^{-n/2} \int e ^{- i \langle \xi , y \rangle} f (y) \dd y \right\} e ^{i \langle x, \xi \rangle} \dd \xi = (2 \pi) ^{-n/2} \int \left\{\int g (\xi) e ^{- i \langle \xi , y - x \rangle} \dd \xi \right\} f (y) \dd y
\\
&= \int \hat {g} (y - x) f (y) \dd y = \int \hat {g} (y) f (x + y) \dd y.
\end{aligned}
\]

If we take $g(\varepsilon \xi)$ for $g(\xi), \varepsilon > 0$, then 
\[
(2 \pi) ^{-n/2} \int e ^{- i \langle y, \xi \rangle} g (\varepsilon \xi) \dd \xi = (2 \pi) ^{-n/2} \varepsilon^{-n} \int g (z) e ^{- i \langle y, z / \varepsilon \rangle} \dd z = \varepsilon^{-n} \hat {g} (y / \varepsilon).
\]

Hence, by \eqref{eq:VI.1.9}, 
\[
\int g (\varepsilon \xi) \hat {f} (\xi) e ^{i \langle x, \xi \rangle} \dd \xi = \int \hat {g} (y) f (x + \varepsilon y) \dd y.
\]

We shall take, following F. \textsc{Riesz}, $g(x) = e^{-|x|^2/2}$ and let $\varepsilon \downarrow 0$. Then 
\[
g (0) \int \hat {f} (\xi) e ^{i \langle x, \xi \rangle} \dd \xi = f (x) \int \hat {g} (y) \dd y.
\]

This proves \eqref{eq:VI.1.8}, since $g(0) = 1$ and $\int \hat{g}(y) \dd y = (2\pi)^{n/2}$ by the well-known facts: 
\begin{align}
(2 \pi) ^{-n/2}& \int e ^ {- | x | ^{2} / 2} e ^{- i \langle y, x \rangle} \dd x = e ^ {- | y | ^{2} / 2}, \label{eq:VI.1.10} \\
(2 \pi) ^{-n/2}& \int e ^ {- | x | ^{2} / 2} \dd x = 1. \tag{10'}\label{eq:VI.1.10'}
\end{align}
\end{proof}


\begin{remarku} For the sake of completeness, we will give the proof of \eqref{eq:VI.1.10}. We have 
\[
(2 \pi) ^{-1/2} \int_{- \lambda} ^{\lambda} e ^ {- t ^{2} / 2} e ^{-iut} \dd t = e ^ {- u ^{2} / 2} (2 \pi) ^{-1/2} \int_{- \lambda} ^{\lambda} e ^ {- (t + i u) ^{2} / 2} \dd t.
\]
\end{remarku}

Let $u > 0$, and integrate the function $e^{-z^2/2}$, which is holomorphic in $z = t + iu$, along the curve consisting of the oriented segments 
\[
\overrightarrow {- \lambda , \lambda}, \overrightarrow {\lambda , \lambda + i u}, \overrightarrow {\lambda + i u , - \lambda + i u} \text{ and } \overrightarrow {- \lambda + i u , - \lambda}
\]
in this order. By Cauchy's integral theorem, the integral vanishes. Thus 
\[
\begin{aligned}
&(2 \pi) ^{-1/2} \int_{- \lambda} ^{\lambda} e ^ {- (t + i u) ^{2} / 2} \dd t = (2 \pi) ^{-1/2} \int_{- \lambda} ^{\lambda} e ^ {- t ^{2} / 2} \dd t + (2 \pi) ^{-1/2} \int_{u} ^{0} e ^ {- (- \lambda + i u) ^{2} / 2} i \dd u
\\
&+ (2 \pi) ^{-1/2} \int_{0} ^{u} e ^ {- (\lambda + i u) ^{2} / 2} i \dd u.
\end{aligned}
\]

The second and third terms on the right tend to $0$ as $\lambda \to \infty$, and so, by \eqref{eq:VI.1.10}, 
\[
(2 \pi) ^{-1/2} \int_{- \infty} ^{\infty} e ^ {- t ^{2} / 2} e ^{-itu} \dd t = e ^ {- u ^{2} / 2} (2 \pi) ^{-1/2} \int_{- \infty} ^{\infty} e ^ {- t ^{2} / 2} \dd t = e ^ {- u ^{2} / 2}.
\]

We have thus proved \eqref{eq:VI.1.10} for the case $n$ = $1$, and it is easy to prove the case of a general $n$ by reducing it to the case $n$ = $1$. 

\begin{corollaryn}[Parseval's relation] We have
\begin{align}
&\int \hat {f} (\xi) g (\xi) \dd \xi = \int f (x) \hat {g} (x) \dd x, \label{eq:VI.1.11} \\
&\int f (\xi) \bar {g} (\xi) \dd \xi = \int \tilde {f} (x) \bar {\tilde {g}} (x) \dd x, \label{eq:VI.1.12} \\
(\widehat {f \star g})& = (2 \pi) ^{n/2} \widehat {f} \cdot \widehat {g} \text{ and } (2 \pi) ^{n/2} (\widehat {f \cdot g}) = \widehat {f} \star \widehat {g}, \label{eq:VI.1.13}
\end{align}
where the convolution $f \star g$ is defined through 
\begin{equation}\label{eq:VI.1.14}
(f \star g) (x) = \int f (x - y) g (y) \dd y = \int g (x - y) f (y) \dd y.
\end{equation}

\end{corollaryn}
\begin{proof}
 \eqref{eq:VI.1.11} is obtained from \eqref{eq:VI.1.9} by putting $x$ = $0$. \eqref{eq:VI.1.12} is obtained from \eqref{eq:VI.1.11} by observing that the Fourier transform of $\bar{g}$ is equal to $\tilde{\bar{g}}$. We next show that 
\begin{equation}\label{eq:VI.1.15}
\begin{gathered}
(2 \pi) ^{-n/2} \int (f * g) (x) e ^{- i \langle \xi , x \rangle} \dd x = (2 \pi) ^{-n/2} \int g (y) e ^{- i \langle \xi , y \rangle} \left\{\int f (x - y) e ^{- i \langle \xi , x - y \rangle} \dd x \right\} \dd y
\\
= (2 \pi) ^{n/2} \hat {f} (\xi) \hat {g} (\xi).
\end{gathered}
\end{equation}

Since the product $\hat{f} \cdot \hat{g}$ of functions $\hat{f}$ and $\hat{g} \in \mathfrak{S}(\mathbb{R}^{n})$ is again a function of $\mathfrak{S}(\mathbb{R}^{n})$, we know that the right hand side of \eqref{eq:VI.1.15} belongs to $\mathfrak{S}(\mathbb{R}^{n})$. 

It is easy to see that, the convolution $f \star g$ of two functions of $\mathfrak{S}(\mathbb{R}^{n})$ belongs also to $\mathfrak{S}(\mathbb{R}^{n})$. Thus we have proved the first formula of \eqref{eq:VI.1.13}. The second formula may be proved similarly to \eqref{eq:VI.1.9} by \eqref{eq:VI.1.15}. 
\end{proof}
\begin{theorem}[Poisson's Summation Formula]\label{thm:VI.1.Theorem.2} Let $\varphi \in \mathfrak{S}(\mathbb{R}^{1})$ and $\hat{\varphi} \in \mathfrak{S}(\mathbb{R}^{1})$ its Fourier transform. Then we have 
\begin{equation}\label{eq:VI.1.16}
\sum\limits_{n = - \infty} ^{\infty} \varphi (2 \pi n) = \sum\limits_{n = - \infty} ^{\infty} \hat {\varphi} (n).
\end{equation}
\end{theorem}

\begin{proof}
 Set $f(x) = \sum\limits_{n=-\infty}^{\infty} \varphi(x + 2\pi n)$. This series is absolutely convergent, $\in C^\infty$ and $f(x + 2\pi) = f(x)$, as may be proved by the fact that $\varphi(x)$ is rapidly decreasing at $\infty$. In particular, both sides of \eqref{eq:VI.1.16} are convergent. We have to prove the equality. 


The Fourier coefficients $c_k$ of $f(x)$ with respect to the complete orthonormal system $\{(2\pi)^{-1/2} e^{-ikx}; k = 0, \pm 1, \pm 2, \ldots\}$ of $L^2(0, 2\pi)$ are given by 
\[
\begin{aligned}
c _{k} &= (2 \pi) ^{-1/2} \int_{0} ^{2 \pi} f (x) e ^{-ikx} \dd x = \sum\limits_{n = - \infty} ^{\infty} (2 \pi) ^{-1/2} \int_{0} ^{2 \pi} \varphi (x + 2 \pi n) e ^{-ikx} \dd x
\\
&= \sum\limits_{n = - \infty} ^{\infty} (2 \pi) ^{-1/2} \int_{2 \pi n} ^{2 \pi (n +1)} \varphi (x) e ^{-ikx} \dd x = \hat {\varphi} (k).
\end{aligned}
\]

Thus, by $f \in L^2(0, 2\pi)$, we have 
\[
f (x) = \sum\limits_{n = - \infty} ^{\infty} \varphi (x + 2 \pi n) = \text{ l.i.m. } \sum\limits_{k=-s} ^{s} \hat {\varphi} (k) e ^{ikx}.
\]

However, since $\hat{\varphi}(x) \in \mathfrak{S}(\mathbb{R}^{n})$, the series $\sum\limits_{k=-\infty}^{\infty} \hat{\varphi}(k) e^{ikx}$ converges absolutely. Hence 
\[
\sum\limits_{n = - \infty} ^{\infty} \varphi (x + 2 \pi n) \equiv \sum\limits_{k = - \infty} ^{\infty} \hat {\varphi} (k) e ^{ikx},
\]
and so we obtain \eqref{eq:VI.1.16} by setting $x = 0$.
\end{proof}

\begin{examplen} We have, by \eqref{eq:VI.1.10},
\[
\begin{aligned} (2 \pi) ^{-1/2} \int_{- \infty} ^{\infty} e ^ {- t x ^{2}} e ^{-ixy} \dd x & = (2 \pi) ^{-1/2} \int_{- \infty} ^{\infty} e ^ {- x ^{2} / 2} e ^ {- i x y / \sqrt {2 t}} (2 t) ^{-1/2} \dd x \\ & = (2 t) ^{-1/2} e ^ {- y ^{2} / 4 t}, t > 0. \end{aligned}
\]

Hence, by \eqref{eq:VI.1.16}, we obtain the so-called $\theta$-formula: 
\begin{equation}\label{eq:VI.1.17}
\sum\limits_{n = - \infty} ^{\infty} e ^ {- 4 t \pi^{2} n ^{2}} = \sum\limits_{n = - \infty} ^{\infty} (2 t) ^{-1/2} e ^ {- n ^{2} / 4 t}, t > 0.
\end{equation}


\end{examplen}
\yossection{2}{The Fourier Transform of Tempered Distributions}
\label{sec:VI.2}

\begin{definition}\label{def:VI.2.Definition.1} $A$ linear functional $T$ defined and continuous on $\mathfrak{S}(\mathbb{R}^{n})$ is called a tempered distribution (in $\mathbb{R}^{n}$). The totality of tempered distributions is denoted by $\mathfrak{S}(\mathbb{R}^{n})'$. As a dual space of $\mathfrak{S}(\mathbb{R}^{n})$, $\mathfrak{S}(\mathbb{R}^{n})'$ is a locally convex linear topological space by the strong dual topology. 
\end{definition}

\begin{proposition}\label{prop:VI.2.Proposition.1} Since $C_0^\infty(\mathbb{R}^{n})$ is contained in $\mathfrak{S}(\mathbb{R}^{n})$ as an abstract set, and since the topology in $\mathfrak{D}(\mathbb{R}^{n})$ is stronger than the topology in $\mathfrak{S}(\mathbb{R}^{n})$, the restriction of a tempered distribution to $C_0^\infty(\mathbb{R}^{n})$ is a distribution in $\mathbb{R}^{n}$. Two different tempered distributions define, when restricted to $C_0^\infty(\mathbb{R}^{n})$, two different distributions in $\mathbb{R}^{n}$, because $C_0^\infty(\mathbb{R}^{n})$ is dense in $\mathfrak{S}(\mathbb{R}^{n})$ with respect to the topology of $\mathfrak{S}(\mathbb{R}^{n})$, and hence a distribution $\in \mathfrak{S}(\mathbb{R}^{n})'$ which vanishes on $C_0^\infty(\mathbb{R}^{n})$ must vanish on $\mathfrak{S}(\mathbb{R}^{n})$. Therefore 
\begin{equation}\label{eq:VI.2.1}
\mathfrak {S} (\mathbb{R}^{n}) ^{\prime} \subseteq \mathfrak {D} (\mathbb{R}^{n}) ^{\prime}.
\end{equation}
\end{proposition}

\begin{example}\label{ex:VI.2.Example.1} $A$ distribution in $\mathbb{R}^{n}$ with a compact support surely belongs to $\mathfrak{S}(\mathbb{R}^{n})'$. Therefore 
\begin{equation}\label{eq:VI.2.2}
\mathfrak {E} (\mathbb{R}^{n}) ^{\prime} \subseteq \mathfrak {S} (\mathbb{R}^{n}) ^{\prime}.
\end{equation}
\end{example}

\begin{example}\label{ex:VI.2.Example.2} $A$ $\sigma$-finite, non-negative measure $\mu(\dd x)$ which is $\sigma$-additive on Baire sets of $\mathbb{R}^{n}$ is called a slowly increasing measure, if, for some nonnegative $k$, 
\begin{equation}\label{eq:VI.2.3}
\int_ {\mathbb{R}^{n}} (1 + | x | ^{2}) ^{-k} \mu (\dd x) < \infty .
\end{equation}
\end{example}

Such a measure $\mu$ defines a tempered distribution by 
\begin{equation}\label{eq:VI.2.4}
T _{\mu} (\varphi) = \int_ {\mathbb{R}^{n}} \varphi (x) \mu (\dd x), \varphi \in \mathfrak {S} (\mathbb{R}^{n}).
\end{equation}

For, by the condition $\varphi \in \mathfrak{S}(\mathbb{R}^{n})$, we have $\varphi(x) = 0((1 + |x|^2)^{-k})$ for large $|x|$. 

\begin{example}\label{ex:VI.2.Example.3} As a special case of Example~\hyperref[ex:VI.2.Example.2]{$2$}, any function $f \in L^{p}(\mathbb{R}^{n})$, $p \geqq 1$, defines a tempered distribution 
\begin{equation}\tag{4'}\label{eq:VI.2.4'}
T _{f} (\varphi) = \int_ {\mathbb{R}^{n}} \varphi (x) f (x) \dd x, \varphi \in \mathfrak {S} (\mathbb{R}^{n}).
\end{equation}
\end{example}

That an $f \in L^{p}(\mathbb{R}^{n})$ gives rise to a slowly increasing measure $\mu(\dd x) = |f(x)| \dd x$ may be proved by applying Hölder's inequality to 
\[
\int_ {\mathbb{R}^{n}} (1 + | x | ^{2}) ^{-k} | f (x) | \dd x.
\]

\begin{definition}\label{def:VI.2.Definition.2} $A$ function $f \in C^{\infty}(\mathbb{R}^{n})$ is called slowly increasing (at $\infty$), if, for any differentiation $D^{j}$, there exists a non-negative integer $N$ such that 
\begin{equation}\label{eq:VI.2.5}
\lim\limits _{| x | \rightarrow \infty} | x | ^{-N} | D ^{j} f (x) | = 0.
\end{equation}
\end{definition}

The totality of slowly increasing functions will be denoted by $\mathfrak{D}_{M}(\mathbb{R}^{n})$. It is a locally convex linear topological space by the algebraic operations of function sum and multiplication of functions by complex numbers, and by the topology defined by the system of semi-norms of the form 
\begin{equation}\label{eq:VI.2.6}
p (f) = p _ {h, D ^{j}} (f) = \sup\limits _ {x \in \mathbb{R}^{n}} | h (x) D ^{j} f (x) |, f \in \mathfrak {D} _{M} (\mathbb{R}^{n}),
\end{equation}
where $h(x)$ is an arbitrary function $\in \mathfrak{S}(\mathbb{R}^{n})$ and $D^j$ an arbitrary differentiation. As may be seen by applying Leibniz' formula of differentiation of the product of functions, $h(x) D^j f(x) \in \mathfrak{S}(\mathbb{R}^{n})$ and so $p_{h,Dj}(f)$ is finite for every $f \in \mathfrak{D}_M(\mathbb{R}^{n})$. Moreover, if $p_{h,Dj}(f) = 0$ for all $h \in \mathfrak{S}(\mathbb{R}^{n})$ and $D^j$, then $f(x) \equiv 0$ as may be seen by taking $D = I$ and $h \in \mathfrak{D}(\mathbb{R}^{n})$. 

\begin{proposition}\label{prop:VI.2.Proposition.2} $C_0^\infty (\mathbb{R}^{n})$ is dense in $\mathfrak{D}_M(\mathbb{R}^{n})$ with respect to the topology of $\mathfrak{D}_M(\mathbb{R}^{n})$. 
\end{proposition}

\begin{proof}
 Let $f \in \mathfrak{D}_M(\mathbb{R}^{n})$, and take $\psi \in C_0^\infty(\mathbb{R}^{n})$ such that $\psi(x) = 1$ for $|x| \leqq 1$. Then $f_{\varepsilon}(x) = f(x) \psi(\varepsilon x) \in C_0^\infty(\mathbb{R}^{n})$ for any $\varepsilon > 0$. As in Proposition~\ref{prop:VI.1.Proposition.3} in Chapter~\ref{chap:VI}, \ref{sec:VI.1}, we easily prove that $f_{\varepsilon}(x)$ tends to $f(x)$ in the topology of $\mathfrak{D}_M(\mathbb{R}^{n})$ when $\varepsilon \downarrow 0$. 
\end{proof}

\begin{proposition}\label{prop:VI.2.Proposition.3} Any function $f \in \mathfrak{D}_M(\mathbb{R}^{n})$ defines a tempered distribution 
\begin{equation}\label{eq:VI.2.7}
T _{f} (\varphi) = \int_ {\mathbb {R} ^{n}} f (x) \varphi (x) \dd x, \varphi \in \mathfrak {S} (\mathbb{R}^{n}).
\end{equation}
\end{proposition}

\begin{definition}\label{def:VI.2.Definition.3} As in the case of a distribution in $\mathbb{R}^{n}$, we can define the generalized derivative of a tempered distribution $T$ by 
\begin{equation}\label{eq:VI.2.8}
D ^{j} T (\varphi) = (- 1) ^{|j|} T (D ^{j} \varphi), \varphi \in \mathfrak {S} (\mathbb{R}^{n}),
\end{equation}
since the mapping $\varphi(x) \to D^{j}\varphi(x)$ of $\mathfrak{S}(\mathbb{R}^{n})$ into $\mathfrak{S}(\mathbb{R}^{n})$ is linear and continuous in the topology of $\mathfrak{S}(\mathbb{R}^{n})$. We can also define a multiplication by a function $f \in \mathfrak{D}_M(\mathbb{R}^{n})$ to a distribution $T \in \mathfrak{S}(\mathbb{R}^{n})'$ through 
\begin{equation}\label{eq:VI.2.9}
(f T) (\varphi) = T (f \varphi), \varphi \in \mathfrak {S} (R) ^{n},
\end{equation}
since the mapping $\varphi(x) \to f(x)\varphi(x)$ of $\mathfrak{S}(\mathbb{R}^{n})$ into $\mathfrak{S}(\mathbb{R}^{n})$ is linear and continuous in the topology of $\mathfrak{S}(\mathbb{R}^{n})$. 
\end{definition}

\yossub{The Fourier Transform of Tempered Distributions}

\begin{definition}\label{def:VI.2.Definition.4} Since the mapping $\varphi(x) \to \hat{\varphi}(x)$ of $\mathfrak{S}(\mathbb{R}^{n})$ onto $\mathfrak{S}(\mathbb{R}^{n})$ is linear and continuous in the topology of $\mathfrak{S}(\mathbb{R}^{n})$, we can define the Fourier transform $\hat{T}$ of a tempered distribution $T$ as the tempered distribution $\hat{T}$ defined through 
\begin{equation}\label{eq:VI.2.10}
\hat {T} (\varphi) = T (\hat {\varphi}), \varphi \in \mathfrak {S} (\mathbb{R}^{n}).
\end{equation}
\end{definition}

\begin{example}\label{ex:VI.2.Example.1a} If $f \in L^{1}(\mathbb{R}^{n})$, then 
\begin{equation}\label{eq:VI.2.11}
\hat {T} _{f} = T _ {\hat {f}}, \text{ where } \hat {f} (x) = (2 \pi) ^{-n/2} \int_ {\mathbb{R}^{n}} e ^{- i \langle x, \xi \rangle} f (\xi) \dd \xi ,
\end{equation}
as may be seen by changing the order of integration in $\hat{T}_f(\varphi) = \int f(x) \cdot \hat{\varphi}(x) \dd x = (2\pi)^{-n/2} \int_{\mathbb{R}^{n}} f(x) \left\{ \int_{\mathbb{R}^{n}} e^{-i\langle x, \xi \rangle} \varphi(\xi) \dd\xi \right\} \dd x.$ 
\end{example}

\begin{remarku} In the above sense, the Fourier transform of a tempered distribution is a generalization of the ordinary Fourier transform of functions. 
\end{remarku}

\begin{proposition}\label{prop:VI.2.Proposition.4} If we define 
\begin{equation}\label{eq:VI.2.12}
\check {f} (x) = f (- x),
\end{equation}
then Fourier's integral theorem in the preceding Section~\hyperref[sec:VI.1]{$1$} is expressed by 
\begin{equation}\label{eq:VI.2.13}
\hat {\tilde {f}} = \check {f}, f \in \mathfrak {S} (\mathbb{R}^{n}).
\end{equation}
\end{proposition}

\begin{corollary}[Fourier's integral theorem]\label{cor:VI.2.Corollary.1} Fourier's integral theorem is generalized to tempered distributions as follows: 
\begin{equation}\label{eq:VI.2.14}
\hat {\tilde {T}} = \check {T}, \text{ where } \check {T} (\varphi) = T (\check {\varphi}).
\end{equation}
\end{corollary}

In particular, the Fourier transform $T \to \hat{T}$ maps $\mathfrak{S}(\mathbb{R}^{n})'$ linearly onto $\mathfrak{S}(\mathbb{R}^{n})'$. 

\begin{proof}
 We have, by definition, 
\[
\hat {\tilde {T}} (\varphi) = T (\hat {\varphi}) = T (\check {\varphi}) = \check {T} (\varphi) \text{ for all } \varphi \in \mathfrak {S} (\mathbb{R}^{n}). \]
\end{proof}

\begin{corollary}\label{cor:VI.2.Corollary.2} The Fourier transform $T\to\tilde{T}$ and its inverse are linear and continuous on $\mathfrak{S}(\mathbb{R}^{n})'$ onto $\mathfrak{S}(\mathbb{R}^{n})'$ with respect to the weak* topology of $\mathfrak{S}(\mathbb{R}^{n})'$: 
\begin{equation}\lim T _{h} (\varphi) = T (\varphi) \text{ for all } \varphi \in \mathfrak {S} (\mathbb{R}^{n}) \text{ implies that } \lim \hat {T} _{h} (\varphi) = \hat {T} (\varphi) \text{ for all } \varphi \in \mathfrak {S} (\mathbb{R}^{n}). \label{eq:VI.2.15}\end{equation}
\end{corollary}

Here the inverse of the mapping $T \rightarrow \tilde{T}$ is defined by the inverse Fourier transform $T \rightarrow \tilde{T}$ given by 
\begin{equation}\tag{10'}\label{eq:VI.2.10'}
\tilde {T} (\varphi) = T (\tilde {\varphi}), \varphi \in \mathfrak {S} (\mathbb{R}^{n}).
\end{equation}

\begin{example}\label{ex:VI.2.Example.2a} 
\begin{equation}\hat {T} _{\delta} = (2 \pi) ^{-n/2} T _{1}, \hat {T} _{1} = (2 \pi) ^{n/2} T _{\delta}.\label{eq:VI.2.16}
\end{equation}
\end{example}

\begin{proof}
\[
 \hat {T} _{\delta} (\varphi) = T _{\delta} (\hat {\varphi}) = \hat {\varphi} (0) = (2 \pi) ^{-n/2} \int_ {\mathbb{R}^{n}} 1 \cdot \varphi (y) \dd y = (2\pi)^{-n/2}T_{1}(\varphi).
\]
and $T_{\delta} = \check{T}_{\delta} = \hat{\tilde{T}}_{\delta} = (2\pi)^{-n/2}\hat{T}_1.$ 
\end{proof}

\begin{example}\label{ex:VI.2.Example.3a} 
\begin{align}
(\widehat {\partial T / \partial x _{j}})& = i x _{j} \hat {T}, \label{eq:VI.2.17} \\
(\widehat {i x _{j} T})& = - (\partial \hat {T} / \partial x _{j}). \label{eq:VI.2.18}
\end{align}
\end{example}

\begin{proof}
 We have, by \eqref{eq:VI.1.5} in Chapter~\ref{chap:VI}, \ref{sec:VI.1}, 
\[
\begin{gathered}
(\widehat {\partial T / \partial x _{j}}) (\varphi) = (\partial T / \partial x _{j}) (\widehat {\varphi}) = - T (\partial \widehat {\varphi} / \partial x _{j}) = - T (- i x _{j} \varphi (x)) = T (\widehat {i x _{j} \varphi}) = (i x _{j} \widehat {T}) (\varphi).
\end{gathered}
\]

We also have, by \eqref{eq:VI.1.6} in Chapter~\ref{chap:VI}, \ref{sec:VI.1}, 
\[
\begin{aligned}
&\widehat {(i x _{j} T)} (\varphi)  = (i x _{j} T) (\hat {\varphi}) = T (i x _{j} \hat {\varphi}) = T (\widehat {\partial \varphi / \partial x _{j}}) = \hat {T} (\partial \varphi / \partial x _{j}) = - (\partial \hat {T} / \partial x _{j}) (\varphi).
\end{aligned}
\]
\end{proof}

\begin{theoremn}[Plancherel's Theorem] If $f \in L^2(\mathbb{R}^{n})$, then the Fourier transform $\hat{T}_f$ of $T_f$ is defined by a function $\hat{f} \in L^2(\mathbb{R}^{n})$, i.e., 
\begin{equation}\label{eq:VI.2.19}
\hat {T} _{f} = T _ {\hat {f}} \text{ with } \hat {f} \in L ^{2} (\mathbb{R}^{n}),
\end{equation}

and 
\begin{equation}\label{eq:VI.2.20}
\left\| \hat {f} \right\| = \left(\int_ {\mathbb{R}^{n}} | \hat {f} (x) | ^{2} \dd x\right) ^{1/2} = \left(\int_ {\mathbb{R}^{n}} | f (x) | ^{2} \dd x\right) ^{1/2} = \left\| f \right\|.
\end{equation}
\end{theoremn}

\begin{proof}
 We have, by Schwarz' inequality, 
\begin{equation}\label{eq:VI.2.21}
\left| \hat {T} _{f} (\varphi) \right| = \left| T _{f} (\hat {\varphi}) \right| = \left| \int_ {\mathbb{R}^{n}} f (x) \hat {\varphi} (x) \dd x \right| \leqq \| f \| \| \hat {\varphi} \| = \| f \| \| \varphi \|. \end{equation}

The equality above $\|\hat{\varphi}\| = \|\varphi\|$ is proved in \eqref{eq:VI.2.12} of the preceding section $1$. Hence, by F. \textsc{Riesz}' representation theorem in the Hilbert space $L^2(\mathbb{R}^{n})$, there exists a uniquely determined $\hat{f} \in L^2(\mathbb{R}^{n})$ such that 
\begin{equation}\label{eq:VI.2.22}
\begin{aligned}
\hat {T} _{f} (\varphi) &= \int_ {\mathbb{R}^{n}} \varphi (x) \hat {f} (x) \dd x = T _ {\hat {f}} (\varphi), \text{ that is, }
\\
\int_ {\mathbb{R}^{n}} \hat {f} (x) \varphi (x) \dd x &= \int_ {\mathbb{R}^{n}} f (x) \hat {\varphi} (x) \dd x \text{ for all } \varphi \in \mathfrak {S} (\mathbb{R}^{n}).
\end{aligned}
\end{equation}

Moreover, we have, from \eqref{eq:VI.2.21}, $\|f\| \geqq \|\hat{f}\|$ since $\mathfrak{S}(\mathbb{R}^{n})$ is dense in $L^2(\mathbb{R}^{n})$ in the topology of $L^2(\mathbb{R}^{n})$. We have thus $\|\hat{\hat{f}}\| \leqq \|\hat{f}\| \leqq \|f\|$. On the other hand, we have, by \eqref{eq:VI.2.13} and \eqref{eq:VI.2.22}, 
\[
\int_{\mathbb{R}^{n}} \hat {\hat {f}} (x) \varphi (x) \dd x = \int_ {\mathbb{R}^{n}} \hat {f} (x) \hat {\varphi} (x) \dd x = \int_ {\mathbb{R}^{n}} f (- x) \varphi (x) \dd x \text{ for all } \varphi \in \mathfrak {S} (\mathbb{R}^{n})
\]
that is, 
\begin{equation}\label{eq:VI.2.23}
\hat {\hat {f}} (x) = f (- x) = \check {f} (x) \text{ a.e. }
\end{equation}

Thus $\|\hat{\hat{f}}\| = \|f\|$ and so, combined with $\|\hat{\hat{f}}\| \leqq \|\hat{f}\| \leqq \|f\|$, we obtain \eqref{eq:VI.2.20}. 
\end{proof}

\begin{definition}\label{def:VI.2.Definition.5} The above obtained $\hat{f}(x) \in L^{2}(\mathbb{R}^{n})$ is called the Fourier transform of the function $f(x) \in L^{2}(\mathbb{R}^{n})$. 
\end{definition}


\begin{corollary}\label{cor:VI.2.Corollary.1a} We have, for any $f \in L^2(\mathbb{R}^{n})$, 
\begin{equation}\label{eq:VI.2.24}
\hat {f} (x) = \underset {h \uparrow \infty} {\text{ l.i.m. }} (2 \pi) ^{-n/2} \int_{| x | \leqq h} e ^{- i \langle x, y \rangle} f (y) \dd y.
\end{equation}
\end{corollary}

\begin{proof}
 Put 
\[
f _{h} (x) = f (x) \text{ or } = 0 \text{ according as } | x | \leqq h \text{ or } | x | > h. \]

Then $\lim\limits_{h\to \infty}\|f_h - f\| = 0$ and so, by \eqref{eq:VI.2.20}, $\lim\limits_{h\to \infty}\|\hat{f}_h - \hat{f} \| = 0$, that is, $\hat{f}(x) = \underset {h\to \infty}{\text{ l.i.m. }}\hat{f}_h(x)$ a.e. But, by \eqref{eq:VI.2.22}, 
\[
\begin{aligned}
&\int_ {\mathbb{R}^{n}} \hat {f} _{h} (x) \varphi (x) \dd x  = \int_ {\mathbb{R}^{n}} f _{h} (x) \hat {\varphi} (x) \dd x = \int_{| x | \leqq h} f (x) \left\{(2 \pi) ^{-n/2} \int_ {\mathbb{R}^{n}} e ^{- i \langle x, y \rangle} \varphi (y) \dd y \right\} \dd x,
\end{aligned}
\]
which is, by changing the order of integration, equal to 
\[
\int_ {\mathbb{R}^{n}} (2 \pi) ^{-n/2} \left\{\int_{| x | \leqq h} e ^{- i \langle x, y \rangle} f (x) \dd x \right\} \varphi (y) \dd y,
\]
since $f_{h}(x)$ is integrable over $|x| \leqq h$ as may be seen by Schwarz' inequality. Thus $\hat{f}_{h}(x) = (2\pi)^{-n/2} \int_{|x| \leqq h} e^{-i\langle x, y \rangle} f(y) \, \dd y$ a.e., and so we obtain \eqref{eq:VI.2.24}. 
\end{proof}

\begin{corollary}\label{cor:VI.2.Corollary.2a} The Fourier transform $f \to \hat{f}$ maps $L^2(\mathbb{R}^{n})$ onto $L^2(\mathbb{R}^{n})$ in a one-one-manner such that 
\begin{equation}\label{eq:VI.2.25}
(f, g) = (\hat {f}, \hat {g}) \text{ for all } f, g \in L ^{2} (\mathbb{R}^{n}).
\end{equation}
\end{corollary}

\begin{proof}
 Like the Fourier transform $f \to \hat{f}$, the inverse Fourier transform $f \to \tilde{f}$ defined by 
\begin{equation}\label{eq:VI.2.26}
\tilde {f} (x) = \underset {h \rightarrow \infty} {\text{ l . i . m . }} (2 \pi) ^{-n/2} \int_{| y | \leqq h} e ^{i \langle x, y \rangle} f (y) \dd y
\end{equation}
maps $L^2(\mathbb{R}^{n})$ into $L^2(\mathbb{R}^{n})$ in such a way that $\|\tilde{f}\| = \|f\|$. Hence we see that the Fourier transform $f \to \hat{f}$ maps $L^2(\mathbb{R}^{n})$ onto $L^2(\mathbb{R}^{n})$ in a one-one way such that $\|\hat{f}\| = \|f\|$. Hence, by the linearity of the Fourier transform and 
\[
(x, y) = 4 ^{-1} \left(\left\| x + y \right\| ^{2} - \left\| x - y \right\| ^{2}\right) + 4 ^{-1} i \left(\left\| x + i y \right\| ^{2} - \left\| x - i y \right\| ^{2}\right),
\]
we obtain \eqref{eq:VI.2.25}. 
\end{proof}

Parseval's Theorem for the Fourier Transform. Let $f_{1}(x)$ and $f_{2}(x)$ both belong to $L^2(\mathbb{R}^{n})$, and let their Fourier transforms be $\hat{f}_1(u)$ and $\hat{f}_2(u)$, respectively. Then 
\begin{equation}\label{eq:VI.2.27}
\int_ {\mathbb{R}^{n}} \hat {f} _{1} (u) \hat {f} _{2} (u) \dd u = \int_ {\mathbb{R}^{n}} f _{1} (x) f _{2} (- x) \dd x,
\end{equation}
and so 
\begin{equation}\label{eq:VI.2.28}
\int_ {\mathbb{R}^{n}} \hat {f} _{1} (u) \hat {f} _{2} (u) e ^{i \langle u, x \rangle} \dd u = \int_ {\mathbb{R}^{n}} f _{1} (y) f _{2} (x - y) \dd y.
\end{equation}

Thus, if $\hat{f}_1(u)\hat{f}_2(u)$ as well as both its factors belong to $L^2 (\mathbb{R}^{n})$, it is the Fourier transform of 
\begin{equation}\label{eq:VI.2.29}
(2 \pi) ^{-n/2} \int_ {\mathbb{R}^{n}} f _{1} (y) f _{2} (x - y) \dd y.
\end{equation}

This will also be true if $f_{1}(x), f_{2}(x)$ and \eqref{eq:VI.2.29} all belong to $L^2(\mathbb{R}^{n})$. 

\begin{proof}
 It is easy to see that 
\[
(2 \pi) ^{-n/2} \int_ {\mathbb{R}^{n}} \overline {{f _{2} (- x)}} e ^{- i \langle u, x \rangle} \dd x = \overline {{\hat {f} _{2} (u)}},
\]
and so, by \eqref{eq:VI.2.25}, we obtain \eqref{eq:VI.2.27}. Next, since the Fourier transform of the function of $y$, $\overline{f_2(x - y)}$, is $\overline{f_2(u)} e^{-i\langle u,x\rangle}$ containing a parameter $x$, we obtain \eqref{eq:VI.2.28} by \eqref{eq:VI.2.25}. The rest of the Theorem is clear from \eqref{eq:VI.2.28} and $\widetilde{f} = f$. 
\end{proof}

The Negative Norm. The Sobolev space $W^{k,2}(\Omega)$ was defined in Chapter~\ref{chap:I}, \ref{sec:I.9}. Let $f(x) \in W^{k,2}(\mathbb{R}^{n})$. Since $f(x) \in L^2(\mathbb{R}^{n})$, $f$ gives rise to a slowly increasing measure $|f(x)| \dd x$ in $\mathbb{R}^{n}$. Thus we can define the Fourier transform $\hat{T}_f$ of the tempered distribution $T_f$. We have, by \eqref{eq:VI.2.17}, 
\[
\widehat {D ^{\alpha} T _{f}} = (i) ^{| \alpha |} \prod\limits_{j=1} ^{n} x _{j} ^ {\alpha_{j}} \hat {T} _{f}.
\]

By the definition of the space $W^{k,2}(\mathbb{R}^{n})$, $D^\alpha T_f \in L^2(\mathbb{R}^{n})$ for $|\alpha| \leqq k$. Thus, by Plancherel's theorem for $L^2(\mathbb{R}^{n})$, 
$\|\widehat{D^{\alpha}T_{f}}\|_{0}=\|D^{\alpha}T_{f}\|_{0}, \text{ where } \|\|_{0} \text{ is the } L^{2}(\mathbb{R}^{n})-\text{ norm }.$ 

Hence we see that $(1 + |x|^2)^{k/2}\hat{T}_f \in L^2(\mathbb{R}^{n})$, and so it is easy to show that the norm $\|f\|_k = \left(\sum\limits_{|\alpha| \leqq k} \int_{\mathbb{R}^{n}} |D^\alpha T_f|^2 \dd x\right)^{1/2}$ is equivalent to the norm 
\begin{equation}\label{eq:VI.2.30}
\left\| (1 + | x | ^{2}) ^{k/2} \hat {T} _{f} \right\| _{0} = \left\| f \right\| _{k} ^{\prime},
\end{equation}
in the sense that there exist two positive constants $c_{1}$ and $c_{2}$ for which we have 
\[
c _{1} \leqq \| f \| _{k} / \| f \| _{k} ^{\prime} \leqq c _{2} \text{ whenever } f \in W ^{k,2} (\mathbb{R}^{n})
\]

We may thus renorm the space $W^{k,2}(\mathbb{R}^{n})$ by $\|f\|_k'$; $W^{k,2}(\mathbb{R}^{n})$ may thus be defined as the totality of $f \in L^2(\mathbb{R}^{n})$ such that $\|f\|_k'$ is finite. One advantage of the new formulation of $W^{k,2}(\mathbb{R}^{n})$ is that we may also consider the case of a negative exponent $k$. Then, as in the case of $L^2(\mathbb{R}^{n})$ pertaining to the ordinary Lebesgue measure $\dd x$, we see that the dual space of the renormed space $W^{k,2}(\mathbb{R}^{n})$ is the space $W^{-k,2}(\mathbb{R}^{n})$ normed by $\|f\|_{-k}'$. This observation due to L. \textsc{Schwartz} \cite{schwartz298} is of an earlier date than the introduction of the negative norm by P. \textsc{Lax} \cite{lewy172}. 


\yossection{3}{Convolutions}
\label{sec:VI.3}

We define the convolution (Faltung) of two functions $f, g$ of $C(\mathbb{R}^{n})$, one of which has a compact support, by (cf. the case in which $f, g \in \mathfrak{S}(\mathbb{R}^{n})$ in Chapter~\ref{chap:VI}, \ref{sec:VI.1}) 
\begin{equation}\label{eq:VI.3.1}
(f * g) (x) = \int_ {\mathbb{R}^{n}} f (x - y) g (y) \dd y = \int_ {\mathbb{R}^{n}} f (y) g (x - y) \dd y = (g * f) (x).
\end{equation}

Suggested by this formula, we define the convolution of a $T \in \mathfrak{D}(\mathbb{R}^{n})'$ and a $\varphi \in \mathfrak{D}(\mathbb{R}^{n})$ (or a $T \in \mathfrak{E}(\mathbb{R}^{n})'$ and a $\varphi \in \mathfrak{E}(\mathbb{R}^{n})$) by 
\begin{equation}\label{eq:VI.3.2}
(T * \varphi) (x) = T _{[y]} (\varphi (x - y)),
\end{equation}
where $T_{[y]}$ indicates that we apply the distribution $T$ on test functions of $y$. 

\begin{proposition}\label{prop:VI.3.Proposition.1} $(T * \varphi)(x) \in C^{\infty}(\mathbb{R}^{n})$ and $\supp(T * \varphi) \subseteq \supp(T) + \supp(\varphi)$, that is, 
\[
\supp (T * \varphi) \subseteq \left\{w \in \mathbb{R}^{n}; w = x + y, x \in \supp (T), y \in \supp (\varphi) \right\}.
\]
\end{proposition}

Moreover, we have 
\begin{equation}\label{eq:VI.3.3}
D ^{\alpha} (T * \varphi) = T * (D ^{\alpha} \varphi) = (D ^{\alpha} T) * \varphi .
\end{equation}

\begin{proof}
 Let $\varphi \in \mathfrak{D}(\mathbb{R}^{n})$ (or $\in \mathfrak{E}(\mathbb{R}^{n})$). If $\lim\limits_{h\to 0}x^h = x$, then, as functions of $y$, $\lim\limits_{h\to 0}\varphi (x^h -y) = \varphi (x - y)$ in $\mathfrak{D}(\mathbb{R}^{n})$ (or in $\mathfrak{E}(\mathbb{R}^{n})$). Hence $T_{[y]}(\varphi (x - y)) = (T*\varphi)(x)$ is continuous in $x$. The inclusion relation of the supports is proved by the fact that $T_{[y]}(\varphi (x - y)) = 0$ unless the support of $T$ and that of $\varphi (x - y)$ as a function of $y$ meet. Next let $e_j$ be the unit vector of $\mathbb{R}^{n}$ along the $x_j$-axis and consider the expression 
\[
T _{[y]} \left(\left(\varphi (x + h e _{j} - y) - \varphi (x - y)\right) / h\right). \]

When $h \to 0$, the function enclosed by the outer parenthesis converges, as a function of $y$, to $\left(\frac{\partial \varphi}{\partial x_j}\right)(x - y)$ in $\mathfrak{D}(\mathbb{R}^{n})$ (or in $\mathfrak{E}(\mathbb{R}^{n})$). Thus we have proved 
\[
\frac {\partial}{\partial x _{j}} (T * \varphi) (x) = \left(T * \frac {\partial \varphi}{\partial x _{j}}\right) (x).
\]

Moreover, we have 
\[
\left(T * \frac {\partial \varphi}{\partial x _{j}}\right) (x) = T _{[y]} \left(- \frac {\partial \varphi (x - y)}{\partial y _{j}}\right) = \frac {\partial T _{[v]}}{\partial y _{j}} (\varphi (x - y)) = \left(\frac {\partial T}{\partial x _{j}} * \varphi\right) (x).
\]
\end{proof}

\begin{proposition}\label{prop:VI.3.Proposition.2} If $\varphi$ and $\psi$ are in $\mathfrak{D}(\mathbb{R}^{n})$ and $T \in \mathfrak{D}(\mathbb{R}^{n})'$ (or $\varphi \in \mathfrak{E}(\mathbb{R}^{n})$, $\psi \in \mathfrak{D}(\mathbb{R}^{n})$ and $T \in \mathfrak{E}(\mathbb{R}^{n})'$), then 
\begin{equation}\label{eq:VI.3.4}
(T * \varphi) * \psi = T * (\varphi * \psi).
\end{equation}
\end{proposition}

\begin{proof}
 We approximate the function $(\varphi * \psi)(x)$ by the Riemann sum 
\[
f _{h} (x) = h ^{n} \sum\limits_{h} \varphi (x - k h) \psi (k h),
\]
where $h$ > $0$ and $k$ ranges over points of $\mathbb{R}^{n}$ with integral coordinates. Then, for every differentiation $D^{\alpha}$ and for every compact set of $x$, 
\[
D ^{\alpha} f _{h} (x) = h ^{n} \sum\limits_{h} D ^{\alpha} \varphi (x - k h) \psi (k h)
\]
converges, as $h \downarrow 0$, to $((D^{\alpha}\varphi) * \psi)(x) = (D^{\alpha}(\varphi * \psi))(x)$ uniformly in $x$. Hence we see that $\lim\limits_{h \downarrow 0} f_h = \varphi * \psi$ in $\mathfrak{D}(\mathbb{R}^{n})$ (or in $\mathfrak{E}(\mathbb{R}^{n})$). Therefore, by the linearity and the continuity of $T$, we have 
\[ \begin{aligned}
&(T * (\varphi * \psi)) (x)  = \lim\limits _{h \downarrow 0} (T * f _{h}) (x) = \lim\limits _{h \downarrow 0} h ^{n} \sum\limits_{k} (T * \varphi) (x - k h) \psi (k h) = ((T * \varphi) * \psi) (x).
\end{aligned}
\]
\end{proof}

\begin{definitionu} Let $\varphi \in \mathfrak{D}(\mathbb{R}^{n})$ be non-negative, $\int_{\mathbb{R}^{n}}\varphi \dd x = 1$ and such that $\supp(\varphi) \subseteq \{x \in \mathbb{R}^{n}; |x| \leqq 1\}$. We may, for instance, take 
\[
\begin{aligned}
&\varphi (x) = \exp (1 / (| x | ^{2} - 1)) / \int_{|x|<1} \exp (1 / (| x | ^{2} - 1)) \dd x \text{ if } | x | < 1, = 0 \text{ if } | x | \geqq 1.
\end{aligned}
\]
\end{definitionu}

We write $\varphi_{\varepsilon}(x)$ for $\varepsilon^{-n}\varphi(x/\varepsilon)$, $\varepsilon > 0$, and call $T * \varphi_{\varepsilon}$ the regularization of $T \in \mathfrak{D}(\mathbb{R}^{n})'$ (or $\in \mathfrak{E}(\mathbb{R}^{n})'$) through $\varphi_{\varepsilon}(x)$'s (Cf. Chapter~\ref{chap:I}, \ref{sec:I.1}). 

\begin{theorem}\label{thm:VI.3.Theorem.1} Let $T \in \mathfrak{D}(\mathbb{R}^{n})'$ (or $\in \mathfrak{E}(\mathbb{R}^{n})'$). Then $\lim\limits_{\varepsilon \downarrow 0} (T * \varphi_\varepsilon) = T$ in the weak* topology of $\mathfrak{D}(\mathbb{R}^{n})'$ (or of $\mathfrak{E}(\mathbb{R}^{n})'$). In this sense, $\varphi_\varepsilon$ is called an approximate identity. 
\end{theorem}

For the proof we prepare a 

\begin{lemmau} For any $\psi \in \mathfrak{D}(\mathbb{R}^{n})$ (or $\in \mathfrak{E}(\mathbb{R}^{n})$), we have $\lim\limits_{\varepsilon \downarrow 0} \psi * \varphi_{\varepsilon} = \psi$ in $\mathfrak{D}(\mathbb{R}^{n})$ (or in $\mathfrak{E}(\mathbb{R}^{n})$). 
\end{lemmau}

\begin{proof}
 We first observe that $\supp(\psi * \varphi_{\varepsilon}) \subseteq \supp(\psi) + \supp(\varphi_{\varepsilon}) = \supp(\psi) + \varepsilon$. We have, by \eqref{eq:VI.3.3}, $D^{\alpha}(\psi * \varphi_{\varepsilon}) = (D^{\alpha}\psi) * \varphi_{\varepsilon}$. Hence we have to show that $\lim\limits_{\varepsilon \downarrow 0} (\psi * \varphi_{\varepsilon})(x) = \psi(x)$ uniformly on any compact set of $x$. But, by $\int \varphi_{\varepsilon}(y) \dd y = 1$, 
\[
(\psi * \varphi_{\varepsilon}) (x) - \psi (x) = \int_ {\mathbb{R}^{n}} \{\psi (x - y) - \psi (x) \} \varphi_{\varepsilon} (y) \dd y. \]

Therefore, by $\varphi_{\varepsilon}(x) \geqq 0$, $\int_{\mathbb{R}^{n}} \varphi_{\varepsilon}(y) \dd y = 1$ and the uniform continuity of $\psi(x)$ on any compact interval of $x$, we obtain the Lemma. 
\end{proof}

\begin{proof}[Proof of Theorem~\ref{thm:VI.3.Theorem.1}.]
We have, by $\check{\psi}(x) = \psi(-x)$,
\begin{equation}\label{eq:VI.3.5}
T (\psi) = (T * \check {\psi}) (0).
\end{equation}

Hence we have to prove that $\lim\limits_{\varepsilon \downarrow 0} ((T * \varphi_{\varepsilon}) * \check{\psi}) (0) = (T * \check{\psi}) (0)$. But, as proved in \eqref{eq:VI.3.4}, $(T * \varphi_{\varepsilon}) * \check{\psi} = T * (\varphi_{\varepsilon} * \check{\psi})$ and so, by \eqref{eq:VI.3.5}, $(T * (\varphi_{\varepsilon} * \check{\psi})) (0) = T((\varphi_{\varepsilon} * \check{\psi})^{\vee})$. Hence we obtain $\lim\limits_{\varepsilon \downarrow 0} ((T * \varphi_{\varepsilon}) * \check{\psi}) (0) = T((\check{\psi})^{\vee}) = T(\psi)$ by the Lemma. 

We next prove a theorem which concerns with a characterization of the operation of convolution. 
\end{proof}

\begin{theorem}[L. \textsc{Schwartz}]\label{thm:VI.3.Theorem.2} Let $L$ be a continuous linear mapping on $\mathfrak{D}(\mathbb{R}^{n})$ into $\mathfrak{F}(\mathbb{R}^{n})$ such that 
\begin{equation}\label{eq:VI.3.6}
L \tau_{h} \varphi = \tau_{h} L \varphi \text{ for any } h \in \mathbb{R}^{n} \text{ and } \varphi \in \mathfrak {D} (\mathbb{R}^{n}),
\end{equation}
where the translation operator $\tau_{h}$ is defined by 
\begin{equation}\label{eq:VI.3.7}
\tau_{h} \varphi (x) = \varphi (x - h).
\end{equation}
\end{theorem}

Then there exists a uniquely determined $T \in \mathfrak{D}(\mathbb{R}^{n})'$ such that $L\varphi = T * \varphi$. Conversely, any $T \in \mathfrak{D}(\mathbb{R}^{n})'$ defines a continuous linear map $L$ on $\mathfrak{D}(\mathbb{R}^{n})$ into $\mathfrak{E}(\mathbb{R}^{n})$ by $L\varphi = T * \varphi$ such that $L$ satisfies \eqref{eq:VI.3.6}. 

\begin{proof}
 Since $\varphi \to \check{\varphi}$ is a continuous linear map of $\mathfrak{D}(\mathbb{R}^{n})$ onto $\mathfrak{D}(\mathbb{R}^{n})$, the linear map $T: \check{\varphi} \to (L\varphi)$ (0) defines a distribution $T \in \mathfrak{D}(\mathbb{R}^{n})'$. Hence, by \eqref{eq:VI.3.5}, $(L\varphi)$ (0) = $T(\check{\varphi}) = (T * \varphi)$ (0). If we replace $\varphi$ by $\tau_h\varphi$ and make use of condition \eqref{eq:VI.3.6}, then we obtain $(L\varphi)$ ( $h) = (T * \varphi)$ ( $h$). The converse part of Theorem~\ref{thm:VI.3.Theorem.2} is easily proved by \eqref{eq:VI.3.2}, Proposition~\ref{prop:VI.3.Proposition.1} and \eqref{eq:VI.3.5}. 
\end{proof}

\begin{corollaryu} Let $T_{1} \in \mathfrak{D}(\mathbb{R}^{n})'$ and $T_{2} \in \mathfrak{E}(\mathbb{R}^{n})'$. Then the convolution $T_{1} * T_{2}$ may be defined, through the continuous linear map $L$ on $\mathfrak{D}(\mathbb{R}^{n})$ into $\mathfrak{E}(\mathbb{R}^{n})$ as follows: 
\begin{equation}\label{eq:VI.3.8}
\left(T _{1} * T _{2}\right) * \varphi = L (\varphi) = T _{1} * \left(T _{2} * \varphi\right), \varphi \in \mathfrak {D} \left(\mathbb{R}^{n}\right).
\end{equation}
\end{corollaryu}

\begin{proof}
 The map $\varphi \to T_2 * \varphi$ is continuous linear on $\mathfrak{D}(\mathbb{R}^{n})$ into $\mathfrak{D}(\mathbb{R}^{n})$, since $\supp(T_2)$ is compact. Hence the map $\varphi \to T_1 * (T_2 * \varphi)$ is continuous linear on $\mathfrak{D}(\mathbb{R}^{n})$ into $\mathfrak{E}(\mathbb{R}^{n})$. It is easily verified that condition \eqref{eq:VI.3.6} is satisfied for the present $L$. 
\end{proof}

\begin{remarku} We see, by \eqref{eq:VI.3.4}, that the above definition of the convolution $T_{1} \ast T_{2}$ agrees with the previous one in the case where $T_{2}$ is defined by a function $\in \mathfrak{D}(\mathbb{R}^{n})$. It is to be noted that we may also define $T_{1} \ast T_{2}$ by 
\begin{equation}\tag{8'}\label{eq:VI.3.8'}
\left(T _{1} * T _{2}\right) (\varphi) = \left(T _{1(x)} \times T _{2(y)}\right) (\varphi (x + y)), \varphi \in \mathfrak {D} \left(\mathbb{R}^{n}\right),
\end{equation}
\end{remarku}

where $T_{1(x)} \times T_{2(y)}$ is the tensor product of $T_{1}$ and $T_{2}$. See L. \textsc{Schwartz} \cite{schwartz294}. 

\begin{theorem}\label{thm:VI.3.Theorem.3} Let $T_{1} \in \mathfrak{D}(\mathbb{R}^{n})'$ and $T_{2} \in \mathfrak{E}(\mathbb{R}^{n})'$. Then we can define another ``convolution'' $T_{2} \boxtimes T_{1}$ through the continuous linear map $L$: 
\[
\varphi \rightarrow T _{2} * (T _{1} * \varphi) \text{ on } \mathfrak {D} (\mathbb{R}^{n}) \text{ into } \mathfrak {E} (\mathbb{R}^{n})
\]
as follows. $(T_{2} \boxtimes T_{1}) (*\varphi) = L(\varphi)$, $\varphi \in \mathfrak{D}(\mathbb{R}^{n})$. Then we can prove $T_{2} \boxtimes T_{1} = T_{1} * T_{2}$ so that the convolution is commutative if it is defined either as $T_{1} * T_{2}$ or as $T_{2} \boxtimes T_{1}$. 
\end{theorem}

\begin{proof}
 The map $\varphi \to T_1 * \varphi$ is linear continuous on $\mathfrak{D}(\mathbb{R}^{n})$ into $\mathfrak{E}(\mathbb{R}^{n})$. Hence the map $\varphi \to T_2 * (T_1 * \varphi)$ is linear continuous on $\mathfrak{D}(\mathbb{R}^{n})$ into $\mathfrak{E}(\mathbb{R}^{n})$. Thus $T_2 \boxplus T_1$ is well-defined. Next, we have for any $\varphi_1, \varphi_2 \in \mathfrak{D}(\mathbb{R}^{n})$, 
\[
\begin{aligned} (T _{1} * T _{2}) * (\varphi_{1} * \varphi_{2}) & = T _{1} * (T _{2} * (\varphi_{1} * \varphi_{2})) = T _{1} * ((T _{2} * \varphi_{1}) * \varphi_{2}) \\ & = T _{1} * (\varphi_{2} * (T _{2} * \varphi_{1})) = (T _{1} * \varphi_{2}) * (T _{2} * \varphi_{1}) \end{aligned}
\]
by the commutativity of the convolution of functions and Proposition~\ref{prop:VI.3.Proposition.2}, observing that $\supp(T_2 * \varphi_1)$ is compact because of $T_2 \in \mathfrak{E}(\mathbb{R}^{n})'$. Similarly we obtain 
\[ \begin{aligned} (T _{2} \boxtimes T _{1}) * (\varphi_{1} * \varphi_{2}) & = T _{2} * (T _{1} * (\varphi_{1} * \varphi_{2})) = T _{2} * ((T _{1} * \varphi_{2}) * \varphi_{1}) \\ & = T _{2} * (\varphi_{1} * (T _{1} * \varphi_{2})) = (T _{2} * \varphi_{1}) * (T _{1} * \varphi_{2}) \end{aligned}
\]

Thus $(T_{1} * T_{2}) * (\varphi_{1} * \varphi_{2}) = (T_{2} \boxtimes T_{1}) * (\varphi_{1} * \varphi_{2})$, and so, by \eqref{eq:VI.3.5} and the Lemma above, we obtain $(T_{1} * T_{2})$ ( $\varphi$) = $(T_{2} \boxtimes T_{1})$ ( $\varphi$) for all $\varphi \in \mathfrak{D}(\mathbb{R}^{n})$, that is, $(T_{1} * T_{2}) = (T_{2} \boxtimes T_{1})$. 
\end{proof}

\begin{corollaryu}
\begin{equation}\label{eq:VI.3.9}
T _{1} * (T _{2} * T _{3}) = (T _{1} * T _{2}) * T _{3}
\end{equation}
if all $T_{j}$ except one have compact support, 

\begin{equation}\label{eq:VI.3.10}
D ^{\alpha} (T _{1} * T _{2}) = (D ^{\alpha} T _{1}) * T _{2} = T _{1} * (D ^{\alpha} T _{2}).
\end{equation}
\end{corollaryu}

\begin{proof}
 We have, by \eqref{eq:VI.3.5} and the definition of $T_{1} * T_{2}$, 
\[
\begin{aligned}
&(T _{1} * (T _{2} * T _{3})) (\varphi)  = \left((T _{1} * (T _{2} * T _{3})) * \check {\varphi}\right) (0) = (T _{1} * ((T _{2} * T _{3}) * \check {\varphi})) (0) \\
&= (T _{1} * (T _{2} * (T _{3} * \check {\varphi}))) (0)
\end{aligned}
\]
and similarly 
\[
\left(\left(T _{1} * T _{2}\right) * T _{3}\right) (\varphi) = \left(T _{1} * \left(T _{2} * \left(T _{3} * \check {\varphi}\right)\right)\right) (0),
\]
so that ($9$) holds. 

The proof of \eqref{eq:VI.3.10} is as follows. We observe that, by \eqref{eq:VI.3.3}, 
\begin{equation}\label{eq:VI.3.11}
(D ^{\alpha} T _{\delta}) * \varphi = T _{\delta} * (D ^{\alpha} \varphi) = D ^{\alpha} (T _{\delta} * \varphi) = D ^{\alpha} \varphi ,
\end{equation}
which implies 
\[
(D ^{\alpha} T) * \varphi = T * (D ^{\alpha} \varphi) = T * ((D ^{\alpha} T _{\delta}) * \varphi) = (T * D ^{\alpha} T _{\delta}) * \varphi ,
\]
that is, by \eqref{eq:VI.3.5}, 
\begin{equation}\label{eq:VI.3.12}
D ^{\alpha} T = (D ^{\alpha} T _{\delta}) * T.
\end{equation}

Hence, by the commutativity (Theorem~\ref{thm:VI.3.Theorem.3}) and the associativity ($9$), 
\[
\begin{aligned} D ^{\alpha} (T _{1} \ast T _{2}) & = (D ^{\alpha} T _{\delta}) \ast (T _{1} \ast T _{2}) = ((D ^{\alpha} T _{\delta}) \ast T _{1}) \ast T _{2} = (D ^{\alpha} T _{1}) \ast T _{2} \\ & = (D ^{\alpha} T _{\delta}) \ast (T _{2} \ast T _{1}) = ((D ^{\alpha} T _{\delta}) \ast T _{2}) \ast T _{1} = (D ^{\alpha} T _{2}) \ast T _{1}. \end{aligned}
\]
\end{proof}

The Fourier Transform of the Convolution. We first prove a theorem which will be sharpened to the Paley-Wiener theorem in the next section. 

\begin{theorem}\label{thm:VI.3.Theorem.4} The Fourier transform of a distribution $T \in \mathfrak{E}(\mathbb{R}^{n})'$ is given by the function 
\begin{equation}\label{eq:VI.3.13}
\hat {T} (\xi) = (2 \pi) ^{-n/2} T _{[x]} (e ^{- i \langle x, \xi \rangle}).
\end{equation}
\end{theorem}

\begin{proof}
 When $\varepsilon \downarrow 0$, the regularization $T_{\varepsilon} = T * \varphi_{\varepsilon}$ tends to $T$ in the weak* topology of $\mathfrak{S}(\mathbb{R}^{n})'$ and a fortiori in the weak* topology of $\mathfrak{S}(\mathbb{R}^{n})'$. This we see from $(T * \varphi_{\varepsilon})(\psi) = (T * (\varphi_{\varepsilon} * \check{\psi})) (0) = T_{[x]}((\varphi_{\varepsilon} * \check{\psi})(-x))$ and the Lemma. Thus, by the continuity of the Fourier transform in the weak* topology of $\mathfrak{S}(\mathbb{R}^{n})'$, $\lim\limits_{\varepsilon \downarrow 0} (\widehat{T * \varphi_{\varepsilon}}) = \hat{T}$ in the weak* topology of $\mathfrak{S}(\mathbb{R}^{n})'$. Now formula \eqref{eq:VI.3.13} is clear for the distribution defined by the function $(T * \varphi_{\varepsilon})(x)$. Hence 
\[
(2 \pi) ^{n/2} (\widehat {T * \varphi_{\varepsilon}}) (\xi) = (T * \varphi_{\varepsilon}) _{[x]} (e ^{- i \langle x, \xi \rangle}),
\]
which is, by \eqref{eq:VI.3.5}, $= (T_{[x]} * (\varphi_{\varepsilon} * e^{-i\langle x,\xi \rangle}))$ (0) $= T_{[x]}(\check{\varphi}_{\varepsilon} * e^{-i\langle x,\xi \rangle})$. The last expression tends, as $\varepsilon \downarrow 0$, to $T_{[x]}(e^{i\langle x,\xi \rangle})$ uniformly in $\xi$ on any bounded set of $\xi$ of the complex $n$-space. This proves Theorem~\ref{thm:VI.3.Theorem.4}. 
\end{proof}

\begin{theorem}\label{thm:VI.3.Theorem.5} If we define the convolution of a distribution $T \in \mathfrak{S}(\mathbb{R}^{n})'$ and a function $\varphi \in \mathfrak{S}(\mathbb{R}^{n})$ by $(T * \varphi)(x) = T_{[y]}(\varphi(x - y))$, then the linear map $L$ on $\mathfrak{S}(\mathbb{R}^{n})$ into $\mathfrak{E}(\mathbb{R}^{n}): \varphi \to T * \varphi$ is characterized by the continuity and the translation invariance $\tau_h L = L\tau_h$. 
\end{theorem}

\begin{proof}
 Similar to the proof of Theorem~\ref{thm:VI.3.Theorem.2}. 
\end{proof}

\begin{theorem}\label{thm:VI.3.Theorem.6} If $T \in \mathfrak{S}(\mathbb{R}^{n})'$ and $\varphi \in \mathfrak{S}(\mathbb{R}^{n})$, then 
\begin{equation}\label{eq:VI.3.14}
(\widehat {T * \varphi}) = (2 \pi) ^{n/2} \hat {\varphi} \hat {T}
\end{equation}
\end{theorem}

If $T_{1} \in \mathfrak{S}(\mathbb{R}^{n})'$ and $T_{2} \in \mathfrak{E}(\mathbb{R}^{n})'$, then 
\begin{equation}\label{eq:VI.3.15}
(\widehat {T _{1} * T _{2}}) = (2 \pi) ^{n/2} \hat {T} _{2} \cdot \hat {T} _{1},
\end{equation}
which has a sense since, as proved above in Theorem~\ref{thm:VI.3.Theorem.4}, $\hat{T}_{2}$ is given by a function. 

\begin{proof}
 Let $\psi \in \mathfrak{S}(\mathbb{R}^{n})$. Then the Fourier transform of $\hat{\varphi} \cdot \psi$ is, by \eqref{eq:VI.3.13} of Chapter~\ref{chap:VI}, \ref{sec:VI.1}, equal to $(2\pi)^{-n/2} \hat{\bar{\varphi}} * \hat{\psi} = (2\pi)^{-n/2} \check{\varphi} * \hat{\psi}$. Thus 
\[
\begin{aligned}
&\widehat {(T * \varphi)} (\psi)  = (T * \varphi) (\hat {\psi}) = ((T * \varphi) * \check {\tilde {\varphi}}) (0) = (T * (\varphi * \check {\tilde {\psi}})) \\
&= T ((\varphi * \check {\tilde {\psi}}) ^{\vee}) = T (\check {\varphi} * \hat {\psi}) = T ((2 \pi) ^{n/2} (\hat {\varphi} \cdot \psi) ^{\wedge}) = (2 \pi) ^{n/2} \hat {T} (\hat {\varphi} \psi) = (2 \pi) ^{n/2} \hat {\varphi} \hat {T} (\psi),
\end{aligned}
\]
which proves \eqref{eq:VI.3.14}. 

Let $\psi_{\varepsilon}$ be the regularization $T_{2} * \varphi_{\varepsilon}$. Then the Fourier transform of $T_{1} * \psi_{\varepsilon} = T_{1} * (T_{2} * \varphi_{\varepsilon}) = (T_{1} * T_{2}) * \varphi_{\varepsilon}$ is, by \eqref{eq:VI.3.14}, equal to 
\[
(2 \pi) ^{n/2} \hat {T} _{1} \cdot \hat {\psi} _{\varepsilon} = (2 \pi) ^{n/2} \hat {T} _{1} \cdot (2 \pi) ^{n/2} \hat {T} _{2} \cdot \hat {\varphi} _{\varepsilon} = (2 \pi) ^{n/2} (\widehat {T _{1} * T _{2}}) \cdot \hat {\varphi} _{\varepsilon}.
\]

Hence we obtain \eqref{eq:VI.3.15}, by letting $\varepsilon \downarrow 0$ and using $\lim\limits_{\varepsilon \downarrow 0} \hat{\varphi}_{\varepsilon}(x) = 1$. 
\end{proof}


\yossection{4}{The Paley-Wiener Theorems. The One-sided Laplace Transform}
\label{sec:VI.4}

The Fourier transform of a function $\in C_0^\infty (\mathbb{R}^{n})$ is characterized by the Paley-Wiener Theorem for Functions. An entire holomorphic function $F(\zeta) = F(\zeta_1,\zeta_2,\dots ,\zeta_n)$ of $n$ complex variables $\zeta_j = \xi_j + i\eta_j$ ( $j = 1,2,\ldots ,n$) is the Fourier-Laplace transform 
\begin{equation}\label{eq:VI.4.1}
F (\zeta) = (2 \pi) ^{-n/2} \int_ {\mathbb{R}^{n}} e ^{- i \langle \zeta , x \rangle} f (x) \dd x
\end{equation}
of a function $f \in C_0^\infty(\mathbb{R}^{n})$ with $\supp(f)$ contained in the sphere $|x| \leqq B$ of $\mathbb{R}^{n}$ iff there exists for every integer $N$ a positive constant $C_N$ such that 
\begin{equation}\label{eq:VI.4.2}
\left| F (\zeta) \right| \leqq C _{N} (1 + | \zeta |) ^{-N} e ^{B | \operatorname{Im} \zeta |}.
\end{equation}

\begin{proof}
 The necessity is clear from 
\[
\prod\limits_{j=1} ^{n} (i \zeta_{j}) ^ {\beta_{j}} F (\zeta) = (2 \pi) ^{-n/2} \int_{| x | \leqq B} e ^{- i \langle \zeta , x \rangle} D ^{\beta} f (x) \dd x
\]
which is obtained by partial integration. 

The sufficiency. Let us define 
\begin{equation}\label{eq:VI.4.3}
f (x) = (2 \pi) ^{-n/2} \int_ {\mathbb{R}^{n}} e ^{i \langle x, \xi \rangle} F (\xi) \dd \xi .
\end{equation}

Then, as for the case of functions $\in \mathfrak{S}(\mathbb{R}^{n})$, we can prove that the Fourier transform $\hat{f}(\xi)$ is equal to $F(\xi)$ and $f \in C^{\infty}(\mathbb{R}^{n})$. The last assertion is proved by differentiation: 
\begin{equation}\label{eq:VI.4.4}
D ^{\beta} f (x) = (2 \pi) ^{-n/2} \int_ {\mathbb{R}^{n}} e ^{i \langle x, \xi \rangle} \prod\limits_{j=1} ^{n} (i \xi_{j}) ^ {\beta_{j}} F (\xi) \dd \xi ,
\end{equation}
by making use of condition \eqref{eq:VI.4.2}. The same condition \eqref{eq:VI.4.2} and Cauchy's integral theorem enable us to shift the real domain of integration in \eqref{eq:VI.4.3} into the complex domain so that we obtain 
\begin{equation}\tag{3'}\label{eq:VI.4.3'}
f (x) = (2 \pi) ^{-n/2} \int_ {\mathbb{R}^{n}} e ^{i \langle x, \xi + i \eta \rangle} F (\xi + i \eta) \dd \xi
\end{equation}
for arbitrary real $\eta$ of the form $\eta = \alpha x / |x|$ with $\alpha > 0$. We thus obtain, by taking $N = n + 1$, 
\[
\left| f (x) \right| \leqq C _{N} e ^{B | \eta | - \langle x, \eta \rangle} (2 \pi) ^{-n/2} \int_ {\mathbb{R}^{n}} (1 + | \xi |) ^{-N} \dd \xi .
\]

If $|x| > B$, we obtain $f(x) = 0$ by letting $\alpha \uparrow + \infty$. Hence $\supp(f) \subseteq \{x \in \mathbb{R}^{n}; |x| \leqq B\}$. 
\end{proof}

The above theorem may be generalized to distributions with compact support. Thus we have 

The Paley-Wiener Theorem for Distributions $\in \mathfrak{E}(\mathbb{R}^{n})'$ (L. \textsc{Schwartz}). An entire holomorphic function $F(\zeta) = F(\zeta_1, \ldots, \zeta_n)$ of $n$ complex variables $\zeta_j = \xi_j + i\,\eta_j$ ( $j = 1, 2, \ldots, n$) is the Fourier-Laplace transform of a distribution $T \in \mathfrak{E}(\mathbb{R}^{n})'$ iff for some positive constants $B$, $N$ and $C$ 
\begin{equation}\label{eq:VI.4.5}
\left| F (\zeta) \right| \leqq C (1 + | \zeta |) ^{N} e ^{B | \operatorname{Im} \zeta |}.
\end{equation}

\begin{proof}
 The necessity is clear from the fact (Theorem~\ref{thm:I.13.Theorem.2} in Chapter~\ref{chap:I}, \ref{sec:I.13}) that, if $T \in \mathfrak{E}(\mathbb{R}^{n})'$, there exist positive constants $C, N$ and $B$ such that 
\[
\left| T (\varphi) \right| \leqq C \sum\limits_{| \beta | \leqq N} \sup\limits _{| x | \leqq B} \left| D ^{\beta} \varphi (x) \right| \quad \text{ whenever } \quad \varphi \in \mathfrak {G} (\mathbb{R}^{n}) \]

For, we have only to take $\varphi(x) = e^{-i\langle x,\xi\rangle}$ and apply \eqref{eq:VI.4.13} of the preceding section. 

The sufficiency. $F(\xi)$ is in $\mathfrak{S}(\mathbb{R}^{n})'$ and so it is the Fourier transform of a distribution $T \in \mathfrak{S}(\mathbb{R}^{n})'$. The Fourier transform of the regularization $T_{\varepsilon} = T * \varphi_{\varepsilon}$ is $(2\pi)^{n/2} \hat{T} \cdot \hat{\varphi}_{\varepsilon}$ by \eqref{eq:VI.4.14} of the preceding section. Since the supp $(\varphi_{\varepsilon})$ is in the sphere $|x| \leqq \varepsilon$ of $\mathbb{R}^{n}$, we have, by the preceding theorem, 
\[
\left| \hat {\varphi} _{\varepsilon} (\xi) \right| \leqq C ^{\prime} \cdot e ^{\varepsilon | \operatorname{Im} \zeta |}.
\]

Moreover, since $\hat{T}$ is defined by the function $F(\xi)$, we see that $(2\pi)^{n/2} \hat{T} \cdot \hat{\varphi}_{\varepsilon}$ is defiend by a function $(2\pi)^{n/2} F(\xi) \cdot \hat{\varphi}_{\varepsilon}(\xi)$ which, when extended to the complex $n$-space analytically, satisfies the estimate of type \eqref{eq:VI.4.5} with $B$ replaced by $B + \varepsilon$. Thus, by the preceding theorem, $T_{\varepsilon} = T * \varphi_{\varepsilon}$ has its support in the sphere $|x| \leqq B + \varepsilon$ of $\mathbb{R}^{n}$. Thus, letting $\varepsilon \downarrow 0$ and making use of the Lemma in the preceding section, we see that the supp $(T)$ is in the sphere $|x| \leqq B$ of $\mathbb{R}^{n}$. 
\end{proof}
\begin{remarku} The formulation and the proofs of the Paley-Wiener theorems given above are adapted from L. \textsc{Hörmander} \cite{hormander114}. 
\end{remarku}

The Fourier Transform and the One-sided Laplace Transform. Let $g(t) \in L^2(0, \infty)$. Then, for $x > 0$, 
\[
g (t) e ^{-tx} \in L ^{1} (0, \infty) \cap L ^{2} (0, \infty)
\]
as may be seen by Schwarz' inequality. Hence, by Plancherel's theorem, we have, for the Fourier transform 
\begin{equation}\label{eq:VI.4.6}
\begin{aligned}
&f (x + i y)  = (2 \pi) ^{-1/2} \int_{0} ^{\infty} g (t) e ^{-tx} e ^{-ity} \dd t = (2 \pi) ^{-1/2} \int_{0} ^{\infty} g (t) e ^{-t(x+iy)} \dd t \quad (x > 0),
\end{aligned}
\end{equation}
the inequality 
\begin{equation}\label{eq:VI.4.7}
\int_{- \infty} ^{\infty} | f (x + i y) | ^{2} \dd y = \int_{0} ^{\infty} | g (t) | ^{2} e ^{-2tx} \dd t \leqq \int_{0} ^{\infty} | g (t) | ^{2} \dd t.
\end{equation}

The function $f(x + iy)$ is holomorphic in $z = x + iy$ in the right half-plane $\operatorname{Re}(z) = x > 0$, as may be seen by differentiating \eqref{eq:VI.4.6} under the integral sign observing that $g(t) te^{-tz}$ as a function of $t$ belongs to $L^{1}(0, \infty)$ and to $L^{2}(0, \infty)$ when $\operatorname{Re}(z) = x > 0$. We have thus proved 

\begin{theorem}\label{thm:VI.4.Theorem.1} Let $g(t) \in L^2(0, \infty)$. Then the one-sided Laplace transform 
\begin{equation}\tag{6'}\label{eq:VI.4.6'}
f (z) = (2 \pi) ^{-1/2} \int_{0} ^{\infty} g (t) e ^{-tz} \dd t \quad (\operatorname{Re} (z) > 0)
\end{equation}
belongs to the so-called Hardy-Lebesgue class $H^2(0)$, that is, \begin{enumerate}[(i)]
\item $f(z)$ is holomorphic in the right half-plane $\operatorname{Re}(z) > 0$,
\item for each fixed $x > 0$, $f(x + iy)$ as a function of $y$ belongs to $L^2(-\infty, \infty)$ in such a way that
\begin{equation}\tag{7'}\label{eq:VI.4.7'}
\sup\limits _{x>0} \left(\int_{- \infty} ^{\infty} | f (x + i y) | ^{2} \dd y\right) < \infty .
\end{equation}
\end{enumerate}
\end{theorem}

This theorem admits a converse, that is, we have 

\begin{theorem}[\textsc{Paley}-\textsc{Wiener}]\label{thm:VI.4.Theorem.2} Let $f(z) \in H^2(0)$. Then the boundary function $f(iy) \in L^2(-\infty, \infty)$ of $f(x + iy)$ exists in the sense that 
\begin{equation}\label{eq:VI.4.8}
\lim\limits _{x \downarrow 0} \int_{- \infty} ^{\infty} | f (i y) - f (x + i y) | ^{2} \dd y = 0
\end{equation}
in such a way that the inverse Fourier transform 
\begin{equation}\label{eq:VI.4.9}
g (t) = (2 \pi) ^{-1/2} \underset {N \rightarrow \infty} {\text{ l.i.m. }} \int_{-N} ^{N} f (i y) e ^{ity} \dd y
\end{equation}
vanishes for $t$ < $0$ and $f(z)$ may be obtained as the one-sided Laplace transform of $g(t)$. 
\end{theorem}

\begin{proof}
 By the local weak compactness of $L^2(-\infty, \infty)$, we see, that there exists a sequence $\{x_n\}$ and an $f(iy) \in L^2(-\infty, \infty)$ such that 
\[
x _{n} \downarrow 0 \quad \text{ and } \quad \operatorname * {w e a k - l i m} _{n \to \infty} f (x _{n} + i y) = f (i y). \]

For any $\delta > 0$, there exists a sequence $\{N_k\}$ such that 
\[
\lim\limits _{k \rightarrow \infty} N _{k} = \infty \quad \text{ and } \quad \lim\limits _{k \rightarrow \infty} \int_{0+} ^{\delta} | f (x \pm i N _{k}) | ^{2} \dd x = 0,
\]
as may be seen from 
\[
\int_{-N} ^{N} \left\{\int_{0+} ^{\delta} | f (x + i y) | ^{2} \dd x \right\} \dd y \leqq \int_{0+} ^{\delta} \left\{\int_{- \infty} ^{\infty} | f (x + i y) | ^{2} \dd y \right\} \dd x < \infty
\]

Thus, by Schwarz' inequality, 
\begin{equation}\label{eq:VI.4.10}
\lim\limits _{k \rightarrow \infty} \int_{0+} ^{\delta} | f (x \pm i N _{k}) | \dd x = 0.
\end{equation}


From this we obtain 
\begin{equation}\label{eq:VI.4.11}
f (z) = (2 \pi) ^{-1} \int_{- \infty} ^{\infty} \frac {f (i t)}{z - i t} \dd t \quad (\operatorname{Re} (z) > 0).
\end{equation}

The proof may be obtained as follows. By Cauchy's integral representation, we obtain 
\begin{equation}\label{eq:VI.4.12}
f (z) = (2 \pi i) ^{-1} \int_{C} \frac {f (\zeta)}{\zeta - z} \dd \zeta \quad (\operatorname{Re} (z) > 0),
\end{equation}
where the path $C$ of integration is composed of the segments 
\[
\begin{aligned}
&\overrightarrow {x _{0} - i N _{k} , x _{1} - i N _{k}}, \overrightarrow {x _{1} - i N _{k} , x _{1} + i N _{k}}, \overrightarrow {x _{1} + i N _{k} , x _{0} + i N _{k}}, \overrightarrow {x _{0} + i N _{k} , x _{0} - i N _{k}}
\\
&(x _{0} < \operatorname{Re} (z) < x _{1}, - N _{k} < \operatorname{Im} (z) < N _{k})
\end{aligned}
\]
so that the closed contour $C$ encloses the point $z$. Hence, by letting $k \to \infty$ and observing \eqref{eq:VI.4.10}, we obtain 
\[
f (z) = (2 \pi) ^{-1} \int_{- \infty} ^{\infty} \frac {f \left(x _{0} + i t\right)}{z - \left(x _{0} + i t\right)} \dd t + (2 \pi) ^{-1} \int_{- \infty} ^{\infty} \frac {f \left(x _{1} + i t\right)}{\left(x _{1} + i t\right) - z} \dd t.
\]

We see that the second term on the right tends to $0$ as $x_{1} \to \infty$, because of \eqref{eq:VI.4.7} and Schwarz' inequality. We set, in the first term on the right, $x_{0} = x_{n}$ and letting $n \to \infty$ we obtain \eqref{eq:VI.4.11}. Similarly, we obtain 
\begin{equation}\label{eq:VI.4.13}
0 = (2 \pi) ^{-1} \int_{- \infty} ^{\infty} \frac {f (i t)}{i t - z} \dd t (\operatorname{Re} (z) < 0).
\end{equation}

Let us put, for $\operatorname{Re}(z) > 0$, 
\[
h (x) = 0 \quad (\text{ for } x < 0); \quad h (x) = e ^{-zx} \quad (\text{ for } x > 0).
\]

We then have 
\[
\int_{- \infty} ^{\infty} h (x) e ^{ixt} \dd x = \int_{0} ^{\infty} e ^{ixt-zx} \dd x = (z - i t) ^{-1},
\]
and so, by Plancherel's theorem, 
\begin{equation}\label{eq:VI.4.14}
\underset {N \rightarrow \infty} {\text{ l . i . m . }} (2 \pi) ^{-1} \int_{-N} ^{N} \frac {e ^{-itx}}{z - i t} \dd t = \begin{cases}0&(\text{ for } x < 0)\\e ^{-zx}&(\text{ for } x > 0)\end{cases} \text{ if } \operatorname{Re} (z) > 0.
\end{equation}

Similarly we have 
\begin{equation}\tag{14'}\label{eq:VI.4.14'}
\underset {N \rightarrow \infty} {\text{ l . i . m . }} (2 \pi) ^{-1} \int_{-N} ^{N} \frac {e ^{-itx}}{z - i t} \dd t = \begin{cases}- e ^{-zx} (\text{ for } x < 0)\\0&(\text{ for } x > 0)\end{cases} \quad \text{ if } \operatorname{Re} (z) < 0.
\end{equation}

Therefore, by applying the Parseval theorem ($27$) of Chapter~\ref{chap:VI}, \ref{sec:VI.2} to \eqref{eq:VI.4.11}, we obtain the result that $f(z)$ is the one-sided Laplace transform of $g(t)$ given by \eqref{eq:VI.4.9}. By applying the Parseval theorem to \eqref{eq:VI.4.13}, we see also that $g(t) = 0$ for $t < 0$. 


We finally shall prove that \eqref{eq:VI.4.8} is true. By adding \eqref{eq:VI.4.11} and \eqref{eq:VI.4.13}, we see that $f(z)$ admits Poisson's integral representation 
\begin{equation}\label{eq:VI.4.15}
f (z) = f (x + i y) = \frac {x}{\pi} \int_{- \infty} ^{\infty} \frac {f (i t)}{(t - y) ^{2} + x ^{2}} \dd t \quad \text{ whenever } \quad x > 0.
\end{equation}

By virtue of 
\begin{equation}\label{eq:VI.4.16}
\frac {x}{\pi} \int_{- \infty} ^{\infty} \frac {\dd t}{(t - y) ^{2} + x ^{2}} = 1,
\end{equation}
we have 
\[
\left| f (x + i y) - f (i y) \right| \leqq \frac {x}{\pi} \int_{- \infty} ^{\infty} \frac {\left| f ^{+} (s + y) - f ^{+} (y) \right|}{s ^{2} + x ^{2}} \dd s, \text{ where } f ^{+} (y) = f (i y),
\]
and so 
\[
\begin{aligned}
&\int_{- \infty} ^{\infty} | f (x + i y) - f (i y) | ^{2} \dd y  \leqq \left(\frac {x}{\pi}\right) ^{2} \int_{- \infty} ^{\infty} \left\{\int_{- \infty} ^{\infty} \frac {| f ^{+} (s + y) - f ^{+} (y) |}{s ^{2} + x ^{2}} \dd s \right\} ^{2} \dd y \\
&\leqq \frac {x ^{2}}{\pi^{2}} \int_{- \infty} ^{\infty} \left(\int_{- \infty} ^{\infty} \frac {\dd s}{s ^{2} + x ^{2}}\right) \left(\int_{- \infty} ^{\infty} \frac {| f ^{+} (s + y) - f ^{+} (y) | ^{2}}{s ^{2} + x ^{2}} \dd s\right) \dd y \\
&= \frac {x}{\pi} \int_{- \infty} ^{\infty} \frac {\dd s}{s ^{2} + x ^{2}} \left\{\int_{- \infty} ^{\infty} | f ^{+} (s + y) - f ^{+} (y) | ^{2} \dd y \right\} = \frac {x}{\pi} \int_{- \infty} ^{\infty} \frac {\mu (f ^{+} ; s)}{s ^{2} + x ^{2}} \dd s,
\end{aligned}
\]
where $0 \leqq \mu(f^{+}; s) \leqq 4 \|f^{+}\|^{2}$ and $\mu(f^{+}; s)$ is continuous in $s$ and has value $0$ at $s = 0$. 

To prove that the right side tends to zero as $x \downarrow 0$, we take, for any $\varepsilon > 0$, a $\delta = \delta(\varepsilon) > 0$ such that $\mu(f^{+}; s) \leqq \varepsilon$ whenever $|s| \leqq \delta$. We decompose the integral 
\[
\frac {x}{\pi} \int_{- \infty} ^{\infty} \frac {\mu (f ^{+} ; s)}{s ^{2} + x ^{2}} \dd s = \frac {x}{\pi} \left\{\int_{- \infty} ^{- \delta} + \int_{- \delta} ^{\delta} + \int_{\delta} ^{\infty} \right\} = I _{1} + I _{2} + I _{3}.
\]

We have $|I_2| \leqq \varepsilon$ by \eqref{eq:VI.4.16}, and $|I_j| \leqq 4\pi^{-1} \|f^+ \|^2 \cdot \cot^{-1}(\delta/x) (j = 1, 2)$. Hence the left integral must tend to zero as $x \downarrow 0$. This proves the Theorem.
\end{proof}

\begin{remarkn}[Remark $1$] For the original Paley-Wiener theorem, see Paley-Wiener \cite{paley243}. For the one-sided Laplace transform of tempered distributions, see L. \textsc{Schwartz} \cite{schwartz295}. 
\end{remarkn}

\begin{remarkn}[Remark $2$] M. \textsc{Sato} \cite{sato284} has introduced a happy idea of defining a ``generalized function'' as the ``boundary value of an analytic function''. His idea may be explained as follows. Let $B$ be the totality of functions $\varphi(z)$ which are defined and regular in the upper-half and the lower-half planes of the complex $z$-plane, and let $R$ be the totality of functions which are regular in the whole complex $z$-plane. Then $B$ is a ring with respect to the function sum and function multiplication, and $R$ is a subring of $B$. Sato calls the residue class (mod $R$) containing $\varphi(z)$ as the ``generalized function'' $\hat{\varphi}(x)$ on the real axis $\mathbb{R}^{1}$ defined through $\varphi(z)$. The ``generalized derivative $\dd\hat{\varphi}(x)/\dd x$ of the generalized function $\hat{\varphi}(x)$ '' is naturally defined as the residue class (mod $R$) containing $\dd\varphi(z)/\dd z$. Thus the ``delta function $\delta(x)$ '' is the residue class (mod $R$) containing $-(2\pi i)^{-1}z^{-1}$. Sato's theory of ``generalized function of many variables'' admits the following very interesting topological interpretation. Let $M$ be an $n$-dimensional real analytic manifold and let $X$ be a complexification of $M$. Then the $n$-th relative cohomology group $H^{n}(X \bmod (X - M))$ with the germ of regular functions in $X$ as coefficients gives rise to a notion of the ``generalized function on $M$ ''. That is, the relative cohomology class is a natural definition of the ``generalized function''. 
\end{remarkn}

\begin{remarkn}[Remark $3$] For more detailed treatment of the Fourier transform of generalized functions, see L. \textsc{Schwartz} \cite{schwartz294} and \textsc{Gelfand}-\textsc{Šilov} \cite{gelfand76}. In the latter book, many interesting classes of basic functions other than $\mathfrak{D}(\mathbb{R}^{n})$, $\mathfrak{S}(\mathbb{R}^{n})$ and $\mathfrak{D}_M(\mathbb{R}^{n})$ are introduced to define generalized functions; the Fourier transform of the corresponding generalized functions are also discussed in \textsc{Gelfand}-\textsc{Šilov} \cite{gelfand76}. Cf. also A. \textsc{Friedman} \cite{friedman68} and L. \textsc{Hörmander} \cite{hormander118}. 
\end{remarkn}


\yossection{5}{Titchmarsh's Theorem}
\label{sec:VI.5}

\begin{theoremn}[E. C. \textsc{Titchmarsh}] Let $f(x)$ and $g(x)$ be real- or complex-valued continuous functions, $0 \leqq x < \infty$, such that
\begin{equation}\label{eq:VI.5.1}
(f * g) (x) = \int_{0} ^{x} f (x - y) g (y) \dd y = \int_{0} ^{x} g (x - y) f (y) \dd y = (g * f) (x)
\end{equation}
vanishes identically. Then either one of $f(x)$ or $g(x)$ must vanish identically. 

There is a variety of proofs of this important theorem, such as by \textsc{Titchmarsh} \cite{titchmarsh323} himself and also by \textsc{Crum} and \textsc{Dufresnoy}. The following proof is elementary in the sense that it does not appeal to the theory of functions of a complex variable. It is due to Ryll-Nardzewski \cite{ryllnardzewski279} and given in the book by J. Mikusiński \cite{mikusinski196}. 

\end{theoremn}
\begin{lemma}[\textsc{Phragmén}$N$]\label{lem:VI.5.Lemma.1} If $g(u)$ is continuous for $0 \leqq u \leqq T$, and $0 \leqq t \leqq T$, then 
\begin{equation}\label{eq:VI.5.2}
\lim\limits _{x \rightarrow \infty} \sum\limits_{k=1} ^{\infty} \frac {(- 1) ^{k-1}}{k !} \int_{0} ^{T} e ^{kx(t-u)} g (u) \dd u = \int_{0} ^{t} g (u) \dd u.
\end{equation}
\end{lemma}

\begin{proof}
 We have $\sum\limits_{k=1}^{\infty} (-1)^{k-1} (k!)^{-1} e^{kx(t-u)} = 1 - \exp(-e^{x(t-u)})$, and, for fixed $x$ and $t$, the series on the left converges uniformly in $u$ for $0 \leqq u \leqq T$. Thus the summation in \eqref{eq:VI.5.2} may be carried out under the integral sign and we obtain \eqref{eq:VI.5.2} by the Lebesgue-Fatou Lemma. 
\end{proof}

\begin{lemma}\label{lem:VI.5.Lemma.2} If $f(t)$ is continuous for $0 \leqq t \leqq T$ and $\left| \int_{0}^{T} e^{nt} f(t) \dd t \right| \leqq M$ for $n = 1; 2, \ldots$, where $M$ is a positive constant which is independent of $n$, then $f(t)$ must be $= 0$ for $0 \leqq t \leqq T$. 
\end{lemma}

\begin{proof}
 We have 
\[ \begin{aligned} & \left| \sum\limits_{k=1} ^{\infty} \int_{0} ^{T} \frac {(- 1) ^{k-1}}{k !} e ^{kn(t-u)} f (T - u) \dd u \right| \leqq \sum\limits_{k=1} ^{\infty} \frac {1}{k !} e ^{kn(t-T)} \Bigg | \int_{0} ^{T} e ^{kn(T-u)} f (T - u) \dd u \\ & \quad \leqq M (\exp (e ^{-n(T-t)}) - 1). \end{aligned}
\]

If $t < T$, the last expression tends to zero as $n \to \infty$. Hence, by Lemma~\ref{lem:VI.5.Lemma.1}, with $g(u) = f(T - u)$, we see that $\int_{0}^{t} f(T - u) \dd u = 0$ for $0 \leqq t \leqq T$. Since $f$ is continuous, it follows that $f(t) = 0$ for $0 \leqq t \leqq T$. 
\end{proof}

\begin{corollary}\label{cor:VI.5.Corollary.1} If $g(x)$ is continuous for $1 \leqq x \leqq X$ and if there exists a positive number $N$ such that $\left|\int_{1}^{X} x^n g(x) \dd x\right| \leqq N (n = 1, 2, \ldots)$, then $g(x) = 0$ for $1 \leqq x \leqq X$. 
\end{corollary}

\begin{proof}
 Putting $x = e^t$, $X = e^T$ and $xg(x) = f(t)$, we obtain, by Lemma~\ref{lem:VI.5.Lemma.2}, that $f(t) = 0$ for $0 \leqq t \leqq T$. Thus $xg(x) = 0$ for $1 \leqq x \leqq X$, and so $g(x) = 0$ for $1 \leqq x \leqq X$. 
\end{proof}

\begin{corollary}[Lerch's theorem]\label{cor:VI.5.Corollary.2} If $f(t)$ is continuous for $0 \leqq t \leqq T$ and $\int_{0}^{T} t^n f(t) \dd t = 0$ ( $n = 1, 2, \ldots$), then $f(t) = 0$ for $0 \leqq t \leqq T$. 
\end{corollary}

\begin{proof}
 Let $t_0$ be any number from the open interval $(0, T)$, and put $t = t_0x$, $T = t_0X$, $f(t) = g(x)$. We then obtain 
\[
t _{0} ^{n+1} \int_{0} ^{X} x ^{n} g (x) \dd x = 0 \quad (n = 1, 2, \dots),
\]
and so $\left|\int_{1}^{X}x^{n}g(x)\dd x\right| = \left|\int_{0}^{1}x^{n}g(x)\dd x\right|\leqq \int_{0}^{1}|g(x)|\dd x = N (n = 1,2,\ldots)$. Hence, by Corollary~\ref{cor:VI.5.Corollary.1}, we obtain $g(x) = 0$ for $1\leqq x\leqq X$, and so $f(t) = 0$ for $t_0\leqq t\leqq T$. Since $t_0$ was an arbitrary point of $(0,T)$, we must have $f(t) = 0$ for $0\leqq t\leqq T$. 
\end{proof}

\begin{proof}[Proof of Titchmarsh's theorem.]
 We shall first prove the Theorem for the special case when $f = g$, that is, if $f(t)$ is continuous and $(f * f)(t) = \int_{0}^{t} f(t - u)f(u)\dd u = 0$ for $0 \leqq t \leqq 2T$, then $f(t) = 0$ for $0 \leqq t \leqq T$. 

We have 
\[
\int_{0} ^{2T} e ^{n(2T-t)} \left(\int_{0} ^{t} f (u) f (t - u) \dd u\right) \dd t = 0,
\]
and, by the change of variables $u = T - v$, $t = 2T - v - w$, we obtain 
\[
\int_{\Delta} \int e ^{n(v+w)} f (T - v) f (T - w) \dd v \dd w = 0
\]
where $\Delta$ is the triangle $v + w \geqq 0$, $v \leqq T$, $w \leqq T$ in the $v - w$ plane. Let $\Delta'$ be the triangle $v + w \leqq 0$, $v \geqq -T$, $w \geqq -T$. Then the join $\Delta + \Delta'$ is the square $-T \leqq v$, $w \leqq T$. The above equality shows that the integral of $e^{n(v+w)} f(T - v)f(T - w)$ over $\Delta + \Delta'$ is equal to the integral over $\Delta'$. The integral over $\Delta + \Delta'$ is the product of two single integrals, and in the integral over $\Delta'$ we have $e^{n(v+w)} \leqq 1$. Thus 
\[
\begin{aligned} \left| \int_{-T} ^{T} e ^{nu} f (T - u) \dd u \right| ^{2} & = \left| \iint_ {\Delta + \Delta^{\prime}} e ^{n(v+w)} f (T - v) f (T - w) \dd v \dd w \right| \\ & \leqq \iint_ {\Delta^{\prime}} | f (T - v) f (T - w) | \dd v \dd w \leqq 2 T ^{2} \cdot A ^{2}, \end{aligned}
\]
where $A$ is the maximum of $|f(t)|$ for $0 \leqq t \leqq 2T$, and $2T^{2}$ is the area of $\Delta'$. We thus have 
\[
\left| \int_{-T} ^{T} e ^{nu} f (T - u) \dd u \right| \leqq \sqrt {2} T \cdot A,
\]
and, moreover, $\left|\int_{-T}^{0}e^{nu}f(T - u)\dd u\right| \leqq TA$. Therefore 
\[
\left| \int_{0} ^{T} e ^{nu} f (T - u) \dd u \right| = \left| \int_{-T} ^{T} - \int_{-T} ^{0} \right| \leqq (1 + \sqrt {2}) T A \quad (n = 1, 2, \dots),
\]
and so, by Lemma~\ref{lem:VI.5.Lemma.2}, $f(t) = 0$ for $0 \leqq t \leqq T$. 


We are now ready to prove Titchmarsh's theorem for the general case. 

Let $\int_0^t f(t - u)g(u)\dd u = 0$ for $0\leqq t < \infty$. Then we have, for $0\leqq t < \infty$ 
\[
\int_{0} ^{t} (t - u) f (t - u) g (u) \dd u + \int_{0} ^{t} f (t - u) u g (u) \dd u = t \int_{0} ^{t} f (t - u) g (u) \dd u = 0.
\]

This may be written as 
\[
\begin{aligned}
&\left(f _{1} * g\right) (t) + \left(f * g _{1}\right) (t) = 0 (0 \leqq t < \infty), \text{ where } f _{1} (t) = t f (t), g _{1} (t) = t g (t).
\end{aligned}
\]

Thus 
\[
[ f * \{g _{1} * (f _{1} * g + f * g _{1}) \} ] (t) = 0
\]
and so 
\[
[ (f * g) * (f _{1} * g _{1}) ] (t) + [ (f * g _{1}) * (f * g _{1}) ] (t) = 0.
\]

Thus, by $(f*g)(t) = 0$, we have $(f*g_1)* (f*g_1) = 0$, and so, by the special case proved above, $(f*g_1)(t) = 0$, that is, 
\[
\int_{0} ^{t} f (t - u) u g (u) \dd u = 0 (0 \leqq t < \infty).
\]

From this we obtain, similarly to above, 
\[
\int_{0} ^{t} f (t - u) u ^{2} g (u) \dd u = 0 (0 \leqq t < \infty).
\]

Repeating the argument, we find that 
\[
\int_{0} ^{t} f (t - u) u ^{n} g (u) \dd u = 0 \quad (0 \leqq t < \infty ; n = 1, 2, \dots).
\]

Hence, by Lerch's theorem proved above, we obtain 
\[
f (t - u) g (u) = 0 \quad \text{ for } \quad 0 \leqq u \leqq t < \infty .
\]

If a $u_0$ exists for which $g(u_0) \neq 0$, then $f(t - u_0) = 0$ for all $t \geqq u_0$, that is, $f(v) = 0$ for all $v \geqq 0$. Therefore, we have either $f(v) = 0$ for all $v \geqq 0$ or $g(v) = 0$ for all $v \geqq 0$. 
\end{proof}

\yossection{6}{Mikusiński's Operational Calculus}
\label{sec:VI.6}

In his ``Electromagnetic Theory'', London ($1899$), the physicist O. \textsc{Heaviside} inaugurated an operational calculus which he successfully applied to ordinary linear differential equations connected with electro-technical problems. In his calculus occurred certain operators whose meaning is not at all obvious. The interpretation of such operators as given by \textsc{Heaviside} himself is difficult to justify. The interpretation given by his successors is unclear with regard to its range of validity, since it is based upon the theory of Laplace transforms. The theory of convolution quotients due to J. Mikusiński provides a clear and simple basis for an operational calculus applicable to ordinary differential equations with constant coefficients, as well as to certain partial differential equations with constant coefficients, difference equations and integral equations. 

We shall give, adapted from K. \textsc{Yosida}-S. \textsc{Okamoto} \cite{yosida383}, a simplified presentation of Mikusiński's theory by introducing the notion of the ring $C_{H}$ in place of his operators (= the convolution quotients). By virtue of this ring we can derive the operational calculus without appealing to Titchmarsh's theorem given in Chapter~\ref{chap:VI}, \ref{sec:VI.5}. 

The ring $C_{H}$. We denote by $C$ the totality of complex-valued continuous functions defined on $[0,\infty)$. In this section, we shall denote such a function by $\{f(t)\}$ or simply by $f$, while $f(t)$ means the value at $t$ of this function $f$. For $f, g \in C$ and $\alpha, \beta \in \mathbb{C}^{1}$ (the complex number field), we define 
\begin{equation}\label{eq:VI.6.1}
f + g = \{f (t) + g (t) \} \quad \text{ and } \quad f g = \left\{\int_{0} ^{t} f (t - u) g (u) \dd u \right\}.
\end{equation}

Then, as proved in Chapter~\ref{chap:VI}, \ref{sec:VI.3}, we have 
\[
f g = g f, f (g h) = (f g) h \quad \text{ and } \quad f (g + h) = f g + f h.
\]

Hence $\mathcal{C}$ is a commutative ring with respect to the above addition and multiplication over the coefficients field $\mathbb{C}^{1}$. 

We shall denote by $h$ ( $l$ in J. Mikusiński \cite{mikusinski196}) the constant function $\{1\} \in \mathcal{C}$ so that we have 
\begin{equation}\label{eq:VI.6.2}
h ^{n} = \left\{\frac {t ^{n-1}}{(n - 1) !} \right\} \quad (n = 1, 2, \dots)
\end{equation}
and 
\begin{equation}\label{eq:VI.6.3}
h f = \left\{\int_{0} ^{t} f (u) \dd u \right\} \text{ for } f \in \mathscr {C},
\end{equation}
i.e., $h$ behaves as an operation of integration. Thus we have the fairly trivial 

\begin{proposition}\label{prop:VI.6.Proposition.1} For $k \in H = \{k; k = h^n (n = 1, 2, \ldots)\}$ and $f \in \mathcal{C}$, the equation $kf = 0$ implies $f = 0$, where $0$ denotes $\{0\} \in \mathcal{C}$. 
\end{proposition}

Therefore, as in the algebra, we can construct the commutative ring $\mathcal{C}_{\mathrm{H}}$ consisting of fractions of the form $f / k$, i.e., 
\begin{equation}\label{eq:VI.6.4}
\mathcal {C} _ {\mathrm{H}} = \left\{\frac {f}{k} = f / k; f \in \mathcal {C} \text{ and } k \in H \right\},
\end{equation}
where the equality of two fractions is defined by 
\begin{equation}\label{eq:VI.6.5}
\frac {f}{k} = \frac {f ^{\prime}}{k ^{\prime}} \quad \text{ iff } \quad f k ^{\prime} = f ^{\prime} k ,
\end{equation}
and the addition and multiplication of two fractions are defined through 
\begin{equation}\label{eq:VI.6.6}
\frac {f}{k} + \frac {f ^{\prime}}{k ^{\prime}} = \frac {f k ^{\prime} + f ^{\prime} k}{k k ^{\prime}} \quad \text{ and } \quad \frac {f}{k} \frac {f ^{\prime}}{k ^{\prime}} = \frac {f f ^{\prime}}{k k ^{\prime}}.
\end{equation}

By identifying $f \in \mathcal{C}$ with $kf / k \in \mathcal{C}_{\mathbf{H}}$, the ring $\mathcal{C}$ can isomorphically be embedded as a subring of $\mathcal{C}_{\mathbf{H}}$. 

We next introduce, after Mikusiński, the important notion of the operation of scalar multiplication. We define, for every $\alpha\in \mathbb{C}^{1}$, 
\begin{equation}\label{eq:VI.6.7}
[ \alpha ] = \frac {\{\alpha \}}{h} = \{\alpha \} / h \in \mathcal {C} _ {\mathrm{H}}.
\end{equation}

Then we have, for $\alpha, \beta \in \mathbb{C}^{1}$, $f \in \mathcal{C}$ and $k \in H$, 
\begin{equation}[ \alpha ] + [ \beta ] = [ \alpha + \beta ], [ \alpha ] [ \beta ] = [ \alpha \beta ], [ \alpha ] f = \alpha f = \{\alpha f (t) \}, [ \alpha ] \frac {f}{k} = \frac {\alpha f}{k}. \label{eq:VI.6.8}\end{equation}

\begin{proof}
 $[\alpha] + [\beta] = [\alpha + \beta]$ is clear. We have 
\begin{align}
[ \alpha ] [ \beta ]& = \frac {\{\alpha \}}{h} \frac {\{\beta \}}{h} = \frac {\{\alpha \beta t \}}{h ^{2}} = \frac {h \{\alpha \beta \}}{h ^{2}} = \frac {\{\alpha \beta \}}{h} = [ \alpha \beta ], \notag \\
[ \alpha ] \frac {\{f (t) \}}{k}& = \frac {\{\alpha \} \{f (t) \}}{h k} = \frac {\left\{\int_{0} ^{t} \alpha f (u) \dd u \right\}}{h k} = \frac {\{\alpha f (t) \}}{k}. \notag
\end{align}

Hence $[\alpha]$ can be identified with the complex number $\alpha$, not with $\{\alpha\}$, and we see that the effect of the multiplication by $[\alpha]$ is exactly the scalar multiplication by $\alpha$. [$1$] may be identified with the multiplicative unit I of the ring $C_{H}$, i.e., 
\begin{equation}\label{eq:VI.6.9}
I = \frac {h ^{n}}{h ^{n}} \quad (n = 1, 2, \dots).
\end{equation}

We next define 
\begin{equation}\label{eq:VI.6.10}
s = \frac {h ^{n}}{h ^{n+1}} \in \mathcal {C} _ {\mathrm{H}} (n = 0, 1, 2, \dots ; h ^{0} = I) \text{ so that } s h = h s = I
\end{equation}
\end{proof}

\begin{proposition}\label{prop:VI.6.Proposition.2} If both $f$ and its derivative $f'$ belong to $C$, then 
\begin{equation}\label{eq:VI.6.11}
f ^{\prime} = s f - f (0), \quad \text{ where } \quad f (0) = [ f (0) ],
\end{equation}
i.e., $s$ behaves as an operation of differentiation. 
\end{proposition}

\begin{proof}
 Clear from \eqref{eq:VI.6.10} and Newton's formula 
\[
h f ^{\prime} = \left\{\int_{0} ^{t} f ^{\prime} (u) \dd u \right\} = \{f (t) - f (0) \} = f - [ f (0) ] h. \]
\end{proof}

\begin{corollaryu} For $n$-times continuously differentiable function $f \in \mathcal{C}$, 
\begin{equation}\label{eq:VI.6.12}
\begin{aligned}
&f ^{(n)} = s ^{n} f - s ^{n-1} f (0) - s ^{n-2} f ^{\prime} (0) - \dots - f ^{(n-1)} (0), \text{ where } f ^{(j)} (0) = [ f ^{(j)} (0) ].
\end{aligned}
\end{equation}
\end{corollaryu}

\begin{proposition}\label{prop:VI.6.Proposition.3} For any $\alpha \in \mathbb{C}^{1}$ and for any positive integer $n$, 
\[
(s - \alpha) ^{n} = (s - [ \alpha ]) ^{n} = \left(\frac {h - [ \alpha ] h ^{2}}{h ^{2}}\right) ^{n} = \frac {(h - [ \alpha ] h ^{2}) ^{n}}{h ^{2n}} \in \mathcal {C} _ {\mathrm{H}}
\]
admits a uniquely determined multiplicative inverse in $\mathcal{C}_{\mathrm{H}}$ given by 
\begin{equation}\label{eq:VI.6.13}
\frac {I}{(s - \alpha) ^{n}} = \left\{\frac {t ^{n-1}}{(n - 1) !} e ^{\alpha t} \right\}
\end{equation}
\end{proposition}

\begin{proof}
 We have, by \eqref{eq:VI.6.11}, 
\[
(s - \alpha) \left\{\mathrm{e} ^{\alpha t} \right\} = \alpha \left\{e ^{\alpha t} \right\} + [ 1 ] - \alpha \left\{e ^{\alpha t} \right\} = I, \quad \text{ i. e., } \quad I / (s - \alpha) = \left\{e ^{\alpha t} \right\}. \]

Hence, we have $I / (s - \alpha)^2 = \left\{\int_{0}^{t} e^{\alpha(t - u)} e^{\alpha u} \dd u\right\} = \{t e^{\alpha t}\}$ and so on. 
\end{proof}

Application to the integration of the Cauchy's problem for linear ordinary differential equations with constant coefficients $\in \mathbb{C}^{1}$: 
\begin{equation}\label{eq:VI.6.14}
\begin{aligned}
\alpha_{n} y ^{(n)} + \alpha_{n-1} y ^{(n-1)} + \ldots + \alpha_{0} y &= f \in \mathcal {C} (\alpha_{n} \neq 0),
\\
y (0) &= \gamma_{0}, y ^{\prime} (0) = \gamma_{1}, \ldots , y ^{(n-1)} (0) = \gamma_{n-1}.
\end{aligned}
\end{equation}

By \eqref{eq:VI.6.12}, we rewrite \eqref{eq:VI.6.14} into 
\begin{align}
\left(\alpha_{n} s ^{n} + \alpha_{n-1} s ^{n-1} + \dots + \alpha_{0}\right) y& = f + \beta_{n-1} s ^{n-1} + \beta_{n-2} s ^{n-2} + \dots + \beta_{0}, \tag{14'}\label{eq:VI.6.14'} \\
\beta_{\nu}& = \alpha_{\nu +1} \gamma_{0} + \alpha_{\nu +2} \gamma_{1} + \dots + \alpha_{n} \gamma_{n - \nu -1} (\nu = 0, 1, 2, \dots , n - 1). \notag
\end{align}

The polynomial ring of polynomials in $s$ with coefficients $\in \mathbb{C}^{1}$ is free from zero factors, because of the fact that 
\[
\left(\delta_{n} s ^{n} + \delta_{n-1} s ^{n-1} + \dots\right) \left(\eta_{m} s ^{m} + \eta_{m-1} s ^{m-1} + \dots\right) = \delta_{n} \eta_{m} s ^{n+m} + \dots .
\]

Hence we can define rational functions in $s$: 
\begin{equation}\label{eq:VI.6.15}
F _{1} = \frac {I}{\alpha_{n} s ^{n} + \dots + \alpha_{0}} \quad \text{ and } \quad F _{2} = \frac {\beta_{n-1} s ^{n-1} + \dots + \beta_{0}}{\alpha_{n} s ^{n} + \dots + \alpha_{0}},
\end{equation}
and obtain their partial fraction expressions: 
\begin{equation}\tag{15'}\label{eq:VI.6.15'}
F _{1} = \sum\limits_{j} \sum\limits_{k=1} ^ {m _{j}} \frac {c _{jk} I}{(s - r _{j}) ^{k}} \quad \text{ and } \quad F _{2} = \sum\limits_{j} \sum\limits_{k=1} ^ {m _{j}} \frac {d _{jk} I}{(s - r _{j}) ^{k}},
\end{equation}
where both $c_{jk}$ and $d_{jk}$ belong to $\mathbb{C}^{1}$ and $r_j's$ are distinct roots of the algebraic equation $\alpha_n z^n + \alpha_{n-1} z^{n-1} + \ldots + \alpha_0 = 0$ so that $\sum\limits_{j} m_j = n$. 

By virtue of \eqref{eq:VI.6.13}, both $F_{1}$ and $F_{2}$ given in \eqref{eq:VI.6.15'} belong to $\mathcal{C} \subseteq \mathcal{C}_{\mathbb{H}}$ so that we obtain the solution of \eqref{eq:VI.6.14'} and hence of \eqref{eq:VI.6.14} as well: 
\begin{equation}\label{eq:VI.6.16}
\begin{aligned}
&\left\{y (t) \right\} = \sum\limits_{j} \sum\limits_{k=1} ^ {m _{j}} c _{jk} \left\{\frac {t ^{k-1}}{(k - 1) !} e ^ {r _{j} t} \right\} \left\{f (t) \right\} + \sum\limits_{j} \sum\limits_{k=1} ^ {m _{j}} d _{jk} \left\{\frac {t ^{k-1}}{(k - 1) !} e ^ {r _{j} t} \right\}.
\end{aligned}
\end{equation}

\begin{example}\label{ex:VI.6.Example.1} Solve the equation 
\[
x ^{\prime \prime} (t) - x ^{\prime} (t) - 6 x (t) = 2, \quad x (0) = 1, \quad x ^{\prime} (0) = 0.
\]
\end{example}

Solution. We rewrite the equation 
\[
\{x ^{\prime \prime} (t) - x ^{\prime} (t) - 6 x (t) \} = 2 / s
\]
by \eqref{eq:VI.6.12}, obtaining 
\[
s ^{2} x - s x (0) - x ^{\prime} (0) - s x + x (0) - 6 x = 2 / s
\]
and so, by substituting the initial condition, we get 
\[
(s ^{2} - s - 6) x = s - 1 + 2 / s.
\]

Hence, by \eqref{eq:VI.6.13}, 
\[
\begin{aligned}
&x = \frac {s ^{2} - s + 2}{s (s - 3) (s + 2)} = - \frac {1}{3} \frac {1}{s} + \frac {8}{1 5} \frac {1}{s - 3} + \frac {4}{5} \frac {1}{s + 2} = \left\{- \frac {1}{3} + \frac {8}{1 5} e ^{3t} + \frac {4}{5} e ^{-2t} \right\}.
\end{aligned}
\]

\begin{example}\label{ex:VI.6.Example.2} Solve the system of equations 
\[
x ^{\prime} (t) - \alpha x (t) - \beta y (t) = \beta e ^{\alpha t}, y ^{\prime} (t) + \beta x (t) - \alpha y (t) = 0,
\]
under the initial condition $x(0)=0$, $y(0)=1$. We here assume that $\alpha$ and $\beta$ are real numbers. 
\end{example}

Solution. Rewrite the equations by \eqref{eq:VI.6.11} and \eqref{eq:VI.6.13}, obtaining 
\[
s x - \alpha x - \beta y = \beta / (s - \alpha), s y - 1 + \beta x - \alpha y = 0.
\]

Hence 
\[
x = \frac {2 \beta}{(s - \alpha) ^{2} + \beta^{2}}, \quad y = \frac {(s - \alpha) ^{2} - \beta^{2}}{(s - \alpha) ((s - \alpha) ^{2} + \beta^{2})},
\]
from which we get 
\begin{gather}x = \frac {1}{i} \left\{\frac {1}{s - \alpha - i \beta} - \frac {1}{s - \alpha + i \beta} \right\} = \frac {1}{i} \left\{e ^{(\alpha + i \beta) t} - e ^{(\alpha - i \beta) t} \right\} = {} \{2 e ^{\alpha t} \sin \beta t \}, \notag
\\\leqq {} y = \frac {2 (s - \alpha)}{(s - \alpha) ^{2} + \beta^{2}} - \frac {1}{(s - \alpha)} = \frac {1}{s - \alpha - i \beta} + \frac {1}{s - \alpha + i \beta} - \frac {1}{s - \alpha} \notag
\\= \left\{e ^{(\alpha + i \beta) t} + e ^{(\alpha - i \beta) t} - e ^{\alpha t} \right\} = \left\{e ^{\alpha t} (2 \cos \beta t - 1) \right\}. \notag
\end{gather}

\begin{remarku} Suggested by \eqref{eq:VI.6.2}, we can define the fractional power integration through 
\begin{equation}\label{eq:VI.6.17}
h ^{\alpha} = h \left\{t ^{\alpha -1} / \Gamma (\alpha) \right\} / h \in \mathcal {C} _ {\mathrm{H}} (0 < \alpha < 1).
\end{equation}
\end{remarku}

Accordingly, we can define the fractional power differentiation through 
\begin{equation}\label{eq:VI.6.18}
s ^{\alpha} = h \left\{t ^{- \alpha} / \Gamma (1 - \alpha) \right\} / h ^{2} \in \mathcal {C} _ {\mathrm{H}} (0 < \alpha < 1),
\end{equation}
because we obtain 
\begin{equation}\label{eq:VI.6.19}
\mathrm{h} ^{\alpha} \mathrm{s} ^{\alpha} = \mathrm{I}
\end{equation}
by the Gamma function formula $B(1 - \alpha, \alpha) = \Gamma(1 - \alpha) \Gamma(\alpha)$. 

For further applications of the operational calculus, we refer the reader to J. Mikusiński \cite{mikusinski196} and also to A. \textsc{Erdélyi} \cite{erdelyi55}. 


\yossection{7}{Sobolev's Lemma}
\label{sec:VI.7}

A generalized function is infinitely differentiable in the sense of the distribution (see Chapter~\ref{chap:I}, \ref{sec:I.8}). So the differentiability in the generalized sense has no bearing on the ordinary differentiability. However, we have the following result which is of fundamental importance in modern treatment of partial differential equations. 

\begin{theoremn}[Sobolev's Lemma]\label{thm:VI.7.Sobolev} Let $G$ be a bounded open domain of $\mathbb{R}^{n}$. Let a function $u(x)$ belong to $W^k(G)$ for $k > 2^{-1}n + \sigma$, where $\sigma$ is an integer $\geqq 0$. We thus assume that the distributional derivatives of $u(x)$ of order up to and including $k$ all belong to $L^2(G)$. Then, for any open subset $G_1$ of $G$ such that the closure $G_1^a$ is a compact subset of $G$, there exists a function $u_1(x) \in C^\sigma(G_1)$ such that $u(x) = u_1(x)$ a.e. in $G_1$.

\end{theoremn}
\begin{proof}
 Let $\alpha(x)$ be a function $\in C_0^\infty(\mathbb{R}^{n})$ such that 
\[
G _{1} \subseteq \supp (\alpha) \subset G, 0 \leqq \alpha (x) \leqq 1 \text{ and } \alpha (x) = 1 \text{ on } G _{1}. \]

We define a function $v(x)$ defined on $\mathbb{R}^{n}$ by 
\[
v (x) = \alpha (x) u (x) \text{ for } x \in G; v (x) = 0 \text{ for } x \in \mathbb{R}^{n} - G.
\]

Then $v(x) = u(x)$ whenever $x \in G_1$. Since $v(x)$ is locally integrable over $\mathbb{R}^{n}$, it defines a distribution $\in \mathfrak{D}(\mathbb{R}^{n})'$. By the assumption that $u \in W^k(G)$, the distributional derivatives $D^s v(x) \in L^2(\mathbb{R}^{n})$ when $|s| \leqq k$. For example, the distributional derivative 
\[
\frac {\partial}{\partial x _{j}} (v) = \frac {\partial}{\partial x _{j}} (\alpha u) = \frac {\partial \alpha}{\partial x _{j}} \cdot u + \alpha \frac {\partial}{\partial x _{j}} u
\]
belongs to $L^2(\mathbb{R}^{n})$ by the fact that $u, \partial u / \partial x_j$ both belong to $L^2(G)$ and that the infinitely differentiable function $\alpha(x)$ has its support contained in a compact subset of the open domain $G$. By the Fourier transformation $v(x) \to \hat{v}(y)$, we obtain 
\[
\widehat {(D ^{s} v)} (y) = (i) ^{|s|} y _{1} ^ {s _{1}} y _{2} ^ {s _{2}} \dots y _{n} ^ {s _{n}} \cdot \hat {v} (y).
\]

Since, by Plancherel's theorem, the $L^2$-norm is preserved by the Fourier transformation, we have $(\widehat{D^s v})(y) \in L^2(\mathbb{R}^{n})$ for $|s| \leqq k$. Thus 
\begin{equation}\label{eq:VI.7.1}
\hat {v} (y) y _{1} ^ {s _{1}} y _{2} ^ {s _{2}} \dots y _{n} ^ {s _{n}} \in L ^{2} (\mathbb{R}^{n}) \quad \text{ for } \quad | s | \leqq k.
\end{equation}

In particular we have 
\begin{equation}\tag{1'}\label{eq:VI.7.1'}
\hat {v} (y) \in L ^{2} (\mathbb{R}^{n}).
\end{equation}

Let $q = (q_{1}, q_{2}, \ldots, q_{n})$ be a system of non-negative integers. Then by \eqref{eq:VI.7.1}, we can prove that 
\begin{equation}\label{eq:VI.7.2}
\hat{v}(y) y_1^{q_1} y_2^{q_2} \ldots y_n^{q_n} \text{ is integrable over } \mathbb{R}^{n} \text{ whenever } |q| + \frac{n}{2} < k.
\end{equation}
For the proof, take any positive number $C$. We have, by Schwarz' inequality, 
\[
\begin{aligned}
&\int_{| y | \leqq C} \left| \hat {v} (y) y _{1} ^ {q _{1}} y _{2} ^ {q _{2}} \dots y _{n} ^ {q _{n}} \right| \dd y \leqq \Big (\int_{| y | \leqq C} \left| y _{1} ^ {q _{1}} y _{2} ^ {q _{2}} \dots y _{n} ^ {q _{n}} \right| ^{2} \dd y \cdot \int_{| y | \leqq C} \left| \hat {\pmb v} (y) \right| ^{2} \dd y \Big) ^{1/2} < \infty , \\
&\quad \int_{|y|>C} \left| \hat {v} (y) y _{1} ^ {q _{1}} y _{2} ^ {q _{2}} \dots y _{n} ^ {q _{n}} \right| \dd y \\
&\leqq \Big (\int_{|y|>C} \left| (1 + | y | ^{2}) ^{-k/2} y _{1} ^ {q _{1}} y _{2} ^ {q _{2}} \dots y _{n} ^ {q _{n}} \right| ^{2} \dd y \cdot \int_{|y|>C} \left| \hat {v} (y) (1 + | y | ^{2}) ^{k/2} \right| ^{2} \dd y \Big) ^{1/2}.
\end{aligned}
\]

The second factor on the right of the last inequality is $< \infty$ by \eqref{eq:VI.7.1}. The first factor is finite by 
\[
\dd y = \dd y _{1} \dd y _{2} \dots \dd y _{n} = r ^{n-1} \dd r \dd \Omega_{n},
\]
where $\dd\Omega_{n}$ is the hypersurface element of the surface of the unit sphere of $\mathbb{R}^{n}$ with the origin $0$ as its centre, provided 
\begin{equation}\label{eq:VI.7.3}
2 | q | - 2 k + n - 1 < - 1, \text{ that is, if } k > \frac {n}{2} + | q |.
\end{equation}

Now, by Plancherel's theorem, we have 
\[
v (x) = \underset {h \rightarrow \infty} {\text{ l.i.m. }} (2 \pi) ^{-n/2} \int_{| y | \leqq h} \hat {v} (y) \exp (i \langle y, x \rangle) \dd y
\]
and so, as in the proof of the completeness of $L^2(\mathbb{R}^{n})$, we can choose a subsequence $\{h'\}$ of positive integers $h$ such that 
\[
v (x) = \lim\limits _ {h ^{\prime} \rightarrow \infty} (2 \pi) ^{-n/2} \int_ {| y | \leqq h ^{\prime}} \hat {v} (y) \exp (i \langle y, x \rangle) \dd y \text{ for a.e. } x \in \mathbb{R}^{n}.
\]

But, since $\hat{v}(y)$ is integrable over $\mathbb{R}^{n}$ as proved above, the right side is equal to 
\[
v _{1} (x) = (2 \pi) ^{-n/2} \int_ {\mathbb{R}^{n}} \hat {v} (y) \exp (i \langle x, y \rangle) \dd y,
\]
that is, $v(x)$ is equal to $v_{1}(x)$ for a.e. $x \in \mathbb{R}^{n}$. By ($2$), the differentiation of $v_{1}(x)$ under the integral sign is justified up to the order $\sigma$; and the result of the differentiation is continuous in $x$. By taking $u_{1}(x) = v_{1}(x)$ for $x \in G_{1}$, we have proved the Theorem. 
\end{proof}

\begin{remarku} For the original proof, see S. L. \textsc{Sobolev} \cite{sobolev306}, \cite{sobolev307} and L. \textsc{Kantorovitch}-G. \textsc{Akilov} \cite{kantorovitch133}.\yosstarfootnote{See also Supplementary Notes, p.~\pageref{notes:VI}.} 
\end{remarku}


\yossection{8}{Gårding's Inequality}
\label{sec:VI.8}

Consider a quadratic integral form defined for $C^\infty$ function $u(x) = u(x_1, x_2, \ldots, x_n)$ with compact support in a bounded domain $G$ of $\mathbb{R}^{n}$: 
\begin{equation}\label{eq:VI.8.1}
B [ u, u ] = \sum\limits_{| s |, | t | \leqq m} (c _{st} D ^{s} u, D ^{t} u) _{0},
\end{equation}
where the complex-valued coefficients $c_{st}$ are continuous on the closure $G^{a}$ of $G$, and $(u,v)_{0}$ denotes the scalar product in $L^{2}(G)$. 

Then we have 

\begin{theoremn}[L. \textsc{Garding} {\cite{garding74}}] A sufficient condition for the existence of positive constants $c$, $C$ so that the inequality
\begin{equation}\label{eq:VI.8.2}
\left\| u \right\| _{m} ^{2} \leqq c \operatorname{Re} B [ u, u ] + C \| u \| _{0} ^{2}
\end{equation}
holds for all $u \in C_0^\infty(G)$, is that, for some positive constant $c_0$, 
$\operatorname{Re} \sum\limits_{|s|,|t|=m} c_{st} \xi^s \xi^t \geqq c_0 |\xi|^2m$ for all $x \in G$ and all real vectors 
\begin{equation}\label{eq:VI.8.3}
\xi = (\xi_{1}, \xi_{2}, \dots , \xi_{m}).
\end{equation}

\end{theoremn}
\begin{remarku} The inequality \eqref{eq:VI.8.2} is called Gårding's inequality. If condition \eqref{eq:VI.8.3} is satisfied, then the differential operator 
\begin{equation}\label{eq:VI.8.4}
L = \sum\limits_{| s | = | t | \leqq m} (- 1) ^{|t|} D ^{t} c _{st} D ^{s}
\end{equation}
\end{remarku}

is said to be strongly elliptic in $G$, assuming that $c_{st}$ is $C^m$ on $G^a$. 

\begin{proof}
 We first prove that, for every $\varepsilon > 0$, there is a constant $C(\varepsilon) > 0$ such that for every $C_0^\infty(G)$ function $u$, 
\begin{equation}\label{eq:VI.8.5}
\left\| u \right\| _{m-1} ^{2} \leqq \varepsilon \left\| u \right\| _{m} ^{2} + C (\varepsilon) \left\| u \right\| _{0} ^{2}. \end{equation}

To this purpose, we consider $u$ to belong to $C_0^\infty (\mathbb{R}^{n})$, defining its value as $0$ outside $G$. By the Fourier transformation we have 
\[
\left\| D ^{s} u \right\| _{0} ^{2} = \left\| (\widehat {D ^{s} u}) \right\| _{0} ^{2} = \int_ {\mathbb{R}^{n}} \left| \prod\limits_{j=1} ^{n} y _{j} ^{sj} \hat {\boldsymbol {u}} (y) \right| ^{2} \dd y,
\]
in virtue of Plancherel's theorem. Thus \eqref{eq:VI.8.5} is a consequence from the fact that 
\[
\left(\sum\limits_{| s | \leqq m -1} \prod\limits_{j=1} ^{n} y _{j} ^ {2 s _{j}}\right) / \left(C + \sum\limits_{| t | \leqq m} \prod\limits_{j=1} ^{n} y _{j} ^ {2 t _{j}}\right) \left(| s | = \sum\limits_{j=1} ^{n} s _{j}, | t | = \sum\limits_{j=1} ^{n} t _{j}\right)
\]
tends to zero uniformly with respect to $y = (y_{1}, y_{2}, \ldots, y_{n})$ as $C \uparrow \infty$. 

Suppose that the coefficients $c_{st}$ are constant and vanish unless $|s| = |t| = m$. By the Fourier transformation $u(x) \to \hat{u}(\xi)$ and Plancherel's theorem, we have, by \eqref{eq:VI.8.3}, 
\[
\begin{aligned}
\operatorname{Re} B [ u, u ]& = \operatorname{Re} \int \sum\limits_{s,t} c _{st} \xi^{s} \xi^{t} | \hat {u} (\xi) | ^{2} \dd \xi
\\
\qquad \qquad \qquad& \geqq \int c _{0} | \xi | ^{2m} | \hat {u} (\xi) | ^{2} \dd \xi \geqq c _{1} (\| u \| _{m} ^{2} - \| u \| _{m-1} ^{2}),
\end{aligned}
\]
where $c_{1} > 0$ is a constant which is independent of $u$. Hence, by \eqref{eq:VI.8.5}, we see that \eqref{eq:VI.8.2} is true for our special case. 

We next consider the case of variable coefficients $c_{st}$. Suppose, first, that the support of $u$ is sufficiently small and contained, say, in a small sphere about the origin. By the preceding case, we have, with a constant $c_{0}^{\prime}>0$ which is independent of $u$, 
\[
\begin{aligned}
c _{0} ^{\prime} \| u \| _{m} ^{2}& \leqq \operatorname{Re} B [ u, u ] + \operatorname{Re} \sum\limits_{|s|=|t|=m} \int (c _{st} (0) - c _{st} (x)) D ^{s} u \cdot D ^{t} \bar {u} \dd x
\\
\quad - \operatorname{Re} \sum\limits_{|s|+|t|<2m}& \int c _{st} (x) D ^{s} u \cdot D ^{t} \bar {u} \dd x + C (\varepsilon) \left\| u \right\| _{0} ^{2}.
\end{aligned}
\]

If the support of $u$ is so small that $c_{st}$ has very small oscillation there, we see that the second term on the right may be bounded by $2^{-1}c_0'\|u\|_m^2$. The third term on the right is bounded by constant times $\|u\|_m\cdot \|u\|_{m-1}$. Hence we find that, constants denoting positive constants, 
\[
2 ^{-1} c _{0} ^{\prime} \| u \| _{m} ^{2} \leqq \operatorname{Re} B [ u, u ] + \text{ constant } \| u \| _{m} \cdot \| u \| _{m-1} + C (\varepsilon) \| u \| _{0} ^{2}.
\]

Thus, by 
\begin{equation}\label{eq:VI.8.6}
\begin{aligned}
2 | \alpha | \cdot | \beta |& \leqq \varepsilon | \alpha | ^{2} + \varepsilon^{-1} | \beta | ^{2} \text{ which is valid for every }
\\
\varepsilon& > 0,
\end{aligned}
\end{equation}
we obtain $\|u\|_{m}^{2}\leqq constant\cdot \operatorname{Re} B[u,u]+\text{ constant }\cdot\|u\|_{m-1}^{2}+C(\varepsilon)\cdot\|u\|_{0}^{2}$, and so, by \eqref{eq:VI.8.5}, we obtain \eqref{eq:VI.8.2}. 


Next we consider the general case. Construct a partition of unity in $G$: 
\[
1 = \sum\limits_{j=1} ^{N} \omega_{j} ^{2}, \omega_{j} \in C _{0} ^{\infty} (G) \quad \text{ and } \quad \omega_{j} (x) \geqq 0 \text{ in } G,
\]
such that the support of each $\omega_{j}$ may be taken as small as we please. Then, by Leibniz' rule of differentiation of the product of functions, Schwarz' inequality and the estimate of the case obtained above, we have 
\[
\begin{aligned}
\operatorname{Re} B [ u, u ] &= \operatorname{Re} \sum\limits_{s,t} \int c _{st} D ^{s} u D ^{t} \overline {{u}} \dd x = \operatorname{Re} \sum\limits_{s,t} \sum\limits_{j} \int \omega_{j} ^{2} c _{st} D ^{s} u D ^{t} \overline {{u}} \dd x
\\
&= \operatorname{Re} \int \sum\limits_{j} \sum\limits_{s,t} c _{st} D ^{s} (\omega_{j} u) D ^{t} \overline {{(\omega_{j} u)}} \dd x + 0 (\| u \| _{m} \| u \| _{m-1})
\\
& \geqq \sum\limits_{j} \text{ constant } (\| \omega_{j} u \| _{m} ^{2} - \| \omega_{j} u \| _{m-1} ^{2}) + 0 (\| u \| _{m} \cdot \| u \| _{m-1})
\\
& \geqq \text{ constant } \| u \| _{m} ^{2} + 0 (\| u \| _{m} \cdot \| u \| _{m-1}).
\end{aligned}
\]

Thus we obtain \eqref{eq:VI.8.2} by \eqref{eq:VI.8.5}. We remark that the constants $c$, $C$ in \eqref{eq:VI.8.2} depend upon $c_{0}$, $c_{st}$ and the domain $G$.
\end{proof}

\yossection{9}{Friedrichs' Theorem}
\label{sec:VI.9}

Let 
\begin{equation}\label{eq:VI.9.1}
L = \sum\limits_{| s |, | t | \leqq m} D ^{s} c _{st} (x) D ^{t}
\end{equation}
be strongly elliptic with real $C^{\infty}$ coefficients $c_{st}(x)$ in a bounded open domain of $\mathbb{R}^{n}$. For a given locally square integrable function $f(x)$ in $G$, a locally square integrable function $u(x)$ in $G$ is called a weak solution of 
\begin{equation}\label{eq:VI.9.2}
L u = f,
\end{equation}
if we have 
\begin{equation}\label{eq:VI.9.3}
(u, L ^{*} \varphi) _{0} = (f, \varphi) _{0}, L ^{*} = \sum\limits_{| s |, | t | \leqq m} (- 1) ^{|s|+|t|} D ^{t} c _{st} (x) D ^{s},
\end{equation}
for every $\varphi \in C_0^\infty(G)$. Here $(f, g)_0$ denotes the scalar product of the Hilbert space $L^2(G)$. Thus a weak solution $u$ of \eqref{eq:VI.9.2} is a solution in the sense of the distribution. Concerning the differentiability of the weak solution $u$, we have the following fundamental result: 

\begin{theoremn}[K. \textsc{Friedrichs} {\cite{friedrichs69}}] Any weak solution $u$ of \eqref{eq:VI.9.2} has square integrable (distributional) derivatives of order up to $(2m + p)$ in the domain $G_{1} \subseteq G$ where $f$ has square integrable (distributional) derivatives of order up to $p$. In other words, any weak solution $u$ of \eqref{eq:VI.9.2} belongs to $W^{p+2m}(G_1)$ whenever $f$ belongs to $W^{p}(G_{1})$.

\end{theoremn}
\begin{corollaryu} If $p = \infty$, then, by Sobolev's lemma, there exists a function $u_0(x) \in C^\infty(G_1)$ such that $u(x) = u_0(x)$ for a.e. $x \in G_1$. Thus, after a correction on a set of measure zero, any weak solution $u(x)$ of \eqref{eq:VI.9.2} is $C^\infty$ in the domain $\subseteq G$ where $f(x)$ is $C^\infty$; the corrected solution is hence a genuine solution of the differential equation \eqref{eq:VI.9.2} in the domain where $f(x)$ is $C^\infty$. 
\end{corollaryu}

\begin{remarku} When $L = \Delta$, the Laplacian, the above Corollary is Weyl's Lemma (see Chapter~\ref{chap:II}, \ref{sec:II.7}). There is extensive literature concerning the extensions of Weyl's Lemma to general elliptic operator $L$; such extensions are sometimes called the Weyl-Schwartz theorem. Among an abundant literature, we refer to the papers by P. \textsc{Lax} \cite{lewy172}, L. \textsc{Nirenberg} \cite{nirenberg236} and L. \textsc{Nirenberg} \cite{nirenberg237}. The proof below is due to the present author (unpublished). A similar proof was given by L. \textsc{Bers} \cite{bers16}. It is to be noted that a non-differentiable, locally integrable function $f(x)$ is a distribution solution of the hyperbolic equation 
\[
\frac {\partial f}{\partial x \partial y} = 0,
\]
\end{remarku}

as may be seen from 
\[
0 = \int_{- \infty} ^{\infty} \left\{\int_{- \infty} ^{\infty} f (x) \frac {\partial \varphi (x , y)}{\partial y \partial x} \dd y \right\} \dd x (\varphi (x, y) \in C _{0} ^{\infty} (\mathbb{R}^{2})).
\]

\begin{proof}[Proof of the Theorem.]
We shall be concerned only with real-valued functions. Replacing, if necessary, $L$ by $I + \alpha L$ with a certain constant $\alpha \neq 0$, we may assume that the strongly elliptic operator $L$ itself satisfies Gårding's inequality
\begin{equation}\label{eq:VI.9.4}
\begin{aligned}
&(\varphi , L ^{*} \varphi) _{0} \geqq \delta \| \varphi \| _{m} ^{2} (\delta > 0), | (\varphi , L ^{*} \psi) _{0} | \leqq \gamma \| \varphi \| _{m} \| \psi \| _{m} (\gamma > 0), \text{ whenever } \varphi , \psi \in C _{0} ^{\infty} (G).
\end{aligned}
\end{equation}

The latter inequality is proved by partial integration. We assume here that each of the derivatives of the coefficients $c_{st}(x)$ up to order $m$ is bounded in $G$, so that the constants $\delta$ and $\gamma$ are independent of the test functions $\psi$, $\varphi \in C_{0}^{\infty}(G)$. 

Suppose that $G_{1}$ is a periodic parallelogram 
\begin{equation}\label{eq:VI.9.5}
0 \leqq x _{j} \leqq 2 \pi \quad (j = 1, 2, \dots , n),
\end{equation}
and that the coefficients of $L$ and $f$ periodic with the period $2\pi$ in each $x_{j}$. Under such assumptions, we are to deal with functions $\varphi(x)$ defined on a compact space without boundary, the $n$-dimensional torus $G_{1}$ given by \eqref{eq:VI.9.5}, and the distribution $\in C^{\infty}(G_{1})'$ associated with the space of test functions $\varphi \in C^{\infty}(G_{1})$ consisting of $C^{\infty}$ functions $\varphi(x) = \varphi(x_{1}, x_{2}, \ldots, x_{n})$ periodic with period $2\pi$ in each of the variables $x_{j}$. It is to be noted that, since $G_{1}$ is without boundary, we need not restrict the supports of our test functions $\varphi(x)$. 

The condition $v \in W^{q}(G_{1})$ thus means that the Fourier coefficients $v_{k}$ of $v(x)$ defined by the Fourier expansion 
\begin{equation}v (x) \sim \sum\limits_{k} v _{k} \exp (i k \cdot x) \left(k = \left(k _{1}, \dots , k _{n}\right), x = \left(x _{1}, \dots , x _{n}\right), \quad \text{ and } \quad k \cdot x = \sum\limits_{j=1} ^{n} k _{j} x _{j}\right) \label{eq:VI.9.6}\end{equation}
satisfy 
\begin{equation}\label{eq:VI.9.7}
\sum\limits_{k} \left| v _{k} \right| ^{2} (1 + | k | ^{2}) ^{q} < \infty \quad \left(| k | ^{2} = \sum\limits_{j=1} ^{n} k _{j} ^{2}\right).
\end{equation}

For, by partial integration, the Fourier coefficients of the distributional derivative $D^s v$ satisfy 
\begin{equation}\label{eq:VI.9.8}
\begin{aligned}
&(D ^{s} v (x), \exp (i k \cdot x)) _{0}  = (- 1) ^{|s|} (v (x), D ^{s} \exp (i k \cdot x)) _{0} = (- i) ^{|s|} \prod\limits_{j=1} ^{n} k _{j} ^{sj} v _{k}, s = (s _{1}, s _{2}, \dots , s _{n}),
\end{aligned}
\end{equation}
and so, by Parseval's relation for the Fourier coefficients of $D^q v \in L^2(G_1)$, we obtain \eqref{eq:VI.9.7}. 

It is convenient to introduce the space $W^q(G_1)$ with integer $q \geqq 0$, by saying that a sequence $\{w_k; k = (k_1, k_2, \ldots, k_n)\}$ of complex numbers $w_k$ with $w_k = \overline{w}_{-k}$ belongs to $W^q(G_1)$ if \eqref{eq:VI.9.7} holds. This space $W^q(G_1)$ is normed by $\|\{w_k\}\|_q = \left(\sum\limits_k |w_k|^2 (1 + |k|^2)^q\right)^{1/2}$. By virtue of the Parseval relation with respect to the complete orthonormal system $\{(2\pi)^{-n/2} \exp(ik \cdot x)\}$ of $L^2(G_1)$, we see that, when $q \geqq 0$, the norm $\|v\|_q = \left(\sum\limits_{|s| \leqq q} \int_{G_1} |D^s v(x)|^2 \dd x\right)^{1/2}$ is equivalent to the norm $\|\{v_k\}\|_q$, where $v(x) \sim \sum\limits_k v_k \exp(ik \cdot x)$. 

The above proof of \eqref{eq:VI.9.7} shows that, if $f \in W^{p}(G_{1})$, then $D^{s}f \in W^{p-|s|}(G_{1})$, and $\varphi f \in W^{p}(G_{1})$ for $\varphi \in C^{\infty}(G_{1})$. Hence 
if $f \in W^{p}(G_{1})$, then, for any differential operator $N$ of 
order $q$ with $C^\infty (G_1)$ coefficients, $Nf\in W^{p-|q|}(G_1)$. 


To prove the Theorem for our periodic case, we first show that we may assume that the weak solution $u \in L^{2}(G_{1}) = W^{0}(G_{1})$ of \eqref{eq:VI.9.2} belongs to $W^{m}(G_{1})$. This is justified as follows. Let 
\[
u (x) \sim \sum\limits_{k} u _{k} \exp (i k \cdot x), v (x) \sim \sum\limits_{k} u _{k} (1 + | k | ^{2}) ^{-m} \exp (i k \cdot x).
\]

Then it is easy to see that $v(x) \in W^{2m}(G_1)$ and that $v$ is a weak solution of $(I - \Delta)^m v = u$, where $\Delta$ denotes the Laplacian $\sum\limits_{j=1}^{n} \partial^2 / \partial x_j^2$. Hence $v$ is a weak solution of the strongly elliptic equation of order 4m: 
\begin{equation}\tag{2'}\label{eq:VI.9.2'}
L (I - \Delta) ^{m} v = f.
\end{equation}

If we can show that this weak solution $v \in W^{2m}(G_1)$ actually belongs to $W^{4m+p}(G_1)$, then, by ($8$), $u =: (I - \Delta)^m v$ belongs to $W^{4m+p-2m}(G_1) =$ 
$W^{p+2m}(G_{1})$. Therefore, without losing the generality, we may assume that the weak solution $u$ of \eqref{eq:VI.9.2} belongs to $W^{m}(G_{1})$, where $m$ is half of the order $2m$ of $L$. 

Next, by Gårding's inequality \eqref{eq:VI.9.4} for $L$ and that for $(I - \Delta)^m$, we may apply the Lax-Milgram theorem (Chapter~\ref{chap:III}, \ref{sec:III.7}) to the following effect. The bilinear forms on $C^\infty(G_1)$: 
\begin{equation}\label{eq:VI.9.9}
(\varphi , \psi) ^{\prime} = (\varphi , L ^{*} \psi) _{0} \text{ and } (\varphi , \psi) ^{\prime \prime} = (\varphi , (I - \Delta) ^{m} \psi) _{0}
\end{equation}
can both be extended to be continuous bilinear forms on $W^{m}(G_{1})$ such that there exist one-one bicontinuous linear mappings $T'$, $T''$ on $W^{m}(G_{1})$ onto $W^{m}(G_{1})$ satisfying the conditions 
\[
(T ^{\prime} \varphi , \psi) ^{\prime} = (\varphi , \psi) _{m}, (T ^{\prime \prime} \varphi , \psi) ^{\prime \prime} = (\varphi , \psi) _{m} \quad \text{ for } \quad \varphi , \psi \in W ^{m} (G _{1}).
\]

Therefore, there exists a one-one bicontinuous linear map $T_{m} = T^{\prime \prime}(T^{\prime})^{-1}$ on $W^{m}(G_{1})$ onto $W^{m}(G_{1})$ such that 
\begin{equation}\label{eq:VI.9.10}
(\varphi , \psi) ^{\prime} = (T _{m} \varphi , \psi) ^{\prime \prime} \text{ whenever } \varphi , \psi \in W ^{m} (G _{1}).
\end{equation}

We can show that 
\begin{equation}
\label{eq:VI.9.11}
\begin{minipage}{0.92\linewidth}
\everymath{}\let\limits\nolimits
for any $j \geqq 1, T _{m}$ maps $W ^{m+j} (G _{1})$ onto $W ^{m+j} (G _{1})$ in a one-one and bicontinuous way.
\end{minipage}
\end{equation}


In fact, we have 
\[
(\varphi , L ^{*} (I - \Delta) ^{j} \psi) _{0} = (T _{m} \varphi , (I - \Delta) ^{m+j} \psi) _{0} \quad \text{ for } \quad \varphi , \psi \in C ^{\infty} (G _{1}).
\]

On the other hand, there exists, by the Lax-Milgram theorem as applied to the strongly elliptic operators $(I-\Delta)^{j}L$ and $(I-\Delta)^{m+j}$, a one-one bicontinuous linear map $T_{m+j}$ of $W^{m+j}(G_{1})$ onto $W^{m+j}(G_{1})$ such that 
\[
(\varphi , L ^{*} (I - \Delta) ^{j} \psi) _{0} = (T _{m+j} \varphi , (I - \Delta) ^{m+j} \psi) _{0} \quad \text{ for } \quad \varphi , \psi \in C ^{\infty} (G _{1}).
\]

Therefore, the function $w = (T_{m+j} - T_m)\varphi$ is, for any $\varphi \in C^\infty(G_1)$, a weak solution of $(I - \Delta)^{m+j}w = 0$. But, such a $w(x)$ is identically zero. For, the Fourier coefficient $w_k$ of $w(x)$ satisfies 
\[
\begin{aligned}
0 &= ((I - \Delta) ^{m+j} w (x), \exp (i k \cdot x)) _{0} = (w (x), (I - \Delta) ^{m+j} \exp (i k \cdot x)) _{0}
\\
&= (1 + | k | ^{2}) ^{m+j} (w (x), \exp (i k \cdot x)) _{0} = (1 + | k | ^{2}) ^{m+j} w _{k},
\end{aligned}
\]
and so $w_{k} = 0$ for all $k$. Thus $(T_{m+j} - T_{m})$ is $0$ on $C^\infty(G_1)$. The space $C^\infty(G_1)$ is dense in $W^{m+j}(G_1) \subseteq W^m(G_1)$, since trigonometric polynomials $\sum\limits_{|k| < \infty} w_k \exp(ik \cdot x)$ are dense in the space $W^{m+j}(G_1)$. Hence $T_{m+j} = T_m$ on $W^{m+j}(G_1)$. 

We are now ready to prove the differentiability theorem for our periodic case. We have, for $\psi\in C^{\infty}(G_{1})$, 
\[
(f, \psi) _{0} = (u, L ^{*} \psi) _{0} = (u, \psi) ^{\prime} = (T _{m} u, \psi) ^{\prime \prime} = (T _{m} u, (I - \Delta) ^{m} \psi) _{0}.
\]

Hence, for 
\[
T _{m} u \sim \sum\limits_{k} c _{k} \exp (i k \cdot x), \psi (x) \sim \sum\limits_{k} \psi_{k} \exp (i k \cdot x),
\]
we obtain, by Parseval's relation, 
\[
(T _{m} u, (I - \Delta) ^{m} \psi) _{0} = \sum\limits_{k} c _{k} (1 + | k | ^{2}) ^{m} \overline {{{\psi}}} _{k} = \sum\limits_{k} f _{k} \overline {{{\psi}}} _{k}.
\]

By the arbitrariness of $\psi \in C^{\infty}(G_{1})$, we have $c_k(1 + |k|^2)^m = f_k$, and so, by $f \in W^{p}(G_{1})$, we must have $T_m u \in W^{p+2m}(G_1)$. Hence, by ($11$), $u \in W^{p+2m}(G_1)$. It is to be noted that the above conclusion $u \in W^{p+2m}(G_1)$ is true even in the case $0 \geqq p \geqq (1-m)$, i.e., the case $\{f_k\} \in W^{p}(G)$ with $0 \geqq p \geqq (1-m)$. For, by $p + 2m \geqq m + 1$, we can apply ($11$). 

We finally shall prove, for the general non-periodic case, our differentiability theorem. The following argument is due to P. \textsc{Lax} \cite{lewy172}. 

We want to prove, for the general non-periodic case, the differentiability theorem in a vicinity of a point $x^{0}$ of $G$. Let $\beta(x) \in C_{0}^{\infty}(G)$ be identically one in a vicinity of $x^{0}$. Denote $\beta u$ by $u'$. $u'$ is a weak solution of 
\begin{equation}\label{eq:VI.9.12}
L u ^{\prime} = \beta f + N u,
\end{equation}
where $N$ is a differential operator of order at most equal to $(2m-1)$ whose coefficients are, together with $\beta(x)$, zero outside some vicinity $V$ of $x^{0}$, and the operator $N$ is to be applied in the sense of the distribution. We denote the distribution $\beta f + Nu$ by $f'$. 

Let the periodic parallelogram $G_{1}$ contain $V$, and imagine the coefficients of $L$ so altered inside $G_{1}$ but outside $V$ that they become periodic without losing their differentiability and ellipticity properties. Denote the so altered $L$ by $L'$. Thus $u'$ is a weak solution in $G$ of 
\begin{equation}\label{eq:VI.9.13}
L ^{\prime} u ^{\prime} = f ^{\prime}, \text{ where } f ^{\prime} = \beta f + N u.
\end{equation}

We can thus apply the result obtained above for the periodic case to our weak solution $u'$. We may assume that the weak solution $u'$ belongs to $W^{m}(G_{1})$. Thus, since $N$ is of order $\leqq(2m-1)$ and with coefficients vanishing onside $V$, $f'=\beta f+Nu$ must satisfy, by ($8$), 
$f^{\prime}\in W^{p^{\prime}}(G_{1})$ with $p^{\prime} = \min (p,m - (2m - 1)) = \min (p,1 - m)\geqq 1 - m.$ 

Therefore, the weak solution $u'$ of \eqref{eq:VI.9.13} must satisfy 
\[
\begin{aligned}
&u ^{\prime} \in W ^ {p ^{\prime \prime}} (G _{1})  \text{ with } p ^{\prime \prime} = \min (p + 2 m, 1 - m + 2 m) = \min (p + 2 m, m + 1).
\end{aligned}
\]

Hence, in a certain vicinity of $x^0$, $u$ has square integrable distributional derivatives up to order $p''$ which is $\geqq (m + 1)$. Thus $f' = \beta f + Nu$ has, in a vicinity of $x^0$, square integrable distributional derivatives of order up to 
\[
p ^{\prime \prime \prime} = \min (p, p ^{\prime \prime} - (2 m - 1)) \geqq \min (p, 2 - m).
\]

Thus again applying the result already obtained, we see that $u'$ has, in a certain vicinity of $x^{0}$, square integrable distributional derivatives up to order 
\[
p ^{(4)} = \min (p + 2 m, 2 - m + 2 m) = \min (p + 2 m, m + 2).
\]

Repeating the process, we see that $u$ has, in a vicinity of $x^{0}$, square integrable distributional derivatives up to order $p + 2m$. 
\end{proof}


\yossection{10}{The Malgrange-Ehrenpreis Theorem}
\label{sec:VI.10}

There is a striking difference between ordinary differential equations and partial differential equations. A classical result of \textsc{Peano} states that the ordinary differential equation $\dd y/\dd x = f(x, y)$ has a solution under a single condition of the continuity of the function $f$. This result has also been extended to equations of higher order or to systems of equations. However, for partial differential equations, the situation is entirely different. H. \textsc{Lewy} \cite{leray176} constructed in $1957$ the equation 
\[
- i \frac {\partial u}{\partial x _{1}} + \frac {\partial u}{\partial x _{2}} - 2 (x _{1} + i x _{2}) \frac {\partial u}{\partial x _{3}} = f (x _{3}),
\]
which has no solution at all if $f$ is not analytic, even if $f$ is $C^{\infty}$. Lewy's example led L. \textsc{Hörmander} \cite{hormander115} to develop a systematic method of constructing linear partial differential equations without solutions. Thus it is important to single out classes of linear partial differential equations with solutions. 

Let $P(\xi)$ be a polynomial in $\xi_{1}, \xi_{2}, \ldots, \xi_{n}$, and let $P(D)$ be the linear differential operator obtained by replacing $\xi_{j}$ by $D_{j} = i^{-1}\partial/\partial x_{j}$. $P(D)$ may be written as 
\begin{equation}P (D) = \sum\limits_{\alpha} c _{\alpha} D _{\alpha}, \text{ where, for } \alpha = (\alpha , \alpha_{2}, \dots , \alpha_{n}), D _{\alpha} = D _{1} ^ {\alpha_{1}} D _{2} ^ {\alpha_{2}} \dots D _{n} ^ {\alpha_{n}}. \notag\end{equation}

\begin{definition}\label{def:VI.10.Definition.1} By a fundamental solution of $P(D)$ we mean a distribution $E$ in $\mathbb{R}^{n}$ such that 
\[
P (D) E = \delta = T _{\delta}.
\]
\end{definition}

The importance of the notion of the fundamental solution is due to the fact that 
\[
u = E * f, \text{ where } f \in C _{0} ^{\infty} (\mathbb{R}^{n}),
\]
gives a solution of the equation 
\[
P (D) u = f.
\]

In fact, the differentiation rule \eqref{eq:VI.3.10} in Chapter~\ref{chap:VI}, \ref{sec:VI.3} implies that $P(D)u = (P(D)E)\star f = \delta \star f = f$. 

\begin{examplen} Let $P(D)$ be the Laplacian $\Delta = \sum\limits_{j=1}^{n} \frac{\partial^2}{\partial x_j^2}$ in $\mathbb{R}^{n}$ with $n \geqq 3$. Then the distribution
$E = T_{g}$, where $g(x) = \frac{1}{(2 - n)S_n} |x|^{2-n}$ and $S_{n} =$ the area of the surface of the unit sphere of $\mathbb{R}^{n}$, 
is a fundamental solution of $\Delta$. 

\end{examplen}
\begin{proof}
 We have, in the polar coordinates, $\dd x = |x|^{n-1} d |x| dS_n$, and so the function $g(x)$ is locally integrable in $\mathbb{R}^{n}$. Hence 
\[
\Delta T _ {| x | ^{2-n}} (\varphi) = \lim\limits _{\varepsilon \downarrow 0} \int_{| x | \geqq \varepsilon} | x | ^{2-n} \cdot \Delta \varphi \dd x, \varphi \in \mathfrak {D} (\mathbb{R}^{n}). \]

Let us take two positive numbers $\varepsilon$ and $R (> \varepsilon)$ such that the $\supp(\varphi)$ is contained in the interior of the sphere $|x| \leqq R$. Consider the domain $G: \varepsilon \leqq |x| \leqq R$ of $\mathbb{R}^{n}$ and apply Green's integral theorem, obtaining 
\[
\int_{G} \left(| x | ^{2-n} \cdot \Delta \varphi - \Delta | x | ^{2-n} \cdot \varphi\right) \dd x = \int_{\partial G} \left(| x | ^{2-n} \cdot \frac {\partial \varphi}{\partial v} - \frac {\partial | x | ^{2-n}}{\partial v} \cdot \varphi\right) \dd S
\]
where $S = \partial G$ is the boundary surfaces given by $|x| = \varepsilon$ and $|x| = R$, and $\nu$ denotes the outwards normal to $S$. Since $\varphi$ vanishes around $|x| = R$, we have, remembering that $\Delta |x|^{2-n} = 0$ for $x \neq 0$ and that $-\frac{\partial}{\partial\nu} = \frac{\partial}{\partial|x|}$ at the points of the inner boundary surface $|x| = \varepsilon$, 
\[
\int_ {\mathbb{R}^{n}} | x | ^{2-n} \Delta \varphi \dd x = - \int_{| x | = \varepsilon} \varepsilon^{2-n} \frac {\partial \varphi}{\partial | x |} \dd S + \int_{| x | = \varepsilon} (2 - n) \varepsilon^{1-n} \varphi \dd S.
\]

When $\varepsilon \downarrow 0$, the expression $\partial \varphi / \partial |x| = \sum\limits_{j=1}^{n} (x_j / |x|) \cdot \partial \varphi / \partial x_j$ is bounded and the area of the boundary surface $|x| = \varepsilon$ is $S_n \varepsilon^{n-1}$. Consequently the first term on the right tends to zero as $\varepsilon \downarrow 0$. By the continuity of $\varphi$ and a similar argument to above, the second term on the right tends, as $\varepsilon \downarrow 0$, to $(2 - n) S_n \cdot \varphi(0)$. Thus $T_g$ is a fundamental solution of $\Delta$. 
\end{proof}

The existence of a fundamental solution for every linear partial differential equation with constant coefficients was proved independently by B. \textsc{Malgrange} \cite{malgrange186} and L. \textsc{Ehrenpreis} \cite{ehrenpreis53} around $1954$---$55$. The exposition of the result given below is due to L. \textsc{Hörmander} \cite{hormander116}. 

\begin{definition}\label{def:VI.10.Definition.2} Set 
\begin{equation}\label{eq:VI.10.1}
\begin{aligned}
&\tilde {P} (\xi) = \left(\sum\limits_{| \alpha | \geqq 0} | P ^{(\alpha)} (\xi) | ^{2}\right) ^{1/2} \quad \text{ where } \quad P ^{(\alpha)} (\xi) = D _{\xi} ^{\alpha} P (\xi), D _{\xi} ^{\alpha} = \frac {\partial^ {\alpha_{1} + \alpha_{2} + \cdots + \alpha_{n}}}{\partial_ {\xi_{1}} ^ {\alpha_{1}} \partial_ {\xi_{2}} ^ {\alpha_{2}} \ldots \partial_ {\xi_{n}} ^ {\alpha_{n}}}.
\end{aligned}
\end{equation}
\end{definition}

We say that a differential operator $Q(D)$ with constant coefficients is weaker than $P(D)$ if 
\begin{equation}\label{eq:VI.10.2}
\tilde {Q} (\xi) \leqq C \tilde {P} (\xi), \xi \in \mathbb{R}^{n}, C \text{ being a positive constant }.
\end{equation}

\begin{theorem}\label{thm:VI.10.Theorem.1} If $\Omega$ is a bounded domain of $\mathbb{R}^{n}$ and $f \in L^2(\Omega)$, then there exists a solution $u$ of $P(D) u = f$ in $\Omega$ such that $Q(D) u \in L^2(\Omega)$ for all $Q$ weaker than $P$. Here the differential operators $P(D)$ and $Q(D)$ are to be applied in the sense of the theory of distributions. 
\end{theorem}

The proof is based on 

\begin{theorem}\label{thm:VI.10.Theorem.2} For every $\varepsilon > 0$, there is a fundamental solution $E$ of $P(D)$ such that, with a constant $C$ independent of $u$, 
\begin{equation}\label{eq:VI.10.3}
\left| (E * u) (0) \right| \leqq C \sup\limits _{| \eta | \leqq \varepsilon} \int_ {R _{n}} \left| \hat {u} (\xi + i \eta) \right| / \tilde {P} (\xi) \dd \xi , u \in C _{0} ^{\infty} (\mathbb{R}^{n}).
\end{equation}
\end{theorem}

Here $\hat{u}$ is the Fourier-Laplace transform of $u$: 
\[
\hat {u} (\zeta) = (2 \pi) ^{-n/2} \int_ {\mathbb{R}^{n}} e ^{- i \langle x, \zeta \rangle} u (x) \dd x, \zeta = \xi + i \eta ,
\]
and the finiteness of the right side of \eqref{eq:VI.10.3} is assured by the Paley-Wiener theorem in Chapter~\ref{chap:VI}, \ref{sec:VI.4}. 

Deduction of Theorem~\ref{thm:VI.10.Theorem.1} from Theorem~\ref{thm:VI.10.Theorem.2}. Replace $u$ in \eqref{eq:VI.10.3} by $Q(D) \, u * v$, where $u$ and $v$ are in $C_0^\infty(\mathbb{R}^{n})$. Then, by \eqref{eq:VI.3.10} in Chapter~\ref{chap:VI}, \ref{sec:VI.3}, 
\[
\left| (Q (D) E * u * v) (0) \right| = \left| (E * Q (D) u * v) (0) \right| \leqq C N (Q (D) u * v),
\]
where $N(u) = \sup\limits_{|\eta| \leqq \varepsilon} \int_{\mathbb{R}^{n}} |\hat{u}(\xi + i\eta)| / \tilde{P}(\xi) \cdot \dd\xi$. 

The Fourier-Laplace transform of $Q(D)$ $u * v$ is, by \eqref{eq:VI.2.17} in Chapter~\ref{chap:VI}, \ref{sec:VI.2} and \eqref{eq:VI.3.15} in Chapter~\ref{chap:VI}, \ref{sec:VI.3}, equal to $(2\pi)^{n/2} Q(\zeta) \hat{u}(\zeta) \hat{v}(\zeta)$. Since, by Taylor's formula, 
\begin{equation}\label{eq:VI.10.4}
Q (\xi + i \eta) = \sum\limits_{\alpha} \frac {1}{\alpha !} (- \eta) ^{\alpha} D _{\alpha} Q (\xi), \text{ where } (- \eta) ^{\alpha} = \prod\limits_{j=1} ^{n} (- \eta_{j}) ^ {\alpha_{j}},
\end{equation}
we have, by \eqref{eq:VI.10.2}, 
\[
\left| Q (\xi + i \eta) \right| / \tilde {P} (\xi) \leqq C ^{\prime} \text{ for } | \eta | \leqq \varepsilon \text{ and } \xi \in \mathbb{R}^{n},
\]
where the constant $C'$ may depend on $\varepsilon$. Thus 
\[
N (Q (D) u * v) \leqq (2 \pi) ^{n/2} C ^{\prime} \cdot \sup\limits _{| \eta | \leqq \varepsilon} \int_ {\mathbb{R}^{n}} | \hat {u} (\xi + i \eta) \hat {v} (\xi + i \eta) | \dd \xi .
\]

By Parseval's theorem for the Fourier transform, we obtain, denoting by $\|\|$ the $L^2(\mathbb{R}^{n})$-norm, 
\[
\int_ {\mathbb{R}^{n}} | \hat {\boldsymbol {u}} (\xi + i \eta) | ^{2} \dd \xi = \int_ {\mathbb{R}^{n}} | u (x) | ^{2} e ^{2 \langle x, \eta \rangle} \dd x \leqq \| u (x) e ^{\varepsilon | x |} \| ^{2} \text{ when } | \eta | \leqq \varepsilon .
\]
and a similar estimate for $\hat{v}$. Thus, by Schwarz' inequality, 
$N(Q(D)u*v)\leqq C''\left\|u(x)e^{\varepsilon|x|}\right\|\left\|v(x)e^{\varepsilon|x|}\right\|$ whenever $u$ and $v\in C_{0}^{\infty}(\mathbb{R}^{n})$, $C''$ denoting a constant which may depend on $\varepsilon$. 

Therefore we have proved 
\begin{equation}\label{eq:VI.10.5}
\left| \int_ {\mathbb{R}^{n}} (Q (D) E * u) (x) v (- x) \dd x \right| \leqq (C C ^{\prime \prime}) \| u e ^{\varepsilon | x |} \| \cdot \| v e ^{\varepsilon | x |} \| \quad (u, v \in C _{0} ^{\infty} (\mathbb{R}^{n})).
\end{equation}

We shall write $L_{\varepsilon}^{2}(\mathbb{R}^{n})$ for the Hilbert space of functions $w(x)$ normed by 
\[
\left(\int_ {\mathbb{R}^{n}} | w (x) | ^{2} e ^{\varepsilon | x |} \dd x\right) ^{1/2} = \left\| w (x) e ^{\varepsilon | x |} \right\|.
\]

Since $C_0^\infty(\mathbb{R}^{n})$ is dense in $L_{\varepsilon}^2(\mathbb{R}^{n})$ and, as may be proved easily, $L_{-\varepsilon}^2(\mathbb{R}^{n})$ is the conjugate space of $L_{\varepsilon}^2(\mathbb{R}^{n})$, we obtain, by dividing \eqref{eq:VI.10.5} by $\|v(x)e^{\varepsilon|x|}\|$ and taking the supremum over $v \in C_0^\infty(\mathbb{R}^{n})$, 
\[
\left\| (Q (D) E * u) (x) e ^{- \varepsilon | x |} \right\| \leqq \left(C C ^{\prime \prime}\right) \left\| u (x) e ^{\varepsilon | x |} \right\|, u \in C _{0} ^{\infty} \left(\mathbb{R}^{n}\right).
\]

Hence the mapping 
\begin{equation}\label{eq:VI.10.6}
u \rightarrow Q (D) E * u
\end{equation}
can be extended by continuity from $C_0^\infty (\mathbb{R}^{n})$ to $L_{\varepsilon}^2 (\mathbb{R}^{n})$, so that it becomes a continuous linear mapping on $L_{\varepsilon}^{2}(\mathbb{R}^{n})$ into $L_{-\varepsilon}^{2}(\mathbb{R}^{n})$. Thus, to prove Theorem~\ref{thm:VI.10.Theorem.1}, we have to take $f_{1} = f$ in $\Omega$ and $f_{1} = 0$ in $\mathbb{R}^{n} - \Omega$ and define $u$ as equal to $u = E * f_{1}$. 

For the proof of Theorem~\ref{thm:VI.10.Theorem.2}, we prepare three Lemmas. 

\begin{lemma}[\textsc{Malgrange}]\label{lem:VI.10.Lemma.1} If $f(z)$ is a holomorphic function of a complex variable $z$ for $|z| \leqq 1$ and $p(z)$ is a polynomial in which the coefficient of the highest order term is $A$, then 
\begin{equation}\label{eq:VI.10.7}
\left| A f (0) \right| \leqq (2 \pi) ^{-1} \int_{- \pi} ^{\pi} \left| f \left(e ^{i \theta}\right) p \left(e ^{i \theta}\right) \right| \dd \theta .
\end{equation}
\end{lemma}

\begin{proof}
 Let $z_{j}$'s be the zeros of $p(z)$ in the unit circle $|z| < 1$ and put 
\[
p (z) = q (z) \prod\limits_{j} \frac {z - z _{j}}{\bar {z} _{j} z - 1}. \]

Then $q(z)$ is regular in the unit circle and $|p(z)| = |q(z)|$ for $|z| = 1$. Hence we have 
\[
\begin{aligned} (2 \pi) ^{-1} \int_{- \pi} ^{\pi} | f (e ^{i \theta}) p (e ^{i \theta}) | \dd \theta & = (2 \pi) ^{-1} \int_{- \pi} ^{\pi} | f (e ^{i \theta}) q (e ^{i \theta}) | \dd \theta \\ & \geqq (2 \pi) ^{-1} \left| \int_{- \pi} ^{\pi} f (e ^{i \theta}) q (e ^{i \theta}) \dd \theta \right| = | f (0) q (0) |. \end{aligned}
\]

Thus Lemma~\ref{lem:VI.10.Lemma.1} is proved since $|q(0)/A|$ is equal to the product of the absolute values of zeros of $p(z)$ outside the unit circle. 
\end{proof}

\begin{lemma}\label{lem:VI.10.Lemma.2} With the notations in Lemma~\ref{lem:VI.10.Lemma.1}, we have, if the degree of $p(z)$ is $\leqq m$, 
\begin{equation}\label{eq:VI.10.8}
\left| f (0) p ^{(k)} (0) \right| \leqq \frac {m !}{(m - k) !} (2 \pi) ^{-1} \int_{- \pi} ^{\pi} \left| f \left(e ^{i \theta}\right) p \left(e ^{i \theta}\right) \right| \dd \theta .
\end{equation}
\end{lemma}

\begin{proof}
 We may assume that the degree of $p(z)$ is $m$ and that 
\[
p (z) = \prod\limits_{j=1} ^{m} (z - z _{j}). \]

Applying the preceding Lemma~\ref{lem:VI.10.Lemma.1} to the polynomial $\prod\limits_{j=1}^{k}(z-z_j)$ and the holomorphic function $f(z) \cdot \prod\limits_{j=k+1}^{m}(z-z_j)$, we obtain 
\[
\left| f (0) \prod\limits_{j=k+1} ^{m} z _{j} \right| \leqq (2 \pi) ^{-1} \int_{- \pi} ^{\pi} | f (e ^{i \theta}) p (e ^{i \theta}) | \dd \theta .
\]

A similar inequality will hold for any product of $(m - k)$ of the numbers $z_{j}$ on the left hand side. Since $p^{(k)}(0)$ is the sum of $m!/(m-k)!$ such terms, multiplied by $(-1)^{m-k}$, we have proved the inequality \eqref{eq:VI.10.8}. 
\end{proof}

\begin{lemma}\label{lem:VI.10.Lemma.3} Let $F(\zeta) = F(\zeta_1, \zeta_2, \ldots, \zeta_n)$ be holomorphic for $|\zeta| = \left(\sum\limits_{j=1}^{n} |\zeta_j|^2\right)^{1/2} < \infty$ and $P(\zeta) = P(\zeta_1, \zeta_2, \ldots, \zeta_n)$. a polynomial of degree $\leqq m$. Let $\Phi(\zeta) = \Phi(\zeta_1, \zeta_2, \ldots, \zeta_n)$ be a non-negative integrable function with compact support depending only on $|\zeta_1|, |\zeta_2|, \ldots, |\zeta_n|$. Then we have 
\begin{equation}\label{eq:VI.10.9}
\left| F (0) D _{\alpha} P (0) \right| \int_{| \zeta | < \infty} \left| \zeta^{\alpha} \right| \Phi (\zeta) \dd \zeta \leqq \frac {m !}{(m - \alpha) !} \int_{| \zeta | < \infty} \left| F (\zeta) P (\zeta) \right| \Phi (\zeta) \dd \zeta ,
\end{equation}
where $\dd\zeta$ is the Lebesgue measure $\dd\xi_1\dd\eta_1\ldots \dd\xi_nd\eta_n$ ( $\zeta_k = \xi_k + i\eta_k$) and $m$ is regarded as multi-index $(m, m, \ldots, m)$. 
\end{lemma}

\begin{proof}
 Let $f(z)$ be an entire holomorphic function and apply \eqref{eq:VI.10.8} to the functions $f(\mathbf{r}z)$ and $p(\mathbf{r}z)$, where $r > 0$. Then we obtain 
\[
\left| f (0) p ^{(k)} (0) \right| \cdot r ^{k} \leqq \frac {m !}{(m - k) !} (2 \pi) ^{-1} \int_{- \pi} ^{\pi} \left| f (r e ^{i \theta}) p (r e ^{i \theta}) \right| \dd \theta . \]

Let $\psi(r)$ be a non-negative, integrable function with compact support. Multiplying the above inequality by $2\pi r\psi(r)$ and integrating with respect to $r$, we obtain 
\begin{equation}\label{eq:VI.10.10}
\left| f (0) p ^{(k)} (0) \right| \int | t | ^{k} \psi (| t |) \dd t \leqq \frac {m !}{(m - k) !} \int | f (t) p (t) | \psi (| t |) \dd t
\end{equation}
where $\dd t = r \dd r \dd\theta$ and the integration is extended over the whole complex $t$-plane. Lemma~\ref{lem:VI.10.Lemma.3} is obtained by applying \eqref{eq:VI.10.10} successively to the variables $\zeta_{1}, \zeta_{2}, \ldots, \zeta_{n}$, one at a time. 
\end{proof}

\begin{proof}[Proof of Theorem~\ref{thm:VI.10.Theorem.2}.]
Put $P(D) \, u = v$, where $u \in C_0^\infty(\mathbb{R}^{n})$. Then $P(\zeta) \, \hat{u}(\zeta) = \hat{v}(\zeta)$. Apply Lemma~\ref{lem:VI.10.Lemma.3} with $F(\zeta) = \hat{u}(\xi + \zeta)$, with $P(\zeta)$ replaced by $P(\xi + \zeta)$ and with $|\Phi(\zeta)| = 1$ when $|\zeta| \leqq \varepsilon$ and $= 0$ otherwise. Since $\tilde{P}(\xi) \leqq \sum\limits_{\alpha} |D^\alpha P(\xi)|$, we obtain, from \eqref{eq:VI.10.9},
\[
\left| \hat {u} (\xi) \tilde {P} (\xi) \right| \leqq C _{1} \int_{| \zeta | \leqq \varepsilon} \left| \hat {u} (\xi + \zeta) P (\xi + \zeta) \right| \dd \zeta = C _{1} \int_{| \zeta | \leqq \varepsilon} \left| \hat {v} (\xi + \xi) \right| \dd \zeta .
\]

Hence, by Fourier's integral theorem, 
\[
\begin{aligned}
\left| u (0) \right|& \leqq (2 \pi) ^{-n/2} \int | \hat {u} (\xi) | \dd \xi \leqq C _{1} ^{\prime} \int_{| \zeta | \leqq e} \left(\int | \hat {v} (\xi + \zeta) | / \tilde {P} (\xi) \dd \xi\right) \dd \zeta
\\
\quad& \leqq C _{1} ^{\prime} \int_ {\left(\xi^{\prime 2} + \eta^{\prime 2} \leqq e\right)} | \hat {v} (\xi + \xi^{\prime} + i \eta^{\prime}) | / \tilde {P} (\xi) \dd \xi^{\prime} \dd \eta^{\prime}\left. \right) \dd \xi .
\end{aligned}
\]

On the other hand, we have 
\[
\tilde {P} (\xi + \xi^{\prime}) / \tilde {P} (\xi) \leqq C _{2} \quad \text{ when } | \xi^{\prime} | \leqq \varepsilon ,
\]
because 
\[
D ^{\alpha} P (\xi + \xi^{\prime}) = \sum\limits_{\beta} \frac {(\xi^{\prime}) ^{\beta}}{\beta !} D ^{\alpha + \beta} P (\xi),
\]
so that $|D^{\alpha}P(\xi + \xi')| / \tilde{P}(\xi)$ is bounded when $|\xi'| \leqq \varepsilon$. Therefore we have 
\begin{equation}\label{eq:VI.10.11}
\begin{aligned}
&\left| u (0) \right| \leqq C _{1} ^{\prime} C _{2} \int \left(\int_ {\xi^{\prime 2} + \eta^{\prime 2} \leqq e ^{2}} \left| \hat {v} (\xi + \xi^{\prime} + i \eta^{\prime}) \right| / \tilde {P} (\xi + \xi^{\prime}) \dd \xi^{\prime} \dd \eta^{\prime}\right) \dd \xi \leqq C _{3} \| v \| ^{\prime}, \text{ where }
\\
&\| v \| ^{\prime} = \int_{| \eta | \leqq e} \left(\int | \hat {v} (\xi + i \eta) | / \tilde {P} (\xi) \dd \xi\right) \dd \eta \quad (u \in C _{0} ^{\infty} (R _{n})),
\end{aligned}
\end{equation}
$C_{3}$ denoting a constant depending only on $\varepsilon$. 

By the way, the finiteness of $\|v\|'$ is a consequence of the Paley-Wiener theorem in Chapter~\ref{chap:VI}, \ref{sec:VI.4}. Consider the space $\tilde{C}_{0}^{\infty}(\mathbb{R}^{n})$ which is the completion of $C_{0}^{\infty}(\mathbb{R}^{n})$ with respect to the norm $\|v\|'$. Then, by the Hahn-Banach extension theorem, the linear functional $L$: 
\[
v = P (D) u \rightarrow u (0) \quad (\text{ where } u \in C _{0} ^{\infty} (\mathbb{R}^{n}))
\]
can be extended to a continuous linear functional $L$ defined on $\tilde{C}_{0}^{\infty}(\mathbb{R}^{n})$. As in the case of the space $L^{1}(\mathbb{R}^{n})$, we see that there exists a Baire function $k(\xi + i\eta)$ bounded a.e. with respect to the measure $\tilde{P}(\xi)^{-1}\dd\xi \dd\eta$ such that the extended linear functional $L$ is represented as 
\begin{equation}\label{eq:VI.10.12}
L (v) = \int_{| \eta | \leqq e} (\int \hat {v} (\xi + i \eta) k (\xi + i \eta) / \tilde {P} (\xi) \cdot \dd \xi) \dd \eta .
\end{equation}

When $v_{h}(x) \in C_0^\infty(\mathbb{R}^{n})$ tends, as $h \to \infty$, to $0$ in the topology of $\mathfrak{D}(\mathbb{R}^{n})$, $v_{h}(x)e^{\langle x,\eta\rangle}$ also tends to zero in the topology of $\mathfrak{D}(\mathbb{R}^{n})$, uniformly with respect to $\eta$ for $|\eta| \leqq \varepsilon$. Hence, as in Chapter~\ref{chap:VI}, \ref{sec:VI.1}, we see easily that $\hat{v}_{h}(\xi + i\eta)$, as function of $\xi$, tends to zero in the topology of $\mathfrak{S}(\mathbb{R}^{n})$, uniformly in $\eta$ for $|\eta| \leqq \varepsilon$. Therefore, by \eqref{eq:VI.10.12}, $L$ defines a distribution $T \in \mathfrak{D}(\mathbb{R}^{n})'$. Thus, by \eqref{eq:VI.3.5} in Chapter~\ref{chap:VI}, \ref{sec:VI.3}. 
\begin{equation}\label{eq:VI.10.13}
L (v) = (T * \check {v}) (0) = (\check {T} * v) (0).
\end{equation}

We have thus proved Theorem~\ref{thm:VI.10.Theorem.2} by taking $E = \check{T}$. \eqref{eq:VI.10.3} is clear from \eqref{eq:VI.10.11}. 
\end{proof}


\yossection{11}{Differential Operators with Uniform Strength}
\label{sec:VI.11}

The existence theorem of the preceding section may be extended to a linear differential operator 
\begin{equation}\label{eq:VI.11.1}
P (x, D) = \sum\limits_{\alpha} a _{\alpha} (x) D _{\alpha}
\end{equation}
whose coefficients $a_{\alpha}(x)$ are continuous in an open bounded domain $\Omega$ of $\mathbb{R}^{n}$. 

\begin{definitionu} $P(x, D)$ is said to be of uniform strength in $\Omega$, if 
\begin{equation}\label{eq:VI.11.2}
\sup\limits _ {x, y \in \Omega , \xi \in \mathbb{R}^{n}} \tilde {P} (x, \xi) / \tilde {P} (y, \xi) < \infty ,
\end{equation}
where $\tilde{P}(x, \xi)$ is defined by $\left(\sum\limits_{\alpha} |P^{(\alpha)}(x, \xi)|^2\right)^{1/2}$ considering $x$ as parameters. 
\end{definitionu}

\begin{examplen} The differential operator $P(x, D) = \sum\limits_{|s|, |t| \leqq m} D^{s} a_{s,t}(x) D^{t}$ with real, bounded $C^{\infty}$ coefficients $a_{s,t}(x) = a_{t,s}(x)$ in $\Omega$ is strongly elliptic in $\Omega$ (see Chapter~\ref{chap:VI}, \ref{sec:VI.9}) if there exists a positive constant $\delta$ such that
\begin{equation}\label{eq:VI.11.3}
\sum\limits_{|s|,|t|=m} \xi^{s} a _{s,t} (x) \xi^{t} \geqq \delta \left(\sum\limits_{j=1} ^{n} \xi_{j} ^{2}\right) ^{m} \text{ in } \Omega .
\end{equation}

In such a case, $P(x, D)$ satisfies the condition \eqref{eq:VI.11.2}. Next let $P(x, D)$ be strongly elliptic in an open bounded domain $\Omega$ of $\mathbb{R}^{n-1}$. Then 
\begin{equation}\label{eq:VI.11.4}
\frac {\partial}{\partial x _{n}} - P (x, D)
\end{equation}
is said to be parabolic in the product space $\Omega\times\{x_{n};0<x_{n}\}$. It is easy to see that the operator \eqref{eq:VI.11.4} is of uniform strength in the above product space. 

\end{examplen}
\begin{theoremn}[\textsc{Hörmander} {\cite{hormander117}}] Let $P(x, D)$ be of uniform strength in an open bounded domain $\Omega$ of $\mathbb{R}^{n}$. For any point $x^0$, there exists an open subdomain $\Omega_1$ of $\Omega$ such that $x^0 \in \Omega_1$ and the equation $P(x, D)$ $u = f$ has, for every $f \in L^2(\Omega_1)$, a distribution solution $u \in L^2(\Omega_1)$ for which, moreover, $Q(D)$ $u \in L^2(\Omega_1)$ for every $Q(D)$ weaker than $P(x, D)$ for any fixed $x \in \Omega_1$.

\end{theoremn}
\begin{proof}
 Write $P(x^{0}, D) = P_{0}(D)$. The set of all the differential operators with constant coefficients weaker than $P_{0}(D)$ is a finite dimensional linear space. For, the degree of such operators $Q(D)$ cannot exceed that of $P_{0}(D)$. Thus there exist $P_{1}(D)$, $P_{2}(D)$, $\ldots$, $P_{N}(D)$ which form a basis for the differential operators weaker than $P_{0}(D)$. Hence we can write 
\begin{equation}\label{eq:VI.11.5}
P (x, D) = P _{0} (D) + \sum\limits_{j=1} ^{N} b _{j} (x) P _{j} (D), b _{j} (\boldsymbol {x} ^{0}) = 0,
\end{equation}
with uniquely determined $b_{j}(x)$ which are continuous in $\Omega$. 


By the result of the preceding section, there exists a bounded linear operator $T$ on $L^2(\Omega_1)$ into $L^2(\Omega_1)$ such that 
\begin{equation}\label{eq:VI.11.6}
P _{0} (D) T f = f \text{ for all } f \in L ^{2} (\Omega_{1}),
\end{equation}
and that the operators $P_j(D)$ $T$ are all bounded as operators on $L^2(\Omega_1)$ into $L^2(\Omega_1)$. Here $\Omega_1$ is any open subdomain of $\Omega$. We have only to take $Tf$ as the restriction to $\Omega_1$ of $E * f_1$ where $f_1 = f$ in $\Omega_1$ and $f_1 = 0$ in $\mathbb{R}^{n} - \Omega_1$. 

The equation $P(x, D)$ $u = f$ is equivalent to 
\begin{equation}\label{eq:VI.11.7}
P _{0} (D) u + \sum\limits_{j=1} ^{N} b _{j} (x) P _{j} (D) u = f.
\end{equation}

We shall seek a solution of the form $u = Tv$. Substituting this in \eqref{eq:VI.11.7}, we obtain, by \eqref{eq:VI.11.6}, 
\begin{equation}\label{eq:VI.11.8}
v + \sum\limits_{j=1} ^{N} b _{j} (x) P _{j} (D) T v = f.
\end{equation}

Let the sum of the norms of the bounded linear operators $P_{j}(D)$ $T$ on $L^{2}(\Omega_{1})$ into $L^{2}(\Omega_{1})$ be denoted by $C$. Since $b_{j}(x)$ is continuous and $b_{j}(x^{0}) = 0$, we may choose $\Omega_{1} \ni x^{0}$ so small that 
\[
C \left| b _{j} (x) \right| < 1 / N \text{ whenever } x \in \Omega_{1} (j = 1, 2, \dots , n).
\]

We may assume that the above inequalities hold whenever $x$ belongs to the compact closure of $\Omega_1$. Thus the norm of the operator $\sum\limits_{j=1}^{N} b_j(x) P_j(D) T$ is less than $1$, and so the equation \eqref{eq:VI.11.8} may be solved by Neumann's series (Theorem~\ref{thm:II.1.Theorem.2} in Chapter~\ref{chap:II}, \ref{sec:II.1}): 
\[
v = \left(I + \sum\limits_{j=1} ^{N} b _{j} P _{j} (D) T\right) ^{-1} f = A f,
\]
where $A$ is a bounded linear operator on $L^2(\Omega_1)$ into $L^2(\Omega_1)$. Hence $u = TAf$ is the required solution of $P(x, D)u = f$.
\end{proof}

\yossection{12}{The Hypoellipticity (Hörmänder's Theorem)}
\label{sec:VI.12}

We have defined in Chapter~\ref{chap:II}, \ref{sec:II.7} the notion of hypoellipticity of $P(D)$ and proved \textsc{Hörmander}'s theorem to the effect that, if $P(D)$ is hypoelliptic, then there exists, for any large positive constant $C_1$, a positive constant $C_2$ such that, for all solutions $\zeta = \xi + i\eta$ of the algebraic equation $P(\zeta) = 0$, 
\begin{equation}\label{eq:VI.12.1}
| \xi | < C _{2} \text{ if } | \eta | < C _{1}.
\end{equation}

To prove conversely that \eqref{eq:VI.12.1} implies the hypoellipticity of $P(D)$, we prepare the 

\begin{lemman}[\textsc{Hörmander} \cite{hormander113}] \eqref{eq:VI.12.1} implies that 
\begin{equation}\label{eq:VI.12.2}
\sum\limits_{| \alpha | > 0} \left| P ^{(\alpha)} (\xi) \right| ^{2} / \sum\limits_{| \alpha | \geqq 0} \left| P ^{(\alpha)} (\xi) \right| ^{2} \rightarrow 0 \text{ as } \xi \in \mathbb{R}^{n}, | \xi | \rightarrow \infty .
\end{equation}
\end{lemman}

\begin{proof}
 We first show that, for any real vector $\Theta\in \mathbb{R}^{n}$, we have 
\begin{equation}\label{eq:VI.12.3}
P (\xi + \Theta) / P (\xi) \rightarrow 1 \text{ as } \xi \in \mathbb{R}^{n}, | \xi | \rightarrow \infty . \end{equation}

We may assume that the coordinates are so chosen that $\Theta = (1,0,0,\ldots ,0)$. We have, by \eqref{eq:VI.12.1}, 
\[
P (\xi + i \eta) \neq 0 \text{ when } | \eta | < C _{1} \text{ and } | \xi | > C _{2}.
\]

Then the inequality $|\xi - \zeta'| \geqq C_1$ holds if $|\xi| \geqq C_1 + C_2$ and $P(\zeta') = 0$. For, setting $\zeta' = \xi' + i\eta'$, we have either $|\eta'| \geqq C_1$ or else $|\xi'| < C_2$ so that $|\xi - \xi'| \geqq C_1$. Giving fixed values to, $\xi_2, \xi_3, \ldots, \xi_n$ we can write 
\[
P (\xi) = C \prod\limits_{k=1} ^{m} \left(\xi_{1} - t _{k}\right), C \neq 0,
\]
where $(t_k, \xi_2, \ldots, \xi_n)$ is a zero of $P$. Hence we have $|t_k - \xi_1| \geqq C_1$ if $|\xi| \geqq C_1 + C_2$. Thus 
\[
\frac {P (\xi + \Theta)}{P (\xi)} = \prod\limits_{k=1} ^{m} \frac {\xi_{1} + . 1 - t _{k}}{\xi_{1} - t _{k}} = \prod\limits_{k=1} ^{m} \left(1 + \frac {1}{\xi_{1} - t _{k}}\right)
\]
satisfies 
\[
\left| \frac {P (\xi + \theta)}{P (\xi)} - 1 \right| \leqq m C _{1} ^{-1} (1 + C _{1} ^{-1}) ^{m-1} \quad \text{ if } \quad | \xi | \geqq C _{1} + C _{2}.
\]

As we may take $C_1$ arbitrarily large by taking $C_2$ sufficiently large, we have proved \eqref{eq:VI.12.3}. 

We have, by Taylor's formula, 
\[
P (\xi + \eta) = \sum\limits_{\alpha} \frac {1}{\alpha !} P ^{(\alpha)} (\xi) \eta^{\alpha}
\]
and so 
\begin{equation}\label{eq:VI.12.4}
\sum\limits_{i=1} ^{k} t _{i} P (\xi + \eta^{(i)}) = \sum\limits_{\alpha} \frac {1}{\alpha !} P ^{(\alpha)} (\xi) \sum\limits_{i=1} ^{k} t _{i} (\eta^{(i)}) ^{\alpha},
\end{equation}
where $\eta^{(i)}$ are arbitrary real vectors and $t_i$ arbitrary complex numbers. The coefficients $\sum\limits_{i=1}^{k} t_i (\eta^{(i)})^\alpha$, $|\alpha| \leqq m$, can be given arbitrary values by a convenient choice of $k$, $t_i$ and $\eta^{(i)}$. If otherwise, there would exist constants $C_\alpha$, $|\alpha| \leqq m$, not all equal to zero and such that $\sum\limits_{\alpha} C_\alpha \eta^\alpha = 0$ for every $\eta$. Thus 
\[
P ^{(\alpha)} (\xi) = \sum\limits_{i=1} ^{k} t _{i} P (\xi + \eta^{(i)}) \text{ with real vectors } \eta^{(i)}.
\]

Since the principal part on the right must cancel out when $|\alpha| \neq 0$, we must have $\sum\limits_{i=1}^{k} t_i = 0$. Hence, by \eqref{eq:VI.12.3}, we obtain \eqref{eq:VI.12.2}. 
\end{proof}

\begin{corollaryu} Suppose that $P_{1}(\xi)$ and $P_{2}(\xi)$ satisfy \eqref{eq:VI.12.2}. Then $P(\xi) = P_{1}(\xi) \cdot P_{2}(\xi)$ also satisfies \eqref{eq:VI.12.2}. Moreover, if $Q_{j}(D)$ is weaker than $P_{j}(D)$ ( $j = 1, 2$), it follows that $Q_{1}(D)$ $Q_{2}(D)$ is weaker than $P(D)$. 
\end{corollaryu}

\begin{proof}
 By applying Leibniz' formula of differentiation of product of functions, we see that $P^{(\alpha)}(\xi)$ is a linear combination of products of derivatives of $P_{1}(\xi)$ and $P_{2}(\xi)$ of the order sum $\leqq (\alpha)$. Hence \eqref{eq:VI.12.2} holds for $P(\xi)$. The latter part of the Corollary may be proved similarly. 
\end{proof}

We are now ready to prove 

\begin{theoremn}[\textsc{Hörmander} {\cite{hormander113}}] $P(D)$ is hypoelliptic iff the condition \eqref{eq:VI.12.2} is satisfied.

\end{theoremn}
\begin{proof}
 The ``only if'' part is already proved (Chapter~\ref{chap:II}, \ref{sec:II.7} and what was proved above). We shall prove the ``if'' part. 

Let $\Omega$ be an open subdomain of $\mathbb{R}^{n}$. A distribution $u \in \mathfrak{D}(\Omega)'$ is said to belong to $H_{\mathrm{loc}}^k(\Omega)$ if, for any $\varphi_0 \in C_0^\infty(\Omega)$, the Fourier transform $\hat{u}_0$ of $u_0 = \varphi_0 u$ satisfies (see Chapter~\ref{chap:VI}, \ref{sec:VI.2}) 
\begin{equation}\label{eq:VI.12.5}
\int_ {\mathbb{R}^{n}} (1 + | \xi | ^{2}) ^{k} | \hat {u _{0}} (\xi) | ^{2} \dd \xi < \infty , \text{ that is, if } u _{0} = \varphi_{0} u \in W ^{k,2} (\mathbb{R}^{n})
\end{equation}

By virtue of Sobolev's lemma in Chapter~\ref{chap:VI}, \ref{sec:VI.7}, the ``if'' part is proved by the following proposition: 

Let $P(\xi)$ satisfy \eqref{eq:VI.12.2}. If a distribution $u \in \mathfrak{D}(\Omega)'$ satisfies 
$P(D) u \in H_{\mathrm{loc}}^{s}(\Omega)$ with a positive $s$, then $u$ belongs to 
\begin{equation}\label{eq:VI.12.6}
H _ {\text{ loc }} ^{s} (\Omega)
\end{equation}

For, if $P(D) u \in C^{\infty}$ in $\Omega$, then $P(D) u \in H_{\mathrm{loc}}^{s}(\Omega)$ for every positive $s$ because of Leibniz' formula of differentiation. 

The following proof of \eqref{eq:VI.12.6} is based upon two Lemmas: 

\begin{lemma}\label{lem:VI.12.Lemma.1} Let $f \in W^{s,2}(\mathbb{R}^{n})$ and $\psi \in C_0^\infty(\mathbb{R}^{n}), s \geqq 0$. Then $\psi f \in W^{s,2}(\mathbb{R}^{n})$. 
\end{lemma}

\begin{lemma}\label{lem:VI.12.Lemma.2} Let $P(\xi)$ satisfy \eqref{eq:VI.12.2}. Then there exists a positive constant $\mu$ such that $|P^{(\alpha)}(\xi)\xi^{\mu}| / |P(\xi)|\to 0$ as $\xi \in \mathbb{R}^{n}$, $|\xi |\to \infty$ for every $\alpha \neq 0$. 
\end{lemma}

The proof of Lemma~\ref{lem:VI.12.Lemma.1} will be given later, and the proof of Lemma~\ref{lem:VI.12.Lemma.2} will not be given here (for the latter, we refer the reader to L. \textsc{Hörmander} \cite{hormander118} or to A. \textsc{Friedman} \cite{friedman68}). 

Next, let $\Omega_1$ and $\Omega_0$ be arbitrary open subdomains of $\Omega$ such that their closures $\Omega_1^a$ and $\Omega_0^a$ are compact and $\Omega_1^a \subseteq \Omega_0$, $\Omega_0^a \subseteq \Omega$. By Schwartz' Theorem in Chapter~\ref{chap:III}, \ref{sec:III.11}, the distribution $u \in \mathfrak{D}(\Omega)'$ is, when considered as a distribution $\in \mathfrak{D}(\Omega_0)'$, a distributional derivative of the form $D^t v$ of a function $v(x) \in L^2(\Omega_0)$. Let $\varphi \in C_0^\infty(\Omega_0)$ be such that $\varphi(x) = 1$ in $\Omega_1$. Then $u = u_0 = D^t \varphi v$ as distributions $\in \mathfrak{D}(\Omega_1)'$. $\varphi v$ being $\in L^2(\mathbb{R}^{n})$, we see that there exists a (possibly) negative integer $k$ such that 
\begin{equation}\label{eq:VI.12.7}
P ^{(\alpha)} (D) u _{0} = P ^{(\alpha)} (D) D ^{t} \varphi v \in W ^{k,2} \left(\mathbb{R}^{n}\right) \text{ for every } \alpha .
\end{equation}

Hence, by Lemma~\ref{lem:VI.12.Lemma.1} and the generalized Leibniz formula (see Chapter~\hyperref[chap:I]{$1$}, $8$) 
\begin{equation}\label{eq:VI.12.8}
P (D) \varphi_{1} u _{0} = \varphi_{1} P (D) u _{0} + \sum\limits_{(\alpha) > 0} \frac {1}{(\alpha) !} D _{\alpha} \varphi_{1} \cdot P ^{(\alpha)} (D) u _{0},
\end{equation}
we see by $P(D)u_0 \in H_{\mathrm{loc}}^s(\Omega)$ that, whenever $\varphi_1 \in C_0^\infty(\Omega_1)$, 
\begin{equation}\label{eq:VI.12.9}
P (D) \varphi_{1} u _{0} \in W ^ {k _{1}, 2} (\mathbb{R}^{n}) \text{ with } k _{1} = \min (s, k).
\end{equation}

Thus the Fourier transform $\hat{u}_1(\xi)$ of $u_1(x) = \varphi_1(x)$ $\dot{u}_0(x)$ satisfies 
\begin{equation}\label{eq:VI.12.10}
\int_ {\mathbb{R}^{n}} | P (\xi) \hat {u} _{1} (\xi) | ^{2} (1 + | \xi | ^{2}) ^ {k _{1}} \dd \xi < \infty
\end{equation}
and hence, by Lemma~\ref{lem:VI.12.Lemma.2}, 
\begin{equation}\label{eq:VI.12.11}
\begin{aligned}
&\int_ {\mathbb{R}^{n}} | P ^{(\alpha)} (\xi) \hat {u} _ {\mathbf {1}} (\xi) | ^{2} (1 + | \xi | ^{2}) ^ {k _{1} + \mu} \dd \xi < \infty , \text{ that is, }
\\
P ^{(\alpha)} (D) u _ {\mathbf {1}}& \in W ^ {k _{1} + \mu , 2} (\mathbb{R}^{n}) \text{ for every } \alpha \neq 0.
\end{aligned}
\end{equation}

Let $\Omega_{2}$ be any open subdomain of $\Omega_{1}$ such that its closure $\Omega_{2}^{a}$ is compact and contained in $\Omega_{1}$. Then, for any $\varphi_{2} \in C_{0}^{\infty}(\Omega_{2})$, we prove, by \eqref{eq:VI.12.8} and \eqref{eq:VI.12.11} as above, that 
\begin{align}
P (D) \varphi_{2} u _{1}& \in W ^ {k _{2}, 2} (\mathbb{R}^{n}) \text{ with } k _{2} = \min (s, k _{1} + \mu) \text{ and hence } \notag \\
P ^{(\alpha)} (D) \varphi_{2} u _{1}& \in W ^ {k _{2} + \mu , 2} (\mathbb{R}^{n}) \text{ for every } \alpha \neq 0. \notag
\end{align}

Repeating the argument a finite number of times, we see that, for any open subdomain $\Omega'$ of $\Omega$ such that its closure is compact and contained in $\Omega$, 
\[
P ^{(\alpha)} (D) \varphi u \in W ^{s,2} (\mathbb{R}^{n}) \text{ for all } \alpha \neq 0 \text{ whenever } \varphi \in C _{0} ^{\infty} (\Omega^{\prime})
\]

Thus, $P^{(\alpha)}(\xi) = \text{ constant } \neq 0$ gives $\varphi u \in W^{s,2}(\mathbb{R}^{n})$. 
\end{proof}

\begin{proof}[Proof of Lemma~\ref{lem:VI.12.Lemma.1}]
 The Fourier transform of $\psi f$ is 
\[
(2 \pi) ^{-n/2} \int_ {\mathbb{R}^{n}} \hat {\psi} (\eta) \hat {f} (\xi - \eta) \dd \eta (\text{ see Theorem~\hyperref[thm:VI.3.Theorem.6]{6} in Chapter~\hyperref[chap:VI]{VI}, 3 })
\]
and thus we have to show that, for $s \geqq 0$, 
\[
\int_ {\mathbb{R}^{n}} (1 + | \xi | ^{2}) ^{s} \left| \int_ {\mathbb{R}^{n}} \hat {\psi} (\eta) \hat {f} (\xi - \eta) \dd \eta \right| ^{2} \dd \xi < \infty .
\]

By Schwarz' inequality, this can be estimated from above by 
\begin{equation}\label{eq:VI.12.12}
\begin{gathered}
\int_ {\mathbb{R}^{n}} (1 + | \xi | ^{2}) ^{s} \left[ \int_ {\mathbb{R}^{n}} | \hat {\psi} (\eta) | \dd \eta \cdot \int_ {\mathbb{R}^{n}} | \hat {\psi} (\eta) | \cdot | \hat {f} (\xi - \eta) | ^{2} \dd \eta \right] \dd \xi
\\
= \int_ {\mathbb{R}^{n}} | \hat {\psi} (\eta) | \dd \eta \left[ \int_ {\mathbb{R}^{n}} \int_ {\mathbb{R}^{n}} (1 + | \xi | ^{2}) ^{s} | \hat {\psi} (\eta) | \cdot | \hat {f} (\xi - \eta) | ^{2} \dd \xi \dd \eta \right].
\end{gathered}
\end{equation}

We then make use of the inequality 
\begin{equation}\label{eq:VI.12.13}
(1 + | \xi | ^{2}) ^{s} \leqq 4 ^{|s|} (1 + | \eta | ^{2}) ^{|s|} (1 + | \xi - \eta | ^{2}) ^{s}
\end{equation}
which may be proved by 
\[
\frac {1 + | \xi | ^{2}}{1 + | \xi - \eta | ^{2}} \leqq 4 (1 + | \eta | ^{2}), 4 (1 + | \xi | ^{2}) \geqq \frac {1 + | \xi - \eta | ^{2}}{1 + | \eta | ^{2}}.
\]

By \eqref{eq:VI.12.13}, the right side of \eqref{eq:VI.12.12} is estimated by $\int_{\mathbb{R}^{n}} |\hat{\psi}(\eta)| \dd\eta$ -times 
\[
4 ^{|s|} \int_ {\mathbb{R}^{n}} | \hat {\psi} (\eta) | (1 + | \eta | ^{2}) ^{|s|} \dd \eta \cdot \left(\int_ {\mathbb{R}^{n}} (1 + | \xi - \eta | ^{2}) ^{s} | \hat {f} (\xi - \eta) | ^{2} \dd \xi\right).
\]

This integral converges since $f \in W^{s,2}(\mathbb{R}^{n})$ and $\hat{\psi}(\eta) \in \mathfrak{S}(\mathbb{R}^{n})$. 

We have thus proved our Theorem.
\end{proof}

\yossub{Further Researches}

$1$. A linear partial differential operator $P(x, D)$ with $C^\infty(\Omega)$ coefficients is said to be formally hypoelliptic in $\Omega \subseteq \mathbb{R}^{n}$ if the following two conditions are satisfied: \begin{enumerate}[i)]
\item $P(x^0, D)$ is hypoelliptic for every fixed $x^0 \in \Omega$
\item $P(x^0, \xi) = 0(P(x', \xi))$ as $\xi \in \mathbb{R}^{n}$, $|\xi| \to \infty$ for every fixed $x^0$ and $x' \in \Omega$.
\end{enumerate}
 L. \textsc{Hörmander} \cite{hormander117} and B. \textsc{Malgrange} \cite{malgrange187} have proved that, for such an operator $P(x, D)$, any distribution solution $u \in \mathfrak{D}(\Omega)'$ of the equation $P(x, D)u = f$ is $C^\infty$ after correction on a set of measure zero in the open subdomain $\subseteq \Omega$ where $f$ is $C^\infty$. The proof above for the constant coefficients case may be modified so as to apply for the formally hypoelliptic case, see. e.g., J. \textsc{Peetre} \cite{peetre246}. 

$2$. It was proved essentially by I. G. \textsc{Petrowsky} \cite{petrowsky248} that all distribution solutions $u \in \mathfrak{D}(\mathbb{R}^{n})$ of $P(D)u = 0$ are analytic functions in $\mathbb{R}^{n}$ iff the homogeneous part $P_m(\xi)$ of $P(\xi)$ of the highest degree $m$ does not vanish for $\xi \in \mathbb{R}^{n}$. If this condition is satisfied then $P(D)$ is said to be (analytically) elliptic. It is proved that in such a case the degree $m$ is even and $P(D)$ is hypoelliptic. It is to be noted that the hypoellipticity of an analytically elliptic operator $P(D)$ can also be proved by Friedrichs' Theorem in Chapter~\ref{chap:VI}, \ref{sec:VI.9}. For, by the non-vanishing of $P_m(\xi)$, we easily see, by the Fourier transformation, that $P(D)$ or $-P(D)$ is strongly elliptic. For the proof of Petrowsky's Theorem, see, e.g., L. \textsc{Hörmander} \cite{hormander118}, F. \textsc{Trèves} \cite{treves325} and C. B. \textsc{Morrey}-L. \textsc{Nirenberg} \cite{morrey214}.\yosstarfootnote{See also Supplementary Notes, p.~\pageref{notes:I-VI}.} 
   % VI
% ==== VII ====
\yoschapter{VII}{Dual Operators}
\label{chap:VII}


\yossection{1}{Dual Operators}
\label{sec:VII.1}

The notion of the transposed matrix may be extended to the notion of dual operator through 

\begin{theorem}\label{thm:VII.1.Theorem.1} Let $X$, $Y$ be locally convex linear topological spaces and $X_{s}^{\prime}, Y_{s}^{\prime}$ their strong dual spaces, respectively. Let $T$ be a linear operator on $D(T) \subseteq X$ into $Y$. Consider the points $\{x^{\prime}, y^{\prime}\}$ of the product space $X_{s}^{\prime} \times Y_{s}^{\prime}$ satisfying the condition 
\begin{equation}\label{eq:VII.1.1}
\langle T x, y ^{\prime} \rangle = \langle x, x ^{\prime} \rangle \quad \text{ for all } \quad x \in D (T).
\end{equation}
\end{theorem}

Then $x'$ is determined uniquely by $y'$ iff $D(T)$ is dense in $X$. 

\begin{proof}
 By the linearity of the problem, we have to consider the condition: 
\[
\langle x, x ^{\prime} \rangle = 0 \quad \text{ for all } \quad x \in D (T) \quad \text{ implies } \quad x ^{\prime} = 0. \]

Thus the ``if'' part is clear from the continuity of the linear functional $x'$. Assume that $D(T)^{a} \neq X$. Then there exists, by the Hahn-Banach theorem, an $x_{0}' \neq 0$ such that $\langle x, x_{0}' \rangle = 0$ for all $x \in D(T)$; consequently, the ``only if'' part is proved. 
\end{proof}

\begin{definition}\label{def:VII.1.Definition.1} $A$ linear operator $T'$ such that $T'y' = x'$ is defined through \eqref{eq:VII.1.1} iff $D(T)^a = X$. $T'$ is called the dual or conjugate operator of $T$; its domain $D(T')$ is the totality of those $y' \in Y_s'$ such that there exists $x' \in X_s'$ satisfying \eqref{eq:VII.1.1}, and $T'y' = x'$. Hence $T'$ is a linear operator defined on $D(T') \subseteq Y_s'$ into $X_s'$ such that 
\begin{equation}\label{eq:VII.1.2}
\langle T x, y ^{\prime} \rangle = \langle x, T ^{\prime} y ^{\prime} \rangle \quad \text{ for all } \quad x \in D (T) \quad \text{ and all } \quad y ^{\prime} \in D (T ^{\prime}).
\end{equation}
\end{definition}

\begin{theorem}\label{thm:VII.1.Theorem.2} If $D(T) = X$ and $T$ is continuous, then $T'$ is a continuous linear operator defined on $Y_s'$ into $X_s'$. 
\end{theorem}

\begin{proof}
 For any $y' \in Y_s'$, $\langle Tx, y' \rangle$ is a continuous linear functional of $x \in X$ and so there exists an $x' \in X_s'$ with $T'y' = x'$. Let $B$ be a bounded set of $X$. Then, by the continuity of $T$, the image $T \cdot B = \{Tx; x \in B\}$ is a bounded set of $Y$. Thus, by the defining relation $\langle Tx, y' \rangle = \langle x, x' \rangle$, the convergence to $0$ of $y'$ in the bounded convergence topology of $Y'$ (given in Chapter~\ref{chap:IV}, \ref{sec:IV.7}) implies the convergence of $x'$ in the bounded convergence topology of $X'$. Thus $T'$ is a continuous linear operator on $Y_s'$ into $X_s'$. 
\end{proof}

\begin{example}\label{ex:VII.1.Example.1} Let $X = Y$ be $n$-dimensional euclidean spaces normed by the $(l^2)$-norm. For any continuous linear operator $T \in L(X, X)$, set 
\[
T x = y, \quad \text{ where } \quad x = (x _{1}, x _{2}, \dots , x _{n}) \quad \text{ and } \quad y = (y _{1}, y _{2}, \dots , y _{n}).
\]
\end{example}

Then $y_{i} = \sum\limits_{j=1}^{n} t_{ij} x_{j} (i = 1, 2, \ldots, n)$ and so, for $z = (z_{1}, z_{2}, \ldots, z_{n})$, 
\[
\langle T x, z \rangle = \langle y, z \rangle = \sum\limits_{j} y _{j} z _{j} = \sum\limits_{i} \left(\sum\limits_{j} t _{ij} x _{j}\right) z _{i} = \sum\limits_{j} x _{j} \left(\sum\limits_{i} t _{ij} z _{i}\right).
\]

Thus $T'z = w$ is given by $w_{j} = \sum\limits_{i=1}^{n} t_{ij} z_{i} (j = 1, 2, \ldots, n)$. This proves that the matrix corresponding to $T'$ is the transposed matrix of the matrix corresponding to $T$. 

\begin{example}\label{ex:VII.1.Example.2} Let $X = Y$ be the real Hilbert space $(l^2)$, and let $T_n \in L(X, X)$ be defined by 
\[
T _{n} (x _{1}, x _{2}, \dots , x _{k}, \dots) = (x _{n}, x _{n+1}, x _{n+2}, \dots).
\]
\end{example}

Then from 
\[
\langle T _{n} (x _{1}, x _{2}, \dots), (z _{1}, z _{2}, \dots) \rangle = x _{n} z _{1} + x _{n+1} z _{2} + x _{n+2} z _{3} + \dots ,
\]
we obtain 
\[
T _{n} ^{\prime} (z _{1}, z _{2}, \dots) = (\overbrace {0 , 0 , \dots , 0}, z _{1}, z _{2}, \dots).
\]

Since $\|T_n(x_1, x_2, \ldots)\| = \left(\sum\limits_{m=n}^{\infty} x_m^2\right)^{1/2} \to 0$ as $n \to \infty$ and $\|T_n'(z_1, z_2, \ldots)\|$ $= \|(z_1, z_2, \ldots)\|$, we have 

\begin{proposition}\label{prop:VII.1.Proposition.1} The mapping $T \to T'$ of $L(X, Y)$ into $L(Y_s', X_s')$ is not, in general, continuous in the simple convergence topology of operators, that is, $\lim\limits_{n \to \infty} T_n x = Tx$ for all $x \in X$ does not necessarily imply $\lim\limits_{n \to \infty} T_n' y' = T'y'$ for all $y' \in Y'$, in the strong dual topology of $X_s'$. 
\end{proposition}

\begin{theoremn}[Theorem~\ref{thm:VII.1.Theorem.2}']\yosnamedtag{2'}\label{thm:VII.1.Theorem.2prime} Let $T$ be a bounded linear operator on a normed linear space $X$ into a normed linear space $Y$. Then the dual operator $T'$ is a bounded linear operator on $Y_s'$ into $X_s'$ such that 
\begin{equation}\label{eq:VII.1.3}
\| T \| = \left\| T ^{\prime} \right\|.
\end{equation}
\end{theoremn}

\begin{proof}
 From the defining relation $\langle Tx, y'\rangle = \langle x, x'\rangle$, we obtain 
\[
\begin{aligned}
&\left\| T ^{\prime} y ^{\prime} \right\| = \left\| x ^{\prime} \right\| = \sup\limits _{\| x \| \leqq 1} \left| \langle x, x ^{\prime} \rangle \right| = \sup\limits _{\| x \| \leqq 1} \left| \langle T x, y ^{\prime} \rangle \right| \leqq \left\| y ^{\prime} \right\| \cdot \sup\limits _{\| x \| \leqq 1} \left\| T x \right\| \leqq \left\| y ^{\prime} \right\| \cdot \left\| T \right\|,
\end{aligned}
\]
and so $\| T' \| \leqq \| T \|$. The reverse inequality is proved as follows. For any $x_0 \in X$, there exists an $f_0 \in Y'$ such that $\| f_0 \| = 1$ and $f_0(Tx_0) = \langle Tx_0, f_0 \rangle = \| Tx_0 \|$. Thus $f_0' = T'f_0$ satisfies $\langle x_0, f_0' \rangle = \| Tx_0 \|$ and so 
\begin{equation}\left\| T x _{0} \right\| = \left\langle x _{0}, T ^{\prime} f _{0} \right\rangle \leqq \left\| T ^{\prime} \right\| \left\| f _{0} \right\| \cdot \left\| x _{0} \right\| = \left\| T ^{\prime} \right\| \cdot \left\| x _{0} \right\|, i. e., \left\| T \right\| \leqq \left\| T ^{\prime} \right\|. \notag\end{equation}
\end{proof}

\begin{theorem}\label{thm:VII.1.Theorem.3} \begin{enumerate}[i)]
\item If $T$ and $S$ are $\in L(X, Y)$, then $(\alpha T + \beta S)' = (\alpha T' + \beta S')$.
\item Let linear operators $T, S$ be such that $D(T), D(S), R(T)$ and $R(S)$ are all contained in $X$. If $S$ is $\in L(X, X)$ and $D(T)^a = X$, then 
\begin{equation}\label{eq:VII.1.4}
(S T) ^{\prime} = T ^{\prime} S ^{\prime}.
\end{equation}
\end{enumerate}
\end{theorem}

If, moreover, $D(TS)^{a} = X$, then 
\begin{equation}\label{eq:VII.1.5}
(T S) ^{\prime} \supseteq S ^{\prime} T ^{\prime}, \text{ i.e., } (T S) ^{\prime} \text{ is an extension of } S ^{\prime} T ^{\prime}.
\end{equation}

\begin{proof}
 \begin{enumerate}[i)]
\item is clear.
\item $D(ST) = D(T)$ is dense in $X$, and so $(ST)'$ exists. If $y \in D((ST)')$, then, for any $x \in D(T) = D(ST)$, $\langle Tx, S'y \rangle = \langle STx, y \rangle = \langle x, (ST)'y \rangle$. This shows that $S'y \in D(T')$ and $T'S'y = (ST)'y$, that is, $(ST)' \subseteq T'S'$. Let, conversely, $y \in D(T'S')$, i.e., $S'y \in D(T')$. Then, for any $x \in D(T) = D(ST)$, $\langle STx, y \rangle = \langle Tx, S'y \rangle = \langle x, T'S'y \rangle$. This shows that $y \in D((ST)')$ and $(ST)'y = T'S'y$, that is, $T'S' \subseteq (ST)'$. We have thus proved \eqref{eq:VII.1.4}.
\end{enumerate}



To prove \eqref{eq:VII.1.5}, let $y \in D(S'T') = D(T')$. Then, for any $x \in D(TS)$, $\langle TSx, y \rangle = \langle Sx, T'y \rangle = \langle x, S'T'y \rangle$. This shows that $y \in D((TS)')$ and $(TS)'y = S'T'y$, that is, $S'T' \subseteq (TS)'$.
\end{proof}

\yossection{2}{Adjoint Operators}
\label{sec:VII.2}

The notion of transposed conjugate matrix may be extended to the notion of adjoint operator through 

\begin{definition}\label{def:VII.2.Definition.1} Let $X$, $Y$ be Hilbert spaces, and $T$ a linear operator defined on $D(T) \subseteq X$ into $Y$. Let $D(T)^{a} = X$ and let $T'$ be the dual operator of $T$. Thus $\langle Tx, y' \rangle = \langle x, T'y' \rangle$ for $x \in D(T)$, $y' \in D(T')$. If we denote by $J_{X}$ the one-to-one norm-preserving conjugate linear correspondence $X_{s}' \ni f \leftrightarrow y_{f} \in X$ (defined in Corollary~\ref{cor:III.6.Corollary.1} in Chapter~\ref{chap:III}, \ref{sec:III.6}), then $\langle Tx, y' \rangle = y'(Tx) = (Tx, J_{Y}y')$, $\langle x, T'y' \rangle = (T'y')(x) = (x, J_{X}T'y')$. 
\end{definition}

We have thus 
\[
(T x, J _{Y} y ^{\prime}) = (x, J _{X} T ^{\prime} y ^{\prime}), \text{ that is, } (T x, y) = (x, J _{X} T ^{\prime} J _{Y} ^{-1} y).
\]

In the special case when $Y = X$, we write 
\[
T ^{*} = J _{X} T ^{\prime} J _{X} ^{-1}
\]
and call $T^*$ the adjoint operator of $T$. 

\begin{remarku} If $X$ is the complex Hilbert space $(l^{2})$, we see, as in the Example in the preceding section, that the matrix corresponding to $T^{*}$ is the transposed conjugate matrix of the matrix corresponding to $T$. 
\end{remarku}

As in the case of dual operators, we can prove 

\begin{theorem}\label{thm:VII.2.Theorem.1} $T^*$ exists iff $D(T)^a = X$. It is defined as follows: $y \in X$ is in the domain of $D(T^*)$ iff there exists a $y^* \in X$ such that 
\begin{equation}\label{eq:VII.2.1}
(T x, y) = (x, y ^{*}) \quad \text{ holds for all } \quad x \in D (T),
\end{equation}
and we define $T^{*}y = y^{*}$. 
\end{theorem}

We can rewrite the above theorem in terms of the graph $G(A)$ of the linear operator $A$ (the graph was introduced in Chapter~\ref{chap:II}, \ref{sec:II.6}): 

\begin{theorem}\label{thm:VII.2.Theorem.2} We introduce a continuous linear operator $V$ on $X \times X$ into $X \times X$ by 
\begin{equation}\label{eq:VII.2.2}
V \{x, y \} = \{- y, x \}.
\end{equation}
\end{theorem}

Then $(VG(T))^{\perp}$ is the graph of a linear operator iff $D(T)^{a} = X$, and, in fact, we have 
\begin{equation}\label{eq:VII.2.3}
G (T ^{*}) = (V G (T)) ^{\perp}.
\end{equation}

\begin{proof}
 The condition $\{-Tx, x\} \perp \{y, y^*\}$ is equivalent to $(Tx, y) = (x, y^*)$. Thus Theorem~\ref{thm:VII.2.Theorem.2} is proved by Theorem~\ref{thm:VII.2.Theorem.1}. 
\end{proof}

\begin{corollaryu} $T^{*}$ is a closed linear operator, since the orthogonal complement of a linear subspace is a closed linear subspace. 
\end{corollaryu}

\begin{theorem}\label{thm:VII.2.Theorem.3} Let $T$ be a linear operator on $D(T) \subseteq X$ into $X$ such that $D(T)^a = X$. Then $T$ admits a closed linear extension iff $T^{**} = (T^*)^*$ exists, i.e., iff $D(T^*)^a = X$. 
\end{theorem}

\begin{proof}
 Sufficiency. We have $T^{**} \supseteq T$ by definition, and $T^{**} = (T^{*})^{*}$ is closed by the above Corollary. 

Necessity. Let $S$ be a closed extension of $T$. Then $G(S)$ contains $G(T)^{a}$ as a closed linear subspace, and so $G(T)^{a}$ is the graph of a linear operator. But $G(T)^{a} = G(T)^{\perp\perp} = (G(T)^{\perp})^{\perp}$ by the continuity of the scalar product, and, moreover, by $VG(T^{*}) = G(T)^{\perp}$, we obtain $(VG(T^{*}))^{\perp} = G(T)^{\perp\perp}$. Therefore, $(VG(T^{*}))^{\perp}$ is the graph of a linear operator. Thus by Theorem~\ref{thm:VII.2.Theorem.2} $D(T^{*})^{a} = X$ and $T^{**}$ exists. 
\end{proof}

\begin{corollaryu} Under the condition that $D(T)^{a}=X$, $T$ is closed linear iff $T=T^{**}$. 
\end{corollaryu}

\begin{proof}
 The sufficiency is clear. Necessity is proved by observing the formula $G(T)^{a} = G(T^{**})$ obtained above. For, $G(T) = G(T)^{a}$ implies that $T = T^{**}$. 
\end{proof}

\begin{theorem}\label{thm:VII.2.Theorem.4} An everywhere defined closed linear operator $T$ is a continuous linear operator. 
\end{theorem}

\begin{proof}
 Clear from the closed graph theorem. 
\end{proof}

\begin{theorem}\label{thm:VII.2.Theorem.5} If $T$ is a bounded linear operator, then $T^*$ is also bounded linear and 
\begin{equation}\label{eq:VII.2.4}
\left\| T \right\| = \left\| T ^{*} \right\|.
\end{equation}
\end{theorem}

\begin{proof}
 Similar to the case of dual operators. 
\end{proof}


\yossection{3}{Symmetric Operators and Self-adjoint Operators}
\label{sec:VII.3}

A Hermitian matrix is a matrix which is equal to its transposed conjugate matrix. It is known that such a matrix can be transformed into a diagonal matrix by a suitable (complex) rotation of the vector space on which the matrix operates as a linear operator. The notion of the Hermitian matrices is extended to the notion of self-adjoint operators in a Hilbert space. 

\begin{definition}\label{def:VII.3.Definition.1} Let $X$ be a Hilbert space. A linear operator on $D(T) \subseteq X$ into $X$ is called symmetric if $T^* \supseteq T$, i.e., if $T^*$ is an extension of $T$. Note that the condition of the existence of $T^*$ implies that $D(T)^a = X$. 
\end{definition}

\begin{proposition}\label{prop:VII.3.Proposition.1} If $T$ is symmetric, then $T^{**}$ is also symmetric. 
\end{proposition}

\begin{proof}
 Since $T$ is symmetric, we have $D(T^{*}) \supseteq D(T)$ and $D(T)^{a} = X$. Hence $D(T^{*})^{a} = X$ and so $T^{**} = (T^{*})^{*}$ exists. $T^{**}$ is surely an extension of $T$ and so $D(T^{**}) \supseteq D(T)$. Thus, again by $D(T)^{a} = X$, we have $D(T^{**})^{a} = X$ and so $T^{***} = (T^{**})^{*}$ exists. We have, from $T^{*} \supseteq T$, $T^{**} \subseteq T^{*}$ and hence $T^{***} \supseteq T^{**}$ which proves that $T^{**}$ is symmetric. Corollary. A symmetric operator $T$ has a closed symmetric extension $T^{**} = (T^{*})^{*}$. 
\end{proof}

\begin{definition}\label{def:VII.3.Definition.2} $A$ linear operator $T$ is called self-adjoint if $T = T^{*}$. 
\end{definition}

\begin{proposition}\label{prop:VII.3.Proposition.2} $A$ self-adjoint operator is closed. An everywhere defined symmetric operator is bounded and self-adjoint. 
\end{proposition}

\begin{proof}
 Being the adjoint of itself, a self-adjoint operator is closed. The last assertion is proved by the fact that an everywhere defined closed operator is bounded (closed graph theorem). 
\end{proof}

\begin{example}[integral operator of the Hilbert-Schmidt type]\label{ex:VII.3.Example.1} Let $-\infty \leqq a < b \leqq \infty$ and consider $L^2(a, b)$. Let $K(s, t)$ be a complexvalued measurable function for $a \leqq s, t \leqq b$ such that 
\[
\int_{a} ^{b} \int_{a} ^{b} | K (s, t) | ^{2} \dd s \dd t < \infty .
\]
\end{example}

For any $x(t) \in L^{2}(a, b)$, we define the operator $K$ by 
\[
(K \cdot x) (s) = \int_{a} ^{b} K (s, t) x (t) \dd t.
\]

We have, by Schwarz' inequality and the Fubini-Tonelli theorem. 
\[
\int_{a} ^{b} | (K \cdot x) (s) | ^{2} \dd s \leqq \int_{a} ^{b} \int_{a} ^{b} | K (s, t) | ^{2} \dd t \dd s \int_{a} ^{b} | x (t) | ^{2} \dd t.
\]

Hence $K$ is a bounded linear operator on $L^2(a, b)$ into $L^2(a, b)$ such that $\|K\| \leqq \left( \int_{a}^{b} \int_{a}^{b} |K(s, t)|^2 \dd s \dd t \right)^{1/2}$. It is easy to see that the operator $K^*$ is defined by $(K^* y)(t) = \int_{a}^{b} \overline{K(t, s)} y(s) \dd s$. Hence $K$ is self-adjoint iff $K(s, t) = \overline{K(t, s)}$ for a.e. $s, t$. 

\begin{example}[the coordinate operator in quantum mechanics]\label{ex:VII.3.Example.2} Let $X = L^2(-\infty, \infty)$. Let $D = \{x(t); x(t) \text{ and } t \cdot x(t) \text{ both } \in L^2(-\infty, \infty)\}$. Then the operator $T$ defined by $Tx(t) = t \cdot x(t)$ on $D$ is self-adjoint. 
\end{example}

\begin{proof}
 It is clear that $D^a = X$, since the linear combinations of defining functions of finite intervals are strongly dense in $L^2(-\infty, \infty)$. Let $y \in D(T^*)$ and set $T^*y = y^*$. Then, for all $x \in D = D(T)$, 
\[
\int_{- \infty} ^{\infty} t x (t) \overline {{y (t)}} \dd t = \int_{- \infty} ^{\infty} x (t) \overline {{y ^{*} (t)}} \dd t. \]

If we take for $x(t)$ the defining function of the interval $[\alpha, t_0]$, we have $\int_{\alpha}^{t_0} t \cdot \overline{y(t)} \dd t = \int_{\alpha}^{t_0} y^*(t) \dd t$, and hence, by differentiation, $t_0 \cdot \overline{y(t_0)} = \overline{y^*(t_0)}$ for a.e. $t_0$. Thus $y \in D$ and $T^* y(t) = t \cdot y(t)$. Conversely, it is clear that $y \in D$ implies that $y \in D(T^*)$ and $T^* y(t) = t \cdot y(t)$. 
\end{proof}

\begin{example}[the momentum operator in quantum mechanics]\label{ex:VII.3.Example.3} Let $X = L^{2}(-\infty, \infty)$. Let $D$ be the totality of $x(t) \in L^{2}(-\infty, \infty)$ such that $x(t)$ is absolutely continuous on every finite interval with the derivative $x'(t) \in L^{2}(-\infty, \infty)$. Then the operator $T$ defined by $Tx(t) = i^{-1}x'(t)$ on $D$ is self-adjoint. 
\end{example}

\begin{proof}
 Let a continuous function $x_{n}(t)$ be defined by 
\begin{equation}x _{n} (t) = 1 \text{ for } t \in [ \alpha , t _{0} ], x _{n} (t) = 0 \text{ for } t \leqq \alpha - n ^{-1} \text{ and for } t \geqq t _{0} + n ^{-1}, \notag\end{equation}
$x_{n}(t)$ is a linear function on $[\alpha - n^{-1}, \alpha]$ and on $[t_0, t_0 + n^{-1}]$. 

Then the linear combinations of functions of the form $x_{n}(t)$ with different values of $\alpha, t_0$ and $n$ are dense in $L^2(-\infty, \infty)$. Thus $D$ is dense in $X$. Let $y \in D(T^*)$ and $T^*y = y^*$. Then for any $x \in D$, 
\[
\int_{- \infty} ^{\infty} i ^{-1} x ^{\prime} (t) \overline {{{y (t)}}} \dd t = \int_{- \infty} ^{\infty} x (t) \overline {{{y ^{*} (t)}}} \dd t.
\]

If we take $x_{n}(t)$ for $x(t)$, we obtain 
\[
n \int_ {x - n ^{-1}} ^{\alpha} i ^{-1} \overline {{{y (t)}}} \dd t - n \int_ {t _{0}} ^ {t _{0} + n ^{-1}} i ^{-1} \overline {{{y (t)}}} \dd t = \int_{- \infty} ^{\infty} x _{n} (t) y ^{*} (t) \dd t,
\]
and so, by letting $n \to \infty$, we obtain $i^{-1}(\overline{y(\alpha)} - \overline{y(t_0)}) = \int_{\alpha}^{t_0} \overline{y^*(t)} \dd t$ for a.e. $\alpha$ and $t_0$. It is clear, by Schwarz' inequality, that $y^*(t)$ is integrable over any finite interval. Thus $y(t_0)$ is absolutely continuous in $t_0$ over any finite interval, and so we have $i^{-1}y'(t_0) = y^*(t_0)$ for a.e. $t_0$. Hence $y \in D$ and $T^*y(t) = i^{-1}y'(t)$. Let, conversely, $y \in D$. Then, by partial integration, 
\[
\int_{a} ^{b} i ^{-1} x ^{\prime} (t) \overline {{{y (t)}}} \dd t = i ^{-1} [ x (t) \overline {{{y (t)}}} ] _{a} ^{b} + \int_{a} ^{b} x (t) \overline {{{(i ^{-1} y ^{\prime} (t))}}} \dd t.
\]

By the integrability of $x(t) \overline{y(t)}$ over $(-\infty, \infty)$, we see that $\lim\limits_{a \downarrow -\infty, b \uparrow \infty} |[x(t)y(t)]_a^b| = 0$, and so $\int_{-\infty}^{\infty} i^{-1} x'(t) \overline{y(t)} \dd t = \int_{-\infty}^{\infty} x(t) (\overline{i^{-1} y'(t)}) \dd t$. Thus $y \in D(T^*)$ and $T^* y(t) = i^{-1} y'(t)$. 
\end{proof}

\begin{theorem}\label{thm:VII.3.Theorem.1} If a self-adjoint operator $T$ admits the inverse $T^{-1}$, then $T^{-1}$ is also self-adjoint. 
\end{theorem}

\begin{proof}
 $T = T^{*}$ is equivalent to $(VG(T))^{\perp} = G(T)$. We have also $G(T^{-1}) = VG(-T)$. Hence, by $(-T)^{*} = -T^{*} = -T$, $(VG(-T))^{\perp} = G(-T)$ and so 
\[
(V G (T ^{-1})) ^{\perp} = G (- T) ^{\perp} = (V G (- T)) ^{\perp \perp} = V G (- T) = G (T ^{-1}),
\]
that is, $(T^{-1})^{*} = T^{-1}$. 


We have used, in the above proof, the fact that $(VG(-T))^a = VG(-T)$ in virtue of the closedness of $(-T)$. 
\end{proof}

\begin{corollaryu} $A$ symmetric operator $T$ in a Hilbert space $X$ is self-adjoint if $D(T) = X$ or if $R(T) = X$. 
\end{corollaryu}

\begin{proof}
 The case $D(T) = X$ was proved already. We shall prove the case $R(T) = X$. $Tx = 0$ implies $0 = (Tx, y) = (x, Ty)$ for all $y \in D(T)$, and so, by $R(T) = X$, we must have $x = 0$. Therefore the inverse $T^{-1}$ exists which is surely symmetric with $T$. $D(T^{-1}) = R(T) = X$, and so the everywhere defined symmetric operator $T^{-1}$ must be self-adjoint. Hence $T = (T^{-1})^{-1}$ is self-adjoint by Theorem~\ref{thm:VII.3.Theorem.1}. 
\end{proof}

We can construct self-adjoint operators from a closed linear operator. More precisely, we have 

\begin{theorem}[J. von \textsc{Neumann} {\cite{neumann232}}]\label{thm:VII.3.Theorem.2} For any closed linear operator $T$ in a Hilbert space $X$ such that $D(T)^{a} = X$, the operators $T^{*}T$ and $IT^{*}$ are self-adjoint, and $(I + T^{*}T)$ and $(I + TT^{*})$ both admit bounded linear inverses. 
\end{theorem}

\begin{proof}
 We know that, in the product space $X \times X$, $G(T)$ and $VG(T^{*})$ are closed linear subspaces orthogonal to each other and spanning the whole product space $X \times X$. Hence, for any $h \in X$, we have the uniquely determined decomposition 
\begin{equation}\label{eq:VII.3.1}
\{h, 0 \} = \{x, T x \} + \{- T ^{*} y, y \} \quad \text{ with } \quad x \in D (T), y \in D (T ^{*}). \end{equation}

Thus $h = x - T^{*}y$, $0 = Tx + y$. Therefore 
\begin{equation}\label{eq:VII.3.2}
x \in D (T ^{*} T) \quad \text{ and } \quad x + T ^{*} T x = h.
\end{equation}

Because of the uniqueness of decomposition \eqref{eq:VII.3.1}, $x$ is uniquely determined by $h$, and so the everywhere defined inverse $(I + T^{*}T)^{-1}$ exists. 

For any $h, k \in X$, let 
\[
x = (I + T ^{*} T) ^{-1} h, y = (I + T ^{*} T) ^{-1} k.
\]

Then $x$ and $y \in D(T^*T)$ and, by the closedness of $T$, $(T^*)^* = T$. Hence 
\[
\begin{aligned}
&(h, (I + T ^{*} T) ^{-1} k)  = ((I + T ^{*} T) x, y) = (x, y) + (T ^{*} T x, y) \\
&= (x, y) + (T x, T y) = (x, y) + (x, T ^{*} T y) = (x, (I + T ^{*} T) y) = ((I + T ^{*} T) ^{-1} h, k),
\end{aligned}
\]
which proves that the operator $(I + T^{*}T)^{-1}$ is self-adjoint. As an everywhere defined self-adjoint operator, $(I + T^{*}T)^{-1}$ is a bounded operator. By Theorem~\ref{thm:VII.3.Theorem.1}, its inverse $(I + T^{*}T)$ and hence $T^{*}T$ are self-adjoint. 

Since $T$ is closed, we have $(T^{*})^{*} = T$, and so, by what was proved above, $TT^{*} = (T^{*})^{*}T^{*}$ is self-adjoint and $(I + TT^{*})$ has a bounded linear inverse. 
\end{proof}

We next give an example of a non-self-adjoint, symmetric operator: 

\begin{example}\label{ex:VII.3.Example.4} Let $X = L^2(0,1)$. Let $D$ be the totality of absolutely continuous functions $x(t) \in L^2(0,1)$ such that $x(0) = x(1) = 0$ and $x'(t) \in L^2(0,1)$. Then the operator $T_1$ defined by $T_1x(t) = i^{-1}x'(t)$ on $D = D(T_1)$ is symmetric but not self-adjoint. 
\end{example}

\begin{proof}
 We shall prove that $T_1^* = T_2$, where $T_2$ is defined by: 
\[
T _{2} x (t) = i ^{-1} x ^{\prime} (t) \quad \text{ on } \quad D (T _{2}) = \{x (t) \in L ^{2} (0, 1); x (t) \text{ is }
\]
absolutely continuous such that $x'(t) \in L^2(0, 1)$. 

Since $D = D(T_{1})$ is dense in $L^2(0,1)$, the operator $T_1^*$ is defined. Let $y \in D(T_1^*)$ and set $T_1^* y = y^*$. Then, for any $x \in D = D(T_1)$, 
\[
\int_{0} ^{1} i ^{-1} x ^{\prime} (t) \overline {{y (t)}} \dd t = \int_{0} ^{1} x (t) \overline {{y ^{*} (t)}} \dd t.
\]

By partial integration, we obtain, remembering $x(0) = x(1) = 0$, 
\[
\int_{0} ^{1} x (t) \overline {{{y ^{*} (t)}}} \dd t = - \int_{0} ^{1} x ^{\prime} (t) \overline {{{Y ^{*} (t)}}} \dd t, \text{ where } Y ^{*} (t) = \int_{0} ^{t} y ^{*} (s) \dd s.
\]

Hence, by $x(1) = \int_{0}^{1} x'(t) \dd t = 0$, we have, for any constant $c$, 
\[
\int_{0} ^{1} x ^{\prime} (t) \overline {{{(Y ^{*} (t))}}} - i ^{-1} \overline {{{y (t)}}} - \overline {{{c}}}) \dd t = 0 \quad \text{ for all } \quad x \in D (T _{1}).
\]

On the other hand, for any $z(t) \in L^2(0, 1)$, the function $Z(t) = \int_0^t z(t) \dd t - t \int_0^1 z(t) \dd t$ surely belongs to $D(T_1)$. Hence, taking $Z(t)$ for the above $x(t)$, we obtain 
\[
\int_{0} ^{1} \left\{z (t) - \int_{0} ^{1} z (t) \dd t \right\} \cdot \overline {{(Y ^{*} (t)}} - i ^{-1} y (t) - \bar {c}) \dd t = 0.
\]

If we take the constant $c$ in such a way that $\int_0^1 (Y^* (t) - i^{-1}y(t) - c)\dd t = 0$, then 
\[
\int_{0} ^{1} z (t) \overline {{(Y ^{*} (t)}} - \overline {{i ^{-1} y (t)}} - \overline {{c}}) \dd t = 0,
\]
and so, by the arbitrariness of $z \in L^2(0,1)$, we must have $Y^*(t) = \int_0^t y^*(t) \dd t = i^{-1} y(t) + c$. Hence $y \in D(T_2)$ and $T_2y = y^*$. This proves that $T_1^* \subseteq T_2$. It is also clear, by partial integration, that $T_2 \subseteq T_1^*$ and so $T_2 = T_1^*$. 
\end{proof}

\begin{theorem}\label{thm:VII.3.Theorem.3} If $H$ is a bounded self-adjoint operator, then 
\begin{equation}\label{eq:VII.3.3}
\left\| H \right\| = \sup\limits _{\left\| x \right\| \leqq 1} \left| (H x, x) \right|.
\end{equation}
\end{theorem}

\begin{proof}
 Set $\sup\limits_{\|x\|\leqq 1}|(Hx,x)| = \gamma$. Then, by $|(Hx,x)|\leqq \|Hx\|\|x\|,$ $\gamma \leqq \|H\|$. For any real number $\lambda$, we have 
\begin{equation}\left| (H (y \pm \lambda z), y \pm \lambda z) \right| = \left| (H y, y) \pm 2 \lambda \operatorname{Re} (H y, z) + \lambda^{2} (H z, z) \right| \leqq \gamma \| y \pm \lambda z \| ^{2}. \notag\end{equation}

Hence 
\[
\left| 4 \lambda \operatorname{Re} (H y, z) \right| \leqq \gamma (\left\| y + \lambda z \right\| ^{2} + \left\| y - \lambda z \right\| ^{2}) = 2 \gamma (\left\| y \right\| ^{2} + \lambda^{2} \left\| z \right\| ^{2}).
\]

By taking $\lambda = \|y\|/\|z\|$, we obtain $|\operatorname{Re}(Hy,z)| \leqq \gamma \|y\|/\|z\|$. Hence, by substituting $ze^{i\theta}$ for $z$, we obtain $|(Hy,z)| \leqq \gamma \|y\|/\|z\|$, and so 
\[
(H y, H y) = \| H y \| ^{2} \leqq \gamma \| y \| \| H y \|, \text{ i.e., } \| H \| \leqq \gamma .
\]
\end{proof}


\yossection{4}{Unitary Operators. The Cayley Transform}
\label{sec:VII.4}

A symmetric operator is not necessarily a bounded operator. Various investigations of a symmetric operator $H$ may be made through the continuous operator $(H - iI)(H + iI)^{-1}$ called the Cayley transform of $H$. We shall begin with the notion of isometric operators. 

\begin{definition}\label{def:VII.4.Definition.1} $A$ bounded linear operator $T$ on a Hilbert space $X$ into $X$ is called (bounded) isometric if $T$ leaves the scalar product invariant: 
\begin{equation}\label{eq:VII.4.1}
(T x, T y) = (x, y) \quad \text{ for all } \quad x, y \in X.
\end{equation}
\end{definition}

If, in particular, $R(T) = X$, then a (bounded) isometric operator $T$ is called a unitary operator. 

\begin{proposition}\label{prop:VII.4.Proposition.1} For a bounded linear operator $T$, condition \eqref{eq:VII.4.1} is equivalent to the condition of the isometry 
\begin{equation}\label{eq:VII.4.2}
\left\| T x \right\| = \left\| x \right\| \quad \text{ for all } \quad x \in X.
\end{equation}
\end{proposition}

\begin{proof}
 It is clear that \eqref{eq:VII.4.1} implies \eqref{eq:VII.4.2}. We have, by \eqref{eq:VII.4.2}, 
\[ \begin{aligned}
&4 \operatorname{Re} (T x, T y)  = \| T (x + y) \| ^{2} - \| T (x - y) \| ^{2} = \| x + y \| ^{2} - \| x - y \| ^{2} = 4 \operatorname{Re} (x, y).
\end{aligned}
\]

By taking $iy$ in place of $y$, we also obtain $4\operatorname{Im}(Tx, Ty) = 4\operatorname{Im}(x, y)$, and so \eqref{eq:VII.4.2} implies \eqref{eq:VII.4.1}. 
\end{proof}

\begin{proposition}\label{prop:VII.4.Proposition.2} $A$ bounded linear operator on a Hilbert space $X$ into $X$ is unitary iff $T^{*} = T^{-1}$. 
\end{proposition}

\begin{proof}
 If $T$ is unitary, then $T^{-1}$ surely exists in virtue of condition \eqref{eq:VII.4.2}, and $D(T^{-1}) = R(T) = X$. Moreover, by \eqref{eq:VII.4.1}, $T^{*}T = I$ and so $T^{*} = T^{-1}$. Conversely, the condition $T^{*} = T^{-1}$ implies $T^{*}T = I$ which is the condition of the invariance of the scalar product. Moreover, $T^{*} = T^{-1}$ implies that $R(T) = D(T^{-1}) = D(T^{*}) = X$ and hence we see that $T$ must be unitary. 
\end{proof}

\begin{example}\label{ex:VII.4.Example.1} Let $X = L^{2}(-\infty, \infty)$. Then, for any real number a, the operator $T$ defined by $Tx(t) = x(t + a)$ on $L^{2}(-\infty, \infty)$ is unitary. 
\end{example}

\begin{example}\label{ex:VII.4.Example.2} The Fourier transform on $L^2(\mathbb{R}^{n})$ onto $L^2(\mathbb{R}^{n})$ is unitary, since it leaves the scalar product $(f, g) = \int f(x) \overline{g(x)} \dd x$ invariant. 
\end{example}

\begin{definition}\label{def:VII.4.Definition.2} Let $X$ be a Hilbert space. A linear operator $T$ defined on $D(T) \subseteq X$ into $X$ such that $D(T)^a = X$ is called normal if 
\begin{equation}\label{eq:VII.4.3}
T T ^{*} = T ^{*} T.
\end{equation}
\end{definition}

Self-adjoint operators and unitary operators are normal. 

\yossub{The Cayley Transform}

\begin{theorem}[J. \textsc{von} \textsc{Neumann} {\cite{neumann228}}]\label{thm:VII.4.Theorem.1} Let $H$ be a closed symmetric operator in a Hilbert space $X$. Then the continuous (but not necessarily everywhere defined) inverse $(H + iI)^{-1}$ exists, and the operator 
\begin{equation}\label{eq:VII.4.4}
U _{H} = (H - iI) (H + iI) ^{-1} \text{ with the domain } D(U_{H}) = D((H + iI) ^{-1})
\end{equation}
is closed isometric $(\|U_H x\| = \|x\|)$, and $(I - U_H)^{-1}$ exists. 
\end{theorem}


We have, moreover, 
\begin{equation}\label{eq:VII.4.5}
H = i (I + U _{H}) (I - U _{H}) ^{-1}.
\end{equation}

Thus, in particular, $D(H) = R(I - U_H)$ is dense in $X$. 

\begin{definition}\label{def:VII.4.Definition.3} $U_{H}$ is called the Cayley transform of $H$. 
\end{definition}

\begin{proof}[Proof of Theorem~\ref{thm:VII.4.Theorem.1}.]
We have, for any $x \in D(H)$,
\[
((H \pm i I) x, (H \pm i I) x) = (H x, H x) \pm (H x, i x) \pm (i x, H x) + (x, x).
\]

The symmetry condition for $H$ implies $(Hx, ix) = -i(Hx, x) = -i(x, Hx) = -(ix, Hx)$ and so 
\begin{equation}\label{eq:VII.4.6}
\left\| (H \pm i I) x \right\| ^{2} = \left\| H x \right\| ^{2} + \left\| x \right\| ^{2}.
\end{equation}

Hence $(H + iI)x = 0$ implies $x = 0$ and so the inverse $(H + iI)^{-1}$ exists. Since $\|(H + iI)x\| \geqq \|x\|$, the inverse $(H + iI)^{-1}$ is continuous. By \eqref{eq:VII.4.6}, it is clear that $\|U_Hy\| = \|y\|$, i.e., $U_H$ is isometric. 

$U_{H}$ is closed. For, let $(H + iI)x_{n} = y_{n}\rightarrow y$ and $(H - iI)x_{n} = z_{n}\rightarrow z$ as $n\to \infty$. Then we have, by \eqref{eq:VII.4.6}, $\| y_n - y_m\|^2 = \| H(x_n - x_m)\|^2 +\| x_n - x_m\|^2$, and so $(x_{n} - x_{m})\to 0$, $H(x_{n} - x_{m})\to 0$ as $n,m\to \infty$. Since $H$ is closed, we must have $x = s\text{- }\lim\limits_{n\to \infty}x_n\in D(H)$ and $s\text{- }\lim\limits_{n\to \infty}Hx_n = Hx$. Thus $(H + iI)x_{n}\to y = (H + iI)x,(H - iI)x_{n}\to z = (H - iI)x$ and so $U_Hy = z$. This proves that $U_{H}$ is closed. 

From $y = (H + iI)x$ and $U_Hy = (H - iI)x$, we obtain $2^{-1}(I - U_H)y = ix$ and $2^{-1}(I + U_H)y = Hx$. Thus $(I - U_H)y = 0$ implies $x = 0$ and so $(I + U_H)y = 2Hx = 0$ which implies $y = 2^{-1}((I - U_H)y + (I + U_H)y) = 0$. Therefore the inverse $(I - U_H)^{-1}$ exists. By the same calculation as above, we obtain 
\[
\begin{gathered}
H x = 2 ^{-1} (I + U _{H}) y = i (I + U _{H}) (I - U _{H}) ^{-1} x, \text{ that is, } H = i (I + U _{H}) (I - U _{H}) ^{-1}.
\end{gathered}
\]
\end{proof}

\begin{theorem}[J. \textsc{von} \textsc{Neumann} {\cite{neumann228}}]\label{thm:VII.4.Theorem.2} Let $U$ be a closed isometric operator such that $R(I - U)^{a} = X$. Then there exists a uniquely determined closed symmetric operator $H$ whose Cayley transform is $U$. 
\end{theorem}

\begin{proof}
 We first show that the inverse $(I-U)^{-1}$ exists. Suppose that $(I-U)y=0$. For any $z=(I-U)w\in R(I-U)$, we have, by the isometric property of $U$, $(y,w)=(Uy,Uw)$ as in Section~\hyperref[sec:VII.1]{$1$}. Hence 
\[
(y, z) = (y, w) - (y, U w) = (U y, U w) - (y, U w) = (U y - y, U w) = 0. \]

Hence, by the condition $R(I - U)^a = X$, $y$ must be $= 0$. Thus $(I - U)^{-1}$ exists. Put $H = i(I + U)(I - U)^{-1}$. Then $D(H) = D((I - U)^{-1}) =$ 
$R(I - U)$ is dense in $X$. We first prove that $H$ is symmetric. Let $x, y \in D(H) = R(I - U)$ and put $x = (I - U)u, y = (I - U)w$. Then $(Uu, Uw) = (u, w)$ implies that 
\[
\begin{aligned}
&(H x, y)  = (i (I + U) u, (I - U) w) = i ((U u, w) - (u, U w)) = ((I - U) u, i (I + U) w) = (x, H y).
\end{aligned}
\]

The proof of $U_H = U$ is obtained as follows. For $x = (I - U)u$, we have $Hx = i(I + U)u$ and so $(H + iI)x = 2iu$, $(H - iI)x = 2iUu$. Thus $D(U_H) = \{2iu; u \in D(U)\} = D(U)$, and $U_H(2iu) = 2i \cdot Uu = U(2iu)$. Hence $U_H = U$. 


To complete the proof of Theorem~\ref{thm:VII.4.Theorem.2}, we show that $H$ is a closed operator. In fact, $H$ is the operator which maps $(I - U) u$ onto $i(I + U) u$. If $(I - U) u_n$ and $i(I + U) u_n$ both converge as $n \to \infty$, then $u_n$ and $U u_n$ both converge as $n \to \infty$. Hence by the closure property of $U$, we must have 
\[
u _{n} \rightarrow u, (I - U) u _{n} \rightarrow (I - U) u, i (I + U) u _{n} \rightarrow i (I + U) u.
\]

This proves that $H$ is a closed operator.
\end{proof}

For the structure of the adjoint operator of a symmetric operator, we have 

\begin{theorem}[J. von Neumann {\cite{neumann228}}]\label{thm:VII.4.Theorem.3} Let $H$ be a closed symmetric operator in a Hilbert space $X$. For the Cayley transform $U_{H} = (H - iI)(H + iI)^{-1}$ of $H$, we set 
\begin{equation}\label{eq:VII.4.7}
X _{H} ^{+} = D (U _{H}) ^{\perp}, X _{H} ^{-} = R (U _{H}) ^{\perp}.
\end{equation}
\end{theorem}

Then we have 
\begin{equation}\label{eq:VII.4.8}
X _{H} ^{+} = \{x \in X; H ^{*} x = i x \}, X _{H} ^{-} = \{x \in X; H ^{*} x = - i x \},
\end{equation}
and the element $x$ of $D(H^{*})$ is uniquely expressed as 
\begin{equation}x = x _{0} + x _{1} + x _{2}, \text{ where } x _{0} \in D (H), x _{1} \in X _{H} ^{+}, x _{2} \in X _{H} ^{-} \text{ so that } H ^{*} x = H x _{0} + i x _{1} + (- i x _{2}). \label{eq:VII.4.9}\end{equation}

\begin{proof}
 $x \in D(U_H)^\perp = D((H + iI)^{-1})^\perp$ implies $(x, (H + iI)y) = 0$ for all $y \in D(H)$. Hence $(x, Hy) = -(x, iy) = (ix, y)$ and so $x \in D(H^*)$, $H^*x = ix$. The last condition implies $(x, (H + iI)y) = 0$ for all $y \in D(H)$, i.e. $x \in D((H + iI)^{-1})^\perp = D(U_H)^\perp$. This proves the first half of \eqref{eq:VII.4.8}; the latter half may be proved similarly. 

Since $U_{H}$ is a closed isometric operator, we see that $D(U_{H})$ and $R(U_{H})$ are closed linear subspaces of $X$. Hence any element $x \in X$ is uniquely decomposed as the sum of an element of $D(U_{H})$ and an element of $D(U_{H})^{\perp}$. If we apply this orthogonal decomposition to the element $(H^{*} + iI)x$, we obtain 
\[
(H ^{*} + i I) x = (H + i I) x _{0} + x ^{\prime} \quad \text{ where } \quad x _{0} \in D (H), x ^{\prime} \in D (U _{H}) ^{\perp}.
\]

But we have $(H + iI)x_0 = (H^* +iI)x_0$ by $x_0\in D(H)$ and $H\subseteq H^{*}$. We have also $H^{*}x^{\prime} = ix^{\prime}$ by $x^{\prime}\in D(U_{H})^{\perp}$ and \eqref{eq:VII.4.8}. Thus 
\[
x ^{\prime} = (H ^{*} + i I) x _{1}, x _{1} = (2 i) ^{-1} x ^{\prime} \in D (U _{H}) ^{\perp},
\]
and so 
\begin{equation}(H ^{*} + i I) x = (H ^{*} + i I) x _{0} + (H ^{*} + i I) x _{1} \text{ where } x _{0} \in D (H), x _{1} \in D (U _{H}) ^{\perp} \notag\end{equation}

Therefore $(x - x_0 - x_1) \in R(U_H)^\perp$ by $H^*(x - x_0 - x_1) = -i(x - x_0 - x_1)$ and \eqref{eq:VII.4.8}. This proves \eqref{eq:VII.4.9}. The uniqueness of the representation \eqref{eq:VII.4.9} is proved as follows. Let $0 = x_0 + x_1 + x_2$ with $x_0 \in D(H)$, $x_1 \in D(U_H)^\perp$, $x_2 \in R(U_H)^\perp$. Then, by $H^* x_0 = Hx_0$, $H^* x_1 = ix_1$, $H^* x_2 = -ix_2$, 
\[
0 = \left(H ^{*} + i I\right) 0 = \left(H ^{*} + i I\right) \left(x _{0} + x _{1} + x _{2}\right) = (H + i I) x _{0} + 2 i x _{1}.
\]

But by the uniqueness of the orthogonal decomposition of $X$ as the sum of $D(U_{H})$ and $D(U_{H})^{\perp}$, we obtain $(H + iI)x_{0} = 0$, $2ix_{1} = 0$. Since the inverse $(H + iI)^{-1}$ exists, we must have $x_{0} = 0$ and so $x_{2} = 0 - x_{0} - x_{1} = 0 - 0 - 0 = 0$. 
\end{proof}

\begin{corollaryu} $A$ closed symmetric operator $H$ in a Hilbert space $X$ is self-adjoint iff its Cayley transform $U_{H}$ is unitary. 
\end{corollaryu}

\begin{proof}
 The condition $D(H) = D(H^{*})$ is equivalent to the condition $D(U_{H})^{\perp} = R(U_{H})^{\perp} = \{0\}$. The last condition in turn is equivalent to the condition that $U_{H}$ is unitary, i.e. the condition that $U_{H}$ maps $X$ onto $X$ one-one and isometrically. 
\end{proof}


\yossection{5}{The Closed Range Theorem}
\label{sec:VII.5}

The closed range theorem of S. \textsc{Banach} \cite{banach9} reads as follows. 

\begin{theoremu} Let $X$ and $Y$ be $B$-spaces, and $T$ a closed linear operator defined in $X$ into $Y$ such that $D(T)^{a} = X$. Then the following propositions are all equivalent: 
\begin{equation}\label{eq:VII.5.1}
R(T) \text{ is closed in } Y,
\end{equation}
\end{theoremu}

\begin{align}
&\quad R \left(T ^{\prime}\right) \text{ is closed in } X ^{\prime}, \label{eq:VII.5.2} \\
R (T)& = N \left(T ^{\prime}\right) ^{\perp} = \{y \in Y; \langle y, y ^{*} \rangle = 0 \text{ for all } y ^{*} \in N \left(T ^{\prime}\right) \}, \label{eq:VII.5.3} \\
R (T ^{\prime})& = N (T) ^{\perp} = \left\{x ^{*} \in X ^{\prime}; \langle x, x ^{*} \rangle = 0 \text{ for all } x \in N (T) \right\}. \label{eq:VII.5.4}
\end{align}

\begin{proof}
 The proof of this theorem requires five steps. 


The first step. The proof of the equivalence $(1) \leftrightarrow \eqref{eq:VII.5.2}$ is reduced to the equivalence $(1) \leftrightarrow \eqref{eq:VII.5.2}$ for the special case when $T$ is a continuous linear operator such that $D(T) = X$. 

The graph $G = G(T)$ of $T$ is a closed linear subspace of $X \times Y$, and so $G$ is a $B$-space by the norm $\|\{x, y\}\| = \|x\| + \|y\|$ of $X \times Y$. Consider a continuous linear operator $S$ on $G$ into $Y$: 
\[
S \{x, T x \} = T x.
\]

Then the dual operator $S'$ of $S$ is a continuous linear operator on $Y'$ into $G'$, and we have 
\[
\begin{aligned}
&\langle \{x, T x \}, S ^{\prime} y ^{*} \rangle  = \langle S \{x, T x \}, y ^{*} \rangle = \langle T x, y ^{*} \rangle = \langle \{x, T x \}, \{0, y ^{*} \} \rangle , x \in D (T), y ^{*} \in Y ^{\prime}.
\end{aligned}
\]

Thus the functional $S' \cdot y^* - \{0, y^*\} \in (X \times Y)' = X' \times Y'$ vanishes at every point of $G$. But, the condition $\langle \{x, Tx\}, \{x^*, y_1^*\} \rangle = 0, x \in D(T)$, is equivalent to the condition $\langle x, x^* \rangle = \langle -Tx, y_1^* \rangle, x \in D(T)$, that is, to the condition $-T' y_1^* = x^*$. Hence 
\[
S ^{\prime} \cdot y ^{*} = \{0, y ^{*} \} + \{- T ^{\prime} y _{1} ^{*}, y _{1} ^{*} \} = \{- T ^{\prime} y _{1} ^{*}, y ^{*} + y _{1} ^{*} \}, y ^{*} \in Y ^{\prime}.
\]

By the arbitrariness of $y^*$, we see that $R(S') = R(-T') \times Y' = R(T') \times Y'$. Therefore $R(S')$ is closed in $X' \times Y'$ iff $R(T')$ is closed in $X'$, and, since $R(S) = R(T)$, $R(S)$ is closed in $Y$ iff $R(T)$ is closed in $Y$. Hence we have only to prove the equivalence \eqref{eq:VII.5.1} $\leftrightarrow$ \eqref{eq:VII.5.2} in the special case of a bounded linear operator $S$, instead of the original $T$. 

The second step. Let $X$ and $Y$ be $B$-spaces, and $T$ a bounded linear operator on $X$ into $Y$. Then $(1) \rightarrow \eqref{eq:VII.5.2}$. 

We consider $T$ as a bounded linear operator $T_{1}$ on $X$ into the $B$-space $Y_{1} = R(T)^{a} = R(T)$. We have to prove that \eqref{eq:VII.5.2} is true. $T_{1}'y_{1}^{*}, y_{1}^{*} \in Y_{1}'$, is defined by 
\[
\langle T _{1} x, y _{1} ^{*} \rangle = \langle T x, y _{1} ^{*} \rangle = \langle x, T _{1} ^{\prime} y _{1} ^{*} \rangle , x \in D (T _{1}) = D (T) = X.
\]

By the Hahn-Banach theorem, the functional $y_1^*$ can be extended to a $y^* \in Y'$ in such a way that $\langle Tx, y_1^* \rangle = \langle Tx, y^* \rangle$, $x \in D(T) = X$. Hence $T_1' y_1^* = T' y_1^*$ and so $R(T_1') = R(T')$. Thus it suffices to assume that $R(T) = Y$. Then, by the open mapping theorem in Chapter~\ref{chap:II}, \ref{sec:II.5}, there exists a $c > 0$ such that for each $y \in Y$, there exists an $x \in X$ with $Tx = y$, $\|x\| \leqq c \|y\|$. Thus, for each $y^*$ in $D(T')$, we have 
\[
| \langle y, y ^{*} \rangle | = | \langle T x, y ^{*} \rangle | = | \langle x, T ^{\prime} y ^{*} \rangle |
\]

Hence 
\begin{align}
&\leqq \left\| x \right\| \cdot \left\| T ^{\prime} y ^{*} \right\| \leqq c \left\| y \right\| \cdot \left\| T ^{\prime} y ^{*} \right\|. \notag \\
\left\| y ^{*} \right\|& = \sup\limits _{\left\| y \right\| \leqq 1} | \langle y, y ^{*} \rangle | \leqq c \| T ^{\prime} y ^{*} \| \notag
\end{align}
and so $(T')^{-1}$ exists and is continuous. Moreover, $(T')^{-1}$ is a closed linear operator as the inverse of a bounded linear operator. Hence we see that the domain $D((T')^{-1}) = R(T')$ must be closed in $X'$. 

The third step. Let $X$ and $Y$ be $B$-spaces, and $T$ a bounded linear operator on $X$ into $Y$. Then \eqref{eq:VII.5.2} $\to$ \eqref{eq:VII.5.1}. 

As in the second step, we consider $T$ as a bounded linear operator $T_{1}$ on $X$ into $Y_{1} = R(T)^{a}$. Then $T_{1}'$ has the inverse, since $T_{1}'y_{1}^{*} = 0$ implies 
\[
\langle T _{1} x, y _{1} ^{*} \rangle = \langle T x, y _{1} ^{*} \rangle = \langle x, T _{1} ^{\prime} y _{1} ^{*} \rangle = 0, x \in D (T _{1}) = D (T) = X,
\]
and so, since $R(T_1) = R(T)$ is dense in $Y_1 = R(T)^a$, $y_1^*$ must be $0$. Therefore, the condition that $R(T') = R(T_1')$ (proved above) is closed, implies that $T_1'$ is a continuous linear operator on the $B$-space $(R(T)^a)' = Y_1$ onto the $B$-space $R(T_1')$ in a one-one way. Hence, by the open mapping theorem, $(T_1')^{-1}$ is continuous. 

We then prove that $R(T)$ is closed. To this purpose, it suffices to derive a contradiction from the condition 
there exists a positive constant $\varepsilon$ such that the image $\{T _{1} x; \| x \| \leqq \varepsilon \}$ is not dense in all the spheres $\| y \| \leqq n ^{-1}   (n = 1, 2, \dots)$ of $Y _{1} = R (T) ^{a} = R (T _{1}) ^{a}.$


For, if otherwise, the proof of the open mapping theorem shows that $R(T_{1}) = R(T) = Y_{1}$. Thus we assume that there exists a sequence $\{y_{n}\} \subseteq Y_{1}$ with 
\[
\slim _{n \to \infty} y _{n} = 0, y _{n} \bar {\in} \{T _{1} x; \| x \| \leqq \varepsilon \} ^{a} (n = 1, 2, \dots).
\]

Since $\{T_{1}x;\|x\|\leqq\varepsilon\}^{a}$ is a closed convex, balanced set of the $B$-space $Y_{1}$ there exists, by Mazur's theorem in Chapter~\ref{chap:IV}, \ref{sec:IV.6}, a continuous linear functional $f_{n}$ on the $B$-space $Y_{1}$ such that 
\[
f _{n} (y _{n}) > \sup\limits _{\| x \| \leqq \varepsilon} | f _{n} (T _{1} x) | \quad (n = 1, 2, \dots).
\]

Hence $\| T_1'f_n\| < \varepsilon^{-1} \| f_n\| \| y_n\|$, and so, by $\slim\limits_{n\to \infty}y_n = 0$, $T_1'$ does not have a continuous inverse. This is a contradiction, and so $R(T)$ must be closed. 

The fourth step. We prove \eqref{eq:VII.5.1} $\rightarrow$ \eqref{eq:VII.5.3}. First, it is clear, from 
\[
\langle T x, y ^{*} \rangle = \langle x, T ^{\prime} y ^{*} \rangle , x \in D (T), y ^{*} \in D \left(T ^{\prime}\right),
\]
that $R(T) \subseteq N(T')^\perp$. We show that \eqref{eq:VII.5.1} implies $N(T')^\perp \subseteq R(T)$. Assume that there exists a $y_0 \in N(T')^\perp$ with $y_0 \bar{\in} R(T)$. Then, by the Hahn-Banach theorem, there exists a $y_0^* \in Y'$ such that $\langle y_0, y_0^* \rangle \neq 0$ and $\langle Tx, y_0^* \rangle = 0$ for all $x \in D(T)$. The latter condition implies $\langle x, T'y_0^* \rangle = 0$, $x \in D(T)$, and hence $T'y_0^* = 0$, i.e., $y_0 \bar{\in} N(T')^\perp$. This is a contradiction and so we must have $N(T')^\perp \subseteq R(T)$. 

The implication $\eqref{eq:VII.5.3} \to (1)$ is clear, since $N(T')^\perp$ is closed by virtue of the continuity in $y$ of $\langle y, y^* \rangle$. 

The fifth step. We prove \eqref{eq:VII.5.2} $\rightarrow$ \eqref{eq:VII.5.4}. The inclusion $R(T') \subseteq N(T)^\perp$ is clear as in the case of \eqref{eq:VII.5.3}. We show that \eqref{eq:VII.5.2} implies that $N(T)^\perp \subseteq R(T')$. To this purpose, let $x^* \in N(T)^\perp$, and define, for $y = Tx$, the functional $f_1(y)$ of $y$ through $f_1(y) = \langle x, x^* \rangle$. It is a one-valued function of $y$, since $Tx = Tx'$ implies $(x - x') \in N(T)$ and so, by $x^* \in N(T)^\perp$, $\langle (x - x'), x^* \rangle = 0$. Thus $f_1(y)$ is a linear functional of $y$. \eqref{eq:VII.5.2} implies \eqref{eq:VII.5.1}, and so, by the open mapping theorem applied to the operator $S$ in the first step, we may choose the solution $x$ of the equation $y = Tx$ in such a way that \slim $y = 0$ implies \slim $x = 0$. Hence $f_{1}(y) = \langle x, x^{*} \rangle$ is a continuous linear functional on $Y_{1} = R(T)$. Let $f \in Y'$ be an extension of $f_{1}$. Then 
\[
f (T x) = f _{1} (T x) = \langle x, x ^{*} \rangle .
\]

This proves that $T'f = x^*$. Hence $N(T)^\perp \subseteq R(T')$. 

That \eqref{eq:VII.5.4} implies \eqref{eq:VII.5.2} is clear, since $\langle x, x^* \rangle$ is a continuous linear functional of $x$.
\end{proof}

\begin{corollary}\label{cor:VII.5.Corollary.1} Let $X$ and $Y$ be $B$-spaces, and $T$ a closed linear operator on $D(T) \subseteq X$ into $Y$ such that $D(T)^a = X$. Then 
\begin{align}
R (T)& = Y \quad \text{ iff } \quad T ^{\prime} \text{ has a continuous inverse }, \label{eq:VII.5.5} \\
R (T ^{\prime})& = X ^{\prime} \text{ iff } T \text{ has a continuous inverse. } \label{eq:VII.5.6}
\end{align}
\end{corollary}

\begin{proof}
 Suppose that $R(T) = Y$. Then, from $\langle Tx, y^* \rangle = \langle x, T'y^* \rangle$, $x \in D(T)$ and $T'y^* = 0$, we obtain $y^* = 0$; that is, $T'$ must have the inverse $(T')^{-1}$. Since, by $R(T) = Y$ and \eqref{eq:VII.5.2}, $R(T')$ is closed, the closed graph theorem implies that $(T')^{-1}$ is continuous. Next let $T'$ admit a continuous inverse. Then $N(T') = \{0\}$ and also \eqref{eq:VII.5.2} holds since $T'$ is closed. Thus, by \eqref{eq:VII.5.3}, $R(T) = Y$. 

Suppose that $R(T') = X'$. Then, from $\langle Tx, y^* \rangle = \langle x, T'y^* \rangle$, $y^* \in D(T')$ and $Tx = 0$, we obtain $x = 0$, i.e., $T$ must have the inverse $T^{-1}$. Since, by $R(T') = X'$ and \eqref{eq:VII.5.1}, $R(T)$ is closed, the closed graph theorem implies that $T^{-1}$ must be continuous. Next let $T$ admit a continuous inverse. Then $N(T) = \{0\}$ and also \eqref{eq:VII.5.1} holds since $T$ is closed. Thus, by \eqref{eq:VII.5.4}, $R(T') = X'$. 
\end{proof}

\begin{corollary}\label{cor:VII.5.Corollary.2} Let $X$ be a Hilbert space with a scalar product $(u, v)$, and $T$ a closed linear operator with dense domain $D(T) \subseteq X$ and range $R(T) \subseteq X$. Suppose that there exists a positive constant $c$ such that 
\begin{equation}\label{eq:VII.5.7}
\operatorname{Re} (T u, u) \geqq c \| u \| ^{2} \quad \text{ for all } \quad u \in D (T).
\end{equation}
\end{corollary}

Then $R(T^{*}) = X$. 

\begin{proof}
 By Schwarz' inequality, we have 
\[
\left\| T u \right\| \cdot \left\| u \right\| \geqq \operatorname{Re} (T u, u) \geqq c \| u \| ^{2} \text{ for all } u \in D (T). \]

Hence $\|Tu\| \geqq c \|u\|, u \in D(T)$, and so $T$ admits a continuous inverse. Thus, by the preceding Corollary, $R(T') = X$. Hence $R(T^*) = R(T') = X$. 
\end{proof}

\begin{remarku} $A$ linear operator $T$ on $D(T) \subseteq X$ into $X$ is called accretive (the terminology is due to K. \textsc{Friedrichs} and T. \textsc{Kato}) if 
\begin{equation}\label{eq:VII.5.8}
\operatorname{Re} (T u, u) \geqq 0 \quad \text{ for all } \quad u \in D (T).
\end{equation}
\end{remarku}

$T$ is called dissipative (the terminology is due to R. S. \textsc{Phillips}) if $-T$ is accretive. 

\yossub{References for Chapter VII}

For a general account concerning Hilbert spaces, see M. H. \textsc{Stone} \cite{stone311}, N. I. \textsc{Achieser}-I. M. \textsc{Glazman} \cite{achieser1} and N. \textsc{Dunford}-J. \textsc{Schwartz} \cite{dunford48}. The closed range theorem is proved essentially in S. \textsc{Banach} \cite{banach9}. 
   % VII
% ==== VIII ====
\yoschapter{VIII}{Resolvent and Spectrum}
\label{chap:VIII}

Let $T$ be a linear operator whose domain $D(T)$ and range $R(T)$ both lie in the same complex linear topological space $X$. We consider the linear operator 
\[
T _{\lambda} = \lambda I - T,
\]
where $\lambda$ is a complex number and $I$ the identity operator. The distribution of the values of $\lambda$ for which $T_{\lambda}$ has an inverse and the properties of the inverse when it exists, are called the spectral theory for the operator $T$. We shall thus discuss the general theory of the inverse of $T_{\lambda}$. 


\yossection{1}{The Resolvent and Spectrum}
\label{sec:VIII.1}

\begin{definitionu} If $\lambda_0$ is such that the range $R(T_{\lambda_0})$ is dense in $X$ and $T_{\lambda_0}$ has a continuous inverse $(\lambda_0I - T)^{-1}$, we say that $\lambda_0$ is in the resolvent set $\varrho(T)$ of $T$, and we denote this inverse $(\lambda_0I - T)^{-1}$ by $R(\lambda_0; T)$ and call it the resolvent (at $\lambda_0$) of $T$. All complex numbers $\lambda$ not in $\varrho(T)$ form a set $\sigma(T)$ called the spectrum of $T$. The spectrum $\sigma(T)$ is decomposed into disjoint sets $P_{\sigma}(T)$, $C_{\sigma}(T)$ and $R_{\sigma}(T)$ with the following properties: $P_{\sigma}(T)$ is the totality of complex numbers $\lambda$ for which $T_{\lambda}$ does not have an inverse; $P_{\sigma}(T)$ is called the point spectrum of $T$. 

$C_{\sigma}(T)$ is the totality of complex numbers $\lambda$ for which $T_{\lambda}$ has a discontinuous inverse with domain dense in $X$; $C_{\sigma}(T)$ is called the continuous spectrum of $T$. 

$R_{\sigma}(T)$ is the totality of complex numbers $\lambda$ for which $T_{\lambda}$ has an inverse whose domain is not dense in $X$; $R_{\sigma}(T)$ is called the residual spectrum of $T$. 
\end{definitionu}

From these definitions and the linearity of $T$ we have the 

\begin{propositionu} $A$ necessary and sufficient condition for $\lambda_{0}\in P_{\sigma}(T)$ is that the equation $Tx=\lambda_{0}x$ has a solution $x\neq0$. In this case $\lambda_{0}$ is called an eigenvalue of $T$, and $x$ the corresponding eigenvector. The null space $N(\lambda_{0}I-T)$ of $T_{\lambda_{0}}$ is called the eigenspace of $T$ corresponding to the eigenvalue $\lambda_{0}$ of $T$. It consists of the vector $0$ and the totality of eigenvectors corresponding to $\lambda_{0}$. The dimension of the eigenspace corresponding to $\lambda_{0}$ is called the multiplicity of the eigenvalue $\lambda_{0}$. 
\end{propositionu}

\begin{theoremu} Let $X$ be a complex $B$-space, and $T$ a closed linear operator with its domain $D(T)$ and range $R(T)$ both in $X$. Then, for any $\lambda_{0} \in \varrho(T)$, the resolvent $(\lambda_{0}I - T)^{-1}$ is an everywhere defined continuous linear operator. 
\end{theoremu}

\begin{proof}
 Since $\lambda_0$ is in the resolvent set $\varrho(T), R(\lambda_0 I - T) = D((\lambda_0 I - T)^{-1})$ is dense in $X$ in such a way that there exists a positive constant $c$ for which 
\[
\left\| \left(\lambda_{0} I - T\right) x \right\| \geqq c \| x \| \text{ whenever } x \in D (T).
\]

We have to show that $R(\lambda_0 I - T) = X$. But, if $s\text{- }\lim\limits_{n \to \infty} (\lambda_0 I - T)x_n = y$ exists, then, by the above inequality, $s\text{- }\lim\limits_{n \to \infty} x_n = x$ exists, and so, by the closure property of $T$, we must have $(\lambda_0 I - T)x = y$. Hence, by the assumption that $R(\lambda_0 I - T)^a = X$, we must have $R(\lambda_0 I - T) = X$. 
\end{proof}

\begin{example}\label{ex:VIII.1.Example.1} If the space $X$ is of finite dimension, then any bounded linear operator $T$ is represented by a matrix $(t_{ij})$. It is known that the eigenvalues of $T$ are obtained as the roots of the algebraic equation, the so-called secular or characteristic equation of the matrix $(t_{ij})$: 
\begin{equation}\label{eq:VIII.1.1}
\det \left(\lambda \delta_{ij} - t _{ij}\right) = 0,
\end{equation}
where $\det(A)$ denotes the determinant of the matrix $A$. 
\end{example}

\begin{example}\label{ex:VIII.1.Example.2} Let $X = L_{2}(-\infty, \infty)$ and let $T$ be defined by 
\[
T \cdot x (t) = t x (t),
\]
that is, $D(T) = \{x(t); x(t) \text{ and } tx(t) \in L^2(-\infty, \infty)\}$ and $Tx(t) = tx(t)$ for $x(t) \in D(T)$. Then every real number $\lambda_0$ is in $C_\sigma(T)$. 
\end{example}

\begin{proof}
 The condition $(\lambda_{0}I-T)x=0$ implies $(\lambda_{0}-t)x(t)=0$ a.e., and so $x(t)=0$ a.e. Thus $(\lambda_{0}I-T)^{-1}$ exists. The domain $D((\lambda_{0}I-T)^{-1})$ comprises those $y(t)\in L^{2}(-\infty,\infty)$ which vanish identically in the neighbourhood of $t=\lambda_{0}$; the neighbourhood may vary with $y(t)$. Hence $D((\lambda_{0}I-T)^{-1})$ is dense in $L^{2}(-\infty,\infty)$. It is easy to see that the operator $(\lambda_{0}I-T)^{-1}$ is not bounded on the totality of such $y(t)$'s. 
\end{proof}

\begin{example}\label{ex:VIII.1.Example.3} Let $X$ be the Hilbert space $(l^2)$, and let $T_0$ be defined by 
\[
T _{0} (\xi_{1}, \xi_{2}, \dots) = (0, \xi_{1}, \xi_{2}, \dots).
\]
\end{example}

Then $0$ is in the residual spectrum of $T$, since $R(T_0)$ is not dense in $X$. 

\begin{example}\label{ex:VIII.1.Example.4} Let $H$ be a self-adjoint operator in a Hilbert space $X$. Then the resolvent set $\varrho(H)$ of $H$ comprises all the complex numbers $\lambda$ with $\operatorname{Im}(\lambda) \neq 0$, and the resolvent $R(\lambda; H)$ is a bounded linear operator with the estimate 
\begin{equation}\label{eq:VIII.1.2}
\left\| R (\lambda ; H) \right\| \leqq 1 / \left| \operatorname{Im} (\lambda) \right|.
\end{equation}
\end{example}

Moreover, 
\begin{equation}\label{eq:VIII.1.3}
\operatorname{Im} \left((\lambda I - H) x, x\right) = \operatorname{Im} (\lambda) \| x \| ^{2}, x \in D (H).
\end{equation}

\begin{proof}
 If $x \in D(H)$, then $(Hx, x)$ is real since $(Hx, x) = (x, Hx) = (\overline{Hx}, x)$. Therefore we have \eqref{eq:VIII.1.3}, and so, by Schwarz' inequality, 
\begin{equation}\label{eq:VIII.1.4}
\left\| (\lambda I - H) x \right\| \cdot \left\| x \right\| \geqq \left| ((\lambda I - H) x, x) \right| \geqq \left| \operatorname{Im} (\lambda) \right| \cdot \left\| x \right\| ^{2}
\end{equation}
which implies that 
\begin{equation}\label{eq:VIII.1.5}
\left\| (\lambda I - H) x \right\| \geqq \left| \operatorname{Im} (\lambda) \right| \cdot \left\| x \right\|, x \in D (H). \end{equation}

Hence the inverse $(\lambda I-H)^{-1}$ exists if $\operatorname{Im}(\lambda) \neq 0$. Moreover, the range $R(\lambda I-H)$ is dense in $X$ if $\operatorname{Im}(\lambda) \neq 0$. If otherwise, there would exist a $y \neq 0$ orthogonal to $R(\lambda I-H)$, i.e., $((\lambda I-H)x,y)=0$ for all $x \in D(H)$ and so $(x,(\overline{\lambda}I-H)y)=0$ for all $x \in D(H)$. Since the domain $D(H)$ of a self-adjoint operator $H$ is dense in $X$, we must have $(\bar{\lambda} I - H)y = 0$, that is, $Hy = \bar{\lambda} y$, contrary to the reality of the value $(Hy, y)$. 

Therefore, by the above Theorem, we see that, for any complex number $\lambda$ with $\operatorname{Im}(\lambda) \neq 0$, the resolvent $R(\lambda; H)$ is a bounded linear operator with the estimate \eqref{eq:VIII.1.2}. 
\end{proof}


\yossection{2}{The Resolvent Equation and Spectral Radius}
\label{sec:VIII.2}

\begin{theorem}\label{thm:VIII.2.Theorem.1} Let $T$ be a closed linear operator with domain and range both in a complex $B$-space $X$. Then the resolvent set $\varrho(T)$ is an open set of the complex plane. In each component (the maximal connected sets) of $\varrho(T)$, $R(\lambda; T)$ is a holomorphic function of $\lambda$. 
\end{theorem}

\begin{proof}
 By the Theorem of the preceding section, $R(\lambda; T)$ for $\lambda \in \varrho(T)$ is an everywhere defined continuous operator. Let $\lambda_0 \in \varrho(T)$ and consider 
\begin{equation}\label{eq:VIII.2.1}
S (\lambda) = R \left(\lambda_{0}; T\right) \left\{I + \sum\limits_{n=1} ^{\infty} \left(\lambda_{0} - \lambda\right) ^{n} R \left(\lambda_{0}; T\right) ^{n} \right\}. \end{equation}

The series is convergent in the operator norm whenever $\left|\lambda_{0}-\lambda\right|$. $\left\|R\left(\lambda_{0};T\right)\right\|<1$, and within this circle of the complex plane, the series defines a holomorphic function of $\lambda$. Multiplication by $(\lambda I-T)=(\lambda-\lambda_{0})I+(\lambda_{0}I-T)$ on the left or right gives $I$ so that the series $S(\lambda)$ actually represents the resolvent $R(\lambda;T)$. Thus we have proved that a circular neighbourhood of $\lambda_{0}$ belongs to $\varrho(T)$ and $R(\lambda;T)$ is holomorphic in this neighbourhood. 
\end{proof}

\begin{theorem}\label{thm:VIII.2.Theorem.2} If $\lambda$ and $\mu$ both belong to $\varrho(T)$, and if $R(\lambda; T)$ and $R(\mu; T)$ are everywhere defined continuous operators, then the resolvent equation holds: 
\begin{equation}\label{eq:VIII.2.2}
R (\lambda ; T) - R (\mu ; T) = (\mu - \lambda) R (\lambda ; T) R (\mu ; T).
\end{equation}
\end{theorem}

\begin{proof}
 We have 
\[
\begin{aligned}
&R (\lambda ; T)  = R (\lambda ; T) (\mu I - T) R (\mu ; T) = R (\lambda ; T) \{(\mu - \lambda) I + (\lambda I - T) \} R (\mu ; T) \\
&= (\mu - \lambda) R (\lambda ; T) R (\mu ; T) + R (\mu ; T).
\end{aligned}
\]
\end{proof}

\begin{theorem}\label{thm:VIII.2.Theorem.3} If $T$ is a bounded linear operator on a complex $B$-space $X$ into $X$, then the following limit exists: 
\begin{equation}\label{eq:VIII.2.3}
\lim\limits _{n \rightarrow \infty} \| T ^{n} \| ^{1/n} = r _{\sigma} (T).
\end{equation}
\end{theorem}

It is called the spectral radius of $T$, and we have 
\begin{equation}\label{eq:VIII.2.4}
\boldsymbol {r} _{\sigma} (T) \leqq \| T \|.
\end{equation}

If $|\lambda| > r_{\sigma}(T)$, then the resolvent $R(\lambda; T)$ exists and is given by the series 
\begin{equation}\label{eq:VIII.2.5}
R (\lambda ; T) = \sum\limits_{n=1} ^{\infty} \lambda^{-n} T ^{n-1}
\end{equation}
which converges in the norm of operators. 

\begin{proof}
 Set $r = \inf\limits_{n \geqq 1} \|T^n\|^{1/n}$. We have to show that $\varlimsup\limits_{n \to \infty} \|T^n\|^{1/n} \leqq r$. For any $\varepsilon > 0$, choose $m$ such that $\|T^m\|^{1/m} \leqq r + \varepsilon$. For arbitrary $n$, write $n = p m + q$ where $0 \leqq q \leqq (m - 1)$. Then, by $\|AB\| \leqq \|A\| \cdot \|B\|$, we obtain 
\[
\left\| T ^{n} \right\| ^{1/n} \leqq \left\| T ^{m} \right\| ^{p/n} \cdot \left\| T \right\| ^{q/n} \leqq (r + \varepsilon) ^{mp/n} \left\| T \right\| ^{q/n}. \]

Since $pm / n \to 1$ and $q / n \to 0$ as $n \to \infty$, we must have $\varlimsup\limits_{n \to \infty} \|T^n\|^{1/n} \leqq r + \varepsilon$. Since $\varepsilon$ was arbitrary, we have proved $\varlimsup\limits_{n \to \infty} \|T^n\|^{1/n} \leqq r$. 

Since $\|T^n\| \leqq \|T\|^n$, we have $\lim\limits_{n \to \infty} \|T^n\|^{1/n} \leqq \|T\|$. The series \eqref{eq:VIII.2.5} is convergent in the norm of operators when $|\lambda| > r_\sigma(T)$. For, if $|\lambda| \geqq r_\sigma(T) + \varepsilon$ with $\varepsilon > 0$, then, by \eqref{eq:VIII.2.3}, $\|\lambda^{-n}T^n\| \leqq (r_\sigma(T) + \varepsilon)^{-n} \cdot (r_\sigma(T) + 2^{-1}\varepsilon)^n$ for large $n$. Multiplication by $(\lambda I - T)$ on the left or right of this series gives $I$ so that the series actually represents the resolvent $R(\lambda; T)$. 
\end{proof}

\begin{corollaryu} The resolvent set $\varrho(T)$ is not empty when $T$ is a bounded linear operator. 
\end{corollaryu}

\begin{theorem}\label{thm:VIII.2.Theorem.4} For a bounded linear operator $T \in L(X, X)$, we have 
\begin{equation}\label{eq:VIII.2.6}
r _{\sigma} (T) = \sup\limits _{\lambda \in \sigma (T)} | \lambda |.
\end{equation}
\end{theorem}

\begin{proof}
 By Theorem~\ref{thm:VIII.2.Theorem.3}, we know that $r_{\sigma}(T) \geqq \sup\limits_{\lambda \in \sigma(T)} |\lambda|$. Hence we have only to show that $r_{\sigma}(T) \leqq \sup\limits_{\lambda \in \sigma(T)} |\lambda|$. 

By Theorem~\ref{thm:VIII.2.Theorem.1}, $R(\lambda; T)$ is holomorphic in $\lambda$ when $|\lambda| > \sup\limits_{\lambda \in \sigma(T)} |\lambda|$. Thus it admits a uniquely determined Laurent expansion in positive and non-positive powers of $\lambda$ convergent in the operator norm for $|\lambda| > \sup |\lambda|$. By Theorem~\ref{thm:VIII.2.Theorem.3}, this Laurent series must coincide with $\sum\limits_{n=1}^{\infty} \lambda^{-n} T^{n-1}$. Hence $\lim\limits_{n \to \infty} \|\lambda^{-n} T^n\| = 0$ if $|\lambda| > \sup\limits_{\lambda \in \sigma(T)} |\lambda|$, and so, for any $\varepsilon > 0$, we must have $\|T^n\| \leqq (\varepsilon + \sup\limits_{\lambda \in \sigma(T)} |\lambda|)^n$ for large $n$. This proves that 
\[
r _{\sigma} (T) = \lim\limits _{n \rightarrow \infty} \| T ^{n} \| ^{1/n} \leqq \sup\limits _{\lambda \in \sigma (T)} | \lambda |.
\]
\end{proof}

\begin{corollaryu} The series $\sum\limits_{n=1}^{\infty} \lambda^{-n} T^{n-1}$ diverges if $|\lambda| < r_{\sigma}(T)$. 
\end{corollaryu}

\begin{proof}
 Let $r$ be the smallest number $\geqq 0$ such that the series $\sum\limits_{n=1}^{\infty} \lambda^{-n} T^{n-1}$ converges in the operator norm for $|\lambda| > r$. The existence of such an $r$ is proved as for ordinary power series in $\lambda^{-1}$. Then, for $|\lambda| > r$, $\lim\limits_{n \to \infty} \|\lambda^{-n} T^n\| = 0$ and so, as in the proof of $r_\sigma(T) \leqq \sup\limits_{\lambda \in \sigma(T)} |\lambda|$, we must have $\lim\limits_{n \to \infty} \|T^n\|^{1/n} \leqq r$. This proves that $r_\sigma(T) \leqq r$. 
\end{proof}


\yossection{3}{The Mean Ergodic Theorem}
\label{sec:VIII.3}

For a particular class of continuous linear operators, the mean ergodic theorem gives a method for obtaining the eigenspace corresponding to the eigenvalue $1$. In this section, we shall state and prove the mean ergodic theorem from the view point of the spectral theory, as was formulated previously by the present author. The historical sketch of the ergodic theory in connection with statistical mechanics will be given in Chapter~\ref{chap:XIII}. 

\begin{theorem}\label{thm:VIII.3.Theorem.1} Let $X$ be a locally convex linear topological space, and $T$ a continuous linear operator on $X$ into $X$. We assume that 
\begin{equation}\label{eq:VIII.3.1}
\parbox{0.92\linewidth}{\everymath{}\let\limits\nolimits the family of operators $\{T^{n}; n=1,2,\ldots\}$ is equicontinuous in the sense that, for any continuous semi-norm $q$ on $X$, there exists a continuous semi-norm $q'$ on $X$ such that $\sup\limits_{n\geqq1}q(T^{n}x)\leqq q'(x)$ for all $x\in X$.}
\end{equation}
Then the closure $R(I - T)^a$ of the range $R(I - T)$ satisfies 
\begin{equation}\label{eq:VIII.3.2}
R (I - T) ^{a} = \left\{x \in X; \lim\limits _{n \rightarrow \infty} T _{n} x = 0, T _{n} = n ^{-1} \sum\limits_{m=1} ^{n} T ^{m} \right\},
\end{equation}
and so, in particular, 
\begin{equation}\label{eq:VIII.3.3}
R (I - T) ^{a} \cap N (I - T) = \{0 \}.
\end{equation}
\end{theorem}

\begin{proof}
 We have $T_{n}(I - T) = n^{-1}(T - T^{n+1})$. Hence, by \eqref{eq:VIII.3.1}, $w \in R(I - T)$ implies that $\lim\limits_{n\to \infty}T_nw = 0$. Next let $z \in R(I - T)^a$. Then, for any continuous semi-norm $q'$ on $X$ and $\varepsilon > 0$, there exists a $w \in R(I - T)$ such that $q'(z - w) < \varepsilon$. Thus, by \eqref{eq:VIII.3.1}, we have $q(T_{n}(z - w)) \leqq n^{-1}\sum\limits_{m=1}^{n}q(T^{m}(z - w)) \leqq q'(z - w) < \varepsilon$. Hence $q(T_{n}z) \leqq q(T_{n}w) + q(T_{n}(z - w)) \leqq q(T_{n}w) + \varepsilon$, and so $\lim\limits_{n\to \infty}T_nz = 0$. This proves that $R(I - T)^a \subseteq \{x \in X; \lim\limits_{n\to \infty}T_nx = 0\}$. 

Let, conversely, $\lim\limits_{n\to \infty}T_nx = 0$. Then, for any continuous semi-norm $q$ on $X$ and $\varepsilon >0$, there exists an $n$ such that $q(x - (x - T_nx)) = q(T_nx) < \varepsilon$. But, by 
\[
\begin{aligned}
&x - T _{n} x  = n ^{-1} \sum\limits_{m=1} ^{n} (I - T ^{m}) x = n ^{-1} \sum\limits_{m=1} ^{n} (I - T) (I + T + T ^{2} + \dots + T ^{m-1}) x,
\end{aligned}
\]
$(x - T_{n}x)\in R(I - T)$. Hence $x$ must belong to $R(I - T)^{a}$ 
\end{proof}

\begin{theorem}[the mean ergodic theorem]\label{thm:VIII.3.Theorem.2} Let condition \eqref{eq:VIII.3.1} be satisfied. Let, for a given $x \in X$, there exist a subsequence $\{n'\}$ of $\{n\}$ 

such that 
\begin{equation}\label{eq:VIII.3.4}
\text{ weak-lim } _ {n ^{\prime} \to \infty} T _ {n ^{\prime}} x = x _{0} \text{ exists. }
\end{equation}
\end{theorem}

Then $Tx_0 = x_0$ and $\lim\limits_{n \to \infty} T_n x = x_0$. 

\begin{proof}
 We have $TT_{n} - T_{n} = n^{-1}(T^{n+1} - T)$, and so, by \eqref{eq:VIII.3.1}, $\lim\limits_{n\to \infty}(TT_nx - T_nx) = 0$. Thus, for any $f\in X'$, $\lim\limits_{n\to \infty}\langle TT_{n'}x,f\rangle = \lim\limits_{n\to \infty}\langle T_{n'}x,T'f\rangle$ exists and $= \lim\limits_{n\to \infty}\langle T_{n'}x,f\rangle = \langle x_0,f\rangle$. Therefore $\langle x_0,f\rangle = \langle Tx_0,f\rangle$ and so, by the arbitrariness of $f\in X'$, we must have $Tx_0 = x_0$. 

We have thus $T^{m}x = T^{m}x_{0} + T^{m}(x - x_{0}) = x_{0} + T^{m}(x - x_{0})$ and so $T_{n}x = x_{0} + T_{n}(x - x_{0})$. But, $(x - x_{0}) = \text{ weak-lim }_{n\to \infty}(x - T_{n'}x)$ and, as proved above, $(x - T_{n'}x)\in R(I - T)$. Therefore, by Theorem~\ref{thm:V.1.Theorem.11} in Chapter~\ref{chap:V}, \ref{sec:V.1}, $(x - x_0)\in R(I - T)^a$. Thus, by Theorem~\ref{thm:VIII.3.Theorem.1}, $\lim\limits_{n\to \infty}T_n(x - x_0) = 0$ and so we have proved that $\lim\limits_{n\to \infty}T_nx = x_0$. 
\end{proof}

\begin{corollaryu} Let condition \eqref{eq:VIII.3.1} be satisfied, and $X$ be locally sequentially weakly compact. Then, for any $x \in X$, $\lim\limits_{n \to \infty} T_n x = x_0$ exists, and the operator $T_0$ defined by $T_0 x = x_0$ is a continuous linear operator such that 
\begin{align}
T _{0}& = T _{0} ^{2} = T T _{0} = T _{0} T, \label{eq:VIII.3.5} \\
R \left(T _{0}\right)& = N (I - T), \label{eq:VIII.3.6} \\
N (T _{0})& = R (I - T) ^{a} = R (I - T _{0}). \label{eq:VIII.3.7}
\end{align}
\end{corollaryu}

Moreover, we have the direct sum decomposition 
\begin{equation}\label{eq:VIII.3.8}
X = R (I - T) ^{a} \oplus N (I - T),
\end{equation}
i.e., any $x \in X$ is represented uniquely as the sum of an element $\in R(I - T)^a$ and an element $\in N(I - T)$. 

\begin{proof}
 The linearity of $T_0$ is clear. The continuity of $T_0$ is proved by the equi-continuity of $\{T_n\}$ implied by \eqref{eq:VIII.3.1}. Next, since $Tx_0 = x_0$, we have $TT_0 = T_0$ and so $T^n T_0 = T_0$, $T_n T_0 = T_0$ which implies that $T_0^2 = T_0$. On the other hand, $T_n - T_n T = n^{-1}(T - T^{n+1})$ and \eqref{eq:VIII.3.1} imply that $T_0 = T_0 T$. The equality \eqref{eq:VIII.3.6} is proved as follows. Let $Tx = x$, then $T^n x = x$, $T_n x = x$. and so $T_0 x = x$, that is, $x \in R(T_0)$. Let conversely, $x \in R(T_0)$. Then, by $T_0^2 = T_0$, we have $T_0 x = x$ and so, by $TT_0 = T_0$, $Tx = TT_0 x = T_0 x = x$. Therefore, the eigenspace of $T$ corresponding to the eigenvalue $1$ of $T$ is precisely the range $R(T_0)$. Hence \eqref{eq:VIII.3.6} is proved. Moreover, we have, by Theorem~\ref{thm:VIII.3.Theorem.1}, $N(T_0) = R(I - T)^a$. But, by $T_0^2 = T_0$, we have $R(I - T_0) \subseteq N(T_0)$, and if $x \in N(T_0)$, then $x = x - T_0 x \in R(I - T_0)$. Thus $N(T_0) = R(I - T_0)$. 

Therefore, by $I = (I - T_0) + T_0$ and \eqref{eq:VIII.3.6} and \eqref{eq:VIII.3.7}, we obtain \eqref{eq:VIII.3.8}. 
\end{proof}

\begin{remarku} The eigenspace $N(\lambda I - T)$ of $T$ belonging to the eigenvalue $\lambda$ with $|\lambda| = 1$ may be obtained as $R(T(\lambda))$, where $T(\lambda)x = \lim\limits_{n\to \infty}n^{-1}\sum\limits_{m=1}^{n}(T / \lambda)^m x$. 
\end{remarku}

The Mean Ergodic Theorem of J. von Neumann. Let $(S, \mathfrak{B}, m)$ be a measure space, and $P$ an equi-measure transformation of $S$, that is, $P$ is a one-one mapping of $S$ onto $S$ such that $P \cdot B \in B$ iff $B$ is $\in B$ and $m(P \cdot B) = m(B)$. Consider the linear operator $T$ on $L^{2}(S, \mathfrak{B}, m)$ onto itself defined by 
\begin{equation}\label{eq:VIII.3.9}
(T x) (s) = x (P s), x \in L ^{2} (S, \mathfrak {B}, m).
\end{equation}

By the equi-measurable property of $P$, we easily see that the operator $T$ is unitary and so the equi-continuity condition \eqref{eq:VIII.3.1} is surely satisfied by $\left\|T^{n}\right\|=1(n=1,2,\ldots)$. Therefore, by the sequential weak compactness of the Hilbert space $L^{2}(S,\mathfrak{B},m)$, we obtain the mean ergodic theorem of J. von Neumann: 
\begin{equation}\label{eq:VIII.3.10}
\begin{aligned}
&\text{ For any } x \in L ^{2} (S, \mathfrak {B}, m), \slim _{n \to \infty} n ^{-1} \sum\limits_{m=1} ^{n} T ^{m} x = x _{0} \in L ^{2} (S, \mathfrak {B}, m) \text{ exists and } T x _{0} = x _{0}.
\end{aligned}
\end{equation}

\begin{remarku} Theorem~\ref{thm:VIII.3.Theorem.1} and Theorem~\ref{thm:VIII.3.Theorem.2} are adapted from K. \textsc{Yosida} \cite{yosida346}. Cf. also S. \textsc{Kakutani} \cite{kakutani125} and F. \textsc{Riesz} \cite{rickart275}. Neumann's mean ergodic theorem was published in J. \textsc{von} \textsc{Neumann} \cite{neumann230}. 
\end{remarku}


\yossection{4}{Ergodic Theorems of the Hille Type Concerning Pseudo-resolvents}
\label{sec:VIII.4}

The notion of resolvent is generalized to that of pseudo-resolvent by E. \textsc{Hille}. We can prove ergodic theorems for pseudo-resolvents by a similar idea to that used in the proof of the mean ergodic theorems in the preceding section. See K. \textsc{Yosida} \cite{yosida347}. Cf. T. \textsc{Kato} \cite{kato134}. These ergodic theorems may be considered as extensions of the abelian ergodic theorems of E. \textsc{Hille} given in E. \textsc{Hille}-R. S. \textsc{Phillips} \cite{hille97}, p. $502$. 

We shall begin with the definition of the pseudo-resolvent. 

\begin{definitionu} Let $X$ be a locally convex complex linear topological space, and $L(X, X)$ the algebra of all continuous linear operators defined on $X$ into $X$. A pseudo-resolvent $J_{\lambda}$ is a function defined on a subset $D(J)$ of the complex $\lambda$-plane with values in $L(X, X)$ such that 
\begin{equation}\label{eq:VIII.4.1}
J _{\lambda} - J _{\mu} = (\mu - \lambda) J _{\lambda} J _{\mu} (\text{ the resolvent equation }).
\end{equation}
\end{definitionu}

\begin{propositionu} All $J_{\lambda}, \lambda \in D(J)$, have a common null space, which we denote by $N(J)$, and a common range which we denote by $R(J)$. Similarly, all $(I - \lambda J_{\lambda}), \lambda \in D(J)$, have a common null space, which we denote by $N(I - J)$, and a common range which we denote by $R(I - J)$. Moreover, we have the commutativity: 
\begin{equation}\label{eq:VIII.4.2}
J _{\lambda} J _{\mu} = J _{\mu} J _{\lambda} (\lambda , \mu \in D (J)).
\end{equation}
\end{propositionu}

\begin{proof}
 By interchanging $\lambda$ and $\mu$ in \eqref{eq:VIII.4.1}, we obtain 
\[
J _{\mu} - J _{\lambda} = (\lambda - \mu) J _{\mu} J _{\lambda} = - (\mu - \lambda) J _{\mu} J _{\lambda},
\]
and hence \eqref{eq:VIII.4.2} is true. The first part of the Proposition is clear from \eqref{eq:VIII.4.1} and \eqref{eq:VIII.4.2}. The second part is also clear from 
\begin{equation}\tag{1'}\label{eq:VIII.4.1'}
(I - \lambda J _{\lambda}) = (I - (\lambda - \mu) J _{\lambda}) (I - \mu J _{\mu})
\end{equation}
which is a variant of \eqref{eq:VIII.4.1}. 
\end{proof}

\begin{theorem}\label{thm:VIII.4.Theorem.1} $A$ pseudo-resolvent $J_{\lambda}$ is a resolvent of a linear operator $A$ iff $N(J) = \{0\}$; and then $R(J)$ coincides with the domain $D(A)$ of $A$. 
\end{theorem}

\begin{proof}
 The ``only if'' part is clear. Suppose $N(J) = \{0\}$. Then, for any $\lambda \in D(J)$, the inverse $J_{\lambda}^{-1}$ exists. We have 
\begin{equation}\label{eq:VIII.4.3}
\lambda I - J _{\lambda} ^{-1} = \mu I - J _{\mu} ^{-1} (\lambda , \mu \in D (J)). \end{equation}

For, by \eqref{eq:VIII.4.1} and \eqref{eq:VIII.4.2}, 
\[
\begin{aligned} J _{\lambda} J _{\mu} (\lambda I - J _{\lambda} ^{-1} - \mu I + J _{\mu} ^{-1}) & = (\lambda - \mu) J _{\lambda} J _{\mu} - J _{\lambda} J _{\mu} (J _{\lambda} ^{-1} - J _{\mu} ^{-1}) \\ & = (\lambda - \mu) J _{\lambda} J _{\mu} - (J _{\mu} - J _{\lambda}) = 0. \end{aligned}
\]

We thus put 
\begin{equation}\label{eq:VIII.4.4}
A = (\lambda I - J _{\lambda} ^{-1}).
\end{equation}

Then $J_{\lambda} = (\lambda I - A)^{-1}$ for $\lambda \in D(J)$. 
\end{proof}
\begin{lemma}\label{lem:VIII.4.Lemma.1} We assume that there exists a sequence $\{\lambda_n\}$ of numbers $\in D(J)$ such that 
\begin{equation}
\label{eq:VIII.4.5}
\begin{minipage}{0.92\linewidth}
\everymath{}\let\limits\nolimits
$\lim\limits_{n\to \infty}\lambda_n = 0$ and the family of operators $\{\lambda_nJ_{\lambda_n}\}$ is equi-continuous.
\end{minipage}
\end{equation}
\end{lemma}


Then we have 
\begin{equation}\label{eq:VIII.4.6}
R (I - J) ^{a} = \left\{x \in X; \lim\limits _{n \rightarrow \infty} \lambda_{n} J _ {\lambda_{n}} x = 0 \right\},
\end{equation}
and hence 
\begin{equation}\label{eq:VIII.4.7}
N (I - J) \cap R (I - J) ^{a} = \{0 \}.
\end{equation}

\begin{proof}
 We have, by \eqref{eq:VIII.4.1}, 
\begin{equation}\label{eq:VIII.4.8}
\lambda J _{\lambda} (I - \mu J _{\mu}) = (1 - \mu (\mu - \lambda) ^{-1}) \lambda J _{\lambda} - \lambda (\lambda - \mu) ^{-1} \mu J _{\mu}. \end{equation}

Hence, by ($5$), the condition $x \in R(I - \mu J_{\mu}) = R(I - J)$ implies that $\lim\limits_{n \to \infty} \lambda_n J_{\lambda_n} x = 0$. Let $y \in R(I - J)^a$. Then, for any continuous semi-norm $q$ on $X$ and $\varepsilon > 0$, there exists an $x \in R(I - J)$ such that $q(y - x) < \varepsilon$. By ($5$), we have, for any continuous semi-norm $q'$ on $X$, 
\[
q ^{\prime} \left(\lambda_{n} J _ {\lambda_{n}} (y - x)\right) \leqq q (y - x) (n = 1, 2, \dots)
\]
with a suitable continuous semi-norm $q$ on $X$. Therefore, by $\lambda_{n}J_{\lambda_{n}}y = \lambda_{n}J_{\lambda_{n}}x + \lambda_{n}J_{\lambda_{n}}(y - x)$, we must have $\lim\limits_{n\to \infty}\lambda_nJ_{\lambda_n}y = 0$. 

Let, conversely, $\lim\limits_{n\to \infty}\lambda_nJ_{\lambda_n}x = 0$. Then, for any continuous semi-norm $q$ on $X$ and $\varepsilon >0$, there exists a $\lambda_{n}$ such that $q(x - (x - \lambda_nJ_{\lambda_n}x)) < \varepsilon$. Hence $x$ must belong to $R(I - \lambda_nJ_{\lambda_n})^a = R(I - J)^a$. 
\end{proof}

\begin{lemman}[Lemma~\ref{lem:VIII.4.Lemma.1}']\yosnamedtag{1'}\label{lem:VIII.4.Lemma.1prime} We assume that there exists a sequence $\{\lambda_n\}$ of numbers $\in D(J)$ such that 
$\lim\limits_{n\to\infty}\left|\lambda_{n}\right|=\infty$ and the family of operators $\{\lambda_{n}J_{\lambda_{n}}\}$ is equi-continuous. ($5$') Then we have 
\begin{equation}\tag{6'}\label{eq:VIII.4.6'}
R (J) ^{a} = \left\{x \in X; \lim\limits _{n \rightarrow \infty} \lambda_{n} J _ {\lambda_{n}} x = x \right\},
\end{equation}
and hence 
\begin{equation}\tag{7'}\label{eq:VIII.4.7'}
N (J) \cap R (J) ^{a} = \{0 \}.
\end{equation}
\end{lemman}

\begin{proof}
 We have, by \eqref{eq:VIII.4.1}, 
\[
\lambda J _{\lambda} J _{\mu} = \frac {\mu}{\lambda} \lambda J _{\lambda} J _{\mu} - \frac {1}{\lambda} \lambda J _{\lambda} + J _{\mu}. \]

Hence, by ($5$'), the condition $x \in R(J_{\mu}) = R(J)$ implies that $\lim\limits_{n \to \infty} \lambda_n J_{\lambda_n} x = x$. Let $y \in R(J)^a$. Then, for any continuous semi-norm $q$ on $X$ and $\varepsilon > 0$, there exists an $x \in R(J)$ such that $q(y - x) < \varepsilon$. By ($5$'), we have, for any continuous semi-norm $q'$ on $X$, 
\[
q ^{\prime} \left(\lambda_{n} J _ {\lambda_{n}} (y - x)\right) \leqq q (y - x) \quad (n = 1, 2, \dots)
\]
with a suitable continuous semi-norm $q$ on $X$. Therefore, by $(5')$ and 
\[
\lambda_{n} J _ {\lambda_{n}} y - y = \left(\lambda_{n} J _ {\lambda_{n}} x - x\right) + (x - y) + \lambda_{n} J _ {\lambda_{n}} (y - x),
\]
we must have $\lim\limits_{n\to \infty}\lambda_nJ_{\lambda_n}y = y.$ 

Let, conversely, $\lim\limits_{n\to \infty}\lambda_nJ_{\lambda_n}x = x$. Then, for any continuous semi-norm $q$ on $X$ and $\varepsilon >0$, there exists a $\lambda_{n}$ such that $q(x - \lambda_nJ_{\lambda_n}x) < \varepsilon$. Hence $x$ must belong to $R(J_{\lambda_n})^a = R(J)^a$. 
\end{proof}

\begin{theorem}\label{thm:VIII.4.Theorem.2} Let ($5$) be satisfied. Let, for a given $x \in X$, there exist a subsequence $\{n'\}$ of $\{n\}$ such that 
\begin{equation}\label{eq:VIII.4.9}
\underset {n \to \infty} {\text{ weak-lim }} \lambda_ {n ^{\prime}} J _ {\lambda_ {n ^{\prime}}} x = x _{h} \text{ exists. }
\end{equation}
\end{theorem}

Then $x_{h} = \lim\limits_{n\to \infty}\lambda_{n}J_{\lambda_{n}}x$ and $x_{h}\in N(I - J), x_{p} = (x - x_{h})\in R(I - J)^{a}$. 

\begin{proof}
 Setting $\mu = \lambda_{n'}$ in \eqref{eq:VIII.4.1} and letting $n' \to \infty$, we see, by ($5$), that $(I - \lambda J_{\lambda}) x = (I - \lambda J_{\lambda})(x - x_h)$, that is, $(I - \lambda J_{\lambda}) x_h = 0$. Hence $x_h \in N(I - J)$ and so 
\begin{equation}\label{eq:VIII.4.10}
\lambda_{n} J _ {\lambda_{n}} x = x _{h} + \lambda_{n} J _ {\lambda_{n}} (x - x _{h}). \end{equation}

Therefore we have only to prove that $\lim\limits_{n\to \infty}\lambda_nJ_{\lambda_n}(x - x_h) = 0$, or, by Lemma~\ref{lem:VIII.4.Lemma.1}, $(x - x_h)\in R(I - J)^a$. But $(x - \lambda_nJ_{\lambda_n}x)\in R(I - J)$, and so, by Theorem~\ref{thm:V.1.Theorem.11} in Chapter~\ref{chap:V}, \ref{sec:V.1}, we must have $(x - x_h)\in R(I - J)^a$. 
\end{proof}

\begin{corollary}\label{cor:VIII.4.Corollary.1} Let ($5$) be satisfied, and let $X$ be locally sequentially weakly compact. Then 
\begin{equation}\label{eq:VIII.4.11}
X = N (I - J) \oplus R (I - J) ^{a}.
\end{equation}
\end{corollary}

\begin{proof}
 For any $x \in X$, let $x_{h} = \lim\limits_{n \to \infty} \lambda_{n} J_{\lambda_{n}} x$ and $x_{p} = (x - x_{h})$ be the components of $x$ in $N(I - J)$ and $R(I - J)^{a}$, respectively. 
\end{proof}

\begin{theoremn}[Theorem~\ref{thm:VIII.4.Theorem.2}']\yosnamedtag{2'}\label{thm:VIII.4.Theorem.2prime} Let $(5')$ be satisfied. Let, for a given $x \in X$, there exist a subsequence $\{n'\}$ of $\{n\}$ such that 
\begin{equation}\tag{9'}\label{eq:VIII.4.9'}
\underset {n \rightarrow \infty} {\text{ weak-lim }} \lambda_ {n ^{\prime}} J _ {\lambda_ {n ^{\prime}}} x = x _ {h ^{\prime}} \text{ exists. }
\end{equation}

Then $x_{h'} = \lim\limits_{n\to \infty}\lambda_nJ_{\lambda_n}x$ and $x_{h'}\in R(J)^a, x_{p'} = (x - x_{h'})\in N(J).$
\end{theoremn}

\begin{proof}
 Setting $\mu = \lambda_{n'}$ in \eqref{eq:VIII.4.8} and letting $n' \to \infty$, we see, by ($5$'), that $\lambda J_{\lambda}(x - x_{h'}) = 0$, that is, $(x - x_{h'}) \in N(J)$. Hence 
\begin{equation}\tag{10'}\label{eq:VIII.4.10'}
\lambda_{n} J _ {\lambda_{n}} x = \lambda_{n} J _ {\lambda_{n}} x _ {h ^{\prime}}. \end{equation}

Therefore we have only to prove that $\lim\limits_{n\to \infty}\lambda_nJ_{\lambda_n}x_{h'} = x_{h'}$, or, by Lemma~\ref{lem:VIII.4.Lemma.1prime}, $x_{h'}\in R(J)^a$. But $\lambda_{n'}J_{\lambda_n},x_{h'}\in R(J)$, and so, by Theorem~\ref{thm:V.1.Theorem.11} in Chapter~\ref{chap:V}, \ref{sec:V.1}, we must have $x_{h'}\in R(J)^a$. 
\end{proof}

\begin{corollaryn}[Corollary~\ref{cor:VIII.4.Corollary.1}']\yosnamedtag{1'}\label{cor:VIII.4.Corollary.1prime} Let $(5')$ be satisfied, and let $X$ be locally sequentially weakly compact. Then 
\begin{equation}\tag{11'}\label{eq:VIII.4.11'}
X = N (J) \oplus R (J) ^{a}.
\end{equation}
\end{corollaryn}

\begin{proof}
 For any $x \in X$, let $x_{h'} = \lim\limits_{n \to \infty} \lambda_n J_{\lambda_n} x$ and $x_{p'} = (x - x_{h'})$ be the components of $X$ in $R(J)^a$ and $N(J)$, respectively. 
\end{proof}

\begin{remarku} As a Corollary we obtain: In a reflexive $B$-space $X$, a pseudo-resolvent $J_{\lambda}$ satisfying $(5')$ is the resolvent of a closed linear operator $A$ with dense domain iff $R(J)^{a} = X$. This result is due to $T$. \textsc{Kato}, loc. cit. The proof is easy, since, by Eberlein's theorem, a $B$-space $X$ is locally sequentially weakly compact iff $X$ is reflexive. 
\end{remarku}


\yossection{5}{The Mean Value of an Almost Periodic Function}
\label{sec:VIII.5}

As an application of the mean ergodic theorem we shall give an existence proof of the mean value of an almost periodic function. 

\begin{definition}\label{def:VIII.5.Definition.1} $A$ set $G$ of elements $g, h, \ldots$ is called a group if in $G$ a product (in general non-commutative) $gh$ of any pair $(g, h)$ of elements $\in G$ is defined satisfying the following conditions: 
\begin{align}
g h& \in G, \label{eq:VIII.5.1} \\
(g h) k& = g (h k) \text{ (the associativity) }, \label{eq:VIII.5.2}
\end{align}
\begin{equation}\label{eq:VIII.5.3}
\begin{minipage}[b]{0.92\linewidth}
there exists a unique element $e$ in $G$ such that $eg = ge = g$ for all $g \in G$; $e$ is called the identity element of the group $G$,
\end{minipage}
\end{equation}

\begin{equation}\label{eq:VIII.5.4}
\begin{minipage}[b]{0.92\linewidth}
for every element $g \in G$, there exists a uniquely determined element in $G$, which is denoted by $g^{-1}$, such that $gg^{-1} = g^{-1}g = e$; the element $g^{-1}$ is called the inverse element of $g$.
\end{minipage}
\end{equation}
\end{definition}


Clearly, $g$ is the inverse of $g^{-1}$ so that $(g^{-1})^{-1} = g$. A group $G$ is said to be commutative if $gh = hg$ for all $g, h \in G$. 

\begin{examplen} The totality of complex matrices of order $n$ with determinants equal to unity is a group with respect to the matrix multiplication; it is called the complex unimodular group of order $n$. The identity of this group is the identity matrix and the inverse element of the matrix $a$ is the inverse matrix $a^{-1}$. The real unimodular group is defined analogously. These groups are non-commutative when $n \geqq 2$.

\end{examplen}
\begin{definition}[J. \textsc{von} \textsc{Neumann} {\cite{neumann231}}]\label{def:VIII.5.Definition.2} Given an abstract group $G$. A complex-valued function $f(g)$ defined on $G$ is called almost periodic on $G$ if the following condition is satisfied: 

\begin{equation}\label{eq:VIII.5.5}
\begin{minipage}[b]{0.92\linewidth}
the set of functions $\{g_s(f, h); s \in G\}$, where $f_s(g, h) = f(gh)$, defined on the direct product $G \times G$ is totally bounded with respect to the topology of uniform convergence on $G \times G$.
\end{minipage}
\end{equation}
\end{definition}


\begin{examplen} Let $G$ be the set $\mathbb{R}^{1}$ of all real numbers in which the group multiplication is defined as the addition of real numbers; this group $\mathbb{R}^{1}$ is called the additive group of real numbers. The function $f(g) = e^{i\alpha g}$, where $\alpha$ is a real number and $i = \sqrt{-1}$, is almost periodic on $\mathbb{R}^{1}$. This we easily see from the addition theorem $f(gsh) = e^{i\alpha g} e^{i\alpha s} e^{i\alpha h}$ and the fact that $\{e^{i\alpha t}; t \in \mathbb{R}^{1}\}$ is totally bounded as a set of complex numbers of absolute value $1$.

\end{examplen}
\begin{proposition}\label{prop:VIII.5.Proposition.1} Suppose $f(g)$ is an almost periodic function on $G$. If we define, following A. \textsc{Weil}, 
\begin{equation}\label{eq:VIII.5.6}
\dist (s, u) = \sup\limits _{g, h \in G} | f (g s h) - f (g u h) |,
\end{equation}
then 
\begin{equation}\label{eq:VIII.5.7}
\dist (s, u) = \dist (a s b, a u b).
\end{equation}
\end{proposition}

\begin{proof}
 Clear from the definition of the group. 
\end{proof}

\begin{corollary}\label{cor:VIII.5.Corollary.1} The set $E$ of all elements $s$ which satisfies $\mathrm{dis}(s, e) = 0$ constitutes an invariant subgroup in $G$, that is, we have 

\begin{equation}\label{eq:VIII.5.8}
\begin{minipage}[b]{0.92\linewidth}
if $e_1, e_2 \in E$, then $e_1e_2 \in E$ and $ae_1a^{-1} \in E$ for every $a \in G$.
\end{minipage}
\end{equation}
\end{corollary}


\begin{proof}
 Let $\dist(e_1, e) = 0$, $\dist(e_2, e) = 0$. Then, by \eqref{eq:VIII.5.7} and the triangle inequality, we obtain 
\[
\dist \left(e _{1} e _{2}, e\right) \leqq \dist \left(e _{1} e _{2}, e _{1} e\right) + \dist \left(e _{1} e, e\right) = 0 + 0 = 0. \]

Similarly we have $\mathrm{dis}(ae_1a^{-1},e) = \mathrm{dis}(ae_1a^{-1},aea^{-1}) = 0$ from $\mathrm{dis}(e_1,e) = 0$. 
\end{proof}

\begin{corollary}\label{cor:VIII.5.Corollary.2} If we write $s \equiv u (\bmod E)$ when $su^{-1} \in E$, then $s \equiv u (\bmod E)$ is equivalent to $\dist(s, u) = 0$. 
\end{corollary}

\begin{proof}
 Clear from $\dist(su^{-1}, e) = \dist(s, eu) = \dist(s, u)$. 
\end{proof}

\begin{corollary}\label{cor:VIII.5.Corollary.3} The concept $s \equiv u \pmod{E}$ has all the general properties of equivalence, namely 
\begin{align}
s& \equiv s (\mathrm{mod} E), \label{eq:VIII.5.9} \\
s& \equiv u (\mathrm{mod} E), \text{ then } u \equiv s (\mathrm{mod} E), \label{eq:VIII.5.10} \\
& \text{ if } s _{1}& \equiv s _{2} (\bmod E) \text{ and } s _{2} \equiv s _{3} (\bmod E), \text{ then } s _{1} \equiv s _{3} \notag
\\
& (\bmod E). \label{eq:VIII.5.11}
\end{align}
\end{corollary}

\begin{proof}
 Clear from Corollary~\ref{cor:VIII.5.Corollary.2} and the triangle inequality for the $\dist(s, u)$. 
\end{proof}

Hence, as in the case of the factor space in a linear space, we can define the factor group or residue class group $G$/$E$ as follows: we shall denote the set of all elements $\in G$ equivalent (mod $E$) to a fixed element $x \in G$ by $\xi_{x}$, the residue class (mod $E$) containing $x$; then the set of all residue classes $\xi_{x}$ constitutes a group $G$/$E$ by the notion of the product 
\begin{equation}\label{eq:VIII.5.12}
\xi_{x} \xi_{y} = \xi_{xy}.
\end{equation}

To justify this definition \eqref{eq:VIII.5.12} of the product, we have to show that 

\begin{equation}\label{eq:VIII.5.13}
\begin{minipage}[b]{0.92\linewidth}
if $x_{1} \equiv x_{2} (\bmod E)$, $y_{1} \equiv y_{2} (\bmod E)$, then $x_{1}y_{1} \equiv x_{2}y_{2} (\bmod E)$.
\end{minipage}
\end{equation}

This is clear, since we have by \eqref{eq:VIII.5.7} and Corollary~\ref{cor:VIII.5.Corollary.2}, 
\[
\begin{aligned}
&\mathrm{dis} (x _{1} y _{1}, x _{2} y _{2})  \leqq \mathrm{dis} (x _{1} y _{1}, x _{2} y _{1}) + \mathrm{dis} (x _{2} y _{1}, x _{2} y _{2}) = \mathrm{dis} (x _{1}, x _{2}) + \mathrm{dis} (y _{1}, y _{2}) = 0 + 0 = 0.
\end{aligned}
\]

Since the function $f(x)$ takes the same constant value on the residue class $\xi_{x}$, we may consider $f(x)$ as a function $F(\xi_{x})$ defined on the residue class group $G$/$E$. 

\begin{corollary}\label{cor:VIII.5.Corollary.4} The residue class group is a metric space by the distance 
\begin{equation}\label{eq:VIII.5.14}
\dist \left(\xi_{x}, \xi_{y}\right) = \dist (x, y).
\end{equation}
\end{corollary}

\begin{proof}
 $x \equiv x_{1} (\bmod E)$ and $y \equiv y_{1} (\bmod E)$ imply 
\[
\dist (x, y) \leqq \dist (x, x _{1}) + \dist (x _{1}, y _{1}) + \dist (y _{1}, y) = 0 + \dist (x _{1}, y _{1}) + 0
\]
and $\dist(x_1, y_1) \leqq \dist(x, y)$ to the effect that $\dist(x, y) = \dist(x_1, y_1)$. Thus \eqref{eq:VIII.5.14} defines a distance in $G/E$. 
\end{proof}

\begin{corollary}\label{cor:VIII.5.Corollary.5} The group $G$/$E$ is a topological group with respect to the distance $\text{ dis }(\xi_{x}, \xi_{y})$, that is, the operation of multiplication $\xi_{x}\xi_{y}$ is continuous as a mapping from the product space $(G/E)\times(G/E)$ onto $G$/$E$, and the operation $\xi_{x}^{-1}$ is continuous as a mapping from $G$/$E$ onto $G$/$E$. 
\end{corollary}

\begin{proof}
 We have, by \eqref{eq:VIII.5.7}, 
\begin{gather}\quad \dist (s u, s ^{\prime} u ^{\prime})\leqq \dist (s u, s ^{\prime} u) + \dist (s ^{\prime} u, s ^{\prime} u ^{\prime}) = \dist (s, s ^{\prime}) + \dist (u, u ^{\prime}) \text{ and } \notag
\\\dist \left(s ^{-1}, u ^{-1}\right) = \dist \left(s s ^{-1} u, s u ^{-1} u\right) = \dist (u, s) = \dist (s, u). \notag
\end{gather}
\end{proof}

We have thus proved the following 

\begin{theorem}[A. \textsc{Weil}]\label{thm:VIII.5.Theorem.1} The topological group $G / E$, metrized by \eqref{eq:VIII.5.14}, is totally bounded, and the function $f(x)$ gives rise to a function $F(\xi_x) (= f(x))$ which is uniformly continuous on this group $G / E$. 
\end{theorem}

\begin{proof}
 The uniform continuity of the function $F(\xi_x)$ is clear from 
\[
\left| F \left(\xi_{x}\right) - F \left(\xi_{y}\right) \right| = \left| f (x) - f (y) \right| \leqq \dist (x, y) = \dist \left(\xi_{x}, \xi_{y}\right). \]

The almost periodicity of the function $f(x)$ implies, by \eqref{eq:VIII.5.7} and \eqref{eq:VIII.5.14}, that the metric space $G / E$ is totally bounded. 
\end{proof}

By the above theorem, the theory of almost periodic functions is reduced to the study of a uniformly continuous function $f(g)$, defined on a totally bounded topological group $G$, metrized by a metric $\text{ dis }(g_{1}, g_{2})$ satisfying condition \eqref{eq:VIII.5.7}. By virtue of this fact, we shall give a proof for the existence of the mean value of an almost periodic function. 

Since $G$ is totally bounded, there exists, for any $\varepsilon > 0$, a finite system of points $g_1, g_2, \ldots, g_n$ such that $\min\limits_{1 \leqq i \leqq n} \text{ dis } (g, g_i) \leqq \varepsilon$ for any $g \in G$. Hence, collecting these finite systems of points corresponding to $\varepsilon = 1, 2^{-1}, 3^{-1}, \ldots$, respectively, we see that there exists a countable system $\{g_j\}$ of points $\in G$ such that $\{g_j\}$ is dense in $G$. We take a sequence of positive numbers $\alpha_j$ such that $\sum\limits_{j=1}^{\infty} \alpha_j = 1$. Let $C(G)$ be the set of all uniformly continuous complex-valued functions $h(g)$ defined on $G$. $C(G)$ is a $B$-space by the operation of the function sum and the norm $\|h\| = \sup |h(g)|$. We define an operator $T$ defined on $C(G)$ into $C(G)$ by $g \in G$. 
\begin{equation}\label{eq:VIII.5.15}
(T \cdot h) (g) = \sum\limits_{j=1} ^{\infty} \alpha_{j} h (g _{j} g).
\end{equation}

By the uniform continuity of $h(g)$ on $G$, there exists, for any $\varepsilon > 0$, a $\delta > 0$ such that $\mathrm{dis}(g, g') \leqq \delta$ implies $|h(g) - h(g')| \leqq \varepsilon$. Thus, by \eqref{eq:VIII.5.7}, $|h(g_j g) - h(g_j g')| \leqq \varepsilon (j = 1, 2, \ldots)$ whenever $\mathrm{dis}(g, g') \leqq \delta$. Hence it is easy to see, by $\alpha_j > 0$ and $\sum\limits_{j=1}^{\infty} \alpha_j = 1$, that $T$ is a bounded linear operator on $C(G)$ into $C(G)$. By the same reasoning we see that the set of functions $h_n(g)$ defined by 
\begin{align}
h _{n} (g)& = n ^{-1} \sum\limits_{m=1} ^{n} \left(T ^{m} h\right) (g), \text{ which is of the form } \notag \\
h _{n} (g)& = \sum\limits_{j=1} ^{\infty} \beta_{j} h \left(g _{j} ^{\prime} g\right) \quad \text{ with } \quad \beta_{j} > 0, \sum\limits_{j=1} ^{\infty} \beta_{j} = 1, \label{eq:VIII.5.16}
\end{align}
is equi-bounded and equi-continuous with respect to $n$. Hence, by the Ascoli-Arzelà theorem, the sequence $\{h_n(g)\}$ contains a subsequence which is uniformly convergent on $G$. 

Therefore, by the mean ergodic theorem, there exists an $h^*(g) \in C(G)$ such that 
\begin{equation}\label{eq:VIII.5.17}
\lim\limits _{n \rightarrow \infty} \sup\limits _{g} | h _{n} (g) - h ^{*} (g) | = 0 \text{ and } T h ^{*} = h ^{*}.
\end{equation}

\begin{proposition}\label{prop:VIII.5.Proposition.2} $h^{*}(g)$ is identically equal to a constant. 
\end{proposition}

\begin{proof}
 We may assume, without losing the generality, that $h(g)$ and $h^{*}(g)$ are real-valued. Suppose that there exist a point $g_{0} \in G$ and a positive constant $\delta$ such that 
\[
h ^{*} (g _{0}) \leqq \beta - 2 \delta , \text{ where } \beta = \sup\limits _{g \in G} h ^{*} (g). \]

By the continuity of $h^*(g)$, there exists a positive number $\varepsilon$ such that $\text{ dis }(g', g'') \leqq \varepsilon$ implies $|h^*(g') - h^*(g'')| \leqq \delta$; in particular, we have $h^*(g'') \leqq \beta - \delta$ whenever $\text{ dis }(g_0, g'') \leqq \varepsilon$. Since the sequence $\{g_j\}$ is dense in $G$, there exists, for any $\varepsilon > 0$, an index $n$ such that, for any $g \in G$, we have $\min\limits_{1 \leqq i \leqq n} \text{ dis }(g, g_j) \leqq \varepsilon$. Hence, by \eqref{eq:VIII.5.7}, we have, for any $g \in G$, 
\[
\min\limits _{1 \leqq j \leqq n} \dist (g _{0}, g _{j} g) \leqq \varepsilon .
\]

Let the minimum be attained at $j = j_0$. Then 
\[
h ^{*} (g) = \left(T h ^{*}\right) (g) = \sum\limits_{j=1} ^{\infty} \alpha_{j} h ^{*} \left(g _{j} g\right) \leqq \alpha_ {j _{0}} (\beta - \delta) + \left(1 - \alpha_ {j _{0}}\right) \beta = \beta - \alpha_ {j _{0}} \delta < \beta ,
\]
contrary to the assumption that $g$ was an arbitrary point of $G$. Therefore $h^{*}(g)$ is identically equal to a constant. 
\end{proof}

\begin{definition}\label{def:VIII.5.Definition.3} We shall call the constant value $h^{*}(g)$ the left mean value of $h(g)$ and denote its value by $M_{g}^{l}(h(g))$: 
\begin{equation}\label{eq:VIII.5.18}
M _{g} ^{l} (h (g)) = \lim\limits _{n \rightarrow \infty} n ^{-1} \sum\limits_{m=1} ^{n} \left(T ^{m} h\right) (g).
\end{equation}
\end{definition}

\begin{theorem}[J. von Neumann]\label{thm:VIII.5.Theorem.2} We have 
\begin{align}
M _{g} ^{l} (\alpha h (g))& = \alpha M _{g} ^{l} (h (g)), \tag{i}\label{eq:VIII.5.i} \\
M _{g} ^{l} (h _{1} (g) + h _{2} (g))& = M _{g} ^{l} (h _{1} (g)) + M _{g} ^{l} (h _{2} (g)), \tag{ii}\label{eq:VIII.5.ii} \\
M _{g} ^{l} \eqref{eq:VIII.5.1}& = 1, \tag{iii}\label{eq:VIII.5.iii} \\
\text{ if } h (g)& \geqq 0 \text{ on } G, \text{ then } M _{g} ^{l} (h (g)) \geqq 0; \notag \\
\text{ if, moreover, } h (g) \not& \equiv 0, \text{ then } M _{g} ^{l} (h (g)) > 0, \tag{iv}\label{eq:VIII.5.iv} \\
\left| M _{g} ^{l} (h (g)) \right|& \leqq M _{g} ^{l} (\left| h (g) \right|), \tag{v}\label{eq:VIII.5.v} \\
M _{g} ^{l} \overline {{(h (g))}}& = \overline {{M _{g} ^{l} (h (g))}}, \tag{vi}\label{eq:VIII.5.vi} \\
M _{g} ^{l} (h (g a))& = M _{g} ^{l} (h (g)), \tag{vii}\label{eq:VIII.5.vii} \\
M _{g} ^{l} (h (a g))& = M _{g} ^{l} (h (g)), \tag{vii'}\label{eq:VIII.5.vii'} \\
M _{g} ^{l} (h (g ^{-1}))& = M _{g} ^{l} (h (g)). \tag{viii}\label{eq:VIII.5.viii}
\end{align}
\end{theorem}

\begin{proof}
 By definition \eqref{eq:VIII.5.18}, it is clear that (i), (ii), (iii), the first part of (iv), (v) and (vi) are true. The truth of (vii) is proved by Proposition~\ref{prop:VIII.5.Proposition.2}. (vii') is proved as follows.


Starting with the linear operator $T'$ defined by 
\[
\left(T ^{\prime} h\right) (g) = \sum\limits_{j=1} ^{n} \alpha_{j} h \left(g g _{j}\right),
\]
we can also define a right mean $M_g^r(h(g))$ which, as a functional of $h(g)$, satisfies (i), (ii), (iii), the first part of (iv), (v), (vi) and (vii'). We thus have to prove that the left mean $M_g^l(h(g))$ coincides with the right mean $M_g^r(h(g))$. By its definition of the left mean, there exist, for any $\varepsilon > 0$, a sequence of elements $\{k_j\} \subseteq G$ and a sequence of positive numbers $\beta^j$ with $\sum\limits_{j=1}^{\infty} \beta_j = 1$ such that
\begin{equation}\label{eq:VIII.5.19}
\sup\limits _{g} \left| \sum\limits_{j=1} ^{\infty} \beta_{j} h (k _{j} g) - M _{g} ^{l} (h (g)) \right| \leqq \varepsilon .
\end{equation}

Similarly, there exist a sequence of elements $\{s_j\} \subseteq G$ and a sequence of positive numbers $\gamma_j$ with $\sum\limits_{j=1}^{\infty} \gamma_j = 1$ such that 
\begin{equation}\label{eq:VIII.5.20}
\sup\limits _{g} \left| \sum\limits_{j=1} ^{\infty} \gamma_{j} h (g s _{j}) - M _{g} ^{\prime} (h (g)) \right| \leqq \varepsilon .
\end{equation}

We have, from \eqref{eq:VIII.5.19} and (vii) 
\[
\sup\limits _{g} \left| \sum\limits_{i,j} \gamma_{i} \beta_{j} h (k _{j} g s _{i}) - M _{g} ^{l} (h (g)) \right| \leqq \varepsilon ,
\]
and, similarly from \eqref{eq:VIII.5.20} and $(\mathrm{vii}^{\prime})$ 
\[
\sup\limits _{g} \left| \sum\limits_{i,j} \gamma_{i} \beta_{j} h (k _{j} g s _{i}) - M _{g} ^{r} (h (g)) \right| \leqq \varepsilon .
\]

Hence we must have $M_g^l(h(g)) = M_g'(h(g))$. 

We next remark that a linear functional $M_{g}(h(g))$ defined on $C(G)$ is uniquely determined by the properties (i), (ii), (iii), the first part of (iv), (v), (vi) and (vii) (or (vii') as well). In fact, we have, by \eqref{eq:VIII.5.20},
\[
M _{g} ^{\prime} (h (g)) - \varepsilon \leqq \sum\limits_{i=1} ^{\infty} \gamma_{i} h (g s _{i}) \leqq M _{g} ^{\prime} (h (g)) + \varepsilon \quad \text{ for real-valued } h (g)
\]

Hence, for real-valued $h(g)$, $M_g(h(g))$ must coincide with the right mean $M_g^r(h(g))$ and hence with the left mean $M_g^l(h(g))$ as well. Therefore we see that we must have $M_g = M_g' = M_g^l$. Being equal to the right mean, $M_g$ must satisfy (vii'). Moreover, since $M_g^l(h(g^{-1}))$ satisfies, as a linear functional, (i), (ii), (iii), the first part of (iv), (v), (vi) and (vii'), we must have $M_g^l(h(g^{-1})) = M_g^r(h(g)) = M_g^l(h(g))$.

Finally we shall prove the last part of (iv). Suppose $h(g_{0}) > 0$. For any $\varepsilon > 0$, there exists, by the total boundedness of $G$, a finite system of elements $s_{1}, s_{2}, \ldots, s_{n}$ such that 
\[
\min\limits _{1 \leqq i \leqq n} \sup\limits _{g} | h (g s _{i}) - h (g s) | < \varepsilon
\]
for all $s \in G$. This we see by the uniform continuity of $h(g)$ and the fact that $\dist(gs_i, gs) = \dist(s_i, s)$. Hence, for $\varepsilon = h(g_0)/2$, we obtain, for any $s \in G$, a suffix $i, 1 \leqq i \leqq n$, such that 
\[
h \left(g _{0} s _{i} ^{-1} s\right) \geqq h \left(g _{0}\right) / 2.
\]

Thus, by the non-negativity of the function $h(g)$, we obtain 
\[
\sum\limits_{i=1} ^{n} h \left(g _{0} s _{i} ^{-1} s\right) \geqq h \left(g _{0}\right) / 2 > 0 \quad \text{ for all } \quad s \in G.
\]

Therefore, by taking the right mean of both sides, we have 
\[
M _{g} ^{r} \left(\sum\limits_{i=1} ^{n} h \left(g _{0} s _{i} ^{-1} s\right)\right) = n M _{s} ^{r} (h (s)) \geqq h \left(g _{0}\right) / 2 > 0.
\]
\end{proof}

\begin{remarku} The happy idea of introducing the distance \eqref{eq:VIII.5.1} is shown in A. \textsc{Weil} \cite{weil334}. The application of the mean ergodic theorem to the existence proof of the mean value is due to the present author. See also W. \textsc{Maak} \cite{lumer185}. 
\end{remarku}


\yossection{6}{The Resolvent of a Dual Operator}
\label{sec:VIII.6}

\begin{lemma}\label{lem:VIII.6.Lemma.1} Let $X$ and $Y$ be complex $B$-spaces. Let $T$ be a linear operator with $D(T)^a = X$ and $R(T) \subseteq Y$. Then $(T')^{-1}$ exists iff $R(T)^a = Y$.
\end{lemma}

\begin{proof}
If $T' \nu_0^* = 0$, then 
\[
\langle x, T ^{\prime} y _{0} ^{*} \rangle = \langle T x, y _{0} ^{*} \rangle = 0 \quad \text{ for all } \quad x \in D (T),
\]
and hence $y_0^*(R(T)^a) = 0$. Thus $R(T)^a = Y$ implies $y_0^* = 0$ and so $T'$ has an inverse. On the other hand, if $y_0 \bar{\in} R(T)^a$, then, the Hahn-Banach theorem asserts that there exists a continuous linear functional $y_0^* \in Y'$ such that $y_0^*(y_0) = 1$ and $y_0^*(R(T)^a) = 0$. Hence $\langle Tx, y_0^* \rangle = 0$ for all $x \in D(T)$, and so $y_0^* \in D(T')$ and $T'y_0^* = 0$, whereas $y_0^*(y_0) \neq 0$, i.e., $y_0^* \neq 0$. Therefore, the condition $R(T)^a \neq Y$ implies that $T'$ cannot have an inverse. 
\end{proof}

\begin{theorem}[R. S. \textsc{Phillips} {\cite{phillips252}}]\label{thm:VIII.6.Theorem.1} Let $T$ be a linear operator with an inverse and such that $D(T)^{a} = X$ and $R(T)^{a} = Y$, where $X$ and $Y$ are $B$-spaces. Then 
\begin{equation}\label{eq:VIII.6.1}
(T ^{\prime}) ^{-1} = (T ^{-1}) ^{\prime}.
\end{equation}
$T^{-1}$ is bounded on $Y$ iff $T$ is closed and $(T')^{-1}$ is bounded on $X_s'$. 
\end{theorem}

\begin{proof}
 $(T^{-1})'$ exists because $D(T^{-1}) = R(T)$ is dense in $Y$. $(T')^{-1}$ exists by Lemma~\ref{lem:VIII.6.Lemma.1}. We have to show the equality \eqref{eq:VIII.6.1}. If $y \in R(T)$ and $y^* \in D(T')$, then 
\[
\langle y, y ^{*} \rangle = \langle T T ^{-1} y, y ^{*} \rangle = \langle T ^{-1} y, T ^{\prime} y ^{*} \rangle . \]

Hence $R(T') \subseteq D((T^{-1})')$ and $(T^{-1})'(T'y^*) = y^*$ for all $y^* \in D(T')$. Thus $(T^{-1})'$ is an extension of $(T')^{-1}$. Next, if $x \in D(T)$, then 
\[
\langle x, x ^{*} \rangle = \langle T ^{-1} T x, x ^{*} \rangle = \langle T x, (T ^{-1}) ^{\prime} x ^{*} \rangle \quad \text{ for all } \quad x ^{*} \in D ((T ^{-1}) ^{\prime}).
\]

Hence $R((T^{-1})') \subseteq D(T')$ and $T'(T^{-1})' x^* = x^*$ for all $x^* \in D((T^{-1})')$. Thus $(T^{-1})'$ is a contraction of $(T')^{-1}$. Therefore we have proved \eqref{eq:VIII.6.1}. 

If, in addition, $T^{-1}$ is bounded on $Y$, then $(T^{-1})'$ is also bounded. Conversely, if $(T')^{-1}$ is bounded on $X_s'$, then, for all $x \in R(T)$ and $x^* \in X'$, we have, by \eqref{eq:VIII.6.1}, 
\[
\begin{aligned}
&| \langle T ^{-1} x, x ^{*} \rangle |  = | \langle x, (T ^{-1}) ^{\prime} x ^{*} \rangle | = | \langle x, (T ^{\prime}) ^{-1} x ^{*} \rangle | \leqq \| (T ^{\prime}) ^{-1} \| \cdot \| x ^{*} \| \cdot \| x \|.
\end{aligned}
\]

Since $T^{-1}$ is closed and $R(T)^{a} = Y$, $T^{-1}$ must be bounded. 
\end{proof}
\begin{lemma}\label{lem:VIII.6.Lemma.2} Let $T$ be a linear operator with $D(T)^a = X$ and $R(T) \subseteq Y$, where $X$ and $Y$ are $B$-spaces. If $R(T')$ is weakly* dense in $X'$, then $T$ has an inverse. 
\end{lemma}

\begin{proof}
 Suppose that there exists an $x_{0} \neq 0$ such that $Tx_{0} = 0$. Then 
\[
\langle x _{0}, T ^{\prime} y ^{*} \rangle = \langle T x _{0}, y ^{*} \rangle = 0 \quad \text{ for all } \quad y ^{*} \in D (T ^{\prime}). \]

This shows that the weak* closure of $R(T')$ is a proper linear subspace of $X'$, contrary to the hypothesis. 
\end{proof}

\begin{theorem}[R. S. \textsc{Phillips} {\cite{phillips252}}]\label{thm:VIII.6.Theorem.2} Let $X$ be a complex $B$-space, and $T$ a closed linear operator with $D(T)^a = X$ and $R(T) \subseteq X$. Then 
\begin{equation}\label{eq:VIII.6.2}
\varrho (T) = \varrho (T ^{\prime}) \quad \text{ and } \quad R (\lambda ; T) ^{\prime} = R (\lambda ; T ^{\prime}) \quad \text{ for } \quad \lambda \in \varrho (T).
\end{equation}
\end{theorem}

\begin{proof}
 If $\lambda \in \varrho(T)$, then, by Theorem~\ref{thm:VIII.6.Theorem.1}, $\lambda \in \varrho(T')$ and $R(\lambda; T)' = R(\lambda; T')$. On the other hand, if $\lambda \in \varrho(T')$, then Lemma~\ref{lem:VIII.6.Lemma.2} shows that $(\lambda I - T)$ has an inverse $(\lambda I - T)^{-1}$ which is closed with $(\lambda I - T)$. Lemma~\ref{lem:VIII.6.Lemma.1} then shows that $D((\lambda I - T)^{-1}) = R(\lambda I - T)$ is strongly dense in $Y$. Hence, by Theorem~\ref{thm:VIII.6.Theorem.1}, $\lambda \in \varrho(T)$. 
\end{proof}


\yossection{7}{Dunford's Integral}
\label{sec:VIII.7}

Let $X$ be a complex $B$-space and $T$ a bounded linear operator $\in L(X, X)$. We shall define a function $f(T)$ of $T$ by Cauchy's type integral 
\[
f (T) = (2 \pi i) ^{-1} \int_{C} f (\lambda) R (\lambda ; T) \dd \lambda .
\]

To this purpose, we denote by $F(T)$ the family of all complex-valued functions $f(\lambda)$ which are holomorphic in some neighbourhood of the spectrum $\sigma(T)$ of $T$; the neighbourhood need not be connected, and can depend on the function $f(\lambda)$. Let $f \in F(T)$, and let an open set $U \supseteq \sigma(T)$ of the complex plane be contained in the domain of holomorphy of $f$, and suppose further that the boundary $\partial U$ of $U$ consists of a finite number of rectifiable Jordan curves, oriented in positive sense. Then the bounded linear operator $f(T)$ will be defined by 
\begin{equation}\label{eq:VIII.7.1}
f (T) = (2 \pi i) ^{-1} \int_{\partial U} f (\lambda) R (\lambda ; T) \dd \lambda ,
\end{equation}
and the integral on the right may be called a Dunford's integral. By Cauchy's integral theorem, the value $f(T)$ depends only on the function $f$ and the operator $T$, but not on the choice of the domain $U$. 

The following operational calculus holds: 

\begin{theoremn}[N. \textsc{Dunford}] If $f$ and $g$ are in $F(T)$, and $\alpha$ and $\beta$ are complex numbers, then
\begin{gather}
\alpha f+\beta g\text{ is in }F(T)\text{ and }\alpha f(T)+\beta g(T)=(\alpha f+\beta g)(T),\label{eq:VIII.7.2}\\
f\cdot g\text{ is in }F(T)\text{ and }f(T)\cdot g(T)=(f\cdot g)(T),\label{eq:VIII.7.3}
\end{gather}
\begin{equation}\label{eq:VIII.7.4}
\left.\begin{gathered}
\parbox{0.86\linewidth}{\setlength{\parindent}{1.2em}\indent
if $f$ has the Taylor expansion $f(\lambda)=\sum\limits_{n=0}^{\infty}\alpha_n\lambda^n$ convergent in a neighbourhood $U$ of $\sigma(T)$, then}\vphantom{\sum}\mathstrut\\[2pt]
f(T)=\sum\limits_{n=0}^{\infty}\alpha_nT^n\text{ (in the operator norm topology),}
\end{gathered}\right\}
\end{equation}
\begin{equation}\label{eq:VIII.7.5}
\left.\begin{gathered}
\parbox{0.86\linewidth}{\setlength{\parindent}{1.2em}\indent
let $f_n\in F(T)$ $(n=1,2,\ldots)$ be holomorphic in a fixed neighbourhood $U$ of $\sigma(T)$. If $f_n(\lambda)$ converges to $f(\lambda)$ uniformly on $U$, then $f_n(T)$ converges to $f(T)$ in the operator norm topology,}
\end{gathered}\right\}
\end{equation}

\begin{equation}\label{eq:VIII.7.6}
\text{ if } f \in F (T), \text{ then } f \in F (T ^{\prime}) \text{ and } f (T ^{\prime}) = f (T) ^{\prime}.
\end{equation}

\end{theoremn}
\begin{proof}
 \eqref{eq:VIII.7.2} is clear. Proof of \eqref{eq:VIII.7.3}. Let $U_{1}$ and $U_{2}$ be open neighbourhoods of $\sigma(T)$ whose boundaries $\partial U_{1}$ and $\partial U_{2}$ consist of a finite number of rectifiable Jordan curves, and assume that $U_{1} + \partial U_{1} \subseteq U_{2}$ and that $U_{2} + \partial U_{2}$ is contained in the domain of holomorphy of $f$ and $g$. Then, by virtue of the resolvent equation and Cauchy's integral theorem, we obtain 
\[
\begin{aligned}
&f (T) g (T) = - (4 \pi^{2}) ^{-1} \int_ {\partial U _{1}} f (\lambda) R (\lambda ; T) \dd \lambda \cdot \int_ {\partial U _{2}} g (\mu) R (\mu ; T) \dd \mu
\\
&= - (4 \pi^{2}) ^{-1} \int_ {\partial U _{1}} \int_ {\partial U _{2}} f (\lambda) g (\mu) (\mu - \lambda) ^{-1} (R (\lambda ; T) - R (\mu ; T)) \dd \lambda \dd \mu
\\
&= (2 \pi i) ^{-1} \int_ {\partial U _{1}} f (\lambda) R (\lambda ; T) \cdot \left\{(2 \pi i) ^{-1} \int_ {\partial U _{2}} (\mu - \lambda) ^{-1} g (\mu) \dd \mu \right\} \dd \lambda
\\
&- (2 \pi i) ^{-1} \int_ {\partial U _{2}} g (\mu) R (\mu ; T) \cdot \left\{(2 \pi i) ^{-1} \int_ {\partial U _{1}} (\mu - \lambda) ^{-1} f (\lambda) \dd \lambda \right\} \dd \mu
\\
&= (2 \pi i) ^{-1} \int_ {\partial U _{1}} f (\lambda) g (\lambda) R (\lambda ; T) \dd \lambda = (f \cdot g) (T).
\end{aligned}
\]

Proof of \eqref{eq:VIII.7.4}. By hypothesis, $U$ must contain a circle $\{\lambda; |\lambda| \leqq r_{\sigma}(T) + \varepsilon\}$ $\varepsilon > 0$, in its interior, where $r_{\sigma}(T)$ is the spectral radius of $T$ (Theorem~\ref{thm:VIII.2.Theorem.4} in Chapter~\ref{chap:VIII}, \ref{sec:VIII.2}). Hence the power series $f(\lambda) = \sum\limits_{n=0}^{\infty} \alpha_n \lambda^n$ converges uniformly on the circle $C = \{\lambda; |\lambda| \leqq r_{\sigma}(T) + \varepsilon\}$, for some $\varepsilon > 0$. Hence, by Cauchy's integral theorem and Laurent's expansion $R(\lambda; T) = \sum\limits_{n=1}^{\infty} \lambda^{-n} T^{n-1}$ of $R(\lambda; T)$ (Chapter~\ref{chap:VIII}, \ref{sec:VIII.2}, \eqref{eq:VIII.2.5}), 
\[
\begin{aligned}
f (T) &= (2 \pi i) ^{-1} \int_{\partial C} \left(\sum\limits_{k=0} ^{\infty} \alpha_{k} \lambda^{k}\right) R (\lambda ; T) \dd \lambda = (2 \pi i) ^{-1} \sum\limits_{k=0} ^{\infty} \alpha_{k} \int_{\partial C} \lambda^{k} R (\lambda ; T) \dd \lambda
\\
&= (2 \pi i) ^{-1} \sum\limits_{k=0} ^{\infty} \alpha_{k} \sum\limits_{n=1} ^{\infty} \int_{\partial C} \lambda^{k-n} T ^{n-1} \dd \lambda = \sum\limits_{k=0} ^{\infty} \alpha_{k} T ^{k}.
\end{aligned}
\]

\eqref{eq:VIII.7.5} is proved by \eqref{eq:VIII.7.1}, and \eqref{eq:VIII.7.6} is also proved by \eqref{eq:VIII.7.1} and formula \eqref{eq:VIII.7.2} of the preceding section. 
\end{proof}

\begin{corollary}[Spectral Mapping Theorem]\label{cor:VIII.7.Corollary.1} If $f$ is in $F(T)$, then $f(\sigma(T)) = \sigma(f(T))$. 
\end{corollary}

\begin{proof}
 Let $\lambda \in \sigma(T)$, and define the function $g$ by $g(\mu) = (f(\lambda) - f(\mu)) / (\lambda - \mu)$. By the Theorem, $f(\lambda)I - f(T) = (\lambda I - T)g(T)$. Hence, if $(f(\lambda)I - f(T))$ had a bounded inverse $B$, then $g(T)B$ would be the bounded inverse of $(\lambda I - T)$. Thus $\lambda \in \sigma(T)$ implies that $f(\lambda) \in \sigma(f(T))$. Let, conversely, $\lambda \in \sigma(f(T))$, and assume that $\lambda \bar{\in} f(\sigma(T))$. Then the function $g(\mu) = (f(\mu) - \lambda)^{-1}$ must belong to $F(T)$, and so, by the preceding Theorem, $g(T)(f(T) - \lambda I) = I$ which contradicts the assumption $\lambda \in \sigma(f(T))$. 
\end{proof}

\begin{corollary}\label{cor:VIII.7.Corollary.2} If $f \in F(T)$, $g \in F(f(T))$ and $h(\lambda) = g(f(\lambda))$, then $h$ is in $F(T)$ and $h(T) = g(f(T))$. 
\end{corollary}

\begin{proof}
 That $h \in F(T)$ follows from Corollary~\ref{cor:VIII.7.Corollary.1}. Let $U_{1}$ be an open neighbourhood of $\sigma(f(T))$ whose boundary $U_{1}$ consists of a finite number of rectifiable Jordan curves such that $U_{1} + \partial U_{1}$ is contained in the domain of holomorphy of $g$. Let $U_{2}$ be a neighbourhood of $\sigma(T)$ whose boundary $\partial U_{2}$ consists of a finite number of rectifiable Jordan curves such that $U_{2} + \partial U_{2}$ is contained in the domain of holomorphy of $f$ and $f(U_{2} + \partial U_{2}) \subseteq U_{1}$. Then we have, for $\lambda \in \partial U_{1}$, 
\[
R (\lambda ; f (T)) = (2 \pi i) ^{-1} \int_ {\partial U _{2}} (\lambda - f (\mu)) ^{-1} R (\mu ; T) \dd \mu ,
\]
since the right hand operator $S$ satisfies, by \eqref{eq:VIII.7.3}, the equation $(\lambda I - f(T)) S = S(\lambda I - f(T)) = I$. Hence, by Cauchy's integral theorem, 
\[
\begin{aligned}
g (f (T)) &= (2 \pi i) ^{-1} \int_ {\partial U _{1}} g (\lambda) R (\lambda ; f (T)) \dd \lambda
\\
\qquad &= (- 4 \pi^{2}) ^{-1} \int_ {\partial U _{1}} \int_ {\partial U _{2}} g (\lambda) (\lambda - f (\mu)) ^{-1} R (\mu ; T) \dd \mu \dd \lambda
\\
\qquad &= (2 \pi i) ^{-1} \int_ {\partial U _{2}} R (\mu ; T) g (f (\mu)) \dd \mu = h (T). \end{aligned}
\]
\end{proof}

\begin{remarku} The introduction of the operational calculus based on a formula like \eqref{eq:VIII.7.1} goes back to the investigations by H. \textsc{Poincare} on  
continuous groups ($1899$). The exposition of the operational calculus in this section is adapted from N. \textsc{Dunford}-J. \textsc{Schwartz} \cite{dunford44}. In the next chapter on semi-groups, we shall frequently make use of Dunford's integral for a closed unbounded operator $T$. 
\end{remarku}


\yossection{8}{The Isolated Singularities of a Resolvent}
\label{sec:VIII.8}

Let $\lambda_{0}$ be an isolated singular point of the resolvent $R(\lambda; T)$ of a closed linear operator $T$ on a complex $B$-space $X$ into $X$. Then $R(\lambda; T)$ can be expanded into Laurent series 
\begin{equation}\label{eq:VIII.8.1}
R (\lambda ; T) = \sum\limits_{n = - \infty} ^{\infty} (\lambda - \lambda_{0}) ^{n} A _{n}, A _{n} = (2 \pi i) ^{-1} \int_{C} (\lambda - \lambda_{0}) ^{-n-1} R (\lambda ; T) \dd \lambda ,
\end{equation}
where $C$ is a circumference of sufficiently small radius: $|\lambda - \lambda_{0}| = \varepsilon$ such that the circle $|\lambda - \lambda_{0}| \leqq \varepsilon$ does not contain other singularities than $\lambda = \lambda_{0}$, and the integration is performed counter-clockwise. By virtue of the resolvent equation we obtain 

\begin{theorem}\label{thm:VIII.8.Theorem.1} $A$'s are mutually commutative bounded linear operators and 
\begin{equation}\label{eq:VIII.8.2}
\begin{gathered}
T A _{k} x = A _{k} T x \text{ for } x \in D (T) (k = 0, \pm 1, \pm 2, \dots), A _{k} A _{m} = 0 \quad \text{ for } \quad k \geqq 0, m \leqq - 1,
\\
A _{n} = (- 1) ^{n} A _{0} ^{n+1} (n \geqq 1),
\\
A _{-p-q+1} = A _{-p} A _{-q} (p, q \geqq 1).
\end{gathered}
\end{equation}
\end{theorem}

\begin{proof}
 The boundedness and the mutual commutativity of $A$'s and the commutativity of $A$'s with $T$ are clear from the integral representation of $A$'s. 


We substitute the expansion of $R(\lambda; T)$ in the resolvent equation $R(\lambda; T) - R(\mu; T) = (\mu - \lambda) R(\lambda; T) R(\mu; T)$, and obtain 
\[
\sum\limits_{k = - \infty} ^{\infty} A _{k} \frac {(\lambda - \lambda_{0}) ^{k} - (\mu - \lambda_{0}) ^{k}}{(\lambda - \lambda_{0}) - (\mu - \lambda_{0})} = - \sum\limits_{k, m = - \infty} ^{\infty} A _{k} A _{m} (\lambda - \lambda_{0}) ^{k} (\mu - \lambda_{0}) ^{m}.
\]

The coefficient of $A_{n}$ on the left is 
\[
\begin{aligned}
(\lambda - \lambda_{0}) ^{n-1} + (\lambda - \lambda_{0}) ^{n-2} (\mu - \lambda_{0}) + \dots + (\mu - \lambda_{0}) ^{n-1}, n& \geqq 1,
\\
& - \{(\lambda - \lambda_{0}) ^{n} (\mu - \lambda_{0}) ^{-1}
\\
&\quad + (\lambda - \lambda_{0}) ^{n+1} (\mu - \lambda_{0}) ^{-2} + \dots
\\
+ (\lambda - \lambda_{0}) ^{-1} (\mu - \lambda_{0}) ^{n} \}, n& < 0.
\end{aligned}
\]

Hence the terms containing $(\lambda - \lambda_0)^k (\mu - \lambda_0)^m$ with $k \geqq 0$ and $m \leqq -1$ are missing so that we must have the orthogonality $A_k A_m = 0$ ( $k \geqq 0$, $m \leqq -1$). Hence 
\[
R ^{+} (\lambda ; T) = \sum\limits_{n=0} ^{\infty} A _{n} (\lambda - \lambda_{0}) ^{n} \quad \text{ and } \quad R ^{-} (\lambda ; T) = \sum\limits_{n = - \infty} ^{-1} A _{n} (\lambda - \lambda_{0}) ^{n}
\]
must both satisfy the resolvent equation. Substituting the expansion of $R^{+}(\lambda; T)$ in the resolvent equation 
\[
R ^{+} (\lambda ; T) - R ^{+} (\lambda ; T) = (\mu - \lambda) R ^{+} (\lambda ; T) R ^{+} (\mu ; T),
\]
we obtain, setting $(\lambda - \lambda_0) = h$, $(\mu - \lambda_0) = k$, 
\[
\sum\limits_{p=1} ^{\infty} A _{p} \left(h ^{p} - k ^{p}\right) = (k - h) \left(\sum\limits_{p=0} ^{\infty} A _{p} h ^{p}\right) \left(\sum\limits_{q=0} ^{\infty} A _{q} k ^{q}\right).
\]

Hence, dividing both sides by $(k - h)$, we obtain 
\[
- \sum\limits_{p=1} ^{\infty} A _{p} (h ^{p-1} + h ^{p-2} k + \dots + k ^{p-1}) = \sum\limits_{p,q=0} ^{\infty} h ^{p} k ^{q} A _{p} A _{q}
\]
so that we have $-A_{p+q+1} = A_pA_q$ ( $p, q \geqq 0$). Thus, in particular, 
\[
A _{1} = - A _{0} ^{2}, A _{2} = - A _{1} A _{0} = (- 1) ^{2} A _{0} ^{3}, \dots , A _{n} = (- 1) ^{n} A _{0} ^{n+1} (n \geqq 1).
\]

Similarly, from the resolvent equation for $R^{-}(\lambda; T)$, we obtain, setting $(\lambda - \lambda_{0})^{-1} = h$, $(\mu - \lambda_{0})^{-1} = k$, 
\[
\sum\limits_{p=1} ^{\infty} A _{-p} \left(h ^{p-1} + h ^{p-2} k + \dots + k ^{p-1}\right) = \sum\limits_{p,q=1} ^{\infty} h ^{p-1} k ^{q-1} A _{-p} A _{-q},
\]
so that we have $A_{-p-q+1} = A_{-p}A_{-q}$ ( $p, q \geqq 1$). In particular, we have 
\[
A _{-1} = A _{-1} ^{2}, A _{-2} = A _{-1} A _{-2}, \dots , A _{-n} = A _{-1} A _{-n} (n \geqq 1).
\]
\end{proof}

\begin{theorem}\label{thm:VIII.8.Theorem.2} We have 
\begin{equation}\label{eq:VIII.8.3}
\begin{gathered}
{A _{n} = (T - \lambda_{0} I) A _{n+1} (n \geqq 0),} {(T - \lambda_{0} I) A _{-n} = A _{-(n+1)} = (T - \lambda_{0} I) ^{n} A _{-1} \quad (n \geqq 1),}
\\
{(T - \lambda_{0} I) A _{0} = A _{-1} - I.}
\end{gathered}
\end{equation}
\end{theorem}

\begin{proof}
 By the integral representation of $A_{k}$, we see that the range $R(A_{k})$ is in the domain of $T$, so that we can multiply $A_{k}$ by $T$ on the left. Thus our Theorem is proved by the identity 
\begin{equation}\label{eq:VIII.8.4}
\begin{aligned}
&I = (\lambda I - T) \sum\limits_{k = - \infty} ^{\infty} A _{k} (\lambda - \lambda_{0}) ^{k} = \{(\lambda - \lambda_{0}) I + (\lambda_{0} I - T) \} \sum\limits_{k = - \infty} ^{\infty} A _{k} (\lambda - \lambda_{0}) ^{k}.
\end{aligned}
\end{equation}
\end{proof}

\begin{theorem}\label{thm:VIII.8.Theorem.3} If $\lambda_0$ is a pole of $R(\lambda; T)$ of order $m$, then $\lambda_0$ is an eigenvalue of $T$. We have 

$R(A_{-1}) = N((\lambda_0I - T)^n)$ and $R(I - A_{-1}) = R((\lambda_0I - T)^n)$ for $n \geqq m$, \eqref{eq:VIII.8.4} so that, in particular, 
\begin{equation}\label{eq:VIII.8.5}
X = N ((\lambda_{0} I - T) ^{n}) \oplus R ((\lambda_{0} I - T) ^{n}) \quad \text{ for } \quad n \geqq m.
\end{equation}
\end{theorem}

\begin{proof}
 Since $A_{-1}$ is a bounded linear operator satisfying $A_{-1}^2 = A_{-1}$, it is easy to see that 
\begin{equation}\label{eq:VIII.8.6}
N (A _{-1}) = R (I - A _{-1}). \end{equation}

We put $X_{1} = N(A_{-1}) = R(I - A_{-1})$, and put also 
\begin{equation}\label{eq:VIII.8.7}
X _{2} = R (A _{-1}), N _{n} = N ((\lambda_{0} I - T) ^{n}) \text{ and } R _{n} = R _{n} ((\lambda_{0} I - T) ^{n})
\end{equation}

Let $x \in N_n$ where $n \geqq 1$. Then we see, by $(T - \lambda_0 I)^n A_{n-1} = (T - \lambda_0 I) A_0 = A_{-1} - I$, that $0 = A_{n-1}(T - \lambda_0 I)^n x = (T - \lambda_0 I)^n A_{n-1} x = (T - \lambda_0 I) A_0 x = A_{-1} x - x$ so that $x = A_{-1} x \in X_2$. Thus $N_n$ with $n \geqq 1$ belongs to $X_2$. Let, conversely, $x \in X_2$. Then we have $x = A_{-1} y$ and so $x = A_{-1} A_{-1} y = A_{-1} x$ by $A_{-1} = A_{-1}^2$; consequently, we have $(T - \lambda_0 I)^n x = A_{-(n+1)} x$ by $(T - \lambda_0 I)^n A_{-1} = A_{-(n+1)}$. Since $A_{-(n+1)} = 0$ for $n \geqq m$ by hypothesis, it follows that $X_2 \subseteq N_n$ for $n \geqq m$ and so 
\begin{equation}\label{eq:VIII.8.8}
N _{n} = X _{2} \quad \text{ if } \quad n \geqq m.
\end{equation}

Because $(T - \lambda_0I)A_{-m} = A_{-(m+1)} = 0$ and $A_{-m}\neq 0$, the number $\lambda_0$ is an eigenvalue of $T$. 


We see that $X_{1} = N(A_{-1}) = R(I - A_{-1}) \subseteq R_{n}$ by $(T - \lambda_0 I)^n A_{n-1} = A_{-1} - I$. If $n \geqq m$, then $x \in R_n \cap N_n$ implies $x = 0$. For, if $x = (\lambda_0 I - T)^n y$ and $(\lambda_0 I - T)^n x = 0$, then, by \eqref{eq:VIII.8.8}, $y \in N_{2n} = N_n$ and therefore $x = 0$. Next suppose $x \in R_n$ with $n \geqq m$, and write $x = x_1 + x_2$ where $x_1 = (I - A_{-1})x \in X_1$, $x_2 = A_{-1}x \in X_2$. Then, since $X_1 \subseteq R_n$, $x_2 = x - x_1 \in R_n$. But $x_2 \in X_2 = N_n$ by \eqref{eq:VIII.8.8}, and so $x_2 \in R_n \cap N_n$, $x_2 = 0$. This proves that $x = x_1 \in X_1$. Therefore we have proved that $R_n = X_1$ if $n \geqq m$.
\end{proof}

\begin{theorem}\label{thm:VIII.8.Theorem.4} If, in particular, $T$ is a bounded linear operator such that $X_{2} = R(A_{-1})$ is a finite dimensional linear subspace of $X$, then $\lambda_{0}$ is a pole of $R(\lambda; T)$. 
\end{theorem}

\begin{proof}
 Let $x_1, x_2, \ldots, x_k$ be a base of the linear space $X_2$. Since $x_1, Tx_1, T^2x_1, \ldots, T^k x_1$ are linearly dependent vectors of $X_2$, there exists a non-zero polynomial $P_1(\lambda)$ such that $P_1(T)x_1 = 0$. Similarly there exist non-zero polynomials $P_2(\lambda), \ldots, P_k(\lambda)$ such that $P_j(T)x_j = 0$ ( $j = 2, 3, \ldots, k$). Then for the polynomial $P(\lambda) = \prod\limits_{j=1}^{k} P_j(\lambda)$, we must have $P(T)x_j = 0$ ( $j = 1, \ldots, k$) and hence $P(T)x = 0$ for every $x \in X_2$. Let 
\[
P (\lambda) = \alpha \prod\limits_{j=0} ^{s} (\lambda - \lambda_{j}) ^ {\nu_{j}} (\alpha \neq 0)
\]
be the factorization of $P(\lambda)$. Then we can prove that $(T - \lambda_0 I)^{\nu_0} x = 0$ for every $x \in X_2$. Assume the contrary, and let $x_0 \in X_2$ be such that $(T - \lambda_0 I)^{\nu_0} x_0 \neq 0$. Then, by $P(T)x_0 = 0$, we see that there exist at least one $\lambda_j (j \neq 0)$ and a polynomial $Q(\lambda)$ such that 
\[
(T - \lambda_{j} I) Q (T) (T - \lambda_{0} I) ^ {\nu_{0}} x _{0} = 0
\]
and $y = Q(T)(T - \lambda_0 I)^{\nu_0} x_0 \neq 0$. Thus $y \in X_2$ is an eigenvector of $T$ corresponding to the eigenvalue $\lambda_j$. Hence $(\lambda I - T)y = (\lambda - \lambda_j)y$ and so, multiplying both sides by $R(\lambda; T)$, we obtain $y = (\lambda - \lambda_j)R(\lambda; T)y$ which implies 
\[
y = A _{-1} y = (2 \pi i) ^{-1} \int_{C} R (\lambda ; T) y \dd \lambda = (2 \pi i) ^{-1} \int_{C} (\lambda - \lambda_{j}) ^{-1} y \dd \lambda = 0,
\]
by taking the circumference $C$ with $\lambda_{0}$ as centre sufficiently small. This is a contradiction, and so there must exist a positive integer $m$ such that $(T-\lambda_{0}I)^{m}X_{2}=0$. Thus, by $X_{2}=R(A_{-1})$ and $(T-\lambda_{0}I)^{n}A_{-1}=A_{-(n+1)}$, we see that $A_{-(n+1)}=0$ for $n\geqq m$. 
\end{proof}

\yossub{Comments and References}

Section~\hyperref[sec:VIII.6]{$6$} is adapted from R. S. \textsc{Phillips} \cite{phillips252}. Section~\hyperref[sec:VIII.8]{$8$} is adapted from M. \textsc{Nagumo} \cite{nagumo215} and A. \textsc{Taylor} \cite{taylor321}. Parts of these sections can easily be extended to the case of a locally convex linear topological space. See, e.g., section $13$ of the following chapter. 
   % VIII
% ==== IX ====
\yoschapter{IX}{Analytical Theory of Semi-groups}
\label{chap:IX}

The analytical theory of semi-groups of bounded linear operators in a $B$-space deals with the exponential functions in infinite dimensional function spaces. It is concerned with the problem of determining the most general bounded linear operator valued function $T(t)$, $t \geqq 0$, which satisfies the equations 
\[
T (t + s) = T (t) \cdot T (s), T (0) = I.
\]

The problem was investigated by E. \textsc{Hille} \cite{hille98} and K. \textsc{Yosida} \cite{yosida348} independently of each other around $1948$. They introduced the notion of the infinitesimal generator $A$ of $T(t)$ defined by 
\[
A = \slim _{t \downarrow 0} t ^{-1} (T (t) - I),
\]
and discussed the generation of $T(t)$ in terms of $A$ and obtained a characterization of the infinitesimal generator $A$ in terms of the spectral property of $A$. 

The basic result of the semi-group theory may be considered as a natural generalization of the theorem of M. H. \textsc{Stone} \cite{stone312} on one-parameter group of unitary operators in a Hilbert space, which will be explained in a later section. Applications of the theory to stochastic processes and to the integration of the equations of evolution, which include diffusion equations, wave equations and Schrödinger equations, will be discussed in Chapter~\ref{chap:XIV}.\yosstarfootnote{See also Supplementary Notes, p.~\pageref{notes:IX-XIV}.} 

In this chapter, we shall develop the theory of semi-groups of continuous linear operators in locally convex linear topological spaces rather than in Banach spaces. 


\yossection{1}{The Semi-group of Class \texorpdfstring{$(C_0)$}{(C0)}}
\label{sec:IX.1}

\begin{propositionn}[E. \textsc{Hille}] Let $X$ be a $B$-space, and $T_{t}, t \geqq 0$, a one-parameter family of bounded linear operators $\in L(X, X)$ satisfying the semi-group property
\begin{equation}\label{eq:IX.1.1}
T _{t} T _{s} = T _{t+s} \quad \text{ for } \quad t, s > 0.
\end{equation}

If $p(t) = \log \|T_t\|$ is bounded from above on the interval $(0, a)$ for each positive $a$, then 
\begin{equation}\label{eq:IX.1.2}
\lim\limits _{t \rightarrow \infty} t ^{-1} \log \| T _{t} \| = \inf\limits _{t>0} t ^{-1} \log \| T _{t} \|.
\end{equation}

\end{propositionn}
\begin{proof}
 We have $p(t + s) \leqq p(t) + p(s)$ from $\| T_{t+s} \| = \| T_t T_s \| \leqq \| T_t \| \cdot \| T_s \|$. Let $\beta = \inf\limits_{t>0} t^{-1} \cdot p(t)$. $\beta$ is either finite or $-\infty$. Suppose that $\beta$ is finite. We choose, for any $\varepsilon > 0$, a number $a > 0$ in such a way that $p(a) \leqq (\beta + \varepsilon)a$. Let $t > a$ and $n$ be an integer such that $na \leqq t < (n + 1)a$. Then 
\[
\begin{aligned}
&\beta \leqq \frac {p (t)}{t} \leqq \frac {p (n a)}{t} + \frac {p (t - n a)}{t} \leqq \frac {n a}{t} \frac {p (a)}{a} + \frac {p (t - n a)}{t} \leqq \frac {n a}{t} (\beta + \varepsilon) + \frac {p (t - n a)}{t}.
\end{aligned}
\]

By hypothesis, $p(t - na)$ is bounded from above as $t \to \infty$. Thus, letting $t \to \infty$ in the above inequality, we obtain $\lim\limits_{t \to \infty} t^{-1} p(t) = \beta$. The case $\beta = -\infty$ may be treated similarly. 
\end{proof}

\begin{definition}\label{def:IX.1.Definition.1} If $\{T_t; t \geqq 0\} \subseteq L(X, X)$ satisfy the conditions 
\begin{align}
T _{t} T _{s}& = T _{t+s} (\text{ for } t, s \geqq 0), \tag{1'}\label{eq:IX.1.1'} \\
T _{0}& = I, \label{eq:IX.1.3} \\
\slim _ {t \to t _{0}} T _{t} x& = T _ {t _{0}} x \text{ for each } t _{0} \geqq 0 \text{ and each } x \in X, \label{eq:IX.1.4}
\end{align}
then $\{T_t\}$ is called a semi-group of class $(C_0)$. 
\end{definition}

In virtue of the Proposition, we see that a semi-group $\{T_t\}$ of class $(C_0)$ satisfies the condition 
\begin{equation}\label{eq:IX.1.5}
\left\| T _{t} \right\| \leqq M e ^{\beta t} \quad (\text{ for } 0 \leqq t < \infty),
\end{equation}
with constants $M > 0$ and $\beta < \infty$. 

The proof is easy. We have only to show that $\|T_t\|$ is, for any interval $(0, a)$ with $\infty > a > 0$, bounded on $(0, a)$. Assume the contrary and let there exist a sequence $\{t_n\} \subseteq (0, a)$ such that $\|T_{t_{n}}\| > n$ and $\lim\limits_{n \to \infty} t_n = t_0 \leqq a$. By the resonance theorem, $\|T_{t_{n}}x\|$ must be unbounded at least for one $x \in X$, which surely contradicts the strong continuity condition \eqref{eq:IX.1.4}. 

\begin{remarku} By multiplying $e^{-\beta t}$, we may assume that a semi-group $\{T_{t}\}$ of class $(C_{0})$ is equi-bounded: 
\begin{equation}\label{eq:IX.1.6}
\left\| T _{t} \right\| \leqq M \quad (\text{ for } 0 \leqq t < \infty).
\end{equation}
\end{remarku}

If, in particular, $M$ is $\leqq 1$, i.e., if 
\begin{equation}\label{eq:IX.1.7}
\left\| T _{t} \right\| \leqq 1 \quad (\text{ for } 0 \leqq t < \infty),
\end{equation}
then the semi-group $\{T_t\}$ is called a contraction semi-group of class $(C_0)$. As for the strong continuity condition \eqref{eq:IX.1.4}, we have the following 

\begin{theoremu} Let a family $\{T_t; t \geqq 0\}$ of operators $\in L(X, X)$ satisfy \eqref{eq:IX.1.1} and \eqref{eq:IX.1.3}. Then condition \eqref{eq:IX.1.4} is equivalent to the condition 
\begin{equation}\label{eq:IX.1.8}
\wlim _{t \downarrow 0} T _{t} x = x \quad \text{ for every } \quad x \in X.
\end{equation}
\end{theoremu}

\begin{proof}
 Suppose that \eqref{eq:IX.1.8} is satisfied. Let $x_0$ be any fixed element of $X$. We shall show that $s\text{- }\lim\limits_{t \to t_0} T_t x_0 = T_{t_0} x_0$ for each $t_0 \geqq 0$. Consider the function $x(t) = T_t x_0$. For each $t_0 \geqq 0$, $x(t)$ is weakly right continuous at $t_0$, because $w\text{- }\lim\limits_{t \downarrow t_0} T_t x = w\text{- }\lim\limits_{h \downarrow 0} T_h T_{t_0} x_0 = T_{t_0} x_0$. We next prove that $\|T_t\|$ is bounded in a vicinity of $t = 0$. For, otherwise, there would exist a sequence $\{t_n\} \in$ such that $t_n \downarrow 0$ and $\lim\limits_{n \to \infty} \|T_{t_{n}} x_0\| = \infty$, contrary to the resonance theorem implied by the weak right continuity of $x(t) = T_t x_0$. Thus by \eqref{eq:IX.1.1}, we see that $T_t x_0 = x(t)$ is bounded on any compact interval of $t$. Moreover, $x(t)$ is weakly measurable. For, a right continuous real-valued function $f(t)$ is Lebesgue measurable, as may be proved from the fact that, for any $\alpha$, the set $\{t; f(t) < \alpha\}$ is representable as the union of intervals of positive length. Next let $\{t_n\}$ be the totality of positive rational numbers, and consider finite linear combinations $\sum\limits_j \beta_j x(t_j)$ where $\beta_j$ are rational numbers (if $X$ is a complex linear space, we take $\beta_j = a_j + ib_j$ with rational coefficients $a_j$ and $b_j$). These elements form a countable set $M = \{x_n\}$ such that $\{x(t); t \geqq 0\}$ is contained in the strong closure of $M$. If otherwise, there would exist a number $t'$ such that $x(t')$ does not belong to $M^a$. But, being a closed linear subspace of $X$, $M^a$ is weakly closed by Theorem~\ref{thm:V.1.Theorem.11} in Chapter~\ref{chap:V}, \ref{sec:V.1}; consequently, the condition $x(t') \in M^a$ is contradictory to the weak right continuity of $x(t)$, i.e., to $x(t') = w\text{- }\lim\limits_{t_n \downarrow t'} x(t_n)$. 

We may thus apply Pettis' theorem in Chapter~\ref{chap:V}, \ref{sec:V.4} to the effect that $x(t)$ is strongly measurable, and so by the boundedness of $\|x(t)\|$ on any compact interval of $t$, we may define the Bochner integral $\int_{\alpha}^{\beta} x(t) \dd t$ and we have $\left\| \int_{\alpha}^{\beta} x(t) \dd t \right\| \leqq \int_{\alpha}^{\beta} \|x(t)\| \dd t$ for $0 \leqq \alpha < \beta < \infty$. By virtue of the strong continuity in $s$ of the integral $\int_{\alpha}^{\beta} x(t + s) \dd t = \int_{\alpha + s}^{\beta + s} x(t) \dd t$, which is implied by the boundedness of $x(t)$ on any compact interval of $t$, N. \textsc{Dunford} \cite{dunford46} proved that $x(t)$ is strongly continuous in $t > 0$. We shall follow his proof. 

Let $0 \leqq \alpha < \eta < \beta < \xi - \varepsilon < \xi$ with $\varepsilon > 0$. Since $x(\xi) = T_{\xi}x_0 = T_{\eta}T_{\xi - \eta}x_0 = T_{\eta}x(\xi - \eta)$, we have 
\[
(\beta - \alpha) x (\xi) = \int_{\alpha} ^{\beta} x (\xi) \dd \eta = \int_{\alpha} ^{\beta} T _{\eta} x (\xi - \eta) \dd \eta ,
\]
and so, by $\sup\limits_{\alpha \leqq \eta \leqq \beta}\|T_{\eta}\| < \infty$ which is implied by \eqref{eq:IX.1.1} and \eqref{eq:IX.1.3}, combined with the boundedness of $\|T_t\|$ near $t = 0$, we obtain 
\begin{align}
(\beta - \alpha) \{x (\xi \pm \varepsilon) - x (\xi) \}& = \int_{\alpha} ^{\beta} T _{\eta} \{x (\xi \pm \varepsilon - \eta) - x (\xi - \eta) \} \dd \eta , \notag \\
(\beta - \alpha) \| x (\xi \pm \varepsilon) - x (\xi) \|& \leqq \sup\limits _{\alpha \leqq \eta \leqq \beta} \| T _{\eta} \| \cdot \int_{\xi - \beta} ^{\xi - \alpha} \| x (\tau \pm \varepsilon) - x (\tau) \| \dd \tau . \notag
\end{align}

The right hand side tends to zero as $\varepsilon \downarrow 0$, as may be seen by approximating $x(\tau)$ by finitely-valued functions. 

We have thus proved that $x(t)$ is strongly continuous in $t > 0$. To prove the strong continuity of $x(t)$ at $t = 0$, we proceed as follows. For any positive rational number $t_n$, we have $T_t x(t_n) = T_t T_{t_{n}} x_0 = T_{t+t_n} x_0 = x(t + t_n)$. Hence, by the strong continuity of $x(t)$ for $t > 0$ proved above, we have $\slim\limits_{t \downarrow 0} T_t x(t_n) = x(t_n)$. Since each $x_m \in M$ is a finite linear combination of $x(t_n)'s$, we have $\slim\limits_{t \downarrow 0} T_t x_m = x_m (m = 1, 2, \ldots)$. On the other hand, we have, for any $t \in [0, 1]$, 
\begin{align}
\left\| x (t) - x _{0} \right\|& \leqq \left\| T _{t} x _{m} - x _{m} \right\| + \left\| x _{m} - x _{0} \right\| + \left\| T _{t} \left(x _{0} - x _{m}\right) \right\| \notag \\
&\leqq \left\| T _{t} x _{m} - x _{m} \right\| + \left\| x _{m} - x _{0} \right\| + \sup\limits _{0<t<1} \left\| T _{t} \right\| \cdot \left\| x _{0} - x _{m} \right\|. \notag
\end{align}

Hence $\varlimsup\limits_{t\downarrow 0}\|x(t) - x_0\|\leqq (1 + \sup\limits_{0\leqq t\leqq 1}\|T_t\|)\|x_m - x_0\|,$ and so $\slim\limits_{t\downarrow 0}x(t)$ $= x_0$ by $\inf\limits_{x_m\in M}\|x_0 - x_m\| = 0.$ 
\end{proof}


\yossection[2. The Equi-continuous Semi-group of Class \texorpdfstring{$(C_0)$}{(C0)}]{2}{The Equi-continuous Semi-group of Class \texorpdfstring{$(C_0)$}{(C0)} in Locally Convex Spaces. Examples of Semi-groups}
\label{sec:IX.2}

Suggested by the preceding section, we shall pass to a more general class of semi-groups. 

\begin{definitionu} Let $X$ be a locally convex linear topological space, and $\{T_t; t \geqq 0\}$ be a one-parameter family of continuous linear operators $\in L(X, X)$ such that 
\begin{align}
T _{t} T _{s}& = T _{t+s}, T _{0} = I, \label{eq:IX.2.1} \\
\lim\limits _ {t \rightarrow t _{0}} T _{t} x& = T _ {t _{0}} x \quad \text{ for any } \quad t _{0} \geqq 0 \quad \text{ and } \quad x \in X, \label{eq:IX.2.2}
\end{align}
the family of mappings $\{T_t\}$ is equi-continuous in $t$, 
i.e., for any continuous semi-norm $\phi$ on $X$, there 
exists a continuous semi-norm $q$ on $X$ such that 
\begin{equation}\label{eq:IX.2.3}
p \left(T _{t} x\right) \leqq q (x) \text{ for all } t \geqq 0 \text{ and all } x \in X.
\end{equation}
\end{definitionu}

Such a family $\{T_t\}$ is called an equi-continuous semi-group of class $(C_0)$. 

The semi-groups $\{T_{t}\}$ satisfying conditions \eqref{eq:IX.2.1}, \eqref{eq:IX.2.3}, ($4$) and ($6$) of Section~\hyperref[sec:IX.1]{$1$} are example of such equi-continuous semi-groups of class $(C_0)$. We shall give concrete examples. 

\begin{example}\label{ex:IX.2.Example.1} Let $C[0, \infty]$ be the space of bounded uniformly continuous real-valued (or complex-valued) functions on the interval $[0, \infty)$, and define $T_t, t \geqq 0$, on $C[0, \infty]$ into $C[0, \infty]$ by 
\[
(T _{t} x) (s) = x (t + s).
\]
\end{example}

Condition \eqref{eq:IX.2.1} is trivially satisfied. \eqref{eq:IX.2.2} follows from the uniform continuity of $x(s)$. Finally $\|T_t\| \leqq 1$ and so $\{T_t\}$ is a contraction semigroup of class $(C_0)$. In this example, we could replace $C[0,\infty]$ by $C[-\infty, \infty]$ or by $L^p(-\infty, \infty)$. 

\begin{example}\label{ex:IX.2.Example.2} Consider the space $C[-\infty, \infty]$. Let 
\[
N _{t} (u) = (2 \pi t) ^{-1/2} e ^ {- u ^{2} / 2 t}, - \infty < u < \infty , t > 0,
\]
which is the Gaussian probability density. Define $T_{t}, t \geqq 0$, on $C[-\infty, \infty]$ into $C[-\infty, \infty]$ by 
\[
\begin{aligned} (T _{t} x) (s) & = \int_{- \infty} ^{\infty} N _{t} (s - u) x (u) \dd u, \text{ for } t > 0, \\ & = x (s), \text{ for } t = 0. \end{aligned}
\]
\end{example}

Each $T_{t}$ is continuous, since, by $\int_{-\infty}^{\infty} N_{t}(s - u) \dd u = 1$, 
\[
\left\| T _{t} x \right\| \leqq \left\| x \right\| \int_{- \infty} ^{\infty} N _{t} (s - u) \dd u = \left\| x \right\|.
\]
$T_{0}=I$ by the definition, and the semi-group property $T_{t}T_{s}=T_{t+s}$ is a consequence from the well-known formula concerning the Gaussian probability distribution: 
\[
\frac {1}{\sqrt {2 \pi (t + t ^{\prime})}} e ^ {- u ^{2} / 2 (t + t ^{\prime})} = \frac {1}{\sqrt {2 \pi t}} \frac {1}{\sqrt {2 \pi t ^{\prime}}} \int_{- \infty} ^{\infty} e ^ {- (u - v) ^{2} / 2 t} e ^ {- v ^{2} / 2 t ^{\prime}} \dd v.
\]

This formula may be proved by applying the Fourier transformation on both sides, remembering ($10$) and \eqref{eq:VI.1.13} in Chapter~\ref{chap:VI}, \ref{sec:VI.1}. To prove the strong continuity in $t$ of $T_{t}$, we observe that $x(s) = \int_{-\infty}^{\infty} N_{t}(s - u) x(s) \dd u$. Thus 
\[
(T _{t} x) (s) - x (s) = \int_{- \infty} ^{\infty} N _{t} (s - u) (x (u) - x (s)) \dd u,
\]
which is, by the change of variable $(s - u) / \sqrt{t} = z$, equal to 
\[
\int_{- \infty} ^{\infty} N _{1} (z) \left(x (s - \sqrt {t} z) - x (s)\right) \dd z.
\]

By the uniform continuity of $x(s)$ on $(-\infty, \infty)$, there exists, for any $\varepsilon > 0$, a number $\delta = \delta(\varepsilon) > 0$ such that $|x(s_1) - x(s_2)| \leqq \varepsilon$ whenever $|s_1 - s_2| \leqq \delta$. Splitting the last integral, we obtain 
\[
\begin{aligned}
&\left| (T _{t} x) (s) - x (s) \right| \leqq \int_ {| \sqrt {t} z | \leqq \delta} N _{1} (z) \left| x (s - \sqrt {t} \cdot z) - x (s) \right| \dd z + \int_ {| \sqrt {t} z | > \delta} (\dots) \dd z
\\
&\leqq \varepsilon \int_ {| \sqrt {t} z | \leqq \delta} N _{1} (z) \dd z + 2 \| x \| \int_ {| \sqrt {t} z | > \delta} N _{1} (z) \dd z \leqq \varepsilon + 2 \| x \| \int_ {| \sqrt {t} z | > \delta} N _{1} (z) \dd z.
\end{aligned}
\]

The second term on the right tends to $0$ as $t \to 0$, since the integral $\int_{-\infty}^{\infty} N_1(z) \dd z$ converges. Thus $\lim\limits_{t \downarrow 0} \sup\limits_s |(T_t x)(s) - x(s)| = 0$ and hence $\slim\limits_{t \downarrow 0} T_t x = x$; consequently, by the Theorem in the preceding section, we have proved \eqref{eq:IX.2.2}. 

In the above example, we can replace $C[-\infty, \infty]$ by $L^p(-\infty, \infty)$. Consider, for example, $L^1(-\infty, \infty)$. In this case, we have 
\[
\left\| T _{t} x \right\| \leqq \iint_{- \infty} ^{\infty} N _{t} (s - u) | x (u) | \dd s \dd u \leqq \left\| x \right\|,
\]
by Fubini's theorem. As for the strong continuity, we have, as above, 
\[
\begin{aligned}
&\| T _{t} x - x \|  = \int_{- \infty} ^{\infty} \int N _{1} (z) (x (s - \sqrt {t} \cdot z) - x (s)) \dd z | \dd s \\
&\leqq \int_{- \infty} ^{\infty} N _{1} (z) \left[ \int_{- \infty} ^{\infty} | x (s - \sqrt {t} \cdot z) - x (s) | \dd s \right] \dd z \leqq 2 \int_{- \infty} ^{\infty} N _{1} (z) \dd z \cdot \| x \|.
\end{aligned}
\]

Hence, by the Lebesgue-Fatou lemma, we obtain 
\[
\varlimsup\limits_{t \downarrow 0} \| T _{t} x - x \| \leqq \int_{- \infty} ^{\infty} N _{1} (z) \left(\varlimsup\limits_{t \downarrow 0} \int_{- \infty} ^{\infty} | x (s - \sqrt {t} \cdot z) - x (s) | \dd s\right) \dd z = 0,
\]
because of the continuity in the mean of the Lebesgue integral, which may be proved by approximating $x(s)$ by finitely-valued functions. 

\begin{example}\label{ex:IX.2.Example.3} Consider $C[-\infty, \infty]$. Let $\lambda > 0$, $\mu > 0$. Define $T_t$, $t \geqq 0$, on $C[-\infty, \infty]$ into $C[-\infty, \infty]$ by 
\[
\left(T _{t} x\right) (s) = e ^{- \lambda t} \sum\limits_{k=0} ^{\infty} \frac {(\lambda t) ^{k}}{k !} x (s - k \mu).
\]
$3$. The Infinitesimal Generator of an Equi-continuous Semi-group We have 
\[
\begin{aligned}
\left(T _{w} (T _{t} x)\right) (s) &= e ^{- \lambda w} \sum\limits_{m=0} ^{\infty} \frac {(\lambda w) ^{m}}{m !} \left[ e ^{- \lambda t} \sum\limits_{k=0} ^{\infty} \frac {(\lambda t) ^{k}}{k !} x (s - k \mu - m \mu) \right]
\\
&= e ^{- \lambda (w + t)} \sum\limits_{p=0} ^{\infty} \frac {1}{p !} \left[ p! \sum\limits_{m=0} ^{p} \frac {(\lambda w) ^{m}}{m !} \frac {(\lambda t) ^{p-m}}{(p - m) !} x (s - p \mu) \right]
\\
&= e ^{- \lambda (w + t)} \sum\limits_{p=0} ^{\infty} \frac {1}{p !} (\lambda w + \lambda t) ^{p} x (s - p \mu) = (T _{w+t} x) (s).
\end{aligned}
\]
\end{example}

Thus it is easy to verify that $T_{t}$ is a contraction semi-group of class $(C_0)$. 


\yossection{3}{The Infinitesimal Generator of an Equi-continuous Semi-group of Class \texorpdfstring{$(C_0)$}{(C0)}}
\label{sec:IX.3}

Let $\{T_t; t \geqq 0\}$ be an equi-continuous semi-group of class $(C_0)$ defined on a locally convex linear topological space $X$ which we assume to be sequentially complete. We define the infinitesimal generator $A$ of $T_t$ by 
\begin{equation}\label{eq:IX.3.1}
A x = \lim\limits _{h \downarrow 0} h ^{-1} (T _{h} - I) x,
\end{equation}
i.e., $A$ is the linear operator whose domain is the set $D(A) = \{x \in X; \lim\limits_{h \downarrow 0} h^{-1}(T_h - I) x \text{ exists in } X\}$, and, for $x \in D(A)$, $Ax = \lim\limits_{h \downarrow 0} h^{-1}(T_h - I) x$. $D(A)$ is non-empty; it contains at least the vector $0$. Actually $D(A)$ is larger. We can prove 

\begin{theorem}\label{thm:IX.3.Theorem.1} $D(A)$ is dense in $X$. 
\end{theorem}

\begin{proof}
 Let $\varphi_{n}(s) = ne^{-ns}$, $n > 0$. Consider the linear operator $C_{\varphi_n}$ which is the Laplace transform of $T_{t}$ multiplied by $n$: 
\begin{equation}\label{eq:IX.3.2}
C _ {\varphi_{n}} x = \int_{0} ^{\infty} \varphi_{n} (s) T _{s} x \dd s \quad \text{ for } \quad x \in X,
\end{equation}
the integral being defined in the sense of Riemann. The ordinary procedure of defining the Riemann integral of numerical functions can be extended to a function with values in a locally convex, sequentially complete space $X$, using continuous semi-norms $p$ on $X$ in place of the absolute value of a number. The convergence of the improper integral is a consequence of the equi-continuity of $T_{t}$, the inequality 
\[
p \left(\varphi_{n} (s) T _{s} x\right) = n e ^{-ns} p \left(T _{s} x\right)
\]
and the sequential completeness of $X$. 


We see, by 
\[
\phi \left(C _ {\varphi_{n}} x\right) \leqq \int_{0} ^{\infty} n e ^{-ns} \phi \left(T _{s} x\right) \dd s \leqq \sup\limits _{s \geqq 0} \phi \left(T _{s} x\right),
\]
that the operator $C_{\varphi_n}$ is a continuous linear operator $\in L(X, X)$. We shall show that 
\begin{equation}\label{eq:IX.3.3}
R \left(C _ {\varphi_{n}}\right) \subseteq D (A) \quad \text{ for each } n > 0,
\end{equation}
and 
\begin{equation}\label{eq:IX.3.4}
\lim\limits _{n \rightarrow \infty} C _ {\varphi_{n}} x = x \quad \text{ for each } \quad x \in X.
\end{equation}

Then $\bigcup\limits_{n=1}^{\infty} R(C_{\varphi_n})$ will be dense in $X$, and a fortiori $D(A)$ will be dense in $X$. To prove \eqref{eq:IX.3.3}, we start with the formula 
\[
h ^{-1} (T _{h} - I) C _ {\varphi_{n}} x = h ^{-1} \int_{0} ^{\infty} \varphi_{n} (s) T _{h} T _{s} x \dd s - h ^{-1} \int_{0} ^{\infty} \varphi_{n} (s) T _{s} x \dd s.
\]

The change of the order: $T_{h}\int\limits_{0}^{\infty}\cdots=\int\limits_{0}^{\infty}T_{h}\cdots$ is justified, using the linearity and the continuity of $T_{h}$. Thus 
\[
\begin{aligned} h ^{-1} (T _{h} - I) C _ {\varphi_{n}} x & = h ^{-1} \int_{0} ^{\infty} \varphi_{n} (s) T _{s+h} x \dd s - h ^{-1} \int_{0} ^{\infty} \varphi_{n} (s) T _{s} x \dd s \\ & = \frac {e ^{nh} - 1}{h} n \int_{h} ^{\infty} e ^{- n \sigma} T _{\sigma} x \dd \sigma - \frac {1}{h} n \int_{0} ^{h} e ^{-ns} T _{s} x \dd s \\ & = \frac {e ^{nh} - 1}{h} \left\{C _ {\varphi_{n}} x - \int_{0} ^{h} n e ^{- n \sigma} T _{\sigma} x \dd \sigma \right\} - \frac {1}{h} \int_{0} ^{h} \varphi_{n} (s) T _{s} x \dd s. \end{aligned}
\]

By the continuity of $\varphi_{n}(s) T_{s}x$ in $s$, the second term on the right tends to $-\varphi_{n}(0) T_{0}x = -nx$ as $h \downarrow 0$. Similarly, the first term on the right tends to $nC_{\varphi_n}x$ as $h \downarrow 0$. Hence we have 
\begin{equation}\label{eq:IX.3.5}
A C _ {\varphi_{n}} x = n \left(C _ {\varphi_{n}} - I\right) x, x \in X.
\end{equation}

We next prove \eqref{eq:IX.3.4}. We have, by $\int_0^\infty ne^{-ns}\dd s = 1$, 
\begin{align}
C _ {\varphi_{n}} x - x& = n \int_{0} ^{\infty} e ^{-ns} (T _{s} x - x) \dd s, \notag \\
p \left(C _ {\varphi_{n}} x - x\right)& \leqq n \int_{0} ^{\infty} e ^{-ns} p \left(T _{s} x - x\right) \dd s = n \int_{0} ^{\delta} + n \int_{\delta} ^{\infty} = I _{1} + I _{2}, \text{ say }, \notag
\end{align}
where $\delta > 0$ is a positive number. For any $\varepsilon > 0$, we can choose, by the continuity in $s$ of $T_s x$, a $\delta > 0$ such that $p(T_s x - x) \leqq \varepsilon$ for $0 \leqq s \leqq \delta$. Then 
\[
I _{1} \leqq \varepsilon n \int_{0} ^{\delta} e ^{-ns} \dd s \leqq \varepsilon n \int_{0} ^{\infty} e ^{-ns} \dd s = \varepsilon .
\]

For a fixed $\delta > 0$, we have, by the equi-boundedness of $\{T_s x\}$ in $s \geqq 0$, 
\[
I _{2} \leqq n \int_{\delta} ^{\infty} e ^{-ns} (p (T _{s} x) + p (x)) \dd s \rightarrow 0 \quad \text{ as } \quad n \uparrow \infty .
\]

Hence we have proved \eqref{eq:IX.3.4}.
\end{proof}

\begin{definitionu} For $x \in X$, we define 
\begin{equation}\label{eq:IX.3.6}
D _{t} T _{t} x = \lim\limits _{h \rightarrow 0} h ^{-1} (T _{t+h} - T _{t}) x
\end{equation}
if the right hand limit exists. 
\end{definitionu}

\begin{theorem}\label{thm:IX.3.Theorem.2} If $x \in D(A)$, then $x \in D(D_t T_t)$ and 
\begin{equation}\label{eq:IX.3.7}
D _{t} T _{t} x = A T _{t} x = T _{t} A x, t \geqq 0.
\end{equation}
\end{theorem}

Thus, in particular, the operator $A$ is commutative with $T_{t}$. 

\begin{proof}
 If $x \in D(A)$, then we have, since $T_{t}$ is continuous linear, 
\[ \begin{aligned} T _{t} A x & = T _{t} \lim\limits _{h \downarrow 0} h ^{-1} (T _{h} - I) x = \lim\limits _{h \downarrow 0} h ^{-1} (T _{t} T _{h} - T _{t}) x = \lim\limits _{h \downarrow 0} h ^{-1} (T _{t+h} - T _{t}) x \\ & = \lim\limits _{h \downarrow 0} h ^{-1} (T _{h} - I) T _{t} x = A T _{t} x. \end{aligned}
\]

Thus, if $x \in D(A)$, then $T_{t}x \in D(A)$ and $T_{t}Ax = AT_{t}x = \lim\limits_{h\downarrow 0} h^{-1}(T_{t+h} - T_{t})x$. We have thus proved that the right derivative of $T_{t}x$ exists for each $x \in D(A)$. We shall show that the left derivative also exists and is equal to the right derivative. 

For this purpose, take any $f_0 \in X'$. Then, for a fixed $x \in D(A)$, the numerical function $f_0(T_t x) = \langle T_t x, f_0 \rangle$ is continuous in $t \geqq 0$ and has right derivative $d^+ f_0(T_t x) / \dd t$ which is equal to $f_0(A T_t x) = f_0(T_t A x)$ by what we have proved above. Hence $d^+ f_0(T_t x) / \dd t$ is continuous in $t$. We shall prove below a well-known Lemma: if one of the Dini derivatives 
\[
\overline {{{D}}} ^{+} f (t), \underline {{{D}}} ^{+} f (t), \overline {{{D}}} ^{-} f (t) \quad \text{ and } \quad \underline {{{D}}} ^{-} f (t)
\]
of a continuous real-valued function $f(t)$ is finite and continuous, then $f(t)$ is differentiable and the derivative is, of course, continuous and equals $\overline{D}^{\pm}f(t)$. Thus $f_0(T_tx)$ is differentiable in $t$ and 
\[
\begin{aligned} f _{0} (T _{t} x - x) & = f _{0} (T _{t} x) - f _{0} (T _{0} x) = \int_{0} ^{t} d ^{+} f _{0} (T _{s} x) / \dd s \cdot \dd s = \int_{0} ^{t} f _{0} (T _{s} A x) \dd s \\ & = f _{0} \left(\int_{0} ^{t} T _{s} A x \dd s\right). \end{aligned}
\]

Since $f_0 \in X'$ was arbitrary, we must have 
\[
T _{t} x - x = \int_{0} ^{t} T _{s} A x \dd s \quad \text{ for each } \quad x \in D (A).
\]

Since $T_{s}Ax$ is continuous in $s$, it follows that $T_{t}x$ is differentiable in $t$ in the topology of $X$ and 
\[
D _{t} T _{t} x = \lim\limits _{h \rightarrow 0} h ^{-1} \int_{t} ^{t+h} T _{s} A x \dd s = T _{t} A x.
\]

Thus we have proved \eqref{eq:IX.3.7}. 
\end{proof}

\begin{proof}[Proof of the Lemma.]
We first prove that the condition: $\overline{D}^{+}f(t) \geqq 0$ for $a \leqq t \leqq b$ implies that $f(b) - f(a) \geqq 0$. Assume the contrary, and let $f(b) - f(a) < -\varepsilon (b - a)$ with some $\varepsilon > 0$. Then, for $g(t) = f(t) - f(a) + \varepsilon (t - a)$, we have $\overline{D}^{+}g(a) = \overline{D}^{+}f(a) + \varepsilon > 0$, and so, by $g(a) = 0$, we must have $g(t_0) > 0$ for some $t_0 > a$ near $a$. By the continuity of $g(t)$ and $g(b) < 0$, there must exist a $t_1$ with $a < t_0 < t_1 < b$ such that $g(t_1) = 0$ and $g(t) < 0$ for $t_1 < t < b$. We have then $\overline{D}^{+}g(t_1) \leqq 0$ which surely contradicts the fact $\overline{D}^{+}g(t_1) = \overline{D}^{+}f(t_1) + \varepsilon > 0$.

By applying a similar argument to $f(t) - \alpha t$ and to $\beta t - f(t)$, we prove the following: if one of the Dini derivatives $Df(t)$ satisfies 
\[
\alpha \leqq D f (t) \leqq \beta \text{ on any interval } [ t _{1}, t _{2} ],
\]
then $\alpha \leqq (f(t_2) - f(t_1)) / (t_2 - t_1) \leqq \beta$. Hence, the suprema (and the infima) on $[t_1, t_2]$ of the four Dini derivatives of a continuous real-valued function $f(t)$ are the same. Thus, in particular, if one of the Dini derivatives of a continuous real-valued function $f(t)$ is continuous on $[t_1, t_2]$, then the four Dini derivatives of $f(t)$ must coincide on $[t_1, t_2]$. 
\end{proof}


\yossection{4}{The Resolvent of the Infinitesimal Generator \texorpdfstring{$A$}{A}}
\label{sec:IX.4}

\begin{theorem}\label{thm:IX.4.Theorem.1} If $n > 0$, then the operator $(nI - A)$ admits an inverse $R(n;A) = (nI - A)^{-1} \in L(X,X)$, and 
\begin{equation}\label{eq:IX.4.1}
R (n; A) x = \int_{0} ^{\infty} e ^{-ns} T _{s} x \dd s \quad \text{ for } \quad x \in X.
\end{equation}
\end{theorem}

In other words, positive real numbers belong to the resolvent set $\varrho(A)$ of $A$. 

\begin{proof}
 We first show that $(nI-A)^{-1}$ exists. Suppose that there exists an $x_{0}\neq0$ such that $(nI-A)x_{0}=0$, that is, $Ax_{0}=nx_{0}$. Let $f_{0}$ be a continuous linear functional $\in X'$ such that $f_{0}(x_{0})=1$, and set $\varphi(t)=f_{0}(T_{t}x_{0})$. Since $x_{0}\in D(A)$, $\varphi(t)$ is differentiable by Theorem~\ref{thm:IX.4.Theorem.2} in the preceding section and 
\[
\dd \varphi (t) / \dd t = f _{0} \left(D _{t} T _{t} x _{0}\right) = f _{0} \left(T _{t} A x _{0}\right) = f _{0} \left(T _{t} n x _{0}\right) = n \varphi (t). \]

If we solve this differential equation under the initial condition $\varphi(0)=f_{0}(x_{0})=1$, we get $\varphi(t)=e^{nt}$. But, $\varphi(t)=f_{0}(T_{t}x_{0})$ is bounded in $t$, because of the equi-boundedness of $T_{t}x_{0}$ in $t\geqq0$ and the continuity of the functional $f_{0}$. This is a contradiction, and so the inverse $(nI-A)^{-1}$ must exist. 

Since, by \eqref{eq:IX.4.5} of the preceding section, $A C_{\varphi_n}x = n(C_{\varphi_n} - I)x$, we have $(nI - A)C_{\varphi_n}x = nx$ for all $x \in X$. Thus $(nI - A)$ maps $R(C_{\varphi_n}) \subseteq D(A)$ onto $X$ in a one-one way. Hence, a fortiori, $(nI - A)$ must map $D(A)$ onto $X$ in a one-one way, in virtue of the existence of $(nI - A)^{-1}$. Therefore, we must have $R(C_{\varphi_n}) = D(A)$ and $(nI - A)^{-1} = n^{-1}C_{\varphi_n}$. 
\end{proof}
\begin{corollary}\label{cor:IX.4.Corollary.1} The right half plane of the complex $\lambda$-plane is in the resolvent set $\varrho(A)$ of $A$, and we have 
\begin{equation}\label{eq:IX.4.2}
R (\lambda ; A) x = (\lambda I - A) ^{-1} x = \int_{0} ^{\infty} e ^{- \lambda t} T _{t} x \dd t \text{ for } \operatorname{Re} (\lambda) > 0 \text{ and } x \in X.
\end{equation}
\end{corollary}

\begin{proof}
 For real fixed $\tau, \{e^{-i\tau t} T_t; t \geqq 0\}$, constitutes an equi-continuous semi-group of class $(C_0)$. It is easy to see that the infinitesimal generator of this semi-group is equal to $(A - i\tau I)$. Thus, for any $\sigma > 0$, the resolvent $R(\sigma + i\tau; A) = ((\sigma + i\tau)I - A)^{-1}$ exists and 
\begin{equation}\tag{2'}\label{eq:IX.4.2'}
R \left((\sigma + i \tau) I - A\right) x = \int_{0} ^{\infty} e ^{-(\sigma + i \tau) s} T _{s} x \dd s \quad \text{ for all } \quad x \in X. \end{equation}
\end{proof}

\begin{corollary}\label{cor:IX.4.Corollary.2} 
\begin{align}
D (A)& = R ((\lambda I - A) ^{-1}) = R (R (\lambda ; A)) \quad \text{ when } \quad \operatorname{Re} (\lambda) > 0, \label{eq:IX.4.3} \\
A R (\lambda ; A) x& = R (\lambda ; A) A x = (\lambda R (\lambda ; A) - I) x \text{ for } x \in D (A) \label{eq:IX.4.4} \\
A R (\lambda ; A) x& = (\lambda R (\lambda ; A) - I) x \quad \text{ for } \quad x \in X, \label{eq:IX.4.5} \\
\lim\limits _{\lambda \uparrow \infty} \lambda R (\lambda ; A) x& = x \quad \text{ for } \quad x \in X. \label{eq:IX.4.6}
\end{align}
\end{corollary}

\begin{proof}
 Clear from $R(\lambda; A) = (\lambda I - A)^{-1}$ and \eqref{eq:IX.4.4} of the preceding section. 
\end{proof}

\begin{corollary}\label{cor:IX.4.Corollary.3} The infinitesimal generator $A$ is a closed operator in the following sense (Cf. Chapter~\ref{chap:II}, \ref{sec:II.6}): 
\begin{align}
\text{ if } x _{h}& \in D (A) \quad \text{ and } \quad \lim\limits _{h \to \infty} x _{h} = x \in X, \lim\limits _{h \to \infty} A x _{h} = y \in X, \text{ then } \notag \\
x& \in D (A) \quad \text{ and } \quad A x = y. \notag
\end{align}
\end{corollary}

\begin{proof}
 Put $(I - A)x_{h} = z_{h}$. Then $\lim\limits_{h\to \infty}z_h = x - y$ and so, by the continuity of $(I - A)^{-1}$, $\lim\limits_{h\to \infty}x_h = \lim\limits_{h\to \infty}(I - A)^{-1}z_h = (I - A)^{-1}(x - y)$ that is, $x = (I - A)^{-1}(x - y)$, $(I - A)x = x - y$. This proves that $y = Ax$. 
\end{proof}

\begin{theorem}\label{thm:IX.4.Theorem.2} The family of operators 
\begin{equation}\label{eq:IX.4.7}
\{(\lambda R (\lambda ; A)) ^{n} \}
\end{equation}
is equi-continuous in $\lambda > 0$ and in $n = 0, 1, 2, \ldots$. 
\end{theorem}

\begin{proof}
 From the resolvent equation (Chapter~\ref{chap:VIII}, \ref{sec:VIII.2}, \eqref{eq:VIII.2.2}) 
\[
R (\mu ; A) - R (\lambda ; A) = (\lambda - \mu) R (\mu ; A) R (\lambda ; A),
\]
we obtain 
\[
\lim\limits _{\mu \rightarrow \lambda} (\mu - \lambda) ^{-1} (R (\mu ; A) - R (\lambda ; a)) x = \dd R (\lambda ; A) x / \dd \lambda = - R (\lambda ; A) ^{2} x, x \in X. \]

To derive the above formula, we have to appeal to \eqref{eq:IX.4.2} in order to show that $\lim\limits_{\mu \to \lambda} R(\mu; A)y = R(\lambda; A)y, y \in X$. 


Therefore, $R(\lambda; A) x$ is infinitely differentiable with respect to $\lambda$ when $\operatorname{Re}(\lambda) > 0$ and 
\begin{equation}\label{eq:IX.4.8}
d ^{n} R (\lambda ; A) x / \dd \lambda^{n} = (- 1) ^{n} n! R (\lambda ; A) ^{n+1} x \quad (n = 0, 1, 2, \dots).
\end{equation}

On the other hand, we have, by differentiating \eqref{eq:IX.4.2} $n$-times with respect to $\lambda$, 
\begin{equation}\label{eq:IX.4.9}
d ^{n} R (\lambda ; A) x / \dd \lambda^{n} = \int_{0} ^{\infty} e ^{- \lambda t} (- t) ^{n} T _{t} x \dd t.
\end{equation}

Here the differentiation under the integral sign is justified since $\{T_t x\}$ is equi-bounded in $t$ and $\int_0^\infty e^{-\lambda t} t^n \dd t = (n!) / \lambda^{n+1}$ when $\operatorname{Re}(\lambda) > 0$. Hence 
\begin{equation}\label{eq:IX.4.10}
(\lambda R (\lambda ; A)) ^{n+1} x = \frac {\lambda^{n+1}}{n !} \int_{0} ^{\infty} e ^{- \lambda t} t ^{n} T _{t} x \dd t \text{ for } x \in X \text{ and } \operatorname{Re} (\lambda) > 0,
\end{equation}
and so, for any continuous semi-norm $p$ on $X$ and $\lambda > 0$, $n > 0$, 
\begin{equation}\label{eq:IX.4.11}
\phi ((\lambda R (\lambda ; A)) ^{n+1} x) \leqq \frac {\lambda^{n+1}}{n !} \int_{0} ^{\infty} e ^{- \lambda t} t ^{n} \dd t \cdot \sup\limits _{t \geqq 0} \phi (T _{t} x) = \sup\limits _{t \geqq 0} \phi (T _{t} x).
\end{equation}

This proves Theorem~\ref{thm:IX.4.Theorem.2} by the equi-continuity of $\{T_t\}$ in $t$. 
\end{proof}


\yossection{5}{Examples of Infinitesimal Generators}
\label{sec:IX.5}

We first define, for $n > 0$, 
\begin{equation}\label{eq:IX.5.1}
J _{n} = (I - n ^{-1} A) ^{-1} = n R (n; A),
\end{equation}
so that 
\begin{equation}\label{eq:IX.5.2}
A J _{n} = n (J _{n} - I).
\end{equation}

\begin{example}\label{ex:IX.5.Example.1} $(T_{t}x)(s) = x(t + s)$ on $C[0,\infty]$. 
\end{example}

Writing $y_{n}(s) = (J_{n}x)(s) = n\int \limits_{0}^{\infty}e^{-nt}x(t + s)\dd t = n\int \limits_{s}^{\infty}e^{-n(t-s)}x(t)\dd t,$ 
we obtain $y_n'(s) = -n e^{-n(s-s)} x(s) + n^2 \int_s^{\infty} e^{-n(t-s)} x(t) \dd t = -n x(s) + n y_n(s)$. Comparing this with the general formula \eqref{eq:IX.5.2}: 
\[
(A J _{n} x) (s) = n ((J _{n} - I) x) (s),
\]
we obtain $A y_{n}(s) = y_{n}^{\prime}(s)$. Since $R(J_{n}) = R(R(n;A)) = D(A)$, we have $A y(s) = y^{\prime}(s)$ for every $y \in D(A)$. 

Conversely, let $y(s)$ and $y'(s)$ both belong to $C[0, \infty]$. We will then show that $y \in D(A)$ and $A y(s) = y'(s)$. To this purpose, define $x(s)$ by 
\[
y ^{\prime} (s) - n y (s) = - n x (s).
\]

Setting $(J_{n}x)(s) = y_{n}(s)$, we obtain, as shown above, 
\[
y _{n} ^{\prime} (s) - n y _{n} (s) = - n x (s).
\]

Hence $w(s) = y(s) - y_{n}(s)$ satisfies $w'(s) = nw(s)$ and so $w(s) = Ce^{ns}$. But, $w$ must belong to $C[0,\infty]$ and this is possible only if $C = 0$. Hence $y(s) = y_{n}(s) \in D(A)$ and $Ay(s) = y'(s)$. 

Therefore, the domain $D(A)$ of $A$ is precisely the set of functions $y \in C[0, \infty]$ whose first derivative also $\in C[0, \infty]$, and for such a function $y$ we have $A y = y'$. We have thus characterized the differential operator $\dd / \dd s$ as the infinitesimal generator of the semi-group associated with the operation of translation by $t$ on the function space $C[0, \infty]$. 

\begin{example}\label{ex:IX.5.Example.2} We shall give a characterization of the second derivative $d^{2}/\dd s^{2}$ as the infinitesimal generator of the semi-group associated with the integral operator by Gaussian kernel. The space is $C[-\infty,\infty]$ and 
\[
(T _{t} x) (s) = (2 \pi t) ^{-1/2} \int_{- \infty} ^{\infty} e ^ {- (s - v) ^{2} / 2 t} x (v) \dd v \text{ if } t > 0, = x (s) \text{ if } t = 0.
\]
\end{example}

We have 
\[
\begin{aligned}
&y _{n} (s) = (J _{n} x) (s) = \int_{- \infty} ^{\infty} x (v) \left[ \int_{0} ^{\infty} n (2 \pi t) ^{-1/2} e ^ {- n t - (s - v) ^{2} / 2 t} \dd t \right] \dd v
\\
&= \int_{- \infty} ^{\infty} x (v) \left[ 2 \sqrt {n} \int_{0} ^{\infty} (2 \pi) ^{-1/2} e ^ {- \sigma^{2} - n (s - v) ^{2} / 2 \sigma^{2}} \dd \sigma \right] \dd v (\text{ by substituting } t = \sigma^{2} / n).
\end{aligned}
\]

Assuming, for the moment, the formula 
\begin{equation}\label{eq:IX.5.3}
\int_{0} ^{\infty} e ^ {- (\sigma^{2} + c / \sigma^{2})} \dd \sigma = \frac {\sqrt {\pi}}{2} e ^{-2c}, c = \sqrt {n} | s - v | / \sqrt {2} > 0,
\end{equation}
we get 
\[
y _{n} (s) = \int_{- \infty} ^{\infty} x (v) (n / 2) ^{1/2} e ^ {- \sqrt {2 n} | s - v |} \dd v.
\]
$x(v)$ being continuous and bounded, we can differentiate twice and obtain 
\[
\begin{aligned}
y _{n} ^{\prime} (s) &= n \int_{s} ^{\infty} x (v) e ^ {- \sqrt {2 n} (v - s)} \dd v - n \int_{- \infty} ^{s} x (v) e ^ {- \sqrt {2 n} (s - v)} \dd v,
\\
y _{n} ^{\prime \prime} (s) &= n \left\{- x (s) - x (s) + \sqrt {2} \sqrt {n} \int_{s} ^{\infty} x (v) e ^ {- \sqrt {2 n} (s - v)} \dd v \right.
\\
&\quad \left. + \sqrt {2} \sqrt {n} \int_{- \infty} ^{s} x (v) e ^ {- \sqrt {2 n} (s - v)} \dd v \right\}
\\
&= - 2 n x (s) + 2 n y _{n} (s).
\end{aligned}
\]

Comparing this with the general formula \eqref{eq:IX.5.2}: 
\[
(A y _{n}) (s) = (A J _{n} x) (s) = n ((J _{n} - I) x) (s) = n (y _{n} (s) - x (s)),
\]
we find that $A y_{n}(s) = y_{n}^{\prime \prime}(s) / 2$. Since $R(J_{n}) = R(R(n;A)) = D(A)$, we  
have proved that 
\[
A y (s) = y ^{\prime \prime} (s) / 2 \quad \text{ for every } \quad y \in D (A).
\]

Conversely, let $y(s)$ and $y''(s)$ both belong to $C[-\infty, \infty]$. Define $x(s)$ by 
\[
y ^{\prime \prime} (s) - 2 n y (s) = - 2 n x (s).
\]

Setting $y_{n}(s) = (J_{n}x)(s)$, we obtain, as shown above, 
\[
y _{n} ^{\prime \prime} (s) - 2 n y _{n} (s) = - 2 n x (s).
\]

Hence $w(s) = y(s) - y_n(s)$ satisfies $w''(s) - 2nw(s) = 0$, and so $w(s) = C_1 e^{\sqrt{2n}s} + C_2 e^{-\sqrt{2n}s}$. This function cannot be bounded unless both $C_1$ and $C_2$ are zero. Thus $y(s) = y_n(s)$, and so $y \in D(A)$, $(Ay)(s) = y''(s)/2$. 

Therefore, the differential operator $\frac{1}{2}\frac{a^{2}}{\dd s^{2}}$ is characterized as the infinitesimal generator of the semi-group associated with the integral transform by the Gaussian kernel on the function space $C[-\infty,\infty]$. 

The proof of \eqref{eq:IX.5.3}. We start with the well-known formula 
\[
\int_{0} ^{\infty} e ^ {- x ^{2}} \dd x = \sqrt {\pi} / 2.
\]

Setting $x = \sigma - c / \sigma$, we obtain 
\[
\begin{aligned}
\frac {\sqrt {\pi}}{2} &= \int_ {\sqrt {c}} ^{\infty} e ^ {- (\sigma - c / \sigma) ^{2}} (1 + c / \sigma^{2}) \dd \sigma = e ^{2c} \int_ {\sqrt {c}} ^{\infty} e ^ {- (\sigma^{2} + c ^{2} / \sigma^{2})} (1 + c / \sigma^{2}) \dd \sigma
\\
&= e ^{2c} \left\{\int_ {\sqrt {c}} ^{\infty} e ^ {- (\sigma^{2} + c ^{2} / \sigma^{2})} \dd \sigma + \int_ {\sqrt {c}} ^{\infty} e ^ {- (\sigma^{2} + c ^{2} / \sigma^{2})} \frac {c}{\sigma^{2}} \dd \sigma \right\}.
\end{aligned}
\]

Setting $\sigma = c / t$ in the last integral, we obtain 
\[
\frac {\sqrt {\pi}}{2} = e ^{2c} \left\{\int_ {\sqrt {c}} ^{\infty} e ^ {- (\sigma^{2} + c ^{2} / \sigma^{2})} \dd \sigma - \int_ {\sqrt {c}} ^{0} e ^ {- (c ^{2} / t ^{2} + t ^{2})} \dd t \right\} = e ^{2c} \int_{0} ^{\infty} e ^ {- (t ^{2} + c ^{2} / t ^{2})} \dd t.
\]

Exercise. Show that the infinitesimal generator $A$ of the semi-group $\{T_t\}$ on $C[-\infty, \infty]$ given by 
\[
(T _{t} x) (s) = e ^{- \lambda t} \sum\limits_{k=0} ^{\infty} \frac {(\lambda t) ^{k}}{k !} x (s - k \mu) \quad (\lambda , \mu > 0),
\]
is the difference operator $A$: 
\[
(A x) (s) = \lambda \left\{x (s - \mu) - x (s) \right\}.
\]


\yossection{6}{The Exponential of a Continuous Linear Operator whose Powers are Equi-continuous}
\label{sec:IX.6}

\begin{propositionu} Let $X$ be a locally convex, sequentially complete, linear topological space. Let $B$ be a continuous linear operator $\in L(X, X)$ such that $\{B^k; k = 1, 2, \ldots\}$ is equi-continuous. Then, for each $x \in X$, the series 
\begin{equation}\label{eq:IX.6.1}
\sum\limits_{k=0} ^{\infty} (k!) ^{-1} (t B) ^{k} x (t \geqq 0)
\end{equation}
converges. 
\end{propositionu}

\begin{proof}
 For any continuous semi-norm $p$ on $X$, there exists, by the equi-continuity of $\{B^k\}$, a continuous semi-norm $q$ on $X$ such that $p(B^k x) \leqq q(x)$ for all $k \geqq 0$ and $x \in X$. Hence 
\[
\phi \left(\sum\limits_{k=n} ^{m} (t B) ^{k} x / k!\right) \leqq \sum\limits_{k=n} ^{m} t ^{k} \phi (B ^{k} x) / k! \leqq q (x) \cdot \sum\limits_{k=n} ^{m} t ^{k} / k!. \]

Therefore $\left\{\sum\limits_{k=0}^{n}(tB)^{k}x/k!\right\}$ is a Cauchy sequence in the sequentially complete space $X$. The limit of this sequence will be denoted by \eqref{eq:IX.6.1}. 
\end{proof}

\begin{corollary}\label{cor:IX.6.Corollary.1} The mapping $x \to \sum\limits_{k=0}^{\infty} (tB)^k x / k!$ defines a continuous linear operator which we shall denote by $\exp(tB)$. 
\end{corollary}

\begin{proof}
 By the equi-continuity of $\{B^k\}$, we can prove that $B_n = \sum\limits_{k=0}^{n} (tB)^k / k!$ ( $n = 0, 1, 2, \ldots$) are equi-continuous when $t$ ranges over any compact interval. In fact, we have 
\[
p (B _{n} x) \leqq \sum\limits_{k=0} ^{n} t ^{k} p (B ^{k} x) / k! \leqq q (x) \cdot \sum\limits_{k=0} ^{n} t ^{k} / k! \leqq e ^{t} \cdot q (x). \]

Hence the limit $\exp (tB)$ satisfies 
\begin{equation}\label{eq:IX.6.2}
p (\exp (t B) x) \leqq \exp (t) \cdot q (x) (t \geqq 0).
\end{equation}
\end{proof}

\begin{corollary}\label{cor:IX.6.Corollary.2} Let $B$ and $C$ be two continuous linear operators $\in L(X, X)$ such that $\{B^k\}$ and $\{C^k\}$ are equi-continuous. If, moreover, $BC = CB$, then we have 
\begin{equation}\label{eq:IX.6.3}
\exp (t B) \exp (t C) = \exp (t (B + C)).
\end{equation}
\end{corollary}

\begin{proof}
 We have 
\[
p \left((B + C) ^{k} x\right) \leqq \sum\limits_{s=0} ^{k} k C _{s} p \left(B ^{k-s} C ^{s} x\right) \leqq \sum\limits_{s=0} ^{k} k C _{s} q \left(C ^{s} x\right) \leqq 2 ^{k} \sup\limits _{0 \leqq s} q \left(C ^{s} x\right) \]

Hence $\{2^{-k}(B + C)^k\}$ is equi-continuous, and we can define $\exp(t(B + C))$. By making use of the commutativity $BC = CB$, we rearrange the series 
\[
\sum\limits_{k=0} ^{\infty} (t (B + C)) ^{k} x / k!
\]
so that we obtain $\left(\sum\limits_{k=0}^{\infty}(tB)^{k}/k!\right)\left(\sum\limits_{k=0}^{\infty}(tC)^{k}x/k!\right)$ as in the case of numerical series $\sum\limits_{k=0}^{\infty}(t(b+c))^{k}/k!$. 
\end{proof}

\begin{corollary}\label{cor:IX.6.Corollary.3} For every $x \in X$, 
\begin{equation}\label{eq:IX.6.4}
\lim\limits _{h \downarrow 0} h ^{-1} (\exp (h B) - I) x = B x,
\end{equation}
and so, by making use of the semi-group property 
\begin{equation}\label{eq:IX.6.5}
\exp ((t + h) B) = \exp (t B) \exp (h B)
\end{equation}
proved above, we obtain 
\begin{equation}\label{eq:IX.6.6}
D _{t} \exp (t B) x = \exp (t B) B x = B \exp (t B) x.
\end{equation}
\end{corollary}

\begin{proof}
 For any continuous semi-norm $\phi$ on $X$, we obtain, as above 
\[
\begin{aligned}
p (h ^{-1} (\exp (h B) - I) x - B x)& \leqq \sum\limits_{k=2} ^{\infty} h ^{k-1} p (B ^{k} x) / k! \leqq q (x) \sum\limits_{k=2} ^{\infty} h ^{k-1} / k!
\\
& \text{ which surely tends to } 0 \text{ as } h \downarrow 0. \end{aligned}
\]
\end{proof}


\yossection[7. The Representation and the Characterization of Equi-continuous Semi-groups]{7}{The Representation and the Characterization of Equi-continuous Semi-groups of Class \texorpdfstring{$(C_{0})$}{(C0)} in Terms of the Corresponding Infinitesimal Generators}
\label{sec:IX.7}

We shall prove the following fundamental 

\begin{theoremu} Let $X$ be a locally convex, sequentially complete, linear topological space. Suppose $A$ is a linear operator with domain $D(A)$ dense in $X$ and range $R(A)$ in $X$ such that the resolvent $R(n; A) = (nI - A)^{-1} \in L(X, X)$ exists for $n = 1, 2, \ldots$. Then $A$ is the infinitesimal generator of a uniquely determined equi-continuous semi-group of class $(C_0)$ iff the operators $\{(I - n^{-1}A)^{-m}\}$ are equi-continuous in $n = 1, 2, \ldots$ and $m = 0, 1, \ldots$. 
\end{theoremu}

\begin{proof}
 The ``only if'' part is already proved. We shall prove the ``if'' part. 

Setting 
\begin{equation}\label{eq:IX.7.1}
J _{n} = (I - n ^{-1} A) ^{-1},
\end{equation}
we shall prove 
\begin{equation}\label{eq:IX.7.2}
\lim\limits _{n \rightarrow \infty} J _{n} x = x \quad \text{ for every } \quad x \in X.
\end{equation}

In fact, we have, for $x \in D(A)$, $A J_n x = J_n A x = n(J_n - I)x$ and so $J_n x - x = n^{-1} J_n(Ax)$ tends to $0$ as $n \to \infty$, in virtue of the equi-boundedness of $\{J_n(Ax)\}$ in $n = 1, 2, \ldots$. Since $D(A)$ is dense in $X$ and $\{J_n\}$ is equi-continuous in $n$, it follows that $\lim\limits_{n \to \infty} J_n x = x$ for every $x \in X$. 


Put 
\begin{equation}\label{eq:IX.7.3}
T _{t} ^{(n)} = \exp (t A J _{n}) = \exp (t n (J _{n} - I)) = \exp (- n t) \exp (n t J _{n}), t \geqq 0.
\end{equation}

Since $\{J_n^k\}$ is equi-bounded in $n$ and $k$, the $\exp(t n J_n)$ can be defined, and we have, as in \eqref{eq:IX.7.2} of the preceding section, 
\[
p \left(\exp (n t J _{n}) x\right) \leqq \sum\limits_{k=0} ^{\infty} (n t) ^{k} (k!) ^{-1} p \left(J _{n} ^{k} \cdot x\right) \leqq \exp (n t) \cdot q (x).
\]

Consequently, the operators $\{T_t^{(n)}\}$ are equi-continuous in $t \geqq 0$ and $n = 1, 2, \ldots$ in such a way that 
\begin{equation}\label{eq:IX.7.4}
p \left(T _{t} ^{(n)} x\right) \leqq q (x).
\end{equation}

We next remark that $J_{n}J_{m}=J_{m}J_{n}$ for $n, m>0$. Thus $J_{n}$ is commutative with $T_{t}^{(m)}$. Hence, by $D_{t}T_{t}^{(n)}x=AJ_{n}T_{t}^{(n)}x=T_{t}^{(n)}AJ_{n}x$, proved in the preceding section, 
\begin{equation}\label{eq:IX.7.5}
\begin{aligned}
&\phi (T _{t} ^{(n)} x - T _{t} ^{(m)} x)  = \phi \left(\int_{0} ^{t} D _{s} (T _{t-s} ^{(m)} T _{s} ^{(n)} x) \dd s\right) = \phi \left(\int_{0} ^{t} T _{t-s} ^{(m)} T _{s} ^{(n)} (A J _{n} - A J _{m}) x \dd s\right).
\end{aligned}
\end{equation}

Hence, if $x \in D(A)$, there exists a continuous semi-norm $\tilde{q}$ on $X$ such that 
\[
p \left(T _{t} ^{(n)} x - T _{t} ^{(m)} x\right) \leqq \int_{0} ^{t} \tilde {q} \left(\left(A J _{n} - A J _{m}\right) x\right) \dd s = t \tilde {q} \left(\left(J _{n} A - J _{m} A\right) x\right).
\]

Therefore, by \eqref{eq:IX.7.2}, we have proved that $\lim\limits_{n,m\to \infty}p(T_t^{(n)}x - T_t^{(m)}x) = 0$ uniformly in $t$ when $t$ varies on every compact interval. Since $D(A)$ is dense in the sequentially complete space $X$, and since the operators $\{T^{(n)}\}$ are equi-continuous in $t\geqq 0$ and in $n$, we see that $\lim\limits_{n\to \infty}T_t^{(n)}x = T_tx$ exists for every $x\in X$ and $t\geqq 0$ uniformly in $t$ on every compact interval of $t$. Thus the operators $\{T_t\}$ are equi-continuous in $t\geqq 0$, and from the uniform convergence in $t$, $T_tx$ is continuous in $t\geqq 0$. 

We next prove that $T_{t}$ satisfies the semi-group property $T_{t}T_{s} = T_{t+s}$. Since $T_{t+s}^{(n)} = T_{t}^{(n)}T_{s}^{(n)}$, we have 
\[
\begin{aligned}
&\phi (T _{t+s} x - T _{t} T _{s} x)\leqq \phi (T _{t+s} x - T _{t+s} ^{(n)} x) + \phi (T _{t+s} ^{(n)} x - T _{t} ^{(n)} T _{s} ^{(n)} x) \\
&\quad + \phi (T _{t} ^{(n)} T _{s} ^{(n)} x - T _{t} ^{(n)} T _{s} x) + \phi (T _{t} ^{(n)} T _{s} x - T _{t} T _{s} x) \leqq \phi (T _{t+s} x - T _{t+s} ^{(n)} x) + q (T _{s} ^{(n)} x - T _{s} x) \\
&\quad + \phi ((T _{t} ^{(n)} - T _{t}) T _{s} x) \rightarrow 0 \text{ as } n \rightarrow \infty .
\end{aligned}
\]

Thus $p(T_{t+s}x - T_tT_s x) = 0$ for every continuous semi-norm $p$ on $X$, which proves that $T_{t+s} = T_tT_s$. 

Let $\hat{A}$ be the infinitesimal generator of this equi-continuous semigroup $\{T_t\}$ of class $(C_0)$. We have to show that $\hat{A} = A$. Let $x \in D(A)$. Then $\lim\limits_{n \to \infty} T_t^{(n)} A J_n x = T_t A x$ uniformly in $t$ on every compact interval of $t$. For, we have, by \eqref{eq:IX.7.4}, 
\[
\begin{aligned} p (T _{t} A x - T _{t} ^{(n)} A J _{n} x) & \leqq p (T _{t} A x - T _{t} ^{(n)} A x) + p (T _{t} ^{(n)} A x - T _{t} ^{(n)} A J _{n} x) \\ & \leqq p ((T _{t} - T _{t} ^{(n)}) A x) + q (A x - J _{n} A x) \end{aligned}
\]
which tends to $0$ as $n \to \infty$, since $\lim\limits_{n \to \infty} J_n A x = A x$. Hence 
\[
\begin{aligned}
&T _{t} x - x= \lim\limits _{n \rightarrow \infty} (T _{t} ^{(n)} x - x) = \lim\limits _{n \rightarrow \infty} \int_{0} ^{t} T _{s} ^{(n)} A J _{n} x \dd s = \int_{0} ^{t} (\lim\limits _{n \rightarrow \infty} T _{s} ^{(n)} A J _{n} x) \dd s = \int_{0} ^{t} T _{s} A x \dd s
\end{aligned}
\]
so that $\lim\limits_{t\downarrow 0}t^{-1}(T_t x - x) = \lim\limits_{t\downarrow 0}t^{-1}\int_0^t T_sAxds$ exists and equals $A x$. We 
have thus proved that $x \in D(A)$ implies $x \in D(\hat{A})$ and $Ax = \hat{A}x$, i.e. $\hat{A}$ is an extension of $A$. $\hat{A}$ being the infinitesimal generator of the semi-group $T_{t}$, we know that, for $n > 0$, $(nI - \hat{A})$ maps $D(\hat{A})$ onto $X$ in one-one way. But, by hypothesis, $(nI - A)$ also maps $D(A)$ onto $X$ in one-one way. Therefore the extension $\hat{A}$ of $A$ must coincide with $A$. 

Finally, the uniqueness of the semi-group $T_{t}$ is proved as follows. Let $\tilde{T}_{t}$ be an equi-continuous semi-group of class $(C_0)$ whose infinitesimal generator is precisely $A$. We construct the semi-group $T_{t}^{(n)}$. Since $A$ is commutative with $\tilde{T}_{t}$, we see that $AJ_{n}$ and $T_{s}^{(n)}$ are commutative with $\tilde{T}_{t}$. Then, for $x \in D(A)$, we obtain, as in \eqref{eq:IX.7.5}, 
\begin{equation}\label{eq:IX.7.6}
\begin{aligned}
&p (T _{t} ^{(n)} x - \tilde {T} _{t} x)  = p \left(\int_{0} ^{t} D _{s} (\tilde {T} _{t-s} T _{s} ^{(n)} x) \dd s\right) = p \left(\int_{0} ^{t} - \tilde {T} _{t-s} T _{s} ^{(n)} (A - A J _{n}) x \dd s\right).
\end{aligned}
\end{equation}

Thus, by virtue of $\lim\limits_{n\to \infty}AJ_nx = Ax$ for all $x\in D(A)$, we prove $\lim\limits_{n\to \infty}T_t^{(n)}x = \tilde{T}_t x$ for all $x\in X$ similarly to the above proof of the existence of $\lim\limits_{n\to \infty}T_t^{(n)}x,x\in X$. 

Therefore, $T_{t}x = \tilde{T}_{t}x$ for all $x \in X$, that is, $T_{t} = \tilde{T}_{t}$.
\end{proof}

\begin{remarku} The above proof shows that, if $A$ is the infinitesimal generator of an equi-continuous semi-group $T_{t}$ of class $(C_{0})$, then 
\begin{equation}\label{eq:IX.7.7}
T _{t} x = \lim\limits _{n \rightarrow \infty} \exp (t A (I - n ^{-1} A) ^{-1}) x, x \in X,
\end{equation}
\end{remarku}

and the convergence in \eqref{eq:IX.7.7} is uniform in $t$ on every compact interval of $t$. This is the representation theorem for semi-groups. 

\begin{corollary}\label{cor:IX.7.Corollary.1} If $X$ is a $B$-space, then the condition of the Theorem reads: $D(A)^{a} = X$ and the resolvent $(I - n^{-1}A)^{-1}$ exists such that 
\begin{equation}\label{eq:IX.7.8}
\left\| (I - n ^{-1} A) ^{-m} \right\| \leqq C \quad (n = 1, 2, \dots ; m = 1, 2, \dots)
\end{equation}
with a positive constant $C$ which is independent of $n$ and $m$. In particular, for the case of a contraction semi-group, the condition reads: $D(A)^{a} = X$ and the resolvent $(I - n^{-1}A)^{-1}$ exists such that 
\begin{equation}\label{eq:IX.7.9}
\left\| (I - n ^{-1} A) ^{-1} \right\| \leqq 1 (n = 1, 2, \dots).
\end{equation}
\end{corollary}

\begin{remarku} The above result \eqref{eq:IX.7.9} was obtained independently by E. \textsc{Hille} \cite{hille98} and by K. \textsc{Yosida} \cite{yosida348}. The result was extended by W. \textsc{Feller} \cite{erdelyi56}, R. S. \textsc{Phillips} \cite{phillips253} and I. \textsc{Miyadera} \cite{miyadera204} and the extension is given in the form \eqref{eq:IX.7.8}. It is to be noted that in condition \eqref{eq:IX.7.8} and \eqref{eq:IX.7.9}, we may replace $(n=1,2,\ldots)$ by (for all sufficiently large $n$). The extension of the semi-group theory to locally convex linear topological spaces, as given in the present book, is suggested by L. \textsc{Schwartz} \cite{schwartz296}. 
\end{remarku}

\begin{corollary}\label{cor:IX.7.Corollary.2} Let $X$ be a $B$-space, and $\{T_{t}; t \geqq 0\}$ be a family of bounded linear operators $L(X, X)$ such that 
\begin{align}
T _{t} T _{s}& = T _{t+s} (t, s \geqq 0), T _{0} = I, \label{eq:IX.7.10} \\
\slim _{t \downarrow 0} T _{t} x& = x \text{ for all } x \in X, \label{eq:IX.7.11} \\
\left\| T _{t} \right\|& \leqq M e ^{\beta t} \text{ for all } t \geqq 0, \text{ where } M > 0 \text{ and } \beta \geqq 0 \text{ are } \label{eq:IX.7.12}
\end{align}
independent of $t$. 
\end{corollary}


Then $(A - \beta I)$ is the infinitesimal generator of the equi-continuous semi-group $S_{t} = e^{-\beta t}T_{t}$ of class $(C_0)$, where $A$ is the operator defined through $Ax = s\lim\limits_{t\downarrow 0}t^{-1}(T_t - I)x$. Thus, by Corollary~\ref{cor:IX.7.Corollary.1}, we see that a closed linear operator $A$ with $D(A)^{a} = X$ and $R(A)\subseteq X$ is the infinitesimal generator of a semi-group $T_{t}$ satisfying \eqref{eq:IX.7.10}, \eqref{eq:IX.7.11} and \eqref{eq:IX.7.12} iff the resolvent $(I - n^{-1}(A - \beta I))^{-1}$ exists such that 
$\left\|(I-n^{-1}(A-\beta I))^{-m}\right\|\leqq M$ (for $m=1,2,\ldots$ and all large $n$). This condition may be rewritten as 

\begin{equation}\tag{13'}\label{eq:IX.7.13'}
\left\| (I - n ^{-1} A) ^{-m} \right\| \leqq M (1 - n ^{-1} \beta) ^{-m} (\text{ for } m = 1, 2, \dots \text{ and all large } n).
\end{equation}

In particular, for those semi-groups $T_{t}$ satisfying \eqref{eq:IX.7.10}, \eqref{eq:IX.7.11} and 
\begin{equation}\label{eq:IX.7.14}
\left\| T _{t} \right\| \leqq e ^{\beta t} \quad \text{ for all } \quad t \geqq 0,
\end{equation}
condition \eqref{eq:IX.7.13'} may be replaced by 
\begin{equation}\tag{13''}\label{eq:IX.7.13}
\left\| (I - n ^{-1} A) ^{-1} \right\| \leqq (1 - n ^{-1} \beta) ^{-1} \quad (\text{ for all large } n)
\end{equation}

An application of the representation theorem to the proof of Weierstrass' polynomial approximation theorem. Consider the semi-group $T_{t}$ defined by $(T_{t}x)(s) = x(t + s)$ on $C[0,\infty]$. The representation theorem gives 
\[
(T _{t} x) (s) = x (t + s) = \slim\limits_{n \rightarrow \infty} \exp (t A J _{n} x) (s) = \slim\limits_{n \rightarrow \infty} \sum\limits_{m=0} ^{\infty} \frac {t ^{m}}{m !} (A J _{n}) ^{m} x (s),
\]
and the above \slim is uniform in $t$ on any compact interval of $t$. From this result we can derive Weierstrass' polynomial approximation theorem. Let $z(s)$ be a continuous function on the closed interval [$0$, $1$]. Let $x(s) \in C[0, \infty]$ be such that $x(s) = z(s)$ for $s \in [0, 1]$. Put $s = 0$ in the above representation of $x(t + s)$. Then we obtain 
\[
(T _{t} x) (0) = x (t) = \slim _{n \to \infty} \sum\limits_{m=0} ^{\infty} t ^{m} (A J _{n}) ^{m} x (0) / m! \quad \text{ in } \quad C [ 0, 1 ],
\]
uniformly in $t$ on $[0, 1]$. Hence $z(t)$ is a uniform limit on $[0, 1]$ of a sequence of polynomials in $t$. 


\yossection{8}{Contraction Semi-groups and Dissipative Operators}
\label{sec:IX.8}

G. \textsc{Lumer} and R. S. \textsc{Phillips} have discussed contraction semi-groups by virtue of the notion of semi-scalar product. The infinitesimal generator of such a semi-group is dissipative in their terminology. 

\begin{propositionu} (\textsc{Lumer}). To each pair $\{x, y\}$ of a complex (or real) normed linear space $X$, we can associate a complex (or real) number $[x, y]$ such that 
\begin{equation}\label{eq:IX.8.1}
\begin{gathered}
{}[ x + y, z ] = [ x, z ] + [ y, z ], \quad [ \lambda x, y ] = \lambda [ x, y ], \quad [ x, x ] = \| x \| ^{2},
\\
| [ x, y ] | \leqq \| x \| \cdot \| y \|.
\end{gathered}
\end{equation}
$[x, y]$ is called a semi-scalar product of the vectors $x$ and $y$. 
\end{propositionu}

\begin{proof}
 According to Corollary~\hyperref[cor:IV.6.Corollary.2]{$2$} of Theorem~\ref{thm:IV.6.Theorem.1} in Chapter~\ref{chap:IV}, \ref{sec:IV.6}, there exists, for each $x_0 \in X$, at least one (and let us choose exactly one) bounded linear functional $f_{x_0} \in X'$ such that $\|f_{x_0}\| = \|x_0\|$ and $\langle x_0, f_{x_0} \rangle = \|x_0\|^2$. Then clearly 
\begin{equation}\label{eq:IX.8.2}
[ x, y ] = \langle x, f _{y} \rangle
\end{equation}
defines a semi-scalar product. 
\end{proof}

\begin{definitionu} Let a complex (or real) $B$-space $X$ be endowed with a semi-scalar product $[x, y]$. A linear operator $A$ with domain $D(A)$ and range $R(A)$ both in $X$ is called dissipative (with respect to $[x, y]$) if 
\begin{equation}\label{eq:IX.8.3}
\operatorname{Re} [ A x, x ] \leqq 0 \quad \text{ whenever } \quad x \in D (A).
\end{equation}
\end{definitionu}

\begin{examplen} Let $X$ be a Hilbert space. Then a symmetric operator $A$ such that $(Ax, x) \leqq 0$ is surely dissipative with respect to the semi-scalar product $[x, y] = (x, y)$, where $(x, y)$ is the ordinary scalar product.

\end{examplen}
\begin{theoremn}[\textsc{Phillips} and \textsc{Lumer}] Let $A$ be a linear operator with domain $D(A)$ and range $R(A)$ both in a complex (or real) $B$-space $X$ such that $D(A)^{a} = X$. Then $A$ generates a contraction semi-group of class $(C_{0})$ in $X$ iff $A$ is dissipative (with respect to any semi-scalar product $[x, y]$) and $R(I - A) = X$.

\end{theoremn}
\begin{proof}
 The ``if'' part. Let $A$ be dissipative and $\lambda > 0$. Then the inverse $(\lambda I - A)^{-1}$ exists and $\|(\lambda I - A)^{-1}y\| \leqq \lambda^{-1}\|y\|$ when $y \in D((\lambda I - A)^{-1})$. 

For, if $y = \lambda x - Ax$, then
\begin{equation}\label{eq:IX.8.4}
\lambda\|x\|^2=\lambda[x,x]\leqq\operatorname{Re}(\lambda[x,x]-[Ax,x])=\operatorname{Re}[y,x]\leqq\|y\|\cdot\|x\|,
\end{equation}

since $A$ is dissipative. By hypothesis, $R(I - A) = X$ so that $\lambda = 1$ is in the resolvent set $\varrho (A)$ of $A$, and we have $\|R(1;A)\|\leqq 1$ by \eqref{eq:IX.8.4}. If $|\lambda -1| < 1$, then the resolvent $R(\lambda ;A)$ exists and is given by 
\[
\begin{aligned}
&R (\lambda ; A)  = R (1; A) \left(I + (\lambda - 1) R (1; A)\right) ^{-1} = R (1; A) \cdot \sum\limits_{n=0} ^{\infty} \big ((1 - \lambda) R (1; A) \big) ^{n},
\end{aligned}
\]

(see Theorem~\ref{thm:VIII.2.Theorem.1} in Chapter~\ref{chap:VIII}, \ref{sec:VIII.2}). Moreover, \eqref{eq:IX.8.4} implies that $\|R(\lambda; A)\| \leqq \lambda^{-1}$ for $\lambda > 0$ with $|\lambda - 1| < 1$. Hence, again by 
\[
R (\mu ; A) = R (\lambda ; A) (I + (\mu - \lambda) R (\lambda ; A)) ^{-1},
\]
which is valid for $\mu > 0$ with $|\mu - \lambda| \|R(\lambda; A)\| < 1$, we prove the existence of $R(\mu; A)$ and $\|R(\mu; A)\| \leqq \mu^{-1}$. Repeating the process, we see that $R(\lambda; A)$ exists for all $\lambda > 0$ and satisfies the estimate $\|R(\lambda; A)\| \leqq \lambda^{-1}$. As $D(A)$ is dense by hypothesis, it now follows, from Corollary~\hyperref[cor:IX.7.Corollary.1]{$1$} of the preceding section, that $A$ generates a contraction semi-group of class $(C_0)$. 


The ``only if'' part. Suppose $\{T_{t}; t \geqq 0\}$ is a contraction semi-group of class $(C_{0})$. Then 
\begin{gather}\operatorname{Re} \left[ T _{t} x - x, x \right] = \operatorname{Re} \left[ T _{t} x, x \right] - \| x \| ^{2} \leqq \| T _{t} x \| \cdot \| x \| - \| x \| ^{2} \leqq 0. x \in D (A) T, \notag \\\operatorname{Re} [ A x, x ] = \lim\limits _{t \downarrow 0} \operatorname{Re} \left\{t ^{-1} \left[ T _{t} x - x, x \right] \right\} \leqq 0 \notag
\end{gather}

Moreover, we know that $R(I - A) = D(R(1; A)) = X$, since $A$ is the infinitesimal generator of a contraction semi-group of class $(C_0)$.
\end{proof}

\begin{corollaryu} If $A$ is a densely defined closed linear operator such that $D(A)$ and $R(A)$ are both in a $B$-space $X$ and if $A$ and its dual operator $A'$ are both dissipative, then $A$ generates a contraction semi-group of class $(C_{0})$. 
\end{corollaryu}

\begin{proof}
 It suffices to show that $R(I - A) = X$. But, since $(I - A)^{-1}$ is closed with $A$ and continuous, $R(I - A)$ is a closed linear subspace of $X$. Thus $R(I - A) \neq X$ implies the existence of a non-trivial $x' \in X'$ such that $\langle (x - Ax), x' \rangle = 0$ for all $x \in D(A)$. 

Hence $x' - A'x' = 0$, contrary to the dissipativity of $A'$ and $x' \neq 0$. 
\end{proof}

\begin{remarku} For further details concerning dissipative operators, see G. \textsc{Lumer}-R. S. \textsc{Phillips} \cite{lumer183}. See also T. \textsc{Kato} \cite{kato139}. 
\end{remarku}


\yossection{9}{Equi-continuous Groups of Class \texorpdfstring{$(C_0)$}{(C0)}. Stone's Theorem}
\label{sec:IX.9}

\begin{definitionu} An equi-continuous semi-group $\{T_t\}$ of class $(C_0)$ is called an equi-continuous group of class $(C_0)$, if there exists an equi-continuous semi-group $\{\hat{T}_t\}$ of class $(C_0)$ with the condition: 

If we define $S_{t}$ by $S_{t} = T_{t}$ for $t \geqq 0$ and $S_{-t} = \hat{T}_{t}$ for $t \geqq 0$, then the family of operators $S_{t}, -\infty < t < \infty$, has the group property 
\end{definitionu}
\begin{equation}\label{eq:IX.9.1}
S _{t} S _{s} = S _{t+s} (- \infty < t, s < \infty), S _{0} = I.
\end{equation}

\begin{theoremu} Let $X$ be a locally convex, sequentially complete linear topological space. Suppose $A$ is a linear operator with domain $D(A)$ dense in $X$ and range $R(A)$ in $X$. Then $A$ is the infinitesimal generator of an equi-continuous group of class $(C_{0})$ of operators $S_{t} \in L(X, X)$ iff the operators $(I - n^{-1}A)^{-m}$ are everywhere defined and equi-continuous in $n = \pm 1, \pm 2, \ldots$ and in $m = 1, 2, \ldots$. 
\end{theoremu}

\begin{proof}
 The ``only if'' part. Let $T_{t}=S_{t}$ for $t\geqq0$ and $\hat{T}_{t}=S_{-t}$ for $t\geqq0$. Let $A$ and $\hat{A}$ be the infinitesimal generator of $\hat{T}_{t}$ and $T_{t}$, respectively. We have to show that $\hat{A}=-A$. If $x\in D(\hat{A})$, then, by putting $x_{h}=h^{-1}(\hat{T}_{h}-I)x$ and making use of the equi-continuity of $T_{h}$, 
\[
\begin{aligned}
&p (T _{h} x _{h} - \hat {A} x)  \leqq p (T _{h} x _{h} - T _{h} \hat {A} x) + p ((T _{h} - I) A x) \leqq q (x _{h} - \hat {A} x) + p ((T _{h} - I) \hat {A} x)
\end{aligned}
\]
where $p$ and $q$ are continuous semi-norms on $X$ such that, for a given $p$, we can choose a $q$ satisfying the above inequality for all $h \geqq 0$ and all $x \in D(\hat{A})$ simultaneously. Thus $\lim\limits_{h \downarrow 0} T_h x_h = \hat{A} x$, and so $x \in D(\hat{A})$ implies $\hat{A} x = \lim\limits_{h \downarrow 0} T_h (h^{-1}(\hat{T}_h - I)) x = \lim\limits_{h \downarrow 0} h^{-1}(I - T_h) x = -Ax$. Hence $-A$ is an extension of $\hat{A}$. In the same way, we can prove that $\hat{A}$ is an extension of $-A$. Therefore $\hat{A} = -A$. 

The ``if'' part. Define, for $t \geqq 0$, 
\begin{align}
T _{t} x& = \lim\limits _{n \rightarrow \infty} T _{t} ^{(n)} x = \lim\limits _{n \rightarrow \infty} \exp (t A (I - n ^{-1} A) ^{-1}) x, \notag \\
\hat {T} _{t} x& = \lim\limits _{n \rightarrow \infty} \hat {T} _{t} ^{(n)} x = \lim\limits _{n \rightarrow \infty} \exp (t \hat {A} (I - n ^{-1} \hat {A}) ^{-1}) x, \text{ where } \hat {A} = - A. \notag
\end{align}

Then $T_{t}$ and $\hat{T}_{t}$ are both equi-continuous semi-groups of class $(C_{0})$. We have 
\[
\begin{aligned} p (T _{t} \hat {T} _{t} x - T _{t} ^{(n)} \hat {T} _{t} ^{(n)} x) & \leqq p (T _{t} \hat {T} _{t} x - T _{t} ^{(n)} \hat {T} _{t} x) + p (T _{t} ^{(n)} \hat {T} _{t} x - T _{t} ^{(n)} \hat {T} _{t} ^{(n)} x) \\ & \leqq p ((T _{t} - T _{t} ^{(n)}) \hat {T} _{t} x) + q (\hat {T} _{t} x - \hat {T} _{t} ^{(n)} x) \end{aligned}
\]
by the equi-continuity of $\{T_t^{(n)}\}$ in $n$ and $t \geqq 0$. Thus we have $\lim\limits_{n \to \infty} T_t^{(n)} \hat{T}_t^{(n)} x = T_t \hat{T}_t x$. We have incidentally proved the equi-continuity of $T_t \hat{T}_t$ by the equi-continuity of $T_t^{(n)} \hat{T}_t^{(n)}$. On the other hand, we have $T_t^{(n)} \hat{T}_s^{(m)} = \hat{T}_s^{(m)} T_t^{(n)}$ by the commutativity of $(I - n^{-1}A)^{-1}$ and $(I - m^{-1}A)^{-1}$. Thus $(T_t^{(n)} \hat{T}_t^{(n)}) (T_s^{(n)} \hat{T}_s^{(n)}) = T_{t+s}^{(n)} \hat{T}_{t+s}^{(n)}$, and so $(T_{t}\hat{T}_{t})$ with $t \geqq 0$ enjoys the semi-group property $(T_{t}\hat{T}_{t})(T_{s}\hat{T}_{s}) = T_{t+s}\hat{T}_{t+s}$. Hence $\{T_{t}\hat{T}_{t}\}$ is an equi-continuous semi-group of class $(C_0)$. If $x \in D(\hat{A}) = D(A)$, then 
\[
\begin{aligned}
&\lim\limits _{h \downarrow 0} h ^{-1} (T _{h} \hat {T} _{h} - I) x  = \lim\limits _{h \downarrow 0} T _{h} h ^{-1} (\hat {T} _{h} - I) x + \lim\limits _{h \downarrow 0} h ^{-1} (T _{h} - I) x = \hat {A} x + A x = 0,
\end{aligned}
\]
so that the infinitesimal generator $A_{1}$ of $\{T_t\hat{T}_t\}$ is $0$ at every $x \in D(\hat{A})$. Since $(I - A_1)$ is the inverse of a continuous linear operator $(I - A_1)^{-1} \in L(X, X)$, we see that $A_{1}$ must be closed. Hence $A_{1}$ must be = $0$, in virtue of the fact that $A_{1}$ vanishes on a dense subset $D(\hat{A}) = D(A)$ of $X$. Therefore $(T_t\hat{T}_t)x = \lim\limits_{n\to \infty}\exp (t\cdot 0\cdot (I - n^{-1}\cdot 0)^{-1})x = x$, that is, $T_{t}\hat{T}_{t} = I$. We have thus proved the group property of $S_{t}, -\infty < t < \infty$, where $S_{t} = T_{t}$ and $S_{-t} = \hat{T}_t$ for $t\geqq 0$. 
\end{proof}

\begin{corollary}\label{cor:IX.9.Corollary.1} For the case of a $B$-space $X$, the condition of the Theorem reads: $D(A)^{a} = X$ and the resolvent $(I - n^{-1}A)^{-1}$ exists such that 
\begin{equation}\label{eq:IX.9.2}
\left\|(I - n ^{-1} A)^{-m}\right\| \leqq M \text{ (for } m = 1, 2, \ldots \text{ and all large } |n|, n \geqq 0).
\end{equation}
For the group $S_{t}$ satisfying $\left\|S_{t}\right\| \leqq M e^{\beta |t|} (\beta \geqq 0)$ for all $t, -\infty < t < \infty$, the condition reads: $D(A)^{a} = X$ and the resolvent $(I - n^{-1}A)^{-1}$ exists such that 
\begin{equation}\label{eq:IX.9.3}
\left\| (I - n ^{-1} A) ^{-m} \right\| \leqq M (1 - | n ^{-1} | \beta) ^{-m} \text{ (for } m = 1, 2, \ldots \text{ and all large } |n|, n \leqq 0).
\end{equation}
\end{corollary}


For the particular case $\| S_t \| \leqq e^{\beta |t|}$ for all $t, -\infty < t < \infty$, the condition reads: $D(A)^a = X$ and the resolvent $(I - n^{-1}A)^{-1}$ exists such that 
\begin{equation}\label{eq:IX.9.4}
\left\| (I - n ^{-1} A) ^{-1} \right\| \leqq (1 - | n ^{-1} | \beta) ^{-1} (\text{ for all large } | n |, n \geqq 0).
\end{equation}

\begin{proof}
 As the proof of Corollary~\ref{cor:IX.9.Corollary.1} and Corollary~\ref{cor:IX.7.Corollary.2} in Chapter~\ref{chap:IX}, \ref{sec:IX.7}. \end{proof}

\begin{corollary}[Stone's theorem]\label{cor:IX.9.Corollary.2} Let $U_{t}, -\infty < t < \infty$, be a group of unitary operators of class $(C_{0})$ in a Hilbert space $X$. Then the infinitesimal generator $A$ of $U_{t}$ is $\sqrt{-1}$ times a self-adjoint operator $H$. 
\end{corollary}

\begin{proof}
 We have $(U_{t}x,y) = (x,U_{t}^{-1}y) = (x,U_{-t}y)$ and so, by differentiation, 
\[
(A x, y) = (x, - A y) \text{ whenever } x \text{ and } y \in D (A) \]

Thus $-iA = H$ is symmetric. $A$ being the infinitesimal generator of $U_{t}$, $(I - n^{-1}A)^{-1} = (I - n^{-1}iH)^{-1}$ must be a bounded linear operator such that $\left\|(I - n^{-1}iH)^{-1}\right\| \leqq 1$ for $n = \pm 1, \pm 2, \ldots$. Hence, taking the case $n = \pm 1$, we see that the Cayley transform of $H$ is unitary. This proves that $H$ is self-adjoint. 
\end{proof}

\begin{remarku} If $A$ is of the form $A = \sqrt{-1} \cdot H$, where $H$ is a self-adjoint operator in a Hilbert space $X$, then condition \eqref{eq:IX.9.4} of Corollary~\ref{cor:IX.9.Corollary.1} is surely satisfied, as may be proved by the theory of Cayley's transform. Therefore, $A$ is the infinitesimal generator of a group of contraction operators $U_t$ in $X$. It is easy to see that such $U_t$ is unitary. For, a contraction operator $U_t$ on a Hilbert space $X$ into $X$ must be unitary, if $U_t^{-1} = U_{-t}$ is also a contraction operator on $X$ into $X$. 
\end{remarku}


\yossection{10}{Holomorphic Semi-groups}
\label{sec:IX.10}

We shall introduce an important class of semi-groups, namely, semi-groups $T_{t}$ which can, as functions of the parameter $t$, be continued holomorphically into a sector of the complex plane containing the positive $t$-axis. We first prove 

\begin{lemmau} Let $X$ be a locally convex, sequentially complete, linear topological space. Let $\{T_{t}; t \geqq 0\} \subseteq L(X, X)$ be an equi-continuous semigroup of class $(C_{0})$. Suppose that, for all $t > 0$, $T_{t}X \subseteq D(A)$, the domain of the infinitesimal generator $A$ of $T_{t}$. Then, for any $x \in X$, $T_{t}x$ is infinitely differentiable in $t > 0$ and we have 
\begin{equation}\label{eq:IX.10.1}
T _{t} ^{(n)} x = (T _{t/n} ^{\prime}) ^{n} x \quad \text{ for all } \quad t > 0,
\end{equation}
where $T_{t}^{\prime} = D_{t}T_{t}, T_{t}^{\prime \prime} = D_{t}T_{t}^{\prime}, \ldots, T_{t}^{(n)} = D_{t}T_{t}^{(n-1)}$. 
\end{lemmau}

\begin{proof}
 If $t > t_0 > 0$, then $T_t'x = AT_tx = T_{t-t_0}AT_{t_0}x$ by the commutativity of $A$ and $T_s, s \geqq 0$. Thus $T_t'X \subseteq T_{t-t_0}X \subseteq D(A)$ when $t > 0$, and so $T_t''x$ exists for all $t > 0$ and $x \in X$. Since $A$ is a closed linear operator, we have 
\[ \begin{aligned}
&T _{t} ^{\prime \prime} x= D _{t} (A T _{t}) x = A \cdot \lim\limits _{n \rightarrow \infty} \cdot n (T _{t+1/n} - T _{t}) x = A (A T _{t}) x = A T _{t/2} A T _{t/2} x = (T _{t/2} ^{\prime}) ^{2} x.
\end{aligned}
\]

Repeating the same argument, we obtain \eqref{eq:IX.10.1}. 
\end{proof}

Let $X$ be a locally convex, sequentially complete, complex linear topological space. Let $\{T_{t}; t \geqq 0\} \subseteq L(X, X)$ be an equi-continuous semi-group of class $(C_{0})$. For such a semi-group, we consider the following three conditions: 

(I) For all $t > 0$, $T_{t}x \in D(A)$, and there exists a positive constant $C$ such that the family of operators $\{(CtT_t)^n\}$ is equi-continuous in $n \geqq 0$ and $t, 0 < t \leqq 1$. 

(II) $T_{t}$ admits a weakly holomorphic extension $T_{\lambda}$ given locally by 
\begin{equation}\label{eq:IX.10.2}
T _{\lambda} x = \sum\limits_{n=0} ^{\infty} (\lambda - t) ^{n} T _{t} ^{(n)} x / n! \quad \text{ for } \quad | \arg \lambda | < \tan^{-1} (C e ^{-1}),
\end{equation}
the family of operators $\{e^{-\lambda}T_{\lambda}\}$ is equi-continuous in $\lambda$ for 
\begin{equation}\label{eq:IX.10.3}
\left| \arg \lambda \right| < \tan^{-1} \left(2 ^{-k} C e ^{-1}\right) \text{ with some positive constant } k.
\end{equation}

(III) Let $A$ be the infinitesimal generator of $T_{t}$. Then there exists a positive constant $C_{1}$ such that the family of operators $\{(C_{1}\lambda R(\lambda ;A))^{n}\}$ is equi-continuous in $n\geqq 0$ and in $\lambda$ with $\operatorname{Re}(\lambda)\geqq 1 + \varepsilon ,\varepsilon >0$. 

We can prove 

\begin{theoremu} The three conditions (I), (II) and (III) are mutually equivalent. 
\end{theoremu}

\begin{proof}
 The implication (I) $\rightarrow$ (II). Let $p$ be any continuous semi-norm on $X$. Then, by hypothesis, there exists a continuous semi-norm $q$ on $X$ such that $p((tT_t')^n x) \leqq C^{-n} q(x)$, $0 < C$ for $1 \geqq t > 0$, $n \geqq 0$ and $x \in X$. Hence, by \eqref{eq:IX.10.1}, we obtain, for any $t > 0$, 
\[ \begin{aligned} \phi ((\lambda - t) ^{n} T _{t} ^{(n)} x / n!) & \leqq \frac {| \lambda - t | ^{n}}{t ^{n}} \frac {n ^{n}}{n !} \frac {1}{C ^{n}} \phi \left(\left(\frac {t}{n} C T _{t/n} ^{\prime}\right) ^{n} x\right) \\ & \leqq \left(\frac {| \lambda - t |}{t} C ^{-1} e\right) ^{n} \cdot q (x), \text{ whenever } 0 \leqq t / n \leqq 1. \end{aligned}
\]

Thus the right side of \eqref{eq:IX.10.2} is convergent for $|\arg \lambda| < \tan^{-1}(Ce^{-1})$, and so, by the sequential completeness of the space $X$, $T_{\lambda}x$ is well defined and weakly holomorphic in $\lambda$ for $|\arg \lambda| < \tan^{-1}(Ce^{-1})$. That is, for any $x \in X$ and $f \in X'$, the numerical function $f(T_t x)$ of $t$, $t > 0$, admits a holomorphic extension $f(T_{\lambda}x)$ for $|\arg \lambda| < \tan^{-1}(Ce^{-1})$; consequently, by the Hahn-Banach theorem, we see that $T_{\lambda}x$ is an extension of $T_t x$ for $|\arg \lambda| < \tan^{-1}(Ce^{-1})$. Next put $S_t = e^{-t} T_t$. Then $S_t' = e^{-t} T_t' - e^{-t} T_t$ and so, by $0 \leqq te^{-t} \leqq 1$ ( $0 \leqq t$) and (I), we easily see that the family of operators $\{(2^{-k} Ct S_t')^n\}$ is equi-continuous in $t > 0$ and $n \geqq 0$, in virtue of the equi-continuity of $\{T_t\}$ in $t > 0$. The equi-continuous semi-group $S_t$ of class $(C_0)$ satisfies the condition that $t > 0$ implies $S_t X \subseteq D(A - I) = D(A)$, where $(A - I)$ is the infinitesimal generator of $S_t$. Therefore, by the same reasoning applied above to $T_t$, we can prove that the weakly holomorphic extension $e^{-\lambda} T_{\lambda}$ of $S_t = e^{-t} T_t$ satisfies the estimate \eqref{eq:IX.10.3}. 
\end{proof}

By the way, we can prove the following 

\begin{corollaryn}[due to E. \textsc{Hille}] If, in particular, $X$ is a complex $B$-space and $\varlimsup\limits_{t\downarrow 0}\|tT_t'| < e^{-1}$, then $X = D(A)$.

\end{corollaryn}
\begin{proof}
 For a fixed $t > 0$, we have $\varlimsup\limits_{n \to \infty} \|(t/n) T_{t/n}'\| < e^{-1}$, and so the series 
\[
\sum\limits_{n=0} ^{\infty} (\lambda - t) ^{n} T _{t} ^{(n)} x / n! = \sum\limits_{n=0} ^{\infty} \frac {(\lambda - t) ^{n}}{t ^{n}} \frac {n ^{n}}{n !} \left(\frac {t}{n} T _{t/n} ^{\prime}\right) ^{n} x
\]
strongly converges in some sector 
\[
\{\lambda ; | \lambda - t | / t < 1 + \delta \quad \text{ with some } \quad \delta > 0 \}
\]
of the complex $\lambda$-plane. This sector surely contains $\lambda = 0$ in its interior. 


The implication (II) $\rightarrow$ (III). We have, by \eqref{eq:IX.4.10} in Chapter~\ref{chap:IX}, \ref{sec:IX.4}, 
\begin{equation}\label{eq:IX.10.4}
(\lambda R (\lambda ; A)) ^{n+1} x = \frac {\lambda^{n+1}}{n !} \int_{0} ^{\infty} e ^{- \lambda t} t ^{n} T _{t} x \dd t \quad \text{ for } \quad \operatorname{Re} (\lambda) > 0, x \in X.
\end{equation}

Hence, putting $S_{t} = e^{-t}T_{t}$, we obtain 
\[
\begin{aligned}
&((\sigma + 1 + i \tau) R (\sigma + 1 + i \tau ; A)) ^{n+1} x = \frac {(\sigma + 1 + i \tau) ^{n+1}}{n !} \int_{0} ^{\infty} e ^{-(\sigma + i \tau) t} t ^{n} S _{t} x \dd t, \sigma > 0.
\end{aligned}
\]

Let $\tau < 0$. Since the integrand is weakly holomorphic, we can deform, by the estimate \eqref{eq:IX.10.3} and Cauchy's integral theorem, the path of integration: $0 \leqq t < \infty$ to the ray: $re^{i\theta} (0 \leqq r < \infty)$ contained in the sector $0 < \arg \lambda < \tan^{-1}(2^{-k} Ce^{-1})$ of the complex $\lambda$-plane. We thus obtain 
\[
\begin{aligned}
&\quad ((\sigma + 1 + i \tau) R (\sigma + 1 + i \tau ; A)) ^{n+1} x = \frac {(\sigma + 1 + i \tau) ^{n+1}}{n !} \int_{0} ^{\infty} e ^ {- (\sigma + i \tau) r e ^{i \theta}} r ^{n} e ^{i n \theta} S _ {r e ^{i \theta}} x e ^{i \theta} \dd r,
\end{aligned}
\]
and so, by \eqref{eq:IX.10.3}, 
\[
\begin{aligned} & p \left(\left((\sigma + 1 + i \tau) R (\sigma + 1 + i \tau ; A)\right) ^{n+1} x\right) \\ & \leqq \sup\limits _{0 < r < \infty} p (S _ {r e ^{i \theta}} x) \frac {| (\sigma + 1 + i \tau) | ^{n+1}}{n !} \int_{0} ^{\infty} e ^{(- \sigma \cos \theta + \tau \sin \theta) r} r ^{n} \dd r \end{aligned}
\]
$\leqq q'(x) \frac{|\sigma + 1 + i\tau|^{n+1}}{|\tau \sin \theta - \sigma \cos \theta|^{n+1}}$, where $q'$ is a continuous semi-norm on $X$. 

A similar estimate is also obtained for the case $\tau > 0$. Hence, combined with \eqref{eq:IX.10.7} in Section~\hyperref[sec:IX.4]{$4$}, we have proved (III). 

The implication (III) $\rightarrow$ (I). For any continuous semi-norm $p$ on $X$, there exists a continuous semi-norm $q$ on $X$ such that $p((C_1\lambda R(\lambda;A))^n x)\leqq q(x)$ whenever $\operatorname{Re}(\lambda)\geqq 1 + \varepsilon, \varepsilon >0$, and $n\geqq 0$. Hence, if $\operatorname{Re}(\lambda_0)\geqq 1 + \varepsilon$, we have 
\[
\phi \left((\lambda - \lambda_{0}) ^{n} R (\lambda_{0}; A) ^{n} x\right) \leqq \frac {| \lambda - \lambda_{0} | ^{n}}{(C _{1} | \lambda_{0} |) ^{n}} q (x) \quad (n = 0, 1, 2, \dots).
\]

Thus, if $|\lambda - \lambda_0| / C_1 |\lambda_0| < 1$, the resolvent $R(\lambda, A)$ exists and is given by 
\begin{align}
R (\lambda ; A) x& = \sum\limits_{n=0} ^{\infty} (\lambda_{0} - \lambda) ^{n} R (\lambda_{0}; A) ^{n+1} x \quad \text{ such that } \notag \\
p (R (\lambda ; A) x)& \leqq (1 - C _{1} ^{-1} | \lambda_{0} | ^{-1} | \lambda - \lambda_{0} |) ^{-1} q (R (\lambda_{0}; A) x). \notag
\end{align}

Therefore, by (III), there exists an angle $\theta_0$ with $\pi / 2 < \theta_0 < \pi$ such that $R(\lambda; A)$ exists and satisfies the estimate 
\begin{equation}\label{eq:IX.10.5}
p(R(\lambda; A)x) \leqq \frac{1}{|\lambda|} q'(x) \text{ with a continuous semi-norm } q' \text{ on } X
\end{equation}
in the sectors $\pi / 2 \leqq \arg \lambda \leqq \theta_0$ and $-\theta_0 \leqq \arg \lambda \leqq -\pi / 2$ and also for $\operatorname{Re}(\lambda) \geqq 0$, when $|\lambda|$ is sufficiently large. 

Hence the integral 
\begin{equation}\label{eq:IX.10.6}
\hat {T} _{t} x = (2 \pi i) ^{-1} \int_ {C _{3}} e ^{\lambda t} R (\lambda ; A) x \dd \lambda \quad (t > 0, x \in X)
\end{equation}
converges if we take the path of integration $C_2 = \lambda(\sigma), -\infty < \sigma < \infty$, in such a way that $\lim\limits_{|\sigma| \uparrow \infty} |\lambda(\sigma)| = \infty$ and, for some $\varepsilon > 0$, 
\[
\pi / 2 + \varepsilon \leqq \arg \lambda (\sigma) \leqq \theta_{0} \quad \text{ and } \quad - \theta_{0} \leqq \arg \lambda (\sigma) \leqq - \pi / 2 - \varepsilon
\]
when $\sigma \uparrow +\infty$ and $\sigma \downarrow -\infty$, respectively; for $|\sigma|$ not large, $\lambda (\sigma)$ lies in the right half plane of the complex $\lambda$-plane. 

We shall show that $\hat{T}_{t}$ coincides with the semi-group $T_{t}$ itself. We first show that $\lim\limits_{t\downarrow0}\hat{T}_{t}x=x$ for all $x\in D(A)$. Let $x_{0}$ be any element $\in D(A)$, and choose any complex number $\lambda_{0}$ to the right of the contour $C_{2}$ of integration, and denote $(\lambda_{0}I-A)x_{0}=y_{0}$. Then, by the resolvent equation, 
\[
\begin{aligned}
&\hat {T} _{t} x _{0} = \hat {T} _{t} R (\lambda_{0}; A) y _{0}
\\
&= (2 \pi i) ^{-1} \int_ {C _{2}} e ^{\lambda t} R (\lambda ; A) R (\lambda_{0}; A) y _{0} \dd \lambda
\\
&= (2 \pi i) ^{-1} \int_ {C _{2}} e ^{\lambda t} (\lambda_{0} - \lambda) ^{-1} R (\lambda ; A) y _{0} \dd \lambda - (2 \pi i) ^{-1} \int_ {C _{2}} e ^{\lambda t} (\lambda_{0} - \lambda) ^{-1} R (\lambda_{0}; A) y _{0} \dd \lambda .
\end{aligned}
\]

The second integral on the right is equal to zero, as may be seen by shifting the path of integration to the left. Hence 
\[
\hat {T} _{t} x _{0} = (2 \pi i) ^{-1} \int_ {C _{3}} e ^{\lambda t} (\lambda_{0} - \lambda) ^{-1} R (\lambda ; A) y _{0} \dd \lambda , y _{0} = (\lambda_{0} I - A) x _{0}.
\]

Because of the estimate \eqref{eq:IX.10.5}, the passage to the limit $t \downarrow 0$ under the integral sign is justified, and so 
\[
\lim\limits _{t \downarrow 0} \hat {T} _{t} x _{0} = (2 \pi i) ^{-1} \int_ {C _{*}} (\lambda_{0} - \lambda) ^{-1} R (\lambda ; A) y _{0} \dd \lambda , y _{0} = (\lambda_{0} I - A) x _{0}.
\]

To evaluate the right hand integral, we make a closed contour out of the original path of integration $C_2$ by adjoining the arc of the circle $|\lambda| = r$ which is to the right of the path $C_2$, and throwing away that portion of the original path $C_2$ which lies outside the circle $|\lambda| = r$. The value of the integral along the new arc and the discarded arc tends to zero as $r \uparrow \infty$, in virtue of \eqref{eq:IX.10.5}. Hence the value of the integral is equal to the residue inside the new closed contour, that is, the value $R(\lambda_0; A) y_0 = x_0$. We have thus proved $\lim\limits_{t \downarrow 0} \hat{T}_t x_0 = x_0$ when $x_0 \in D(A)$. 

We next show that $\hat{T}'x = A\hat{T}_t x$ for $t > 0$ and $x \in X$. We have $R(\lambda; A)X = D(A)$ and $AR(\lambda; A) = \lambda R(\lambda; A) - I$, so that, by the
convergence factor $e^{\lambda t}$, the integral $(2\pi i)^{-1} \int_{C_2} e^{\lambda t} A R(\lambda; A)x \, \dd\lambda$ has a sense. This integral is equal to $A \hat{T}_t x$, as may be seen by approximating the integral \eqref{eq:IX.10.6} by Riemann sums and using the fact that $A$ is closed: $\lim\limits_{n \to \infty} x_n = x$ and $\lim\limits_{n \to \infty} A x_n = y$ imply $x \in D(A)$ and $Ax = y$. Therefore 
\[
A \hat {T} _{t} x = (2 \pi i) ^{-1} \int_ {C _{3}} e ^{\mu} A R (\lambda ; A) x \dd \lambda , t > 0.
\]

On the other hand, by differentiating \eqref{eq:IX.10.6} under the integral sign, we obtain 
\begin{equation}\label{eq:IX.10.7}
\hat {T} _{t} ^{\prime} x = (2 \pi i) ^{-1} \int_ {C _{2}} e ^{\lambda t} \lambda R (\lambda ; A) x \dd \lambda , t > 0.
\end{equation}

The difference of these two integrals is $(2\pi i)^{-1}\int_{C_{s}}e^{\lambda t}xd\lambda$, and the value of the last integral is zero, as may be seen by shifting the path of integration to the left. 

Thus we have proved that $\hat{x}(t) = \hat{T}_t x_0$, $x_0 \in D(A)$, satisfies \begin{enumerate}[i)]
\item $\lim\limits_{t \downarrow 0} \hat{x}(t) = x_0$,
\item $\dd\hat{x}(t)/\dd t = A\hat{x}(t)$ for $t > 0$
\item $\{\hat{x}(t)\}$ is by \eqref{eq:IX.10.6} of exponential growth when $t \uparrow \infty$.
\end{enumerate}
 On the other hand, since $x_0 \in D(A)$ and since $\{T_t\}$ is equi-continuous in $t \geqq 0$, we see that $x(t) = T_t x_0$ also satisfies $\lim\limits_{t \downarrow 0} x(t) = x_0$, $\dd x(t)/\dd t = Ax(t)$ for $t \geqq 0$, and $\{x(t)\}$ is bounded when $t \geqq 0$. Let us put $\hat{x}(t) - x(t) = y(t)$. Then $\lim\limits_{t \downarrow 0} y(t) = 0$, $\dd y(t)/\dd t = Ay(t)$ for $t > 0$ and $\{y(t)\}$ is of exponential growth when $t \uparrow \infty$. Hence we may consider the Laplace transform 
\[
L (\lambda ; y) = \int_{0} ^{\infty} e ^{- \lambda t} y (t) \dd t \text{ for large positive } \operatorname{Re} (\lambda).
\]

We have 
\[
\int_{\alpha} ^{\beta} e ^{- \lambda t} y ^{\prime} (t) \dd t = \int_{\alpha} ^{\beta} e ^{- \lambda t} A y (t) \dd t = A \int_{\alpha} ^{\beta} e ^{- \lambda t} y (t) \dd t, 0 \leqq \alpha < \beta < \infty ,
\]
by approximating the integral by Riemann sums and using the fact that $A$ is closed. By partial integration, we obtain 
\[
\int_{\alpha} ^{\beta} e ^{- \lambda t} y ^{\prime} (t) \dd t = e ^{- \lambda \beta} y (\beta) - e ^{- \lambda \alpha} y (\alpha) + \lambda \int_{\alpha} ^{\beta} e ^{- \lambda t} y (t) \dd t,
\]
which tends to $\lambda L(\lambda; y)$ as $\alpha \downarrow 0$, $\beta \uparrow \infty$. For, $y(0) = 0$ and $\{y(\beta)\}$ is of exponential growth as $\beta \uparrow \infty$. Thus again, by using the closure property of $A$, we obtain 
\[
A L (\lambda ; y) = \lambda L (\lambda ; y) \quad \text{ for large positive } \quad \operatorname{Re} (\lambda).
\]

Since the inverse $(\lambda I - A)^{-1}$ exists for $\operatorname{Re}(\lambda) > 0$, we must have $L(\lambda; y) = 0$ when $\operatorname{Re}(\lambda)$ is large positive. Thus, for any continuous linear functional 
$f\in X^{\prime}$, we have 
\[
\int_{0} ^{\infty} e ^{- \lambda t} f (y (t)) \dd t = 0 \quad \text{ when } \quad \operatorname{Re} (\lambda) \text{ is large positive. }
\]

We set $\lambda = \sigma + i\tau$ and put 
\[
g _{\sigma} (t) = e ^{- \sigma t} f (y (t)) \text{ or } = 0 \quad \text{ according as } \quad t \geqq 0 \quad \text{ or } \quad t < 0.
\]

Then, the above equality shows that the Fourier transform $(2\pi)^{-1}\int_{-\infty}^{\infty}e^{-i\tau t}g_{\sigma}(t)\dd t$ vanishes identically in $\tau$, $-\infty < \tau < \infty$, so that, by Fourier's integral theorem, $g_{\sigma}(t) = 0$ identically. Thus $f(y(t)) = 0$ and so we must have $y(t) = 0$ identically, in virtue of the Hahn-Banach theorem. 

Therefore $\hat{T}_t x = T_t x$ for all $t > 0$ and $x \in D(A)$. $D(A)$ being dense in $X$ and $\hat{T}_t, T_t$ both belonging to $L(X, X)$, we easily conclude that $\hat{T}_t x = T_t x$ for all $x \in X$ and $t > 0$. Hence, by defining $\hat{T}_0 = I$, we have $\hat{T}_t = T_t$ for all $t \geqq 0$. Hence, by \eqref{eq:IX.10.7}, $T_t' x = (2\pi i)^{-1} \int_{C_1} e^{\lambda t} \lambda R(\lambda; A) x \dd\lambda$, $t > 0$, and so, by \eqref{eq:IX.10.1} and \eqref{eq:IX.10.5}, we obtain 
\[
(T _{t/n} ^{\prime}) ^{n} x = T _{t} ^{(n)} x = (2 \pi i) ^{-1} \int_ {C _{3}} e ^{\lambda t} \lambda^{n} R (\lambda ; A) x \dd \lambda , t > 0.
\]

Hence 
\[
(t T _{t} ^{\prime}) ^{n}. x = (2 \pi i) ^{-1} \int_ {C _{s}} e ^{n \lambda t} (t \lambda) ^{n} R (\lambda ; A) x \dd \lambda , t > 0.
\]

Therefore, by (III), 
\[
\phi \left((t T _{t} ^{\prime}) ^{n} x\right) \leqq (2 \pi) ^{-1} q (x) \int_ {C _{s}} | e ^{n \lambda t} | t ^{n} | \lambda | ^{n-1} d | \lambda |.
\]

If $0 < t \leqq 1$, then the last integral is majorized by $C_{3}^{n}$ where $C_{3}$ is a certain positive constant. This we see, by splitting the integration path $C_{2}$ into the sum of that in the right half-plane $\operatorname{Re}(\lambda) \geqq 0$ and those in the left half-plane $\operatorname{Re}(\lambda) < 0$, and remembering the integral representation of the $\Gamma$-function.
\end{proof}

References. The result of the present section is due to K. \textsc{Yosida} \cite{yosida349}. See also E. \textsc{Hille} \cite{hille99} and E. \textsc{Hille}-R. S. \textsc{Phillips} \cite{hille97}. 


\yossection{11}{Fractional Powers of Closed Operators}
\label{sec:IX.11}

Let $X$ be a $B$-space, and $\{T_t; t \geqq 0\} \subseteq L(X, X)$ an equi-continuous semi-group of class $(C_0)$. We introduce 
\begin{equation}\label{eq:IX.11.1}
\begin{aligned}
&f _{t, \alpha} (\lambda)  = \frac {1}{2 \pi i} \int_{\sigma - i \infty} ^{\sigma + i \infty} e ^ {i \lambda - t s ^{\alpha}} \dd z (\sigma > 0, t > 0, \lambda \geqq 0, 0 < \alpha < 1), = 0 \quad (\text{ when } \lambda < 0),
\end{aligned}
\end{equation}
where the branch of $z^{\alpha}$ is so taken that $\operatorname{Re}(z^{\alpha}) > 0$ for $\operatorname{Re}(z) > 0$. This branch is a one-valued function in the $z$-plane cut along the negative real axis. The convergence of the integral \eqref{eq:IX.11.1} is apparent in virtue of the convergence factor $e^{-tz^{\alpha}}$. Following S. \textsc{Bochner} \cite{bochner21} and R. S. \textsc{Phillips} \cite{phillips255}, we can show that the operators 
\begin{equation}\label{eq:IX.11.2}
\begin{aligned} \hat {T} _{t, \alpha} x & = \hat {T} _{t} x = \int_{0} ^{\infty} f _{t, \alpha} (s) T _{s} x \dd s (t > 0), \\ & = x (t = 0), \end{aligned}
\end{equation}
constitute an equi-continuous semi-group of class $(C_{0})$. Moreover, we can show that $\{\hat{T}_{t}\}$ is a holomorphic semi-group (K. \textsc{Yosida} \cite{yosida351} and V. \textsc{Balakrishnan} \cite{balakrishnan8}). The infinitesimal generator $\hat{A} = \hat{A}_{\alpha}$ of $\hat{T}_{t}$ is connected with the infinitesimal generator $A$ of $T_{t}$ by 
\begin{equation}\label{eq:IX.11.3}
\hat {A} _{\alpha} x = - (- A) ^{\alpha} x \quad \text{ for } \quad x \in D (A),
\end{equation}
where the fractional power $(-A)^{\alpha}$ of $(-A)$ is given by 
\begin{equation}\label{eq:IX.11.4}
(- A) ^{\alpha} x = \frac {\sin \alpha \pi}{\pi} \int_{0} ^{\infty} \lambda^{\alpha -1} (\lambda I - A) ^{-1} (- A x) \dd \lambda \quad \text{ for } \quad x \in D (A)
\end{equation}
and also by 
\begin{equation}\label{eq:IX.11.5}
(- A) ^{\alpha} x = \Gamma (- \alpha) ^{-1} \int_{0} ^{\infty} \lambda^{- \alpha -1} (T _{\lambda} - I) x \dd \lambda \quad \text{ for } \quad x \in D (A).
\end{equation}

Formula \eqref{eq:IX.11.4} and \eqref{eq:IX.11.5} were obtained by V. \textsc{Balakrishnan}. For the resolvent of $\hat{A}_{\alpha}$, we have the following formula due to T. \textsc{Kato}: 
\begin{equation}\label{eq:IX.11.6}
(\mu I - \hat {A} _{\alpha}) ^{-1} = \frac {\sin \alpha \pi}{\pi} \int_{0} ^{\infty} (r I - A) ^{-1} \frac {r ^{\alpha}}{\mu^{2} - 2 \mu r ^{\alpha} \cos \alpha \pi + r ^{2 \alpha}} \dd r.
\end{equation}

In this way, we see the abundance of holomorphic semi-groups among the class of equi-continuous semi-groups of class $(C_{0})$. 

To prove the above result, we have to investigate the properties of the function $f_{t,\alpha}(\lambda)$ in a series of propositions. 

\begin{proposition}\label{prop:IX.11.Proposition.1} We have 
\begin{equation}\label{eq:IX.11.7}
e ^ {- t a ^{\alpha}} = \int_{0} ^{\infty} e ^{- \lambda a} f _{t, \alpha} (\lambda) \dd \lambda (t > 0, a > 0).
\end{equation}
\end{proposition}

\begin{proof}
 It is easy to see, by the convergence factor $e^{-z^{\alpha}t}$, that the function $f_{t,\alpha}(\lambda)$ is of exponential growth in $\lambda$. By Cauchy's integral theorem, integral \eqref{eq:IX.11.1} is independent of $\sigma > 0$. Let $a > \sigma = \operatorname{Re}(z) > 0$. Then, by Cauchy's theorem of residue, 
\[ \begin{aligned} \int_{0} ^{\infty} e ^{- \lambda a} f _{t, \alpha} (\lambda) \dd \lambda & = \frac {1}{2 \pi i} \int_{\sigma - i \infty} ^{\sigma + i \infty} \left[ \frac {e ^{\lambda (z - a)}}{z - a} \right] _{\lambda = 0} ^{\lambda = \infty} e ^ {- z ^{\alpha} t} \dd z \\ & = \frac {- 1}{2 \pi i} \int_{\sigma - i \infty} ^{\sigma + i \infty} \frac {1}{z - a} e ^ {- z ^{\alpha} t} \dd z = e ^ {- t a ^{\alpha}}. \end{aligned}
\]
\end{proof}

\begin{proposition}\label{prop:IX.11.Proposition.2} We have 
\begin{equation}\label{eq:IX.11.8}
f _{t, \alpha} (\lambda) \geqq 0 \quad \text{ for all } \quad \lambda > 0.
\end{equation}
\end{proposition}

\begin{proof}
 If we set $a^{\alpha} = g(a)$, $e^{-tx} = h(x)$, then 
\[
\begin{aligned}
(-1)^{n-1} g^{(n)}(a) &\geqq 0 \quad (n = 1, 2, \ldots), \quad g(a) \geqq 0, \\
\text{and } (-1)^{n} h^{(n)}(x) &\geqq 0 \quad (n = 0, 1, \ldots), \quad \text{when } a \geqq 0 \text{ and } x \geqq 0.
\end{aligned}
\]

Hence $k(a) = h(g(a)) = e^{-ta^{\alpha}}$ satisfies 
\begin{equation}\label{eq:IX.11.9}
\begin{yoslines}& (- 1) ^{n} k ^{(n)} (a) = (- 1) h ^{\prime} (x) (- 1) ^{n-1} g ^{(n)} (a)\\& + \sum\limits_{(p)} C _ {(p _{0} p _{1} \dots p _{v})} ^{(n)} (- 1) ^ {p _{0}} h ^ {(p _{0})} (x) (- 1) ^ {p _{1}} g ^ {(p _{1} + 1)} (a) \dots (- 1) ^ {p _{v}} g ^ {(p _{v} + 1)} (a)\\& \left(C _{(p)} ^{(n)} \geqq 0, p _{0} \geqq 2, p _{1} \geqq 0, \dots , p _{v} \geqq 0 \text{ with } p _{0} \leqq \sum\limits_{i=1} ^{v} p _{i} = n, \text{ and } v \text{ arbitrary }\right)\\& \geqq 0, (n = 0, 1, \dots).\end{yoslines}
\end{equation}

That is, the function $k(a) = e^{-ta^{\alpha}}$ is completely monotone in $a \geqq 0$. 

We next prove the Post-Widder inversion formula 
\begin{equation}\label{eq:IX.11.10}
f _{t, \alpha} (\lambda) = \lim\limits _{n \rightarrow \infty} \frac {(- 1) ^{n}}{n !} \left(\frac {n}{\lambda}\right) ^{n+1} k ^{(n)} \left(\frac {n}{\lambda}\right), \quad \lambda > 0
\end{equation}
so that, by \eqref{eq:IX.11.9}, $f_{t,\alpha}(\lambda) \geqq 0$. The proof of \eqref{eq:IX.11.10} is obtained as follows. We find, differentiating \eqref{eq:IX.11.7} $n$-times, 
\[
k ^{(n)} \left(\frac {n}{\lambda}\right) = (- 1) ^{n} \int_{0} ^{\infty} s ^{n} e ^{- s n / \lambda} f _{t, \alpha} (s) \dd s.
\]

Substituting this into the right side of \eqref{eq:IX.11.10}, we get 
\[
\lim\limits _{n \rightarrow \infty} \frac {n ^{n+1}}{e ^{n} n !} \frac {1}{\lambda} \int_{0} ^{\infty} \left[ \frac {s}{\lambda} \cdot \exp \left(1 - \frac {s}{\lambda}\right)\right] ^{n} f _{t, \alpha} (s) \dd s.
\]

Since 
\[
\lim\limits _{n \rightarrow \infty} n ^{n} / \sqrt {2 \pi n} e ^{n} n! = 1 \quad (\text{ Stirling's formula }),
\]
we have to prove that 
\begin{equation}\label{eq:IX.11.11}
f _{t, \alpha} (\lambda_{0}) = \lim\limits _{n \rightarrow \infty} \frac {\sqrt {2 \pi}}{\lambda_{0}} \int_{0} ^{\infty} n ^{3/2} \left[ \frac {s}{\lambda_{0}} \cdot \exp \left(1 - \frac {s}{\lambda_{0}}\right)\right] ^{n} f _{t, \alpha} (s) \dd s, \lambda_{0} > 0.
\end{equation}

Let $\eta$ be a fixed positive number such that $\eta < \lambda_0$. We decompose the last integral into three parts, 
\[
\int_{0} ^{\infty} = \int_{0} ^ {\lambda_{0} - \eta} + \int_ {\lambda_{0} - \eta} ^ {\lambda_{0} + \eta} + \int_ {\lambda_{0} + \eta} ^{\infty} = J _{1} + J _{2} + J _{3}.
\]

Since $x \cdot \exp(1 - x)$ increases monotonically in [$0$, $1$] from $0$ to $1$, we see, by the boundedness of $f_{t,\alpha}(s)$ in $s$, that $\lim\limits_{n\to \infty}J_1 = 0$. Next, since $x\cdot \exp (1 - x)$ decreases monotonically in $[1,\infty]$ from $1$ to $0$, we have 
\[
\frac {\lambda_{0} + \eta}{\lambda_{0}} \cdot \exp \left(1 - \frac {\lambda_{0} + \eta}{\lambda_{0}}\right) < \beta < 1,
\]
and so, since $f_{t,\alpha}(s)$ is of exponential growth as $s \uparrow \infty$, 
\[
\left| J _{3} \right| \leqq n ^{3/2} e ^ {n _{0}} \beta^ {n - n _{0}} \int_{0} ^{\infty} \left(\frac {s}{\lambda_{0}}\right) ^ {n _{0}} \exp \left(- \frac {n _{0} s}{\lambda_{0}}\right)\left| f _{t, \alpha} (s) \right| \dd s \rightarrow 0 \text{ as } n \rightarrow \infty .
\]

By the continuity of $f_{t,\alpha}(s)$ in $s$, we have, for any positive number $\varepsilon$, $f_{t,\alpha}(\lambda_0) - \varepsilon \leqq f_{t,\alpha}(s) \leqq f_{t,\alpha}(\lambda_0) + \varepsilon$ whenever $\lambda_0 - \eta \leqq s \leqq \lambda_0 + \eta$, if we take $\eta > 0$ sufficiently small. Thus 
\begin{equation}\label{eq:IX.11.12}
\left(f _{t, \alpha} (\lambda_{0}) - \varepsilon\right) J _{0} \leqq J _{2} \leqq \left(f _{t, \alpha} (\lambda_{0}) + \varepsilon\right) J _{0},
\end{equation}
where 
\begin{equation}\label{eq:IX.11.13}
J _{0} = \int_ {\lambda_{0} - \eta} ^ {\lambda_{0} + \eta} n ^{3/2} \left[ \frac {s}{\lambda_{0}} \cdot \exp \left(1 - \frac {s}{\lambda_{0}}\right) \right] ^{n} \dd s.
\end{equation}

The whole preceding argument is true for the particular case of the completely monotone function 
\[
k (a) = a ^{-1} = \int_{0} ^{\infty} e ^{- \lambda a} \dd \lambda .
\]

In this case, $k^{(n)}(n / \lambda_0) = (-1)^n n!(\lambda_0 / n)^{n+1}$. Substituting this in \eqref{eq:IX.11.10}, we find that \eqref{eq:IX.11.10} holds for $f_{t,\alpha}(\lambda) \equiv 1$. Since \eqref{eq:IX.11.10} and \eqref{eq:IX.11.11} are equivalent, \eqref{eq:IX.11.11} must also hold for $f_{t,\alpha}(\lambda) \equiv 1$. Thus, since $\lim\limits_{n\to \infty}J_1 = 0$ and $\lim\limits_{n\to \infty}J_3 = 0$ for a general $f_{t,\alpha}$, we obtain 
\[
1 = \lim\limits _{n \rightarrow \infty} \sqrt {2 \pi} \lambda_{0} ^{-1} J _{0}.
\]

Therefore, by \eqref{eq:IX.11.12}, we get \eqref{eq:IX.11.11} and the equivalent formula \eqref{eq:IX.11.10} is proved. 
\end{proof}

\begin{proposition}\label{prop:IX.11.Proposition.3} 
\begin{align}
&\int_{0} ^{\infty} f _{t, \alpha} (\lambda) \dd \lambda = 1, \label{eq:IX.11.14} \\
f _{t + s, \alpha} (\lambda)& = \int_{0} ^{\infty} f _{t, \alpha} (\lambda - \mu) f _{s, \alpha} (\mu) \dd \mu . \label{eq:IX.11.15}
\end{align}
\end{proposition}

\begin{proof}
 Since the function $f_{t,\alpha}(\lambda)$ is non-negative, we have, by the Lebesgue-Fatou lemma and \eqref{eq:IX.11.7}, 
\[
\int_{0} ^{\infty} \varliminf\limits_{a \downarrow 0} \left(e ^{- \lambda a} f _{t, \alpha} (\lambda)\right) \dd \lambda \leqq \varliminf\limits_{a \downarrow 0} e ^ {- t a ^{\alpha}} = 1. \]

Thus $f_{t,\alpha}(\lambda)$ is integrable with respect to $\lambda$ over $(0,\infty)$ and so, again by 
the Lebesgue-Fatou lemma and \eqref{eq:IX.11.7}, we obtain \eqref{eq:IX.11.14}. We next have, by \eqref{eq:IX.11.7}, 
\[
\begin{aligned}
&\int e ^{- \lambda a} \left\{\int f _{t, \alpha} (\lambda - \mu) f _{s, \alpha} (\mu) \dd \mu \right\} \dd \lambda = \int e ^{-(\lambda - \mu) a} f _{t, \alpha} (\lambda - \mu) d (\lambda - \mu) \cdot \int e ^{- \mu a} f _{s, \alpha} (\mu) \dd \mu
\\
&= e ^ {- t a ^{\alpha}} e ^ {- s a ^{\alpha}} = e ^ {- (t + s) a ^{\alpha}} = \int e ^{- \lambda a} f _{t + s, \alpha} (\lambda) \dd \lambda , a > 0.
\end{aligned}
\]

Hence, by inverting the Laplace transform as in the preceding section, we get \eqref{eq:IX.11.15}. 
\end{proof}

\begin{proposition}\label{prop:IX.11.Proposition.4} We have 
\begin{equation}\label{eq:IX.11.16}
\int_{0} ^{\infty} \partial f _{t, \alpha} (\lambda) / \partial t \cdot \dd \lambda = 0, \quad t > 0.
\end{equation}
\end{proposition}

\begin{proof}
 By deforming the path of integration in \eqref{eq:IX.11.1} to the union of two paths $re^{-i\theta}(-\infty<-r<0)$ and $re^{i\theta}(0<r<\infty)$, where $\pi/2\leqq\theta\leqq\pi$, we obtain 
\begin{equation}\label{eq:IX.11.17}
\begin{aligned}
&f _{t, \alpha} (s) = \frac {1}{\pi} \int_{0} ^{\infty} \exp (s r \cdot \cos \theta - t r ^{\alpha} \cos \alpha \theta) \times \sin (s r \cdot \sin \theta - t r ^{\alpha} \sin \alpha \theta + \theta) \dd r.
\end{aligned}
\end{equation}

Similarly, by deforming the path of integration in 
\[
\partial f _{t, \alpha} (\lambda) / \partial t = \frac {1}{2 \pi i} \int_{\sigma - i \infty} ^{\sigma + i \infty} e ^ {z \lambda - z ^{\alpha} t} (- z ^{\alpha}) \dd z
\]
to the union of two paths, $re^{-i\theta}(-\infty < -r < 0)$ and $re^{i\theta}(0 < r < \infty)$, we obtain 
\begin{equation}\label{eq:IX.11.18}
\begin{aligned} f _{t, \alpha} ^{\prime} (s) & = \partial f _{t, \alpha} (s) / \partial t = \frac {(- 1)}{\pi} \int_{0} ^{\infty} \exp (s r \cdot \cos \theta - t r ^{\alpha} \cos \alpha \theta) \\ & \times \sin (s r \cdot \sin \theta - t r ^{\alpha} \sin \alpha \theta + \alpha \theta + \theta) r ^{\alpha} \dd r. \end{aligned}
\end{equation}

If we take 
\[
\theta = \theta_{\alpha} = \pi / (1 + \alpha),
\]
then 
\begin{equation}\label{eq:IX.11.19}
f _{t, \alpha} ^{\prime} (s) = \frac {1}{\pi} \int_{0} ^{\infty} \exp \left((s r + t r ^{\alpha}) \cos \theta_{\alpha}\right) \cdot \sin \left((s r - t r ^{\alpha}) \sin \theta_{\alpha}\right) r ^{\alpha} \dd r.
\end{equation}

Thus we see, by the factor $r^{\alpha}$ ( $0 < \alpha < 1$), that $f_{t,\alpha}'(s)$ is integrable with respect to $s$ over $(0, \infty)$. Hence, by differentiating \eqref{eq:IX.11.14} with respect to $t$, we get \eqref{eq:IX.11.16}. 
\end{proof}

We are now able to prove 

\begin{theorem}\label{thm:IX.11.Theorem.1} $\{\hat{T}_t\}$ is a holomorphic semi-group. 
\end{theorem}

\begin{proof}
 That $\{\hat{T}_t\}$ enjoys the semi-group property $\hat{T}_t\hat{T}_s = \hat{T}_{t+s}$ ( $t, s > 0$) is clear from \eqref{eq:IX.11.2} and \eqref{eq:IX.11.15}. We have, by \eqref{eq:IX.11.2} and \eqref{eq:IX.11.17} with $\theta = \theta_{\alpha}$, 
\begin{equation}\label{eq:IX.11.20}
\begin{aligned}
&\hat {T} _{t} x = \frac {1}{\pi} \int_{0} ^{\infty} T _{s} x \dd s \int_{0} ^{\infty} \exp ((s r + t r ^{\alpha}) \cos \theta_{\alpha}) \times \sin ((s r - t r ^{\alpha}) \sin \theta_{\alpha} + \theta_{\alpha}) \dd r,
\end{aligned}
\end{equation}
which gives by the change of variables 
\begin{gather}s = v t ^{1 / \alpha}, r = u t ^{-1 / \alpha}, \hat {T} _{t} x= \frac {1}{\pi} \int_{0} ^{\infty} T _ {v t ^{1 / \alpha}} \cdot x \dd v \int_{0} ^{\infty} \exp \left((u v + u ^{\alpha}) \cos \theta_{\alpha}\right) \label{eq:IX.11.21}
\\\quad \times \sin \left((u v - u ^{\alpha}) \sin \theta_{\alpha} + \theta_{\alpha}\right) \dd u.  \tag{20'} \label{eq:IX.11.20'}
\end{gather}

The second integral on the right is exactly $\pi \cdot f_{1,\alpha}(v)$, and so, by the equi-boundedness of $\{\|T_t x\|\}$ in $t \geqq 0$, we see, by \eqref{eq:IX.11.14}, that 
\begin{equation}\label{eq:IX.11.22}
\left\| \hat {T} _{t} x \right\| \leqq \sup\limits _{t \geqq 0} \left\| T _{t} x \right\| \int_{0} ^{\infty} f _{1, \alpha} (v) \dd v = \sup\limits _{t \geqq 0} \left\| T _{t} x \right\|.
\end{equation}

By passing to the limit $t \downarrow 0$ in $\eqref{eq:IX.11.20}$, which is justified by the integrability of $f_{1,\alpha}(v)$ over $[0, \infty)$, we obtain, by \eqref{eq:IX.11.14}, 
\[
\slim _{t \downarrow 0} \hat {T} _{t} x = \int_{0} ^{\infty} f _{1, \alpha} (v) \dd v \cdot x = x.
\]

Hence $\{\hat{T}_t\}$ is an equi-continuous semi-group of class $(C_0)$ such that \eqref{eq:IX.11.22} holds. 


By the integrability of $f_{t,\alpha}'(s) = \partial f_{t,\alpha}(s) / \partial t$ over $[0, \infty)$ and the equicontinuity of $\{T_t\}$, we obtain, by differentiating \eqref{eq:IX.11.2} with respect to $t$ under the integral sign, 
\begin{equation}\label{eq:IX.11.23}
\begin{aligned}
\hat {T} _{t} ^{\prime} x &= \int_{0} ^{\infty} f _{t, \alpha} ^{\prime} (s) T _{s} x \dd s
\\
&= \frac {1}{\pi} \int_{0} ^{\infty} T _{s} x \dd s \int_{0} ^{\infty} \exp ((s r + t r ^{\alpha}) \cos \theta_{\alpha}) \cdot \sin ((s r - t r ^{\alpha}) \sin \theta_{\alpha}) r ^{\alpha} \dd r,
\end{aligned}
\end{equation}
which is, by the change of variables \eqref{eq:IX.11.21}, 
\[
= \int_{0} ^{\infty} \left(T _ {v t ^{1 / \alpha}} \cdot x\right) \cdot f _{1, \alpha} ^{\prime} (v) \dd v \cdot t ^{-1}.
\]

Thus, by the integrability of $f_{1,\alpha}'(v)$ over $[0, \infty)$ and the equi-continuity of $\{T_t\}$ in $t \geqq 0$, we see that 
\[
\varlimsup\limits_{t \downarrow 0} \| t \hat {T} _{t} ^{\prime} \| < \infty ,
\]
that is, $\{\hat{T}_t\}$ is a holomorphic semi-group.
\end{proof}

\begin{theorem}\label{thm:IX.11.Theorem.2} The infinitesimal generator $\hat{A}_{\alpha}$ of $\hat{T}_t$ is connected with the infinitesimal generator $A$ of $T$ by \eqref{eq:IX.11.3}, where $(-A)^{\alpha}$ is defined by \eqref{eq:IX.11.4} and also by \eqref{eq:IX.11.5}. We have, moreover, \eqref{eq:IX.11.6}. 
\end{theorem}

\begin{proof}
 By \eqref{eq:IX.11.16} and \eqref{eq:IX.11.23}, we obtain 
\begin{equation}\label{eq:IX.11.24}
\begin{aligned}
&\hat {T} _{t} ^{\prime} x = \frac {1}{\pi} \int_{0} ^{\infty} (T _{s} - I) x \dd s \int_{0} ^{\infty} \exp ((s r + t r ^{\alpha}) \cos \theta_{\alpha}) \times \sin ((s r - t r ^{\alpha}) \sin \theta_{\alpha}) r ^{\alpha} \dd r.
\end{aligned}
\end{equation}

If $x \in D(A)$, then $\slim\limits_{s \downarrow 0} s^{-1}(T_s - I)x = Ax$ and $\|(T_s - I) \cdot x\|$ is bounded in $s \geqq 0$. Thus, we obtain, letting $t \downarrow 0$ in \eqref{eq:IX.11.24}, 
\[
\begin{aligned} \slim _{t \downarrow 0} \hat {T} _{t} ^{\prime} x & = \frac {1}{\pi} \int_{0} ^{\infty} (T _{s} - I) x \dd s \int_{0} ^{\infty} \exp (s r \cdot \cos \theta_{\alpha}) \cdot \sin (s r \cdot \sin \theta_{\alpha}) r ^{\alpha} \dd r \\ & = (- \Gamma (- \alpha)) ^{-1} \int_{0} ^{\infty} s ^{- \alpha -1} (T _{s} - I) x \dd s, \end{aligned}
\]
because, by the $\Gamma$-function formulae 
\begin{equation}\label{eq:IX.11.25}
\Gamma (z) = c ^{z} \int_{0} ^{\infty} e ^{-cr} r ^{z-1} \dd r \left(\operatorname{Re} (z) > 0, \operatorname{Re} (c) > 0\right)
\end{equation}
and 
\begin{equation}\label{eq:IX.11.26}
\Gamma (z) \Gamma (1 - z) = \pi / \sin \pi z,
\end{equation}
we obtain, by $(\alpha + 1)\theta_{\alpha} = \pi$, 
\[
\begin{aligned}
&\frac {1}{\pi} \int_{0} ^{\infty} \exp (s r \cdot \cos \theta_{\alpha}) \cdot \sin (s r \cdot \sin \theta_{\alpha}) r ^{\alpha} \dd r = (\pi i) ^{-1} \operatorname{Im} \left\{\int_{0} ^{\infty} e ^ {- r (- s e ^ {i \theta_{\alpha}})} r ^{\alpha} \dd r \right\}
\\
&= (\pi i) ^{-1} \operatorname{Im} ((- s e ^ {i \theta_{\alpha}}) ^{- \alpha -1}) \Gamma (1 + \alpha) = s ^{- \alpha -1} \pi^{-1} \sin (\alpha \pi) \Gamma (1 + \alpha) = s ^{- \alpha -1} \frac {- \Gamma (1 + \alpha)}{\Gamma (- \alpha) \Gamma (1 + \alpha)}
\\
&= (- \Gamma (- \alpha)) ^{-1} s ^{- \alpha -1} .
\end{aligned}
\]

Thus, by $\hat{T}_t' x = \hat{A}_{\alpha} \hat{T}_t x$ (when $t > 0$), the continuity at $t = 0$ of $\hat{T}_t x$ and the closure property of the infinitesimal generator $\hat{A}_{\alpha}$, we obtain 
\[
\hat {A} _{\alpha} x = (- \Gamma (- \alpha)) ^{-1} \int_{0} ^{\infty} s ^{- \alpha -1} (T _{s} - I) x \dd s \quad \text{ when } \quad x \in D (A).
\]

Hence, by \eqref{eq:IX.11.25}, \eqref{eq:IX.11.26} and $(tI - A)^{-1} = \int_{0}^{\infty} e^{-ts} T_s \dd s$, we obtain 
\[
\begin{aligned} \hat {A} _{\alpha} x & = \Gamma (- \alpha) ^{-1} \Gamma (1 + \alpha) ^{-1} \int_{0} ^{\infty} \left\{\int_{0} ^{\infty} e ^{-st} t ^{\alpha} \dd t \right\} (I - T _{s}) x \dd s \\ & = \frac {\sin \alpha \pi}{\pi} \int_{0} ^{\infty} t ^{\alpha} ((t I - A) ^{-1} - t ^{-1} I) x \dd t \\ & = \frac {\sin \alpha \pi}{\pi} \int_{0} ^{\infty} t ^{\alpha -1} (t I - A) ^{-1} A x \dd t \quad \text{ for } \quad x \in D (A). \end{aligned}
\]

Finally we have, by taking $\theta = \pi$ in \eqref{eq:IX.11.17} and \eqref{eq:IX.11.2}, 
\[
\begin{gathered}
(\mu I - \hat {A} _{\alpha}) ^{-1} = \int_{0} ^{\infty} e ^{- \mu t} \hat {T} _{t} \dd t
\\
= \pi^{-1} \int_{0} ^{\infty} \dd r \int_{0} ^{\infty} e ^{-sr} T _{s} \dd s \int_{0} ^{\infty} \exp (- \mu t - t r ^{\alpha} \cos \alpha \pi) \cdot \sin (t r ^{\alpha} \sin \alpha \pi) \dd t
\\
= \pi^{-1} \int_{0} ^{\infty} (r I - A) ^{-1} \left\{\int_{0} ^{\infty} \exp (- \mu t - t r ^{\alpha} \cos \alpha \pi) \cdot \sin (t r ^{\alpha} \sin \alpha \pi) \dd t \right\} \dd r
\\
= \frac {\sin \alpha \pi}{\pi} \int_{0} ^{\infty} (r I - A) ^{-1} \frac {r ^{\alpha}}{\mu^{2} - 2 r ^{\alpha} \mu \cos \alpha \pi + r ^{2 \alpha}} \dd r.
\end{gathered}
\]
\end{proof}

\begin{remarku} Formula \eqref{eq:IX.11.2} was devised by S. \textsc{Bochner} \cite{bochner21} without detailed proof. Cf. R. S. \textsc{Phillips} \cite{phillips255}. That $\hat{T}_t$ is a holomorphic semi-group was proved by K. \textsc{Yosida} \cite{yosida351}, V. \textsc{Balakrishnan} \cite{balakrishnan8} and T. \textsc{Kato} \cite{kato135}. Formulae \eqref{eq:IX.11.4} and \eqref{eq:IX.11.5} are due to V. \textsc{Balakrishnan} \cite{balakrishnan8}, who, by virtue of \eqref{eq:IX.11.4}, defined the fractional power $(-A)^{\alpha}$ of a closed linear operator $A$ satisfying the condition: 
\begin{equation}
\label{eq:IX.11.27}
\begin{minipage}{0.92\linewidth}
\everymath{}\let\limits\nolimits
the resolvent $R (\lambda; A) = (\lambda I - A) ^{-1}$ exists for $\operatorname{Re} (\lambda) > 0$ and $\sup\limits _{\operatorname{Re} (\lambda) > 0} | \operatorname{Re} (\lambda) | \cdot \| R (\lambda; A) \| < \infty.$
\end{minipage}
\end{equation}
\end{remarku}

He also proved that $(-A)^{a}$ enjoys properties to be demanded for the fractional power. In fact we have 

\begin{theorem}\label{thm:IX.11.Theorem.3} Let a closed linear operator $A$ satisfy condition \eqref{eq:IX.11.27}. Then, by \eqref{eq:IX.11.4}, a linear operator $(-A)^{\alpha}$ is defined, and we have 
\begin{align}
(- A) ^{\alpha} (- A) ^{\beta} x& = (- A) ^{\alpha + \beta} x \quad \text{ if } \quad x \in D (A ^{2}) \quad \text{ and } \quad 0 < \alpha , \beta \notag \\
\text{ with } \quad \alpha + \beta& < 1, \label{eq:IX.11.28} \\
\slim _{\alpha \uparrow 1} (- A) ^{\alpha} x& = - A x \quad \text{ if } \quad x \in D (A), \label{eq:IX.11.29} \\
\slim _{\alpha \downarrow 0} (- A) ^{\alpha} x& = x \quad \text{ if } \quad \slim _{\lambda \downarrow 0} \lambda R (\lambda ; A) x = 0, \label{eq:IX.11.30}
\end{align}
and, if $A$ is the infinitesimal generator of an equi-continuous semi-group $T_{t}$ of class $(C_0)$, 
\begin{equation}(A _{\alpha}) _{\beta} = A _{\alpha \beta}, \text{ where } A _{\alpha} \text{ is the operator } \hat {A} _{\alpha} \text{ defined through } \text{ Kato's formula } \eqref{eq:IX.11.6}. \label{eq:IX.11.31}\end{equation}
\end{theorem}

\begin{remarku} The last formula \eqref{eq:IX.11.31} is due to J. \textsc{Watanabe} \cite{watanabe331}. 
\end{remarku}

\begin{proof}
 $\left\|r^{\alpha-1}(rI-A)^{-1}(-Ax)\right\|$ is, by \eqref{eq:IX.11.27}, of order $O(r^{\alpha-2})$ when $r\uparrow\infty$, and, by $(rI-A)^{-1}(-Ax)=x-r(rI-A)^{-1}x$ and \eqref{eq:IX.11.27}, it is of order $O(r^{\alpha-1})$ when $r\downarrow0$. Thus the right side of \eqref{eq:IX.11.4} is convergent. 


It is clear that $x \in D(A^{2})$ implies $(-A)^{\beta}x \in D(A)$. For, by approximating the integral by Riemann sums and making use of the closure property of $A$, we have $(-A)^{\beta}x \subseteq D(A)$. We can thus define $(-A)^{\alpha}(-A)^{\beta}x$. 
\[
(- A) ^{\alpha} (- A) ^{\beta} x = \frac {\sin \alpha \pi}{\pi} \frac {\sin \beta \pi}{\pi} \int_{0} ^{\infty} \int_{0} ^{\infty} \lambda^{\beta -1} \mu^{\alpha -1} R (\lambda ; A) R (\mu ; A) A ^{2} x \dd \lambda \dd \mu
\]
can be rewritten, decomposing the domain of integration into those for which $\lambda \geqq \mu$ and $\lambda < \mu$, as follows. 
\[
\frac {\sin \alpha \pi}{\pi} \frac {\sin \beta \pi}{\pi} \int_{0} ^{1} \left(\sigma^{\beta -1} + \sigma^{\alpha -1}\right) \dd \sigma \int_{0} ^{\infty} \lambda^{\alpha + \beta -1} R (\lambda \sigma ; A) R (\lambda ; A) A ^{2} x \dd \lambda .
\]

By the resolvent equation $R(\lambda; A) - R(\mu; A) = (\mu - \lambda)R(\lambda; A)R(\mu; A)$ and $R(\lambda; A)(-A) = I - \lambda R(\lambda; A)$ valid on $D(A)$, we obtain 
$R(\lambda \sigma ;A)R(\lambda ;A)A^{2}x = (1 - \sigma)^{-1}\{-\sigma R(\lambda \sigma ;A) + R(\lambda ;A)\} (-Ax),$ and hence 
\[
\begin{aligned}
&(- A) ^{\alpha} (- A) ^{\beta} x = \frac {\sin \alpha \pi}{\pi} \frac {\sin \beta \pi}{\pi} \slim _{t \uparrow 1} \int_{0} ^{t} (\sigma^{\beta -1} + \sigma^{\alpha -1}) (1 - \sigma) ^{-1} \dd \sigma \\
&\times \int_{0} ^{\infty} \lambda^{\alpha + \beta -1} (- \sigma R (\lambda \sigma ; A) + R (\lambda ; A)) (- A x) \dd \lambda \\
&= \int_{0} ^{\infty} \left(\frac {\sin \alpha \pi}{\pi} \frac {\sin \beta \pi}{\pi} \int_{0} ^{1} \frac {\sigma^{\beta -1} + \sigma^{\alpha -1} - \sigma^{- \alpha} - \sigma^{- \beta}}{1 - \sigma} \dd \sigma\right) \times \lambda^{\alpha + \beta -1} R (\lambda ; A) (- A x) \dd \lambda .
\end{aligned}
\]

The coefficient () above is evaluated as $\pi^{-1}\sin\pi(\alpha+\beta)$, as may be seen by expanding $(1-\sigma)^{-1}$ into powers of $\sigma$. We have thus proved \eqref{eq:IX.11.28}. 

To prove \eqref{eq:IX.11.29}, we shall make use of $\int_{0}^{\infty} \lambda^{\alpha -1}(1 + \lambda)^{-1} \dd\lambda = \pi / \sin \alpha \pi$. Thus 
\[
(- A) ^{\alpha} x - (- A) x = \frac {\sin \alpha \pi}{\pi} \int_{0} ^{\infty} \lambda^{\alpha -1} (R (\lambda ; A) - \frac {1}{\lambda + 1}) (- A x) \dd \lambda .
\]

We split the integral into two parts, one from $0$ to $C$ and the other from $C$ to $\infty$. For a fixed $C$, the first part goes to zero as $\alpha\uparrow1$, since $R(\lambda;A)(-Ax)=x-\lambda R(\lambda;A)x$ is bounded in $\lambda>0$. The second part is, in norm, 
\[
\leqq \frac {\sin \alpha \pi}{\pi (1 - \alpha)} C ^{\alpha -1} \sup\limits _{\lambda \geqq C} \left\| (\lambda R (\lambda ; A) - \frac {\lambda}{\lambda + 1}) A x \right\|.
\]

We have $\slim\limits_{\lambda \uparrow \infty} \lambda R(\lambda; A) x = x$ by $x - \lambda R(\lambda; A) x = R(\lambda; A) (-A x)$ and \eqref{eq:IX.11.27}. Hence $\slim\limits_{\alpha \uparrow 1}$ of the second part is arbitrarily near $0$ if we take $C$ sufficiently large. This prove \eqref{eq:IX.11.29}. 

To prove \eqref{eq:IX.11.30}, we split the integral into two parts, one from $0$ to $C$ and the other from $C$ to $\infty$. Because of \eqref{eq:IX.11.27}, the second part tends to zero as $\alpha \downarrow 0$. By $R(\lambda; A)(-Ax) = x - \lambda R(\lambda; A)x$ and the hypothesis $\slim\limits_{\lambda \downarrow 0} \lambda R(\lambda; A)x = 0$, the first is, for sufficiently small $C$, arbitrarily near the value $(\alpha\pi)^{-1} \sin \alpha\pi \cdot C^{\alpha}x$, which tends to $x$ when $\alpha \downarrow 0$. This proves \eqref{eq:IX.11.30}. 

We shall prove \eqref{eq:IX.11.31}. By virtue of the representation \eqref{eq:IX.11.6}, we obtain 
\[
\begin{aligned}
(\mu I - (A _{\alpha}) _{\beta}) ^{-1}& = \int_{0} ^{\infty} \int_{0} ^{\infty} (2 \pi i) ^{-2} \left(\frac {1}{\mu - \lambda^{\beta} \cdot e ^{- i \pi \beta}} - \frac {1}{\mu - \lambda^{\beta} \cdot e ^{i \pi \beta}}\right)
\\
&\quad \times \left(\frac {1}{\lambda - \zeta^{\alpha} \cdot e ^{- i \pi \alpha}} - \frac {1}{\lambda - \zeta^{\alpha} \cdot e ^{i \pi \alpha}}\right) (\zeta I - A) ^{-1} \dd \lambda \dd \zeta .
\end{aligned}
\]

This double integral is absolutely convergent in norm, and so we may interchange the order of integration. Hence we obtain \eqref{eq:IX.11.31}, because the inner integral is 
\[
\begin{aligned}
&= (2 \pi i) ^{-2} \int_{C} \frac {- 1}{\mu - z ^{\beta}} \left(\frac {1}{z - \zeta^{\alpha} e ^{- i \pi \alpha}} - \frac {1}{z - \zeta^{\alpha} e ^{i \pi \alpha}}\right) \dd z
\\
&= (2 \pi i) ^{-1} \left(\frac {- 1}{\mu - \zeta^{\alpha \beta} e ^{- i \pi \alpha \beta}} - \frac {- 1}{\mu - \zeta^{\alpha \beta} e ^{i \pi \alpha \beta}}\right)
\end{aligned}
\]
where the path of integration $C$ runs from $\infty e^{i\pi}$ to $0$ and from $0$ to $\infty e^{-i\pi}$.
\end{proof}

An example of the fractional power. If $\alpha = 1/2$, then we have, by taking $\theta = \pi$ in \eqref{eq:IX.11.17}, 
\begin{equation}\label{eq:IX.11.32}
f _{t,1/2} (s) = \pi^{-1} \int_{0} ^{\infty} e ^{-sr} \sin \left(t r ^{1/2}\right) \dd r = \pi^{-1} \sqrt {\pi} t \left(2 ^{3} \sqrt {s}\right) ^{-3} \cdot e ^ {- t ^{2} / 4 s}.
\end{equation}

Thus, if we take the semi-group $\{T_t\}$ associated with the Gaussian kernel: 
\[
\left(T _{t} x\right) (u) = \frac {1}{2 \sqrt {\pi s}} \int_{- \infty} ^{\infty} e ^ {- (u - v) ^{2} / 4 s} x (v) \dd v, x \in C [ - \infty , \infty ],
\]
then 
\[
\begin{aligned}
&\left(\hat {T} _{t,1/2} x\right) (u)  = \int_{- \infty} ^{\infty} \left\{\int_{0} ^{\infty} x (v) \frac {t}{4 \pi s ^{2}} e ^ {- ((u - v) ^{2} + t ^{2}) / 4 s} \dd s \right\} \dd v = \frac {t}{\pi} \int_{- \infty} ^{\infty} \frac {1}{t ^{2} + (u - v) ^{2}} x (v) \dd v,
\end{aligned}
\]
that is, the semi-group $\{\hat{T}_t\}$ is associated with the Poisson kernel. In this case, the infinitesimal generator $A$ of $T_{t}$ is given by the differential operator $d^2 /\dd s^2$, while the infinitesimal generator $\hat{A}$ of $\hat{T}_t$ is given by the singular integral operator 
\[
(\hat {A} _{1/2} x) (s) = \slim _{h \downarrow 0} \frac {1}{\pi} \int_{- \infty} ^{\infty} \frac {x (s - v) - x (s)}{v ^{2} + h ^{2}} \dd v,
\]
and not by the differential operator $\dd/\dd s$. For another example, see K. \textsc{Yosida} \cite{yosida373}. 


\yossection{12}{The Convergence of Semi-groups. The Trotter-Kato Theorem}
\label{sec:IX.12}

We shall denote by $\exp(tA)$ the semi-group of class $(C_{0})$ with the infinitesimal generator $A$. Concerning the convergence of semi-groups, we have 

\begin{theorem}\label{thm:IX.12.Theorem.1} Let $X$ be a locally convex, sequentially complete complex linear space. Let $\{\exp(tA_{n})\} \subseteq L(X, X)$ be a sequence of equi-continuous semi-groups of class $(C_{0})$ such that the family of operators $\{\exp(tA_{n})\}$ is equi-continuous in $t \geqq 0$ and in $n = 1, 2, \ldots$. Thus we assume that, for any continuous semi-norm $p(x)$ on $X$, there exists a continuous semi-norm $q(x)$ on $X$ such that 
\begin{equation}\label{eq:IX.12.1}
p(\exp(tA_{n})x) \leqq q(x) \text{ for all } t \geqq 0, x \in X \text{ and } n = 1, 2, \ldots.
\end{equation}
Suppose that, for some $\lambda_0$ with $\operatorname{Re}(\lambda_0) > 0$, 
\begin{equation}
\label{eq:IX.12.2}
\begin{minipage}{0.92\linewidth}
\everymath{}\let\limits\nolimits
$\lim\limits_{n\to \infty}R(\lambda_0;A_n)x = J(\lambda_0)x$ exists for all $x\in X$ in such a way that the range $R(J(\lambda_0))$ is dense in $X.$
\end{minipage}
\end{equation}
Then $J(\lambda_0)$ is the resolvent of the infinitesimal generator $A$ of an equicontinuous semi-group $\exp(tA)$ of class $(C_0)$ in $X$ and 
\begin{equation}\label{eq:IX.12.3}
\lim\limits _{n \rightarrow \infty} \exp (t A _{n}) x = \exp (t A) x \text{ for every } x \in X.
\end{equation}
\end{theorem}

Moreover, the convergence in \eqref{eq:IX.12.3} is uniform in $t$ on every compact interval of $t$. 

For the proof, we prepare 

\begin{lemmau} Let $T_{t} = \exp(tA)$ be an equi-continuous semi-group of class $(C_0)$ in $X$. Then, for any continuous semi-norm $p(x)$ on $X$, there exists a continuous semi-norm $q(x)$ on $X$ such that 
\begin{equation}\label{eq:IX.12.4}
p\left(T_{t}x - (I - t n^{-1}A)^{-n}x\right) \leqq (2n)^{-1}t^{2}q\left(A^{2}x\right) \quad (n = 1, 2, \ldots)
\end{equation}
whenever $x \in D(A^2)$. 
\end{lemmau}


\begin{proof}
 Set $T(t, n) = (I - n^{-1}tA)^{-n}$. Then we know (Chapter~\ref{chap:IX}, \ref{sec:IX.7}) that $\{T(t, n)\}$ is equi-continuous in $t \geqq 0$ and $n = 1, 2, \ldots$. Moreover, we have (Chapter~\ref{chap:IX}, \ref{sec:IX.4}), for any $x \in D(A)$, 
\begin{equation}D _{t} T (t, n) = (I - n ^{-1} t A) ^{-n-1} A x = A (I - n ^{-1} t A) ^{-n-1} x, D _{t} T x = T _{t} A x = A T _{t} x. \notag\end{equation}

Thus, by the commutativity of $T_{t}$ and $T(t, n)$, 
\begin{equation}\label{eq:IX.12.5}
\begin{aligned} T _{t} x - T (t, n) & x = \int_{0} ^{t} [ D _{s} T (t - s, n) T _{s} x ] \dd s \\ & = \int_{0} ^{t} T (t - s, n) T _{s} \left(A x - \left(I - \frac {t - s}{n} A\right) ^{-1} A x\right) \dd s, x \in D (A). \end{aligned}
\end{equation}

Hence, if $x \in D(A^2)$, we have, by $(I - m^{-1}A)^{-1}Ax = -m(I - (I - m^{-1}A)^{-1})x$, 
\[
p (T _{t} x - T (t, n) x) \leqq \int_{0} ^{t} p [ T (t - s, n) T _{s} (I - n ^{-1} (t - s) A) ^{-1} \frac {s - t}{n} A ^{2} x ] \dd s,
\]
and so, by a continuous semi-norm $q(x)$ on $X$ which is independent of $x$ and $n$, 
\[
p \left(T _{t} x - T (t, n) x\right) \leqq (2 n) ^{-1} t ^{2} q \left(A ^{2} x\right).
\]
\end{proof}

\begin{corollaryu} For any $x \in D(A^2)$, $s > 0$ and $t \geqq 0$, 
\begin{equation}\label{eq:IX.12.6}
p \left(T _{t} x - (I - s A) ^{-[t/s]} x\right) \leqq s q _{1} (A x) + \frac {t s}{2} q \left(A ^{2} x\right), \text{ where }
\end{equation}
$q_{1}(x)$ is a continuous semi-norm on $X$ which is independent of $x, t$ and $s$, and $[t / s]$ is the largest integer $\leqq t / s$. 
\end{corollaryu}

\begin{proof}
 We have, for $t = ns$, 
\[
p \left(T _{ns} x - (I - s A) ^{-n} x\right) \leqq 2 ^{-1} s t q \left(A ^{2} x\right). \]

If $t = ns + u$ with $0 \leqq u < s$ and $n = [t / s]$, 
\[
p \left(T _{t} x - T _{ns} x\right) = p \left(\int_{ns} ^{t} T _{\sigma} ^{\prime} x \dd \sigma\right) \leqq \int_{ns} ^{t} p \left(T _{\sigma} A x\right) \dd \sigma \leqq s q _{1} (A x).
\]
\end{proof}

\begin{proof}[Proof of Theorem~\ref{thm:IX.12.Theorem.1}.]
By \eqref{eq:IX.12.1} and by \eqref{eq:IX.4.11} in Chapter~\ref{chap:IX}, \ref{sec:IX.4}, we see that $\{(\operatorname{Re}(\lambda) R(\lambda; A_n)^m)\}$ is equi-continuous in $\operatorname{Re}(\lambda) > 0$, in $n = 1, 2, \ldots$ and in $m = 0, 1, 2, \ldots$. From this and by \eqref{eq:IX.12.2}, we can prove that $J(\lambda_0) = (\lambda_0 I - A)^{-1}$ with some $A$ and
\begin{equation}
\label{eq:IX.12.7}
\begin{minipage}{0.92\linewidth}
\everymath{}\let\limits\nolimits
$\lim\limits _{n \rightarrow \infty} R (\lambda; A _{n}) x = R (\lambda; A) x$ whenever $\operatorname{Re} (\lambda) > 0,$ and the convergence is uniform with respect to $\lambda$ on every compact subset of the right half-plane $\operatorname{Re} (\lambda) > 0.$
\end{minipage}
\end{equation}

To this purpose, we observe that 
\[
R (\lambda ; A _{n}) x = \sum\limits_{m=0} ^{\infty} (\lambda_{0} - \lambda) ^{m} R (\lambda_{0}; A _{n}) ^{m+1} x \text{ (for } | \lambda - \lambda_{0} | / \operatorname{Re} (\lambda_{0}) < 1)
\]
and that the series is, by the equi-continuity of $\{(\operatorname{Re}(\lambda_0)R(\lambda_0^{\prime};A_n))^m\}$ in $n = 1,2,\ldots$ and $m = 0,1,2,\ldots$, uniformly convergent for $|\lambda -\lambda_0| / \operatorname{Re}(\lambda_0)\leqq 1 - \varepsilon$ and $n = 1,2,\ldots$. Here $\varepsilon$ is a fixed positive number. Therefore, for any $\delta >0$, there exists an $m_0$ and a continuous semi-norm $q(x)$ on $X$ such that, for $|\lambda -\lambda_0| / \operatorname{Re}(\lambda_0)\leqq 1 - \varepsilon$, 
\begin{align}
p \left(R (\lambda ; A _{n}) x - R (\lambda ; A _ {n ^{\prime}}) x\right)& \leqq \sum\limits_{m=0} ^ {m _{0}} | \lambda_{0} - \lambda | ^{m}. \notag \\
p \left(R \left(\lambda_{0}; A _{n}\right) ^ {m _{0} + 1} x - R \left(\lambda_{0}; A _ {n ^{\prime}}\right) ^ {m _{0} + 1} x\right) + 2 \delta q (x) \text{ for all } x& \in X. \notag
\end{align}

Hence, by \eqref{eq:IX.12.2}, we see that $\lim\limits_{n\to \infty}R(\lambda ;A_n)x = J(\lambda)x$ exists uniformly for $|\lambda -\lambda_0| / \operatorname{Re}(\lambda_0)\leqq 1 - \varepsilon$. In this way, extending the convergence domain of the sequence $\{R(\lambda ;A_n)\}$, we see that 
$\lim\limits_{n\to \infty}R(\lambda ;A)x = J(\lambda)x$ exists and the convergence is uniform 
on any compact set of $\lambda$ in the right half-plane $\operatorname{Re}(\lambda) > 0$. 

Thus $J(\lambda)$ is a pseudo-resolvent, because $J(\lambda)$ satisfies the resolvent equation with $R(\lambda; A_n)$ ( $n = 1, 2, \ldots$). However, by $R(J(\lambda_0))^a = X$ and the ergodic theorem for pseudo-resolvents in Chapter~\ref{chap:VIII}, \ref{sec:VIII.4}, we see that $J(\lambda)$ is the resolvent of a closed linear operator $A$ in such a way that $J(\lambda) = R(\lambda; A)$ and $D(A) = R(R(\lambda; A))$ is dense in $X$. 

Thus we see that $\exp(tA)$ is an equi-continuous semi-group of class $(C_{0})$ in $X$. We have to show that \eqref{eq:IX.12.3} is true. But we have, by \eqref{eq:IX.12.6}, 
\[
\begin{aligned} & p \left((\exp (t A _{n}) - (I - s A _{n}) ^{-[t/s]}) (I - A _{n}) ^{-2} x\right) \\ & \leqq s q _{1} (A _{n} (I - A _{n}) ^{-2} x) + 2 ^{-1} t s q (A _{n} ^{2} (I - A _{n}) ^{-2} x), \end{aligned}
\]
for any $x \in X$, $s > 0$ and $t \geqq 0$. 

The operators 
\[
\begin{aligned} A _{n} (I - A _{n}) ^{-1} & = (I - A _{n}) ^{-1} - I, A _{n} (I - A) ^{-2} = A _{n} (I - A _{n}) ^{-1} (I - A _{n}) ^{-1} \\ & \text{ and } A _{n} ^{2} (I - A _{n}) ^{-2} = (A _{n} (I - A _{n}) ^{-1}) ^{2} \end{aligned}
\]
are equi-continuous in $n = 1,2,\ldots$. On the other hand, by \eqref{eq:IX.12.7}, 
\[
\lim\limits _{n \rightarrow \infty} (I - s A _{n}) ^{-[t/s]} (I - A _{n}) ^{-2} x = (I - s A) ^{-[t/s]} (I - A) ^{-2} x
\]
uniformly in $s$ and $t$, if $s > 0$ is bounded away from $0$ and $\infty$, and if $t$ runs over a compact interval of $[0, \infty)$. Moreover, we have, by \eqref{eq:IX.12.6}, 
\[
\begin{aligned}
&\phi \left((\exp (t A) - (I - s A) ^{-[t/s]}) (I - A) ^{-2} x\right) \leqq s q _{1} (A (I - A) ^{-2} x) + 2 ^{-1} t s q (A ^{2} (I - A) ^{-2} x)
\end{aligned}
\]
for every $x \in X$, $s > 0$ and $t \geqq 0$. Thus, by taking $s > 0$ sufficiently small, we see that 
\[
\lim\limits _{n \rightarrow \infty} \exp (t A _{n}) y = \exp (t A) y \text{ for any } y \in R (1; A) ^{2} \cdot X,
\]
and the convergence is uniform in $t$ on every compact interval of $t$. $R(1;A)^{2} \cdot X$ being dense in $X$, we see, by the equi-continuity of $\exp(tA)$ and $\exp(tA_{n})$ in $t \geqq 0$ and in $n = 1, 2, \ldots$, that \eqref{eq:IX.12.3} holds. 
\end{proof}

\begin{theorem}\label{thm:IX.12.Theorem.2} Let a sequence $\{\exp (tA_n)\}$ of equi-continuous semi-groups of class $(C_0)$ in $X$ be such that $\{\exp (tA_n)\}$ is equi-continuous in $t\geqq 0$ and in $n = 1,2,\ldots$. If, for each $x\in X$, 
\[
\lim\limits _{n \rightarrow \infty} \exp (t A _{n}) x = \exp (t A) x
\]
uniformly with respect to $t$ on every compact interval of $t$, then 
$\lim\limits_{n \to \infty} R(\lambda; A_n) x = R(\lambda; A) x$ for each $x \in X$ and $\operatorname{Re}(\lambda) > 0$ and the convergence is uniform on every compact set of $\lambda$ in the right half-plane $\operatorname{Re}(\lambda) > 0$. 
\end{theorem}

\begin{proof}
 We have 
\[
R (\lambda ; A) x - R (\lambda ; A _{n}) x = \int_{0} ^{\infty} e ^{- \lambda t} (\exp (t A) - \exp (t A _{n})) x \dd t. \]

Hence, splitting the integral into two parts, one from $0$ to $C$ and the other from $C$ to $\infty$, we obtain the result. 
\end{proof}

\begin{remarku} For the case of a Banach space $X$, Theorem~\ref{thm:IX.12.Theorem.1} was first proved by H. F. \textsc{Trotter} \cite{trotter327}. In his paper, the proof that $J(\lambda)$ is the resolvent $R(\lambda; A)$ is somewhat unclear. This was pointed out by T. \textsc{Kato}. The proof given above is adapted from Kato's modification of Trotter's proof. For perturbation of semi-groups, see E. \textsc{Hille}-R. S. \textsc{Phillips} \cite{hille97}, T. \textsc{Kato} \cite{kato142} and K. \textsc{Yosida} \cite{yosida374}. 
\end{remarku}


\yossection{13}{Dual Semi-groups. Phillips' Theorem}
\label{sec:IX.13}

Let $X$ be a locally convex sequentially complete linear topological space, and $\{T_t; t \geqq 0\} \subseteq L(X, X)$ an equi-continuous semi-group of class $(C_0)$. Then the family $\{T_t^*; t \geqq 0\}$ of operators $\in L(X', X')$ where (*) denotes the dual operators in this section, satisfies the semi-group property: $T_t^* T_s^* = T_{t+s}^*$, $T_0^* = I^* =$ the identity in $X_s'$ (see Theorem~\ref{thm:VII.1.Theorem.3} in Chapter~\ref{chap:VII}, \ref{sec:VII.1}). However, it is not of class $(C_0)$ in general. For, the mapping $T_t \to T_t^*$ does not necessarily conserve the continuity in $t$ (see Proposition~\ref{prop:VII.1.Proposition.1} in Chapter~\ref{chap:VII}, \ref{sec:VII.1}). But we can show that $\{T^*\}$ is equicontinuous in $t \geqq 0$. For, we can prove 

\begin{proposition}\label{prop:IX.13.Proposition.1} If $\{S_t; t \geqq 0\} \subseteq L(X, X)$ is equi-continuous in $t \geqq 0$, then $\{S_t^*; t \geqq 0\} \subseteq L(X', X')$ is also equi-continuous in $t \geqq 0$. 
\end{proposition}

\begin{proof}
 For any bounded set $B$ of $X$, the set $\bigcup\limits_{t\geqq 0}S_t\cdot B$ is by hypothesis a bounded set of $X$. Let $U'$ and $V'$ be the polar sets of $B$ and $\bigcup\limits_{t\geqq 0}S_t\cdot B$: $U' = \{x' \in X'; \sup\limits_{b\in B}|\langle b,x'\rangle |\leqq 1\}$, $V' = \{x' \in X'; \sup\limits_{b\in B,t\geqq 0}|\langle S_t\cdot b,x'\rangle |\leqq 1\}$. 

Then (see Chapter~\ref{chap:IV}, \ref{sec:IV.7}), $U'$ and $V'$ are neighbourhoods of $0$ of $X_s'$. From 
\[
\left| \langle S _{t} \cdot b, x ^{\prime} \rangle \right| = | < b, S _{t} ^{*} x ^{\prime} \rangle | \leqq 1 (\text{ when } b \in B, x ^{\prime} \in V ^{\prime})
\]
we see that $S_{t}^{*} \cdot V' \subseteq U'$ for all $t \geqq 0$. This proves that $\{S_{t}^{*}\}$ is equicontinuous in $t \geqq 0$. 

Let $A$ be the infinitesimal generator of the semi-group $T_{t}$. Then $D(A)^{a} = X$, $R(A) \subseteq X$, and, for $\lambda > 0$, the resolvent $(\lambda I - A)^{-1} \in L(X, X)$ exists in such a way that 
\begin{equation}
\label{eq:IX.13.1}
\begin{minipage}{0.92\linewidth}
\everymath{}\let\limits\nolimits
$\{\lambda^m (\lambda I - A)^{-m}\}$ is equi-continuous in $\lambda >0$ and in $m = 0,1\ldots$
\end{minipage}
\end{equation}
We can prove (cf. Theorem~\ref{thm:VIII.6.Theorem.2} in Chapter~\ref{chap:VIII}, \ref{sec:VIII.6}) 
\end{proof}


\begin{proposition}\label{prop:IX.13.Proposition.2} For $\lambda > 0$, the resolvent $(\lambda I^{*} - A^{*})^{-1}$ exists and 
\begin{equation}\label{eq:IX.13.2}
(\lambda I ^{*} - A ^{*}) ^{-1} = ((\lambda I - A) ^{-1}) ^{*}.
\end{equation}
\end{proposition}

\begin{proof}
 We have $(\lambda I - A)^{*} = \lambda I^{*} - A^{*}$. Since $(\lambda I - A)^{-1} \in L(X, X)$, the operator $((\lambda I - A)^{-1})^{*} \in L(X', X')$ exists. We shall prove that $(\lambda I^{*} - A^{*})^{-1}$ exists and equals $((\lambda I - A)^{-1})^{*}$. Suppose there exists an $x' \in X'$ such that $(\lambda I^{*} - A^{*}) x' = 0$. Then $0 = \langle x, (\lambda I^{*} - A^{*}) x' \rangle = \langle (\lambda I - A) x, x' \rangle$ for all $x \in D(A)$. But, since $R(\lambda I - A) = X$, we must have $x' = 0$. Hence the inverse $(\lambda I^{*} - A^{*})^{-1}$ must exist. We have, for $x \in X$, $x' \in D(A^{*})$, 
\[
\langle x, x ^{\prime} \rangle = \left\langle (\lambda I - A) (\lambda I - A) ^{-1} x, x ^{\prime} \right\rangle = \left\langle (\lambda I - A) ^{-1} x, (\lambda I ^{*} - A ^{*}) x ^{\prime} \right\rangle . \]

Thus $D(((\lambda I - A)^{-1})^*) \supseteq R(\lambda I^* - A^*)$ and $((\lambda I - A)^{-1})^* \cdot (\lambda I^* - A^*) x' = x'$ for every $x' \in D(A^*)$. This proves that $((\lambda I - A)^{-1})^* \supseteq (\lambda I^* - A^*)^{-1}$. On the other hand, if $x \in D(A)$ and $x' \in D\left((\lambda I - A)^{-1})^*\right)$, then 
$\langle x, x' \rangle = \langle (\lambda I - A)^{-1} (\lambda I - A) x, x' \rangle = \langle (\lambda I - A) x, ((\lambda I - A)^{-1})^* x' \rangle$. This proves that $D(A^*) = D((\lambda I - A)^*) \supseteq R(((\lambda I - A)^{-1})^*)$ and $(\lambda I - A)^* \cdot ((\lambda I - A)^{-1})^* x' = x'$ for every $x' \in D(((\lambda I - A)^{-1})^*)$, that is, $((\lambda I - A)^{-1})^* \subseteq (\lambda I^* - A^*)^{-1}$. We have thus proved \eqref{eq:IX.13.2}. 
\end{proof}

We are now ready to prove 

\begin{theoremu} Let $X$ be a locally convex sequentially complete linear topological space such that its strong dual space $X'$ is also sequentially complete. Let $\{T_{t}\} \subseteq L(X, X)$, $t \geqq 0$, be an equi-continuous semi-group of class $(C_{0})$ with the infinitesimal generator $A$. Let us denote by $X^{+}$ the closure $D(A^{*})^{a}$ of the domain $D(A^{*})$ in the strong topology of $X'$. Let $T_{t}^{+}$ be the restriction of $T_{t}^{*}$ to $X^{+}$. Then $T_{t}^{+} \in L(X^{+}, X^{+})$ and $\{T_{t}^{+}; t \geqq 0\}$ is an equi-continuous semi-group of class $(C_{0})$ such that its infinitesimal generator $A^{+}$ is the largest restriction of $A^{*}$ with domain and range in $X^{+}$. 
\end{theoremu}

\begin{remarku} The above Theorem was proved by R. S. \textsc{Phillips} \cite{phillips252} for the special case of a $B$-space $X$. The extension given above is due to H. \textsc{Komatsu} \cite{komatsu156}. 
\end{remarku}

\begin{proof}[Proof of the Theorem.]
We have the resolvent equation $R(\lambda; A) - R(\mu; A) = (\mu - \lambda)R(\lambda; A)R(\mu; A)$ and the equi-continuity of $\{\lambda^m R(\lambda; A)^m\}$ in $\lambda > 0$ and in $m = 0, 1, 2, \ldots$. Thus, by Proposition~\ref{prop:IX.13.Proposition.1} and $2$,
\begin{equation}\label{eq:IX.13.3}
(\lambda I ^{*} - A ^{*}) ^{-1} - (\mu I ^{*} - A ^{*}) ^{-1} = (\mu - \lambda) (\lambda I ^{*} - A ^{*}) ^{-1} (\mu I ^{*} - A ^{*}) ^{-1}
\end{equation}
$\{\lambda^m (\lambda I^* - A^*)^{-m}\}$ is equi-continuous in $\lambda > 0$ and in $m = 0,1,2,\ldots$ ($4$) 

Thus, if we denote by $J(\lambda)$ the restriction to $X^{+}$ of $(\lambda I^{*} - A^{*})^{-1}$, we have 
\begin{equation}\tag{3'}\label{eq:IX.13.3'}
J (\lambda) - J (\mu) = (\mu - \lambda) J (\lambda) J (\mu),
\end{equation}
$\{\lambda^m J(\lambda)^m\}$ is equi-continuous in $\lambda > 0$ and in $m = 0, 1, \ldots$ ($4$') 

Since $D(A^{*})$ is dense in $X^{+}$ and since $(4')$ holds, we see, as in Chapter~\ref{chap:IX}, \ref{sec:IX.7}, that $\lim\limits_{\lambda \to \infty} \lambda J(\lambda) x = x$ for $x \in X^{+}$. Thus we have $R(J(\lambda))^{a} = X^{+}$ and so, by $(7')$ in Chapter~\ref{chap:VIII}, \ref{sec:VIII.4}, that $N(J(\lambda)) = \{0\}$. Thus the pseudo-resolvent $J(\lambda)$ must be the resolvent of a closed linear operator $A^{+}$ in $X^{+}$. Hence, by the sequential completeness of $X^{+}$ and $(4')$, $A^{+}$ is the infinitesimal generator of an equi-continuous semi-group of class $(C_0)$ of operators $T_t^+ \in L(X^+, X^+)$. For any $x \in X$ and $y' \in X^{+}$, we have 
\[
\langle (I - m ^{-1} t A) ^{-m} x, y ^{\prime} \rangle = \langle x, (I ^{*} - m ^{-1} t A ^{+}) ^{-m} y ^{\prime} \rangle ,
\]
and so, by the result of the preceding section, we obtain by letting $m \to \infty$ the equality $\langle T_{t}x, y' \rangle = \langle x, T_{t}^{+}y' \rangle$. Hence $T_{t}^{*}y' = T_{t}^{+}y'$, that is, $T_{t}^{+}$ is the restriction to $X^{+}$ of $T_{t}^{*}$. 

We finally show that $A^{+}$ is the largest restriction of $A^{*}$ with domain and range in $X^{+}$. It is clear, by the above derivation of the operator $A^{+}$, that $A^{+}$ is a restriction of $A^{*}$. Suppose that $x' \in D(A^{*})$ and that $x' \in X^{+}$, $A^{*}x' \in X^{+}$. Then $(\lambda I^{*} - A^{*})x' \in X^{+}$ and hence $(\lambda I^{*} - A^{+})^{-1}(\lambda I^{*} - A^{*})x' = x'$. Thus, applying $(\lambda I^{*} - A^{+})$ from the left on both sides, we obtain $A^{*}x' = A^{+}x'$. This proves that $A^{+}$ is the largest restriction of $A^{*}$ with domain as well as range in $X^{+}$. 
\end{proof}
   % IX
% ==== $X$ ====
\yoschapter{X}{Compact Operators}
\label{chap:X}

Let $X$ and $Y$ be complex $B$-spaces, and let $S$ be the unit sphere in $X$. An operator $T \in L(X, Y)$ is said to be compact or completely continuous if the image $T \cdot S$ is relatively compact in $Y$. For a compact operator $T \in L(X, X)$, the eigenvalue problem can be treated fairly completely, in the sense that the classical theory of Fredholm concerning linear integral equations may be extended to the linear functional equation $Tx - \lambda x = y$ with a complex parameter $\lambda$. This result is known as the Riesz-Schauder theory. F. \textsc{Riesz} \cite{rickart273} and J. \textsc{Schauder} \cite{schauder289}. 


\yossection{1}{Compact Sets in \texorpdfstring{$B$}{B}-spaces}
\label{sec:X.1}

A compact set in a linear topological space must be bounded. The converse is, however, not true in general; we know (Chapter~\ref{chap:III}, \ref{sec:III.2}) that the closed unit sphere of a normed linear space $X$ is strongly compact iff $X$ is of finite dimension. Let $S$ be a compact metric space and $C(S)$ the $B$-space of real- or complex-valued continuous functions $x(s)$ on $S$, normed by $\|x\| = \sup\limits_{s \in S} |x(s)|$. We know (Chapter~\ref{chap:III}, \ref{sec:III.3}) that a subset $\{x_{\alpha}(s)\}$ of $C(S)$ is strongly relatively compact in $C(S)$ iff $\{x_{\alpha}(s)\}$ is equi-bounded and equi-continuous in $\alpha$. For the case of the space $L^p(S, \mathfrak{B}, m)$, $1 \leqq p < \infty$, we have 

\begin{theoremn}[\textsc{Fréchet}-\textsc{Kolmogorov}] Let $S$ be the real line, $\mathfrak{B}$ the $\sigma$-ring of Baire subsets $B$ of $S$ and $m(B) = \int_{B} \dd x$ the ordinary Lebesgue measure of $B$. Then a subset $K$ of $L^{p}(S, \mathfrak{B}, m)$, $1 \leqq p < \infty$, is strongly pre-compact iff it satisfies the conditions:
\begin{align}
\sup\limits _{x \in K} \| x \|& = \sup\limits _{x \in K} \left(\int_{S} | x (s) | ^{p} \dd s\right) ^{1/p} < \infty , \label{eq:X.1.1} \\
\lim\limits _{t \to 0}& \int_{S} | x (t + s) - x (s) | ^{p} \dd s = 0 \text{ uniformly in } x \in K, \label{eq:X.1.2} \\
\lim\limits _{\alpha \uparrow \infty}& \int_{| s | > \alpha} | x (s) | ^{p} \dd s = 0 \text{ uniformly in } x \in K. \label{eq:X.1.3}
\end{align}

\end{theoremn}
\begin{proof}
 Let $K$ be strongly relatively compact. Then $K$ is bounded and so \eqref{eq:X.1.1} is true. Let $\varepsilon > 0$ be given. Then there exists a finite number of functions $\in L^p: f_1, f_2, \ldots, f_n$ such that, for each $f \in K$, there is a $j$ with $\|f - f_j\| \leqq \varepsilon$. Otherwise, we would have an infinite sequence $\{f_j\} \subseteq K$ with $\|f_j - f_i\| > \varepsilon$ for $j \neq i$, contrary to the relative compactness of $K$. We then find, by the definition of the Lebesgue integral, finitely-valued functions $g_1, g_2, \ldots, g_n$ such that $\|f_j - g_j\| \leqq \varepsilon$ ( $j = 1, 2, \ldots, n$). Since each finitely-valued function $g_j(x)$ vanishes outside some sufficiently large interval, we have, for large $\alpha$, 
\[
\begin{gathered}
\left(\int_{\alpha} ^{\infty} + \int_{- \infty} ^{- \alpha} | f (s) | ^{p} \dd s\right) ^{1/p} \leqq \left(\int_{\alpha} ^{\infty} + \int_{- \infty} ^{- \alpha} | f (s) - g _{j} (s) | ^{p} \dd s\right) ^{1/p}
\\
+ \left(\int_{\alpha} ^{\infty} + \int_{- \infty} ^{- \alpha} | g _{j} (s) | ^{p} \dd s\right) ^{1/p} \leqq \| f - g _{j} \| + \left(\int_{\alpha} ^{\infty} + \int_{- \infty} ^{- \alpha} | g _{j} (s) | ^{p} \dd s\right) ^{1/p}. \end{gathered}
\]

This proves \eqref{eq:X.1.3} by $\| f - g_j\| \leqq \| f - f_j\| + \| f_j - g_j\| \leqq 2\varepsilon$. 

The proof of \eqref{eq:X.1.2} is based on the fact that, for the defining function $C_I(s)$ of a finite interval $I$, $\lim\limits_{t \to 0} \int_{-\infty}^{\infty} |C_I(s + t) - C_I(s)|^p \dd s = 0$ (see Chapter~\ref{chap:0}, \ref{sec:0.3}). Thus \eqref{eq:X.1.2} holds for finitely-valued functions $g_j(s)$ ( $j = 1, 2, \ldots, n$). Hence we have, for any $f \in K$, 
\[
\begin{aligned}
&\overline {{\lim\limits _{t \to 0}}} \left(\int_{- \infty} ^{\infty} | f (s + t) - f (s) | ^{p} \dd s\right) ^{1/p} \leqq {} \overline {{\lim\limits _{t \to 0}}} \left(\int_{- \infty} ^{\infty} | f (s + t) - f _{j} (s + t) | ^{p} \dd s\right) ^{1/p}
\\
&\quad + \overline {{\lim\limits _{t \to 0}}} \left(\int_{- \infty} ^{\infty} | f _{j} (s + t) - g _{j} (s + t) | ^{p} \dd s\right) ^{1/p} + \overline {{\lim\limits _{t \to 0}}} \left(\int_{- \infty} ^{\infty} | g _{j} (s + t) - g _{j} (s) | ^{p} \dd s\right) ^{1/p}
\\
&\quad + \left(\int_{- \infty} ^{\infty} | g _{j} (s) - f _{j} (s) | ^{p} \dd s\right) ^{1/p} + \left(\int_{- \infty} ^{\infty} | f _{j} (s) - f (s) | ^{p} \dd s\right) ^{1/p} \leqq {} \varepsilon + \varepsilon + 0 + \varepsilon + \varepsilon ,
\end{aligned}
\]
by taking $f_{j}$ in such a way that $\| f - f_{j}\| \leqq \varepsilon$. This proves \eqref{eq:X.1.2}. 
\end{proof}

\begin{theoremn}[We next prove the converse part of the Theorem] We define the translation operator $T_{t}$ by $(T_{t}f)(s) = f(t + s)$. Condition \eqref{eq:X.1.2} says that  
$s\text{- }\lim\limits_{t \to 0} T_t f = f$ uniformly in $f \in K$. We next define the mean value $(M_a f)(s) = (2a)^{-1} \int_{-a}^{a} (T_t f)(s) \dd t$. Then, by Hölder's inequality and the Fubini-Tonelli theorem, 
\[
\begin{aligned}
\left\| M _{a} f - f \right\|& \leqq \left(\int_{- \infty} ^{\infty} \left\{\int_{-a} ^{a} (2 a) ^{-1} | f (s + t) - f (s) | \dd t \right\} ^{p} \dd s\right) ^{1/p}
\\
\quad& \leqq (2 a) ^{-1} \left(\int_{- \infty} ^{\infty} \int_{-a} ^{a} | f (s + t) - f (s) | ^{p} \dd t \cdot (2 a) ^ {p / p ^{\prime}} \dd s\right) ^{1/p}
\\
\quad& \leqq \left((2 a) ^{-1} \int_{-a} ^{a} \dd t \int_{- \infty} ^{\infty} | f (s + t) - f (s) | ^{p} \dd s\right) ^{1/p} \text{ if } 1 \leqq p < \infty .
\end{aligned}
\]
\end{theoremn}

Thus we have 
\[
\left\| M _{a} f - f \right\| \leqq \sup\limits _{| t | \leqq a} \left\| T _{i} f - f \right\|,
\]
so that $s\text{- }\lim\limits_{a \downarrow 0} M_a f = f$ uniformly in $f \in K$. Therefore, we have to prove the relative compactness of the set $\{M_a f; f \in K\}$ for a sufficiently small fixed $a > 0$. 

We will show that, for a fixed $a > 0$, the set of functions $\{(M_{a}f)(s); f \in K\}$ is equi-bounded and equi-continuous. In fact, we have, as above, 
\[
\begin{aligned} \left| (M _{a} f) (s _{1}) - (M _{a} f) (s _{2}) \right| & \leqq (2 a) ^{-1} \int_{-a} ^{a} | f (s _{1} + t) - f (s _{2} + t) | \dd t \\ & \leqq \left((2 a) ^{-1} \int_{-a} ^{a} | f (s _{1} + t) - f (s _{2} + t) | ^{p} \dd t\right) ^{1/p}. \end{aligned}
\]

Thus, by \eqref{eq:X.1.2}, we have proved the equi-continuity of the set of functions $\{(M_{a}f)(s); f \in K\}$ for a fixed $a > 0$. The equi-boundedness of the set may be proved similarly. Thus, by the Ascoli-Arzelà theorem, there exists, for any positive $a > 0$, a finite number of functions $M_{a}f_{1}, M_{a}f_{2}, \ldots, M_{a}f_{n}$ with $f_{j} \in K$ ( $j = 1, 2, \ldots, n$) such that, for any $f \in K$, there exists some $j$ for which $\sup |(M_{a}f)(s) - (M_{a}f_{j})(s)| \leqq \varepsilon$. Therefore 
\begin{equation}\label{eq:X.1.4}
\begin{aligned}
&\left\| M _{a} f - M _{a} f _{j} \right\| ^{p}  \leqq \int_{- \alpha} ^{\alpha} | (M _{a} f) (s) - (M _{a} f _{j}) (s) | ^{p} \dd s + \int_{| s | > \alpha} | (M _{a} f) (s) - (M _{a} f _{j}) (s) | ^{p} \dd s.
\end{aligned}
\end{equation}

The second term on the right is, by Minkowski's inequality, smaller than $\left(\|M_a f - f\| + \left(\int_{|s| > \alpha} |f(s) - f_j(s)|^p \dd s\right)^{1/p} + \left(\int_{|s| > \alpha} |f_j(s) - (M_a f_j)(s)|^p \dd s\right)^{1/p}\right)^p$. 

The term $\|M_{a}f - f\|$ is small for sufficiently small $a > 0$, and, by virtue of \eqref{eq:X.1.3}, $\int_{|s| > \alpha}|f(s) - f_j(s)|^p \dd s$ and $\int_{|s| > \alpha}|f_j(s) - (M_{a}f_j)(s)|^p \dd s$ are both small for sufficiently large $\alpha > 0$, if $a > 0$ is bounded. Also the first term on the right of \eqref{eq:X.1.4} is $\leqq 2\alpha\varepsilon^{p}$ for an appropriate choice of $j$. These estimates are valid uniformly with respect to $f \in K$. Thus we have proved the relative compactness in $L^{p}$ of the set $\{M_{a}f; f \in K\}$ for sufficiently small $a > 0$.\yosstarfootnote{See also Supplementary Notes, p.~\pageref{notes:X}.} 


\yossection{2}{Compact Operators and Nuclear Operators}
\label{sec:X.2}

\begin{definition}\label{def:X.2.Definition.1} Let $X$ and $Y$ be $B$-spaces, and let $S$ be the unit sphere of $X$. An operator $T \in L(X, Y)$ is said to be compact or completely continuous if the image $T \cdot S$ is relatively compact in $Y$. 
\end{definition}

\begin{example}\label{ex:X.2.Example.1} Let $K(x, y)$ be a real- or complex-valued continuous function defined for $-\infty < a \leqq x, y \leqq b < \infty$. Then the integral operator $K$ defined by 
\begin{equation}\label{eq:X.2.1}
(K f) (x) = \int_{a} ^{b} K (x, y) f (y) \dd y
\end{equation}
is compact as an operator $\in L(C[a,b],C[a,b])$. 
\end{example}

\begin{proof}
 Clearly $K$ maps $C[a, b]$ into $C[a, b]$. Set $\sup |K(x, y)| = M$. 

Then $\| K \cdot f \| \leqq (b - a) M \| f \|$ so that $K \cdot S$ is equi-bounded. By Schwarz' inequality, we have 
\[
\left| (K f) \left(x _{1}\right) - (K f) \left(x _{2}\right) \right| ^{2} \leqq \int_{a} ^{b} \left| K \left(x _{1}, y\right) - K \left(x _{2}, y\right) \right| ^{2} \dd y \cdot \int_{a} ^{b} | f (y) | ^{2} \dd y,
\]
and hence $K \cdot S$ is equi-continuous, that is, 
\[
\lim\limits _{\delta \downarrow 0} \sup\limits _ {| x _{1} - x _{2} | \leqq \delta} | (K f) (x _{1}) - (K f) (x _{2}) | = 0 \text{ uniformly in } f \in S.
\]

Therefore, by the Ascoli-Arzelà theorem (Chapter~\ref{chap:III}, \ref{sec:III.3}), the set $K \cdot S$ is relatively compact in $C[a, b]$. 
\end{proof}

\begin{example}\label{ex:X.2.Example.2} Let $K(x, y)$ be a real- or complex-valued $\mathfrak{B}$-measurable function on a measure space $(S, \mathfrak{B}, m)$ such that 
\begin{equation}\label{eq:X.2.2}
\iint_{S} \left| K (x, y) \right| ^{2} m (\dd x) m (\dd y) < \infty .
\end{equation}
\end{example}

Then the integral operator $K$ defined by the kernel $K(x, y)$: 
\begin{equation}\label{eq:X.2.3}
(K f) (x) = \int_{S} K (x, y) f (y) m (\dd y), f \in L ^{2} (S) = L ^{2} (S, \mathfrak {B}, m),
\end{equation}
is compact as an operator $\in L(L^2(S), L^2(S))$. The kernel $K(x, y)$ satisfying \eqref{eq:X.2.2} is said to be of the Hilbert-Schmidt type. 

\begin{proof}
 Take any sequence $\{f_n\}$ from the unit sphere of $L^2(S)$. We have to show that the sequence $\{K \cdot f_n\}$ is relatively compact in $L^2(S)$. Since a Hilbert space $L^2(S)$ is locally sequentially weakly compact, we may assume that $\{f_n\}$ converges weakly to an element $f \in L^2(S)$; otherwise, we choose a suitable subsequence. By \eqref{eq:X.2.2} and the Fubini-Tonelli theorem, $\int_{S} |K(x, y)|^2 m(\dd y) < \infty$ for $m$-a.e. $x$. Hence, for such an $x$, $\lim\limits_{n \to \infty} (Kf_n)(x) = \lim\limits_{n \to \infty} \int_{S} K(x, y) f_n(y) m(\dd y) = \lim\limits_{n \to \infty} (f_n(.), \overline{K(x,.)}) = (f(.), \overline{K(x,.)}) = \int_{S} K(x, y) f(y) m(\dd y)$. On the other hand, we have, by Schwarz' inequality, 
\begin{equation}\label{eq:X.2.4}
\begin{aligned}
&\left| (K f _{n}) (x) \right| ^{2} \leqq \int_{S} | K (x, y) | ^{2} m (\dd y) \cdot \int_{S} \left| f _{n} (y) \right| ^{2} m (\dd y) \leqq \int_{S} | K (x, y) | ^{2} m (\dd y) \text{ for } m\text{-a.e.} x.
\end{aligned}
\end{equation}

Hence, by the Lebesgue-Fatou theorem, $\lim\limits_{n\to \infty}\int_S|(Kf_n)(x)|^2 m(\dd x) = \int_S|(Kf)(x)|^2 m(\dd x)$. This result, if combined with $w - \lim\limits_{n\to \infty}K\cdot f_n = K\cdot f$, implies $\slim\limits_{n\to \infty}K\cdot f_n = K\cdot f$ by Theorem~\ref{thm:V.1.Theorem.8} in Chapter~\ref{chap:V}, \ref{sec:V.1}. But, as proved above in \eqref{eq:X.2.4}, we have 
\[
\int_{S} | (K h) (x) | ^{2} m (\dd x) \leqq \iint_{S} | K (x, y) | ^{2} m (\dd y) m (\dd x) \cdot \int_{S} | h (y) | ^{2} m (\dd y),
\]
and so 
\begin{equation}\label{eq:X.2.5}
\left\| K \right\| \leqq \left(\iint_{S} \left| K (x, y) \right| ^{2} m (\dd x) m (\dd y)\right) ^{1/2}.
\end{equation}

Hence, from \wlim $f_n = f$, we obtain \wlim $K \cdot f_n = K \cdot f$, because, for any $g \in L^2(S)$, $\lim\limits_{n \to \infty} (K \cdot f_n, g) = \lim\limits_{n \to \infty} (f_n, K^*g) = (f, K^*g) = (K \cdot f, g)$. 
\end{proof}

\begin{theoremu}
 \begin{enumerate}[(i)]
\item A linear combination of compact operators is compact.
\item The product of a compact operator with a bounded linear operator is compact; thus the set of compact operators $\in L(X, X)$ constitutes a closed two-sided ideal of the algebra $L(X, X)$ of operators.
\item Let a sequence $\{T_n\}$ of compact operators $\in L(X, Y)$ converge to an operator $T$ in the sense of the uniform operator topology, i.e., $\lim\limits_{n \to \infty} \|T - T_n\| = 0$. Then $T$ is also compact.
\end{enumerate}
\end{theoremu}

\begin{proof}
 (i) and (ii) are clear from the definition of compact operators. The closedness, in the sense of the uniform operator topology, of the ideal of compact operators in the algebra $L(X, X)$ is implied by (iii). 

We shall prove (iii). Let $\{x_h\}$ be a sequence from the closed unit sphere $S$ of $X$. By the compact property of each $T_n$, we can choose, by the diagonal method, a subsequence $\{x_{h'}\}$ such that $s\lim\limits_{h\to \infty}T_n x_{h'}$ exists for every fixed $n$. We have 
\[
\begin{aligned} & {\left\| T x _{h ^{\prime}} - T x _{k ^{\prime}} \right\| \leqq \left\| T x _{h ^{\prime}} - T _{n} x _{h ^{\prime}} \right\| + \left\| T _{n} x _{h ^{\prime}} - T _{n} x _{k ^{\prime}} \right\| + \left\| T _{n} x _{k ^{\prime}} - T x _{k ^{\prime}} \right\|} \\ & {\quad \leqq \| T - T _{n} \| + \left\| T _{n} x _{h ^{\prime}} - T _{n} x _{k ^{\prime}} \right\| + \| T _{n} - T \|,} \end{aligned}
\]
and so $\varlimsup\limits_{h,k\to\infty}\left\|T\cdot x_{h'}-T\cdot x_{k'}\right\|\leqq 2\|T-T_n\|$. Hence $\{Tx_{k'}\}$ is a Cauchy sequence in the $B$-space $Y$. 
\end{proof}

Nuclear Operator. As an application of the Theorem, we shall consider the nuclear operator introduced by A. \textsc{Grothendieck} \cite{grothendieck88}. 

\begin{definition}\label{def:X.2.Definition.2} Let $X, Y$ be $B$-spaces and $T \in L(X, Y)$. If there exist a sequence $\{f_n'\} \subseteq X'$, a sequence $\{y_n\} \subseteq Y$ and a sequence $\{c_n\}$ of numbers such that 
\begin{align}
\sup\limits _{n} \left\| f _{n} ^{\prime} \right\|& < \infty , \sup\limits _{n} \left\| y _{n} \right\| < \infty , \sum\limits_{n} \left| c _{n} \right| < \infty \text{ and } \notag \\
T \cdot x& = \slim _{m \rightarrow \infty} \sum\limits_{n=1} ^{m} c _{n} \left\langle x, f _{n} ^{\prime} \right\rangle y _{n} \text{ in } Y \text{ for every } x \in X, \label{eq:X.2.6}
\end{align}
then $T$ is called a nuclear operator on $X$ into $Y$. 
\end{definition}

\begin{remarku} The existence of the \slim in \eqref{eq:X.2.6} is clear, since 
\[
\left\| \sum\limits_{j=n} ^{m} c _{j} \langle x, f _{j} ^{\prime} \rangle y _{j} \right\| \leqq \sum\limits_{j=n} ^{m} | c _{j} | \cdot \| x \| \cdot \| f _{j} ^{\prime} \| \cdot \| y _{j} \| \leqq \text{ constant } \cdot \sum\limits_{j=n} ^{m} | c _{j} | \cdot \| x \|.
\]
\end{remarku}

The nuclear condition says that the \slim is equal to $T \cdot x$ for every $x \in X$. 

\begin{propositionu} $A$ nuclear operator $T$ is compact. 
\end{propositionu}

\begin{proof}
 Define the operator $T_{n}$ by 
\begin{equation}\label{eq:X.2.7}
T _{n} x = \sum\limits_{j=1} ^{n} c _{j} \langle x, f _{j} ^{\prime} \rangle y _{j}. \end{equation}

Since the range $R(T_{n})$ is of finite dimension, $T_{n}$ is compact as may be proved by the Bolzano-Weierstrass theorem. Moreover, by \eqref{eq:X.2.6} and 
\[
\left\| T x - T _{n} x \right\| = \left\| \sum\limits_{j=n+1} ^{\infty} c _{j} \langle x, f _{j} ^{\prime} \rangle y _{j} \right\| \leqq \text{ constant } \sum\limits_{j=n+1} ^{\infty} \left| c _{j} \right| \cdot \left\| x \right\|,
\]
we have $\lim\limits_{n\to \infty}\left\|T - T_n\right| = 0$ and so $T$ must be compact. 
\end{proof}

An Example of the Nuclear Operator. Let $G$ be a bounded open domain of $\mathbb{R}^{n}$, and consider the Hilbert space $H_{0}^{k}(G)$. Suppose $(k-j)>n$. Then the mapping $T$ 
\begin{equation}\label{eq:X.2.8}
H _{0} ^{k} (G) \ni \varphi \rightarrow \varphi \in H _{0} ^{j} (G)
\end{equation}
is a nuclear operator $\in L(H_0^k (G),H_0^j (G))$ 

\begin{proof}
 We may assume that the bounded domain $G$ is contained in the interior of the parallelogram $P$: 
\[
0 \leqq x _{j} \leqq 2 \pi (j = 1, 2, \dots , n). \]

We recall that $H_0^k (G)$ is the completion of $\hat{H}_0^k (G) = C_0^k (G)$ with respect to the norm $\|\varphi \|_k = \left(\sum\limits_{|s|\leqq k}\int_G |D^s\varphi (x)|^2 \dd x\right)^{1/2}$ (see Chapter~\ref{chap:I}, \ref{sec:I.10}). We 
extend the functions $\in \tilde{H}_0^k (G)$ to be periodic with period $2\pi$ in each variable $x_{s}$ by defining the function values as $0$ in $P - G$. The functions 
% 按正文段落排: 左对齐到版心、首行缩进, 编号排在末行右端。
% 用 \notag 关掉居中编号, 改在段内用 \hfill 把 (\theequation) 推到右边,
% 这样公式体就是整版心宽的段落, 不再是居中排的块。
\begin{equation}\notag
\refstepcounter{equation}\label{eq:X.2.9}%
\parbox[b]{\linewidth}{\setlength{\parindent}{1.2em}\setlength{\parfillskip}{0pt}$f _{\beta} (x) = (2 \pi) ^{-n/2} \exp (i \beta \cdot x)$, where $\beta = (\beta_{1}, \beta_{2}, \dots , \beta_{n})$ is an $n$-tuple of integers and \mbox{$\beta \cdot x = \sum\limits_{s=1}^{n} \beta_s x_s,$}\hfill(\theequation)}
\end{equation}


form a complete orthonormal system of $L^2(P) = H_0^0(P)$. Thus, denoting by $D^s$ the distributional derivative, we have, for $|s| \leqq k$, the Fourier expansion in $L^2(P)$ of functions $D^s\varphi(x)$ where $\varphi \in \hat{H}_0^k(G)$: 
\begin{equation}\label{eq:X.2.10}
D ^{s} \varphi (x) = \sum\limits_{\beta} \left(D ^{s} \varphi , f _{\beta}\right) _{0} f _{\beta}, \text{ where } (\psi , f _{\beta}) _{0} = \int_{P} \psi (x) \overline {{f _{\beta} (x)}} \dd x.
\end{equation}

We have, by 
\[
(D ^{s} \varphi , f _{\beta}) _{0} = (- 1) ^{|s|} (\varphi , D ^{s} f _{\beta}) _{0} = \prod\limits_{m=1} ^{n} (i \beta_{m}) ^{sm} (\varphi , f _{\beta}) _{0}
\]
and the Parseval relation 
\[
\sum\limits_{\beta} \left| (D ^{s} \varphi , f _{\beta}) _{0} \right| ^{2} = \int_{P} \left| D ^{s} \varphi (x) \right| ^{2} \dd x \leqq \left\| \varphi \right\| _{k} ^{2} (| s | \leqq k),
\]
the inequality 
\[
\left| \left(\varphi , (1 + | \beta | ^{2}) ^{k/2} f _{\beta}\right) _{0} \right| ^{2} \leqq \text{ constant } \sum\limits_{| s | \leqq k} \left| \left(D ^{s} \varphi , f _{\beta}\right) _{0} \right| ^{2} \leqq \text{ constant } \| \varphi \| _{k} ^{2}.
\]

Therefore the functional $f_{\beta}^{\prime} \in H_0^k (G)'$ defined by 
\[
\langle \varphi , f _{\beta} ^{\prime} \rangle = (\varphi , (1 + | \beta | ^{2}) ^{k/2} f _{\beta}) _{0}
\]
satisfies $\sup\limits_{\beta}\|f_{\beta}'\| < \infty$. Moreover, 
\[
y _{\beta} = (1 + | \beta | ^{2}) ^{-j/2} f _{\beta}
\]
satisfies $\sup\limits_{\beta}\|y_{\beta}\|_j < \infty$ by $D^{s}f_{\beta} = \prod\limits_{t=1}^{n}(i\beta_{t})^{s_{t}}f_{\beta}$. We also have 
\[
\sum\limits_{\beta} | c _{\beta} | < \infty , \text{ where } c _{\beta} = (1 + | \beta | ^{2}) ^{(j-k)/2},
\]
because, for positive integers $\beta_{s}$, 
\[
\begin{aligned}
\sum\limits_{\beta} \frac {1}{(\beta_{1} + \beta_{2} + \cdots + \beta_{n}) ^{k-j}} &= \sum\limits_{\beta} \left(\frac {1}{(\beta_{1} + \beta_{2} + \cdots + \beta_{n}) ^{n}}\right) ^{(k-j)/n}
\\
\leqq \sum\limits_{\beta} \left(\frac {1}{\beta_{1} \beta_{2} \cdots \beta_{n}}\right) ^{(k-j)/n} &= \sum\limits_ {\beta_{1}} \left(\frac {1}{\beta_{1}}\right) ^{(k-j)/n} \cdot \sum\limits_ {\beta_{2}} \left(\frac {1}{\beta_{2}}\right) ^{(k-j)/n}
\\
\dots \sum\limits_ {\beta_{n}} \left(\frac {1}{\beta_{n}}\right) ^{(k-j)/n}& < \infty \text{ by } \frac {(k - j)}{n} > 1.
\end{aligned}
\]

Therefore, we have proved the (Fourier) expansion 
\[
\varphi = \sum\limits_{\beta} c _{\beta} \left< \varphi , f _{\beta} ^{\prime} \right> y _{\beta}.
\]
\end{proof}

\begin{remarku} If there exists a complete orthonormal system $\{\varphi_{j}\}$ of eigenfunctions of a given bounded linear operator $K \in L(L^{2}(S), L^{2}(S))$ such that $K\varphi_{j} = \lambda_{j}\varphi_{j} (j = 1, 2, \ldots)$, then, from the Fourier expansion 
\[
f = \sum\limits_{j=1} ^{\infty} (f, \varphi_{j}) \varphi_{j}, f \in L ^{2} (S),
\]
\end{remarku}

we obtain 
\[
K f = \sum\limits_{j=1} ^{\infty} \lambda_{j} (f, \varphi_{j}) \varphi_{j}.
\]

We have $\lambda_{j} = (K\varphi_{j},\varphi_{j})$, and so, if the eigenvalues $\lambda_{j}$ are all $>0$ and $\sum\limits_{j=1}^{\infty}\lambda_{j}<\infty$, then the operator $K$ is nuclear. The condition $\sum\limits_{j=1}^{\infty}|(K\varphi_{j},\varphi_{j})|<\infty$ is surely satisfied, if the operator $K$ is defined by a kernel 
\[
K(x,y)=\int_S\overline{K_2(z,x)}K_1(z,y)m(\dd z),
\]
where the kernels $K_1(x,y)$ and $K_2(x,y)$ are of the Hilbert-Schmidt type.
For,
\[
\sum\limits_{j=1} ^{\infty} \left| \left(K _{2} ^{*} K _{1} \varphi_{j}, \varphi_{j}\right) \right| = \sum\limits_{j=1} ^{\infty} \left| \left(K _{1} \varphi_{j}, K _{2} \varphi_{j}\right) \right| \leqq \left(\sum\limits_{j=1} ^{\infty} \left\| K _{1} \varphi_{j} \right\| ^{2} \cdot \sum\limits_{j=1} ^{\infty} \left\| K _{2} \varphi_{j} \right\| ^{2}\right) ^{1/2}
\]
and, by Parseval's relation, we have 
\[
\begin{aligned}
&\sum\limits_{j=1} ^{\infty} \| K _{1} \varphi_{j} \| ^{2}  = \sum\limits_{j=1} ^{\infty} \int_{S} \left| \int_{S} K _{1} (z, y) \varphi_{j} (y) m (\dd y) \right| ^{2} m (\dd z) \\
&= \int_{S} \sum\limits_{j=1} ^{\infty} \left| \int_{S} K _{1} (z, y) \varphi_{j} (y) m (\dd y) \right| ^{2} m (\dd z) = \int_{S} \left\{\int_{S} | K _{1} (z, y) | ^{2} m (\dd y) \right\} m (\dd z) < \infty ,
\end{aligned}
\]
and similar for $\sum\limits_{j=1}^{\infty} \|K_2\varphi_j\|^2$. A bounded linear operator $K$ in a separable Hilbert space $X$ is said to be of the trace class if $\sum\limits_{j=1}^{\infty} |(K\varphi_j, \psi_j)| < \infty$ for arbitrary complete orthonormal systems $\{\varphi_j\}$ and $\{\psi_j\}$ of $X$. For a general account concerning the trace class and the nuclear operator, see R. \textsc{Schatten} \cite{schatten288}, and I. M. \textsc{Gelfand}-N. Y. \textsc{Vilenkin} \cite{gelfand77}. 


\yossection{3}{The Rellich-Gårding Theorem}
\label{sec:X.3}

\begin{theoremn}[\textsc{Garding} {\cite{garding74}}] Let $G$ be a bounded open domain of $\mathbb{R}^{n}$. If an operator $T \in L(H_0^k(G), H_0^k(G))$ satisfies, for $j < k$,
\begin{equation}\label{eq:X.3.1}
\|T\varphi\|_{k} \leqq C \|\varphi\|_{j} \quad \text{for all } \varphi \in H_{0}^{k}(G) \text{, where } C \text{ is a constant,}
\end{equation}
then $T$ is compact as an operator $\in L(H_{0}^{k}(G), H_{0}^{k}(G))$. 

\end{theoremn}
\begin{proof}
 By the definition of the space $H_0^k(G)$ (see Chapter~\ref{chap:I}, \ref{sec:I.10}), it would be sufficient to show the following: Let a sequence $\{\varphi_\nu\} \subseteq \hat{H}_0^k(G) = C_0^k(G)$ be such that $\|\varphi_\nu\|_k \leqq 1 (\nu = 1,2,\ldots)$. Then the sequence $\{T\varphi_\nu\}$ contains a subsequence strongly convergent in $H_0^k(G)$. The Fourier transform $\hat{\varphi}_\nu(\xi) = (2\pi)^{-n/2} \int_G \varphi_\nu(x) \exp(-ix\xi) \dd x$ satisfies, by Schwarz' inequality, 
\[
\left| \hat {\varphi} _{\nu} (\xi) \right| ^{2} \leqq (2 \pi) ^{-n} \int_{G} \dd x \int_{G} \left| \varphi_{\nu} (x) \right| ^{2} \dd x \leqq (2 \pi) ^{-n} \int_{G} \dd x,
\]
and hence $\{\hat{\varphi}_{\nu}(\xi)\}$ is equi-bounded in $\xi\in \mathbb{R}^{n}$ and in $\pmb{\nu}$. We may assume, by the boundedness of $\|\varphi_{\nu}\|_{0}$, that a subsequence $\{\varphi_{\nu}\}$ is weakly convergent in $L^2(G) = H_0^0(G)$. Since, for each $\xi$, the function $\exp(-ix\xi)$ belongs to $L^2(G)$, we know that the sequence of bounded functions $\hat{\varphi}_{\nu'}(\xi) = (\varphi_{\nu'}, (2\pi)^{-n/2} \exp(-ix\xi))_0$ converges at every $\xi$. Thus, by \eqref{eq:X.3.1} and the Parseval relation for the Fourier transform (Chapter~\ref{chap:VI}, \ref{sec:VI.2}), 
\[
\begin{aligned}
&\left\| T \varphi_ {\nu^{\prime}} - T \varphi_ {\mu^{\prime}} \right\| _{k} ^{2} = \left\| T (\varphi_ {\nu^{\prime}} - \varphi_ {\mu^{\prime}}) \right\| _{k} ^{2} \leqq C ^{2} \left\| \varphi_ {\nu^{\prime}} - \varphi_ {\mu^{\prime}} \right\| _{j} ^{2}
\\
&= C ^{2} \sum\limits_{| s | \leqq j} \left\| D ^{s} (\varphi_ {\nu^{\prime}} - \varphi_ {\mu^{\prime}}) \right\| _{0} ^{2} = C ^{2} \sum\limits_{| s | \leqq j} \left\| (D ^{s} (\varphi_ {\nu^{\prime}} - \varphi_ {\mu^{\prime}})) \right\| _{0} ^{2}
\\
&= C ^{2} \sum\limits_{| s | \leqq j} \left\| \prod\limits_{t=1} ^{n} (i \xi_{t}) ^{st} (\hat {\varphi} _ {\nu^{\prime}} - \hat {\varphi} _ {\mu^{\prime}}) (\xi) \right\| _{0} ^{2} \leqq C ^{2} \sum\limits_{| s | \leqq j} \int_{| \xi | \leqq r} \left| \prod\limits_{t=1} ^{n} \xi_{t} ^{st} (\hat {\varphi} _ {\nu^{\prime}} (\xi) - \hat {\varphi} _ {\mu^{\prime}} (\xi)) \right| ^{2} \dd \xi
\\
&\quad + C ^{2} C _{1} \int_{| \xi | > r} | \xi | ^{2j} | \hat {\varphi} _ {\nu^{\prime}} (\xi) - \hat {\varphi} _ {\mu^{\prime}} (\xi) | ^{2} \dd \xi ,
\end{aligned}
\]
where $C_1$ is a positive constant. 

The first term on the right converges, for fixed $r$, to $0$ as $v'$ and $\mu' \to \infty$. This we see by the Lebesgue-Fatou Lemma. The second term on the right is, for $r > 1$. 
\[
\begin{aligned}
&\leqq C ^{2} C _{1} r ^{2j-2k} \int_{| \xi | > r} | \xi | ^{2k} | \hat {\varphi} _ {\nu^{\prime}} (\xi) - \hat {\varphi} _ {\mu^{\prime}} (\xi) | ^{2} \dd \xi \leqq C ^{2} C _{1} r ^{2j-2k} \int_ {\mathbb{R}^{n}} | \xi | ^{2k} | \hat {\varphi} _ {\nu^{\prime}} (\xi) - \hat {\varphi} _ {\mu^{\prime}} (\xi) | ^{2} \dd \xi
\\
&\leqq C ^{2} C _{1} C _{2} r ^{2j-2k} \sum\limits_{| s | \leqq k} \| (D ^{s} \widehat {\varphi_ {\nu^{\prime}} - D ^{s}} \varphi_ {\mu^{\prime}}) \| _{0} ^{2} = C ^{2} C _{1} C _{2} r ^{2j-2k} \sum\limits_{| s | \leqq k} \| D ^{s} (\varphi_ {\nu^{\prime}} - \varphi_ {\mu^{\prime}}) \| _{0} ^{2}
\\
&= C ^{2} C _{1} C _{2} r ^{2j-2k} \| \varphi_ {\nu^{\prime}} - \varphi_ {\mu^{\prime}} \| _{k} ^{2} \leqq 4 C ^{2} C _{1} C _{2} r ^{2j-2k}
\end{aligned}
\]
with a constant $C_2$. 

The last term converges, by $j < k$, to $0$ as $r \to \infty$. Therefore, $\lim\limits_{\nu, \mu \to \infty} \| T \varphi_{\nu'} - T \varphi_{\mu'} \|_k = 0$. 
\end{proof}


\yossection{4}{Schauder's Theorem}
\label{sec:X.4}

\begin{theoremn}[Schauder] An operator $T \in L(X, Y)$ is compact iff its dual operator $T'$ is compact.

\end{theoremn}
\begin{proof}
 Let $S, S'$ be the closed unit sphere in $X, Y'$, respectively. Let $T \in L(X, Y)$ be compact. Let $\{y_j'\}$ be an arbitrary sequence in $S'$. The functions $F_j(y) = \langle y, y_j' \rangle$ are equi-continuous in the sense that 
\[
\left| F _{j} (y) - F _{j} (z) \right| = \left| \langle y - z, y _{j} ^{\prime} \rangle \right| \leqq \left\| y - z \right\|. \]

Moreover, $\{F_j(y)\}$ is equi-bounded in $j$ on any bounded set of $y$, since $|F_j(y)| \leqq \|y\|$. Therefore, by the Ascoli-Arzelà theorem, as applied to the functions $\{F_j(y)\}$ defined on the compact set $(T\cdot S)^a$, we see that some subsequence $\{F_{j'}(y)\}$ converges uniformly in $y\in (T\cdot S)^a$. Hence $\langle Tx,y_j'\rangle = \langle x,T'y_j'\rangle$ converges uniformly in $x\in S$, and so $\{T'\cdot y_j'\}$ converges in the strong topology of $X'$. This proves that $T'$ is compact. 

Conversely, let $T'$ be compact. Then, by what we have proved above, $T''$ is compact. Hence, if $S''$ is the closed unit sphere in $X''$, $(T'' \cdot S'')$ is relatively compact. We know that $Y$ is isometrically embedded in $Y''$ (see Theorem~\ref{thm:IV.8.Theorem.2} in Chapter~\ref{chap:IV}, \ref{sec:IV.8}). Hence $T \cdot S \subseteq T'' \cdot S''$ and so $T \cdot S$ is relatively compact in the strong topology of $Y''$ and so in the strong topology of $Y$. Therefore, $T$ is compact. 
\end{proof}


\yossection{5}{The Riesz-Schauder Theory}
\label{sec:X.5}

\begin{lemman}[F. \textsc{Riesz} \cite{rickart273}] Let $V$ be a compact operator $\in L(X, X)$, where $X$ is a $B$-space. Then, for any complex number $\lambda_{0} \neq 0$, the range $R(\lambda_{0}I - V)$ is strongly closed.

\end{lemman}
\begin{proof}
 We may assume that $\lambda_{0}=1$. Let $\{x_{n}\}$ be a sequence of $X$ such that $y_{n}=(I-V)x_{n}$ converges strongly to $y$. If $\{x_{n}\}$ is bounded, then by, the compactness of the operator $V$, there exists a subsequence $\{x_{n'}\}$ such that $\{Vx_{n}\}$ converges strongly. Since $x_{n'}=y_{n'}+Vx_{n'},\{x_{n'}\}$ converges to some $x$, and so $y=(I-V)x$. 

We next assume that $\{\| x_n \| \}$ is unbounded. Set $T = (I - V)$ and put $\alpha_n = \text{ dis }(x_n, N(T))$, where $N(T) = \{x; Tx = 0\}$. Take a $w_n \in N(T)$ such that $\alpha_n \leqq \| x_n - w_n \| \leqq (1 + n^{-1}) \alpha_n$. Then $T(x_n - w_n) = Tx_n$, and so in the case when $\{\alpha_n\}$ is bounded, we can prove, as above, that $y \in R(T) = R(I - V)$. Suppose that $\lim\limits_{n \to \infty} \alpha_n = \infty$. Since $z_n = (x_n - w_n) / \| x_n - w_n \| $ satisfies $ \| z_n \| = 1$ and $\slim\limits_{n \to \infty} T z_n = 0$, we can prove, as above, that there exists a subsequence $\{z_{n'}\}$ such that $\slim\limits_{n \to \infty} z_{n'} = w_0$, $\slim\limits_{n \to \infty} T z_{n'} = 0$. Hence $w_0 \in N(T)$. But, if we put $z_n - w_0 = u_n$, then, in 
\[
x _{n} - w _{n} - w _{0} \left\| x _{n} - w _{n} \right\| = u _{n} \left\| x _{n} - w _{n} \right\|,
\]
the second and the third terms on the left belong to $N(T)$ so that we must have $\| u_n \| \cdot \| x_n - w_n \| \geqq \alpha_n$. This is a contradiction, since $s \text{- } \lim\limits_{n \to \infty} u_n = 0$, $\| x_n - w_n \| \leqq (1 + n^{-1}) \alpha_n$ and $\lim\limits_{n \to \infty} \alpha_n = \infty$. 

We are now able to prove the Riesz-Schauder theory; for convenience sake, we shall state the theory in a series of three theorems. 
\end{proof}
\begin{theorem}\label{thm:X.5.Theorem.1} Let $V$ be a compact operator $\in L(X, X)$. If $\lambda_0 \neq 0$ is not an eigenvalue of $V$, then $\lambda_0$ is in the resolvent set of $V$. 
\end{theorem}

\begin{proof}
 By the preceding Lemma and the hypothesis, the operator $T_{\lambda_{0}} = (\lambda_{0}I - V)$ gives a one-one map of $X$ onto the set $R(T_{\lambda_{0}})$ which is strongly closed in $X$. Hence, by the Corollary of the open mapping theorem in Chapter~\ref{chap:II}, \ref{sec:II.5}, $T_{\lambda_{0}}$ has a continuous inverse. We have to show that 

$R(T_{\lambda_0}) = X$. If not, the topological image $X_1 = T_{\lambda_0}X$ of $X$ is a proper closed subspace of $X$. Hence, if we set $X_2 = T_{\lambda_0}X_1$, $X_3 = T_{\lambda_0}X_2, \ldots$, then $X_{n+1}$ is a proper closed subspace of $X_n$ ( $n = 0, 1, 2, \ldots; X_0 = X$). By F. \textsc{Riesz}' theorem in Chapter~\ref{chap:III}, \ref{sec:III.2}, there exists a sequence $\{y_n\}$ such that $y_n \in X_n$, $\|y_n\| = 1$ and $\text{ dis }(y_n, X_{n+1}) \geqq 1/2$. Thus, if $n > m$, 
\[
\lambda_{0} ^{-1} \left(V y _{m} - V y _{n}\right) = y _{m} + \left\{- y _{n} - \left(T _ {\lambda_{0}} y _{m} - T _ {\lambda_{0}} y _{n}\right) / \lambda_{0} \right\} = y _{m} - y
\]
with some $y \in X_{m+1}$. 

Hence $\|Vy_n - Vy_m\| \geqq |\lambda_0| / 2$, contrary to the compactness of the operator $V$. 
\end{proof}

\begin{theorem}\label{thm:X.5.Theorem.2} Let $V$ be a compact operator $\in L(X, X)$. Then, \begin{enumerate}[(i)]
\item its spectrum consists of an at most countable set of points of the complex plane which has no point of accumulation except possibly $\lambda = 0$;
\item every non-zero number in the spectrum of $V$ is an eigenvalue of $V$ of finite multiplicity;
\item a non-zero number is an eigenvalue of $V$ iff it is an eigenvalue of $V'$.
\end{enumerate}
\end{theorem}

\begin{proof}
 By Theorem~\ref{thm:X.5.Theorem.1}, a non-zero number in the spectrum of $V$ is an eigenvalue of $V$. The same is also true of $V'$, since, by Schauder's theorem, $V'$ is compact when $V$ is. But the resolvent sets are the same for $V$ and $V'$ (see Chapter~\ref{chap:VIII}, \ref{sec:VIII.6}). Hence (iii) is proved. Since the eigenvectors belonging respectively to different eigenvalues of $V$ are linearly independent, the proofs of:
 (i) and (ii) are completed if we derive a contradiction from the following situation:
There exists a sequence $\{x _{n} \}$ of linearly independent vectors such that $V x _{n} = \lambda_{n} x (n = 1, 2, \dots)$ and $\lim\limits _{n \to \infty} \lambda_{n} = \lambda \neq 0.$


To derive a contradiction, we consider the closed subspace $X_{n}$ spanned by $x_{1}, x_{2}, \ldots, x_{n}$. By F. \textsc{Riesz}' theorem in Chapter~\ref{chap:III}, \ref{sec:III.2}, there exists a sequence $\{y_{n}\}$ such that $y_{n} \in X_{n}$, $\| y_{n} \| = 1$ and $\text{ dis }(y_{n}, X_{n-1}) \geqq 1/2$ ( $n = 2, 3, \ldots$). If $n > m$, then 
\[
\begin{aligned}
&\lambda_{n} ^{-1} V y _{n} - \lambda_{m} ^{-1} V y _{m} = y _{n} + (- y _{m} - \lambda_{n} ^{-1} T _ {\lambda_{n}} y _{n} + \lambda_{m} ^{-1} T _ {\lambda_{m}} y _{m}) = y _{n} - z, \text{ where } z \in X _{n-1}.
\end{aligned}
\]

For, if $y_{n} = \sum\limits_{j=1}^{n} \beta_{j} x_{j}$, then we have $y_{n} - \lambda_{n}^{-1} V y_{n} = \sum\limits_{j=1}^{n} \beta_{j} x_{j} - \sum\limits_{j=1}^{n} \beta_{j} \lambda_{n}^{-1} \lambda_{j} x_{j} \in X_{n-1}$ and similarly $T_{\lambda_m} y_m \in X_m$. Therefore $\|\lambda_n^{-1} V y_n - \lambda_m^{-1} V y_m\| \geqq 1/2$. This contradicts the compactness of $V$ combined with the hypothesis $\lim\limits_{n \to \infty} \lambda_n \neq 0$. 
\end{proof}
\begin{theorem}\label{thm:X.5.Theorem.3} Let $\lambda_0 \neq 0$ be an eigenvalue of a compact operator $V \in L(X, X)$. Then $\lambda_0$ is also an eigenvalue of $V'$ by the preceding theorem. We can prove: \begin{enumerate}[(i)]
\item the multiplicities for the eigenvalue $\lambda_0$ are the same for $V$ and $V'$.
\item The equation $(\lambda_0 I - V)x = y$ admits a solution $x$ iff $y \in N(\lambda_0 I' - V')^\perp$, that is, iff $V'f = \lambda_0 f$ implies $\langle y, f \rangle = 0$.
\item The equation $(\lambda_0 I' - V'f) = g$ admits a solution $f$ iff $g \in N(\lambda_0 I - V)^\perp$, that is, iff $Vx = \lambda_0 x$ implies $\langle x, g \rangle = 0$.
\end{enumerate}
\end{theorem}

\begin{proof}
 Since the eigenvalue $\lambda_{0} \neq 0$ is an isolated singularity of the resolvent $R(\lambda; V) = (\lambda I - V)^{-1}$, we can expand $R(\lambda; V)$ in Laurent series 
\[
R (\lambda ; V) = \sum\limits_{n = - \infty} ^{\infty} (\lambda - \lambda_{0}) ^{n} A _{n}. \]

We are particularly interested in the residue $A_{-1} = (2\pi i)^{-1} \int_{|\lambda - \lambda_0| = \varepsilon} R(\lambda; V) \dd\lambda$. As was proved in Chapter~\ref{chap:VIII}, \ref{sec:VIII.8}, $A_{-1}$ is an idempotent, i.e., $A_{-1}^2 = A_{-1}$. If we set $(\lambda I - V)^{-1} = \lambda^{-1}I + V_\lambda$, then from $(\lambda I - V)(\lambda^{-1}I + V_\lambda) = I$ we obtain $V_\lambda = V(\lambda^{-1}V_\lambda + \lambda^{-2}I)$, and so $V_\lambda$ is compact when $V$ is. Hence, by 
\[
\begin{aligned}
A _{-1} &= (2 \pi i) ^{-1} \int_ {| \lambda - \lambda_{0} | = \varepsilon} R (\lambda ; V) \dd \lambda = (2 \pi i) ^{-1} \int_ {| \lambda - \lambda_{0} | = \varepsilon} \lambda^{-1} \dd \lambda \cdot I
\\
+ (2 \pi i) ^{-1} \int_ {| \lambda - \lambda_{0} | = \varepsilon} V _{\lambda} \dd \lambda &= (2 \pi i) ^{-1} \int_ {| \lambda - \lambda_{0} | = \varepsilon} V _{\lambda} \dd \lambda .
\end{aligned}
\]

Thus, by the Theorem in Chapter~\ref{chap:X}, \ref{sec:X.2}, $A_{-1}$ is a compact operator. 

Therefore, by $A_{-1}X = A_{-1}(A_{-1}X)$ and the compactness of $A_{-1}$, the unit sphere of the normed linear space $A_{-1}X$ is relatively compact. Hence, by F. \textsc{Riesz}' theorem in Chapter~\ref{chap:III}, \ref{sec:III.2}, the range $R(A_{-1})$ is of finite dimension. On the other hand, $Vx = \lambda_0x$, $x \neq 0$, implies that $(\lambda I - V)^{-1}x = (\lambda - \lambda_0)^{-1}x$ by $(\lambda I - V)x = (\lambda - \lambda_0)x$, and so $A_{-1}x = (2\pi i)^{-1}\int_{|\lambda - \lambda_0| = \varepsilon}(\lambda - \lambda_0)^{-1}\dd\lambda \cdot x = x$. Therefore, the eigenvalue equation $Vx = \lambda_0x$ is equivalent to $Vx = \lambda_0x$, $x \in R(A_{-1})$. In the same way, we can prove that the eigenvalue equation $V'f = \lambda_0f$ is equivalent to $V'f = \lambda_0f$, $f \in R(A_{-1}')$. But $R(A_{-1})$ and $R(A_{-1}')$ are of the same dimension. For, $A_{-1}'f = g$ satisfies $A_{-1}'g = A_{-1}'(A_{-1}'f) = g$, and this is equivalent to $\langle x, g \rangle = \langle A_{-1}x, g \rangle$ for all $x \in X$ and so the functional $g$ may be considered as a functional defined on the finite dimensional space $R(A_{-1})$. 

Now, by the well-known theorem in matrix theory, the eigenvalue equation $Vx = \lambda_{0}x$ (in $R(A_{-1})$) and its transposed equation $V'f = \lambda_{0}f$ (in $R(A_{-1}')$) both have the same number of linearly independent solutions. We have thus proved:
 (i). The propositions (ii) and (iii) are already proved by the Lemma and the closed range theorem (Chapter~\ref{chap:VII}, \ref{sec:VII.5}).
\end{proof}

Extension of the Riesz-Schauder Theory. Let a power $V^{n}$ of $V \in L(X, X)$ be compact for some positive integer $n$. Then, by the spectral mapping theorem in Chapter~\ref{chap:VIII}, \ref{sec:VIII.7}, $\sigma(V^{n}) = \sigma(V)^{n}$ and, by the compactness of $V^{n}$, $\sigma(V^{n})$ is either a finite set or else a countable set accumulating only at $0$. Therefore, $\sigma(V)$ is either a finite set or a countable set accumulating 
only at $0$. $V^n$ being compact, 
\[
(2 \pi i) ^{-1} \int_ {| \lambda - \lambda_{0} | = \varepsilon} R (\lambda ; V ^{n}) \dd \lambda
\]
is, for any $\lambda_0 \neq 0$ of $\sigma(V^n)$ and sufficiently small $\varepsilon > 0$, of finite-dimensional range. Hence $\lambda_0$ is a pole of $R(\lambda; V^n)$ (see Chapter~\ref{chap:VIII}, \ref{sec:VIII.8}). But $(\lambda^n I - V^n) = (\lambda - V)(\lambda^{n-1}I + \lambda^{n-2}V + \cdots + V^{n-1})$ and hence 
\[
(\lambda^{n} I - V ^{n}) ^{-1} (\lambda^{n-1} I + \dots + V ^{n-1}) = (\lambda I - V) ^{-1},
\]
which proves that any $\lambda_0 \neq 0$ of $\sigma(V^n)$ is a pole of $R(\lambda; V)$ and so is an eigenvalue of $V$. These facts enable us to extend the Riesz-Schauder theory to operators $V$ for which some power $V^n$ is compact. This extension is highly important in view of its application to concrete problems of integral equations, such as the Dirichlet problem pertaining to potentials. See, e.g., O. D. \textsc{Kellogg} \cite{kellogg148}. It can be proved that the Riesz-Schauder theory for $\lambda_0 = 1$ is valid also for an operator $V \in L(X, X)$ if there exist a positive integer $m$ and a compact operator $K \in L(X, X)$ such that $\|K - V^m\| < 1$. See K. \textsc{Yosida} \cite{yosida352}. It is to be noted here that, if $K_1(s, t)$ and $K_2(s, t)$ are bounded measurable for $0 \leqq s$, $t \leqq 1$, then the integral operator $T$ defined by 
\[
x (s) \rightarrow (T x) (s) = \left(K _{1} K _{2} x\right) (s), \text{ where } \left(K _{j} x\right) (s) = \int_{0} ^{1} K _{j} (s, t) x (t) \dd t,
\]
is compact as an operator $L(L^{1}(0,1),L^{1}(0,1))$. See K. \textsc{Yosida}-Y. \textsc{Mimura}-S. \textsc{Kakutani} \cite{yosida353}. 


\yossection{6}{Dirichlet's Problem}
\label{sec:X.6}

Let $G$ be a bounded open domain of $\mathbb{R}^{n}$, and 
\[
L = \sum\limits_{| s |, | t | \leqq m} D ^{s} c _{st} (x) D ^{t}
\]
a strongly elliptic differential operator with real $C^\infty(G^a)$ coefficients $c_{st}(x) = c_{ts}(x)$. We shall deal only with real-valued functions. Let $f \in L^2(G)$ and $u_1 \in H^m(G)$ be given. Consider a distribution solution $u_0 \in L^2(G)$ of 
\begin{equation}\label{eq:X.6.1}
L u = f \text{ such that } (u _{0} - u _{1}) \in H _{0} ^{m} (G).
\end{equation}

The latter condition $(u_0 - u_1) \in H_0^m(G)$ means that each of the distributional derivatives 
\begin{equation}\label{eq:X.6.2}
\left(D ^{j} u _{0} - D ^{j} u _{1}\right) \text{ for } | j | \leqq m
\end{equation}
is the $L^{2}(G)$-limit of a sequence $\{D^{j}\varphi_{h,j}\}$, where $\varphi_{h,j} \in C_{0}^{\infty}(G)$ (see Chapter~\ref{chap:I}, \ref{sec:I.10}). Thus it gives roughly the boundary conditions: 
\begin{equation}\label{eq:X.6.3}
D ^{j} u _{0} = D ^{j} u _{1} \text{ on the boundary } \partial G \text{ of } G \text{ for } | j | < m.
\end{equation}

In such a sense, \eqref{eq:X.6.1} will be called a Dirichlet's problem for the operator 

$L$. We shall follow the treatment of the problem as formulated and solved by L. \textsc{Garding} \cite{garding74}. 

We first solve 
\begin{equation}\label{eq:X.6.4}
u + \alpha L u = f, (u - u _{1}) \in H _{0} ^{m} (G),
\end{equation}
where the positive constant $\alpha$ is so chosen that Gårding's inequality 
\begin{equation}\label{eq:X.6.5}
(\varphi + \alpha L ^{*} \varphi , \varphi) _{0} \geqq \delta \| \varphi \| _{m} ^{2} \text{ holds whenever } \varphi \in C _{0} ^{\infty} (G)
\end{equation}

Here $L^{*} = \sum\limits_{|s|, |t| \leqq m} (-1)^{|s|+|t|} D^{t} c_{st}(x) D^{s}$ and $\delta$ is a positive constant. The existence of such an $\alpha$ is guaranteed if the coefficients $c_{st}(x)$ are continuous on the closure $G^{a}$ of $G$. We also have, by $m$-times partial differentiation, the inequality 
\begin{equation}\label{eq:X.6.6}
\left| \left(\varphi + \alpha L ^{*} \varphi , \psi\right) _{0} \right| \leqq \gamma \left\| \varphi \right\| _{m} \cdot \left\| \psi \right\| _{m} \text{ whenever } \varphi , \psi \in C _{0} ^{\infty} (G),
\end{equation}
where $\gamma$ is another positive constant independent of $\varphi$ and $\psi$. 

We have, for $u_{1} \in H^{m}(G)$ and $\varphi \in C_0^\infty(G)$, 
\[
(L ^{*} \varphi , u _{1}) _{0} = \sum\limits_{s,t} ((- 1) ^{|s|+|t|} D ^{t} c _{st} D ^{s} \varphi , u _{1}) _{0} = \sum\limits_{s,t} (- 1) ^{|s|} (c _{st} D ^{s} \varphi , D ^{t} u _{1}) _{0},
\]
by partial differentiation. By Schwarz' inequality, we obtain, remembering that the coefficients $c_{st}$ are bounded on $G^{a}$, 
\[
\left| \left(L ^{*} \varphi , u _{1}\right) _{0} \right| \leqq \eta \sum\limits_{| s |, | t | \leqq m} \left\| D ^{s} \varphi \right\| _{0} \left\| D ^{t} u _{1} \right\| _{0} \quad \left(\sup\limits _{s,t;x} \left| c _{st} (x) \right| = \eta\right).
\]

The right hand side is smaller than constant times $\|\varphi\|_m$. 

Thus the linear functional 
\[
F (\varphi) = \left(\varphi + \alpha L ^{*} \varphi , u _{1}\right) _{0}, \varphi \in C _{0} ^{\infty} (G),
\]
can be extended to a bounded linear functional defined on $H_{0}^{m}(G)$ which is the completion of $C_{0}^{\infty}(G)$ with respect to the norm $\left\|\varphi\right\|_{m}$. Similarly, we see, from 
\[
\left| (\varphi , f) _{0} \right| \leqq \left\| \varphi \right\| _{0} \cdot \left\| f \right\| _{0} \leqq \left\| \varphi \right\| _{m} \cdot \left\| f \right\| _{0},
\]
that the linear functional $(\varphi, f)_0$ of $\varphi \in C_0^\infty(G)$ can be extended to a bounded linear functional of $\varphi \in H_0^m(G)$. Hence, by F. \textsc{Riesz}' representation theorem as applied to the Hilbert space $H_0^m(G)$, there exists an $f' = f'(f, u_1) \in H_0^m(G)$ such that 
\[
(\varphi , f) _{0} - (\varphi + \alpha L ^{*} \varphi , u _{1}) _{0} = (\varphi , f ^{\prime}) _{m} \text{ whenever } \varphi \in C _{0} ^{\infty} (G).
\]

Hence, by the Milgram-Lax theorem in Chapter~\ref{chap:III}, \ref{sec:III.7}, applied to the Hilbert space $H_0^m(G)$, we have 
\begin{equation}\label{eq:X.6.7}
(\varphi , f) _{0} - (\varphi + \alpha L ^{*} \varphi , u _{1}) _{0} = (\varphi , f ^{\prime}) _{m} = B (\varphi , S f ^{\prime}), S f ^{\prime} \in H _{0} ^{m} (G),
\end{equation}
where 
\begin{equation}\label{eq:X.6.8}
B (\varphi , \psi) = (\varphi + \alpha L ^{*} \varphi , \psi) _{0} \quad \text{ for } \quad \varphi \in C _{0} ^{\infty} (G), \psi \in H _{0} ^{m} (G).
\end{equation}

Thus 
\[
(\varphi , f) _{0} = (\varphi + \alpha L ^{*} \varphi , u _{1} + S f ^{\prime}) _{0} \quad \text{ whenever } \quad \varphi \in C _{0} ^{\infty} (G),
\]
and so $u_0 = u_1 + Sf'$ is the desired solution $\in L^2(G)$ of \eqref{eq:X.6.4}. 

We shall next discuss the original equation \eqref{eq:X.6.1}. If $u_0 \in L^2(G)$ satisfies \eqref{eq:X.6.1}, then $u_2 = u_0 - u_1 \in H_0^m(G)$ satisfies 
\[
(u _{0}, L ^{*} \varphi) _{0} = (u _{1}, L ^{*} \varphi) _{0} + (u _{2}, L ^{*} \varphi) _{0} = (f, \varphi) _{0}, \varphi \in C _{0} ^{\infty} (G).
\]

We obtain, by partial integration as above, 
$\left|(u_1, L^*\varphi)_0\right| \leqq a$ constant times $\left\|\varphi\right\|_m$, 
\[
\left| (f, \varphi) _{0} \right| \leqq \left\| f \| _{0} \right. \left\| \varphi \| _{0} \right. \leqq \left\| f \| _{0} \cdot \right. \left\| \varphi \| _{m} \cdot \right.
\]

We may thus apply F. \textsc{Riesz}' representation theorem in $H_0^m(G)$ to the linear functional $(f, \varphi)_0 - (u_1, L^*\varphi)_0$ of $\varphi \in C_0^\infty(G)$. Hence there exists a uniquely determined $v \in H_0^m(G)$ such that 
\[
(f, \varphi) _{0} - (u _{1}, L ^{*} \varphi) _{0} = (v, \varphi) _{m} \quad \text{ whenever } \quad \varphi \in C _{0} ^{\infty} (G).
\]

By the Milgram-Lax theorem, applied to $(v,\varphi)_m$, we obtain an $S_1v\in H_0^m (G)$ such that 
\[
(v, \varphi) _{m} = B (S _{1} v, \varphi) \quad \text{ whenever } \quad \varphi \in C _{0} ^{\infty} (G), v \in H _{0} ^{m} (G).
\]

Thus the Dirichlet problem \eqref{eq:X.6.1} is equivalent to the problem: For a given $S_{1}v\in H_{0}^{m}(G)$, find a solution $\pmb{u}_{2}\in H_{0}^{m}(G)$ of 
\begin{equation}\tag{1'}\label{eq:X.6.1'}
(u _{2}, L ^{*} \varphi) _{0} = B (S _{1} v, \varphi), \varphi \in C _{0} ^{\infty} (G).
\end{equation}

Now, for a given $u \in L^2(G) = H_0^0(G)$, 
\[
\left| (u, \varphi) _{0} \right| \leqq \left\| u \| _{0} \cdot \| \varphi \| _{0} \right. \leqq \left\| u \| _{0} \cdot \| \varphi \| _{m} \right.
\]
so that, by F. \textsc{Riesz}' representation theorem in the Hilbert space $H_0^m(G)$, there exists a uniquely determined $u' = Tu \in H_0^m(G)$ such that, whenever $\varphi \in C_0^\infty(G)$, 
\[
(u, \varphi) _{0} = (u ^{\prime}, \varphi) _{m} \quad \text{ and } \quad \| u ^{\prime} \| _{m} \leqq \| u \| _{0}.
\]

Hence, by the Milgram-Lax theorem, we obtain 
\begin{equation}\label{eq:X.6.9}
(u, \varphi) _{0} = (u ^{\prime}, \varphi) _{m} = B \left(S _{1} u ^{\prime}, \varphi\right) = B \left(S _{1} T u, \varphi\right), \| S _{1} T u \| _{m} \leqq \delta^{-1} \| u \| _{0}.
\end{equation}

Therefore, by \eqref{eq:X.6.1}, we have, whenever $\varphi \in C_0^\infty(G)$, 
\[
\begin{aligned}
&B (u _{2}, \varphi)  = (u _{2}, \varphi + \alpha L ^{*} \varphi) _{0} = (u _{2}, \varphi) _{0} + \alpha (u _{2}, L ^{*} \varphi) _{0} = B (S _{1} T u _{2}, \varphi) + \alpha B (S _{1} v, \varphi),
\end{aligned}
\]
that is, 
\[
B \left(u _{2} - S _{1} T u _{2} - \alpha S _{1} v, \varphi\right) = 0.
\]

Because of the positivity $B(\varphi, \varphi) > 0$ of $B$, we must have 
\begin{equation}\tag{1''}\label{eq:X.6.1''}
u _{2} - S _{1} T u _{2} = \alpha S _{1} v.
\end{equation}

The right hand term $\alpha S_1 v \in H_0^m(G)$ is a known function. By $\|S_1 T u\|_m \leqq \delta^{-1} \|u\|_0$, we see that the operator $S_1 T$ defined on $H_0^m(G)$ into $H_0^m(G)$ is compact (the Rellich-Gårding theorem in Chapter~\ref{chap:X}, \ref{sec:X.3}). Therefore, we may apply the Riesz-Schauder theory to the effect that one of the following alternatives holds: 

Either the homogeneous equation $u - S_1 Tu = 0$ has a non-trivial solution $u \in H_0^m(G)$, or the inhomogeneous equation $u - S_1 Tu = w$ has, for every given $w \in H_0^m(G)$, a uniquely determined solution $u \in H_0^m(G)$. 

The first alternative corresponds to the case $(u, \varphi + \alpha L^{*}\varphi)_{0} = (u, \varphi)_{0}$, that is, to the case $Lu = 0$. Hence, returning to the original equation \eqref{eq:X.6.1}, we have 

\begin{theoremu} One of the following alternatives holds: Either \begin{enumerate}[i)]
\item the homogeneous equation $L u = 0$ has a non-trivial solution $u \in H_0^m(G)$
\item for any $f \in L^2(G)$ and any $u_1 \in H^m(G)$, there exists a uniquely determined solution $u_0 \in L^2(G)$ of $L u = f$, $u - u_1 \in H_0^m(G)$.
\end{enumerate}
\end{theoremu}
   % X
% ==== AX ====
\yoschapterstar{Appendix to Chapter X. The Nuclear Space of A. Grothendieck}
\label{chap:AX}

The nuclear operator defined in Chapter~\ref{chap:X}, \ref{sec:X.2} may be extended to locally convex spaces as follows. 

\begin{proposition}\label{prop:AX.0.Proposition.1} Let $X$ be a locally convex linear topological space, and $Y$ a $B$-space. Suppose that there exist an equi-continuous sequence $\{f_{j}^{\prime}\}$ of continuous linear functionals on $X$, a bounded sequence $\{y_{j}\}$ of elements $\in Y$ and a sequence of non-negative numbers $\{c_{j}\}$ with $\sum\limits_{j=1}^{n} c_{j} < \infty$. Then 
\begin{equation}\label{eq:AX.0.1}
T \cdot x = \slim _{n \rightarrow \infty} \sum\limits_{j=1} ^{n} c _{j} \langle x, f _{j} ^{\prime} \rangle y _{j}
\end{equation}
defines a continuous linear operator on $X$ into $Y$. 
\end{proposition}

\begin{proof}
 By the equi-continuity of $\{f_j'\}$, there exists a continuous seminorm $p$ on $X$ such that $\sup |\langle x, f_j' \rangle| \leqq p(x)$ for $x \in X$. Hence, for $m > n$, 
\[
\left\| \sum\limits_{j=n} ^{m} c _{j} \langle x, f _{j} ^{\prime} \rangle y _{j} \right\| \leqq p (x) \sup\limits _{j \geqq 1} \| y _{j} \| \cdot \sum\limits_{j=n} ^{m} c _{j}. \]

This proves that the right hand side of \eqref{eq:AX.0.1} exists and defines a continuous linear operator $T$ on $X$ into the $B$-space $Y$. 
\end{proof}

\begin{definition}\label{def:AX.0.Definition.1} An operator $T$ of the form \eqref{eq:AX.0.1} is said to be a nuclear operator on $X$ into the $B$-space $Y$. 
\end{definition}

\begin{corollaryu} $A$ nuclear operator $T$ is a compact operator in the sense that it maps a neighbourhood of $O$ of $X$ into a relatively compact set of $Y$. 

\end{corollaryu}

\begin{proof}
 We define 
\[
T _{n} \cdot x = \sum\limits_{j=1} ^{n} c _{j} \langle x, f _{j} ^{\prime} \rangle y _{j}.
\]
$T_{n}$ is compact, since the image by $T_{n}$ of the set $V = \{x; p(x) \leqq 1\}$ of $X$ is relatively compact in $Y$. On the other hand, we have 
\[
\left\| T x - T _{n} x \right\| = \left\| \sum\limits_{j=n+1} ^{\infty} c _{j} \langle x, f _{j} ^{\prime} \rangle y _{j} \right\| \leqq p (x) \sup\limits _{j \geqq 1} \left\| y _{j} \right\| \sum\limits_{j=n+1} ^{\infty} c _{j},
\]
and so $T_{n}x$ converges to $Tx$ strongly and uniformly on $V$. Hence the operator $T$ is compact. 
\end{proof}

As was proved in Chapter~\ref{chap:X}, \ref{sec:X.2}, we have a typical example of a nuclear operator: 

\begin{examplen} Let $K$ be a compact subset of $\mathbb{R}$. Then, for $(k-j)>n$, the identity mapping $T$ of $H_{0}^{k}(K)$ into $H_{0}^{j}(K)$ is a nuclear operator.

\end{examplen}
\begin{proposition}\label{prop:AX.0.Proposition.2} Let $X$ be a locally convex linear topological space, and $V$ a convex balanced neighbourhood of $0$ of $X$. Let $p_{V}(x) = \inf\limits_{x/\lambda \in V, \lambda > 0} \lambda$ be the Minkowski functional of $V$. $p_{V}$ is a continuous semi-norm on $X$. Set 
\[
N _{V} = \{x \in X; p _{V} (x) = 0 \} = \{x \in X; \lambda x \in V \text{ for all } \lambda > 0 \}.
\]
\end{proposition}

Then $N_V$ is a closed linear subspace of $X$, and the quotient space $X_V = X / N_V$ is a normed linear space by the norm 
\begin{equation}\left\| \tilde {x} \right\| _{V} = p _{V} (x), \text{ where } \tilde {x} \text{ is the residue class mod } N _{V} \text{ containing the element } x. \label{eq:AX.0.2}\end{equation}

\begin{proof}
 Let $(x - x_1) \in N_V$. Then $p_V(x_1) \leqq p_V(x) + p_V(x_1 - x) = p_V(x)$, and similarly $p_V(x) \leqq p_V(x_1)$. Thus $p_V(x) = p_V(x_1)$ if $x$ and $x_1$ are in the same residue class mod $N_V$. We have $\|\tilde{x}\|_V \geqq 0$ and $\|\tilde{0}\|_V = 0$. If $\|\tilde{x}\|_V = 0$, then $x \in \tilde{x}$ implies $x \in N_V$ and so $\tilde{x} = \tilde{0}$. The triangle inequality is proved by $\|\tilde{x} + \tilde{y}\|_V = p_V(x + y) \leqq p_V(x) + p_V(y) = \|\tilde{x}\|_V + \|\tilde{y}\|_V$. We have also $\|\alpha \tilde{x}\|_V = p_V(\alpha x) = |\alpha | p_V(x) = |\alpha \|\tilde{x}\|_V$. 
\end{proof}

\begin{corollaryu} By the equivalence 
\begin{equation}\label{eq:AX.0.3}
\left(p _ {V _{1}} \leqq p _ {V _{2}}\right) \leftrightarrow \left(V _{2} \subseteq V _{1}\right),
\end{equation}
we can define the canonical mapping 
\[
X _ {V _{2}} \rightarrow X _ {V _{1}} (\text{ when } V _{2} \subseteq V _{1})
\]
by associating the residue class $\tilde{x}_{V_{1}}\pmod{N_{V_{1}}}$ containing $x$ to the residue class $\tilde{x}_{V_{1}}\pmod{N_{V_{1}}}$ containing $\tilde{x}$. The mapping thus obtained is continuous, since 
\[
\left\| \tilde {x} _ {V _{1}} \right\| _ {V _{1}} = p _ {V _{1}} (x) \geqq p _ {V _{1}} (x) = \left\| \tilde {x} _ {V _{1}} \right\| _ {V _{1}}.
\]
\end{corollaryu}

We are now ready to give the notion of a nuclear space, introduced in analysis by A. \textsc{Grothendieck} \cite{grothendieck88}. 

\begin{definition}\label{def:AX.0.Definition.2} $A$ locally convex linear topological space $X$ is said to be a nuclear space, if, for any convex balanced neighbourhood $V$ of $0$, there exists another convex balanced neighbourhood $U \subseteq V$ of $0$ such that the canonical mapping 
\begin{equation}\label{eq:AX.0.4}
T \colon X _{U} \to \hat {X} _{V}
\end{equation}
is nuclear. Here $\hat{X}_{V}$ is the completion of the normed linear space $X_{V}$. 
\end{definition}

\begin{example}\label{ex:AX.0.Example.1} Let $\mathbb{R}^{A}$ be the topological product of real number field $\mathbb{R}$ in such a way that $\mathbb{R}^{A}$ is the totality of real-valued finite functions $x(a)$ defined on $A$ and topologized by the system of semi-norms 
\begin{equation}\label{eq:AX.0.5}
p_{a} (x) = | x (a) |, a \in A.
\end{equation}
\end{example}

Then $\mathbb{R}^{A}$ is a nuclear space. 

\begin{proof}
 $N_{V}$ is the totality of functions $x(a) \in \mathbb{R}^{A}$ such that, for some finite set $\{a_{j} \in A; j = 1, 2, \ldots, n\}$, $x(a_{j}) = 0$ ( $j = 1, 2, \ldots, n$). Hence $X_{V} = \mathbb{R}^{A}/N_{V}$ is equivalent to the space of functions $x_{V}(a)$ such that $x_{V}(a) = 0$ for $a \neq a_{j}$ ( $j = 1, 2, \ldots, n$) and normed by 
\[
\left\| x _{V} (a) \right\| _{V} = \sup\limits _{1 \leqq j \leqq n} | x (a _{j}) |. \]

We take $N_U$ the totality of functions $x(a) \in \mathbb{R}^{A}$ such that $x(a_\alpha) = 0$ for $\alpha \in A'$ where $A'$ is any finite set of integers containing $1, 2, \ldots, n$. Thus, for $U \subseteq V$, the canonical mapping $X_U = \mathbb{R}^{A} / N_U \to \mathbb{R}^{A} / N_V = X_V$ is nuclear. For the mapping is a continuous linear mapping with a finite-dimensional range. 
\end{proof}

\begin{example}\label{ex:AX.0.Example.2} $A$ nuclear $B$-space $X$ must be of finite dimension. 
\end{example}

\begin{proof}
 Since $X = X_{V}$ for any convex balanced neighbourhood $V$ of $0$ of a $B$-space, the compactness of the identity mapping $X \to X$ implies that $X$ is of finite dimension by F. \textsc{Riesz}' theorem in Chapter~\ref{chap:III}, \ref{sec:III.2}. 
\end{proof}

\begin{example}\label{ex:AX.0.Example.3} Let $K$ be a compact subset of $\mathbb{R}^{n}$. Then the space $\mathfrak{D}_K(\mathbb{R}^{n})$ introduced in Chapter~\ref{chap:I}, \ref{sec:I.1} is a nuclear space. 
\end{example}

\begin{proof}
 As in Chapter~\ref{chap:I}, \ref{sec:I.1}, let 
\[
\phi_{K,k} (f) = \sup\limits _{x \in K, | s | \leqq k} | D ^{s} f (x) |
\]
be one of the semi-norms which define the topology of $\mathfrak{D}_K(\mathbb{R}^{n})$. Let $V_k = \{f \in \mathfrak{D}_K(\mathbb{R}^{n}); p_{K,k}(f) \leqq 1\}$. Then $N_{V_{k}}$ is $\{0\}$, and $X_{V_{k}} = X / N_{V_{k}} = \mathfrak{D}_K(\mathbb{R}^{n}) / N_{V_{k}}$ is precisely the space $\mathfrak{D}_K(\mathbb{R}^{n})$ normed by $p_{K,k}$. If $(k-j) > n$, then it is easy to prove, as in the example following the Corollary of Definition~\ref{def:AX.0.Definition.1} above, that the canonical mapping of $X_{V_{k}}$ into $X_{V_{j}}$ is a nuclear transformation. Hence $\mathfrak{D}_K(\mathbb{R}^{n})$ is a nuclear space. 
\end{proof}

\begin{theorem}\label{thm:AX.0.Theorem.1} $A$ locally convex linear topological space $X$ is nuclear, iff, for any convex balanced neighbourhood $V$ of $0$, the canonical mapping $X \rightarrow \hat{X}_{V}$ is nuclear. 
\end{theorem}

\begin{proof}
 Necessity. Let $U \subseteq V$ be a convex balanced neighbourhood of $0$ of $X$ such that the canonical mapping $X_{U} \to \hat{X}_{V}$ is nuclear. The canonical mapping $T: X \to \hat{X}_{V}$ is the product of the canonical mapping $X \to \hat{X}_{U}$ and the canonical nuclear transformation $X_{U} \to \hat{X}_{V}$. Hence $T$ must be a nuclear transformation. 

Sufficiency. Let the canonical mapping $T: X \to \hat{X}_{V}$ be given by a nuclear transformation 
\[
T x = \sum\limits_{j=1} ^{\infty} c _{j} \langle x, f _{j} ^{\prime} \rangle y _{j}.
\]

For any $\alpha > 0$, the set $\{x \in X; |\langle x, f_j' \rangle| \leqq \alpha$ for $j = 1, 2, \ldots\}$ is a convex balanced neighbourhood $U_{\alpha}$ of $0$ of $X$, because of the equi-continuity of $\{f_j'\} \subseteq X'$. Moreover, 
$\| Tx\| _V = \left\| \sum\limits_{j}c_j\langle x,f_j'\rangle y_j\right\| _V\leqq \alpha \sup\limits_{j}\| y_j\| _V\sum\limits_{j}c_j$ whenever $x\in U_{\alpha}$ 

Let $\alpha$ be so small that the right hand side is $< 1$. Then $\|Tx\|_V < 1$ and $U_{\alpha} \subseteq V$. Each $f_j'$ may be considered as belonging to the dual space $X_{U_{\alpha}}'$, and so 
\[
T x = T z = \sum\limits_{j} c _{j} \left\langle x, f _{j} ^{\prime} \right\rangle y _{j} \quad \text{ whenever } \quad (x - z) \in N _ {U _{\alpha}}.
\]

Thus the canonical mapping $X_{U_{\alpha}} \to \hat{X}_V$ is given by a nuclear transformation 
\[
\tilde {x} _ {U _{\alpha}} \rightarrow \sum\limits_{j} c _{j} \left\langle \tilde {x} _ {U _{\alpha}}, f _{j} ^{\prime} \right\rangle y _{j}.
\]
\end{proof}

\begin{theorem}\label{thm:AX.0.Theorem.2} Let a locally convex linear topological space $X$ be nuclear. Then, for any convex balanced neighbourhood $V$ of $0$ of $X$, there exists a convex balanced neighbourhood $W \subseteq V$ of $0$ of $X$ such that $\hat{X}_{W}$ is a Hilbert space. 
\end{theorem}

\begin{proof}
 The nuclear canonical mapping $X_{U} \to \hat{X}_{V} (U \subseteq V)$ defined by 
\[
T \tilde {x} _{U} = \sum\limits_{j} c _{j} \left\langle \tilde {x} _{U}, f _{j} ^{\prime} \right\rangle y _{j}
\]
is factored as the product of the two mappings 
$\alpha : X_U \to (l^2)$ and $\beta : (l^2) \to \hat{X}_V$, where $\alpha$ is given by 
\[
\tilde {x} _{U} \rightarrow \left\{c _{j} ^{1/2} \langle \tilde {x} _{U}, f _{j} ^{\prime} \rangle \right\} \quad \text{ and } \quad \beta \text{ is given by } \{\xi_{j} \} \rightarrow \sum\limits_{j} c _{j} ^{1/2} \xi_{j} y _{j}. \]

The continuity of $\alpha$ is clear from 
\[
\sum\limits_{j} \left| c _{j} ^{1/2} \langle \tilde {x} _{U}, f _{j} ^{\prime} \rangle \right| ^{2} \leqq \left(\sup\limits _{j} \left\| f _{j} ^{\prime} \right\| \cdot \left\| \tilde {x} _{U} \right\| _{U}\right) ^{2} \cdot \sum\limits_{j} c _{j},
\]
and that of $\beta$ is proved by 
\[
\left\| \sum\limits_{j} c _{j} ^{1/2} \xi_{j} y _{j} \right\| _{V} ^{2} \leqq \sum\limits_{j} c _{j} \| y _{j} \| _{V} ^{2} \cdot \sum\limits_{j} | \xi_{j} | ^{2} \leqq \sup\limits _{j} \| y _{j} \| _{V} ^{2} \cdot \| \{\xi_{j} \} \| _ {l ^{n}} ^{2} \cdot \sum\limits_{j} c _{j}.
\]

Let $U_{2}$ be the inverse image in $(l^{2})$ by $\beta$ of the unit sphere of $\hat{X}_{V}$. Then $U_{2}$ is a neighbourhood of $0$ of $(l^{2})$ and so contains a sphere $S$ of centre $0$ of $(l^{2})$. 

Let $W$ be the inverse image in $X$ of $S$ by the continuous mapping $\tilde{\alpha}$ defined as the product of continuous canonical mapping $X \to X_U$ and the continuous mapping $\alpha: X_U \to (l^2)$. Then clearly $W \subseteq V$ and, for any $\tilde{x}_W \in X_W$, 
\[
\left\| \tilde {x} _{W} \right\| _{W} = \inf\limits _{x / \lambda \in W, \lambda > 0} \lambda = \inf\limits _ {\tilde {\alpha} x / \lambda \in S, \lambda > 0} \lambda = \left\| \tilde {\alpha} x \right\| _ {L ^{2}} (\text{ the radius of } S).
\]

Since $\|\|_{l^2}$ is the norm in the Hilbert space $(l^2)$, $X_W$ is a pre-Hilbert space. 
\end{proof}

\begin{corollaryu} Let $X$ be a locally convex nuclear space. Then, for any convex balanced neighbourhood $V$ of $0$ of $X$, there exist convex balanced neighbourhoods $W_{1}$ and $W_{2}$ of $0$ of $X$ with the properties: 

$W_{2} \subseteq W_{1} \subseteq V, \hat{X}_{W_{1}}$ and $\hat{X}_{W_{1}}$ are Hilbert spaces and the canonical 
mappings $X \to \hat{X}_{W_2}, \hat{X}_{W_2} \to \hat{X}_{W_1}, \hat{X}_{W_1} \to \hat{X}_V$ are all nuclear. 
\end{corollaryu}

Therefore, a nuclear space $X$ has a fundamental system $\{V_{\alpha}\}$ of neighbourhoods of $0$ such that the spaces $\hat{X}_{V_{\alpha}}$ are Hilbert spaces. 

Further Properties of Nuclear Spaces. It can be proved that: 

$1$. A linear subspace and a factor space of a nuclear space are also nuclear. 

$2$. The topological vector product of a family of nuclear spaces and the inductive limit of a sequence of nuclear spaces are also nuclear. 

$3$. The strong dual of the inductive limit of a sequence of nuclear spaces, each of which is an $F$-space, is also nuclear. 

For the proof, see the book by \textsc{Grothendieck} \cite{grothendieck87} referred to above, p. $47$. As a consequence of $2$., the space $\mathfrak{D}(\mathbb{R}^{n})$, which is the inductive limit of the sequence $\{\mathfrak{D}_{K_{r}}(\mathbb{R}^{n}); r = 1, 2, \ldots\}$ (here $K_r$ is the sphere $|x| \leqq r$ of $\mathbb{R}^{n}$), is nuclear. Hence, by $3$., the space $\mathfrak{D}(\mathbb{R}^{n})'$ is also nuclear. The spaces $\mathfrak{E}(\mathbb{R}^{n}), \mathfrak{E}(\mathbb{R}^{n})'$, $\mathfrak{S}(\mathbb{R}^{n})$ and $\mathfrak{S}(\mathbb{R}^{n})'$ are also nuclear spaces. 

The importance of the notion of the nuclear space has recently been stressed by R. A. \textsc{Minlos} \cite{minlos202}. He has proved the following generalization of Kolmogorov's extension theorem of measures: 

Let $X$ be a nuclear space whose topology is defined through a countable system of convex balanced neighbourhoods of $0$. Let $X'$ be the strong dual space of $X$. A cylinder set of $X'$ is defined as a set of the form 
\[
Z ^{\prime} = \{f ^{\prime} \in X ^{\prime}; a _{i} < \langle x _{i}, f ^{\prime} \rangle < b _{i} (i = 1, 2, \dots , n) \}.
\]

Suppose there is given a set function $\mu_{0}$, defined and $\geqq 0$ for all cylinder sets. Let $\mu_{0}$ be $\sigma$-additive for those cylinder sets $Z'$ with fixed $x_{1}, x_{2}, \ldots, x_{n}$. Then, under a compatibility condition and a continuity condition, there exists a uniquely determined extension of $\mu_{0}$ which is $\sigma$-additive and $\geqq 0$ for all sets of the smallest $\sigma$-additive family of sets of $X'$ containing all the cylinder sets of $X'$. 

For a detailed proof and applications of this result, see I. M. \textsc{Gelfand}-N. Y. \textsc{Vilenkin} \cite{gelfand77}.\yosstarfootnote{See also Supplementary Notes, p.~\pageref{notes:AX}.} 
   % Appendix to Chapter X
% ==== XI ====
\yoschapter{XI}{Normed Rings and Spectral Representation}
\label{chap:XI}

A linear space $A$ over a scalar field $(F)$ is said to be an algebra or a ring over $(F)$, if to each pair of elements $x, y \in A$ a unique product $xy \in A$ is defined with the properties: 
\begin{equation}\label{eq:XI.0.1}
\begin{yoslines}
&(x y) z = x (y z) (\text{ associativity }), x (y + z) = x y + x z (\text{ distributivity }), \alpha \beta (x y) = (\alpha x) (\beta y).
\end{yoslines}
\end{equation}

If there exists a unit element $e$ such that $ex = xe = x$ for every $x \in A$, then $A$ is said to be an algebra with a unit. A unit $e$ of $A$, if it exists, is uniquely determined. For, if $e'$ be another unit of $A$, then we must have $ee' = e = e'$. If the multiplication is commutative, i.e., $xy = yx$ for every pair $x, y \in A$, then $A$ is called a commutative algebra. Let $A$ be an algebra with a unit $e$. If, for an $x \in A$, there exists an $x' \in A$ such that $xx' = x'x = e$, then $x'$ is called an inverse of $x$. An inverse $x'$ of $x$, if it exists, is uniquely determined. For, if $x''$ be another inverse of $x$, then we must have 
\[
x ^{\prime \prime} (x x ^{\prime}) = x ^{\prime \prime} e = x ^{\prime \prime} = (x ^{\prime \prime} x) x ^{\prime} = e x ^{\prime} = x ^{\prime}.
\]

Thus we shall denote by $x^{-1}$ the inverse of $x$ if $x$ has an inverse. 

An algebra is called a Banach algebra, or in short a $B$-algebra if it is a $B$-space and satisfies 
\begin{equation}\label{eq:XI.0.2}
\left\| x y \right\| \leqq \left\| x \right\| \left\| y \right\|.
\end{equation}

The inequality 
\[
\begin{aligned}
&\left\| x _{n} y _{n} - x y \right\|  \leqq \left\| x _{n} (y _{n} - y) \right\| + \left\| (x _{n} - x) y \right\| \leqq \left\| x _{n} \right\| \left\| (y _{n} - y) \right\| + \left\| (x _{n} - x) \right\| \left\| y \right\|
\end{aligned}
\]
shows that $xy$ is a continuous function of both variables together. 

\begin{example}\label{ex:XI.0.Example.1} Let $X$ be a $B$-space. Then $L(X, X)$ is a $B$-algebra with a unit by the operator sum $T + S$ and operator product $TS$; the identity operator $I$ is the unit of this algebra $L(X, X)$, and the operator norm $\|T\|$ is the norm of the element $T$ of this algebra $L(X, X)$. 
\end{example}

\begin{example}\label{ex:XI.0.Example.2} Let $S$ be a compact topological space. Then $C(S)$ is a $B$-algebra by $(x_1 + x_2)(s) = x_1(s) + x_2(s), (\alpha x)(s) = \alpha x(s), (x_1x_2)(s) = x_1(s)x_2(s)$ and $\|x\| = \sup\limits_{s\in S}|x(s)|$. 
\end{example}

\begin{example}\label{ex:XI.0.Example.3} Let $B$ be the totality of continuous functions $x(s)$, $0 \leqq s \leqq 1$, which are representable as absolutely convergent Fourier series: 
\begin{equation}\label{eq:XI.0.3}
x (s) = \sum\limits_{n = - \infty} ^{\infty} c _{n} e ^{2 \pi i n s} \quad \text{ with } \quad \sum\limits_{n = - \infty} ^{\infty} | c _{n} | < \infty .
\end{equation}
\end{example}

Then it is easy to see that $B$ is a commutative $B$-algebra with a unit by the ordinary function sum and function multiplication, normed by 
\begin{equation}\label{eq:XI.0.4}
\left\| x \right\| = \sum\limits_{n = - \infty} ^{\infty} \left| c _{n} \right|.
\end{equation}

In the last two examples, the unit is given by the function $e(s) \equiv 1$ and $\|e\| = 1$. In the following sections, we shall be concerned with the commutative $B$-algebra with the unit $e$ such that 
\begin{equation}\label{eq:XI.0.5}
\left\| e \right\| = 1.
\end{equation}

Such an algebra is called a normed ring. 

A Historical Sketch. The notion of Banach algebras was introduced in analysis by M. \textsc{Nagumo} \cite{nagumo215}. He proved that Cauchy's complex function theory can be extended to functions with values in such an algebra, and applied it to the investigation of the resolvent of a bounded linear operator around an isolated singular point. The result is an abstract treatment of those given in our Chapter~\ref{chap:VIII}, \ref{sec:VIII.8}. K. \textsc{Yosida} \cite{yosida354} proved that a connected group embedded in a $B$-algebra is a Lie group iff the group is locally compact. This result is an extension of a result due to J. \textsc{von} \textsc{Neumann} \cite{neumann233} concerning matrix groups. Cf. E. \textsc{Hille}-R. S. \textsc{Phillips} \cite{hille97}, in which the result of K. \textsc{Yosida} \cite{yosida354} is reproduced. 

The ideal theory of normed rings was initiated by I. M. \textsc{Gelfand} \cite{gelfand76b}. He has shown that such a ring can be represented as the ring of continuous functions defined on the space of maximal ideals of the ring. By virtue of this representation, we can give an integration free treatment of the spectral resolution of bounded normal operators in a Hilbert space; see K. \textsc{Yosida} \cite{yosida355}. This result will be exposed in the following sections. The Gelfand representation may also be applied to a new proof of the Tauberian theorem of N. \textsc{Wiener} \cite{wiener340}. We shall expose this application in the last section of this chapter. For further details about $B$-algebras, see N. A. \textsc{Naimark} \cite{naimark221}, C. E. \textsc{Rickart} \cite{rickart271} and I. M. \textsc{Gelfand}-D. A. \textsc{Raikov}-G. E. \textsc{Šilov} \cite{gelfand79}.\yosstarfootnote{See also Supplementary Notes, p.~\pageref{notes:XI}.} 


\yossection{1}{Maximal Ideals of a Normed Ring}
\label{sec:XI.1}

We shall be concerned with a commutative $B$-algebra $B$ with a unit $e$ such that $\| e \| = 1$. 

\begin{definition}\label{def:XI.1.Definition.1} $A$ subset $J$ of $B$ is called an ideal of $B$ if $x, y \in J$ implies that $(\alpha x + \beta y) \in J$ and $zx \in J$ for every $z \in B$. $B$ itself and $\{0\}$ are ideals of $B$. Ideals other than $B$ and $\{0\}$ are called non-trivial ideals. A non-trivial ideal $J$ is said to be a maximal ideal if there exists no non-trivial ideal containing $J$ as a proper subset. 
\end{definition}

\begin{proposition}\label{prop:XI.1.Proposition.1} Each non-trivial ideal $J_0$ of $B$ is contained in a maximal ideal $J$. 
\end{proposition}

\begin{proof}
 Let $[J_0]$ be the set of all non-trivial ideals containing $J_0$. We order the ideals of $[J_0]$ by inclusion relation, that is, we denote $J_1 \prec J_2$ if $J_1$ is a subset of $J_2$. Suppose that $\{J_\alpha\}$ is a linearly ordered subset of $[J_0]$ and put $J_\beta = \bigcup\limits_{J_\alpha \in \{J_\alpha\}} J_\alpha$. We shall show that $J_\beta$ is an upper bound of $\{J_\alpha\}$. For, if $x, y \in J_\beta$, then there exist ideals $J_{\alpha_1}$ and $J_{\alpha_2}$ such that $x \in J_{\alpha_1}$ and $y \in J_{\alpha_2}$. Since $\{J_\alpha\}$ is linearly ordered, $J_{\alpha_1} \prec J_{\alpha_2}$ (or $J_{\alpha_1} > J_{\alpha_2}$) and so $x$ and $y$ both belong to $J_{\alpha_2}$; consequently $(x - y) \in J_{\alpha_2} \subseteq J_\beta$ and $zx \in J_{\alpha_2} \subseteq J_\beta$ for any $z \in B$. This proves that $J_\beta$ is an ideal. Since the unit element $e$ is not contained in any $J_\alpha, e$ is not contained in $J_\beta = \bigcup\limits_{J_\alpha \in \{J_\alpha\}} J_\alpha$. 

Thus $J_{\beta}$ is a non-trivial ideal containing every $J_{\alpha}$. Therefore, by Zorn's Lemma, there exists at least one maximal ideal which contains $J_{0}$. 
\end{proof}

\begin{corollaryu} An element $x$ of $B$ has the inverse $x^{-1} \in B$ such that $x^{-1}x = xx^{-1} = e$ iff $x$ is contained in no maximal ideal. 
\end{corollaryu}

\begin{proof}
 If $x^{-1} \in B$ exists, then any ideal $J \ni x$ must contain $e = xx^{-1}$ so that $J$ must coincide with $B$ itself. Let conversely, $x$ be contained in no maximal ideal. Then the ideal $xB = \{xb; b \in B\} \neq \{0\}$ must coincide with $B$ itself, since, otherwise, there exists at least one maximal ideal containing $xB \ni x = xe$. It follows that $xB = B$, and so there must exist an element $b \in B$ such that $xb = e$. By the commutativity of $B$, we have $xb = bx = e$, that is, $b = x^{-1}$. 
\end{proof}

\begin{proposition}\label{prop:XI.1.Proposition.2} $A$ maximal ideal $J$ is a closed linear subspace of $B$. 
\end{proposition}

\begin{proof}
 By the continuity of the algebraic operations (addition, multiplication and scalar multiplication) in $B$, the strong closure $J^a$ is also an ideal containing $J$. Suppose $J^a \neq J$. Then $J^a = B$, because of the maximality of the ideal $J$. Thus $e \in J^a$, and so there exists an $x \in J$ such that $\| e - x \| < 1$. $x$ has the inverse $x^{-1} \in B$ which is given by Neumann's series $e + (e - x) + (e - x)^2 + \cdots$ 
\[
e + (e - x) + (e - x) ^{2} + \dots \]

For, by $\|(e-x)^{n}\|\leqq\|e-x\|^{n}$, the series converges to an element $\in B$ which is the inverse of $x$, as may be seen by multiplying the series by $x=e-(e-x)$. Hence $e=x^{-1}x\in J$ and so $J$ cannot be a maximal ideal. 
\end{proof}

\begin{proposition}\label{prop:XI.1.Proposition.3} For any ideal $J$ of $B$, we write 

$x \equiv y (\bmod J)$ or $x \sim y (\bmod J)$ or in short $x \sim y$, if $(x - y) \in J$. Then $x \sim y$ is an equivalence relation, that is, we have 
\begin{equation}\label{eq:XI.1.1}
\begin{yoslines}
&x \sim x (\mathrm{reflexivity}), x \sim y \mathrm{implies} y \sim x (\mathrm{symmetry}), x \sim y \mathrm{and} y \sim z \mathrm{implies} x \sim z (\mathrm{transitivity}).
\end{yoslines}
\end{equation}
\end{proposition}

We denote by $\bar{x}$ the set $\{y; (y-x) \in J\}$; it is called the class (mod $J$) containing $x$. Then the classes $(\overline{x+y})$, $\alpha x$ and $(\overline{xy})$ are determined independently of the choice of elements $x$ and $y$ from the classes $\bar{x}$ and $\bar{y}$, respectively. 

Proof We have to show that $x \sim x'$, $y \sim y'$ implies that $(x + y) \sim (x' + y')$, $\alpha x \sim \alpha x'$ and $xy \sim x'y'$. These are clear from the condition that $J$ is an ideal. For instance, we have $xy - x'y' = (x - x')y + x'(y - y') \in J$ by $(x - x') \in J$ and $(y - y') \in J$. 

\begin{corollaryu} The set of classes $\overline{x}$ (mod $J$) thus constitutes an algebra by 
\begin{equation}\label{eq:XI.1.2}
\overline {{{x}}} + \overline {{{y}}} = \overline {{{x + y}}}, \alpha \overline {{{x}}} = \overline {{{\alpha x}}}, \overline {{{x}}} \overline {{{y}}} = \overline {{{x y}}}.
\end{equation}
\end{corollaryu}

\begin{definition}\label{def:XI.1.Definition.2} The above obtained algebra is called the residue class algebra of $B \pmod{J}$ and is denoted by $B / J$. Thus the mapping $x \to \overline{x}$ of the algebra $B$ onto $\overline{B} = B / J$ is a homomorphism, that is, relation \eqref{eq:XI.1.2} holds. 
\end{definition}

\begin{proposition}\label{prop:XI.1.Proposition.4} Let $J$ be a maximal ideal of $B$. Then $\overline{B}=B/J$ is a field, that is, each non-zero element $\bar{x}\in\overline{B}$ has an inverse $\bar{x}^{-1}\in\overline{B}$ such that $\bar{x}^{-1}\bar{x}=\bar{x}\bar{x}^{-1}=\bar{e}$. 
\end{proposition}

\begin{proof}
 Suppose the inverse $\bar{x}^{-1}$ does not exist. Then the set $\bar{x}\overline{B}=\{\bar{x}\bar{b};\bar{b}\in\overline{B}\}$ is an ideal of $\overline{B}$. It is non-trivial since it does not contain $\bar{e}$, but does contain $\bar{x}\neq0$. The inverse image of an ideal by the homomorphism is an ideal. Therefore, $B$ contains a non-trivial ideal containing $J$ as a proper subset, contrary to the maximality of the ideal $J$. 
\end{proof}

We are now able to prove 

\begin{theoremu} Let $B$ be a normed ring over the field of complex numbers, and $J$ a maximal ideal of $B$. Then the residue class algebra $\overline{B}=B/J$ is isomorphic to the complex number field, in the sense that each $\overline{x}\in\overline{B}$ is represented uniquely as $\overline{x}=\xi\overline{e}$, where $\xi$ is a complex number. 
\end{theoremu}

\begin{proof}
 We shall prove that $\overline{B}=B/J$ is a normed ring by the norm 
\begin{equation}\label{eq:XI.1.3}
\left\| \overline {{x}} \right\| = \inf\limits _ {x \in \overline {{x}}} \left\| x \right\|. \end{equation}

If this is proved, then $B / J$ is a normed field and so, by the Gelfand-Mazur theorem in Chapter~\ref{chap:V}, \ref{sec:V.3}, $\overline{B} = B / J$ is isomorphic to the complex number field. 

Now we have $\|\alpha \overline{x}\| = |\alpha \|\overline{x}\|$, and $\|\overline{x} + \overline{y}\| = \inf\limits_{x \in \overline{x}, y \in \overline{y}} \|x + y\| \leqq \inf\limits_{x \in \overline{x}} \|x\| + \inf\limits_{y \in \overline{y}} \|y\| = \|\overline{x}\| + \|\overline{y}\|; \|\overline{x} \overline{y}\| \leqq \|\overline{x}\| \|\overline{y}\|$ is proved similarly. If $\|\overline{x}\| = 0$, then there exists a sequence $\{x_n\} \subseteq \overline{x}$ such that $\slim\limits_{n \to \infty} x_n = 0$. Hence, for any $x \in \overline{x}$, $(x - x_n) \in J$ and so $\slim\limits_{n \to \infty} (x - x_n) = x$ which proves that $x \in J^a = J$, that is, $\overline{x} = \overline{0}$. Hence $\|\overline{x}\| = 0$ is equivalent to $\overline{x} = \overline{0}$. We have $\|\overline{e}\| \leqq \|e\| = 1$. If $\|\overline{e}\| < 1$, then there exists an element $x \in J$ such that $\|e - x\| < 1$. As in the proof of Proposition~\ref{prop:XI.1.Proposition.2}, the inverse $x^{-1}$ exists, which is contradictory to the Corollary of Proposition~\ref{prop:XI.1.Proposition.1}. Thus we must have $\|\overline{e}\| = 1$. Finally, since $B$ is a $B$-space and $J$ is a closed linear subspace by Proposition~\ref{prop:XI.1.Proposition.2}, the factor space $\overline{B} = B / J$ is complete with respect to the norm \eqref{eq:XI.1.3} (see Chapter~\ref{chap:I}, \ref{sec:I.11}). We have thus proved the Theorem. 
\end{proof}

\begin{corollaryu} We shall denote by $x(J)$ the number $\xi$ in the representation $\overline{x} = \xi \overline{e}$. Thus, for each $x \in B$, we obtain a complex-valued function $x(J)$ defined on the set $\{J\}$ of all the maximal ideals of $B$. Then we have 
\begin{equation}\label{eq:XI.1.4}
\begin{gathered}
(x + y) (J) = x (J) + y (J), (\alpha x) (J) = \alpha x (J), (x y) (J) = x (J) y (J), \text{ and } e (J) \equiv 1.
\end{gathered}
\end{equation}
\end{corollaryu}

We have, moreover, 
\begin{equation}\label{eq:XI.1.5}
\sup\limits _ {J \in \{J \}} | x (J) | \leqq \| x \|,
\end{equation}
and 
\begin{equation}\label{eq:XI.1.6}
\sup\limits _ {J \in \{J \}} | x (J) | = 0 \text{ implies } x = 0 \text{ iff } \bigcap\limits_ {J \in \{J \}} J = \{0 \}.
\end{equation}

\begin{proof}
 The mapping $x \to \bar{x} = x(J) \bar{e}$ of the algebra $B$ onto the residue class algebra $\overline{B} = B/J$ is a homomorphism, that is, relation \eqref{eq:XI.1.2} holds. Hence we have \eqref{eq:XI.1.4}. Inequality \eqref{eq:XI.1.5} is proved by 
\[
| \xi | = | \xi | \| \overline {{e}} \| = \| \overline {{x}} \| = \inf\limits _ {x \in \overline {{x}}} \| x \| \leqq \| x \|. \]

Property \eqref{eq:XI.1.6} is clear, since $x(J) = 0$ identically on $\{J\}$ iff $x \in \bigcap\limits_{J \in \{J\}} J$. 
\end{proof}

\begin{definition}\label{def:XI.1.Definition.3} The representation 
\begin{equation}\label{eq:XI.1.7}
x \rightarrow x (J)
\end{equation}
of the normed ring $B$, by the ring of functions $x(J)$ defined on the set $\{J\}$ of all the maximal ideals $J$ of $B$, is called the Gelfand-representation of $B$. 
\end{definition}


\yossection{2}{The Radical. The Semi-simplicity}
\label{sec:XI.2}

\begin{definition}\label{def:XI.2.Definition.1} Let $B$ be a normed ring over the complex number field, and $\{J\}$ the totality of the maximal ideals $J$ of $B$. Then the ideal $\bigcap\limits_{J \in \{J\}} J$ is called the radical of the ring $B$. $B$ is said to be semi-simple if its radical $R = \bigcap\limits_{J \in \{J\}} J$ reduces to the zero ideal $\{0\}$. 
\end{definition}

\begin{theorem}\label{thm:XI.2.Theorem.1} For any $x \in B$, $\lim\limits_{n \to \infty} \| x^n \|^{1/n}$ exists and we have 
\begin{equation}\label{eq:XI.2.1}
\lim\limits _{n \rightarrow \infty} \| x ^{n} \| ^{1/n} = \sup\limits _ {J \in \{J \}} | x (J) |.
\end{equation}
\end{theorem}

\begin{proof}
 Set $\alpha = \sup\limits_{J \in \{J\}} |x(J)|$. Then, by $\|x^n\| \geqq |x^n(J)| = |x(J)|^n$, we have $\|x^n\| \geqq \alpha^n$, and so $\lim\limits_{n \to \infty} \|x^n\|^{1/n} \geqq \alpha$. We have thus to prove $\varlimsup\limits_{n \to \infty} \|x^n\|^{1/n} \leqq \alpha$. 

Let $|\beta| > \alpha$. Then, for any $J \in \{J\}$, $x(J) - \beta \neq 0$, i.e., $(x - \beta e) \bar{\in} J$. Hence the inverse $(\beta e - x)^{-1}$ exists. Setting $\beta^{-1} = \lambda$, we see that the inverse $(\beta e - x)^{-1} = \lambda (e - \lambda x)^{-1}$ exists whenever $|\lambda| < \alpha^{-1}$. Moreover, as in Theorem~\ref{thm:VIII.2.Theorem.1} in Chapter~\ref{chap:VIII}, \ref{sec:VIII.2}, we see that $\lambda (e - \lambda x)^{-1}$ is, for $|\lambda| < \alpha^{-1}$, holomorphic in $\lambda$. Hence we have the Taylor expansion 
\[
\lambda (e - \lambda x) ^{-1} = \lambda (e + \lambda x _{1} + \lambda^{2} x _{2} + \dots + \lambda^{n} x _{n} + \dots).
\]

That $x_{n} = x^{n}$ may be seen by Neumann's series $(e - \lambda x)^{-1} = \sum\limits_{n=0}^{\infty} \lambda^{n} x^{n}$ which is valid for $\|\lambda x\| < 1$. By the convergence of the above Taylor series, we see that 
\[
\lim\limits _{n \rightarrow \infty} \| \lambda^{n} x ^{n} \| = 0 \text{ if } | \lambda | < \alpha^{-1}.
\]

Thus $\|x^n \| = |\lambda|^{-n}\|\lambda^n x^n \| < |\lambda|^{-n}$ for large $n$ when $|\lambda| < \alpha^{-1}$, and so $\varlimsup\limits_{n \to \infty} \|x^n \|^{1/n} \leqq |\lambda|^{-1}$ when $|\lambda|^{-1} > \alpha$, that is, $\varlimsup\limits_{n \to \infty} \|x^n \|^{1/n} \leqq \alpha$. 
\end{proof}

\begin{corollaryu} The radical $R = \bigcap\limits_{J \in \{J\}} J$ of $B$ coincides with the totality of the generalized nilpotent elements $x \in B$ which are defined by 
\begin{equation}\label{eq:XI.2.2}
\lim\limits _{n \rightarrow \infty} \| x ^{n} \| ^{1/n} = 0.
\end{equation}
\end{corollaryu}

\begin{definition}\label{def:XI.2.Definition.2} $A$ complex number $\lambda$ is said to belong to the spectrum of $x \in B$, if the inverse $(x - \lambda e)^{-1}$ does not exist in $B$. 
\end{definition}

If $\lambda$ belongs to the spectrum of $x$, then there exists a maximal ideal $J$ such that $(x - \lambda e) \in J$. Conversely, if $(x - \lambda e)$ belongs to a maximal ideal $J$, then the inverse $(x - \lambda e)^{-1}$ does not exist. Hence we obtain 

\begin{theorem}\label{thm:XI.2.Theorem.2} The spectrum of $x \in B$ coincides with the totality of values taken by the function $x(J)$ on the space $\{J\}$ consisting of all the maximal ideals $J$ of $B$. 
\end{theorem}

Application of Tychonov's theorem. We define, for any $J_0 \in \{J\}$, a fundamental system of neighbourhoods of $J_0$ by 
\begin{equation}\label{eq:XI.2.3}
\{J \in \{J \}; | x _{i} (J) - x _{i} (J _{0}) | < \varepsilon_{i} (i = 1, 2, \dots , n) \},
\end{equation}
where $\varepsilon_{i} > 0$, $n$ and $x_{i} \in B$ are arbitrary. Then $\{J\}$ becomes a topological space and each $x(J)$, $x \in B$, becomes a continuous function on $\{J\}$. We have only to verify that if $J_{0} \neq J_{1}$, then there exist a neighbourhood $V_{0}$ of $J_{0}$ and a neighbourhood $V_{1}$ of $J_{1}$ with empty intersection. This may be done as follows. Let $x_{0} \in J_{0}$ and $x_{0} \bar{\in} J_{1}$ so that $x_{0}(J_{0}) = 0$ and $x_{0}(J_{1}) = \alpha \neq 0$. Then $V_{0} = \{J \in \{J\}; |x_{0}(J)| < |\alpha| / 2\}$ and $V_{1} = \{J \in \{J\}; |x_{0}(J) - x_{0}(J_{1})| < |\alpha| / 2\}$ have an empty intersection. 

\begin{theorem}\label{thm:XI.2.Theorem.3} The space $\{J\}$, topologized as above, is a compact space. 
\end{theorem}

\begin{proof}
 We attach, to each $x \in B$, the compact set 
\[
K _{x} = \{z; | z | \leqq \| x \| \}
\]
of the complex $z$-plane. Then the topological product 
\[
S = \prod\limits_{x \in B} K _{x}
\]
is a compact space, by virtue of Tychonov's theorem. See Chapter~\ref{chap:0}. To any maximal ideal $J_0 \in \{J\}$, we assign the point 
\[
\prod\limits_{x \in B} x (J _{0}) = s (J _{0}) \in S.
\]
$\{J\}$ is in one-one correspondence with a subset $S_{1}$ of $S$ by the above correspondence $J_{0} \rightarrow s(J_{0})$. Moreover, the topology of $\{J\}$ is the same as the relative topology of $S_{1}$ as a subset of $S$. Hence, if we can show that $S_{1}$ is a closed subset of the compact space $S$, then its topological image $\{J\}$ is compact. 

To prove that $S_{1}$ is a closed set, we consider an accumulation point $\omega = \prod\limits_{x \in B} \lambda_x \in S$ of the set $S_{1}$ in $S$. We shall show that the mapping $x \to \lambda_x$ is a homomorphism of the algebra $B$ into the complex number field $(K)$. Then $J_0 = \{x; \lambda_x = 0\}$ becomes, as may be seen from the isomorphism of $B / J_0$ with $(K)$, a maximal ideal of $B$ and $(x - \lambda_x e) \in J_0$, that is, $x(J_0) = \lambda_x$. This proves that the point $\omega = \prod\limits_{x \in B} \lambda_x = \prod\limits_{x \in B} x(J_0)$ belongs to $S_{1}$. 

We thus have to show that 
\[
\lambda_{x+y} = \lambda_{x} + \lambda_{y}, \lambda_{\alpha x} = \alpha \lambda_{x}, \lambda_{xy} = \lambda_{x} \lambda_{y}, \lambda_{e} = 1.
\]

We shall, for instance, prove that $\lambda_{x+y} = \lambda_x + \lambda_y$. Since $\omega = \prod\limits_{x \in B} \lambda_x$ is an accumulation point of $S_1$, there exists, for any $\varepsilon > 0$, a maximal ideal $J$ such that 
\begin{equation}\label{eq:XI.2.4}
\left| \lambda_{x} - x (J) \right| < \varepsilon , \left| \lambda_{y} - y (J) \right| < \varepsilon , \left| \lambda_{x+y} - (x + y) (J) \right| < \varepsilon .
\end{equation}

By $(x + y)(J) = x(J) + y(J)$ and the arbitrariness of $\varepsilon > 0$, we easily see that $\lambda_{x+y} = \lambda_x + \lambda_y$ is true. 
\end{proof}

We can now state the fundamental facts about the Gelfand representation $x \to x(J)$ of the normed ring $B$ in the form of 

\begin{theorem}\label{thm:XI.2.Theorem.4} $A$ normed ring $B$ over the complex number field is represented homomorphically by the ring of functions $x(J)$ on the compact space $\{J\}$ of all the maximal ideals $J$ of $B$. The radical $R$ of $B$ consists of those and only those elements which are represented by functions identically equal zero on $\{J\}$. The representation $x \to x(J)$ is isomorphic iff the ring $B$ is semi-simple. 
\end{theorem}

Application of the Stone-Weierstrass Theorem (of Chapter~\ref{chap:0}). The above obtained ring of functions is dense in the space of all complex-valued continuous functions on $\{J\}$ with uniform convergence topology if the ring $B$ is symmetric (or involutive) in the following sense: 

For any $x \in B$, there exists an $x^{*} \in B$ such that $x^{*}(J) = \overline{x\langle J\rangle}$ on $\{J\}$. ($4$) 

\yossub{Examples of Gelfand Representations}

\begin{example}\label{ex:XI.2.Example.1} Let $B = C(S)$ where $S$ is a compact topological space, and $J_0$ a maximal ideal of $C(S)$. Then there exists a point $s_0 \in S$ such that $x(s_0) = 0$ for all $x \in J_0$. Otherwise, for any $s_\alpha \in S$, there exists an $x_\alpha \in J_0$ such that $x_\alpha(s_\alpha) \neq 0$. $x_\alpha(s)$ being a continuous function, there exists a neighbourhood $V_\alpha$ of $s_\alpha$ such that $x_\alpha(s) \neq 0$ in $V_\alpha$. Since $S$ is compact, there exists a finite system, say, $V_{\alpha_1}, V_{\alpha_2}, \ldots, V_{\alpha_n}$ such that $\bigcup\limits_{j=1}^{n} V_{\alpha_j} = S$. Hence the function 
\[
x (s) = \sum\limits_{i=1} ^{n} \overline {{{{x _ {\alpha_{i}} (s)}}}} x _ {\alpha_{i}} (s) \in J _{0}
\]
does not vanish over $S$, and the inverse $x^{-1}$, $x^{-1}(s) = x(s)^{-1}$, of $x \in J_0$ exists, contrary to the maximality of the ideal $J_0$. Thus we see that $J_0$ is contained in the maximal ideal $J' = \{x \in B; x(s_0) = 0\}$. By the maximality of $J_0$, we must have $J_0 = J'$. In this way we see that the space $\{J\}$ of the maximal ideals $J$ of $B$ is in one-one correspondence with the points $s$ of $S$. 
\end{example}

\begin{example}\label{ex:XI.2.Example.2} Let $B$ be the totality of functions $x(s)$, $0 \leqq s \leqq 1$, which can be represented by absolutely convergent Fourier series: 
\[
x (s) = \sum\limits_{n = - \infty} ^{\infty} c _{n} e ^{2 \pi i s n}, \sum\limits_{n = - \infty} ^{\infty} | c _{n} | < \infty .
\]
\end{example}

$B$ is a normed ring by $(x + y)(s) = x(s) + y(s), (xy)(s) = x(s)y(s)$ and $\|x\| = \sum\limits_{j} |c_j|$. Let $J_0$ be a maximal ideal of $B$. Set $e^{2\pi is} = x_1$. Then $x_1^{-1} = e^{-2\pi is}$, and so, by $|x_1(J_0)| \leqq \|x_1\| = 1, |x_1^{-1}(J_0)| = |x_1(J_0)^{-1}| \leqq \|x_1^{-1}\| = 1$, we see that $|x_1(J_0)| = 1$. Hence there exists a point $s_0, 0 \leqq s_0 \leqq 1$, such that $x_1(J_0) = e^{2\pi is_0}$. Thus $x_n = e^{2\pi isn} = x_1^n$ satisfies $x_n(J_0) = e^{2\pi is_n n}$, and so $x(J_0) = \sum\limits_{n=-\infty}^{\infty} c_n e^{2\pi is_n n} = x(s_0)$. In this way, we see that for any maximal ideal $J_0$ of $B$, there exists a point $s_0, 0 \leqq s_0 \leqq 1$, such that the homomorphism $x \to x(J_0)$ is given by $x(J_0) = x(s_0)$, for all $x \in B$. It is also clear that the mapping $x \to x(s_0)$ gives a homomorphism of the algebra $B$ into the complex number field. Therefore we see that the maximal ideal space of $B$ coincides with $\{e^{2\pi is}; 0 \leqq s \leqq 1\}$. 

\begin{corollaryn}[N. \textsc{Wiener}'s theorem] If an absolutely convergent Fourier series $x(s) = \sum\limits_{n=-\infty}^{\infty} c_n e^{2\pi isn}$ does not vanish on $[0,1]$, then the function $1/x(s)$ is also representable as an absolutely convergent Fourier series. For, $x$ does not belong to any of the maximal ideals of the normed ring of the above Example~\hyperref[ex:XI.2.Example.2]{$2$}.

\end{corollaryn}
\begin{example}\label{ex:XI.2.Example.3} We take $B_{1} = C[0,1]$, and define, for $x, y \in B_{1}$, 
\[
\begin{aligned}
(x + y) (s) &= x (s) + y (s), (\alpha x) (s) = \alpha x (s), (x y) (s) =
\\
\int_{0} ^{s} x (s - t) y (t) \dd t \text{ and } \| x \| &= \sup\limits _{s \in [ 0, 1 ]} | x (s) |.
\end{aligned}
\]
$B_{1}$ is a commutative $B$-algebra without unit. Adjoining formally a unit $e$ by the rule $ex = xe = x$, $\|e\| = 1$, the set $B = \{z = \lambda e + x; x \in B_{1}\}$ becomes a normed ring by the operations 
\[
\begin{yoslines}
&(\lambda_{1} e + x _{1}) + (\lambda_{2} e + x _{2}) = (\lambda_{1} + \lambda_{2}) e + (x _{1} + x _{2}), \alpha (\lambda e + x) \\
&= \alpha \lambda e + \alpha x, (\lambda_{1} e + x _{1}) (\lambda_{2} e + x _{2}) = \lambda_{1} \lambda_{2} e + \lambda_{1} x _{2} + \lambda_{2} x _{1} + x _{1} x _{2} \text{ and } \| \lambda e + x \| = | \lambda | + \| x \|.
\end{yoslines}
\]
\end{example}

We have, by induction, 
\[
\left| x ^{2} (s) \right| \leqq M ^{2} s, \left| x ^{3} (s) \right| \leqq M ^{3} \frac {s ^{2}}{2}, \dots , \left| x ^{n} (s) \right| \leqq M ^{n} \frac {s ^{n-1}}{(n - 1) !}, \dots
\]
where $M = \sup\limits_{s \in S} |x(s)| = \|x\|$. Thus every $x \in B_1$ is a generalized nilpotent element of $B$, due to the fact that $\lim\limits_{n \to \infty} (n!)^{1/n} = \infty$. 


\yossection{3}{The Spectral Resolution of Bounded Normal Operators}
\label{sec:XI.3}

Let $X$ be a Hilbert space, and let a system $M$ of bounded normal operators $\in L(X, X)$ satisfy the conditions: 
\begin{align}
T, S& \in M \text{ implies } T S = S T (\text{ commutativity }), \label{eq:XI.3.1} \\
T& \in M \text{ implies } T ^{*} \in M. \label{eq:XI.3.2}
\end{align}

A system $M$ consisting of a bounded normal operator $T \in L(X, X)$ and its adjoint $T^{*}$ surely satisfies \eqref{eq:XI.3.1} and \eqref{eq:XI.3.2}. 

Let $M'$ be the totality of operators $\in L(X, X)$ which commute with every $T \in M$, and let $B = M'' = (M')'$ be the totality of operators $\in L(X, X)$ commutative with every operator $S \in M'$. 

\begin{proposition}\label{prop:XI.3.Proposition.1} Every element of $B$ is a normal operator. $B$ is a normed ring over the complex number field by the operator sum, the operator product, the unit $I$ (the identity operator) and the operator norm $\|T\|$. 
\end{proposition}

\begin{proof}
 $M \subseteq M'$ by \eqref{eq:XI.3.1}, and so $M' \supseteq M''$. Hence $M''' = (M'')' \supseteq M''$ and so $B = M''$ is a commutative ring. The identity operator $I$ belongs to $B$ and is the unit of this algebra $B$. By \eqref{eq:XI.3.2}, we easily see that every operator $\in B$ is normal. Since the multiplication $TS$ and the adjoint formation $T \to T^*$ in the algebra $B$ are continuous with respect to the norm of the operator, it is easy to see that the ring $B$ is complete with respect to the operator norm. 
\end{proof}

\begin{theorem}\label{thm:XI.3.Theorem.1} By the Gelfand representation 
\begin{equation}\label{eq:XI.3.3}
B \ni T \rightarrow T (J),
\end{equation}
the ring $B$ is represented isomorphically by the algebra $C(\{J\})$ of all continuous functions $T(J)$ on the compact space $\{J\}$ of all the maximal 
ideals $J$ of $B$ in such a way that 
\begin{equation}\label{eq:XI.3.4}
\|T\|=\sup_{J\in\{J\}}|T(J)|,
\end{equation}
\begin{equation}\label{eq:XI.3.5}
\begin{minipage}{0.92\linewidth}
\everymath{}\let\limits\nolimits
$T(J)$ is real-valued on $\{J\}$ iff $T$ is self-adjoint,
\end{minipage}
\end{equation}
\begin{equation}\label{eq:XI.3.6}
\begin{minipage}{0.92\linewidth}
\everymath{}\let\limits\nolimits
$T(J)\geqq0$ on $\{J\}$ iff $T$ is self-adjoint and positive, that is, $(Tx,x)\geqq0$ for all $x\in X$.
\end{minipage}
\end{equation}

\end{theorem}

\begin{proof}
 We first show that, for any bounded normal operator $T$, 
\begin{equation}\label{eq:XI.3.7}
\left\| T ^{2} \right\| = \left\| T \right\| ^{2}. \end{equation}

By the normality of $T$, we see that $H = TT^{*} = T^{*}T$ is self-adjoint. Hence, by Theorem~\ref{thm:VII.3.Theorem.3} in Chapter~\ref{chap:VII}, \ref{sec:VII.3}, 
\[
\left\| T \right\| ^{2} = \sup\limits _{\| x \| \leqq 1} (T x, T x) = \sup\limits _{\| x \| \leqq 1} \left| \left(T ^{*} T x, x\right) \right| = \| H \| = \| T ^{*} T \| = \| T T ^{*} \|.
\]

Since $(T^{*})^{2} = (T^{2})^{*}$, $T^{2}$ is normal with $T$. Thus, as above, we obtain $\| T^2\|^2 = \| T^{*2}T^2\|$, which is, by the commutativity $TT^{*} = T^{*}T$, equal to $\| (T^{*}T)^{2}\| = \| H^{2}\|$. Since $H^{2}$ is self-adjoint, we obtain, again by Theorem~\ref{thm:VII.3.Theorem.3} in Chapter~\ref{chap:VII}, \ref{sec:VII.3}, 
\[
\left\| H \right\| ^{2} = \sup\limits _{\| x \| \leqq 1} (H x, H x) = \sup\limits _{\| x \| \leqq 1} \left| \left(H ^{2} x, x\right) \right| = \left\| H ^{2} \right\|.
\]

Therefore, $\|T^2\|^2 = \|H^2\| = \|H\|^2 = (\|T\|^2)^2$, that is, $\|T^2\| = \|T\|^2$. We have, by \eqref{eq:XI.3.7}, $\|T\| = \lim\limits_{n\to \infty} \|T^n\|^{1/n}$, because we know already that the right hand limit exists (see \eqref{eq:VIII.2.3} in Chapter~\ref{chap:VIII}, \ref{sec:VIII.2}). Hence, by Theorem~\hyperref[thm:XI.2.Theorem.4]{$4$} of the preceding section, the representation \eqref{eq:XI.3.3} is isomorphic and \eqref{eq:XI.3.4} is true. 
\end{proof}

\begin{proof}[Proof of \eqref{eq:XI.3.5}.]
Let a self-adjoint $T \in B$ satisfy, for a certain $J_0 \in \{J\}$, $T(J_0) = a + ib$ with $b \neq 0$. Then the self-adjoint operator $S = (T - aI)/b \in B$ satisfies $(I + S^2)(J_0) = 1 + i^2 = 0$, and so $(I + S^2)$ does not have an inverse in $B$. But, by Theorem~\ref{thm:VII.3.Theorem.2} in Chapter~\ref{chap:VII}, \ref{sec:VII.3}, $(I + S^2)$ has an inverse which surely belongs to $B$. Thus, if $T \in B$ is self-adjoint, $T(J)$ must be real-valued. Let $T \in B$ be not self-adjoint, and put
\[
T = \frac {T + T ^{*}}{2} + i \frac {T - T ^{*}}{2 i}.
\]

Then, since the first term on the right is self-adjoint, the self-adjoint operator $(T - T^{*}) / 2i$ must be $\neq 0$. Thus, by the isomorphism of representation \eqref{eq:XI.3.3}, there must exist a $J_0 \in \{J\}$ such that $\frac{T - T^*}{2i}(J_0) \neq 0$. Hence $T(J_0) = \frac{T + T^*}{2}(J_0) + i\frac{T - T^*}{2i}(J_0)$ is not real. For, as proved above, self-adjoint operators are represented by real-valued functions. 
\end{proof}

\begin{proof}[Proof of \eqref{eq:XI.3.6}.]
We first show that
\begin{equation}\label{eq:XI.3.8}
T ^{*} (J) = \overline {{{{T (J)}}}} \text{ on } \{J \}.
\end{equation}

This is clear, since self-adjoint operators $(T + T^{*})/2$ and $(T - T^{*})/2i$ are represented by real-valued functions. Therefore, by \eqref{eq:XI.3.4} and the result of the preceding section, the ring $B$ is represented by the ring of all continuous complex-valued functions on $\{J\}$ satisfying \eqref{eq:XI.3.5} and \eqref{eq:XI.3.8}. Let $T(J) \geqq 0$ over $\{J\}$. Then $S(J) = T(J)^{1/2}$ is a continuous function on $\{J\}$. Hence, by the isomorphism of the representation \eqref{eq:XI.3.3}, $S^{2} = T$. By \eqref{eq:XI.3.5}, we have $S = S^{*}$. Hence $(Tx, x) = (S^{2}x, x) = (Sx, Sx) \geqq 0$. To prove, conversely, that the condition $(Tx, x) \geqq 0$ for all $x \in X$ implies $T(J) \geqq 0$ over $\{J\}$, we set $T_{1}(J) = \max(T(J), 0)$ and $T_{2}(J) = T_{1}(J) - T(J)$. Then, by what we have proved above, $T_{1}$ and $T_{2}$ are both $\in B$, self-adjoint and positive: $(T_{j}x, x) \geqq 0$ for all $x \in X$ ($j = 1$, $2$). Moreover, we have $T_{1}T_{2} = 0$ and $T_{2} = T_{1} - T$. The former equality is implied by $T_{1}(J)T_{2}(J) = 0$. 

Therefore, we have 
$0 \leqq (TT_{2}x, T_{2}x) = (-T_{2}^{2}x, T_{2}x) = -(T_{2}^{3}x, x) = -(T_{2}T_{2}x, T_{2}x) \leqq 0$. Thus $(T_{2}^{3}x, x) = 0$ and so, by Theorem~\ref{thm:VII.3.Theorem.3} in Chapter~\ref{chap:VII}, \ref{sec:VII.3}, we must have $T_{2}^{3} = 0$. Hence, by $\|T_{2}\| = \lim\limits_{n \to \infty} \|T_{2}^{n}\|^{1/n}$, we obtain $T_{2} = 0$. We have thus proved $T = T_{1}$ and hence $T(J) \geqq 0$ on $\{J\}$. 

We shall thus write $T \geqq 0$ if $T$ is self-adjoint and positive. We also write $S \geqq T$ if $(S - T) \geqq 0$. 
\end{proof}

\begin{theorem}\label{thm:XI.3.Theorem.2} Let $\{T_n\} \subseteq B$ be a sequence of self-adjoint operators such that 
\begin{equation}\label{eq:XI.3.9}
0 \leqq T _{1} \leqq T _{2} \leqq \dots \leqq T _{n} \leqq \dots \leqq S \in B.
\end{equation}
\end{theorem}

Then, for any $x \in X$, $\slim\limits_{n \to \infty} T_n x = Tx$ exists, i.e., $\slim\limits_{n \to \infty} T_n = T$ exists and $T \in B$, $S \geqq T \geqq T_n$ ( $n = 1, 2, \ldots$). 

\begin{proof}
 We first remark that, by \eqref{eq:XI.3.6}, 
\begin{equation}\label{eq:XI.3.10}
E, F \in B \text{ and } E \geqq 0, F \geqq 0 \text{ imply } E + F \geqq 0 \text{ and } EF \geqq 0.
\end{equation}
Thus $0 \leqq T_1^2 \leqq T_2^2 \leqq \cdots \leqq T_n^2 \leqq \cdots \leqq S^2$. Hence, for any $x \in X$, a finite $\lim\limits_{n \to \infty} (T_n^2 x, x)$ exists. Since, by \eqref{eq:XI.3.6}, $T_{n+k}^2 \geqq T_{n+k} T_n \geqq T_n^2$, we also have 
\[
\lim\limits _{n, k \rightarrow \infty} (T _{n+k} ^{2} x, x) = \lim\limits _{n, k \rightarrow \infty} (T _{n+k} T _{n} x, x) = \lim\limits _{n \rightarrow \infty} (T _{n} ^{2} x, x). \]

Therefore, $\lim\limits_{n,m\to \infty}((T_n - T_m)^2 x,x) = \lim\limits_{n,m\to \infty}\|T_nx - T_mx\|^2 = 0$ so that $s\text{- }\lim\limits_{n\to \infty}T_nx = Tx$ exists. That $T\in B$ and $S\geqq T\geqq T_n$ is clear from the process of the proof. 
\end{proof}

\begin{theorem}\label{thm:XI.3.Theorem.3} Let a sequence of real-valued functions $\{T_n(J)\}$, where $T_n \in B$, satisfy the condition 
\begin{equation}\label{eq:XI.3.11}
0 \leqq T_{1}(J) \leqq T_{2}(J) \leqq \cdots \leqq T_{n}(J) \leqq \cdots \leqq a \text{ finite constant on } \{J\}.
\end{equation}
\end{theorem}

Then, by \eqref{eq:XI.3.6} and Theorem~\ref{thm:XI.3.Theorem.2}, $\slim\limits_{n \to \infty} T_n = T$ exists. In such a case, we can prove that 
\[
D = \{J \in \{J \}; T (J) \neq \lim\limits _{n \rightarrow \infty} T _{n} (J) \}
\]
is a set of the first category, and so $D^{C} = \{J\} - D$ is dense in $\{J\}$. 

\begin{proof}
 By Theorem~\ref{thm:XI.3.Theorem.2}, $T \geqq T_n$ and so $T(J) \geqq \lim\limits_{n \to \infty} T_n(J)$ on $\{J\}$. By Baire's theorem in Chapter~\ref{chap:0}, \ref{sec:0.2}, the set of points of discontinuity of the function $\lim\limits_{n \to \infty} T_n(J)$ is of the first category. Hence, if the set $D$ is not of the first category, then there exists at least one point $J_0 \in D$ at which $\lim\limits_{n \to \infty} T_n(J)$ is continuous. In other words, there exists a positive number $\delta$ and an open set $V(J_0) \ni J_0$ of $\{J\}$ such that 
\[
T (J) \geqq \delta + \lim\limits _{n \rightarrow \infty} T _{n} (J) \text{ whenever } J \in V (J _{0}). \]

Since the compact space $\{J\}$ is normal, and since $T(J) \geqq \lim\limits_{n \to \infty} T_n(J)$ on $\{J\}$, we may construct, by Urysohn's theorem an open set $V_1(J_0) \ni J_0$ and a function $W(J) \in C(\{J\})$ such that $0 \leqq W(J) \leqq \delta$ on $\{J\}$, $V_1(J_0)^a \subseteq V(J_0)$, $W(J) = \delta/2$ on $V_1(J_0)$ and $W(J) = 0$ on $V(J_0)^C$. Hence $T(J) - W(J) \geqq \lim\limits_{n \to \infty} T_n(J)$ on $\{J\}$, and so, by \eqref{eq:XI.3.6}, $T - W \geqq T_n$ ( $n = 1, 2, \ldots$). $W(J) \not\equiv 0$ implies, by the isomorphism \eqref{eq:XI.3.3}, that $W \neq 0$, $W \geqq 0$. Thus, again by \eqref{eq:XI.3.6}, $T - W \geqq s\text{- }\lim\limits_{n \to \infty} T_n$, contrary to $T = s\text{- }\lim\limits_{n \to \infty} T_n$. Finally, since $\{J\}$ is a compact space, the complement $D^C = \{J\} - D$ of the set $D$ of the first category must be dense in $\{J\}$. 
\end{proof}

We are now able to prove (K. \textsc{Yosida} \cite{yosida355}) 

\yossub{The spectral resolution or the spectral representation of operators $\in B$.}

Consider the set $C'(\{J\})$ of all complex-valued bounded functions $T'(J)$ on $\{J\}$ such that $T'(J)$ is different from a continuous function $T(J)$ only on a set of the first category. We identify two functions from $C'(\{J\})$ if they differ only on a set of the first category. Then $C'(\{J\})$ is divided into classes. Since the complement of a set of the first category is dense in the compact space $\{J\}$, each class $T'$ contains exactly one continuous function $T(J)$ which corresponds, by the isomorphism $B \leftrightarrow C(\{J\})$, to an element $T \in B$. 

For any $T \in B$ and for any complex number $z = \lambda + i\mu$, we put $E_z =$ the element $\in B$ which corresponds to the class $E_z'$ containing the defining function $E_z'(J)$ of the set $\{J \in \{J\}; \operatorname{Re} T(J) < \lambda, \operatorname{Im} T(J) < \mu\}$. It is clear that there exists a monotone increasing sequence of continuous functions $f_n$ of complex argument such that $E_z'(J) = \lim\limits_{n \to \infty} f_n(T(J))$ and 
so $E_{z}^{\prime}(J)\in C^{\prime}(\{J\})$. We have then 
\[
\begin{gathered}
| T (J) - \sum\limits_{j=2} ^{n} (\lambda_{j} + i \mu_{j}) \left(E _ {\lambda_{j} + i \mu_{j}} ^{\prime} (J) + E _ {\lambda_{j-1} + i \mu_{j-1}} ^{\prime} (J) - E _ {\lambda_{j-1} + i \mu_{j}} ^{\prime} (J) \right.
\\
\left. - E _ {\lambda_{j} + i \mu_{j-1}} ^{\prime} (J)\right) | \leqq \varepsilon \text{ on } \{J \} \text{ if }
\end{gathered}
\]
\[
\begin{gathered}
\lambda_{1} = - \alpha - \frac {\varepsilon}{\sqrt {2}} \leqq \lambda_{2} \leqq \dots \leqq \lambda_{n} = \alpha = \sup\limits _ {J \in \{J \}} | \operatorname{Re} T (J) |,
\\
\mu_{1} = - \beta - \frac {\varepsilon}{\sqrt {2}} \leqq \mu_{2} \leqq \dots \leqq \mu_{n} = \beta = \sup\limits _ {J \in \{J \}} | \operatorname{Im} T (J) |,
\\
\left(\sup\limits _{j} (\lambda_{j} - \lambda_{j-1}) ^{2} + \sup\limits _{j} (\mu_{j} - \mu_{j-1}) ^{2}\right) ^{1/2} \leqq \varepsilon .
\end{gathered}
\]

Thus, by the definition of $E_{z}(J)$, we have 
\[
\begin{gathered}
\left| T (J) - \sum\limits_{j=2} ^{n} (\lambda_{j} + i \mu_{j}) \left(E _ {\lambda_{j} + i \mu_{j}} (J) + E _ {\lambda_{j-1} + i \mu_{j-1}} (J) \right. \right. \left. \left. - E _ {\lambda_{j-1} + i \mu_{j}} (J) - E _ {\lambda_{j} + i \mu_{j-1}} (J)\right) \right| \leqq \varepsilon
\end{gathered}
\]
on $\{J\}$, since the complement of a set of the first category is dense in the compact space $\{J\}$. Therefore, by \eqref{eq:XI.3.4}, we have 
\[
\left\| T - \sum\limits_{j=2} ^{n} \left(\lambda_{j} + i \mu_{j}\right) \left(E _ {\lambda_{j} + i \mu_{j}} + E _ {\lambda_{j-1} + i \mu_{j-1}} - E _ {\lambda_{j-1} + i \mu_{j}} - E _ {\lambda_{j} + i \mu_{j-1}}\right) \right\| \leqq \varepsilon ,
\]
which we shall write as 
\begin{equation}\label{eq:XI.3.12}
T = \iint z \dd E _{z},
\end{equation}
and it is called the spectral resolution of the normal operator $T$. 


\yossection{4}{The Spectral Resolution of a Unitary Operator}
\label{sec:XI.4}

If $T$ is a unitary operator $\in B$, then, by 
\begin{equation}\label{eq:XI.4.1}
T (J) T ^{*} (J) = T (J) \overline {{T (J)}} = 1,
\end{equation}
we see that the values taken by the function $T(J)$ on $\{J\}$ are complex numbers of absolute value $1$. From this fact, we can simplify the spectral resolution $\iint z dE_z$ of $T$. 

The defining function $E_{\theta}'(J)$ of the set $\{J \in \{J\}; \arg(T(J)) \in (0, \theta)\}$, $0 < \theta < 2\pi$, belongs to $C'(\{J\})$, and we have, by setting $E_0'(J) = 0$, $E_{2\pi}'(J) = 1$, 
\[
\begin{gathered}
\left| T (J) - \sum\limits_{j=1} ^{n} e ^ {i \theta_{j}} (E _ {\theta_{j}} ^{\prime} (J) - E _ {\theta_{j-1}} ^{\prime} (J)) \right| \leqq \max\limits _{j} \left| e ^ {i \theta_{j}} - e ^ {i \theta_{j-1}} \right| (0 = \theta_{0} < \theta_{1} < \dots < \theta_{n} = 2 \pi).
\end{gathered}
\]

Let $E_{\theta}(J)$ be the continuous function on $\{J\}$ which is different from $E_{\theta}'(J)$ only on a set of the first category, and let $E_{\theta}$ be the operator $\in B$ which corresponds to $E_{\theta}(J)$ by the isomorphic representation $B \ni T \leftrightarrow T(J)$. Then, as in the preceding section, 
\[
\left\| T - \sum\limits_{j=1} ^ {\dot {n}} e ^ {i \theta_{j}} (E _ {\theta_{j}} - E _ {\theta_{j-1}}) \right\| \leqq \max\limits _{j} \left| e ^ {i \theta_{j}} - e ^ {i \theta_{j-1}} \right|.
\]

This we write, in view of the fact $e^{2\pi i} = 1$, 
\begin{equation}T = \int_{0} ^{2 \pi} e ^{i \theta} \dd F (\theta), \text{ where } F (\theta) = E _{\theta +0} - E _{+0} \text{ for } 0 < \theta < 2 \pi , F (0) = 0, F (2 \pi) = I. \label{eq:XI.4.2}\end{equation}

Here $E_{\theta +0}$ is defined by $E_{\theta +0}x = \slim\limits_{\theta \downarrow 0} E_{\theta}x$, the existence of this limit will be proved below. 

\begin{theorem}\label{thm:XI.4.Theorem.1} The system of operators $F(\theta), 0 \leqq \theta \leqq 2\pi$, satisfies the conditions: 
\begin{equation}\label{eq:XI.4.3}
\parbox{0.90\linewidth}{each $F(\theta)$ is a projection operator, commutative with every bounded linear operator commutative with $T$,}
\end{equation}
\end{theorem}

\begin{gather}F (\theta) F \left(\theta^{\prime}\right) = F \left(\min \left(\theta , \theta^{\prime}\right)\right), \label{eq:XI.4.4} \\F (0) = 0, F (2 \pi) = I, \label{eq:XI.4.5} \\F (\theta + 0) = F (\theta), \quad 0 \leqq \theta < 2 \pi , \text{ in the sense that } \slim _ {\theta^{\prime} \downarrow \theta} F (\theta^{\prime}) x = F (\theta) x \text{ for every } x \in X. \label{eq:XI.4.6}
\end{gather}

\begin{proof}
 It is sufficient to prove that $E_{\theta}, 0 \leqq \theta \leqq 2\pi$, satisfies the conditions: 
\begin{align}
\text{ each } E _{\theta} \text{ is a projection operator }& \in B , \tag{3'}\label{eq:XI.4.3'} \\
E _{\theta} E _ {\theta^{\prime}}& = E _ {\min (\theta , \theta^{\prime})}, \tag{4'}\label{eq:XI.4.4'} \\
E _{0}& = 0, E _{2 \pi} = I, \tag{5'}\label{eq:XI.4.5'} \\
E _{\theta} x& = \slim _ {\theta^{\prime} \downarrow \theta} E _ {\theta^{\prime}} x \text{ for every } x \in X \text{ and } 0 \leqq \theta < 2 \pi .  \tag{6'}\label{eq:XI.4.6'} \end{align}

We have $E_{\theta}^{\prime}(J) = \overline{E_{\theta}^{\prime}(J)}$ and $E_{\theta}^{\prime}(J)^{2} = E_{\theta}^{\prime}(J)$. Hence, by the result of the preceding section, we obtain $E_{\theta} = E_{\theta}^{*}$ and $E_{\theta}^{2} = E_{\theta}$. This proves ($3$'). \eqref{eq:XI.4.4} is proved similarly from $E_{\theta}^{\prime}(J)E_{\theta^{\prime}}^{\prime}(J) = E_{\min (\theta ,\theta^{\prime})}^{\prime}(J)$, and \eqref{eq:XI.4.5} is proved similarly. Next let $\theta_{n}\downarrow \theta$. Then $E_{\theta_n}^{\prime}(J)\geqq E_{\theta_{n+1}}^{\prime}(J)\geqq E_{\theta}^{\prime}(J)$, and so, by the result in the preceding section, $s\text{- }\lim\limits_{n\to \infty}E_{\theta_n} = E$ exists and $E(J) = E_{\theta}^{\prime}(J) = \lim\limits_{\theta_{n}\downarrow \theta}E_{\theta_{n}}^{\prime}(J)$ on $\{J\}$ except possibly on a set of the first category. Thus $E = E_{\theta}$. 
\end{proof}

\begin{example}\label{ex:XI.4.Example.1} Let a linear operator $T$ defined by 
\[
T y (s) = e ^{is} y (s), \text{ where } y (s) \in L ^{2} (- \infty , \infty).
\]
$T$ is unitary. We define, when $2n\pi < s \leqq 2(n + 1)\pi$, 
\begin{equation}F (\theta) y (s) = y (s) \quad \text{ for } \quad s \leqq \theta + 2 n \pi \leqq 2 (n + 1) \pi , F (\theta) y (s) = 0 \quad \text{ for } \quad \theta + 2 n \pi < s. \notag\end{equation}
\end{example}

It is easy to see that $T = \int_{0}^{2\pi} e^{i\theta} dF(\theta)$. 


\begin{example}\label{ex:XI.4.Example.2} Let a linear operator $T_{1}$ be defined by 
\[
T _{1} x (t) = x (t + 1) \quad \text{ in } \quad L ^{2} (- \infty , \infty).
\]
$T_{1}$ is unitary. By the Fourier transformation 
\[
y (s) = U x (t) = \underset {n \rightarrow \infty} {\text{ l . i . m . }} (2 \pi) ^{-1/2} \int_{-n} ^{n} e ^{-ist} x (t) \dd t,
\]
we obtain $Ux(t + 1) = e^{is}Ux(t) = e^{is}y(s)$. Thus 
\[
T _{1} x (t) = x (t + 1) = U ^{-1} e ^{is} y (s) = U ^{-1} T y (s) = U ^{-1} T U x (t),
\]
that is, $T_{1} = U^{-1}TU$. Therefore, we have 
\[
T _{1} = \int_{0} ^{2 \pi} e ^{i \theta} \dd F _{1} (\theta), \text{ where } F _{1} (\theta) = U ^{-1} F (\theta) U.
\]
\end{example}

The uniqueness of the spectral resolution. Since $T^{-1} = T^{*}$ and $T^{-1}(J) = T^{*}(J) = T(J)^{-1}$, we easily see that 
\begin{equation}\label{eq:XI.4.7}
T ^{-1} = \int_{0} ^{2 \pi} e ^{- i \theta} \dd F (\theta).
\end{equation}

Let $\max\limits_{j}|e^{i\theta_j} - e^{i\theta_{j-1}}| < \varepsilon$. Then, from 
\[
T = \sum\limits_{j} e ^ {i \theta_{j}} (F (\theta_{j}) - F (\theta_{j-1})) + \delta , \| \delta \| < \varepsilon ,
\]
we obtain, by \eqref{eq:XI.4.4}, 
\[
\begin{gathered}
T ^{2} = \sum\limits_{j} e ^ {2 i \theta_{j}} (F (\theta_{j}) - F (\theta_{j-1})) + \delta^{\prime}, \text{ where } \| \delta^{\prime} \| \leqq \| (T - \delta) \delta \| + \| \delta (T - \delta) \| + \| \delta^{2} \|
\\
\leqq (\| T \| + \varepsilon) \varepsilon + \varepsilon (\| T \| + \varepsilon) + \varepsilon^{2}.
\end{gathered}
\]

Hence we have $T^2 = \int_0^{2\pi} e^{2i\theta} dF(\theta)$, and, more generally, 
\begin{equation}\label{eq:XI.4.8}
T ^{n} = \int_{0} ^{2 \pi} e ^{i n \theta} \dd F (\theta) \quad (n = 0, \pm 1, \pm 2, \dots).
\end{equation}

Therefore, if there exists another spectral resolution $T = \int_{0}^{2\pi} e^{i\theta} dF_1(\theta)$ satisfying ($3$) to \eqref{eq:XI.4.6}, then, for any polynomial $p(\theta)$ in $e^{i\theta}$ and $e^{-i\theta}$, 
\[
\int_{0} ^{2 \pi} p (\theta) d ((F (\theta) x, y) - (F _{1} (\theta) x, y)) = 0, (x, y \in X).
\]

Thus, by continuity, the above equality holds for any continuous function $p(\theta)$ with $p(0) = p(2\pi)$. Let $0 < \theta_{0} < \theta_{1} < 2\pi$, and take 
\[
\begin{aligned}
&p _{n} (\theta)  = 0 \text{ for } 0 \leqq \theta \leqq \theta_{0} \text{ and for } \theta_{1} + \frac {1}{n} \leqq \theta \leqq 2 \pi , = 1 \text{ for } \theta_{0} + \frac {1}{n} \leqq \theta \leqq \theta_{1}, \\
&= \text{ linear for } \theta_{0} \leqq \theta \leqq \theta_{0} + \frac {1}{n} \text{ and for } \theta_{1} \leqq \theta \leqq \theta_{1} + \frac {1}{n}.
\end{aligned}
\]

Then letting $n \to \infty$, we obtain, by \eqref{eq:XI.4.6}, 
\[
\lim\limits _{n \rightarrow \infty} \int_{0} ^{2 \pi} p _{n} (\theta) d [ (F (\theta) x, y) - (F _{1} (\theta) x, y) ] = [ (F (\theta) x, y) - (F _{1} (\theta) x, y) ] _ {\theta_{0}} ^ {\theta_{1}} = 0,
\]
which is valid for all $x, y \in X$. Hence, letting $\theta_{0} \downarrow 0$ and making use of conditions \eqref{eq:XI.4.5} and \eqref{eq:XI.4.6}, we see that $F(\theta_{1}) = F_{1}(\theta)$. Therefore, the spectral resolution for a unitary operator is uniquely determined. 


\yossection{5}{The Resolution of the Identity}
\label{sec:XI.5}

\begin{definition}\label{def:XI.5.Definition.1} $A$ family of projections $E(\lambda), -\infty < \lambda < \infty$, in a Hilbert space $X$ is called a (real) resolution of the identity if it satisfies the conditions: 
\begin{align}
E (\lambda) E (\mu)& = E (\min (\lambda , \mu)), \label{eq:XI.5.1} \\
E (-& \infty) = 0, E (+ \infty) = I, \text{ where } E (- \infty) x = \slim _{\lambda \downarrow - \infty} E (\lambda) x \text{ and } \notag \\
E (+& \infty) x = \slim _{\lambda \uparrow \infty} E (\lambda) x, \label{eq:XI.5.2} \\
E (\lambda + 0)& = E (\lambda), \text{ where } E (\lambda + 0) x = s \underset {\mu \downarrow \lambda} {- \lim} E (\mu) x. \label{eq:XI.5.3}
\end{align}
\end{definition}

\begin{proposition}\label{prop:XI.5.Proposition.1} For any $x, y \in X$, the function $(E(\lambda)x, y)$ is, as a function of $\lambda$, of bounded variation. 
\end{proposition}

\begin{proof}
 Let $\lambda_1 < \lambda_2 < \cdots < \lambda_n$. Then, by \eqref{eq:XI.5.1}, $E(\alpha, \beta] = E(\beta) - E(\alpha)$ is a projection. Thus we have, by Schwarz' inequality, 
\[
\begin{aligned}
&\sum\limits_{j} | (E (\lambda_{j-1}, \lambda_{j} ] x, y) | = \sum\limits_{j} | (E (\lambda_{j-1}, \lambda_{j} ] x, E (\lambda_{j-1}, \lambda_{j} ] y) |
\\
&\quad \leqq \sum\limits_{j} \| E (\lambda_{j-1}, \lambda_{j} ] x \| \cdot \| E (\lambda_{j-1}, \lambda_{j} ] y \| \leqq \left(\sum\limits_{j} \| E (\lambda_{j-1}, \lambda_{j} ] x \| ^{2}\right) ^{1/2}
\\
&\quad \cdot \left(\sum\limits_{j} \| E (\lambda_{j-1}, \lambda_{j} ] y \| ^{2}\right) ^{1/2} = (\| E (\lambda_{1}, \lambda_{n} ] x \| ^{2}) ^{1/2} \cdot (\| E (\lambda_{1}, \lambda_{n} ] y \| ^{2}) ^{1/2} \leqq \| x \| \cdot \| y \|.
\end{aligned}
\]

For, by the orthogonality 
\begin{equation}\label{eq:XI.5.4}
E \left(\lambda_{j-1}, \lambda_{j} \right] \cdot E \left(\lambda_{i-1}, \lambda_{i} \right] = 0 \quad (i \neq j)
\end{equation}
implied by \eqref{eq:XI.5.1}, we have, for $m > n$, 
\begin{equation}\label{eq:XI.5.5}
\left\| x \right\| ^{2} \geqq \left\| E \left(\lambda_{n}, \lambda_{m} \right] x \right\| ^{2} = \sum\limits_{i=n} ^{m-1} \left\| E \left(\lambda_{i}, \lambda_{i+1} \right] x \right\| ^{2}.
\end{equation}
\end{proof}

\begin{corollaryu} For any $\lambda, -\infty < \lambda < \infty$, the operators $E(\lambda + 0) = s\text{- }\lim\limits_{\lambda' \downarrow \lambda} E(\lambda')$ and $E(\lambda - 0) = s\text{- }\lim\limits_{\lambda' \uparrow \lambda} E(\lambda')$ do exist. 
\end{corollaryu}

\begin{proof}
 From \eqref{eq:XI.5.5}, we see that, if $\lambda_{n} \uparrow \lambda$, then 
\[
\lim\limits _{j, k \rightarrow \infty} \| E (\lambda_{j}, \lambda_{k} ] x \| ^{2} = 0,
\]
and the same is true for the case $\lambda_{n} \downarrow \lambda$. 
\end{proof}

\begin{proposition}\label{prop:XI.5.Proposition.2} Let $f(\lambda)$ be a complex-valued continuous function on $(-\infty, \infty)$, and let $x \in X$. Then we can define, for $-\infty < \alpha < \beta < \infty$, 
\[
\int_{\alpha} ^{\beta} f (\lambda) \dd E (\lambda) x
\]
as the \slim of the Riemann sums 
$\sum\limits_{j}f(\lambda_{j}^{\prime})E(\lambda_{j},\lambda_{j+1}]x,\text{ where }\alpha=\lambda_{1}<\lambda_{2}<\cdots<\lambda_{n}=\beta,x_{j}\in(\lambda_{j},\lambda_{j+1}],\text{ when the max }|\lambda_{j+1}-\lambda_{j}| \text{ tends to zero. }$ 
\end{proposition}

\begin{proof}
 $f(\lambda)$ is uniformly continuous on the compact interval $[\alpha, \beta]$. Let $|f(\lambda) - f(\lambda')| \leqq \varepsilon$ whenever $|\lambda - \lambda'| \leqq \delta$. We consider two partitions of $[\alpha, \beta]$: 
\begin{equation}\alpha = \lambda_{1} < \dots < \lambda_{n} = \beta , \max\limits _{j} | \lambda_{j+1} - \lambda_{j} | \leqq \delta , \alpha = \mu_{1} < \dots < \mu_{m} = \beta , \max\limits _{j} | \mu_{j+1} - \mu_{j} | \leqq \delta , \notag\end{equation}
and let 
\[
\alpha = \nu_{1} < \dots < \nu_{p} = \beta , p \leqq m + n,
\]
be the superposition of these two partitions. Then, if $\mu_k' \in (\mu_k, \mu_{k+1}]$, we have 
\[
\begin{aligned}
&\quad \sum\limits_{j} f (\lambda_{j} ^{\prime}) E (\lambda_{j}, \lambda_{j+1} ] x - \sum\limits_{k} f (\mu_{k} ^{\prime}) E (\mu_{k}, \mu_{k+1} ] x = \sum\limits_{s} \varepsilon_{s} E (\nu_{s}, \nu_{s+1} ] x, \text{ with } | \varepsilon_{s} | \leqq 2 \varepsilon ,
\end{aligned}
\]
and so the square of the norm of the left side is, as in \eqref{eq:XI.5.5}, 
\[
\leqq \varepsilon^{2} \left\| \sum\limits_{s} E (\nu_{s}, \nu_{s+1} ] x \right\| ^{2} = \varepsilon^{2} \left\| E (\alpha , \beta ] x \right\| ^{2} \leqq \varepsilon^{2} \left\| x \right\| ^{2}. \]
\end{proof}

\begin{corollaryu} We may define $\int_{\infty}^{\infty}f(\lambda)dE(\lambda)x$ as the $\slim\limits_{\alpha \downarrow -\infty ,\beta \uparrow \infty}\int_{\alpha}^{\beta}f(\lambda)dE(\lambda)x,$ 
when the right side limit exists. 
\end{corollaryu}

\begin{theorem}\label{thm:XI.5.Theorem.1} For a given $x \in X$, the following three conditions are mutually equivalent: 
\begin{align}
&\int_{- \infty} ^{\infty} f (\lambda) \dd E (\lambda) x \text{ exists }, \label{eq:XI.5.6} \\
&\int_{- \infty} ^{\infty} | f (\lambda) | ^{2} \dd \| E (\lambda) x \| ^{2} < \infty , \label{eq:XI.5.7} \\
F (y)& = \int_{- \infty} ^{\infty} f (\lambda) \dd (E (\lambda) y, x) \text{ defines a bounded linear functional. } \label{eq:XI.5.8}
\end{align}
\end{theorem}

\begin{proof}
 We shall prove the implications $\eqref{eq:XI.5.6} \to \eqref{eq:XI.5.8} \to \eqref{eq:XI.5.7} \to \eqref{eq:XI.5.6}$. 


\eqref{eq:XI.5.6} $\rightarrow$ \eqref{eq:XI.5.8}. The scalar product of $y$ with the approximate Riemann sum of $\int_{-\infty}^{\infty} f(\lambda) dE(\lambda)x$ is a bounded linear functional of $y$. Hence, by $(y, E(\lambda)x) = (E(\lambda)y, x)$ and the resonance theorem, we obtain \eqref{eq:XI.5.8}. 

\eqref{eq:XI.5.8} $\rightarrow$ \eqref{eq:XI.5.7}. We apply the operator $E(\alpha, \beta]$ to the approximate Riemann sum of $y = \int_{\alpha}^{\beta} f(\overline{\lambda}) dE(\lambda)x$. We then see, by \eqref{eq:XI.5.1}, that $y = E(\alpha, \beta]y$. Thus, again by \eqref{eq:XI.5.1}, 
\[
\begin{aligned}
&\overline {{F (y)}} = \int_{- \infty} ^{\infty} \overline {{f (\lambda)}} \dd (E (\lambda) x, y) = \lim\limits _ {\alpha^{\prime} \downarrow - \infty , \beta^{\prime} \uparrow \infty} \int_ {\alpha^{\prime}} ^ {\beta^{\prime}} \overline {{f (\lambda)}} \dd (E (\lambda) x, y)
\\
&= \lim\limits _ {\alpha^{\prime} \downarrow - \infty , \beta^{\prime} \uparrow \infty} \int_ {\alpha^{\prime}} ^ {\beta^{\prime}} \overline {{f (\lambda)}} \dd (E (\lambda) x, E (\alpha , \beta ] y) = \lim\limits _ {\alpha^{\prime} \downarrow - \infty , \beta^{\prime} \uparrow \infty} \int_ {\alpha^{\prime}} ^ {\beta^{\prime}} \overline {{f (\lambda)}} \dd (E (\alpha , \beta ] E (\lambda) x, y)
\\
&= \int_{\alpha} ^{\beta} \overline {{f (\lambda)}} \dd (E (\lambda) x, y) = \| y \| ^{2}.
\end{aligned}
\]

Hence $\| y\| ^2 \leqq \| F\| \cdot \| y\|$, i.e., $\| y\| \leqq \| F\|$. On the other hand, by approximating $y = \int_{\alpha}^{\beta}\overline{f(\lambda)} dE(\lambda)x$ by Riemann sums, we obtain, by \eqref{eq:XI.5.1}, 
\[
\left\| y \right\| ^{2} = \left\| \int_{\alpha} ^{\beta} \overline {{f (\lambda)}} \dd E (\lambda) x \right\| ^{2} = \int_{\alpha} ^{\beta} | f (\lambda) | ^{2} d \left\| E (\lambda) x \right\| ^{2},
\]
so that $\int_{\alpha}^{\beta}|f(\lambda)|^2 \dd\|E(\lambda)x\|^2 \leqq \|F\|^2$. Therefore, by letting $\alpha \downarrow -\infty$, $\beta \downarrow \infty$, we see that \eqref{eq:XI.5.7} is true. 

\eqref{eq:XI.5.7} $\rightarrow$ \eqref{eq:XI.5.6}. We have, for $\alpha' < \alpha < \beta < \beta'$, 
\[
\begin{aligned}
&\quad \left\| \int_ {\alpha^{\prime}} ^ {\beta^{\prime}} f (\lambda) \dd E (\lambda) x - \int_{\alpha} ^{\beta} f (\lambda) \dd E (\lambda) x \right\| ^{2} = \int_ {\alpha^{\prime}} ^{\alpha} | f (\lambda) | ^{2} \dd \| E (\lambda) x \| ^{2} + \int_{\beta} ^ {\beta^{\prime}} | f (\lambda) | ^{2} \dd \| E (\lambda) x \| ^{2}
\end{aligned}
\]
as above. Hence \eqref{eq:XI.5.7} implies \eqref{eq:XI.5.6}.
\end{proof}

\begin{theorem}\label{thm:XI.5.Theorem.2} Let $f(\lambda)$ be a real-valued continuous function. Then, a self-adjoint operator $H$ with $D(H) = D$ is defined through 
\begin{equation}\label{eq:XI.5.9}
\begin{gathered}
(H x, y) = \int_{- \infty} ^{\infty} f (\lambda) \dd (E (\lambda) x, y), \text{ where }
\\
x \in D = \left\{x; \int_{- \infty} ^{\infty} | f (\lambda) | ^{2} \dd \| E (\lambda) x \| ^{2} < \infty \right\} \text{ and any } y \in X,
\end{gathered}
\end{equation}
and we have $HE(\lambda) \supseteq E(\lambda)H$, that is, $HE(\lambda)$ is an extension of $E(\lambda)H$. 
\end{theorem}

\begin{proof}
 For any $y \in X$ and for any $\varepsilon > 0$, there exist $\alpha$ and $\beta$ with $-\infty < \alpha < \beta < \infty$ such that $\|y - E(\alpha, \beta] y\| < \varepsilon$. Moreover, we have 
\[
\int_{- \infty} ^{\infty} | f (\lambda) | ^{2} \dd \| E (\lambda) E (\alpha , \beta ] y \| ^{2} = \int_{\alpha} ^{\beta} | f (\lambda) | ^{2} \dd \| E (\lambda) y \| ^{2}. \]

Hence $E(\alpha, \beta] y \in D$ and so, by \eqref{eq:XI.5.2}, $D^a = X$. $H$ is symmetric by 
\[
f (\lambda) = \overline {{f (\lambda)}}, (E (\lambda) x, y) = \overline {{(E (\lambda) y , x)}}.
\]

If $y \in D(H^{*})$ and $H^{*}y = y^{*}$, then, by $E(\alpha, \beta]z \in D$ and \eqref{eq:XI.5.1}, 
\[
(z, E (\alpha , \beta ] y ^{*}) = (E (\alpha , \beta ] z, H ^{*} y) = (H E (\alpha , \beta ] z, y) = \int_{\alpha} ^{\beta} f (\lambda) \dd (E (\lambda) z, y).
\]

Thus, by the resonance theorem, 
\[
\lim\limits _{\alpha \downarrow - \infty , \beta \uparrow \infty} (z, E (\alpha , \beta ] y ^{*}) = \int_{- \infty} ^{\infty} f (\lambda) \dd (E (\lambda) z, y) = F (z)
\]
is a bounded linear functional. Hence, by the preceding theorem, 
\[
\int_{- \infty} ^{\infty} | f (\lambda) | ^{2} \dd \| E (\lambda) y \| ^{2} < \infty , \text{ that is, } y \in D.
\]

Therefore, $D = D(H) \supseteq D(H^*)$. Since $H$ is a symmetric operator, we have $H \subseteq H^*$ and so $H$ must be self-adjoint, i.e., $H = H^*$. 


Finally, let $x \in D(H)$. Then, by applying $E(\mu)$ to the approximate Rie-
mann sums of $Hx = \int_{-\infty}^{\infty} f(\lambda) dE(\lambda)x$, we obtain, by \eqref{eq:XI.5.1}, 
\[
\begin{aligned}
&E (\mu) H x  = \int_{- \infty} ^{\infty} f (\lambda) \dd (E (\mu) E (\lambda) x) = \int_{- \infty} ^{\infty} f (\lambda) \dd (E (\lambda) E (\mu) x) = H E (\mu) x.
\end{aligned}
\]
\end{proof}

\begin{corollary}\label{cor:XI.5.Corollary.1} In the particular case $f(\lambda) = \lambda$, we have 
\begin{equation}\label{eq:XI.5.10}
(H x, y) = \int_{- \infty} ^{\infty} \lambda \dd (E (\lambda) x, y), \quad \text{ for } \quad x \in D (H), y \in X.
\end{equation}
\end{corollary}

We shall write it symbolically 
\[
H = \int_{- \infty} ^{\infty} \lambda \dd E (\lambda),
\]
and call it the spectral resolution or the spectral representation of the self-adjoint operator $H$. 

\begin{corollary}\label{cor:XI.5.Corollary.2} We have, for $H = \int_{-\infty}^{\infty} f(\lambda) dE(\lambda)$ given by \eqref{eq:XI.5.9}, 
\begin{equation}\label{eq:XI.5.11}
\left\| H x \right\| ^{2} = \int_{- \infty} ^{\infty} | f (\lambda) | ^{2} d \left\| E (\lambda) x \right\| ^{2} \quad \text{ whenever } \quad x \in D (H).
\end{equation}
\end{corollary}

In particular, if $H$ is a bounded self-adjoint operator, then 
\begin{equation}\label{eq:XI.5.12}
(H ^{n} x, y) = \int_{- \infty} ^{\infty} f (\lambda) ^{n} \dd (E (\lambda) x, y) \quad \text{ for } \quad x, y \in X (n = 0, 1, 2, \dots).
\end{equation}

\begin{proof}
 Since $E(\lambda) Hx = HE(\lambda)x$ for $x \in D(H)$, we have, by \eqref{eq:XI.5.1}, 
\[ \begin{aligned} (H x, H x) & = \int f (\lambda) \dd (E (\lambda) x, H x) = \int f (\lambda) d (H E (\lambda) x, x) \\ & = \int f (\lambda) d _{\lambda} \left\{\int f (\mu) d _{\mu} (E (\mu) E (\lambda) x, x) \right\} \\ & = \int f (\lambda) d _{\lambda} \left\{\int_{- \infty} ^{\lambda} f (\mu) \dd (E (\mu) x, x) \right\} = \int f (\lambda) ^{2} \dd \| E (\lambda) x \| ^{2}. \end{aligned}
\]


The last part of the Corollary may be proved similarly.
\end{proof}

\begin{examplen} It is easy to see that the multiplication operator
\[
H x (t) = t x (t) \quad \text{ in } \quad L ^{2} (- \infty , \infty)
\]
admits the spectral resolution $H = \int_{-\infty}^{\infty}\lambda dE(\lambda)$, where 
\begin{equation}\label{eq:XI.5.13}
\begin{aligned} E (\lambda) x (t) & = x (t) \quad \text{ for } \quad t \leqq \lambda , \\ & = 0 \quad \text{ for } \quad t > \lambda . \end{aligned}
\end{equation}

For, 
\[
\begin{aligned}
\int_{- \infty} ^{\infty} \lambda^{2} \dd \| E (\lambda) x \| ^{2} &= \int_{- \infty} ^{\infty} \lambda^{2} d _{\lambda} \int_{- \infty} ^{\lambda} | x (t) | ^{2} \dd t = \int_{- \infty} ^{\infty} t ^{2} | x (t) | ^{2} \dd t = \| H x \| ^{2},
\\
\int_{- \infty} ^{\infty} \lambda \dd (E (\lambda) x, y) &= \int_{- \infty} ^{\infty} \lambda d _{\lambda} \int_{- \infty} ^{\lambda} x (t) \overline {{y (t)}} \dd t = \int_{- \infty} ^{\infty} t \cdot x (t) \overline {{y (t)}} \dd t = (H x, y).
\end{aligned}
\]


\end{examplen}
\yossection{6}{The Spectral Resolution of a Self-adjoint Operator}
\label{sec:XI.6}

\begin{theorem}\label{thm:XI.6.Theorem.1} $A$ self-adjoint operator $H$ in a Hilbert space $X$ admits a uniquely determined spectral resolution. 
\end{theorem}

\begin{proof}
 The Cayley transform $U = U_{H} = (H - iI)(H + iI)^{-1}$ of the self-adjoint operator $H$ is unitary (see Chapter~\ref{chap:VII}, \ref{sec:VII.4}). Let $U = \int_{0}^{2\pi} e^{i\theta} dF(\theta)$ be the spectral resolution of $U$. Then we have 
\[
F (2 \pi - 0) = \slim _{\theta \downarrow 0} F (2 \pi - \theta) = F (2 \pi) = I. \]

If otherwise, the projection $F(2\pi)-F(2\pi-0)$ would not be equal to the zero operator. Thus there exists an element $y \neq 0$ such that 
\[
(F (2 \pi) - F (2 \pi - 0)) y = y.
\]

Hence, by $F(\theta)F(\theta') = F(\min (\theta ,\theta'))$, 
\[
\begin{aligned}
&U y  = \int_{0} ^{2 \pi} e ^{i \theta} \dd (F (\theta) (F (2 \pi) - F (2 \pi - 0))) y = (F (2 \pi) - F (2 \pi - 0)) y = y.
\end{aligned}
\]

Thus $(y, z) = (Uy, Uz) = (y, Uz)$ and so $(y, z - Uz) = 0$ for every $z \in X$. $U$ being the Cayley transform of a self-adjoint operator $H$, we know (see Chapter~\ref{chap:VII}, \ref{sec:VII.4}) that the range $R(I - U)$ is dense in $X$. Hence we must have $y = 0$, which is a contradiction. 

Thus, if we set 
\[
\lambda = - \cot \theta , E (\lambda) = F (\theta),
\]
then $0 < \theta < 2\pi$ and $-\infty < \lambda < \infty$ are in a topological correspondence. Hence $E(\lambda)$ is a resolution of the identity with $F(\theta)$. We shall show that the self-adjoint operator 
\[
H ^{\prime} = \int_{- \infty} ^{\infty} \lambda \dd E (\lambda)
\]
is equal to $H$. Since $H = i(I + U)(I - U)^{-1}$, we have only to show that 
\[
(H ^{\prime} (y - U y), x) = (i (y + U y), x) \quad \text{ for all } \quad x, y \in X.
\]

But, since $D(H')^a = X$, we may restrict $x$ to be in the domain $D(H')$. Now, by $F(\theta) \cdot F(\theta') = F(\min(\theta, \theta'))$, 
\[
\begin{aligned}
&(y - U y, F (\theta) x)  = \int_{0} ^{2 \pi} (1 - e ^ {i \theta^{\prime}}) d _ {\theta^{\prime}} (F (\theta^{\prime}) y, F (\theta) x) = \int_{0} ^{2 \pi} (1 - e ^ {i \theta^{\prime}}) d _ {\theta^{\prime}} (F (\theta) F (\theta^{\prime}) y, x) \\
&= \int_{0} ^{\theta} (1 - e ^ {i \theta^{\prime}}) \dd (F (\theta^{\prime}) y, x).
\end{aligned}
\]

Hence 
\[
\begin{aligned}
&(y - U y, H ^{\prime} x)  = \int_{- \infty} ^{\infty} \lambda d (y - U y, E (\lambda) x) = \int_{0} ^{2 \pi} - \cot \theta d \left\{\int_{0} ^{\theta} (1 - e ^ {i \theta^{\prime}}) \dd (F (\theta^{\prime}) y, x) \right\} \\
&= \int_{0} ^{2 \pi} i (1 + e ^{i \theta}) \dd (F (\theta) y, x) = (i (y + U y), x).
\end{aligned}
\]

The uniqueness of the spectral representation. Suppose $H = \int_{-\infty}^{\infty} \lambda dE(\lambda)$ admits another spectral representation $H = \int_{-\infty}^{\infty} \lambda dE'(\lambda)$ such that 
$E^{\prime}(\lambda_{0})\neq E(\lambda_{0})$ for some $\lambda_0$. Then, by setting 
\[
\lambda = - \cot \theta , E ^{\prime} (\lambda) = F ^{\prime} (\theta),
\]
we have $F'(\theta_0) \neq F(\theta_0)$, where $\lambda_0 = -\cot \theta_0$. By a similar calculation to the above, we can prove that the Cayley transform of $\int_{-\infty}^{\infty} \lambda dE'(\lambda)$ is equal to $\int_{0}^{2\pi} e^{i\theta} dF'(\theta)$. Hence the unitary operator $U$ admits two different spectral representations $U = \int_{0}^{2\pi} e^{i\theta} dF(\theta)$ and $U = \int_{0}^{2\pi} e^{i\theta} dF'(\theta)$, contrary to what we have proved in Chapter~\ref{chap:XI}, \ref{sec:XI.4}. 
\end{proof}

We have thus proved (see Chapter~\ref{chap:VII}, \ref{sec:VII.3} and $4$) the fundamental result due to J. \textsc{von} \textsc{Neumann} \cite{neumann228}: 


\begin{theorem}\label{thm:XI.6.Theorem.2} $A$ symmetric operator $H$ has a closed symmetric extension $H^{**}$. A closed symmetric operator $H$ admits a uniquely determined spectral representation iff $H$ is self-adjoint. $H$ is self-adjoint iff its Cayley transform is unitary. 
\end{theorem}

\begin{remarku} It sometimes happens, in applications, that $H$ is not self-adjoint but $H^{*}$ is self-adjoint. In such a case, $H$ is said to be essentially self-adjoint. In this connection, see T. \textsc{Kato} \cite{kato140} concerning Schrödinger's operators in quantum mechanics. 
\end{remarku}

The spectral representation of the momentum operator $H_{1}$: 
\[
H _{1} x (t) = \frac {1}{i} \frac {d}{\dd t} x (t) \quad \text{ in } \quad L ^{2} (- \infty , \infty).
\]

The Fourier transform $U$ defined by 
\[
x (t) = U \cdot y (s) = \underset {n \rightarrow \infty} {\text{ l . i . m . }} (2 \pi) ^{-1/2} \int_{-n} ^{n} e ^{ist} y (s) \dd s
\]
is unitary and $U^{-1}x(t) = U^{*}x(t) = Ux(-t)$. Hence, denoting by $E(\lambda)$ the resolution of the identity given by \eqref{eq:XI.5.13} in Chapter~\ref{chap:XI}, \ref{sec:XI.5}, we obtain a resolution of the identity $\{E'(\lambda)\}, E'(\lambda) = UE(\lambda)U^{-1}$. If both $y(s)$ and $sy(s)$ belong to $L^2(-\infty, \infty) \cap L^1(-\infty, \infty)$, then 
\[
\begin{aligned}
\frac {1}{i} \frac {d}{\dd t} x (t) &= \frac {1}{i} \frac {d}{\dd t} \left((2 \pi) ^{-1/2} \int_{- \infty} ^{\infty} e ^{ist} y (s) \dd s\right)
\\
&= (2 \pi) ^{-1/2} \int_{- \infty} ^{\infty} e ^{ist} s y (s) \dd s = U (s y (s)) = U s U ^{-1} x (t),
\end{aligned}
\]
or, symbolically, 
\begin{equation}\label{eq:XI.6.1}
\frac {1}{i} \frac {d}{\dd t} = U s U ^{-1}.
\end{equation}

Hence, for the self-adjoint operator $H = s \cdot = \int_{-\infty}^{\infty} \lambda dE(\lambda)$, we have 
\begin{equation}\int U ^{-1} H _{1} U \cdot y (s) = s \cdot y (s) = H y (s) \text{ whenever } y (s), L ^{2} (- \infty , \infty) \cap L ^{1} (- \infty , \infty) \notag\end{equation}

For any $y(s) \in D(H) = D(s \cdot)$, set 
\[
y _{n} (s) = y (s) \text{ or } = 0 \text{ according as } | s | \leqq n \text{ or } | s | > n.
\]

Then surely $y_{n}(s)$, $s y_{n}(s)$ both $\in L^{2}(-\infty, \infty) \cap L^{1}(-\infty, \infty)$ and, moreover, \slim $y_{n} = y$, \slim $H y_{n} = H y$. Thus, since the self-adjoint operators $U^{-1}H_{1}U$ and $H$ are closed, we have, by $(U^{-1}H_{1}U)y_{n} = Hy_{n}$, 
\[
\left(U ^{-1} H _{1} U\right) y = H y \quad \text{ whenever } \quad y \in D (H),
\]
that is, $U^{-1}H_{1}U$ is a self-adjoint extension of the self-adjoint operator $H$. Hence, by taking the adjoint, we see that $H^{*}=H$ is also an extension of $(U^{-1}H_{1}U)^{*}=U^{-1}H_{1}U$; consequently $U^{-1}H_{1}U=H$ and so 
\[
H _{1} = U H U ^{-1} = \int_{- \infty} ^{\infty} \lambda d (U E (\lambda) U ^{-1}) = \int_{- \infty} ^{\infty} \lambda \dd E ^{\prime} (\lambda).
\]


\yossection{7}{Real Operators and Semi-bounded Operators. Friedrichs' Theorem}
\label{sec:XI.7}

Real operators and semi-bounded operators, defined below, have self-adjoint extensions. Thus we can apply von Neumann's theorem to these extensions to the effect that they admit spectral resolutions. 

\begin{definition}\label{def:XI.7.Definition.1} Let $X = L^2(S, \mathfrak{B}, m)$ and let $H$ be a symmetric operator defined in $X$ into $X$. $H$ is said to be a real operator, if \begin{enumerate}[i)]
\item $x(s) \in D(H)$ then $\overline{x(s)} \in D(H)$
\item $H$ maps real-valued functions into real-valued functions.
\end{enumerate}
\end{definition}

\begin{examplen} Let $f(s)$ be a real-valued continuous function in $(-\infty, \infty)$. Then, for $X = L^2(-\infty, \infty)$, the operator of multiplication by $f(s)$ is a real operator.

\end{examplen}
\begin{theorem}[J. von Neumann {\cite{neumann228}}]\label{thm:XI.7.Theorem.1} $A$ real operator $H$ admits a self-adjoint extension. 
\end{theorem}

\begin{proof}
 Let $U = U_H$ be the Cayley transform of $H$. Then $D(U) = \{(H + iI)x; x \in D(H)\}$ consists of the functions obtained by taking the complex-conjugate of the functions of $R(U) = \{(H - iI)x; x \in D(H)\}$. Thus, if we define an extension $U_1$ of $U$ through 
\begin{equation}U _{1} = U \text{ in } D (U), U _{1} \left(\sum\limits_{\alpha} c _{\alpha} \varphi_{\alpha}\right) = \sum\limits_{\alpha} c _{\alpha} \bar {\varphi} _{\alpha}, \text{ where } \{\varphi_{\alpha} \} \text{ is a complete } \notag\end{equation}
orthonormal system of the Hilbert space $D(U)^{\perp}$, 
then $U_{1}$ is a unitary extension of $U$. Therefore the self-adjoint extension $H_{1}$ of $H$ exists such that $U_{1} = U_{H_{1}}$ (see Chapter~\ref{chap:VII}, \ref{sec:VII.4}). 
\end{proof}

\begin{definition}\label{def:XI.7.Definition.2} $A$ symmetric operator $H$ is said to be upper semi-bounded (or lower semi-bounded) if there exists a real constant $\alpha$ such that 
\[
(H x, x) \leqq \alpha \| x \| ^{2} (\text{ or } (H x, x) \geqq \alpha \| x \| ^{2}) \text{ for all } x \in D (H)
\]
\end{definition}

If $(H x, x) \geqq 0$ for all $x \in D(H)$, then $H$ is said to be a positive operator. 

\begin{examplen} Let $q(s)$ be continuous and non-negative in $(-\infty, \infty)$. Consider the operator $H$ defined for $C^{2}$ functions $x(s)$ with compact support by
\[
(H x) (s) = - x ^{\prime \prime} (s) + q (s) x (s).
\]

Then $H$ is a positive operator in the Hilbert space $L^{2}(-\infty,\infty)$, as may be verified by partial integration. 

\end{examplen}
\begin{theorem}[K. \textsc{Friedrichs} {\cite{friedrichs71}}]\label{thm:XI.7.Theorem.2} $A$ semi-bounded operator $H$ admits a self-adjoint extension. 
\end{theorem}

Proof (due to H. \textsc{Freudenthal} \cite{freudenthal66}). $-H$ is lower semi-bounded if $H$ is upper semi-bounded. If $H$ is lower semi-bounded as above, then $H_{1} = H + (1 - \alpha)I$ satisfies the condition that $(H_{1}x, x) \geqq \|x\|^{2}$ for all $x \in D(H_{1})$. Therefore, since $\alpha I$ is self-adjoint, we may assume that the symmetric operator $H$ satisfies the condition 
\begin{equation}\label{eq:XI.7.1}
(H x, x) \geqq \| x \| ^{2} \text{ for all } x \in D (H).
\end{equation}

We introduce a new norm $\|x\|'$ and the associated new scalar product $(x, y)'$ in $D(H)$ through 
\begin{equation}\label{eq:XI.7.2}
\left\| x \right\| ^{\prime} = (H x, x), (x, y) ^{\prime} = (H x, y).
\end{equation}

Since $H$ is symmetric and satisfies \eqref{eq:XI.7.1}, it is easy to see that $D(H)$ becomes a pre-Hilbert space with respect to $\|x\|'$ and $(x, y)'$. We denote by $D(H)'$ the completion of this pre-Hilbert space. 

We shall show that $D(H)'$ is, as an abstract set without topology, a subset of the set $X$ which is the original Hilbert space. Proof: A Cauchy sequence $\{x_n\}'$ of the pre-Hilbert space $D(H)$ satisfies $\| x_n - x_m \|' \geqq \| x_n - x_m \|$ and $\lim\limits_{n,m \to \infty} \| x_n - x_m \|' = 0$; consequently $\{x_n\}$ is also a Cauchy sequence of the original Hilbert space $X$. If we can show, for a Cauchy sequence $\{y_n\}$ of $D(H)'$, that 
\begin{equation}\label{eq:XI.7.3}
\lim\limits _{n \rightarrow \infty} \| y _{n} \| ^{\prime} \neq 0 \text{ does not imply } \lim\limits _{n \rightarrow \infty} \| y _{n} \| = 0,
\end{equation}
then the correspondence 
\begin{equation}\label{eq:XI.7.4}
\left\{x _{n} \right\} ^{\prime} \rightarrow \left\{x _{n} \right\}
\end{equation}
is one-one from the Cauchy sequences of $D(H)$ into the Cauchy sequences of $X$. Two Cauchy sequences $\{x_n\}', \{z_n\}'$ of $D(H)$ (of $X$) being identified if $\lim\limits_{n \to \infty} \|x_n - z_n\|' = 0$ (if $\lim\limits_{n \to \infty} \|x_n - z_n\| = 0$). Since $X$ is complete, we may thus identify its Cauchy sequence $\{x_n\}$ with the element $x \in X$ such that $\lim\limits_{n\to \infty}\|x_n - x\| = 0$. Therefore, $D(H)'$ may, as an abstract set without topology, be identified with a subset of $X$ by the correspondence \eqref{eq:XI.7.4}. The proof of \eqref{eq:XI.7.3} is obtained, remembering the continuity of the scalar product in $D(H)'$ and in $X$, as follows. $\lim\limits_{n,m\to \infty}\|x_n - x_m\|' = 0$, $\lim\limits_{n\to \infty}\|x_n\|' = \alpha > 0$ and $\lim\limits_{n\to \infty}\|x_n\| = 0$ imply a contradiction 
\[
\alpha^{2} = \lim\limits _{n, m \rightarrow \infty} (x _{n}, x _{m}) ^{\prime} = \lim\limits _{n, m \rightarrow \infty} (H x _{n}, x _{m}) = \lim\limits _{n \rightarrow \infty} (H x _{n}, 0) = 0.
\]

We next set 
\begin{equation}\label{eq:XI.7.5}
\tilde {D} = D (H ^{*}) \cap D (H) ^{\prime}.
\end{equation}

Since $D(H) \subseteq D(H^{*})$, we must have $D(H) \subseteq \tilde{D} \subseteq D(H^{*})$. Hence we can define an extension $\tilde{H}$ of $H$ by restricting $H^{*}$ to the domain $\tilde{D} = D(\tilde{H})$. We have to show that $\tilde{H}$ is self-adjoint. 

We first show that $\tilde{H}$ is symmetric. Suppose $x, y \in \tilde{D}$; there exist two sequences $\{x_n\}'$, $\{y_n\}'$ of $D(H)$ such that $\|x - x_n\|' \to 0$, $\|y - y_n\|' \to 0$ as $n \to \infty$. Hence, by the continuity of the scalar product in $D(H)'$, we see that a finite $\lim\limits_{n,m \to \infty} (x_n, y_m)' = \lim\limits_{n,m \to \infty} (Hx_ny_m), = \lim\limits_{n,m \to \infty} (x_n, \tilde{H}y_m)$ exists. This limit is equal to 
\[
\lim\limits _{n \rightarrow \infty , m \rightarrow \infty} (H x _{n}, y _{m}) = \lim\limits _{n \rightarrow \infty} (H x _{n}, y) = \lim\limits _{n \rightarrow \infty} (x _{n}, \tilde {H} y) = (x, \tilde {H} y)
\]
and also to 
\[
\lim\limits _{m \rightarrow \infty} \lim\limits _{n \rightarrow \infty} (H x _{n}, y _{m}) = \lim\limits _{m \rightarrow \infty} (x, H y _{m}) = \lim\limits _{m \rightarrow \infty} (\tilde {H} x, y _{m}) = (\tilde {H} x, y).
\]

Hence $\tilde{H}$ is symmetric, that is, $\tilde{H} \subseteq (\tilde{H})^*$. 

Next let $x \in D(H)$, $y \in X$. Then, by 
\[
\left| (x, y) \right| \leqq \left\| x \| \right. \left\| y \right\| \leqq \left\| x \| ^{\prime} \cdot \right. \left\| y \right\|,
\]
we see that $f(x) = (x, y)$ is a bounded linear functional on the pre-Hilbert space $D(H)$. Hence $f(x)$ can, by continuity, be extended to a bounded linear functional on the Hilbert space $D(H)'$. Hence, by F. \textsc{Riesz}' representation theorem as applied to the Hilbert space $D(H)'$, there exists a uniquely determined $y' \in D(H)'$ such that 
\[
f (x) = (x, y) = \left(x, y ^{\prime}\right) ^{\prime} = \left(H x, y ^{\prime}\right) \quad \text{ for all } \quad x \in D (H).
\]

This proves that $y' \in D(H^*)$ and $H^*y' = y$. Hence $y' \in \tilde{D}$ and $\tilde{H}y' = y$. We have thus proved that $R(\tilde{H}) = X$, and so, by the Corollary of Theorem~\ref{thm:VII.3.Theorem.1} in Chapter~\ref{chap:VII}, \ref{sec:VII.3}, $\tilde{H}$ must be self-adjoint. 


\yossection[8. The Spectrum of a Self-adjoint Operator]{8}{The Spectrum of a Self-adjoint Operator. Rayleigh's Principle and the Krylov-Weinstein Theorem. The Multiplicity of the Spectrum}
\label{sec:XI.8}

\begin{theorem}\label{thm:XI.8.Theorem.1} Let $H = \int \lambda dE(\lambda)$ be a self-adjoint operator in a Hilbert space $X$. Let $\sigma(H), P_{\sigma}(H), C_{\sigma}(H)$ and $R_{\sigma}(H)$ be the spectrum, the point spectrum, the continuous spectrum and the residual spectrum of $H$, respectively. Then \begin{enumerate}[(i)]
\item $\sigma(H)$ is a set on the real line;
\item $\lambda_0 \in P_{\sigma}(H)$ is equivalent to the condition $E(\lambda_0) \neq E(\lambda_0 - 0)$ and the eigenspace of $H$ corresponding to the eigenvalue $\lambda_0$ is $R(E(\lambda_0) - E(\lambda_0 - 0))$;
\item $\lambda_0 \in C_{\sigma}(H)$ is equivalent to the condition $E(\lambda_0) = E(\lambda_0 - 0)$ in such a way that $E(\lambda_1) \neq E(\lambda_2)$ whenever $\lambda_1 < \lambda_0 < \lambda_2$;
\item $R_{\sigma}(H)$ is void.
\end{enumerate}
\end{theorem}

\begin{proof}
 We know already that, for a self-adjoint operator $H$, the resolvent set $\varrho(H)$ of $H$ comprises all the complex numbers $\lambda$ with $\operatorname{Im}(\lambda) \neq 0$ (see Chapter~\ref{chap:VIII}, \ref{sec:VIII.1}). Hence (i) is clear. We have $\lambda_0 I = \lambda_0 \int_{-\infty}^{\infty} dE(\lambda)$ by the definition of the resolution of the identity $\{E(\lambda)\}$ and so $(H - \lambda_0 I) = \int_{-\infty}^{\infty} (\lambda - \lambda_0) dE(\lambda)$. Hence, as in Corollary~\hyperref[cor:XI.5.Corollary.2]{$2$} of Theorem~\ref{thm:XI.5.Theorem.2} in Chapter~\ref{chap:XI}, \ref{sec:XI.5}, we obtain 
\begin{equation}\label{eq:XI.8.1}
\left\| (H - \lambda_{0} I) x \right\| ^{2} = \int_{- \infty} ^{\infty} (\lambda - \lambda_{0}) ^{2} d \left\| E (\lambda) x \right\| ^{2}, x \in D (H). \end{equation}

Thus, by $E(-\infty) = 0$ and the right continuity of $\|E(\lambda)x\|^2$ in $\lambda$, we see that $Hx = \lambda_0x$ iff 
\[
\begin{yoslines}
&E (\lambda) x = E (\lambda_{0} + 0) x = E (\lambda_{0}) x \quad \text{ for } \quad \lambda \geqq \lambda_{0} E (\lambda) x = E (\lambda_{0} - 0) x = 0 \quad \text{ for } \quad \lambda < \lambda_{0},
\end{yoslines}
\]
that is, $Hx = \lambda_0 x$ iff $(E(\lambda_0) - E(\lambda_0 - 0))x = x$. This proves (ii). Next we shall prove (iv). If $\lambda_0 \in R_\sigma(H)$, then, by (i), $\lambda_0$ is a real number. By the condition $R(H - \lambda_0 I)^a = D((H - \lambda_0 I)^{-1})^a \neq X$, we see that there exists a $y \neq 0$ which is orthogonal to $R(H - \lambda_0 I)$, i.e., $((H - \lambda_0 I)x, y) = 0$ for all $x \in D(H)$. Hence $(Hx, y) = (\lambda_0 x, y) = (x, \lambda_0 y)$ and so $y \in D(H^*)$, $H^*y = \lambda_0 y$. This proves that $Hy = \lambda_0 y$, i.e., $\lambda_0$ is an eigenvalue of $H$. Hence we have obtained a contradiction $\lambda_0 \in R_\sigma(H) \cap P_\sigma(H)$, and so $R_\sigma(H)$ must be a void set. 

Let $\lambda_0$ be a real number not belonging to $\sigma(H)$. Then the resolvent $(\lambda_0 I - H)^{-1}$ exists. Hence $H_{\lambda_0} = (H - \lambda_0 I)$ has a continuous inverse $(H - \lambda_0 I)^{-1}$. The last condition is, by (iv), equivalent to $\lambda_0 \in \varrho(H)$ and to the condition that there exists a positive number $\alpha$ such that 
\[
\left\| (H - \lambda_{0} I) x \right\| \geqq \alpha \left\| x \right\| \quad \text{ for all } \quad x \in D (H).
\]

This condition is by \eqref{eq:XI.8.1} equivalent to 
\begin{equation}\label{eq:XI.8.2}
\int_{- \infty} ^{\infty} (\lambda - \lambda_{0}) ^{2} \dd \| E (\lambda) x \| ^{2} \geqq \alpha^{2} \| x \| ^{2} \quad \text{ for all } \quad x \in D (H).
\end{equation}

Suppose that we have, for $\lambda_1 < \lambda_0 < \lambda_2$ with $\lambda_0 - \lambda_1 = \lambda_2 - \lambda_0 < \alpha$, $E(\lambda_1) = E(\lambda_2)$. Then, contrary to \eqref{eq:XI.8.2}, we obtain 
\[
\int_{- \infty} ^{\infty} (\lambda - \lambda_{0}) ^{2} \dd \| E (\lambda) x \| ^{2} < \alpha^{2} \int_{- \infty} ^{\infty} \dd \| E (\lambda) x \| ^{2} = \alpha^{2} \| x \| ^{2}.
\]

This proves (iii) by (i), (ii), (iv) and \eqref{eq:XI.8.2}.
\end{proof}

\begin{remarku} The Example in Chapter~\ref{chap:XI}, \ref{sec:XI.5} gives a self-adjoint operator $H$ for which all real numbers are in the continuous spectrum of $H$. 
\end{remarku}

\begin{theorem}\label{thm:XI.8.Theorem.2} Let $H$ be a bounded self-adjoint operator. Then 
\begin{equation}\label{eq:XI.8.3}
\sup\limits _{\lambda \in \sigma (H)} \lambda = \sup\limits _{\| x \| \leqq 1} (H x, x), \inf\limits _{\lambda \in \sigma (H)} \lambda = \inf\limits _{\| x \| \leqq 1} (H x, x).
\end{equation}
\end{theorem}

\begin{proof}
 Since $(Hx, x) = (x, Hx) = (\overline{Hx, x}) = \text{ real number }$, we can consider 
\[
\alpha_{1} = \inf\limits _{\| x \| \leqq 1} (H x, x) \quad \text{ and } \quad \alpha_{2} = \sup\limits _{\| x \| \leqq 1} (H x, x). \]

Let $\lambda_0 \in \sigma(H)$. Then, by Theorem~\ref{thm:XI.8.Theorem.1}, there exists, for every pair $(\lambda_1, \lambda_2)$ of real numbers with $\lambda_1 < \lambda_0 < \lambda_2$, a $y = y_{\lambda_1, \lambda_2} \neq 0$ such that $(E(\lambda_2) - E(\lambda_1))y = y$. We may assume that $\|y\| = 1$. Hence 
\[
\begin{aligned}
&(H y, y)  = \int \lambda \dd (E (\lambda) y, y) = \int \lambda \dd \| E (\lambda) y \| ^{2} = \int \lambda \dd \| E (\lambda) (E (\lambda_{2}) - E (\lambda_{1})) y \| ^{2} \\
&= \int_ {\lambda_{1}} ^ {\lambda_{2}} \lambda \dd \| (E (\lambda) - E (\lambda_{1})) y \| ^{2}.
\end{aligned}
\]

Thus, by letting $\lambda_1 \uparrow \lambda_0$, $\lambda_2 \downarrow \lambda_0$, we obtain $\lim (H y_{\lambda_1, \lambda_2}, y_{\lambda_1, \lambda_2}) = \lambda_0$. This proves that $\sup\limits_{\lambda \in g(H)} \lambda = \sup\limits_{\lambda_0} \leqq \alpha_2$. 

Let us assume that $\alpha_{2} \in \sigma(H)$. Then, by Theorem~\ref{thm:XI.8.Theorem.1}, there exists a pair $(\lambda_{1}, \lambda_{2})$ of real numbers such that $\lambda_{1} < \alpha_{2} < \lambda_{2}$ and $E(\lambda_{2}) = E(\lambda_{1})$. Hence $I = I - E(\lambda_{2}) + E(\lambda_{1})$, $(I - E(\lambda_{2}))E(\lambda_{1}) = E(\lambda_{1})(I - E(\lambda_{2})) = 0$ and so either $(I - E(\lambda_{2}))$ or $E(\lambda_{1})$ is not equal to the zero operator. If $(I - E(\lambda_{2})) \neq 0$, then there exists a $y$ with $\|y\| = 1$, $(I - E(\lambda_{2}))y = y$. In this case, we have 
\[
\begin{aligned}
&(H y, y)  = \int \lambda \dd \| E (\lambda) y \| ^{2} = \int \lambda \dd \| E (\lambda) (I - E (\lambda_{2})) y \| ^{2} = \int_ {\lambda_{1}} ^{\infty} \lambda \dd \| E (\lambda) y \| ^{2} \geqq \lambda_{2} > \alpha_{2};
\end{aligned}
\]
and in the case $E(\lambda_1) \neq 0$, we obtain 
\[
(H z, z) \leqq \lambda_{1} < \alpha_{2} \text{ for a } z \text{ with } \| z \| = 1, E (\lambda_{1}) z = z.
\]

Therefore the assumption $\alpha_{2} \bar{\in} \sigma(H)$ is absurd and so we have proved that $\sup\limits_{\lambda = \sup\limits_{\lambda \in \sigma(H)} H(x, x)} \inf\limits_{\|x\| \leqq 1} \inf\limits_{\lambda \in \sigma(H)} \lambda = \inf\limits_{\|x\| \leqq 1} (Hx, x)$. 
\end{proof}

\begin{theorem}[\textsc{Krylov}-\textsc{Weinstein}]\label{thm:XI.8.Theorem.3} Let $H$ be self-adjoint, and define, for any $x \in D(H)$ with $\| x \| = 1$, 
\begin{equation}\label{eq:XI.8.4}
\alpha_{x} = (H x, x), \beta_{x} = \| H x \|.
\end{equation}
\end{theorem}

Then, for any $\varepsilon > 0$, we can find a $\lambda_0 \in \sigma(H)$ satisfying the inequality 
\begin{equation}\label{eq:XI.8.5}
\alpha_{x} - (\beta_{x} ^{2} - \alpha_{x} ^{2}) ^{1/2} - \varepsilon \leqq \lambda_{0} \leqq \alpha_{x} + (\beta_{x} ^{2} - \alpha_{x} ^{2}) ^{1/2} + \varepsilon .
\end{equation}

\begin{proof}
 We have 
\[
\begin{aligned}
&\beta_{x} ^{2} = (H x, H x) = (H ^{2} x, x) = \int \lambda^{2} \dd \| E (\lambda) x \| ^{2}, \alpha_{x} = (H x, x) = \int \lambda \dd \| E (\lambda) x \| ^{2},
\\
&\| x \| ^{2} = \int \dd \| E (\lambda) x \| ^{2},
\end{aligned}
\]
and so 
\[ \begin{aligned} \beta_{x} ^{2} - \alpha_{x} ^{2} & = \int \lambda^{2} \dd \| E (\lambda) x \| ^{2} - 2 \alpha_{x} \int \lambda \dd \| E (\lambda) x \| ^{2} + \alpha_{x} ^{2} \int \dd \| E (\lambda) x \| ^{2} \\ & = \int (\lambda - \alpha_{x}) ^{2} \dd \| E (\lambda) x \| ^{2}. \end{aligned}
\]

Therefore, if $\left\|E(\lambda)x\right\|^{2}$ does not vary in the interval given by \eqref{eq:XI.8.5}, we would obtain a contradiction 
\[
\beta_{x} ^{2} - \alpha_{x} ^{2} \geqq ((\beta_{x} ^{2} - \alpha_{x} ^{2}) ^{1/2} + \varepsilon) ^{2} > \beta_{x} ^{2} - \alpha_{x} ^{2}.
\]
\end{proof}
\begin{remarku} The so-called Rayleigh principle consists in taking $\alpha_{x}$ as an approximation to the spectrum of the operator $H$. If we calculate $\beta_{x}$, then Theorem~\ref{thm:XI.8.Theorem.3} gives an upper bound of the error when we take $\alpha_{x}$ as an approximation of the spectrum of $H$. For a concrete application of such error estimate, we refer the reader to K. \textsc{Yosida} \cite{yosida344}. 
\end{remarku}

The Multiplicity of the Spectrum. We shall begin with the case of a self-adjoint matrix $H = \int \lambda dE(\lambda)$ in an $n$-dimensional Hilbert space $X_{n}$. Let $\lambda_1, \lambda_2, \ldots, \lambda_p (p \leqq n)$ be the eigenvalues of $H$ with the multiplicity $m_1, m_2, \ldots, m_p$, respectively $\left( \sum\limits_{j=1}^{p} m_j = n \right)$. Let $x_{j_1}, x_{j_2}, \ldots, x_{j_{m_j}}$ be orthonormal eigenvectors of $H$ belonging to the eigenvalue $\lambda_j (Hx_{j_{s}} = \lambda_j x_{j_{s}})$ so that $\{x_{j_{s}}; s = 1, 2, \ldots, m_j\}$ spans the eigenspace $E_{\lambda_j} = R(E(\lambda_j) - E(\lambda_j - 0))$ of $H$ belonging to the eigenvalue $\lambda_j$. Then the set $\{x_{j_{s}}; j = 1, 2, \ldots, p \text{ and } s = 1, 2, \ldots, m_j\}$ is a complete orthonormal system of vectors of the space $X_n$, and hence every vector $y$ of $X_n$ is represented uniquely as the linear combination of $x_{j_{s}}$'s: 
\begin{equation}\label{eq:XI.8.6}
y = \sum\limits_{j=1} ^{p} \sum\limits_{s=1} ^ {m _{j}} x _ {j _{s}} x _ {j _{s}}.
\end{equation}

Thus, denoting by $P_{\lambda_{j}}$ the projection $(E(\lambda_{j})-E(\lambda_{j}-0))$ upon the eigenspace $E_{\lambda_{j}}$, we have, for any $\alpha<\beta$, 
\begin{equation}\label{eq:XI.8.7}
(E (\beta) - E (\alpha)) y = \sum\limits_ {\alpha < \lambda_{j} \leqq \beta} \left(\sum\limits_{s=1} ^ {m _{j}} \alpha_ {j _{s}} x _ {j _{s}}\right) = \sum\limits_ {\alpha < \lambda_{j} \leqq \beta} P _ {\lambda_{j}} y,
\end{equation}
\begin{equation}\label{eq:XI.8.8}
P_{\lambda_j}(E(\beta) - E(\alpha))y = \sum\limits_{s=1}^{m_j}\alpha_{j_{s}}x_{j_{s}} \text{ or } = 0 \text{ according as } \alpha < \lambda_j\leqq \beta \text{ or not.}
\end{equation}

Hence, for fixed $\alpha < \beta$ and a fixed linear subspace $M$ of $X$, the set 
\[
\{(E (\beta) - E (\alpha)) y; y \in M \}
\]
does not contain $E_{\lambda_j}$ if the dimension of $M$, $\dim(M)$, is $< m_j$. Moreover, for a suitable $M$ with $\dim(M) = m_j$, the set $\{(E(\beta) - E(\alpha))y; y \in M\}$ with $\alpha < \lambda_j \leqq \beta$ contains $E_{\lambda_j}$. In fact, the statement is true for $M$ containing $x_{j_1}, x_{j_2}, \ldots, x_{j_{m_j}}$. In particular, $m_1 = m_2 = \cdots = m_p = 1$ with $p = n$ iff there exists a fixed vector $y \in X_n$ such that the set of vectors 
\[
\{(E (\beta) - E (\alpha)) y; \alpha < \beta \}
\]
spans the whole space $X_{n}$. 

These considerations lead to the following definitions: 

\begin{definition}\label{def:XI.8.Definition.1} The spectrum of a self-adjoint operator $H = \int \lambda dE(\lambda)$ in a Hilbert space $X$ is said to be simple, if there exists a fixed vector $y \in X$ such that the set of vectors $\{(E(\beta) - E(\alpha))y; \alpha < \beta\}$ spans a linear subspace strongly dense in $X$. 
\end{definition}

\begin{definition}\label{def:XI.8.Definition.2} Let $H = \int \lambda dE(\lambda)$ be a self-adjoint operator in a Hilbert space $X$. For fixed $\alpha < \beta$, consider linear subspaces $M$ of $(E(\beta) - E(\alpha)) \cdot X$ such that 
\begin{equation}\label{eq:XI.8.9}
(E (\beta) - E (\alpha)) \cdot M = (E (\beta) - E (\alpha)) \cdot X.
\end{equation}
\end{definition}

We may take, for example, $M = (E(\beta) - E(\alpha)) \cdot X$ to meet condition \eqref{eq:XI.8.9}. The minimum of the set of values $\dim(M)$, where $M$ satisfies condition \eqref{eq:XI.8.9}, will be called the total multiplicity of the spectrum of $H$ contained in the interval $(\alpha, \beta]$. 

\begin{definition}\label{def:XI.8.Definition.3} The multiplicity of the spectrum of a self-adjoint operator $H = \int \lambda dE(\lambda)$ at $\lambda = \lambda_0$ is defined as the limit, as $n \to \infty$, of the total multiplicity of the spectrum of $H$ contained in the interval $(\lambda_0 - n^{-1}, \lambda_0 + n^{-1}]$. 
\end{definition}

\begin{examplen} The coordinate operator $H$, i.e., the operator $H$ defined by $H \cdot x(t) = t \cdot x(t)$ in $L^2(-\infty, \infty)$, is of the simple spectrum.

\end{examplen}
\begin{proof}
 We know that the spectral resolution $H = \int \lambda dE(\lambda)$ is given by 
\[
E (\lambda) x (t) = x (t) \text{ or } = 0 \text{ according as } t \leqq \lambda \text{ or } t > \lambda . \]

Let $y(t)$ be defined by 
\[
y (t) = c _{k} > 0 \text{ for } k - 1 < t \leqq k \quad (k = 0, \pm 1, \pm 2, \dots)
\]
with $\sum\limits_{k} c_k^2 < \infty$ so that $y(t) \in L^2(-\infty, \infty)$. 

Then it is easy to see that the linear combinations of vectors of the form $(E(\beta)-E(\alpha))y,\alpha<\beta$, are strongly dense in the totality of step functions with compact support and consequently strongly dense in $L^{2}(-\infty,\infty)$. 
\end{proof}

The Problem of the Unitary Equivalence. Two self-adjoint operators $H_{1}$ and $H_{2}$ in an $n$-dimensional Hilbert space $X_{n}$ are said to be unitarily equivalent to each other if there exists a unitary matrix $U$ in $X_{n}$ such that $H_{1}=UH_{2}U^{-1}$. It is well-known that $H_{1}$ and $H_{2}$ are unitarily equivalent to each other iff they have the same system of eigenvalues with respectively the same multiplicities. Thus the eigenvalues together with the respective multiplicities are the unitary invariants of a self-adjoint matrix. 

The investigation of the unitary invariants for a self-adjoint operator in an infinite dimensional Hilbert space goes back to a paper by E. \textsc{Hellinger} \cite{hausdorff85} published in $1909$. See, e.g., M. H. \textsc{Stone} \cite{stone311}. There the Hilbert space is assumed to be separable. For the non-separable Hilbert space case, see F. \textsc{Wecken} \cite{wecken333} and H. \textsc{Nakano} \cite{nakano224} and also P. R. \textsc{Halmos} \cite{halmos93}. K. \textsc{Yosida} \cite{yosida356} proved the following theorem: 

Let $H$ be a self-adjoint operator in a Hilbert space $X$, and let us denote by $(H)'$ the totality of bounded linear operators $\in L(X, X)$ which are commutative with $H$. Then two self-adjoint operators $H_1$ and $H_2$ in $X$ are unitarily equivalent to each other iff there exists a one-to-one mapping $T$ of $(H_1)'$ onto $(H_2)'$ such that $T$ defines a ring-isomorphism of the ring $(H_1)'$ with the ring $(H_2)'$ in such a way that $(T \cdot B)^* = T \cdot B^*$ for every $B \in (H_1)'$. 

Thus the algebraic structure of the ring $(H_{\mathbf{1}})'$ is the unitary invariant of $H_{\mathbf{1}}$. 


\yossection{9}{The General Expansion Theorem. A Condition for the Absence of the Continuous Spectrum}
\label{sec:XI.9}

Let $H = \int \lambda dE(\lambda)$ be a self-adjoint operator in a Hilbert space $X$. Then, by $E(+\infty) = I$ and $E(-\infty) = 0$, we have the representation 
\begin{equation}\label{eq:XI.9.1}
x = \underset {\lambda_{1} \downarrow - \infty , \lambda_{2} \uparrow \infty} {\slim} \int_ {\lambda_{1}} ^ {\lambda_{2}} \dd E (\lambda) x = \underset {\lambda_{1} \downarrow - \infty , \lambda_{2} \uparrow \infty} {\slim} (E (\lambda_{2}) - E (\lambda_{1})) x \text{ for every } x \in X.
\end{equation}

We shall call \eqref{eq:XI.9.1} the general expansion theorem associated with the self-adjoint operator $H$. In concrete cases, it sometimes happens that the resolvent $(\lambda I - H)^{-1}$ is obtainable more easily than the spectral resolution $H = \int \lambda dE(\lambda)$. In such a case, the general expansion theorem \eqref{eq:XI.9.1}  
may be replaced by 
\begin{equation}\tag{1'}\label{eq:XI.9.1'}
\begin{aligned}
&x = \underset {\alpha \downarrow - \infty , \beta \uparrow \infty} {\slim} s - \underset {v \downarrow 0} {\lim} \frac {1}{2 \pi i} \left[ \int_{\alpha} ^{\beta} ((u - i v) I - H) ^{-1} x \dd u \right. \left. + \int_{\beta} ^{\alpha} ((u + i v) I - H) ^{-1} x \dd u \right], \quad x \in X.
\end{aligned}
\end{equation}

\begin{proof}
 If $v \neq 0$, then 
\[
((u + i v) I - H) ^{-1} x = \int_{- \infty} ^{\infty} \frac {1}{u + i v - \lambda} \dd E (\lambda) x, \text{ for every } x \in X. \]

For, by approximating the integral by Riemann sums and remembering the relation $E(\lambda)E(\lambda') = E(\min (\lambda ,\lambda'))$, we obtain, for $\operatorname{Im}(\mu)\neq 0$, 
\[
\begin{aligned}
&\int_{- \infty} ^{\infty} (\lambda - \mu) \dd E (\lambda) \left\{\int_{- \infty} ^{\infty} \frac {1}{\lambda^{\prime} - \mu} \dd E (\lambda^{\prime}) \right\} = \int_{- \infty} ^{\infty} (\lambda - \mu) d _{\lambda} \left\{\int_{- \infty} ^{\infty} \frac {1}{(\lambda^{\prime} - \mu)} d _ {\lambda^{\prime}} (E (\lambda) E (\lambda^{\prime})) \right\}
\\
&= \int_{- \infty} ^{\infty} (\lambda - \mu) d _{\lambda} \left\{\int_{- \infty} ^{\lambda} \frac {1}{\lambda^{\prime} - \mu} \dd E (\lambda^{\prime}) \right\} = \int_{- \infty} ^{\infty} \dd E (\lambda) = I,
\\
&\int_{- \infty} ^{\infty} \frac {1}{\lambda - \mu} d _{\lambda} \left\{E (\lambda) \int_{- \infty} ^{\infty} (\lambda^{\prime} - \mu) \dd E (\lambda^{\prime}) \right\} = \int_{- \infty} ^{\infty} \frac {1}{\lambda - \mu} d _{\lambda} \left\{\int_{- \infty} ^{\infty} (\lambda^{\prime} - \mu) d _ {\lambda^{\prime}} (E (\lambda) E (\lambda^{\prime})) \right\}
\\
&= \int_{- \infty} ^{\infty} \frac {1}{\lambda - \mu} d _{\lambda} \left\{\int_{- \infty} ^{\lambda} (\lambda^{\prime} - \mu) \dd E (\lambda^{\prime}) \right\} = \int_{- \infty} ^{\infty} \dd E (\lambda) = I.
\end{aligned}
\]

Therefore, we have 
\[
\begin{aligned}
& \int_{\alpha} ^{\beta} ((u - i v) I - H) ^{-1} x \dd u + \int_{\beta} ^{\alpha} ((u + i v) I - H) ^{-1} x \dd u
\\
&= \int_{- \infty} ^{\infty} \dd E (\lambda) x \left\{\int_{\alpha} ^{\beta} \frac {\dd u}{u - i v - \lambda} + \int_{\beta} ^{\alpha} \frac {\dd u}{u + i v - \lambda} \right\}
\\
&= \int_{- \infty} ^{\infty} \dd E (\lambda) x \left\{\int_{\alpha} ^{\beta} d _{u} \log (u - i v - \lambda) + \int_{\beta} ^{\alpha} d _{u} \log (u + i v - \lambda) \right\}
\end{aligned}
\]
which tends, when $v \downarrow 0$, strongly to 
\[
\begin{aligned}
&\int_{\alpha +0} ^{\beta -0} 2 \pi i \dd E (\lambda) x + \pi i (E (\beta) - E (\beta - 0)) x + \pi i (E (\alpha) - E (\alpha - 0)) x
\\
&= \pi i (E (\beta) + E (\beta - 0)) x - \pi i (E (\alpha) + E (\alpha - 0)) x.
\end{aligned}
\]

This proves formula \eqref{eq:XI.9.1}. 
\end{proof}

\begin{remarku} The eigenfunction expansion associated with the second order differential operator 
\[
- \frac {d ^{2}}{\dd x ^{2}} + q (x)
\]
\end{remarku}

with real continuous $q(x)$ in an open interval $(a, b)$ was inaugurated by H. \textsc{Weyl} \cite{weyl337}, further developed by M. H. \textsc{Stone} \cite{stone311}, and completed by E. C. \textsc{Titchmarsh} \cite{titchmarsh324} and K. \textsc{Kodaira} \cite{kodaira151} who gave a formula which determines the expansion explicitly. The expansion is exactly a concrete application of \eqref{eq:XI.9.1}. The crucial point in their theory is to give the possible boundary conditions at $x = a$ and $x = b$ so that the operator 
\[
- \frac {d ^{2}}{\dd x ^{2}} + q (x)
\]
becomes a self-adjoint operator $H$ in the Hilbert space $L^{2}(a,b)$. Their theory is very important in that it gives a unified treatment of the classical expansions in terms of special functions, such as the Fourier series expansion, the Fourier integral representation, the Hermite polynomials expansion, the Laguerre polynomials expansion and the Bessel functions expansion. We do not go into details, and refer the reader to the above cited book by \textsc{Titchmarsh} and the paper by \textsc{Kodaira}. Cf. also N.A.\textsc{Naimark} \cite{naimark222}, N.\textsc{Dunford}-J.\textsc{Schwartz} \cite{dunford48} and K.\textsc{Yosida} \cite{yosida344}. The last cited book gives an elementary treatment of the theory. 

If the continuous spectrum $C_{\sigma}(H)$ is absent, then the expansion \eqref{eq:XI.9.1} will be replaced by a series rather than the integral. We have, for instance, the following 

\begin{theorem}\label{thm:XI.9.Theorem.1} Let $H = \int \lambda dE(\lambda)$ be a self-adjoint compact operator in a Hilbert space $X$. Then \begin{enumerate}[(i)]
\item $C_{\sigma}(H)$ contains no real number except possibly $0$;
\item the eigenvalues of $H$ constitute at most a countable system of real numbers accumulating only at $0$;
\item for any eigenvalue $\lambda_{0} \neq 0$ of $H$, the corresponding eigenspace $E_{\lambda_{0}}$ is of finite dimension.
\end{enumerate}
\end{theorem}

\begin{proof}
 Suppose a closed interval $[\lambda', \lambda'']$ on the real line does not contain the number $0$. Then the range $R(E(\lambda'') - E(\lambda'))$ is of finite dimension. If otherwise, there exists, by E. \textsc{Schmidt}'s orthogonalization in Chapter~\ref{chap:III}, \ref{sec:III.5}, a countable orthonormal system $\{x_j\}$ contained in $R(E(\lambda'') - E(\lambda'))$. We have \wlim $x_j = 0$, since, by Bessel's inequality, 
\[
\sum\limits_{j=1} ^{\infty} | (f, x _{j}) | ^{2} \leqq \| f \| ^{2} \quad \text{ for any } \quad f \in X. \]

Hence, by the compactness of the operator $H$, there exists a subsequence $\{x_{j'}\}$ such that $\slim\limits_{j \to \infty} H x_{j'} = w - \lim\limits_{j \to \infty} H x_{j'} = 0$. On the other hand, we have 
\[
\begin{aligned}
&\| H x _{i} \| ^{2}  = \int \lambda^{2} \dd \| E (\lambda) x _{j} \| ^{2} = \int \lambda^{2} \dd \| E (\lambda) (E (\lambda^{\prime \prime}) - E (\lambda^{\prime})) x _{j} \| ^{2} \\
&= \int_ {\lambda^{\prime}} ^ {\lambda^{\prime \prime}} \lambda^{2} \dd \| (E (\lambda) - E (\lambda^{\prime})) x _{j} \| ^{2} \geqq \| x _{j} \| ^{2} \cdot \min (| \lambda^{\prime} | ^{2}, | \lambda^{\prime \prime} | ^{2}) = \min (| \lambda^{\prime} | ^{2}, | \lambda^{\prime \prime} | ^{2}),
\end{aligned}
\]
which is a contradiction. 

If $C_{\sigma}(H)$ contains a number $\lambda_0 \neq 0$, then, by Theorem~\ref{thm:XI.9.Theorem.1} in the preceding section, $\slim\limits_{\varepsilon \downarrow 0} (E(\lambda_0 - \varepsilon) - E(\lambda_0 - \varepsilon)) x = 0$ for any $x \in X$. As proved above, the range $R(E(\lambda_0 + \varepsilon) - E(\lambda_0 - \varepsilon))$ is of finite dimension and this dimension number is monotone decreasing as $\varepsilon \downarrow 0$. Hence we see that $(E(\lambda_0 + \varepsilon) - E(\lambda_0 - \varepsilon)) = 0$ for sufficiently small $\varepsilon > 0$, and so, by Theorem~\ref{thm:XI.9.Theorem.1} in the preceding section, $\lambda_0$ cannot be contained in $C_{\sigma}(H)$. This proves our Theorem. 
\end{proof}

\begin{corollary}\label{cor:XI.9.Corollary.1} Let $\{\lambda_j\}$ be the system of all eigenvalues of $H$ different from $0$. Then, for any $x \in X$, we have 
\begin{equation}\label{eq:XI.9.2}
x = (E (0) - E (0 - 0)) x + \slim _{n \rightarrow \infty} \sum\limits_{j=1} ^{n} (E (\lambda_{j}) - E (\lambda_{j} - 0)) x.
\end{equation}
\end{corollary}

\begin{proof}
 Clear from \eqref{eq:XI.9.1}. 
\end{proof}

\begin{corollary}[The Hilbert-Schmidt Expansion Theorem]\label{cor:XI.9.Corollary.2} For any $x \in X$, we have 
\begin{equation}\label{eq:XI.9.3}
H x = \underset {n \rightarrow \infty} {\text{ \slim }} \sum\limits_{j=1} ^{n} \lambda_{j} (E (\lambda_{j}) - E (\lambda_{j} - 0)) x.
\end{equation}
\end{corollary}

\begin{proof}
 Clear from the continuity of the operator $H$ and the fact that $H(E(0)-E(0-0))=0$ and $H(E(\lambda_{j})-E(\lambda_{j}-0))x=\lambda_{j}(E(\lambda_{j})-E(\lambda_{j}-0))x$, the latter being implied by $R(E(\lambda_{j})-E(\lambda_{j}-0))=E_{\lambda_{j}}$, the eigenspace of $H$ belonging to the eigenvalue $\lambda_{j}$. 
\end{proof}

\begin{remarku} The strong convergence in \eqref{eq:XI.9.3} is uniform in the unit circle $\{x; \| x \| \leqq 1\}$. For, we have 
\[
\begin{aligned}
\left\| H x - \int_{| \lambda | > \varepsilon} \lambda \dd E (\lambda) x \right\| ^{2} &= \left\| \int_{| \lambda | \leqq \varepsilon} \lambda \dd E (\lambda) x \right\| ^{2} = \int_{- \varepsilon} ^{\varepsilon} \lambda^{2} \dd \| E (\lambda) x \| ^{2}
\\
\quad \leqq \varepsilon^{2} \int_{- \varepsilon} ^{\varepsilon} \dd \| E (\lambda) x \| ^{2} \leqq \varepsilon^{2} \int \dd \| E (\lambda) x \| ^{2} &= \varepsilon^{2} \| x \| ^{2} \leqq \varepsilon^{2}.
\end{aligned}
\]
\end{remarku}


\yossection{10}{The Peter-Weyl-Neumann Theory}
\label{sec:XI.10}

Let $G$ be a totally bounded topological group, metrized by a distance satisfying the condition (see Chapter~\ref{chap:VIII}, \ref{sec:VIII.5}) 
\begin{equation}\label{eq:XI.10.1}
\dist (x, y) = \dist (a x b, a y b) \text{ for every } x, y, a \text{ and } b \in G.
\end{equation}

Let $f(g)$ be a complex-valued bounded uniformly continuous function defined on $G$. For any $\varepsilon > 0$, set 
\begin{equation}\label{eq:XI.10.2}
V = \left\{y \in G; \sup\limits _{x \in G} | f (x) - f (y ^{-1} x) | < \varepsilon \right\}.
\end{equation}

Then, by the continuity of $f$, the set $V$ is an open set containing the unit $e$ of the group $G$. Hence the set $U = V \cap V^{-1}$, where $V^{-1} = \{y^{-1}; y \in V\}$, is also open and $\ni e$. If we put 
\begin{align}
k (x)& = 2 ^{-1} \left(k _{1} (x) + k _{1} \left(x ^{-1}\right)\right), \text{ where } \notag \\
k _{1} (x)& = \frac {\dist \left(x , U ^{c}\right)}{\dist (x , e) + \dist \left(x , U ^{c}\right)} \quad \left(\dist \left(x, U ^{c}\right) = \inf\limits _ {y \in U ^{c}} \dist (x, y)\right), \label{eq:XI.10.3}
\end{align}
then we have the results: 

$k(x)$ is bounded uniformly continuous on $G$, and 
\begin{equation}k (x) = k \left(x ^{-1}\right), \quad 0 \leqq k (x) \leqq 1 \text{ on } G, \quad k (e) = 1 \text{ and } k (x) = 0 \text{ whenever } x \in U ^{C}. \label{eq:XI.10.4}\end{equation}

Hence, for all $x, y \in G$, we obtain 
\[
\left| k (y) \left(f (x) - f \left(y ^{-1} x\right)\right) \right| \leqq \varepsilon k (y).
\]

By taking the mean value (see Chapter~\ref{chap:VIII}, \ref{sec:VIII.5}) of both sides, we obtain 
\[
\left| M _{y} (k (y)) f (x) - M _{y} (k (y) f \left(y ^{-1} x\right)) \right| \leqq \varepsilon M _{y} (k (y)).
\]

We have $M_y(k(y)) > 0$ by $k(y) \geqq 0$ and $k(y) \not\equiv 0$. Hence 
\begin{equation}\label{eq:XI.10.5}
\left| f (x) - M _{y} \left(k _{0} (y) f \left(y ^{-1} x\right)\right) \right| \leqq \varepsilon , \quad \text{ where } \quad k _{0} (x) = k (x) / M _{x} (k (x)).
\end{equation}

Thus, by virtue of the invariance $M_y(g(y^{-1})) = M_y(g(y)) = M_y(g(ay)) = M_y(g(ya))$ of the mean value, we obtain, from \eqref{eq:XI.10.5}, 
\begin{equation}\label{eq:XI.10.6}
\left| f (x) - M _{y} \left(k _{0} \left(x y ^{-1}\right) f (y)\right) \right| \leqq \varepsilon \quad \text{ for all } \quad x \in G.
\end{equation}

\begin{proposition}\label{prop:XI.10.Proposition.1} We shall denote by $C(G)$ the set of all complex-valued bounded uniformly continuous functions $h(g)$ defined on $G$. Then $C(G)$ is a $B$-space by the norm $\|h\|_0 = \sup |h(g)|$. Then, for any $b$ and $h \in C(G)$, 
\begin{equation}\label{eq:XI.10.7}
(b \times h) (x) = M _{y} (b (x y ^{-1}) h (y))
\end{equation}
also belongs to $C(G)$. 
\end{proposition}

\begin{proof}
 By $\dist(x,z)=\dist(axc,azc)$ and the uniform continuity of the function $b(g)$, there exists, for any $\delta>0$, an $\eta=\eta(\delta)>0$ such that 
\[
\sup\limits _{y} \left| b (x y ^{-1}) - b \left(x ^{\prime} y ^{-1}\right) \right| \leqq \delta \quad \text{ whenever } \quad \dist (x, x ^{\prime}) \leqq \eta . \]

Hence, as in the case of Schwarz' inequality, we obtain 
\begin{equation}\label{eq:XI.10.8}
\begin{aligned}
& \left| M _{y} (b (x y ^{-1}) h (y)) - M _{y} (b (x ^{\prime} y ^{-1}) h (y)) \right| ^{2}
\\
&\leqq M _{y} ((b (x y ^{-1}) - b (x ^{\prime} y ^{-1})) ^{2}) \cdot M _{y} (| h (y) | ^{2}) \leqq \delta^{2} M _{y} (| h (y) | ^{2})
\end{aligned}
\end{equation}
whenever $\dist(x, x') \leqq \eta$. This proves our Proposition~\ref{prop:XI.10.Proposition.1}. 
\end{proof}

\begin{proposition}\label{prop:XI.10.Proposition.2} $C(G)$ is a pre-Hilbert space by the operation of the function sum and the scalar product 
\begin{equation}\label{eq:XI.10.9}
(b, h) = (b \times h ^{*}) (e) = M _{y} (b (y ^{-1}) \overline {{h (y ^{-1})}}) = M _{y} (b (y) \overline {{h (y)}}),
\end{equation}
$h^{*}(y) = \overline{h(y^{-1})}$. We shall denote this pre-Hilbert space by $\hat{C}(G)$. 
\end{proposition}

\begin{proof}
 Easy. 
\end{proof}

\begin{proposition}\label{prop:XI.10.Proposition.3} We shall denote the completion of the pre-Hilbert space $\hat{C}(G)$ by $\tilde{C}(G)$, and the norm in the Hilbert space $\tilde{C}(G)$ by $\|h\| = (h, h)^{1/2}$. Then the linear mapping $T$ on $\hat{C}(G)$ into $\hat{C}(G)$ defined by 
\begin{equation}\label{eq:XI.10.10}
(T h) (x) = \left(k _{0} \times h\right) (x), x \in G,
\end{equation}
can, by the continuity in $\hat{C}(G)$, be extended to a compact linear operator $\tilde{T}$ on $\tilde{C}(G)$ into $\tilde{C}(G)$. 
\end{proposition}

\begin{proof}
 By virtue of $M_y(1) = 1$, we obtain 
\begin{equation}\label{eq:XI.10.11}
\left\| h \right\| = (h, h) ^{1/2} = M _{y} (h (y) \overline {{{{h (y)}}}}) ^{1/2} \leqq \sup\limits _{y} | h (y) | = \left\| h \right\| _{0}. \end{equation}

The continuity of the operator $T$ in $\hat{C}(G)$ is clear from the Schwarz inequality for \eqref{eq:XI.10.7}, and so, by the denseness of $\hat{C}(G)$ in $\tilde{C}(G)$, we can extend $T$ to a bounded linear operator $\tilde{T}$ in $\tilde{C}(G)$. On the other hand, we see, by \eqref{eq:XI.10.8}, that $T$ is a compact operator on $\hat{C}(G)$ into $C(G)$. We prove this by the Ascoli-Arzelà theorem. Thus, by \eqref{eq:XI.10.11}, we easily see that $\tilde{T}$ is a compact operator on $\hat{C}(G)$ into $\hat{C}(G)$. 

Therefore, by the denseness of $\hat{C}(G)$ in $\tilde{C}(G)$, we see that the extended operator $\tilde{T}$ is also compact as an operator on $\tilde{C}(G)$ into $\tilde{C}(G)$. 
\end{proof}

We are now ready to prove the Peter-Weyl-Neumann theory on the representation of almost periodic functions. 

It is easy to see that, by $k_{0}(xy^{-1}) = k_{0}(yx^{-1})$, the compact operator $\tilde{T}$ is self-adjoint in the Hilbert space $\tilde{C}(G)$. Hence, by the Hilbert-Schmidt expansion theorem in the preceding section, we obtain 
\begin{equation}\label{eq:XI.10.12}
\tilde {T} h = \slim _{n \to \infty} \sum\limits_{m=1} ^{n} \lambda_{m} P _ {\lambda_{m}} h \text{ uniformly in } h \text{ satisfying } \| h \| \leqq 1,
\end{equation}
if we denote by $\{\lambda_{m}\}$ the system of all eigenvalues of $\tilde{T}$ different from $0$, and by $P_{\lambda_{m}}$ the projection upon the eigenspace of $\tilde{T}$ belonging to the eigenvalue $\lambda_{m}$. 

Since $f(g)$, introduced at the beginning of this section, belongs to $C(G)$, we have $\tilde{T}f = Tf \in C(G)$. Since the eigenspace $R(P_{\lambda_m}) = P_{\lambda_m} \cdot \tilde{C}(G)$ is of finite dimension, there exists, for every eigenvalue $\lambda_m$, a finite system $\{h_{mj}\}_{j=1,\ldots,n_m}$ of elements $\in \tilde{C}(G)$ such that each $h \in R(P_{\lambda_m}) =$ 
$P_{\lambda_m} \cdot \tilde{C}(G)$ can be represented as a uniquely determined linear combination of $h_{mj}$'s ( $j = 1, 2, \ldots, n_m$). Let 
\begin{equation}\label{eq:XI.10.13}
P _ {\lambda_{m}} h = \sum\limits_{j=1} ^ {n _{m}} c _{j} h _{mj}, \text{ where } c ^{\prime} \text{ s are complex numbers. }
\end{equation}

Then, since $h_{mj} \in R(P_{\lambda_m})$, we have $\tilde{T}h_{mj} = \lambda_m h_{mj}$, and so, by applying \eqref{eq:XI.10.8} to the operator $T$ given by \eqref{eq:XI.10.10}, we see that $h_{mj} = \lambda_m^{-1}(\tilde{T}h_{mj})$ must belong to $C(G)$. Hence, by \eqref{eq:XI.10.13}, we see that, for each eigenvalue $\lambda_m$ of $\tilde{T}$, the eigenspace $R(P_{\lambda_m}) = P_{\lambda_m} \cdot \tilde{C}(G)$ is spanned by the functions $h_{mj} \in C(G)$. 

Therefore, we have, by \eqref{eq:XI.10.12}, 
\[
(Tf)(x)=\slim_{n\to\infty}\sum_{m=1}^n\lambda_m f_m(x)
\]
in the strong topology of $\tilde C(G)$, where $f_m=P_{\lambda_m}\cdot f\in C(G)$ for each $m$.

By applying \eqref{eq:XI.10.8} to the operator $T$ given by \eqref{eq:XI.10.10}, we see that 
\begin{equation}\label{eq:XI.10.14}
(T ^{2} f) (x) = \lim\limits _{n \rightarrow \infty} M _{y} \left(k _{0} (x y ^{-1}) \cdot \sum\limits_{m=1} ^{n} \lambda_{m} f _{m} (y)\right) \text{ uniformly in } x.
\end{equation}

On the other hand, by \eqref{eq:XI.10.6} and $M_y(k_0(y)) = 1$, we obtain 
\[
\left| M _{z} \left(k _{0} (x z ^{-1}) f (z)\right) - M _{z} \left(k _{0} (x z ^{-1}) M _{y} \left(k _{0} (z y ^{-1}) f (y)\right) \right. \right| \leqq \varepsilon ,
\]
and so, combined with \eqref{eq:XI.10.6}, we have 
\begin{equation}\label{eq:XI.10.15}
\left| f (x) - M _{z} \left(k _{0} \left(x z ^{-1}\right) M _{y} \left(k _{0} \left(z y ^{-1}\right) f (y)\right)\right) \right| \leqq 2 \varepsilon .
\end{equation}

The left hand side is precisely $|f(x) - (T^2 f)(x)| \leqq 2\varepsilon$. 

Since $T \cdot R(P_{\lambda m}) \subseteq R(P_{\lambda m})$, we have proved 

\begin{theorem}\label{thm:XI.10.Theorem.1} The function $f(x)$ can be approximated uniformly on $G$ by linear combinations of the eigenfunctions of $\tilde{T}$ belonging to the eigenvalues which are different from $0$. 
\end{theorem}

We shall take a fixed eigenvalue $\lambda \neq 0$ of $\tilde{T}$, and shall denote the base $\{h_j\} \subseteq C(G)$ of the corresponding eigenspace $P_{\lambda} \cdot \tilde{C}(G)$ by $e_1(x)$, $e_2(x), \ldots, e_k(x)$. Then, by the invariance of the mean value, we obtain, for any $a \in G$, 
\[
\begin{aligned} M _{y} \left(k _{0} (x y ^{-1}) e _{j} (y a)\right) & = M _{y} \left(k _{0} (x a \cdot a ^{-1} y ^{-1}) e _{j} (y a)\right) = M _{z} \left(k _{0} (x a \cdot z ^{-1}) e _{j} (z)\right) \\ & = \left(T e _{j}\right) (x a) = \lambda e _{j} (x a). \end{aligned}
\]

Since the left hand side is equal to the result of applying the operator $T$ upon the function $e_j(ya)$ of $y$, we see that, for any given $a \in G$, the function $e_j(xa)$ of $x$ must be uniquely represented as a linear combination of the functions $e_1(x), e_2(x), \ldots, e_k(x)$. We have thus 
\begin{equation}\label{eq:XI.10.16}
e _{j} (x a) = \sum\limits_{i=1} ^{k} d _{ji} (a) e _{i} (x) \quad (j = 1, 2, \dots , k),
\end{equation}
or, in vectorial notation, 
\begin{equation}\tag{16'}\label{eq:XI.10.16'}
e (x a) = D (a) e (x).
\end{equation}

By $e(x \cdot ab) = D(ab)e(x), e(xa \cdot b) = D(b)e(xa) = D(b)D(a)e(x)$, we see, remembering the linear independence of $e_1(x), e_2(x), \ldots, e_n(x)$, that 
\begin{equation}\label{eq:XI.10.17}
D (a b) = D (b) D (a), D (e) = \text{ the unit matrix of degree } k.
\end{equation}

By applying E. \textsc{Schmidt}'s orthogonalization, we may assume that $\{e_j(x)\}$ constitutes an orthonormal system in $\tilde{C}(G)$. Then, by \eqref{eq:XI.10.16}, 
\begin{equation}\label{eq:XI.10.18}
M _{x} (e _{j} (x a) e _{i} ^{*} (x)) = d _{ji} (a)
\end{equation}
and so the elements $d_{ji}(a)$ of the matrix $D(a)$ belong to $C(G)$. By the invariance of the mean value, we see that 
\[
M _{y} \left(e _{j} (y a) \overline {{e _{i} (y a)}}\right) = M _{y} \left(e _{j} (y) \overline {{e _{i} (y)}}\right) = \delta_{ij}.
\]

Hence the matrix $D(a)$ gives a linear mapping transforming the orthonormal system $\{e_j(x)\}$ onto the orthonormal system $\{e_j(xa)\}$. Therefore $D(a)$ must be a unitary matrix. Hence the tranposed matrix $D(a)'$ of $D(a)$ is also unitary and we have 
\begin{equation}\tag{17'}\label{eq:XI.10.17prime}
D (a b) ^{\prime} = D (a) ^{\prime} D (b) ^{\prime}, D (e) ^{\prime} = \text{ the unit matrix of degree } k.
\end{equation}

Hence $D(a)'$ gives a unitary matrix representation of the group $G$ such that its matrix elements are continuous functions of $a$. Letting $x = e$ in \eqref{eq:XI.10.16}, we see that each $e_j(a)$ is a linear combination of the matrix elements of the representation $D(a)'$. 

We have thus proved 

\begin{theorem}[\textsc{Peter}-\textsc{Weyl}-\textsc{Neumann}]\label{thm:XI.10.Theorem.2} Let $G$ be a totally bounded topological group, metrized by a distance satisfying $\text{ dis }(x, y) = \text{ dis }(axb, ayb)$. Let $f(g)$ be any complex-valued bounded uniformly continuous function defined on $G$. Then, $f(g)$ can be approximated uniformly on $G$ by linear combinations of the matrix elements of unitary uniformly continuous matrix representations $D(g)'$ of the group $G$. 
\end{theorem}

Referring to A. Weil's reduction given in Chapter~\ref{chap:VIII}, \ref{sec:VIII.5}, we obtain the following 

\begin{corollaryu} Let $G$ be an abstract group, and $f(g)$ an almost periodic function on $G$. Then $f(g)$ can be approximated uniformly on $G$ by linear combinations of the matrix elements of unitary matrix representations $D(g)'$ of the group $G$. 
\end{corollaryu}

\begin{remarkn}[Remark $1$] Let the degree of a unitary matrix representation $D(g)'$ of the group $G$ be $d$, i.e., the degree of the matrix $D(g)'$ be $d$. Then each $D(g)'$ gives a linear mapping of a fixed $d$-dimensional complex Hilbert space $X_d$ onto itself. The representation $D(g)'$ is said to be irreducible if there is no proper linear subspace $\neq \{0\}$ of $X_d$ which is invariant by applying the mappings $D(g)'$, $g \in G$. Otherwise, the representation 
$D(g)'$ is called reducible, and there exists a proper linear subspace $X_{d,1} \neq \{0\}$ invariant by every $D(g)'$, $g \in G$. Then the orthogonal complement $X_{d,1}^{\perp}$ of $X_{d,1}$ in $X_d$ is, by the unitarity of the representation $D(g)'$, also invariant by every $D(g)$, $g \in G$. If we take for the base of the linear space $X_d$ the orthonormal system of vectors composed of one orthonormal system of vectors of $X_{d,1}$ and one of $X_{d,1}^{\perp}$, then, by this choice of an orthonormal base of $X_d$, the representation $D(g)'$ will be transformed into the form 
\[
UD(g)'U^{-1}=\begin{pmatrix}D_1(g)'&0\\0&D_2(g)'\end{pmatrix},
\]
where $U$ is a fixed unitary matrix.
\end{remarkn}

Hence a reducible unitary representation $D(g)'$ is completely reducible into the sum of two unitary representations $D_1(g)'$ and $D_2(g)'$ of the group $G$ acting respectively upon $X_{d,1}$ and $X_{d,1}'$. In this way, we finally can choose a fixed unitary matrix $U_d$ such that the representation $U_dD(g)'U_d^{-1}$ is the sum of irreducible unitary representations of the group $G$. Therefore, in the statements of Theorem~\ref{thm:XI.10.Theorem.2} and its Corollary, we can impose the condition that the matrix representations $D(g)'$ are all irreducible. 

\begin{remarkn}[Remark $2$] In the particular case when $G$ is the additive group of real numbers, a unitary irreducible representation $D(g)'$ is given by 
\begin{equation}\label{eq:XI.10.19}
D (g) ^{\prime} = e ^{i \alpha g}, \text{ where } \alpha \text{ is a real number and } i = \sqrt {- 1}.
\end{equation}
\end{remarkn}

For, by the commutativity of the unitary matrices $D(g)'$, $g \in G$, the representation $D(g)'$ is completely reducible into the sum of one-dimensional unitary representations $\chi(g)$, that is, complex-valued solutions $\chi(g)$ of 
\begin{equation}\label{eq:XI.10.20}
\chi \left(g _{1} + g _{2}\right) = \chi \left(g _{1}\right) \cdot \chi \left(g _{2}\right), | \chi (g _{1}) | = 1 (g _{1}, g _{2} \in G), \chi (0) = 1.
\end{equation}

It is well known that any continuous solution of \eqref{eq:XI.10.20} is of the form $\chi(g) = e^{i\alpha g}$. Hence, any continuous almost periodic function $f(g)$ on the additive group $G$ of real numbers can be approximated uniformly on $G$ by linear combinations of $e^{i\alpha g}$ with real $\alpha$'s. This constitutes the fundamental theorem in H. \textsc{Bohr}'s theory of almost periodic functions. According to the original definition by Bohr, a continuous function $f(x)$ ( $-\infty < x < \infty$) is called almost periodic if for each $\varepsilon > 0$ there exists a positive number $p = p(\varepsilon)$ such that any interval of the form $(t, t + p)$ contains at least one $\tau$ such that 
\[
\left| f (x + \tau) - f (x) \right| \leqq \varepsilon \quad \text{ for } \quad - \infty < x < \infty .
\]

See H. \textsc{Bohr} \cite{bohr25}. S. \textsc{Bochner} \cite{bochner23} has shown that a continuous function $f(x) (-\infty < x < \infty)$ is almost periodic in Bohr's sense iff the following condition is satisfied: For any sequence of real numbers $\{a_n\}$, the system of functions $\{f_{a_{n}}(x); f_{a_{n}}(x) = f(x + a_n)\}$ is totally bounded in the topology of the uniform convergence on $(-\infty, \infty)$. It was extended by J. von \textsc{Neumann} \cite{neumann231} to almost periodic functions in a group. Neumann's result contains as a special case, the Peter-Weyl theory of continuous representation of a compact Lie group (\textsc{Peter}-\textsc{Weyl} \cite{peter247}). According to our treatment, Bohr's result is easily proved by observing that $\lim |s - t| = 0$ implies $\lim \left[\sup\limits_{a,b} |f(asb) - f(atb)|\right] = 0$. 


\yossection{11}{Tannaka's Duality Theorem for Non-commutative Compact Groups}
\label{sec:XI.11}

Let $G$ be a compact (topological) group. This means that $G$ is a compact topological space as well as a group in such a way that the mapping 
\[
(\stackrel {\bullet} {x}, y) \rightarrow x y ^{-1}
\]
of the product space $G \times G$ onto $G$ is continuous. 

\begin{proposition}\label{prop:XI.11.Proposition.1} $A$ complex-valued continuous function $f(g)$ defined on a compact group $G$ is uniformly continuous in the following sense: 
\begin{equation}\label{eq:XI.11.1}
\parbox{0.92\linewidth}{\everymath{}\let\limits\nolimits for any $\varepsilon > 0$, there exists a neighbourhood $U(e)$ of the unit element $e$ of $G$ such that $|f(x) - f(y)| < \varepsilon$ whenever $xy^{-1} \in U(e)$ and also whenever $x^{-1}y \in U(e)$.}
\end{equation}
\end{proposition}


\begin{proof}
 Since $f(x)$ is continuous at every point $a \in G$, there exists a neighbourhood $V_{a}$ of $a$ such that $x \in V_{a}$ implies $|f(x) - f(a)| < \varepsilon / 2$. If we denote by $U_{a}$ the neighbourhood of $e$ defined by $U_{a} = V_{a}a^{-1} = \{va^{-1}; v \in V\}$, then $xa^{-1} \in U_{a}$ implies $|f(x) - f(a)| < \varepsilon / 2$. Let us denote by $W_{a}$ the neighbourhood of $e$ such that $W_{a}^{2} \subseteq U_{a}$, where $W_{a}^{2} = \{w_{1}w_{2}; w_{i} \in W_{a}(i = 1,2)\}$. Obviously, the system of all open sets of the form $W_{a} \cdot a$, where $a$ is an arbitrary element of $G$, covers the whole space $G$. $G$ being compact, there exists a finite set $\{a_{i}; i = 1,2,\ldots,n\}$ such that the system of open sets $W_{a_{i}} \cdot a_{i} (i = 1,2,\ldots,n)$ covers $G$. We denote by $U(e)$ the intersection of all open sets of the system $\{W_{a_{i}}\}$. Then $U(e)$ is a neighbourhood of $e$. We shall show that if $xy^{-1} \in U(e)$, then $|f(x) - f(y)| \leqq \varepsilon$. Since the system $W_{a_{i}} \cdot a_{i}$ covers $G$, there exists a number $k$ such that $ya_{k}^{-1} \in W_{a_{k}} \subseteq U_{a_{k}}$ and therefore $|f(y) - f(a_{k})| < \varepsilon / 2$. Furthermore we have $xa_{k}^{-1} = xy^{-1}ya_{k}^{-1} \in U(e) W_{a_{k}} \subseteq W_{a_{k}}^{2} \subseteq U_{a_{k}}$ so that $|f(x) - f(a_{k})| < \varepsilon / 2$. Combining these two inequalities we get $|f(x) - f(y)| < \varepsilon$. 

If we start with a neighbourhood $U_{a}$ of $e$ such that $x^{-1}a \in U_{a}$ implies $|f(x) - f(a)| < \varepsilon / 2$, we would obtain a neighbourhood $U(e)$ of $e$ such that $|f(x) - f(y)| < \varepsilon$ whenever $x^{-1}y \in U(e)$. Thus taking the intersection of these two $U(e)$'s as the $U(e)$ in the statement of the Proposition, we complete the proof. 
\end{proof}

\begin{corollaryu} $A$ complex-valued continuous function $f(x)$ on a compact group $G$ is almost periodic on $G$. 
\end{corollaryu}

\begin{proof}
 Let $U(e)$ be the neighbourhood of $e$ given in Proposition~\ref{prop:XI.11.Proposition.1}. For any $a \in G$, $U(e)a$ is a neighbourhood of $a$. The compact space $G$ is covered by the system of open sets $U(e)a$, $a \in G$, and therefore some finite subsystem $\{U(e)a_i; i = 1, 2, \ldots, n\}$ covers $G$. That is, for any $a \in G$, there exists some $a_k$ with $k \leqq n$ such that $aa_k^{-1} \in U(e)$. Hence, by $(ax)(a_kx)^{-1} = aa_k^{-1}$, we have $\sup\limits_{x}|f(ax) - f(a_kx)| < \varepsilon$. Similarly, we can find a finite system $\{b_j, j = 1, 2, \ldots, m\}$ such that, for any $b \in G$, there exists some $b_j$ with $j \leqq m$ satisfying the inequality 
\[
\sup\limits _{i,x} | f (a _{i} x b) - f (a _{i} x b _{j}) | < \varepsilon . \]

Therefore, for any pair $a, b$ of elements $\in G$, we can find $a_{k}$ and $b_{j}$ ( $k \leqq n, j \leqq m$) such that 
\[
\sup\limits _{x} | f (a x b) - f \left(a _{k} x b _{j}\right) | < 2 \varepsilon .
\]

This proves that the system of functions $\{f_{a,b}(x); f_{a,b}(x) = f(axb), a \text{ and } b \in G\}$ is totally bounded by the maximum norm $\|h\| = \sup\limits_x |h(x)|$. Hence $f(x)$ is almost periodic on $G$. 
\end{proof}

We are now able to extend the Peter-Weyl-Neumann Theory of the preceding section to complex-valued continuous functions $f(x)$ on a compact group $G$. For such a function $f(x)$ and $\varepsilon > 0$, we set 
\[
V = \left\{y \in G; \sup\limits _{x} | f (x) - f (y ^{-1} x) | < \varepsilon \right\}.
\]

Then, by the continuity of $f$, the set $V$ is an open set containing $e$. By Urysohn's theorem as applied to the normal space $G$, we see that there exists a continuous function $k_{1}(x)$ defined on $G$ such that 
\[
0 \leqq k (x) \leqq 1 \text{ on } G, k _{1} (e) = 1 \quad \text{ and } \quad k _{1} (x) = 0 \text{ whenever } x \in V ^{C}.
\]

Then the continuous function 
\begin{equation}\label{eq:XI.11.2}
k (x) = 2 ^{-1} \left(k _{1} (x) + k _{1} \left(x ^{-1}\right)\right)
\end{equation}
satisfies the condition $k(x^{-1}) = k(x)$ and 
\[
\begin{yoslines}
&0 \leqq k (x) \leqq 1 \text{ on } G, k (e) = 1 \text{ and } k (x) = 0 \text{ whenever } x \in U ^{C}, \text{ where } U = V \cup V ^{-1}.
\end{yoslines}
\]

Therefore, if we denote by $C(G)$ the $B$-space of all complex-valued continuous functions $h(x)$ defined on $G$ normed by the maximum norm, we can define a linear operator $T$ on $\hat{C}(G)$ into $\hat{C}(G)$ by 
\begin{equation}\label{eq:XI.11.3}
(T h) (x) = \left(k _{0} \times h\right) (x), x \in G.
\end{equation}

Here, as in the preceding section, 
\begin{equation}\label{eq:XI.11.4}
k _{0} (x) = k (x) / M _{x} (k (x))
\end{equation}
and $\hat{C}(G)$ is the space $C(G)$ endowed with the scalar product 
\begin{equation}\label{eq:XI.11.5}
(b, h) = M _{y} (b (y) h (y)) = (b \times h ^{*}) (e), h ^{*} (y) = \overline {{{h (y ^{-1})}}}.
\end{equation}

Thus, as in the preceding section, we obtain 

\begin{theorem}\label{thm:XI.11.Theorem.1} Any complex-valued continuous function $f(g)$ defined on a compact group $G$ is almost periodic, and $f(g)$ can be approximated uniformly on $G$ by linear combinations of the matrix elements of unitary, continuous, irreducible matrix representations of $G$. 
\end{theorem}

We shall say that two matrix representations $A_{1}(g)$ and $A_{2}(g)$ of a group $G$ are equivalent if there exists a fixed non-singular matrix $B$ such that $B^{-1}A_{1}(g)B = A_{2}(g)$ for all $g \in G$. 

\begin{proposition}[I. \textsc{Schur}'s lemma]\label{prop:XI.11.Proposition.2} If the representations $A_{1}(g), A_{2}(g)$ are irreducible and inequivalent, then there is no matrix $B$ such that 
\begin{equation}\label{eq:XI.11.6}
A _{1} (g) B = B A _{2} (g)
\end{equation}
holds identically in $g$, except $B = 0$. In \eqref{eq:XI.11.6}, the matrix $B$ is assumed to be of $n_1$ rows and $n_2$ columns, where $n_1, n_2$ are the degrees of $A_1(g), A_2(g)$, respectively. 
\end{proposition}

\begin{proof}
 Let $X_{1}$ and $X_{2}$ be the linear spaces subject to the linear transformations $A_{1}(g)$ and $A_{2}(g)$, respectively. $B$ in \eqref{eq:XI.11.6} can be interpreted as a linear mapping $x_{2} \to x_{1} = Bx_{2}$ of $X_{2}$ onto $X_{1}$. The linear subspace of $X_{1}$ consisting of all vectors $x_{1}$ of the form $Bx_{2}$ is invariant, for $A_{1}(g)x_{1} = Bx_{2}'$ with $x_{2}' = A_{2}(g)x_{2}$. By virtue of the irreducibility of $A_{1}(g)$, there are only two possibilities: either $Bx_{2} = 0$ for all $x_{2}$ in $X_{2}$, i.e. $B = 0$, or $X_{1} = BX_{2}$. On the other hand, the set of all vectors $x_{2}$ in $X_{2}$ such that $Bx_{2} = 0$ is an invariant subspace of $X_{2}$, for $BA_{2}(g)x_{2} = A_{1}(g)Bx_{2} = 0$. From the irreducibility of $A_{2}(g)$ we conclude: either $Bx_{2} = 0$ for all $x_{2}$ in $X_{2}$, i.e. $B = 0$, or $x_{2} = 0$ is the only vector in $X_{2}$ such that $Bx_{2} = 0$, so that distinct vectors in $X_{2}$ go into distinct vectors in $X_{1}$ under the linear mapping $B$. Hence if $B \neq 0$ we conclude that $B$ defines a one-one linear mapping of $X_{2}$ onto $X_{1}$. But this means that $B$ is a non-singular matrix ( $n_1 = n_2$) and so $A_{1}(g)$ and $A_{2}(g)$ would be equivalent. 
\end{proof}

\begin{proposition}[the orthogonality relations]\label{prop:XI.11.Proposition.3} Let $A_{1}(g) = (a_{ij}^{1}(g))$ and $A_{2}(g) = (a_{kl}^{2}(g))$ be unitary, continuous, irreducible matrix representations of a compact group $G$. Then we have the orthogonality relations: 
\begin{align}
M _{g} (a _{ij} ^{1} (g) \overline {{a _{kl} ^{2} (g)}})& = 0, \text{ if } A _{1} (g) \text{ is inequivalent to } A _{2} (g), \label{eq:XI.11.7} \\
M _{g} (a _{ij} ^{1} (g) \overline {{a _{kl} ^{1} (g)}})& = n _{1} ^{-1} \delta_{ik} \delta_{jl} \text{ where } n _{1} \text{ is the degree of } A _{1} (G). \notag
\end{align}
\end{proposition}

\begin{proof}
 Let $n_{1}, n_{2}$ be the degrees of $A_{1}(g), A_{2}(g)$ respectively. Take any matrix $B$ of $n_{1}$ rows and $n_{2}$ columns, and set $A(g) = A_{1}(g) BA_{2}(g^{-1})$. 

Then the matrix $A = M_g(A(g))$ satisfies $A_1(g)A = AA_2(g)$. For, we have, by the invariance of the mean value, 
\[
\begin{gathered}
A _{1} (y) A A _{2} (y ^{-1}) = M _{g} (A _{1} (y) A _{1} (g) B A _{2} (g ^{-1}) A _{2} (y ^{-1})) = M _{g} (A _{1} (y g) B A _{2} ((y g) ^{-1}) = A.
\end{gathered}
\]

By Schur's lemma, $A$ must be equal to a zero matrix. If we take $B = (b_{jl})$ in such a way that only the element $b_{jl}$ is not zero, then, by the unitarity condition $A_{2}(g^{-1}) = \overline{A_{2}(g)}'$, we obtain 
\[
M _{g} (a _{ij} ^{1} (g) \overline {{a _{kl} ^{2} (g)}}) = 0.
\]


Next we have, as above, $A_{1}(g)A = AA_{1}(g)$ for $A = M_{g}(A_{1}(g)BA_{1}(g^{-1}))$. Let $\alpha$ be any one of the eigenvalues of the matrix $A$. Then the matrix $(A - \alpha I_{n_1})$, where $I_{n_1}$ denotes the unit matrix of degree $n_1$, satisfies 
\[
A _{1} (g) \left(A - \alpha I _ {n _{1}}\right) = \left(A - \alpha I _ {n _{1}}\right) A _{1} (g).
\]

Therefore, by Schur's lemma, the matrix $(A - \alpha I_{n_1})$ is either non-singular, or $(A - \alpha I_{n_1}) = 0$. The first possibility is excluded since $\alpha$ is an eigenvalue of $A$. Thus $A = \alpha I_{n_1}$. By taking the trace (the sum of the diagonal elements) of both sides of $A = M_g(A_1(g)BA_1(g^{-1}))$, we obtain 
\[
\begin{aligned} n _{1} \alpha & = \tr (A) = M _{g} (\tr (A _{1} (g) B A _{1} (g) ^{-1}) = M _{g} (\tr (B)) \\ & = \tr (B). \end{aligned}
\]

Therefore, if we take $B = (b_{jl})$ in such a way that $b_{jl} = 1$ while the other elements are all zero, then, from $M_g(A_1(g)BA_1(g^{-1})) = n_1^{-1}$ trace $(B)I_{n_1}$, we obtain 
\[
M _{g} \left(a _{ij} ^{1} (g) \overline {{{a _{kl} ^{1} (g)}}}\right) = n _{1} ^{-1} \delta_{ik} \delta_{jl}.
\]
\end{proof}

\begin{corollaryu} There exists a set $U$ of mutually inequivalent, continuous, unitary, irreducible matrix representations $U(g) = (u_{ij}(g))$ of $G$ satisfying the following three conditions: 

\begin{enumerate}[i)]
\item for any pair of distinct points $g_1, g_2$ of $G$, there exists an $U(g) \in \mathfrak{U}$ such that $U(g_1) \neq U(g_2)$,
\item if $U(g) \in \mathfrak{U}$, then the complex conjugate representation $\bar{U}(g)$ of $U(g)$ also belongs to $\mathfrak{U}$,
\item if $U_{1}(g), U_{2}(g)$ are in $\mathfrak{U}$, then the product representation $U_{1}(g) \times U_{2}(g)$, explained below, is completely reducible to a sum of a finite number of representations $\in \mathfrak{U}$.
\end{enumerate}
\end{corollaryu}

\begin{proof}
 The definition of the product representation $U_{1}(g) \times U_{2}(g)$ is given as follows. Let $(e_{1}^{1}, e_{2}^{1}, \ldots, e_{n}^{1})$ and $(e_{1}^{2}, e_{2}^{2}, \ldots, e_{m}^{2})$ be orthonormal bases respectively of the finite dimensional complex Hilbert spaces subject to the linear mappings $U_{1}(g)$ and $U_{2}(g)$, respectively. The product space of these two Hilbert spaces is spanned by the product base consisting of $nm$ vectors $e_i^1 \times e_j^2$ ( $i = 1, 2, \ldots, n; j = 1, 2, \ldots, m$). Referring to this base, the product representation $U_1(g) \times U_2(g)$ of $U_1(g) = (u_{ij}^1(g))$ and $U_2(g) = (u_{kl}^2(g))$ is given by 
\[
\left(U _{1} (g) \times U _{2} (g)\right) \left(e _{i} ^{1} \times e _{j} ^{2}\right) = \sum\limits_{s,t} u _{si} ^{1} (g) u _{tj} ^{2} (g) \left(e _{i} ^{1} \times e _{j} ^{2}\right). \]

Let us take a maximal set $U$ of mutually inequivalent, continuous, unitary, irreducible matrix representations $U(g)$ satisfying the condition ii). Then, by Theorem~\ref{thm:XI.11.Theorem.1}, condition i) is satisfied. Condition iii) is also satisfied since the product representation of two unitary representations is also unitary and hence is completely reducible. 
\end{proof}

We are now ready to formulate T. \textsc{Tannaka}'s duality theorem. Let $R$ be the set of all Fourier polynomials: 
\begin{equation}\label{eq:XI.11.8}
x (g) = \sum \gamma_{ij} ^{(\alpha)} u _{ij} ^{(\alpha)} (g),
\end{equation}
that is, finite linear combinations of $u_{ij}^{(\alpha)}(g)$, where $(u_{ij}^{(\alpha)}(g)) \in \mathfrak{U}$ and $\gamma_{ij}^{(\alpha)}$ denote complex numbers. Then $\operatorname{Re}$ is a ring with the multiplicative unit $u(u(g) \equiv 1 \text{ on } G)$ and with complex multipliers; the sum and the multiplication in the ring $\operatorname{Re}$ being understood as the function sum and the function multiplication, respectively. Let $\mathfrak{T}$ be the set of all linear homomorphisms $T$ of the ring $\operatorname{Re}$ onto the field of complex numbers such that 
$Tu = 1, T\bar{x} = \overline{Tx}$, the bar indicating the complex-conjugate. 


$\mathfrak{T}$ is not void since each $g \in G$ induces such a homomorphism $T_g$: 
\begin{equation}\label{eq:XI.11.9}
T _{g} x = x (g).
\end{equation}

By the condition i) for $\mathfrak{U}$, we see that 
\begin{equation}\label{eq:XI.11.10}
g _{1} \neq g _{2} \quad \text{ implies } \quad T _ {g _{1}} \neq T _ {g _{2}}.
\end{equation}

\begin{proposition}\label{prop:XI.11.Proposition.4} $\mathfrak{T}$ may be considered as a group which contains $G$ as a subgroup. 
\end{proposition}

\begin{proof}
 We shall define a product $T_{1} \cdot T_{2} = T$ in $T$ as follows. Let 
\[
U ^{(\alpha)} (g) = \left(u _{ij} ^{(\alpha)} (g)\right) \quad (i, j = 1, 2, \dots , n)
\]
be members of $\mathfrak{U}$. We put, for $u_{ij}^{(\alpha)}(gh) = \sum\limits_{k} u_{ik}^{(\alpha)}(g) u_{kj}^{(\alpha)}(h)$, 
\begin{equation}\label{eq:XI.11.11}
T u _{ij} ^{(\alpha)} = \sum\limits_{k} T _{1} u _{ik} ^{(\alpha)} \cdot T _{2} u _{kj} ^{(\alpha)}. \end{equation}

By the orthogonal relations \eqref{eq:XI.11.7}, we see that the functions $u_{ij}^{(\alpha)}(g)$'s, where $(u_{ij}^{(\alpha)}(g)) \in \mathfrak{U}$, are linearly independent on $G$. Hence, we see that $T$ may be extended linearly on the whole $\mathfrak{R}$. It is easy to see that the extension $T$ is also a member of $\mathfrak{T}$ and that $\mathfrak{T}$ is a group with the unit $T_e$ ( $e =$ the unit of $G$) and $T_g^{-1} = T_{g^{-1}}$. $G$ is isomorphically embedded in this group $\mathfrak{T}$ by the correspondence $g \leftrightarrow T_g$, as will be seen from \eqref{eq:XI.11.11} and \eqref{eq:XI.11.10}. 
\end{proof}

In truth, we have 

\begin{theorem}[T. \textsc{Tannaka}]\label{thm:XI.11.Theorem.2} $\mathfrak{T} = G$, viz. every $T \in \mathfrak{T}$ is equal to a certain $T_g$: 
\begin{equation}\label{eq:XI.11.12}
T x = x (g) \quad \text{ for all } \quad x \in \mathfrak {R}.
\end{equation}
\end{theorem}

\begin{proof}
 We introduce a weak topology in the group $T$ by taking the sets of the form 
\begin{equation}\label{eq:XI.11.13}
\{T \in \mathfrak {T}; | T x _{i} - T _{e} x _{i} | < \varepsilon_{i} (i = 1, 2, \dots , n) \}
\end{equation}
as neighbourhoods of the unit $T_{e}$ of the group $\mathfrak{T}$. Then $\mathfrak{T}$ is a compact space. For, by $\sum\limits_{s}|T u_{si}^{(\alpha)}(g)|^2 = (T \cdot T)(1) = 1$ implied by ($8$) and \eqref{eq:XI.11.11}, 
$\mathfrak{T}$ is a closed subset of the topological product of compact spaces: 
\[
\prod\limits_ {x \in \mathfrak {R}} \left\{z; | z | \leqq \sup\limits _ {T \in \mathfrak {T}} | T x | \right\},
\]
and so we can apply Tychonov's theorem. It is easily seen that the isomorphic embedding $g \leftrightarrow T_{g}$ is also a topological one, because a one-one continuous map of a compact space onto a compact space is a topological map. Therefore $G$ may be considered as a closed subgroup of the compact group $T$. 

By the above weak topology of $\mathfrak{T}$, each $x(g) \in \operatorname{Re}$ gives rise to a continuous function $x(T)$ on the compact group $\mathfrak{T}$ such that $x(T_g) = x(g)$. We have only to set $x(T) = T \cdot x$. The set of all these continuous functions $x(T)$ constitutes a ring $\operatorname{Re}(\mathfrak{T})$ of complex-valued continuous functions on $\mathfrak{T}$ satisfying the conditions: 

$1$) $1 = u(T)\in \operatorname{Re} (\mathfrak{T})$ 
$2$) for any pair of distinct points $T_{1}, T_{2}$ of $\mathfrak{T}$, there exists an $x(T) \in \operatorname{Re}(\mathfrak{T})$ such that $x(T_{1}) \neq x(T_{2})$, 
$3$) for any $x(T) \in \operatorname{Re}(\mathfrak{T})$, there exists the complex-conjugate function $\overline{x}(T) = \overline{x(T)}$ which belongs to $\operatorname{Re}(\mathfrak{T})$. 

Now let us suppose that $\mathfrak{T} - G$ is not void. Since a compact space $\mathfrak{T}$ is normal, we may apply Urysohn's theorem to the effect that there exist a point $T_0 \in (\mathfrak{T} - G)$ and a continuous function $y(T)$ on $\mathfrak{T}$ such that 
\begin{equation}\label{eq:XI.11.14}
y (T) \geqq 0 \text{ on } \mathfrak {T}, y (g) = 0 \text{ on } G \text{ and } y (T _{0}) = 1.
\end{equation}

By the Stone-Weierstrass theorem in Chapter~\ref{chap:0}, as applied to the ring $\operatorname{Re}(\mathfrak{T})$ satisfying $1$), $2$) and $3$), we see that there exists, for any $\varepsilon > 0$, a function $x(g) = \sum \gamma_{ij}^{(\alpha)} u_{ji}^{(\alpha)}(g) \in \operatorname{Re}$ such that $|y(T) - \sum \gamma_{ij}^{(\alpha)} u_{ij}^{(\alpha)}(T)| < \varepsilon$ on $\mathfrak{T}$, and, in particular, $|y(g) - \sum \gamma_{ij}^{(\alpha)} u_{ij}^{(\alpha)}(g)| < \varepsilon$ on $G$. Let $u_{11}^{(\alpha_n)}(g) = u(g) = 1$. Then, by taking the mean value and remembering the orthogonality relation \eqref{eq:XI.11.7}, as applied to the compact groups $\mathfrak{T}$ and $G$, we obtain 
\begin{equation}\label{eq:XI.11.15}
\left| M _{T} (y (T)) - \gamma_{11} ^ {\left(\alpha_{0}\right)} \right| < \varepsilon , \left| M _{g} (y (g)) - \gamma_{11} ^ {\left(\alpha_{0}\right)} \right| < \varepsilon .
\end{equation}

We have thus arrived at a contradiction, since $M_T(y(T)) > 0$ and $M_g(y(g)) = 0$ by \eqref{eq:XI.11.14}. 
\end{proof}
\begin{remarkn}[Remark $1$] The above proof of Tannaka's theorem is adapted from K. \textsc{Yosida} \cite{yosida357}. The original proof is given in T. \textsc{Tannaka} \cite{tannaka320}. Since a continuous, unitary and irreducible matrix representation of a compact abelian group $G$ is precisely a continuous function $\chi(g)$ on $G$ satisfying 
\[
\chi \left(g _{1}\right) \chi \left(g _{2}\right) = \chi \left(g _{1} g _{2}\right) \quad \text{ and } \quad | \chi (g) | = 1,
\]
\end{remarkn}

Tannaka's theorem contains, as a special case, the duality theorem of L. \textsc{Pontrjagin} \cite{pontrjagin263}. For further references, see M. A. \textsc{Naimark} \cite{naimark221}. 

\begin{remarkn}[Remark $2$] To define mean values of continuous functions $f_{k}(g)$ ( $k = 1, 2, \ldots, m$) on $G$ by the method given in Chapter~\ref{chap:VIII}, \ref{sec:VIII.5}, we have only to replace the $\text{ dis }(g_1, g_2)$ there by $\text{ dis }(g_1, g_2) = \sup\limits_{g, h \in G; k=1,2,\ldots,m} |f_k(g g_1 h) - f_k(g g_2 h)|$. 
\end{remarkn}


\yossection{12}{Functions of a Self-adjoint Operator}
\label{sec:XI.12}

Let $H = \int \lambda dE(\lambda)$ be the spectral resolution of a self-adjoint operator $H$ in a Hilbert space $X$. For a complex-valued Baire function $f(\lambda)$, we consider the set 
\begin{equation}\label{eq:XI.12.1}
D (f (H)) = \left\{x \in X; \int_{- \infty} ^{\infty} | f (\lambda) | ^{2} \dd \| E (\lambda) x \| ^{2} < \infty \right\},
\end{equation}
where the integral is taken with respect to the Baire measure $m$ determined by $m((\lambda_{1}, \lambda_{2}]) = \|E(\lambda_{2})x\|^{2} - \|E(\lambda_{1})x\|^{2}$. As in the case of a continuous function $f(\lambda)$, treated in Chapter~\ref{chap:XI}, \ref{sec:XI.5}, we see that the integral 
\begin{equation}\label{eq:XI.12.2}
\int_{- \infty} ^{\infty} f (\lambda) \dd (E (\lambda) x, y), x \in D (f (H)), y \in X,
\end{equation}
exists and is finite with respect to the Baire measure $m$ determined by $m((\lambda_{1}, \lambda_{2}]) = (E(\lambda_{2})x, y) - (E(\lambda_{1})x, y)$. Further we see that \eqref{eq:XI.12.2} gives the complex conjugate of a bounded linear functional of $y$. Hence we may write \eqref{eq:XI.12.2} as $(f(H)x, y)$ by virtue of F. \textsc{Riesz}' representation theorem in Chapter~\ref{chap:III}, \ref{sec:III.6}. In this way, we can define functions of $H$: 
\begin{equation}\label{eq:XI.12.3}
f (H) = \int_{- \infty} ^{\infty} f (\lambda) \dd E (\lambda)
\end{equation}
through \eqref{eq:XI.12.1} and \eqref{eq:XI.12.2}. 

\begin{example}\label{ex:XI.12.Example.1} If $H$ is bounded self-adjoint and $f(\lambda) = \sum\limits_{j=1}^{n} \alpha_j \lambda^j$, then, as in Chapter~\ref{chap:XI}, \ref{sec:XI.5}. 
\[
f (H) = \int_{- \infty} ^{\infty} f (\lambda) \dd E (\lambda) = \sum\limits_{j=1} ^{n} \alpha_{j} H ^{j}.
\]
\end{example}

\begin{example}\label{ex:XI.12.Example.2} If $f(\lambda) = (\lambda - i)(\lambda + i)^{-1}$, then $f(H)$ is equal to the Cayley transform $U_H$ of $H$. We have, in this case, $D(f(H)) = X$ since $|f(\lambda)| = 1$ for real $\lambda$. It is easy to see, by $E(\lambda_1)E(\lambda_2) = E(\min (\lambda_1, \lambda_2))$, that if we apply the bounded self-adjoint operator $(H - iI)^{-1} = \int (\lambda - i)^{-1} dE(\lambda)$ to the operator $f(H) = \int f(\lambda) dE(\lambda)$ then the result is equal to $(H + iI)^{-1} = \int (\lambda + i)^{-1} dE(\lambda)$. Hence $f(H) = U_H$. 
\end{example}

\begin{example}\label{ex:XI.12.Example.3} As in Chapter~\ref{chap:XI}, \ref{sec:XI.5}, we have 
\begin{equation}\label{eq:XI.12.4}
\left\| f (H) x \right\| ^{2} = \int_{- \infty} ^{\infty} | f (\lambda) | ^{2} d \left\| E (\lambda) x \right\| ^{2} \text{ whenever } x \in D (f (H)).
\end{equation}
\end{example}

\begin{definitionu} Let $A$ be a not necessarily bounded linear operator, and $B$ a bounded linear operator in a Hilbert space. If 
\begin{equation}\label{eq:XI.12.5}
x \in D (A) \quad \text{ implies } \quad B x \in D (A) \quad \text{ and } \quad A B x = B A x,
\end{equation}
that is, if $AB \supseteq BA$, then we shall write $B \in (A)'$ and say that $B$ is commutative with $A$. We shall thus write the totality of bounded linear operators $B$ commutative with $A$ by $(A)'$. 
\end{definitionu}

\begin{theorem}\label{thm:XI.12.Theorem.1} For a function $f(H)$ of a self-adjoint operator $H = \int \lambda dE(\lambda)$ in a Hilbert space $X$, we have 
\begin{equation}\label{eq:XI.12.6}
(f (H)) ^{\prime} \supseteq (H) ^{\prime},
\end{equation}
that is, $f(H)$ is commutative with every bounded linear operator which commutes with $H$. (Since $E(\lambda) \in (H)'$ by Theorem~\ref{thm:XI.5.Theorem.2} in Chapter~\ref{chap:XI}, \ref{sec:XI.5}, we see, in particular, that $f(H)$ is commutative with every $E(\lambda)$.) 
\end{theorem}

\begin{proof}
 Suppose $S \in (H)'$. Then we can show that $S$ is commutative with every $E(\lambda)$. We first show that $S$ is commutative with the Cayley transform $U_{H}$ of $H$. For, if $x \in D(H)$, then we have, by $S \in (H)'$, 
\[
S (H + i I) x = (H + i I) S x, (H - i I) S x = S (H - i I) x. \]

By putting $(H + iI)x = y$, we see, from the first of the above relations, that 
\[
(H + i I) ^{-1} S y = S (H + i I) ^{-1} y \quad \text{ for all } \quad y \in X = R (H + i I).
\]

Hence 
\[
S (H - i I) (H + i I) ^{-1} = (H - i I) (H + i I) ^{-1} S, \text{ that is, } S U _{H} = U _{H} S.
\]

Therefore, $S$ is commutative with $(U_H)^n = \int_0^{2\pi} e^{ni\theta} dF(\theta) (n = 0, \pm 1, \ldots)$ and so 
\[
\int_{0} ^{2 \pi} e ^{n i \theta} d (S F (\theta) x, y) = \left(S \int_{0} ^{2 \pi} e ^{n i \theta} \dd F (\theta) x, y\right) = \int_{0} ^{2 \pi} e ^{n i \theta} \dd (F (\theta) S x, y).
\]

Hence, as in the proof of the uniqueness of the spectral resolution of a unitary operator, given in Chapter~\ref{chap:XI}, \ref{sec:XI.4}, we obtain $SF(\theta) = F(\theta)S$. This proves $SE(\lambda) = E(\lambda)S$ since $E(-\cot \theta) = F(\theta)$. Thus, for $x \in D(f(H))$,  
we obtain 
\[
(S \int f (\lambda) \dd E (\lambda) x, y) = \int f (\lambda) d (S E (\lambda) x, y) = \int f (\lambda) \dd (E (\lambda) S x, y),
\]
that is, $Sf(H) \subseteq f(H) S$. 
\end{proof}

The above Theorem~\ref{thm:XI.12.Theorem.1} admits a converse in the form of 

\begin{theorem}[\textsc{Neumann}-\textsc{Riesz}-\textsc{Mimura}]\label{thm:XI.12.Theorem.2} Let $H$ be a self-adjoint operator in a separable Hilbert space $X$. Let $T$ be a closed linear operator in $X$ such that $D(T)^{a} = X$. Then a necessary and sufficient condition in order that $T$ be a function $f(H)$ of $H$ with an everywhere finite Baire function $f(\lambda)$ is that 
\begin{equation}\label{eq:XI.12.7}
(T) ^{\prime} \supseteq (H) ^{\prime}.
\end{equation}
\end{theorem}

\begin{proof}
 We have only to show that condition \eqref{eq:XI.12.7} is sufficient. We may assume that $H$ is a bounded self-adjoint operator. If $H$ is not bounded, then we consider $H_{1} = \tan^{-1}H$. Since $|\tan^{-1}\lambda| \leqq \pi/2$, we easily see that $H_{1}$ is bounded self-adjoint. By Theorem~\ref{thm:XI.12.Theorem.1}, the operator $H = \tan H_{1}$ is commutative with every operator $\in (H_{1})'$. Thus, by hypothesis, we have 
\[
(T) ^{\prime} \supseteq (H) ^{\prime} \supseteq (H _{1}) ^{\prime}. \]

Therefore, if Theorem~\ref{thm:XI.12.Theorem.2} is proved for bounded $H_{1}$, then $T = f_{1}(H_{1}) = f_{1}(\tan^{-1}H) = f_{2}(H)$, where $f_{2}(\lambda) = f_{1}(\tan^{-1}\lambda)$. 

We may thus assume that $H$ is bounded and self-adjoint. 

The first step. For any fixed $x_0 \in D(T)$, we can find a Baire function $F(\lambda)$ such that $Tx_0 = F(H)x_0$. This may be proved as follows. Let $M(x_0)$ be the smallest closed linear subspace of $X$ spanned by $x_0, Hx_0, H^2x_0, \ldots, H^n x_0, \ldots$. We denote by $L$ the projection upon $M(x_0)$. Then $(T)' \ni L$. For, by $HM(x_0) \subseteq M(x_0)$, we obtain $HL = LHL$ and so $LH = (HL)^* = (LHL)^* = LHL = HL$, that is, $L \in (H)'$. Hence, by hypothesis, $(T)' \ni L$. 

Therefore $Tx_{0}=TLx_{0}=LTx_{0}\in M(x_{0})$, and so there exists a sequence $\{\phi_{n}(\lambda)\}$ of polynomials such that 
\begin{equation}\label{eq:XI.12.8}
T x _{0} = \slim\limits_{n \rightarrow \infty} p _{n} (H) x _{0}.
\end{equation}

Hence, by \eqref{eq:XI.12.4}, we obtain 
\[
\left\| p _{n} (H) x _{0} - p _{m} (H) x _{0} \right\| ^{2} = \int_{- \infty} ^{\infty} \left| p _{n} (\lambda) - p _{m} (\lambda) \right| ^{2} d \left\| E (\lambda) x _{0} \right\| ^{2}.
\]

As in the proof of the completeness of $L^{2}(S, \mathfrak{B}, m)$, we see that there exists a Baire function $F(\lambda)$ which is square-integrable with respect to the measure $m$ determined by $m((\lambda_{1}, \lambda_{2}]) = \|E(\lambda_{2}) x_{0}\|^{2} - \|E(\lambda_{1}) x_{0}\|^{2}$ such that 
\[
\lim\limits _{n \rightarrow \infty} \int_{- \infty} ^{\infty} | F (\lambda) - p _{n} (\lambda) | ^{2} \dd \| E (\lambda) x _{0} \| ^{2} = 0.
\]

Hence, for $F(H)$, we have 
\[
\lim\limits _{n \rightarrow \infty} \| (F (H) - p _{n} (H)) x _{0} \| ^{2} = \lim\limits _{n \rightarrow \infty} \int_{- \infty} ^{\infty} | F (\lambda) - p _{n} (\lambda) | ^{2} \dd \| E (\lambda) x _{0} \| ^{2} = 0.
\]

This proves that $Tx_0 = F(H)x_0$. Since $F(\lambda)$ is finite almost everywhere with respect to the measure $m$ determined by $m((\lambda_1, \lambda_2]) = \|E(\lambda_2)x_0\|^2 - \|E(\lambda_1)x_0\|^2$, we may assume that $F(\lambda)$ is a Baire function which takes a finite value at every $\lambda$; we may put $F(\lambda) = 0$ for those $\lambda$ for which $|F(\lambda)| = \infty$. 

The second step. Since $X$ is separable and $D(T)^a = X$, we may choose a countable sequence $\{g_n\} \subseteq D(T)$ such that $\{g_n\}$ is strongly dense in $X$. We set 
\begin{equation}\label{eq:XI.12.9}
f _{1} = g _{1}, f _{2} = g _{2} - L _{1} g _{2}, \dots , f _{n} = g _{n} - \sum\limits_{k=1} ^{n-1} L _{k} g _{n}, \text{ where } L _{k}
\end{equation}
is the projection on the closed linear subspace $M(f_k)$. 


By the first step, we have $(T)' \ni L_k$ and so 
\begin{equation}\label{eq:XI.12.10}
L _{k} g _{n} \in D (T) \text{ which implies } f _{n} \in D (T) \quad (n = 1, 2, \dots).
\end{equation}

We may show that 
\begin{equation}\label{eq:XI.12.11}
L _{i} L _{k} = 0 \quad (\text{ for } i \neq k)
\end{equation}
and 
\begin{equation}\label{eq:XI.12.12}
I = \sum\limits_{k=1} ^{\infty} L _{k}.
\end{equation}

Suppose that \eqref{eq:XI.12.11} is proved for $i, k < n$. Then, for $i < n$, 
\begin{equation}L _{i} f _{n} = L _{i} g _{n} - L _{i} \left(\sum\limits_{k=1} ^{n-1} L _{k} g _{n}\right) = L _{i} g _{n} - L _{i} ^{2} g _{n} = L _{i} g _{n} - L _{i} g _{n} = 0, L _{i} H ^ {k ^{\prime}} f _{n} = H ^ {k ^{\prime}} L _{i} f _{n} = 0. \notag\end{equation}

Thus $M(f_{n})$ is orthogonal to $M(f_{i})$. This proves $L_{i}L_{n}=L_{n}L_{i}=0$. 

Next we put $\sum\limits_{k=1}^{\infty} L_k = P$ and show that $P g_n = g_n$ ( $n = 1, 2, \ldots$). Then, since $\{g_n\}$ is dense in $X$, we obtain $P = I$. But, by ($9$), we have 
\[
P g _{n} = P f _{n} + \sum\limits_{k=1} ^{n-1} P L _{k} g _{n}.
\]

We also have $Pf_{n} = f_{n}$ by $f_{n} \in M(f_{n})$. Hence, by $PL_{k} = L_{k}$ implied by \eqref{eq:XI.12.11}, we obtain $Pg_{n} = g_{n}$ ( $n = 1, 2, \ldots$). 

The third step. Take a sequence $\{c_n\}$ of positive numbers such that 
\[
\slim _{k \to \infty} \sum\limits_{n=1} ^{k} c _{n} f _{n}, \slim _{k \to \infty} \sum\limits_{n=1} ^{k} c _{n} T f _{n}
\]
both exist. For instance, we may take $c_{n} = 2^{-n} (\|f_{n}\| + \|Tf_{n}\|)^{-1}$. Since $T$ is a closed operator, we have 
\begin{equation}\label{eq:XI.12.13}
x _{0} = \sum\limits_{n=1} ^{\infty} c _{n} f _{n} \in D (T) \text{ and } T x _{0} = y _{0} = \sum\limits_{n=1} ^{\infty} c _{n} T f _{n}.
\end{equation}

Hence, by the first step, 
\begin{equation}\label{eq:XI.12.14}
T x _{0} = F (H) x _{0}.
\end{equation}

Let $B \in (H)'$ be a bounded self-adjoint operator. Then, by hypothesis, $B \in (T)'$. By Theorem~\ref{thm:XI.12.Theorem.1}, $F(H)$ is commutative with $B$. Hence 
\begin{equation}\label{eq:XI.12.15}
F (H) B x _{0} = B F (H) x _{0} = B T x _{0} = T B x _{0}.
\end{equation}

Let $e_n(\lambda)$ be the defining function of the set $\{\lambda; |F(\lambda)| \leqq n\}$, and put 
\[
B = c _{m} ^{-1} P _{n} H ^{k} L _{m}, \quad \text{ where } \quad P _{n} = e _{n} (H).
\]

Then we can show that 
\begin{equation}\label{eq:XI.12.16}
T P _{n} = F (H) P _{n}.
\end{equation}

In fact, we have $L_{m}x_{0} = \sum\limits_{n=1}^{\infty} c_{n}L_{m}f_{n} = c_{m}f_{m}$, by \eqref{eq:XI.12.11} and $f_{m} \in M(f_{m})$. Hence, by \eqref{eq:XI.12.15}, 
\[
\begin{aligned}
&F (H) P _{n} H ^{k} f _{m}  = F (H) c _{m} ^{-1} P _{n} H ^{k} L _{m} x _{0} = F (H) B x _{0} = T B x _{0} = T c _{m} ^{-1} P _{n} H ^{k} L _{m} x _{0} = T P _{n} H ^{k} f _{m},
\end{aligned}
\]
that is, for $h$ spanned by $H^{k}f_{m}$ with fixed $m$, we have 
\begin{equation}\tag{16'}\label{eq:XI.12.16'}
F (H) P _{n} h = T P _{n} h.
\end{equation}

But such $h$'s are dense in $M(f_m)$ and so, by \eqref{eq:XI.12.12}, we see that if we let $m$ take all positive integers then the $h$'s are dense in $X$. Hence we have proved \eqref{eq:XI.12.16} for those $h$'s which are dense in $X$. 

Now $P_{n}$ is bounded by \eqref{eq:XI.12.4}. By the operational calculus given below, the operator $F(H)P_{n}$ is equal to the function $F_{n}(H)$ where 
\[
\begin{aligned}
&F _{n} (\lambda)  = F (\lambda) e _{n} (\lambda) = F (\lambda), \text{ for } | F (\lambda) | \leqq n, = 0, \text{ for } | F (\lambda) | > n.
\end{aligned}
\]

Thus $F_{n}(H) = F(H)P_{n}$ is bounded. 

Let $h^{*} \in X$ be arbitrary, and let $h^{*} = \slim\limits_{j \to \infty} h_{j}$ where $h_{j}^{\prime}s$ are linear combinations of elements of the form $H^{h}f_{m}$. Such a choice of $h_{j}$ is possible as proved above. By the continuity of the operator $F(H)P_{n}$, we have 
\[
F (H) P _{n} h ^{*} = \slim _{j \rightarrow \infty} F (H) P _{n} h _{j}.
\]

Hence, by $\slim\limits_{j \to \infty} P_n h_j = P_n h^*$ and \eqref{eq:XI.12.16}, we see that the closed operator $T$ satisfies \eqref{eq:XI.12.16}. 

The fourth step. Let $y \in D(F(H))$ and set $y_n = P_n y$. Since $F(\lambda)$ is finite everywhere, we must have $\slim\limits_{n \to \infty} P_n = I$. Thus $\slim\limits_{n \to \infty} y_n = \slim\limits_{n \to \infty} P_n y = y$. Hence, by \eqref{eq:XI.12.16}, we obtain $T \supseteq F(H)$. For, $\slim\limits_{n \to \infty} F(H) P_n y = \slim\limits_{n \to \infty} F_n(H) y = F(H) y$ whenever $y \in D(F(H))$. 

Let $y \in D(T)$ and set $y_{n} = P_{n}y$. Then 
\[
\begin{gathered}
T y _{n} = T P _{n} y = F (H) P _{n} y (\text{ by } (1 6)), T P _{n} y = P _{n} T y (\text{ by } P _{n} = e _{n} (H) \in (H) ^{\prime}).
\end{gathered}
\]

By the operational calculus given below, the function $F(H)$ of $H$ is a closed operator. Hence, letting $n \to \infty$ in the above relations, we see that $F(H) \geqq T$. We have thus proved that $T = F(H)$. 
\end{proof}

\begin{theorem}\label{thm:XI.12.Theorem.3} \begin{enumerate}[(i)]
\item Let $\bar{f}(\lambda)$ be the complex-conjugate function of $f(\lambda)$. Then $D(\bar{f}(H)) = D(f(H))$ and, for $x, y \in D(f(H)) = D(\bar{f}(H))$, we have 
\begin{equation}\label{eq:XI.12.17}
(f (H) x, y) = (x, \bar {f} (H) y).
\end{equation}
\item If $x \in D(f(H))$, $y \in D(g(H))$, then 
\begin{equation}\label{eq:XI.12.18}
(f (H) x, g (H) y) = \int_{- \infty} ^{\infty} f (\lambda) \bar {g} (\lambda) \dd (E (\lambda) x, y).
\end{equation}
\item $(\alpha f)(H)x = \alpha f(H)x$ if $x\in D(f(H))$. If $x\in D(f(H))\cap D(g(H))$, then 
\begin{equation}\label{eq:XI.12.19}
(f + g) (H) x = f (H) x + g (H) x.
\end{equation}
\item If $x \in D(f(H))$, then the condition $f(H)x \in D(g(H))$ is equivalent to the condition $x \in D(f \cdot g(H))$, where $f \cdot g(\lambda) = f(\lambda)g(\lambda)$, and we have 
\begin{equation}\label{eq:XI.12.20}
g (H) f (H) x = (g \cdot f) (H) x.
\end{equation}
\item If $f(\lambda)$ is finite everywhere, then $f(H)$ is a normal operator and 
\begin{equation}\label{eq:XI.12.21}
f (H) ^{*} = \bar {f} (H).
\end{equation}

In particular, $f(H)$ is self-adjoint if $f(\lambda)$ is real-valued and finite everywhere.
\end{enumerate}
\end{theorem}

\begin{proof}
\begin{enumerate}[(i)]
\item $D(f(H)) = D(\bar{f}(H))$ is clear and 
\[ \begin{aligned}
&(f (H) x, y)  = \int_{- \infty} ^{\infty} f (\lambda) \dd (E (\lambda) x, y) = \int_{- \infty} ^{\infty} f (\lambda) d (x, E (\lambda) y) = \overline {{(f (H) y , x)}} = (x, \overline {{f}} (H) y).
\end{aligned}
\]
\item We know, by Theorem~\ref{thm:XI.12.Theorem.1}, that $E(\lambda)$ is commutative with $g(H)$. Thus 
\[
\begin{aligned}
(f (H) x, g (H) y) &= \int_{- \infty} ^{\infty} f (\lambda) \dd (E (\lambda) x, g (H) y) = \int_{- \infty} ^{\infty} f (\lambda) d (x, E (\lambda) g (H) y)
\\
&= \int_{- \infty} ^{\infty} f (\lambda) \dd \overline {{(g (H) E (\lambda) y , x)}} = \int_{- \infty} ^{\infty} f (\lambda) \dd \left(\int_{- \infty} ^{\infty} \overline {{g (\mu)}} \dd \overline {{(E (\mu) E (\lambda) y , x)}}\right)
\\
&= \int_{- \infty} ^{\infty} f (\lambda) \dd \left(\int_{- \infty} ^{\lambda} \overline {{g (\mu)}} \dd \overline {{(y , E (\mu) x)}}\right) = \int_{- \infty} ^{\infty} f (\lambda) \overline {{g (\lambda)}} \dd (E (\lambda) x, y).
\end{aligned}
\]
\item is clear.
\item Let $x$ satisfy $\int_{-\infty}^{\infty}|f(\lambda)|^2\dd\|E(\lambda)x\|^2 < \infty$. Then, by $E(\lambda)E(\mu) = E(\min (\lambda ,\mu))$, the condition $\int_{-\infty}^{\infty}|g(\lambda)|^2\dd\|E(\lambda)f(H)x\|^2 < \infty$ implies, by virtue of the commutativity of $E(\lambda)$ with $f(H)$, 
\[
\begin{gathered}
\infty > \int_{- \infty} ^{\infty} | g (\lambda) | ^{2} \dd \| E (\lambda) f (H) x \| ^{2} = \int_{- \infty} ^{\infty} | g (\lambda) | ^{2} \dd \| f (H) E (\lambda) x \| ^{2}
\\
= \int_{- \infty} ^{\infty} | g (\lambda) | ^{2} \dd \left(\int_{- \infty} ^{\infty} | f (\mu) | ^{2} \dd \| E (\mu) E (\lambda) x \| ^{2}\right)
\\
= \int_{- \infty} ^{\infty} | g (\lambda) | ^{2} \dd \left(\int_{- \infty} ^{\lambda} | f (\mu) | ^{2} \dd \| E (\mu) x \| ^{2}\right) = \int_{- \infty} ^{\infty} | g (\lambda) f (\lambda) | ^{2} \dd \| E (\lambda) x \| ^{2}.
\end{gathered}
\]

Since the above calculation may be traced conversely, we see that, under the hypothesis $x \in D(f(H))$, the two conditions $f(H)$ $x \in D(g(H))$ and $x \in D(f \cdot g(H))$ are equivalent and we have 
\[
\begin{aligned}
&(g (H) f (H) x, y)  = \int_{- \infty} ^{\infty} g (\lambda) \dd (E (\lambda) f (H) x, y) = \int_{- \infty} ^{\infty} g (\lambda) \dd \left(\int_{- \infty} ^{\lambda} f (\mu) \dd (E (\mu) x, y)\right) \\
&= \int_{- \infty} ^{\infty} g (\lambda) f (\lambda) \dd (E (\lambda) x, y) = ((g \cdot f) (H) x, y).
\end{aligned}
\]
\item Let us put $h(\lambda)=|f(\lambda)|+\alpha$, $k(\lambda)=h(\lambda)^{-1}$, $g(\lambda)=f(\lambda)h(\lambda)^{-1}$, where $\alpha$ is any positive integer. Then $k(\lambda)$ and $g(\lambda)$ are both bounded functions. Hence $D(k(H))=D(g(H))=X$. Thus, by (iv), 
\begin{equation}\label{eq:XI.12.22}
f (H) = h (H) g (H) = g (H) h (H).
\end{equation}

We have, by (i) and $D(k(H)) = X$, $k(H)^* = k(H)$, i.e., $k(H)$ is self-adjoint. By (iv) we have $x = h(H)k(H)x$ for all $x \in X$ and $x = k(H)h(H)x$ for all $x \in D(h(H))$. Hence $h(H) = k(H)^{-1}$. Thus, by Theorem~\ref{thm:VII.3.Theorem.1} in Chapter~\ref{chap:VII}, \ref{sec:VII.3}, $h(H)$ is self-adjoint. Therefore $D(f(H)) = D(h(H))$ is dense in $X$ and so we may define $f(H)^*$. We shall show that $f(H)^* = \overline{f}(H)$. Let a pair $\{y, y^*\}$ of elements $\in X$ be such that $(f(H)x, y) = (x, y^*)$ for all $x \in D(f(H))$. Then, by $g(H)^* = \overline{g}(H)$ (implied by (i)) and \eqref{eq:XI.12.22}, 
\[
(f (H) x, y) = (g (H) h (H) x, y) = (h (H) x, \bar {g} (H) y).
\]

Hence, by $x \in D(f(H)) = D(h(H))$ and the self-adjointness of $h(H)$, we obtain 
\[
\bar {g} (H) y \in D (h (H)) \quad \text{ and } \quad h (H) \bar {g} (H) y = y ^{*}.
\]

Thus, again by \eqref{eq:XI.12.22}, we obtain $\bar{f}(H)y=y^{*}$ so that $f(H)^{*}=\bar{f}(H)$. Therefore, by (iv), we see that $f(H)$ is normal, that is, $f(H)^{*}f(H)=f(H)f(H)^{*}$. \qedhere
\end{enumerate}
\end{proof}

\begin{corollaryu} If $f(\lambda)$ is finite everywhere, then $f(H)$ is closed. 
\end{corollaryu}

\begin{proof}
 Clear from $f(H)^{**} = \bar{f}(H)^* = \bar{\bar{f}}(H) = f(H)$. 
\end{proof}

A historical note. Theorem~\ref{thm:XI.12.Theorem.2} was first proved by J. \textsc{von} \textsc{Neumann} \cite{neumann234} for the case of a bounded self-adjoint operator $T$. Cf. F. \textsc{Riesz} \cite{rickart276}. The general case of a closed linear operator $T$ was proved by Y. \textsc{Mimura} \cite{mimura201}. The exposition given above is adapted from Y. \textsc{Mimura} \cite{mimura201} and B. Sz. \textsc{Nagy} \cite{nagy217}. 


\yossection{13}{Stone's Theorem and Bochner's Theorem}
\label{sec:XI.13}

As an example of functions of a self-adjoint operator, we give 

\begin{theorem}[M. H. \textsc{Stone}]\label{thm:XI.13.Theorem.1} Let $\{U_t\}$, $-\infty < t < \infty$, be a one-parameter group of class $(C_0)$ of unitary operators in a Hilbert space $X$. Then 
\begin{equation}\label{eq:XI.13.1}
U _{t} = f _{t} (H), \text{ where } f _{t} (\lambda) = \exp (i t \lambda) \text{ and } i H = A, H ^{*} = H,
\end{equation}
is the infinitesimal generator of the group $U_{t}$. 
\end{theorem}


Conversely, for any self-adjoint operator $H$ in $X$, $U_{t} = f_{t}(H)$ defines a one-parameter group of class $(C_0)$ of unitary operators. 

\begin{proof}
 We have, by the representation theorem of the semi-group theory, 
\[
U _{t} x = \slim _{n \to \infty} \exp (t i H (I - n ^{-1} i H) ^{-1}) x. \]

Since the function $g(t) = \exp (ti\lambda (1 - n^{-1}i\lambda)^{-1})$ is smaller than $\exp ((-nt\lambda^2) / (n^2 +\lambda^2))$ in absolute value, we have, for $H = \int \lambda dE(\lambda)$, 
\[
\exp (t i H (I - n ^{-1} i H) ^{-1}) = \int_{- \infty} ^{\infty} \exp \left(\frac {t i \lambda}{1 - n ^{-1} i \lambda}\right) \dd E (\lambda)
\]
and, moreover, 
\[
\begin{aligned}
\lim\limits _{n \rightarrow \infty}& \int_{- \infty} ^{\infty} \left| \exp \left(\frac {t i \lambda}{1 - n ^{-1} i \lambda}\right) - \exp (t i \lambda) \right| ^{2} \dd \| E (\lambda) x \| ^{2}
\\
&= \lim\limits _{n \rightarrow \infty} \int_{- \infty} ^{\infty} \left| \exp \left(\frac {- t \lambda^{2}}{n - i \lambda}\right) - 1 \right| ^{2} \dd \| E (\lambda) x \| ^{2} = 0.
\end{aligned}
\]

This proves that $U_{t} = f_{t}(H) = \int_{-\infty}^{\infty}\exp (it\lambda)dE(\lambda).$ 

For the converse part of Theorem~\ref{thm:XI.13.Theorem.1}, we observe that, by the operational calculus of the preceding section, 
\[
f _{t} (H) ^{*} = f _{-t} (H) \quad \text{ and } \quad f _{t} (H) f _{s} (H) = f _{t+s} (H), f _{0} (H) = I.
\]

We also have the strong continuity of $f_{t}(H)$ at $t = 0$ by 
\[
\left\| f _{t} (H) x - x \right\| ^{2} = \int_{- \infty} ^{\infty} | \exp (i t \lambda) - 1 | ^{2} \dd \| E (\lambda) x \| ^{2} \rightarrow 0 \text{ as } t \rightarrow 0.
\]

Hence $U_{t} = f_{t}(H)$ is a one-parameter group of class $(C_0)$ of unitary operators. 
\end{proof}

\begin{remarku} For the original proof, see M. H. \textsc{Stone} \cite{stone312}, Cf. also J. \textsc{von} \textsc{Neumann} \cite{neumann235}. Another proof given by E. \textsc{Hopf} \cite{hirsch109} is based on a theorem due to S. \textsc{Bochner}: 
\end{remarku}

\begin{theorem}[\textsc{Bochner}]\label{thm:XI.13.Theorem.2} $A$ complex-valued continuous function $f(t)$, $-\infty < t < \infty$, is representable as 
\begin{equation}\label{eq:XI.13.2}
f (t) = \int_{- \infty} ^{\infty} e ^{i t \lambda} \dd v (\lambda) \text{ with a non-decreasing right-continuous bounded function } v(\lambda),
\end{equation}
right-continuous bounded function $v(\lambda)$, 


iff $f(t)$ is positive definite in the following sense: 
\begin{equation}\label{eq:XI.13.3}
\int_{- \infty} ^{\infty} \int_{- \infty} ^{\infty} f (t - s) \varphi (t) \overline {{\varphi (s)}} \dd t \dd s \geqq 0
\end{equation}
for every continuous function $\varphi$ with compact support. 
\end{theorem}


The proof of Theorem~\ref{thm:XI.13.Theorem.1} given by E. \textsc{Hopf} starts with the fact that $f(t) = (U_{t}x, x)$ satisfies condition ($3$) as may be seen from 
\[
\begin{aligned} & \int_{- \infty} ^{\infty} \int_{- \infty} ^{\infty} (U _{t-s} x, x) \varphi (t) \overline {{{{\varphi (s)}}}} \dd t \dd s = \int_{- \infty} ^{\infty} \int_{- \infty} ^{\infty} (U _{t} x, U _{s} x) \varphi (t) \overline {{{{\varphi (s)}}}} \dd t \dd s \\ & \qquad = \left(\int_{- \infty} ^{\infty} \varphi (t) U _{t} x \dd t, \int_{- \infty} ^{\infty} \varphi (s) U _{s} x \dd s\right) \geqq 0. \end{aligned}
\]

We shall show that Bochner's theorem is a consequence of Stone's theorem. 

Deduction of Theorem~\ref{thm:XI.13.Theorem.2} from Theorem~\ref{thm:XI.13.Theorem.1}. Consider the totality $\mathfrak{F}$ of complex-valued functions $x(t), -\infty < t < \infty$, such that $x(t) = 0$ except possibly for a finite set of values of $t$; the finite set may vary with $x$. $\mathfrak{F}$ is a pre-Hilbert space by 
\begin{equation}\label{eq:XI.13.4}
\begin{aligned}
&(x + y) (t) = x (t) + y (t), (\alpha x) (t) = \alpha x (t) \text{ and } (x, y) = \sum\limits_{- \infty < t, s < \infty} f (t - s) x (t) \overline {{y (s)}} \text{ for } x, y \in \mathfrak {F},
\end{aligned}
\end{equation}
excepting the axiom that $(x, x) = 0$ implies $x = 0$. That $(x, x) \geqq 0$ for any $x \in F$ is a simple consequence of the positive definiteness of the function $f(t)$. 

Let us set $\mathfrak{N} = \{x \in \mathfrak{F}; (x, x) = 0\}$. Then the factor space $\mathfrak{F}/\mathfrak{N}$ is a pre-Hilbert space with respect to the scalar product $(\bar{x}, \bar{y}) = (x, y)$ where $\bar{x}$ is the residue class mod $\mathfrak{N}$ containing $x \in \mathfrak{F}$. Let $X$ be the completion of the pre-Hilbert space $X = \mathfrak{F}/\mathfrak{N}$. The operator $U_{\tau}$ defined by 
\begin{equation}\label{eq:XI.13.5}
(U _{\tau} x) (t) = x (t - \tau), x \in \mathfrak {F},
\end{equation}
surely satisfies the conditions 
\begin{equation}\label{eq:XI.13.6}
(U _{\tau} x, U _{\tau} y) = (x, y), U _{\tau} U _{\sigma} = U _{\tau + \sigma} \text{ and } U _{0} = I.
\end{equation}

Therefore, it is easy to see that $\{U_{\tau}\}$ naturally defines a unitary operator $\hat{U}_{\tau}$ in $X$ in such a way that $\{\hat{U}_{\tau}\}, -\infty < \tau < \infty$, constitutes a one-parameter semi-group of class $(C_{0})$ of unitary operators in $X$; the strong continuity in $t$ of $\hat{U}_{t}$ follows from the continuity of the function $f(t)$. 

Hence, by Stone's theorem, $\hat{U}_t = \int_{-\infty}^{\infty} e^{it\lambda} dE(\lambda)$. Let $x_0(t) \in \mathfrak{F}$ be defined by $x_0(\tau) = 1$ and $x_0(t) = 0$ whenever $t \neq \tau$. Then, by \eqref{eq:XI.13.4} and \eqref{eq:XI.13.5}, $f(\tau) = (U_\tau x_0, x_0)$. Therefore, 
\[
f (\tau) = \int_{- \infty} ^{\infty} e ^{i \tau \lambda} \dd \| E (\lambda) \bar {x} _{0} \| ^{2},
\]
which proves Bochner's theorem. 

\begin{remarku} The idea of using a positive definite function to define a pre-Hilbert space as in \eqref{eq:XI.13.4} was systematically applied by B. Sz. \textsc{Nagy} \cite{nagy219} to various interesting problems concerning the Hilbert space. 
\end{remarku}


\yossection{14}{A Canonical Form of a Self-adjoint Operator with Simple Spectrum}
\label{sec:XI.14}

Let $H = \int \lambda dE(\lambda)$ be a self-adjoint operator in a Hilbert space $X$ with simple spectrum as defined in Chapter~\ref{chap:XI}, \ref{sec:XI.8}. Thus there exists an element $y \in X$ such that the set $\{(E(\beta) - E(\alpha))y; \alpha < \beta\}$ spans a dense linear subspace of $X$. We put 
\begin{equation}\label{eq:XI.14.1}
\sigma (\lambda) = (E (\lambda) y, y).
\end{equation}

Then $\sigma(\lambda)$ is monotone non-decreasing, right-continuous and bounded. We shall denote by $\sigma(B)$ the Baire measure determined on Baire sets on $\mathbb{R}^{1}$ from $\sigma((a,b]) = \sigma(b) - \sigma(a)$. We denote by $L_{\sigma}^{2}(-\infty,\infty)$ the totality of complex-valued Baire-measurable functions $f(\lambda), -\infty < \lambda < \infty$, such that $\|f\|_{\sigma} = \left(\int_{-\infty}^{\infty}|f(\lambda)|^{2}\sigma(\dd\lambda)\right)^{1/2} < \infty$. Then $L_{\sigma}^{2}(-\infty,\infty)$ is a Hilbert space by the scalar product $(f,g)_{\sigma} = \int_{-\infty}^{\infty}f(\lambda)\overline{g(\lambda)}\sigma(\dd\lambda)$ with the convention to consider $f = g$ in $L_{\sigma}^{2}$ iff $f(\lambda) = g(\lambda)\sigma$ -a.e. 

\begin{theoremu} With any $f(\lambda) \in L_{\sigma}^{2}(-\infty, \infty)$ we associate a vector $\hat{f}$ of $X$ defined by 
\begin{equation}\label{eq:XI.14.2}
\hat {f} = \int_{- \infty} ^{\infty} f (\lambda) \dd E (\lambda) y.
\end{equation}
\end{theoremu}

Then the correspondence $f(\lambda) \to \hat{f}$ is a one-one linear isometric mapping of $L_{\sigma}^{2}(-\infty, \infty)$ onto $X$. Let this mapping be denoted by $V$, i.e., $\hat{f} = Vf$. Then the operator $H_{1} = V^{-1}HV$ in $L_{\sigma}^{2}(-\infty, \infty)$ is precisely the 
operator of multiplication by $\lambda$: 
\begin{equation}\label{eq:XI.14.3}
\begin{aligned}
D (H _{1}) &= D (V ^{-1} H V) = \{f (\lambda); f (\lambda) \text{ and } \lambda f (\lambda) \text{ both } \in L _{\sigma} ^{2} (- \infty , \infty) \}
\\
\text{ and } (H _{1} f) (\lambda) &= \lambda f (\lambda) \text{ whenever } f (\lambda) \in D (H _{1}).
\end{aligned}
\end{equation}

\begin{proof}
 We have, by $E(\lambda)E(\mu) = E(\min (\lambda ,\mu))$, 
\[
\begin{aligned} (E (\lambda) y, \hat {f}) & = \int_{- \infty} ^{\infty} \overline {{f (\mu)}} d _{\mu} (E (\lambda) y, E (\mu) y) = \int_{- \infty} ^{\infty} \overline {{f (\mu)}} d _{\mu} (E (\mu) E (\lambda) y, y) \\ & = \int_{- \infty} ^{\lambda} \overline {{f (\mu)}} \dd (E (\mu) y, y) = \int_{- \infty} ^{\lambda} \overline {{f (\mu)}} \sigma (\dd \mu), \end{aligned}
\]
and so 
\begin{equation}\label{eq:XI.14.4}
(\hat {f}, \hat {g}) = \int_{- \infty} ^{\infty} f (\lambda) \dd (E (\lambda) y, \hat {g}) = \int_{- \infty} ^{\infty} f (\lambda) \overline {{g (\lambda)}} \sigma (\dd \lambda) = (f, g) _{\sigma}. \end{equation}

Therefore, $V$ maps $D(V) = L_{\sigma}^{2}(-\infty, \infty)$ onto $\{\hat{f}; \hat{f} = \int_{-\infty}^{\infty} f(\lambda) dE(\lambda) y, f \in L_{\sigma}^{2}(-\infty, \infty)\}$ one-to-one, linearly and isometrically. Hence, in particular, $R(V)$ is a closed linear subspace of $X$. But $R(V)$ surely contains the elements of the form $\int_{\alpha}^{\beta} dE(\lambda) y = (E(\beta) - E(\alpha)) y, -\infty < \alpha < \beta < \infty$, and so, by the hypothesis that the spectrum of $H$ is simple, we see that $R(V) = R(V)^a = X$. Thus the first half of the Theorem is proved. 

Next we have 
\[
\begin{aligned}
&E (\lambda) \hat {f}  = E (\lambda) \int_{- \infty} ^{\infty} f (\mu) \dd E (\mu) y = \int_{- \infty} ^{\infty} f (\mu) d _{\mu} (E (\lambda) E (\mu) y) = \int_{- \infty} ^{\lambda} f (\mu) \dd E (\mu) y,
\end{aligned}
\]
and so, by \eqref{eq:XI.14.4}, 
\begin{equation}\label{eq:XI.14.5}
(E (\lambda) \hat {f}, \hat {g}) = \int_{- \infty} ^{\lambda} f (\mu) \overline {{g (\mu)}} \sigma (\dd \mu).
\end{equation}

Hence the condition $\int_{-\infty}^{\infty}\lambda^{2}\dd(E(\lambda)\hat{f},\hat{f})<\infty$, equivalent to $\hat{f}\in D(H)$, is equivalent to the condition $\int_{-\infty}^{\infty}\lambda^{2}|f(\lambda)|^{2}\sigma(\dd\lambda)<\infty$. Moreover, in the last case, we have, by \eqref{eq:XI.12.20} in Chapter~\ref{chap:XI}, \ref{sec:XI.12}, 
\[
H V f = H \hat {f} = \int_{- \infty} ^{\infty} \lambda \dd E (\lambda) \hat {f},
\]
and hence, by \eqref{eq:XI.14.4} and \eqref{eq:XI.14.5}, 
\[
\begin{aligned} & {(H _{1} f, g) _{\sigma} = (V ^{-1} H V f, g) _{\sigma} = (H V \cdot f, V \cdot g) = (H \cdot \hat {f}, \hat {g})} \\ & {\qquad = \int_{- \infty} ^{\infty} \lambda \dd (E (\lambda) \hat {f}, \hat {g}) = \int_{- \infty} ^{\infty} \lambda f (\lambda) \overline {{g (\lambda)}} \sigma (\dd \lambda).} \end{aligned}
\]

On the other hand, we have 
\[
(H _{1} f, g) _{\sigma} = \int_{- \infty} ^{\infty} (H _{1} f) (\lambda) g (\lambda) \sigma (\dd \lambda).
\]

Therefore we must have 
\[
(H _{1} f) (\lambda) = \lambda f (\lambda) \sigma \text{-almost everywhere. }
\]
\end{proof}

\begin{remarku} There is a close connection between self-adjoint operators with simple spectrum and the Jacobi matrices. See M. H. \textsc{Stone} \cite{stone311}, p. $275$. For a canonical form of a self-adjoint operator with not necessarily simple spectrum, see J. von Neumann's reduction theory: \textsc{Neumann} [$9$]. 
\end{remarku}


\yossection{15}{The Defect Indices of a Symmetric Operator. The Generalized Resolution of the Identity}
\label{sec:XI.15}

\begin{definition}\label{def:XI.15.Definition.1} Let $U = U_H = (H - iI)(H + iI)^{-1}$ be the Cayley transform of a closed symmetric operator $H$ in a Hilbert space $X$. Let $X_H^+ = D(U_H)^\perp$, $X_H^- = R(U_H)^\perp$, and let $m = \dim(X_H^+)$, $n = \dim(X_H^-)$ be the dimension numbers of $X_H^+$, $X_H^-$ respectively. Then $H$ is said to be of the defect indices $(m, n)$. $H$ is self-adjoint iff it is of the defect indices $(0, 0)$ (see Chapter~\ref{chap:VII}, \ref{sec:VII.4}). 
\end{definition}

\begin{proposition}\label{prop:XI.15.Proposition.1} The defect indices $(m, n)$ of a closed symmetric operator $H$ may be defined as follows: $m$ is the dimension number of the linear subspace $\{x \in X; H^{*}x = ix\}$; $n$ is the dimension number of the linear subspace $\{x \in X; H^{*}x = -ix\}$. 
\end{proposition}

\begin{proof}
 Clear from Theorem~\ref{thm:VII.4.Theorem.3} in Chapter~\ref{chap:VII}, \ref{sec:VII.4}. 
\end{proof}

\begin{example}\label{ex:XI.15.Example.1} Let $X = L^2(0,1)$. Let $D$ be the totality of absolutely continuous functions $x(t) \in L^2(0,1)$ such that $x(0) = x(1) = 0$ and $x'(t) \in L^2(0,1)$. Then the operator $T_1$ defined by $T_1x(t) = i^{-1}x'(t)$ on $D = D(T_1)$ is of the defect indices ($1$, $1$). 
\end{example}

\begin{proof}
 As was shown in Example~\hyperref[ex:VII.3.Example.4]{$4$} of Chapter~\ref{chap:VII}, \ref{sec:VII.3}, $T_1^* = T_2$ is defined by 
\[
T _{2} x (t) = i ^{-1} x ^{\prime} (t) \text{ on } D (T _{2}) = \{x (t) \in L ^{2} (0, 1); x (t) \text{ is }
\]
absolutely continuous such that $x'(t) \in L^2(0, 1)$. 

Thus the solution $y \in L^2(0,1)$ of $T_1^*y = T_2y = iy$ is a distribution solution of the differential equation 
\begin{equation}\label{eq:XI.15.1}
y ^{\prime} (t) = - y (t) \left(y, y ^{\prime} \in L ^{2} (0, 1)\right).
\end{equation}

Then $z(t) = y(t) \exp(t)$ is a distribution solution of the differential equation 
\begin{equation}\label{eq:XI.15.2}
z ^{\prime} (t) = 0 \left(z, z ^{\prime} \in L ^{2} (0, 1)\right).
\end{equation}

We shall show that there exists a constant $C$ such that $z(t) = C$ for a.e. $t \in (0,1)$. To this purpose, take any function $x_0(t) \in C_0^\infty(0,1)$ such that 
\[
\begin{aligned}
&\int_{0} ^{1} x _{0} (t) \dd t  = 1 \text{ and put } x (t) - x _{0} (t) \int_{0} ^{1} x (t) \dd t  = u (t), w (t) = \int_{0} ^{t} u (s) \dd s,
\end{aligned}
\]
where $x(t)$ is an arbitrary function from $C_0^\infty(0,1)$. We have $w \in C_0^\infty(0,1)$ by $\int_{0}^{1} u(s) \dd s = 0$. Therefore, by \eqref{eq:XI.15.2}, we have 
\[
- \int_{0} ^{1} z (t) w ^{\prime} (t) \dd t = - \int_{0} ^{1} z (t) u (t) \dd t = 0,
\]
that is, 
\[
\int_{0} ^{1} z (t) x (t) \dd t = C \int_{0} ^{1} x (t) \dd t, \text{ where } C = \int_{0} ^{1} z (t) x _{0} (t) \dd t.
\]

This proves, by the arbitrariness of $x(t) \in C_0^\infty(0,1)$, that $z(t) = C$ for a.e. $t \in (0,1)$. 

Hence any solution of $T^{*}y = iy$ is of the form $y(t) = C\exp (-t)$. In the same manner we see that any solution of $T^{*}y = -iy$ is of the form $y(t) = C\exp (t)$. Thus $T$ is of the defect indices ($1$, $1$). 
\end{proof}

\begin{definition}\label{def:XI.15.Definition.2} $A$ symmetric operator $H$ in a Hilbert space $X$ is called maximal symmetric if there is no proper symmetric extension of $H$. 
\end{definition}

\begin{proposition}\label{prop:XI.15.Proposition.2} $A$ maximal symmetric operator $H$ is closed and $H = H^{**}$. A self-adjoint operator $H$ is maximal symmetric. 
\end{proposition}

\begin{proof}
 By Proposition~\ref{prop:VII.3.Proposition.1} in Chapter~\ref{chap:VII}, \ref{sec:VII.3}, $H^{**}$ is a closed symmetric extension of $H$. Thus the first half of Proposition~\ref{prop:XI.15.Proposition.2} is clear. Let $H_0$ be a symmetric extension of a self-adjoint operator $H$. Then, from $H \subseteq H_0$, $H_0 \subseteq H_0^*$, we obtain $H_0 \subseteq H_0^* \subseteq H^*$ and so, by $H = H^*$, $H \subseteq H_0 \subseteq H$. This proves that a self-adjoint operator $H$ is maximal symmetric. 
\end{proof}

\begin{corollary}\label{cor:XI.15.Corollary.1} Every maximal symmetric extension $H_0$ of a given symmetric operator $H$ is also an extension of $H^{**}$. 
\end{corollary}

\begin{proof}
 The relation $H \subseteq H_0$ implies the relation $H_0^* \subseteq H^*$ and $H^{**} \subseteq H_0^{**}$. Since, by the preceding proposition, $H_0 = H_0^{**}$, Corollary~\ref{cor:XI.15.Corollary.1} is true. 
\end{proof}

\begin{corollary}\label{cor:XI.15.Corollary.2} If $H$ is a symmetric operator such that $H^{*} = H^{**}$, then the self-adjoint operator $H^{*}$ is the only maximal symmetric extension of $H$. 
\end{corollary}

\begin{proof}
 Being self-adjoint, $H^{**}$ is maximal symmetric. Thus any maximal symmetric extension $H_{0}$ of $H$, which is also a symmetric extension of $H^{**}$ by Corollary~\ref{cor:XI.15.Corollary.1}, is identical to $H^{**} = H^{*}$. 
\end{proof}

We are thus in a position to state 

\begin{definition}\label{def:XI.15.Definition.3} $A$ symmetric operator $H$ such that $H^{*} = H^{**}$ is said to be essentially self-adjoint. A self-adjoint operator $H$ is said to be hypermaximal. The latter terminology is due to J. \textsc{von} \textsc{Neumann}. 
\end{definition}

\begin{example}\label{ex:XI.15.Example.2} Let $X = L^2(-\infty, \infty)$, and define an operator $H$ by $Hx(t) = t \cdot x(t)$ for $x(t) \in C_0^0(-\infty, \infty)$. $H$ is surely a symmetric operator in $X$. It is easy to see that $H^*$ is the coordinate operator defined in Example~\hyperref[ex:VII.3.Example.2]{$2$} in Chapter~\ref{chap:VII}, \ref{sec:VII.3} so that $H$ is essentially self-adjoint. The operator $H$ defined by $Hx(t) = i^{-1}x'(t)$ for $x(t) \in C_0^1(-\infty, \infty)$ is also essentially self-adjoint in $X = L^2(-\infty, \infty)$. For, in this case, $H^*$ is the momentum operator defined in Example~\hyperref[ex:VII.3.Example.3]{$3$} in Chapter~\ref{chap:VII}, \ref{sec:VII.3}. 
\end{example}

\begin{theorem}\label{thm:XI.15.Theorem.1} Let the defect indices $(m, n)$ of a closed symmetric operator $H$ satisfy 
\[
m = m ^{\prime} + p, n = n ^{\prime} + p \quad (p > 0).
\]
\end{theorem}

Then there exists a closed symmetric extension $H'$ of $H$ with the defect indices $(m', n')$. 

\begin{proof}
 Let $\{\varphi_1, \varphi_2, \ldots, \varphi_p, \varphi_{p+1}, \ldots, \varphi_{p+m}\}, \{\psi_1, \psi_2, \ldots, \psi_p, \psi_{p+1}, \ldots, \psi_{p+n'}\}$ be complete orthonormal systems of $X_H^+ = D(U_H)^\perp$, $X_H^- = R(U_H)^\perp$, respectively. Define an isometric extension $V$ of $U_H$ by 
\begin{align}
V x& = U _{H} x \quad \text{ for } \quad x \in D (U _{H}), \notag \\
V \cdot \sum\limits_{i=1} ^{p} \alpha_{i} \varphi_{i}& = \sum\limits_{i=1} ^{p} \alpha_{i} \psi_{i}. \notag
\end{align}

We have $R(I - U_{H})^{a} = X$ by Theorem~\ref{thm:VII.4.Theorem.1} in Chapter~\ref{chap:VII}, \ref{sec:VII.4}. Thus, by $R(I - V) \supseteq R(I - U_{H})$ and Theorem~\ref{thm:VII.4.Theorem.2} in Chapter~\ref{chap:VII}, \ref{sec:VII.4}, there exists a uniquely determined closed symmetric extension $H'$ of $H$ such that $V = (H' - iI)(H' + iI)^{-1}$. The defect indices of $H'$ are $(m', n')$ as may be seen from $\dim(D(V)^{\perp}) = m'$, $\dim(R(V)^{\perp}) = n'$. 
\end{proof}

\begin{corollaryu} $A$ closed symmetric operator $H$ with the defect indices $(m, n)$ is maximal symmetric iff either $m = 0$ or $n = 0$. 
\end{corollaryu}

\begin{proof}
 The ``only if'' part is clear from Theorem~\ref{thm:XI.15.Theorem.1}. If, for instance, $m = 0$, then, by $D(U_{H}) = X$, we must have $U_{H_{0}} = U_{H}$ for a closed symmetric extension $H_{0}$ of $H$. This proves $H_{0} = i(I + U_{H_{0}})(I - U_{H_{0}})^{-1} = i(I + U_{H})(I - U_{H})^{-1} = H$. Also, in the case $n = 0$, we see that there is no proper closed symmetric extension of $H$. 
\end{proof}

\begin{example}\label{ex:XI.15.Example.3} Suppose that $\{\varphi_1, \varphi_2, \ldots, \varphi_n, \ldots\}$ is a complete ortho-normal system of a (separable) Hilbert space $X$. Then, by 
\[
U \cdot \sum\limits_{i=1} ^{\infty} \alpha_{i} \varphi_{i} = \sum\limits_{i=1} ^{\infty} \alpha_{i} \varphi_{i+1}, \sum\limits_{i=1} ^{\infty} | \alpha_{i} | ^{2} < \infty
\]
we define a closed isometric operator $U$ such that $D(U) = X$ and $\dim (R(U)^\perp) = 1$. If $R(I - U)^a \neq X$, then there exists an $x \neq 0$ such that $x \in R(I - U)^\perp$. Consequently $((I - U)x, x) = 0$ and so $(Ux, x) = \|x\|^2 = \|Ux\|^2$. This implies that 
\begin{equation}\left\| (I - U) x \right\| ^{2} = \left\| x \right\| ^{2} - (U x, x) - (x, U x) + \left\| U x \right\| ^{2} = \| x \| ^{2} - \| x \| ^{2} - \| x \| ^{2} + \| x \| ^{2} = 0, \notag\end{equation}
that is, $Ux = x$ and so, by the above definition of $U$, $x$ must be zero. This contradiction shows that $R(I - U)^a = X$. Hence, by Theorem~\ref{thm:VII.4.Theorem.2} in Chapter~\ref{chap:VII}, \ref{sec:VII.4}, $U$ is the Cayley transform of a closed symmetric operator $H$. $H$ is of the defect indices ($0$, $1$) since $D(U) = X$ and $R(U)^\perp$ is spanned by $\varphi_1$. Thus $H$ is maximal symmetric without being self-adjoint. 
\end{example}

\begin{theorem}[M. A. \textsc{Naimark} {\cite{naimark223}}]\label{thm:XI.15.Theorem.2} Let a closed symmetric operator $H_{1}$ in a Hilbert space $X_{1}$ be of the defect indices $(m, n)$. Then we can construct a Hilbert space $X$, containing $X_{1}$ as a closed linear subspace, and a closed symmetric operator $H$ in $X$ with the defect indices $(m + n, m + n)$ such that 

$H_{1} = P(X_{1})HP(X_{1})$, where $P(X_{1})$ is the projection of $X$ onto $X_{1}$ 
\end{theorem}

\begin{proof}
 Consider a Hilbert space $X_{2}$ of the same dimensionality as $X_{1}$. We construct a closed symmetric operator $H_{2}$ in $X_{2}$ with the defect indices $(n, m)$. For instance, we may take $H_{2} = -H_{1}$ supposing that $X_{2}$ coincides with $X_{1}$. Then we would have $\{x \in X_{2}; H_{2}^{*}x = ix\} = \{x \in X_{1}; H_{1}^{*}x = -ix\}, \{x \in X_{2}; H_{2}^{*}x = -ix\} = \{x \in X_{1}; H_{1}^{*}x = ix\}$ and so the defect indices of $H_{2}$ must be $(n, m)$. 

We then consider an operator $H$ defined by 
\[
H \{x, y \} = \left\{H _{1} x, H _{2} y \right\} \text{ for } \{x, y \} \in X _{1} \times X _{2} \text{ with } x \in D (H _{1}), y \in D (H _{2}).
\]

It is easy to see that $H$ is a closed symmetric operator in the product Hilbert space $X = X_{1} \times X_{2}$. The condition 
\[
H ^{*} \{x, y \} = i \{x, y \} \quad (\text{ or the condition } H ^{*} \{x, y \} = - i \{x, y \})
\]
means that $H_1^* x = ix, H_2^* y = iy$ (or $H_1^* x = -ix, H_2^* y = -iy$). Hence we see that the defect indices of $H$ are $(m + n, m + n)$. 
\end{proof}

\begin{corollaryu} Let, by Theorem~\ref{thm:XI.15.Theorem.1}, $\hat{H}$ be a self-adjoint extension of $H$, and let $\hat{H} = \int \lambda \dd\hat{E} (\lambda)$ be the spectral resolution of $\hat{H}$. Since $H$ and a fortiori $\hat{H}$ are extensions of $H_{1}$ when $H_{1}$ is considered as an operator in $X$, we have the result: 
\begin{gather}\text{ if } x \in D (H _{1}) \subseteq X _{1} = P (X _{1}) X, \text{ then } x = P (X _{1}) x \in D (\hat {H}) \notag \\H _{1} x = P (X _{1}) \hat {H} x = P (X _{1}) \hat {H} P (X _{1}) x = \int_{- \infty} ^{\infty} \lambda \dd F (\lambda) x, \text{ where } F (\lambda) = P (X _{1}) \hat {E} (\lambda) P (X _{1}). \label{eq:XI.15.3}
\end{gather}
\end{corollaryu}

The system $\{F(\lambda); -\infty < \lambda < \infty\}$ surely satisfies the conditions: 

$F(\lambda)$ is a self-adjoint operator in $X_{1}$, 
\begin{gather}\lambda_{1} < \lambda_{2} \text{ implies that } (F (\lambda_{1}) x, x) \leqq (F (\lambda_{2}) x, x) \text{ for every } x \in X _{1}, F (\lambda + 0) = F (\lambda), \label{eq:XI.15.4} \\F (- \infty) x = \slim _{\lambda \downarrow - \infty} F (\lambda) x = 0, F (\infty) x = \slim _{\lambda \uparrow \infty} F (\lambda) x = x \text{ for all } x \in X. \notag
\end{gather}

\begin{remarkn}[Remark $1$] For the closed symmetric operator $H_{1}$, we have 
\[
H _{1} x = \int_{- \infty} ^{\infty} \lambda \dd F (\lambda) x \text{ for all } x \in D (H _{1}),
\]
\end{remarkn}

where $\{F(\lambda);-\infty<\lambda<\infty\}$ satisfies \eqref{eq:XI.15.4}. In this sense, $H_{1}$ admits a generalized spectral resolution. 

\begin{example}\label{ex:XI.15.Example.4} Let $X_{1} = L^{2}(-\infty,0)$, and let $D_{1}$ be the totality of absolutely continuous functions $x(t) \in L^{2}(-\infty,0)$ such that $x(0) = 0$ and $x'(t) \in L^{2}(-\infty,0)$. Then the operator $H_{1}$ defined by $H_{1}x(t) = i^{-1}x'(t)$ on $D_{1} = D(H_{1})$ is a maximal symmetric operator with the defect indices ($0$, $1$). This we see as in the above Example~\hyperref[ex:XI.15.Example.1]{$1$}. Let $X_{2} = L^{2}(0,\infty)$, and let $D_{2}$ be the totality of absolutely continuous functions $x(t) \in L^{2}(0,\infty)$ such that $x(0) = 0$ and $x'(t) \in L^{2}(0,\infty)$. Then the operator $H_{2}$ defined by $H_{2}x(t) = i^{-1}x'(t)$ on $D_{2} = D(H_{2})$ is a maximal symmetric operator with the defect indices ($1$, $0$). In this case, $X = X_{1} \times X_{2} = L^{2}(-\infty,\infty)$, and the operator $H$ in Theorem~\ref{thm:XI.15.Theorem.2} is given concretely by $Hx(t) = i^{-1}x'(t)$ for those $x(t) \in L^{2}(-\infty,\infty)$ for which $x(0) = 0$ and $x'(t) \in L^{2}(-\infty,\infty)$. 
\end{example}

\begin{remarkn}[Remark $2$] Since $\hat{H} = \int \lambda \dd\hat{E} (\lambda)$ is an extension of $H_{1}$, we have, for $x \in D(H_{1})$, 
\begin{equation}\label{eq:XI.15.5}
\begin{aligned} \| H _{1} x \| ^{2} & = \| \hat {H} x \| ^{2} = \int_{- \infty} ^{\infty} \lambda^{2} \dd \| \hat {E} (\lambda) x \| ^{2} = \int_{- \infty} ^{\infty} \lambda^{2} \dd (\hat {E} (\lambda) x, x) \\ & = \int_{- \infty} ^{\infty} \lambda^{2} \dd (\hat {E} (\lambda) P (X _{1}) x, P (X _{1}) x) = \int_{- \infty} ^{\infty} \lambda^{2} \dd (F (\lambda) x, x). \end{aligned}
\end{equation}
\end{remarkn}

However, the condition $\int_{-\infty}^{\infty} \lambda^2 \dd(F(\lambda)x, x) < \infty$ does not necessarily imply that $x \in D(H_1)$. Concerning this point we have 

\begin{theorem}\label{thm:XI.15.Theorem.3} In the case of a maximal symmetric operator $H_{1}$, we have, for the corresponding generalized spectral resolution $\int \lambda dF(\lambda)$ in $X_{1}$, 

$x \in D(H_1)$ is equivalent to the condition $\int_{-\infty}^{\infty} \lambda^2 \dd(F(\lambda)x, x) < \infty$. 
\end{theorem}


\begin{proof}
 The reasoning in Remark $2$ shows that $\int_{-\infty}^{\infty} \lambda^2 \dd(F(\lambda)x, x) < \infty$ implies that $\int_{-\infty}^{\infty} \lambda^2 \dd\|\hat{E}(\lambda)x\|^2 < \infty$, i.e., $x \in D(\hat{H})$. The operator 
\[
H ^{\prime} = P (X _{1}) \hat {H} P (X _{1}),
\]
when considered as an operator in $X_{1}$, is a symmetric extension of $H_{1}$ and $D(H') = D(\hat{H}) \cap X_{1}$. Thus, by the maximality of $H_{1}$, we must have $H_{1} = H'$. Thus $D(H_{1}) = D(H') = D(\hat{H}) \cap X_{1}$ and so, by $H' = P(X_{1})\hat{H}P(X_{1})$, the condition $\int_{-\infty}^{\infty}\lambda^{2}\dd\|\hat{E}(\lambda)x\|^{2} < \infty$ with $x \in X_{1}$ implies 
$\int_{-\infty}^{\infty}\lambda^{2}\dd(F(\lambda)x,x) < \infty$; and conversely, the condition $x\in D(H_1)$ implies $\int_{-\infty}^{\infty}\lambda^{2}\dd(F(\lambda)x,x) = \int_{-\infty}^{\infty}\lambda^{2}\dd\|\hat{E} (\lambda)x\|^{2} < \infty .$ 
\end{proof}

\begin{remarkn}[Remark $3$] Since, by Theorem~\ref{thm:XI.15.Theorem.1}, any closed symmetric operator can be extended to a maximal symmetric operator $H_{1}$, we can apply our Theorem~\ref{thm:XI.15.Theorem.3} to the effect that ($5$) is valid. For a detailed exposition of the generalized spectral resolution, see N. I. \textsc{Achieser}-I. M. \textsc{Glazman} \cite{achieser1} or B. Sz. \textsc{Nagy} \cite{nagy219}. The spectral representation of a self-adjoint operator in a Hilbert space can be extended to a certain class of linear operators in a Banach space with an appropriate modification. This result is due to N. \textsc{Dunford} and may be considered as an ``elementary divisor theory'' in infinite dimensional spaces. N. \textsc{Dunford}-J. \textsc{Schwartz} \cite{dunford49}. 
\end{remarkn}


\yossection{16}{The Group-ring \texorpdfstring{$L^1$}{L1} and Wiener's Tauberian Theorem}
\label{sec:XI.16}

The Gelfand representation admits another important application in functional analysis, namely, an operator-theoretical treatment of the Tauberian theorem of N. \textsc{Wiener}. 

The linear space $L^{1}(-\infty,\infty)$ is a ring with respect to the function sum and the product $\times$ defined by 
\begin{equation}\label{eq:XI.16.1}
(f \times g) (t) = (f * g) (t) = \int_{- \infty} ^{\infty} f (t - s) g (s) \dd s.
\end{equation}

For, by the Fubini-Tonelli theorem, 
\[
\begin{aligned}
&\int_{- \infty} ^{\infty} \left| \int_{- \infty} ^{\infty} f (t - s) g (s) \dd s \right| \dd t  \leqq \int_{- \infty} ^{\infty} | f (t - s) | \dd t \int_{- \infty} ^{\infty} | g (s) | \dd s = \int_{- \infty} ^{\infty} | f (t) | \dd t \int_{- \infty} ^{\infty} | g (t) | \dd t.
\end{aligned}
\]

We have thus proved 
\begin{equation}\label{eq:XI.16.2}
\left\| f \times g \right\| \leqq \left\| f \right\| \cdot \left\| g \right\|, \text{ where } \| \| \text{ is the norm in } L ^{1} (- \infty , \infty).
\end{equation}

Therefore we have proved 

\begin{proposition}\label{prop:XI.16.Proposition.1} We can introduce formally a multiplicative unit $e$ in the ring $L^{1}$ consisting of all elements $\tilde{z}$ given formally by 
\begin{equation}\label{eq:XI.16.3}
\tilde {z} = \lambda e + x, x \in L ^{1} (- \infty , \infty).
\end{equation}
\end{proposition}

In fact $L^1$ is a normed ring by the following rule: 
\begin{equation}\label{eq:XI.16.4}
\begin{gathered}
(\lambda_{1} e + x _{1}) + (\lambda_{2} e + x _{2}) = (\lambda_{1} + \lambda_{2}) e + (x _{1} + x _{2}), \alpha (\lambda e + x) = \alpha \lambda e + \alpha x,
\\
(\lambda_{1} e + x _{1}) \times (\lambda_{2} e + x _{2}) = \lambda_{1} \lambda_{2} e + \lambda_{1} x _{2} + \lambda_{2} x _{1} + x _{1} \times x _{2}, \| \lambda e + x \| = | \lambda | + \| x \|.
\end{gathered}
\end{equation}

This normed ring $L^{1}$ is called the group-ring of the group $\mathbb{R}^{1}$ of real numbers written additively. We shall find all maximal ideals of this ring $L^{1}$. One is $I_{0}=L^{1}(-\infty,\infty)$ from which $L^{1}$ is obtained by adjunction of the unit $e$. We shall find all maximal ideals $I\neq I_{0}$. 

For any maximal ideal $I$ of the normed ring $L^1$, we shall denote by $(\tilde{z}, I)$ the complex number which corresponds to the element $\tilde{z}$ by the ring homomorphic mapping $L^1 \to L^1 / I$. Thus $(\tilde{z}, I_0) = \lambda$ for $\tilde{z} = \lambda e + x, x \in I_0$. 

Let $I$ be a maximal ideal of $L^1$ different from $I_0$. Then there exists a function $x \in L^1(-\infty, \infty) = I_0$ such that $(x, I) \neq 0$. We set 
\begin{equation}\label{eq:XI.16.5}
\chi (\alpha) = (x _{\alpha}, I) / (x, I), \text{ where } x _{\alpha} (t) = x (t + \alpha).
\end{equation}

Then $\chi(0) = 1$, and, by $|(x_{\alpha}, I)| \leqq \|x_{\alpha}\| = \|x\|$, we see that $|\chi(\alpha)| \leqq \|x\| / |(x, I)|$. Thus the function $\chi(\alpha)$ is bounded in $\alpha$. Moreover, by 
\[
\left| \chi (\alpha + \delta) - \chi (\alpha) \right| \leqq \left\| x _{\alpha + \delta} - x _{\alpha} \right\| / |(x, I)|,
\]
the function $\chi (\alpha)$ is continuous in $\alpha$. For, by \eqref{eq:X.1.2} in Chapter~\ref{chap:X}, \ref{sec:X.1}, we have 
\[
\lim\limits _{\delta \rightarrow 0} \int_{- \infty} ^{\infty} | x (\alpha + \delta + t) - x (\alpha + t) | \dd t = 0.
\]

On the other hand, we have, by $(x_{\alpha + \beta} \times x)(t) = (x_{\alpha} \times x_{\beta})(t)$, 
\[
\left(x _{\alpha + \beta}, I\right) (x, I) = \left(x _{\alpha}, I\right) \left(x _{\beta}, I\right)
\]
so that 
\begin{equation}\label{eq:XI.16.6}
\chi (\alpha + \beta) = \chi (\alpha) \chi (\beta).
\end{equation}

From this we can prove that there exists a uniquely determined real number $\xi = \xi(I)$ such that 
\begin{equation}\label{eq:XI.16.7}
\chi (\alpha) = \exp (i \cdot \xi (I) \cdot \alpha).
\end{equation}

In fact, by $\chi(n\alpha) = \chi(\alpha)^n$ and the boundedness of the function $\chi$, we obtain $|\chi(\alpha)| \leqq 1$. Consequently, by $\chi(\alpha) \chi(-\alpha) = \chi(0) = 1$, we must have $|\chi(\alpha)| = 1$. Thus, as a continuous solution of \eqref{eq:XI.16.6} whose absolute value is one, the function $\chi(\alpha)$ must be given in form \eqref{eq:XI.16.7}. That the value $\xi(I)$ is determined by $I$ independently of the choice of $x \in L^1(-\infty, \infty)$ may be seen from $\chi(\alpha)(y, I) = (y_\alpha, I)$ which is implied by $x_\alpha \times y = x \times y_\alpha$. 

Every continuous solution of \eqref{eq:XI.16.6} with the absolute value identically equal to one is called a continuous unitary character of the group $\mathbb{R}^{1}$ of real numbers written additively. 

We have thus constructed the (continuous unitary) character $\chi(\alpha)$ of the group $\mathbb{R}^{1}$ with respect to the given maximal ideal $I \neq I_{0}$. 

We next show how to reconstruct the ideal $I$ with respect to this character or, what amounts to the same thing, to reconstruct the value $(\tilde{z}, I)$ from $\chi$. 

For any $y(t) \in L^{1}(-\infty, \infty)$, we have 
\[
(x \times y) (t) = \int_{- \infty} ^{\infty} x (t - s) y (s) \dd s = \int_{- \infty} ^{\infty} x _{-s} (t) y (s) \dd s.
\]

Hence, by \eqref{eq:XI.16.5} and the continuity of the ring homomorphism $L^1 \to L^1 / I$, we obtain 
\[
(x \times y, I) = (x, I) (y, I) = (x, I) \int_{- \infty} ^{\infty} \chi (- s) y (s) \dd s.
\]

Thus, by $(x, I) \neq 0$, we obtain 
\begin{equation}\label{eq:XI.16.8}
(y, I) = \int_{- \infty} ^{\infty} y (s) \exp (- i \cdot \xi (I) \cdot s) \dd s.
\end{equation}

Therefore, for any $\tilde{z} = \lambda e + x$ with $x \in L^{1}(-\infty, \infty)$, 
\begin{equation}\label{eq:XI.16.9}
(\tilde {z}, I) = (\lambda e, I) + (x, I) = \lambda + \int_{- \infty} ^{\infty} x (s) \exp (- i \cdot \xi (I) \cdot s) \dd s.
\end{equation}

Conversely, any (continuous, unitary) character $\chi(\alpha) = \exp(i \cdot \xi \cdot \alpha)$ of the group $\mathbb{R}^{1}$ defines a ring homomorphism 
\begin{equation}\label{eq:XI.16.10}
\tilde {z} \rightarrow \lambda + \int_{- \infty} ^{\infty} x (s) \exp (- i \cdot \xi \cdot s) \dd s (\tilde {z} = \lambda e + x, x \in L ^{1} (- \infty , \infty))
\end{equation}
of $L^{1}$ onto the ring of complex numbers. For, we have, by the Fubini-Tonelli theorem. 
\begin{equation}\label{eq:XI.16.11}
\begin{aligned}
&\int_{- \infty} ^{\infty} (x _{1} \times x _{2}) (t) \exp (- i \cdot \xi \cdot t) \dd t = \int_{- \infty} ^{\infty} x _{1} (t) \exp (- i \cdot \xi \cdot t) \dd t \int_{- \infty} ^{\infty} x _{2} (t) \exp (- i \cdot \xi \cdot t) \dd t.
\end{aligned}
\end{equation}

We have thus proved 

\begin{theorem}[\textsc{Gelfand} {\cite{gelfand78}} and \textsc{Raikov} {\cite{raikov267}}]\label{thm:XI.16.Theorem.1} There exists a one-to-one correspondence between the set of all maximal ideals $I \neq L^{1}(-\infty, \infty)$ of the group ring $L^{1}$ of the group $\mathbb{R}^{1}$ and the set of all continuous, unitary characters $\chi(\alpha)$ of this group $\mathbb{R}^{1}$. This correspondence is defined by formula \eqref{eq:XI.16.9}. 
\end{theorem}

We shall show that the normed ring $L^{1}$ is semi-simple or, what amounts to the same thing (see Chapter~\ref{chap:XI}, \ref{sec:XI.2}), that the following theorem is true. 

\begin{theorem}\label{thm:XI.16.Theorem.2} The normed ring $L^{1}$ has no generalized nilpotent elements other than $0$. 
\end{theorem}

\begin{proof}
 Let $x \in L^{1}(-\infty, \infty)$ and $y \in L^{2}(-\infty, \infty)$. Then, by Schwarz' inequality, 
\[
\left| \int_{- \infty} ^{\infty} x (t - s) y (s) \dd s \right| \leqq \left(\int_{- \infty} ^{\infty} | x (t - s) | \dd s \int_{- \infty} ^{\infty} | x (t - s) | | y (s) | ^{2} \dd s\right) ^{1/2},
\]
and so, by the Fubini-Tonelli theorem, we see that the left hand side belongs to $L^2(-\infty, \infty)$ and 
$\|x \times y\|_{2} \leqq \|x\| \|y\|_{2}, \text{ where } \|\|_{2} \text{ is the norm in } L^{2_1}(-\infty, \infty).$ 


Hence we can define a bounded linear operator $T_{x}$ on $L^{2}(-\infty, \infty)$ into $L^{2}(-\infty, \infty)$ by 
\begin{equation}\label{eq:XI.16.12}
(T _{x} y) (t) = \int_{- \infty} ^{\infty} x (t - s) y (s) \dd s \quad \text{ whenever } \quad x \in L ^{1} (- \infty , \infty);
\end{equation}
and, moreover, we have 
\begin{align}
\left\| T _{x} \right\| _{2}& \leqq \left\| x \right\|, \label{eq:XI.16.13} \\
T _{x} ^{*}& = T _ {x ^{*}} \text{ where } x ^{*} (t) = \overline {{x (- t)}}. \label{eq:XI.16.14}
\end{align}

Thus, by applying the Fubini-Tonelli theorem again, we obtain 
\begin{equation}\label{eq:XI.16.15}
T_{x}T_{x}^{*} = T_{x\times x^{*}} = T_{x^{*}\times x} = T_{x}^{*}T_{x}, \text{ i.e., } T_{x} \text{ is a normal operator.}
\end{equation}

Hence, by Chapter~\ref{chap:XI}, \ref{sec:XI.3}, we have $\| T_x \|_2 = \lim\limits_{n \to \infty} (\|T_x^n \|_2)^{1/n}$. Therefore, by 
\[
\left| \underbrace {\left| x \times x \times x \times \cdots \times x \right|} _ {n \text{ factor }} \right| \geqq \left\| T _{x} ^{n} \right\| _{2},
\]
we see that, if $x$ is a generalized nilpotent element, then $\| T_x \|_2 = 0$. From this fact we easily prove that $x$ must be a zero vector of $L^1(-\infty, \infty)$. 

Let now $\tilde{z} = \lambda e + x$ with $x \in L^{1}(-\infty, \infty)$ be a generalized nilpotent element of the normed ring $L^{1}$. Then, for any maximal ideal $I$ of $L^{1}$, we must have $(\tilde{z}, I) = \lambda + (x, I) = 0$. Hence the Fourier transform 
\[
(2 \pi) ^{-1/2} \int_{- \infty} ^{\infty} x (t) \exp (- i \cdot \xi \cdot t) \dd t
\]
must be identically equal to $-(2\pi)^{-1/2}\lambda$. From this we prove that $\lambda$ must be $0$. For, we have (the Riemann-Lebesgue theorem) 
\begin{align}
\left| \int_{- \infty} ^{\infty} x (t) \exp (- i \cdot \xi \cdot t) \dd t \right|& = \left| 2 ^{-1} \int_{- \infty} ^{\infty} \left[ x (t) - x \left(t + \frac {\pi}{\xi}\right) \right] \exp (- i \cdot \xi \cdot t) \dd t \right| \notag \\
&\leqq 2 ^{-1} \int_{- \infty} ^{\infty} \left| x (t) - x \left(t + \frac {\pi}{\xi}\right)\right| \dd t \rightarrow 0 \text{ as } \xi \rightarrow \infty . \notag
\end{align}

Thus a generalized nilpotent element $\tilde{z} \in L^1$ must be of the form $\tilde{z} = x$ with $x \in L^1(-\infty, \infty)$ and so, by what we have proved above, $\tilde{z} = x = 0$. 
\end{proof}

We are now in a position to state and prove the Tauberian theorem of N. \textsc{Wiener} \cite{wiener340}: 

\begin{theorem}\label{thm:XI.16.Theorem.3} Let $x(t) \in L^{1}(-\infty, \infty)$ be such that its Fourier transform 
\[
(2 \pi) ^{-1/2} \int_{- \infty} ^{\infty} x (t) \exp (- i \cdot \xi \cdot t) \dd t
\]
does not vanish for any real $\xi$. Then, for any $y(t) \in L^{1}(-\infty, \infty)$ and $\varepsilon > 0$, we can find the real numbers $\beta$'s, the complex numbers $\alpha$'s and a positive integer $N$ in such a way that 
\begin{equation}\label{eq:XI.16.16}
\int_{- \infty} ^{\infty} \left| y (t) - \sum\limits_{j=1} ^{N} \alpha_{j} x (t - \beta_{j}) \right| \dd t < \varepsilon .
\end{equation}
\end{theorem}

\begin{proof}
 It suffices to find $\tilde{z} \in L^1$ in such a way that 
\begin{equation}\label{eq:XI.16.17}
\left\| y - x \times \tilde {z} \right\| \leqq \varepsilon / 2. \end{equation}

We first show that 
\begin{equation}\lim\limits _{\alpha \rightarrow \infty} \int_{- \infty} ^{\infty} | y (t) - y ^{(\alpha)} (t) | \dd t = 0, \text{ where } y ^{(\alpha)} (t) = \frac {1}{\pi} \int_{- \infty} ^{\infty} y (t - s) \frac {1 - \cos \alpha s}{\alpha s ^{2}} \dd s. \label{eq:XI.16.18}\end{equation}

For, by $\int_{-\infty}^{\infty}(1 - \cos \alpha s)(\alpha s^2)^{-1}\dd s = \pi (\alpha >0)$, the left side is 
\[
\leqq (\pi) ^{-1} \lim\limits _{\alpha \rightarrow \infty} \int_{- \infty} ^{\infty} \frac {(1 - \cos s)}{s ^{2}} \dd s \left\{\int_{- \infty} ^{\infty} \left| y (t) - y \left(t - \frac {s}{\alpha}\right)\right| \dd t \right\} = 0.
\]


Next we have 
\begin{equation}\label{eq:XI.16.19}
(2 \pi) ^{-1/2} \int_{- \infty} ^{\infty} \left(\frac {2}{\pi}\right) ^{1/2} \frac {1 - \cos \alpha u}{\alpha u ^{2}} e ^{- i u \xi} \dd u = \begin{cases}1 - | \xi | / \alpha & (| \xi | < \alpha), \\ 0 & (| \xi | \geqq \alpha).\end{cases}
\end{equation}

For, we have 
\[
\begin{aligned}
&(2 \pi) ^{-1/2} \int_{- \alpha} ^{\alpha} (1 - | \xi | / \alpha) e ^{i \xi u} \dd \xi  = \left(\frac {2}{\pi}\right) ^{1/2} \int_{0} ^{\alpha} \left(1 - \frac {\xi}{\alpha}\right) \cos u \xi \cdot \dd \xi \\
&= \left(\frac {2}{\pi}\right) ^{1/2} \frac {1}{u} \int_{0} ^{\alpha} \left(1 - \frac {\xi}{\alpha}\right) d (\sin u \xi) = \left(\frac {2}{\pi}\right) ^{1/2} \int_{0} ^{\alpha} \frac {\sin u \xi}{u \alpha} \dd \xi = \left(\frac {2}{\pi}\right) ^{1/2} \frac {1 - \cos u \alpha}{u ^{2} \alpha},
\end{aligned}
\]
and so we have only to apply Plancherel's theorem in Chapter~\ref{chap:VI}, \ref{sec:VI.2}. Hence, by the Parseval relation of the Fourier transform, 
\begin{equation}\label{eq:XI.16.20}
y ^{(\alpha)} (t) \text{ satisfies } \int_{- \infty} ^{\infty} y ^{(\alpha)} (t) \exp (- i \cdot \xi \cdot t) \dd t = 0 \text{ if } | \xi | \geqq \alpha .
\end{equation}

Therefore, we may assume that $y \in L^{1}(-\infty, \infty)$ in \eqref{eq:XI.16.17} satisfies the condition 
\begin{equation}\label{eq:XI.16.21}
\int_{- \infty} ^{\infty} y (t) \exp (- i \cdot \xi \cdot t) \dd t = 0 \text{ whenever } | \xi | \geqq \alpha .
\end{equation}

Next we choose a positive number $\beta$ and a sufficiently large positive number $\gamma$ such that $[-\beta - \gamma, -\beta + \gamma]$ and $[\beta - \gamma, \beta + \gamma]$ both contain $[- \alpha, \alpha]$. Let $C_1(\xi)$ and $C_2(\xi)$ be the defining functions of the intervals $[- \gamma, \gamma]$ and $[- \beta, \beta]$, respectively. Then 
\begin{equation}\label{eq:XI.16.22}
\begin{yoslines}
&u (\xi) = (2 \beta) ^{-1} \int_{- \infty} ^{\infty} C _{1} (\xi - \eta) C _{2} (\eta) \dd \eta = 1 \text{ for } \xi \in [ - \alpha , \alpha ], \\
&\quad = 0 \text{ for sufficiently large } | \xi |, 0 \leqq u (\xi) \leqq 1 \text{ for all real } \xi .
\end{yoslines}
\end{equation}

By the Parseval relation of the Fourier transform, we have 
\begin{equation}\label{eq:XI.16.23}
\hat {\boldsymbol {u}} (t) = \frac {1}{2 \beta} (2 \pi) ^{1/2} \hat {C} _{1} (t) \hat {C} _{2} (t).
\end{equation}

Thus, by Plancherel's theorem, we see that $\hat{\pmb{u}}(t)$ belongs to $L^1(-\infty, \infty) \cap L^2(-\infty, \infty)$. Hence we may apply Theorem~\ref{thm:XI.16.Theorem.1} so that 
$f(t) = \hat{u}(t) = (2\pi)^{-1/2}(u,I_t)$ with a certain maximal ideal $I_{t}\neq I_{0}$ 


Moreover, by Plancherel's theorem, the inverse Fourier transform of $f(t)$ is equal to $u(\xi)$, i.e., we have 
\begin{equation}\label{eq:XI.16.24}
u (\xi) = (2 \pi) ^{-1/2} (f, I _{- \xi}).
\end{equation}

We next set $\tilde{g} = e - (2\pi)^{-1/2}f$. Then, by what was proved above, 
\begin{equation}\label{eq:XI.16.25}
\begin{aligned}
&0 \leqq (\tilde {g}, I _{\xi}) \leqq 1; \quad (\tilde {g}, I _{\xi}) = 0 \text{ whenever } \xi \in [ - \alpha , \alpha ]; \\
&\qquad (\tilde{g}, I_{\xi}) = 1 \text{ for all sufficiently large } |\xi|.
\end{aligned}
\end{equation}


On the other hand, we have, for $x^{*}(t) = \overline{x(-t)}$, the relation $(\overline{x,I_{\xi}}) = (x^{*},I_{\xi})$. Thus, by hypothesis, we have $(x^{*}\times x,I_{\xi}) = |(x,I_{\xi})|^2 > 0$ for all real $\xi$. Hence the element 
\[
\tilde {g} + x ^{*} \times x \in L ^{1}
\]
satisfies the condition that $(\tilde{g} + x^{*}\times x, I) > 0$ for all maximal ideals $I$ of $L^1$. Consequently, the inverse $(\tilde{g} + x^{*}\times x)^{-1}$ in $L^1$ does exist. 

We define 
\begin{equation}\label{eq:XI.16.26}
\tilde {z} = (\tilde {g} + x \times x ^{*}) ^{-1} \times x ^{*} \times y.
\end{equation}

Then, by \eqref{eq:XI.16.4}, the element $x \times \tilde{z}$ belongs to $L^1(-\infty, \infty)$. Moreover, we have for every real number $\xi$, 
\[
(x \times \tilde {z}, I _{\xi}) = (x, I _{\xi}) (\tilde {z}, I _{\xi}) = (x, I _{\xi}) \frac {(x ^{*} , I _{\xi}) (y , I _{\xi})}{(\tilde {g} , I _{\xi}) + (x , I _{\xi}) (x ^{*} , I _{\xi})}.
\]

Hence, by \eqref{eq:XI.16.25} and the hypothesis that $(y, I_{\xi}) = 0$ whenever $|\xi| \geqq \alpha$, we obtain 
\[
(x \times \tilde {z}, I _{\xi}) = (y, I _{\xi}) \text{ for all real } \xi .
\]

The normed ring $L^1$ is semi-simple by Theorem~\ref{thm:XI.16.Theorem.2}. Hence we must have $x \times \tilde{z} = y$. We have thus proved Theorem~\ref{thm:XI.16.Theorem.3}.
\end{proof}

\begin{corollaryu} Let $k_{1}(t)$ belong to $L^{1}(-\infty, \infty)$, and let its Fourier transform vanish for no real argument. Let $k_{2}(t)$ belong to $L^{1}(-\infty, \infty)$. Let $f(t)$ be Baire measurable and bounded over $(-\infty, \infty)$. Let, for a constant $C$, 
\begin{equation}\label{eq:XI.16.27}
\lim\limits _{t \rightarrow \infty} \int_{- \infty} ^{\infty} k _{1} (t - s) f (s) \dd s = C \int_{- \infty} ^{\infty} k _{1} (t) \dd t.
\end{equation}
\end{corollaryu}

Then 
\begin{equation}\label{eq:XI.16.28}
\lim\limits _{t \rightarrow \infty} \int_{- \infty} ^{\infty} k _{2} (t - s) f (s) \dd s = C \int_{- \infty} ^{\infty} k _{2} (t) \dd t.
\end{equation}

\begin{proof}
 We may plainly suppose that $C$ = $0$. And \eqref{eq:XI.16.28} holds for $k_{2}(t)$ of the form $k_{2}(t) = (x \times k_{1})(t)$ with $x(t) \in L^{1}(-\infty, \infty)$. It is easy to prove that \eqref{eq:XI.16.28} holds for $k_{2}(t)$ if $k_{2}(t) = \slim\limits_{n \to \infty} k^{(n)}(t)$ in $L^{1}(-\infty, \infty)$ with $k^{(n)}(t) \in L^{1}(-\infty, \infty)$ for which \eqref{eq:XI.16.28} holds $(n = 1, 2, \ldots)$. Hence, by Theorem~\ref{thm:XI.16.Theorem.3}, we see that \eqref{eq:XI.16.28} holds for every $k_{2}(t) \in L^{1}(-\infty, \infty)$. 
\end{proof}

\begin{remarku} N. \textsc{Wiener} (\cite{wiener339}, \cite{wiener340} and \cite{wiener341}) has applied the above Corollary to a unified treatment of classical results concerning the limit relations in series and integrals which include a new proof of the prime number theorem. Cf. also H. R. \textsc{Pitt} \cite{pitt262}. The above proof of Theorem~\ref{thm:XI.16.Theorem.3} is adapted from M. \textsc{Fukamiya} \cite{fukamiya72} and I. E. \textsc{Segal} \cite{segal301}. Cf. also M. A. \textsc{Naimark} \cite{naimark221} and C. E. \textsc{Rickart} \cite{rickart271} and the bibliographies cited in these books. For the sake of comprehension of the scope of the above Corollary, we shall reproduce Wiener's deduction of the special Tauberian theorem. The special Tauberian theorem, as formulated by J. E. \textsc{Littlewood} reads: 
\end{remarku}

\begin{theorem}\label{thm:XI.16.Theorem.4} Let $\sum\limits_{n=0}^{\infty} a_n x^n$ converge to $s(x)$ for $|x| < 1$, and let 
\begin{equation}\label{eq:XI.16.29}
\lim\limits _{x \rightarrow 1-0} s (x) = C.
\end{equation}
\end{theorem}

Let, moreover, 
\begin{equation}\label{eq:XI.16.30}
\sup\limits _{n \geqq 1} n | a _{n} | = K < \infty .
\end{equation}

Then 
\begin{equation}\label{eq:XI.16.31}
\sum\limits_{n=0} ^{\infty} a _{n} = C.
\end{equation}

\begin{proof}
 Put $f(x) = \sum\limits_{n=0}^{[x]} a_n$. Then we see, by 
\begin{align}
\left| f (x) - s \left(e ^{-1/x}\right) \right|& = \left| \sum\limits_{n=1} ^{[x]} a _{n} \left(1 - e ^{-n/x}\right) - \sum\limits_{[x]+1} ^{\infty} a _{n} e ^{-n/x} \right| \notag \\
&\leqq \sum\limits_{n=1} ^{[x]} \frac {K}{n} \frac {n}{x} + \sum\limits_{[x]+1} ^{\infty} \frac {K}{n} e ^{-n/x} \leqq 2 K + K \int_{[x]} ^{\infty} e ^{-u/x} u ^{-1} \dd u \notag \\
&\leqq 2 K + K \int_{1} ^{\infty} e ^{-u} u ^{-1} \dd u = \text{ constant }, \notag
\end{align}
that $f(x)$ is bounded. 

Hence, by partial integration, 
\[
s \left(e ^{-x}\right) = \sum\limits_{n=0} ^{\infty} a _{n} e ^{-nx} = \int_{-0} ^{\infty} e ^{-ux} \dd f (u) = \int_{0} ^{\infty} x e ^{-ux} f (u) \dd u.
\]

Thus we have 
\begin{equation}\tag{31'}\label{eq:XI.16.31'}
C = \lim\limits _{x \rightarrow 0} \int_{0} ^{\infty} x e ^{-ux} f (u) \dd u = \lim\limits _{\xi \rightarrow \infty} \int_{- \infty} ^{\infty} e ^{- \xi} e ^ {- e ^{\eta - \xi}} f (e ^{\eta}) e ^{\eta} \dd \eta .
\end{equation}

This may be written as 
\begin{equation}\label{eq:XI.16.32}
\lim\limits _{t \rightarrow \infty} \int_{- \infty} ^{\infty} k _{1} (t - s) f \left(e ^{s}\right) \dd s = C \int_{- \infty} ^{\infty} k _{1} (t) \dd t, \text{ where } k _{1} (t) = e ^{-t} e ^ {- e ^{-t}},
\end{equation}
since 
\[
\int_{- \infty} ^{\infty} k _{1} (t) \dd t = \int_{- \infty} ^{\infty} e ^{-t} e ^ {- e ^{-t}} \dd t = \int_{0} ^{\infty} e ^{-x} \dd x = 1.
\]

Furthermore, we have 
\[
\int_{- \infty} ^{\infty} k _{1} (t) e ^{-iut} \dd t = \int_{0} ^{\infty} x ^{iu} e ^{-x} \dd x = \Gamma (1 + i u) \neq 0.
\]

Thus we can apply the Corollary and obtain, for 
\[
k _{2} (t) = 0 (t < 0); k _{2} (t) = e ^{-t} (t > 0),
\]
the limit relation 
\[
\begin{aligned}
&C= C \int_{0} ^{\infty} e ^{-t} \dd t = C \int_{- \infty} ^{\infty} k _{2} (t) \dd t = \lim\limits _{t \rightarrow \infty} \int_{- \infty} ^{\infty} k _{2} (t - s) f (e ^{s}) \dd s = \lim\limits _{x \rightarrow \infty} x ^{-1} \int_{0} ^{x} f (y) \dd y.
\end{aligned}
\]

It follows that, if $\lambda > 0$, 
\begin{equation}\label{eq:XI.16.33}
\begin{aligned}
C &= \frac {(1 + \lambda) C - C}{\lambda} = \lim\limits _{x \rightarrow \infty} \frac {1}{\lambda x} \left\{\int_{0} ^{(1+ \lambda) x} f (y) \dd y - \int_{0} ^{x} f (y) \dd y \right\}
\\
&= \lim\limits _{x \rightarrow \infty} \frac {1}{\lambda x} \int_{x} ^{(1+ \lambda) x} f (y) \dd y = \lim\limits _{x \rightarrow \infty} \left\{f (x) + \frac {1}{\lambda x} \int_{x} ^{(1+ \lambda) x} [ f (y) - f (x) ] \dd y \right\}.
\end{aligned}
\end{equation}

On the other hand, we have, by \eqref{eq:XI.16.30}, 
\[
\begin{aligned}
&\left| \frac {1}{\lambda x} \int_{x} ^{(1+ \lambda) x} [ f (y) - f (x) ] \dd y \right|  \leqq \frac {1}{\lambda x} \int_{x} ^{(1+ \lambda) x} \sum\limits_{[x]+1} ^{[y]} \frac {K}{n} \dd y \leqq \sum\limits_{[x]+1} ^{[ (1+ \lambda) x ]} \frac {K}{[ x ]} \leqq \frac {[ \lambda x ] K}{[ x ]} \leqq 2 \lambda K
\end{aligned}
\]
for sufficiently large value of $x$. Hence, by \eqref{eq:XI.16.33}, 
\[
\varlimsup\limits_{x \rightarrow \infty} | f (x) - C | \leqq 2 \lambda K,
\]
and since $\lambda$ is any positive number, we get 
\[
\lim\limits _{x \rightarrow \infty} f (x) = C.
\]

Thus we have proved \eqref{eq:XI.16.31}. 
\end{proof}
   % XI
% ==== XII ====
\yoschapter{XII}{Other Representation Theorems in Linear Spaces}
\label{chap:XII}

In this chapter, we shall prove three representation theorems in linear spaces. The first one, the Krein-Milman theorem says that a non-void convex compact subset $K$ of a locally convex linear topological space is equal to the closure of the convex hull of the extremal points of $K$. The other two theorems concern the representations of a vector lattice as point functions and as set functions. 


\yossection{1}{Extremal Points. The Krein-Milman Theorem}
\label{sec:XII.1}

\begin{definitionu} Let $K$ be a subset of a real- or complex-linear space $X$. A non-void subset $M \subseteq K$ is said to be an extremal subset of $K$, if a proper convex combination $\alpha k_{1} + (1 - \alpha) k_{2}, 0 < \alpha < 1$, of two points $k_{1}$ and $k_{2}$ of $K$ lies in $M$ only if both $k_{1}$ and $k_{2}$ are in $M$. An extremal set of $K$ consisting of just one point is called an extremal point of $K$. 
\end{definitionu}

\begin{examplen} In a three dimensional Euclidean space, the surface of a solid sphere is an extremal subset of the sphere, and every point of the surface is an extremal point of the surface.

\end{examplen}
\begin{theoremn}[\textsc{Krein}-\textsc{Milman}] A non-void compact convex subset $K$ of a locally convex linear topological space $X$ has at least one extremal point.

\end{theoremn}
\begin{proof}
 The set $K$ is itself an extremal set of $K$. Let $M$ be the totality of compact extremal subsets $M$ of $K$. Order $M$ by inclusion relation. It is easy to see that, if $M_{1}$ is a linearly ordered subfamily of $M$, the non-void set $\bigcap\limits_{M\in M_{1}}M$ is a compact extremal subset of $K$ which is a lower bound for the subfamily $M_{1}$. 

Thus, by Zorn's lemma, $\mathfrak{M}$ contains a minimal element $M_0$. Suppose that $M_0$ contains two distinct points $x_0$ and $y_0$. Then there exists a continuous linear functional $f$ on $X$ such that $f(x_0) \neq f(y_0)$. We may assume that $\operatorname{Re} f(x_0) \neq \operatorname{Re} f(y_0)$. $M_0$ being compact, the set $M_1 = \{x \in M_0; \operatorname{Re} f(x) = \inf\limits_{y \in M_0} \operatorname{Re} f(y)\}$ is a proper subset of $M_0$. On the other hand, if $k_1$ and $k_2$ are points of $K$ such that $\alpha k_1 + (1 - \alpha) k_2 \in M_1$ for some $\alpha$ with $0 < \alpha < 1$, then, by the extremal property of $M_{0}$, $k_{1}$ and $k_{2}$ both $\in M_{0}$. It follows from the definition of $M_{1}$ that $k_{1}$ and $k_{2}$ both $\in M_{1}$. Hence $M_{1}$ is a closed extremal subset properly contained in $M_{0}$. Since $M_{0}$ is a minimal element of $M$, we have arrived at a contradiction. Therefore $M_{0}$ must consist of only one point which is thus an extremal point of $K$. 
\end{proof}
\begin{corollaryu} Let $K$ be a non-void compact convex subset of a locally convex real linear topological space $X$. Let $E$ be the totality of the extremal points of $K$. Then $K$ coincides with the smallest closed set containing every convex combination $\sum\limits_{i} \alpha_{i} e_{i} (\alpha_{i} \geqq 0, \sum\limits_{i} \alpha_{i} = 1)$ of points $e_{i} \in E$, i.e., $K$ is equal to the closure of the convex hull Conv $(E)$ of $E$. 
\end{corollaryu}

\begin{proof}
 The inclusion $E \subseteq K$ and the convexity of $K$ imply that $\Conv(E)^a$ is contained in the compact set $K$. Suppose that there exists a point $k_0$ contained in $(K - \Conv(E)^a)$. We can then take a point $c \in \Conv(E)^a$ so that $(k_0 - c) \bar{\in} (\Conv(E)^a - c)$. The set $(\Conv(E)^a - c)$ being compact convex and $\ni 0$, there exists, by Theorem~\ref{thm:IV.6.Theorem.3prime} in Chapter~\ref{chap:IV}, \ref{sec:IV.6}, a continuous real linear functional $f$ on $X$ such that 
\[
f (k _{0} - c) > 1 \text{ and } f (k - c) \leqq 1 \text{ for } (k - c) \in (\Conv (E) ^{a} - c). \]

Let $K_{1} = \{x \in K; f(x) = \sup\limits_{y \in K} f(y)\}$. Then, since $k_{0} \in K$, the set $K_{1} \cap E$ must be void. Moreover, since $K$ is compact, $K_{1}$ is a closed extremal subset of $K$. On the other hand, any extremal subset of $K_{1}$ is also an extremal subset of $K$ and hence any extremal point of $K_{1}$, which surely exists by the preceding Theorem, is also an extremal point of $K$. Since $K_{1} \cap E$ is void, we have arrived at a contradiction. 
\end{proof}
\begin{remarku} The above Theorem and the Corollary were first proved by M. \textsc{Krein}-D. \textsc{Milman} \cite{krein165}. The proof given above is adapted from J. L. \textsc{Kelley} \cite{kelley147}. It is to be noted that for the unit sphere $S = \{x \in X; \| x \| \leqq 1\}$ in a Hilbert space $X$ the extremal points of $S$ are precisely those on the surface of $S$, i.e., those of norm $1$. This we easily see from \eqref{eq:I.5.1} in Chapter~\ref{chap:I}, \ref{sec:I.5}. For applications of the notion of the extremal points to concrete function spaces, see, e.g., K. \textsc{Hoffman} \cite{hirsch108}. 
\end{remarku}

A simple example. Let $C[0,1]$ be the space of real-valued continuous functions $x(t)$ defined on $[0,1]$ normed by $x = \max\limits_t |x(t)|$. The dual space $X = C[0,1]'$ is the space of real Baire measures on $[0,1]$ of bounded total variations. The unit sphere $K$ of $X$ is compact in the weak* topology of $X$ (see Theorem~\hyperref[thm:AV.1.Theorem.1]{$1$} in the Appendix to Chapter~\ref{chap:IV}). It is easy to see that extremal points of $K$ are in one-to-one correspondence with the linear functionals $f_{t_0} \in X$ of the form $\langle x, f_{t_0} \rangle = x(t_0), t_0 \in [0,1]$. The above Corollary says that any linear functional $f \in X$ is given as the weak* limit of the functionals of the form 
\[
\sum\limits_{j=1} ^{n} \alpha_{j} x (t _{j}), \quad \text{ where } \quad \alpha_{j} > 0, \sum\limits_{j=1} ^{n} \alpha_{j} = 1 \quad \text{ and } \quad t _{j} \in [ 0, 1 ].
\]

Recently, G. \textsc{Choquet} \cite{choquet33} proved a more precise result: If $X$ is a metric space, then $E$ is a $G_{\delta}$ set and, for every $x \in K$, there exists a non-negative Baire measure $\mu_x(B)$ defined for Baire sets $B$ of $X$ such that $\mu_x(X - E) = 0$, $\mu_x(E) = 1$ and $x = \int_E y \mu_x(\dd y)$. As for the uniqueness of $\mu_x$ and further literature, see G. \textsc{Choquet} and P. A. \textsc{Meyer} \cite{choquet34}.\yosstarfootnote{See also Supplementary Notes, p.~\pageref{notes:XII}.} 


\yossection{2}{Vector Lattices}
\label{sec:XII.2}

The notion of ``positivity'' in concrete function spaces is very important, in theory as well as in application. A systematic abstract treatment of the ``positivity'' in linear spaces was introduced by F. \textsc{Riesz} \cite{rickart277}, and further developed by H. \textsc{Freudenthal} \cite{freudenthal67}, G. \textsc{Birkhoff} \cite{birkhoff17} and many other authors. These results are called the theory of vector lattice. We shall begin with the definition of the vector lattice. 

\begin{definition}\label{def:XII.2.Definition.1} $A$ real linear space $X$ is said to be a vector lattice if $X$ is a lattice by a partial order relation $x \leqq y$ satisfying the conditions: 
\begin{equation}\label{eq:XII.2.1}
x \leqq y \text{ implies } x + z \leqq y + z,
\end{equation}
\begin{equation}\label{eq:XII.2.2}
x \leqq y \text{ implies } \alpha x \leqq \alpha y \text{ (or } \alpha x \geqq \alpha y\text{) for every } \alpha \geqq 0 \text{ (or } \alpha \leqq 0\text{)}.
\end{equation}
\end{definition}


\begin{proposition}\label{prop:XII.2.Proposition.1} If, in a vector lattice $X$, we define 
\begin{equation}\label{eq:XII.2.3}
x ^{+} = x \vee 0 \text{ and } x ^{-} = x \wedge 0,
\end{equation}
then we have 
\begin{equation}\label{eq:XII.2.4}
x \vee y = (x - y) ^{+} + y, x \wedge y = - ((- x) \vee (- y)).
\end{equation}
\end{proposition}

\begin{proof}
 The one-one mappings $x \rightarrow x + z$ and $x \rightarrow \alpha x (\alpha > 0)$ of $X$ onto $X$ both preserve the partial order in $X$. 
\end{proof}

\begin{examplen} The totality $A(S, \mathfrak{B})$ of real-valued, $\sigma$-additive set functions $x(B)$ defined and finite on a $\sigma$-additive family $(S, \mathfrak{B})$ of sets $B \subseteq S$ is a vector lattice by
\[
(x + y) (B) = x (B) + y (B), (\alpha x) (B) = \alpha x (B)
\]
and the partial order $x \leqq y$ defined by $x(B) \leqq y(B)$ on $\mathfrak{B}$. In fact, we have, in this case, 
\begin{equation}\label{eq:XII.2.5}
x ^{+} (B) = \sup\limits _{N \subseteq B} x (N) = \text{ the positive variation } \overline {{V}} (x; B) \text{ of } x \text{ on } B.
\end{equation}

\end{examplen}
\begin{proof}
 We have to show that $\overline{V}(x; B) = (x \vee 0)(B)$. It is clear that $\overline{V}(x; B) \geqq 0$ and $x(B) \leqq \overline{V}(x; B)$ on $\mathfrak{B}$. If $0 \leqq y(B)$ and $x(B) \leqq y(B)$ on $\mathfrak{B}$, then, for any $N \subseteq B$, $y(B) = y(N) + y(B - N) \geqq x(N)$ and so $y(B) \geqq \overline{V}(x; B)$ on $\mathfrak{B}$. 
\end{proof}

\begin{proposition}\label{prop:XII.2.Proposition.2} In a vector lattice $X$, we have 
\begin{align}
x \underset {\wedge} {\vee} y + z& = (x + z) \underset {\wedge} {\vee} (y + z), \label{eq:XII.2.6} \\
\alpha (x \underset {\wedge} {\vee} y)& = (\alpha x) \underset {\wedge} {\vee} (\alpha y) \quad \text{ for } \quad \alpha > 0, \label{eq:XII.2.7} \\
\alpha (x \underset {\wedge} {\vee} y)& = (\alpha x) \underset {\vee} {\wedge} (\alpha y) \quad \text{ for } \quad \alpha < 0, \label{eq:XII.2.8} \\
x \wedge y& = - (- x) \vee (- y), x ^{-} = - (- x) ^{+}, x ^{+} = - (- x) ^{-}. \label{eq:XII.2.9}
\end{align}
\end{proposition}

\begin{proof}
 Clear from \eqref{eq:XII.2.1} and \eqref{eq:XII.2.2}. 
\end{proof}

\begin{corollaryu} 
\begin{equation}\label{eq:XII.2.10}
x+y=x\vee y+x\wedge y,\qquad\text{in particular, }x=x^++x^-.
\end{equation}
\end{corollaryu}

\begin{proof}
\[
x\vee y-x-y=0\vee(y-x)-y=(-y)\vee(-x)=-(y\wedge x).
\]
\end{proof}



\begin{proposition}\label{prop:XII.2.Proposition.3} We have 
\begin{align}
x \underset {\wedge} {\vee} y& = y \underset {\wedge} {\vee} x (\text{ commutativity }), \label{eq:XII.2.11} \\
x \vee (y \vee z)& = (x \vee y) \vee z = \sup (x, y, z) \Big \} _ {\text{ (associativity) }} \label{eq:XII.2.12} \\
x \wedge (y \wedge z)& = (x \wedge y) \wedge z = \inf (x, y, z) \left\} ^ {(\text{ associativity })}, \right. \label{eq:XII.2.13} \\
&\quad (x \wedge y) \vee z = (x \vee z) \wedge (y \vee z) \quad (\text{ distributivity }), \label{eq:XII.2.14} \\
&\quad (x \vee y) \wedge z = (x \wedge z) \vee (y \wedge z) \quad (\text{ distributivity }). \label{eq:XII.2.15}
\end{align}
\end{proposition}


\begin{proof}
 We have only to prove the distributivity. It is clear that $(x \wedge y) \vee z \leqq x \vee z, y \vee z$. Let $w \leqq x \vee z, y \vee z$. Then $w \leqq x \vee z = x + z - x \wedge z$ and so $x + z \geqq w + x \wedge z$. Similarly we have $y + z \geqq w + y \wedge z$. Hence 
\[
\begin{aligned}
&w + (x \wedge z) \wedge (y \wedge z) = (w + x \wedge z) \wedge (w + y \wedge z) \leqq (x + z) \wedge (y + z) = x \wedge y + z,
\end{aligned}
\]
and so 
\[
w \leqq (x \land y) + z - (x \land y) \land z = (x \land y) \lor z. \]

We have thus proved \eqref{eq:XII.2.14}. \eqref{eq:XII.2.15} is proved by substituting $-x$, $-y$, $-z$ for $x$, $y$, $z$, respectively in \eqref{eq:XII.2.14}. 
\end{proof}

\begin{remarkn}[Remark $1$] In a lattice, the distributive identity 
\begin{equation}\label{eq:XII.2.16}
x \wedge (y \vee z) = (x \wedge y) \vee (x \wedge z)
\end{equation}
\end{remarkn}

does not hold in general. The modular identity: 
\begin{equation}\label{eq:XII.2.17}
x \leqq z \text{ implies } x \vee (y \wedge z) = (x \vee y) \wedge z
\end{equation}
is weaker than the distributive identity \eqref{eq:XII.2.15}. Let $G$ be a group. Then the totality of invariant subgroups $N$ of $G$ constitutes a modular lattice, that is, a lattice satisfying the modular identity \eqref{eq:XII.2.17}, if we define $N_{1} \vee N_{2}$ and $N_{1} \wedge N_{2}$ as the invariant subgroup generated by $N_{1}$ and $N_{2}$ and the invariant subgroup $N_{1} \cap N_{2}$, respectively. 

\begin{remarkn}[Remark $2$] $A$ typical example of a distributive lattice is given by the Boolean algebra; a distributive lattice $B$ is called a Boolean algebra if it satisfies the conditions: \begin{enumerate}[(i)]
\item there exist elements $I$ and $0$ such that $0 \leqq x \leqq I$ for every $x \in B$,
\item for any $x \in B$, there exists a uniquely determined complement $x' \in B$ which satisfies $x \vee x' = I$, $x \wedge x' = 0$.
\end{enumerate}
 The totality $B$ of subsets of a fixed set is a Boolean algebra by defining the partial order in $B$ by the inclusion relation. 
\end{remarkn}

\begin{proposition}\label{prop:XII.2.Proposition.4} We define, in a vector lattice $X$, the absolute value 
\begin{equation}\label{eq:XII.2.18}
| x | = x \vee (- x).
\end{equation}
\end{proposition}

Then 
\begin{align}
| x |& \geqq 0, \text{ and } | x | = 0 \text{ iff } x = 0, \label{eq:XII.2.19} \\
| x + y |& \leqq | x | + | y |, | \alpha x | = | \alpha | | x |. \label{eq:XII.2.20}
\end{align}

\begin{proof}
 We have 
\begin{equation}\label{eq:XII.2.21}
x ^{+} \wedge (- x ^{-}) = x ^{+} \wedge (- x) ^{+} = 0. \end{equation}

For, by $0 = x - x = x \vee (-x) + x \wedge (-x) \geqq 2(x \wedge (-x))$ and the distributive identity \eqref{eq:XII.2.14}, $0 = (x \wedge (-x)) \vee 0 = x^{+} \wedge (-x)^{+}$. Thus 
\[
x ^{+} - x ^{-} = x ^{+} + (- x) ^{+} = x ^{+} \vee (- x) ^{+}.
\]

On the other hand, from $x \vee (-x) \geqq x \wedge (-x) \geqq -((-x) \vee x)$, we have $x \vee (-x) \geqq 0$ and so by \eqref{eq:XII.2.21} 
\[
x \vee (- x) = (x \vee (- x)) \vee 0 = x ^{+} \vee (- x) ^{+} = x ^{+} + (- x) ^{+}.
\]

We have thus proved 
\begin{equation}\label{eq:XII.2.22}
x ^{+} - x ^{-} = x ^{+} + (- x) ^{+} = x ^{+} \vee (- x) ^{+} = x \vee (- x).
\end{equation}

We now prove \eqref{eq:XII.2.19} and \eqref{eq:XII.2.20}. If $x = x^{+} + x^{-} \neq 0$, then $x^{+}$ or $(-x^{-})$ is $> 0$ so that $|x| = x^{+} \vee (-x^{-}) > 0$. $|\alpha x| = (\alpha x) \vee (-\alpha x) = |\alpha | (x \vee (-x)) = |\alpha | |x|$. From $|x| + |y| \geqq x + y, -x - y$, we obtain $|x| + |y| \geqq (x + y) \vee (-x - y) = |x + y|$. 
\end{proof}

\begin{remarku} The decomposition $x = x^{+} + x^{-}$, $x^{+} \wedge (-x^{-}) = 0$, is called the Jordan decomposition of $x; x^{+}, x^{-}$ and $|x|$ correspond to the positive variation, the negative variation and the total variation, of a function $x(t)$ of bounded variation, respectively. 
\end{remarku}

\begin{proposition}\label{prop:XII.2.Proposition.5} For any $y \in X$, we have 
\begin{equation}\label{eq:XII.2.23}
\left| x - x _{1} \right| = \left| x \vee y - x _{1} \vee y \right| + \left| x \wedge y - x _{1} \wedge y \right|.
\end{equation}
\end{proposition}

\begin{proof}
 We have 
\[ \begin{aligned}
&| a - b |  = (a - b) ^{+} - (a - b) ^{-} = a \vee b - b - (a \wedge b - b) = a \vee b - a \wedge b.
\end{aligned}
\]

Hence the right side of \eqref{eq:XII.2.23} is, by \eqref{eq:XII.2.10}, \eqref{eq:XII.2.14} and \eqref{eq:XII.2.15}, 
\[
\begin{aligned}
&= (x \lor y) \lor x _{1} - (x \lor y) \land (x _{1} \lor y) + (x \land y) \lor (x _{1} \land y) - (x \land y) \land x _{1}
\\
&= (x \lor x _{1}) \lor y - (x \land x _{1}) \lor y + (x \lor x _{1}) \land y - (x \land x _{1}) \land y
\\
&= x \lor x _{1} + y - (x \land x _{1} + y) = x \lor x _{1} - x \land x _{1} = | x - x _{1} |.
\end{aligned}
\]
\end{proof}

\begin{definition}\label{def:XII.2.Definition.2} $A$ sequence $\{x_{n}\}$ of a vector lattice $X$ is said to $O$-converge to an element $x \in X$, in symbol, $\Olim_{n \to \infty} x_{n} = x$ if there exists a sequence $\{w_{n}\}$ such that $|x - x_{n}| \leqq w_{n}$ and $w_{n} \downarrow 0$. Here $w_{n} \downarrow 0$ means that $w_{1} \geqq w_{2} \geqq \cdots$ and $\bigwedge\limits_{n \geqq 1} w_{n} = 0$. The \Olim $x_{n}$, if it exists, is uniquely determined. For, let \Olim $x_{n} = x$ and \Olim $x_{n} = y$. Then $|x - x_{n}| \leqq w_{n}$, $w_{n} \downarrow 0$ and $|y - x_{n}| \leqq u_{n}, u_{n} \downarrow 0$. Thus $|x - y| \leqq |x - x_{n}| + |x_{n} - y| \leqq w_{n} + u_{n}$ and $(w_{n} + u_{n}) \downarrow 0$ as may be seen from $w_{n} + \bigwedge\limits_{n \geqq 1} u_{n} \geqq \bigwedge\limits_{n \geqq 1} (w_{n} + u_{n})$. This proves that $x = y$. 
\end{definition}

\begin{proposition}\label{prop:XII.2.Proposition.6} With respect to the notion of \Olim, the operations $x + y, x \vee y$ and $x \wedge y$ are continuous in $x$ and $y$. 
\end{proposition}

\begin{proof}
 Let $\Olim_{n \to \infty} x_n = x$, $\Olim_{n \to \infty} y_n = y$. Then $|x + y - x_n - y_n| \leqq |x - x_n| + |y - y_n|$ implies $\Olim_{n \to \infty} (x_n + y_n) = x + y$. By \eqref{eq:XII.2.23}, we have 
\[
\begin{aligned}
&| x \underset {\wedge} {\vee} y - x _{n} \underset {\wedge} {\vee} y _{n} |  \leqq | x \underset {\wedge} {\vee} y - x _{n} \underset {\wedge} {\vee} y | + | x _{n} \underset {\wedge} {\vee} y - x _{n} \underset {\wedge} {\vee} y _{n} | \leqq | x - x _{n} | + | y - y _{n} |,
\end{aligned}
\]
and so \Olim $(x_{n} \wedge^{V} y_{n}) = (\Olim_{n \to \infty} x_{n}) \wedge^{V} (\Olim_{n \to \infty} y_{n})$. 
\end{proof}

\begin{remarku} We have 
\[
\Olim _{n \to \infty} \alpha \cdot x _{n} = \alpha \cdot \Olim _{n \to \infty} x _{n}.
\]
\end{remarku}

But, in general, 
\[
\Olim _{n \to \infty} \alpha_{n} x \neq \left(\lim\limits _{n \to \infty} \alpha_{n}\right) x.
\]

The former relation is clear from $\left|\alpha x - \alpha x_{n}\right| = \left|\alpha\right| \left|x - x_{n}\right|$. The latter inequality is proved by the following counter example: We introduce in the two-dimensional vector space the lexicographic partial order in which $(\xi_{1}, \eta_{1}) \geqq (\xi_{2}, \eta_{2})$ means that either $\xi_{1} > \xi_{2}$ or $\xi_{1} = \xi_{2}$, $\eta_{1} \geqq \eta_{2}$. It is easy to see that we obtain a vector lattice. We have, in this lattice, 
\[
n ^{-1} (1, 0) \geqq (0, 1) > 0 = (0, 0) (n = 1, 2, \dots).
\]

Hence \Olim $n^{-1}(1,0) \neq 0$. A necessary and sufficient condition for the validity of the equation 
\begin{equation}\label{eq:XII.2.24}
\Olim _ {\alpha_{n} \rightarrow \alpha} \alpha_{n} x = \alpha x
\end{equation}
is the so-called Archimedean axiom 
\begin{equation}\label{eq:XII.2.25}
\Olim _{n \to \infty} n ^{-1} x = 0 \quad \text{ for every } \quad x \geqq 0.
\end{equation}

This we see by the Jordan decomposition $y = y^{+} + y^{-}$. 

\begin{definition}\label{def:XII.2.Definition.3} $A$ subset $\{x_{\alpha}\}$ of a vector lattice $X$ is said to be bounded if there exist $y$ and $z$ such that $y \leqq x_{\alpha} \leqq z$ for all $x_{\alpha}$. $X$ is said to be complete if, for any bounded set $\{x_{\alpha}\}$ of $X$, $\sup\limits_{\alpha} x_{\alpha}$ and $\inf\limits_{\alpha} x_{\alpha}$ exist in $X$. Here $\sup\limits_{\alpha} x_{\alpha}$ is the least upper bound in the sense of the partial order in $X$, and $\inf\limits_{\alpha} x_{\alpha}$ is the largest lower bound in the sense of the partial order in $X$. A vector lattice $X$ is said to be $\sigma$-complete, if, for any bounded sequence $\{x_{n}\}$ of $X$, $\sup\limits_{n\geqq 1}x_n$ and $\inf\limits_{n\geqq 1}x_n$ exist in $X$. We define, in a $\sigma$-complete vector lattice $X$, 
\begin{equation}\label{eq:XII.2.26}
O _{n \to \infty} ^{- \lim} x _{n} = \inf\limits _{m} \left(\sup\limits _{n \geqq m} x _{n}\right), O _{n \to \infty} ^{- \lim} x _{n} = \sup\limits _{m} \left(\inf\limits _{n \geqq m} x _{n}\right).
\end{equation}
\end{definition}

\begin{proposition}\label{prop:XII.2.Proposition.7} \Olim $x_{n} = x$ iff \Olim $x_{n} = O$-lim $x_{n} = x$. 
\end{proposition}

\begin{proof}
 Suppose $|x - x_n| \leqq w_n, w_n \downarrow 0$. Then $x - w_n \leqq x_n \leqq x + w_n$ and so we obtain $O\text{- }\varlimsup\limits_{n \to \infty} (x - w_n) = x \leqq O\text{- }\varliminf\limits_{n \to \infty} x_n \leqq O\text{- }\varliminf\limits_{n \to \infty} (x + w_n) = x$, that is, $O\text{- }\varliminf\limits_{n \to \infty} x_n = x$. Similarly, we obtain $O\text{- }\varliminf\limits_{n \to \infty} x_n = x$. 

We next prove the sufficiency. Put $u_{n} = \sup\limits_{m\geqq n}x_{m}, v_{n} = \inf\limits_{m\geqq n}x_{m}, u_{n} - v_{n} = w_{n}$. Then, by hypothesis, $w_{n}\downarrow 0$. Also, by $x_{n}\leqq u_{n} = x + (u_{n} - x)\leqq x + (u_{n} - v_{n}) = x + w_{n}$ and $x_{n}\geqq x - w_{n}$ obtained similarly, we prove $|x - x_n|\leqq w_n$. Thus $\Olim_{n\to \infty}x_n = x$. 
\end{proof}

\begin{proposition}\label{prop:XII.2.Proposition.8} In a $\sigma$-complete vector lattice $X$, $\alpha x$ is continuous in $\alpha$, $x$ with respect to \Olim. 
\end{proposition}

\begin{proof}
 We have $|\alpha x - \alpha_n x_n| \leqq |\alpha x - \alpha x_n| + |\alpha x_n - \alpha_n x_n| = |\alpha| |x - x_n| + |\alpha - \alpha_n| |x_n|$. Hence, if $\Olim_{n \to \infty} x_n = x$, $\lim\limits_{n \to \infty} \alpha_n = \alpha$, then the $\Olim$ of the first term on the right is $0$. Therefore, by putting $\sup\limits_{n \geqq 1} |x_n| = y$, $\sup\limits_{m \geqq n} |\alpha - \alpha_m| = \beta_n$, we have to prove $\Olim_{n \to \infty} \beta_n y = 0$. But, by $y \geqq 0$ and $\beta_n \downarrow 0$, we see that $\Olim_{n \to \infty} \beta_n y = z$ exists and $\Olim_{n \to \infty} 2^{-1} \beta_n y = 2^{-1} z$. Since there exists, for any $n$, an $n_0$ such that $\beta_{n_0} \leqq 2^{-1} \beta_n$, we must have $z = 2^{-1} z$, that is, $z = 0$. 
\end{proof}

\begin{proposition}\label{prop:XII.2.Proposition.9} In a $\sigma$-complete vector lattice $X$, \Olim $x_{n}$ exists iff 
\begin{equation}\label{eq:XII.2.27}
\underset {n, m \rightarrow \infty} {O - \lim} | x _{n} - x _{m} | = 0.
\end{equation}
\end{proposition}

\begin{proof}
 The necessity is clear from $|x_{n} - x_{m}| \leqq |x_{n} - x| + |x - x_{m}|$. If we set $|x_{n} - x_{m}| = y_{nm}$, then $O\text{- }\varlimsup\limits_{n \to \infty} x_{n} \leqq x_{m} + O\text{- }\varliminf\limits_{n \to \infty} y_{nm}, O\text{- }\varliminf\limits_{n \to \infty} x_{n} \geqq x_{m} - O\text{- }\varliminf\limits_{n \to \infty} y_{nm}$. Hence 
\[
0 \leqq O - \varlimsup\limits_{n \rightarrow \infty} x _{n} - O - \varliminf\limits_{n \rightarrow \infty} x _{n} \leqq O - \lim\limits _{m \rightarrow \infty} \left(O - \varlimsup\limits_{n \rightarrow \infty} y _{nm} - O - \varliminf\limits_{n \rightarrow \infty} y _{nm}\right) = 0. \]

Thus we have proved the sufficiency. 
\end{proof}

\begin{proposition}\label{prop:XII.2.Proposition.10} $A$ vector lattice $X$ is $\sigma$-complete iff every monotone increasing, bounded sequence $\{x_{n}\} \subseteq X$ has $\sup x_{n}$ in $X$. 
\end{proposition}

\begin{proof}
 We have only to prove the sufficiency. Let $\{z_n\}$ be any bounded sequence in $X$, and set $x_{n} = \sup\limits_{m\leqq n}z_{m}$. Then, by hypothesis, $\sup\limits_{n\geqq 1}x_{n} = z$ exists in $X$, and $z = \sup\limits_{n\geqq 1}z_{n}$. Similarly, we see that $\inf\limits_{n\geqq 1}z_n = \inf\limits_{n\geqq 1}\left(\inf\limits_{m\leqq n}z_m\right)$ also exists in $\dot{X}$. 
\end{proof}


\yossection{3}{\texorpdfstring{$B$}{B}-lattices and \texorpdfstring{$F$}{F}-lattices}
\label{sec:XII.3}

\begin{definitionu} $A$ real $B$-space ( $F$-space) $X$ is said to be a $B$-lattice ( $F$-lattice) if it is a vector lattice such that 
\begin{equation}\label{eq:XII.3.1}
| x | \leqq | y | \quad \text{ implies } \quad \| x \| \leqq \| y \|.
\end{equation}
\end{definitionu}

\begin{examplen} $C(S), L^p(S)$ are $B$-lattices by the natural partial order $x \leqq y$ which means $x(s) \leqq y(s)$ on $S(x(s) \leqq y(s))$ a.e. on $S$ for the case $L^p(S)$). $M(S, \mathfrak{B}, m)$ with $m(S) < \infty$ is an $F$-lattice by the natural partial order as in the case of $L^p(S)$. $A(S, \mathfrak{B})$ is a $B$-lattice by
$\left\| x \right\| = | x | (S) =$ the total variation of $x$ over $S.$


We have, by \eqref{eq:XII.3.1} and $|x| = |(|x|)|$, 
\begin{equation}\label{eq:XII.3.2}
\left\| x \right\| = \left\| (| x |) \right\|.
\end{equation}

In $A(S, \mathfrak{B})$, we have, moreover, 
\begin{equation}\label{eq:XII.3.3}
x \geqq 0, y \geqq 0 \quad \text{ imply } \quad \| x + y \| = \| x \| + \| y \|.
\end{equation}

S. \textsc{Kakutani} called a $B$-lattice satisfying \eqref{eq:XII.3.3} an abstract $L^1$-space. From \eqref{eq:XII.3.3} we have 
\begin{equation}\label{eq:XII.3.4}
| x | < | y | \quad \text{ implies } \quad \| x \| < \| y \|.
\end{equation}

The norm in $A(S, \mathfrak{B})$ is continuous in \Olim, that is 
\begin{equation}\label{eq:XII.3.5}
\Olim _{n \to \infty} x _{n} = x \quad \text{ implies } \quad \lim\limits _{n \to \infty} \| x _{n} \| = \| x \|.
\end{equation}

For, in $A(S, \mathfrak{B})$, $\Olim_{n \to \infty} x_n = x$ is equivalent to the existence of $y_n \in A(S, \mathfrak{B})$ such that $|x - x_n| (S) \leqq y_n(S)$ with $y_n(S) \downarrow 0$. It is easy to see that $M(S, \mathfrak{B}, m)$ satisfies \eqref{eq:XII.3.4} and \eqref{eq:XII.3.5}. 

\end{examplen}
\begin{proposition}\label{prop:XII.3.Proposition.1} $A$ $\sigma$-complete $F$-lattice $X$ satisfying \eqref{eq:XII.3.4} and \eqref{eq:XII.3.5} is a complete lattice. In particular, $A(S, \mathfrak{B})$ and $L^p(S)$ are complete lattices. 
\end{proposition}

\begin{proof}
 Let $\{x_{\alpha}\} \subseteq X$ be bounded. We may assume $0 \leqq x_{\alpha} \leqq x$ for all $\alpha$, and we shall show that $\sup\limits_{\alpha} x_{\alpha}$ exists. Consider the totality $\{z_{\beta}\}$ of 

$z_{\beta}$ obtained as the sup of a finite number of $x_{\alpha}$'s: $z_{\beta} = \bigvee\limits_{j=1}^{n} x_{\alpha_j}$. Set $\gamma = \sup\limits_{\beta} \|z_{\beta}\|$. Then there exists a sequence $\{z_{\beta_j}\}$ such that $\lim\limits_{j \to \infty} \|z_{\beta_j}\| = \gamma$. If we put $z_n = \sup\limits_{j \leqq n} z_{\beta_j}$, then $O - \lim\limits_{n \to \infty} z_n = w$ exists, and, by \eqref{eq:XII.3.5} and the definition of $\gamma$, we have $\|w\| = \gamma$. We shall prove that $w = \sup\limits_{\alpha} x_{\alpha}$. Suppose that $x_{\alpha} \vee w > w$ for a certain $x_{\alpha}$. Then, by \eqref{eq:XII.3.4}, $\|x_{\alpha} \vee w\| > \|w\| = \gamma$. But, by $x_{\alpha} \vee w = O - \lim\limits_{n \to \infty} (x_{\alpha} \vee z_n)$, $x_{\alpha} \vee z_n \in \{z_{\beta}\}$ and \eqref{eq:XII.3.5}, we have $\|x_{\alpha} \vee w\| = \lim\limits_{n \to \infty} \|x_{\alpha} \vee z_n\| \leqq \gamma$, which is a contradiction. Hence we must have $w \geqq x_{\alpha}$ for all $x_{\alpha}$. Let $x_{\alpha} \leqq u$ for all $x_{\alpha}$. Suppose $w \wedge u < w$.


Then, by \eqref{eq:XII.3.4}, $\|w \wedge u\| < \gamma$, contrary to the fact that $w \wedge u \geqq z_{\beta}$ for all $z_{\beta}$. Hence we must have $w = \sup x_{\alpha}$. 
\end{proof}

\begin{remarku} In $C(S)$, $\Olim_{n \to \infty} x_n = x$ does not necessarily imply $\slim\limits_{n \to \infty} x_n = x$. In $M(S, \mathfrak{B}, m)$, $\slim\limits_{n \to \infty} x_n = x$ does not necessarily imply $\Olim_{n \to \infty} x_n = x$. To see this, let $x_1(s), x_2(s), \ldots$ be the defining functions of the intervals of [$0$, $1$]: 
\[
\left[ 0, \frac {1}{2} \right], \left[ \frac {1}{2}, \frac {2}{2} \right], \left[ 0, \frac {1}{4} \right], \left[ \frac {1}{4}, \frac {2}{4} \right], \left[ \frac {2}{4}, \frac {3}{4} \right], \left[ \frac {3}{4}, \frac {4}{4} \right], \left[ 0, \frac {1}{8} \right], \left[ \frac {1}{8}, \frac {2}{8} \right], \dots .
\]
\end{remarku}

Then the sequence $\{x_{n}(s)\} \subseteq M([0,1])$ converges to $0$ asymptotically but does not converge to $0$ a.e., that is, we have $s\text{- }\lim\limits_{n\to \infty}x_n = 0$ but $O\text{- }\lim\limits_{n\to \infty}x_n = 0$ does not hold. 

\begin{proposition}\label{prop:XII.3.Proposition.2} Let $X$ be an $F$-lattice. Let a sequences $\{x_{n}\} \subseteq X$ satisfy $\slim\limits_{n \to \infty} x_{n} = x$. Then $\{x_{n}\}$ *-converges to $x$ relative uniformly. This means that, from any subsequence $\{y_{n}\}$ of $\{x_{n}\}$, we can choose a subsequence $\{y_{n(k)}\}$ and a $z \in X$ such that 
\begin{equation}\label{eq:XII.3.6}
\left| y _{n(k)} - x \right| \leqq k ^{-1} z \quad (k = 1, 2, \dots).
\end{equation}
\end{proposition}

Conversely, if $\{x_{n}\}^{*}$-converges relative uniformly to $x$, then \slim $x_{n} = x$. 

\begin{proof}
 We may restrict ourselves to the case $x = 0$. From $\lim\limits_{n \to \infty} \| y_n \| = 0$, we see that there exists a sequence $\{n(k)\}$ of positive integers such that $\| k y_{n(k)} \| \leqq k^{-2}$. Then \eqref{eq:XII.3.6} holds by taking $z = \sum\limits_{k=1}^{\infty} |k y_{n(k)}|$. Conversely, let condition \eqref{eq:XII.3.6} with $x = 0$ be satisfied. Then we have $\| y_{n(k)} \| \leqq \| k^{-1} z \|$ from $|y_{n(k)}| \leqq k^{-1} z$. Hence $s \cdot \lim\limits_{k \to \infty} y_{n(k)} = 0$. Therefore, there cannot exist a subsequence $\{y_n\}$ of $\{x_n\}$ such that $\lim\limits_{n \to \infty} \| y_n \| > 0$. 
\end{proof}

\begin{remarku} The above proposition is an abstraction of the fact, that an asymptotically convergent sequence of $M(S, \mathfrak{B}, m)$ with $m(S) < \infty$ contains a subsequence which converges $m$-a.e. 
\end{remarku}


\yossection{4}{A Convergence Theorem of Banach}
\label{sec:XII.4}

This theorem is concerned with the almost everywhere convergence of a sequence of linear operators whose ranges are measurable functions. See S. \textsc{Banach} \cite{banach10}. A lattice-theoretic formulation of the theorem reads as follows (K. \textsc{Yosida} \cite{yosida358}): 

\begin{theoremu} Let $X$ be a real $B$-space with the norm $\|\|$ and $Y$ a $\sigma$-complete $F$-lattice with the quasi-norm $\|\|_1$ such that 
\begin{equation}\label{eq:XII.4.1}
\Olim _{n \to \infty} y _{n} = y \quad \text{ implies } \quad \lim\limits _{n \to \infty} \| y _{n} \| _{1} = \| y \| _{1}.
\end{equation}
\end{theoremu}

Let $\{T_n\}$ be a sequence of bounded linear operators $\in L(X, Y)$. Suppose that 
\begin{equation}\label{eq:XII.4.2}
\Olim _{n \to \infty} | T _{n} x | \text{ exists for those } x ^{\prime} s \in X \text{ which form a set }
\end{equation}

$G$ of the second category. 

Then, for any $x \in X$, $O - \varliminf\limits_{n \to \infty} T_n x$ and $O - \varliminf\limits_{n \to \infty} T_n x$ both exist and the (not necessarily linear) operator $\tilde{T}$ defined by 
\begin{equation}\label{eq:XII.4.3}
\tilde {T} x = \left(O - \varlimsup\limits_{n \rightarrow \infty} T _{n} x\right) - \left(O - \varliminf\limits_{n \rightarrow \infty} T _{n} x\right)
\end{equation}
is continuous as an operator defined on $X$ into $Y$. 

\begin{remarku} The space $M(S, \mathfrak{B}, m)$ with $m(S) < \infty$ satisfies \eqref{eq:XII.4.1} if we write $y_1 \leqq y_2$ when and only when $y_1(s) \leqq y_2(s)$ $m$-a.e., the quasi-norm $\|y\|_1$ being defined by $\|y\|_1 = \int_S |y(s)|(1 + |y(s)|)^{-1} m(\dd s)$. Likewise $L^p(S, \mathfrak{B}, m)$ with $m(S) \leqq \infty$ also satisfies \eqref{eq:XII.4.1} by the same semi-order. 
\end{remarku}

\begin{proof}[Proof of the Theorem.]
Set $T_{n}x = y_{n}, y_{n}' = \sup\limits_{n\geqq m}|y_{m}|, y' = \sup\limits_{n\geqq 1}|y_{n}|$, and consider the operators $V_{n}x = y_{n}'$ and $Vx = y'$ defined at least on $G$ into $Y$. By ($23$) of the preceding section $2$, each $V_{n}$ is strongly continuous with $T_{k}$. Since $\lim\limits_{n\to \infty}\|V_n x - V x\|_1 = 0$ by \eqref{eq:XII.4.1}, we have $\lim\limits_{n\to \infty}\|k^{-1}V_n x\| = \|k^{-1}Vx\|$ and further $\lim\limits_{k\to \infty}\|k^{-1}Vx\|_1 = 0$. These are implied by the continuity in $\alpha, y$ of $\alpha y$ in the $F$-space $Y$. Hence for any $\varepsilon > 0$, 
\begin{equation}\label{eq:XII.4.4}
G \subseteq \bigcup\limits_{k=1}^{\infty} G_k \text{ where } G_k = \{x \in X; \sup\limits_{n\geqq 1}\|k^{-1}V_n x\|_1 \leqq \varepsilon\}.
\end{equation}

By the strong continuity of $V_{n}$, each $G_{k}$ is a strongly closed set of $X$. Thus some $G_{k_0}$ contains a sphere of $X$, in virtue of the hypothesis that $G$ is of the second category. That is, there exist an $x_0 \in X$ and a $\delta > 0$ such that $\|x_0 - x\| \leqq \delta$ implies $\sup\limits_{n \geqq 1} \|k_0^{-1} V_n x\|_1 \leqq \varepsilon$. Hence, by putting $z = x_0 - x$, we see that 
\[
\sup\limits _{n \geqq 1} \| k _{0} ^{-1} V _{n} z \| _{1} \leqq \sup\limits _{n \geqq 1} \| k _{0} ^{-1} V _{n} x _{0} \| _{1} + \sup\limits _{n \geqq 1} \| k _{0} ^{-1} V _{n} x \| _{1} \leqq 2 \varepsilon ,
\]
that is, since $V_{n}(k_{0}^{-1}z) = k_{0}^{-1}V_{n}z$, we have 
\[
\sup\limits _{n \geqq 1} \| V _{n} z \| _{1} \leqq 2 \varepsilon \text{ whenever } \| z \| \leqq \delta / k _{0}.
\]

This proves that $\slim\limits_{\|z\| \to 0} V_n \cdot z = 0$ uniformly in $n$. 

$G$ being dense in $X$, $V \cdot x$ is defined for all $x \in X$ and $V \cdot x$ is strongly continuous at $x = 0$ with $V \cdot 0 = 0$. Hence, by 
\[
| \tilde {T} \cdot x | \leqq 2 V \cdot x \quad \text{ and } \quad \| \tilde {T} x _{1} - \tilde {T} x _{2} \| _{1} \leqq \| \tilde {T} (x _{1} - x _{2}) \| _{1},
\]
we see that $\tilde{T} \cdot x$ is strongly continuous at every $x \in X$. 
\end{proof}

\begin{corollaryu} Under condition \eqref{eq:XII.4.1}, the set $G = \{x \in X; O - \lim\limits_{n \to \infty} T_n x \text{ exists }\}$ either coincides with $X$ or is a set of the first category. 
\end{corollaryu}

\begin{proof}
 Suppose the set $G$ be of the second category. Then, by our Theorem, the operator $\tilde{T}$ is a strongly continuous operator defined on $X$ into $Y$. Thus $G = \{x \in X; \tilde{T}y = 0\}$ is strongly closed in $X$. Moreover, $G$ is a linear subspace of $X$. Hence $G$ must coincide with $X$. Otherwise, $G$ would be non dense in $X$. 
\end{proof}


\yossection{5}{The Representation of a Vector Lattice as Point Functions}
\label{sec:XII.5}

Let a vector lattice $X$ contain a ``unit'' $I$ with the properties: 

$I > 0$, and, for any $f \in X$, there exists an $\alpha > 0$ such 
\begin{equation}\label{eq:XII.5.1}
\text{ that } - \alpha I \leqq f \leqq \alpha I.
\end{equation}

For such a vector lattice $X$, we can give an analogous representation similar to the representation of a normed ring as point functions. 

An element $f \in X$ is called ``nilpotent'' if $n |f| \leqq I$ ( $n = 1, 2, \ldots$). The set $R$ of all nilpotent elements $f \in X$ is called the ``radical'' of $X$. By ($20$) of Chapter~\ref{chap:XII}, \ref{sec:XII.2}, $R$ constitutes a linear subspace of $X$. Moreover, $R$ is an ``ideal'' of $X$ in the sense that 
\begin{equation}\label{eq:XII.5.2}
f \in R \text{ and } | g | \leqq | f | \text{ imply } g \in R.
\end{equation}

\begin{lemmau} Let $X_{1}$ and $X_{2}$ be vector lattices. A linear operator $T$ defined on $X_{1}$ onto $X_{2}$ is called a lattice homomorphism if 
\begin{equation}\label{eq:XII.5.3}
T (x \underset {\wedge} {\vee} y) = (T x) \underset {\wedge} {\vee} (T y).
\end{equation}
\end{lemmau}

Then \eqref{eq:XII.5.3} implies that $N = \{x \in X_1; Tx = 0\}$ is an ideal of $X_1$. Conversely, if a linear subspace $N$ of $X_1$ is an ideal of $X_1$, then the mapping $x \to Tx = \overline{x} = (\text{ the residue class mod } N \text{ containing } x)$ satisfies \eqref{eq:XII.5.3} by defining $\overline{x}_{\wedge}^{\vee}\overline{y}$ through $\overline{x}_{\wedge}^{\vee}\overline{y}$. 

\begin{proof}
 Let $T$ be a lattice homomorphism. Let $x \in N$ and $|y| \leqq |x|$. Then, by $T(|x|) = T(x \vee -x) = (Tx) \vee (T(-x)) = 0$, we obtain, $0 \leqq Ty^{+} = T(y^{+} \wedge |x|) = Ty^{+} \wedge T|x| = 0$ and so $y^{+} \in N$. Similarly we obtain $y^{-} \in N$ and thus $y = y^{+} + y^{-} \in N$. 


Next let a linear subspace $N = \{x \in X_1; Tx = 0\}$ be an ideal of $X_1$. If $\overline{y} = \overline{z}$, i.e., if $y - z \in N$, then, by \eqref{eq:XII.2.23} in Chapter~\ref{chap:XII}, \ref{sec:XII.2}, we have 
\[
\left| x _{\wedge} ^{\vee} y - x _{\wedge} ^{\vee} z \right| \leqq | y - z | \in N
\]
so that the residue class $(x \stackrel{\vee}{\wedge} y)$ is determined independently of the choice of the representative elements $x$, $y$ from the residue classes $\overline{x}, \overline{y}$, respectively. Hence we can put $\overline{x} \stackrel{\vee}{\wedge} \overline{y} = \overline{x} \stackrel{\vee}{\wedge} y$ by defining $\overline{z} \geqq \overline{0}$ through $z \vee 0 = \overline{z}$.
\end{proof}

\begin{remarku} The above Lemma may be phrased as follows. Let $N$ be a linear subspace of a vector lattice. Then the linear-congruence $a \equiv b \pmod{N}$ is also a lattice-congruence: 
\[
a \equiv b, a ^{\prime} \equiv b ^{\prime} (\mathrm{mod} N) \quad \text{ implies } \quad a \underset {\wedge} {\vee} b \equiv a ^{\prime} \underset {\wedge} {\vee} b ^{\prime} (\mathrm{mod} N),
\]
\end{remarku}

iff $N$ is an ideal of $X$. 

Now an ideal $N$ is called ``non-trivial'' if $N \neq \{0\}$, $X$. A non-trivial ideal $N$ is called ``maximal'' if it is contained in no other ideal $\neq X$. Denote by $M$ the set of all maximal ideals $N$ of $X$. The residual class $X$/$N$ of $X$ mod any ideal $N \in M$ is ``simple'', that is, $X$/$N$ does not contain non-trivial ideals. It will be proved below that a simple vector lattice with a unit is linear-lattice isomorphic to the vector lattice of real numbers, the non-negative elements and the unit $I$ being represented by non-negative numbers and the number $1$. We denote by $f(N)$ the real number which corresponds to $f \in X$ by the linear-lattice-homomorphism $X \to X/N$, $N \in M$. 

After these preliminaries we may state 

\begin{theorem}\label{thm:XII.5.Theorem.1} The radical $R$ coincides with the intersection ideal $\bigcap\limits_{N \in \mathfrak{M}} N$. 
\end{theorem}

\begin{proof}
 The first step. Let $X$ be a simple vector lattice with a unit $I$. Then we must have $X = \{\alpha I; -\infty < \alpha < \infty\}$. 


\begin{proof}
 $X$ does not contain a nilpotent element $f \neq 0$, for otherwise $X$ would contain a non-trivial ideal $N = \{g; |g| \leqq \eta |f|$ with some $\eta < \infty\}$. Hence, by \eqref{eq:XII.5.1}, $X$ satisfies the Archimedean axiom: 
\begin{equation}\label{eq:XII.5.4}
\underset {n \rightarrow \infty} {\text{ order-lim }} n ^{-1} | x | = 0 \quad \text{ for all } \quad x \in X. \end{equation}

Suppose there exists an $f_0 \in X$ such that $f_0 \neq \gamma I$ for any real number $\gamma$. Let 
\[
\alpha = \inf\limits _ {f _{0} \leqq \alpha^{\prime} I} \alpha^{\prime}, \beta = \sup\limits _ {\beta^{\prime} I \leqq f _{0}} \beta^{\prime}.
\]

Then, by \eqref{eq:XII.5.4}, $\beta I \leqq f_0 \leqq \alpha I$ and $\beta < \alpha$. Hence $(f_0 - \delta I)^+ \neq 0$, $(f_0 - \delta I)^- \neq 0$ for $\beta < \delta < \alpha$. Thus, by $x^+ \wedge (-x^-) = 0$, the set $N_0 = \{g; |g| \leqq \eta(f_0 - \delta I)^+$ with some $\eta < \infty\}$ is a non-trivial ideal, contrary to the hypothesis. 
\end{proof}

The second step. For any non-trivial ideal $N_0$, there exists a maximal ideal $N_1$ containing $N_0$. 

\begin{proof}
 Let $\{N_0\}$ be the totality of non-trivial ideals containing $N_0$. We order the ideals of $\{N_0\}$ by inclusion relation, that is, we denote $N_{\alpha_1} \leqq N_{\alpha_2}$ if $N_{\alpha_1}$ is a subset of $N_{\alpha_2}$. Suppose that $\{N_\alpha\}$ is a linearly ordered subset of $\{N_0\}$ and set $N_\beta = \bigcup\limits_{N_\alpha \in \{N_\alpha\}} N_\alpha$. We shall then show that $N_\beta$ is an upper bound of $\{N_\alpha\}$. For, if $x, y \in N_\beta$, there exist ideals $N_{\alpha_1}$ and $N_{\alpha_2}$ such that $x \in N_{\alpha_1}$ and $y \in N_{\alpha_2}$. Since $\{N_\alpha\}$ is linearly ordered, $N_{\alpha_1} \subseteq N_{\alpha_2}$ (or $N_{\alpha_1} \supseteq N_{\alpha_2}$) and so $x$ and $y$ both belong to $N_{\alpha_2}$. This proves that $(\gamma x + \delta y) \in N_{\alpha_1} \subseteq N_\beta$ and that $|z| \leqq |x|$ implies $z \in N_{\alpha_1} \subseteq N_\beta$. Since the unit $I$ is contained in no $N_\alpha$, $I$ is not contained in $N_\beta$. Therefore $N_\beta$ is a non-trivial ideal containing every $N_\alpha$, i.e., $N_\beta$ is an upper bound of $\{N_\alpha\}$. Thus, by Zorn's Lemma, there exists at least one maximal ideal containing $N_0$. 
\end{proof}

The third step. $R \subseteq \bigcap\limits_{N \in \mathfrak{M}} N$. Let $f > 0$ and $nf \leqq I (n = 1, 2, \ldots)$. Then; for any $N \in \mathfrak{M}$, we have $nf(N) \leqq I(N) = 1 (n = 1, 2, \ldots)$, and hence $f(N) = 0$, that is, $f \in N$. 

The fourth step. $R \supseteq \bigcap\limits_{N \in \mathfrak{M}} N$. Let $f > 0$ be not nilpotent. Then we have to show that there exists an ideal $N \in \mathfrak{M}$ such that $f \bar{\in} N$. This may be proved as follows. 

Since $f > 0$ is not nilpotent, there exists an integer $n$ such that $nf \not\leqq I$. We may assume that $nf \not\leqq I$, since otherwise $f \not\in N$ for any $N \in \mathfrak{M}$ and so we have nothing to prove. Thus suppose $p = I - (n \cdot f) \cap I > 0$. Then, for any positive integer $m$, we do not have $m \cdot p \geqq I$. If otherwise, we would have $m^{-1}I \leqq I - (n \cdot f) \cap I$ and hence 
\[
(n \cdot f) \wedge I = (n \cdot f) \wedge (1 - m ^{-1}) I.
\]

Thus, by \eqref{eq:XII.2.6} in Chapter~\ref{chap:XII}, \ref{sec:XII.2}, 
\[
(n \cdot f - (1 - m ^{-1}) I) \wedge m ^{-1} I = (n \cdot f - (1 - m ^{-1}) I) \wedge 0 \leqq 0,
\]
and so, by the distributivity of the vector lattice, 
\[
0 = \{(n \cdot f - (1 - m ^{-1}) I) \wedge m ^{-1} I \} \vee 0 = (n \cdot f - (1 - m ^{-1}) I) ^{+} \wedge m ^{-1} I,
\]
that is, $(n \cdot f - (1 - m^{-1}) I)^{+} \wedge I = 0$. Put $b = (n \cdot f - (1 - m^{-1}) I)^{+}$ and assume that $b > 0$. By hypothesis \eqref{eq:XII.5.1}, we have $b < \alpha I$ with some $\alpha > 1$. Then $0 < b = b \wedge \alpha I$ and so $0 < (\alpha^{-1} b) \wedge I \leqq b \wedge I$, contrary to $b \wedge I = 0$. Therefore $b = 0$, i.e., $n \cdot f \leqq (1 - m^{-1}) I$. This contradicts the fact that $n \cdot f \not\leqq I$. Hence the set $N_{0} = \{g; |g| \leqq \eta |p|$ for some $\eta < \infty\}$ is a non-trivial ideal. $N_{0}$ is contained in at least one maximal ideal $N$, by the second step. Then $0 = p(N) = 1 - (n \cdot f(N)) \wedge 1$ which shows that $f(N) > 0$, that is, $f \not\in N$. 

We have thus proved our Theorem~\ref{thm:XII.5.Theorem.1}.
\end{proof} 

The vector lattice $\overline{X} = X / R$ is again a vector lattice with a unit $\overline{I}$. By Theorem~\ref{thm:XII.5.Theorem.1}, the intersection ideal $\bigcap\limits_{\overline{N}} \overline{N}$ of all the maximal ideals $\overline{N}$ of $\overline{X}$ is the zero ideal and $\overline{X}$ contains no nilpotent element $\neq 0$. Hence $\overline{X}$ satisfies the Archimedean axiom 
\begin{equation}\label{eq:XII.5.5}
\underset {n \uparrow \infty} {\text{ order-limit }} n ^{-1} | \bar {f} | = 0 \quad \text{ for all } \quad \bar {f} \in \overline {{X}}.
\end{equation}

Let $\bar{N}$ be any maximal ideal of $\bar{X}$. Then the factor space $\bar{X}/\bar{N}$ is a simple vector lattice, and so, by the first step in the proof of Theorem~\ref{thm:XII.5.Theorem.1}, $\bar{X}/\bar{N}$ is linear-lattice-isomorphic to the vector lattice of real numbers; the non-negative elements and the unit being represented by non-negative numbers and $1$. We denote by $\overline{f}(\overline{N})$ the real number which corresponds to $\overline{f}$ by this homomorphism $\overline{X} \to \overline{X}/\overline{N}$. We also denote by $\overline{\mathfrak{M}}$ the set of all maximal ideals of $\overline{X}$. We have thus 

\begin{theorem}\label{thm:XII.5.Theorem.2} By the correspondence $\overline{f} \to \overline{f}(\overline{N})$, $\overline{X}$ is linear-lattice-isomorphically mapped on the vector lattice $F(\mathfrak{M})$ of real-valued bounded functions on $\overline{\mathfrak{M}}$ such that \begin{enumerate}[(i)]
\item $|\overline{f}| \to |\overline{f}(\overline{N})|$,
\item $\overline{I}(\overline{N}) \equiv 1$ on $\overline{\mathfrak{M}}$
\item $F(\overline{\mathfrak{M}})$ separates the points of $\overline{\mathfrak{M}}$ in the sense that 

\begin{equation}\label{eq:XII.5.6}
\begin{minipage}[b]{0.92\linewidth}
for two different points $\overline{N}_1, \overline{N}_2$ of $\overline{\mathfrak{M}}$, there exists at least one $\overline{f} \in \overline{X}$ such that $\overline{f}(\overline{N}_1) \neq \overline{f}(\overline{N}_2)$.
\end{minipage}
\end{equation}
\end{enumerate}
\end{theorem}


\begin{remarku} We introduce a topology in $\overline{M}$ by calling the sets of the form 
\begin{align}
\{\bar {N}& \in \overline {{\mathfrak {M}}}; | \bar {f} _{i} (\bar {N}) - \bar {f} _{i} (\bar {N} _{0}) | < \varepsilon_{i} (i = 1, 2, \dots , n), \notag \\
\text{ where } - \bar {I}& \leqq \bar {f} \leqq \bar {I} \quad \text{ for all } i \} \notag
\end{align}
\end{remarku}

neighbourhoods of $N_0$. Then $\overline{\mathfrak{M}}$ is compact since it may be identified with a closed subset of a topological product (of the same potency as the cardinal number of the set of elements $\bar{f} \in X$ which satisfy $-\bar{I} \leqq \bar{f} \leqq \bar{I}$) of the closed intervals [---$1$, $1$]. The proof is entirely similar to the case of the set of all maximal ideals of a normed ring in Chapter~\ref{chap:XI}, \ref{sec:XI.2}. Moreover, each function $\bar{f}(\overline{N}) \in F(\overline{\mathfrak{M}})$ is continuous on the compact space $\overline{\mathfrak{M}}$ topologized in this way. Thus, by the Kakutani-Krein theorem in Chapter~\ref{chap:0}, \ref{sec:0.2}, we see that $F(\overline{\mathfrak{M}})$ is dense in the $B$-space $C(\overline{\mathfrak{M}})$. The above two theorems are adapted from K. \textsc{Yosida}-M. \textsc{Fukamiya} \cite{yosida359}. Cf. also S. \textsc{Kakutani} \cite{kakutani128} and M. \textsc{Krein}-S. \textsc{Krein} \cite{krein166}. 


\yossection{6}{The Representation of a Vector Lattice as Set Functions}
\label{sec:XII.6}

Let $X$ be a $\sigma$-complete vector lattice. Choose any positive element $x$ of $X$ and call it a ``unit'' of $X$ and write $1$ for $x$; when not ambiguous we also write $\alpha$ for $\alpha \cdot 1$. A non-negative element $e \in X$ is called a ``quasi-unit'' if $e \wedge (1 - e) = 0$. A finite linear combination $\sum\limits_{i} \alpha_{i} e_{i}$ of quasi-units $e_{i}$ is called a ``step-element'', and we call the element $y \in X$ ``absolutely continuous'' (with respect to the unit $1$) if $y$ can be expressed as the \Olim of a sequence of step-elements. An element $z \in X$ is called ``singular'' (with respect to the unit $1$) if $|z| \wedge 1 = 0$. 

We shall give an abstract formulation of the Radon-Nikodym theorem in integration theory. 

\begin{theoremu} Any element of $X$ is uniquely expressed as the sum of an absolutely continuous element and a singular element. 
\end{theoremu}

\begin{proof}
 The first step. If $f > 0$ and $f \wedge 1 \neq 0$, then there exist a positive number $\alpha$ and a quasi-unit $e_{\alpha} \neq 0$ such that $f \geqq \alpha e_{\alpha}$. We can, 

in fact, take 
\begin{equation}\label{eq:XII.6.1}
e _{\alpha} = \bigvee\limits_{n \geqq 1} \left\{n \left(\alpha^{-1} f - \alpha^{-1} f \wedge 1\right) \wedge 1 \right\}. \end{equation}


\begin{proof}
 Put $y_{\alpha} = \alpha^{-1}f - \alpha^{-1}; \wedge 1$. Then we obtain 
\[
2 e _{a} \wedge 1 = \left\{\bigvee\limits_{n \geqq 1} (2 n y _{a} \wedge 2) \right\} \wedge 1 = e _{\alpha},
\]
and so $e_{\alpha}$ is a quasi-unit. We obtain $f \geqq \alpha e_{\alpha}$ from 
\[ \begin{aligned} n y _{\alpha} \wedge 1 & = n \alpha^{-1} f \wedge [ 1 + n (\alpha^{-1} f \wedge 1) ] - n (\alpha^{-1} f \wedge 1) \\ & \leqq (n + 1) \alpha^{-1} f \wedge (n + 1) - n (\alpha^{-1} f \wedge 1) \leqq \alpha^{-1} f \wedge 1 \leqq \alpha^{-1} f. \end{aligned}
\]

If we can show that $y_{\alpha} \wedge 1 > 0$ for some $\alpha > 0$, then $e_{\alpha} > 0$ for such $\alpha$. Suppose that such positive $\alpha$ does not exist. Then, for every $\alpha$ with $0 < \alpha < 1$, we have 
\[
\alpha^{-1} (\alpha^{-1} f - \alpha^{-1} f \wedge 1) \wedge \alpha^{-1} = 0.
\]

Hence $(f - f \wedge \alpha) \wedge 1 = 0$ and, letting $\alpha \downarrow 0$, we obtain $f \wedge 1 = 0$, contrary to the hypothesis $f \wedge 1 \neq 0$. 
\end{proof}

The second step. Let $f \geqq 0$, and $f \geqq \alpha e$ where $\alpha > 0$ and $e$ is a quasi-unit. Then, for $0 < \alpha' < \alpha$, $e_{\alpha'} \geqq e$ and $f \geqq \alpha' e_{\alpha'}$ where $e_{\alpha'}$ is defined by \eqref{eq:XII.6.1}. 

\begin{proof}
 For the sake of simplicity, we assume that $\alpha = 1$. For $0 < \delta < 1$ 
\[ \begin{aligned}
&\frac {f}{1 - \delta} + \frac {e}{1 - \delta} \wedge 1  = \left(\frac {f}{1 - \delta} + \frac {e}{1 - \delta}\right) \wedge \left(1 + \frac {f}{1 - \delta}\right) \geqq \left(\frac {f}{1 - \delta} + \frac {e}{1 - \delta}\right) \wedge \left(1 + \frac {e}{1 - \delta}\right) \\
&= \frac {f}{1 - \delta} \wedge 1 + \frac {e}{1 - \delta}.
\end{aligned}
\]

Since $e$ is a quasi-unit, we have $2e \wedge 1 = e$, $e \wedge 1 = e$. Hence $me \wedge 1 = e$ when $m \geqq 1$. Thus, by $1 < (1 - \delta)^{-1}$, we have $(1 - \delta)^{-1}e \wedge 1 = e$. Therefore, we obtain, from the above inequality, 
\[
\frac {\delta}{1 - \delta} e = \frac {e}{1 - \delta} - \frac {e}{1 - \delta} \wedge 1 \leqq \frac {f}{1 - \delta} - \frac {f}{1 - \delta} \wedge 1 = y _{1- \delta},
\]
and so, by \eqref{eq:XII.6.1}, $e \leqq e_{1- \delta}$. 
\end{proof}

The third step. The set of all quasi-units constitutes a Boolean algebra, that is, if $e_{1}$ and $e_{2}$ are quasi-units, then $e_{1} \vee e_{2}$ and $e_{1} \wedge e_{2}$ are also quasi-units and $0 \leqq e_{i} \leqq 1$. The quasi-unit $(1 - e)$ is the complement of $e$, and $0$,$1$ are the least and the greatest elements, respectively in the totality of the quasi-units. 

\begin{proof}
 The condition $e \wedge (1 - e) = 0$ is equivalent to $2e \wedge 1 = e$. Hence, if $e_1, e_2$ are quasi-units, then 
\begin{align}
2 \left(e _{1} \wedge e _{2}\right) \wedge 1& = \left(2 e _{1} \wedge 1\right) \wedge \left(2 e _{2} \wedge 1\right) = e _{1} \wedge e _{2}, \notag \\
2 \left(e _{1} \vee e _{2}\right) \wedge 1& = \left(2 e _{1} \wedge 1\right) \vee \left(2 e _{2} \wedge 1\right) = e _{1} \vee e _{2}, \notag
\end{align}
so that $e_1 \wedge e_2$ and $e_1 \vee e_2$ are also quasi-units. 
\end{proof}

The fourth step. Let $f > 0$, and set $\overline{f} = \sup \beta e_{\beta}$ where sup is taken for all positive rational numbers $\beta$. Because of the third step, the sup of a finite number of elements of the form $\beta_i e_{\beta_i}$ is a step-element. Hence $\overline{f}$ is absolutely continuous with respect to the unit element $1$. We have to show that $g = f - \overline{f}$ is singular with respect to the unit element $1$. Suppose that $g$ is not singular. Then, by the first step, there exist a positive number $\alpha$ and a quasi-unit $e$ such that $g \geqq \alpha e$. Hence $f \geqq \alpha e$, and so, by the second step, there exists, for $0 < \alpha_1 < \alpha$, a quasi-unit $e_{a_1} \geqq e$ such that $f \geqq \alpha_1 e_{\alpha_1}$. We may assume that $\alpha_1$ is a rational number. Thus $\overline{f} \geqq \alpha_1 e_{\alpha_1}$ and so $f = \overline{f} + g \geqq 2\alpha_1 e_{\alpha_1}$. Again, by the second step, there exists, for $0 < \alpha_1' < \alpha_1$, a quasi-unit $e_{2\alpha_1'} \geqq e_{\alpha_1}$ such that $f \geqq 2\alpha_1' e_{2\alpha_1'}$. We may assume that $2\alpha_1'$ is a rational number so that $\overline{f} \geqq 2\alpha_1' e_{2\alpha_1'}$. Hence $f = \overline{f} + g \geqq 3\alpha_1' e$. Repeating the process, we can prove that, for any rational number $\alpha_n$ with $0 < \alpha_n < \alpha$, 
\[
f \geqq (n + 1) \alpha_{n} e \quad (n = 1, 2, \dots).
\]

If we take $\alpha_{n} \geqq \alpha / 2$, we have $(n + 1)\alpha_{n}e \geqq 2^{-1}n\alpha e$. Hence $f \geqq n\alpha e$ ( $n = 1, 2, \ldots$) with $\alpha > 0, e > 0$. This is a contradiction. For, by the $\sigma$-completeness of $X$, the Archimedean axiom holds in $X$. 

The fifth step. Let $f = f^{+} + f^{-}$ be the Jordan decomposition of a general element $f \in X$. Applying the fourth step to $f^{+}$ and $f^{-}$ separately, we see that $f$ is decomposed as the sum of an absolutely continuous element and a singular element. The uniqueness of the decomposition is proved if we can show that an element $h \in X$ is $= 0$ if $h$ is absolutely continuous as well as singular. But, since $h$ is absolutely continuous, we have $h = \Olim_{n \to \infty} h_{n}$ where $h_{n}$ are step-elements. $h_{n}$ being a step-element, there exists a positive number $\alpha_{n}$ such that $|h_{n}| \leqq \alpha_{n} \cdot 1$. Since $h$ is singular, we have $|h| \wedge |h_{n}| = 0$. Therefore $|h| = |h| \wedge |h| = \Olim_{n \to \infty} (|h| \wedge |h_{n}|) = 0$. 
\end{proof}

Application to the Radon-Nikodym theorem. Consider the case $X = A(S, \mathfrak{B})$. We know already (Proposition~\ref{prop:XII.3.Proposition.1} in Chapter~\ref{chap:XII}, \ref{sec:XII.3}) that it is a complete lattice, and that in $A(S, \mathfrak{B})$, 
\begin{equation}\label{eq:XII.6.2}
x^{+}(B) = \sup\limits_{B^{\prime}\subseteq B}x(B^{\prime}) = \text{ the positive variation } \overline{V} (x;B) \text{ of } x \text{ on } B.
\end{equation}

We shall prepare a 


\begin{propositionu} In $A(S, \mathfrak{B})$, let $x > 0$, $z \geqq 0$ in such a way that $x \wedge z = 0$. Then there exists a set $B_0 \in \mathfrak{B}$ such that $x(B_0) = 0$ and $z(S - B_0) = 0$. 
\end{propositionu}


\begin{proof}
 Since $x \wedge z = (x - z)^{-} + z$, we know, from \eqref{eq:XII.6.2}, that 
\begin{equation}\label{eq:XII.6.3}
(x \wedge z) (B) = \inf\limits _ {B ^{\prime} \subseteq B} (x - z) \left(B ^{\prime}\right) + z (B) = \inf\limits _ {B ^{\prime} \subseteq B} \left[ x \left(B ^{\prime}\right) + z \left(B - B ^{\prime}\right) \right]. \end{equation}

Hence, by the hypothesis $x \wedge z = 0$, we have 
\[
\inf\limits _ {B \in \mathfrak {B}} [ x (B) + z (S - B) ] = (x \wedge z) (S) = 0.
\]

Hence, for any $\varepsilon > 0$, there exists a $B_{\varepsilon} \in \mathfrak{B}$ such that $x(B_{\varepsilon}) \leqq \varepsilon$, $z(S - B_{\varepsilon}) \leqq \varepsilon$. Put $B_0 = \bigwedge\limits_{k \geqq 1} \left( \bigvee\limits_{n \geqq k} B_2^{-n} \right)$. Then, by the $\sigma$-additivity of $x(B)$ and $z(B)$, 
\begin{align}
0& \leqq x (B _{0}) = \lim\limits _{k \rightarrow \infty} x \left(\bigvee\limits_{n \geqq k} B _ {2 ^{-n}}\right) \leqq \lim\limits _{k \rightarrow \infty} \sum\limits_{n=k} ^{\infty} 2 ^{-n} = 0, \notag \\
0& \leqq z (S - B _{0}) = \lim\limits _{k \rightarrow \infty} z \left(S - \bigvee\limits_{n \geqq k} B _ {2 ^{-n}}\right) \leqq \lim\limits _{k \rightarrow \infty} z (S - B _ {2 ^{-k}}) = 0. \notag
\end{align}
\end{proof}

\begin{corollaryu} Let $e$ be a quasi-unit with respect to $x > 0$ in $A(S, \mathfrak{B})$. Then there exists a set $B_{1} \in \mathfrak{B}$ such that 
\begin{equation}\label{eq:XII.6.4}
e (B) = x \left(B \cap B _{1}\right) \quad \text{ for all } \quad B \in \mathfrak {B}.
\end{equation}
\end{corollaryu}

\begin{proof}
 Since $(x - e) \wedge e = 0$, there exists a set $B_0 \in \mathfrak{B}$ such that $e(B_0) = 0$, $(x - e)(S - B_0) = 0$. Hence $e(S - B_0) = x(S - B_0) = e(S)$, and so $e(B) = x(B - B_0) = x(B \cap B_1)$ where $B_1 = S - B_0$. 
\end{proof}

We are now able to prove the Radon-Nikodym theorem in integration theory. In virtue of the above Corollary, a quasi-unit $e$ with respect to $x > 0$ is, in $A(S, \mathfrak{B})$, a contracted measure $e(B) = x(B \cap B_e)$. Hence a step-element in $A(S, \mathfrak{B})$ is the integral of the form: 
\[
\sum\limits_{i} \lambda_{i} \int_{B \cap B i} x (\dd s),
\]
that is, indefinite integral of a step function (= finitely-valued function). Hence an absolutely continuous element in $A(S, \mathfrak{B})$ is an indefinite integral with respect to the measure $x(B)$. The Proposition above says that a singular element $g$ (with respect to the unit $x$) is associated with a set $B_0 \in \mathfrak{B}$ such that $x(B_0) = 0$ and $g(B) = g(B \cap B_0)$ for all $B \in \mathfrak{B}$. Such a measure $g(B)$ is the so-called singular measure (with respect to $x(B)$). Thus any element $f \in A(S, \mathfrak{B})$ is expressed as the sum of an indefinite integral (with respect to $x(B)$) and a measure $g(B)$ which is singular (with respect to $x(B)$). This decomposition is unique. The obtained result is exactly the Radon-Nikodym theorem. 

\begin{remarku} The above Theorem is adapted from K. \textsc{Yosida} \cite{yosida345}. Cf. F. \textsc{Riesz} \cite{rickart277}, H. \textsc{Freudenthal} \cite{freudenthal67} and S. \textsc{Kakutani} \cite{kakutani129}. For further references, see G. \textsc{Birkhoff} \cite{birkhoff17}. 
\end{remarku}
   % XII
% ==== XIII ====
\yoschapter{XIII}{Ergodic Theory and Diffusion Theory}
\label{chap:XIII}

These theories constitute fascinating fields of application of the analytical theory of semi-groups. Mathematically speaking, the ergodic theory is concerned with the ``time average'' $\lim\limits_{t\uparrow\infty}t^{-1}\int\limits_{0}^{t}T_{s}\dd s$ of a semi-group $T_{t}$, and the diffusion theory is concerned with the investigation of a stochastic process in terms of the infinitesimal generator of the semi-group intrinsically associated with the stochastic process. 


\yossection{1}{The Markov Process with an Invariant Measure*}
\label{sec:XIII.1}
\yosstarfootnotetext{See also Supplementary Notes, p.~\pageref{notes:XIII}.}

In $1862$, an English botanist R. \textsc{Brown} observed under a microscope that small particles, pollen of some flower, suspended in a liquid move chaotically, changing position and direction incessantly. To describe such a phenomenon, we shall consider the transition probability $P(t,x;s,E)$ that a particle starting from the position $x$ at time $t$ belongs to the set $E$ at a later time $s$. The introduction of the transition probability $P(t,x;s,E)$ is based upon the fundamental hypothesis that the chaotic motion of the particle after the time moment $t$ is entirely independent of its past history before the time moment $t$. That is, the future history of the particle after the time moment $t$ is entirely determined chaotically if we know the position $x$ of the particle at time $t$. The hypothesis that the particle has no memory of the past implies that the transition probability $P$ satisfies the equation 
\begin{equation}\label{eq:XIII.1.1}
P (t, x; s, E) = \int_{S} P (t, x; u, \dd y) P (u, y; s, E) \text{ for } t < u < s,
\end{equation}
where the integration is performed over the entire space $S$ of the chaotic movement of the particle. 

The process of evolution in time governed by a transition probability satisfying \eqref{eq:XIII.1.1} is called a Markov process, and the equation \eqref{eq:XIII.1.1} is called the Chapman-Kolmogorov equation. The Markov process is a natural generalization of the deterministic process for which $P(t, x; s, E) = 1 \text{ or } = 0$ according as $y \in E$ or $y \bar{\in} E$; that is, the process in which the particle at the position $x$ at the time moment $t$ moves to a definite position $y = y(x, t, s)$ with probability $1$ at every fixed later time moment $s$. The Markov process $P$ is said to be temporally homogeneous if $P(t, x; s, E)$ is a function of $(s - t)$ independently of $t$. In such a case, we shall be concerned with the transition probability $P(t, x, E)$ that a particle at the position $x$ is transferred into the set $E$ after the lapse of $t$ units of time. The equation \eqref{eq:XIII.1.1} then becomes 
\begin{equation}\label{eq:XIII.1.2}
P (t + s, x, E) = \int_{S} P (t, x, \dd y) P (s, y, E) \text{ for } t, s > 0.
\end{equation}

In a suitable function space $X$, $P(t, x, E)$ gives rise to a linear transformation $T_t$: 
\begin{equation}\label{eq:XIII.1.3}
(T_t f)(x) = \int_{S} P(t, x, \dd y) f(y), f \in X,
\end{equation}

such that, by \eqref{eq:XIII.1.2}, the semi-group property holds: 
\begin{equation}\label{eq:XIII.1.4}
T _{t+s} = T _{t} T _{s} \quad (t, s > 0).
\end{equation}

A fundamental mathematical question in statistical mechanics is concerned with the existence of the time average 
\begin{equation}\label{eq:XIII.1.5}
\lim\limits _{t \uparrow \infty} t ^{-1} \int_{0} ^{t} T _{s} f \dd s.
\end{equation}

In fact, let $S$ be the phase space of a mechanical system governed by the classical Hamiltonian equations whose Hamiltonian does not contain the time variable explicitly. Then a point $x$ of $S$ is moved to the point $y_{t}(x)$ of $S$ after the lapse of $t$ units of time in such a way that, by a classical theorem due to \textsc{Liouville}, the mapping $x \to y_{t}(x)$ of $S$ onto $S$, for each fixed $t$, is an equi-measure transformation, that is, the mapping $x \to y_{t}(x)$ leaves the ``phase volume'' of $S$ invariant. In such a deterministic case, 
we have 
\begin{equation}\label{eq:XIII.1.6}
\left(T _{t} f\right) (x) = f \left(y _{t} (x)\right),
\end{equation}
and hence the ergodic hypothesis of \textsc{Boltzmann} that 
the time average of any physical quantity = the space 
average of this physical quantity 
is expressed, assuming $\int_{S} \dd x < \infty$, by 
\begin{equation}\label{eq:XIII.1.7}
\lim\limits _{t \uparrow \infty} t ^{-1} \int_{0} ^{t} f (y _{s} (x)) \dd s = \int_{S} f (x) \dd x \Big | \int_{S} \dd x \text{ for all } f \in X,
\end{equation}
$\dd x$ denoting the phase volume element of $S$. 


A natural generalization of the equi-measure transformation $x \to y_t(x)$ to the case of a Markov process $P(t, x, E)$ is the condition of the existence of an invariant measure $m(\dd x)$: 
\begin{equation}\label{eq:XIII.1.8}
\int_{S} m (\dd x) P (t, x, E) = m (E) \text{ for all } t > 0 \text{ and all } E.
\end{equation}

We are thus lead to the 

\begin{definitionu} Let $\mathfrak{B}$ be a $\sigma$-additive family of subsets $B$ of a set $S$ such that $S$ itself $\in \mathfrak{B}$. For every $t > 0, x \in S$ and $E \in \mathfrak{B}$, let there be associated a function $P(t, x, E)$ such that 
\begin{equation}\label{eq:XIII.1.9}
P(t,x,E)\geqq0,\qquad P(t,x,S)=1,
\end{equation}
\begin{equation}\label{eq:XIII.1.10}
\begin{minipage}{0.92\linewidth}
\everymath{}\let\limits\nolimits
For fixed $t$ and $x$, $P(t,x,E)$ is $\sigma$-additive in $E\in\mathfrak{B}$,
\end{minipage}
\end{equation}
\begin{equation}\label{eq:XIII.1.11}
\begin{minipage}{0.92\linewidth}
\everymath{}\let\limits\nolimits
For fixed $t$ and $E$, $P(t,x,E)$ is $\mathfrak{B}$-measurable in $x$,
\end{minipage}
\end{equation}
\begin{equation}\label{eq:XIII.1.12}
P(t+s,x,E)=\int_S P(t,x,\dd y)P(s,y,E)\qquad\text{(the Chapman-Kolmogorov equation).}
\end{equation}

\end{definitionu}

Such a system $P(t, x, E)$ is said to define a Markov process on the phase space $(S, \mathfrak{B})$. If we further assume that $(S, \mathfrak{B}, m)$ is a measure space in such a way that 
\begin{equation}\label{eq:XIII.1.13}
\int_{S} m (\dd x) P (t, x, E) = m (E) \text{ for all } E \in \mathfrak {B},
\end{equation}
then $P(t, x, E)$ is said to be a Markov process with an invariant measure $m(E)$. 

\begin{theorem}\label{thm:XIII.1.Theorem.1} Let $P(t, x, E)$ be a Markov process with an invariant measure $m$ such that $m(S) < \infty$. Let the norm in $X_{p} = L^{p}(S, \mathfrak{B}, m)$ be denoted by $\|f\|_p$, $p \geqq 1$. Then, by \eqref{eq:XIII.1.3}, a bounded linear operator $T_{t} \in L(X_{p}, X_{p})$ is defined such that $T_{t+s} = T_{t}T_{s} (t, s > 0)$ and 

\begin{equation}\label{eq:XIII.1.14}
\parbox{0.92\linewidth}{\everymath{}\let\limits\nolimits $T_{t}$ is positive, i.e., $(T_{t}f)(x)\geqq 0$ on $S$ $m$-a.e. if $f(x)\geqq 0$ on $S$ $m$-a.e.,}
\end{equation}
\end{theorem}

\begin{align}
T _{t} \cdot 1& = 1, \label{eq:XIII.1.15} \\
\left\| T _{t} f \right\| _{p}& \leqq \left\| f \right\| _{p} \text{ for } f \in X _{p} = L ^{p} (S, \mathfrak {B}, m) \text{ with } p = 1, 2 \text{ and } \infty . \label{eq:XIII.1.16}
\end{align}

\begin{proof}
 \eqref{eq:XIII.1.14} and \eqref{eq:XIII.1.15} are clear. Let $f \in L^{\infty}(S, \mathfrak{B}, m)$. Then, by \eqref{eq:XIII.1.9}, \eqref{eq:XIII.1.10} and \eqref{eq:XIII.1.11}, we see that $f_t(x) = (T_t f)(x) \in L^{\infty}(S, \mathfrak{B}, m)$ is defined and $\|f_t\|_{\infty} \leqq \|f\|_{\infty}$. Hence, by \eqref{eq:XIII.1.9} and \eqref{eq:XIII.1.13}, we obtain, for $t \in L^{\infty}(S, \mathfrak{B}, m)$ with $p = 1$ or $p = 2$, 
\[
\begin{aligned} & {\| f _{t} \| _{p} = \left\{\int_ {\mathcal {S}} m (\dd x) \left| \int_ {\mathcal {S}} P (t, x, \dd y) f (y) \right| ^{p} \right\} ^{1/p}} \\ & {\quad \leqq \left\{\int_ {\mathcal {S}} m (\dd x) \left[ \int_ {\mathcal {S}} P (t, x, \dd y) | f (y) | ^{p} \cdot \int_ {\mathcal {S}} P (t, x, \dd y) 1 ^{p} \right] \right\} ^{1/p}} \\ & {\quad = \left\{\int_ {\mathcal {S}} m (\dd y) | f (y) | ^{p} \right\} ^{1/p} = \| f \| _{p}.} \end{aligned}
\]

We put, for a non-negative $f \in L^p(S, \mathfrak{B}, m)$ with $p = 1$ or $p = 2$, 
$_{n}f(s)=\min(f(s),n)$, where $n$ is a positive integer. 

Then, by the above, we obtain $0 \leqq (_n f(s))_t \leqq (_{n+1} f(s))_t$ and $\|(n f)_t\|_p \leqq \|_n f\|_p \leqq \|f\|_p$. Thus, if we put $f_t(s) = \lim\limits_{n \to \infty} (_n f(s))_t$, then, by the Lebesgue-Fatou lemma, $\|f_t\|_p \leqq \|f\|_p$, that is, $f_t \in L^p(S, \mathfrak{B}, m)$. Again by the Lebesgue-Fatou lemma, we obtain 
\[
\begin{aligned}
&f _{t} (x) = \lim\limits _{n \rightarrow \infty} \int_{S} P (t, x, \dd y) (n f (y)) \geqq \int_{S} P (t, x, \dd y) (\lim\limits _{n \rightarrow \infty} (n f (y))) = \int_{S} P (t, x, \dd y) f (y).
\end{aligned}
\]

Thus $f(y)$ is integrable with respect to the measure $P(t, x_0, \dd y)$ for those $x_0$'s for which $f_t(x_0) \neq \infty$, that is, for $m$-a.e. $x_0$. Hence, by the Lebesgue-Fatou lemma, we finally obtain 
\[
\int_{S} P (t, x _{0}, \dd y) \left(\lim\limits _{n \rightarrow \infty} (_{n} f (y))\right) = \lim\limits _{n \rightarrow \infty} \int_{S} P (t, x _{0}, \dd y) (_{n} f (y)).
\]

Therefore, $f_{t}(x_0) = \int \limits_{S} P(t, x_0, \dd y) f(y)$ $m$-a.e. and $\|f_t\|_p \leqq \|f\|_p$. For a general $f \in L^p(S, \mathfrak{B}, m)$, we obtain the same result by applying the positive operator $T_t$ to $f^+$ and $f^-$, separately. 
\end{proof}

\begin{theorem}[K. \textsc{Yosida}]\label{thm:XIII.1.Theorem.2} Let $P(t, x, E)$ be a Markoff process with an invariant measure $m$ such that $m(S) < \infty$. Then, for any $f \in L^p(S, \mathfrak{B}, m)$ with $p = 1$ or $p = 2$, the mean ergodic theorem holds: 
\begin{equation}\label{eq:XIII.1.17}
\begin{aligned}
&\slim _{n \to \infty} n ^{-1} \sum\limits_{k=1} ^{n} T _{k} f = f ^{*} \text{ exists in } L ^{p} (S, \mathfrak {B}, m) \text{ and } T _{1} f ^{*} = f \text{whenever } f \in L^{p}(S, \mathfrak{B}, m),
\end{aligned}
\end{equation}


and we have 
\begin{equation}\label{eq:XIII.1.18}
\int_{S} f (s) m (\dd s) = \int_{S} f ^{*} (s) m (\dd s).
\end{equation}
\end{theorem}

\begin{proof}
 In the Hilbert space $L^2(S, \mathfrak{B}, m)$, the mean ergodic theorem \eqref{eq:XIII.1.17} holds by \eqref{eq:XIII.1.16} and the general mean ergodic theorem in Chapter~\ref{chap:VIII}, \ref{sec:VIII.3}. Since $m(S) < \infty$, we see, by Schwarz' inequality, that any $f \in L^2(S, \mathfrak{B}, m)$ belongs to $L^1(S, \mathfrak{B}, m)$ and $\|f\|_1 \leqq \|f\|_2 \cdot m(S)^{1/2}$. Hence the mean ergodic theorem $\lim\limits_{n \to \infty} \int_S |f^*(s) - n^{-1} \sum\limits_{k=1}^{n} (T_k f)(s)|m(\dd s) = 0$, together with $T_1 f^* = f^*$, hold for any $f \in L^2(S, \mathfrak{B}, m)$. Again, by $m(S) < \infty$ and $0 = \lim\limits_{n \to \infty} \|f - _n f\|_1 = \lim\limits_{n \to \infty} \int_S |f(s) - _n f(s)|m(\dd s)$ with $_n f(s) = \min(f(s), n)$, we see that $L^2(S, \mathfrak{B}, m)$ is $L^1$-dense in $L^1(S, \mathfrak{B}, m)$. That is, for any $f \in L^1(S, \mathfrak{B}, m)$ and $\varepsilon > 0$, there exists an $f_\varepsilon \in L^2(S, \mathfrak{B}, m) \cap L^1(S, \mathfrak{B}, m)$ such that $\|f - f_\varepsilon\|_1 < \varepsilon$. Hence, by \eqref{eq:XIII.1.16}, we obtain 
\[
\left\| n ^{-1} \sum\limits_{k=1} ^{n} T _{k} f - n ^{-1} \sum\limits_{k=1} ^{n} T _{k} f _{\varepsilon} \right\| _{1} \leqq \| f - f _{\varepsilon} \| _{1} \leqq \varepsilon . \]

The mean ergodic theorem ($17$) in $L^{1}(S, \mathfrak{B}, m)$ holds for $f_{\varepsilon}$ and so, by the above inequality, we see that ($17$) in $L^{1}(S, \mathfrak{B}, m)$ must hold also for $f$. 

Since the strong convergence implies the weak convergence, \eqref{eq:XIII.1.18} is a consequence of \eqref{eq:XIII.1.17}. 
\end{proof}

\begin{remarkn}[Remark $1$] The above Theorem~\ref{thm:XIII.1.Theorem.2} is due to K. \textsc{Yosida} \cite{yosida360}. Cf. S. \textsc{Kakutani} \cite{kakutani130}, where, moreover, the $m$-a.e. convergence $\lim\limits_{n\to \infty}n^{-1}\sum\limits_{k=1}^{n}(T_kf)(x)$ is proved for every $f\in L^{\infty}(S,\mathfrak{B},m)$. We can prove, when the semi-group $T_{t}$ is strongly continuous in $t$, that we may replace \slim $n^{-1}\sum\limits_{k=1}^{n}T_{k}f$ in ($17$) by \slim $t^{-1}\int_{0}^{t}T_{s}fds$. We shall not go into the details, since they are given in the books on ergodic theory due to E. \textsc{Hopf} \cite{hirsch109} and K. \textsc{Jacobs} \cite{ito123}. We must mention a fine report by S. \textsc{Kakutani} \cite{kakutani132} on the development of ergodic theory between \textsc{Hopf}'s report in $1937$ and the $1950$ International Congress of Mathematicians held at Cambridge. 
\end{remarkn}

In order to deal with the ergodic hypothesis ($7$), we must prove the individual ergodic theorem to the effect that 
\[
\lim\limits _{n \rightarrow \infty} n ^{-1} \sum\limits_{k=1} ^{n} (T _{k} f) (x) = f ^{*} (x) m\text{-a.e.}
\]

In the next section, we shall be concerned with the $m$-a.e. convergence of the sequence $n^{-1} \sum\limits_{k=1}^{n} (T_k f)(x)$. Our aim is to derive the $m$-a.e. convergence from the mean convergence using Banach's convergence theorem in Chapter~\ref{chap:XII}, \ref{sec:XII.4}. 


\yossection{2}{An Individual Ergodic Theorem and Its Applications}
\label{sec:XIII.2}

We first prove 

\begin{theorem}[K. \textsc{Yosida}]\label{thm:XIII.2.Theorem.1} Let $X_{1}$ be a real, $\sigma$-complete $F$-lattice provided with a quasi-norm $\| x \|_{1}$ such that 
\begin{equation}\label{eq:XIII.2.1}
\Olim _{n \to \infty} x _{n} = x \text{ implies } \lim\limits _{n \to \infty} \| x _{n} \| _{1} = \| x \| _{1}.
\end{equation}
\end{theorem}

Let a linear subspace $X$ of $X_{1}$ be a real $B$-space provided with a norm $\| x \|$ such that 
\begin{equation}\label{eq:XIII.2.2}
\slim _{n \to \infty} x _{n} = x \quad \text{ in } X \text{ implies } \quad \slim _{n \to \infty} x _{n} = x \text{ in } X _{1}.
\end{equation}

Let $\{T_n\}$ be a sequence of bounded linear operators defined on $X$ into $X$ such that 
\begin{equation}\label{eq:XIII.2.3}
\begin{minipage}[b]{0.92\linewidth}
$\underset {n \rightarrow \infty} {\mathcal {O} - \varlimsup} | T _{n} x |$ exists for those $x$'s which form a set $S$ of the second category in $X$.
\end{minipage}
\end{equation}


Suppose that, to a certain $z \in X$, there corresponds a $\bar{z} \in X$ such that 
\begin{align}
\lim\limits _{n \rightarrow \infty} \| T _{n} z - \bar {z} \|& = 0, \label{eq:XIII.2.4} \\
T _{n} \bar {z}& = \bar {z} (n = 1, 2, \dots), \label{eq:XIII.2.5} \\
\Olim _{n \to \infty} (T _{n} z - T _{n} T _{k} z)& = 0 \quad \text{ for } \quad k = 1, 2, \dots \label{eq:XIII.2.6}
\end{align}

Then 
\begin{equation}\label{eq:XIII.2.7}
\Olim _{n \to \infty} T _{n} z = \bar {z}.
\end{equation}

\begin{proof}
 Put $z = \bar{z} + (z - \bar{z})$. We define an operator $\tilde{T}$ on $S$ into $X_{1}$ by 
\begin{equation}\label{eq:XIII.2.8}
\tilde {T} x = O \underset {n \rightarrow \infty} {\overline {{\lim}}} T _{n} x - O \underset {n \rightarrow \infty} {\underline {{\lim}}} T _{n} x. \end{equation}

Then, by \eqref{eq:XIII.2.5}, 
\begin{equation*}
0 \leqq \tilde {T} z \leqq \tilde {T} (z - \bar {z}).
\end{equation*}

We have $\tilde{T}(z - T_k z) = 0$ ( $k = 1, 2, \ldots$) by \eqref{eq:XIII.2.6}, and $\lim\limits_{k \to \infty} \|(z - \bar{z}) - (z - T_k z)\| = 0$ by \eqref{eq:XIII.2.4}. Hence, by Banach's theorem in Chapter~\ref{chap:XII}, \ref{sec:XII.4}, we have $\tilde{T}(z - \bar{z}) = 0$. Hence $0 \leqq \tilde{T} z \leqq 0$, that is, $O$ - $\lim\limits_{n \to \infty} T_n z = w$ exists. 

We have to show that $w = \bar{z}$. We have, by \eqref{eq:XIII.2.4} and \eqref{eq:XIII.2.2}, $\lim\limits_{n\to \infty}\|T_n z - \bar{z}\|_1 = 0$. Also we have, by \eqref{eq:XIII.2.1} and $O\text{- }\lim\limits_{n\to \infty}T_n z = w$, $\lim\limits_{n\to \infty}\|T_n z - w\|_1 = 0$. Hence $w = \bar{z}$. 
\end{proof}

Specializing $X_{1}$ as the real space $M(S, \mathfrak{B}, m)$ and $X$ as the real space $L^{1}(S, \mathfrak{B}, m)$, respectively, we can prove the following individual ergodic theorem. 

\begin{theorem}[K. \textsc{Yosida}]\label{thm:XIII.2.Theorem.2} Let $T$ be a bounded linear operator defined on $L^{1}(S, \mathfrak{B}, m)$ into $L^{1}(S, \mathfrak{B}, m)$ where $m(S) < \infty$. Suppose that 
\begin{align}
\left\| T ^{n} \right\|& \leqq C < \infty (n = 1, 2, \dots), \label{eq:XIII.2.9} \\
\varlimsup\limits_{n \rightarrow \infty} \left| n ^{-1} \sum\limits_{m=1} ^{n} (T ^{m} x) (s) \right|& < \infty \quad m\text{-a.e.} \label{eq:XIII.2.10}
\end{align}
\end{theorem}

Suppose that, for a certain $z \in L^{1}(S, \mathfrak{B}, m)$, 
\begin{equation}\label{eq:XIII.2.11}
\lim\limits _{n \rightarrow \infty} n ^{-1} (T ^{n} z) (s) = 0 \quad m\text{-a.e.},
\end{equation}
\begin{equation}\label{eq:XIII.2.12}
\parbox{0.92\linewidth}{\everymath{}\let\limits\nolimits and the sequence $\left\{n^{-1}\sum\limits_{m=1}^{n}T^{m}z\right\}$ contains a subsequence which converges weakly to an element $\bar{z} \in L^{1}(S, \mathfrak{B}, m)$.}
\end{equation}


Then 
\begin{equation}\label{eq:XIII.2.13}
\slim _{n \to \infty} n ^{-1} \sum\limits_{m=1} ^{n} T ^{m} z = \bar {z}, T \bar {z} = \bar {z},
\end{equation}
and 
\begin{equation}\label{eq:XIII.2.14}
\lim\limits _{n \rightarrow \infty} n ^{-1} \sum\limits_{m=1} ^{n} (T ^{m} z) (s) = \bar {z} (s) m\text{-a.e.}
\end{equation}

\begin{proof}
 Consider the space $M(S, \mathfrak{B}, m)$, and take the $F$-lattice $X_1 = M(S, \mathfrak{B}; m)$ with $\|x\|_1 = \int_S |x(s)| (1 + |x(s)|)^{-1} m(\dd s)$, the $B$-lattice $X = L^1(S, \mathfrak{B}, m)$ with $\|x\| = \int_S |x(s)| m(\dd s)$ and $T_n = n^{-1} \sum\limits_{m=1}^{n} T^m$. Then the conditions of Theorem~\ref{thm:XIII.2.Theorem.1} are satisfied. We shall, for instance, verify that \eqref{eq:XIII.2.6} is satisfied. We have 
\[
\begin{gathered}
T _{n} z - T _{n} T ^{k} z = n ^{-1} (T + T ^{2} + \dots + T ^{n}) z - n ^{-1} (T ^{k+1} + T ^{k+2} + \dots + T ^{k+n}) z
\end{gathered}
\]
and so, by \eqref{eq:XIII.2.11}, $\lim\limits_{n\to \infty}(T_nz - T_nt^k z)(s) = 0$ $m$-a.e. for $k = 1,2,\ldots$. Hence, by taking the arithmetic mean with respect to $k$, we obtain $\lim\limits_{n\to \infty}(T_nz - T_nt_kz)(s) = 0$ $m$-a.e. for $k = 1,2,\ldots$. Conditions \eqref{eq:XIII.2.4} and \eqref{eq:XIII.2.5}, that is condition \eqref{eq:XIII.2.13} is a consequence of the mean ergodic theorem in Chapter~\ref{chap:VIII}, \ref{sec:VIII.3}. 
\end{proof}

\begin{remarku} The above two theorems are adapted from K. \textsc{Yosida} \cite{yosida358} and \cite{yosida361}. In these papers there are given other ergodic theorems implied by Theorem~\ref{thm:XIII.2.Theorem.1}. 
\end{remarku}

As for condition \eqref{eq:XIII.2.11} above, we have the following result due to E. \textsc{Hopf} \cite{hirsch111}: 

\begin{theorem}\label{thm:XIII.2.Theorem.3} Let $T$ be a positive linear operator on the real $L^1(S, \mathfrak{B}, m)$ into itself with the $L^1$-norm $\| T \| \leqq 1$. If $f \in L^1(S, \mathfrak{B}, m)$ and $p \in L^1(S, \mathfrak{B}, m)$ be such that $p(s) \geqq 0m$-a.e., then 

$\lim\limits_{n \to \infty} (T^n f)(s) / \sum\limits_{j=0}^{n-1} (T^j p)(s) = 0$ $m$-a.e. on the set where $p(s) > 0$. 
\end{theorem}

If $m(S) < \infty$ and $T \cdot 1 = 1$, then, by taking $p(s) = 1$, we obtain \eqref{eq:XIII.2.11}. \begin{proof}
It suffices to prove the Theorem for the case $f \geqq 0$. Choose $\varepsilon > 0$ arbitrarily and consider the functions 
\[
g _{n} = T ^{n} f - \varepsilon \cdot \sum\limits_{j=0} ^{n-1} T ^{j} p, g _{0} = f.
\]

Let $x_{n}(s)$ be the defining function of the set $\{s \in S; g_{n}(s) \geqq 0\}$. From $x_{n}g_{n} = g_{n}^{+} = \max(g, 0)$ and $g_{n+1} + \varepsilon p = Tg_{n}$, we obtain, by the positivity of the operator $T$ and $\|T\| \leqq 1$, 
\begin{gather}\int_{S} g _{n+1} ^{+} \cdot m (\dd s) + \varepsilon \int_{S} x _{n+1} p \cdot m (\dd s) = \int_{S} x _{n+1} \left(g _{n+1} + \varepsilon p\right) m (\dd s) \notag \\= \int_{S} x _{n+1} T g _{n} \cdot m (\dd s) \leqq \int_{S} x _{n+1} T g _{n} ^{+} \cdot m (\dd s) \leqq \int_{S} T g _{n} ^{+} \cdot m (\dd s) \leqq \int_{S} g _{n} ^{+} \cdot m (\dd s). \notag
\end{gather}

Summing up these inequalities from $n = 0$ on, we obtain 
\[
\int_{S} g _{n} ^{+} \cdot m (\dd s) + \varepsilon \int_{S} p \cdot \sum\limits_{k=1} ^{n} x _{k} \cdot m (\dd s) \leqq \int_{S} g _{0} ^{+} \cdot m (\dd s),
\]
and hence 
\[
\int_{S} p \cdot \sum\limits_{k=1} ^{\infty} x _{k} \cdot m (\dd s) \leqq \varepsilon^{-1} \int_{S} g _{0} ^{+} \cdot m (\dd s).
\]

Hence $\sum\limits_{k=1}^{\infty} x_k(s)$ converges $m$-a.e. on the set where $p(s) > 0$. Thus, on the set where $p(s) > 0$, we must have $g_n(s) < 0$ $m$-a.e. for all large $n$. 

Therefore, we have proved that $(T^n f)(s) < \varepsilon \sum\limits_{k=0}^{n-1} (T^k p)(s)$ $m$-a.e. for all large $n$ on the set where $p(s) > 0$. As $\varepsilon > 0$ was arbitrary, we have proved our Theorem. 
\end{proof}

As for condition \eqref{eq:XIII.2.10}, we have the following Theorem due to R. V. \textsc{Chacon}-D. S. \textsc{Ornstein} \cite{chacon32}: 

\begin{theorem}[R. V. \textsc{Chacon}-D. S. \textsc{Ornstein}]\label{thm:XIII.2.Theorem.4} Let $T$ be a positive linear operator on the real $L^{1}(S, \mathfrak{B}, m)$ into itself with the $L^{1}$-norm

$\| T \|_{1} \leqq 1$. If $f$ and $p$ are functions in $L^{1}(S, \mathfrak{B}, m)$ and $p(s) \geqq 0$, then 
\[
\varlimsup\limits_{n \rightarrow \infty} \left\{\sum\limits_{k=0} ^{n} (T ^{k} f) (s) / \sum\limits_{k=0} ^{n} (T ^{k} p) (s) \right\} \text{ is finite } m\text{-a.e.}
\]
on the set where $\sum\limits_{n=0}^{\infty}\left(T^{n}p\right)(s)>0$. 
\end{theorem}

If $m(S) < \infty$ and $T \cdot 1 = 1$, then, by taking $p(s) = 1$, we obtain \eqref{eq:XIII.2.10}. For the proof, we need a 

\begin{lemman}[\textsc{Chacon}-\textsc{Ornstein}\cite{chacon32}] If $f = f^{+} + f^{-}$, and if $\varlimsup\limits_{n\to \infty}\sum\limits_{k=0}^{n}(T^{k}f)(s) > 0$ on a set $B$, then there exist sequences $\{d_k\}$ and $\{f_k\}$ of non-negative functions such that, for every $N$,
\begin{align}
&\int_{S} \sum\limits_{k=0} ^{N} d _{k} \cdot m (\dd s) + \int_{S} f _{N} \cdot m (\dd s) \leqq \int_{S} f ^{+} \cdot m (\dd s), \label{eq:XIII.2.15} \\
\sum\limits_{k=0} ^{\infty} d _{k} (s)& = - f (s) \text{ on } B, \label{eq:XIII.2.16} \\
T ^{N} f ^{+}& = \sum\limits_{k=0} ^{N} T ^{N-k} d _{k} + f _{N}. \label{eq:XIII.2.17}
\end{align}

\end{lemman}
\begin{proof}
 Define inductively 
\begin{equation}d _{0} = 0, f _{0} = f ^{+}, f _{-1} = 0, f _{i+1} = \left(T f _{i} + f ^{-} + d _{0} + \dots + d _{i}\right) ^{+}, d _{i+1} \doteq T f _{i} - f _{i+1}.  \label{eq:XIII.2.18}\end{equation}

Note that 
\begin{equation}\label{eq:XIII.2.19}
+ f ^{-} + d _{0} + \dots + d _{i} \leqq 0,
\end{equation}
and that the equality holds on the set where $f_{i}(s) > 0$, for 
\[
\begin{aligned}
&{f _{i} = (T f _{i-1} + f ^{-} + d _{0} + \cdots + d _{i-1}) ^{+}} {\quad = (T f _{i-1} - f _{i} + f ^{-} + d _{0} + \cdots + d _{i-1} + f _{i}) ^{+}} \\
&{\quad = (d _{i} + f ^{-} + d _{0} + \cdots + d _{i-1} + f _{i}) ^{+}.}
\end{aligned}
\]

It follows from \eqref{eq:XIII.2.18} that 
\begin{equation}\label{eq:XIII.2.20}
T ^{j} f ^{+} = \sum\limits_{k=0} ^{j} T ^{j-k} d _{k} + f _{j}.
\end{equation}

By the definition, $f_{i}$ is non-negative, and so is $d_{i}$ by the last two equations of \eqref{eq:XIII.2.18} and \eqref{eq:XIII.2.19}. From \eqref{eq:XIII.2.20}, we have 
\begin{equation}\label{eq:XIII.2.21}
\sum\limits_{j=0} ^{n} T ^{j} f ^{+} = \sum\limits_{j=0} ^{n} \sum\limits_{k=0} ^{j} T ^{j-k} d _{k} + \sum\limits_{j=0} ^{n} f _{j}.
\end{equation}

We then prove 
\begin{equation}\label{eq:XIII.2.22}
\sum\limits_{j=0} ^{n} T ^{j} f ^{+} \leqq \sum\limits_{j=0} ^{n} d _{j} - \sum\limits_{j=1} ^{n} T ^{j} f ^{-} + \sum\limits_{j=0} ^{n} f _{j}.
\end{equation}

To this purpose, we note that 
\[
\sum\limits_{j=0} ^{n} \sum\limits_{k=0} ^{j} T ^{j-k} d _{k} = \sum\limits_{j=0} ^{n} T ^{j} \left(\sum\limits_{k=0} ^{n-j} d _{k}\right)
\]
and that, by \eqref{eq:XIII.2.19} and the positivity of $T$, 
\begin{equation}\label{eq:XIII.2.23}
- T ^{j} f ^{-} \geqq T ^{j} \left(\sum\limits_{k=0} ^{n-j} d _{k}\right) \quad \text{ for } \quad 1 \leqq j \leqq n.
\end{equation}

Rewriting \eqref{eq:XIII.2.22}, we have 
\begin{equation}\label{eq:XIII.2.24}
\sum\limits_{j=0} ^{n} T ^{j} (f ^{+} + f ^{-}) \leqq \sum\limits_{j=0} ^{n} (d _{j} + f _{j}) + f ^{-}.
\end{equation}

We will now prove that 
\begin{equation}\label{eq:XIII.2.25}
\sum\limits_{j=0} ^{\infty} d _{j} (s) + f ^{-} (s) \geqq 0 m \text{-a.e. on } B.
\end{equation}

It is clear from the remark after \eqref{eq:XIII.2.19} that \eqref{eq:XIII.2.25} holds with equality $m$-a.e. on the set $C = \{s \in S; f_j(s) > 0 \text{ for some } j \geqq 0\}$. It remains to show that \eqref{eq:XIII.2.25} holds on the set $B$ - $C$. This is proved by noting that \eqref{eq:XIII.2.24} implies that on $B$, and on $B$ - $C$ in particular, we have the inequality 
\[
\sum\limits_{j=0} ^{\infty} \left(d _{j} + f _{j}\right) (s) + f ^{-} (s) \geqq 0.
\]

Now we note that \eqref{eq:XIII.2.17} is exactly \eqref{eq:XIII.2.20} and that \eqref{eq:XIII.2.16} is implied from \eqref{eq:XIII.2.19} and \eqref{eq:XIII.2.25}. To see that \eqref{eq:XIII.2.15} holds, note that we have, by the assumptions on $T$ and \eqref{eq:XIII.2.18}, 
\[
\int_{S} \left(\sum\limits_{k=0} ^{j} d _{k} + f _{j}\right) m (\dd s) \geqq \int_{S} \left(\sum\limits_{k=0} ^{j} d _{k} + T \cdot f _{j}\right) m (\dd s) = \int_{S} \left(\sum\limits_{k=0} ^{j+1} d _{k} + f _{j+1}\right) m (\dd s).
\]

Hence we obtain \eqref{eq:XIII.2.15} by induction on $j$, since $d_{0} + f_{0} = f^{+}$. 
\end{proof}

\begin{proof}[Proof of Theorem~\ref{thm:XIII.2.Theorem.4}.]
It is sufficient to prove the Theorem under the hypothesis that $f(s) \geqq 0$ on $S$, and only to prove the finiteness of the indicated supremum when $p(s) > 0$. The first remark is obvious, and the second is proved as follows. By the hypothesis that the indicated supremum is finite at point $s$ where $p(s) > 0$, we have
\begin{equation}\label{eq:XIII.2.26}
\varlimsup\limits_{n \rightarrow \infty} \left\{\sum\limits_{j=0} ^{n} (T ^{j+k} f) (s) / \sum\limits_{j=0} ^{n} (T ^{j+k} p) (s) \right\} \text{ is finite } m\text{-a.e.}
\end{equation}
on the set where $(T^k p)(s) > 0$, 


and this implies that $\varlimsup\limits_{n\to \infty}\left\{\sum\limits_{j=0}^{n}(T^j f)(s)\Big/\sum\limits_{j=0}^{n}(T^j p)(s)\right\}$ is finite $m$-a.e. on the set where $(T^k p)(s) > 0$. 

Now assume the contrary to the hypothesis. Then 
\[
\varlimsup\limits_{n \rightarrow \infty} \left\{\sum\limits_{j=0} ^{n} (T ^{j} f) (s) / \sum\limits_{j=0} ^{n} (T ^{j} p) (s) \right\}
\]
is infinite $m$-a.e. on a set $E$ of positive $m$-measure, and $p(s) > \beta > 0$ on $E$ for some positive constant $\beta$. We have therefore, that for any positive constant $\alpha$, 
\[
\sup\limits _{n} \left\{\sum\limits_{j=0} ^{n} \left(T ^{j} ((f - \alpha p) ^{+} + (f - \alpha p) ^{-})\right) (s) \right\} > 0
\]
$m$-a.e. on $E$. Applying the Lemma with $f$ replaced by $(f - \alpha p)$, we obtain from \eqref{eq:XIII.2.15} and \eqref{eq:XIII.2.16} that 
\[
\int_{S} (f - \alpha p) ^{+} m (\dd s) \geqq \int_{S} \sum\limits_{k=0} ^{\infty} d _{k} m (\dd s) \geqq - \int_{E} (f - \alpha p) ^{-} m (\dd s).
\]

However, the extreme right term tends to $\infty$ as $\alpha\uparrow\infty$, and the extreme left term is bounded as $\alpha\uparrow\infty$. This is a contradiction, and so we have proved the Theorem. 
\end{proof}

We have thus proved 


\begin{theorem}\label{thm:XIII.2.Theorem.5} Let $T$ be a positive linear operator on $L^1(S, \mathfrak{B}, m)$ into itself with the $L^1$-norm $\| T \|_1 \leqq 1$. If $m(S) < \infty$ and $T \cdot 1 = 1$, then, for any $f \in L^1(S, \mathfrak{B}, m)$, the mean convergence of the sequence $\left\{ n^{-1} \sum\limits_{j=1}^{n} T^j f \right\}$ implies its $m$-a.e. convergence. 
\end{theorem}

As a corollary we obtain 

\begin{theorem}\label{thm:XIII.2.Theorem.6} Let $P(t, x, E)$ be a Markov process with an invariant measure $m$ on a measure space $(S, \mathfrak{B}, m)$ such that $m(S) < \infty$. Then, for the linear operator $T_t$ defined by $(T_t f)(x) = \int_S P(t, x, \dd y)f(y)$, we have
\begin{enumerate}[i)]
\item the mean ergodic theorem: 
\begin{equation}
\label{eq:XIII.2.27}
\begin{minipage}{0.92\linewidth}
\everymath{}\let\limits\nolimits
for any $f \in L ^{p} (S, \mathfrak {B}, m), \slim _{n \to \infty} n ^{-1} \sum\limits_{k=1} ^{n} T _{k} f = f ^{*}$ exists in $L^p(S, \mathfrak{B}, m)$ and $T_1 f^* = f^*   (p = 1, 2),$
\end{minipage}
\end{equation}


\item the individual ergodic theorem: 

\begin{equation}\label{eq:XIII.2.28}
\parbox{0.92\linewidth}{\everymath{}\let\limits\nolimits for any $f \in L^{p}(S, \mathfrak{B}, m)$ with $p = 1$ or $p = 2$, finite $\lim\limits_{n\to \infty}n^{-1}\sum\limits_{k=1}^{n}(T_kf)(s)$ exists $m$-a.e. and is $= f^{*}(s)$ $m$-a.e., and moreover, $\int_{S}f(s)m(ds) = \int_{S}f^{*}(s)m(ds)$.}
\end{equation}
\end{enumerate}
\end{theorem}

\begin{proof}
 By Theorem~\ref{thm:XIII.2.Theorem.2} in the preceding section, the mean ergodic theorem \eqref{eq:XIII.2.27} holds and so we may apply Theorem~\ref{thm:XIII.2.Theorem.5}. 
\end{proof}

\begin{remarku} If $T_{t}$ is given by an equi-measure transformation $x \to y_{t}(x)$ of $S$ onto $S$, the result \eqref{eq:XIII.2.27} is precisely the mean ergodic theorem of J. \textsc{von} \textsc{Neumann} \cite{neumann230} and the result \eqref{eq:XIII.2.28} is precisely the individual ergodic theorem of G. D. \textsc{Birkhoff} \cite{birkhoff18} and A. \textsc{Khintchine} \cite{khintchine149}. 
\end{remarku}

A Historical Sketch. The first operator-theoretical generalization of the individual ergodic theorem of the Birkhoff-Khintchine type was given by J. L. Doob \cite{doob40}. He proved that $n^{-1} \sum\limits_{k=1}^{n} (T_k f)(x)$ converges $m$-a.e. when $T_t$ is defined by a Markov process $P(t, x, E)$ with an invariant measure $m$ on a measure space $(S, \mathfrak{B}, m)$ such that $m(S) = 1$ and $f$ is the defining function of a set $\in \mathfrak{B}$. It was remarked by S. \textsc{Kakutani} \cite{kakutani130} that 

Doob's method is applicable and gives the same result for $f$ merely a bounded $\mathfrak{B}$-measurable function. E. \textsc{Hopf} \cite{hirsch110} then proved the theorem assuming merely that $f$ is $m$-integrable. N. \textsc{Dunford}-J. \textsc{Schwartz} \cite{dunford47} extended Hopf's result by proving \eqref{eq:XIII.2.28} for a linear operator $T_{t}$ which increases neither the $L^{1}$ - nor $L^{\infty}$-norm, without assuming that $T_{t}$ is positive but assuming that $T_{k} = T_{1}^{k}$ and $T_{1} \cdot 1 = 1$. It is noted that Hopf's and the Dunford-Schwartz arguments use the idea of our Theorem~\ref{thm:XIII.2.Theorem.1}. R. V. Chacon-D. S. \textsc{Ornstein} \cite{chacon32} proved \eqref{eq:XIII.2.28} for a positive linear operator $T_{t}$ of $L^{1}$-norm $\leqq 1$, without assuming that $T_{t}$ does not increase the $L^{\infty}$-norm, and, moreover, without appealing to our Theorem~\ref{thm:XIII.2.Theorem.1}. Here, of course, it is assumed that $T_{k} = T_{1}^{k}$ and $T_{1} \cdot 1 = 1$. We will not go into details, since, for the exposition of this Chapter which deals with Markov processes, it suffices to base the ergodic theory upon Theorem~\ref{thm:XIII.2.Theorem.6} which is a consequence of our Theorem~\ref{thm:XIII.2.Theorem.1}. 


\yossection{3}{The Ergodic Hypothesis and the \texorpdfstring{$H$}{H}-theorem}
\label{sec:XIII.3}

Let $P(t, x, E)$ be a Markov process with an invariant measure $m$ on a measure space $(S, \mathfrak{B}, m)$ such that $m(S) = 1$. We shall define the ergodicity of the process $P(t, x, E)$ by the condition: 
\begin{equation}\label{eq:XIII.3.1}
\begin{aligned}
&\text{ the time average }  f ^{*} (x) = \lim\limits _{n \to \infty} n ^{-1} \sum\limits_{k=1} ^{n} (T _{k} f) (x) = \lim\limits _{n \to \infty} n ^{-1} \sum\limits_{k=1} ^{n} \int_{S} P (k, x, \dd y) f (y) \\
&= \text{ the space average } \int_{S} f (x) m (\dd x) m\text{-a.e.}
\end{aligned}
\end{equation}
for every $f \in L^{p}(S, \mathfrak{B}, m) (\phi = 1, 2)$. 

Since $\int_{S} f^{*}(x) m(\dd x) = \int_{S} f(x) m(\dd x)$, \eqref{eq:XIII.3.1} may be rewritten as 
\begin{equation}\tag{1'}\label{eq:XIII.3.1'}
f ^{*} (x) = \text{ a constant } m \text{-a.e. for every } f \in L ^{p} (S, \mathfrak {B}, m) (p = 1, 2).
\end{equation}

We shall give three different interpretations of the ergodic hypothesis \eqref{eq:XIII.3.1} and $\eqref{eq:XIII.3.1}$. 

\begin{enumerate}[1.]
\item Let $\chi_B(x)$ be the defining function of the set $B \in \mathfrak{B}$. Then, for any two sets $B_1, B_2 \in \mathfrak{B}$, the time average of the probability that the points of $B_1$ will be transferred into the set $B_2$ after $k$ units of time is equal to the product $m(B_1)m(B_2)$. That is, we have 
\begin{equation}\label{eq:XIII.3.2}
\left(\chi_ {B _{1}} ^{*}, \chi_ {B _{1}}\right) = m \left(B _{1}\right) m \left(B _{2}\right).
\end{equation}

\begin{proof}
 If $f, g \in L^2(S, \mathfrak{B}, m)$, then the strong convergence $n^{-1} \sum\limits_{k=1}^{n} T_k f \to f^*$ in $L^2(S, \mathfrak{B}, m)$ implies, by \eqref{eq:XIII.3.1}, 
\[ \begin{aligned}\lim\limits _{n \rightarrow \infty} n ^{-1} \sum\limits_{k=1} ^{n} (T _{k} f, g)&= (f ^{*}, g) = f ^{*} (x) \int_{S} g (x) m (\dd x)\\&= \int_{S} f (x) m (\dd x) \cdot \int_{S} g (x) m (\dd x) \text{ for } m\text{-a.e.} x.\end{aligned}
\]

By taking $f(x) = \chi_{B_2}(x)$, $g(x) = \chi_{B_1}(x)$, we obtain \eqref{eq:XIII.3.2}. 
\end{proof}

\begin{remarku} Since linear combinations of the defining functions $\chi_{B}(x)$ are dense in the space $L^{p}(S, \mathfrak{B}, m)$ for $p = 1$ and $p = 2$, we easily see that \eqref{eq:XIII.3.2} is equivalent to the ergodic hypothesis \eqref{eq:XIII.3.1}. \eqref{eq:XIII.3.2} says that every part of $S$ will be transferred uniformly into every part of $S$ in the sense of the time average. 
\end{remarku}

\item The mean ergodic theorem in Chapter~\ref{chap:VIII}, \ref{sec:VIII.3} says that the mapping $f \to T^*f = f^*$ gives the eigenspace of $T_1$ belonging to the eigenvalue $1$ of $T_1$; this eigenspace is given by the range $R(T^*)$. Hence the ergodic hypothesis \eqref{eq:XIII.3.1'} means precisely the hypothesis that $R(T^*)$ is of one-dimension. Thus the ergodic hypothesis of $P(t,x,E)$ is interpreted in terms of the spectrum of the operator $T_1$. 

\item The Markov process $P(t, x, E)$ is said to be metrically transitive or indecomposable if the following conditions are satisfied: 

$S$ cannot be decomposed into the sum of disjoint sets 
\begin{align}
B _{1}, B _{2}& \in \mathfrak {B} \text{ such that } m (B _{1}) > 0, m (B _{2}) > 0 \notag \\
\text{ and } P (1, x, E)& = 0 \text{ for every } x \in B _{i}, E \subseteq B _{j} \text{ with } i \neq j. \label{eq:XIII.3.3}
\end{align}

\begin{proof}
 Suppose $P(t, x, E)$ is ergodic and that $S$ is decomposed as in \eqref{eq:XIII.3.3}. The defining function $\chi_{B_1}(x)$ satisfies $T_1\chi_{B_1} = \chi_{B_1}$ by the condition 
\[ \begin{aligned}
&\int_{S} P (1, x, \dd y) \chi_ {B _{1}} (y) = P (1, x, B _{1}) = 1 = \chi_ {B _{1}} (x),  \text{ for } x \in B _{1}, = 0 = \chi_ {B _{1}} (x),  \text{ for } x \bar {\in} B _{1}.
\end{aligned}
\]

But, by $m(B_1)m(B_2) > 0$, the function $\chi_{B_1}(x) = \chi_{B_1}^*(x)$ cannot be equal to a constant $m$-a.e. 


Next suppose that the process $P(t,x,E)$ is indecomposable. Let $T_{1}f=f$ and we shall show that $f^{*}(x)=f(x)$ equals a constant $m$-a.e. Since $T_{1}$ maps real-valued functions into real-valued functions, we may assume that the function $f(x)$ is real-valued. If $f(x)$ is not equal to a constant $m$-a.e., then there exists a constant $\alpha$ such that both the sets 
\[
B _{1} = \{s \in S; f (s) > \alpha \} \quad \text{ and } \quad B _{2} = \{s \in S; f (s) \leqq \alpha \}
\]
are of $m$-measure $> 0$. Since $T_{1}(f - \alpha) = f - \alpha$, the argument under The Angle Variable (p.~\pageref{sec:XIII.3}) implies that $T_{1}(f - \alpha)^{+} = (f - \alpha)^{+}$, $T_{1}(f - \alpha)^{-} = (f - \alpha)^{-}$. 

Hence we would obtain $P(1, x, E) = 0$ if $x \in B_i$, $E \subseteq B_j$ ( $i \neq j$).
\end{proof}

\begin{remarku} The notion of the metric transitivity was introduced by G. D. \textsc{Birkhoff}-P. A. \textsc{Smith} \cite{birkhoff19} for the case of an equi-measure transformation $x \rightarrow y_{t}(x)$ of $S$ onto $S$. 
\end{remarku}

An Example of the Ergodic Equi-measure Transformation. Let $S$ be a torus. That is, $S$ is the set of all pairs $s = \{x, y\}$ of real numbers $x$, $y$ such that $s = \{x, y\}$ and $s' = \{x', y'\}$ are identified iff $x \equiv x' \pmod{1}$ and $y \equiv y' \pmod{1}$; $S$ is topologized by the real number topology of its coordinates $x$, $y$. We consider a mapping 
\[
s = \{x, y \} \rightarrow T _{t} s = s _{t} = \{x + t \alpha , y + t \beta \}
\]
of $S$ onto $S$ which leaves invariant the measure $\dd x \, \dd y$ on $S$. Here we assume that the real numbers $\alpha, \beta$ are integrally linearly independent mod $1$ in the following sense: if integers $n, k$ satisfy $n\alpha + k\beta \equiv 0 \pmod{1}$, then $n = k = 0$. Then the mapping $s \to T_t^s$ is ergodic. 

\begin{proof}
 Let $f(s) \in L^2(S)$ be invariant by $T_1$, that is, let $f(s) = f(s_1) \dd x \dd y$ -a.e. on $S$. We have to show that $f(s) = f^*(s) = a \text{ constant } \dd x \dd y$ -a.e. on $S$. Let us consider the Fourier coefficients of $f(s)$ and $f(s_1) = f(T_1 s)$, respectively. 
\[
\begin{aligned}
&\int_{0} ^{1} \int_{0} ^{1} f (s) \exp (- 2 \pi i (k _{1} x + k _{2} y)) \dd x \dd y, \int_{0} ^{1} \int_{0} ^{1} f (T _{1} s) \exp (- 2 \pi i (k _{1} x + k _{2} y)) \dd x \dd y.
\end{aligned}
\]

By $T_{1}s = \{x + \alpha, y + \beta\}$ and the invariance of the measure $\dd x \, \dd y$ by the mapping $s \to T_{1}s$, the latter integral is equal to 
\[
\int_{0} ^{1} \int_{0} ^{1} f (s) \exp \left(- 2 \pi i (k _{1} x + k _{2} y)\right) \exp \left(2 \pi i (k _{1} \alpha + k _{2} \beta)\right) \dd x \dd y.
\]

Hence, by the uniqueness of the corresponding Fourier coefficients of $f(s)$ and $f(T_{1}s)$, we must have 
\[
\begin{gathered}
\exp (2 \pi i (k _{1} \alpha + k _{2} \beta)) = 1 \quad \text{ whenever } \int_{0} ^{1} \int_{0} ^{1} f (s) \exp (- 2 \pi i (k _{1} x + k _{2} y)) \dd x \dd y \neq 0.
\end{gathered}
\]

Hence, by the hypothesis concerning $\alpha$, $\beta$, we must have 
\[
\int_{0} ^{1} \int_{0} ^{1} f (s) \exp (- 2 \pi i (k _{1} x + k _{2} y)) \dd x \dd y = 0 \text{ unless } k _{1} = k _{2} = 0.
\]

Therefore $f(s) = f^{*}(s)$ must reduce to a constant $\dd x \dd y$-a.e. 
\end{proof}

The Angle Variable. Let $T_{t}$ be defined by the Markov process $P(t, x, E)$ with an invariant measure $m$ on $(S, \mathfrak{B}, m)$ such that $m(S) = 1$. Let $f(s) \in L^{2}(S, \mathfrak{B}, m)$ be the eigenvector of $T_{1}$ pertaining to the eigenvalue $\lambda$ of absolute value $1$ of $T_{1}$: 
\[
T _{1} f = \lambda f, | \lambda | = 1.
\]

Then we have $T_{1}|f| = |f|$. For, we have, by the positivity of the operator $T_{1}, (T_{1}|f|)(x) \geqq |(T_{1}f)(x)| = |f(x)|$, that is, $\int_{\xi} P(1, x, \dd y) |f(y)| \geqq$ 
$|f(x)|$. In this inequality, we must have the equality for $m$-a.e. $x$, as may be seen by integrating both sides with respect to $m(\dd x)$ and remembering the invariance of the measure $m$. Hence, if $P(t, x, E)$ is ergodic, then $|f(x)|$ must reduce to a constant $m$-a.e. Therefore, if we put 
\[
f (x) = | f (x) | \exp (i \Theta (x)), 0 \leqq \Theta (x) < 2 \pi ,
\]
we must have 
\[
T _{1} \exp (i \Theta (x)) = \lambda \exp (i \Theta (x)).
\]

In case $\lambda \neq 1$, such $\Theta(x)$ is called an angle variable of the ergodic Markov process $P(t, x, E)$. 

The Mixing Hypothesis. Let $T_{t}$ be defined by a Markov process $P(t, x, E)$ with an invariant measure $m$ on $(S, \mathfrak{B}, m)$ such that $m(S) = 1$. A stronger condition than the ergodic hypothesis of $P(t, x, E)$ is given by 
\begin{equation}\label{eq:XIII.3.4}
\lim\limits _{k \rightarrow \infty} (T _{k} f, g) = (f ^{*}, g) = \int_{S} f (x) m (\dd x) \cdot \int_{S} g (x) m (\dd x)
\end{equation}
for every pair $\{f, g\}$ of vectors $\in L^2(S, \mathfrak{B}, m)$. 


This condition is called the mixing hypothesis of the Markov process $P(t, x, E)$. As in the case of the ergodic hypothesis, it may be interpreted as follows: every part of $S$ will be transferred into every part of $S$ uniformly in the long run. As for the examples of the mixing equi-measure transformation $x \to y_t(x)$, we refer the reader to the above cited book by E. \textsc{Hopf}. 

$H$-theorem. Let $P(t, x, E)$ be a Markov process with an invariant measure $m$ on $(S, \mathfrak{B}, m)$ such that $m(S) = 1$. Let us consider the function 
\begin{equation}\label{eq:XIII.3.5}
H (z) = - z \log z, z \geqq 0.
\end{equation}

We can prove 

\begin{theoremn}[K. \textsc{Yosida} {\cite{yosida360}}] Let a non-negative function $f(x) \in L^{1}(S, \mathfrak{B}, m)$ belong to the Zygmund class, that is, let us assume that $\int_{S} f(x) \log^{+} f(x) m(\dd x) < \infty$, where $\log^{+}|z| = \log |z|$ or $= 0$ according as $|z| \geqq 1$ or $|z| < 1$. Then we have
\begin{equation}\label{eq:XIII.3.6}
\int_{S} H (f (x)) m (\dd x) \leqq \int_{S} H ((T _{t} f) (x)) m (\dd x).
\end{equation}

\end{theoremn}
\begin{proof}
 Since $H(z) = -z \log z$ satisfies $H''(z) = -1/z < 0$ for $z > 0$, $H(z)$ is a concave function. Hence we obtain 

the weighted mean of $H(f(x)) \leqq H$ (the weighted mean of $f(x)$), and so 
\[
\int_{S} P (t, x, \dd y) H (f (y)) \leqq H \left(\int_{S} P (t, x, \dd y) f (y)\right). \]

Integrating with respect to $m(\dd x)$ and remembering the invariance of the measure $m(E)$, we obtain \eqref{eq:XIII.3.6}. 
\end{proof}

\begin{remarku} By virtue of the semi-group property $T_{t+s} = T_t T_s$ and \eqref{eq:XIII.3.6}, we easily obtain 
\begin{equation}\tag{6'}\label{eq:XIII.3.6'}
\int_{S} H \left(\left(T _ {t _{1}} f\right) (x)\right) m (\dd x) \leqq \int_{S} H \left(\left(T _ {t _{1}} f\right) (x)\right) m (\dd x) \text{ whenever } t _{1} < t _{2}.
\end{equation}
\end{remarku}

This may be considered as an analogue of the classical $H$-theorem in statistical mechanics. 


\end{enumerate}

\yossection{4}{The Ergodic Decomposition of a Markov Process with a Locally Compact Phase Space}
\label{sec:XIII.4}

Let $S$ be a separable metric space whose bounded closed sets are compact. Let $B$ be the set of all Baire subsets of $S$, and consider a Markov process $P(t, x, E)$ on $(S, \mathfrak{B})$. We assume that 
\begin{equation}\label{eq:XIII.4.1}
f (x) \in C _{0} ^{0} (S) \text{ implies } f _{t} (x) = \int_{S} P (t, x, \dd y) f (y) \in C _{0} ^{0} (S).
\end{equation}

The purpose of this section is to give the decomposition of $S$ into ergodic parts and dissipative part as an extension of the Krylov-Bogolioubov case of the deterministic, reversible transition process in a compact metric space $S$ (see N. \textsc{Krylov}-N. \textsc{Bogolioubov} \cite{krein169}). The possibility of such an extension, to the case where $S$ is a compact metric space and with condition \eqref{eq:XIII.4.1}, was observed by K. \textsc{Yosida} \cite{yosida360} and carried out by N. \textsc{Beboutov} \cite{beboutov12}, independently of \textsc{Yosida}. The extension to the case of a locally compact space $S$ was given in K. \textsc{Yosida} \cite{yosida362}. We shall follow the last cited paper. 

\begin{lemma}\label{lem:XIII.4.Lemma.1} Let a linear functional $L(f)$ on the normed linear space $C_0^0(S)$, normed by the maximal norm $\|f\| = \sup\limits_{x \in S} |f(x)|$, be non-negative, viz., $L(f) \geqq 0$ when $f(x) \geqq 0$ on $S$. Then $L(f)$ is represented by a uniquely determined regular measure $\varphi(E)$, which is $\sigma$-additive and $\geqq 0$ for Baire subsets $E$ of $S$, as follows: 
\begin{equation}\label{eq:XIII.4.2}
L (f) = \int_{S} f (x) \varphi (\dd x) \text{ for all } f \in C _{0} ^{0} (S).
\end{equation}
\end{lemma}

Here the regularity of the measure $\varphi$ means that 
\begin{equation}\label{eq:XIII.4.3}
\varphi (E) = \inf \varphi (G) \text{ when } G \text{ ranges over all open sets } \supseteq E.
\end{equation}

\begin{proof}
 We refer the reader to P. R. \textsc{Halmos} \cite{halmos92}. 
\end{proof}

We shall call a non-negative, $\sigma$-additive regular measure $\varphi(E)$ defined on $\mathfrak{B}$ such that $\varphi(S) \leqq 1$ an invariant measure for the Markov process if 
\begin{equation}\label{eq:XIII.4.4}
\varphi (E) = \int_{S} \varphi (\dd x) P (t, x, E) \text{ for all } t > 0 \text{ and } E \in \mathfrak {B}.
\end{equation}

We have then 

\begin{lemma}\label{lem:XIII.4.Lemma.2} For any $f \in C_0^0(S)$ and for any invariant measure $\varphi(E)$, 
\begin{equation}\label{eq:XIII.4.5}
\lim\limits _{n \rightarrow \infty} n ^{-1} \sum\limits_{k=1} ^{n} f _{k} (x) = f ^{*} (x) \left(f _{k} (x) = \int_{S} P (k, x, \dd y) f (y)\right) \text{ exists }
\end{equation}
$\varphi$-a.e.and 
\begin{equation}\label{eq:XIII.4.6}
\int_{S} f ^{*} (x) \varphi (\dd x) = \int_{S} f (x) \varphi (\dd x).
\end{equation}
\end{lemma}

This is exactly a corollary of Theorem $6$ in the preceding section. 

Now, let $_1f, _2f, _3f, \ldots$ be dense in the normed linear space $C_0^0(S)$. The existence of such a sequence is guaranteed by the hypothesis concerning the space $S$. Applying Lemma~\ref{lem:XIII.4.Lemma.2} to $_1f, _2f, _3f, \ldots$ and summing up the exceptional sets of $\varphi$-measure zero, we see that there exists a set $N$ of $\varphi$-measure zero with the property: 

\begin{equation}
\label{eq:XIII.4.7}
\begin{minipage}{0.92\linewidth}
\everymath{}\let\limits\nolimits
for any $x \in N$ and for any $f \in C_0^0(S), \lim\limits_{n \to \infty} n^{-1} \sum\limits_{k=1}^{n} f_k(x) = f^*(x)$ exists.
\end{minipage}
\end{equation}

A Baire set $S' \subseteq S$ is said to be of maximal probability if $\varphi(S - S') = 0$ for every invariant measure $\varphi$. Thus the set $S'$ of all $x$ for which $f^*(x) = \lim\limits_{n \to \infty} n^{-1} \sum\limits_{k=1}^{n} f_k(x)$ exists for every $f \in C_0^0(S)$ is of maximal probability. Hence if there exists an invariant measure with $\varphi(S) > 0$, then there exist a $g \in C_0^0(S)$ and a point $x_0$ such that 
\begin{equation}\label{eq:XIII.4.8}
g ^{**} (x _{0}) = \varlimsup\limits_{n \rightarrow \infty} n ^{-1} \sum\limits_{k=1} ^{n} g _{k} (x _{0}) > 0.
\end{equation}

For, if otherwise, we would have $f^{*}(x) = 0$ on $S$ for all $f \in C_0^0(S)$ and hence, by \eqref{eq:XIII.4.6}, $\int_{S} f(x) \varphi(\dd x) = 0$. 

Let, conversely, \eqref{eq:XIII.4.8} be satisfied for a certain $g \in C_0^0(S)$ and for a certain $x_0$. Let a subsequence $\{n'\}$ of natural numbers be chosen in such a way that $\lim\limits_{n' \to \infty} (n')^{-1} \sum\limits_{k=1}^{n'} g_k(x_0) = g^{**}(x_0)$. By a diagonal method, we may choose a subsequence $\{n''\}$ of $\{n'\}$ such that $\lim\limits_{n'' \to \infty} (n'')^{-1} \sum\limits_{k=1}^{n''} (jf)_k(x_0)$ exists for $j = 1, 2, \ldots$. By the denseness of $\{jf\}$ in $C_0^0(S)$, we easily see that 
\[
\lim\limits _ {n ^{\prime \prime} \rightarrow \infty} (n ^{\prime \prime}) ^{-1} \sum\limits_{k=1} ^ {n ^{\prime \prime}} f _{k} (x _{0}) = f ^{***} (x _{0}) \text{ exists for every } f \in C _{0} ^{0} (S).
\]

If we put $f^{***}(x_0) = L_{x_0}(f)$, then 
\begin{equation}\label{eq:XIII.4.9}
L _ {x _{0}} (f) = L _ {x _{0}} \left(f _{1}\right).
\end{equation}

For, by condition \eqref{eq:XIII.4.1}, we have $f_{1} \in C_{0}^{0}(S)$ and so 
\[
L _ {x _{0}} (f) = \lim\limits _ {n ^{\prime \prime} \rightarrow \infty} (n ^{\prime \prime}) ^{-1} \sum\limits_{k=1} ^ {n ^{\prime \prime}} f _{k+1} (x _{0}) = f _{1} ^{***} (x _{0}) = L _ {x _{0}} (f _{1}).
\]

By Lemma~\ref{lem:XIII.4.Lemma.1}, there exists a regular measure $\varphi_{x_0}(E)$ such that 
\begin{equation}\label{eq:XIII.4.10}
L _ {x _{0}} (f) = f ^{***} \left(x _{0}\right) = \int_{S} f (x) \varphi_ {x _{0}} (\dd x).
\end{equation}

Surely we have 
\begin{equation}\label{eq:XIII.4.11}
0 \leqq \varphi_ {x _{0}} (E) \leqq 1,
\end{equation}
and, by \eqref{eq:XIII.4.9}, 
\[
\int_{S} f (x) \varphi_ {x _{0}} (\dd x) = \int_{S} \left(\int_{S} P (1, x, \dd y) f (y)\right) \varphi_ {x _{0}} (\dd x).
\]

By letting $f(x)$ tend to the defining function of the Baire set $E$, we see that $\varphi_{x_0}(E)$ is an invariant measure. We have $\varphi_{x_0}(S) > 0$, since, by \eqref{eq:XIII.4.8} and \eqref{eq:XIII.4.10}, we have $L_{x_0}(g) = g^{***}(x_0) = \int_S g(x) \varphi_{x_0}(\dd x) > 0$. We have thus proved 

\begin{theorem}\label{thm:XIII.4.Theorem.1} $A$ necessary and sufficient condition for the non-existence of non-trivial invariant measures is 
\begin{equation}\label{eq:XIII.4.12}
\lim\limits _{n \rightarrow \infty} n ^{-1} \sum\limits_{k=1} ^{n} f _{k} (x) = f ^{*} (x) = 0 \text{ on } S \text{ for any } f \in C _{0} ^{0} (S).
\end{equation}
\end{theorem}

\begin{definitionu} We will call the process $P(t, x, E)$ on $(S, \mathfrak{B})$ dissipative if condition \eqref{eq:XIII.4.12} is satisfied. 
\end{definitionu}

\begin{examplen} Let $S$ be the half line $(0, \infty)$ and $\mathfrak{B}$ the set of all Baire sets of $(0, \infty)$. Then the Markov process $P(t, x, E)$ defined by
\[
P (t, x, E) = 1 \text{ if } (x + t) \in E \text{ and } = 0 \text{ if } (x + t) \bar {\in} E
\]
is dissipative. 

We shall assume that $P(t, x, E)$ is not dissipative. Let $D$ denote the set 
\[
\left\{x \in S; f ^{*} (x) = \lim\limits _{n \rightarrow \infty} n ^{-1} \sum\limits_{k=1} ^{n} f _{k} (x) = 0 \quad \text{ for all } \quad f \in C _{0} ^{0} (S) \right\}
\]

Since it is equal to 
\[
\{x \in S; (j f) ^{*} (x) = 0 \quad \text{ for } \quad j = 1, 2, \dots \},
\]

$D$ is a Baire set. We shall call $D$ the dissipative part of $S$. We can prove below that $\varphi(D)=0$ for any invariant measure $\varphi$, and so, since the process $P(t,x,E)$ is assumed to be not dissipative, $S_{0}=S-D$ is not void. We already know that there exists a Baire set $S_{1}\subseteq S_{0}$ with the property: 

\begin{equation}\label{eq:XIII.4.13}
\parbox{0.92\linewidth}{\everymath{}\let\limits\nolimits to any $x \in S_1$, there corresponds a non-trivial invariant measure $\varphi_{x} (E)$ such that $f ^{*} (x) = \lim\limits_{n \to \infty} n ^{-1} \sum\limits_{k=1} ^{n} f _{k} (x) = \int_{S} f (z) \varphi_{x} (d z)$}
\end{equation}
and $\varphi(S_0 - S_1) = 0$ for any invariant measure $\varphi$. 


For any invariant measure $\varphi$, we have \eqref{eq:XIII.4.6} and hence, by \eqref{eq:XIII.4.13}, 
\[
\int_{S} f (y) \varphi (\dd y) = \int_ {S _{1}} \left(\int_{S} f (z) \varphi_{y} (\dd z)\right) \varphi (\dd y).
\]

Thus we obtain 
\begin{equation}\label{eq:XIII.4.14}
\varphi (E) = \int_ {S _{1}} \varphi_{y} (E) \varphi (\dd y),
\end{equation}
and so we have proved 

\end{examplen}
\begin{theorem}\label{thm:XIII.4.Theorem.2} Any invariant measure $\varphi$ may be obtained as a convex combination of the invariant measure $\varphi_y(E)$ with $y$ as a parameter. 
\end{theorem}

We see, from \eqref{eq:XIII.4.11} and \eqref{eq:XIII.4.14}, that the set 
\[
S _{2} = \{x \in S; x \in S _{1}, \varphi_{x} (S _{1}) = 1 \}
\]
is of maximal probability. Hence $S_0$ is of maximal probability. 

For any $f \in C_0^0(S)$ and for any invariant measure $\varphi$, we have 
\begin{equation}\label{eq:XIII.4.15}
\int_ {S _{3}} \varphi (\dd x) \left(\int_ {S _{3}} \left(f ^{*} (y) - f ^{*} (x)\right) ^{2} \varphi_{x} (\dd y)\right) = 0,
\end{equation}
since the left hand side is equal to 
\[
\begin{aligned}
&\int_ {S _{2}} \varphi (\dd x) \left(\int_ {S _{2}} f ^{*} (y) ^{2} \varphi_{x} (\dd y)\right)  - 2 \int_ {S _{2}} f ^{*} (x) \varphi (\dd x) \left(\int_ {S _{2}} f ^{*} (y) \varphi_{x} (\dd y)\right) \\
&+ \int_ {S _{2}} f ^{*} (x) ^{2} \varphi (\dd x) \cdot \int_ {S _{2}} \varphi_{x} (\dd y) = \int_ {S _{2}} f ^{*} (y) ^{2} \varphi (\dd y) - 2 \int_ {S _{2}} f ^{*} (x) ^{2} \varphi (\dd x) + \int_ {S _{2}} f ^{*} (x) ^{2} \varphi (\dd x),
\end{aligned}
\]
by ($13$), \eqref{eq:XIII.4.14}, \eqref{eq:XIII.4.6} and the definition of $S_{2}$. By applying \eqref{eq:XIII.4.15} to $_{1}f, _{2}f, _{3}f, \ldots$, we see that the set 
\[
S _{3} = \left\{x \in S _{2}; \int_ {S _{1}} (f ^{*} (y) - f ^{*} (x)) ^{2} \varphi_{x} (\dd y) = 0 \text{ for all } f \in C _{0} ^{0} (S) \right\}
\]
is of maximal probability. 

We shall give the ergodic decomposition of $S$. For any $x \in S_3$, put 
\begin{equation}\label{eq:XIII.4.16}
E _{x} = \{y \in S _{3}; f ^{*} (y) = f ^{*} (x) \quad \text{ for all } \quad f \in C _{0} ^{0} (S) \}.
\end{equation}

We can then show that each $E_x$ contains a set $\hat{E}_x$ with the property: 
\begin{equation}\label{eq:XIII.4.17}
\varphi_{x} (E _{x}) = \varphi_{x} (\hat {E} _{x}) \text{ and } P (1, y, \hat {E} _{x}) = 1 \text{ for any } y \in \hat {E} _{x}.
\end{equation}

\begin{proof}
 By the definition of $S_{3}$, we see that $f^{*}(y) = f^{*}(x)$ if the measure $\varphi_{x}(E)$ has variation at the point $y$. Hence $\varphi_{x}(E_{x}) = \varphi_{x}(S_{3}) = 1$. Thus, by the invariance of the measure $\varphi_{x}$, we obtain 
\[
1 = \varphi_{x} (E _{x}) = \int_ {S _{3}} P (1, z, E _{x}) \varphi_{x} (\dd z) = \int_ {E _{x}} P (1, z, E _{x}) \varphi_{x} (\dd z). \]

Since $0 \leqq P(1, z, E_x) \leqq 1$, there exists a Baire set $E^1 \subseteq E_x$ such that 
$\varphi_{x}(E^{1}) = \varphi_{x}(E_{x})$ and $z\in E^{1}$ implies $P(1,z,E_x) = 1$ 

Put 
\[
E ^{2} = \{z \in E ^{1}, P (1, z, E ^{1}) = 1 \}.
\]

Since 
\[
\int_ {E ^{1}} P (1, z, E ^{1}) \varphi_{x} (\dd z) = \varphi_{x} (E ^{1}) = \varphi_{x} (E _{x}) = \int_ {E ^{1}} P (1, z, E _{x}) \varphi_{x} (\dd z),
\]
we must have 
\[
\int_ {E ^{1}} \left(P (1, z, E _{x}) - P (1, z, E ^{1})\right) \varphi_{x} (\dd z) = 0.
\]

As $z \in E^2$ implies $P(1, z, E^1) = 1$ by the definition of $E^2$, we obtain $\varphi_x(E^1 - E^2) = 0$. 

Next put 
\[
E ^{3} = \{z \in E ^{2}; P (1, z, E ^{2}) = 1 \}.
\]

Then, as above, $\varphi_{x}(E^{3}) = \varphi_{x}(E_{x})$. In this way, we obtain a sequence $\{E^n\}$ such that 
\begin{align}
E _{x}& \supseteq E ^{1} \supseteq E ^{2} \supseteq \dots , \varphi_{x} (E ^{n}) = \varphi_{x} (E _{x}) \text{ and } \notag \\
P (1, z, E ^{n})& = 1 \text{ if } z \in E ^{n+k} (n \geqq 0, k \geqq 1, E ^{0} = E _{x}). \notag
\end{align}

We have thus obtained the set $\hat{E}_x = \bigcap\limits_{n=1}^{\infty} E^n$ with the demanded property. 
\end{proof}
Therefore, we have 

\begin{theorem}\label{thm:XIII.4.Theorem.3} $P(t, y, E)$ defines a Markov process in each $\hat{E}_x$ which is ergodic, in the sense that 

\begin{equation}\label{eq:XIII.4.18}
\parbox{0.92\linewidth}{\everymath{}\let\limits\nolimits $\hat{E}_x$ is not decomposable in two parts $A$ and $B$ such that $\varphi(A) \cdot \varphi(B) > 0$ for a certain invariant measure $\varphi$ and $P(1, a, B) = 0$ when $a \in A$ and $P(1, b, A) = 0$ when $b \in B$.}
\end{equation} 
\end{theorem}


\begin{proof}
 Let $\varphi$ be any invariant measure in $\hat{E}_x$ with $\varphi(\hat{E}_x) = 1$. Then, by \eqref{eq:XIII.4.14}, $\varphi(E) = \int_{\hat{E}_x} \varphi(\dd z) \varphi_z(E)$ whenever $E \subseteq \hat{E}_x$. Since $\varphi_z(E) = \varphi_x(E)$ for every $z \in \hat{E}_x$ by the definition of $\hat{E}_x$ we have 
\[
\varphi (E) = \varphi_{x} (E) \int_ {\hat {E} x} \varphi (\dd z) = \varphi_{x} (E). \]

Thus there is essentially only one invariant measure $\varphi_{x}(E)$ in $\hat{E}_{x}$. This proves the ergodicity. For, let $\hat{E}_{x}$ be decomposed as in ($18$). Then, by the invariance of $\varphi$, 
\[
\varphi (C) = \int_ {\hat {E} _{z}} P (1, z, C) \varphi (\dd z) \quad \text{ whenever } \quad C \subseteq \hat {E} _{x}.
\]

Thus the measure $\psi$ defined by 
\[
\begin{aligned} \psi (C) & = \varphi (C) / \varphi (A) \quad \text{ if } \quad C \subseteq A \\ & = 0 \text{ if } C \subseteq B \end{aligned}
\]
is invariant in $\hat{E}_x$ and differs from the unique invariant measure $\varphi$. 
\end{proof}


\yossection{5}{The Brownian Motion on a Homogeneous Riemannian Space}
\label{sec:XIII.5}

There is an interesting interplay between Markov processes and differential equations. Already in the early thirties, A. \textsc{Kolmogorov} proved, under a certain regularity hypothesis concerning the transition probability $P(t, x, E)$, that 
\[
u (t, x) = \int_{S} P (t, x, \dd y) f (y)
\]
satisfies the equation of the diffusion type: 
\begin{equation}\label{eq:XIII.5.1}
\frac {\partial u}{\partial t} = b ^{ij} (x) \frac {\partial^{2} u}{\partial x _{i} \partial x _{j}} + a ^{i} (x) \frac {\partial u}{\partial x _{i}} = A u, t > 0,
\end{equation}
where the differential operator $A$ is elliptic in the local coordinates $(x_1, x_2, \ldots, x_n)$ of the point $x$ of the phase space $S$. Following Einstein's convention of the tensor notation, it is understood that, e.g., $a^i(x) \partial/\partial x_i$ means $\sum\limits_{i=1}^{n} a^i(x) \partial/\partial x_i$. For the derivation of equation \eqref{eq:XIII.5.1}, see A. \textsc{Kolmogorov} \cite{kolmogorov152}. The leit-motif in his research was the investigation of local characteristics of Markov processes. 

Following work done by E. \textsc{Hille} at Scandinavian Congress in $1949$ (Cf. E. \textsc{Hille} \cite{hille105}) and K. \textsc{Yosida} \cite{yosida372} in $1948$, W. \textsc{Feller} \cite{erdelyi57} began in $1952$ a systematic research of this new field in probability theory using the analytical theory of semi-groups. His investigations were further developed by E. B. \textsc{Dynkin} \cite{dunford50}, \cite{dunford51}, K. \textsc{Ito}-H. McKean \cite{hunt120}, D. \textsc{Ray}, \cite{ray270} G. A. \textsc{Hunt} \cite{hunt119}, A. A. \textsc{Yushkevitch} \cite{yushkevitch384}, G. \textsc{Maruyama} \cite{maruyama189} and many younger scholars, especially in Japan, USSR and USA. These investigations are called the theory of diffusion. A comprehensive treatise on the diffusion theory by K. \textsc{Ito}-H. McKean \cite{hunt120} is in the course of printing\footnote{Berlin/G\"{o}ttingen/Heidelberg: Springer.} We shall sketch some of the salient analytical features of the theory. 

Let $S$ be a locally compact space such that $B$ is the totality of the Baire sets in $S$. To define the spatial homogeneity of a Markov process $P(t, x, E)$ on the space $S$, we suppose that $S$ is an $n$-dimensional, orientable connected $C^{\infty}$ Riemannian space such that the full group $G$ of the isometries of $S$, which is a Lie group, is transitive on $S$. That is, for each pair $\{x, y\}$ of points $\in S$, there exists an isometry $M \in G$ such that $M \cdot x = y$. Then the process $P(t, x, E)$ is spatially homogeneous if 
\begin{equation}\label{eq:XIII.5.2}
P (t, x, E) = P (t, M \cdot x, M \cdot E) \text{ for each } x \in S, E \in \mathfrak {B} \text{ and } M \in \mathfrak {G}.
\end{equation}

A temporally and spatially homogeneous Markov process on $S$ is called a Brownian motion on $S$, if the following condition, known as the continuity condition of Lindeberg's type, is satisfied: 
\begin{equation}\label{eq:XIII.5.3}
\lim\limits _{t \downarrow 0} t ^{-1} \int_ {\dist (x, y) > \varepsilon} P (t, x, \dd y) = 0 \quad \text{ for every } \quad \varepsilon > 0 \text{ and } x \in S.
\end{equation}

\begin{propositionu} Let $C(S)$ be the $B$-space of all bounded uniformly continuous real-valued functions $f(x)$ on $S$ normed by $\|f\| = \sup |f(x)|$. 
\end{propositionu}

Define 
\begin{equation}\label{eq:XIII.5.4}
\begin{aligned} (T _{t} f) (x) & = \int_{S} P (t, x, \dd y) f (y), \quad \text{ if } \quad t > 0, \\ & = f (x), \text{ if } t = 0. \end{aligned}
\end{equation}

Then $\{T_t\}$ constitutes a contraction semi-group of class $(C_0)$ in $C(S)$. 

\begin{proof}
 By $P(t, x, E) \geqq 0$ and $P(t, x, S) = 1$, we obtain 
\[
\left| \left(T _{t} f\right) (x) \right| \leqq \sup\limits _{y} | f (y) |
\]
and the positivity of the operator $T_{t}$. The semi-group property $T_{t+s} = T_{t}T_{s}$ is a consequence from the Chapman-Kolmogorov equation. If we define a linear operator $M'$ by $(M'f)(x) = f(M\cdot x), M \in \mathfrak{G}$, we obtain 
\[
T _{t} M ^{\prime} = M ^{\prime} T _{t}, t \geqq 0. \]

For, 
\begin{equation}\label{eq:XIII.5.5}
\begin{aligned}
&(M ^{\prime} T _{t} f) (x)  = (T _{t} f) (M \cdot x) = \int_{S} P (t, M \cdot x, \dd y) f (y) = \int_{S} P (t, M \cdot x, d (M \cdot y)) f (M \cdot y) \\
&= \int_{S} P (t, x, \dd y) f (M \cdot y) = (T _{t} M ^{\prime} f) (x).
\end{aligned}
\end{equation}

If $M \in \mathfrak{G}$ be such that $M \cdot x = x'$, we have 
\[
\left(T _{t} f\right) (x) - \left(T _{t} f\right) \left(x ^{\prime}\right) = \left(T _{t} f\right) (x) - \left(M ^{\prime} T _{t} f\right) (x) = T _{t} \left(f - M ^{\prime} f\right) (x).
\]

Hence by the uniform continuity of $f(x)$, we see that $(T_{t}f)(x)$ is uniformly continuous and bounded. 


To prove the strong continuity in $t$ of $T_{t}$, it suffices by the Theorem in Chapter~\ref{chap:IX}, \ref{sec:IX.1}, to verify weak right continuity of $T_{t}$ at $t = 0$. Therefore, it is surely enough to show that $\lim\limits_{t\downarrow0}(T_{t}f)(x)=f(x)$ uniformly in $x$. Now 
\[
\begin{aligned} & {\left| \left(T _{t} f\right) (x) - f (x) \right| = \left| \int_{S} P (t, x, \dd y) [ f (y) - f (x) ] \right|} \\ & {\leqq \left| \int_{d (x, y) \leqq \varepsilon} P (t, x, \dd y) [ f (y) - f (x) ] \right| + \left| \int_{d (x, y) > \varepsilon} P (t, x, \dd y) [ f (y) - f (x) ] \right|} \\ & {\leqq \left| \int_{d (x, y) \leqq \varepsilon} P (t, x, \dd y) [ f (y) - f (x) ] \right| + 2 \| f \| \int_{d (x, y) > \varepsilon} P (t, x, \dd y).} \end{aligned}
\]

The first term on the right tends to zero uniformly in $x$ as $\varepsilon \downarrow 0$ and, for fixed $\varepsilon > 0$, the second term on the right tends to zero uniformly in $x$ as $t \downarrow 0$. The latter fact is implied by \eqref{eq:XIII.5.3} and the spatial homogeneity. Thus $\lim\limits_{t \downarrow 0} (T_t f)(x) = f(x)$ uniformly in $x$.
\end{proof}

\begin{theoremu} Let $x_0$ be a fixed point of $S$. Let us assume that the group of isotropy $\mathfrak{G}_0 = \{M \in G; M \cdot x_0 = x_0\}$ is compact. Being a closed subgroup of a Lie group $\mathfrak{G}, \mathfrak{G}_0$ is a Lie group by a theorem due to E. \textsc{Cartan}. Let $A$ be the infinitesimal generator of the semi-group $T_t$. Then we have the following results. 


\begin{enumerate}[{[}1{]}]
\item If $f \in D(A) \cap C^2(S)$, then, for a coordinate system $(x_0^1, x_0^2, \ldots, x_0^n)$ at $x_0$, 
\begin{equation}\label{eq:XIII.5.6}
(A f) (x _{0}) = a ^{i} (x _{0}) \frac {\partial f}{\partial x _{0} ^{i}} + b ^{ij} (x _{0}) \frac {\partial^{2} f}{\partial x _{0} ^{i} \partial x _{0} ^{j}},
\end{equation}
where 
\begin{align}
a ^{i} \left(x _{0}\right)& = \lim\limits _{t \downarrow 0} t ^{-1} \int_ {d \left(x _{0}, x _{t}\right) \leqq \varepsilon} \left(x ^{i} - x _{0} ^{i}\right) P (t, x _{0}, \dd x), \label{eq:XIII.5.7} \\
b ^{ij} (x _{0})& = \lim\limits _{t \downarrow 0} t ^{-1} \int_ {d (x _{0}, x) \leqq \varepsilon} \left(x ^{i} - x _{0} ^{i}\right) \left(x ^{j} - x _{0} ^{j}\right) P (t, x _{0}, \dd x), \label{eq:XIII.5.8}
\end{align}
the limits existing independently of sufficiently small $\varepsilon > 0$. 

\item The set $D(A) \cap C^2(S)$ is ``big'' in the sense that, for any $C^\infty(S)$ function $h(x)$ with compact support there exists an $f(x) \in D(A) \cap C^2(S)$ such that $f(x_0), \partial f / \partial x_0^i, \partial^2 f / \partial x_0^i \partial x_0^j$ are arbitrarily near to $h(x_0), \partial h / \partial x_0^i, \partial^2 h / \partial x_0^i \partial x_0^j$, respectively. 

\end{enumerate}
\end{theoremu}

\begin{proof}
 The first step. Let $h(x)$ be a $C^\infty$ function with compact support. If $f \in D(A)$, then the ``convolution'' 
\begin{equation}\label{eq:XIII.5.9}
(f \otimes h) (x) = \int_ {\mathfrak {G}} f (M _{y} \cdot x) h (M _{y} \cdot x _{0}) \dd y
\end{equation}
$(M_{y}$ denotes a generic element of $G$ and $\dd y$ a fixed right invariant Haar measure on $G$ such that $\dd y = d(y \cdot M)$ for every $M \in \mathfrak{G}$) is $C^{\infty}$ and belongs to $D(A)$. The above integral exists since the isotropy group $\mathfrak{G}_{0}$ is compact and $h$ has a compact support. By the uniform continuity of $f$ and the compactness of the support of $h$ we can approximate the integral by Riemann sums $\sum\limits_{i=1}^{k} f(M_{y_{i}} \cdot x) C_{i}$ uniformly in $x$: 
\[
(f \otimes h) (x) = \slim _{k \rightarrow \infty} \sum\limits_{i=1} ^{k} f \left(M _ {y _{i}} \cdot x\right) C _{i}. \]

Since $T_{t}M' = M'T_{t}$, $M'$ commutes with $A$, i.e., if $f \in D(A)$, then $M'f \in D(A)$ and $AM'f = M'Af$. Putting $g(x) = (Af)(x)$, we obtain $g \in C(S)$ and 
\[
\begin{aligned}
&A \left(\sum\limits_{i=1} ^{k} f (M _ {y _{i}} \cdot x) C _{i}\right)  = \sum\limits_{i=1} ^{k} (A M _ {y _{i}} ^{\prime} \cdot f) (x) C _{i} = \sum\limits_{i=1} ^{k} (M _ {y _{i}} ^{\prime} A f) (x) C _{i} = \sum\limits_{i=1} ^{k} g (M _ {y _{i}} \cdot x) C _{i},
\end{aligned}
\]
and the right hand side tends to $(g\otimes h)(x) = (Af\otimes h)(x)$. Since the infinitesimal generator $A$ is closed, it follows that $f\otimes h\in D(A)$ and $A(f\otimes h) = Af\otimes h$. Since $S$ is a homogeneous space of the Lie group $\mathfrak{G}$, i.e., $S = \mathfrak{G} / \mathfrak{G}_0$, we can find a coordinate neighbourhood $U$ of $x_0$ such that for each $x\in U$ there exists an element $M = M(x)\in \mathfrak{G}$ satisfying the conditions: 

\begin{enumerate}[i)]
\item $M(x)\cdot x = x_0,$
\item $M(x) x_0$ depends analytically on the coordinates $(x^1, x^2, \ldots, x^n)$ of the point $x \in S$.
\end{enumerate}


This is so, since the set $\{M_y \in \mathfrak{G}; M_y \cdot x = x_0\}$ forms an analytic submanifold of $\mathfrak{G}$; it is one of the cosets of $\mathfrak{G}$ with respect to the Lie subgroup $\mathfrak{G}_0$. Hence, by the right invariance of $\dd y$, we have 
\[
\begin{gathered}
(f \otimes h) (x) = \int_{G} f (M _{y} M (x) x) h (M _{y} M (x) x _{0}) \dd y = \int_{G} f (M _{y} \cdot x _{0}) h (M _{y} M (x) \cdot x _{0}) \dd y.
\end{gathered}
\]

The right side is infinitely differentiable in the vicinity of $x_{0}$, and 
\begin{equation}\label{eq:XIII.5.10}
D _{x} ^{s} (f \otimes h) (x) = \int_{G} f \left(M _{y} \cdot x _{0}\right) D _{x} ^{s} h \left(M _{y} M (x) \cdot x _{0}\right) \dd y.
\end{equation}

The second step. Remembering that $D(A)$ is dense in $C(S)$ and choosing $f$ and $h$ appropriately, we obtain: 

\begin{equation}\label{eq:XIII.5.11}
\parbox{0.92\linewidth}{\everymath{}\let\limits\nolimits there exist $C^\infty$ functions $F^1(x), F^2(x), \ldots, F^n(x) \in D(A)$ such that the Jacobian $\frac{\partial(F^1(x),\ldots,F^n(x))}{\partial(x^1,\ldots,x^n)} >0$ at $x_0$, and, moreover there exists a $C^\infty$ function $F_0(x) \in D(A)$ such that}
\end{equation}
\begin{equation}\label{eq:XIII.5.12}
\left(x ^{i} - x _{0} ^{i}\right) \left(x ^{j} - x _{0} ^{j}\right) \frac {\partial^{2} F _{0}}{\partial x _{0} ^{i} \partial x _{0} ^{j}} \geqq \sum\limits_{i=1} ^{n} \left(x ^{i} - x _{0} ^{i}\right) ^{2}.
\end{equation}

We can use $F^1(x), F^2(x), \ldots, F^n(x)$ as coordinate functions in a neighbourhood $\{x; \text{ dis }(x_0, x) < \varepsilon\}$ of the point $x_0$. We denote these new local coordinates by $(x^1, x^2, \ldots, x^n)$. Since $F^j(x) \in D(A)$, we have 
\[
\begin{aligned}
\underset {t \downarrow 0} {\lim} t ^{-1} \int_{S} P (t, x _{0}, \dd x) (F ^{j} (x) - F ^{j} (x _{0})) &= (A F ^{j}) (x _{0})
\\
&= \underset {t \downarrow 0} {\lim} t ^{-1} \int_ {\mathrm{dis} (x, x _{0}) \leqq \varepsilon} P (t, x _{0}, \dd x) (F ^{j} (x) - F ^{j} (x _{0}))
\end{aligned}
\]
independently of sufficiently small $\varepsilon > 0$, by virtue of Lindeberg's condition \eqref{eq:XIII.5.3}. Hence there exists, for the coordinate function $x^1, x^2, \ldots, x^n$ ( $x^j = F^j$), a finite 
\begin{equation}\label{eq:XIII.5.13}
\lim\limits _{t \downarrow 0} t ^{-1} \int_ {\dist (x, x _{0}) \leqq \varepsilon} \left(x ^{j} - x _{0} ^{j}\right) P (t, x _{0}, \dd x) = a ^{j} \left(x _{0}\right)
\end{equation}
which is independent of sufficiently small $\varepsilon > 0$. Since $F_0 \in D(A)$, we have, again using Lindeberg's condition \eqref{eq:XIII.5.3}, 
\[
\begin{aligned}
(A F _{0}) (x) &= \lim\limits _{t \downarrow 0} t ^{-1} \int_{S} P (t, x _{0}, \dd x) (F _{0} (x) - F _{0} (x _{0}))
\\
&= \lim\limits _{t \downarrow 0} t ^{-1} \int_ {\mathrm{dis} (x _{0}, x) \leqq \varepsilon} P (t, x _{0}, \dd x) (F _{0} (x) - F _{0} (x _{0}))
\\
&= \lim\limits _{t \downarrow 0} \left[ t ^{-1} \int_ {\mathrm{dis} (x _{0}, x) \leqq \varepsilon} (x ^{i} - x _{0} ^{i}) \frac {\partial F _{0}}{\partial x _{0} ^{i}} P (t, x _{0}, \dd x) \right.
\\
& \qquad + \left. t ^{-1} \int_ {\mathrm{dis} (x _{0}, x) \leqq \varepsilon} (x ^{i} - x _{0} ^{i}) (x ^{j} - x _{0} ^{j}) \left(\frac {\partial^{2} F _{0}}{\partial x ^{i} \partial x ^{j}}\right) _ {x = x _{0} + \theta (x - x _{0})} P (t, x _{0}, \dd x) \right],
\end{aligned}
\]
where $0 < \theta < 1$. The first term on the right has a limit $a^i(x_0) \frac{\partial F_0}{\partial x_0^i}$. Hence, by the positivity of $P(t, x, E)$ and \eqref{eq:XIII.5.12}, we see that 
\begin{equation}\label{eq:XIII.5.14}
\varlimsup\limits_{t \downarrow 0} t ^{-1} \int_ {\dist (x _{0}, x) < \varepsilon} \sum\limits_{i=1} ^{n} (x ^{i} - x _{0} ^{i}) ^{2} P (t, x _{0}, \dd x) < \infty .
\end{equation}

The third step. Let $f \in D(A) \cap C^2$. Then, by expanding $f(x) - f(x_0)$, we obtain 
\begin{equation}\label{eq:XIII.5.15}
\begin{aligned}
&\frac {(T _{t} f) (x _{0}) - f (x _{0})}{t} = t ^{-1} \int_{S} (f (x) - f (x _{0})) P (t, x _{0}, \dd x)
\\
&= t ^{-1} \int_ {\mathrm{dis} (x, x _{0}) > \varepsilon} (f (x) - f (x _{0})) P (t, x _{0}, \dd x) + t ^{-1} \int_ {\mathrm{dis} (x, x _{0}) \leqq \varepsilon} (x ^{i} - x _{0} ^{i}) \frac {\partial f}{\partial x _{0} ^{i}} P (t, x _{0}, \dd x)
\\
&\quad + t ^{-1} \int_ {\mathrm{dis} (x, x _{0}) \leqq \varepsilon} (x ^{i} - x _{0} ^{i}) (x ^{j} - x _{0} ^{j}) \frac {\partial^{2} f}{\partial x _{0} ^{i} \partial x _{0} ^{j}} P (t, x _{0}, \dd x)
\\
&\quad + t ^{-1} \int_ {\mathrm{dis} (x, x _{0}) \leqq \varepsilon} (x ^{i} - x _{0} ^{i}) (x ^{j} - x _{0} ^{j}) C _{ij} (\varepsilon) P (t, x _{0}, \dd x)
\\
&= C _{1} (t, \varepsilon) + C _{2} (t, \varepsilon) + C _{3} (t, \varepsilon) + C _{4} (t, \varepsilon). \text{ say },
\end{aligned}
\end{equation}
where $C_{ij}(\varepsilon) \to 0$ as $\varepsilon \downarrow 0$. We know, by \eqref{eq:XIII.5.3}, that $\lim\limits_{t \downarrow 0} C_1(t, \varepsilon) = 0$ for fixed $\varepsilon > 0$, and that $\lim\limits_{t \downarrow 0} C_2(t, \varepsilon) = a^i(x_0) \frac{\partial f}{\partial x_0^i}$, independently of sufficiently small $\varepsilon > 0$. By \eqref{eq:XIII.5.14} and Schwarz' inequality, we see that $\lim\limits_{\varepsilon \downarrow 0} C_4(t, \varepsilon) = 0$, boundedly in $t > 0$. The left side also has a finite limit $(Af)(x_0)$ as $t \downarrow 0$. So the difference 
\[
\varlimsup\limits_{t \downarrow 0} C _{3} (t, \varepsilon) - \varliminf\limits_{t \downarrow 0} C _{3} (t, \varepsilon)
\]
can be made arbitrarily small by taking $\varepsilon > 0$ sufficiently small. But, by \eqref{eq:XIII.5.14}, Schwarz' inequality and \eqref{eq:XIII.5.3}, the difference is independent of sufficiently small $\varepsilon > 0$. Thus a finite limit $\lim\limits_{t \downarrow 0} C_3(t, \varepsilon)$ exists independently of sufficiently small $\varepsilon > 0$. Since we may choose $F \in D(A) \cap C^\infty(S)$ in such a way that $\partial^2 F / \partial x_0^i \partial x_0^j$ ( $i, j = 1, 2, \ldots, n$) is arbitrarily near $\alpha_{ij}$, where $\alpha_{ij}$ are arbitrarily given constants, it follows, by an argument similar to the above, that 
a finite $\lim\limits_{t\downarrow 0}t^{-1}\int_{\mathrm{dis}(x,x_0)\leqq s}(x^i -x_0^i)(x^j -x_0^j)P(t,x_0,\dd x) = b^{ij}(x_0)$ exists and 
\begin{equation}\label{eq:XIII.5.16}
\lim\limits _{t \downarrow 0} C _{3} (t, \varepsilon) = b ^{ij} \left(x _{0}\right) \frac {\partial^{2} f}{\partial x _{0} ^{i} \partial x _{0} ^{j}}.
\end{equation}

This completes the proof of our Theorem.
\end{proof}

\begin{remarku} The above Theorem and proof are adapted from K. \textsc{Yosida} \cite{yosida363}. It is to be noted that $b^{ij}(x) = b^{ji}(x)$ and 
$b^{ij}(x_0)\xi_i\xi_j\geqq 0$ for every real vector $(\xi_1,\xi_2,\ldots ,\xi_n)$ 
\begin{equation}\label{eq:XIII.5.17}
\text{ because } (x ^{i} - x _{0} ^{i}) (x ^{j} - x _{0} ^{j}) \xi_{i} \xi_{j} = \left(\sum\limits_{i=1} ^{n} (x ^{i} - x _{0} ^{i}) \xi_{i}\right) ^{2}.
\end{equation}
\end{remarku}

The Brownian Motion on the Surface of the $3$-sphere. In the special case when $S$ is the surface of the $3$-sphere $S^{3}$ and $G$ is the group of rotations of $S^{3}$, the infinitesimal generator $A$ of the semi-group induced by a Brownian motion on $S$ is proved to be of the form 
$A = C\Lambda$, where $C$ is a positive constant and $\pmb{\Lambda}$ is the Laplacian 
$= \frac{1}{\sin\theta}\frac{\partial}{\partial\theta}\sin \theta \frac{\partial}{\partial\theta} +\frac{1}{\sin^2\theta}\frac{\partial^2}{\partial\varphi^2}$ on the surface of $S^3$ 


Thus there exists essentially one Brownian motion on the surface of $S^{3}$. For the detail, see K. \textsc{Yosida} \cite{yosida370}. 


\yossection{6}{The Generalized Laplacian of W. Feller}
\label{sec:XIII.6}

Let $S$ be an open interval $(r_1, r_2)$, finite or infinite, on the real line, and $\mathfrak{B}$ be the set of all Baire sets in $(r_1, r_2)$. Consider a Markov process $P(t, x, E)$ on $(S, \mathfrak{B})$ satisfying Lindeberg's condition 
\begin{equation}\label{eq:XIII.6.1}
\lim\limits_{t \downarrow 0} t^{-1} \int_{|x - y| > \varepsilon} P(t, x, \dd y) = 0 \text{ for every } x \in (r_1, r_2) \text{ and } \varepsilon > 0.
\end{equation}


Let $C[r_1, r_2]$ be the $B$-space of real-valued, bounded uniformly continuous functions $f(x)$ defined in $(r_1, r_2)$ and normed by $\|f\| = \sup |f(x)|$. Then 
\begin{equation}\label{eq:XIII.6.2}
(T _{t} f) (x) = \int_ {r _{1}} ^ {r _{2}} P (t, x, \dd y) f (y) (t > 0); = f (x) (t = 0)
\end{equation}
defines a positive contraction semi-group on $C[r_1, r_2]$. We prove by ($1$) the weak right continuity of $T_t$ at $t = 0$ so that $T_t$ is of class $(C_0)$ by the Theorem in Chapter~\ref{chap:IX}, \ref{sec:IX.1}. 

Concerning the infinitesimal generator $A$ of the semi-group $T_{t}$, we have the following two theorems due to W. \textsc{Feller}. 

\begin{theorem}\label{thm:XIII.6.Theorem.1} 
\begin{equation}\label{eq:XIII.6.3}
A \cdot 1 = 0;
\end{equation}
$A$ is of local character, in the sense that, if $f \in D(A)$ 
vanishes in a neighbourhood of a point $x_0$ then 
\begin{equation}\label{eq:XIII.6.4}
(A f) (x _{0}) = 0;
\end{equation}
\end{theorem}

If $f \in D(A)$ has a local maximum at $x_0$, then 
\begin{equation}\label{eq:XIII.6.5}
(A f) (x _{0}) \leqq 0.
\end{equation}

\begin{proof}
 \eqref{eq:XIII.6.3} is clear from $T_{t} \cdot 1 = 1$. We have, by ($1$), 
\[
(A f) \left(x _{0}\right) = \lim\limits _{t \downarrow 0} t ^{-1} \int_ {\left| x _{0} - y \right| \leqq \varepsilon} P (t, x _{0}, \dd y) (f (y) - f \left(x _{0}\right))
\]
independently of sufficiently small $\varepsilon > 0$. Hence, by $P(t, x, E) \geqq 0$, we easily obtain \eqref{eq:XIII.6.4} and \eqref{eq:XIII.6.5}. 
\end{proof}

\begin{theorem}\label{thm:XIII.6.Theorem.2} Let a linear operator $A$, defined on a linear subspace $D(A)$ of $C[r_{1}, r_{2}]$ into $C[r_{1}, r_{2}]$, satisfy conditions \eqref{eq:XIII.6.3}, \eqref{eq:XIII.6.4} and \eqref{eq:XIII.6.5}. Let us assume that $A$ does not degenerate in the sense that the following two conditions are satisfied: 

\begin{equation}
\label{eq:XIII.6.6}
\begin{minipage}{0.92\linewidth}
\everymath{}\let\limits\nolimits
there exists at least one $f_0 \in D(A)$ such that $(Af_0)(x)$
\end{minipage}
\end{equation}
\begin{equation}\label{eq:XIII.6.7}
> 0 \text{ for all } x \in (r _{1}, r _{2}),
\end{equation}
there exists a solution $v$ of $A \cdot v = 0$ which is linearly 


independent of the function $1$ in every subinterval 
\[
\left(x _{1}, x _{2}\right) \text{ of the interval } \left(r _{1}, r _{2}\right).
\]
\end{theorem}

Then there exists a strictly increasing continuous solution $s = s(x)$ of $A \cdot s = 0$ in $(r_1, r_2)$ such that, if we define a strictly increasing function $m = m(x)$, not necessarily continuous nor bounded in $(r_1, r_2)$, by 
\begin{equation}\label{eq:XIII.6.8}
m (x) = \int^{x} \frac {1}{(A f _{0}) (t)} d \left(D _{s} ^{+} f _{0} (t)\right),
\end{equation}
we obtain the representation: 
\begin{equation}\label{eq:XIII.6.9}
(A f) (x) = D _{m} ^{+} D _{s} ^{+} f (x) \text{ in } (r _{1}, r _{2}) \text{ for any } f \in D (A).
\end{equation}

Here the right derivative $D_{p}^{+}$ with respect to a strictly increasing function $p = p(x)$ is defined by 
\begin{equation}\label{eq:XIII.6.10}
\begin{aligned} D _{p} ^{+} g (x) & = \lim\limits _{y \downarrow x} \frac {g (y + 0) - g (x - 0)}{p (y + 0) - p (x - 0)} \text{ at the point of continuity of } p, \\ & = \frac {g (x + 0) - g (x - 0)}{p (x + 0) - p (x - 0)} \text{ at the point of discontinuity of } p. \end{aligned}
\end{equation}

\begin{proof}
 The first step. Let $u \in D(A)$ satisfy $A \cdot u = 0$ and $u(x_1) = u(x_2)$, where $r_1 < x_1 < x_2 < r_2$. Then $u(x)$ must reduce to a constant in $(x_1, x_2)$. If otherwise, $u(x)$ would, for example, assume a maximum at an interior point $x_0$ of $(x_1, x_2)$ in such a way that $u(x_0) - \varepsilon \geqq u(x_1) = u(x_2)$ with some $\varepsilon > 0$. Let $f_0$ be the function $\in D(A)$ given by condition \eqref{eq:XIII.6.6}. Then, for a sufficiently large $\delta > 0$, the function $F(x) = f_0(x) + \delta u(x)$ satisfies $F(x_0) > F(x_i)$ ( $i = 1, 2$). Hence the maximum of $F(x)$ in $[x_1, x_2]$ is attained at an interior point $x_0'$ although $(AF)(x_0') = (Af_0)(x_0') > 0$. This is a contradiction to condition \eqref{eq:XIII.6.5}. 


The second step. By \eqref{eq:XIII.6.7} and the first step, we see that there exists a strictly increasing continuous solution $s(x)$ of $A \cdot s = 0$. We reparametrize the interval by $s$ so that we may assume that 
\begin{equation}
\label{eq:XIII.6.11}
\begin{minipage}{0.92\linewidth}
\everymath{}\let\limits\nolimits
the functions $1$ and $x$ are both solutions of $A \cdot f = 0.$
\end{minipage}
\end{equation}
We can then prove that 
\begin{equation}\label{eq:XIII.6.12}
\parbox{0.92\linewidth}{\everymath{}\let\limits\nolimits if $(Ah)(x) > 0$ for all $x$ in a subinterval $(x_{1}, x_{2})$ of the interval $(r_1, r_2)$, then $h(x)$ is convex downwards in $(x_1, x_2)$.}
\end{equation}


For, the function $u(x) = h(x) - \alpha x - \beta$, which satisfies $(Au)(x) = (Ah)(x) > 0$ for all $x$ in $(x_1, x_2)$, can have no local maximum at an interior point of $(x_1, x_2)$. 

The third step. Let $f \in D(A)$. Then, by \eqref{eq:XIII.6.12} and \eqref{eq:XIII.6.7}, we see that, for sufficiently large $\delta > 0$, the two functions $f_{1}(x) = f(x) + \delta f_{0}(x)$ and $f_{2}(x) = \delta f_{0}(x)$ are both convex downwards. Being the difference of two convex functions, $f(x) = f_{1}(x) - f_{2}(x)$ has a right derivative at every point $x$ of $(r_{1}, r_{2})$. Let us put 
\begin{equation}\label{eq:XIII.6.13}
A \cdot f = \varphi , A \cdot f _{0} = \varphi_{0}, D _{x} ^{+} f _{0} (x) = \mu (x).
\end{equation}

For $g_{t} = f - tf_{0}$, we have $A \cdot g_{t} = Af - tA \cdot f_{0} = \varphi - t\varphi_{0}$. Hence, by ($12$), 
$t < \min\limits_{x_1<x<x_2} \varphi(x) / \varphi_0(x)$ implies that $g_t(t)$ is convex downwards in $(x_1, x_2)$ and so $D_x^+ g_t(x)$ is increasing in $(x_1, x_2)$. 

Applying the same argument for $t > \max\limits_{x_1<x<x_2} \varphi(x) / \varphi_0(x)$, we obtain: 

for any sub-interval $(x_{1}, x_{2})$ of $(r_1, r_2)$, we have 
\[
\begin{aligned}
\min\limits _ {x _{1} < x < x _{2}} \frac {\varphi (x)}{\varphi_{0} (x)} \times (D _{x} ^{+} f _{0} (x _{2}) - D _{x} ^{+} f _{0} (x _{1}))& \leqq D _{x} ^{+} f (x _{2}) - D _{x} ^{+} f (x _{1})
\\
\quad& \leqq \max\limits _ {x _{1} < x < x _{2}} \frac {\varphi (x)}{\varphi_{0} (x)} \times (D _{x} ^{+} f _{0} (x _{2}) - D _{x} ^{+} f _{0} (x _{1})).
\end{aligned}
\]

The continuity of the function $\varphi(x)/\varphi_{0}(x)$ implies, by the above inequality, that 
\[
D _{x} ^{+} f (x _{2}) - D _{x} ^{+} f (x _{1}) = \int_ {x _{1}} ^ {x _{3}} \frac {\varphi (x)}{\varphi_{0} (x)} d \left(D _{x} ^{+} f _{0} (x)\right),
\]
which is precisely the integrated version of $A \cdot f = D_{m}^{+}D_{x}^{+}f$.
\end{proof}

\begin{remarkn}[Remark $1$] We may consider the operator $D_{m}^{+}D_{s}^{+}$ as a generalized Laplacian in the sense that $D_{s}^{+}$ and $D_{m}^{+}$ correspond to the generalized gradient and the generalized divergence in one dimension, respectively. Feller called $s = s(x)$ the canonical scale and $m = m(x)$ the canonical measure of the Markov process in question. We easily see, from the first step above, that the function $1$ and $s$ constitute a fundamental system of linearly independent solutions of $A\cdot v = 0$ to the effect that any solution of $A\cdot v = 0$ can be expressed uniquely as a linear combination of $1$ and $s(x)$. Thus the canonical scale of $P(t,x,E)$ is determined up to a linear transformation, i.e., another canonical scale $s_1$ must be of the form $s_1 = \alpha s + \beta$ with $\alpha >0$; hence the canonical measure $m_1$ which corresponds to $s_1$ must be of the form $m_1 = \alpha^{-1}m$. 
\end{remarkn}

\begin{remarkn}[Remark $2$] Theorem~\ref{thm:XIII.6.Theorem.2} gives the representation of the infinitesimal generator $A$ at the interior point $x$ of $(r_{1}, r_{2})$. To determine the operator $A$ as the infinitesimal generator of a positive contraction semi-group $T_{t}$ of class $(C_{0})$ on $C[r_{1}, r_{2}]$ into $C[r_{1}, r_{2}]$, we must consider the lateral condition, that is, the boundary condition of the operator $A$ at both boundary points $r_{1}$ and $r_{2}$ in order to describe the domain $D(A)$ of $A$ concretely and completely. According to \textsc{Feller} \cite{erdelyi57} and \cite{erdelyi61} (cf. E. \textsc{Hille} \cite{hille102}), the boundary points $r_{1}$ (or $r_{2}$) are classified into the regular-boundary, the exit-boundary, the entrance-boundary and the natural-boundary. To this purpose, we introduce four quantities: 
\begin{align}
\sigma_{1}& = \iint_ {r _{1} < y < x < r _{1} ^{\prime}} \dd m (x) \dd s (y), \mu_{1} = \iint_ {r _{1} < y < x < r _{1} ^{\prime}} \dd s (x) \dd m (y), \notag \\
\sigma_{2}& = \iint_ {r _{1} > y > x > r _{1} ^{\prime}} \dd m (x) \dd s (y), \quad \mu_{2} = \iint_ {r _{1} > y > x > r _{1} ^{\prime}} \dd s (x) \dd m (y). \notag
\end{align}
\end{remarkn}

The boundary point $r_i$ ( $i = 1, 2$) is called 
\[
\sigma_{i} < \infty , \mu_{i} < \infty
\]
exit in case 
\[
\sigma_{i} = \infty , \mu_{i} < \infty
\]
natural in case 
\begin{gather}r _{1} ^{\prime} \notag \\\left. r _{2} ^{\prime}\right) \notag
\end{gather}

We shall illustrate by simple examples. 

\begin{example}\label{ex:XIII.6.Example.1} $D_{m}^{+}D_{s}^{+} = d^{2} / \dd x^{2}, S = (-\infty, \infty)$. We can take $s = x$, $m = x$. Hence 
\begin{equation}\sigma_{2} = \iint_ {\infty > y > x > r _{3} ^{\prime}} \dd x \dd y = \int_ {r _{3} ^{\prime}} ^{\infty} (y - r _{2} ^{\prime}) \dd y = \infty , \mu_{2} = \iint_ {\infty > y > x > r _{1} ^{\prime}} \dd x \dd y = \infty , \notag\end{equation}
and so $\infty$ is a natural boundary. Similarly $-\infty$ is also a natural boundary. 
\end{example}

\begin{example}\label{ex:XIII.6.Example.2} $D_{m}^{+}D_{s}^{+} = d^{2} / \dd x^{2}, S = (-\infty,0)$. We can take $s = x$, $m = x$. In this case, $-\infty$ is a natural boundary and $0$ is a regular boundary. 
\end{example}

\begin{example}\label{ex:XIII.6.Example.3} $D_{m}^{+}D_{s}^{+} = x^{2}\frac{d^{2}}{\dd x^{2}} -\frac{d}{\dd x}, S = (0,\infty)$. A strictly increasing continuous solution $s = s(x)$ of $D_{m}^{+}D_{s}^{+}$ $s = 0$ is given by $s(x) = \int e^{-1/t}\dd t$ so that 
\[
D _{s} = e ^{1/x} \frac {d}{\dd x}, D _{m} ^{+} D _{s} ^{+} = x ^{2} e ^{-1/x} \frac {d}{\dd x} e ^{1/x} \frac {d}{\dd x}.
\]
\end{example}

Therefore $\dd s = e^{-1/x} \dd x$, $\dd m = x^{-2} e^{1/x} \dd x$. Hence 
\begin{align}
\sigma_{1}& = \iint_{0<y<x<1} x ^{-2} e ^{1/x} \dd x e ^{-1/y} \dd y = \int_{0} ^{1} [ - e ^{1/x} ] _{y} ^{1} e ^{-1/y} \dd y < \infty , \notag \\
\mu_{1}& = \iint_{0<y<z<1} e ^{-1/x} \dd x y ^{-2} e ^{1/y} \dd y = \int_{0} ^{1} e ^{-1/x} [ - e ^{1/y} ] _{0} ^{x} \dd x = \infty . \notag
\end{align}

Thus $0$ is an exit boundary. Similarly we see 
\begin{align}
\sigma_{2}& = \iint_{\infty > y > x > 1} e ^{1/x} x ^{-2} \dd x e ^{-1/y} \dd y = \int_{1} ^{\infty} [ - e ^{1/x} ] _{1} ^{y} \cdot e ^{-1/y} \dd y = \infty , \notag \\
\mu_{2}& = \iint_{\infty > y > x > 1} e ^{-1/x} \dd x e ^{1/y} y ^{-2} \dd y = \int_{1} ^{\infty} e ^{-1/x} [ - e ^{1/y} ] _{x} ^{\infty} \dd x = \infty . \notag
\end{align}

That $\infty$ is a natural boundary. 

\begin{example}\label{ex:XIII.6.Example.4} $D_{m}^{+}D_{s}^{+} = x^{2}\frac{d^{2}}{\dd x^{2}} +\frac{d}{\dd x}, S = (0,2)$. As above, we see that $\dd s = e^{1/x}\dd x, \dd m = x^{-2}e^{-1/x}\dd x$, and so we can verify that $0$ is an entrance boundary and $2$ is a regular boundary. 
\end{example}

Feller's probabilistic interpretation of the above classification is as follows: 

The probability that a particle, located at first in the interior of the open interval $(r_{1}, r_{2})$, will, after a finite lapse of time, reach a regular boundary or an exit boundary, is positive; while, the particle can, after a finite lapse of time, neither reach an entrance boundary nor a natural boundary. 

\yossub{References}

\textsc{Feller}'s original paper was published in $1952$: W. \textsc{Feller} \cite{erdelyi57}. The proof of Theorem~\ref{thm:XIII.6.Theorem.2} given above is adapted from W. \textsc{Feller} \cite{erdelyi58} and \cite{erdelyi59}, who also delivered an inspiring address \cite{erdelyi60}. We also refer the readers to the forthcoming book by K. \textsc{Ito}-H. McKean \cite{hunt120} and E. B. \textsc{Dynkin} [$3$], and the references given in these books. 


\yossection{7}{An Extension of the Diffusion Operator}
\label{sec:XIII.7}

Let the possible state of a (not necessarily temporally homogeneous) Markov process be represented by the point $x$ of an $n$-dimensional $C^{\infty}$ Riemannian space $R$. We denote by $P(s, x, t, E)$, $s \leqq t$, the transition probability that a state $x$ at a time moment $s$ is transferred into a Baire set $E \subseteq R$ at a later time moment $t$. 

We shall be concerned with the possible form of the operator $A_{s}$ defined by 
\begin{equation}\label{eq:XIII.7.1}
\left(A _{s} f\right) (x) = \lim\limits _{t \downarrow 0} t ^{-1} \int_{R} P (s, x, s + t, \dd y) (f (y) - f (x)), f \in C _{0} ^{0} (R),
\end{equation}
without assuming the Lindeberg type condition: 

\begin{equation}\label{eq:XIII.7.2}
\begin{minipage}[b]{0.92\linewidth}
$\lim\limits_{t \downarrow 0} t^{-1} \int_{d(x, y) \geqq \varepsilon} P(s, x, s + t, \dd y) = 0$ for all positive constants $\varepsilon$ \\
$(d(x,y) =$ the geodesic distance between two points $x$ and $y)$
\end{minipage}
\end{equation}


We can prove 

\begin{theoremu} Let there exist an increasing sequence $\{k\}$ of positive integers such that, for a fixed pair $\{s, x\}$, 
\begin{align}
\lim\limits _{a \uparrow \infty} k \cdot& \int_{d (x, y) \geqq a} P (s, x, s + k ^{-1}, \dd y) = 0 \text{ uniformly in } k, \label{eq:XIII.7.3} \\
k& \int_ {\dot {R}} \frac {d (x , y) ^{2}}{1 + d (x , y) ^{2}} P (s, x, s + k ^{-1}, \dd y) \text{ is uniformly bounded in } k. \label{eq:XIII.7.4}
\end{align}
\end{theoremu}

Suppose that, for a function $f(x) \in C_0^2(R)$, 
\begin{equation}\label{eq:XIII.7.5}
\text{ a finite } \lim\limits _{k \to \infty} k \left\{\int_{R} P (s, x, s + k ^{-1}, \dd y) f (y) - f (x) \right\} \text{ exists. }
\end{equation}

We have then, in any fixed local coordinates $(x_{1}, x_{2}, \ldots, x_{n})$ of the point $x \in R$, 
\begin{equation}\label{eq:XIII.7.6}
\begin{aligned}
(A _{s} f) (x)& = a _{j} (s, x) \frac {\partial f}{\partial x _{j}} + b _{ij} (s, x) \frac {\partial^{2} f}{\partial x _{i} \partial x _{j}}
\\
\qquad + \varlimsup\limits_{\varepsilon \downarrow 0}& \int_{d (x, y) \geqq \varepsilon} \left\{f (y) - f (x) - \frac {\varrho (x , y)}{1 + d (x , y) ^{2}} (y _{j} - x _{j}) \frac {\partial f}{\partial x _{j}} \right\}
\\
&\quad \qquad \times \frac {1 + d (x , y) ^{2}}{d (x , y) ^{2}} G (s, x, \dd y),
\end{aligned}
\end{equation}
where 
$G(s, x, E)$ is non-negative and $\sigma$-additive for Baire sets 
\begin{equation}\label{eq:XIII.7.7}
E \subseteq R \text{ and } G (s, x, R) < \infty ,
\end{equation}
$\varrho(x, y)$ is continuous in $x, y$ such that $\varrho(x, y) = 1$ or 
\begin{gather}0 \text{ according as } d (x, y) \leqq \delta / 2 \text{ or } > \delta (\delta \text{ is a fixed constant } > 0), \label{eq:XIII.7.8}
\\\text{ the quadratic form } b _{ij} (s, x) \xi_{i} \xi_{j} \text{ is } \geqq 0. \label{eq:XIII.7.9}
\end{gather}

\begin{remarku} Formula \eqref{eq:XIII.7.6} is given in K. \textsc{Yosida} \cite{yosida369}. It is an extension of the diffusion operator of the form of the second order elliptic differential operator discussed in the preceding sections. The third term on the right of \eqref{eq:XIII.7.6} is the sum of infinitely many difference operators. The appearance of such a term is due to the fact that we do not assume the condition ($2$) of Lindeberg's type. Formula \eqref{eq:XIII.7.6} is an operator-theoretical counterpart of the P. Lévy-A. Khintchine-K. Itô infinitely divisible law in probability theory. About this point, see E. \textsc{Hille}-\textsc{Phillips} \cite{hille97}, p. $652$. 
\end{remarku}

\begin{proof}[Proof of Theorem.]
Consider a sequence of non-negative, $\sigma$-additive measures given by
\begin{equation}\label{eq:XIII.7.10}
G _{k} (s, x, E) = k \int_{E} \frac {d (x , y) ^{2}}{1 + d (x , y) ^{2}} P (s, x, s + k ^{-1}, \dd y).
\end{equation}

We have, by \eqref{eq:XIII.7.3} and \eqref{eq:XIII.7.4}, 
\begin{align}
&\quad G _{k} (s, x, E) \text{ is uniformly bounded in } E \text{ and } k, \label{eq:XIII.7.11} \\
\lim\limits _{a \uparrow \infty}& \int_{d (x, y) \geqq a} G _{k} (s, x, \dd y) = 0 \text{ uniformly in } k. \label{eq:XIII.7.12}
\end{align}

Hence, for fixed $\{s, x\}$, the linear functional $L_{k}$ given by 
\[
L _{k} (g) = \int_{R} G _{k} (s, x, \dd y) g (y), g \in C _{0} ^{0} (R),
\]
is non-negative and continuous on the normed linear space $C_{0}^{0}(R)$ with the norm $\|g\|=\sup |g(x)|$; and the norm $L_{k}$ is uniformly bounded in $k$. 

Therefore, by the separability of the normed linear space $C_0^0 (R)$, we can choose a subsequence $\{k'\}$ of $\{k\}$ such that $\lim\limits_{k\to \infty}L_{k'}(g) = L(g)$ exists as a non-negative linear functional on $C_0^0 (R)$. By virtue of a lemma in measure theory (P. R. \textsc{Halmos} \cite{halmos92}) there exists a non-negative, $\sigma$-additive measure $G(s,x,E)$ such that $G(s,x,R) < \infty$ and 
\begin{equation}\label{eq:XIII.7.13}
\lim\limits_{k = k' \to \infty} \int_{R} G_k(s, x, \dd y) g(y) = \int_{R} G(s, x, \dd y) g(y) \quad \text{ for all } \quad g \in C_0^0(R).
\end{equation}

\begin{equation}\label{eq:XIII.7.14}
\begin{aligned}
&\quad \text{ We have }
\\
&\quad k \left[ \int_ {\tilde {R}} P (s, x, s + k ^{-1}, \dd y) f (y) - f (x) \right]
\\
&= \int_ {\tilde {R}} \left\{\left[ f (y) - f (x) - \frac {\varrho (y , x)}{1 + d (y , x) ^{2}} (y _{j} - x _{j}) \frac {\partial f}{\partial x _{j}} \right] \frac {1 + d (y , x) ^{2}}{d (y , x) ^{2}} \right\} G _{k} (s, x, \dd y)
\\
+& \int_ {\tilde {R}} \frac {\varrho (y , x)}{d (y , x) ^{2}} (y _{j} - x _{j}) \frac {\partial f}{\partial x _{j}} G _{k} (s, x, \dd y).
\end{aligned}
\end{equation}

The term $\{\}$ in the first integral on the right is, for sufficiently small $d(y, x)$, 
\[
\begin{aligned}
&= (y _{j} - x _{j}) \frac {\partial f}{\partial x _{j}} + (y _{i} - x _{i}) (y _{j} - x _{j}) \left(\frac {\partial^{2} f}{\partial X _{i} \partial X _{j}}\right) \frac {1 + d (y , x) ^{2}}{d (y , x) ^{2}},
\\
\text{ where } X _{j} &= x _{j} + \theta (y _{j} - x _{j}), 0 < \theta < 1.
\end{aligned}
\]

Thus $\{\}$ is bounded and continuous in $y$. Hence, by \eqref{eq:XIII.7.12} and ($13$), the first term on the right of \eqref{eq:XIII.7.14} tends, as $k = k' \to \infty$, to $\int_{R} \{\}\ G(s, x, \dd y)$. Therefore, by \eqref{eq:XIII.7.5}, 
\[
\text{ a finite } \lim\limits _ {k = k ^{\prime} \rightarrow \infty} \int_{R} \frac {\varrho (y , x)}{d (y , x) ^{2}} \left(y _{j} - x _{j}\right) \frac {\partial f}{\partial x _{j}} G _{k} (s, x, \dd y) = a _{j} (s, x) \frac {\partial f}{\partial x _{j}}
\]
exists. Hence, by \eqref{eq:XIII.7.3}, we obtain \eqref{eq:XIII.7.6} by taking 
\[
b _{ij} (s, x) = \lim\limits _{\varepsilon \downarrow 0} \lim\limits _ {k = k ^{\prime} \rightarrow \infty} k \int_{d (y, x) \leqq \varepsilon} \left(y _{i} - x _{i}\right)\left(y _{j} - x _{j}\right) P (s, x, s + k ^{-1}, \dd y).
\]
\end{proof}


\yossection{8}{Markov Processes and Potentials}
\label{sec:XIII.8}

Let $\{T_t\}$ be an equi-continuous semi-group of class $(C_0)$ on a function space $X$ defined by $(T_t f)(x) = \int_S P(t, x, \dd y) f(y)$ where $P(t, x, E)$ is a Markov process on a measure space $(S, \mathfrak{B})$. Let $A$ be the infinitesimal generator of $T_t$. Suggested by the special case where $A$ is the Laplacian, an element $f \in X$ is called harmonic if $A \cdot f = 0$. Then $f$ is harmonic iff $\lambda (\lambda I - A)^{-1} f = f$ for every $\lambda > 0$. An element $f \in X$ is called a potential if $f = \lim\limits_{\lambda \downarrow 0} (\lambda I - A)^{-1} g$ with some $g$. Because, for such an element $f$, there holds by the closedness of $A$ the Poisson equation 
\[
A \cdot f = \lim\limits _{\lambda \downarrow 0} A (\lambda I - A) ^{-1} g = \lim\limits _{\lambda \downarrow 0} \{- g + \lambda (\lambda I - A) ^{-1} g \} = - g.
\]

Suppose that $X$ is a vector lattice which is also a locally convex linear topological space such that every monotone increasing bounded sequence of elements $\in X$ converges weakly to an element of $X$ which is greater than the elements of the sequence. We also assume that the resolvent $J_{\lambda} = (\lambda I - A)^{-1}$ is positive in the sense that $f \geqq 0$ implies $J_{\lambda} f \geqq 0$. The situation is suggested by the special case where $A$ is the Laplacian considered in a suitable function space. 

We may thus call subharmonic those elements $f \in X$ for which the inequality $A \cdot f \geqq 0$ holds. By virtue of the positivity of $J_{\lambda}$, a subharmonic element $f$ satisfies $\lambda J_{\lambda} f \geqq f$ for all $\lambda > 0$. We shall prove an analogue of a well-known theorem of F. \textsc{Riesz} concerning ordinary subharmonic functions (see T. \textsc{Rado} \cite{rado266}): 

\begin{theoremu} Any subharmonic element $x$ is decomposed as the sum of a harmonic element $x_{h}$ and a potential $x_{p}$, where the harmonic part $x_{h}$ of $x$ is given by $x_{h} = \lim\limits_{\lambda \downarrow 0} \lambda J_{\lambda} x$ and $x_{h}$ is the least harmonic majorant of $x$ in the sense that any harmonic element $x_{H} \geqq x$ satisfies $x_{H} \geqq x_{h}$. 
\end{theoremu}

Proof (K. \textsc{Yosida} \cite{yosida347}). By the resolvent equation 
\begin{equation}\label{eq:XIII.8.1}
J _{\lambda} - J _{\mu} = (\mu - \lambda) J _{\lambda} J _{\mu},
\end{equation}
we obtain 
\begin{equation}\label{eq:XIII.8.2}
(I - \lambda J _{\lambda}) = (I + (\mu - \lambda) J _{\lambda}) (I - \mu J _{\mu}).
\end{equation}

Since $x$ is subharmonic, we see, by the positivity of $J_{\lambda}$, that 
\[
\lambda > \mu \quad \text{ implies } \quad \lambda J _{\lambda} x \geqq \mu J _{\mu} x \geqq x.
\]

Hence the weak-lim $\lambda J_{\lambda}x = x_{h}$ exists by virtue of the hypothesis concerning bounded monotone sequences in $X$. Therefore, by the ergodic theorem in Chapter~\ref{chap:VIII}, \ref{sec:VIII.4}, we see that $x_{h} = \lim\limits_{\lambda \downarrow 0} \lambda J_{\lambda}x$ exists and $x_{h}$ is harmonic, i.e., $\lambda J_{\lambda}x_{h} = x_{h}$ for all $\lambda > 0$. We also have $x_{p} = (x - x_{h}) = \lim\limits_{\lambda \downarrow 0} (I - \lambda J_{\lambda})x = \lim\limits_{\lambda \downarrow 0} (-A)(\lambda I - A)^{-1}x = \lim\limits_{\lambda \downarrow 0} (\lambda I - A)^{-1}(-Ax)$ which shows that $x_{p}$ is a potential. 

Let a harmonic element $x_{H}$ satisfy $x_{H} \geqq x$. Then, by the positivity of $\lambda J_{\lambda}$ and the harmonic property of $x_{H}$, we have 
\[
x _{H} = \lambda J _{\lambda} x _{H} \geqq \lambda J _{\lambda} x \quad \text{ and hence } \quad x _{H} \geqq \lim\limits _{\lambda \downarrow 0} \lambda J _{\lambda} x = x _{h}.
\]


\yossection{9}{Abstract Potential Operators and Semi-groups}
\label{sec:XIII.9}

Let $\{T_{t}; t \geqq 0\}$ be an equi-continuous semi-group of class $(C_0)$ of bounded linear operators on a Banach space $X$ into $X$, and let $A$ be its infinitesimal generator. Thus $D(A)^a = X$ and the resolvent $(\lambda I - A)^{-1}$ of $A$ exists for $\lambda > 0$ as a bounded linear operator on $X$ into $X$ satisfying the condition 
\begin{equation}\label{eq:XIII.9.1}
\sup\limits _{\lambda > 0} \| \lambda (\lambda I - A) ^{-1} \| < \infty
\end{equation}
so that, by applying abelian ergodic theorems in Chapter~\ref{chap:VIII}, \ref{sec:VIII.4}, we obtain 
\begin{equation}\label{eq:XIII.9.2}
\begin{aligned}
&R (A) ^{a}  = R (A (\mu I - A) ^{-1}) ^{a} = \{x \in X; \slim _{\lambda \downarrow 0} \lambda (\lambda I - A) ^{-1} x = 0 \} \text{ for all } \mu > 0
\end{aligned}
\end{equation}
and 
\begin{equation}\label{eq:XIII.9.3}
\begin{aligned}
&D (A) ^{a}  = R ((\mu I - A) ^{-1}) ^{a} = \{x \in X; \slim _{\lambda \uparrow \infty} \lambda (\lambda I - A) ^{-1} x = x \} = X \text{ for all } \mu > 0.
\end{aligned}
\end{equation}

We shall be concerned with the situation in which 
\begin{equation}
\label{eq:XIII.9.4}
\begin{minipage}{0.92\linewidth}
\everymath{}\let\limits\nolimits
the inverse $A^{-1}$ exists with its domain $D(A^{-1}) = R(A)$ strongly dense in $X,$
\end{minipage}
\end{equation}
and in case \eqref{eq:XIII.9.4} we shall call $V = -A^{-1}$ as the ``abstract potential operator'' associated with the semi-group $T_{t}$ so that the ``abstract potential $V$ $x$'' satisfies the ``abstract Poisson equation'' 
\begin{equation}\label{eq:XIII.9.5}
A V x = - x
\end{equation}
with respect to the ``abstract Laplacian operator'' $A$. 

\begin{proposition}[K. \textsc{Yosida} {\cite{yosida377}}]\label{prop:XIII.9.Proposition.1} \eqref{eq:XIII.9.4} is equivalent to the condition 
\begin{equation}\label{eq:XIII.9.6}
\slim _{\lambda \downarrow 0} \lambda (\lambda I - A) ^{-1} x = 0 \quad \text{ for every } \quad x \in X.
\end{equation}
\end{proposition}

\begin{proof}
 The condition $A$ $x$ = $0$ is equivalent to $\lambda(\lambda I - A)^{-1}$ $x$ = $x$ and so \eqref{eq:XIII.9.6} implies the existence of $A^{-1}$. Moreover, by \eqref{eq:XIII.9.2}, the denseness of $R(A)$ is equivalent to \eqref{eq:XIII.9.6} so that ($4$) must be equivalent to \eqref{eq:XIII.9.6}. 
\end{proof}

\begin{corollaryu} In case \eqref{eq:XIII.9.6}, the abstract potential operator $V = -A^{-1}$ is defined also by 
\begin{equation}\label{eq:XIII.9.7}
V x = \slim _{\lambda \downarrow 0} (\lambda I - A) ^{-1} x.
\end{equation}
\end{corollaryu}

\begin{proof}
 For $x \in D(A^{-1})$, we have 
\[
\begin{aligned}
&- A ^{-1} x - (\lambda I - A) ^{-1} x  = - A ^{-1} x - A (\lambda I - A) ^{-1} A ^{-1} x = \lambda (\lambda I - A) ^{-1} A ^{-1} x
\end{aligned}
\]
so that, by \eqref{eq:XIII.9.6}, we obtain $\slim\limits_{\lambda \downarrow 0} (\lambda I - A)^{-1} x = -A^{-1}x$. On the other hand, let $\slim\limits_{\lambda \downarrow 0} (\lambda I - A)^{-1} x = y$ exists. Then $A(\lambda I - A)^{-1} x = -x + \lambda (\lambda I - A)^{-1} x$ implies, by the closure property of the infinitesimal generator $A$, the equality $A y = -x$, that is, $y = -A^{-1}x$ and so we have proved \eqref{eq:XIII.9.7}. 
\end{proof}

From now on we shall be concerned with the case 
\begin{equation}\label{eq:XIII.9.8}
D (A) ^{a} = X \quad \text{ and } \quad \sup\limits _{\lambda > 0} \| \lambda (\lambda I - A) ^{-1} \| \leqq 1
\end{equation}
which characterises the infinitesimal generator $A$ of a contraction semi-group $\{T_{t}; t \geqq 0\}$ of class $(C_{0})$. We can prove 

\begin{theorem}\label{thm:XIII.9.Theorem.1} The abstract potential operator $V$ and its dual operator $V^{*}$ must satisfy the following inequalities, respectively: 
\begin{align}
\| \lambda V x + x \|& \geqq \| \lambda V x \| \quad \text{ for all } \quad x \in D (V) \quad \text{ and } \quad \lambda > 0, \label{eq:XIII.9.9} \\
\| \lambda V ^{*} f ^{*} + f ^{*} \|& \geqq \| \lambda V ^{*} f ^{*} \| \quad \text{ for all } \quad f ^{*} \in D (V ^{*}) \quad \text{ and } \quad \lambda > 0. \label{eq:XIII.9.10}
\end{align}
\end{theorem}

\begin{proof}
 We put $J_{\lambda} = (\lambda I - A)^{-1}$. Then $J_{\lambda}$ satisfies the resolvent equation 
\begin{equation}\label{eq:XIII.9.11}
J _{\lambda} - J _{\mu} = (\mu - \lambda) J _{\lambda} J _{\mu}. \end{equation}

Letting $\mu \downarrow 0$ in \eqref{eq:XIII.9.11}, we obtain 
\begin{equation}\label{eq:XIII.9.12}
J _{\lambda} x - V x = - \lambda J _{\lambda} V x \quad \text{ for all } \quad x \in D (V) \quad \text{ and all } \quad \lambda > 0.
\end{equation}

Hence $J_{\lambda}(\lambda V x + x) = V x$ and so, by $\| \lambda J_{\lambda} \| \leqq 1$, we obtain \eqref{eq:XIII.9.9}. 

Since $D(A)$ and $R(A)$ are both strongly dense in $X$, we have 
\begin{equation}\label{eq:XIII.9.13}
(A ^{*}) ^{-1} = (A ^{-1}) ^{*} \quad \text{ or equivalently } \quad V ^{*} = (- A ^{-1}) ^{*} = (- A ^{*}) ^{-1}
\end{equation}
by R. S. \textsc{Phillips}' theorem in Chapter~\ref{chap:VIII}, \ref{sec:VIII.6}. Therefore, we have, when $f^{*} \in D(V^{*})$, 
\begin{align}
V ^{*} f ^{*} - (\lambda I ^{*} - A ^{*}) ^{-1} f ^{*}& = (- A ^{*}) ^{-1} f ^{*} - (\lambda I ^{*} - A ^{*}) ^{-1} f ^{*} \notag \\
&= (- A ^{*}) ^{-1} f ^{*} - A ^{*} (\lambda I ^{*} - A ^{*}) ^{-1} (A ^{*}) ^{-1} f ^{*} = \lambda (\lambda I ^{*} - A ^{*}) ^{-1} (A ^{*}) ^{-1} f ^{*}. \notag
\end{align}

Hence we obtain 
\begin{equation}\label{eq:XIII.9.14}
V ^{*} f ^{*} = w ^{*} - \lim\limits _{\lambda \downarrow 0} J _{\lambda} ^{*} f ^{*}
\end{equation}
since $w^{*} - \lim\limits_{\lambda \downarrow 0} \lambda J_{\lambda}^{*} g^{*} = 0$ for all $g^{*} \in X^{*}$ by \eqref{eq:XIII.9.6}. Thus, as in the case of $V$, we prove \eqref{eq:XIII.9.10} by 
\begin{equation}\tag{11'}\label{eq:XIII.9.11'}
J _{\lambda} ^{*} - J _{\mu} ^{*} = (\mu - \lambda) J _{\lambda} ^{*} J _{\mu} ^{*} \quad \text{ and } \quad \sup\limits _{\lambda > 0} \| \lambda J _{\lambda} ^{*} \| \leqq 1.
\end{equation}
\end{proof}

\begin{corollaryu} The inverse $(\lambda V + I)^{-1}$ and the inverse $(\lambda V^{*} + I^{*})^{-1}$ both exist as bounded linear operators. 
\end{corollaryu}

\begin{proof}
 The existence of $(\lambda V + I)^{-1}$ and of $(\lambda V^{*} + I^{*})^{-1}$ are clear from \eqref{eq:XIII.9.9} and \eqref{eq:XIII.9.10}, respectively. We shall prove 
\begin{equation}\label{eq:XIII.9.15}
R (\lambda V + I) = X \quad \text{ for all } \quad \lambda > 0. \end{equation}

The existence of the inverse $(\lambda V^{*} + I^{*})^{-1}$ implies, by the Hahn-Banach theorem, that 
\[
R (\lambda V + I)
\]
is strongly dense in $X$.\yosparagraphtag{15'}\label{eq:XIII.9.15'}

Hence, for any $y \in X$, there exists a sequence $\{x_n\} \subseteq D(V)$ such that $\slim\limits_{n \to \infty} (\lambda V x_n + x_n) = y$. Thus, by \eqref{eq:XIII.9.9}, we have $\| \lambda V(x_n - x_m) + (x_n - x_m) \| \geqq \| \lambda V(x_n - x_m) \|$ and so $\slim\limits_{n \to \infty} V x_n = z$ exists, proving that $\slim\limits_{n \to \infty} x_n = x$ exists. Since an abstract potential operator is closed by definition, the operator $V$ must be closed. Hence we must have $Vx = z$, that is, $y = \lambda Vx + x$. This proves \eqref{eq:XIII.9.15}, and thus $(\lambda V + I)^{-1}$ is a bounded linear operator by the closed graph theorem. Hence, by the closed range theorem in Chapter~\ref{chap:VIII}, \ref{sec:VIII.5}, we have 
\begin{equation}\label{eq:XIII.9.16}
R (\lambda V ^{*} + I ^{*}) = X ^{*} \quad \text{ for all } \quad \lambda > 0,
\end{equation}
and so, again by the closed range theorem and \eqref{eq:XIII.9.15}, $(\lambda V^{*} + I^{*})^{-1}$ is a bounded linear operator on the dual space $X^{*}$ into $X^{*}$.
\end{proof}

We can now give a characterization of abstract potential operators: 

\begin{theorem}[K. \textsc{Yosida} {\cite{yosida378}}]\label{thm:XIII.9.Theorem.2} Let a closed linear operator $V$ with domain $D(V)$ and range $R(V)$ both strongly dense in $X$ be such that $V$ and its dual operator $V^{*}$ satisfy \eqref{eq:XIII.9.9} and \eqref{eq:XIII.9.10}, respectively. Then $V$ is an abstract potential operator, i.e., -$V$ is the inverse of an operator $A$ which is the infinitesimal generator of a contraction semi-group of class $(C_{0})$ in $X$. 
\end{theorem}

\begin{proof}
 We first remark that the Corollary above holds good for $V$ and for $V^{*}$, and hence $(\lambda V + I)^{-1}$ is a bounded linear operator on $X$ into $X$. Moreover, the linear operator $J_{\lambda}$ defined through 
\begin{equation}\label{eq:XIII.9.17}
\hat {J} _{\lambda} (\lambda V x + x) = V x \quad [ \text{ for all } x \in D (V) \text{ and all } \lambda > 0 ]
\end{equation}
is a bounded linear operator on $X$ into $X$ such that 
\begin{equation}\label{eq:XIII.9.18}
\hat {J} _{\lambda} = V (\lambda V + I) ^{-1} \quad \text{ and } \quad \sup\limits _{\lambda > 0} \| \lambda \hat {J} _{\lambda} \| \leqq 1. \end{equation}

We can prove, by \eqref{eq:XIII.9.17}, that $J_{\lambda}$ is a pseudo-resolvent. In fact, we have 
\[
\begin{aligned} \hat {J} _{\lambda} (\lambda V x + x) - \hat {J} _{\mu} (\lambda V x + x) & = V x - \hat {J} _{\mu} \left(\frac {\lambda}{\mu} (\mu V x + x) + \left(1 - \frac {\lambda}{\mu}\right) x\right) \\ & = V x - \frac {\lambda}{\mu} V x - \left(1 - \frac {\lambda}{\mu}\right) \hat {J} _{\mu} x = \left(1 - \frac {\lambda}{\mu}\right) (V x - \hat {J} _{\mu} x) \end{aligned}
\]
and 
\[
\begin{aligned}
&(\mu - \lambda) \hat {J} _{\mu} \hat {J} _{\lambda} (\lambda V x + x) = (\mu - \lambda) \hat {J} _{\mu} V x = (\mu - \lambda) \hat {J} _{\mu} \left(V x + \frac {1}{\mu} x - \frac {1}{\mu} x\right)
\\
&= (\mu - \lambda) \frac {1}{\mu} V x - (\mu - \lambda) \frac {1}{\mu} \hat {J} _{\mu} x = \left(1 - \frac {\lambda}{\mu}\right) (V x - \hat {J} _{\mu} x).
\end{aligned}
\]

Because of \eqref{eq:XIII.9.18}, we can apply abelian ergodic theorems: 
\begin{align}
R (\hat {J} _{\mu}) ^{a}& = \{x \in X; \slim _{\lambda \uparrow \infty} \lambda \hat {J} _{\lambda} x = x \} \quad \text{ for all } \quad \mu > 0, \label{eq:XIII.9.19} \\
R (I - \mu \hat {J} _{\mu}) ^{a}& = \{x \in X; \slim _{\lambda \downarrow 0} \lambda \hat {J} _{\lambda} x = 0 \} \quad \text{ for all } \quad \mu > 0 \label{eq:XIII.9.20}
\end{align}

Since $R(V)^{a} = X$, we see $R(\hat{\mathcal{J}}_{\mu})^{a} = X$ by \eqref{eq:XIII.9.17}. On the other hand, the null space of $\hat{\mathcal{J}}_{\lambda}$ is independent of $\lambda$ (see Proposition in Chapter~\ref{chap:VIII}, \ref{sec:VIII.4}). Hence, by \eqref{eq:XIII.9.19} and $R(\hat{\mathcal{J}}_{\mu})^{a} = X$, the null space of $\hat{\mathcal{J}}_{\lambda}$ consists of zero vector only. Therefore $\hat{\mathcal{J}}_{\lambda}$ is the resolvent of a linear operator $A$: 
\begin{equation}\label{eq:XIII.9.21}
\hat {J} _{\lambda} = (\lambda I - A) ^{-1}, \quad \text{ where } \quad A = \lambda I - \hat {J} _{\lambda} ^{-1}.
\end{equation}

Hence we obtain 
\begin{align}
D (A) ^{a}& = R (\hat {J} _{\mu}) ^{a} = \{x \in X; \slim _{\lambda \uparrow \infty} \lambda \hat {J} _{\lambda} x = x \} \quad \text{ for all } \quad \mu > 0 \tag{19'}\label{eq:XIII.9.19'} \\
R (A) ^{a}& = R (I - \mu \hat {J} _{\mu}) ^{a} = \{x \in X; \slim _{\lambda \downarrow 0} \lambda \hat {J} _{\lambda} x = 0 \} \quad \text{ for all } \quad \mu > 0 \tag{20'}\label{eq:XIII.9.20'}
\end{align}

By \eqref{eq:XIII.9.19'} and $R(\hat{J}_{\mu})^a = X$, we see that $D(A)^a = X$. We can also prove that $R(A)^a = X$. In fact, we have, by \eqref{eq:XIII.9.17} and \eqref{eq:XIII.9.21}, 
\begin{gather}\quad (\lambda I - A) \hat {J} _{\lambda} (\lambda V x + x)= \lambda V x + x = (\lambda I - A) V x = \lambda V x - A V x, \text{ that is }, \notag
\\- A V x = x \quad \text{ whenever } \quad x \in D (V) . \label{eq:XIII.9.22}
\end{gather}

Hence $R(A) \supseteq D(V)$ is strongly dense in $X$, that is, $R(A)^{a} = X$ and so, by \eqref{eq:XIII.9.20'}, $\slim\limits_{\lambda \downarrow 0} \lambda \hat{J}_{\lambda} x = 0$ for all $x \in X$. Therefore, by Proposition~\ref{prop:XIII.9.Proposition.1}, $-A^{-1}$ is an abstract potential operator. 


That $V = -A^{-1}$ may be proved as follows. Firstly, the inverse $V^{-1}$ exists since, by \eqref{eq:XIII.9.22}, $V$ $x$ = $0$ implies $x$ = $0$. Thus, by \eqref{eq:XIII.9.18} and \eqref{eq:XIII.9.21}, we obtain 
\[
\lambda I - A = \hat {J} _{\lambda} ^{-1} = (\lambda V + I) V ^{-1} = \lambda I + V ^{-1}
\]
which proves that $-A^{-1}=V$.
\end{proof}

\begin{remarku} The above proof shows that \eqref{eq:XIII.9.10} is used only to prove \eqref{eq:XIII.9.15'} and hence \eqref{eq:XIII.9.15}. 
\end{remarku}

As in Chapter~\ref{chap:IX}, \ref{sec:IX.8}, we introduce a semi-scalar product $[x, y]$ in $X$ satisfying conditions: 
\begin{equation}\label{eq:XIII.9.23}
\begin{aligned}
&{[ x + y, z ] = [ x, z ] + [ y, z ], [ \lambda x, y ] = \lambda [ x, y ],} {[ x, x ] = \| x \| ^{2} \quad \text{ and } \quad | [ x, y ] | \leqq \| x \| \cdot \| y \|.}
\end{aligned}
\end{equation}

Then a linear operator $V$, with domain $D(V)$ and range $R(V)$ both in $X$, is called accretive (with respect to $[x, y]$) if 
\begin{equation}\label{eq:XIII.9.24}
\operatorname{Re} [ x, V x ] \geqq 0 \quad \text{ whenever } \quad x \in D (V) .
\end{equation}

We shall prepare three propositions. 

\begin{proposition}\label{prop:XIII.9.Proposition.2} An abstract potential operator $V$ must be accretive. 
\end{proposition}

\begin{proof}
 Let $\{T_t; t \geqq 0\}$ be the contraction semi-group of class $(C_0)$ whose infinitesimal generator $A$ is given by $A = -V^{-1}$. We have, by $\| T_t \| \leqq 1$, 
\[
\operatorname{Re} \left[ T _{t} x - x, x \right] = \operatorname{Re} [ T x, x ] - \| x \| ^{2} \leqq \| T _{t} x \| \cdot \| x \| - \| x \| ^{2} \leqq 0. \]

Hence, for $x \in D(A)$, we obtain 
\[
\operatorname{Re} [ A x, x ] = \lim\limits _{t \downarrow 0} [ t ^{-1} (T _{t} x - x), x ] \leqq 0
\]
so that, for $x_0 \in D(V) = D(A^{-1})$, 
\[
\operatorname{Re} \left[ A V x _{0}, V x _{0} \right] = \operatorname{Re} \left[ - x _{0}, V x _{0} \right] \leqq 0, \quad \text{ that is }, \quad \operatorname{Re} \left[ x _{0}, V x _{0} \right] \geqq 0.
\]
\end{proof}

\begin{proposition}\label{prop:XIII.9.Proposition.3} $V$ satisfies \eqref{eq:XIII.9.9} if $V$ is accretive. 
\end{proposition}

\begin{proof}
 We have, by \eqref{eq:XIII.9.23} and \eqref{eq:XIII.9.24}, 
\[
\begin{aligned}
&{\| \lambda V x \| ^{2} = [ \lambda V x, \lambda V x ] \leqq \operatorname{Re} \{[ \lambda V x, \lambda V x ] \}} {\qquad = \operatorname{Re} [ \lambda V x + x, \lambda V x ] \leqq \| \lambda V x + x \| \cdot \| \lambda V x \|,}
\end{aligned}
\]
that is, we have proved \eqref{eq:XIII.9.9}. 
\end{proof}

\begin{proposition}\label{prop:XIII.9.Proposition.4} Let $V$ be an abstract potential operator associated with a contraction semi-group $\{T_{t}; t \geqq 0\}$ of class $(C_{0})$ with the infinitesimal generator $A$ so that $V = -A^{-1}$. Then the dual operator $V^{*}$ must satisfy, for any semi-scalar product $[f^{*}, g^{*}]$ in the dual space $X^{*}$ of $X$, the inequality 
\begin{equation}\label{eq:XIII.9.25}
[ f ^{*}, V ^{+} f ^{*} ] \geqq 0
\end{equation}
where 
$V^{+}$ is the largest restriction of $V^{*}$ with domain and range both in 
\begin{equation}\label{eq:XIII.9.26}
R (V ^{*}) ^{a}
\end{equation}
\end{proposition}

\begin{proof}
 Let $\{T_t^+; t \geqq 0\}$ be the dual semi-group (see Chapter~\ref{chap:IX}, \ref{sec:IX.13}) of $\{T_t; t \geqq 0\}$, so that, $\{T_t^+; t \geqq 0\}$ be a contraction semi-group of class $(C_0)$ in $X^+ = D(A^*)^a = R(V^*)^a$. Thus the infinitesimal generator $A^+$ of the semi-group $T_t^+$ is the largest restriction of $A^*$ with domain and range both in $X^+ = D(A^*)^a$. Then $V^+ = (-A^+)^{-1}$, and so \eqref{eq:XIII.9.25} is proved as in the case of Proposition~\ref{prop:XIII.9.Proposition.2}. 
\end{proof}

We are now able to prove 

\begin{theorem}\label{thm:XIII.9.Theorem.3} Let $V$ be a closed linear operator with domain and range both strongly dense in $X$. Then $V$ is an abstract potential operator iff $V$ is accretive and $V^{+}$ is accretive. 
\end{theorem}

\begin{proof}
 The ``only if part'' is proved by Proposition~\ref{prop:XIII.9.Proposition.2} and Proposition~\ref{prop:XIII.9.Proposition.4}. We shall prove the ``if part''. Firstly, $V$ satisfies \eqref{eq:XIII.9.9} by Proposition~\ref{prop:XIII.9.Proposition.3}. Next let $\lambda V^{*}f^{*} + f^{*} = 0$. Then $f^{*} \in R(V^{*})$ and so $V^{*}f^{*} = -\lambda^{-1}f^{*} \in R(V^{*})$. Thus $f^{*} \in D(V^{+})$ and $\lambda V^{+}f^{*} = -f^{*}$. Therefore, by the accretive property of $V^{+}$, we obtain 
\[
[ \lambda f ^{*}, \lambda V ^{+} f ^{*} ] = \lambda [ f ^{*}, - f ^{*} ] = - \lambda \| f ^{*} \| ^{2} \geqq 0, \quad \text{ that is }, \quad f ^{*} = 0. \]

Hence the inverse $(\lambda V^{*} + I^{*})^{-1}$ exists and so, by the Hahn-Banach theorem, \eqref{eq:XIII.9.15'} and consequently \eqref{eq:XIII.9.15} both hold good. This proves that $V$ is an abstract potential operator. 
\end{proof}

Comparison with G. A. \textsc{Hunt}'s theory of potentials. Consider a special case in which $X$ is the completion $C_{\infty}(S)$, with respect to the maximum norm, of the space $C_{0}(S)$ of real- or complex-valued continuous functions $x(s)$ with compact support in a locally compact, non-compact separable Hausdorff space $S$. In this case, the semi-scalar product may be defined through 
$[x, y] = x(s_0) \cdot y(s_0)$, where $s_0$ is any fixed point of $S$ such that 
\begin{equation}\label{eq:XIII.9.27}
| y (s _{0}) | = \sup\limits _{s \in S} | y (s) |.
\end{equation}

\begin{equation}\label{eq:XIII.9.28}
\parbox{0.92\linewidth}{\everymath{}\let\limits\nolimits Therefore the accretive property of the linear operator $V$ is given by $\operatorname{Re}(x(s_0) \cdot (\overline{Vx}(s_0)) \geqq 0$ whenever $|(Vx)(s_0)| = \sup\limits_{s \in S} |(Vx)(s)|$.}
\end{equation}

This may be compared with the principle of positive maximum for $V$ in Hunt's theory of potentials (G. A. \textsc{Hunt} \cite{hunt119}, cf. P. A. \textsc{Meyer} \cite{meyer194} and K. \textsc{Yosida} \cite{yosida375} and the bibliography referred to in these). Assuming that the space $C_{\infty}(S)$ consists of real-valued functions only, \textsc{Hunt} introduced the notion of the potential operator $U$ as a positive linear operator on the domain $D(U) \subseteq C_{\infty}(S)$ into $C_{\infty}(S)$ satisfying three conditions: \begin{enumerate}[i)]
\item $D(U) \supseteq C_0(S)$,
\item $U \cdot C_0(S)$ is strongly dense in $C_{\infty}(S)$
\item the principle of positive maximum given by 
\begin{equation}\label{eq:XIII.9.29}
\parbox{0.92\linewidth}{\everymath{}\let\limits\nolimits for each $x(s) \in C_0(S)$, the value $x(s_0) \geqq 0$ whenever $(U x)(s)$ attains its positive supremum at $s = s_{0}$.}
\end{equation}
\end{enumerate}



(In the above, the positivity for $U$ means that $U$ maps non-negative functions into non-negative functions.) \textsc{Hunt} then proved that there exists a uniquely determined semi-group $\{T_{t}; t \geqq 0\}$ of class $(C_{0})$ of positive contraction operators $T_{t}$ on $C_{\infty}(S)$ into $C_{\infty}(S)$ satisfying the following two conditions: 

\begin{equation}\label{eq:XIII.9.30}
\parbox{0.92\linewidth}{\everymath{}\let\limits\nolimits $A U x = -x$ for every $x(s) \in C_{0}(S)$, where $A$ is the infinitesimal generator of $T_{t}$,}
\end{equation}

\begin{equation}\label{eq:XIII.9.31}
(U x) (s) = \int_{0} ^{\infty} \left(T _{t} x\right) (s) \dd t \text{ for every } x (s) \in C _{0} (S), \text{ the integral being }
\end{equation}
taken in the strong topology of the $B$-space $C_{\infty}(S)$. 

It is to be noted that our abstract potential operator $V$ is defined without assuming the positivity of the operator $V$. Moreover, in our formulation, ($30$) can be replaced by the true Poisson equation 
\begin{equation}\tag{30'}\label{eq:XIII.9.30'}
V = - A ^{-1}.
\end{equation}

In this connection, it is to be noted that the smallest closed extension $V$ of the operator $U$ restricted to $C_0(S)$ satisfies ($30$) $'$. See K. \textsc{Yosida}, T. \textsc{Watanabe}, H. \textsc{Tanaka} \cite{yosida379} and K. \textsc{Yosida} \cite{yosida380}. Furthermore, in our formulation of the abstract potential operator $V$, \eqref{eq:XIII.9.31} is replaced by a
more general $(Vx)$ $(s)$ = $s$- $\lim\limits_{n\to\infty}\left((\lambda I-A)^{-1}x\right)$ $(s)$ which is based upon \eqref{eq:XIII.9.7}. This formulation has the advantage that it can be applied to essentially more general class of semi-groups than \eqref{eq:XIII.9.31}. See, e.g., K. \textsc{Yosida} \cite{yosida380}, \cite{yosida381}, A. \textsc{Yamada} \cite{yamada343}, and K. \textsc{Sato} \cite{sato282}, \cite{sato283}. It is to be noted here that F. \textsc{Hirsch} \cite{hirsch106}, \cite{hirsch107} has also developed an abstract potential theory directed by essentially the same idea as the present author. 
   % XIII
% ==== XIV ====
\yoschapter{XIV}{The Integration of the Equation of Evolution*}
\label{chap:XIV}
\yosstarfootnotetext{See also Supplementary Notes, p.~\pageref{notes:IX-XIV}.}

The ordinary exponential function solves the initial value problem 
\[
\dd y / \dd x = \alpha y, y (0) = 1.
\]

We consider the diffusion equation 
\[
\partial u/\partial t=\Delta u,
\]
where $\Delta=\sum_{j=1}^m\partial^2/\partial x_j^2$ is the Laplacian in $\mathbb R^m$.

We wish to find a solution $u = u(x, t)$, $t > 0$, of this equation satisfying the initial condition $u(x, 0) = f(x)$, where $f(x) = f(x_{1}, \ldots, x_{m})$ is a given function of $x$. We shall also study the wave equation 
\[
\partial^{2} u / \partial t ^{2} = \Delta u, - \infty < t < \infty ,
\]
with the initial data 
\[
u (x, 0) = f (x) \quad \text{ and } \quad (\partial u / \partial t) _{t=0} = g (x),
\]
$f$ and $g$ being given functions. This may be written in vector form as follows: 
\[
\frac {\partial}{\partial t} \begin{pmatrix}u \\ v\end{pmatrix} = \begin{pmatrix}0 & I \\ \Delta & 0\end{pmatrix} \begin{pmatrix}u \\ v\end{pmatrix}, v = \frac {\partial u}{\partial t}
\]
with the initial condition 
\[
\begin{pmatrix}u (x, 0) \\ v (x, 0)\end{pmatrix} = \begin{pmatrix}f (x) \\ g (x)\end{pmatrix}.
\]

So in a suitable function space, the wave equation is of the same form as the diffusion (or heat) equation---differentiation with respect to the time parameter on the left and another operator on the right---or again similar to the equation $\dd y/\dd t = \alpha y$. Since the solution in the last case is the exponential function, it is suggested that the heat equation and the wave equation may be solved by properly defining the exponential functions of the operators 
\[
\Delta \text{ and } \begin{pmatrix}0 & I \\ \Delta & 0\end{pmatrix}
\]
in suitable function spaces. This is the motivation for the application of the semi-group theory to Cauchy's problem. It is to be noted that the 

Schrödinger equation 
$i^{-1}\partial u / \partial t = Hu = (\Delta +U(x))u$, where $U(x)$ is a given function, 
gives another example of the equation of evolution of the form 
\begin{equation}\label{eq:XIV.0.1}
\partial u / \partial t = A u, t > 0,
\end{equation}
where $A$ is a not necessarily continuous, linear operator in a function space. 

The equation of the form \eqref{eq:XIV.0.1} may be called a temporally homogeneous equation of evolution. We may integrate such an equation by the semi-group theory. In the following three sections, we shall give typical examples of such an integration. We shall then expound the integration theory of the temporally inhomogeneous equation of evolution 
\begin{equation}\label{eq:XIV.0.2}
\partial u / \partial t = A (t) u, a < t < b.
\end{equation}


\yossection{1}{Integration of Diffusion Equations in \texorpdfstring{$L^{2}(\mathbb{R}^{m})$}{L2(Rm)}}
\label{sec:XIV.1}

Consider a diffusion equation 
\begin{equation}\label{eq:XIV.1.1}
\partial u / \partial t = A u, t > 0,
\end{equation}
where the differential operator 
\begin{equation}\label{eq:XIV.1.2}
A = a ^{ij} (x) \frac {\partial^{2}}{\partial x _{i} \partial x _{j}} + b ^{i} (x) \frac {\partial}{\partial x _{i}} + c (x) \left(a ^{ij} (x) = a ^{ji} (x)\right)
\end{equation}
is strictly elliptic in an $m$-dimensional euclidean space $\mathbb{R}^{m}$. We assume that the real-valued coefficients $a, b$ and $c$ are $C^{\infty}(\mathbb{R}^{m})$ functions and that 
\begin{equation}\label{eq:XIV.1.3}
\begin{gathered}
\max \left(\sup\limits _{x} | a ^{ij} (x) |, \sup\limits _{x} | b ^{i} (x) |, \sup\limits _{x} | c (x) |, \sup\limits _{x} | a _ {x _{k}} ^{ij} (x) |, \right. \left. \sup\limits _{x} | b _ {x _{k}} ^{i} (x) |, \sup\limits _{x} | a _ {x _{k} x _{8}} ^{ij} (x) |\right) = \eta < \infty .
\end{gathered}
\end{equation}

The strictly elliptic hypothesis concerning $A$ means that positive constants $\lambda_{0}$ and $\mu_{0}$ exist such that 
\begin{equation}\mu_{0} \sum\limits_{j=1} ^{m} \xi_{j} ^{2} \geqq a ^{ij} (x) \xi_{i} \xi_{j} \geqq \lambda_{0} \sum\limits_{i=1} ^{m} \xi_{i} ^{2} \text{ on } \mathbb{R}^{m} \text{ for any real } \text{ vector } \xi = (\xi_{1}, \xi_{2}, \dots , \xi_{m}) \label{eq:XIV.1.4}\end{equation}

Let $\hat{H}_0^1$ be the space of all real-valued $C_0^\infty (\mathbb{R}^{m})$ functions $f(x)$ normed by 
\begin{equation}\label{eq:XIV.1.5}
\left\| f \right\| _{1} = \left(\int_ {\mathbb{R}^{m}} f ^{2} \dd x + \sum\limits_{j=1} ^{m} \int_ {\mathbb{R}^{m}} f _ {x _{j}} ^{2} \dd x\right) ^{1/2},
\end{equation}
and let $H_0^1$ be the completion of $\hat{H}_0^1$ with respect to the norm $\|f\|_1$. Let similarly $H_0^0$ be the completion of $\hat{H}_0^1$ with respect to the norm 
\begin{equation}\label{eq:XIV.1.6}
\left\| f \right\| _{0} = \left(\int_ {\mathbb{R}^{m}} f ^{2} \dd x\right) ^{1/2}.
\end{equation}


We have thus introduced two real Hilbert spaces $H_0^1$ and $H_0^0$, and $H_0^1$ and $\hat{H}_0^1$ are $\|\|_0$-dense in $H_0^0$. We know, from the Proposition in Chapter~\ref{chap:I}, \ref{sec:I.10}, that $H_0^1$ coincides with the real Sobolev space $W^1(\mathbb{R}^{m})$; we know also that $H_0^0$ coincides with the real Hilbert space $L^2(\mathbb{R}^{m})$. We shall denote the scalar product in the Hilbert space $H_0^1$ (or in the Hilbert space $H_0^0$) by $(f,g)_1$ (or by $(f,g)_0$). 

To integrate equation \eqref{eq:XIV.1.1} in the complex Hilbert space $L^{2}(\mathbb{R}^{m})$ under condition \eqref{eq:XIV.1.3} and \eqref{eq:XIV.1.4}, we shall prepare some lemmas. These lemmas will also play important roles in the following sections. 

\begin{lemma}[concerning partial integration]\label{lem:XIV.1.Lemma.1} Let $f, g \in \tilde{H}_0^1$. We have then 
\begin{equation}\label{eq:XIV.1.7}
\begin{aligned}
&(A f, g) _{0} =  - \int_ {\mathbb{R}^{m}} a ^{ij} f _ {x _{i}} g _ {x _{j}} \dd x - \int_ {\mathbb{R}^{m}} a _ {x _{j}} ^{ij} f _ {x _{i}} g \dd x + \int_ {\mathbb{R}^{m}} b ^{i} f _ {x _{i}} g \dd x + \int_ {\mathbb{R}^{m}} c f g \dd x,
\end{aligned}
\end{equation}
that is, we may partially integrate, in $(Af, g)_{0}$, the terms containing the second order derivatives as if the integrated terms are zero. 
\end{lemma}

\begin{proof}
 By \eqref{eq:XIV.1.3} and the fact that $f$ and $g$ both belong to $H_0^1$, we see that $a^{ij}f_{x_ix_j}g$ is integrable over $\mathbb{R}^{m}$. Hence, by the fact that $f$ and $g$ are both of compact support, 
\[
\int_ {\mathbb{R}^{m}} a ^{ij} f _ {x _{i} x _{j}} g \dd x = - \int_ {\mathbb{R}^{m}} a ^{ij} f _ {x _{i}} g _ {x _{j}} \dd x - \int_ {\mathbb{R}^{m}} a _ {x _{j}} ^{ij} f _ {x _{i}} g \dd x. \]
\end{proof}

\begin{remarku} The formal adjoint $A^{*}$ of $A$ is defined by 
\begin{equation}\label{eq:XIV.1.8}
(A ^{*} f) (x) = \frac {\partial^{2}}{\partial x _{i} \partial x _{j}} \left(a ^{ij} (x) f (x)\right) - \frac {\partial}{\partial x _{i}} \left(b ^{i} (x) f (x)\right) + c (x) f (x).
\end{equation}
\end{remarku}

Then, as above, we have the result: If $f, g \in \tilde{H}_0^1$ then we may partially integrate, in $(A^*f, g)_0$, the terms containing the first and the second order derivatives as if the integrated terms are zero. That is, we have 
\begin{equation}\tag{7'}\label{eq:XIV.1.7'}
\begin{aligned}
&(A ^{*} f, g) _{0} =  - \int_ {\mathbb{R}^{m}} a ^{ij} f _ {x _{i}} g _ {x _{j}} \dd x - \int_ {\mathbb{R}^{m}} a _ {x _{i}} ^{ij} f g _ {x _{j}} \dd x - \int_ {\mathbb{R}^{m}} b ^{i} f g _ {x _{i}} \dd x + \int_ {\mathbb{R}^{m}} c f g \dd x.
\end{aligned}
\end{equation}

\begin{corollaryu} There exist positive constants $x$, $y$ and $\delta$ such that, for all sufficiently small positive constants $\alpha$, 
\begin{align}
\alpha \delta \| f \| _{1} ^{2}& \leqq (f - \alpha A f, f) _{0} \leqq (1 + \alpha \gamma) \| f \| _{1} ^{2} \quad \text{ when } \quad f \in \hat {H} _{0} ^{1}, \label{eq:XIV.1.9} \\
\alpha \delta \| f \| _{1} ^{2}& \leqq (f - \alpha A ^{*} f, f) _{0} \leqq (1 + \alpha \gamma) \| f \| _{1} ^{2} \quad \text{ when } \quad f \in \hat {H} _{0} ^{1}, \notag \\
\left| (f - \alpha A f, g) _{0} \right|& \leqq (1 + \alpha \gamma) \left\| f \right\| _{1} \left\| g \right\| _{1} \quad \text{ when } \quad f, g \in \hat {H} _{0} ^{1}, \label{eq:XIV.1.10} \\
\left| (f - \alpha A ^{*} f, g) _{0} \right|& \leqq (1 + \alpha \gamma) \left\| f \| _{1} \right. \| g \| _{1} \quad \text{ when } \quad f, g \in \hat {H} _{0} ^{1}, \notag \\
\left| (A f, g) _{0} - (f, A g) _{0} \right|& \leqq \varkappa \| f \| _{1} \| g \| _{0} \quad \text{ when } \quad f, g \in \hat {H} _{0} ^{1}. \label{eq:XIV.1.11}
\end{align}
\end{corollaryu}

\begin{proof}
 \eqref{eq:XIV.1.9} and \eqref{eq:XIV.1.10} may be proved by \eqref{eq:XIV.1.3}, \eqref{eq:XIV.1.4}, \eqref{eq:XIV.1.7} and \eqref{eq:XIV.1.7} remembering the inequality 
\begin{equation}\label{eq:XIV.1.12}
2 \alpha | a b | \leqq \alpha (\nu | a | ^{2} + \nu^{-1} | b | ^{2})
\end{equation}
which is valid whenever $\alpha$ and $\nu$ are $>0$. In fact, we can use an estimate such as 
\[
\left| \int_ {\mathbb{R}^{m}} a _ {x _{i}} ^{ij} f _ {x _{i}} g \dd x \right| \leqq \sum\limits_{i=1} ^{m} m \eta (v \| f _ {x _{i}} \| _{0} ^{2} + v ^{-1} \| g \| _{0} ^{2}). \]

We also obtain \eqref{eq:XIV.1.11} from 
\[
(A f, g) _{0} - (f, A g) _{0} = - \int_ {\mathbb{R}^{m}} \left(2 a _ {x _{i}} ^{ij} f _ {x _{j}} g + a _ {x _{i} x _{j}} ^{ij} f g - 2 b ^{i} f _ {x _{i}} g - b _ {x _{i}} ^{i} f g\right) \dd x.
\]
\end{proof}

\begin{lemma}[concerning the existence of solutions of $u - \alpha A u = f$]\label{lem:XIV.1.Lemma.2} Let a positive number $\alpha_0$ be so chosen that the above Corollary is valid for $0 < \alpha \leqq \alpha_0$. Then, for any function $f(x) \in \hat{H}_0^1$, the equation 
\begin{equation}\label{eq:XIV.1.13}
u - \alpha A u = f (0 < \alpha \leqq \alpha_{0})
\end{equation}
admits a uniquely determined solution $u \in H_{0}^{1} \cap C^{\infty}(\mathbb{R}^{m})$. 
\end{lemma}

\begin{proof}
 Let us define a bilinear functional $\hat{B}(u,v) = (u - \alpha A^{*}u,v)_{0}$ defined for functions $u, v \in \hat{H}_{0}^{1}$. From the above Corollary, we have 
\begin{equation}\label{eq:XIV.1.14}
\left| \hat {B} (u, v) \right| \leqq (1 + \alpha \gamma) \left\| u \right\| _{1} \left\| v \right\| _{1}, \alpha \delta_{r} \| u \| _{1} ^{2} \leqq \hat {B} (u, u). \end{equation}

Hence we may extend $\dot{B}(u,v)$, by continuity, to a bilinear functional $B(u,v)$ defined for $u,v\in H_{0}^{1}$ and satisfying 
\begin{equation}\tag{14'}\label{eq:XIV.1.14'}
\left| B (u, v) \right| \leqq (1 + \alpha \gamma) \left\| u \| _{1} \right\| v \| _{1}, \alpha \delta \| u \| _{1} ^{2} \leqq B (u, u).
\end{equation}

The linear functional $F(u) = (u, f)_0$ defined on $H_0^1$, is a bounded linear functional by $|(u, f)_0| \leqq \|u\|_0 \|f\|_0 \leqq \|u\|_1 \|f\|_0$. Hence, by F. \textsc{Riesz}' representation theorem in the Hilbert space $H_0^1$, there exists a uniquely determined $v = v(f) \in H_0^1$ such that $(u, f)_0 = (u, v(f))_1$. Thus, by the Milgram-Lax theorem applied to the Hilbert space $H_0^1$, we have 
\begin{equation}\label{eq:XIV.1.15}
(u, f) _{0} = (u, v (f)) _{1} = B (u, S v (f)) \text{ for all } u \in H _{0} ^{1}, \text{ where }
\end{equation}

$S$ is a bounded linear operator on $H_{0}^{1}$ onto $H_{0}^{1}$. 


Let $u$ run over $C_0^\infty(\mathbb{R}^{m})$, and let $v_n \in \hat{H}_0^1$ be such that $\lim\limits_{n \to \infty} \|v_n - Sv(f)\|_1 = 0$. Then 
\[
\begin{aligned}
&B (u, S v (f))= \lim\limits _{n \rightarrow \infty} B (u, v _{n}) = \lim\limits _{n \rightarrow \infty} \hat {B} (u, v _{n}) = \lim\limits _{n \rightarrow \infty} (u - \alpha A ^{*} u, v _{n}) _{0} = (u - \alpha A ^{*} u, S v (f)) _{0},
\end{aligned}
\]
because the norm $\|\|_1$ is larger than the norm $\|\|_0$. Hence 
\begin{equation}\tag{15'}\label{eq:XIV.1.15'}
(u, f) _{0} = (u - \alpha A ^{*} u, S v (f)) _{0},
\end{equation}
that is, $Sv(f) \in H_0^1$ is a distribution solution of the equation \eqref{eq:XIV.1.13}. Hence, by the strict ellipticity of $(I - \alpha A)$ and by the fact $f \in C_0^\infty(\mathbb{R}^{m})$, we see, from the Corollary of Friedrichs' theorem in Chapter~\ref{chap:VI}, \ref{sec:VI.9}, that we may consider $u = Sv(f) \in H_0^1$ to be a $C^\infty(\mathbb{R}^{m})$ solution of \eqref{eq:XIV.1.13}. Hence $u = Sv(f) \in H_0^1 \cap C^\infty(\mathbb{R}^{m})$. 


The uniqueness of such a solution $u$ of \eqref{eq:XIV.1.13} may be proved as follows. Let a function $u \in H_0^1 \cap C^\infty(\mathbb{R}^{m})$ satisfy $u - \alpha A u = 0$. Thus $Au \in H_0^1 \cap C^\infty(\mathbb{R}^{m}) \subseteq H_0^0$ and so the expression $(u - \alpha A u, u)_0$ is defined and $= 0$. Let $u_n \in \hat{H}_0^1$ be such that $\lim\limits_{n \to \infty} \|u - u_n\|_1 = 0$. We obtain, by partial integration as in \eqref{eq:XIV.1.9}, 
\[
0 = (u - \alpha A u, u) _{0} = \lim\limits _{n \rightarrow \infty} (u - \alpha A u, u _{n}) _{0} \geqq \alpha \delta \| u \| _{1} ^{2}, \text{ that is, } u = 0.
\]
\end{proof}

\begin{corollary}\label{cor:XIV.1.Corollary.1} Positive constants $\hat{\alpha}_{0}$ and $\eta_{0}$ exist such that, for any $f\in\hat{H}_{0}^{1}$, the equation 
\begin{equation}\label{eq:XIV.1.16}
\alpha u - A u = f \left(0 < \hat {\alpha} _{0} + \lambda_{0} + \eta_{0} \leqq \alpha\right)
\end{equation}
admits a uniquely determined solution $u = u_{f} \in H_{0}^{1} \cap C^{\infty}(\mathbb{R}^{m})$, and we have the estimate 
\begin{equation}\label{eq:XIV.1.17}
\left\| u _{f} \right\| _{0} \leqq (\alpha - \lambda_{0} - \eta_{0}) ^{-1} \left\| f \right\| _{0}.
\end{equation}
\end{corollary}

\begin{proof}
 By Schwarz' inequality we have 
\begin{equation}\label{eq:XIV.1.18}
\left\| (\alpha I - A) u \right\| _{0} \cdot \left\| u \right\| _{0} \geqq \left| \left((\alpha I - A) u, u\right) _{0} \right| \quad \text{ for } \quad u \in \hat {H} _{0} ^{1}. \end{equation}

By partial integration, we have 
\[
\begin{aligned}
&\left((\alpha I - A) u, u\right) _{0} = \alpha \| u \| _{0} ^{2} + \int_ {\mathbb{R}^{m}} a ^{ij} u _ {x _{i}} u _ {x _{j}} \dd x + \int_ {\mathbb{R}^{m}} a _ {x _{j}} ^{ij} u _ {x _{j}} u \dd x - \int_ {\mathbb{R}^{m}} b ^{i} u _ {x _{i}} u \dd x - \int_ {\mathbb{R}^{m}} c u u \dd x.
\end{aligned}
\]

Hence, we have, by \eqref{eq:XIV.1.3}, \eqref{eq:XIV.1.4} and \eqref{eq:XIV.1.12}, 
\[
\begin{aligned}
((\alpha I - A) u, u) _{0}& \geqq \alpha \| u \| _{0} ^{2} + \lambda_{0} (\| u \| _{1} ^{2} - \| u \| _{0} ^{2})
\\
&\quad \quad - m ^{2} \eta [ v (\| u \| _{1} ^{2} - \| u \| _{0} ^{2}) + v ^{-1} \| u \| _{0} ^{2} + m ^{-2} \| u \| _{0} ^{2} ]
\\
&= (\alpha - \lambda_{0} - m ^{2} \eta (v ^{-1} - v + m ^{-2})) \| u \| _{0} ^{2} + (\lambda_{0} - m ^{2} \eta v) \| u \| _{1} ^{2}.
\end{aligned}
\]

Thus by \eqref{eq:XIV.1.18} we have, for $\eta_0 = m^2\eta (\nu^{-1} - \nu +m^{-2})$ 
\begin{equation}\tag{17'}\label{eq:XIV.1.17'}
\|(\alpha I-A)u\|_{0}\geqq(\alpha-\lambda_{0}-\eta_{0})\|u\|_{0}\quad\text{whenever }u\in\hat{H}_{0}^{1}.
\end{equation}
by taking $\nu>0$ so small that $(\lambda_{0}-m^{2}\eta\nu)$ and $\eta_{0}$ are >$0$. We then take $\hat{\alpha}_{0}$ so large that for $\alpha\geqq\hat{\alpha}_{0}+\lambda_{0}+\eta_{0}$ we can apply Lemma~\ref{lem:XIV.1.Lemma.2} in solving \eqref{eq:XIV.1.16}. 

Since the solution $u = u_f \in H_0^1 \cap C^\infty(\mathbb{R}^{m})$ of \eqref{eq:XIV.1.16} is approximated by $\|\|_1$-norm by a sequence of functions $\in \hat{H}_0^1$, we obtain the estimate \eqref{eq:XIV.1.17} from \eqref{eq:XIV.1.18} and \eqref{eq:XIV.1.17'}. 
\end{proof}

\begin{corollary}\label{cor:XIV.1.Corollary.2} Consider $A$ as an operator on $D(A) = (\alpha I - A)^{-1} \hat{H}_0^1 \subseteq H_0^0$ into $H_0^0$. Then the smallest closed extension $\hat{A}$ in $H_0^0$ of $A$ admits, for $\alpha > \hat{\alpha}_0 + \lambda_0 + \eta_0$, the resolvent $(\alpha I - \hat{A})^{-1}$ defined on $H_0^0$ into $H_0^0$ such that 
\begin{equation}\label{eq:XIV.1.19}
\left\| (\alpha I - \hat {A}) ^{-1} \right\| _{0} \leqq (\alpha - \lambda_{0} - \eta_{0}) ^{-1}.
\end{equation}
\end{corollary}

\begin{proof}
 Clear from Corollary~\ref{cor:XIV.1.Corollary.1} remembering the fact that $D(A)$ and $\hat{H}_0^1$ both are $\|\|_0$-dense in $H_0^0$. 
\end{proof}

Thus, by Corollary~\ref{cor:IX.7.Corollary.2} in Chapter~\ref{chap:IX}, \ref{sec:IX.7}, $\hat{A}$ is the infinitesimal generator of a semi-group $T_{t}$ of class $(C_0)$ in the $B$-space $H_0^0$ such that $\| T_t \|_0 \leqq e^{(\lambda_0+ \eta_0)t}$ for $t \geqq 0$. 

Actually we can prove 

\begin{theorem}\label{thm:XIV.1.Theorem.1} Let Complex- $\hat{H}_0^1$ be the space of all complex-valued functions $f \in C_0^\infty(\mathbb{R}^{m})$ with $\|f\|_1 = \left( \int_{\mathbb{R}^{m}} |f|^2 \dd x + \sum\limits_{j=1}^{m} \int_{\mathbb{R}^{m}} |f_{x_{j}}|^2 \dd x \right)^{1/2} < \infty$. 
\end{theorem}

Let $\tilde{H}_0^1$ and $\tilde{H}_0^0$ be the completions of Complex- $\hat{H}_0^1$ with respect to the norm $\|f\|_1$ and $\|f\|_0$, respectively. We know that $\tilde{H}_0^1 = (\text{ complex })$ Sobolev space $W^1(\mathbb{R}^{m})$ (see Chapter~\ref{chap:I}, \ref{sec:I.10}). It is also clear that $\tilde{H}_0^0 = (\text{ complex })$ Hilbert space $L^2(\mathbb{R}^{m})$. We consider $A$ as an operator defined on $D(A) = (\alpha I - A)^{-1}\tilde{H}_0^1 \subseteq L^2(\mathbb{R}^{m})$ into $L^2(\mathbb{R}^{m})$. Then the smallest closed extension $\tilde{A}$ in $L^2(\mathbb{R}^{m})$ of $A$ is the infinitesimal generator of a holomorphic semi-group $T_t$ of class $(C_0)$ in $L^2(\mathbb{R}^{m})$ such that $\|T_t\|_0 \leqq e^{(\lambda_0+ \eta_0)t}$ for $t \geqq 0$. 

\begin{proof}
 By the preceding Corollary~\ref{cor:XIV.1.Corollary.2} and the reality of the coefficients in the differential operator $A$, we see that the range $R(\alpha I - A) = (\alpha I - A) \cdot D(A)$ is $\|\|_0$-dense in $L^2(\mathbb{R}^{m})$ for $\alpha > \hat{\alpha}_0 + \lambda_0 + \eta_0$. Moreover, we have, for $(u + iv) \in L^2(\mathbb{R}^{m})$ such that $(u + iv) \in D(A)$, 
\begin{equation}\left\| (\alpha I - A) (u + i v) \right\| _{0} ^{2} = \left\| (\alpha I - A) u \right\| _{0} ^{2} + \left\| (\alpha I - A) v \right\| _{0} ^{2} \geqq (\alpha - \lambda_{0} - \eta_{0}) ^{2} (\| u \| _{0} ^{2} + \| v \| _{0} ^{2}). \notag\end{equation}

Hence the inverse $(\alpha I - \tilde{A})^{-1}$ is a bounded linear operator on $L^2 (\mathbb{R}^{m})$ into $L^2 (\mathbb{R}^{m})$ such that $\|(\alpha I - \tilde{A})^{-1}\|_0 \leqq (\alpha - \lambda_0 - \eta_0)^{-1}$ for $\alpha > \hat{\alpha}_0 + \lambda_0 + \eta_0$. 

Therefore, by the Theorem in Chapter~\ref{chap:IX}, \ref{sec:IX.10}, we have only to show that 
\begin{equation}\label{eq:XIV.1.20}
\varlimsup\limits_{| \tau | \uparrow \infty} | \tau | \| ((\alpha + i \tau) I - \tilde {A}) ^{-1} \| _{0} < \infty .
\end{equation}

We have, for $w \in D(A)$, 
\begin{equation}\label{eq:XIV.1.21}
\left\| ((\alpha + i \tau) I - A) w \right\| _{0} \left\| w \right\| _{0} \geqq \left| \left(\left((\alpha + i \tau) I - A\right) w, w\right) _{0} \right|.
\end{equation}

By partial integration, we obtain, as in the proof of \eqref{eq:XIV.1.17}, 
\begin{gather}\quad \left| \text{ Real part } \left(\left((\alpha + i \tau) I - A\right) w, w\right) _{0} \right| = \bigl\| \alpha | w | _{0} ^{2} + \text{ Real part } \left(\int_ {\mathbb{R}^{m}} a ^{ij} w _ {x _{i}} \overline {{w}} _ {x _{j}} \dd x \right.
\\\quad \left. + a _ {x _{i}} ^{ij} w _ {x _{j}} \overline {{w}} \dd x - \int_ {\mathbb{R}^{m}} b ^{i} w _ {x _{i}} \overline {{w}} \dd x - \int_ {\mathbb{R}^{m}} c w \overline {{w}} \dd x\right) \bigr| \geqq (\alpha - \lambda_{0} - \eta_{0}) \| w \| _{0} ^{2} \notag
\\\quad + (\lambda_{0} - m ^{2} \eta \nu) \| w \| _{1} ^{2}. \notag
\end{gather}

Similarly we have 
\begin{gather}\left| \text{ Imaginary part } \left(\left((\alpha + i \tau) I - A\right) w, w\right) _{0} \right| \notag \\\geqq | \tau | \| w \| _{0} ^{2} - m ^{2} \eta (\| w \| _{1} ^{2} + m ^{-2} \| w \| _{0} ^{2}) | \geqq (| \tau | - \eta) \| w \| _{0} ^{2} - m ^{2} \eta \| w \| _{1} ^{2}. \notag
\end{gather}

Suppose there exists a $w_0 \in D(A)$, $\|w_0\|_0 \neq 0$, such that 
$\left|\operatorname{Im}\left(\left(\left(\alpha I+i\tau\right)I-A\right)w_{0},w_{0}\right)_{0}\right|<2^{-1}\left(\left|\tau\right|-\eta\right)\left\|w_{0}\right\|_{0}^{2}$ for sufficiently large $\tau$ (or for sufficiently large $-\tau$). Then, for such large $\tau$ (or $-\tau$), we must have 
\[
m ^{2} \eta \| w _{0} \| _{1} ^{2} > 2 ^{-1} (| \tau | - \eta) \| w _{0} \| _{0} ^{2}
\]
and so 
$\left|\left(\text{ Real part } \left(\left(\alpha+i\tau\right)I-A\right)w_{0},w_{0}\right)_{0}\right|\geqq\left(\lambda_{0}-m^{2}\eta\nu\right)\frac{\left(\left|\tau\right|-\eta\right)}{2m^{2}\eta}\left\|w_{0}\right\|_{0}^{2}.$ Hence we have proved Theorem~\ref{thm:XIV.1.Theorem.1} by \eqref{eq:XIV.1.21}. 
\end{proof}

\begin{theorem}\label{thm:XIV.1.Theorem.2} For any $f \in L^2(\mathbb{R}^{m})$, $u(t, x) = (T_t f)(x)$ is infinitely differentiable in $t > 0$ and in $x \in \mathbb{R}^{m}$, and $u(t, x)$ satisfies the diffusion equation \eqref{eq:XIV.1.1} as well as the initial condition $\lim\limits_{t \downarrow 0} \|u(t, x) - f(x)\|_0 = 0$. 
\end{theorem}

\begin{proof}
 We denote by $T_{t}^{(k)}$ the $k$-th strong derivative in $L^2(\mathbb{R}^{m})$ of $T_{t}$ with respect to $t$. $T_{t}$ being a holomorphic semi-group of class $(C_0)$ in $L^2(\mathbb{R}^{m})$, we have $T_{t}^{(k)}f = \tilde{A}^{k}T_{t}f\in L^{2}(\mathbb{R}^{m})$ if $t > 0$ ( $k = 0,1,\ldots$). Since $\tilde{A}$ is the smallest closed extension in $L^2(\mathbb{R}^{m})$ of $A$, we see that $A^{k}T_{t}f\in L^{2}(\mathbb{R}^{m})$ for fixed $t > 0$ ( $k = 0,1,2,\ldots$) if we apply the differential operator $A^{k}$ in the distributional sense. Thus, by the Corollary of Friedrichs' theorem in Chapter~\ref{chap:VI}, \ref{sec:VI.9}, $u(t,x)$ is, for fixed $t > 0$, equal to a $C^\infty(\mathbb{R}^{m})$ function after a correction on a set of measure zero. 


Because of the estimate $\left\|T_{t}\right\|_{0}\leqq e^{(\lambda_{0}+\eta_{0})t}$, we easily see that, if we apply, in the distributional sense, the elliptic differential operator 
\[
\left(\frac {\partial^{2}}{\partial t ^{2}} + A\right)
\]
any number of times to $u(t, x)$, then the result is locally square integrable in the product space $\{t; 0 < t < \infty\} \times \mathbb{R}^{m}$. Thus again, by the Corollary of Friedrichs' theorem in Chapter~\ref{chap:VI}, \ref{sec:VI.9}, we see that $u(t, x)$ is equal to a function which is $C^\infty$ in $(t, x)$, $t > 0$ and $x \in \mathbb{R}^{m}$, after a correction on a set of measure zero of the product space. Thus we may consider that $u(t, x)$ is a genuine solution of \eqref{eq:XIV.1.1} satisfying the initial condition $\lim\limits_{t \downarrow 0} \| u(t, x) - f(x) \|_0 = 0$.
\end{proof}

\begin{remarku} The above obtained solution $u(t, x)$ satisfies the ``forward and backward unique continuation property'': 
\end{remarku}

if, for a fixed $t_0 > 0$, $u(t_0, x) \equiv 0$ on an open set $G \subseteq \mathbb{R}^{m}$; 


then $u(t, x) = 0$ for every $t > 0$ and every $x \in G$. 

\begin{proof}
 Since $T_{t}$ is a holomorphic semi-group of class $(C_0)$, we have 
\[
\lim\limits _{n \rightarrow \infty} \left\| T _ {t _{0} + h} f - \sum\limits_{k=0} ^{n} (k!) ^{-1} h ^{k} A ^{k} T _ {t _{0}} f \right\| _{0} = 0
\]
for a sufficiently small $h$. Hence, as in the proof of the completeness of the space $L^2(\mathbb{R}^{m})$, there exists a sequence $\{n'\}$ of natural numbers such that 
\[
u \left(t _{0} + h, x\right) = \lim\limits _ {n ^{\prime} \rightarrow \infty} \sum\limits_{k=0} ^ {n ^{\prime}} (k!) ^{-1} h ^{k} A ^{k} u \left(t _{0}, x\right) \text{ for a.e. } x \in \mathbb{R}^{m}. \]

By hypothesis given in ($22$), we have $A^{k}u(t_{0}, x) = 0$ in $G$ and so we must have $u(t_{0} + h, x) = 0$ in $G$ for sufficiently small $h$. Repeating the process we see that conclusion ($22$) is true. 
\end{proof}

\yossub{References}

The result of this section is adapted from K. \textsc{Yosida} \cite{yosida364}. As for the ``forward and backward unique continuation property'' ($22$), we have the more precise result that $u(t, x) = 0$ for every $t > 0$ and every $x \in \mathbb{R}^{m}$. For, S. \textsc{Mizohata} \cite{mizohata210} has proved a ``space-like unique continuation property'' of solutions $u(t, x)$ of a diffusion equation to the effect that $u(t, x) = 0$ for all $t > 0$ and all $x \in \mathbb{R}^{m}$ if $u(t, x) = 0$ for all $t > 0$ and all $x \in G$. Concerning the holomorphic character of the semi-group $T_{t}$ obtained above, see also R. S. \textsc{Phillips} \cite{phillips256}. It is to be noted here that the unique continuation property of the solutions of a heat equation $\partial u/\partial t = \Delta u$ was first proposed and solved by H. \textsc{Yamabe}-S. ITô \cite{yamabe342}. 

There is a fairly complete discussion of parabolic equations from the view point of the theory of dissipative operators. See R. S. \textsc{Phillips} \cite{phillips257}. 


\yossection{2}{Integration of Diffusion Equations in a Compact Riemannian Space}
\label{sec:XIV.2}

Let $R$ be a connected, orientable $m$-dimensional $C^\infty$ Riemannian space with the metric 
\begin{equation}\label{eq:XIV.2.1}
\dd s ^{2} = g _{ij} (x) \dd x ^{i} \dd x ^{j}.
\end{equation}

Let $A$ be a second order linear partial differential operator in $R$ with real-valued $C^\infty$ coefficients: 
\begin{equation}\label{eq:XIV.2.2}
A = a ^{ij} (x) \frac {\partial^{2}}{\partial x ^{i} \partial x ^{j}} + b ^{i} (x) \frac {\partial}{\partial x ^{i}}.
\end{equation}

We assume that $a^{ij}$ is a symmetric contravariant tensor and that $b^{i}(x)$ satisfies the transformation rule 
\begin{equation}\label{eq:XIV.2.3}
\bar {b} ^{i} = b ^{k} \frac {\partial \bar {x} ^{i}}{\partial x ^{k}} + a ^{kj} \frac {\partial^{2} \bar {x} ^{i}}{\partial x ^{k} \partial x ^{j}}
\end{equation}
by the coordinate change $(x^{1}, x^{2}, \ldots, x^{m}) \to (\bar{x}^{1}, \bar{x}^{2}, \ldots, \bar{x}^{m})$ so that the value $(A f)(x)$ is determined independently of the choice of the local coordinates. We further assume that $A$ is strictly elliptic in the sense that there exist positive constants $\lambda_{0}$ and $\mu_{0}$ such that 
\begin{equation}\mu_{0} \sum\limits_{j=1} ^{m} \xi_{j} ^{2} \geqq a ^{ij} (x) \xi_{i} \xi_{j} \geqq \lambda_{0} \sum\limits_{j=1} ^{m} \xi_{j} ^{2} \text{ for every real vector } \left(\xi_{1}, \dots , \xi_{m}\right) \text{ and every } x \in \mathbb{R}^{m}. \label{eq:XIV.2.4}\end{equation}

We consider the Cauchy problem in the large on $R$ for the diffusion equation: to find solution $u(t, x)$ such that 
\begin{equation}\label{eq:XIV.2.5}
\begin{aligned}
&\partial u / \partial t  = A u, t > 0, u (0, x) = f (x) \text{ where } f (x) \text{ is a given } \text{ function on } R.
\end{aligned}
\end{equation}

We shall prove 

\begin{theoremu} If $R$ is compact so that $R$ is without boundary, then the equation \eqref{eq:XIV.2.5} admits, for any initial function $f \in C^{\infty}(R)$, a uniquely determined solution $u(t, x)$ which is $C^{\infty}$ in $(t, x)$, $t > 0$ and $x \in R$. This solution can be represented in the form 
\begin{equation}\label{eq:XIV.2.6}
u (t, x) = \int_{R} P (t, x, \dd y) f (y)
\end{equation}
where $P(t, x, E)$ is the transition probability of a Markov process on $R$. \end{theoremu}

\begin{proof}
Let us denote by $C(R)$ the $B$-space of real-valued continuous functions $f(x)$ on $R$ normed by $\|f\| = \sup |f(x)|$. We first prove 
\begin{align}
\text{ for any } f& \in C ^{\infty} (R) \text{ and any } n > 0, \text{ we have } \notag \\
\max\limits _{x} h (x)& \geqq f (x) \geqq \min\limits _{x} h (x) \text{ where } h (x) = f (x) - n ^{-1} (A f) (x). \label{eq:XIV.2.7}
\end{align}


Suppose that $f(x)$ attains its maximum at $x = x_0$. We choose a local coordinate system at $x_0$ such that $a^{ij}(x_0) = \delta_{ij} (= 1 \text{ or } 0 \text{ according as } i = j \text{ or } i \neq j)$. Such a choice is possible owing to condition \eqref{eq:XIV.2.4}. Then 
\[
\begin{aligned}
&h (x _{0})  = f (x _{0}) - n ^{-1} (A f) (x _{0}) = f (x _{0}) - n ^{-1} b ^{i} (x _{0}) \frac {\partial f}{\partial x _{0} ^{i}} - n ^{-1} \sum\limits_{i=1} ^{m} \frac {\partial^{2} f}{\partial (x _{0} ^{i}) ^{2}} \geqq f (x _{0}),
\end{aligned}
\]
since we have, at the maximum point $x_0$, 
\[
\frac {\partial f}{\partial x _{0} ^{i}} = 0 \quad \text{ and } \quad \frac {\partial^{2} f}{\partial (x _{0} ^{i}) ^{2}} \leqq 0.
\]

Thus $\max h(x) \geqq f(x)$. Similarly we have $f(x) \geqq \min h(x)$. 

We shall consider $A$ as an operator on $D(A) = C^{\infty}(R) \subseteq C(R)$ into $C(R)$. Then, by \eqref{eq:XIV.2.7}, we see that the inverse $(I - n^{-1}A)^{-1}$ exists for $n > 0$ and $\|(I - n^{-1}A)^{-1}g\| \leqq \|g\|$ for $g$ in the range $R(I - n^{-1}A) =$ $(I - n^{-1}A) \cdot D(A)$. This range is, for sufficiently large $n$, strongly dense in $C(R)$. For, as in the preceding section, we have the result: For any $g \in C^{\infty}(R)$ and for sufficiently large $n > 0$, the equation $u - n^{-1}Au = g$ admits a uniquely determined solution $u \in C^{\infty}(R)$. Because, by the compactness of the Riemann space $R$, Lemma~\hyperref[lem:XIV.1.Lemma.1]{$1$} concerning the partial integration in the preceding section may be adapted to our Riemannian space $R$ without boundary. Moreover, $C^{\infty}(R)$ is strongly dense in $C(R)$ as may be seen by the regularization of functions in $C(R)$ (see Proposition~\ref{prop:I.1.Proposition.8} in Chapter~\ref{chap:I}, \ref{sec:I.1}). 

Hence the smallest closed extension $\overline{A}$ in $C(R)$ of the operator $A$ satisfies the conditions: 

for sufficiently large $n > 0$, the resolvent $(I - n^{-1}\overline{A})^{-1}$ 
exists as a bounded linear operator on $C(R)$ into $C(R)$ 
\begin{align}
\text{ such that } \left\| (I - n ^{-1} \overline {{A}}) ^{-1} \right\|& \leqq 1, \label{eq:XIV.2.8} \\
((I - n ^{-1} \bar {A}) ^{-1} h) (x)& \geqq 0 \text{ on } R \text{ whenever } h (x) \geqq 0 \text{ on } R, \label{eq:XIV.2.9} \\
(I - n ^{-1} \overline {{{A}}}) ^{-1} \cdot 1& = 1. \label{eq:XIV.2.10}
\end{align}

The positivity of the operator $(I - n^{-1}\overline{A})^{-1}$ given in \eqref{eq:XIV.2.9} is clear from \eqref{eq:XIV.2.7}. The equation \eqref{eq:XIV.2.10} follows from $A\cdot 1 = 0$. 

Therefore $\overline{A}$ is the infinitesimal generator of a contraction semigroup $T_{t}$ in $C(R)$ of class $(C_0)$. As in the preceding section, we see, by the strict ellipticity of $A$, that, for any $f \in C^{\infty}(R)$, the function $u(t,x) = (T_{t}f)(x)$ is a $C^{\infty}$ function in $(t,x)$ for $t > 0$ and $x \in R$ so that $u(t,x)$ is a genuine solution of \eqref{eq:XIV.2.5}. 

Since the dual space of the space $C(R)$ is the space of Baire measures in $R$, we easily prove the latter part of the Theorem remembering \eqref{eq:XIV.2.9} and \eqref{eq:XIV.2.10}. 
\end{proof}

\yossection{3}{Integration of Wave Equations in a Euclidean Space \texorpdfstring{$\mathbb{R}^{m}$}{Rm}}
\label{sec:XIV.3}

Consider a wave equation 
\begin{equation}\label{eq:XIV.3.1}
\partial^{2} u / \partial t ^{2} = A u, - \infty < t < \infty ,
\end{equation}
where the differential operator 
\begin{equation}\label{eq:XIV.3.2}
A = a ^{ij} (x) \frac {\partial^{2}}{\partial x _{i} \partial x _{j}} + b ^{i} (x) \frac {\partial}{\partial x ^{i}} + c (x) \quad \left(a ^{ij} (x) = a ^{ji} (x)\right)
\end{equation}
is strictly elliptic in an $m$-dimensional Euclidean space $\mathbb{R}^{m}$. We assume that the real-valued $C^{\infty}$ coefficients $a, b$, and $c$ satisfy conditions \eqref{eq:XIV.3.3} and \eqref{eq:XIV.1.4} in Chapter~\ref{chap:XIV}, \ref{sec:XIV.1}. As done there, we shall denote by $\hat{H}_{0}^{1}$ the space of all real-valued $C_{0}^{\infty}(\mathbb{R}^{m})$ functions $f(x)$ normed by 
\[
\left\| f \right\| _{1} = \left(\int_ {\mathbb{R}^{m}} f ^{2} \dd x + \sum\limits_{j=1} ^{m} \int_ {\mathbb{R}^{m}} f _ {x _{j}} ^{2} \dd x\right) ^{1/2}
\]
and let $H_0^1$ (and $H_0^0$) be the completion of $\hat{H}_0^1$ with respect to the norm $\|f\|_1$ (and with respect to the norm $\|f\|_0 = \left( \int_{\mathbb{R}^{m}} f^2 \dd x \right)^{1/2}$). 

\begin{lemmau} For any pair $\{f, g\}$ of elements $\in \hat{H}_0^{1'}$ the equation 
\begin{equation}\label{eq:XIV.3.3}
\left(\begin{pmatrix}I & 0 \\ 0 & I\end{pmatrix} - n ^{-1} \begin{pmatrix}0 & I \\ A & 0\end{pmatrix}\right) \begin{pmatrix}u \\ v\end{pmatrix} = \begin{pmatrix}f \\ g\end{pmatrix}
\end{equation}
admits, if the integer $n$ be such that $|n^{-1}|$ is sufficiently small, a uniquely determined solution $\{u, v\}$, $u$ and $v \in H_0^1 \cap C^\infty(\mathbb{R}^{m})$, satisfying the estimate 
\begin{equation}\label{eq:XIV.3.4}
\begin{gathered}
\left((u - \alpha_{0} A u, u) _{0} + \alpha_{0} (v, v) _{0}\right) ^{1/2} \leqq (1 - \beta | n | ^{-1}) ^{-1} \left((f - \alpha_{0} A f, f) _{0} + \alpha_{0} (g, g) _{0}\right) ^{1/2},
\end{gathered}
\end{equation}
with positive constants $\alpha_0$ and $\beta$ which are independent of $n$ and $\{f, g\}$. \end{lemmau}

\begin{proof}
Let $u_1 \in H_0^1 \cap C^\infty(\mathbb{R}^{m})$ and $v_1 \in H_0^1 \cap C^\infty(\mathbb{R}^{m})$ be respectively the solutions of 
\begin{equation}\label{eq:XIV.3.5}
u _{1} - n ^{-2} A u _{1} = f \quad \text{ and } \quad v _{1} - n ^{-2} A v _{1} = g.
\end{equation}


The existence of such solutions for sufficiently small $|n^{-1}|$ was proved in Lemma~\ref{lem:XIV.1.Lemma.2} in Chapter~\ref{chap:XIV}, \ref{sec:XIV.1}. Then 
\begin{equation}\label{eq:XIV.3.6}
u = u _{1} + n ^{-1} v _{1}, v = n ^{-1} A u _{1} + v _{1}
\end{equation}
satisfies \eqref{eq:XIV.3.3}, i.e., we have $u - n^{-1}v = f$, $v - n^{-1}Au = g$. 

We next prove \eqref{eq:XIV.3.4}. We remark that 
$A u = n (v - g) \in H_0^1 \cap C^\infty(\mathbb{R}^{m}) \subseteq H_0^0$ and hence, by $f, g \in C_0^\infty(\mathbb{R}^{m})$, 
\[
A v = n (A u - A f) \in H _{0} ^{1} \cap C ^{\infty} (\mathbb{R}^{m}) \subseteq H _{0} ^{0}.
\]

We have thus, by \eqref{eq:XIV.3.3}, 
and 
\begin{gather}(f - \alpha_{0} A f, f) _{0}  = (u - n ^{-1} v - \alpha_{0} A (u - n ^{-1} v), u - n ^{-1} v) _{0} \notag
\\= (u - \alpha_{0} A u, u) _{0} - 2 n ^{-1} (u, v) _{0} + \alpha_{0} n ^{-1} (A u, v) _{0} + \alpha_{0} n ^{-1} (A v, u) _{0} + n ^{-2} (v - \alpha_{0} A v, v) _{0} \notag
\\\alpha_{0} (g, g) _{0}  = \alpha_{0} (v - n ^{-1} A u, v - n ^{-1} A u) _{0} \notag
\\= \alpha_{0} (v, v) _{0} - \alpha_{0} n ^{-1} (v, A u) _{0} - \alpha_{0} n ^{-1} (A u, v) _{0} + \alpha_{0} n ^{-2} (A u, A u) _{0}. \notag
\end{gather}

By a limiting process, we can prove that \eqref{eq:XIV.3.9}, \eqref{eq:XIV.3.10} and \eqref{eq:XIV.1.11} in Chapter~\ref{chap:XIV}, \ref{sec:XIV.1} are valid for $f = u$ and $g = v$. Thus, by \eqref{eq:XIV.1.12} in Chapter~\ref{chap:XIV}, \ref{sec:XIV.1}, there exists a positive constant $\beta$ satisfying 
\[
\begin{aligned} ((f - \alpha_{0} A f, f) _{0} + \alpha_{0} (g, g) _{0}) ^{1/2} & \geqq ((u - \alpha_{0} A u, u) _{0} + \alpha_{0} (v, v) _{0} \\ & - \alpha_{0} | n ^{-1} | | (A u, v) _{0} - (A v, u) _{0} | - 2 | n ^{-1} | | (u, v) _{0} |) ^{1/2} \\ & \geqq (1 - \beta | n ^{-1} |) ((u - \alpha_{0} A u, u) _{0} + \alpha_{0} (v, v) _{0}) ^{1/2} \end{aligned}
\]
for sufficiently large $|\pmb{n}|$. 

The above estimate for the solutions $\{u, v\}$ belonging to $H_0^1 \cap C^\infty(\mathbb{R}^{m})$ shows that such solutions are uniquely determined by $\{f, g\}$. 
\end{proof}

\begin{corollaryu} The product space $H_{0}^{1} \cap H_{0}^{0}$ of vectors 
\begin{equation}\label{eq:XIV.3.7}
\begin{pmatrix}u \\ v\end{pmatrix} = \{u, v \} ^{\prime}, \text{ where } u \in H _{0} ^{1} \text{ and } v \in H _{0} ^{0},
\end{equation}
is a $B$-space by the norm 
\begin{equation}\label{eq:XIV.3.8}
\left\| \begin{pmatrix}u \\ v\end{pmatrix} \right\| = \| \{u, v \} ^{\prime} \| = (B (u, u) + \alpha_{0} (v, v) _{0}) ^{1/2},
\end{equation}
where $B(f, g)$ is the extension by continuity with respect to the norm $\|\|_1$ of the bilinear functional 
\[
\hat {B} (f, g) = (f - \alpha_{0} A f, g) _{0} \text{ defined for } f, g \in \hat {H} _{0} ^{1}.
\]
\end{corollaryu}

We know that $B(u, u)^{1/2}$ is equivalent to the norm $\|u\|_1$ (see Chapter~\ref{chap:XIV}, \ref{sec:XIV.1}): 
\begin{equation}\label{eq:XIV.3.9}
\alpha_{0} \delta \| u \| _{1} ^{2} \leqq B (u, u) \leqq (1 + \alpha_{0} \gamma) \| u \| _{1} ^{2}.
\end{equation}

Let the domain $D(\mathfrak{A})$ of the operator 
\begin{equation}\label{eq:XIV.3.10}
\mathfrak {A} = \begin{pmatrix}0 & I \\ A & 0\end{pmatrix}
\end{equation}
be the vectors $\{u, v\}' \in H_0^1 \times H_0^0$ such that $u, v \in H_0^0$ are given by \eqref{eq:XIV.3.6}. Then the Lemma shows that the range of the operator 
\[
\operatorname{Im} - n ^{-1} \mathfrak {A}, \text{ where } \operatorname{Im} = \begin{pmatrix}I & 0 \\ 0 & I\end{pmatrix},
\]
contains all the vectors $\{f, g\}'$ such that $f, g \in H_0^1$. Thus the smallest closed extension $\overline{\mathfrak{A}}$ in $H_0^1 \times H_0^0$ of the operator $\mathfrak{A}$ is such that the operator $(\mathfrak{I} - n^{-1} \overline{\mathfrak{A}})$ with integral parameter $n$ admits, for sufficiently large $|n|$, an everywhere in $H_0^1 \times H_0^0$ defined inverse $(\mathfrak{I} - n^{-1} \overline{\mathfrak{A}})^{-1}$ satisfying 
\begin{equation}\label{eq:XIV.3.11}
\left\| (\operatorname{Im} - n ^{-1} \overline {{{\mathfrak {A}}}}) ^{-1} \right\| \leqq (1 - \beta | n ^{-1} |) ^{-1}.
\end{equation}

We are now prepared to prove 

\begin{theoremu} For any pair $\{f(x), g(x)\}$ of $C_0^\infty(\mathbb{R}^{m})$ functions, the equation \eqref{eq:XIV.3.1} admits a $C^\infty$ solution $u(t, x)$ satisfying the initial condition 
\begin{equation}\label{eq:XIV.3.12}
u (0, x) = f (x), u _{t} (0, x) = g (x)
\end{equation}
and the estimate 
\begin{equation}\label{eq:XIV.3.13}
(B (u, u) + \alpha_{0} (u _{t}, u _{t}) _{0}) ^{1/2} \leqq \exp (\beta | t |) (B (f, f) + \alpha_{0} (g, g) _{0}) ^{1/2}.
\end{equation}
\end{theoremu}

\begin{remarku} Formula \eqref{eq:XIV.3.9} shows that $B(u, u)$ is comparable to the potential energy of the wave (= the solution of \eqref{eq:XIV.3.1}) $u(t, x)$, and $(u_t, u_t)_0$ is comparable to the kinetic energy of the wave $u(t, x)$. Thus \eqref{eq:XIV.3.13} means that the total energy of the wave $u(t, x)$ does not increase more rapidly than $\exp(\beta |t|)$ when the time $t$ tends to $\pm\infty$. This is a kind of energy inequality which governs wave equations in general. 
\end{remarku}

\begin{proof}[Proof of Theorem.]
Estimate \eqref{eq:XIV.3.11} shows that $\overline{\mathfrak{A}}$ is the infinitesimal generator of a group $T_{t}$ of class $(C_0)$ in $H_0^1 \times H_0^0$ such that
\begin{equation}\label{eq:XIV.3.14}
\left\| T _{t} \right\| \leqq \exp (\beta | t |), - \infty < t < \infty .
\end{equation}

By hypothesis we have, for $k = 0, 1, 2, \ldots$, 
\[
\overline {{\mathfrak {A}}} ^{k} \begin{pmatrix}f \\ g\end{pmatrix} = \mathfrak {A} ^{k} \begin{pmatrix}f \\ g\end{pmatrix} \in C _{0} ^{\infty} (\mathbb{R}^{m}) \times C _{0} ^{\infty} (\mathbb{R}^{m}) \subseteq H _{0} ^{1} \times H _{0} ^{0}.
\]

Hence, if we put 
\[
\begin{pmatrix}u (t, x) \\ v (t, x)\end{pmatrix} = T _{t} \begin{pmatrix}f (x) \\ g (x)\end{pmatrix},
\]
then, by the commutativity of $\overline{A}$ with $T_{t}$, we have 
\[
\frac {\partial^{k}}{\partial t ^{k}} \begin{pmatrix}u (t, x) \\ v (t, x)\end{pmatrix} = \frac {\partial^{k} T _{t}}{\partial t ^{k}} \begin{pmatrix}f (x) \\ g (x)\end{pmatrix} = \mathfrak {A} ^{k} \begin{pmatrix}u (t, x) \\ v (t, x)\end{pmatrix} \in H _{0} ^{1} \times H _{0} ^{0}
\]
for $k = 0, 1, 2, \ldots$. Here we denote by $\partial^k T_t / \partial t^k$ the $k$ th strong derivative in $H_0^1 \times H_0^0$. Therefore, by $H_0^1 \subseteq H_0^0 = L^2(\mathbb{R}^{m})$ and the strict ellipticity of the operator $A$, we see, as in the proof of Theorem~\ref{thm:XIV.1.Theorem.2} in Chapter~\ref{chap:XIV}, \ref{sec:XIV.1}, that $u(t, x)$ is $C^\infty$ in $(t, x)$ for $-\infty < t < \infty$, $x \in \mathbb{R}^{m}$, and satisfies equation \eqref{eq:XIV.3.1} with \eqref{eq:XIV.3.12} and estimate \eqref{eq:XIV.3.13}. 
\end{proof}

\begin{remarku} The result of the present section is adapted from K. \textsc{Yosida} \cite{yosida365}. Cf. J. L. \textsc{Lions} \cite{liggett178}. P. D. \textsc{Lax} has kindly communicated to the present author that the method of integration given in this section is very similar to that announced by him in Abstract $180$, Bull. Amer. Math. Soc. $58$, $192$ ($1952$). It is to be noted here that our method can be modified to the integration of wave equations in an open domain of a Riemannian space. There is another approach to the integration of wave equations based upon the theory of dissipative semi-groups. See R. S. \textsc{Phillips} \cite{phillips258} and \cite{phillips259}. The method is closely connected with the theory of symmetric positive systems by K. \textsc{Friedrichs} \cite{friedrichs70}. Cf. also P. \textsc{Lax}-R. S. \textsc{Phillips} \cite{lewy173}. 
\end{remarku}


\yossection{4}{Integration of Temporally Inhomogeneous Equations of Evolution in a \texorpdfstring{$B$}{B}-space}
\label{sec:XIV.4}

We shall be concerned with the integration of the equation 
\begin{equation}\label{eq:XIV.4.1}
\dd x (t) / \dd t = A (t) x (t), a \leqq t \leqq b.
\end{equation}

Here the unknown $x(t)$ is an element of a $B$-space $X$, depending on a real parameter $t$, while $A(t)$ is a given, in general unbounded, linear operator with domain $D(A(t))$ and range $R(A(t))$, both in $X$, depending also on $t$. 

T. \textsc{Kato} \cite{kato136}, \cite{kato137} was the first to make a successful attack on the problem of integration of \eqref{eq:XIV.4.1}. He assumed the following four conditions: 

\begin{enumerate}[(i)]
\item The domain $D(A(t))$ is independent of $t$ and is strongly dense in $X$ such that, for $\alpha > 0$, the resolvent $(I - \alpha A(t))^{-1}$ exists as a bounded linear operator $\in L(X, X)$ with the norm $\leqq 1$.
\item The operator $B(t, s) = (I - A(t))(I - A(s))^{-1}$ is uniformly bounded in norm for $t \geqq s$.
\item $B(t, s)$ is, at least for some $s$, of bounded variation in $t$ in norm, that is, for every partition $s = t_{0} < t_{1} < \cdots < t_{n} = t$ of $[a, b]$, 
\[
\sum\limits_{j=0} ^{n-1} \left\| B (t _{j+1}, s) - B (t _{j}, s) \right\| \leqq N (a, b) < \infty .
\]
\item $B(t, s)$ is, at least for some $s$, weakly differentiable in $t$ and $\partial B(t, s)/\partial t$ is strongly continuous in $t$.
\end{enumerate}
Under these conditions, \textsc{Kato} proved that the limit 
\[
U (t, s) x _{0} = \underset {\max \left| t _{j+1} - t _{j} \right|\rightarrow 0} {\slim} \prod\limits_{j=n-1} ^{0} \exp \left(\left(t _{j+1} - t _{j}\right) A (t _{j})\right) x _{0}
\]
exists for every $x_0 \in X$ and gives the unique solution of \eqref{eq:XIV.4.1} with the initial condition $x(s) = x_0$ at least when $x_0 \in D(A(s))$. 

Kato's method is thus an abstraction of the classical polygon method of Cauchy for ordinary differential equation $\dd x(t) / \dd t = a(t)x(t)$. Although his method is very simple and natural in its idea, the proof is somewhat lengthy because it is concerned with a general partition of the interval $[s,t]$. \textsc{Kato} \cite{kato136} has shown that the proof can be simplified when the space $X$ is reflexive. For the reflexive $B$-space case, see also K. \textsc{Yosida} \cite{yosida371}. 

In this section, we shall be concerned with an equi-partition of a fixed interval, say, [$0$, $1$], independently of $s$ and $t$, to the effect that \textsc{Kato}'s original method be modified so as to yield a fairly simplified presentation. 

We shall assume the following four conditions which are essentially the same as \textsc{Kato}'s conditions (i) through (iv) above: 
\begin{equation}
\label{eq:XIV.4.2}
\begin{minipage}{0.92\linewidth}
\everymath{}\let\limits\nolimits
$D (A (t))$ is independent of $t$ and it is dense in $X.$
\end{minipage}
\end{equation}

\begin{equation}\label{eq:XIV.4.3}
\parbox{0.92\linewidth}{\everymath{}\let\limits\nolimits For every $\lambda \geqq 0$ and $t, 0 \leqq t \leqq 1$, the resolvent $(\lambda I - A(t))^{-1}$ exists as an operator $\in L(X, X)$ such that $\| (\lambda I - A(t))^{-1} \| \leqq \lambda^{-1}$ for $\lambda > 0$.}
\end{equation}

\begin{equation}\label{eq:XIV.4.4}
A (t) A (s) ^{-1} \in L (X, X) \quad \text{ for } \quad 0 \leqq s, t \leqq 1.
\end{equation}

\begin{equation}\label{eq:XIV.4.5}
\parbox{0.92\linewidth}{\everymath{}\let\limits\nolimits For every $x \in X$, $(t - s)^{-1} C(t, s) x = (t - s)^{-1}(A(t) A(s)^{-1} - I)x$ is bounded and uniformly strongly continuous in $t$ and $s$, $t \neq s$, and \slim $k \cdot C\left(t, t - \frac{1}{k}\right)x = C(t)x$ exists uniformly in $t$ so that $C(t) \in L(X, X)$ is strongly continuous in $t$.}
\end{equation}


\begin{remarku} Condition \eqref{eq:XIV.4.4} is stated for the sake of convenience. It is implied by (2) and the closed graph theorem, since $A(s)$ is a closed linear operator by ($3$). Moreover, conditions (2) and ($3$) imply that $A(s)$ is the infinitesimal generator of a contraction semi-group $\{\exp(tA(s)); t \geqq 0\}$ of class $(C_0)$. Hence we have (see Chapter~\ref{chap:IX}): 
\begin{align}
    \exp (t A (s)) x = & \slim _{n \rightarrow \infty} \exp \left(t A (s) (I - n ^{-1} A (s)) ^{-1}\right) x \text{ uniformly in } t, 0 \leqq t \leqq 1; \label{eq:XIV.4.6} \\
    \exp (t _{1} A (s)) \cdot \exp (t _{2} A (s)) = & \exp ((t _{1} + t _{2}) A (s)); \label{eq:XIV.4.7} \\
    \frac {\dd \exp (t A (s)) y}{\dd t} = & A (s) \exp (t A (s)) y = \exp (t A (s)) \cdot A (s) y, y\in D (A (s)), \notag \\  \text{where } \dd / \dd t & \text{ means differentiation in the strong topology of } X; \label{eq:XIV.4.8} \\
    \slim _ {t \to t _{0}} \exp (t A (s)) x = & \exp (t _{0} A (s)) x. \label{eq:XIV.4.9}
\end{align}
\end{remarku}

We are now able to state our result. 

\begin{theorem}\label{thm:XIV.4.Theorem.1} For any positive integer $k$ and $0 \leqq s \leqq t \leqq 1$, define the operator $U_k(t, s) \in L(X, X)$ through 
\begin{equation}\label{eq:XIV.4.1a}
\left.\begin{aligned}
    U _{k} (t, s)& = \exp \left((t - s) A \left(\frac {i - 1}{k}\right)\right) \quad \text{ for } \quad \frac {i - 1}{k} \leqq s \leqq t \leqq \frac {i}{k}, (1 \leqq i \leqq k), \\
    U _{k} (t, r)& = U _{k} (t, s) U _{k} (s, r) \quad \text{ for } \quad 0 \leqq r \leqq s \leqq t \leqq 1. 
\end{aligned}\right\}
\end{equation}

Then, for every $x \in X$ and $0 \leqq s \leqq t \leqq 1$, 
\begin{equation}\label{eq:XIV.4.11}
\slim _{k \to \infty} U _{k} (t, s) x = U (t, s) x \text{ exists uniformly in } t \text{ and } s.
\end{equation}

Moreover, if $y \in D(A(0))$ then the Cauchy problem 
\begin{equation}\tag{1'}\label{eq:XIV.4.1'}
\frac {\dd x (t)}{\dd t} = A (t) x (t), x (0) = y \text{ and } x (t) \in D (A (t)), 0 \leqq t \leqq 1
\end{equation}
is solved by $x(t) = U(t, 0)y$ which satisfies the estimate $\| x(t) \| \leqq \| y \|$.
\end{theorem}

\begin{proof} By ($3$), \eqref{eq:XIV.4.6} and \eqref{eq:XIV.4.1a}, we obtain 
\begin{equation}\label{eq:XIV.4.12}
\| U _{k} (t, s) x \| \leqq \| x \| (k = 1, 2, \dots ,; 0 \leqq s \leqq t \leqq 1; x \in X).
\end{equation}

We also need, for the operator 
\begin{equation}\label{eq:XIV.4.13}
W _{k} (t, s) = A (t) U _{k} (t, s) A (s) ^{-1},
\end{equation}
the following estimate which is the key for the whole proof: 
\begin{equation}\label{eq:XIV.4.14}
\begin{yoslines}
&\| W _{k} (t, s) x \| \leqq (1 + k ^{-1} N) ^{2} \cdot \exp (N (t - s)) \cdot \| x \|, N = \sup\limits _{0 \leqq s, t \leqq 1, s \neq t} \| (t - s) ^{-1} C (t, s) \|.
\end{yoslines}
\end{equation}

To prove this, we shall rewrite $W_{k}(t,s)$, remembering \eqref{eq:XIV.4.1a} and the commutativity $A(s)^{-1} \exp((t - s)A(s)) = \exp((t - s)A(s))A(s)^{-1}$: 
\[
\begin{aligned}
&\quad W _{k} (t, s) = A (t) A \left(\frac {[ k t ]}{k}\right) ^{-1} U _{k} \left(t, \frac {[ k t ]}{k}\right) A \left(\frac {[ k t ]}{k}\right) A \left(\frac {[ k t ] - 1}{k}\right) ^{-1} U _{k} \left(\frac {[ k t ]}{k}, \frac {[ k t ] - 1}{k}\right) \dots
\\
&\quad \dots A \left(\frac {[ k s ] + 2}{k}\right) A \left(\frac {[ k s ] + 1}{k}\right) ^{-1} U _{k} \left(\frac {[ k s ] + 2}{k}, \frac {[ k s ] + 1}{k}\right) A \left(\frac {[ k s ] + 1}{k}\right) A \left(\frac {[ k s ]}{k}\right) ^{-1} \times
\\
&\quad \times U _{k} \left(\frac {[ k s ] + 1}{k}, s\right) A \left(\frac {[ k s ]}{k}\right) A (s) ^{-1}.
\end{aligned}
\]

Expanding the right side and again remembering \eqref{eq:XIV.4.1a}, we obtain 
\[
\begin{aligned}
W _{k} (t, s)& = \left(I + C \left(t, \frac {[ k t ]}{k}\right)\right) \left\{U _{k} (t, s) + \sum\limits_{ku=[ks]+1} ^{[kt]} U _{k} (t, u) C \left(u, u - \frac {1}{k}\right) U _{k} (u, s) + \right.
\\
&\quad \left. + \sum\limits_{kv=[ks]+1} ^{[kt]} U _{k} (t, v) C \left(v, v - \frac {1}{k}\right) \sum\limits_{ku=[ks]+1} ^{[kv]} U _{k} (v, u) C \left(u, u - \frac {1}{k}\right) U _{k} (u, s) + \dots \right\} \times
\\
&\quad \times \left(I + C \left(\frac {[ k s ]}{k}, s\right)\right),
\end{aligned}
\]
that is, 
\begin{equation}\label{eq:XIV.4.15}
\begin{aligned} 
    W _{k} (t, s) & = \left(I + C \left(t, \frac {[ k t ]}{k}\right)\right) \left\{U _{k} (t, s) + W _{k} ^{(1)} (t, s) + W _{k} ^{(2)} (t, s) + \dots \right\} \\
&\quad \times \left(I + C \left(\frac {[ k s ]}{k}, s\right)\right), \\ 
    W _{k} ^{(1)} (t, s) & = \sum\limits_{ku=[ks]+1} ^{[kt]} U _{k} (t, u) C \left(u, u - \frac {1}{k}\right) U _{k} (u, s), \\ 
    W _{k} ^{(m+1)} (t, s) & = \sum\limits_{ku=[ks]+1} ^{[kt]} U _{k} (t, u) C \left(u, u - \frac {1}{k}\right) W _{k} ^{(m)} (u, s), (m = 1, 2, \dots , [ k t ] - 1). 
\end{aligned}
\end{equation}

We have $\left\| C\left(u, u - \frac{1}{k}\right)x\right\| \leqq \frac{1}{k} N \| x\|$ by \eqref{eq:XIV.4.14}. Hence from \eqref{eq:XIV.4.12} we conclude that 
\begin{equation}\label{eq:XIV.4.16}
\begin{aligned}
&\left\| W _{k} ^{(1)} (t, s) x \right\| \leqq (t - s) N \| x \|, \dots \left\| W _{k} ^{(m)} (t, s) x \right\| \leqq \frac {(t - s) ^{m}}{m !} N ^{m} \| x \|.
\end{aligned}
\end{equation}

Combined with the definition of $N$, we have proved \eqref{eq:XIV.4.14} by \eqref{eq:XIV.4.12}.

Incidentally, we have proved that 
\begin{equation}\label{eq:XIV.4.17}
U _{k} (t, s) y \in D (A (t)) \quad \text{ when } \quad y \in D (A (s)).
\end{equation}

Hence 
\[
U _{k} (t, s) y = \exp \left(\left(t - \frac {[ k t ]}{k}\right) A \left(\frac {[ k t ]}{k}\right)\right) U _{k} \left(\frac {[ k t ]}{k}, s\right) y
\]
is differentiable at $t \neq \frac{i}{k} (i = 0, 1, \ldots, k)$ and 
\begin{equation}\label{eq:XIV.4.18}
\frac {\dd U _{k} (t , s) y}{\dd t} = A \left(\frac {[ k t ]}{k}\right) U _{k} (t, s) y, y \in D (A (0)).
\end{equation}

Similarly, $U_{k}(t,s)y$ is differentiable at $s\neq \frac{i}{k} (i = 0,1,\dots ,k)$ and 
\begin{equation}\label{eq:XIV.4.19}
\frac {\dd U _{k} (t , s) y}{\dd s} = - U _{k} (t, s) A \left(\frac {[ k s ]}{k}\right) y, y \in D (A (0)).
\end{equation}

These derivatives are bounded and are strongly continuous in $t$ and $s$ except at $t = \frac{i}{k}$ and $s = \frac{i}{k}$. For the proof, rewrite the right side of \eqref{eq:XIV.4.18} as $= C\left(\frac{[kt]}{k}, t\right) W_k(t, s) A(s) A(0)^{-1}x, y = A(0)^{-1}x$, and make use of ($5$), \eqref{eq:XIV.4.9}, \eqref{eq:XIV.4.12}, \eqref{eq:XIV.4.14} and \eqref{eq:XIV.4.15}; similarly for \eqref{eq:XIV.4.19}. 

Since $U_{n}(t,s)U_{k}(s,0)A(0)^{-1}x$ is strongly continuous in $s$ by \eqref{eq:XIV.4.9}, we have, remembering the fact $U_{k}(s,0)A(0)^{-1}x\in D(A(s))$ together with \eqref{eq:XIV.4.18} and \eqref{eq:XIV.4.19}, 
\[
\begin{aligned}
&\left(U _{k} (t, 0) - U _{n} (t, 0)\right) A (0) ^{-1} x = [ U _{n} (t, s) U _{k} (s, 0) A (0) ^{-1} x ] _{s=0} ^{s=t}
\\
&= \int_{0} ^{t} \frac {d}{\dd s} \left\{U _{n} (t, s) U _{k} (s, 0) A (0) ^{-1} x \right\} \dd s = \int_{0} ^{t} U _{n} (t, s) \left\{A \left(\frac {[ k s ]}{k}\right) - A \left(\frac {[ n s ]}{n}\right) \right\} A \left(\frac {[ k s ]}{k}\right) ^{-1}
\\
&\quad A \left(\frac {[ k s ]}{k}\right) U _{k} (s, 0) A (0) ^{-1} x \dd s
\\
&= \int_{0} ^{t} - U _{n} (t, s) C \left(\frac {[ n s ]}{n}, \frac {[ k s ]}{k}\right) A \left(\frac {[ k s ]}{k}\right) A (s) ^{-1} W _{k} (s, 0) x \dd s.
\end{aligned}
\]

Hence, \eqref{eq:XIV.4.12} and \eqref{eq:XIV.4.14} yield 
\begin{gather*}{\| U _{k} (t, 0) A (0) ^{-1} x - U _{n} (t, 0) A (0) ^{-1} x \|} \leqq  \int_{0} ^{t} \left\| C \left(\frac {[ n s ]}{n}, \frac {[ k s ]}{k}\right) \left(I + C \left(\frac {[ k s ]}{k}, s\right)\right) W _{k} (s, 0) x \right\| d s \\\leqq  N \int_{0} ^{t} \left| \frac {[ k s ]}{k} - \frac {[ n s ]}{n} \right| \cdot \left\| \left(I + C \left(\frac {[ k s ]}{k}, s\right)\right) \right. \cdot \left. \cdot W _{k} (s, 0) x \right\| d s \\\leqq  N \int_{0} ^{t} \left| \frac {[ k s ]}{k} - \frac {[ n s ]}{n} \right| \cdot \left(1 + N \left| s - \frac {[ k s ]}{k} \right|\right) g \cdot \left(1 + \frac {1}{k} N\right) \exp (N s) \cdot \| x \| d s.
\end{gather*}

This proves that $\slim\limits_{k \to \infty} U_k(t, 0) A(0)^{-1} x$ exists uniformly in $t, 0 \leqq t \leqq 1$, 
and so, by (2) and \eqref{eq:XIV.4.12}, we prove that, for every $x \in X$, $\slim\limits_{k \to \infty} U_k(t, 0) x = U(t, 0) x$ exists uniformly in $t, 0 \leqq t \leqq 1$. 

Similar argument shows that, for every $x \in X$,
\begin{equation}\label{eq:XIV.4.20}
\parbox{0.92\linewidth}{\everymath{}\let\limits\nolimits \slim $U_{k}(t,s)x = U(t,s)x$ exists uniformly in $t$ and $s,0 \leqq s \leqq t \leqq 1$.}
\end{equation} 

Therefore, by \eqref{eq:XIV.4.9}, we see that $U(t, s)$ is uniformly strongly continuous in $t$ and $s$. 

We also prove easily, remembering \eqref{eq:XIV.4.15}, \eqref{eq:XIV.4.16} and ($20$), that as $k \to \infty$ $W_{k}(t, 0)x$ strongly converges boundedly in $t$ such that 
\begin{equation}\label{eq:XIV.4.21}
\begin{yoslines}& \slim _{k \rightarrow \infty} W _{k} (t, 0) x = W (t, 0) x = U (t, 0) x + W ^{(1)} (t, 0) x + W ^{(2)} (t, 0) x \dots ,\\& \text{ where } W ^{(1)} (t, 0) x = \int_{0} ^{t} U (t, s) C (s) U (s, 0) x \dd s,\\& \dots\\& W ^{(m+1)} (t, 0) x = \int_{0} ^{t} U (t, s) C (s) W ^{(m)} (s, 0) x \dd s \quad (m = 1, 2, \dots).\end{yoslines}
\end{equation}

Therefore, if $y \in D(A(0))$ then the limits 
$\slim\limits_{k \to \infty} U_k(t,0) y = U(t,0) y$ and $\slim\limits_{k \to \infty} A(t) U_k(t,0) y = W(t,0) A(0) y$ both exist boundedly and uniformly in $t$ and $W(t,0) A(0) y$ is uniformly strongly continuous in $t$. $A(t)$ being a closed linear operator by ($3$), we have proved that, for every $y \in D(A(0))$, 
\begin{equation}\label{eq:XIV.4.22}
\begin{yoslines}& U (t, 0) y \in D (A (t)) \quad \text{ and }\\& A (t) U (t, 0) y = \slim\limits_{k \to \infty} A (t) A \left(\frac {[ k t ]}{k}\right) ^{-1} A \left(\frac {[ k t ]}{k}\right) U _{k} (t, 0) y = W (t, 0) A (0) y,\\& \text{ where the \slim \, exists boundedly and uniformly in } t, 0 \leqq t \leqq 1.\end{yoslines}
\end{equation}
Hence, by letting $k \to \infty$ in 
\[
U _{k} (t, 0) y - y = \int_{0} ^{t} \left(\frac {d}{\dd s} U _{k} (s, 0) y\right) \dd s = \int_{0} ^{t} A \left(\frac {[ k s ]}{k}\right) U _{k} (s, 0) y \dd s
\]

We obtain 
\[
U (t, 0) y - y = \int_{0} ^{t} A (s) U (s, 0) y \dd s.
\]


As proved in \eqref{eq:XIV.4.22}, the integrand is strongly continuous in $s$ so that we have proved that $x(t) = U(t,0)y$ solves Cauchy problem \eqref{eq:XIV.4.1'}.
\end{proof}

\begin{remarku} Theorem~\ref{thm:XIV.4.Theorem.1} and its proof are adapted from K. \textsc{Yosida} \cite{yosida373}. It is to be noted here that an inequality similar to our \eqref{eq:XIV.4.14} was obtained in T. \textsc{Kato} \cite{kato136} by solving a Volterra type integral equation. The appropriateness of the equi-partition of the interval [$0$,$1$] as revealed in \eqref{eq:XIV.4.1a} is suggested by a paper by J. \textsc{Kisyński} \cite{kisynski150}. 
\end{remarku}

The approximation $x_{k}(t) = U_{k}(t,0)y$ to the true solution $x(t) = U(t,0)y$ of \eqref{eq:XIV.4.1'} is a kind of difference approximation to the differential equation given in \eqref{eq:XIV.4.1'}. However, the true difference approximation would be to use the scheme 
\begin{equation}\label{eq:XIV.4.23}
\frac {x \left(t _{j}\right) - x \left(t _{j-1}\right)}{t _{j} - t _{j-1}} = A \left(t _{j}\right) x \left(t _{j}\right),
\end{equation}
choosing the backward difference. This scheme gives 
\[
x \left(t _{j}\right) = \left(I - \left(t _{j} - t _{j-1}\right) A \left(t _{j}\right)\right) ^{-1} x \left(t _{j-1}\right).
\]

Therefore, suggested by \eqref{eq:XIV.4.1a}, we are lead to the approximation 
\begin{equation}\tag{10'}\label{eq:XIV.4.10'}
\begin{yoslines}
&V _{k} (t, s) = \left(I - \left(t - \frac {[ k t ]}{k}\right) A \left(\frac {[ k t ]}{k}\right)\right) ^{-1} \left(I - \frac {1}{k} A \left(\frac {[ k t ] - 1}{k}\right)\right) ^{-1} \dots \\
&\dots \left(I - \frac {1}{k} A \left(\frac {[ k s ] + 1}{k}\right)\right) ^{-1} \left(I - \left(\frac {[ k s ] + 1}{k} - s\right) A \left(\frac {[ k s ]}{k}\right)\right) ^{-1}, (0 \leqq s \leqq t \leqq 1).
\end{yoslines}
\end{equation}

From the point of view of the numerical analysis, this approximation $V_{k}(t,s)$ would be much more practical than the approximation $U_{k}(t,s)$, since in the construction of the approximation $U_{k}(t,s)$ we have to appeal to the construction of $\exp(tA(s))$. 

We shall prove (Cf. the last section in T. \textsc{Kato} \cite{kato143}). 

\begin{theorem}\label{thm:XIV.4.Theorem.2} Under the same conditions (2) through ($5$), the Cauchy problem given in \eqref{eq:XIV.4.1'} is, for every initial data $y \in D(A(0))$, solved by 
\[
x (t) = \slim\limits_{k \rightarrow \infty} V _{k} (t, 0) y.
\]
\end{theorem}

\begin{proof}
 From ($3$) it follows that 
\begin{equation}\tag{12'}\label{eq:XIV.4.12'}
\| V _{k} (t, s) x \| \leqq \| x \| \quad (k = 1, 2, \dots ; 0 \leqq s \leqq t \leqq 1; x \in X). \end{equation}

We next see, from 
\[
V _{k} (t, s) = V _{k} \left(t, \frac {[ k s ] + 1}{k}\right) \left(I - \left(\frac {[ k s ] + 1}{k} - s\right) A \left(\frac {[ k s ]}{k}\right)\right) ^{-1},
\]
that $V_{k}(t, s)$ is differentiable at $s \neq \frac{i}{k} (i = 0, 1, \ldots, k)$ and 
\begin{equation}\tag{19'}\label{eq:XIV.4.19'}
\frac {\dd V _{k} (t , s) y}{\dd s} = - V _{k} (t, s) \left(I - \left(\frac {[ k s ] + 1}{k} - s\right) A \left(\frac {[ k s ]}{k}\right)\right) ^{-1} A \left(\frac {[ k s ]}{k}\right) y.
\end{equation}

We see, by $y \in D(A(0)) = D(A(s))$, ($3$) and \eqref{eq:XIV.4.12'}, that this derivative is bounded and strongly continuous in $s$ except at $s = \frac{i}{k}$. Since 
$V_{k}(t,s)U_{k}(s,0)A(0)^{-1}x$ is strongly continuous in $s$, we obtain, remembering the fact $U_{k}(s,0)A(0)^{-1}x\in D(A(0))$, 
\[
\begin{aligned}
& \left(U _{k} (t, 0) - V _{k} (t, 0)\right) A (0) ^{-1} x = \left[ V _{k} (t, s) U _{k} (s, 0) A (0) ^{-1} x \right] _{s=0} ^{s=t}
\\
= & \int_{0} ^{t} \frac {d}{\dd s} \left(V _{k} (t, s) U _{k} (s, 0) A (0) ^{-1} x\right) \dd s
\\
= & \int_{0} ^{t} V _{k} (t, s) \left(A \left(\frac {[ k s ]}{k}\right)\right. \left.- A \left(\frac {[ k s ]}{k}\right) \left(I - \left(\frac {[ k s ] + 1}{k} - s\right) A \left(\frac {[ k s ]}{k}\right)\right) ^{-1}\right) U _{k} (s, 0) A (0) ^{-1} x \dd s
\\
= & \int_{0} ^{t} V _{k} (t, s) \left(I - \left(I - \left(\frac {[ k s ] + 1}{k} - s\right) A \left(\frac {[ k s ]}{k}\right)\right) ^{-1}\right)\left(I + C \left(\frac {[ k s ]}{k}, s\right)\right) W _{k} (s, 0) x \dd s.
\end{aligned}
\]

Thus we have 
\begin{equation}\label{eq:XIV.4.24}
\begin{aligned}
& \left\| \left(U _{k} (t, 0) - V _{k} (t, 0)\right) A (0) ^{-1} x \right\|
\\
&\leqq \int_{0} ^{t} \left\| \left(I - \left(I - \left(\frac {[ k s ] + 1}{k} - s\right) A \left(\frac {[ k s ]}{k}\right)\right) ^{-1}\right) \left(I + C \left(\frac {[ k s ]}{k}, s\right)\right) W _{k} (s, 0) x \right\| \dd s.
\end{aligned}
\end{equation}

From ($5$) and \eqref{eq:XIV.4.21}, it follows that 
\[
\slim _{k \rightarrow \infty} \left(I + C \left(\frac {[ k s ]}{k}, s\right)\right) W _{k} (s, 0) x = W (s, 0) x \text{ uniformly in } s, 0 \leqq s \leqq 1.
\]

On the other hand, we have, for $z \in D(A(s))$, 
\begin{equation}\label{eq:XIV.4.25}
\begin{aligned} & - \left(I - \left(I - \left(\frac {[ k s ]}{k} - s\right) A \left(\frac {[ k s ]}{k}\right)\right) ^{-1}\right) z \\ & = \left(\frac {[ k s ]}{k} - s\right) \left(I - \left(\frac {[ k s ]}{k} - s\right) A \left(\frac {[ k s ]}{k}\right)\right) ^{-1} A \left(\frac {[ k s ]}{k}\right) A (s) ^{-1} A (s) z. \end{aligned}
\end{equation}

By ($3$) and ($5$), we see that this tends strongly to $0$ as $k \to \infty$. The integrand in \eqref{eq:XIV.4.24} is bounded in $k$ and $s$. Thus, since $D(A(s)) = D(A(0))$ is dense in $X$, we see from \eqref{eq:XIV.4.25} that 
\[
\slim _{k \to \infty} (U _{k} (t, 0) - V _{k} (t, 0)) A (0) ^{-1} x = 0.
\]

Therefore, by virtue of Theorem~\ref{thm:XIV.4.Theorem.1}, we have proved Theorem~\ref{thm:XIV.4.Theorem.2}. 
\end{proof}

\begin{remarku} The above proof is communicated from H. \textsc{Fujita} who is suggested by the last section of T. \textsc{Kato} \cite{kato143}. 
\end{remarku}

Other approaches. K. \textsc{Yosida} \cite{yosida366} and \cite{yosida371} devised the idea of approximating Cauchy problem \eqref{eq:XIV.4.1'} by 
\begin{equation}\tag{1''}\label{eq:XIV.4.1''}
\frac {\dd x _{k} (t)}{\dd t} = A (t) \left(I - \frac {1}{k} A (t)\right) ^{-1} x _{k} (t), x _{k} (0) = y \in D (A (0)), 0 \leqq t \leqq 1.
\end{equation}

Applying the Cauchy polygon method to equation \eqref{eq:XIV.4.1''}, J. \textsc{Kisyński} \cite{kisynski150} proved the strong convergence of $\{x_{k}(t)\}$ to the solution of \eqref{eq:XIV.4.1'}. Another method of integration of \eqref{eq:XIV.4.1} has been devised by J. L. \textsc{Lions} \cite{liggett179}. He assumes the operator $A(t)$ to be an elliptic differential operator with smooth coefficients depending on $t$, and seeks distributional solutions by transforming equation \eqref{eq:XIV.4.1} into an integrated form in concrete function spaces such as the Sobolev space $W^{k,p}(\Omega)$ and its variants. We also refer the reader to a paper by O. A. \textsc{Ladyzhenskaya}-I. M. \textsc{Visik} \cite{ladyzhenskaya170} which is motivated by a similar idea as Lions'.\yosstarfootnote{See also Supplementary Notes, p.~\pageref{notes:XIV}.} 


\yossection{5}{The Method of Tanabe and Sobolevski}
\label{sec:XIV.5}

Let $X$ be a complex $B$-space, and consider the equation of evolution in $X$ with a given \emph{inhomogeneous term} $f(t)$: 
\begin{equation}\label{eq:XIV.5.1}
\dd x (t) / \dd t = A x (t) + f (t), a \leqq t \leqq b.
\end{equation}

Then the solution $x(t) \in X$ with the initial condition $x(a) = x_0 \in X$ is given formally by the so-called \emph{Duhamel Principle} from the solution $\exp((t - a)A)x$ of the \emph{homogeneous equation} $\dd x/\dd t = Ax$: 
\begin{equation}\label{eq:XIV.5.2}
x (t) = \exp ((t - a) A) x _{0} + \int_{a} ^{t} \exp ((t - s) A) \cdot f (s) \dd s.
\end{equation}

This suggests that a temporally inhomogeneous equation in $X$: 
\begin{equation}\label{eq:XIV.5.3}
\dd x (t) / \dd t = A (t) x (t), a \leqq t \leqq b,
\end{equation}
may be solved formally as follows. We rewrite equation \eqref{eq:XIV.5.3} in the form 
\begin{equation}\label{eq:XIV.5.4}
\dd x (t) / \dd t = A (a) x (t) + (A (t) - A (a)) x (t).
\end{equation}

By virtue of the formalism \eqref{eq:XIV.5.2}, the solution $x(t)$ of \eqref{eq:XIV.5.4} with the initial condition $x(a) = x_{0}$ will be given as the solution of an abstract integral equation 
\begin{equation}\label{eq:XIV.5.5}
\begin{aligned}
&x (t) = \exp ((t - a) A (a)) x _{0} + \int_{a} ^{t} \exp ((t - s) A (s)) (A (s) - A (a)) x (s) \dd s.
\end{aligned}
\end{equation}

Solving \eqref{eq:XIV.5.5} formally by successive approximation, we obtain approximate solutions: 
\[
\begin{aligned}
&x _{1} (t) = \exp ((t - a) A (a)) x _{0}, \dots x _{n+1} (t) = \exp ((t - a) A (a)) x _{0}
\\
&+ \int_{a} ^{t} \exp ((t - s) A (s)) (A (s) - A (a)) x _{n} (s) \dd s.
\end{aligned}
\]

Hence the solution $x(t)$ of \eqref{eq:XIV.5.5} would be given formally by 
\begin{equation}\label{eq:XIV.5.6}
x (t) = \exp ((t - a) A (a)) x _{0} + \int_{a} ^{t} \exp ((t - s) A (s)) R (s, a) x _{0} \dd s,
\end{equation}
where 
\begin{equation}\label{eq:XIV.5.7}
\begin{aligned}
&R (t, s) = \sum\limits_{m=1} ^{\infty} R _{m} (t, s), R _{1} (t, s) = (A (t) - A (s)) \exp ((t - s) A (s)), s < t, = 0, s \geqq t,
\\
&R _{m} (t, s) = \int_{s} ^{t} R _{1} (t, \sigma) R _{m-1} (\sigma , s) \dd \sigma \quad (m = 2, 3, \dots).
\end{aligned}
\end{equation}

Justification for the above formal procedure of integration has been given by H. \textsc{Tanabe} \cite{tanabe315} using the theory of holomorphic semi-groups as given in Chapter~\ref{chap:IX}, \ref{sec:IX.10}. We shall follow Tanabe's approach, and assume the following conditions: 

\begin{equation}\label{eq:XIV.5.8}
\parbox{0.92\linewidth}{\everymath{}\let\limits\nolimits For each $t \in [a, b]$, $A(t)$ is a closed linear operator with domain dense in $X$ and range $X$ such that the resolvent set $\varrho(A(t))$ of $A(t)$ contains a fixed angular domain $\Theta$ of the complex $\lambda$-plane consisting of the origin $0$ plus the set $\{\lambda; -\theta < \arg \lambda < \theta \text{ with } \theta > \pi/2\}$. The resolvent $(\lambda I - A(t))^{-1}$ is strongly continuous in $t$ uniformly in $\lambda$ on any compact set $\subseteq \Theta$.}
\end{equation}

\begin{equation}\label{eq:XIV.5.9}
\parbox{0.92\linewidth}{\everymath{}\let\limits\nolimits There exists positive constants $M$ and $N$ such that, for $\lambda \in \Theta$ and $t \in [a, b]$, we have $\|(\lambda I - A(t))^{-1}\| \leqq N(|\lambda| - M)^{-1}$ whenever $|\lambda| > M$ with $N = 1$ for real $\lambda$.}
\end{equation}

\begin{equation}\label{eq:XIV.5.10}
\parbox{0.92\linewidth}{\everymath{}\let\limits\nolimits The domain $D(A(t))$ of $A(t)$ is independent of $t$ so that, by the closed graph theorem in Chapter~\ref{chap:II}, \ref{sec:II.6}, the operator $A(t)A(s)^{-1}$ is in $L(X,X)$. It is assumed that there exists a positive constant $K$ such that $\|A(t)A(s)^{-1} - A(r)A(s)^{-1}\| \leqq K|t - r|$ for $s, t$ and $r \in [a,b]$.}
\end{equation}



Under these conditions, we can prove 

\begin{theoremu} For any $x_0 \in X$ and $s$ with $a \leqq s < b$, the equation 
\begin{equation}\tag{3'}\label{eq:XIV.5.3'}
\dd x (t) / \dd t = A (t) x (t), x (s) = x _{0}, s < t \leqq b
\end{equation}
admits a uniquely determined solution $x(t) \in \mathbf{X}$. This solution is given by 
\begin{align}
x (t)& = U (t, s) x (s) = U (t, s) x _{0}, \text{ where } \label{eq:XIV.5.11} \\
&\left.\begin{yoslines}& U (t, s) = \exp ((t - s) A (s)) + W (t, s),\\& W (t, s) = \int_{s} ^{t} \exp ((t - \sigma) A (\sigma)) R (\sigma , s) \dd \sigma \text{ with } R (t, s) \quad \text{ given  by \eqref{eq:XIV.5.7}. }\end{yoslines}\right\} \label{eq:XIV.5.12}
\end{align}
\end{theoremu}

For the proof, we prepare three lemmas. 

\begin{lemma}\label{lem:XIV.5.Lemma.1} $R(t,s)$ satisfies, with a constant $C$, 
\begin{equation}\label{eq:XIV.5.13}
\left\| R (t, s) \right\| \leqq K C \cdot \exp (K C (t - s)),
\end{equation}
and $R(t, s)$ is strongly continuous in $a \leqq s < t \leqq b$. 
\end{lemma}

\begin{proof}
 By \eqref{eq:XIV.5.8} and \eqref{eq:XIV.5.9}, we see that each $A(s)$ generates a holomorphic semi-group which is given by (see Chapter~\ref{chap:IX}, \ref{sec:IX.10}) 
\begin{equation}\label{eq:XIV.5.14}
\exp (t A (s)) = (2 \pi i) ^{-1} \int_ {C ^{\prime}} e ^{\lambda t} (\lambda I - A (s)) ^{-1} \dd \lambda , \text{ where } C ^{\prime} \text{ is }
\end{equation}
a smooth contour running from $\infty e^{-i\theta}$ to $\infty e^{i\theta}$ in $\Theta$. 

Hence, by $A(s)(\lambda I - A(s)) = \lambda (\lambda I - A(s))^{-1} - I$, we have, for $(b - a) > t > 0$, 
\begin{equation}\label{eq:XIV.5.15}
\begin{aligned}
&\left\| \exp (t A (s)) \right\| \leqq C \quad \text{ and } \quad \left\| A (s) \exp (t A (s)) \right\| \leqq C t ^{-1}, \\
&\qquad \text{where the positive constant } C \text{ is independent of } t > 0 \text{ and } s\in [a,b].
\end{aligned}
\end{equation}


We have, by \eqref{eq:XIV.5.7}, 
\[
R _{1} (t, s) = \left(A (t) - A (s)\right) A (s) ^{-1} A (s) \exp ((t - s) A (s)), t > s,
\]
and so, by \eqref{eq:XIV.5.10} and \eqref{eq:XIV.5.15}, 
\begin{equation}\label{eq:XIV.5.16}
\left\| R _{1} (t, s) \right\| \leqq K C.
\end{equation}

It is also clear, from \eqref{eq:XIV.5.8} and \eqref{eq:XIV.5.14}, that $R_{1}(t,s)$ is strongly continuous in $a \leqq s < t \leqq b$. Next, by induction, we obtain 
\[
\begin{aligned}
&\left\| R _{m} (t, s) \right\|  \leqq \int_{s} ^{t} \left\| R _{1} (t, \sigma) \right\| \cdot \left\| R _{m-1} (\sigma , s) \right\| \dd \sigma \leqq \int_{s} ^{t} (K C) ^{m} (\sigma - s) ^{m-2} (\underline {{m - 2}}) ^{-1} \dd \sigma \\
&= (K C) ^{m} (t - s) ^{m-1} (| m - 1) ^{-1},
\end{aligned}
\]
and hence \eqref{eq:XIV.5.13} is obtained. In the same way, we see that $R(t, s)$ is strongly continuous in $a \leqq s < t \leqq b$. 
\end{proof}

\begin{lemma}\label{lem:XIV.5.Lemma.2} For $s < \tau < t$, we have 
\begin{equation}\label{eq:XIV.5.17}
\left\| R (t, s) - R (\tau , s) \right\| \leqq C _{1} \left(\frac {t - \tau}{t - s} + (t - \tau) \log \frac {t - s}{t - \tau}\right),
\end{equation}
where $C_1$ is a positive constant which is independent of $s, \tau$ and $t$. 
\end{lemma}

\begin{proof}
 We have, by \eqref{eq:XIV.5.7}, 
\[
\begin{aligned}
R _{1} (t, s) - R _{1} (\tau , s)& = \left(A (t) - A (\tau)\right) \exp \left((t - s) A (s)\right)
\\
&\quad \quad + \left(A (\tau) - A (s)\right) \left[ \exp \left((t - s) A (s)\right) - \exp \left((\tau - s) A (s)\right) \right]. \end{aligned}
\]

By \eqref{eq:XIV.5.10} and \eqref{eq:XIV.5.15}, the norm of the first term on the right is dominated by $KC(t - \tau)(t - s)^{-1}$. The second term on the right is 
\[
\begin{aligned}
&= (A (\tau) - A (s)) \int_{\tau - s} ^{t-s} \frac {d}{\dd \sigma} \exp (\sigma A (s)) = (A (\tau) - A (s)) A (s) ^{-1} \times \int_{\tau - s} ^{t-s} A (s) ^{2} \exp (\sigma A (s)) \dd \sigma ,
\end{aligned}
\]
and we have, by \eqref{eq:XIV.5.15}, 
\[
\begin{aligned}
&\left\| \int_{\tau - s} ^{t-s} A (s) ^{2} \exp (\sigma A (s)) \dd \sigma \right\| \leqq \int_{\tau - s} ^{t-s} \| (A (s) \exp (2 ^{-1} \sigma A (s))) ^{2} \| \dd \sigma \leqq \int_{\tau - s} ^{t-s} (2 C / \sigma) ^{2} \dd \sigma
\\
&\quad = 4 C ^{2} \left[ \frac {- 1}{\sigma} \right] _{\tau - s} ^{t-s} = 4 C ^{2} (t - \tau) (t - s) ^{-1} (\tau - s) ^{-1}.
\end{aligned}
\]

Therefore 
\begin{equation}\label{eq:XIV.5.18}
\left\| R _{1} (t, s) - R _{1} (\tau , s) \right\| \leqq K C (1 + 4 C) \frac {t - \tau}{t - s}.
\end{equation}

On the other hand, by \eqref{eq:XIV.5.7}, 
\[
\begin{aligned}
&\sum\limits_{m=2} ^{\infty} R _{m} (t, s) - \sum\limits_{m=2} ^{\infty} R _{m} (\tau , s) = \int_{s} ^{t} R _{1} (t, \sigma) R (\sigma , s) \dd \sigma
\\
&- \int_{s} ^{\tau} R _{1} (\tau , \sigma) R (\sigma , s) \dd \sigma = \int_{\tau} ^{t} R _{1} (t, \sigma) R (\sigma , s) \dd \sigma + \int_{s} ^{\tau} (R _{1} (t, \sigma) - R _{1} (\tau , \sigma)) R (\sigma , s) \dd \sigma .
\end{aligned}
\]

The norm of the first term on the right is dominated by 
\[
\int_{\tau} ^{t} \left\| R _{1} (t, \sigma) \right\| \left\| R (\sigma , s) \right\| \dd \sigma \leqq K ^{2} C ^{2} \exp (K C (b - a)) (t - \tau).
\]

The norm of the second term on the right is, by \eqref{eq:XIV.5.13} and \eqref{eq:XIV.5.18}, dominated by 
\[
\begin{aligned}
&\int_{s} ^{\tau} \left\| R _{1} (t, \sigma) - R _{1} (\tau , \sigma) \right\| \left\| R (\sigma , s) \right\| \dd \sigma
\\
&\leqq K ^{2} C ^{2} (1 + 4 C) \exp (K C (b - a)) \int_{s} ^{\tau} (t - \tau) (t - \sigma) ^{-1} \dd \sigma = K _{1} (t - \tau) \log \frac {t - s}{t - \tau}.
\end{aligned}
\]

Therefore we obtain \eqref{eq:XIV.5.17}. 
\end{proof}

\begin{lemma}\label{lem:XIV.5.Lemma.3} For $s < t$, we have 
\begin{equation}\label{eq:XIV.5.19}
\begin{aligned} & {\left\| A (t) \left\{\exp ((t - s) A (t)) - \exp ((t - s) A (s)) \right\} \right\| \leqq C _{2}, \text{ where }} \\ & {\quad C _{2} \text{ is a positive constant independent of } s \text{ and } t.} \end{aligned}
\end{equation}
\end{lemma}

\begin{proof}
 We obtain, from \eqref{eq:XIV.5.14}, 
\[
\begin{gathered}
A (t) \left\{\exp ((t - s) A (t)) - \exp ((t - s) A (s)) \right\}
\\
= (2 \pi i) ^{-1} \int_{C} \dot {e} ^{\lambda (t - s)} A (t) (\lambda I - A (t)) ^{-1} (A (t) - A (s)) (\lambda I - A (s)) ^{-1} \dd \lambda . \end{gathered}
\]

On the other hand, we have $A(t)(\lambda I - A(t))^{-1} = \lambda (\lambda I - A(t))^{-1} - I$, and so, by \eqref{eq:XIV.5.9}, 
\begin{equation}\label{eq:XIV.5.20}
\left\| A (t) (\lambda I - A (t)) ^{-1} \right\| \leqq \frac {| \lambda |}{| \lambda | - M} + 1 \text{ for } \lambda \in \Theta \text{ and } t \in [ a, b ].
\end{equation}

Hence we obtain \eqref{eq:XIV.5.19} by \eqref{eq:XIV.5.10} and 
\[
\begin{aligned}
&{\left\| (A (t) - A (s)) (\lambda I - A (s)) ^{-1} \right\|} {\qquad \leqq \left\| (A (t) - A (s)) A (s) ^{-1} \right\| \left\| A (s) (\lambda I - A (s)) ^{-1} \right\|.}
\end{aligned}
\]
\end{proof}

\begin{proof}[Proof of the Theorem.]
We rewrite $W(t, s)$ given in \eqref{eq:XIV.5.12} as follows:
\[
\begin{aligned}
W (t, s)& = \int_{s} ^{t} \exp ((t - \tau) A (t)) R (t, s) \dd \tau
\\
\qquad +& \int_{s} ^{t} \left\{\exp ((t - \tau) A (\tau)) - \exp ((t - \tau) A (t)) \right\} R (\tau , s) \dd \tau
\\
\qquad +& \int_{s} ^{t} \exp ((t - \tau) A (t)) (R (\tau , s) - R (t, s)) \dd \tau .
\end{aligned}
\]

By approximating the integrals by Riemann sums and making use of the closure property of the operator $A(t)$, we see that we can apply $A(t)$ to each term of the right side of the above equality. For, by \eqref{eq:XIV.5.19}, we can apply $A(t)$ to the second term of the right side; and also to the third term of the right side by \eqref{eq:XIV.5.15} and \eqref{eq:XIV.5.17}; we also have, by $A(t) \exp((t - \tau)A(t)) = -\dd\exp((t - \tau)A(t))/\dd\tau$, 
\[
A (t) \int_{s} ^{t} \exp ((t - \tau) A (t)) R (t, s) \dd \tau = \left\{\exp ((t - s) A (t)) - I \right\} R (t, s).
\]

Hence we obtain 
\begin{equation}\label{eq:XIV.5.21}
\begin{aligned}
A (t) U (t, s)& = A (t) \exp ((t - s) A (s)) + \{\exp ((t - s) A (t)) - I \} R (t, s)
\\
\quad +& \int_{s} ^{t} A (t) \left\{\exp ((t - \tau) A (\tau)) - \exp ((t - \tau) A (t)) \right\} R (\tau , s) \dd \tau
\\
\quad +& \int_{s} ^{t} A (t) \exp ((t - \tau) A (t)) (R (\tau , s) - R (t, s)) \dd \tau .
\end{aligned}
\end{equation}

The above proof shows that $A(t)U(t,s)$ is strongly continuous in $a \leqq s < t \leqq b$ and 
\begin{equation}\label{eq:XIV.5.22}
\left\| A (t) W (t, s) \right\| \leqq C _{3} \text{ and } \left\| A (t) U (t, s) \right\| \leqq C _{3} (t - s) ^{-1},
\end{equation}
where $C_{3}$ is a positive constant independent of $s$ and $t$. 

We next define, for $s < (t - h) < t$, 
\begin{equation}\label{eq:XIV.5.23}
U _{h} (t, s) = \exp ((t - s) A (s)) + \int_{s} ^{t-h} \exp ((t - \tau) A (\tau)) R (\tau , s) \dd \tau .
\end{equation}

Since a holomorphic semi-group $\exp(tA(u))$ is differentiable in $t > 0$, we have 
\[
\begin{aligned}
\frac {\partial}{\partial t} U _{h} (t, s)& = A (s) \exp ((t - s) A (s)) + \exp (h A (t - h)) R (t - h, s)
\\
\qquad +& \int_{s} ^{t-h} A (\tau) \exp ((t - \tau) A (\tau)) R (\tau , s) \dd \tau .
\end{aligned}
\]

Hence we have, by \eqref{eq:XIV.5.7}, 
\begin{equation}\label{eq:XIV.5.24}
\begin{aligned}
&\frac {\partial}{\partial t} U _{h} (t, s) - A (t) U _{h} (t, s)  = \exp (h A (t - h)) R (t - h, s) - R _{1} (t, s) - \int_{s} ^{t-h} R _{1} (t, \tau) R (\tau , s) \dd \tau .
\end{aligned}
\end{equation}

By \eqref{eq:XIV.5.8}, \eqref{eq:XIV.5.13} and \eqref{eq:XIV.5.14}, $\exp(hA(t-h))R(t-h,s)$ tends strongly to $R(t,s)$ as $h\downarrow 0$. Thus we have 
\begin{equation}\label{eq:XIV.5.25}
\begin{gathered}
\slim _{h \downarrow 0} \left(\frac {\partial}{\partial t} U _{h} (t, s) - A (t) U _{h} (t, s)\right) x _{0}
\\
= \left(R (t, s) - R _{1} (t, s) - \int_{s} ^{t} R _{1} (t, \sigma) R (\sigma , s) \dd \sigma\right) x _{0}, x _{0} \in X.
\end{gathered}
\end{equation}

The right side must be $0$ as may easily be proved by \eqref{eq:XIV.5.7}. Since we have 
\[
\slim _{h \downarrow 0} A (t) U _{h} (t, s) x _{0} = A (t) U (t, s) x _{0}
\]
by the reasoning used in proving \eqref{eq:XIV.5.21}, we obtain from \eqref{eq:XIV.5.25} 
\begin{equation}\label{eq:XIV.5.26}
\slim _{h \downarrow 0} \frac {\partial}{\partial t} U _{h} (t, s) x _{0} = A (t) U (t, s) x _{0} \text{ for } t > s \text{ and } x _{0} \in X.
\end{equation}

The right side of \eqref{eq:XIV.5.26} being strongly continuous in $t > s$, we see, by integrating \eqref{eq:XIV.5.26} and remembering $\slim\limits_{h\downarrow 0}U_h(t,s)x_0 = U(t,s)x_0$, that 
\begin{equation}\label{eq:XIV.5.27}
\frac {\partial}{\partial t} U (t, s) x _{0} = A (t) U (t, s) x _{0} \text{ for } t > s \text{ and } x _{0} \in X.
\end{equation}

Therefore, $x(t) = U(t, s)$ $x_{0}$ is the desired solution of $\eqref{eq:XIV.5.3}$. The uniqueness of the solution may be proved as in the preceding section. 
\end{proof}

\yossub{Comments and References}

The above Theorem and proof are adapted from H. \textsc{Tanabe} \cite{tanabe315}. To illustrate his idea we have made Tanabe's conditions somewhat stronger. Thus, e.g., the condition \eqref{eq:XIV.5.9} may be replaced by a weaker one: 
\[
\left\| A (t) A (s) ^{-1} - A (r) A (s) ^{-1} \right\| \leqq K _{1} | t - r | ^{\varrho} \text{ with } 0 < \varrho < 1.
\]

For the details, see the above cited paper by H. \textsc{Tanabe} which is a refinement of H. \textsc{Tanabe} \cite{tanabe316} and \cite{tanabe317}. It is to be noted that the Russian school independently has developed a similar method. See, e.g., P. E. \textsc{Sobolevski} \cite{sobolevski308} and the reference cited in this paper. Cf. E. T. \textsc{Poulsen} \cite{poulsen264}. 

Komatsu's work. H. \textsc{Komatsu} \cite{komatsu153} made an important remark regarding Tanabe's result given above. Let $\Delta$ be a convex complex neighbourhood of the real segment $[a, b]$ considered as embedded in the complex plane. Suppose that $A(t)$ is defined for $t \in \Delta$ and satisfies \eqref{eq:XIV.5.8} and \eqref{eq:XIV.5.9} in which the phrase `` $t \in [a, b]$ '' is replaced by `` $t \in \Delta$ ''. Assume, moreover, that there exists a bounded linear operator $A_0$ which maps $X$ onto $D$, the domain of $A(t)$ which is assumed to be independent of $t \in \Delta$, in a one-to-one manner and such that $B(t) = A(t)A_0$ is strongly holomorphic in $t \in \Delta$. Under these assumptions, \textsc{Komatsu} proved that the operator $U(t, s)$ constructed as above is strongly holomorphic in $t \in \Delta$ if 
\[
\left| \arg (t - s) \right| < \theta_{0} \text{ with a certain } \theta_{0} \text{ satisfying } 0 < \theta_{0} < \pi / 2.
\]

The result may be applied to the ``forward and backward unique continuation'' of solutions of temporally inhomogeneous diffusion equations as in Chapter~\ref{chap:XIV}, \ref{sec:XIV.1}. In this connection, we refer to H. \textsc{Komatsu} \cite{komatsu154}, \cite{komatsu155} and T. \textsc{Kotaké}-M. \textsc{Narasimhan} \cite{konishi162}. 

Kato's work. To get rid of the assumption that the domain $D(A(t))$ is independent of $t$, T. \textsc{Kato} \cite{kato139} proved that, in the above Theorem, we may replace condition \eqref{eq:XIV.5.10} by the following: 

For a certain positive integer $k$, the domain $D\left((-A(t))^{1/k}\right)$ is independent of $t$. (Here $(-A(t))^{1/k}$ is the fractional power as defined in Chapter~\ref{chap:IX}, \ref{sec:IX.11}.) Further, there exist constants $K_{2} > 0$ and $\gamma$ with $1 - k^{-1} < \gamma \leqq 1$ such that, for $s, t \in [a, b]$, $\left\| (-A(t))^{1/k}(-A(s))^{1/k} - I \right\| \leqq K_{2}|t - s|^{\gamma}$. 

\yosparagraphtag{10'}\label{eq:XIV.5.10'} 

Tanabe's and the Kato-Tanabe recent work. With the same motivation as \textsc{Kato}, H. \textsc{Tanabe} \cite{kato134} devised a method to replace condition \eqref{eq:XIV.5.10} by 
$A(t)^{-1}$ is once strongly differentiable in $a \leqq t \leqq b$ 
and such that, for positive constants $K_{3}$ and $\alpha$, 
\begin{equation}\tag{10''}\label{eq:XIV.5.10''}
\left\| \frac {\dd A (t) ^{-1}}{\dd t} - \frac {\dd A (s) ^{-1}}{\dd s} \right\| \leqq K _{3} | t - s | ^{\alpha}.
\end{equation}

Further, there exist positive constants $N$ and $\varrho$ with $0 < \varrho \leqq 1$ such that 
\[
\left\| \frac {\partial}{\partial t} (\lambda I - A (t)) ^{-1} \right\| \leqq N | \lambda | ^{q-1}, q < 1.
\]

For details, see \textsc{Kato}-\textsc{Tanabe} \cite{kato141}. Their point is to start with the first approximation $\exp((t-a)A(t))x_0$ instead of $\exp((t-a)A(a))x_0$. 

Nelson's work on Feynman integrals. The semi-group method of integration of Schrödinger equations gives an interpretation of Feynman integrals. See E. \textsc{Nelson} \cite{nelson227}. 

The Agmon-Nirenberg work. Agmon-Nirenberg \cite{agmon3} discussed the behaviour as $t \uparrow \infty$ of solutions of the equation $\frac{1}{i} \frac{\dd u}{\dd t} - Au = 0$ in some $B$-spaces. 


\yossection{6}{Non-linear Evolution Equations 1 (The Kōmura-Kato Approach)}
\label{sec:XIV.6}

Let $X$ be a real or complex Banach space, endowed with a semi-scalar product $[x, y]$ such that (see Chapter~\ref{chap:IX}, \ref{sec:IX.8}) 
\begin{equation}\label{eq:XIV.6.1}
\begin{gathered}
\left[ \alpha_{1} x _{1} + \alpha_{2} x _{2}, y \right] = \alpha_{1} \left[ x _{1}, y \right] + \alpha_{2} \left[ x _{2}, y \right], | [ x, y ] | \leqq \| x \| \cdot \| y \| \quad \text{ and }
\\
{}[ x, x ] = \| x \| ^{2},
\end{gathered}
\end{equation}
and consider a family $\{T_t; t \geqq 0\}$ of non-linear mappings of $X$ into $X$ satisfying conditions: 

$\int T_{t}T_{s} = T_{t+s}, T_{0} =$ the identity mapping $I$ (semi-group property), 
$\left\|T_{t}x-T_{t}y\right\|\leqq\left\|x-y\right\|\text{ (contraction property) }$, 
and strong continuity in $t$ of $T_{t}$. As in the case of linear mappings, we define the infinitesimal generator $A$ of $\{T_{t}; t \geqq 0\}$ by 
\[
A \cdot x = \slim _{h \downarrow 0} h ^{-1} \left(T _{t} x - x\right).
\]

Then $A$ must be dissipative in the following sense: 
\begin{equation}\label{eq:XIV.6.2}
\operatorname{Re} [ A x - A y, x - y ] \leqq 0.
\end{equation}

The proof is easy, since we have 
\begin{gather}\operatorname{Re} \left[ h ^{-1} (T _{h} x - x) - h ^{-1} (T _{h} y - y), x - y \right] = {} h ^{-1} \operatorname{Re} \left[ T _{h} x - T _{h} y, x - y \right] - h ^{-1} [ x - y, x - y ] \notag \\\leqq {} h ^{-1} \| T _{h} x - T _{h} y \| \cdot \| x - y \| - h ^{-1} \| x - y \| ^{2} \leqq {} h ^{-1} \| x - y \| ^{2} - h ^{-1} \| x - y \| ^{2} = 0. \notag
\end{gather}

A celebrated example due to Y. \textsc{Kōmura} \cite{komura158} is this: Let $X = \mathbb{R}^{1}$ with $[x, y] = x \cdot y$ and $\|x\| = |x|$, and 
\[
T _{t} x = \begin{cases}\max (x - t, 0) & \text{ for } \quad x > 0, \\ x & \text{ for } \quad x \leqq 0,\end{cases} \quad A x = \begin{cases}- 1 & \text{ for } \quad x > 0, \\ 0 & \text{ for } \quad x \leqq 0.\end{cases}
\]

By virtue of \eqref{eq:XIV.6.2} we can prove that, for all $\lambda > 0$, the mapping $(I - \lambda A)$ of $D(A)$ into $X$ has an inverse $J_{\lambda} = (I - \lambda A)^{-1}$, because 
\[
x _{1} - \lambda A x _{1} = z = x _{2} - \lambda A x _{2}
\]
implies 
\[
\begin{aligned}
0 &= [ (x _{1} - \lambda A x _{1}) - (x _{2} - \lambda A x _{2}), x _{1} - x _{2} ]
\\
\quad &= [ x _{1} - x _{2}, x _{1} - x _{2} ] - \lambda [ A x _{1} - A x _{2}, x _{1} - x _{2} ] \text{ so that }
\\
0 &= \| x _{1} - x _{2} \| ^{2} - \lambda \operatorname{Re} [ A x _{1} - A x _{2}, x _{1} - x _{2} ] \geqq \| x _{1} - x _{2} \| ^{2}, \text{ i.e., } x _{1} = x _{2}.
\end{aligned}
\]

Therefore, the theory of linear contraction semi-groups (Chapter~\ref{chap:IX}, \ref{sec:IX.8}) suggests us to approximate the non-linear evolution equation 
$\frac{\dd u(t)}{\dd t} = A u(t)$ for $t \geqq 0$ with $u(0) = x_0 \in D(A)$ by equations $(\lambda > 0)$ $\frac{\dd u^{(\lambda)}(t)}{\dd t} = A_{\lambda} u^{(\lambda)}(t)$ for $t \geqq 0$ with $u^{(\lambda)}(0) = x_0$, where $A_{\lambda} = \lambda^{-1}(J_{\lambda} - 1)$, under the assumption that $D(J_{\lambda}) = R(I - \lambda A) = X$, expecting that we shall have $u(t) = \slim\limits_{\lambda \downarrow 0} u^{(\lambda)}(t)$. 

In the above example of \textsc{Kōmura}, we have $D(J_{\lambda}) = (-\infty, 0] \cup (\lambda, \infty)$ which does not coincide with $X = \mathbb{R}^{1}$. To obtain $D(J_{\lambda}) = X = \mathbb{R}^{1}$, it will be necessary to extend $A$ to a multi-valued mapping $\hat{A}$ by 
\[
\hat {A} x = \begin{cases}- 1 & \text{ for } x > 0, \\ - [ 0, 1 ] & \text{ for } x = 0, \\ 0 & \text{ for } x < 0,\end{cases}
\]
for which the dissipative property is preserved in the following sense: $\operatorname{Re}\langle y_1 - y_2, x_1 - x_2 \rangle \leqq 0$ for any $\{x_i, y_i\}$ with $y_i \in \hat{A} x_i (i = 1, 2)$. \yosparagraphtag{2'}\label{eq:XIV.6.2'} 

In this way, we are led to the following setup for our further discussions. Let $X$ be a Banach space whose dual space is denoted by $X'$. An element of $X \times X$ will be written in the form $\{x, y\}$ where $x$ and $y$ are both $\in X$. For a subset $A$ of $X \times X$, the following notations will conveniently be used: 

$1$) $D(A) = \{x; \{x, y\} \in A \text{ for some } y\}$, 
$2$) $R(A) = \{y; \{x, y\} \in A \text{ for some } x\}$, 
$3$) $A^{-1} = \{\{y, x\}; \{x, y\} \in A\}$, 
$4$) $\lambda A = \{\{x, \lambda y\}; \{x, y\} \in A$ and for real $\lambda\}$, 
$5$) $A + B = \{\{x, y + z\}; \{x, y\} \in A \text{ and } \{x, z\} \in B\}$, 
$6$) $A x = \{y; \{x, y\} \in A$ for a fixed $x \in D(A)\}$, 
$7$) $\|Ax\| = \inf \{|y\|; y \in Ax\}$, 
$8$) $J_{\lambda} = (I - \lambda A)^{-1} = \{\{x - \lambda y, x\} ; \{x, y\} \in A\}$ for real $\lambda$ and $I = \{\{x, x\} ; x \in X\}$, 
$9$) $A_{\lambda} = \{\{x - \lambda y, y\}; \{x, y\} \in A\}$; 

we have 
\begin{equation}\label{eq:XIV.6.3}
\lambda A _{\lambda} = J _{\lambda} - I.
\end{equation}

If $A x$ consists of a single element for all $x \in D(A)$, then $A$ is the graph of a uniquely determined function ( $= \text{ single-valued mapping }$) from $D(A)$ with values in $X$; in this case, we may identify the set $A$ and the above function. 

To state the notion of dissipative sets in $X \times X$, we give 

\begin{definition}\label{def:XIV.6.Definition.1} By the duality map of $X$ into $X'$ we mean the multi-valued mapping $F$ from $X$ into $X'$ defined by 
\begin{equation}\label{eq:XIV.6.4}
F (x) = \{f \in X ^{\prime}; \langle x, f \rangle = \| x \| ^{2} = \| f \| ^{2} \}.
\end{equation}
\end{definition}

That $F(x)$ is not void is clear from the \textsc{Hahn}-\textsc{Banach} theorem. If $X$ is a Hilbert space, then F. \textsc{Riesz}' representation theorem asserts that $F(x)$ consists of $x$ alone and $\langle y, F(x) \rangle = (y, x)$, the scalar product of $y$ and $x$. 

\begin{definition}\label{def:XIV.6.Definition.2} $A$ set $A$ in $X \times X$ is called a dissipative set if, for arbitrary two points $\{x_{1}, y_{1}\}$ and $\{x_{2}, y_{2}\}$ of $A$, there exists an $f \in F(x_{1} - x_{2})$ such that $\operatorname{Re}\langle y_{1} - y_{2}, f \rangle \leqq 0$. 
\end{definition}

\begin{remarku} $A$ is a dissipative set iff $-A$ is a monotonic set in the sense of G. \textsc{Minty} \cite{minty203}. 
\end{remarku}

We also give 

\begin{definition}\label{def:XIV.6.Definition.3} Let $D$ be a subset of $X$, and $T$ be a function from $D$ into $X$. $T$ is called a Lipschitzian mapping (function) with Lipschitz constant $k > 0$ if $\|Tx_{1} - Tx_{2}\| \leqq k\|x_{1} - x_{2}\|$ for all $x_{1}, x_{2} \in D$. If we can take $k = 1$, then $T$ is called a contraction mapping (function). We have 
\end{definition}

\begin{lemma}[T. \textsc{Kato} {\cite{kato144}}]\label{lem:XIV.6.Lemma.1} Let $x, y \in X$. Then $\|x - \lambda y\| \geqq \|x\|$ for all $\lambda > 0$ iff there exists an $f \in F(x)$ such that $\operatorname{Re}\langle y, f \rangle \leqq 0$. 
\end{lemma}

\begin{proof}
 The assertion is trivial if $x = 0$. So we shall assume $x \neq 0$ in the following. If $\operatorname{Re}\langle y, f \rangle \leqq 0$ for some $f \in F(x)$, then $\|x\|^2 = \langle x, f \rangle = \operatorname{Re}\langle x, f \rangle \leqq \operatorname{Re}\langle x - \lambda y, f \rangle \leqq \|x - \lambda y\| \cdot \|f\|$. Since $\|x\| = \|f\|$, we obtain $\|x\| \leqq \|x - \lambda y\|$. 

Suppose, conversely, that $\|x\| \leqq \|x - \lambda y\|$ for all $\lambda > 0$. For each $\lambda > 0$, let $f_{\lambda} \in F(x - \lambda y)$ and put $g_{\lambda} = f_{\lambda} / \|f_{\lambda}\|$ so that $\|g_{\lambda}\| = 1$. Then $\|x\| \leqq \|x - \lambda y\| = \langle x - \lambda y, g_{\lambda} \rangle = \operatorname{Re}\langle x, g_{\lambda} \rangle - \lambda \operatorname{Re}\langle y, g_{\lambda} \rangle \leqq \|x\| - \lambda \operatorname{Re}\langle y, g; \rangle$ by $\|g_{\lambda}\| = 1$. Hence 
\[
\liminf\limits _{\lambda \downarrow 0} \operatorname{Re} \left\langle x, g _{\lambda} \right\rangle = \| x \| \quad \text{ and } \quad - \lambda \operatorname{Re} \left\langle y, g _{\lambda} \right\rangle \geqq 0.
\]

Since the closed unit sphere of the dual space $X'$ is weakly* compact (see Theorem~\hyperref[thm:AV.1.Theorem.1]{$1$} in the Appendix to Chapter~\ref{chap:V}, \ref{sec:AV.1}) the sequence $\{g_{1/n}\}$ has a weak* accumulation point $g \in X'$ with $\|g\| \leqq 1$. Thus we see that $g$ must satisfy $\operatorname{Re}\langle x, g \rangle \geqq \|x\|$ and $\operatorname{Re}\langle y, g \rangle \leqq 0$, that is, we must have $\|g\| = 1$ and $\langle x, g \rangle = \|x\|$. Therefore $f = \|x\| \cdot g$ satisfies $f \in F(x)$ and $\operatorname{Re}\langle y, f \rangle \leqq 0$. 
\end{proof}

\begin{corollaryu} $A$ is a dissipative set iff 
\begin{equation}\label{eq:XIV.6.5}
\begin{aligned}
&{\| (x _{1} - \lambda y _{1}) - (x _{2} - \lambda y _{2}) \| \geqq \| x _{1} - x _{2} \| \text{ whenever } \lambda > 0} {\quad \text{ and } \quad \{x _{i}, y _{i} \} \in A (i = 1, 2).}
\end{aligned}
\end{equation}
\end{corollaryu}

We have the following 

\begin{proposition}\label{prop:XIV.6.Proposition.1} If $A$ is a dissipative set and $\lambda > 0$, then $J_{\lambda}$ and $A_{\lambda}$ are both single-valued mappings, and 
\begin{align}
\left\| J _{\lambda} x _{1} - J _{\lambda} x _{2} \right\|& \leqq \left\| x _{1} - x _{2} \right\| \quad \text{ for } \quad x _{1}, x _{2} \in D (J _{\lambda}), \label{eq:XIV.6.6} \\
\left\| A _{\lambda} x _{1} - A _{\lambda} x _{2} \right\|& \leqq \frac {2}{\lambda} \left\| x _{1} - x _{2} \right\| \quad \text{ for } \quad x _{1}, x _{2} \in D (A _{\lambda}) = D (J _{\lambda}). \label{eq:XIV.6.7}
\end{align}
moreover, $A_{\lambda}$ is dissipative and 
\begin{align}
A J _{\lambda} x& = A (J _{\lambda} x) \ni A _{\lambda} x \quad \text{ for } \quad x \in D (J _{\lambda}), \label{eq:XIV.6.8} \\
\left\| A _{\lambda} x \right\|& \leqq \left\| A x \right\| \quad \text{ for all } \quad x \in D (A) \cap D (J _{\lambda}). \label{eq:XIV.6.9}
\end{align}
\end{proposition}

\begin{proof}
 That $J_{\lambda}$ and $A_{\lambda}$ are both single-valued is clear from \eqref{eq:XIV.6.5}. We also have \eqref{eq:XIV.6.6} by \eqref{eq:XIV.6.5}, and thus \eqref{eq:XIV.6.7} is proved by \eqref{eq:XIV.6.3} and \eqref{eq:XIV.6.6}. Next let $f \in F(x_1 - x_2)$. Then we obtain, by \eqref{eq:XIV.6.3}, \eqref{eq:XIV.6.4}, and \eqref{eq:XIV.6.6}, 
\[
\begin{gathered}
\operatorname{Re} \left\langle A _{\lambda} x _{1} - A _{\lambda} x _{2}, f \right\rangle = \lambda^{-1} \operatorname{Re} \left\langle (J _{\lambda} x _{1} - x _{1}) - (J _{\lambda} x _{2} - x _{2}), f \right\rangle
\\
= \lambda^{-1} \operatorname{Re} \left\langle J _{\lambda} x _{1} - J _{\lambda} x _{2}, f \right\rangle - \lambda^{-1} \left\langle x _{1} - x _{2}, f \right\rangle \leqq \lambda^{-1} \| J _{\lambda} x _{1} - J _{\lambda} x _{2} \| \cdot \| f \|
\\
- \lambda^{-1} \left\langle x _{1} - x _{2}, f \right\rangle \leqq \lambda^{-1} \| x _{1} - x _{2} \| ^{2} - \lambda^{-1} \| x _{1} - x _{2} \| ^{2} = 0. \end{gathered}
\]

This proves that $A_{\lambda}$ is dissipative. \eqref{eq:XIV.6.8} is clear. \eqref{eq:XIV.6.9} is proved as follows. By \eqref{eq:XIV.6.3} and \eqref{eq:XIV.6.6}, we obtain, for any $y \in Ax$, 
\[
\lambda \| A _{\lambda} x \| = \| J _{\lambda} x - x \| = \| J _{\lambda} x - J _{\lambda} (x - \lambda y) \| \leqq \| x - (x - \lambda y) \| = \lambda \| y \|.
\]
\end{proof}

\begin{lemma}[Y. \textsc{Kōmura} {\cite{komura158}}]\label{lem:XIV.6.Lemma.2} Let $A$ be a dissipative set $\subseteq X \times X$, and assume that $D(J_{\lambda}) = X$ for a $\lambda > 0$. Then $D(J_{\mu}) = X$ for every $\mu > 0$ with $0 \leqq |(\mu - \lambda)/\mu| < 1$. 
\end{lemma}

\begin{proof}
 Take any point $x \in X$, and consider a single-valued mapping $T$ defined by 
\[
X \ni z \rightarrow T z = J _{\lambda} \left(\frac {\lambda}{\mu} x + \frac {\mu - \lambda}{\mu} z\right). \]

Since $J_{\lambda}$ is a contraction mapping by \eqref{eq:XIV.6.6}, we obtain 
\[
\left\| T z - T w \right\| \leqq \left\| \left(\frac {\lambda}{\mu} x + \frac {\mu - \lambda}{\mu} z\right) - \left(\frac {\lambda}{\mu} x + \frac {\mu - \lambda}{\mu} w\right) \right\| = \left| \frac {\mu - \lambda}{\mu} \right| \cdot \| z - w \|.
\]
that is, $T$ is a Lipschitzian mapping with Lipschitz constant 
\[
\alpha = | (\mu - \lambda) / \mu | < 1.
\]

Thus, for $n > m$ and for any point $z \in X$, we have 
\[
\begin{aligned} & {\| T ^{n} z - T ^{m} z \| \leqq \alpha^{m} \| T ^{n-m} z - z \| \leqq \alpha^{m} (\| T z - z \| + \| T ^{2} z - T z \| + \dots)} \\ & {\quad \leqq \alpha^{m} (1 + \alpha + \alpha^{2} + \dots) \cdot \| T z - z \| \leqq \alpha^{m} (1 - \alpha) ^{-1} \| T z - z \|.} \end{aligned}
\]

Hence, by the completeness of the space $X$, $\slim\limits_{n\to 0}T^n z = y$ exists in $X$. $T$ being continuous as a Lipschitzian mapping, we obtain $T\cdot y = y$ by $y = \slim\limits_{n\to 0}T^{n+1}z = \slim\limits_{n\to 0}T(T^n z)$. Hence 
\[ y = J _{\lambda} \left(\frac {\lambda}{\mu} x + \frac {\mu - \lambda}{\mu} y\right) = J _{\lambda} \left(y - \lambda \left(\frac {1}{\mu} y - \frac {1}{\mu} x\right)\right),
\]
that is, $\left(\frac {1}{\mu} y - \frac {1}{\mu} x\right) \in A y$, and so $y - \mu z = x$ with $z \in Ay$. This proves that $J_{\mu}x = y$. Since $x$ was arbitrary, we must have $D(J_{\mu}) = X$. 
\end{proof}

\begin{corollaryu} Repeating the above argument, we can prove that $D(J_{\mu}) = X$ for all $\mu > 0$. 
\end{corollaryu}

We are now able to give 

\begin{definition}\label{def:XIV.6.Definition.4} $A$ dissipative set $A \subseteq X \times X$ is called a hyperdissipative set if $D(J_{\lambda}) = X$ for some $\lambda > 0$ and hence for all $\lambda > 0$. 
\end{definition}

\begin{proposition}\label{prop:XIV.6.Proposition.2} $A$ hyperdissipative set $A \subseteq X \times X$ must be maximally dissipative in the sense that there does not exist a dissipative set $B \subseteq X \times X$ which contains $A$ as a proper subset. 
\end{proposition}

\begin{proof}
 Let us assume that a dissipative set $B \subseteq X \times X$ contains $A$ as a subset. Let $y \in Bx$. Then, since $A$ is a hyperdissipative set, there exists a point $\{x_{1}, y_{1}\} \in A$ such that $x - y = x_{1} - y_{1}$. Since $A \subseteq B$, we must have $\{x_{1}, y_{1}\} \in B$ and hence we obtain $x = x_{1}$ and $y = y_{1}$ by \eqref{eq:XIV.6.5} applied to the dissipative set $B$. 
\end{proof}

We have thus prepared tools for the exposition of \textsc{Kōmura}'s approach to the integration in Hilbert space of non-linear evolution equations. Let $H$ be a real or complex Hilbert space with the scalar product $(x, z)$, and $A \subseteq H \times H$ be a hyperdissipative set. We shall be concerned with strong solutions of the initial value problem: 
\begin{equation}\label{eq:XIV.6.10}
\begin{yoslines}
&\frac {\dd u (t)}{\dd t} \in A u (t) \quad \text{ for almost every } t \text{ on the interval } [ 0, \infty), u (0) = x _{0} \in D (A),
\end{yoslines}
\end{equation}
where $u(t)$ defined on $[0, \infty)$ with values in $H$ is called a strong solution of \eqref{eq:XIV.6.10} if $u(t)$ is strongly absolutely continuous in $t$, $u(t)$ is strongly differentiable almost everywhere on $[0, \infty)$ such that the strong differential quotient $\dd u(t)/\dd t$ is Bochner integrable on any compact interval of $t$ and $u$ satisfies \eqref{eq:XIV.6.10}. We shall approximate this problem \eqref{eq:XIV.6.10} by 
\begin{equation}\tag{10'}\label{eq:XIV.6.10'}
\begin{yoslines}& \frac {\dd u ^{(\lambda)} (t)}{\dd t} - A _{\lambda} u ^{(\lambda)} (t) = 0 \quad \text{ for all } t \quad \text{ on } [ 0, \infty) \quad (\lambda > 0),\\& u ^{(\lambda)} (0) = x _{\lambda} = x _{0} - \lambda y _{0} \quad \text{ with a fixed } \quad y _{0} \in A x _{0}.\end{yoslines}
\end{equation}

Since $A_{\lambda}$ with $\lambda > 0$ is a Lipschitzian mapping, we see that equation \eqref{eq:XIV.6.10'} admits a uniquely determined solution $u^{(\lambda)}(t)$ which can be obtained by E. \textsc{Picard}'$S$ successive approximation: 
\begin{equation}\tag{10''}\label{eq:XIV.6.10''}
u _{n+1} ^{(\lambda)} (t) = x _{\lambda} + \int_{0} ^{t} A _{\lambda} u _{n} ^{(\lambda)} (s) \dd s \quad (n = 0, 1 \dots ; u _{0} ^{(\lambda)} (t) = x _{\lambda}).
\end{equation}

\begin{lemma}[Y. \textsc{Kōmura} {\cite{komura158}}]\label{lem:XIV.6.Lemma.3} For the solution $u^{(\lambda)}(t)$ of $\eqref{eq:XIV.6.10'}$, we have the following estimates: 
\begin{equation}\label{eq:XIV.6.11}
\left\| \frac {d}{\dd t} u ^{(\lambda)} (t) \right\| \leqq \left\| \frac {d}{\dd s} u ^{(\lambda)} (s) \right\| \quad \text{ for } 0 \leqq s < t,
\end{equation}
and 
\begin{equation}\label{eq:XIV.6.12}
\left\| u ^{(\lambda)} (t) - u ^{(\mu)} (t) \right\| ^{2} \leqq \left((\lambda - \mu) ^{2} + 4 t (\lambda + \mu)\right) \left\| y _{0} \right\| ^{2}.
\end{equation}
\end{lemma}

\begin{proof}
 Putting $u^{(\lambda)}(t + h) = v^{(\lambda)}(t)$ with $h > 0$, we obtain 
\[
\begin{aligned} \frac {d}{\dd t} \| v ^{(\lambda)} (t) - u ^{(\lambda)} (t) \| ^{2} & = 2 \operatorname{Re} \left(\frac {d}{\dd t} v ^{(\lambda)} (t) - \frac {d}{\dd t} u ^{(\lambda)} (t), v ^{(\lambda)} (t) - u ^{(\lambda)} (t)\right) \\ & = 2 \operatorname{Re} \left(A _{\lambda} v ^{(\lambda)} (t) - A _{\lambda} u ^{(\lambda)} (t), v ^{(\lambda)} (t) - u ^{(\lambda)} (t)\right) \leqq 0 \end{aligned}
\]
by the dissipative property of $A_{\lambda}$. Thus $\|u^{(\lambda)}(t + h) - u^{(\lambda)}(t)\|$ is monoton decreasing in $t$ and so we obtain \eqref{eq:XIV.6.11}. Therefore we have 
\begin{equation}\label{eq:XIV.6.13}
\left\| \frac {d}{\dd t} u ^{(\lambda)} (t) \right\| = \| A _{\lambda} u ^{(\lambda)} (t) \| \leqq \| A _{\lambda} u ^{(\lambda)} (0) \| = \| A _{\lambda} (x _{0} - \lambda y _{0}) \| = \| y _{0} \|. \end{equation}

By a similar argument as above, we obtain, by $\lambda A_{\lambda} = J_{\lambda} - I$, 
\begin{gather}\left\| u ^{(\lambda)} (t) - u ^{(\mu)} (t) \right\| ^{2} - \left\| u ^{(\lambda)} (0) - u ^{(\mu)} (0) \right\| ^{2} = \left\| u ^{(\lambda)} (t) - u ^{(\mu)} (t) \right\| ^{2} - \left\| (\lambda - \mu) y _{0} \right\| ^{2} \notag
\\= \int_{0} ^{t} \frac {d}{\dd s} \left\| u ^{(\lambda)} (s) - u ^{(\mu)} (s) \right\| ^{2} \dd s = 2 \int_{0} ^{t} \operatorname{Re} \left(A _{\lambda} u ^{(\lambda)} (s) - A _{\mu} u ^{(\mu)} (s), \right. \left. u ^{(\lambda)} (s) - u ^{(\mu)} (s)\right) \dd s \notag
\\= 2 \int_{0} ^{t} \operatorname{Re} \left(A _{\lambda} u ^{(\lambda)} (s) - A _{\mu} u ^{(\mu)} (s), \right. \left. J _{\lambda} u ^{(\lambda)} (s) - J _{\mu} u ^{(\mu)} (s)\right) \dd s \notag
\\\quad - 2 \int_{0} ^{t} \operatorname{Re} \left(A _{\lambda} u ^{(\lambda)} (s) - A _{\mu} u ^{(\mu)} (s), \right. \left. \lambda A _{\lambda} u ^{(\lambda)} (s) - \mu A _{\mu} u ^{(\mu)} (s)\right) \dd s. \notag
\end{gather}

The first term on the extreme right is non-negative, because $A$ is a dissipative set in $H \times H$ and $A_{\lambda}u^{(\lambda)}(s) \in A J_{\lambda}u^{(\lambda)}(s)$ by \eqref{eq:XIV.6.8}. The second term on the extreme right is, by \eqref{eq:XIV.6.13} smaller than $4t(\lambda + \mu)\|y_{0}\|^{2}$ so that 
\[
\left\| u ^{(\lambda)} (t) - u ^{(\mu)} (t) \right\| ^{2} - (\mu - \lambda) ^{2} \left\| y _{0} \right\| ^{2} \leqq 4 t (\lambda + \mu) \left\| y _{0} \right\| ^{2}.
\]

This proves \eqref{eq:XIV.6.12}. 
\end{proof}

\begin{corollaryu} 
\begin{equation}
\label{eq:XIV.6.14}
\begin{minipage}{0.92\linewidth}
\everymath{}\let\limits\nolimits
$\slim_{\lambda \downarrow 0} u^{(\lambda)}(t) = u(t)$ exists uniformly on every compact set of $t.$
\end{minipage}
\end{equation}
\end{corollaryu}

\begin{equation}
\label{eq:XIV.6.15}
\begin{minipage}{0.92\linewidth}
\everymath{}\let\limits\nolimits
$\slim _{\lambda \downarrow 0} J _{\lambda} u ^{(\lambda)} (t) = u (t)$ uniformly on every compact set of $t.$
\end{minipage}
\end{equation}

\begin{proof}
 ($14$) is clear from \eqref{eq:XIV.6.12}. We also have \eqref{eq:XIV.6.15} from ($14$), since we have 
\[
\left\| J _{\lambda} u ^{(\lambda)} (t) - u ^{(\lambda)} (t) \right\| = \lambda \left\| A _{\lambda} u ^{(\lambda)} (t) \right\| \leqq \lambda \left\| y _{0} \right\| \text{ by } (1 3). \]
\end{proof}

\begin{lemma}[Y. \textsc{Kōmura} {\cite{komura158}}]\label{lem:XIV.6.Lemma.4} $u(t) \in D(A)$ for all $t \geqq 0$. 
\end{lemma}

\begin{proof}
 By \eqref{eq:XIV.6.8} we have, for fixed $t > 0$, 
\begin{equation}\label{eq:XIV.6.16}
\left\{J _{\lambda} u ^{(\lambda)} (t), A _{\lambda} u ^{(\lambda)} (t) \right\} \in A \quad \text{ for every } \quad \lambda > 0. \end{equation}

Thus, setting $A_{\lambda}u^{(\lambda)}(t) = w^{(\lambda)}(t)$, we obtain $\|w^{(\lambda)}(t)\| \leqq \|y_0\|$ by \eqref{eq:XIV.6.13}. Hence, by the local weak compactness of the Hilbert space $H$, there exists, for fixed $t > 0$, a sequence $\{\lambda_n\}$ of positive numbers satisfying $\lambda_n \downarrow 0$ and such that 
\begin{equation}\label{eq:XIV.6.17}
\wlim _{n \to \infty} w ^ {(\lambda_{n})} (t) = w (t) \quad H \quad \text{ with } \quad \| w (t) \| \leqq \| y _{0} \|.
\end{equation}

Since $A$ is a dissipative set, we have, for every point $\{x, y\} \in A$, the inequality $\operatorname{Re}(y - w^{(\lambda_n)}(t), x - J_{\lambda_n}u^{(\lambda_n)}(t)) \leqq 0$. Hence, by letting $n \uparrow \infty$, we obtain $\operatorname{Re}(y - w(t), x - u(t)) \leqq 0$ by \eqref{eq:XIV.6.15}. This proves that 
\begin{equation}\label{eq:XIV.6.18}
\{u (t), w (t) \} \in A,
\end{equation}
because $A$ is a maximally dissipative set $\subseteq H \times H$ as proved in Proposition~\ref{prop:XIV.6.Proposition.2}. 
\end{proof}

\begin{lemma}[Y. \textsc{Kōmura} {\cite{komura158}}]\label{lem:XIV.6.Lemma.5} $u(t)$ is strongly absolutely continuous in $t$, and it is strongly differentiable at almost every $t \geqq 0$. 
\end{lemma}

\begin{proof}
 By \eqref{eq:XIV.6.13}, we obtain, for $0 \leqq t_{1} < t_{2} < \cdots$ with $\sum\limits_{i} (t_{i+1} - t_{i}) < \infty$, 
\[
\left\| u ^{\lambda} \left(t _{i+1}\right) - u ^{(\lambda)} \left(t _{i}\right) \right\| \leqq \int_ {t _{i}} ^ {t _{i+1}} \left\| \frac {d}{\dd s} u ^{(\lambda)} (s) \right\| \dd s \leqq \left(t _{i+1} - t _{i}\right) \cdot \left\| y _{0} \right\|. \]

Therefore, by letting $\lambda \downarrow 0$, we obtain 
\begin{equation}\label{eq:XIV.6.19}
\sum\limits_{i} \left\| u \left(t _{i+1}\right) - u \left(t _{i}\right) \right\| \leqq \sum\limits_{i} \left(t _{i+1} - t _{i}\right) \cdot \left\| y _{0} \right\|.
\end{equation}

This proves the strong absolute continuity in $t$ of $u(t)$. 

Next let $0 < t_0 < \infty$. Then, by \eqref{eq:XIV.6.19}, the set $\{u(t); 0 \leqq t \leqq t_0\}$ is compact in $H$, hence it is separable. Therefore, without losing the generality, we may assume that $H$ is separable. Thus let $\{x_k\}$ be a strongly dense countable sequence of $H$. Each numerical function $v_k(t) = (u(t), x_k)$ is absolutely continuous by \eqref{eq:XIV.6.19}, and so there exists a set $N_k$ of measure zero such that $v_k(t)$ is differentiable on $[0, t_0] - N_k$. In view of \eqref{eq:XIV.6.19}, we see that $u(t)$ is weakly differentiable on $[0, t_0] - \bigcup\limits_{k=1}^{\infty} N_k$, and the weak derivative $u'(t)$ satisfies $\|u'(t)\| \leqq \|y_0\|$. $H$ being assumed separable, the weakly measurable function $u'(t)$ is strongly measurable by \textsc{Pettis}' theorem (Chapter~\ref{chap:V}, \ref{sec:V.4}) and so, by the boundedness condition $\|u'(t)\| \leqq \|y_0\|$, $u'(t)$ is Bochner integrable. Therefore, by \textsc{Bochner}'$S$ theorem (Chapter~\ref{chap:V}, \ref{sec:V.5}), 
\begin{equation}\label{eq:XIV.6.20}
u (t) - u (0) = \int_{0} ^{t} u ^{\prime} (s) \dd s
\end{equation}
is strongly differentiable for almost every $t \in [0, t_0]$ with $u'(t)$ as the strong derivative. 
\end{proof}
We are now able to prove Y. \textsc{Kōmura}'s 

\begin{theoremu} $u(t)$ is a strong solution of the initial value problem \eqref{eq:XIV.6.10}. 
\end{theoremu}

\begin{proof}
 For a fixed positive number $t_0$, we define the space $L^2([0, t_0], H)$ of $H$-valued strongly measurable function $\tilde{x} = x(t)$ such that $\int_{0}^{t_0} \|x(s)\|^2 \dd s < \infty$, and normed by $\|\tilde{x}\|^2 = \int_{0}^{t_0} \|x(s)\|^2 \dd s$. It is easy to show that $L^2([0, t_0], H)$ is a Hilbert space with the scalar product 
\begin{equation}\label{eq:XIV.6.21}
(\tilde {x}, \tilde {y}) = \int_{0} ^ {t _{0}} (x (s), y (s)) \dd s,
\end{equation}
where $(x(s), y(s))$ is the scalar product in $H$ of $x(s)$ and $y(s)$. 

Thus our $\tilde{u} = u(t)$ can be considered as an element of $L^2([0, t_0], H)$, because $u(t)$ is a strongly continuous $H$-valued function. Moreover, the hyperdissipative set $A$ can naturally be extended to a hyperdissipative set $\tilde{A} \subseteq L^2([0, t_0], H) \times L^2(0, t_0], H$ by defining 
\begin{equation}\label{eq:XIV.6.22}
\tilde {A} \tilde {x} = \{\tilde {y} \in L ^{2} ([ 0, t _{0} ], H); y (t) \in A x (t) \text{ for almost every } t \in [ 0, t _{0} ] \}
\end{equation}

We have, by \eqref{eq:XIV.6.16} and \eqref{eq:XIV.6.13}, 
\begin{equation}\tag{16'}\label{eq:XIV.6.16'}
\{\widetilde {J} _{\lambda} \widetilde {u} ^{(\lambda)}, \widetilde {A} _{\lambda} \widetilde {u} ^{(\lambda)} \} \in \tilde {A} \quad \text{ for every } \lambda > 0.
\end{equation}

Again by \eqref{eq:XIV.6.13}, $\{\widetilde{A}_{\lambda}\widetilde{u}^{(\lambda)};\lambda >0\}$ is a norm bounded subset of the Hilbert space $L^2 ([0,t_0],H)$. Hence there exists a sequence $\{\lambda_n\}$ of positive numbers satisfying $\lambda_{n}\downarrow 0$ and such that 
\begin{equation}\tag{17'}\label{eq:XIV.6.17'}
\wlim \widetilde {A} _ {\lambda_{n}} \widetilde {u} ^ {(\lambda_{n})} = \tilde {w} \in L ^{2} ([ 0, t _{0} ], H) .
\end{equation}

Moreover, we have, by \eqref{eq:XIV.6.15}, 
\begin{equation}\tag{15'}\label{eq:XIV.6.15'}
\slim _{n \rightarrow \infty} \widetilde {J} _ {\lambda_{n}} \widetilde {u} ^ {(\lambda_{n})} = \tilde {u} \in L ^{2} ([ 0, t _{0} ], H).
\end{equation}
$\tilde{A}$ is a hyperdissipative set by the fact that $A$ is a hyperdissipative set. Thus, by Proposition~\ref{prop:XIV.6.Proposition.2}, $\tilde{A}$ is a maximally dissipative set. Therefore, we have, as in the proof of Lemma~\ref{lem:XIV.6.Lemma.4}, 
\begin{equation}\tag{18'}\label{eq:XIV.6.18'}
w (t) \in A u (t) \quad \text{ for almost every } \quad t \in [ 0, t _{0} ] .
\end{equation}

On the other hand, we have 
\[
u ^ {(\lambda_{n})} (t) - u ^ {(\lambda_{n})} (0) = \int_{0} ^{t} \frac {d}{\dd s} u ^ {(\lambda_{n})} (s) \dd s \quad (\text{ for } 0 \leqq t \leqq t _{0}),
\]
and so, for every $x \in H$, 
\[
(u ^ {(\lambda_{n})} (t), x) - \langle x _ {\lambda_{n}}, x \rangle = \int_{0} ^{t} \left(\frac {d}{\dd s} u ^ {(\lambda_{n})} (s), x\right) \dd s \quad (\text{ for } \quad 0 \leqq t \leqq t _{0})
\]
by Colollary $2$ in Chapter~\ref{chap:V}, \ref{sec:V.5}. Hence, we obtain, by ($14$) and \eqref{eq:XIV.6.17'}, 
\begin{equation}\label{eq:XIV.6.23}
(u (t), x) - (x _{0}, x) = \int_{0} ^{t} (w (s), x) \dd s = (\int_{0} ^{t} w (s) \dd s, x).
\end{equation}

Since $x \in H$ and $t_0 > 0$ were arbitrary, we obtain, by \eqref{eq:XIV.6.23}, \eqref{eq:XIV.6.18'} and \eqref{eq:XIV.6.20}, that $\frac{\dd u(t)}{\dd t} = u'(t) = w(t) \in Au(t)$ for almost every $t \geqq 0$. 
\end{proof}

\begin{remarku} T. \textsc{Kato} \cite{kato144} extended the above Theorem to the case where $X$ is a Banach space whose dual space is uniformly convex. An almost complete version of the Hille-Yosida theorem for non-linear contraction semi-groups was devised and proved in Hilbert space by Y. \textsc{Kōmura} \cite{komura159}. Cf. M. G. \textsc{Crandall}-A. \textsc{Pazy} \cite{clarkson36}, T. \textsc{Kato} \cite{kato145} and J. R. \textsc{Dorroh} \cite{doob43} cited in § $5$ of Y. \textsc{Kōmura} \cite{komura159}. In this paper of \textsc{Kōmura}, the most crucial point is the proof that the infinitesimal generator is densely defined, and this proof was later fairly simplified by T. \textsc{Kato} \cite{kato145}. Recently, H. \textsc{Brezis} has published a comprehensive treatise on nonlinear semi-groups in Hilbert space. This book (H. \textsc{Brezis} \cite{brezis28}) contains an extensive bibliography which is not necessarily restricted to the Hilbert space. 
\end{remarku}


\yossection[7. Non-linear Evolution Equations 2]{7}{Non-linear Evolution Equations 2 (The Approach through the Crandall-Liggett Convergence Theorem)}
\label{sec:XIV.7}

Let $X$ be a real or complex Banach space, and $A$ the infinitesimal generator of an equi-continuous semi-group $\{T_t; t \geqq 0\}$ of linear operators $\in L(X, X)$ of class $(C_0)$. Then the Lemma in Chapter~\ref{chap:IX}, \ref{sec:IX.12} gives the proof of $E$. Hille's approximation of $T_t$: 

\begin{equation}
\label{eq:XIV.7.1}
\begin{minipage}{0.92\linewidth}
\everymath{}\let\limits\nolimits
For every $x_0 \in X,  T_t x_0 = \slim\limits_{n \to \infty} (I - n^{-1} t A)^{-n} \cdot x_0,$ and the convergence is uniform on every compact interval of $t.$
\end{minipage}
\end{equation}

The titled convergence theorem is suggested by \eqref{eq:XIV.7.1}. It reads as follows. 

\begin{theorem}[M. \textsc{Crandall}-T. \textsc{Liggett} {\cite{clarkson37}}]\label{thm:XIV.7.Theorem.1} Let $A$ be a hyperdissipative set $\subseteq X \times X$. Then, for all $x_0 \in D(A)$, 

\begin{equation}
\tag{1'}\label{eq:XIV.7.1'}
\begin{minipage}{0.92\linewidth}
\everymath{}\let\limits\nolimits
$\slim\limits_{n\to \infty}(J_{t/n})^n \cdot x_0$ exists uniformly on every compact interval of $t,$ where $J_{\lambda} = (I - \lambda A)^{-1}$ as in the preceding Section~\hyperref[sec:XIV.6]{6}.
\end{minipage}
\end{equation}
\end{theorem}

The following proof is adapted from S. \textsc{Rasmussen} \cite{rasmussen269} which seems to give a good modification of the original proof by \textsc{Crandall}-\textsc{Liggett}: 

\begin{proof}
 The first step. We shall prepare the following ($2$)-($4$): 

\begin{equation}\label{eq:XIV.7.2}
\|A_{\lambda}x\| \leqq \|Ax\| \text{ for all } x \in D(A), \text{ and for all } \lambda > 0,
\end{equation}

\begin{equation}\label{eq:XIV.7.3}
\left\| \prod\limits_{i=1} ^{n} J _ {\lambda_{i}} x - x \right\| = \| J _ {\lambda_{n}} J _ {\lambda_{n-1}} \dots J _ {\lambda_{1}} x - x \| \leqq \left(\sum\limits_{i=1} ^{n} \lambda_{i}\right) \| A x \|
\end{equation}
for all $x \in D(A)$ and for all $\lambda_i > 0$ ( $i = 1, 2, \ldots, \overline{n}$). 

\begin{equation}\label{eq:XIV.7.4}
J_{\mu}\left(\frac{\mu}{\lambda} x + \frac{\lambda - \mu}{\lambda} J_{\lambda}x\right) = J_{\lambda}x \text{ for all } x \in X \text{ and for all } \lambda, \mu > 0.
\end{equation}


($2$) was proved in the preceding Section~\hyperref[sec:XIV.6]{$6$}. The proof of ($3$) is obtained by the contraction property of $J_{\lambda}$, ($2$) and $\lambda A_{\lambda} = J_{\lambda} - I$ as follows. 
\[
\begin{aligned}
\left\| \prod\limits_{i=1} ^{n} J _ {\lambda_{i}} x - x \right\|& = \left\| \sum\limits_{i=1} ^{n} \left(\prod\limits_{j=i} ^{n} J _ {\lambda_{j}} x - \prod\limits_{j=i+1} ^{n} J _ {\lambda_{j}} x\right) \right\| \leqq \sum\limits_{i=1} ^{n} \left\| \prod\limits_{j=i} ^{n} J _ {\lambda_{j}} x - \prod\limits_{j=i+1} ^{n} J _ {\lambda_{j}} x \right\|
\\
&\leqq \sum\limits_{i=1} ^{n} \| J _ {\lambda_{i}} x - x \| \leqq \sum\limits_{i=1} ^{n} \| \lambda_{i} A _ {\lambda_{i}} x \| \leqq \left(\sum\limits_{i=1} ^{n} \lambda_{i}\right) \| A x \|.
\end{aligned}
\]

We next prove ($4$). Since $A$ is a hyperdissipative set, there exists, for any $x \in X$ and $\lambda > 0$, a point $\{x_1, y_1\} \in A$ such that $x = x_1 - \lambda y_1$. Thus $J_\lambda x = x_1$ and so 
\[
\frac {\mu}{\lambda} x + \frac {\lambda - \mu}{\lambda} J _{\lambda} x = \frac {\mu}{\lambda} (x _{1} - \lambda y _{1}) + \frac {\lambda - \mu}{\lambda} x _{1} = x _{1} - \mu y _{1},
\]
proving that $J_{\mu}(x_1 - \mu y_1) = x_1 = J_{\lambda}x$. 

The second step. Let $x \in D(A)$ and $\lambda > 0$. Let furthermore $\{\mu_n; n \geqq 1\}$ be a sequence such that $0 < \mu_n \leqq \lambda$ for all $n$. We then define 
\begin{align}
A _{n,m}& = \left\| \prod\limits_{i=j} ^{n} J _ {\mu_{i}} x - J _{\lambda} ^{m} x \right\| \text{ for } n, m = 0, 1, 2, \dots \label{eq:XIV.7.5} \\
&\quad \left(\text{ We set } \prod\limits_{i=1} ^{0} J _ {\mu_{i}} = I = J _{\lambda} ^{0} \text{ for convenience. }\right) \notag
\end{align}

Moreover, let 
\begin{equation}\label{eq:XIV.7.6}
t _{n} = \sum\limits_{i=1} ^{n} \mu_{i} (n = 0, 1, 2, \dots ; t _{0} = 0)
\end{equation}
and 
\begin{equation}\label{eq:XIV.7.7}
\alpha_{n} = \mu_{n} / \lambda \quad \text{ and } \quad \beta_{n} = 1 - \alpha_{n} = (\lambda - \mu_{n}) / \lambda \quad \text{ for } \quad n = 1, 2, \dots
\end{equation}

Then we have 
\begin{equation}\label{eq:XIV.7.8}
A _{n,m} \leqq \left\{\left[ (m \lambda - t _{n}) ^{2} + m \lambda^{2} \right] ^{1/2} + \left[ (m \lambda - t _{n}) ^{2} + \lambda t _{n} \right] ^{1/2} \right\} \| A x \|.
\end{equation}

For the proof, we first show 
\begin{equation}\label{eq:XIV.7.9}
A _{n,m} \leqq \alpha_{n} A _{n-1,m-1} + \beta_{n} A _{n-1,m} \quad \text{ for } \quad n, m = 1, 2, \dots .
\end{equation}

In fact we have, by ($4$) and the contraction property of $J_{\mu_{n}}$, 
\[
\begin{aligned}
&A _{n,m} = \left\| \prod\limits_{i=1} ^{n} J _ {\mu_{i}} x - J _{\lambda} ^{m} x \right\| = \left\| \prod\limits_{i=1} ^{n} J _ {\mu_{i}} x - J _ {\mu_{n}} \left(\frac {\mu_{n}}{\lambda} J _{\lambda} ^{m-1} x + \frac {\lambda - \mu_{n}}{\lambda} J _{\lambda} ^{m} x\right) \right\|
\\
&\leqq \left\| \prod\limits_{i=1} ^{n-1} J _ {\mu_{i}} x - \left(\frac {\mu_{n}}{\lambda} J _{\lambda} ^{m-1} x + \frac {\lambda - \mu_{n}}{\lambda} J _{\lambda} ^{m} x\right) \right\| \leqq \alpha_{n} \left\| \prod\limits_{i=1} ^{n-1} J _ {\mu_{i}} x - J _{\lambda} ^{m-1} x \right\| + \beta_{n} \left\| \prod\limits_{i=1} ^{n-1} J _ {\mu_{i}} x - J _{\lambda} ^{m} x \right\|
\\
&= \alpha_{n} A _{n -1 \cdot m -1} + \beta_{n} A _{n-1,m}.
\end{aligned}
\]

We shall prove \eqref{eq:XIV.7.8} by \eqref{eq:XIV.7.9} and induction. We start with the proof that $A_{0,m}$ satisfies \eqref{eq:XIV.7.8} for all $m = 0, 1, \ldots$. In fact, we have 
\[
\begin{aligned} A _{0,m} & = \| x - J _{\lambda} ^{m} x \| \leqq m \lambda \| A x \| \quad \text{ (by \eqref{eq:XIV.7.3} and \eqref{eq:XIV.7.5}) } \\ & \leqq \left\{\left[ (m \lambda - t _{0}) ^{2} + m \lambda^{2} \right] ^{1/2} + \left[ \left(m \lambda - t _{0}\right) ^{2} + \lambda t _{0} \right] ^{1/2} \right\} \| A x \| (\text{since $t_{0}$ = 0}) \end{aligned}
\]

Let now $n$ and $m$ be arbitrary and assume that $A_{n,m}$ and $A_{n,m-1}$ both satisfy \eqref{eq:XIV.7.8}. We shall then show that $A_{n+1,m}$ satisfies \eqref{eq:XIV.7.8}. In fact, we have 
\[
\begin{aligned}
A_{n+1,m}&\leqq\alpha_{n+1}A_{n,m-1}+\beta_{n+1}A_{n,m}
  \quad\text{(by \eqref{eq:XIV.7.9})}\\
&\leqq\alpha_{n+1}\bigl\{[((m-1)\lambda-t_n)^2+(m-1)\lambda^2]^{1/2}\\
&\qquad+[((m-1)\lambda-t_n)^2+\lambda t_n]^{1/2}\bigr\}\|Ax\|\\
&\quad+\beta_{n+1}\bigl\{[(m\lambda-t_n)^2+m\lambda^2]^{1/2}
  +[(m\lambda-t_n)^2+\lambda t_n]^{1/2}\bigr\}\|Ax\|\\
&\leqq(\alpha_{n+1}+\beta_{n+1})^{1/2}
  \bigl\{\alpha_{n+1}[((m-1)\lambda-t_n)^2+(m-1)\lambda^2]\\
&\qquad+\beta_{n+1}[(m\lambda-t_n)^2+m\lambda^2]\bigr\}^{1/2}\|Ax\|\\
&\quad+(\alpha_{n+1}+\beta_{n+1})^{1/2}
  \bigl\{\alpha_{n+1}[((m-1)\lambda-t_n)^2+\lambda t_n]\\
&\qquad+\beta_{n+1}[(m\lambda-t_n)^2+\lambda t_n]\bigr\}^{1/2}\|Ax\|
  \quad\text{(by Schwarz' inequality)}\\
&=\bigl\{m^2\lambda^2-2\alpha_{n+1}m\lambda^2+\alpha_{n+1}\lambda^2
  -2m\lambda t_n+2\alpha_{n+1}\lambda t_n+t_n^2\\
&\qquad+m\lambda^2-\alpha_{n+1}\lambda^2\bigr\}^{1/2}\|Ax\|\\
&\quad+\bigl\{m^2\lambda^2-2\alpha_{n+1}m\lambda^2+\alpha_{n+1}\lambda^2
  -2m\lambda t_n+2\alpha_{n+1}\lambda t_n+t_n^2+\lambda t_n\bigr\}^{1/2}\|Ax\|\\
&\hspace{7em}\text{(by $\beta_{n-1}=1-\alpha_{n-1}$)}\\
&=\bigl\{m^2\lambda^2-2\mu_{n+1}m\lambda+\mu_{n+1}\lambda
  -2m\lambda t_n+2\mu_{n+1}t_n+t_n^2+m\lambda^2-\mu_{n+1}\lambda\bigr\}^{1/2}\|Ax\|\\
&\quad+\bigl\{m^2\lambda^2-2\mu_{n+1}m\lambda+\mu_{n+1}\lambda
  -2m\lambda t_n+2\mu_{n+1}t_n+t_n^2+\lambda t_n\bigr\}^{1/2}\|Ax\|\\
&=\bigl\{(m\lambda-(\mu_{n+1}+t_n))^2-\mu_{n+1}^2+m\lambda^2\bigr\}^{1/2}\|Ax\|\\
&\quad+\bigl\{(m\lambda-(\mu_{n+1}+t_n))^2-\mu_{n+1}^2
  +\lambda(t_n+\mu_{n+1})\bigr\}^{1/2}\|Ax\|\\
&\leqq\bigl\{[(m\lambda-t_{n+1})^2+m\lambda^2]^{1/2}
  +[(m\lambda-t_{n+1})^2+\lambda t_{n+1}]^{1/2}\bigr\}\|Ax\|,
\end{aligned}
\]
and we have thus proved that $A_{n+1,m}$ satisfies \eqref{eq:XIV.7.8}. We can also prove that $A_{n,0}$ satisfies \eqref{eq:XIV.7.8} for all $n = 1, 2, \ldots$. In fact, we obtain by \eqref{eq:XIV.7.3} 
\[
\begin{aligned}
A _{n,0}& = \left\| \prod\limits_{i=1} ^{n} J _ {u _{i}} x - x \right\| \leqq \sum\limits_{i=1} ^{n} \mu_{i} \cdot \| A x \| = t _{n} \| A x \|
\\
\quad& \leqq \{[ (0 \cdot \lambda - t _{n}) ^{2} + 0 \cdot \lambda^{2} ] ^{1/2} + [ (0 \cdot \lambda - t _{n}) ^{2} + \lambda t _{n} ] ^{1/2} \} \| A x \|.
\end{aligned}
\]

Therefore, we have completed the induction and hence \eqref{eq:XIV.7.8} is proved. 

The third step. Let $t$ be a positive number, and consider a partition $\Delta$ of the closed interval $[0, t]$. 
\[
\Delta : 0 = t _{0} < t _{1} < \dots < t _{i-1} < t _{i} < t _{i+1} < \dots < t _{n} = t.
\]

For this partition, we define 
\[
\Delta_{i} = \left(t _{i} - t _{i-1}\right) \quad (i = 1, 2, \dots , n) \quad \text{ and } \quad | \Delta | = \max\limits _{1 \leqq i \leqq n} \left(t _{i} - t _{i-1}\right).
\]

Take another partition 
\[
\Delta^{\prime}: 0 = t _{0} ^{\prime} < t _{1} ^{\prime} < \dots < t _{j-1} ^{\prime} < t _{j} ^{\prime} < t _{j+1} ^{\prime} < \dots < t _{k} ^{\prime} = t
\]
and define 
\[
\Delta_{j} ^{\prime} = \left(t _{j} ^{\prime} - t _{j-1} ^{\prime}\right) \quad (j = 1, 2, \dots , k) \quad \text{ and } \quad | \Delta^{\prime} | = \max\limits _{1 \leqq j \leqq k} \left(t _{j} ^{\prime} - t _{j-1} ^{\prime}\right)
\]

Then, for $x_0 \in D(A)$, we obtain, by taking 
$\lambda = \max (|\Delta |, |\Delta '|)$ and $m =$ the largest integer $\leqq t / \lambda$, 
the following inequality implied from \eqref{eq:XIV.7.8}: 
\begin{equation}\label{eq:XIV.7.10}
\begin{yoslines}& \left\| \prod\limits_{i=1} ^{n} J _ {\Delta_{i}} x _{0} - \prod\limits_{j=1} ^{k} J _ {\Delta_ {j ^{\prime}}} x _{0} \right\| \leqq \left\| \prod\limits_{i=1} ^{n} J _ {\Delta_{i}} x _{0} - J _{\lambda} ^{m} x _{0} \right\| + \left\| J _{\lambda} ^{m} x _{0} - \prod\limits_{j=1} ^{k} J _ {\Delta_ {j ^{\prime}}} x _{0} \right\|\\& \leqq 2 \left\{\left[ (m \lambda - t) ^{2} + m \lambda^{2} \right] ^{1/2} + \left[ (m \lambda - t) ^{2} + \lambda t \right] ^{1/2} \right\} \| A x _{0} \|\\& \leqq 2 \left\{\left[ \lambda^{2} + t \lambda \right] ^{1/2} + \left[ \lambda^{2} + \lambda t \right] ^{1/2} \right\} \| A x \|.\end{yoslines}
\end{equation}

Therefore, by taking $\Delta_{i} = t / n$ ( $i = 1,2,\ldots,n$) and $\Delta j = t / k$ ( $j = 1,2,\ldots,k$) in \eqref{eq:XIV.7.10}, we obtain \eqref{eq:XIV.7.1'}. 
\end{proof}

\begin{remarku} We have incidentally proved that 
\begin{equation}
\tag{1''}\label{eq:XIV.7.1''}
\begin{minipage}{0.92\linewidth}
\everymath{}\let\limits\nolimits
$T _{t} x _{0} = \slim _{n \rightarrow \infty} (I - n ^{-1} t A) ^{-n} x _{0} = \slim (I - n ^{-1} A) ^{-[nt]} x _{0}$ uniformly on every compact interval of $t,$ where $[ n t ] =$ the largest integer $\leqq n t.$
\end{minipage}
\end{equation}
\end{remarku}

The scope of this Theorem will be seen from Theorem~\ref{thm:XIV.7.Theorem.2}, Theorem~\ref{thm:XIV.7.Theorem.3} and Theorem~\ref{thm:XIV.7.Theorem.4} given below. 

\begin{theorem}\label{thm:XIV.7.Theorem.2} Let $A$ be a hyperdissipative closed set of $X \times X$. Then, for every $x_{0} \in D(A)$, the following two conditions are mutually equivalent.
\begin{enumerate}[i)]
\item $u(t)$ is a strong solution of the initial value problem 
\begin{equation}\label{eq:XIV.7.11}
\frac {\dd u (t)}{\dd t} \in A u (t) \quad \text{ for almost every } \quad t \geqq 0 \quad \text{ and } \quad u (0) = x _{0},
\end{equation}
as formulated in the preceding Section~\hyperref[sec:XIV.6]{$6$}. 

\item 
\begin{equation}
\label{eq:XIV.7.12}
\begin{minipage}{0.92\linewidth}
\everymath{}\let\limits\nolimits
$u (t) = \slim _{n \rightarrow \infty} (I - n ^{-1} A) ^{-[nt]} x _{0}$ and $u (t)$ is strongly differentiable at almost every $t \geqq 0.$
\end{minipage}
\end{equation}
\end{enumerate}
\end{theorem}

We first give the \textsc{Brezis}-\textsc{Pazy} \cite{brezis29} proof of the assertion i) $\to$ ii), and the proof of the assertion ii) $\to$ i) due to \textsc{Crandall}-\textsc{Liggett} \cite{clarkson37} will be given later; after the proof of Theorem~\ref{thm:XIV.7.Theorem.3}. We shall begin with a key Lemma due to T. \textsc{Kato} \cite{kato144}. 

\begin{lemmau} Let $S$ be the set of those $s \geqq 0$ at which the strong solution $u(s)$ of \eqref{eq:XIV.7.11} is strongly differentiable. Then, we have, at almost every $s \in S$, 
\begin{equation}\label{eq:XIV.7.13}
2 ^{-1} \frac {d}{\dd s} \| u (s) \| ^{2} = \| u (s) \| \frac {\dd \| u (s) \|}{\dd s} = \operatorname{Re} \left\langle \frac {\dd u (s)}{\dd s}, f \right\rangle \text{ whenever } f \in F (u (s)).
\end{equation}
\end{lemmau}

\begin{proof}
 $u(s)$ being a strong solution of \eqref{eq:XIV.7.11}, the strong derivative $\dd u(s)/\dd s$ is Bochner integrable on every compact interval of $s$ so that 
\[
u (s) - u (0) = \int_{0} ^{s} \frac {\dd u (s)}{\dd s} \dd s
\]
by \textsc{Bochner}'s theorem in Chapter~\ref{chap:V}, \ref{sec:V.5}. Hence $\|u(s)\|$ is of bounded variation on every compact interval of $s$. Therefore, $\|u(s)\|$ is differentiable at almost every $s \geqq 0$. 

Since $\operatorname{Re}\langle u(t), f \rangle \leqq \|u(t)\| \cdot \|u(s)\|$ and $\operatorname{Re}\langle u(s), f \rangle = \|u(s)\|^2$ by $f \in F(u(s))$, we obtain 
\[
\operatorname{Re} \left\langle u (t) - u (s), f \right\rangle \leqq \| u (s) \| \left(\| u (t) \| - \| u (s) \|\right).
\]

Dividing the both side by $(t - s)$ and letting $t \rightarrow s$ from above and from below, we obtain at almost every $s \in S$, 
\[
\operatorname{Re} \left\langle \frac {\dd u (s)}{\dd s}, f \right\rangle \leqq \| u (s) \| \frac {\dd \| u (s) \|}{\dd s} \text{ and } \operatorname{Re} \left\langle \frac {\dd u (s)}{\dd s}, f \right\rangle \geqq \| u (s) \| \frac {\dd \| u (s) \|}{\dd s}.
\]

We have thus proved \eqref{eq:XIV.7.13}. 
\end{proof}

\begin{corollaryn}[H. \textsc{Brezis}-A. \textsc{Pazy} \cite{brezis28}] For the strong solution $u(t)$ of \eqref{eq:XIV.7.11}, we have
\begin{equation}\label{eq:XIV.7.14}
\left\| \frac {\dd u (t)}{\dd t} \right\| \leqq \| A x _{0} \| = \inf\limits _ {x \in A x _{0}} \| z \| \quad \text{ at almost every } t \geqq 0.
\end{equation}

\end{corollaryn}
\begin{proof}
 We put $\dd u(t) / \dd t = y(t)$ so that, by \eqref{eq:XIV.7.11}, $y(t) \in Au(t)$ at almost every $t \geqq 0$. For such $t$ and for every $z \in Ax_0$, there exists, by the dissipative property of $A$, an $f_0 \in F(u(t) - x_0)$ such that $\operatorname{Re}\langle y(t) - z, f_0 \rangle \leqq 0$. Hence, by the same argument which enabled us to prove \eqref{eq:XIV.7.13}, we obtain 
\[
2 ^{-1} \frac {d}{\dd t} \| u (t) - x _{0} \| ^{2} = \| u (t) - x _{0} \| \frac {d}{\dd t} \| u (t) - x _{0} \| = \operatorname{Re} \left\langle \frac {\dd u (t)}{\dd t} - 0, f _{0} \right\rangle
\]
$= \operatorname{Re}\langle y(t) - z,f_0\rangle +\operatorname{Re}\langle z,f_0\rangle \leqq \operatorname{Re}\langle z,f_0\rangle$ at almost every $t\geqq 0$ 

Thus, at almost every $t \geqq 0$, we have 
\[
\left\| u (t) - x _{0} \right\| \frac {d}{\dd t} \left\| u (t) - x _{0} \right\| \leqq | \langle z, f _{0} \rangle | \leqq \| z \| \cdot \| u (t) - x _{0} \|
\]
and so, since $z \in Ax_0$ was arbitrary, 
\begin{equation}\label{eq:XIV.7.15}
\left\| u (t) - x _{0} \right\| \leqq t \left\| A x _{0} \right\| \text{ at every } t \geqq 0.
\end{equation}

Next, let $h > 0$ and put $u(t + h) = v(t)$. Then 
$\frac{\dd v(t)}{\dd t} \in A v(t)$ for almost every $t \geqq 0$ and $v(0) = u(h)$. \yosparagraphtag{11'}\label{eq:XIV.7.11'} 

Thus, similarly as above, we obtain, for a certain $f \in F(v(t) - u(t))$, 
\[
2 ^{-1} \frac {d}{\dd t} \| v (t) - u (t) \| ^{2} = \operatorname{Re} \left\langle \frac {\dd v (t)}{\dd t} - \frac {\dd u (t)}{\dd t}, f \right\rangle \leqq 0 \text{ at almost every } t \geqq 0.
\]

Therefore, we have 
\[
\left\| u (t + h) - u (t) \right\| \leqq \left\| u (s + h) - u (s) \right\| \text{ whenever } t \geqq s \geqq 0,
\]
and so, combined with \eqref{eq:XIV.7.15}, we finally obtain 
\begin{equation}\tag{15'}\label{eq:XIV.7.15'}
\left\| u (t + h) - u (t) \right\| \leqq \left\| u (h) - u (0) \right\| \leqq h \left\| A x _{0} \right\|.
\end{equation}

This proves \eqref{eq:XIV.7.14}. 
\end{proof}

\begin{proof}[Proof of the assertion i) $\to$ ii) in Theorem~\ref{thm:XIV.7.Theorem.2}.]
The existence of the right hand limit of \eqref{eq:XIV.7.12} is guaranteed by Theorem~\ref{thm:XIV.7.Theorem.1}. Let $T$ be an arbitrary positive constant. Besides step functions $u_{n}(t) = (I - n^{-1}A)^{-[nt]}x_{0}$, we define on [$0$, $T$] another sequence of functions $v_{n}(t)$ by
\begin{equation}\label{eq:XIV.7.16}
\begin{aligned}
&v _{n} (t)  = u _{n} (j / n) + (t - j / n) \cdot n \cdot [ u _{n} ((j + 1) / n) - u _{n} (j / n) ] \\
&\quad (\text{ for } \quad j / n \leqq t \leqq (j + 1) / n; j = 0, 1, \dots , [ n T ] - 1) = u _{n} (t) \quad (\text{ for } \quad n ^{-1} [ n T ] \leqq t \leqq T).
\end{aligned}
\end{equation}

Clearly, $v_{n}(t)$ is strongly differentiable on $[0, T]$ except for finite number of points of the form $t = j / n$, and for $j / n < t < (j + 1) / n$ we have 
\begin{equation}\label{eq:XIV.7.17}
\begin{yoslines}
&\frac {\dd v _{n} (t)}{\dd t}  = n [ u _{n} ((j + 1) / n) - u _{n} (j / n) ] = n [ (I - n ^{-1} A) ^{-1} u _{n} (j / n) - u _{n} (j / n) ] \\
&= A _{1/n} u _{n} (j / n) \quad (\text{ by } \quad A _{\lambda} = \lambda^{-1} (J _{\lambda} - I)) = A _{1/n} u _{n} (t) = A _{1/n} v _{n} (j / n).
\end{yoslines}
\end{equation}

By \eqref{eq:XIV.6.7}, \eqref{eq:XIV.6.8} and \eqref{eq:XIV.6.9} of the preceding Section and by 
\begin{align}
A u _{n} (j / n)& = A J _{1/n} u _{n} ((j - 1 / n) \ni A _{1/n} u _{n} ((j - 1) / n), \text{ we obtain } \notag \\
\left\| A _{1/n} u _{n} (j / n) \right\|& \leqq \left\| A u _{n} (j / n) \right\| \leqq \left\| A _{1/n} u _{n} ((j - 1) / n) \right\| \leqq \left\| A u _{n} ((j - 1) / n) \right\|. \label{eq:XIV.7.18}
\end{align}

Thus, by \eqref{eq:XIV.7.17}, we have 
\begin{equation}\label{eq:XIV.7.19}
\left\| \frac {\dd v _{n} (t)}{\dd t} \right\| = \| A _{1/n} u _{n} (j / n) \| \leqq \| A u _{n} (0) \| = \| A x _{0} \| \text{ for a.e. } t \geqq 0.
\end{equation}

We have incidentally proved from the proof of \eqref{eq:XIV.7.17}-\eqref{eq:XIV.7.19} that, whenever $j / n \leqq t \leqq (j + 1) / n$, 
\[
\begin{aligned} & {\| v _{n} (t) - u _{n} (t) \| = (t - j / n) \cdot n \cdot \| u _{n} ((j + 1) / n) - u _{n} (j / n) \|} \\ & {\qquad \leqq n ^{-1} \| A _{1/n} u _{n} (j / n) \| \leqq n ^{-1} \| A x _{0} \|,} \end{aligned}
\]
and hence 
\begin{equation}\label{eq:XIV.7.20}
\left\| v _{n} (t) - u _{n} (t) \right\| \leqq n ^{-1} \left\| A x _{0} \right\| \quad \text{ for all } \quad t \geqq 0.
\end{equation}

We put $\dd u(t) / \dd t = y(t)$ and $n[u_n(t) - u_n(t - 1/n)] = y_n(t)$. Then, by \eqref{eq:XIV.7.11}, we have $y(t) \in Au(t)$ for almost every $t \geqq 0$. We have also $y_n(t) \in Au_n(t)$ when $t \neq j/n$ ( $j = 0, 1, \ldots$), because 
\[
y _{n} (t) = A _{1/n} u _{n} (t - 1 / n) \in A (I - n ^{-1} A) ^{-1} u _{n} (t - 1 / n) = A u _{n} (t)
\]
by \eqref{eq:XIV.6.8} of the preceding Section~\hyperref[sec:XIV.6]{$6$}. A being a dissipative set, we obtain, by \eqref{eq:XIV.6.5} of the preceding Section~\hyperref[sec:XIV.6]{$6$}, that for almost every $t \geqq 0$ the inequality 
\[
\left\| [ n \cdot u _{n} (t) - y _{n} (t) ] - [ n \cdot u (t) - y (t) ] \right\| \geqq \left\| n \cdot u _{n} (t) - n \cdot u (t) \right\|.
\]

Therefore we have, for almost every $t \geqq 0$, 
\begin{equation}\label{eq:XIV.7.21}
\begin{aligned}
&\left\|\frac{\dd u(t)}{\dd t}-n[u(t)-u(t-1/n)]\right\|\\
&=\bigl\|n[u_n(t)-u_n(t-1/n)]-n[u(t)-u(t-1/n)]+y(t)-y_n(t)\bigr\|\\
&\geqq\bigl\|[n u_n(t)-y_n(t)]-[n u(t)-y(t)]\bigr\|
  -\bigl\|n u_n(t-1/n)-n u(t-1/n)\bigr\|\\
&\geqq n\|u_n(t)-u(t)\|-n\|u_n(t-1/n)-u(t-1/n)\|.
\end{aligned}
\end{equation}

We here extend $u_{n}(t)$ and $u(t)$ for negative value of $t$ by 
\begin{equation}\label{eq:XIV.7.22}
\begin{yoslines}
&u _{n} (t) = x _{0} - n ^{-1} z _{0}, \text{ where } z _{0} \text{ is any fixed point } A x _{0}, u (t) = x _{0}.
\end{yoslines}
\end{equation}

Then, integrating the resulting inequality \eqref{eq:XIV.7.21} on $[0, \theta]$ where $n^{-1} \leqq \theta \leqq T$, we obtain 
\[
\begin{aligned}
-& \int_{-1/n} ^{0} n \| u _{n} (t) - u (t) \| \dd t + \int_{\theta -1 / n} ^{\theta} n \| u _{n} (t) - u (t) \| \dd t
\\
&\leqq \int_{0} ^{\theta} \left\| \frac {\dd u (t)}{\dd t} - n [ u (t) - u (t - 1 / n) ] \right\| \dd t,
\end{aligned}
\]
that is, 
\begin{equation*}
\begin{aligned}
&n \int_{\theta -1 / n} ^{\theta} \| u _{n} (t) - u (t) \| \dd t  \leqq \int_{0} ^{\theta} \left\| \frac {\dd u (t)}{\dd t} - n [ u (t) - u (t - 1 / n) ] \right\| \dd t + n \int_{-1/n} ^{0} \| u _{n} (t) - x _{0} \| \dd t \\
&\leqq \int_{0} ^{\theta} \left\| \frac {\dd u (t)}{\dd t} - n [ u (t) - u (t - 1 / n) ] \right\| \dd t + n ^{-1} \| z _{0} \| . \quad \text{ (by \eqref{eq:XIV.7.22}) }
\end{aligned}
\end{equation*}

Adding these inequalities for $\theta = 1 / n, 2 / n, \ldots, N / n$ with $N = [nT]$ yields 
\[
\begin{aligned}
&n \int_{0} ^{N/n} \| u _{n} (t) - u (t) \| \dd t  \leqq N \int_{0} ^{T} \left\| \frac {\dd u (t)}{\dd t} - n [ u (t) - u (t - 1 / n) ] \right\| \dd t + n ^{-1} N \| z _{0} \|
\end{aligned}
\]
and therefore 
\begin{equation}\label{eq:XIV.7.23}
\begin{aligned}
&\int_{0} ^{N/n} \| u _{n} (t) - u (t) \| \dd t  \leqq T \int_{0} ^{T} \left\| \frac {\dd u (t)}{\dd t} - n [ u (t) - u (t - 1 / n) ] \right\| \dd t + n ^{-1} T \cdot \| z _{0} \| .
\end{aligned}
\end{equation}

Since, by the assumption, $\slim\limits_{n\to \infty}n[u(t) - u(t - 1 / n)] = \dd u(t) / \dd t$ for almost every $t$ on $[0,T]$, we obtain $\lim\limits_{n\to \infty}\int_0^T \|u_n(t) - u(t)\|\dd t = 0$ by \eqref{eq:XIV.7.23} 
and by 
\[
\left\| \frac {\dd u (t)}{\dd t} - n [ u (t) - u (t - 1 / n) ] \right\| \leqq 2 \cdot \| z _{0} \| ,
\]
which is implied by \eqref{eq:XIV.7.14} and $n[u(t) - u(t - 1/n)] = \int_{t-1/n}^{t}\frac{\dd u(s)}{\dd s} \quad \dd s$. Therefore, combined with \eqref{eq:XIV.7.20}, we obtain 
\begin{equation}\label{eq:XIV.7.24}
\lim\limits _{n \rightarrow \infty} \int_{0} ^{T} \| v _{n} (t) - u (t) \| \dd t = 0.
\end{equation}

On the other hand, we have, by \eqref{eq:XIV.7.14}, \eqref{eq:XIV.7.19} and by the same argument which enabled us to prove \eqref{eq:XIV.7.13}, 
\[
2 ^{-1} \frac {d}{\dd t} \| v _{n} (t) - u (t) \| ^{2} \leqq \| v _{n} (t) - u (t) \| \cdot \left\| \frac {\dd v _{n} (t)}{\dd t} - \frac {\dd u (t)}{\dd t} \right\|
\]
$\leqq \|v_n(t) - u(t)\| \cdot 2 \cdot \|z_0\|$ for almost every $t \in [0, T]$. 

Therefore, by $v_{n}(0) = u(0) = x_{0}$, we obtain 
\[
\left\| v _{n} (t) - u (t) \right\| ^{2} \leqq 4 \int_{0} ^{t} \left\| v _{n} (s) - u (s) \right\| \dd s \cdot \left\| z _{0} \right\| \quad \text{ for } \quad 0 \leqq t \leqq T.
\]

Hence, by \eqref{eq:XIV.7.24}, we obtain $\slim\limits_{n\to \infty}v_n(t) = u(t)$ uniformly on $[0,T]$. Thus, by \eqref{eq:XIV.7.20}, we have proved that $\slim\limits_{n\to \infty}u_n(t) = u(t)$ uniformly on $[0,T]$. 
\end{proof}

\begin{theorem}\label{thm:XIV.7.Theorem.3} Let $A$ be a hyperdissipative set $\subseteq X \times X$, and let $x_0 \in D(A)$ and $\lambda > 0$. As in the preceding Section~\hyperref[sec:XIV.6]{$6$}, the initial value problem 
\begin{equation}\label{eq:XIV.7.25}
\frac {\dd u ^{(\lambda)} (t)}{\dd t} = A _{\lambda} u ^{(\lambda)} (t) \quad \text{ for } \quad t \geqq 0 \quad \text{ and } \quad u ^{(\lambda)} (0) = x _{0}
\end{equation}
has a uniquely determined solution $u^{(\lambda)}(t)$, because $A_{\lambda}$ is a Lipschitzian mapping. We have then 
\begin{equation}\label{eq:XIV.7.26}
\begin{aligned}
\slim_{\lambda \downarrow 0} u^{(\lambda)}(t) &\ \text{ exists uniformly on every compact interval of } t \text{ and }
\\
\slim_{\lambda \downarrow 0} u^{(\lambda)}(t) &= \slim_{n \rightarrow \infty} (I - n^{-1} t A)^{-n} x _{0}.
\end{aligned}
\end{equation}
\end{theorem}

\begin{remarku} If the strong solution $u(t)$ of \eqref{eq:XIV.7.11} exists, then, as in the linear case, $u(t)$ can be given by Hille's type approximation $\slim\limits_{n\to \infty}(I - n^{-1}tA)^{-n}\cdot x_0$ as well as by Yosida's type approximation $\slim\limits_{\lambda \downarrow \infty}u^{(\lambda)}(t)$. 
\end{remarku}

\begin{proof}[Proof of Theorem~\ref{thm:XIV.7.Theorem.3}.]
Since $A_{1} = (J_{1} - I)$ is dissipative, we obtain, as in the case of \eqref{eq:XIV.7.15'},
\begin{equation}\label{eq:XIV.7.27}
\left\| u ^{(1)} (t) - u ^{(1)} (0) \right\| \leqq t \cdot \left\| A _{1} u ^{(1)} (0) \right\| = t \cdot \left\| A _{1} x _{0} \right\| = t \cdot \left\| J _{1} x _{0} - x _{0} \right\|.
\end{equation}

On the other hand, it is easy to verify that the solution $u^{(1)}(t)$ of \eqref{eq:XIV.7.25} with $\lambda = 1$ is given by 
\begin{equation}\label{eq:XIV.7.28}
u ^{(1)} (t) = e ^{-t} x _{0} + \int_{0} ^{t} e ^{s-t} J _{1} u ^{(1)} (s) \dd s.
\end{equation}

Hence 
\[
u ^{(1)} (t) - J _{1} ^{n} x _{0} = e ^{-t} \left(x _{0} - J _{1} ^{n} x _{0}\right) + \int_{0} ^{t} e ^{s-t} \left[ J _{1} u ^{(1)} (s) - J _{1} ^{n} x _{0} \right] \dd s
\]
and so, by the fact that $J_{1}$ is a contraction mapping, we obtain 
\begin{equation}\label{eq:XIV.7.29}
\left\| u ^{(1)} (t) - J _{1} ^{n} x _{0} \right\| \leqq n e ^{-t} \left\| \left(J _{1} - I\right) x _{0} \right\| + \int_{0} ^{t} e ^{s-t} \left\| u ^{(1)} (s) - J _{1} ^{n-1} x _{0} \right\| \dd s.
\end{equation}

From this we shall derive 
\begin{equation}\label{eq:XIV.7.30}
\left\| u ^{(1)} (n) - J _{1} ^{n} x _{0} \right\| \leqq \sqrt {n} \left\| (J _{1} - I) x _{0} \right\|.
\end{equation}

We may assume that $(J_{1}-I)x_{0}=A_{1}x_{0}\neq0$. For, if otherwise, the solution $u^{(1)}(t)$ of \eqref{eq:XIV.7.25} with $\lambda=1$ is given by $u^{(1)}(t)\equiv x_{0}$ and hence \eqref{eq:XIV.7.30} would be trivial. The assumption $(J_{1}-I)x_{0}\neq0$, combined with \eqref{eq:XIV.7.29}, gives 
\begin{equation}\varphi_{n} (t) \leqq n e ^{-t} + \int_{0} ^{t} e ^{s-t} \varphi_{n-1} (s) \dd s, \text{ where } \varphi_{n} (t) = \left\| u ^{(1)} (t) - J _{1} ^{n} x _{0} \right\| \cdot \left\| \left(J _{1} - I\right) x _{0} \right\| ^{-1}. \label{eq:XIV.7.31}\end{equation}

We shall prove, by induction with respect to $n$, that 
\begin{equation}\label{eq:XIV.7.32}
\varphi_{n} (t) \leqq \{(n - t) ^{2} + t \} ^{1/2}.
\end{equation}

This gives \eqref{eq:XIV.7.30} by putting $t = n$. Inequality \eqref{eq:XIV.7.32} is clear for $n = 0$. If \eqref{eq:XIV.7.32} 
is true for $(n - 1)$, then $\varphi_n(t) \leqq n e^{-t} + \int_0^t e^{s-t}\{ (n - 1 - s\}^2 + s\}^{1/2} \dd s$. Denoting 
the right hand side by $\psi_{n}(t)e^{-t}$, we have to show that 
\[
\psi_{n} (t) \leqq e ^{t} \left\{(n - t) ^{2} + t \right\} ^{1/2}.
\]

Since $\psi_{n}(0) = n$, it is sufficient to prove that 
\[
\psi_{n} ^{\prime} (t) = e ^{t} \left\{(n - 1 - t) ^{2} + t \right\} ^{1/2} \leqq \frac {d}{\dd t} e ^{t} \left\{(n - t) ^{2} + t \right\} ^{1/2}.
\]

The last inequality is easily checked and so we have proved \eqref{eq:XIV.7.30}. 

We are now prepared to prove \eqref{eq:XIV.7.26}. Put $v^{(\lambda)}(t) = u^{(\lambda)}(\lambda t)$. Then, by $A_{\lambda} = \lambda^{-1}(J_{\lambda} - I)$, we have 
\begin{equation}\tag{25'}\label{eq:XIV.7.25'}
\frac {\dd v ^{(\lambda)} (t)}{\dd t} = \left(J _{\lambda} - I\right) v ^{(\lambda)} \text{ for all } t \geqq 0 \text{ with } v ^{(\lambda)} (0) = x _{0}.
\end{equation}

Taking $J_{\lambda}$ for $J_{1}$ and $v^{(\lambda)}$ for $u^{(\lambda)}$ in \eqref{eq:XIV.7.30} and making use of \eqref{eq:XIV.6.9} of the preceding Section~\hyperref[sec:XIV.6]{$6$}, we obtain 
\begin{equation}\label{eq:XIV.7.33}
\begin{aligned}
&\| v ^{(\lambda)} (n) - J _{\lambda} ^{n} x _{0} \|  = \| u ^{(\lambda)} (n \lambda) - J _{\lambda} ^{n} x _{0} \| \leqq \sqrt {n} \| (J _{\lambda} - I) x _{0} \| \leqq \sqrt {n} \lambda \| A _{\lambda} x _{0} \| \leqq \sqrt {n} \lambda \| A x _{0} \|.
\end{aligned}
\end{equation}

Let $\lambda_{k} \downarrow 0$ and let $n_{k} = [t / \lambda_{k}]$ so that $n_{k} \lambda_{k} \uparrow t$. Hence, by \eqref{eq:XIV.7.1'} and \eqref{eq:XIV.7.10}, we obtain $\slim\limits_{k \to \infty} J_{\lambda_k}^{n_k} x_0 = \slim\limits_{n \to \infty} (I - n^{-1} tA)^{-n} x_0$. Therefore, combined with \eqref{eq:XIV.7.33}, we have 
\[
\left\| u ^ {\left(\lambda_{k}\right)} \left(n _{k} \lambda_{k}\right) - J _ {\lambda_{k}} ^ {n _{k}} x _{0} \right\| \leqq \sqrt {n _{k}} \lambda_{k} \left\| A x _{0} \right\| \leqq n _{k} ^{-1/2} t \left\| A x _{0} \right\|.
\]

Since $\{\lambda_k\}$ was arbitrary, we have thus proved \eqref{eq:XIV.7.26}. 
\end{proof}

\begin{remarku} The above proof is adapted from H. \textsc{Brezis}-A. \textsc{Pazy} \cite{brezis29}. It is to be noted here that inequality \eqref{eq:XIV.7.30} is due to I. \textsc{Miyadera}-S. \textsc{Oharu}\cite{miyadera204}. 
\end{remarku}

\begin{proof}[Proof of the assertion ii) $\rightarrow$ i) in Theorem~\ref{thm:XIV.7.Theorem.2}.]
Let $u(t) = s\text{- }\lim\limits_{n\to \infty}(I - n^{-1}tA)^{-n}x_0$ be strongly differentiable at $t_0 > 0$ with strong derivative $y = \dd u(t) / \dd t|_{t=t_0}$ so that
\begin{equation}\label{eq:XIV.7.34}
u \left(t _{0} + h\right) = u \left(t _{0}\right) + h y + o (h), \text{ where } s \underset {h \downarrow 0} {\lim} o (h) / h = 0.
\end{equation}

Since \eqref{eq:XIV.7.19} and \eqref{eq:XIV.7.20} implies that $t \to u(t)$ is a Lipschitzian mapping with Lipschitz constant $\|Ax_0\|$, we obtain i) if we can prove 
\begin{equation}\label{eq:XIV.7.35}
\left\{u \left(t _{0}\right), y \right\} \in A.
\end{equation}

The proof of \eqref{eq:XIV.7.35} reads as follows. Let $0 < \lambda < t_{0}$. Then, by the strong differentiability of $u(t)$ at $t = t_{0}$ and by the hyperdissipative property of $A$, there exists a uniquely determined point $\{x_{\lambda}, y_{\lambda}\} \in A$ such that 
\begin{equation}\label{eq:XIV.7.36}
x _{\lambda} - \lambda y _{\lambda} = u (t _{0} - \lambda) = u (t _{0}) - \lambda y + o (\lambda).
\end{equation}

We obtain \eqref{eq:XIV.7.35} if we can show 
\begin{equation}\label{eq:XIV.7.37}
\slim _{\lambda \downarrow 0} x _{\lambda} = u (t _{0}) \quad \text{ and } \quad \slim _{\lambda \downarrow 0} y _{\lambda} = y,
\end{equation}
because $A$ is a closed set in $X \times X$ by the assumption. For the proof of \eqref{eq:XIV.7.37}, we put 
\begin{equation}\label{eq:XIV.7.38}
\parbox{0.92\linewidth}{\everymath{}\let\limits\nolimits $\hat{x}_{\lambda} = \hat{x} - \lambda \hat{y}$, where $\{\hat{x}, \hat{y}\}$ is a point of $A$ to be determined later.}
\end{equation} 


Then, by \eqref{eq:XIV.7.13} and \eqref{eq:XIV.7.25}, we obtain a certain $f \in F(u^{(\lambda)}(t) - \hat{x}_{\lambda})$ such that 
\begin{align}
2 ^{-1} \frac {d}{\dd t} \| u ^{(\lambda)} (t) - \hat {x} _{\lambda} \| ^{2}& = \operatorname{Re} \left\langle A _{\lambda} u ^{(\lambda)} (t), f \right\rangle = \operatorname{Re} \left\langle A _{\lambda} u ^{(\lambda)} (t) - A _{\lambda} \hat {x} _{\lambda}, f \right\rangle + \notag \\
+ \operatorname{Re} \left\langle A _{\lambda} \hat {x} _{\lambda}, f \right\rangle& \leqq \operatorname{Re} \left\langle \hat {y}, f \right\rangle , \notag
\end{align}
because $A_{\lambda} \hat{x}_{\lambda} = \hat{y}$ and, moreover, $A_{\lambda}$ is dissipative by Proposition~\hyperref[prop:XIV.6.Proposition.1]{$1$} in the preceding Section~\hyperref[sec:XIV.6]{$6$}. Thus 
\begin{equation}\label{eq:XIV.7.39}
\left\| u ^{(\lambda)} \left(t _{0} + h\right) - \hat {x} _{\lambda} \right\| ^{2} - \left\| u ^{(\lambda)} \left(t _{0}\right) - \hat {x} _{\lambda} \right\| ^{2} \leqq 2 \int_{0} ^{h} \left\langle \hat {y}, u ^{(\lambda)} \left(t _{0} + \tau\right) - \hat {x} _{\lambda} \right\rangle_{s} \cdot \dd \tau ,
\end{equation}
where 
\begin{equation}\label{eq:XIV.7.40}
\left\langle w, z \right\rangle_{s} = \sup\limits _{k \in F (z)} \operatorname{Re} \left\langle w, k \right\rangle .
\end{equation}

It will be proved later that 
\begin{equation}
\label{eq:XIV.7.41}
\begin{minipage}{0.92\linewidth}
\everymath{}\let\limits\nolimits
$\langle w,z\rangle_s $ is upper semi-continuous, i.e., $\slim_{n\to\infty}w_n=w_\infty$ and $\slim_{n\to\infty}z_n =z_\infty$ implies that $\limsup\limits_{n\to\infty}\langle w_n,z_n\rangle_s\leqq\langle w_\infty,z_\infty\rangle_s.$
\end{minipage}
\end{equation} 

Hence, by $\slim\limits_{\lambda \downarrow \infty} u^{(\lambda)}(t) = u(t) = \slim\limits_{n \to \infty} (I - n^{-1}tA)^{-n} x_0$ proved in \eqref{eq:XIV.7.26} and by $\slim\limits_{n \to \infty} \hat{x}_\lambda = \hat{x}$, we obtain 
\begin{equation}\label{eq:XIV.7.42}
\left\| u \left(t _{0} + h\right) - \hat {x} \right\| ^{2} - \left\| u \left(t _{0}\right) - \hat {x} \right\| ^{2} \leqq 2 \int_{0} ^{h} \left\langle \hat {y}, u \left(t _{0} + \tau\right) - \hat {x} \right\rangle_{s} \cdot \dd \tau .
\end{equation}

For every $f \in F(u(t_0) - \hat{x})$, we have 
\begin{gather}\operatorname{Re} \left\langle u (t _{0} + h) - u (t _{0}), f \right\rangle - \operatorname{Re} \left\langle \hat {x} - u (t _{0}), f \right\rangle = {} \operatorname{Re} \left\langle u (t _{0} + h) - \hat {x}, f \right\rangle \notag \\\leqq {} \left\| u \left(t _{0} + h\right) - \hat {x} \right\| \cdot \left\| \hat {x} - u \left(t _{0}\right) \right\| \leqq {} 2 ^{-1} \left\{\left\| u \left(t _{0} + h\right) - \hat {x} \right\| ^{2} + \left\| \hat {x} - u \left(t _{0}\right) \right\| ^{2} \right\} \notag
\end{gather}

Therefore, by $\operatorname{Re}\langle \hat{x} - u(t_0), f \rangle = -\|\hat{x} - u(t_0)\|^2$ and \eqref{eq:XIV.7.42}, we obtain 
\[
2 \cdot \operatorname{Re} \left\langle \frac {u (t _{0} + h) - u (t _{0})}{h}, f \right\rangle \leqq 2 \int_{0} ^{1} \left\langle \hat {y}, u (t _{0} + h \tau) - \hat {x} \right\rangle_{s} \cdot \dd \tau .
\]

Hence, by the strong differentiability of $u(t)$ at $t = t_0$ and \eqref{eq:XIV.7.41}, we have 
\begin{equation}\tag{42'}\label{eq:XIV.7.42'}
\operatorname{Re} \langle y, f \rangle \leqq \langle \hat {y}, u (t _{0}) - \hat {x} \rangle_{s}.
\end{equation}

We may replace the supremum on the right hand side of \eqref{eq:XIV.7.42'} by the maximum, because the set $F(u(t_0) - \hat{x})$ contained in the strongly closed sphere of $X'$ is $w^*$-compact in $X'$ (see Theorem~\hyperref[thm:AV.1.Theorem.1]{$1$} in Appendix to Chapter~\ref{chap:V}, \ref{sec:AV.1}). Thus there exists an $\hat{f} \in F(u(t_0) - \hat{x})$ such that 
\[
\operatorname{Re} \langle y, f _{1} \rangle \leqq \operatorname{Re} \langle \hat {y}, \hat {f} \rangle \quad \text{ for all } \quad f _{1} \in F (u (t _{0}) - \hat {x}).
\]

By taking $\hat{x} = x_{\lambda},\hat{y} = y_{\lambda}$ and $f_{1} = f$, we obtain 
\begin{equation}\label{eq:XIV.7.43}
\operatorname{Re} \langle y, f \rangle \leqq \operatorname{Re} \langle y _{\lambda}, f \rangle , \text{ that is, } \operatorname{Re} \langle y _{\lambda} - y, f \rangle \geqq 0.
\end{equation}

On the other hand, we have, from \eqref{eq:XIV.7.36}, 
\begin{equation}\label{eq:XIV.7.44}
\frac {u \left(t _{o}\right) - x _{\lambda}}{\lambda} + \frac {o (\lambda)}{\lambda} = y - y _{\lambda},
\end{equation}
and hence, by \eqref{eq:XIV.7.43}, we obtain 
\[
\left\| u \left(t _{0}\right) - x _{\lambda} \right\| ^{2} = \operatorname{Re} \left\langle u \left(t _{0}\right) - x _{\lambda}, \hat {f} \right\rangle \leqq \operatorname{Re} \left\langle o (\lambda), \hat {f} \right\rangle \leqq o (\lambda) \cdot \left\| u \left(t _{0}\right) - x _{\lambda} \right\|.
\]

Hence $\|u(t_0) - x_{\lambda}\| \leqq o(\lambda)$ and so we obtain \eqref{eq:XIV.7.37} by \eqref{eq:XIV.7.44}. 

Finally, the proof of \eqref{eq:XIV.7.41} is given as follows. Since the set $\{f\in X';\|f\|\leqq \|z_m\|\}$ is $w^*$-compact (Theorem~\ref{thm:AV.1.Theorem.1} in Appendix to Chapter~\ref{chap:V}, \ref{sec:AV.1}), we can find $f_{m}\in X'$ such that 
\[
\left\langle w _{m}, z _{m} \right\rangle_{s} = \operatorname{Re} \left\langle w _{m}, f _{m} \right\rangle \quad \text{ and } \quad \| f _{m} \| \leqq \| z _{m} \|.
\]

We may assume, without loss of generality, that $\lim\limits_{m\to \infty}\langle w_m,z_m\rangle_s$ exists. 

Let $f_{\infty} \in X'$ be any one of the $w^*$-accumulation points of $\{f_m\}$. Then we have 
\[
\lim\limits _{m \rightarrow \infty} \left\langle w _{m}, z _{m} \right\rangle_{s} = \lim\limits _{m \rightarrow \infty} \operatorname{Re} \left\langle w _{m}, f _{m} \right\rangle = \operatorname{Re} \left\langle w _{\infty}, f _{\infty} \right\rangle , \| f _{\infty} \| \leqq \| z _{\infty} \|,
\]
because $s$ - $\lim\limits_{m \to \infty} w_m = w_\infty$ and $w^*$ - $\lim\limits_{m' \to \infty} f_{m'} = f_\infty$ for a suitable subsequence 
$\{m'\}$ of the sequence $\{m\}$. Therefore we must have \eqref{eq:XIV.7.41}. 
\end{proof}

\begin{remarku} It is an unfortunate and a fortiori an interesting situation that, although $\slim\limits_{n\to \infty}(I - n^{-1}tA)^{-n}x_0$ and $\slim\limits_{\lambda \downarrow 0}u^{(\lambda)}(t)$ both exist and coincide each other under a general condition, this limit need not be strongly differentiable in $t$ anywhere. Such an example is given in M. \textsc{Crandall}-T. \textsc{Liggett} \cite{clarkson37}. Recently, P. \textsc{Bénilan} \cite{benilan13} devised the following remedy for the above situation. Let $u(t)$ be a strong solution of initial value problem \eqref{eq:XIV.7.11}. Let, further, $\{z,w\}$ be any fixed point of the hyperdissipative subset $A$ of $X\times X$. Then, as in the case of \eqref{eq:XIV.7.42}, we obtain 
\[
\|u(t)-z\|^2\leqq\|u(s)-z\|^2+2\int_s^t\langle w,u(\tau)-z\rangle_s\cdot\dd\tau
\]
for all $0\leqq s\leqq t$ and with a fixed initial condition $u(0)=x_0\in D(A)$.
\end{remarku}

P. \textsc{Bénilan} \cite{benilan13} called a strongly continuous $X$-valued function $u(t)$ as an integral solution of initial value problem \eqref{eq:XIV.7.11}, if $u(t)$ satisfies the above condition for all $\{z, w\} \in A$ simultaneously, and he proved 

\begin{theorem}\label{thm:XIV.7.Theorem.4} Let $A$ be a hyperdissipative set $\subseteq X \times X$. Then $u(t) = s\text{- }\lim\limits_{\lambda \downarrow 0} u^{(\lambda)}(t) = s\text{- }\lim\limits_{n \to \infty} (I - n^{-1}tA)^{-n} \cdot x_0$ is the uniquely determined integral solution of \eqref{eq:XIV.7.11}. 
\end{theorem}

For the proof of this theorem and further interesting extensions as well as comments, we refer the reader to P. \textsc{Bénilan} \cite{benilan14}. 

\begin{remarku} The scope of applications of Theorems in the present and the preceding Sections may well be seen, e.g., by reading S. \textsc{Aizawa} \cite{aizawa4}, M. G. \textsc{Crandall} \cite{clarkson38}, Y. \textsc{Konishi} \cite{konishi160} and \cite{konishi161}, and B. K. \textsc{Quinn} \cite{quinn265}. 
\end{remarku}

\yoschapterstar{Supplementary Notes}\label{chap:SupplementaryNotes}

\yossub{Chapter I and Chapter VI}\phantomsection\label{notes:I-VI}

\begin{enumerate}[1.]
\item For the distributions or the generalized functions due to S. L. \textsc{Sobolev}, L. \textsc{Schwartz} and I. M. \textsc{Gelfand}, see the new and enlarged edition L. \textsc{Schwartz} \cite{schwartz299} and the English edition I. M. \textsc{Gelfand} \cite{gelfand80}--\cite{gelfand84}.
\item For the recent development of the hyperfunctions due to M. \textsc{Sato}, see M. \textsc{Sato} \cite{sato285}--\cite{sato286} and P. \textsc{Schapira} \cite{schapira287}.
\end{enumerate}

\yossub{Chapter VI}\phantomsection\label{notes:VI}

\begin{enumerate}[1.]
\item (Section~\hyperref[sec:VI.7]{$7$}) Sobolev Spaces. The theorem (Sobolev's Lemma) given on p.~\pageref{thm:VI.7.Sobolev} is a very special case of the so-called ``Sobolev Imbedding Theorems'' in S. L. \textsc{Sobolev} \cite{sobolev306}--\cite{sobolev307}. A comprehensive description of the development of this imbedding theorems is given in R. A. \textsc{Adams} \cite{adams2}. See also E. M. \textsc{Stein} \cite{stein310}.
\end{enumerate}

\yossub{Chapter X}\phantomsection\label{notes:X}

\begin{enumerate}[1.]
\item (Section~\hyperref[sec:X.1]{$1$}) The trace operator or the generalized boundary values. We can prove 

\begin{theoremu} Let $\Omega$ be a bounded open domain in $\mathbb{R}^{n}$ such that its boundary hypersurface $\partial\Omega$ is $C^2$. Then, for any $f(x) \in W^{1,2}(\Omega)$ and for almost all $\xi \in \partial\Omega$, the limiting value $\varphi(\xi) = \lim\limits_{x \to \xi} f(x)$ exists if we let $x$ tend to $\xi$ along the normal at the point $\xi \in \partial\Omega$. Moreover, this boundary value $\varphi(\xi)$ of $f(x)$ belongs to $L^2(\partial\Omega)$ and satisfies the inequality: 
\[
\left\| \varphi (\xi) \right\| _ {L^{2} (\partial \Omega)} \leqq C \left\{\left\| f (x) \right\| _ {L^{2} (\Omega)} + \sum\limits_{j=1} ^{n} \left\| \partial f / \partial x _{j} \right\| _ {L^{2} (\Omega)} \right\},
\]
where the positive constant $C$ is independent of the individual choice of $f$ and the notation $\partial f/\partial x_{j}$ means differentiation in the sense of the ``distribution''. 
\end{theoremu}

The linear map $f \to \varphi$ is called the trace operator. For the proof as well as generalizations of the above theorem, see, e.g., S. \textsc{Mizohata} \cite{mizohata213}, R. A. \textsc{Adams} \cite{adams2} and F. \textsc{Trèves} \cite{trotter326}. 

\item \phantomsection\label{notes:AX} (Appendix to Chapter~\ref{chap:X}) For extensions with a unified treatment of the result of R. A. \textsc{Minlos} \cite{minlos202}, see L. \textsc{Schwartz} \cite{schwartz300} and the bibliography cited there.
\end{enumerate}

\yossub{Chapter XI}\phantomsection\label{notes:XI}

\begin{enumerate}[1.]
\item For recent development of the theory of contraction operators in Hilbert spaces, see B. Sz. \textsc{Nagy}-C. \textsc{Foias} \cite{nagy220}.
\end{enumerate}

\yossub{Chapter XII}\phantomsection\label{notes:XII}

\begin{enumerate}[1.]
\item (Section~\hyperref[sec:XII.1]{$1$}) For Choquet's refinement of the Krein-Milman Theorem, see R. R. \textsc{Phelps} \cite{phelps250}.
\end{enumerate}

\yossub{Chapter XIII}\phantomsection\label{notes:XIII}
\setcounter{equation}{18}

\begin{enumerate}[1.]
\item (Section~\hyperref[sec:XIII.1]{$1$}) We can prove (K. \textsc{Yosida} \cite{yosida360}) 

\begin{theoremn}[Theorem~\ref{thm:XIII.1.Theorem.1}']\yosnamedtag{1'}\label{thm:notes.XIII.1.1prime} Let $\mathrm{P}(t, x, E)$ be a Markov process on the phase space $(S, \mathfrak{B})$ with an invariant measure $m$ such that $m(S) = 1$. Then, for any $f \in L^1(S, \mathfrak{B}, m)$ and $t > 0$, there corresponds an $f_{-t} \in L^1(S, \mathfrak{B}, m)$ with $\|f_{-t}\|_1 \leqq \|f\|_1$ in such a way that 
\begin{equation}\label{eq:XIV.7.19a}
\int_{S} m (\dd x) f (x) P (t, x, E) = \int_ {\boldsymbol {E}} f _{-t} (x) m (\dd x) \quad \text{ for all } \quad E \in \mathfrak {B}.
\end{equation}
\end{theoremn}

\begin{proof}
 Since $P(t, x, E)$ is a bounded function, $\int_{S} m(\dd x) f(x)$ $P(t, x, E)$ does exist for any $f \in L^{1}(S, \mathfrak{B}, m)$ and we have 
\begin{equation}\label{eq:XIV.7.20a}
\sup\limits _ {\boldsymbol {E}} | \int_{S} m (\dd x) f (x) P (t, x, E) | \leqq \sup\limits _ {\boldsymbol {E}} \int_{S} m (\dd x) | f (x) | P (t, x, E) \leqq \| f \| _ {\mathbf {1}}. \end{equation}

If $f \in L^{1}(S, \mathfrak{B}, m)$ is $\in L^{\infty}(S, \mathfrak{B}, m)$, then, by \eqref{eq:XIII.1.8}, the $\sigma$-additive function $\int_{S} m(\dd x)f(x)P(t, x, E)$ of $E \in \mathfrak{B}$ is $m$-absolutely continuous such that $|\int_{S} m(\dd x)f(x)P(t, x, E)| \leqq \|f\|_{\mathbb{L}^{\infty}} \cdot m(E)$, where $\|f\|_{\mathbb{L}^{\infty}} = \text{ essential sup } |f(x)|$. Next, let $f \in L^{1}(S, \mathfrak{B}, m)$ be real-valued non-negative and put $f^{(n)}(x) = \min(f(x), n)$. Since 
\[
\int_{S} m (\dd x) f (x) P (t, x, E) = \lim\limits _{n \rightarrow \infty} \int_{S} m (\dd x) f ^{(n)} (x) P (t, x, E) \quad \text{ for every } E \in \mathfrak {B}
\]
the left hand side above is $m$-absolutely continuous by the Vitali-Hahn-Saks theorem on p.~\pageref{thm:II.2.VitaliHahnSaks}. Thus, in this special case for $f$, we can define $f_{-t} \in L^{1}(S, \mathfrak{B}, m)$ by \eqref{eq:XIV.7.19a}, and we have $\|f_{-t}\|_{1} \leqq \|f\|_{1}$ by \eqref{eq:XIV.7.20a}. This proves our Theorem~\ref{thm:notes.XIII.1.1prime}. 
\end{proof}

\begin{remarku} Consider the deterministic case given by $P(t, x, E) = C_E(y_t(x))$, defined through a one-parameter family $\{y_t(x); -\infty < t < \infty\}$ of one-one transformation $x \to y_t(x)$ of $S$ onto $S$ such that $m(E) = m(y_t \cdot E)$. Then we see that the mapping $f \to f_t$ in Theorem~\ref{thm:XIII.1.Theorem.1} (p.~\pageref{thm:XIII.1.Theorem.1}) and the mapping $f \to f_{-t}$ in Theorem~\ref{thm:notes.XIII.1.1prime} are respectively given by $f_t(x) = f(y_t(x))$. 
and $f_{-t}(x) = f(y_{-t}(x))$. In this sense, a Markov process with an invariant measure $m$ is ``time reversible''.
\end{remarku}
\item (Section~\hyperref[sec:XIII.1]{$1$}) As a corollary of Theorem~\ref{thm:notes.XIII.1.1prime} above, we obtain, as in the case of Theorem~\ref{thm:XIII.1.Theorem.2} (p.~\pageref{thm:XIII.1.Theorem.2}), the following 

\begin{theoremn}[Theorem~\ref{thm:XIII.1.Theorem.2}' (K. \textsc{Yosida} {\cite{yosida360}})]\yosnamedtag{2'}\label{thm:notes.XIII.1.2prime} Under the assumption as in Theorem~\ref{thm:notes.XIII.1.1prime}, the following mean ergodic theorem holds good: 
\begin{align}
\underset {n \rightarrow \infty} {\slim} n ^{-1} \sum\limits_{k=1} ^{n} f _{-k}& = f ^{-*} \text{ exists in } L ^{1} (S, \mathfrak {B}, m), \label{eq:XIV.7.21a} \\
&\int_{S} f (s) m (\dd s) = \int_{S} f ^{-*} (s) m (\dd s). \label{eq:XIV.7.22b}
\end{align}
\end{theoremn}
\end{enumerate}

\yossub{Chapter XIV}\phantomsection\label{notes:XIV}

\begin{enumerate}[1.]
\item (Section~\hyperref[sec:XIV.4]{$4$}) We can prove 

\begin{theoremn}[Theorem~\ref{thm:XIV.4.Theorem.1}']\yosnamedtag{1'}\label{thm:notes.XIV.4.1prime} Under the assumption $x(0) = y \in D(A(0))$ and \eqref{eq:XIV.4.3} (p.~\pageref{eq:XIV.4.3}), the solution of \eqref{eq:XIV.4.1'} (p.~\pageref{eq:XIV.4.1'}), if it exists, is uniquely determined.
\end{theoremn}

\begin{proof}
 We follow an argument given in T. \textsc{Kato} \cite{kato136}. Thus we have, for $\delta > 0$, 
\[ \begin{aligned}
&x (t + \delta)  = x (t) + \delta A (t) x (t) + o (\delta) = (I + \delta A (t)) (I - \delta A (t)) (I - \delta A (t)) ^{-1} x (t) + o (\delta) \\
&= (I - \delta A (t)) ^{-1} x (t) - \delta^{2} A (t) (I - \delta A (t)) ^{-1} A (t) x (t) + o (\delta) \\
&= (I - \delta A (t)) ^{-1} x (t) - \delta ((I - \delta A (t)) ^{-1} - I) A (t) x (t) + o (\delta)
\end{aligned}
\]

Because of \eqref{eq:XIV.4.3}, we have $s\text{- }\lim\limits_{\delta \to 0} (I - \delta A(t))^{-1} z = z$ for any $z \in X$ [see \eqref{eq:IX.7.2} in Chapter~\ref{chap:IX}, \ref{sec:IX.7}]. Hence, again by \eqref{eq:XIV.4.3}, we have 
\[
\left\| x (t + \delta) \right\| \leqq \left\| x (t) \right\| + o (\delta)
\]
and so $d^{+}\|x(t)\| / \dd t \leqq 0$ which implies that $\|x(t)\| \leqq \|x(0)\|$. 

Therefore, the difference between two solutions of \eqref{eq:XIV.4.1'} each with the same initial condition $\in D(A(0))$ must be $0$. 
\end{proof}
\end{enumerate}

\yossub{Chapter IX and Chapter XIV}\phantomsection\label{notes:IX-XIV}

\begin{enumerate}[1.]
\item For Linear evolution equations, see S. \textsc{Krein} \cite{krein168}, P. \textsc{Lax}-R. S. \textsc{Phillips} \cite{lewy174}, J. L. \textsc{Lions} \cite{liggett182} and F. \textsc{Trèves} \cite{trotter326}.
\item For Linear as well as nonlinear evolution equations, see: V. \textsc{Barbu} \cite{barbu11}, F. \textsc{Browder} \cite{browder31}, J. L. \textsc{Lions} \cite{liggett182}, R. H. \textsc{Martin} \cite{martin188}, K. \textsc{Masuda} \cite{masuda190}, I. \textsc{Miyadera} \cite{miyadera207} and H. \textsc{Tanabe} \cite{tanabe319}. For the last cited book of Tanabe, an English translation is in preparation.
\item Currently, convex analysis is playing an important rôle in nonlinear evolution equations. For this analysis, see, e.g., R. S. \textsc{Rockafellar} \cite{rockafeller278} and I. \textsc{Ekeland}-R. \textsc{Temam} \cite{ekeland54}.
\end{enumerate}

   % XIV
\yoschapterstar{Bibliography}\label{chap:Bibliography}
\vspace{1.8mm}
\nocite{*}
\printbibliography[heading=none]

\yoschapterstar{Index}\label{chap:Index}
\setlength{\columnsep}{8mm}\setlength{\multicolsep}{0pt}
\begin{multicols}{2}\raggedcolumns\footnotesize\raggedright\setlength{\parindent}{0pt}\setlength{\parskip}{0pt}
\noindent absolutely continuous \hyperlink{page.16}{16}, \hyperlink{page.89}{89}, \hyperlink{page.131}{131}, \hyperlink{page.360}{360}\par
\noindent absorb \hyperlink{page.43}{43}\par
\noindent absorbing \hyperlink{page.24}{24}\par
\noindent abstract $L^{1}$-space \hyperlink{page.354}{354}\par
\noindent accretive operator \hyperlink{page.203}{203}\par
\noindent accumulation point \hyperlink{page.3}{3}\par
\noindent \textsc{Achieser}, N. I. \hyperlink{page.339}{339}, \hyperlink{page.418}{418}, \hyperlink{page.447}{447}\par
\noindent adjoint operator \hyperlink{page.191}{191}\par
\noindent a.e. \hyperlink{page.14}{14}\par
\noindent \textsc{Agmon}, S. \hyperlink{page.422}{422}, \hyperlink{page.447}{447}\par
\noindent \textsc{Aizawa}, S. \hyperlink{page.440}{440}, \hyperlink{page.447}{447}\par
\noindent \textsc{Akilov}, G. \hyperlink{page.170}{170}, \hyperlink{page.447}{447}\par
\noindent \textsc{Alexandrov}, P. \hyperlink{page.447}{447}\par
\noindent algebra \hyperlink{page.285}{285}\par
\hspace*{1.1em}Banach-- \hyperlink{page.285}{285}\par
\hspace*{1.1em}Boolean-- \hyperlink{page.350}{350}\par
\noindent almost everywhere \hyperlink{page.14}{14}\par
\noindent almost periodic function \hyperlink{page.214}{214}, \hyperlink{page.318}{318}\par
\hspace*{1.1em}Bohr's theory \hyperlink{page.318}{318}\par
\hspace*{1.1em}mean value of \hyperlink{page.217}{217}\par
\noindent angle variable \hyperlink{page.374}{374}\par
\noindent approximate identity \hyperlink{page.153}{153}\par
\noindent approximation\par
\hspace*{1.1em}Stone-Weierstrass theorem \hyperlink{page.8}{8}\par
\hspace*{1.1em}Weierstrass' polynomial - theorem \hyperlink{page.7}{7}, \hyperlink{page.243}{243}\par
\hspace*{1.1em}Weierstrass' trigonometric - theorem \hyperlink{page.10}{10}\par
\hspace*{1.1em}of $C_{0}^{0}$-functions by $C_{0}^{\infty}$-functions \hyperlink{page.29}{29}\par
\noindent \textsc{Aronszajn}, N. \hyperlink{page.92}{92}, \hyperlink{page.447}{447}\par
\hspace*{1.1em}--Bergman reproducing kernel \hyperlink{page.92}{92}\par
\noindent Ascoli-Arzelà theorem \hyperlink{page.82}{82}\par
\noindent asymptotic convergence \hyperlink{page.37}{37}, \hyperlink{page.119}{119}\par
\noindent \textsc{Baire}, R.\par
\hspace*{1.1em}category argument \hyperlink{page.11}{11}\par
\hspace*{1.1em}function \hyperlink{page.17}{17}\par
\hspace*{1.1em}--Hausdorff theorem \hyperlink{page.11}{11}\par
\hspace*{1.1em}measure \hyperlink{page.16}{16}\par
\hspace*{1.1em}set \hyperlink{page.17}{17}\par
\hspace*{1.1em}theorem \hyperlink{page.11}{11}\par
\noindent \textsc{Balakrishnan}, V. \hyperlink{page.252}{252}, \hyperlink{page.257}{257}, \hyperlink{page.447}{447}\par
\noindent balanced \hyperlink{page.24}{24}\par
\noindent $B$-algebra \hyperlink{page.285}{285}\par
\noindent \textsc{Banach}, S. \hyperlink{page.23}{23}, \hyperlink{page.65}{65}, \hyperlink{page.69}{69}, \hyperlink{page.71}{71}, \hyperlink{page.74}{74}, \hyperlink{page.99}{99}, \hyperlink{page.101}{101}, \hyperlink{page.103}{103}, \hyperlink{page.109}{109}, \hyperlink{page.138}{138}, \hyperlink{page.200}{200}, \hyperlink{page.355}{355}, \hyperlink{page.447}{447}\par
\hspace*{1.1em}--algebra(or $B$--) \hyperlink{page.285}{285}\par
\hspace*{1.1em}--lattice (or $B$--) \hyperlink{page.354}{354}\par
\hspace*{1.1em}--space (or $B$--) \hyperlink{page.23}{23}, \hyperlink{page.50}{50}\par
\hspace*{2.2em}locally compact \hyperlink{page.82}{82}\par
\hspace*{2.2em}reflexive \hyperlink{page.123}{123}, \hyperlink{page.137}{137}\par
\hspace*{1.1em}closed graph theorem \hyperlink{page.74}{74}\par
\hspace*{1.1em}generalized limit \hyperlink{page.101}{101}\par
\hspace*{1.1em}open mapping theorem \hyperlink{page.71}{71}\par
\hspace*{1.1em}theorem on a.e. convergence of linear operators \hyperlink{page.355}{355}\par
\noindent barrel space \hyperlink{page.134}{134}\par
\noindent base, open \hyperlink{page.57}{57}\par
\noindent basis \hyperlink{page.19}{19}\par
\noindent \textsc{Beboutov}, M. \hyperlink{page.376}{376}, \hyperlink{page.447}{447}\par
\noindent \textsc{Bénilan}, P. \hyperlink{page.440}{440}, \hyperlink{page.447}{447}\par
\noindent \textsc{Bergman}, S. \hyperlink{page.92}{92}, \hyperlink{page.447}{447}\par
\hspace*{1.1em}kernel \hyperlink{page.92}{92}\par
\noindent \textsc{Bernstein}, S. \hyperlink{page.7}{7}\par
\noindent \textsc{Bers}, L. \hyperlink{page.172}{172}, \hyperlink{page.447}{447}\par
\noindent Bessel's inequality \hyperlink{page.83}{83}\par
\noindent bidual space \hyperlink{page.109}{109}\par
\noindent \textsc{Birkhoff}, G. \hyperlink{page.348}{348}, \hyperlink{page.362}{362}, \hyperlink{page.447}{447}\par
\noindent \textsc{Birkhoff}, G. D. \hyperlink{page.371}{371}, \hyperlink{page.373}{373}, \hyperlink{page.447}{447}\par
\noindent $B$-lattice \hyperlink{page.354}{354}\par
\noindent \textsc{Bochner}, S. \hyperlink{page.128}{128}, \hyperlink{page.252}{252}, \hyperlink{page.257}{257}, \hyperlink{page.319}{319}, \hyperlink{page.332}{332}, \hyperlink{page.448}{448}\par
\hspace*{1.1em}integral \hyperlink{page.128}{128}\par
\hspace*{1.1em}theorem \hyperlink{page.332}{332}\par
\noindent \textsc{Bogolioubov}, N. \hyperlink{page.376}{376}, \hyperlink{page.448}{448}\par
\noindent \textsc{Bohnenblust}, H. F. \hyperlink{page.102}{102}\par
\noindent \textsc{Bohr}, H. \hyperlink{page.319}{319}, \hyperlink{page.448}{448}\par
\noindent Boolean algebra \hyperlink{page.350}{350}\par
\noindent \textsc{Borel}, E. \hyperlink{page.17}{17}\par
\hspace*{1.1em}measure \hyperlink{page.16}{16}\par
\hspace*{1.1em}set \hyperlink{page.17}{17}\par
\noindent bornologic space \hyperlink{page.43}{43}\par
\noindent bound\par
\hspace*{1.1em}greatest lower \hyperlink{page.2}{2}, \hyperlink{page.352}{352}\par
\hspace*{1.1em}least upper \hyperlink{page.2}{2}, \hyperlink{page.352}{352}\par
\noindent bound\par
\hspace*{1.1em}lower \hyperlink{page.2}{2}\par
\hspace*{1.1em}upper \hyperlink{page.2}{2}\par
\noindent boundary\par
\hspace*{1.1em}condition \hyperlink{page.276}{276}, \hyperlink{page.387}{387}\par
\hspace*{1.1em}function \hyperlink{page.159}{159}\par
\noindent bounded\par
\hspace*{1.1em}convergence topology \hyperlink{page.107}{107}\par
\hspace*{1.1em}essentially -- function \hyperlink{page.33}{33}\par
\hspace*{1.1em}linear operator \hyperlink{page.41}{41}\par
\hspace*{1.1em}set \hyperlink{page.5}{5}, \hyperlink{page.43}{43}, \hyperlink{page.352}{352}\par
\hspace*{1.1em}totally -- set \hyperlink{page.12}{12}\par
\noindent boundedness\par
\hspace*{1.1em}uniform -- theorem \hyperlink{page.65}{65}\par
\noindent \textsc{Bourbaki}, N. \hyperlink{page.23}{23}, \hyperlink{page.50}{50}, \hyperlink{page.141}{141}, \hyperlink{page.448}{448}\par
\noindent \textsc{Brezis}, H. \hyperlink{page.430}{430}, \hyperlink{page.433}{433}, \hyperlink{page.434}{434}, \hyperlink{page.438}{438}, \hyperlink{page.448}{448}\par
\noindent \textsc{Browder}, F. E. \hyperlink{page.448}{448}\par
\noindent \textsc{Brown}, R. \hyperlink{page.363}{363}\par
\noindent Brownian motion \hyperlink{page.381}{381}, \hyperlink{page.385}{385}\par
\noindent $B$-space \hyperlink{page.23}{23}, \hyperlink{page.50}{50}, \hyperlink{page.82}{82}, \hyperlink{page.122}{122}, \hyperlink{page.137}{137}\par
\noindent canonical\par
\hspace*{1.1em}scale \hyperlink{page.387}{387}\par
\hspace*{1.1em}measure \hyperlink{page.387}{387}\par
\noindent carrier (= support) \hyperlink{page.26}{26}\par
\noindent \textsc{Cartan}, E. \hyperlink{page.382}{382}\par
\noindent Cartesian product \hyperlink{page.3}{3}\par
\noindent category\par
\hspace*{1.1em}first \hyperlink{page.11}{11}\par
\hspace*{1.1em}second \hyperlink{page.11}{11}\par
\noindent Cauchy\par
\hspace*{1.1em}convergence condition \hyperlink{page.10}{10}, \hyperlink{page.50}{50}\par
\hspace*{1.1em}integral representation \hyperlink{page.124}{124}\par
\hspace*{1.1em}integral theorem \hyperlink{page.124}{124}\par
\hspace*{1.1em}sequence \hyperlink{page.10}{10}\par
\noindent Cayley transform \hyperlink{page.197}{197}\par
\noindent \textsc{Chacon}, R. V. \hyperlink{page.369}{369}, \hyperlink{page.372}{372}, \hyperlink{page.448}{448}\par
\hspace*{1.1em}--Ornstein theorem \hyperlink{page.369}{369}\par
\noindent Chapman-Kolmogorov equation \hyperlink{page.363}{363}\par
\noindent character, continuous unitary \hyperlink{page.340}{340}\par
\noindent \textsc{Choquet}, G. \hyperlink{page.348}{348}, \hyperlink{page.448}{448}\par
\noindent \textsc{Clarkson}, J. A. \hyperlink{page.123}{123}, \hyperlink{page.448}{448}\par
\noindent closable operator \hyperlink{page.73}{73}\par
\noindent closed\par
\hspace*{1.1em}graph theorem \hyperlink{page.74}{74}\par
\hspace*{1.1em}linear operator \hyperlink{page.73}{73}\par
\hspace*{1.1em}range theorem \hyperlink{page.200}{200}\par
\hspace*{1.1em}set \hyperlink{page.3}{3}\par
\noindent closure\par
\hspace*{1.1em}convex \hyperlink{page.27}{27}\par
\hspace*{1.1em}strong \hyperlink{page.122}{122}\par
\hspace*{1.1em}weak \hyperlink{page.122}{122}\par
\noindent commutative \hyperlink{page.326}{326}\par
\noindent compact \hyperlink{page.4}{4}\par
\hspace*{1.1em}group \hyperlink{page.319}{319}\par
\hspace*{1.1em}in $B$-spaces \hyperlink{page.265}{265}\par
\hspace*{1.1em}locally -- \hyperlink{page.4}{4}\par
\noindent compact\par
\hspace*{1.1em}locally sequentially weakly -- \hyperlink{page.122}{122}, \hyperlink{page.138}{138}\par
\hspace*{1.1em}operator \hyperlink{page.265}{265}, \hyperlink{page.267}{267}, \hyperlink{page.313}{313}\par
\hspace*{1.1em}relatively -- \hyperlink{page.4}{4}\par
\noindent compactification, one point -- \hyperlink{page.4}{4}\par
\noindent complement \hyperlink{page.3}{3}, \hyperlink{page.350}{350}\par
\hspace*{1.1em}orthogonal \hyperlink{page.79}{79}\par
\noindent complete \hyperlink{page.10}{10}, \hyperlink{page.50}{50}, \hyperlink{page.102}{102}, \hyperlink{page.352}{352}\par
\hspace*{1.1em}orthonormal system \hyperlink{page.83}{83}\par
\noindent $\sigma-367$ \hyperlink{page.352}{352}\par
\hspace*{1.1em}sequentially -- \hyperlink{page.102}{102}\par
\hspace*{1.1em}sequentially weakly -- \hyperlink{page.117}{117}, \hyperlink{page.121}{121}\par
\noindent completely continuous operator \hyperlink{page.265}{265}\par
\noindent completion \hyperlink{page.53}{53}, \hyperlink{page.102}{102}, \hyperlink{page.110}{110}\par
\noindent complex measure \hyperlink{page.34}{34}\par
\noindent condensation of singularities \hyperlink{page.69}{69}\par
\noindent conjugate\par
\hspace*{1.1em}exponent \hyperlink{page.32}{32}\par
\hspace*{1.1em}operator \hyperlink{page.189}{189}\par
\noindent continuous\par
\hspace*{1.1em}absolutely \hyperlink{page.16}{16}\par
\hspace*{1.1em}mapping \hyperlink{page.3}{3}\par
\hspace*{1.1em}spectrum \hyperlink{page.205}{205}, \hyperlink{page.311}{311}\par
\hspace*{1.1em}unitary character \hyperlink{page.340}{340}\par
\noindent contracted measure \hyperlink{page.362}{362}\par
\noindent contraction semi-group \hyperlink{page.228}{228}\par
\noindent convergence \hyperlink{page.101}{101}\par
\hspace*{1.1em}asymptotic -- \hyperlink{page.37}{37}, \hyperlink{page.119}{119}\par
\hspace*{1.1em}in measure \hyperlink{page.37}{37}, \hyperlink{page.119}{119}\par
\hspace*{1.1em}mean -- \hyperlink{page.33}{33}\par
\hspace*{1.1em}$O$-- \hyperlink{page.351}{351}\par
\hspace*{1.1em}relatively uniform *-- \hyperlink{page.355}{355}\par
\hspace*{1.1em}strong -- \hyperlink{page.30}{30}, \hyperlink{page.31}{31}\par
\hspace*{1.1em}weak -- \hyperlink{page.117}{117}\par
\hspace*{1.1em}weak* -- \hyperlink{page.122}{122}\par
\noindent convex \hyperlink{page.24}{24}\par
\hspace*{1.1em}closure \hyperlink{page.27}{27}\par
\hspace*{1.1em}hull \hyperlink{page.347}{347}\par
\hspace*{1.1em}uniformly - space \hyperlink{page.123}{123}\par
\noindent convolution \hyperlink{page.145}{145}, \hyperlink{page.152}{152}, \hyperlink{page.339}{339}\par
\noindent coordinate operator \hyperlink{page.193}{193}, \hyperlink{page.336}{336}\par
\noindent countably additive (= $\sigma$-additive) \hyperlink{page.14}{14}\par
\noindent covering \hyperlink{page.4}{4}, \hyperlink{page.57}{57}\par
\hspace*{1.1em}scattered open -- \hyperlink{page.57}{57}\par
\noindent \textsc{Crandall}, M. G. \hyperlink{page.430}{430}, \hyperlink{page.434}{434}, \hyperlink{page.440}{440}, \hyperlink{page.448}{448}\par
\noindent \textsc{Crumm}, M. M. \hyperlink{page.161}{161}\par
\noindent $D^{\alpha}$ \hyperlink{page.27}{27}\par
\noindent $D_{\alpha}$ \hyperlink{page.176}{176}\par
\noindent decomposition\par
\hspace*{1.1em}Hahn -- \hyperlink{page.35}{35}\par
\hspace*{1.1em}Jordan -- \hyperlink{page.35}{35}, \hyperlink{page.351}{351}\par
\hspace*{1.1em}of a subharmonic element \hyperlink{page.392}{392}\par
\noindent defect indices \hyperlink{page.334}{334}\par
\noindent defining function \hyperlink{page.16}{16}\par
\noindent dense \hyperlink{page.11}{11}\par
\noindent derivative\par
\hspace*{1.1em}Dini -- \hyperlink{page.234}{234}\par
\hspace*{1.1em}distributional -- \hyperlink{page.47}{47}, \hyperlink{page.68}{68}\par
\noindent \textsc{Dieudonné}, J. \hyperlink{page.102}{102}, \hyperlink{page.448}{448}\par
\noindent difference\par
\hspace*{1.1em}operator \hyperlink{page.238}{238}\par
\hspace*{1.1em}symmetric -- \hyperlink{page.18}{18}\par
\noindent differential operator\par
\hspace*{1.1em}analytically elliptic -- \hyperlink{page.187}{187}\par
\hspace*{1.1em}as infinitesimal generator \hyperlink{page.237}{237}\par
\hspace*{1.1em}formally hypoelliptic -- \hyperlink{page.186}{186}\par
\hspace*{1.1em}hypoelliptic -- \hyperlink{page.75}{75}, \hyperlink{page.183}{183}\par
\hspace*{1.1em}parabolic -- \hyperlink{page.182}{182}\par
\hspace*{1.1em}strictly elliptic -- \hyperlink{page.394}{394}\par
\hspace*{1.1em}strongly elliptic -- \hyperlink{page.170}{170}, \hyperlink{page.182}{182}, \hyperlink{page.276}{276}\par
\hspace*{1.1em}uniformly strong -- \hyperlink{page.182}{182}\par
\hspace*{1.1em}weaker \hyperlink{page.178}{178}\par
\noindent diffusion \hyperlink{page.380}{380}, \hyperlink{page.390}{390}\par
\hspace*{1.1em}equation \hyperlink{page.380}{380}, \hyperlink{page.400}{400}, \hyperlink{page.405}{405}\par
\noindent dimension \hyperlink{page.19}{19}\par
\noindent Dini derivative \hyperlink{page.234}{234}\par
\noindent Dirac distribution \hyperlink{page.46}{46}, \hyperlink{page.48}{48}, \hyperlink{page.149}{149}\par
\noindent directed set \hyperlink{page.100}{100}\par
\hspace*{1.1em}fundamental \hyperlink{page.102}{102}\par
\noindent Dirichlet's problem \hyperlink{page.275}{275}\par
\noindent discrete topological space \hyperlink{page.33}{33}\par
\noindent dissipative\par
\hspace*{1.1em}Markov process \hyperlink{page.377}{377}\par
\hspace*{1.1em}operator \hyperlink{page.203}{203}, \hyperlink{page.244}{244}\par
\hspace*{1.1em}part \hyperlink{page.378}{378}\par
\hspace*{1.1em}set \hyperlink{page.424}{424}\par
\noindent distance function \hyperlink{page.3}{3}, \hyperlink{page.29}{29}\par
\noindent distribution (= generalized function) \hyperlink{page.45}{45}\par
\hspace*{1.1em}convergence of -- \hyperlink{page.119}{119}\par
\hspace*{1.1em}Dirac -- \hyperlink{page.46}{46}, \hyperlink{page.48}{48}, \hyperlink{page.149}{149}\par
\hspace*{1.1em}direct (= tensor) product of -- \hyperlink{page.63}{63}\par
\hspace*{1.1em}local structure of a -- \hyperlink{page.96}{96}\par
\hspace*{1.1em}product with a function \hyperlink{page.49}{49}\par
\hspace*{1.1em}tempered -- \hyperlink{page.146}{146}\par
\hspace*{1.1em}with compact support \hyperlink{page.59}{59}, \hyperlink{page.61}{61}\par
\noindent distributional derivative \hyperlink{page.47}{47}, \hyperlink{page.68}{68}\par
\noindent distributive\par
\hspace*{1.1em}identity \hyperlink{page.350}{350}\par
\hspace*{1.1em}lattice \hyperlink{page.350}{350}\par
\noindent divergence, generalized \hyperlink{page.387}{387}\par
\noindent domain 1, \hyperlink{page.1}{1}, \hyperlink{page.19}{19}\par
\noindent \textsc{Doob}, J. L. \hyperlink{page.371}{371}, \hyperlink{page.449}{449}\par
\noindent \textsc{Dorroh}, J. R. \hyperlink{page.430}{430}, \hyperlink{page.449}{449}\par
\noindent duality map \hyperlink{page.424}{424}\par
\noindent duality theorem of Tannaka \hyperlink{page.323}{323}\par
\noindent dual operator \hyperlink{page.189}{189}\par
\hspace*{1.1em}inverse of the -- \hyperlink{page.220}{220}\par
\hspace*{1.1em}resolvent of the -- \hyperlink{page.220}{220}\par
\noindent dual semi-group \hyperlink{page.262}{262}\par
\noindent dual space \hyperlink{page.46}{46}, \hyperlink{page.88}{88}, \hyperlink{page.108}{108}\par
\hspace*{1.1em}strong topology \hyperlink{page.108}{108}\par
\hspace*{1.1em}weak* topology \hyperlink{page.108}{108}\par
\noindent \textsc{Dufresnoy}, J. \hyperlink{page.162}{162}\par
\noindent Duhamel principle \hyperlink{page.411}{411}\par
\noindent \textsc{Dunford}, N. \hyperlink{page.124}{124}, \hyperlink{page.221}{221}, \hyperlink{page.229}{229}, \hyperlink{page.313}{313}, \hyperlink{page.372}{372}, \hyperlink{page.449}{449}\par
\hspace*{1.1em}integral \hyperlink{page.221}{221}\par
\hspace*{1.1em}theorem \hyperlink{page.124}{124}\par
\noindent \textsc{Dynkin}, E. B. \hyperlink{page.380}{380}, \hyperlink{page.389}{389}, \hyperlink{page.449}{449}\par
\noindent \textsc{Eberlein}, W. F. \hyperlink{page.137}{137}, \hyperlink{page.449}{449}\par
\hspace*{1.1em}--Shmulyan theorem \hyperlink{page.137}{137}\par
\noindent Egorov's theorem \hyperlink{page.14}{14}\par
\noindent \textsc{Ehrenpreis}, L. \hyperlink{page.177}{177}, \hyperlink{page.449}{449}\par
\noindent eigenspace \hyperlink{page.205}{205}\par
\noindent eigenvalue \hyperlink{page.205}{205}\par
\noindent eigenvector \hyperlink{page.205}{205}\par
\noindent elliptic differential equation\par
\hspace*{1.1em}analytically -- \hyperlink{page.187}{187}\par
\hspace*{1.1em}formally -- \hyperlink{page.186}{186}\par
\hspace*{1.1em}strictly -- \hyperlink{page.400}{400}\par
\hspace*{1.1em}strongly -- \hyperlink{page.170}{170}, \hyperlink{page.182}{182}, \hyperlink{page.276}{276}\par
\noindent energy inequality \hyperlink{page.409}{409}\par
\noindent equation\par
\hspace*{1.1em}diffusion -- \hyperlink{page.380}{380}, \hyperlink{page.400}{400}, \hyperlink{page.405}{405}\par
\hspace*{1.1em}of evolution \hyperlink{page.399}{399}, \hyperlink{page.410}{410}, \hyperlink{page.417}{417}\par
\hspace*{1.1em}wave -- \hyperlink{page.407}{407}\par
\noindent equi-bounded \hyperlink{page.82}{82}\par
\hspace*{1.1em}semi-group \hyperlink{page.228}{228}\par
\noindent equi-continuous \hyperlink{page.82}{82}, \hyperlink{page.209}{209}\par
\hspace*{1.1em}semi-group of class (Co) \hyperlink{page.229}{229}\par
\hspace*{1.1em}group of class (Co) \hyperlink{page.245}{245}\par
\noindent equi-measure transformation \hyperlink{page.210}{210}, \hyperlink{page.364}{364}\par
\noindent \textsc{Erdélyi}, A. \hyperlink{page.168}{168}, \hyperlink{page.449}{449}\par
\noindent ergodic\par
\hspace*{1.1em}decomposition of a Markov process \hyperlink{page.375}{375}\par
\noindent ergodic\par
\hspace*{1.1em}hypothesis \hyperlink{page.372}{372}\par
\hspace*{2.2em}of Boltzmann \hyperlink{page.364}{364}\par
\noindent ergodic theorem\par
\hspace*{1.1em}individual -- \hyperlink{page.366}{366}, \hyperlink{page.371}{371}\par
\hspace*{1.1em}mean -- \hyperlink{page.209}{209}, \hyperlink{page.366}{366}\par
\hspace*{1.1em}of the Hille type \hyperlink{page.211}{211}\par
\noindent essentially\par
\hspace*{1.1em}bounded function \hyperlink{page.33}{33}\par
\hspace*{1.1em}self-adjoint operator \hyperlink{page.304}{304}\par
\noindent evolution equation \hyperlink{page.399}{399}, \hyperlink{page.410}{410}, \hyperlink{page.417}{417}\par
\hspace*{1.1em}non--linear -- \hyperlink{page.422}{422}, \hyperlink{page.424}{424}, \hyperlink{page.430}{430}\par
\noindent exponential of a bounded operator \hyperlink{page.238}{238}\par
\noindent extension of an operator \hyperlink{page.1}{1}, \hyperlink{page.19}{19}\par
\hspace*{1.1em}smallest closed -- \hyperlink{page.73}{73}\par
\noindent extremal\par
\hspace*{1.1em}point \hyperlink{page.347}{347}\par
\hspace*{1.1em}subset \hyperlink{page.347}{347}\par
\noindent factor\par
\hspace*{1.1em}group \hyperlink{page.215}{215}\par
\hspace*{1.1em}space \hyperlink{page.20}{20}, \hyperlink{page.56}{56}\par
\noindent Fatou, Lebesgue-- -- lemma \hyperlink{page.16}{16}\par
\noindent \textsc{Feller}, W. \hyperlink{page.242}{242}, \hyperlink{page.380}{380}, \hyperlink{page.385}{385}, \hyperlink{page.387}{387}, \hyperlink{page.449}{449}\par
\noindent \textsc{Feynman}, R. P., integral \hyperlink{page.440}{440}\par
\noindent field, normed -- \hyperlink{page.125}{125}\par
\noindent finite intersection property \hyperlink{page.5}{5}\par
\noindent finitely additive \hyperlink{page.68}{68}\par
\hspace*{1.1em}measure \hyperlink{page.115}{115}\par
\noindent finitely valued function \hyperlink{page.15}{15}, \hyperlink{page.126}{126}\par
\noindent $F$-lattice \hyperlink{page.354}{354}\par
\noindent \textsc{Foias}, C. \hyperlink{page.450}{450}\par
\noindent formally hypoelliptic \hyperlink{page.186}{186}\par
\noindent forward and backward unique continuation property \hyperlink{page.405}{405}\par
\noindent \textsc{Fourier}, J. B. J.\par
\hspace*{1.1em}absolutely convergent -- series \hyperlink{page.285}{285}, \hyperlink{page.292}{292}\par
\hspace*{1.1em}coefficient \hyperlink{page.83}{83}\par
\hspace*{1.1em}expansion \hyperlink{page.70}{70}, \hyperlink{page.84}{84}\par
\hspace*{1.1em}inverse Fourier transform \hyperlink{page.144}{144}\par
\hspace*{1.1em}--Laplace transform \hyperlink{page.157}{157}\par
\hspace*{1.1em}transform \hyperlink{page.143}{143}, \hyperlink{page.148}{148}, \hyperlink{page.149}{149}\par
\noindent fractional power of a closed operator \hyperlink{page.252}{252}\par
\noindent \textsc{Fréchet}, M. \hyperlink{page.23}{23}, \hyperlink{page.37}{37}, \hyperlink{page.50}{50}, \hyperlink{page.265}{265}\par
\hspace*{1.1em}--Kolmogorov theorem \hyperlink{page.265}{265}\par
\hspace*{1.1em}--space (or F--) \hyperlink{page.50}{50}\par
\noindent Fredholm theory (Riesz-Schauder theory) \hyperlink{page.265}{265}, \hyperlink{page.273}{273}\par
\noindent \textsc{Freudenthal}, H. \hyperlink{page.306}{306}, \hyperlink{page.348}{348}, \hyperlink{page.450}{450}\par
\noindent \textsc{Friedmann}, A. \hyperlink{page.184}{184}, \hyperlink{page.450}{450}\par
\noindent \textsc{Friedrichs}, K. \hyperlink{page.172}{172}, \hyperlink{page.203}{203}, \hyperlink{page.306}{306}, \hyperlink{page.410}{410}\par
\hspace*{1.1em}theorem on semi-bounded operators \hyperlink{page.305}{305}\par
\hspace*{1.1em}theorem on strongly elliptic differential equations \hyperlink{page.172}{172}\par
\noindent $F$-space \hyperlink{page.50}{50}\par
\noindent Fubini-Tonelli theorem \hyperlink{page.16}{16}\par
\noindent \textsc{Fujita}, H. \hyperlink{page.411}{411}\par
\noindent FUKAMiya, M. \hyperlink{page.344}{344}, \hyperlink{page.359}{359}, \hyperlink{page.450}{450}\par
\noindent function \hyperlink{page.1}{1}\par
\hspace*{1.1em}almost periodic -- \hyperlink{page.214}{214}, \hyperlink{page.318}{318}\par
\hspace*{1.1em}Baire -- \hyperlink{page.17}{17}\par
\hspace*{1.1em}Bochner m--integrable -- \hyperlink{page.128}{128}\par
\hspace*{1.1em}boundary -- \hyperlink{page.159}{159}\par
\hspace*{1.1em}continuous -- \hyperlink{page.4}{4}\par
\hspace*{1.1em}defining -- \hyperlink{page.16}{16}\par
\hspace*{1.1em}finitely valued -- \hyperlink{page.15}{15}\par
\hspace*{1.1em}generalized -- \hyperlink{page.45}{45}\par
\hspace*{1.1em}harmonic -- \hyperlink{page.391}{391}\par
\hspace*{1.1em}Heaviside -- \hyperlink{page.47}{47}\par
\hspace*{1.1em}holomorphic -- \hyperlink{page.39}{39}\par
\hspace*{1.1em}ideal -- \hyperlink{page.45}{45}\par
\hspace*{1.1em}integrable -- \hyperlink{page.15}{15}, \hyperlink{page.16}{16}\par
\hspace*{1.1em}locally integrable -- \hyperlink{page.46}{46}\par
\hspace*{1.1em}measurable -- \hyperlink{page.14}{14}\par
\hspace*{1.1em}positive definite -- \hyperlink{page.332}{332}\par
\hspace*{1.1em}potential -- \hyperlink{page.391}{391}\par
\hspace*{1.1em}rapidly decreasing -- \hyperlink{page.143}{143}\par
\hspace*{1.1em}slowly increasing -- \hyperlink{page.147}{147}\par
\hspace*{1.1em}testing -- \hyperlink{page.45}{45}\par
\hspace*{1.1em}uniformly continuous -- \hyperlink{page.319}{319}\par
\noindent fundamental\par
\hspace*{1.1em}directed set \hyperlink{page.102}{102}\par
\hspace*{1.1em}solution \hyperlink{page.176}{176}\par
\hspace*{1.1em}system of neighbourhoods \hyperlink{page.27}{27}\par
\noindent \textsc{Garding}, L. \hyperlink{page.170}{170}, \hyperlink{page.271}{271}, \hyperlink{page.276}{276}, \hyperlink{page.450}{450}\par
\hspace*{1.1em}inequality \hyperlink{page.170}{170}\par
\noindent Gaussian probability density \hyperlink{page.230}{230}, \hyperlink{page.237}{237}, \hyperlink{page.259}{259}\par
\noindent G$\delta$-set \hyperlink{page.17}{17}\par
\noindent \textsc{Gelfand}, I. M. \hyperlink{page.124}{124}, \hyperlink{page.161}{161}, \hyperlink{page.271}{271}, \hyperlink{page.283}{283}, \hyperlink{page.286}{286}, \hyperlink{page.289}{289}, \hyperlink{page.341}{341}, \hyperlink{page.450}{450}\par
\hspace*{1.1em}--Mazur theorem \hyperlink{page.124}{124}\par
\hspace*{1.1em}--Raikov theorem \hyperlink{page.341}{341}\par
\hspace*{1.1em}representation \hyperlink{page.289}{289}\par
\noindent general expansion theorem \hyperlink{page.311}{311}\par
\noindent generalized derivative \hyperlink{page.47}{47}\par
\hspace*{1.1em}divergence \hyperlink{page.387}{387}\par
\noindent generalized\par
\hspace*{1.1em}function \hyperlink{page.45}{45}\par
\hspace*{1.1em}function with compact support \hyperlink{page.59}{59}\par
\hspace*{1.1em}gradient \hyperlink{page.387}{387}\par
\hspace*{1.1em}Laplacian \hyperlink{page.385}{385}\par
\hspace*{1.1em}Leibniz' formula \hyperlink{page.49}{49}\par
\hspace*{1.1em}limit \hyperlink{page.101}{101}\par
\hspace*{1.1em}nilpotent element \hyperlink{page.289}{289}, \hyperlink{page.292}{292}\par
\hspace*{1.1em}spectral resolution \hyperlink{page.338}{338}\par
\noindent generator, infinitesimal -- \hyperlink{page.227}{227}, \hyperlink{page.232}{232}\par
\noindent \textsc{Glazman}, I. M. \hyperlink{page.339}{339}, \hyperlink{page.451}{451}\par
\noindent gradient, generalized -- \hyperlink{page.387}{387}\par
\noindent graph \hyperlink{page.73}{73}\par
\noindent Green's integral theorem \hyperlink{page.48}{48}\par
\noindent \textsc{Grothendieck}, A. \hyperlink{page.141}{141}, \hyperlink{page.268}{268}, \hyperlink{page.279}{279}, \hyperlink{page.283}{283}, \hyperlink{page.451}{451}\par
\noindent group \hyperlink{page.214}{214}, \hyperlink{page.319}{319}\par
\hspace*{1.1em}of isotropy \hyperlink{page.382}{382}\par
\hspace*{1.1em}--ring \hyperlink{page.339}{339}\par
\noindent group, equicontinuous -- of class (C$_{0}$) \hyperlink{page.245}{245}\par
\hspace*{1.1em}representation of a -- \hyperlink{page.318}{318}\par
\hspace*{1.1em}of unitary operators (Stone's theorem) \hyperlink{page.331}{331}\par
\noindent \textsc{Hahn}, H. \hyperlink{page.35}{35}, \hyperlink{page.66}{66}, \hyperlink{page.99}{99}, \hyperlink{page.102}{102}, \hyperlink{page.103}{103}\par
\hspace*{1.1em}--Banach theorem \hyperlink{page.99}{99}, \hyperlink{page.102}{102}, \hyperlink{page.103}{103}\par
\hspace*{1.1em}decomposition \hyperlink{page.35}{35}\par
\noindent \textsc{Halmos}, P. R. \hyperlink{page.91}{91}, \hyperlink{page.115}{115}, \hyperlink{page.311}{311}, \hyperlink{page.376}{376}, \hyperlink{page.390}{390}, \hyperlink{page.451}{451}\par
\noindent \textsc{Hardy}, G. H. \hyperlink{page.39}{39}\par
\hspace*{1.1em}--Lebesgue class \hyperlink{page.39}{39}, \hyperlink{page.50}{50}, \hyperlink{page.159}{159}\par
\noindent harmonic \hyperlink{page.391}{391}\par
\hspace*{1.1em}least -- majorant \hyperlink{page.392}{392}\par
\noindent \textsc{Hausdorff}, F. \hyperlink{page.2}{2}, \hyperlink{page.11}{11}, \hyperlink{page.71}{71}, \hyperlink{page.451}{451}\par
\hspace*{1.1em}axiom of separation \hyperlink{page.2}{2}\par
\hspace*{1.1em}Baire-- theorem \hyperlink{page.11}{11}\par
\noindent \textsc{Heaviside}, O. \hyperlink{page.164}{164}\par
\hspace*{1.1em}function \hyperlink{page.48}{48}\par
\hspace*{1.1em}operational calculus \hyperlink{page.164}{164}\par
\noindent \textsc{Hellinger}, E. \hyperlink{page.311}{311}, \hyperlink{page.451}{451}\par
\noindent \textsc{Helly}, E. \hyperlink{page.106}{106}, \hyperlink{page.110}{110}, \hyperlink{page.451}{451}\par
\noindent Hermite polynomials \hyperlink{page.86}{86}\par
\noindent \textsc{Hilbert}, D.\par
\hspace*{1.1em}--Schmidt integral operator \hyperlink{page.193}{193}, \hyperlink{page.267}{267}\par
\hspace*{1.1em}--space \hyperlink{page.50}{50}, \hyperlink{page.76}{76}, \hyperlink{page.83}{83}\par
\noindent \textsc{Hille}, E. \hyperlink{page.211}{211}, \hyperlink{page.227}{227}, \hyperlink{page.242}{242}, \hyperlink{page.248}{248}, \hyperlink{page.286}{286}, \hyperlink{page.380}{380}, \hyperlink{page.430}{430}, \hyperlink{page.437}{437}, \hyperlink{page.451}{451}\par
\hspace*{1.1em}--Yosida theorem \hyperlink{page.242}{242}\par
\hspace*{1.1em}\parbox[t]{\dimexpr\linewidth-1.1em\relax}{(The result pertaining to (9) of Corollary~1 on p.~\pageref{cor:IX.7.Corollary.1} is sometimes called the Hille--Yosida theorem.)}\par
\noindent \textsc{Hirsch}, F. \hyperlink{page.398}{398}, \hyperlink{page.452}{452}\par
\noindent \textsc{Hoffmann}, K. \hyperlink{page.348}{348}, \hyperlink{page.452}{452}\par
\noindent Hölder's inequality \hyperlink{page.32}{32}\par
\noindent Holomorphic \hyperlink{page.39}{39}\par
\hspace*{1.1em}semi--group \hyperlink{page.247}{247}\par
\hspace*{1.1em}strongly \hyperlink{page.125}{125}\par
\hspace*{1.1em}weakly \hyperlink{page.124}{124}\par
\noindent homeomorphic \hyperlink{page.72}{72}\par
\noindent \textsc{Hopf}, E. \hyperlink{page.332}{332}, \hyperlink{page.366}{366}, \hyperlink{page.368}{368}, \hyperlink{page.372}{372}, \hyperlink{page.452}{452}\par
\noindent \textsc{Hörmander}, L. \hyperlink{page.49}{49}, \hyperlink{page.75}{75}, \hyperlink{page.176}{176}, \hyperlink{page.177}{177}, \hyperlink{page.182}{182}, \hyperlink{page.183}{183}, \hyperlink{page.186}{186}, \hyperlink{page.452}{452}\par
\hspace*{1.1em}theorem \hyperlink{page.75}{75}, \hyperlink{page.184}{184}\par
\noindent $H$-theorem \hyperlink{page.375}{375}\par
\noindent \textsc{Hunt}, G. A. \hyperlink{page.380}{380}, \hyperlink{page.397}{397}, \hyperlink{page.452}{452}\par
\noindent hyperdissipative set \hyperlink{page.426}{426}, \hyperlink{page.429}{429}, \hyperlink{page.430}{430}, \hyperlink{page.433}{433}, \hyperlink{page.437}{437}, \hyperlink{page.440}{440}\par
\noindent hypermaximal operator \hyperlink{page.336}{336}\par
\noindent hypoelliptic \hyperlink{page.75}{75}, \hyperlink{page.183}{183}\par
\hspace*{1.1em}formally \hyperlink{page.186}{186}\par
\noindent ideal \hyperlink{page.286}{286}, \hyperlink{page.356}{356}\par
\hspace*{1.1em}function \hyperlink{page.45}{45}\par
\hspace*{1.1em}maximal -- \hyperlink{page.286}{286}\par
\hspace*{1.1em}theory in normed rings \hyperlink{page.286}{286}\par
\noindent idempotent \hyperlink{page.80}{80}\par
\noindent identity\par
\hspace*{1.1em}distributive -- \hyperlink{page.350}{350}\par
\hspace*{1.1em}modular -- \hyperlink{page.350}{350}\par
\hspace*{1.1em}resolution of the -- \hyperlink{page.298}{298}\par
\noindent image \hyperlink{page.1}{1}\par
\noindent indecomposable Markov process \hyperlink{page.373}{373}\par
\noindent individual ergodic theorem \hyperlink{page.367}{367}, \hyperlink{page.371}{371}\par
\noindent (strict) inductive limit \hyperlink{page.27}{27}, \hyperlink{page.136}{136}\par
\noindent inequality\par
\hspace*{1.1em}Gärding's -- \hyperlink{page.170}{170}\par
\hspace*{1.1em}Hölder's -- \hyperlink{page.32}{32}\par
\noindent infimum (= greatest lower bound) \hyperlink{page.2}{2}\par
\noindent infinitely divisible law \hyperlink{page.390}{390}\par
\noindent infinitesimal generator \hyperlink{page.227}{227}, \hyperlink{page.232}{232}\par
\hspace*{1.1em}of non--linear evolution equation \hyperlink{page.423}{423}, \hyperlink{page.429}{429}\par
\hspace*{1.1em}resolvent of the -- \hyperlink{page.235}{235}\par
\noindent inner product \hyperlink{page.39}{39}\par
\noindent integrable function \hyperlink{page.15}{15}, \hyperlink{page.16}{16}\par
\hspace*{1.1em}locally -- \hyperlink{page.46}{46}\par
\noindent integral \hyperlink{page.15}{15}\par
\hspace*{1.1em}Bochner -- \hyperlink{page.128}{128}\par
\hspace*{1.1em}Dunford -- \hyperlink{page.221}{221}\par
\hspace*{1.1em}Feynman -- \hyperlink{page.417}{417}\par
\noindent integral \hyperlink{page.15}{15}\par
\hspace*{1.1em}Lebesgue -- \hyperlink{page.17}{17}\par
\hspace*{1.1em}Radon -- \hyperlink{page.115}{115}\par
\hspace*{1.1em}solution of the initial value problem \hyperlink{page.440}{440}\par
\noindent integral operator \hyperlink{page.193}{193}, \hyperlink{page.267}{267}\par
\noindent interior \hyperlink{page.12}{12}\par
\noindent invariant measure with respect to a Markov process \hyperlink{page.365}{365}, \hyperlink{page.375}{375}, \hyperlink{page.378}{378}\par
\hspace*{1.1em}right invariant Haar measure \hyperlink{page.382}{382}\par
\noindent inverse \hyperlink{page.1}{1}, \hyperlink{page.19}{19}, \hyperlink{page.41}{41}, \hyperlink{page.72}{72}\par
\noindent involutive ring \hyperlink{page.291}{291}\par
\noindent isometric \hyperlink{page.110}{110}, \hyperlink{page.197}{197}, \hyperlink{page.198}{198}\par
\noindent isotropy \hyperlink{page.382}{382}\par
\noindent ITô, K. \hyperlink{page.380}{380}, \hyperlink{page.389}{389}, \hyperlink{page.452}{452}\par
\noindent ITô, S. \hyperlink{page.400}{400}, \hyperlink{page.452}{452}\par
\noindent Jacobi matrix \hyperlink{page.334}{334}\par
\noindent \textsc{Jacobs}, K. \hyperlink{page.366}{366}, \hyperlink{page.453}{453}\par
\noindent \textsc{John}, F. \hyperlink{page.453}{453}\par
\noindent \textsc{Jordan}, C. decomposition \hyperlink{page.35}{35}, \hyperlink{page.351}{351}\par
\noindent \textsc{Jordan}, P. \hyperlink{page.37}{37}\par
\noindent \textsc{Kakutani}, S. \hyperlink{page.9}{9}, \hyperlink{page.30}{30}, \hyperlink{page.115}{115}, \hyperlink{page.123}{123}, \hyperlink{page.210}{210}, \hyperlink{page.275}{275}, \hyperlink{page.354}{354}, \hyperlink{page.359}{359}, \hyperlink{page.362}{362}, \hyperlink{page.366}{366}, \hyperlink{page.372}{372}, \hyperlink{page.453}{453}\par
\noindent Riesz-Markov-Theorem \hyperlink{page.115}{115}\par
\noindent \textsc{Kantorovitch}, L. \hyperlink{page.170}{170}, \hyperlink{page.453}{453}\par
\noindent \textsc{Kato}, T. \hyperlink{page.203}{203}, \hyperlink{page.211}{211}, \hyperlink{page.214}{214}, \hyperlink{page.245}{245}, \hyperlink{page.252}{252}, \hyperlink{page.257}{257}, \hyperlink{page.262}{262}, \hyperlink{page.304}{304}, \hyperlink{page.411}{411}, \hyperlink{page.415}{415}, \hyperlink{page.416}{416}, \hyperlink{page.422}{422}, \hyperlink{page.424}{424}, \hyperlink{page.430}{430}, \hyperlink{page.434}{434}, \hyperlink{page.454}{454}\par
\noindent \textsc{Kelley}, J. L. \hyperlink{page.348}{348}, \hyperlink{page.454}{454}\par
\noindent \textsc{Kellog}, O. D. \hyperlink{page.275}{275}, \hyperlink{page.454}{454}\par
\noindent kernel\par
\hspace*{1.1em}Aronszajn--Bergman -- \hyperlink{page.92}{92}\par
\hspace*{1.1em}Hilbert--Schmidt -- \hyperlink{page.193}{193}, \hyperlink{page.267}{267}\par
\hspace*{1.1em}reproducing -- \hyperlink{page.91}{91}\par
\noindent \textsc{Khintchine}, A. \hyperlink{page.371}{371}, \hyperlink{page.454}{454}\par
\noindent \textsc{Kisyński}, J. \hyperlink{page.409}{409}, \hyperlink{page.454}{454}\par
\noindent \textsc{Kodaira}, K. \hyperlink{page.312}{312}, \hyperlink{page.454}{454}\par
\noindent \textsc{Kolmogorov}, A. \hyperlink{page.265}{265}, \hyperlink{page.380}{380}, \hyperlink{page.454}{454}\par
\hspace*{1.1em}Chapman-equation \hyperlink{page.363}{363}\par
\hspace*{1.1em}extension theorem of measures \hyperlink{page.283}{283}\par
\noindent \textsc{Komatsu}, H. \hyperlink{page.264}{264}, \hyperlink{page.422}{422}, \hyperlink{page.454}{454}\par
\noindent \textsc{Komura}, Y. \hyperlink{page.423}{423}, \hyperlink{page.425}{425}, \hyperlink{page.426}{426}, \hyperlink{page.427}{427}, \hyperlink{page.428}{428}, \hyperlink{page.429}{429}, \hyperlink{page.430}{430}, \hyperlink{page.454}{454}\par
\noindent \textsc{Konishi}, Y. \hyperlink{page.440}{440}, \hyperlink{page.455}{455}\par
\noindent \textsc{Kotaké}, T. \hyperlink{page.416}{416}, \hyperlink{page.455}{455}\par
\noindent \textsc{Köthe}, G. \hyperlink{page.64}{64}, \hyperlink{page.102}{102}, \hyperlink{page.455}{455}\par
\noindent \textsc{Krein}, M. \hyperlink{page.9}{9}, \hyperlink{page.359}{359}, \hyperlink{page.455}{455}\par
\hspace*{1.1em}Milman theorem \hyperlink{page.347}{347}\par
\noindent \textsc{Krein}, S. \hyperlink{page.9}{9}, \hyperlink{page.359}{359}, \hyperlink{page.455}{455}\par
\noindent \textsc{Krylov}, N. \hyperlink{page.309}{309}, \hyperlink{page.376}{376}, \hyperlink{page.455}{455}\par
\hspace*{1.1em}--Weinstein theorem \hyperlink{page.309}{309}\par
\noindent \textsc{Ladyzhenskaya}, O. A. \hyperlink{page.411}{411}, \hyperlink{page.455}{455}\par
\noindent Laguerre polynomials \hyperlink{page.86}{86}\par
\noindent Laplace transform \hyperlink{page.158}{158}, \hyperlink{page.232}{232}\par
\noindent Laplacian, generalized -- of Feller \hyperlink{page.385}{385}\par
\noindent lateral condition \hyperlink{page.387}{387}\par
\noindent lattice \hyperlink{page.2}{2}\par
\hspace*{1.1em}B-- -- \hyperlink{page.354}{354}\par
\hspace*{1.1em}distributive -- \hyperlink{page.350}{350}\par
\hspace*{1.1em}F-- -- \hyperlink{page.354}{354}\par
\hspace*{1.1em}homomorphism \hyperlink{page.357}{357}\par
\hspace*{1.1em}modular -- \hyperlink{page.350}{350}\par
\hspace*{1.1em}vector -- \hyperlink{page.349}{349}\par
\noindent \textsc{Lax}, P. D. \hyperlink{page.88}{88}, \hyperlink{page.94}{94}, \hyperlink{page.410}{410}, \hyperlink{page.455}{455}\par
\hspace*{1.1em}--Milgram theorem \hyperlink{page.88}{88}\par
\hspace*{1.1em}negative norm \hyperlink{page.94}{94}, \hyperlink{page.151}{151}\par
\noindent \textsc{Lebesgue}, H. \hyperlink{page.39}{39}\par
\hspace*{1.1em}--Fatou lemma \hyperlink{page.16}{16}\par
\hspace*{1.1em}measure \hyperlink{page.17}{17}\par
\hspace*{1.1em}--Nikodym theorem \hyperlink{page.89}{89}\par
\noindent Legendre polynomial \hyperlink{page.86}{86}\par
\noindent Leibniz' formula \hyperlink{page.49}{49}\par
\noindent \textsc{Leray}, J. \hyperlink{page.96}{96}, \hyperlink{page.455}{455}\par
\noindent Lerch's theorem \hyperlink{page.162}{162}\par
\noindent \textsc{Lewy}, H. \hyperlink{page.176}{176}, \hyperlink{page.455}{455}\par
\noindent lexicographic partial order \hyperlink{page.352}{352}\par
\noindent \textsc{Liggett}, T. \hyperlink{page.430}{430}, \hyperlink{page.434}{434}, \hyperlink{page.440}{440}, \hyperlink{page.455}{455}\par
\noindent limit\par
\hspace*{1.1em}generalized -- \hyperlink{page.101}{101}\par
\hspace*{1.1em}point \hyperlink{page.3}{3}\par
\noindent Lindeberg's condition \hyperlink{page.381}{381}\par
\noindent linear\par
\hspace*{1.1em}continuous -- functional \hyperlink{page.41}{41}\par
\hspace*{1.1em}functional \hyperlink{page.19}{19}\par
\hspace*{1.1em}independence \hyperlink{page.18}{18}\par
\hspace*{1.1em}operator \hyperlink{page.19}{19}\par
\hspace*{1.1em}( = total) ordering \hyperlink{page.2}{2}\par
\hspace*{1.1em}space \hyperlink{page.18}{18}\par
\hspace*{1.1em}subspace (or, in short, subspace) \hyperlink{page.19}{19}\par
\hspace*{1.1em}topological space \hyperlink{page.25}{25}\par
\noindent \textsc{Lions}, J. L. \hyperlink{page.410}{410}, \hyperlink{page.416}{416}, \hyperlink{page.455}{455}\par
\noindent Liouville's theorem \hyperlink{page.125}{125}, \hyperlink{page.364}{364}\par
\noindent Lipschitz condition \hyperlink{page.130}{130}\par
\noindent Lipschitzian mapping \hyperlink{page.424}{424}, \hyperlink{page.426}{426}, \hyperlink{page.427}{427}, \hyperlink{page.437}{437}, \hyperlink{page.438}{438}\par
\noindent \textsc{Littlewood}, J. E. \hyperlink{page.344}{344}\par
\noindent locally\par
\hspace*{1.1em}compact space \hyperlink{page.4}{4}\par
\hspace*{1.1em}convex space \hyperlink{page.26}{26}\par
\noindent locally\par
\hspace*{1.1em}integrable \hyperlink{page.46}{46}\par
\hspace*{1.1em}sequentially weakly compact\par
\hspace*{2.2em}space \hyperlink{page.122}{122}, \hyperlink{page.137}{137}\par
\noindent lower bound \hyperlink{page.2}{2}\par
\noindent L$^{1}$--space, abstract -- \hyperlink{page.354}{354}\par
\noindent \textsc{Lumer}, G. \hyperlink{page.243}{243}, \hyperlink{page.456}{456}\par
\hspace*{1.1em}Phillips--theorem \hyperlink{page.243}{243}\par
\noindent \textsc{Maak}, W. \hyperlink{page.219}{219}, \hyperlink{page.456}{456}\par
\noindent \textsc{Malgrange}, B. \hyperlink{page.177}{177}, \hyperlink{page.186}{186}, \hyperlink{page.456}{456}\par
\hspace*{1.1em}--Ehrenpreis theorem \hyperlink{page.177}{177}\par
\noindent mapping \hyperlink{page.1}{1}\par
\hspace*{1.1em}continuous -- \hyperlink{page.4}{4}\par
\hspace*{1.1em}finitely-valued \hyperlink{page.126}{126}\par
\hspace*{1.1em}homeomorphic -- \hyperlink{page.72}{72}\par
\hspace*{1.1em}isometric -- \hyperlink{page.110}{110}\par
\hspace*{1.1em}m--almost separately valued -- \hyperlink{page.126}{126}\par
\hspace*{1.1em}one--(to--)one -- \hyperlink{page.2}{2}\par
\hspace*{1.1em}separably valued -- \hyperlink{page.126}{126}\par
\hspace*{1.1em}strongly $\exists $--measurable -- \hyperlink{page.126}{126}\par
\hspace*{1.1em}weakly $\exists $--measurable -- \hyperlink{page.126}{126}\par
\noindent \textsc{Markov}, A. \hyperlink{page.115}{115}\par
\noindent Markov process \hyperlink{page.363}{363}\par
\hspace*{1.1em}canonical measure of a -- \hyperlink{page.387}{387}\par
\hspace*{1.1em}dissipative -- \hyperlink{page.378}{378}\par
\hspace*{1.1em}ergodic decomposition of a -- \hyperlink{page.376}{376}\par
\hspace*{1.1em}indecomposable -- \hyperlink{page.373}{373}\par
\hspace*{1.1em}metrically transitive -- \hyperlink{page.373}{373}\par
\hspace*{1.1em}spatially homogeneous -- \hyperlink{page.381}{381}\par
\hspace*{1.1em}temporally homogeneous -- \hyperlink{page.363}{363}\par
\noindent \textsc{Maruyama}, G. \hyperlink{page.380}{380}, \hyperlink{page.456}{456}\par
\noindent maximal \hyperlink{page.2}{2}\par
\hspace*{1.1em}ideal \hyperlink{page.286}{286}\par
\hspace*{1.1em}ideal (in vector lattice) \hyperlink{page.357}{357}\par
\hspace*{1.1em}set of probability \hyperlink{page.376}{376}\par
\hspace*{1.1em}symmetric operator \hyperlink{page.336}{336}\par
\noindent maximally dissipative \hyperlink{page.426}{426}, \hyperlink{page.428}{428}, \hyperlink{page.429}{429}\par
\noindent \textsc{Mazur}, S. \hyperlink{page.105}{105}, \hyperlink{page.106}{106}, \hyperlink{page.117}{117}, \hyperlink{page.456}{456}\par
\hspace*{1.1em}Gelfand--theorem \hyperlink{page.125}{125}\par
\hspace*{1.1em}theorem \hyperlink{page.106}{106}, \hyperlink{page.117}{117}\par
\noindent McKean, H. \hyperlink{page.380}{380}, \hyperlink{page.389}{389}, \hyperlink{page.456}{456}\par
\noindent mean\par
\hspace*{1.1em}convergence of p--th order \hyperlink{page.33}{33}\par
\hspace*{1.1em}ergodic theorem \hyperlink{page.209}{209}, \hyperlink{page.211}{211}, \hyperlink{page.366}{366}, \hyperlink{page.371}{371}\par
\hspace*{1.1em}value of an almost periodic function \hyperlink{page.217}{217}\par
\noindent measure \hyperlink{page.14}{14}\par
\hspace*{1.1em}absolutely continuous -- \hyperlink{page.16}{16}, \hyperlink{page.89}{89}, \hyperlink{page.131}{131}, \hyperlink{page.360}{360}\par
\hspace*{1.1em}Baire -- \hyperlink{page.17}{17}\par
\hspace*{1.1em}Borel -- \hyperlink{page.17}{17}\par
\noindent measure \hyperlink{page.14}{14}\par
\hspace*{1.1em}canonical -- of a Markov process\par
\hspace*{1.1em}complex -- \hyperlink{page.34}{34}\par
\hspace*{1.1em}contracted -- \hyperlink{page.362}{362}\par
\hspace*{1.1em}convergence in -- \hyperlink{page.37}{37}\par
\hspace*{1.1em}invariant -- of a Markov process\par
\noindent \hyperlink{page.364}{364}, \hyperlink{page.378}{378}\par
\hspace*{1.1em}Lebesgue -- \hyperlink{page.17}{17}\par
\hspace*{1.1em}product -- \hyperlink{page.16}{16}\par
\hspace*{1.1em}regular Borel -- \hyperlink{page.17}{17}, \hyperlink{page.376}{376}\par
\hspace*{1.1em}signed -- \hyperlink{page.34}{34}\par
\hspace*{1.1em}singular -- \hyperlink{page.362}{362}\par
\hspace*{1.1em}slowly increasing -- \hyperlink{page.147}{147}\par
\hspace*{1.1em}space \hyperlink{page.14}{14}\par
\hspace*{1.1em}topological -- \hyperlink{page.17}{17}\par
\noindent measurable\par
\hspace*{1.1em}function \hyperlink{page.14}{14}\par
\hspace*{1.1em}set \hyperlink{page.14}{14}, \hyperlink{page.18}{18}\par
\hspace*{1.1em}strongly -- \hyperlink{page.126}{126}\par
\hspace*{1.1em}weakly -- \hyperlink{page.126}{126}\par
\noindent metric \hyperlink{page.3}{3}\par
\hspace*{1.1em}space \hyperlink{page.3}{3}\par
\noindent metrically transitive Markov process \hyperlink{page.373}{373}\par
\noindent \textsc{Meyer}, P. A. \hyperlink{page.348}{348}, \hyperlink{page.397}{397}, \hyperlink{page.456}{456}\par
\noindent \textsc{Mikusiński}, J. \hyperlink{page.162}{162}, \hyperlink{page.164}{164}, \hyperlink{page.456}{456}\par
\noindent \textsc{Milgram}, A. N. \hyperlink{page.88}{88}, \hyperlink{page.456}{456}\par
\hspace*{1.1em}Lax-- -- theorem \hyperlink{page.88}{88}\par
\noindent \textsc{Milman}, D. \hyperlink{page.123}{123}, \hyperlink{page.347}{347}, \hyperlink{page.456}{456}\par
\hspace*{1.1em}Krein-- -- theorem \hyperlink{page.347}{347}\par
\noindent \textsc{Mimura}, Y. \hyperlink{page.115}{115}, \hyperlink{page.275}{275}, \hyperlink{page.326}{326}, \hyperlink{page.331}{331}, \hyperlink{page.456}{456}\par
\noindent \textsc{Minkowski}, H. \hyperlink{page.24}{24}, \hyperlink{page.32}{32}\par
\hspace*{1.1em}functional \hyperlink{page.24}{24}, \hyperlink{page.25}{25}\par
\hspace*{1.1em}inequality \hyperlink{page.32}{32}\par
\noindent \textsc{Minlos}, R. A. \hyperlink{page.283}{283}, \hyperlink{page.456}{456}\par
\noindent \textsc{Minty}, G. \hyperlink{page.424}{424}, \hyperlink{page.457}{457}\par
\noindent mixing hypothesis \hyperlink{page.375}{375}\par
\noindent \textsc{Miyadera}, I. \hyperlink{page.242}{242}, \hyperlink{page.430}{430}, \hyperlink{page.438}{438}, \hyperlink{page.457}{457}\par
\noindent \textsc{Mizohata}, S. \hyperlink{page.400}{400}, \hyperlink{page.457}{457}\par
\noindent modular\par
\hspace*{1.1em}identity \hyperlink{page.350}{350}\par
\hspace*{1.1em}lattice \hyperlink{page.350}{350}\par
\noindent moment problem \hyperlink{page.103}{103}\par
\noindent momentum operator \hyperlink{page.194}{194}, \hyperlink{page.304}{304}, \hyperlink{page.336}{336}\par
\noindent Moore-Smith limit \hyperlink{page.101}{101}\par
\noindent \textsc{Morrey}, C. B. \hyperlink{page.187}{187}, \hyperlink{page.457}{457}\par
\noindent multiplicity\par
\hspace*{1.1em}of an eigenvalue \hyperlink{page.205}{205}\par
\hspace*{1.1em}of the spectrum \hyperlink{page.309}{309}\par
\noindent multi-valued mapping \hyperlink{page.424}{424}\par
\noindent \textsc{Nagumo}, M. \hyperlink{page.226}{226}, \hyperlink{page.286}{286}, \hyperlink{page.457}{457}\par
\noindent \textsc{Nagy}, B. \textsc{von} Sz. \hyperlink{page.331}{331}, \hyperlink{page.339}{339}, \hyperlink{page.457}{457}\par
\noindent \textsc{Naimark}, M. A. \hyperlink{page.286}{286}, \hyperlink{page.313}{313}, \hyperlink{page.337}{337}, \hyperlink{page.344}{344}, \hyperlink{page.457}{457}\par
\noindent \textsc{Nakano}, H. \hyperlink{page.311}{311}, \hyperlink{page.458}{458}\par
\noindent \textsc{Narasinham}, M. S. \hyperlink{page.416}{416}, \hyperlink{page.458}{458}\par
\noindent nearly orthogonal element \hyperlink{page.81}{81}\par
\noindent negative\par
\hspace*{1.1em}norm \hyperlink{page.94}{94}, \hyperlink{page.152}{152}\par
\hspace*{1.1em}variation \hyperlink{page.34}{34}, \hyperlink{page.351}{351}\par
\noindent neighbourhood \hyperlink{page.3}{3}\par
\noindent \textsc{Nelson}, E. \hyperlink{page.417}{417}, \hyperlink{page.458}{458}\par
\noindent \textsc{Neumann}, C. series \hyperlink{page.66}{66}\par
\noindent \textsc{Neumann}, J. \textsc{von} \hyperlink{page.37}{37}, \hyperlink{page.88}{88}, \hyperlink{page.90}{90}, \hyperlink{page.195}{195}, \hyperlink{page.198}{198}, \hyperlink{page.199}{199}, \hyperlink{page.210}{210}, \hyperlink{page.218}{218}, \hyperlink{page.286}{286}, \hyperlink{page.304}{304}, \hyperlink{page.319}{319}, \hyperlink{page.326}{326}, \hyperlink{page.331}{331}, \hyperlink{page.336}{336}, \hyperlink{page.458}{458}\par
\hspace*{1.1em}--Riesz--Mimura theorem \hyperlink{page.326}{326}\par
\hspace*{1.1em}theorem on\par
\hspace*{2.2em}Cayley transform \hyperlink{page.197}{197}\par
\hspace*{2.2em}real operators \hyperlink{page.305}{305}\par
\hspace*{2.2em}spectral resolution \hyperlink{page.303}{303}\par
\noindent \textsc{Nikodym}, O. \hyperlink{page.89}{89}\par
\hspace*{1.1em}Radon- - theorem \hyperlink{page.362}{362}\par
\noindent nilpotent \hyperlink{page.356}{356}\par
\hspace*{1.1em}generalized -- \hyperlink{page.289}{289}, \hyperlink{page.292}{292}\par
\noindent \textsc{Nirenberg}, L. \hyperlink{page.172}{172}, \hyperlink{page.187}{187}, \hyperlink{page.417}{417}, \hyperlink{page.458}{458}\par
\noindent non-dense \hyperlink{page.11}{11}\par
\noindent non-negative functional \hyperlink{page.376}{376}\par
\noindent norm \hyperlink{page.29}{29}, \hyperlink{page.41}{41}\par
\hspace*{1.1em}of a bounded linear operator \hyperlink{page.42}{42}\par
\hspace*{1.1em}negative -- \hyperlink{page.94}{94}, \hyperlink{page.152}{152}\par
\hspace*{1.1em}quasi-- -- \hyperlink{page.23}{23}, \hyperlink{page.30}{30}\par
\hspace*{1.1em}semi-- -- \hyperlink{page.23}{23}\par
\noindent normal\par
\hspace*{1.1em}operator \hyperlink{page.197}{197}\par
\hspace*{2.2em}spectral resolution of a -- -- \hyperlink{page.296}{296}\par
\hspace*{1.1em}space \hyperlink{page.6}{6}\par
\noindent normed\par
\hspace*{1.1em}factor space \hyperlink{page.57}{57}\par
\hspace*{1.1em}field \hyperlink{page.125}{125}\par
\hspace*{1.1em}linear space \hyperlink{page.29}{29}\par
\hspace*{1.1em}ring \hyperlink{page.286}{286}\par
\hspace*{2.2em}semi--simple -- -- \hyperlink{page.289}{289}\par
\noindent nuclear\par
\hspace*{1.1em}operator \hyperlink{page.269}{269}, \hyperlink{page.279}{279}\par
\hspace*{1.1em}space \hyperlink{page.281}{281}\par
\noindent null space \hyperlink{page.19}{19}\par
\noindent $O$-convergence \hyperlink{page.351}{351}\par
\noindent \textsc{Oharu}, S. \hyperlink{page.438}{438}, \hyperlink{page.458}{458}\par
\noindent \textsc{Okamoto}, S. \hyperlink{page.164}{164}\par
\noindent one point compactification \hyperlink{page.4}{4}\par
\noindent one-to-one mapping \hyperlink{page.1}{1}\par
\noindent open\par
\hspace*{1.1em}base \hyperlink{page.57}{57}\par
\hspace*{1.1em}mapping theorem \hyperlink{page.71}{71}\par
\hspace*{1.1em}set \hyperlink{page.3}{3}\par
\hspace*{1.1em}sphere \hyperlink{page.3}{3}\par
\noindent operational calculus\par
\hspace*{1.1em}Dunford's -- \hyperlink{page.221}{221}\par
\hspace*{1.1em}on functions of a self-adjoint operator \hyperlink{page.329}{329}\par
\hspace*{1.1em}Heaviside's -- \hyperlink{page.164}{164}\par
\hspace*{1.1em}Mikusiński's -- \hyperlink{page.164}{164}\par
\noindent operator \hyperlink{page.19}{19}\par
\hspace*{1.1em}accretive -- \hyperlink{page.203}{203}\par
\hspace*{1.1em}adjoint -- \hyperlink{page.191}{191}\par
\hspace*{1.1em}adjoint differential -- \hyperlink{page.76}{76}\par
\hspace*{1.1em}bounded linear -- \hyperlink{page.42}{42}\par
\hspace*{1.1em}closable -- \hyperlink{page.73}{73}\par
\hspace*{1.1em}closed linear -- \hyperlink{page.73}{73}\par
\hspace*{1.1em}generalized spectral resolution of a closed symmetric -- \hyperlink{page.338}{338}\par
\hspace*{1.1em}compact -- \hyperlink{page.265}{265}, \hyperlink{page.267}{267}, \hyperlink{page.313}{313}\par
\hspace*{1.1em}comparison of -s (Hörmander's theorem) \hyperlink{page.75}{75}\par
\hspace*{1.1em}completely continuous -- \hyperlink{page.265}{265}, \hyperlink{page.267}{267}\par
\hspace*{1.1em}conjugate -- \hyperlink{page.189}{189}\par
\hspace*{1.1em}continuous linear -- \hyperlink{page.41}{41}\par
\hspace*{1.1em}coordinate -- \hyperlink{page.193}{193}, \hyperlink{page.336}{336}\par
\hspace*{1.1em}$ \delta $ -- \hyperlink{page.165}{165}\par
\hspace*{1.1em}dissipative -- \hyperlink{page.203}{203}, \hyperlink{page.244}{244}\par
\hspace*{1.1em}dual -- \hyperlink{page.189}{189}\par
\hspace*{2.2em}inverse of the -- \hyperlink{page.220}{220}\par
\hspace*{2.2em}resolvent of the -- \hyperlink{page.220}{220}\par
\noindent operator\par
\hspace*{1.1em}essentially self--adjoint -- \hyperlink{page.304}{304}, \hyperlink{page.336}{336}\par
\hspace*{1.1em}exponential of a bounded -- \hyperlink{page.238}{238}\par
\hspace*{1.1em}extension of an -- \hyperlink{page.1}{1}, \hyperlink{page.19}{19}\par
\hspace*{1.1em}fractional powers of a closed -- \hyperlink{page.252}{252}\par
\hspace*{1.1em}Hilbert--Schmidt integral -- \hyperlink{page.193}{193}, \hyperlink{page.267}{267}\par
\hspace*{1.1em}hypermaximal (= self--adjoint) -- \hyperlink{page.336}{336}\par
\hspace*{1.1em}idempotent -- \hyperlink{page.80}{80}\par
\hspace*{1.1em}identity -- \hyperlink{page.42}{42}\par
\hspace*{1.1em}integral -- \hyperlink{page.193}{193}, \hyperlink{page.267}{267}\par
\hspace*{1.1em}inverse -- \hyperlink{page.19}{19}, \hyperlink{page.72}{72}\par
\hspace*{1.1em}isometric -- \hyperlink{page.197}{197}, \hyperlink{page.198}{198}\par
\hspace*{1.1em}linear -- \hyperlink{page.19}{19}, \hyperlink{page.41}{41}\par
\hspace*{1.1em}maximal symmetric -- \hyperlink{page.336}{336}\par
\noindent operator\par
\hspace*{1.1em}momentum -- \hyperlink{page.194}{194}, \hyperlink{page.304}{304}, \hyperlink{page.336}{336}\par
\hspace*{1.1em}norm of a bounded -- \hyperlink{page.42}{42}\par
\hspace*{1.1em}normal -- \hyperlink{page.197}{197}\par
\hspace*{2.2em}spectral resolution of a - -\par
\hspace*{1.1em}nuclear -- \hyperlink{page.269}{269}, \hyperlink{page.279}{279}\par
\hspace*{1.1em}positive -- \hyperlink{page.305}{305}\par
\hspace*{1.1em}pre--closed -- \hyperlink{page.73}{73}\par
\hspace*{1.1em}projection -- \hyperlink{page.79}{79}\par
\hspace*{1.1em}real -- \hyperlink{page.305}{305}\par
\hspace*{1.1em}self--adjoint -- \hyperlink{page.80}{80}, \hyperlink{page.193}{193}, \hyperlink{page.329}{329}, \hyperlink{page.336}{336}\par
\hspace*{2.2em}functions of a -- -- \hyperlink{page.325}{325}\par
\hspace*{2.2em}spectral resolution of a -- -- \hyperlink{page.303}{303}\par
\hspace*{2.2em}spectrum of a -- -- \hyperlink{page.307}{307}\par
\hspace*{1.1em}semi--bounded -- \hyperlink{page.306}{306}\par
\hspace*{1.1em}symmetric -- \hyperlink{page.193}{193}, \hyperlink{page.334}{334}, \hyperlink{page.338}{338}\par
\hspace*{1.1em}of the trace class \hyperlink{page.271}{271}\par
\hspace*{1.1em}unitary -- \hyperlink{page.197}{197}\par
\hspace*{2.2em}spectral resolution of a -- -- \hyperlink{page.296}{296}\par
\noindent ordered set\par
\hspace*{1.1em}linearily -- \hyperlink{page.2}{2}\par
\hspace*{1.1em}partially -- \hyperlink{page.2}{2}\par
\hspace*{1.1em}semi-- -- \hyperlink{page.2}{2}\par
\hspace*{1.1em}totally -- \hyperlink{page.2}{2}\par
\noindent \textsc{Ornstein}, D. S. \hyperlink{page.369}{369}, \hyperlink{page.372}{372}, \hyperlink{page.458}{458}\par
\noindent orthogonal \hyperlink{page.79}{79}\par
\hspace*{1.1em}base \hyperlink{page.83}{83}\par
\hspace*{1.1em}complement \hyperlink{page.79}{79}\par
\hspace*{1.1em}nearly -- \hyperlink{page.81}{81}\par
\hspace*{1.1em}projection \hyperlink{page.79}{79}\par
\noindent orthogonalization (E. \textsc{Schmidt}) \hyperlink{page.85}{85}\par
\noindent orthonormal system \hyperlink{page.83}{83}\par
\noindent \textsc{Paley}, E. R. A. C. \hyperlink{page.157}{157}, \hyperlink{page.458}{458}\par
\hspace*{1.1em}--Wiener theorem \hyperlink{page.157}{157}, \hyperlink{page.158}{158}, \hyperlink{page.159}{159}\par
\noindent parabolic differential equation \hyperlink{page.182}{182}\par
\noindent \textsc{Parseval}(-\textsc{Deschènes}, M. A.) relation \hyperlink{page.83}{83}, \hyperlink{page.145}{145}\par
\hspace*{1.1em}theorem for the Fourier transform \hyperlink{page.151}{151}\par
\noindent partial order \hyperlink{page.2}{2}\par
\noindent partition of unity \hyperlink{page.57}{57}\par
\noindent \textsc{Pazy}, A. \hyperlink{page.430}{430}, \hyperlink{page.433}{433}, \hyperlink{page.434}{434}, \hyperlink{page.438}{438}, \hyperlink{page.459}{459}\par
\noindent \textsc{Peetre}, J. \hyperlink{page.186}{186}, \hyperlink{page.459}{459}\par
\noindent \textsc{Peter}, F. \hyperlink{page.317}{317}, \hyperlink{page.459}{459}\par
\hspace*{1.1em}--Weyl--Neumann theory \hyperlink{page.317}{317}\par
\noindent \textsc{Petrowsky}, I. G. \hyperlink{page.187}{187}, \hyperlink{page.459}{459}\par
\noindent \textsc{Pettis}, B. J. \hyperlink{page.126}{126}, \hyperlink{page.459}{459}\par
\noindent phase space \hyperlink{page.364}{364}\par
\noindent \textsc{Phillips}, R. S. \hyperlink{page.211}{211}, \hyperlink{page.220}{220}, \hyperlink{page.226}{226}, \hyperlink{page.242}{242}, \hyperlink{page.243}{243}, \hyperlink{page.252}{252}, \hyperlink{page.262}{262}, \hyperlink{page.286}{286}, \hyperlink{page.390}{390}, \hyperlink{page.405}{405}, \hyperlink{page.410}{410}, \hyperlink{page.459}{459}\par
\hspace*{1.1em}--Lumer theorem \hyperlink{page.243}{243}\par
\hspace*{1.1em}theorem \hyperlink{page.264}{264}\par
\noindent \textsc{Pitt}, H. R. \hyperlink{page.344}{344}, \hyperlink{page.459}{459}\par
\noindent Plancherel's theorem for the Fourier transform \hyperlink{page.150}{150}\par
\noindent \textsc{Poincare}, H. \hyperlink{page.223}{223}\par
\noindent point spectrum \hyperlink{page.205}{205}\par
\noindent \textsc{Poisson}, S. D.\par
\hspace*{1.1em}equation \hyperlink{page.391}{391}\par
\hspace*{1.1em}integral representation \hyperlink{page.160}{160}\par
\hspace*{1.1em}kernel \hyperlink{page.259}{259}\par
\hspace*{1.1em}summation formula \hyperlink{page.146}{146}\par
\noindent polar set \hyperlink{page.133}{133}\par
\noindent polynomial\par
\hspace*{1.1em}approximation theorem \hyperlink{page.7}{7}, \hyperlink{page.243}{243}\par
\hspace*{1.1em}Hermite -- \hyperlink{page.86}{86}\par
\hspace*{1.1em}Laguerre -- \hyperlink{page.86}{86}\par
\hspace*{1.1em}Legendre -- \hyperlink{page.86}{86}\par
\hspace*{1.1em}Tchebyshev -- \hyperlink{page.86}{86}\par
\noindent \textsc{Pontrjagin}, L. \hyperlink{page.324}{324}, \hyperlink{page.459}{459}\par
\noindent positive\par
\hspace*{1.1em}definite function \hyperlink{page.332}{332}\par
\hspace*{1.1em}functional \hyperlink{page.88}{88}\par
\hspace*{1.1em}operator \hyperlink{page.306}{306}\par
\hspace*{1.1em}variation \hyperlink{page.34}{34}, \hyperlink{page.351}{351}\par
\noindent \textsc{Post}, E. L.-Widder inversion formula \hyperlink{page.253}{253}\par
\noindent potential \hyperlink{page.391}{391}\par
\noindent \textsc{Poulsen}, E. T. \hyperlink{page.416}{416}, \hyperlink{page.459}{459}\par
\noindent pre-closed operator \hyperlink{page.73}{73}\par
\noindent pre-Hilbert space \hyperlink{page.37}{37}\par
\noindent principle of the condensation of singularities \hyperlink{page.69}{69}\par
\noindent projection \hyperlink{page.79}{79}\par
\noindent product\par
\hspace*{1.1em}direct (= tensor) -- of distributions \hyperlink{page.63}{63}\par
\hspace*{1.1em}inner (= scalar) -- \hyperlink{page.39}{39}\par
\hspace*{1.1em}measure \hyperlink{page.16}{16}\par
\noindent pseudo-resolvent \hyperlink{page.211}{211}\par
\noindent quasi-norm \hyperlink{page.23}{23}, \hyperlink{page.30}{30}\par
\noindent quasi-normed linear space \hyperlink{page.30}{30}\par
\hspace*{1.1em}completion of a -- \hyperlink{page.53}{53}\par
\noindent quasi-unit \hyperlink{page.360}{360}\par
\noindent \textsc{Quinn}, B. K. \hyperlink{page.440}{440}, \hyperlink{page.459}{459}\par
\noindent radical of a\par
\hspace*{1.1em}normed ring \hyperlink{page.289}{289}\par
\hspace*{1.1em}vector lattice \hyperlink{page.356}{356}\par
\noindent \textsc{Rado}, T. \hyperlink{page.391}{391}, \hyperlink{page.459}{459}\par
\noindent \textsc{Radon}, J.\par
\hspace*{1.1em}integral \hyperlink{page.115}{115}\par
\hspace*{1.1em}--Nikodým theorem \hyperlink{page.362}{362}\par
\noindent \textsc{Raikov}, D. A. \hyperlink{page.286}{286}, \hyperlink{page.341}{341}, \hyperlink{page.460}{460}\par
\noindent range 1, \hyperlink{page.1}{1}, \hyperlink{page.19}{19}\par
\hspace*{1.1em}closed -- theorem \hyperlink{page.200}{200}\par
\noindent rapidly decreasing function \hyperlink{page.143}{143}\par
\noindent \textsc{Rasmussen}, S. \hyperlink{page.430}{430}, \hyperlink{page.460}{460}\par
\noindent \textsc{Ray}, D. \hyperlink{page.380}{380}, \hyperlink{page.460}{460}\par
\noindent Rayleigh's principle \hyperlink{page.309}{309}\par
\noindent real operator \hyperlink{page.305}{305}\par
\noindent reduction theory \hyperlink{page.334}{334}\par
\noindent reflexive space \hyperlink{page.88}{88}, \hyperlink{page.110}{110}, \hyperlink{page.123}{123}, \hyperlink{page.136}{136}\par
\hspace*{1.1em}semi-- -- \hyperlink{page.136}{136}\par
\noindent regular space \hyperlink{page.6}{6}\par
\noindent regularization \hyperlink{page.29}{29}\par
\noindent relative topology \hyperlink{page.3}{3}\par
\noindent Rellich-Gärding theorem \hyperlink{page.271}{271}\par
\noindent representation\par
\hspace*{1.1em}Riesz' -- theorem \hyperlink{page.86}{86}\par
\hspace*{1.1em}of a semi--group \hyperlink{page.240}{240}\par
\hspace*{1.1em}spectral -- for a\par
\hspace*{2.2em}normal operator \hyperlink{page.296}{296}\par
\hspace*{2.2em}self--adjoint operator \hyperlink{page.303}{303}\par
\hspace*{2.2em}unitary operator \hyperlink{page.296}{296}\par
\hspace*{1.1em}of a vector lattice as\par
\hspace*{2.2em}point functions \hyperlink{page.356}{356}\par
\hspace*{2.2em}set functions \hyperlink{page.359}{359}\par
\noindent reproducing kernel \hyperlink{page.91}{91}\par
\noindent residual spectrum \hyperlink{page.205}{205}\par
\noindent residue class \hyperlink{page.215}{215}, \hyperlink{page.288}{288}\par
\noindent resolution of the identity \hyperlink{page.298}{298}\par
\noindent resolvent \hyperlink{page.205}{205}\par
\hspace*{1.1em}equation \hyperlink{page.207}{207}\par
\hspace*{1.1em}isolated singularities of a -- \hyperlink{page.223}{223}\par
\hspace*{1.1em}Laurent expansion of a -- \hyperlink{page.208}{208}\par
\hspace*{1.1em}pseudo-- -- \hyperlink{page.211}{211}\par
\hspace*{1.1em}set \hyperlink{page.205}{205}\par
\noindent resonance theorem \hyperlink{page.65}{65}\par
\noindent restriction \hyperlink{page.1}{1}, \hyperlink{page.19}{19}\par
\noindent \textsc{Rickart}, C. E. \hyperlink{page.286}{286}, \hyperlink{page.344}{344}, \hyperlink{page.460}{460}\par
\noindent \textsc{Riemann}, B.\par
\hspace*{1.1em}--Lebesgue theorem \hyperlink{page.342}{342}\par
\hspace*{1.1em}sum \hyperlink{page.101}{101}\par
\noindent \textsc{Riesz}, F. \hyperlink{page.81}{81}, \hyperlink{page.86}{86}, \hyperlink{page.88}{88}, \hyperlink{page.115}{115}, \hyperlink{page.210}{210}, \hyperlink{page.265}{265}, \hyperlink{page.273}{273}, \hyperlink{page.326}{326}, \hyperlink{page.331}{331}, \hyperlink{page.348}{348}, \hyperlink{page.362}{362}, \hyperlink{page.460}{460}\par
\hspace*{1.1em}--Markov--Kakutani theorem \hyperlink{page.115}{115}\par
\hspace*{1.1em}representation theorem \hyperlink{page.86}{86}\par
\hspace*{1.1em}--Schauder theory \hyperlink{page.272}{272}\par
\noindent ring \hyperlink{page.285}{285}\par
\hspace*{1.1em}involutive -- \hyperlink{page.291}{291}\par
\noindent ring \hyperlink{page.285}{285}\par
\hspace*{1.1em}normed -- \hyperlink{page.286}{286}\par
\hspace*{2.2em}semi--simple -- -- \hyperlink{page.289}{289}\par
\hspace*{1.1em}symmetric -- \hyperlink{page.291}{291}\par
\noindent \textsc{Ryll-nardzewski}, C. \hyperlink{page.162}{162}, \hyperlink{page.460}{460}\par
\noindent $\sigma $-additive \hyperlink{page.14}{14}\par
\noindent \textsc{Saks}, S. \hyperlink{page.66}{66}, \hyperlink{page.91}{91}, \hyperlink{page.460}{460}\par
\noindent \textsc{Sato}, M. \hyperlink{page.161}{161}, \hyperlink{page.460}{460}\par
\noindent scalar product \hyperlink{page.39}{39}\par
\noindent \textsc{Schatten}, R. \hyperlink{page.271}{271}, \hyperlink{page.460}{460}\par
\noindent \textsc{Schauder}, J. \hyperlink{page.265}{265}, \hyperlink{page.272}{272}, \hyperlink{page.460}{460}\par
\hspace*{1.1em}Riesz theory \hyperlink{page.273}{273}\par
\hspace*{1.1em}theorem \hyperlink{page.272}{272}\par
\noindent \textsc{Schmidt}, E. orthogonalization \hyperlink{page.85}{85}\par
\noindent \textsc{Schur}, I. \hyperlink{page.119}{119}, \hyperlink{page.321}{321}\par
\hspace*{1.1em}lemma \hyperlink{page.321}{321}\par
\noindent \textsc{Schwartz}, J. \hyperlink{page.313}{313}, \hyperlink{page.372}{372}, \hyperlink{page.461}{461}\par
\noindent \textsc{Schwartz}, L. \hyperlink{page.23}{23}, \hyperlink{page.44}{44}, \hyperlink{page.96}{96}, \hyperlink{page.154}{154}, \hyperlink{page.161}{161}, \hyperlink{page.172}{172}, \hyperlink{page.242}{242}, \hyperlink{page.461}{461}\par
\noindent $\sigma $-complete \hyperlink{page.352}{352}\par
\noindent \textsc{Segal}, I. E. \hyperlink{page.344}{344}, \hyperlink{page.461}{461}\par
\noindent self-adjoint operator \hyperlink{page.193}{193}, \hyperlink{page.329}{329}, \hyperlink{page.336}{336}\par
\hspace*{1.1em}essentially -- \hyperlink{page.304}{304}, \hyperlink{page.336}{336}\par
\hspace*{1.1em}functions of a -- \hyperlink{page.325}{325}\par
\hspace*{1.1em}spectral resolution of a -- \hyperlink{page.303}{303}\par
\hspace*{1.1em}spectrum of a -- \hyperlink{page.307}{307}\par
\noindent semi-bounded operator \hyperlink{page.306}{306}\par
\noindent semi-group (of class $ (C_{0}) $) \hyperlink{page.228}{228}\par
\hspace*{1.1em}contraction -- \hyperlink{page.243}{243}\par
\hspace*{1.1em}convergence of a sequence of -- \hyperlink{page.260}{260}\par
\hspace*{1.1em}differentiability of a -- \hyperlink{page.232}{232}\par
\hspace*{1.1em}dual -- \hyperlink{page.262}{262}\par
\hspace*{1.1em}equi--bounded -- \hyperlink{page.228}{228}\par
\hspace*{1.1em}equi--continuous -- \hyperlink{page.230}{230}\par
\hspace*{1.1em}holomorphic -- \hyperlink{page.247}{247}\par
\hspace*{1.1em}infinitesimal generator of a -- \hyperlink{page.232}{232}\par
\hspace*{1.1em}representation of a -- \hyperlink{page.240}{240}\par
\noindent semi-norm \hyperlink{page.23}{23}\par
\noindent semi-reflexive space \hyperlink{page.136}{136}\par
\noindent semi-scalar product \hyperlink{page.244}{244}\par
\noindent semi-simple normed ring \hyperlink{page.289}{289}\par
\noindent separably valued mapping \hyperlink{page.126}{126}\par
\noindent separating points \hyperlink{page.9}{9}\par
\noindent separation, axiom of -- \hyperlink{page.2}{2}\par
\noindent sesqui-linear functional \hyperlink{page.88}{88}\par
\noindent set \hyperlink{page.1}{1}\par
\hspace*{1.1em}Baire -- \hyperlink{page.17}{17}\par
\hspace*{1.1em}Borel -- \hyperlink{page.17}{17}\par
\hspace*{1.1em}closed -- \hyperlink{page.3}{3}\par
\noindent set \hyperlink{page.1}{1}\par
\hspace*{1.1em}convex -- \hyperlink{page.24}{24}\par
\hspace*{1.1em}$ G_{\delta^{*}} -- $ \hyperlink{page.17}{17}\par
\hspace*{1.1em}open -- \hyperlink{page.3}{3}\par
\hspace*{1.1em}ordered -- \hyperlink{page.2}{2}\par
\noindent $ \sigma $-finite \hyperlink{page.14}{14}\par
\noindent \textsc{Shmulyan}, V. L. \hyperlink{page.137}{137}, \hyperlink{page.461}{461}\par
\noindent signed measure \hyperlink{page.34}{34}\par
\noindent \textsc{Šilov}, G. \hyperlink{page.161}{161}, \hyperlink{page.286}{286}, \hyperlink{page.461}{461}\par
\noindent simple convergence topology \hyperlink{page.107}{107}\par
\noindent simple spectrum \hyperlink{page.310}{310}, \hyperlink{page.333}{333}\par
\noindent singular\par
\hspace*{1.1em}element of a vector lattice \hyperlink{page.359}{359}\par
\hspace*{1.1em}measure \hyperlink{page.362}{362}\par
\noindent slowly increasing\par
\hspace*{1.1em}function \hyperlink{page.147}{147}\par
\hspace*{1.1em}measure \hyperlink{page.147}{147}\par
\noindent \textsc{Smith}, P. A. \hyperlink{page.373}{373}, \hyperlink{page.461}{461}\par
\noindent \textsc{Sobczyk}, A. \hyperlink{page.102}{102}\par
\noindent \textsc{Sobolev}, S. \hyperlink{page.23}{23}, \hyperlink{page.52}{52}, \hyperlink{page.151}{151}, \hyperlink{page.168}{168}, \hyperlink{page.461}{461}\par
\hspace*{1.1em}lemma \hyperlink{page.168}{168}\par
\hspace*{1.1em}space \hyperlink{page.52}{52}\par
\noindent \textsc{Sobolevski}, P. E. \hyperlink{page.411}{411}, \hyperlink{page.416}{416}, \hyperlink{page.461}{461}\par
\noindent solution of a partial differential equation\par
\hspace*{1.1em}fundamental -- \hyperlink{page.177}{177}\par
\hspace*{1.1em}genuine -- \hyperlink{page.172}{172}\par
\hspace*{1.1em}weak -- \hyperlink{page.172}{172}\par
\noindent space\par
\hspace*{1.1em}Banach (or B--) -- \hyperlink{page.23}{23}, \hyperlink{page.50}{50}\par
\hspace*{1.1em}barrel -- \hyperlink{page.134}{134}\par
\hspace*{1.1em}bidual -- \hyperlink{page.88}{88}, \hyperlink{page.110}{110}, \hyperlink{page.136}{136}\par
\hspace*{1.1em}bornologic -- \hyperlink{page.43}{43}\par
\hspace*{1.1em}complete -- \hyperlink{page.10}{10}, \hyperlink{page.50}{50}, \hyperlink{page.102}{102}\par
\hspace*{1.1em}discrete topological -- \hyperlink{page.33}{33}\par
\hspace*{1.1em}dual -- \hyperlink{page.46}{46}, \hyperlink{page.88}{88}, \hyperlink{page.108}{108}\par
\hspace*{2.2em}strong -- -- \hyperlink{page.108}{108}\par
\hspace*{2.2em}weak* -- -- \hyperlink{page.108}{108}\par
\hspace*{1.1em}eigen -- \hyperlink{page.205}{205}\par
\hspace*{1.1em}factor -- \hyperlink{page.20}{20}, \hyperlink{page.56}{56}\par
\hspace*{1.1em}Fréchet (or F--) -- \hyperlink{page.50}{50}\par
\hspace*{1.1em}Hilbert -- \hyperlink{page.50}{50}, \hyperlink{page.77}{77}, \hyperlink{page.83}{83}\par
\noindent space\par
\hspace*{1.1em}homogeneous -- \hyperlink{page.383}{383}\par
\hspace*{1.1em}linear -- \hyperlink{page.18}{18}\par
\hspace*{1.1em}linear topological -- \hyperlink{page.25}{25}\par
\hspace*{1.1em}locally compact -- \hyperlink{page.4}{4}\par
\hspace*{1.1em}locally convex -- \hyperlink{page.26}{26}\par
\hspace*{1.1em}locally sequentially weakly compact -- \hyperlink{page.122}{122}, \hyperlink{page.138}{138}\par
\hspace*{1.1em}measure -- \hyperlink{page.14}{14}\par
\hspace*{1.1em}metric -- \hyperlink{page.3}{3}\par
\noindent space\par
\hspace*{1.1em}normal -- \hyperlink{page.7}{7}\par
\hspace*{1.1em}normed linear -- \hyperlink{page.29}{29}\par
\hspace*{1.1em}nuclear -- \hyperlink{page.281}{281}\par
\hspace*{1.1em}null -- \hyperlink{page.19}{19}\par
\hspace*{1.1em}phase -- \hyperlink{page.364}{364}\par
\hspace*{1.1em}pre--Hilbert -- \hyperlink{page.37}{37}\par
\hspace*{1.1em}product -- \hyperlink{page.5}{5}\par
\hspace*{1.1em}quasi--normed linear -- \hyperlink{page.30}{30}\par
\hspace*{1.1em}reflexive -- \hyperlink{page.88}{88}, \hyperlink{page.110}{110}, \hyperlink{page.123}{123}, \hyperlink{page.136}{136}\par
\hspace*{1.1em}regular -- \hyperlink{page.6}{6}\par
\hspace*{1.1em}semi--reflexive -- \hyperlink{page.136}{136}\par
\hspace*{1.1em}sequentially complete -- \hyperlink{page.102}{102}\par
\hspace*{1.1em}Sobolew -- \hyperlink{page.52}{52}\par
\hspace*{1.1em}topological -- \hyperlink{page.2}{2}\par
\hspace*{1.1em}uniformly convex -- \hyperlink{page.123}{123}\par
\noindent space-like unique continuation property \hyperlink{page.400}{400}\par
\noindent spatially homogeneous Markov process \hyperlink{page.381}{381}\par
\noindent spectral\par
\hspace*{1.1em}mapping theorem \hyperlink{page.222}{222}\par
\hspace*{1.1em}radius \hyperlink{page.207}{207}\par
\hspace*{1.1em}resolution for a\par
\hspace*{2.2em}closed symmetric operator \hyperlink{page.338}{338}\par
\hspace*{2.2em}normal operator \hyperlink{page.296}{296}\par
\hspace*{2.2em}self--adjoint operator \hyperlink{page.303}{303}\par
\hspace*{2.2em}unitary operator \hyperlink{page.296}{296}\par
\hspace*{1.1em}theory \hyperlink{page.205}{205}\par
\noindent spectrum \hyperlink{page.205}{205}, \hyperlink{page.290}{290}\par
\hspace*{1.1em}continuous -- \hyperlink{page.205}{205}, \hyperlink{page.311}{311}\par
\hspace*{1.1em}multiplicity of the -- \hyperlink{page.309}{309}\par
\hspace*{1.1em}point -- \hyperlink{page.205}{205}\par
\hspace*{1.1em}residual -- \hyperlink{page.205}{205}\par
\hspace*{1.1em}of a self-adjoint operator \hyperlink{page.307}{307}\par
\hspace*{1.1em}simple -- \hyperlink{page.310}{310}, \hyperlink{page.333}{333}\par
\noindent sphere, open -- \hyperlink{page.3}{3}\par
\noindent $ \sigma $-ring \hyperlink{page.14}{14}\par
\noindent \textsc{Steinhaus}, H. \hyperlink{page.69}{69}\par
\noindent \textsc{Stone}, M. H. \hyperlink{page.8}{8}, \hyperlink{page.227}{227}, \hyperlink{page.246}{246}, \hyperlink{page.311}{311}, \hyperlink{page.312}{312}\par
\noindent \hyperlink{page.331}{331}, \hyperlink{page.461}{461}\par
\hspace*{1.1em}theorem \hyperlink{page.246}{246}, \hyperlink{page.331}{331}\par
\hspace*{1.1em}--Weierstrass theorem \hyperlink{page.8}{8}\par
\noindent strictly elliptic differential operator \hyperlink{page.394}{394}\par
\noindent strong convergence \hyperlink{page.30}{30}, \hyperlink{page.31}{31}\par
\hspace*{1.1em}of operators \hyperlink{page.65}{65}\par
\noindent strong dual space \hyperlink{page.108}{108}\par
\noindent strong topology of operators \hyperlink{page.107}{107}\par
\noindent strongly elliptic differential operator \hyperlink{page.170}{170}, \hyperlink{page.276}{276}\par
\noindent strongly measurable mapping \hyperlink{page.126}{126}\par
\noindent subadditive \hyperlink{page.25}{25}\par
\noindent subharmonic \hyperlink{page.392}{392}\par
\noindent subspace \hyperlink{page.19}{19}\par
\noindent support (= carrier) of a\par
\hspace*{1.1em}distribution \hyperlink{page.59}{59}\par
\hspace*{1.1em}function \hyperlink{page.26}{26}\par
\noindent supremum (= least upper bound) \hyperlink{page.2}{2}\par
\noindent symmetric\par
\hspace*{1.1em}difference \hyperlink{page.18}{18}\par
\hspace*{1.1em}operator \hyperlink{page.193}{193}, \hyperlink{page.334}{334}, \hyperlink{page.338}{338}\par
\hspace*{1.1em}ring \hyperlink{page.291}{291}\par
\noindent \textsc{Szegő}, G. \hyperlink{page.86}{86}, \hyperlink{page.462}{462}\par
\noindent \textsc{Tanabe}, H. \hyperlink{page.411}{411}, \hyperlink{page.416}{416}, \hyperlink{page.417}{417}, \hyperlink{page.462}{462}\par
\noindent \textsc{Tannaka}, T. \hyperlink{page.323}{323}, \hyperlink{page.462}{462}\par
\hspace*{1.1em}duality theorem \hyperlink{page.323}{323}\par
\noindent Tauberian theorem\par
\hspace*{1.1em}special -- \hyperlink{page.344}{344}\par
\hspace*{1.1em}Wiener's -- \hyperlink{page.342}{342}\par
\noindent \textsc{Taylor}, A. \hyperlink{page.226}{226}, \hyperlink{page.462}{462}\par
\noindent Taylor's expansion \hyperlink{page.125}{125}\par
\noindent Tchebyshev polynomials \hyperlink{page.86}{86}\par
\noindent tempered distribution \hyperlink{page.146}{146}\par
\noindent temporally homogeneous\par
\hspace*{1.1em}equation of evolution \hyperlink{page.400}{400}\par
\hspace*{1.1em}Markov process \hyperlink{page.363}{363}\par
\noindent temporally inhomogeneous equation\par
\hspace*{1.1em}of evolution \hyperlink{page.400}{400}, \hyperlink{page.410}{410}\par
\noindent tensor product of distributions \hyperlink{page.63}{63}\par
\noindent testing function \hyperlink{page.45}{45}\par
\noindent $\theta $-formula \hyperlink{page.146}{146}\par
\noindent theorem\par
\hspace*{1.1em}Ascoli--Arzelà -- \hyperlink{page.82}{82}\par
\hspace*{1.1em}Baire -- \hyperlink{page.11}{11}\par
\hspace*{1.1em}Baire--Hausdorff -- \hyperlink{page.11}{11}\par
\hspace*{1.1em}Banach on\par
\hspace*{2.2em}a.e. convergence of linear\par
\hspace*{3.3em}operators \hyperlink{page.355}{355}\par
\hspace*{2.2em}generalized limits \hyperlink{page.101}{101}\par
\hspace*{1.1em}Bochner -- \hyperlink{page.128}{128}, \hyperlink{page.130}{130}\par
\hspace*{2.2em}on positive definite functions\par
\hspace*{1.1em}Cauchy's integral -- \hyperlink{page.124}{124}\par
\noindent theorem\par
\hspace*{1.1em}Chacon--Ornstein -- \hyperlink{page.369}{369}\par
\hspace*{1.1em}closed graph -- \hyperlink{page.74}{74}\par
\hspace*{1.1em}closed range -- \hyperlink{page.200}{200}\par
\hspace*{1.1em}Crandall--Liggett -- \hyperlink{page.430}{430}, \hyperlink{page.434}{434}, \hyperlink{page.448}{448}\par
\hspace*{1.1em}Dunford -- \hyperlink{page.124}{124}\par
\hspace*{1.1em}Eberlein--Shmulyan -- \hyperlink{page.137}{137}\par
\hspace*{1.1em}Egorov -- \hyperlink{page.14}{14}\par
\hspace*{1.1em}ergodic of Hille type \hyperlink{page.211}{211}\par
\noindent Feller -- \hyperlink{page.385}{385}\par
\noindent Fourier's inversion -- \hyperlink{page.144}{144}\par
\noindent Fréchet--Kolmogorov -- \hyperlink{page.265}{265}\par
\noindent Friedrichs on\par
\hspace*{1.1em}semi--bounded operators \hyperlink{page.305}{305}\par
\hspace*{1.1em}strongly elliptic differential equations \hyperlink{page.172}{172}\par
\noindent Fubini--Tonelli -- \hyperlink{page.16}{16}\par
\noindent Gelfand--Mazur -- \hyperlink{page.125}{125}\par
\noindent Gelfand--Raikov -- \hyperlink{page.341}{341}\par
\noindent general expansion -- \hyperlink{page.311}{311}\par
\noindent H.-- -- \hyperlink{page.375}{375}\par
\noindent Hahn--Banach -- \hyperlink{page.99}{99}, \hyperlink{page.102}{102}\par
\noindent Hilbert--Schmidt expansion -- \hyperlink{page.314}{314}\par
\noindent Hille--Yosida -- \hyperlink{page.242}{242}\par
\noindent Hormander -- \hyperlink{page.75}{75}, \hyperlink{page.184}{184}\par
\noindent individual ergodic -- \hyperlink{page.367}{367}, \hyperlink{page.371}{371}\par
\noindent Kolmogorov's extension of measures \hyperlink{page.283}{283}\par
\noindent Komura -- \hyperlink{page.429}{429}\par
\noindent Krein--Milman -- \hyperlink{page.347}{347}\par
\noindent Krylov--Weinstein -- \hyperlink{page.309}{309}\par
\noindent Lax--Milgram -- \hyperlink{page.88}{88}\par
\noindent Lebesgue--Fatou -- \hyperlink{page.16}{16}\par
\noindent Lebesgue--Nikodym -- \hyperlink{page.89}{89}\par
\noindent Lerch -- \hyperlink{page.162}{162}\par
\noindent Liouville -- \hyperlink{page.125}{125}, \hyperlink{page.364}{364}\par
\noindent Malgrange--Ehrenpreis -- \hyperlink{page.176}{176}\par
\noindent Mazur on\par
\hspace*{1.1em}existence of continuous linear functionals \hyperlink{page.105}{105}\par
\hspace*{1.1em}weak limits \hyperlink{page.117}{117}\par
\noindent mean ergodic -- \hyperlink{page.209}{209}, \hyperlink{page.210}{210}, \hyperlink{page.366}{366}, \hyperlink{page.371}{371}\par
\noindent Milman -- \hyperlink{page.123}{123}\par
\noindent Naimark -- \hyperlink{page.337}{337}\par
\noindent von Neumann on\par
\hspace*{1.1em}Cayley transform \hyperlink{page.197}{197}\par
\hspace*{1.1em}real operators \hyperlink{page.305}{305}\par
\hspace*{1.1em}spectral representation \hyperlink{page.304}{304}\par
\noindent Neumann--Riesz--Mimura -- \hyperlink{page.326}{326}\par
\noindent open mapping -- \hyperlink{page.71}{71}\par
\noindent Paley--Wiener -- \hyperlink{page.157}{157}, \hyperlink{page.158}{158}, \hyperlink{page.159}{159}\par
\noindent Peter--Weyl--Neumann -- \hyperlink{page.317}{317}\par
\noindent Pettis -- \hyperlink{page.126}{126}\par
\noindent Phillips -- \hyperlink{page.220}{220}, \hyperlink{page.264}{264}\par
\noindent Phillips--Lumer -- \hyperlink{page.243}{243}\par
\noindent Plancherel -- \hyperlink{page.150}{150}\par
\noindent Radon--Nikodym -- \hyperlink{page.360}{360}, \hyperlink{page.362}{362}\par
\noindent Rellich--Gärding -- \hyperlink{page.271}{271}\par
\noindent theorem\par
\hspace*{1.1em}representation -- for equicontinuous groups of class $ (C_{0}) $ \hyperlink{page.245}{245}\par
\hspace*{2.2em}semi-groups of class $ (C_{0}) $ \hyperlink{page.240}{240}\par
\hspace*{1.1em}resonance -- \hyperlink{page.65}{65}\par
\hspace*{1.1em}Riemann--Lebesgue -- \hyperlink{page.342}{342}\par
\hspace*{1.1em}Riesz--Markov--Kakutani -- \hyperlink{page.115}{115}\par
\hspace*{1.1em}Riesz' representation -- \hyperlink{page.86}{86}\par
\hspace*{1.1em}Schauder -- \hyperlink{page.272}{272}\par
\hspace*{1.1em}Schwartz -- 100, 158, 178 (Weyl--Schwartz --)\par
\hspace*{1.1em}special Tauberian -- \hyperlink{page.344}{344}\par
\hspace*{1.1em}spectral--mapping -- \hyperlink{page.222}{222}\par
\hspace*{1.1em}Stone -- \hyperlink{page.246}{246}, \hyperlink{page.331}{331}\par
\hspace*{1.1em}Stone--Weierstrass -- \hyperlink{page.8}{8}\par
\hspace*{1.1em}Tannaka's duality -- \hyperlink{page.323}{323}\par
\hspace*{1.1em}Tauberian -- of Wiener \hyperlink{page.342}{342}\par
\hspace*{1.1em}Titchmarsh -- \hyperlink{page.162}{162}\par
\hspace*{1.1em}Trotter--Kato -- \hyperlink{page.259}{259}\par
\hspace*{1.1em}Tychonov -- \hyperlink{page.5}{5}\par
\hspace*{1.1em}uniform boundedness -- \hyperlink{page.65}{65}\par
\hspace*{1.1em}Urysohn -- \hyperlink{page.7}{7}\par
\hspace*{1.1em}Vitali--Hahn--Saks -- \hyperlink{page.66}{66}\par
\hspace*{1.1em}Weierstrass -- on\par
\hspace*{2.2em}the existence of nowhere differentiable functions \hyperlink{page.68}{68}\par
\hspace*{2.2em}polynomial approximation \hyperlink{page.7}{7}, \hyperlink{page.243}{243}\par
\hspace*{2.2em}trigonometric approximation \hyperlink{page.10}{10}\par
\hspace*{1.1em}Weyl--Schwartz -- \hyperlink{page.172}{172}\par
\hspace*{1.1em}Wiener -- \hyperlink{page.292}{292}\par
\hspace*{1.1em}Wiener's Tauberian -- \hyperlink{page.342}{342}\par
\noindent \textsc{Titchmarsh}, E.C. \hyperlink{page.162}{162}, \hyperlink{page.312}{312}, \hyperlink{page.462}{462}\par
\hspace*{1.1em}theorem \hyperlink{page.161}{161}\par
\noindent topology \hyperlink{page.3}{3}\par
\hspace*{1.1em}bounded convergence -- \hyperlink{page.107}{107}\par
\hspace*{1.1em}relative -- \hyperlink{page.3}{3}\par
\hspace*{1.1em}simple convergence -- \hyperlink{page.107}{107}\par
\hspace*{1.1em}strong -- \hyperlink{page.108}{108}\par
\hspace*{1.1em}uniform -- of operators \hyperlink{page.108}{108}\par
\hspace*{1.1em}weak -- \hyperlink{page.108}{108}\par
\hspace*{1.1em}weak* -- \hyperlink{page.108}{108}\par
\noindent total\par
\hspace*{1.1em}ordering \hyperlink{page.2}{2}\par
\hspace*{1.1em}variation \hyperlink{page.34}{34}, \hyperlink{page.35}{35}, \hyperlink{page.36}{36}, \hyperlink{page.67}{67}, \hyperlink{page.114}{114}, \hyperlink{page.351}{351}\par
\noindent totally bounded \hyperlink{page.12}{12}\par
\noindent trace class, operators of the -- \hyperlink{page.271}{271}\par
\noindent transition probability \hyperlink{page.363}{363}\par
\noindent translation operator \hyperlink{page.154}{154}, \hyperlink{page.237}{237}\par
\noindent transposed matrix \hyperlink{page.190}{190}\par
\noindent \textsc{Trèves}, F. \hyperlink{page.187}{187}, \hyperlink{page.462}{462}\par
\noindent triangle inequality \hyperlink{page.29}{29}\par
\noindent trigonometric approximation theorem (Weierstrass) \hyperlink{page.10}{10}\par
\noindent \textsc{Trotter}, H. F. \hyperlink{page.262}{262}, \hyperlink{page.462}{462}\par
\hspace*{1.1em}--Kato theorem \hyperlink{page.259}{259}\par
\noindent \textsc{Tychonov}, A. theorem \hyperlink{page.5}{5}\par
\noindent uniform\par
\hspace*{1.1em}boundedness theorem \hyperlink{page.65}{65}\par
\hspace*{1.1em}strength \hyperlink{page.182}{182}\par
\hspace*{1.1em}topology of operators \hyperlink{page.107}{107}\par
\noindent uniformly\par
\hspace*{1.1em}continuous function \hyperlink{page.319}{319}\par
\hspace*{1.1em}convex $B$-space \hyperlink{page.123}{123}\par
\noindent unimodular group \hyperlink{page.214}{214}\par
\noindent unit \hyperlink{page.285}{285}, \hyperlink{page.356}{356}\par
\noindent unit sphere (or disk) \hyperlink{page.42}{42}\par
\noindent unitary\par
\hspace*{1.1em}character \hyperlink{page.340}{340}\par
\hspace*{1.1em}equivalence \hyperlink{page.311}{311}\par
\hspace*{1.1em}invariant \hyperlink{page.311}{311}\par
\hspace*{1.1em}matrix representation \hyperlink{page.317}{317}\par
\hspace*{1.1em}operator \hyperlink{page.197}{197}, \hyperlink{page.246}{246}, \hyperlink{page.331}{331}\par
\hspace*{2.2em}spectral resolution of a -- -- \hyperlink{page.296}{296}\par
\noindent unity, partition of -- \hyperlink{page.57}{57}\par
\noindent upper bound \hyperlink{page.2}{2}\par
\noindent \textsc{Urysohn}, P. theorem \hyperlink{page.7}{7}\par
\noindent variation\par
\hspace*{1.1em}negative \hyperlink{page.34}{34}, \hyperlink{page.35}{35}, \hyperlink{page.36}{36}, \hyperlink{page.351}{351}\par
\hspace*{1.1em}positive \hyperlink{page.34}{34}, \hyperlink{page.35}{35}, \hyperlink{page.351}{351}\par
\hspace*{1.1em}total \hyperlink{page.34}{34}, \hyperlink{page.35}{35}, \hyperlink{page.36}{36}, \hyperlink{page.67}{67}, \hyperlink{page.68}{68}, \hyperlink{page.114}{114}, \hyperlink{page.115}{115}, \hyperlink{page.350}{350}\par
\noindent vector lattice \hyperlink{page.349}{349}\par
\hspace*{1.1em}with unit \hyperlink{page.356}{356}\par
\noindent \textsc{Vilenkin}, N. Y. \hyperlink{page.271}{271}, \hyperlink{page.283}{283}, \hyperlink{page.462}{462}\par
\noindent \textsc{Visik}, I. M. \hyperlink{page.411}{411}, \hyperlink{page.462}{462}\par
\noindent \textsc{Vitali}, G. \hyperlink{page.66}{66}, \hyperlink{page.462}{462}\par
\hspace*{1.1em}--Hahn--Saks theorem \hyperlink{page.66}{66}\par
\noindent \textsc{Watanabe}, J. \hyperlink{page.258}{258}, \hyperlink{page.462}{462}\par
\noindent \textsc{Watanabe}, T. \hyperlink{page.398}{398}\par
\noindent wave equation \hyperlink{page.402}{402}\par
\noindent weak\par
\hspace*{1.1em}convergence \hyperlink{page.117}{117}\par
\noindent in $L^{1}(S, \mathfrak{B}, m)$ \hyperlink{page.118}{118}\par
\hspace*{1.1em}topology \hyperlink{page.109}{109}\par
\noindent weak*\par
\hspace*{1.1em}convergence \hyperlink{page.121}{121}\par
\hspace*{1.1em}dual space \hyperlink{page.108}{108}\par
\hspace*{1.1em}topology \hyperlink{page.108}{108}\par
\noindent weakly measurable mapping \hyperlink{page.126}{126}\par
\noindent \textsc{Wecken}, F. J. \hyperlink{page.311}{311}, \hyperlink{page.462}{462}\par
\noindent \textsc{Webb}, G. \hyperlink{page.462}{462}\par
\noindent \textsc{Weierstrass}, C. \hyperlink{page.7}{7}, \hyperlink{page.8}{8}, \hyperlink{page.69}{69}, \hyperlink{page.243}{243}\par
\hspace*{1.1em}Stone-- -- theorem \hyperlink{page.8}{8}\par
\hspace*{1.1em}theorem on\par
\hspace*{2.2em}the existence of nowhere differentiable functions \hyperlink{page.68}{68}\par
\hspace*{2.2em}polynomial approximation \hyperlink{page.7}{7}, \hyperlink{page.243}{243}\par
\hspace*{2.2em}trigonometric approximation \hyperlink{page.10}{10}\par
\noindent \textsc{Weil}, A. \hyperlink{page.215}{215}, \hyperlink{page.462}{462}\par
\noindent \textsc{Weinstein}, A. \hyperlink{page.309}{309}\par
\noindent \textsc{Weyl}, H. \hyperlink{page.75}{75}, \hyperlink{page.172}{172}, \hyperlink{page.312}{312}, \hyperlink{page.317}{317}, \hyperlink{page.463}{463}\par
\hspace*{1.1em}lemma \hyperlink{page.75}{75}\par
\noindent \textsc{Widder}, D. V. \hyperlink{page.253}{253}\par
\hspace*{1.1em}Post-- inversion formula \hyperlink{page.253}{253}\par
\noindent \textsc{Wiener}, N. \hyperlink{page.23}{23}, \hyperlink{page.157}{157}, \hyperlink{page.158}{158}, \hyperlink{page.159}{159}, \hyperlink{page.161}{161}, \hyperlink{page.286}{286}, \hyperlink{page.292}{292}, \hyperlink{page.342}{342}, \hyperlink{page.463}{463}\par
\hspace*{1.1em}Tauberian theorem \hyperlink{page.342}{342}\par
\hspace*{1.1em}theorem \hyperlink{page.291}{291}\par
\noindent \textsc{Yamabe}, H. \hyperlink{page.405}{405}, \hyperlink{page.463}{463}\par
\noindent \textsc{Yamada}, A. \hyperlink{page.398}{398}, \hyperlink{page.463}{463}\par
\noindent \textsc{Yosida}, K. \hyperlink{page.86}{86}, \hyperlink{page.91}{91}, \hyperlink{page.164}{164}, \hyperlink{page.173}{173}, \hyperlink{page.211}{211}, \hyperlink{page.227}{227}, \hyperlink{page.242}{242}, \hyperlink{page.252}{252}, \hyperlink{page.257}{257}, \hyperlink{page.259}{259}, \hyperlink{page.262}{262}, \hyperlink{page.275}{275}, \hyperlink{page.286}{286}, \hyperlink{page.295}{295}, \hyperlink{page.309}{309}, \hyperlink{page.311}{311}, \hyperlink{page.313}{313}, \hyperlink{page.324}{324}, \hyperlink{page.355}{355}, \hyperlink{page.359}{359}, \hyperlink{page.362}{362}, \hyperlink{page.366}{366}, \hyperlink{page.367}{367}, \hyperlink{page.375}{375}, \hyperlink{page.376}{376}, \hyperlink{page.380}{380}, \hyperlink{page.385}{385}, \hyperlink{page.392}{392}, \hyperlink{page.405}{405}, \hyperlink{page.410}{410}, \hyperlink{page.411}{411}, \hyperlink{page.415}{415}, \hyperlink{page.416}{416}, \hyperlink{page.437}{437}, \hyperlink{page.463}{463}\par
\hspace*{1.1em}Hille-- -- theorem \hyperlink{page.242}{242}\par
\hspace*{1.1em}individual ergodic theorem \hyperlink{page.367}{367}, \hyperlink{page.371}{371}\par
\hspace*{1.1em}mean ergodic theorem \hyperlink{page.209}{209}, \hyperlink{page.366}{366}\par
\noindent \textsc{Yushkevitch}, A. A. \hyperlink{page.380}{380}, \hyperlink{page.465}{465}\par
\noindent \textsc{Zorn}, M. lemma \hyperlink{page.2}{2}\par
\noindent \textsc{Zygmund}, A. -class \hyperlink{page.375}{375}\par
\end{multicols}
\yoschapterstar{Notation of Spaces}
\begin{multicols}{2}\raggedcolumns\footnotesize\raggedright\setlength{\parindent}{0pt}\setlength{\parskip}{0pt}
\noindent $ A^{2}(G) $ \hyperlink{page.39}{39}, \hyperlink{page.51}{51}, \hyperlink{page.92}{92}\par
\noindent $A(S, \mathfrak{B})$ \hyperlink{page.34}{34}, \hyperlink{page.354}{354}\par
\noindent $(c)$ \hyperlink{page.34}{34}, \hyperlink{page.112}{112}\par
\noindent $ (c_{0}) $ \hyperlink{page.33}{33}, \hyperlink{page.110}{110}\par
\noindent $C(S)$ \hyperlink{page.9}{9}, \hyperlink{page.31}{31}, \hyperlink{page.285}{285}, \hyperlink{page.291}{291}, \hyperlink{page.354}{354}\par
\noindent $ C^{k}(\Omega) $ \hyperlink{page.26}{26}, \hyperlink{page.39}{39}\par
\noindent $ C_0^k(\Omega) $ \hyperlink{page.26}{26}, \hyperlink{page.39}{39}, \hyperlink{page.54}{54}\par
\noindent $ \mathfrak{D}_{K}(\Omega) $ \hyperlink{page.28}{28}, \hyperlink{page.135}{135}, \hyperlink{page.281}{281}\par
\noindent $ \mathfrak{D}(\Omega) $ \hyperlink{page.28}{28}, \hyperlink{page.44}{44}, \hyperlink{page.283}{283}\par
\noindent $ \mathfrak{E}^{k}(\Omega) $ \hyperlink{page.27}{27}, \hyperlink{page.36}{36}\par
\noindent $ \mathfrak{E}(\Omega) = \mathfrak{E}^\infty(\Omega) $ \hyperlink{page.27}{27}, \hyperlink{page.50}{50}, \hyperlink{page.59}{59}, \hyperlink{page.283}{283}\par
\noindent $\exp(tA)$ \hyperlink{page.260}{260}\par
\noindent $H$-$L^{2}$ \hyperlink{page.39}{39}, \hyperlink{page.52}{52}\par
\noindent $ \hat{H}^{k}(\Omega) $ \hyperlink{page.39}{39}, \hyperlink{page.54}{54}\par
\noindent $ \hat{H}_0^k(\Omega) $ \hyperlink{page.39}{39}, \hyperlink{page.54}{54}\par
\noindent $ H^k(\Omega) $ \hyperlink{page.54}{54}\par
\noindent $ H_0^k(\Omega) $ \hyperlink{page.54}{54}, \hyperlink{page.94}{94}, \hyperlink{page.269}{269}, \hyperlink{page.271}{271}\par
\noindent $ H_0^{-k}(\Omega) $ \hyperlink{page.95}{95}\par
\hspace*{1.1em}$(l^{p})$ \hyperlink{page.34}{34}, \hyperlink{page.113}{113}\par
\noindent $L^{1}$ \hyperlink{page.340}{340}\par
\noindent $L(X, Y)$ \hyperlink{page.107}{107}, \hyperlink{page.285}{285}\par
\noindent $ L^{p}(S), L^{p}(S, \mathfrak{B}, m) $ \hyperlink{page.31}{31}, \hyperlink{page.51}{51}, \hyperlink{page.112}{112}, \hyperlink{page.265}{265}, \hyperlink{page.354}{354}\par
\noindent $ L^{\infty}(S), L^{\infty}(S, \mathfrak{B}, m) $ \hyperlink{page.33}{33}, \hyperlink{page.51}{51}, \hyperlink{page.114}{114}\par
\noindent $(m)$ \hyperlink{page.34}{34}\par
\noindent $ M(S, \mathfrak{B}, m) $ \hyperlink{page.36}{36}, \hyperlink{page.43}{43}, \hyperlink{page.52}{52}, \hyperlink{page.113}{113}, \hyperlink{page.354}{354}\par
\noindent $ \mathfrak{O}_{M}(\mathbb{R}^{n}) $ \hyperlink{page.147}{147}\par
\hspace*{1.1em}$(s)$ \hyperlink{page.52}{52}\par
\noindent $ \mathfrak{S}(\mathbb{R}^{n}) $ \hyperlink{page.143}{143}, \hyperlink{page.283}{283}\par
\hspace*{1.1em}$(S, \mathfrak{B}, m)$ \hyperlink{page.14}{14}, \hyperlink{page.67}{67}\par
\noindent $ W^{k,p}(\Omega) $ \hyperlink{page.52}{52}\par
\noindent $ W^{k}(\Omega) $ \hyperlink{page.52}{52}, \hyperlink{page.55}{55}, \hyperlink{page.152}{152}\par
\end{multicols}


\end{document}
